% Auto-generated from data/segments.json + layout.json (+ transforms) — do not edit by hand.
% volume: 16An02
% mode: reading
% segments: 3292
% transform_rules: 0
% toc_marks: 564

% Volume layout defaults (layout.json ``layout``).
\csromanlayoutapply{1}{1.5}{21.6pt}{5pt}{6.3pt}{65pt}{2.5em}%

% Reading physical-page layout (layout.json ``page_layout_reading_mode``).
\csromanreadingpagelayoutvolume{1}{1.5}{21.6pt}{5pt}{6.3pt}{65pt}{2.5em}%
\csromanreadingpagelayoutenable

\csromanheader{2}{\textbf{องฺคุตฺตรนิกาย}}
\csromanmatikaanchor{mtk.0}
\csromantocmarknopage{chapter}{ปญฺจกนิปาตปาฬิ}{mtk.0}
\bookmark[dest={mtk.0},level=0]{ปญฺจกนิปาตปาฬิ}
\gambhira{\textbf{ปญฺจกนิปาตปาฬิ}}
\namakkaram{นโม ตสฺส ภควโต อรหโต สมฺมาสมฺพุทฺธสฺส.}
\csromanmatikaanchor{mtk.1}
\csromanmatikaanchor{mtk.126}
\csromanmatikaanchor{mtk.127}
\csromantocmarknopage{section}{(10) 5. กกุธวคฺค}{mtk.1}
\bookmark[dest={mtk.1},level=1]{(10) 5. กกุธวคฺค}
\chapterhead{1. ปฐมปณฺณาสก}
\csromanheader{2}{1. สํขิตฺตสุตฺต}
\proseitem{1}{\csromanfolio{1}เอวํ เม สุตํ\csromandash{}เอกํ สมยํ ภควา สาวตฺถิยํ วิหรติ เชตวเน อนาถปิณฺฑิกสฺส อาราเม.\csromanspacer{} ตตฺร โข ภควา ภิกฺขู อามนฺเตสิ “ภิกฺขโว”ติ. “ภทนฺเต”ติ เต ภิกฺขู ภควโต ปจฺจสฺโสสุํ.\csromanspacer{} ภควา เอตทโวจ\csromandash{}}

\prose{ปญฺจิมานิ ภิกฺขเว เสขพลานิ\footnote{เสกฺขพลานิ (ก)}.\csromanspacer{} กตมานิ ปญฺจ? สทฺธาพลํ หิรีพลํ\footnote{หิริพลํ (สี, อิ)} โอตฺตปฺปพลํ วีริยพลํ\footnote{วิริยพลํ (สี, สฺยา, กํ, อิ)} ปญฺญาพลํ.\csromanspacer{} อิมานิ โข ภิกฺขเว ปญฺจ เสขพลานิ.}

\prose{ตสฺมาติห ภิกฺขเว เอวํ สิกฺขิตพฺพํ “สทฺธาพเลน สมนฺนาคตา ภวิสฺสาม เสขพเลน, หิรีพเลน สมนฺนาคตา ภวิสฺสาม เสขพเลน, โอตฺตปฺปพเลน สมนฺนาคตา ภวิสฺสาม เสขพเลน, วีริยพเลน สมนฺนาคตา ภวิสฺสาม เสขพเลน, ปญฺญาพเลน สมนฺนาคตา ภวิสฺสาม เสขพเลนา”ติ.\csromanspacer{} เอวํ หิ โว ภิกฺขเว สิกฺขิตพฺพนฺติ.\csromanspacer{} อิทมโวจ ภควา.\csromanspacer{} อตฺตมนา เต ภิกฺขู ภควโต ภาสิตํ อภินนฺทุนฺติ.\csromanspacer{}\csromanspacer{}ปฐมํ.}

\csromanheader{2}{2. วิตฺถตสุตฺต}
\proseitem{2}{\csromanfolio{2}ปญฺจิมานิ ภิกฺขเว เสขพลานิ.\csromanspacer{} กตมานิ ปญฺจ? สทฺธาพลํ หิรีพลํ โอตฺตปฺปพลํ วีริยพลํ ปญฺญาพลํ.\csromanspacer{} กตมญฺจ ภิกฺขเว สทฺธาพลํ? อิธ ภิกฺขเว อริยสาวโก สทฺโธ โหติ, สทฺทหติ ตถาคตสฺส โพธิํ “อิติปิ โส ภควา อรหํ สมฺมาสมฺพุทฺโธ วิชฺชาจรณสมฺปนฺโน สุคโต โลกวิทู อนุตฺตโร ปุริสทมฺมสารถิ สตฺถา เทวมนุสฺสานํ พุทฺโธ ภควา”ติ.\csromanspacer{} อิทํ วุจฺจติ ภิกฺขเว สทฺธาพลํ.}

\prose{กตมญฺจ ภิกฺขเว หิรีพลํ? อิธ ภิกฺขเว อริยสาวโก หิรีมา โหติ, หิรียติ กายทุจฺจริเตน วจีทุจฺจริเตน มโนทุจฺจริเตน, หิรียติ ปาปกานํ อกุสลานํ ธมฺมานํ สมาปตฺติยา.\csromanspacer{} อิทํ วุจฺจติ ภิกฺขเว หิรีพลํ.}

\prose{กตมญฺจ ภิกฺขเว โอตฺตปฺปพลํ? อิธ ภิกฺขเว อริยสาวโก โอตฺตปฺปี โหติ, โอตฺตปฺปติ กายทุจฺจริเตน วจีทุจฺจริเตน มโนทุจฺจริเตน, โอตฺตปฺปติ ปาปกานํ อกุสลานํ ธมฺมานํ สมาปตฺติยา.\csromanspacer{} อิทํ วุจฺจติ ภิกฺขเว โอตฺตปฺปพลํ.}

\prose{กตมญฺจ ภิกฺขเว วีริยพลํ? อิธ ภิกฺขเว อริยสาวโก อารทฺธวีริโย วิหรติ อกุสลานํ ธมฺมานํ ปหานาย กุสลานํ ธมฺมานํ อุปสมฺปทาย, ถามวา ทฬฺหปรกฺกโม อนิกฺขิตฺตธุโร กุสเลสุ ธมฺเมสุ.\csromanspacer{} อิทํ วุจฺจติ ภิกฺขเว วีริยพลํ.}

\prose{กตมญฺจ ภิกฺขเว ปญฺญาพลํ? อิธ ภิกฺขเว อริยสาวโก ปญฺญวา โหติ, อุทยตฺถคามินิยา ปญฺญาย สมนฺนาคโต อริยาย นิพฺเพธิกาย สมฺมา ทุกฺขกฺขยคามินิยา, อิทํ วุจฺจติ ภิกฺขเว ปญฺญาพลํ.\csromanspacer{} อิมานิ โข ภิกฺขเว ปญฺจ เสขพลานิ.}

\prose{ตสฺมาติห ภิกฺขเว เอวํ สิกฺขิตพฺพํ “สทฺธาพเลน สมนฺนาคตา ภวิสฺสาม เสขพเลน, หิรีพเลน.\csromanspacer{} โอตฺตปฺปพเลน.\csromanspacer{} วีริยพเลน.\csromanspacer{} ปญฺญาพเลน สมนฺนาคตา ภวิสฺสาม เสขพเลนา”ติ.\csromanspacer{} เอวํ หิ โข ภิกฺขเว สิกฺขิตพฺพนฺติ.\csromanspacer{}\csromanspacer{}ทุติยํ.}

\csromanheader{2}{3. ทุกฺขสุตฺต}
\proseitem{3}{ปญฺจหิ ภิกฺขเว ธมฺเมหิ สมนฺนาคโต ภิกฺขุ ทิฏฺเฐว ธมฺเม ทุกฺขํ วิหรติ สวิฆาตํ ส-อุปายาสํ สปริฬาหํ, กายสฺส จ เภทา ปรํ \csromanfolio{3}มรณา ทุคฺคติ ปาฏิกงฺขา.\csromanspacer{} กตเมหิ ปญฺจหิ? อิธ ภิกฺขเว ภิกฺขุ อสฺสทฺโธ โหติ, อหิริโก โหติ.\csromanspacer{} อโนตฺตปฺปี โหติ, กุสีโต โหติ, ทุปฺปญฺโญ โหติ.\csromanspacer{} อิเมหิ โข ภิกฺขเว ปญฺจหิ ธมฺเมหิ สมนฺนาคโต ภิกฺขุ ทิฏฺเฐว ธมฺเม ทุกฺขํ วิหรติ สวิฆาตํ ส-อุปายาสํ สปริฬาหํ, กายสฺส จ เภทา ปรํ มรณา ทุคฺคติ ปาฏิกงฺขา.}

\prose{ปญฺจหิ ภิกฺขเว ธมฺเมหิ สมนฺนาคโต ภิกฺขุ ทิฏฺเฐว ธมฺเม สุขํ วิหรติ อวิฆาตํ อนุปายาสํ อปริฬาหํ, กายสฺส จ เภทา ปรํ มรณา สุคติ ปาฏิกงฺขา.\csromanspacer{} กตเมหิ ปญฺจหิ? อิธ ภิกฺขเว ภิกฺขุ สทฺโธ โหติ, หิรีมา โหติ, โอตฺตปฺปี โหติ, อารทฺธวีริโย โหติ, ปญฺญวา โหติ.\csromanspacer{} อิเมหิ โข ภิกฺขเว ปญฺจหิ ธมฺเมหิ สมนฺนาคโต ภิกฺขุ ทิฏฺเฐว ธมฺเม สุขํ วิหรติ อวิฆาตํ อนุปายาสํ อปริฬาหํ, กายสฺส จ เภทา ปรํ มรณา สุคติ ปาฏิกงฺขาติ.\csromanspacer{}\csromanspacer{}ตติยํ.}

\csromanheader{2}{4. ยถาภตสุตฺต}
\proseitem{4}{ปญฺจหิ ภิกฺขเว ธมฺเมหิ สมนฺนาคโต ภิกฺขุ ยถาภตํ นิกฺขิตฺโต เอวํ นิรเย.\csromanspacer{} กตเมหิ ปญฺจหิ? อิธ ภิกฺขเว ภิกฺขุ อสฺสทฺโธ โหติ, อหิริโก โหติ, อโนตฺตปฺปี โหติ, กุสีโต โหติ, ทุปฺปญฺโญ โหติ.\csromanspacer{} อิเมหิ โข ภิกฺขเว ปญฺจหิ ธมฺเมหิ สมนฺนาคโต ภิกฺขุ ยถาภตํ นิกฺขิตฺโต เอวํ นิรเย.}

\prose{ปญฺจหิ ภิกฺขเว ธมฺเมหิ สมนฺนาคโต ภิกฺขุ ยถาภตํ นิกฺขิตฺโต เอวํ สคฺเค.\csromanspacer{} กตเมหิ ปญฺจหิ? อิธ ภิกฺขเว ภิกฺขุ สทฺโธ โหติ, หิรีมา โหติ, โอตฺตปฺปี โหติ, อารทฺธวีริโย โหติ, ปญฺญวา โหติ.\csromanspacer{} อิเมหิ โข ภิกฺขเว ปญฺจหิ ธมฺเมหิ สมนฺนาคโต ภิกฺขุ ยถาภตํ นิกฺขิตฺโต เอวํ สคฺเคติ.\csromanspacer{}\csromanspacer{}จตุตฺถํ.}

\csromanheader{2}{5. สิกฺขาสุตฺต}
\proseitem{5}{โย หิ โกจิ ภิกฺขเว ภิกฺขุ วา ภิกฺขุนี วา สิกฺขํ ปจฺจกฺขาย หีนายาวตฺตติ, ตสฺส ทิฏฺเฐว\footnote{ทิฏฺเฐ เจว (สี)} ธมฺเม ปญฺจ สหธมฺมิกา วาทานุปาตา\footnote{วาทานุวาทา (อํ 3. 26; อํ 1. 161 ปิฏฺเฐปิ.)} คารยฺหา ฐานา อาคจฺฉนฺติ.\csromanspacer{} กตเม ปญฺจ? สทฺธาปิ นาม เต นาโหสิ กุสเลสุ ธมฺเมสุ, หิรีปิ นาม เต นาโหสิ กุสเลสุ ธมฺเมสุ, \csromanfolio{4}โอตฺตปฺปมฺปิ นาม เต นาโหสิ กุสเลสุ ธมฺเมสุ, วีริยมฺปิ นาม เต นาโหสิ กุสเลสุ ธมฺเมสุ, ปญฺญาปิ นาม เต นาโหสิ กุสเลสุ ธมฺเมสุ.\csromanspacer{} โย หิ โกจิ ภิกฺขเว ภิกฺขุ วา ภิกฺขุนี วา สิกฺขํ ปจฺจกฺขาย หีนายาวตฺตติ, ตสฺส ทิฏฺเฐว ธมฺเม อิเม ปญฺจ สหธมฺมิกา วาทานุปาตา คารยฺหา ฐานา อาคจฺฉนฺติ.}

\prose{โย หิ โกจิ ภิกฺขเว ภิกฺขุ วา ภิกฺขุนี วา สหาปิ ทุกฺเขน สหาปิ โทมนสฺเสน อสฺสุมุโข\footnote{อสฺสุมุโขปิ (สฺยา)} รุทมาโน ปริปุณฺณํ ปริสุทฺธํ พฺรหฺมจริยํ จรติ, ตสฺส ทิฏฺเฐว ธมฺเม ปญฺจ สหธมฺมิกา ปาสํสา ฐานา\footnote{ปาสํสํ ฐานํ (สฺยา)} อาคจฺฉนฺติ.\csromanspacer{} กตเม ปญฺจ? สทฺธาปิ นาม เต อโหสิ กุสเลสุ ธมฺเมสุ, หิรีปิ นาม เต อโหสิ กุสเลสุ ธมฺเมสุ, โอตฺตปฺปมฺปิ นาม เต อโหสิ กุสเลสุ ธมฺเมสุ, วีริยมฺปิ นาม เต อโหสิ กุสเลสุ ธมฺเมสุ, ปญฺญาปิ นาม เต อโหสิ กุสเลสุ ธมฺเมสุ.\csromanspacer{} โย หิ โกจิ ภิกฺขเว ภิกฺขุ วา ภิกฺขุนี วา สหาปิ ทุกฺเขน สหาปิ โทมนสฺเสน อสฺสุมุโข รุทมาโน ปริปุณฺณํ ปริสุทฺธํ พฺรหฺมจริยํ จรติ, ตสฺส ทิฏฺเฐว ธมฺเม อิเม ปญฺจ สหธมฺมิกา ปาสํสา ฐานา อาคจฺฉนฺตีติ.\csromanspacer{}\csromanspacer{}ปญฺจมํ.}

\csromanheader{2}{6. สมาปตฺติสุตฺต}
\proseitem{6}{น ตาว ภิกฺขเว อกุสลสฺส สมาปตฺติ โหติ, ยาว สทฺธา ปจฺจุปฏฺฐิตา โหติ กุสเลสุ ธมฺเมสุ.\csromanspacer{} ยโต จ โข ภิกฺขเว สทฺธา อนฺตรหิตา โหติ อสฺสทฺธิยํ ปริยุฏฺฐาย ติฏฺฐติ, อถ อกุสลสฺส สมาปตฺติ}

\prose{น ตาว ภิกฺขเว อกุสลสฺส สมาปตฺติ โหติ, ยาว หิรี ปจฺจุปฏฺฐิตา โหติ กุสเลสุ ธมฺเมสุ.\csromanspacer{} ยโต จ โข ภิกฺขเว หิรี อนฺตรหิตา โหติ อหิริกํ ปริยุฏฺฐาย ติฏฺฐติ, อถ อกุสลสฺส สมาปตฺติ โหติ.}

\prose{น ตาว ภิกฺขเว อกุสลสฺส สมาปตฺติ โหติ, ยาว โอตฺตปฺปํ ปจฺจุปฏฺฐิตํ โหติ กุสเลสุ ธมฺเมสุ.\csromanspacer{} ยโต จ โข ภิกฺขเว โอตฺตปฺปํ อนฺตรหิตํ โหติ อโนตฺตปฺปํ ปริยุฏฺฐาย ติฏฺฐติ, อถ อกุสลสฺส สมาปตฺติ}

\prose{\csromanfolio{5}น ตาว ภิกฺขเว อกุสลสฺส สมาปตฺติ โหติ, ยาว วีริยํ ปจฺจุปฏฺฐิตํ โหติ กุสเลสุ ธมฺเมสุ.\csromanspacer{} ยโต จ โข ภิกฺขเว วีริยํ อนฺตรหิตํ โหติ โกสชฺชํ ปริยุฏฺฐาย ติฏฺฐติ, อถ อกุสลสฺส สมาปตฺติ}

\prose{น ตาว ภิกฺขเว อกุสลสฺส สมาปตฺติ โหติ, ยาว ปญฺญา ปจฺจุปฏฺฐิตา โหติ กุสเลสุ ธมฺเมสุ.\csromanspacer{} ยโต จ โข ภิกฺขเว ปญฺญา อนฺตรหิตา โหติ ทุปฺปญฺญา\footnote{ทุปฺปญฺญํ (ก)} ปริยุฏฺฐาย ติฏฺฐติ, อถ อกุสลสฺส สมาปตฺติ โหตีติ.\csromanspacer{}\csromanspacer{}ฉฏฺฐํ.}

\csromanheader{2}{7. กามสุตฺต}
\proseitem{7}{เยภุยฺเยน ภิกฺขเว สตฺตา กาเมสุ ลฬิตา\footnote{ปลาฬิตา (สี)}.\csromanspacer{} อสิตพฺยาภงฺคิํ\footnote{อสิตพฺยาภงฺคิํ เจปิ (?)} ภิกฺขเว กุลปุตฺโต โอหาย อคารสฺมา อนคาริยํ ปพฺพชิโต โหติ, “สทฺธาปพฺพชิโต กุลปุตฺโต”ติ อลํ วจนาย.\csromanspacer{} ตํ กิสฺส เหตุ? ลพฺภา\footnote{ลพฺภา หิ (สฺยา)} ภิกฺขเว โยพฺพเนน กามา, เต จ โข ยาทิสา วา ตาทิสา วา.\csromanspacer{} เย จ ภิกฺขเว หีนา กามา, เย จ มชฺฌิมา กามา, เย จ ปณีตา กามา, สพฺเพ กามา “กามา”เตฺวว สงฺขํ คจฺฉนฺติ.\csromanspacer{} เสยฺยถาปิ ภิกฺขเว ทหโร กุมาโร มนฺโท อุตฺตานเสยฺยโก ธาติยา ปมาทมนฺวาย กฏฺฐํ วา กฐลํ\footnote{กถลํ (ก)} วา มุเข อาหเรยฺย, ตเมนํ ธาติ สีฆํ สีฆํ\footnote{สีฆสีฆํ (สี)} มนสิ กเรยฺย, สีฆํ สีฆํ มนสิ กริตฺวา สีฆํ สีฆํ อาหเรยฺย.\csromanspacer{} โน เจ สกฺกุเณยฺย สีฆํ สีฆํ อาหริตุํ, วาเมน หตฺเถน สีสํ ปริคฺคเหตฺวา ทกฺขิเณน หตฺเถน วงฺกงฺคุลิํ กริตฺวา สโลหิตมฺปิ อาหเรยฺย.\csromanspacer{} ตํ กิสฺส เหตุ? อตฺเถสา ภิกฺขเว กุมารสฺส วิเหสา, เนสา นตฺถีติ วทามิ.\csromanspacer{} กรณียญฺจ โข เอตํ\footnote{เอวํ (ก)} ภิกฺขเว ธาติยา อตฺถกามาย หิเตสินิยา อนุกมฺปิกาย อนุกมฺปํ อุปาทาย.\csromanspacer{} ยโต จ โข ภิกฺขเว โส กุมาโร วุทฺโธ โหติ อลํปญฺโญ, อนเปกฺขา ทานิ\footnote{อนเปกฺขา ปน (สี, สฺยา, กํ)} ภิกฺขเว ธาติ ตสฺมิํ กุมาเร โหติ “อตฺตคุตฺโต ทานิ กุมาโร, นาลํ ปมาทายา”ติ.}

\prose{เอวเมวํ โข ภิกฺขเว ยาวกีวญฺจ ภิกฺขุโน สทฺธาย อกตํ โหติ กุสเลสุ ธมฺเมสุ, หิริยา อกตํ โหติ กุสเลสุ ธมฺเมสุ, โอตฺตปฺเปน อกตํ โหติ กุสเลสุ ธมฺเมสุ, วีริเยน อกตํ \csromanfolio{6}โหติ กุสเลสุ ธมฺเมสุ, ปญฺญาย อกตํ โหติ กุสเลสุ ธมฺเมสุ.\csromanspacer{} อนุรกฺขิตพฺโพ ตาว เม โส ภิกฺขเว ภิกฺขุ โหติ.\csromanspacer{} ยโต จ โข ภิกฺขเว ภิกฺขุโน สทฺธาย กตํ โหติ กุสเลสุ ธมฺเมสุ, หิริยา กตํ โหติ กุสเลสุ ธมฺเมสุ, โอตฺตปฺเปน กตํ โหติ กุสเลสุ ธมฺเมสุ, วีริเยน กตํ โหติ กุสเลสุ ธมฺเมสุ, ปญฺญาย กตํ โหติ กุสเลสุ ธมฺเมสุ.\csromanspacer{} อนเปกฺโข ทานาหํ\footnote{ปนาหํ (สี, สฺยา, กํ)} ภิกฺขเว ตสฺมิํ ภิกฺขุสฺมิํ โหมิ “อตฺตคุตฺโต ทานิ ภิกฺขุ, นาลํ ปมาทายา”ติ.\csromanspacer{}\csromanspacer{}สตฺตมํ.}

\csromanheader{2}{8. จวนสุตฺต}
\proseitem{8}{ปญฺจหิ ภิกฺขเว ธมฺเมหิ สมนฺนาคโต ภิกฺขุ จวติ นปฺปติฏฺฐาติ สทฺธมฺเม.\csromanspacer{} กตเมหิ ปญฺจหิ? อสฺสทฺโธ ภิกฺขเว ภิกฺขุ จวติ นปฺปติฏฺฐาติ สทฺธมฺเม, อหิริโก ภิกฺขเว ภิกฺขุ จวติ นปฺปติฏฺฐาติ สทฺธมฺเม, อโนตฺตปฺปี ภิกฺขเว ภิกฺขุ จวติ นปฺปติฏฺฐาติ สทฺธมฺเม, กุสีโต ภิกฺขเว ภิกฺขุ จวติ นปฺปติฏฺฐาติ สทฺธมฺเม, ทุปฺปญฺโญ ภิกฺขเว ภิกฺขุ จวติ นปฺปติฏฺฐาติ สทฺธมฺเม.\csromanspacer{} อิเมหิ โข ภิกฺขเว ปญฺจหิ ธมฺเมหิ สมนฺนาคโต ภิกฺขุ จวติ นปฺปติฏฺฐาติ สทฺธมฺเม.}

\prose{ปญฺจหิ ภิกฺขเว ธมฺเมหิ สมนฺนาคโต ภิกฺขุ น จวติ ปติฏฺฐาติ สทฺธมฺเม.\csromanspacer{} กตเมหิ ปญฺจหิ? สทฺโธ ภิกฺขเว ภิกฺขุ น จวติ ปติฏฺฐาติ สทฺธมฺเม, หิรีมา ภิกฺขเว ภิกฺขุ น จวติ ปติฏฺฐาติ สทฺธมฺเม, โอตฺตปฺปี ภิกฺขเว ภิกฺขุ น จวติ ปติฏฺฐาติ สทฺธมฺเม, อารทฺธวีริโย ภิกฺขเว ภิกฺขุ น จวติ ปติฏฺฐาติ สทฺธมฺเม, ปญฺญวา ภิกฺขเว ภิกฺขุ น จวติ ปติฏฺฐาติ สทฺธมฺเม.\csromanspacer{} อิเมหิ โข ภิกฺขเว ปญฺจหิ ธมฺเมหิ สมนฺนาคโต ภิกฺขุ น จวติ ปติฏฺฐาติ สทฺธมฺเมติ.\csromanspacer{}\csromanspacer{}อฏฺฐมํ.}

\csromanheader{2}{9. ปฐม อคารวสุตฺต}
\proseitem{9}{ปญฺจหิ ภิกฺขเว ธมฺเมหิ สมนฺนาคโต ภิกฺขุ อคารโว อปฺปติสฺโส จวติ นปฺปติฏฺฐาติ สทฺธมฺเม.\csromanspacer{} กตเมหิ ปญฺจหิ? อสฺสทฺโธ ภิกฺขเว ภิกฺขุ อคารโว อปฺปติสฺโส จวติ นปฺปติฏฺฐาติ สทฺธมฺเม, อหิริโก ภิกฺขเว ภิกฺขุ อคารโว อปฺปติสฺโส จวติ นปฺปติฏฺฐาติ สทฺธมฺเม, อโนตฺตปฺปี ภิกฺขเว ภิกฺขุ อคารโว อปฺปติสฺโส จวติ นปฺปติฏฺฐาติ สทฺธมฺเม, \csromanfolio{7}กุสีโต ภิกฺขเว ภิกฺขุ อคารโว อปฺปติสฺโส จวติ นปฺปติฏฺฐาติ สทฺธมฺเม, ทุปฺปญฺโญ ภิกฺขเว ภิกฺขุ อคารโว อปฺปติสฺโส จวติ นปฺปติฏฺฐาติ สทฺธมฺเม.\csromanspacer{} อิเมหิ โข ภิกฺขเว ปญฺจหิ ธมฺเมหิ สมนฺนาคโต ภิกฺขุ อคารโว อปฺปติสฺโส จวติ นปฺปติฏฺฐาติ สทฺธมฺเม.}

\prose{ปญฺจหิ ภิกฺขเว ธมฺเมหิ สมนฺนาคโต ภิกฺขุ สคารโว สปฺปติสฺโส น จวติ ปติฏฺฐาติ สทฺธมฺเม.\csromanspacer{} กตเมหิ ปญฺจหิ? สทฺโธ ภิกฺขเว ภิกฺขุ สคารโว สปฺปติสฺโส น จวติ ปติฏฺฐาติ สทฺธมฺเม, หิรีมา ภิกฺขเว ภิกฺขุ สคารโว สปฺปติสฺโส น จวติ ปติฏฺฐาติ สทฺธมฺเม, โอตฺตปฺปี ภิกฺขเว ภิกฺขุ สคารโว สปฺปติสฺโส น จวติ ปติฏฺฐาติ สทฺธมฺเม, อารทฺธวีริโย ภิกฺขเว ภิกฺขุ สคารโว สปฺปติสฺโส น จวติ ปติฏฺฐาติ สทฺธมฺเม, ปญฺญวา ภิกฺขเว ภิกฺขุ สคารโว สปฺปติสฺโส น จวติ ปติฏฺฐาติ สทฺธมฺเม.\csromanspacer{} อิเมหิ โข ภิกฺขเว ปญฺจหิ ธมฺเมหิ สมนฺนาคโต ภิกฺขุ สคารโว สปฺปติสฺโส น จวติ ปติฏฺฐาติ สทฺธมฺเมติ.\csromanspacer{}\csromanspacer{}นวมํ.}

\csromanheader{2}{10. ทุติย อคารวสุตฺต}
\proseitem{10}{ปญฺจหิ ภิกฺขเว ธมฺเมหิ สมนฺนาคโต ภิกฺขุ อคารโว อปฺปติสฺโส อภพฺโพ อิมสฺมิํ ธมฺมวินเย วุทฺธิํ วิรูฬฺหิํ เวปุลฺลํ อาปชฺชิตุํ.\csromanspacer{} กตเมหิ ปญฺจหิ? อสฺสทฺโธ ภิกฺขเว ภิกฺขุ อคารโว อปฺปติสฺโส อภพฺโพ อิมสฺมิํ ธมฺมวินเย วุทฺธิํ วิรูฬฺหิํ เวปุลฺลํ อาปชฺชิตุํ, อหิริโก ภิกฺขเว ภิกฺขุ อคารโว อปฺปติสฺโส อภพฺโพ อิมสฺมิํ ธมฺมวินเย วุทฺธิํ วิรูฬฺหิํ เวปุลฺลํ อาปชฺชิตุํ, อโนตฺตปฺปี ภิกฺขเว ภิกฺขุ อคารโว อปฺปติสฺโส อภพฺโพ อิมสฺมิํ ธมฺมวินเย วุทฺธิํ วิรูฬฺหิํ เวปุลฺลํ อาปชฺชิตุํ, กุสีโต ภิกฺขเว ภิกฺขุ อคารโว อปฺปติสฺโส อภพฺโพ อิมสฺมิํ ธมฺมวินเย วุทฺธิํ วิรูฬฺหิํ เวปุลฺลํ อาปชฺชิตุํ, ทุปฺปญฺโญ ภิกฺขเว ภิกฺขุ อคารโว อปฺปติสฺโส อภพฺโพ อิมสฺมิํ ธมฺมวินเย วุทฺธิํ วิรูฬฺหิํ เวปุลฺลํ อาปชฺชิตุํ.\csromanspacer{} อิเมหิ โข ภิกฺขเว ปญฺจหิ ธมฺเมหิ สมนฺนาคโต ภิกฺขุ อคารโว อปฺปติสฺโส อภพฺโพ อิมสฺมิํ ธมฺมวินเย วุทฺธิํ วิรูฬฺหิํ เวปุลฺลํ อาปชฺชิตุํ.}

\prose{ปญฺจหิ ภิกฺขเว ธมฺเมหิ สมนฺนาคโต ภิกฺขุ สคารโว สปฺปติสฺโส ภพฺโพ อิมสฺมิํ ธมฺมวินเย วุทฺธิํ วิรูฬฺหิํ เวปุลฺลํ อาปชฺชิตุํ.\csromanspacer{} กตเมหิ ปญฺจหิ? สทฺโธ ภิกฺขเว ภิกฺขุ สคารโว สปฺปติสฺโส ภพฺโพ อิมสฺมิํ \csromanfolio{8}ธมฺมวินเย วุทฺธิํ วิรูฬฺหิํ เวปุลฺลํ อาปชฺชิตุํ, หิรีมา ภิกฺขเว ภิกฺขุ ฯเปฯ โอตฺตปฺปี ภิกฺขเว ภิกฺขุ ฯเปฯ อารทฺธวีริโย ภิกฺขเว ภิกฺขุ ฯเปฯ ปญฺญวา ภิกฺขเว ภิกฺขุ สคารโว สปฺปติสฺโส ภพฺโพ อิมสฺมิํ ธมฺมวินเย วุทฺธิํ วิรูฬฺหิํ เวปุลฺลํ อาปชฺชิตุํ.\csromanspacer{} อิเมหิ โข ภิกฺขเว ปญฺจหิ ธมฺเมหิ สมนฺนาคโต ภิกฺขุ สคารโว สปฺปติสฺโส ภพฺโพ อิมสฺมิํ ธมฺมวินเย วุทฺธิํ วิรูฬฺหิํ เวปุลฺลํ อาปชฺชิตุนฺติ.\csromanspacer{}\csromanspacer{}ทสมํ.\csromanspacer{} เสขพลวคฺโค ปฏฺฐโม.}

\titlehead{\textbf{ตสฺสุทฺทานํ}}
\csromangathameasure{สํขิตฺตํ วิตฺถตํ ทุกฺขา, ภตํ สิกฺขาย ปญฺจมํ. \\ สมาปตฺติ จ กาเมสุ, จวนา เทฺว อคารวาติ.}
\csromangathagroup{\csromangathastanza{สํขิตฺตํ วิตฺถตํ ทุกฺขา, ภตํ สิกฺขาย ปญฺจมํ. \\ สมาปตฺติ จ กาเมสุ, จวนา เทฺว อคารวาติ.}}

\csromanheader{2}{1. อนนุสฺสุตสุตฺต}
\proseitem{11}{ปุพฺพาหํ ภิกฺขเว อนนุสฺสุเตสุ ธมฺเมสุ อภิญฺญาโวสานปารมิปฺปตฺโต ปฏิชานามิ.\csromanspacer{} ปญฺจิมานิ ภิกฺขเว ตถาคตสฺส ตถาคตพลานิ, เยหิ พเลหิ สมนฺนาคโต ตถาคโต อาสภํ ฐานํ ปฏิชานาติ, ปริสาสุ สีหนาทํ นทติ, พฺรหฺมจกฺกํ ปวตฺเตติ.\csromanspacer{} กตมานิ ปญฺจ? สทฺธาพลํ หิรีพลํ โอตฺตปฺปพลํ วีริยพลํ ปญฺญาพลํ.\csromanspacer{} อิมานิ โข ภิกฺขเว ปญฺจ ตถาคตสฺส ตถาคตพลานิ, เยหิ พเลหิ สมนฺนาคโต ตถาคโต อาสภํ ฐานํ ปฏิชานาติ, ปริสาสุ สีหนาทํ นทติ, พฺรหฺมจกฺกํ ปวตฺเตตีติ.\csromanspacer{}\csromanspacer{}ปฐมํ.}

\csromanheader{2}{2. กูฏสุตฺต}
\proseitem{12}{ปญฺจิมานิ ภิกฺขเว เสขพลานิ.\csromanspacer{} กตมานิ ปญฺจ? สทฺธาพลํ หิรีพลํ โอตฺตปฺปพลํ วีริยพลํ ปญฺญาพลํ.\csromanspacer{} อิมานิ โข ภิกฺขเว ปญฺจ เสขพลานิ.\csromanspacer{} อิเมสํ โข ภิกฺขเว ปญฺจนฺนํ เสขพลานํ เอตํ อคฺคํ เอตํ สงฺคาหิกํ เอตํ สงฺฆาตนิยํ, ยทิทํ ปญฺญาพลํ.}

\prose{\csromanfolio{9}เสยฺยถาปิ ภิกฺขเว กูฏาคารสฺส เอตํ อคฺคํ เอตํ สงฺคาหิกํ เอตํ สงฺฆาตนิยํ, ยทิทํ กูฏํ.\csromanspacer{} เอวเมวํ โข ภิกฺขเว อิเมสํ ปญฺจนฺนํ เสขพลานํ เอตํ อคฺคํ เอตํ สงฺคาหิกํ เอตํ สงฺฆาตนิยํ, ยทิทํ ปญฺญาพลํ.}

\prose{ตสฺมาติห ภิกฺขเว เอวํ สิกฺขิตพฺพํ “สทฺธาพเลน สมนฺนาคตา ภวิสฺสาม เสขพเลน, หิรีพเลน.\csromanspacer{} โอตฺตปฺปพเลน.\csromanspacer{} วีริยพเลน.\csromanspacer{} ปญฺญาพเลน สมนฺนาคตา ภวิสฺสาม เสขพเลนา”ติ.\csromanspacer{} เอวํ หิ โว ภิกฺขเว สิกฺขิตพฺพนฺติ.\csromanspacer{}\csromanspacer{}ทุติยํ.}

\csromanheader{2}{3. สํขิตฺตสุตฺต}
\proseitem{13}{ปญฺจิมานิ ภิกฺขเว พลานิ.\csromanspacer{} กตมานิ ปญฺจ? สทฺธาพลํ วีริยพลํ สติพลํ สมาธิพลํ ปญฺญาพลํ.\csromanspacer{} อิมานิ โข ภิกฺขเว ปญฺจ พลานีติ.\csromanspacer{}\csromanspacer{}ตติยํ.}

\csromanheader{2}{4. วิตฺถตสุตฺต}
\proseitem{14}{ปญฺจิมานิ ภิกฺขเว พลานิ.\csromanspacer{} กตมานิ ปญฺจ? สทฺธาพลํ วีริยพลํ สติพลํ สมาธิพลํ ปญฺญาพลํ.}

\prose{กตมญฺจ ภิกฺขเว สทฺธาพลํ? อิธ ภิกฺขเว อริยสาวโก สทฺโธ โหติ, สทฺทหติ ตถาคตสฺส โพธิํ “อิติปิ โส ภควา อรหํ สมฺมาสมฺพุทฺโธ วิชฺชาจรณสมฺปนฺโน สุคโต โลกวิทู อนุตฺตโร ปุริสทมฺมสารถิ สตฺถา เทวมนุสฺสานํ พุทฺโธ ภควา”ติ.\csromanspacer{} อิทํ วุจฺจติ ภิกฺขเว สทฺธาพลํ.}

\prose{กตมญฺจ ภิกฺขเว วีริยพลํ? อิธ ภิกฺขเว อริยสาวโก อารทฺธวีริโย วิหรติ อกุสลานํ ธมฺมานํ ปหานาย กุสลานํ ธมฺมานํ อุปสมฺปทาย, ถามวา ทฬฺหปรกฺกโม อนิกฺขิตฺตธุโร กุสเลสุ ธมฺเมสุ.\csromanspacer{} อิทํ วุจฺจติ ภิกฺขเว วีริยพลํ.}

\prose{กตมญฺจ ภิกฺขเว สติพลํ? อิธ ภิกฺขเว อริยสาวโก สติมา โหติ ปรเมน สติเนปกฺเกน สมนฺนาคโต จิรกตมฺปิ จิรภาสิตมฺปิ สริตา อนุสฺสริตา.\csromanspacer{} อิทํ วุจฺจติ ภิกฺขเว สติพลํ.}

\prose{กตมญฺจ ภิกฺขเว สมาธิพลํ? อิธ ภิกฺขเว อริยสาวโก วิวิจฺเจว กาเมหิ วิวิจฺจ อกุสเลหิ ธมฺเมหิ สวิตกฺกํ สวิจารํ วิเวกชํ \csromanfolio{10}ปีติสุขํ ปฐมํ ฌานํ อุปสมฺปชฺช วิหรติ.\csromanspacer{} วิตกฺกวิจารานํ วูปสมา อชฺฌตฺตํ สมฺปสาทนํ เจตโส เอโกทิภาวํ อวิตกฺกํ อวิจารํ สมาธิชํ ปีติสุขํ ทุติยํ ฌานํ อุปสมฺปชฺช วิหรติ.\csromanspacer{} ปีติยา จ วิราคา อุเปกฺขโก จ วิหรติ, สโต จ สมฺปชาโน, สุขญฺจ กาเยน ปฏิสํเวเทติ, ยํ ตํ อริยา อาจิกฺขนฺติ “อุเปกฺขโก สติมา สุขวิหารี”ติ, ตติยํ ฌานํ อุปสมฺปชฺช วิหรติ.\csromanspacer{} สุขสฺส จ ปหานา ทุกฺขสฺส จ ปหานา ปุพฺเพว โสมนสฺสโทมนสฺสานํ อตฺถงฺคมา อทุกฺขมสุขํ อุเปกฺขาสติปาริสุทฺธิํ จตุตฺถํ ฌานํ อุปสมฺปชฺช วิหรติ.\csromanspacer{} อิทํ วุจฺจติ ภิกฺขเว สมาธิพลํ.}

\prose{กตมญฺจ ภิกฺขเว ปญฺญาพลํ? อิธ ภิกฺขเว อริยสาวโก ปญฺญวา โหติ อุทยตฺถคามินิยา ปญฺญาย สมนฺนาคโต อริยาย นิพฺเพธิกาย สมฺมา ทุกฺขกฺขย คามินิยา.\csromanspacer{} อิทํ วุจฺจติ ภิกฺขเว ปญฺญาพลํ.\csromanspacer{} อิมานิ โข ภิกฺขเว ปญฺจ พลานีติ.\csromanspacer{}\csromanspacer{}จตุตฺถํ.}

\csromanheader{2}{5. ทฏฺฐพฺพสุตฺต}
\proseitem{15}{ปญฺจิมานิ ภิกฺขเว พลานิ.\csromanspacer{} กตมานิ ปญฺจ? สทฺธาพลํ วีริยพลํ สติพลํ สมาธิพลํ ปญฺญาพลํ.\csromanspacer{} กตฺถ จ ภิกฺขเว สทฺธาพลํ ทฏฺฐพฺพํ? จตูสุ โสตาปตฺติยงฺเคสุ, เอตฺถ สทฺธาพลํ ทฏฺฐพฺพํ.\csromanspacer{} กตฺถ จ ภิกฺขเว วีริยพลํ ทฏฺฐพฺพํ? จตูสุ สมฺมปฺปธาเนสุ, เอตฺถ วีริยพลํ ทฏฺฐพฺพํ.\csromanspacer{} กตฺถ จ ภิกฺขเว สติพลํ ทฏฺฐพฺพํ? จตูสุ สติปฏฺฐาเนสุ, เอตฺถ สติพลํ ทฏฺฐพฺพํ.\csromanspacer{} กตฺถ จ ภิกฺขเว สมาธิพลํ ทฏฺฐพฺพํ? จตูสุ ฌาเนสุ, เอตฺถ สมาธิพลํ ทฏฺฐพฺพํ.\csromanspacer{} กตฺถ จ ภิกฺขเว ปญฺญาพลํ ทฏฺฐพฺพํ? จตูสุ อริยสจฺเจสุ, เอตฺถ ปญฺญาพลํ ทฏฺฐพฺพํ.\csromanspacer{} อิมานิ โข ภิกฺขเว ปญฺจ พลานีติ.\csromanspacer{}\csromanspacer{}ปญฺจมํ.}

\csromanheader{2}{6. ปุนกูฏสุตฺต}
\proseitem{16}{ปญฺจิมานิ ภิกฺขเว พลานิ.\csromanspacer{} กตมานิ ปญฺจ? สทฺธาพลํ วีริยพลํ สติพลํ สมาธิพลํ ปญฺญาพลํ.\csromanspacer{} อิมานิ โข ภิกฺขเว ปญฺจ พลานิ.\csromanspacer{} อิเมสํ โข ภิกฺขเว ปญฺจนฺนํ พลานํ เอตํ อคฺคํ เอตํ สงฺคาหิกํ เอตํ สงฺฆาตนิยํ, ยทิทํ ปญฺญาพลํ.\csromanspacer{} เสยฺยถาปิ ภิกฺขเว กูฏาคารสฺส เอตํ อคฺคํ เอตํ สงฺคาหิกํ เอตํ สงฺฆาตนิยํ, ยทิทํ กูฏํ.\csromanspacer{} เอวเมวํ โข ภิกฺขเว อิเมสํ ปญฺจนฺนํ พลานํ เอตํ อคฺคํ เอตํ สงฺคาหิกํ เอตํ สงฺฆาตนิยํ, ยทิทํ ปญฺญาพลนฺติ.\csromanspacer{}\csromanspacer{}ฉฏฺฐํ.}

\csromanheader{2}{7. ปฐมหิตสุตฺต}
\proseitem{17}{\csromanfolio{11}ปญฺจหิ ภิกฺขเว ธมฺเมหิ สมนฺนาคโต ภิกฺขุ อตฺตหิตาย ปฏิปนฺโน โหติ โน ปรหิตาย.\csromanspacer{} กตเมหิ ปญฺจหิ? อิธ ภิกฺขเว ภิกฺขุ อตฺตนา สีลสมฺปนฺโน โหติ, โน ปรํ สีลสมฺปทาย สมาทเปติ.\csromanspacer{} อตฺตนา สมาธิสมฺปนฺโน โหติ, โน ปรํ สมาธิสมฺปทาย สมาทเปติ.\csromanspacer{} อตฺตนา ปญฺญาสมฺปนฺโน โหติ, โน ปรํ ปญฺญาสมฺป ทาย สมาทเปติ.\csromanspacer{} อตฺตนา วิมุตฺติสมฺปนฺโน โหติ, โน ปรํ วิมุตฺติสมฺปทาย สมาท เปติ.\csromanspacer{} อตฺตนา วิมุตฺติญาณทสฺสนสมฺปนฺโน โหติ, โน ปรํ วิมุตฺติญาณทสฺสน สมฺปทาย สมาทเปติ.\csromanspacer{} อิเมหิ โข ภิกฺขเว ปญฺจหิ องฺเคหิ สมนฺนาคโต ภิกฺขุ อตฺตหิตาย ปฏิปนฺโน โหติ โน ปรหิตายาติ.\csromanspacer{}\csromanspacer{}สตฺตมํ.}

\csromanheader{2}{8. ทุติยหิตสุตฺต}
\proseitem{18}{ปญฺจหิ ภิกฺขเว ธมฺเมหิ สมนฺนาคโต ภิกฺขุ ปรหิตาย ปฏิปนฺโน โหติ โน อตฺตหิตาย.\csromanspacer{} กตเมหิ ปญฺจหิ? อิธ ภิกฺขเว ภิกฺขุ อตฺตนา น สีลสมฺปนฺโน โหติ, ปรํ สีลสมฺปทาย สมาทเปติ.\csromanspacer{} อตฺตนา น สมาธิสมฺปนฺโน โหติ, ปรํ สมาธิสมฺปทาย สมาทเปติ.\csromanspacer{} อตฺตนา น ปญฺญาสมฺปนฺโน โหติ, ปรํ ปญฺญาสมฺปทาย สมาทเปติ.\csromanspacer{} อตฺตนา น วิมุตฺติสมฺปนฺโน โหติ, ปรํ วิมุตฺติสมฺปทาย สมาทเปติ.\csromanspacer{} อตฺตนา น วิมุตฺติญาณทสฺสนสมฺปนฺโน โหติ, ปรํ วิมุตฺติญาณทสฺสนสมฺปทาย สมาทเปติ.\csromanspacer{} อิเมหิ โข ภิกฺขเว ปญฺจหิ ธมฺเมหิ สมนฺนาคโต ภิกฺขุ ปรหิตาย ปฏิปนฺโน โหติ โน อตฺตหิตายาติ.\csromanspacer{}\csromanspacer{}อฏฺฐมํ.}

\csromanheader{2}{9. ตติยหิตสุตฺต}
\proseitem{19}{ปญฺจหิ ภิกฺขเว ธมฺเมหิ สมนฺนาคโต ภิกฺขุ เนว อตฺตหิตาย ปฏิปนฺโน โหติ โน ปรหิตาย.\csromanspacer{} กตเมหิ ปญฺจหิ? อิธ ภิกฺขเว ภิกฺขุ อตฺตนา น สีลสมฺปนฺโน โหติ, โน ปรํ สีลสมฺปทาย สมาทเปติ.\csromanspacer{} อตฺตนา น สมาธิสมฺปนฺโน โหติ, โน ปรํ สมาธิสมฺปทาย สมาทเปติ.\csromanspacer{} อตฺตนา น ปญฺญาสมฺปนฺโน โหติ, โน ปรํ ปญฺญาสมฺปทาย สมาทเปติ.\csromanspacer{} อตฺตนา น วิมุตฺติสมฺปนฺโน โหติ, โน ปรํ วิมุตฺติสมฺปทาย สมาทเปติ.\csromanspacer{} อตฺตนา น วิมุตฺติญาณทสฺสนสมฺปนฺโน โหติ, โน ปรํ วิมุตฺติญาณทสฺสนสมฺปทาย สมาทเปติ.\csromanspacer{} อิเมหิ \csromanfolio{12}โข ภิกฺขเว ปญฺจหิ ธมฺเมหิ สมนฺนาคโต ภิกฺขุ เนว อตฺตหิตาย ปฏิปนฺโน โหติ โน ปรหิตายาติ.\csromanspacer{}\csromanspacer{}นวมํ.}

\csromanheader{2}{10. จตุตฺถหิตสุตฺต}
\proseitem{20}{ปญฺจหิ ภิกฺขเว ธมฺเมหิ สมนฺนาคโต ภิกฺขุ อตฺตหิตาย จ ปฏิปนฺโน โหติ ปรหิตาย จ.\csromanspacer{} กตเมหิ ปญฺจหิ? อิธ ภิกฺขเว ภิกฺขุ อตฺตนา จ สีลสมฺปนฺโน โหติ, ปรญฺจ สีลสมฺปทาย สมาทเปติ.\csromanspacer{} อตฺตนา จ สมาธิสมฺปนฺโน โหติ, ปรญฺจ สมาธิสมฺปทาย สมาทเปติ.\csromanspacer{} อตฺตนา จ ปญฺญาสมฺปนฺโน โหติ, ปรญฺจ ปญฺญาสมฺปทาย สมาทเปติ.\csromanspacer{} อตฺตนา จ วิมุตฺติสมฺปนฺโน โหติ, ปรญฺจ วิมุตฺติสมฺปทาย สมาทเปติ.\csromanspacer{} อตฺตนา จ วิมุตฺติญาณทสฺสนสมฺปนฺโน โหติ, ปรญฺจ วิมุตฺติญาณทสฺสนสมฺปทาย สมาทเปติ.\csromanspacer{} อิเมหิ โข ภิกฺขเว ปญฺจหิ ธมฺเมหิ สมนฺนาคโต ภิกฺขุ อตฺตหิตาย จ ปฏิปนฺโน โหติ ปรหิตาย จาติ.\csromanspacer{}\csromanspacer{}ทสมํ.\csromanspacer{} พลวคฺโค ทุติโย.}

\titlehead{\textbf{ตสฺสุทฺทานํ}}
\csromangathameasure{อนนุสฺสุตกูฏญฺจ, สํขิตฺตํ วิตฺถเตน จ. \\ ทฏฺฐพฺพญฺจ ปุน กูฏํ, จตฺตาโรปิ หิเตน จาติ.}
\csromangathagroup{\csromangathastanza{อนนุสฺสุตกูฏญฺจ, สํขิตฺตํ วิตฺถเตน จ. \\ ทฏฺฐพฺพญฺจ ปุน กูฏํ, จตฺตาโรปิ หิเตน จาติ.}}

\chapterhead{3. ปญฺจงฺคิกวคฺค}
\csromanheader{2}{1. ปฐม อคารวสุตฺต}
\proseitem{21}{โส วต ภิกฺขเว ภิกฺขุ อคารโว อปฺปติสฺโส อสภาควุตฺติโก สพฺรหฺมจารีสุ อาภิสมาจาริกํ ธมฺมํ ปริปูเรสฺสตีติ เนตํ ฐานํ วิชฺชติ, อาภิสมาจาริกํ ธมฺมํ อปริปูเรตฺวา เสขํ\footnote{เสกฺขํ (ก)} ธมฺมํ ปริปูเรสฺสตีติ เนตํ ฐานํ วิชฺชติ, เสขํ ธมฺมํ อปริปูเรตฺวา สีลานิ ปริปูเรสฺสตีติ เนตํ ฐานํ วิชฺชติ, สีลานิ อปริปูเรตฺวา สมฺมาทิฏฺฐิํ ปริปูเรสฺสตีติ เนตํ ฐานํ วิชฺชติ, สมฺมาทิฏฺฐิํ อปริปูเรตฺวา สมฺมาสมาธิํ ปริปูเรสฺสตีติ เนตํ ฐานํ วิชฺชติ.}

\chapterheadpageread{3. ปญฺจงฺคิกวคฺค}
\prose{\csromanfolio{13}โส วต ภิกฺขเว ภิกฺขุ สคารโว สปฺปติสฺโส สภาควุตฺติโก สพฺรหฺมจารีสุ อาภิสมาจาริกํ ธมฺมํ ปริปูเรสฺสตีติ ฐานเมตํ วิชฺชติ, อาภิสมาจาริกํ ธมฺมํ ปริปูเรตฺวา เสขํ ธมฺมํ ปริปูเรสฺสตีติ ฐานเมตํ วิชฺชติ, เสขํ ธมฺมํ ปริปูเรตฺวา สีลานิ ปริปูเรสฺสตีติ ฐานเมตํ วิชฺชติ, สีลานิ ปริปูเรตฺวา สมฺมาทิฏฺฐิํ ปริปูเรสฺสตีติ ฐานเมตํ วิชฺชติ, สมฺมาทิฏฺฐิํ ปริปูเรตฺวา สมฺมาสมาธิํ ปริปูเรสฺสตีติ ฐานเมตํ วิชฺชตีติ.\csromanspacer{}\csromanspacer{}ปฐมํ.}

\csromanheader{2}{2. ทุติย อคารวสุตฺต}
\proseitem{22}{โส วต ภิกฺขเว ภิกฺขุ อคารโว อปฺปติสฺโส อสภาควุตฺติโก สพฺรหฺมจารีสุ อาภิสมาจาริกํ ธมฺมํ ปริปูเรสฺสตีติ เนตํ ฐานํ วิชฺชติ, อาภิสมาจาริกํ ธมฺมํ อปริปูเรตฺวา เสขํ ธมฺมํ ปริปูเรสฺสตีติ เนตํ ฐานํ วิชฺชติ, เสขํ ธมฺมํ อปริปูเรตฺวา สีลกฺขนฺธํ ปริปูเรสฺสตีติ เนตํ ฐานํ วิชฺชติ, สีลกฺขนฺธํ อปริปูเรตฺวา สมาธิกฺขนฺธํ ปริปูเรสฺสตีติ เนตํ ฐานํ วิชฺชติ, สมาธิกฺขนฺธํ อปริปูเรตฺวา ปญฺญากฺขนฺธํ ปริปูเรสฺสตีติ เนตํ ฐานํ วิชฺชติ.}

\prose{โส วต ภิกฺขเว ภิกฺขุ สคารโว สปฺปติสฺโส สภาควุตฺติโก สพฺรหฺมจารีสุ อาภิสมาจาริกํ ธมฺมํ ปริปูเรสฺสตีติ ฐานเมตํ วิชฺชติ, อาภิสมาจาริกํ ธมฺมํ ปริปูเรตฺวา เสขํ ธมฺมํ ปริปูเรสฺสตีติ ฐานเมตํ วิชฺชติ, เสขํ ธมฺมํ ปริปูเรตฺวา สีลกฺขนฺธํ ปริปูเรสฺสตีติ ฐานเมตํ วิชฺชติ, สีลกฺขนฺธํ ปริปูเรตฺวา สมาธิกฺขนฺธํ ปริปูเรสฺสตีติ ฐานเมตํ วิชฺชติ, สมาธิกฺขนฺธํ ปริปูเรตฺวา ปญฺญากฺขนฺธํ ปริปูเรสฺสตีติ ฐานเมตํ วิชฺชตีติ.\csromanspacer{}\csromanspacer{}ทุติยํ.}

\csromanheader{2}{3. อุปกฺกิเลสสุตฺต}
\proseitem{23}{ปญฺจิเม ภิกฺขเว ชาตรูปสฺส อุปกฺกิเลสา, เยหิ อุกฺกิเลเสหิ อุปกฺกิลิฏฺฐํ ชาตรูปํ น เจว มุทุ โหติ น จ กมฺมนิยํ น จ ปภสฺสรํ, ปภงฺคุ จ, น จ สมฺมา อุเปติ กมฺมาย.\csromanspacer{} กตเม ปญฺจ? อโย โลหํ ติปุ สีสํ สชฺฌํ\footnote{สชฺฌุ (สี)}.\csromanspacer{} อิเม โข ภิกฺขเว ปญฺจ ชาตรูปสฺส อุปกฺกิ \csromanfolio{14}เลสา, เยหิ อุปกฺกิเลเสหิ อุปกฺกิลิฏฺฐํ ชาตรูปํ น เจว มุทุ โหติ น จ กมฺมนิยํ น จ ปภสฺสรํ, ปภงฺคุ จ, น จ สมฺมา อุเปติ กมฺมาย.\csromanspacer{} ยโต จ โข ภิกฺขเว ชาตรูปํ อิเมหิ ปญฺจหิ อุปกฺกิเลเสหิ วิมุตฺตํ\footnote{วิปฺปมุตฺตํ (สี)} โหติ, ตํ โหติ ชาตรูปํ มุทุ จ กมฺมนิยญฺจ ปภสฺสรญฺจ, น จ ปภงฺคุ, สมฺมา อุเปติ กมฺมาย, ยสฺสา ยสฺสา จ\footnote{ยสฺส กสฺสจิ (สฺยา, อิ)} ปิฬนฺธนวิกติยา อากงฺขติ, ยทิ มุทฺทิกาย ยทิ กุณฺฑลาย ยทิ คีเวยฺยกาย\footnote{คีเวยฺยเกน (สฺยา, กํ, ก; อํ 1. 255 ปิฏฺเฐปิ.)} ยทิ สุวณฺณมาลาย, ตญฺจสฺส อตฺถํ อนุโภติ.}

\prose{เอวเมวํ โข ภิกฺขเว ปญฺจิเม จิตฺตสฺส อุปกฺกิเลสา, เยหิ อุปกฺกิเลเสหิ อุปกฺกิลิฏฺฐํ จิตฺตํ น เจว มุทุ โหติ น จ กมฺมนิยํ น จ ปภสฺสรํ, ปภงฺคุ จ, น จ สมฺมา สมาธิยติ อาสวานํ ขยาย.\csromanspacer{} กตเม ปญฺจ? กามจฺฉนฺโท พฺยาปาโท ถินมิทฺธํ\footnote{ถีนมิทฺธํ (สี, สฺยา, กํ, อิ)} อุทฺธจฺจกุกฺกุจฺจํ วิจิกิจฺฉา.\csromanspacer{} อิเม โข ภิกฺขเว ปญฺจ จิตฺตสฺส อุปกฺกิเลสา เยหิ อุปกฺกิเลเสหิ อุปกฺกิลิฏฺฐํ จิตฺตํ น เจว มุทุ โหติ น จ กมฺมนิยํ น จ ปภสฺสรํ, ปภงฺคุ จ, น จ สมฺมา สมาธิยติ อาสวานํ ขยาย.\csromanspacer{} ยโต จ โข ภิกฺขเว จิตฺตํ อิเมหิ ปญฺจหิ อุปกฺกิเลเสหิ วิมุตฺตํ โหติ, ตํ โหติ จิตฺตํ มุทุ จ กมฺมนิยญฺจ ปภสฺสรญฺจ, น จ ปภงฺคุ, สมฺมา สมาธิยติ อาสวานํ ขยาย, ยสฺส ยสฺส จ อภิญฺญาสจฺฉิ กรณียสฺส ธมฺมสฺส จิตฺตํ อภินินฺนาเมติ อภิญฺญาสจฺฉิกิริยาย.\csromanspacer{} ตตฺร ตเตฺรว สกฺขิภพฺพตํ ปาปุณาติ สติ สติ อายตเน.}

\prose{โส สเจ อากงฺขติ “อเนกวิหิตํ อิทฺธิวิธํ ปจฺจนุภเวยฺยํ, เอโกปิ หุตฺวา พหุธา อสฺสํ, พหุธาปิ หุตฺวา เอโก อสฺสํ, อาวิภาวํ ติโรภาวํ ติโรกุฏฺฏํ ติโรปาการํ ติโรปพฺพตํ อสชฺชมาโน คจฺเฉยฺยํ เสยฺยถาปิ อากาเส, ปถวิยาปิ อุมฺมุชฺชนิมุชฺชํ กเรยฺยํ เสยฺยถาปิ อุทเก, อุทเกปิ อภิชฺชมาโน คจฺเฉยฺยํ เสยฺยถาปิ ปถวิยํ, อากาเสปิ ปลฺลงฺเกน กเมยฺยํ เสยฺยถาปิ ปกฺขี สกุโณ, อิเมปิ จนฺทิมสูริเย\footnote{จนฺทิมสุริเย (สี, สฺยา, กํ, อิ)} เอวํมหิทฺธิเก เอวํมหานุภาเว ปาณินา ปริมเสยฺยํ\footnote{ปรามเสยฺยํ (ก)} ปริมชฺเชยฺยํ, ยาว พฺรหฺมโลกาปิ กาเยน วสํ วตฺเตยฺยนฺ”ติ, ตตฺร ตเตฺรว สกฺขิภพฺพตํ ปาปุณาติ สติ สติ อายตเน.}

\chapterheadpageread{3. ปญฺจงฺคิกวคฺค}
\prose{\csromanfolio{15}โส สเจ อากงฺขติ “ทิพฺพาย โสตธาตุยา วิสุทฺธาย อติกฺกนฺตมานุสิกาย อุโภ สทฺเท สุเณยฺยํ ทิพฺเพ จ มานุเส จ เย ทูเร สนฺติเก จา”ติ, ตตฺร ตเตฺรว สกฺขิ ภพฺพตํ ปาปุณาติ สติ สติ อายตเน.}

\prose{โส สเจ อากงฺขติ “ปรสตฺตานํ ปรปุคฺคลานํ เจตสา เจโต ปริจฺจ ปชาเนยฺยํ, สราคํ วา จิตฺตํ ‘สราคํ จิตฺตนฺ’ติ ปชาเนยฺยํ, วีตราคํ วา จิตฺตํ ‘วีตราคํ จิตฺตนฺ’ติ ปชาเนยฺยํ, สโทสํ วา จิตฺตํ ‘สโทสํ จิตฺตนฺ’ติ ปชาเนยฺยํ, วีตโทสํ วา จิตฺตํ ‘วีตโทสํ จิตฺตนฺ’ติ ปชาเนยฺยํ, สโมหํ วา จิตฺตํ ‘สโมหํ จิตฺตนฺ’ติ ปชาเนยฺยํ, วีตโมหํ วา จิตฺตํ ‘วีตโมหํ จิตฺตนฺ’ติ ปชาเนยฺยํ, สํขิตฺตํ วา จิตฺตํ ‘สํขิตฺตํ จิตฺตนฺ’ติ ปชาเนยฺยํ, วิกฺขิตฺตํ วา จิตฺตํ ‘วิกฺขิตฺตํ จิตฺตนฺ’ติ ปชาเนยฺยํ, มหคฺคตํ วา จิตฺตํ ‘มหคฺคตํ จิตฺตนฺ’ติ ปชาเนยฺยํ, อมหคฺคตํ วา จิตฺตํ ‘อมหคฺคตํ จิตฺตนฺ’ติ ปชาเนยฺยํ, ส-อุตฺตรํ วา จิตฺตํ ‘ส-อุตฺตรํ จิตฺตนฺ’ติ ปชาเนยฺยํ, อนุตฺตรํ วา จิตฺตํ ‘อนุตฺตรํ จิตฺตนฺ’ติ ปชาเนยฺยํ, สมาหิตํ วา จิตฺตํ ‘สมาหิตํ จิตฺตนฺ’ติ ปชาเนยฺยํ, อสมาหิตํ วา จิตฺตํ ‘อสมาหิตํ จิตฺตนฺ’ติ ปชาเนยฺยํ, วิมุตฺตํ วา จิตฺตํ ‘วิมุตฺตํ จิตฺตนฺ’ติ ปชาเนยฺยํ, อวิมุตฺตํ วา จิตฺตํ ‘อวิมุตฺตํ จิตฺตนฺ’ติ ปชาเนยฺยนฺ’ติ, ตตฺร ตเตฺรว สกฺขิภพฺพตํ ปาปุณาติ สติ สติ อายตเน.}

\prose{โส สเจ อากงฺขติ “อเนกวิหิตํ ปุพฺเพนิวาสํ อนุสฺสเรยฺยํ.\csromanspacer{} เสยฺยถิทํ\footnote{เสยฺยถีทํ (สี, สฺยา, กํ, อิ)}? เอกมฺปิ ชาติํ เทฺวปิ ชาติโย ติสฺโสปิ ชาติโย จตสฺโสปิ ชาติโย ปญฺจปิ ชาติโย ทสปิ ชาติโย วีสมฺปิ ชาติโย ติํสมฺปิ ชาติโย จตฺตารีสมฺปิ ชาติโย ปญฺญาสมฺปิ ชาติโย ชาติสตมฺปิ ชาติสหสฺสมฺปิ ชาติสตสหสฺสมฺปิ อเนเกปิ สํวฏฺฏกปฺเป อเนเกปิ วิวฏฺฏกปฺเป อเนเกปิ สํวฏฺฏวิวฏฺฏกปฺเป ‘อมุตฺราสิํ เอวํนาโม เอวํโคตฺโต เอวํวณฺโณ เอวมาหาโร เอวํสุขทุกฺขปฺปฏิสํเวที เอวมายุปริยนฺโต, โส ตโต จุโต อมุตฺร อุทปาทิํ, ตตฺราปาสิํ เอวํนาโม เอวํโคตฺโต เอวํวณฺโณ เอวมาหาโร เอวํสุขทุกฺขปฺปฏิสํเวที เอวมายุปริยนฺโต, โส ตโต จุโต อิธูปปนฺโน’ติ.\csromanspacer{} อิติ สาการํ ส-อุทฺเทสํ อเนกวิหิตํ ปุพฺเพนิวาสํ อนุสฺสเรยฺยนฺ”ติ, ตตฺร ตเตฺรว สกฺขิภพฺพตํ ปาปุณาติ สติ สติ อายตเน.}

\prose{\csromanfolio{16}โส สเจ อากงฺขติ “ทิพฺเพน จกฺขุนา วิสุทฺเธน อติกฺกนฺตมานุสเกน สตฺเต ปสฺเสยฺยํ จวมาเน อุปปชฺชมาเน หีเน ปณีเต สุวณฺเณ ทุพฺพณฺเณ สุคเต ทุคฺคเต, ยถากมฺมูปเค สตฺเต ปชาเนยฺยํ ‘อิเม วต โภนฺโต สตฺตา กายทุจฺจริเตน สมนฺนาคตา วจีทุจฺจริเตน สมนฺนาคตา มโนทุจฺจริเตน สมนฺนาคตา อริยานํ อุปวาทกา มิจฺฉาทิฏฺฐิกา มิจฺฉาทิฏฺฐิกมฺมสมาทานา, เต กายสฺส เภทา ปรํ มรณา อปายํ ทุคฺคติํ วินิปาตํ นิรยํ อุปปนฺนา.\csromanspacer{} อิเม วา ปน โภนฺโต สตฺตา กายสุจริเตน สมนฺนาคตา วจีสุจริเตน สมนฺนาคตา มโนสุจริเตน สมนฺนาคตา อริยานํ อนุปวาทกา สมฺมาทิฏฺฐิกา สมฺมาทิฏฺฐิกมฺมสมาทานา, เต กายสฺส เภทา ปรํ มรณา สุคติํ สคฺคํ โลกํ อุปปนฺนา’ติ.\csromanspacer{} อิติ ทิพฺเพน จกฺขุนา วิสุทฺเธน อติกฺกนฺตมานุสเกน สตฺเต ปสฺเสยฺยํ จวมาเน อุปปชฺชมาเน หีเน ปณีเต สุวณฺเณ ทุพฺพณฺเณ สุคเต ทุคฺคเต, ยถากมฺมูปเค สตฺเต ปชาเนยฺยนฺ”ติ, ตตฺร ตเตฺรว สกฺขิภพฺพตํ ปาปุณาติ สติ สติ อายตเน.}

\prose{โส สเจ อากงฺขติ “อาสวานํ ขยา อนาสวํ เจโตวิมุตฺติํ ปญฺญาวิมุตฺติํ ทิฏฺเฐว ธมฺเม สยํ อภิญฺญา สจฺฉิกตฺวา อุปสมฺปชฺชวิหเรยฺยนฺ”ติ, ตตฺร ตเตฺรว สกฺขิภพฺพตํ ปาปุณาติ สติ สติ อายตเนติ.\csromanspacer{}\csromanspacer{}ตติยํ.}

\csromanheader{2}{4. ทุสฺสีลสุตฺต}
\proseitem{24}{\csromansymbolfootnote{*}{อุปริ 175, 316, 472 ปิฏฺเฐปิ.}ทุสฺสีลสฺส ภิกฺขเว สีลวิปนฺนสฺส หตูปนิโส โหติ สมฺมาสมาธิ, สมฺมาสมาธิมฺหิ อสติ สมฺมาสมาธิวิปนฺนสฺส หตูปนิสํ โหติ ยถาภูตญาณทสฺสนํ, ยถาภูตญาณทสฺสเน อสติ ยถาภูตญาณทสฺสนวิปนฺนสฺส หตูปนิโส โหติ นิพฺพิทาวิราโค, นิพฺพิทาวิราเค อสติ นิพฺพิทาวิราควิปนฺนสฺส หตูปนิสํ โหติ วิมุตฺติญาณทสฺสนํ.\csromanspacer{} เสยฺยถาปิ ภิกฺขเว รุกฺโข สาขาปลาสวิปนฺโน, ตสฺส ปปฏิกาปิ น ปาริปูริํ คจฺฉติ, ตโจปิ น ปาริปูริํ คจฺฉติ, เผคฺคุปิ น ปาริปูริํ คจฺฉติ, สาโรปิ น ปาริปูริํ คจฺฉติ.\csromanspacer{} เอวเมวํ โข ภิกฺขเว ทุสฺสีลสฺส สีลวิปนฺนสฺส หตูปนิโส โหติ สมฺมาสมาธิ, สมฺมาสมาธิมฺหิ อสติ สมฺมาสมาธิวิปนฺนสฺส หตูปนิสํ โหติ ยถาภูตญาณทสฺสนํ, ยถาภูตญาณทสฺสเน อสติ ยถาภูตญาณทสฺสน วิปนฺนสฺส หตูปนิโส}

\vspace{\tipitakaparskip}
\csromancenter{ปญฺจงฺคิกวคฺค โหติ นิพฺพิทาวิราโค, นิพฺพิทาวิราเค อสติ นิพฺพิทาวิราควิปนฺนสฺส หตูปนิสํ โหติ วิมุตฺติญาณทสฺสนํ.}
\prose{\csromanfolio{17}สีลวโต ภิกฺขเว สีลสมฺปนฺนสฺส อุปนิสสมฺปนฺโน โหติ สมฺมาสมาธิ, สมฺมาสมาธิมฺหิ สติ สมฺมาสมาธิสมฺปนฺนสฺส อุปนิสสมฺปนฺนํ โหติ ยถาภูตญาณทสฺสนํ, ยถาภูตญาณทสฺสเน สติ ยถาภูตญาณทสฺสนสมฺปนฺนสฺส อุปนิสสมฺปนฺโน โหติ นิพฺพิทาวิราโค, นิพฺพิทาวิราเค สติ นิพฺพิทาวิราคสมฺปนฺนสฺส อุปนิสสมฺปนฺนํ โหติ วิมุตฺติญาณทสฺสนํ.\csromanspacer{} เสยฺยถาปิ ภิกฺขเว รุกฺโข สาขาปลาสสมฺปนฺโน, ตสฺส ปปฏิกาปิ ปาริปูริํ คจฺฉติ, ตโจปิ ปาริปูริํ คจฺฉติ, เผคฺคุปิ ปาริปูริํ คจฺฉติ, สาโรปิ ปาริปูริํ คจฺฉติ.\csromanspacer{} เอวเมวํ โข ภิกฺขเว สีลวโต สีลสมฺปนฺนสฺส อุปนิสสมฺปนฺโน โหติ สมฺมาสมาธิ, สมฺมาสมาธิมฺหิ สติ สมฺมาสมาธิสมฺปนฺนสฺส อุปนิสสมฺปนฺนํ โหติ ยถาภูตญาณทสฺสนํ, ยถาภูตญาณทสฺสเน สติ ยถาภูตญาณทสฺสน สมฺปนฺนสฺส อุปนิสสมฺปนฺโน โหติ นิพฺพิทาวิราโค, นิพฺพิทาวิราเค สติ นิพฺพิทาวิราค สมฺปนฺนสฺส อุปนิสสมฺปนฺนํ โหติ วิมุตฺติญาณทสฺสนนฺติ.\csromanspacer{}\csromanspacer{}จตุตฺถํ.}

\csromanheader{2}{5. อนุคฺคหิตสุตฺต}
\proseitem{25}{ปญฺจหิ ภิกฺขเว องฺเคหิ อนุคฺคหิตา สมฺมาทิฏฺฐิ เจโตวิมุตฺติผลา จ โหติ เจโตวิมุตฺติผลานิสํสา จ, ปญฺญาวิมุตฺติผลา จ โหติ ปญฺญาวิมุตฺติผลานิสํสา จ.}

\prose{กตเมหิ ปญฺจหิ? อิธ ภิกฺขเว สมฺมาทิฏฺฐิ สีลานุคฺคหิตา จ โหติ, สุตานุคฺคหิตา จ โหติ, สากจฺฉานุคฺคหิตา จ โหติ, สมถานุคฺคหิตา จ โหติ, วิปสฺสนานุคฺคหิตา จ โหติ.\csromanspacer{} อิเมหิ โข ภิกฺขเว ปญฺจหิ องฺเคหิ อนุคฺคหิตา สมฺมาทิฏฺฐิ เจโตวิมุตฺติผลา จ โหติ เจโตวิมุตฺติผลานิสํสา จ, ปญฺญาวิมุตฺติผลา จ โหติ ปญฺญาวิมุตฺติผลานิสํสา จาติ.\csromanspacer{}\csromanspacer{}ปญฺจมํ.}

\csromanheader{2}{6. วิมุตฺตายตนสุตฺต}
\proseitem{26}{ปญฺจิมานิ ภิกฺขเว วิมุตฺตายตนานิ, ยตฺถ ภิกฺขุโน อปฺปมตฺตสฺส อาตาปิโน ปหิตตฺตสฺส วิหรโต อวิมุตฺตํ วา จิตฺตํ วิมุจฺจติ, \csromanfolio{18}อปริกฺขีณา วา อาสวา ปริกฺขยํ คจฺฉนฺติ, อนนุปฺปตฺตํ วา อนุตฺตรํ โยคกฺเขมํ อนุปาปุณาติ.}

\prose{กตมานิ ปญฺจ? อิธ ภิกฺขเว ภิกฺขุโน สตฺถา ธมฺมํ เทเสติ อญฺญตโร วา ครุฏฺฐานิโย สพฺรหฺมจารี, ยถา ยถา ภิกฺขเว ตสฺส ภิกฺขุโน สตฺถา ธมฺมํ เทเสติ อญฺญตโร วา ครุฏฺฐานิโย สพฺรหฺมจารี, ตถา ตถา โส ตสฺมิํ ธมฺเม อตฺถปฏิสํเวที จ โหติ ธมฺมปฏิสํเวที จ, ตสฺส อตฺถปฏิสํเวทิโน ธมฺมปฏิสํเวทิโน ปาโมชฺชํ\footnote{ปามุชฺชํ (สี, สฺยา, กํ) ที 3. 201 ปิฏฺเฐปิ ปสฺสิตพฺพํ.} ชายติ, ปมุทิตสฺส ปีติ ชายติ, ปีติมนสฺส กาโย ปสฺสมฺภติ, ปสฺสทฺธกาโย สุขํ เวเทติ\footnote{เวทิยติ (สี)} สุขิโน จิตฺตํ สมาธิยติ.\csromanspacer{} อิทํ ภิกฺขเว ปฐมํ วิมุตฺตายตนํ, ยตฺถ ภิกฺขุโน อปฺปมตฺตสฺส อาตาปิโน ปหิตตฺตสฺส วิหรโต อวิมุตฺตํ วา จิตฺตํ วิมุจฺจติ อปริกฺขีณา วา อาสวา ปริกฺขยํ คจฺฉนฺติ, อนนุปฺปตฺตํ วา อนุตฺตรํ โยคกฺเขมํ อนุปาปุณาติ.}

\prose{ปุน จปรํ ภิกฺขเว ภิกฺขุโน น เหว โข สตฺถา ธมฺมํ เทเสติ อญฺญตโร วา ครุฏฺฐานิโย สพฺรหฺมจารี, อปิ จ โข ยถาสุตํ ยถาปริยตฺตํ ธมฺมํ วิตฺถาเรน ปเรสํ เทเสติ.\csromanspacer{} ยถา ยถา ภิกฺขเว ภิกฺขุ ยถาสุตํ ยถาปริยตฺตํ ธมฺมํ วิตฺถาเรน ปเรสํ เทเสติ, ตถา ตถา โส ตสฺมิํ ธมฺเม อตฺถปฏิสํเวที จ โหติ ธมฺมปฏิสํเวที จ.\csromanspacer{} ตสฺส อตฺถปฏิสํเวทิโน ธมฺมปฏิสํเวทิโน ปาโมชฺชํ ชายติ, ปมุทิตสฺส ปีติ ชายติ, ปีติมนสฺส กาโย ปสฺสมฺภติ, ปสฺสทฺธกาโย สุขํ เวเทติ, สุขิโน จิตฺตํ สมาธิยติ.\csromanspacer{} อิทํ ภิกฺขเว ทุติยํ วิมุตฺตายตนํ, ยตฺถ ภิกฺขุโน อปฺปมตฺตสฺส อาตาปิโน ปหิตตฺตสฺส วิหรโต อวิมุตฺตํ วา จิตฺตํ วิมุจฺจติ, อปริกฺขีณา วา อาสวา ปริกฺขยํ คจฺฉนฺติ, อนนุปฺปตฺตํ วา อนุตฺตรํ โยคกฺเขมํ อนุปาปุณาติ.}

\prose{ปุน จปรํ ภิกฺขเว ภิกฺขุโน น เหว โข สตฺถา ธมฺมํ เทเสติ, อญฺญตโร วา ครุฏฺฐานิโย สพฺรหฺมจารี, นาปิ ยถาสุตํ ยถาปริยตฺตํ ธมฺมํ วิตฺถาเรน ปเรสํ เทเสติ, อปิ จ โข ยถาสุตํ ยถาปริยตฺตํ ธมฺมํ วิตฺถาเรน สชฺฌายํ กโรติ.\csromanspacer{} ยถา ยถา ภิกฺขเว ภิกฺขุ ยถาสุตํ ยถาปริยตฺตํ ธมฺมํ วิตฺถาเรน สชฺฌายํ กโรติ, ตถา ตถา โส ตสฺมิํ ธมฺเม อตฺถปฏิสํเวที จ โหติ ธมฺมปฏิสํเวที จ.\csromanspacer{} ตสฺส}

\vspace{\tipitakaparskip}
\csromancenter{ปญฺจงฺคิกวคฺค อตฺถปฏิสํเวทิโน ธมฺมปฏิสํเวทิโน ปาโมชฺชํ ชายติ, ปมุทิตสฺส ปีติ ชายติ, ปีติมนสฺส กาโย ปสฺสมฺภติ, ปสฺสทฺธกาโย สุขํ เวเทติ, สุขิโน จิตฺตํ สมาธิยติ.\csromanspacer{} อิทํ ภิกฺขเว ตติยํ วิมุตฺตายตนํ, ยตฺถ ภิกฺขุโน อปฺปมตฺตสฺส อาตาปิโน ฯเปฯ โยคกฺเขมํ อนุปาปุณาติ.}
\prose{\csromanfolio{19}ปุน จปรํ ภิกฺขเว ภิกฺขุโน น เหว โข สตฺถา ธมฺมํ เทเสติ, อญฺญตโร วา ครุฏฺฐานิโย สพฺรหฺมจารี, นาปิ ยถาสุตํ ยถาปริยตฺตํ ธมฺมํ วิตฺถาเรน ปเรสํ เทเสติ, นาปิ ยถาสุตํ ยถาปริยตฺตํ ธมฺมํ วิตฺถาเรน สชฺฌายํ กโรติ, อปิ จ โข ยถาสุตํ ยถาปริยตฺตํ ธมฺมํ เจตสา อนุวิตกฺเกติ อนุวิจาเรติ มนสานุเปกฺขติ.\csromanspacer{} ยถา ยถา ภิกฺขเว ภิกฺขุ ยถาสุตํ ยถาปริยตฺตํ ธมฺมํ เจตสา อนุวิตกฺเกติ อนุวิจาเรติ มนสานุเปกฺขติ, ตถา ตถา โส ตสฺมิํ ธมฺเม อตฺถปฏิสํเวที จ โหติ ธมฺมปฏิสํเวที จ.\csromanspacer{} ตสฺส อตฺถปฏิสํเวทิโน ธมฺมปฏิสํเวทิโน ปาโมชฺชํ ชายติ, ปมุทิตสฺส ปีติ ชายติ, ปีติมนสฺส กาโย ปสฺสมฺภติ, ปสฺสทฺธกาโย สุขํ เวเทติ, สุขิโน จิตฺตํ สมาธิยติ.\csromanspacer{} อิทํ ภิกฺขเว จตุตฺถํ วิมุตฺตายตนํ, ยตฺถ ภิกฺขุโน อปฺปมตฺตสฺส อาตาปิโน ปหิตตฺตสฺส วิหรโต อวิมุตฺตํ วา จิตฺตํ วิมุจฺจติ, อปริกฺขีณา วา อาสวา ปริกฺขยํ คจฺฉนฺติ, อนนุปฺปตฺตํ วา อนุตฺตรํ โยคกฺเขมํ อนุปาปุณาติ.}

\prose{ปุน จปรํ ภิกฺขเว ภิกฺขุโน น เหว โข สตฺถา ธมฺมํ เทเสติ, อญฺญตโร วา ครุฏฺฐานิโย สพฺรหฺมจารี, นาปิ ยถาสุตํ ยถาปริยตฺตํ ธมฺมํ วิตฺถาเรน ปเรสํ เทเสติ, นาปิ ยถาสุตํ ยถาปริยตฺตํ ธมฺมํ วิตฺถาเรน สชฺฌายํ กโรติ, นาปิ ยถาสุตํ ยถาปริยตฺตํ ธมฺมํ เจตสา อนุวิตกฺเกติ อนุวิจาเรติ มนสานุเปกฺขติ, อปิ จ ขฺวสฺส อญฺญตรํ สมาธินิมิตฺตํ สุคฺคหิตํ โหติ สุมนสิกตํ สูปธาริตํ สุปฺปฏิวิทฺธํ ปญฺญาย.\csromanspacer{} ยถา ยถา ภิกฺขเว ภิกฺขุโน อญฺญตรํ สมาธินิมิตฺตํ สุคฺคหิตํ โหติ สุมนสิกตํ สูปธาริตํ สุปฺปฏิวิทฺธํ ปญฺญาย, ตถา ตถา โส ตสฺมิํ ธมฺเม อตฺถปฏิสํเวที จ โหติ ธมฺมปฏิสํเวที จ.\csromanspacer{} ตสฺส อตฺถปฏิสํเวทิโน ธมฺมปฏิสํเวทิโน ปาโมชฺชํ ชายติ, ปมุทิตสฺส ปีติ ชายติ, ปีติมนสฺส กาโย ปสฺสมฺภติ, ปสฺสทฺธกาโย สุขํ เวเทติ, สุขิโน จิตฺตํ สมาธิยติ.\csromanspacer{} อิทํ ภิกฺขเว ปญฺจมํ วิมุตฺตายตนํ, ยตฺถ ภิกฺขุโน อปฺปมตฺตสฺส อาตาปิโน ปหิตตฺตสฺส วิหรโต อวิมุตฺตํ วา \csromanfolio{20}จิตฺตํ วิมุจฺจติ, อปริกฺขีณา วา อาสวา ปริกฺขยํ คจฺฉนฺติ, อนนุปฺปตฺตํ วา อนุตฺตรํ โยคกฺเขมํ อนุปาปุณาติ.}

\prose{อิมานิ โข ภิกฺขเว ปญฺจ วิมุตฺตายตนานิ, ยตฺถ ภิกฺขุโน อปฺปมตฺตสฺส อาตาปิโน ปหิตตฺตสฺส วิหรโต อวิมุตฺตํ วา จิตฺตํ วิมุจฺจติ, อปริกฺขีณา วา อาสวา ปริกฺขยํ คจฺฉนฺติ, อนนุปฺปตฺตํ วา อนุตฺตรํ โยคกฺเขมํ อนุปาปุณาตีติ.\csromanspacer{}\csromanspacer{}ฉฏฺฐํ.}

\csromanheader{2}{7. สมาธิสุตฺต}
\proseitem{27}{สมาธิํ ภิกฺขเว ภาเวถ อปฺปมาณํ นิปกา ปติสฺสตา, สมาธิํ ภิกฺขเว ภาวยตํ อปฺปมาณํ นิปกานํ ปติสฺสตานํ ปญฺจ ญาณานิ ปจฺจตฺตญฺเญว อุปฺปชฺชนฺติ.\csromanspacer{} กตมานิ ปญฺจ? “อยํ สมาธิ ปจฺจุปฺปนฺนสุโข เจว อายติํ จ สุขวิปาโก”ติ ปจฺจตฺตญฺเญว ญาณํ อุปฺปชฺชติ, “อยํ สมาธิ อริโย นิรามิโส”ติ ปจฺจตฺตญฺเญว ญาณํ อุปฺปชฺชติ, “อยํ สมาธิ อกาปุริสเสวิโต”ติ\footnote{มหาปุริสเสวิโตติ (ก)} ปจฺจตฺตญฺเญว ญาณํ อุปฺปชฺชติ, “อยํ สมาธิ สนฺโต ปณีโต ปฏิปฺปสฺสทฺธลทฺโธ เอโกทิภาวาธิคโต น สสงฺขารนิคฺคยฺหวาริตคโต”ติ\footnote{น จ สสงฺขารนิคฺคยฺหวาริตปฺปติโตติ (สี), น จ สสงฺขารนิคฺคยฺหวาริตปตฺโตติ (สฺยา), น จ สสงฺขารนิคฺคยฺหวาริวาวโฏติ (ก), น สสงฺขารนิคฺคยฺหวาริยาธิคโตติ (?) ที 3. 234; อํ 1. 256; อํ 3. 225 ปิฏฺเฐปิ.} ปจฺจตฺตญฺเญว ญาณํ อุปฺปชฺชติ, “สโต โข ปนาหํ อิมํ สมาปชฺชามิ สโต วุฏฺฐหามี”ติ\footnote{โส โข ปนาหํ อิมํ สมาธิํ สโตว สมาปชฺชามิ, สโต อุฏฺฐหามีติ (สี, สฺยา, กํ)} ปจฺจตฺตญฺเญว ญาณํ อุปฺปชฺชติ.}

\prose{สมาธิํ ภิกฺขเว ภาเวถ อปฺปมาณํ นิปกา ปติสฺสตา, สมาธิํ ภิกฺขเว ภาวยตํ อปฺปมาณํ นิปกานํ ปติสฺสตานํ อิมานิ ปญฺจ ญาณานิ ปจฺจตฺตญฺเญว อุปฺปชฺชนฺตีติ.\csromanspacer{}\csromanspacer{}สตฺตมํ.}

\csromanheader{2}{8. ปญฺจงฺคิกสุตฺต}
\proseitem{28}{อริยสฺส ภิกฺขเว ปญฺจงฺคิกสฺส สมฺมาสมาธิสฺส ภาวนํ เทเสสฺสามิ, ตํ สุณาถ สาธุกํ มนสิ กโรถ ภาสิสฺสามีติ. “เอวํ ภนฺเต”ติ โข เต ภิกฺขู ภควโต ปจฺจสฺโสสุํ.\csromanspacer{} ภควา เอตทโวจ\csromandash{}}

\chapterheadpageread{3. ปญฺจงฺคิกวคฺค}
\prose{\csromanfolio{21}กตมา จ ภิกฺขเว อริยสฺส ปญฺจงฺคิกสฺส สมฺมาสมาธิสฺส ภาวนา? อิธ ภิกฺขเว ภิกฺขุ วิวิจฺเจว กาเมหิ ฯเปฯ ปฐมํ ฌานํ อุปสมฺปชฺช วิหรติ, โส อิมเมว กายํ วิเวกเชน ปีติสุเขน อภิสนฺเทติ ปริสนฺเทติ ปริปูเรติ ปริปฺผรติ, นาสฺส กิญฺจิ สพฺพาวโต กายสฺส วิเวกเชน ปีติสุเขน อปฺผุฏํ โหติ.\csromanspacer{} เสยฺยถาปิ ภิกฺขเว ทกฺโข นฺหาปโก\footnote{นหาปโก (สี, อิ)} วา นฺหาปกนฺเตวาสี วา กํสถาเล นฺหานียจุณฺณานิ\footnote{นหานียจุณฺณานิ (สี, อิ)} อากิริตฺวา อุทเกน ปริปฺโผสกํ ปริปฺโผสกํ สนฺเนยฺย, สายํ นฺหานียปิณฺฑิ\footnote{สาสฺส นหานียปิณฺฑี (สี, สฺยา, กํ)} สฺเนหานุคตา สฺเนหปเรตา สนฺตรพาหิรา ผุฏา สฺเนเหน น จ ปคฺฆรินี.\csromanspacer{} เอวเมวํ โข ภิกฺขเว ภิกฺขุ อิมเมว กายํ วิเวกเชน ปีติสุเขน อภิสนฺเทติ ปริสนฺเทติ ปริปูเรติ ปริปฺผรติ, นาสฺส กิญฺจิ สพฺพาวโต กายสฺส วิเวกเชน ปีติสุเขน อปฺผุฏํ โหติ, อริยสฺส ภิกฺขเว ปญฺจงฺคิกสฺส สมฺมาสมาธิสฺส อยํ ปฐมา ภาวนา.}

\prose{ปุน จปรํ ภิกฺขเว ภิกฺขุ วิตกฺกวิจารานํ วูปสมา ฯเปฯ ทุติยํ ฌานํ อุปสมฺปชฺช วิหรติ, โส อิมเมว กายํ สมาธิเชน ปีติสุเขน อภิสนฺเทติ ปริสนฺเทติ ปริปูเรติ ปริปฺผรติ, นาสฺส กิญฺจิ สพฺพาวโต กายสฺส สมาธิเชน ปีติสุเขน อปฺผุฏํ โหติ.\csromanspacer{} เสยฺยถาปิ ภิกฺขเว อุทกรหโท คมฺภีโร อุพฺภิโททโก\footnote{อุพฺภิโตทโก (สฺยา, กํ, ก)}, ตสฺส เนวสฺส ปุรตฺถิมาย ทิสาย อุทกสฺส อายมุขํ, น ปจฺฉิมาย ทิสาย อุทกสฺส อายมุขํ, น อุตฺตราย ทิสาย อุทกสฺส อายมุขํ, น ทกฺขิณาย ทิสาย อุทกสฺส อายมุขํ, เทโว จ กาเลน กาลํ สมฺมา ธารํ นานุปฺปเวจฺเฉยฺย\footnote{เทโว จ น กาเลน ... อนุปเวจฺเฉยฺย (ที 1. 70; ม 2. 208 ปิฏฺเฐสุ)}, อถ โข ตมฺหาว อุทกรหทา สีตา วาริธารา อุพฺภิชฺชิตฺวา ตเมว อุทกรหทํ สีเตน วารินา อภิสนฺเทยฺย ปริสนฺเทยฺย ปริปูเรยฺย ปริปฺผเรยฺย, นาสฺส กิญฺจิ สพฺพาวโต อุทกรหทสฺส สีเตน วารินา อปฺผุฏํ อสฺส.\csromanspacer{} เอวเมวํ โข ภิกฺขเว ภิกฺขุ อิมเมว กายํ สมาธิเชน ปีติสุเขน อภิสนฺเทติ ปริสนฺเทติ ปริปูเรติ ปริปฺผรติ, นาสฺส กิญฺจิ สพฺพาวโต กายสฺส สมาธิเชน ปีติสุเขน อปฺผุฏํ โหติ.\csromanspacer{} อริยสฺส ภิกฺขเว ปญฺจงฺคิกสฺส สมฺมาสมาธิสฺส อยํ ทุติยา ภาวนา.}

\prose{\csromanfolio{22}ปุน จปรํ ภิกฺขเว ภิกฺขุ ปีติยา จ วิราคา ฯเปฯ ตติยํ ฌานํ อุปสมฺปชฺช วิหรติ, โส อิมเมว กายํ นิปฺปีติเกน สุเขน อภิสนฺเทติ ปริสนฺเทติ ปริปูเรติ ปริปฺผรติ, นาสฺส กิญฺจิ สพฺพาวโต กายสฺส นิปฺปีติเกน สุเขน อปฺผุฏํ โหติ.\csromanspacer{} เสยฺยถาปิ ภิกฺขเว อุปฺปลินิยํ วา ปทุมินิยํ วา ปุณฺฑรีกินิยํ วา อปฺเปกจฺจานิ อุปฺปลานิ วา ปทุมานิ วา ปุณฺฑรีกานิ วา อุทเก ชาตานิ อุทเก สํวฑฺฒานิ อุทกานุคฺคตานิ อนฺโต นิมุคฺคโปสีนิ, ตานิ ยาว จคฺคา ยาว จ มูลา สีเตน วารินา อภิสนฺนานิ ปริสนฺนานิ ปริปูรานิ ปริปฺผุฏานิ, นาสฺส กิญฺจิ สพฺพาวตํ อุปฺปลานํ วา ปทุมานํ วา ปุณฺฑรีกานํ วา สีเตน วารินา อปฺผุฏํ อสฺส.\csromanspacer{} เอวเมว โข ภิกฺขเว ภิกฺขุ อิมเมว กายํ นิปฺปีติเกน สุเขน อภิสนฺเทติ ปริสนฺเทติ ปริปูเรติ ปริปฺผรติ, นาสฺส กิญฺจิ สพฺพาวโต กายสฺส นิปฺปีติเกน สุเขน อปฺผุฏํ โหติ.\csromanspacer{} อริยสฺส ภิกฺขเว ปญฺจงฺคิกสฺส สมฺมาสมาธิสฺส อยํ ตติยา ภาวนา.}

\prose{ปุน จปรํ ภิกฺขเว ภิกฺขุ สุขสฺส จ ปหานา ฯเปฯ จตุตฺถํ ฌานํ อุปสมฺปชฺช วิหรติ, โส อิมเมว กายํ ปริสุทฺเธน เจตสา ปริโยทาเตน ผริตฺวา นิสินฺโน โหติ, นาสฺส กิญฺจิ สพฺพาวโต กายสฺส ปริสุทฺเธน เจตสา ปริโยทาเตน อปฺผุฏํ โหติ.\csromanspacer{} เสยฺยถาปิ ภิกฺขเว ปุริโส โอทาเตน วตฺเถน สสีสํ ปารุปิตฺวา นิสินฺโน อสฺส, นาสฺส กิญฺจิ สพฺพาวโต กายสฺส โอทาเตน วตฺเถน อปฺผุฏํ อสฺส.\csromanspacer{} เอวเมว โข ภิกฺขเว ภิกฺขุ อิมเมว กายํ ปริสุทฺเธน เจตสา ปริโยทาเตน ผริตฺวา นิสินฺโน โหติ, นาสฺส กิญฺจิ สพฺพาวโต กายสฺส ปริสุทฺเธน เจตสา ปริโยทาเตน อปฺผุฏํ โหติ.\csromanspacer{} อริยสฺส ภิกฺขเว ปญฺจงฺคิกสฺส สมฺมาสมาธิสฺส อยํ จตุตฺถา ภาวนา.}

\prose{ปุน จปรํ ภิกฺขเว ภิกฺขุโน ปจฺจเวกฺขณานิมิตฺตํ สุคฺคหิตํ โหติ สุมนสิกตํ สูปธาริตํ สุปฺปฏิวิทฺธํ ปญฺญาย.\csromanspacer{} เสยฺยถาปิ ภิกฺขเว อญฺโญว อญฺญํ\footnote{อญฺโญ วา อญฺญํ วา (สี), อญฺโญ วา อญฺญํ (สฺยา, กํ), อญฺโญ อญฺญํ (?)} ปจฺจเวกฺเขยฺย, ฐิโต วา นิสินฺนํ ปจฺจเวกฺเขยฺย, นิสินฺโน วา นิปนฺนํ ปจฺจเวกฺเขยฺย.\csromanspacer{} เอวเมว โข ภิกฺขเว ภิกฺขุโน ปจฺจเวกฺขณานิมิตฺตํ สุคฺคหิตํ โหติ สุมนสิกตํ สูปธาริตํ สุปฺปฏิวิทฺธํ ปญฺญาย.\csromanspacer{} อริยสฺส ภิกฺขเว ปญฺจงฺคิกสฺส สมฺมาสมาธิสฺส อยํ ปญฺจมา\footnote{ปญฺจมี (สี)} ภาวนา.\csromanspacer{} เอวํ ภาวิเต}

\vspace{\tipitakaparskip}
\csromancenter{ปญฺจงฺคิกวคฺค โข ภิกฺขเว ภิกฺขุ\footnote{เอวํ ภาวิเต โข ภิกฺขเว (สี)} อริเย ปญฺจงฺคิเก สมฺมาสมาธิมฺหิ เอวํ พหุลีกเต ยสฺส ยสฺส อภิญฺญาสจฺฉิกรณียสฺส ธมฺมสฺส จิตฺตํ อภินินฺนาเมติ อภิญฺญาสจฺฉิกิริยาย, ตตฺร ตเตฺรว สกฺขิภพฺพตํ ปาปุณาติ สติ สติ อายตเน.}
\prose{\csromanfolio{23}เสยฺยถาปิ ภิกฺขเว อุทกมณิโก อาธาเร ฐปิโต ปูโร อุทกสฺส สมติตฺติโก กากเปยฺโย, ตเมนํ พลวา ปุริโส ยโต ยโต อาวชฺเชยฺย\footnote{อาวฏฺเฏยฺย (สฺยา, กํ)} “อาคจฺเฉยฺย อุทกนฺ”ติ.\csromanspacer{} เอวํ ภนฺเต.\csromanspacer{} เอวเมวํ โข ภิกฺขเว ภิกฺขุ เอวํ ภาวิเต อริเย ปญฺจงฺคิเก สมฺมาสมาธิมฺหิ เอวํ พหุลีกเต ยสฺส ยสฺส อภิญฺญาสจฺฉิกรณียสฺส ธมฺมสฺส จิตฺตํ อภินินฺนาเมติ อภิญฺญาสจฺฉิกิริยาย, ตตฺร ตเตฺรว สกฺขิภพฺพตํ ปาปุณาติ สติ สติ อายตเน.}

\prose{เสยฺยถาปิ ภิกฺขเว สเม ภูมิภาเค โปกฺขรณี จตุรํสา อาลิพทฺธา ปูรา อุทกสฺส สมติตฺติกา กากเปยฺยา, ตเมนํ พลวา ปุริโส ยโต ยโต อาลิํ มุญฺเจยฺย “อาคจฺเฉยฺย อุทกนฺ”ติ, เอวํ ภนฺเต.\csromanspacer{} เอวเมวํ โข ภิกฺขเว ภิกฺขุ เอวํ ภาวิเต อริเย ปญฺจงฺคิเก สมฺมาสมาธิมฺหิ เอวํ พหุลีกเต ยสฺส ยสฺส อภิญฺญาสจฺฉิกรณียสฺส ธมฺมสฺส ฯเปฯ สติ สติ อายตเน.}

\prose{เสยฺยถาปิ ภิกฺขเว สุภูมิยํ จตุมหาปเถ\footnote{จาตุมฺมหาปเถ (สี, อิ), จตุมฺมหาปเถ (สฺยา, กํ)} อาชญฺญรโถ ยุตฺโต อสฺส ฐิโต โอธสฺตปโตโท.\csromanspacer{} ตเมนํ ทกฺโข โยคฺคาจริโย อสฺสทมฺมสารถิ อภิรุหิตฺวา วาเมน หตฺเถน รสฺมิโย คเหตฺวา ทกฺขิเณน หตฺเถน ปโตทํ คเหตฺวา เยนิจฺฉกํ ยทิจฺฉกํ สาเรยฺยปิ ปจฺจาสาเรยฺยปิ.\csromanspacer{} เอวเมวํ โข ภิกฺขเว ภิกฺขุ เอวํภาวิเต อริเย ปญฺจงฺคิเก สมฺมาสมาธิมฺหิ เอวํ พหุลีกเต ยสฺส ยสฺส อภิญฺญาสจฺฉิกรณียสฺส ธมฺมสฺส จิตฺตํ อภินินฺนาเมติ อภิญฺญาสจฺฉิกิริยาย, ตตฺร ตเตฺรว สกฺขิภพฺพตํ ปาปุณาติ สติ สติ อายตเน.}

\prose{โส สเจ อากงฺขติ “อเนกวิหิตํ อิทฺธิวิธํ ปจฺจนุภเวยฺยํ, เอโกปิ หุตฺวา พหุธา อสฺสํ ฯเปฯ ยาว พฺรหฺมโลกาปิ กาเยน วสํ วตฺเตยฺยนฺ”ติ, ตตฺร ตเตฺรว สกฺขิภพฺพตํ ปาปุณาติ สติ สติ อายตเน.}

\prose{\csromanfolio{24}โส สเจ อากงฺขติ “ทิพฺพาย โสตธาตุยา วิสุทฺธาย อติกฺกนฺตมานุสิกาย อุโภ สทฺเท สุเณยฺยํ ทิพฺเพ จ มานุเส จ เย ทูเร สนฺติเก จา”ติ, ตตฺร ตเตฺรว สกฺขิ ภพฺพตํ ปาปุณาติ สติ สติ อายตเน.}

\prose{โส สเจ อากงฺขติ “ปรสตฺตานํ ปรปุคฺคลานํ เจตสา เจโต ปริจฺจ ปชาเนยฺยํ, สราคํ วา จิตฺตํ ‘สราคํ จิตฺตนฺ’ติ ปชาเนยฺยํ, วีตราคํ วา จิตฺตํ ‘วีตราคํ จิตฺตนฺ’ติ ปชาเนยฺยํ.\csromanspacer{} สโทสํ วา จิตฺตํ.\csromanspacer{} วีตโทสํ วา จิตฺตํ.\csromanspacer{} สโมหํ วา จิตฺตํ.\csromanspacer{} วีตโมหํ วา จิตฺตํ.\csromanspacer{} สํขิตฺตํ วา จิตฺตํ.\csromanspacer{} วิกฺขิตฺตํ วา จิตฺตํ.\csromanspacer{} มหคฺคตํ วา จิตฺตํ.\csromanspacer{} อมหคฺคตํ วา จิตฺตํ.\csromanspacer{} ส-อุตฺตรํ วา จิตฺตํ.\csromanspacer{} อนุตฺตรํ วา จิตฺตํ.\csromanspacer{} สมาหิตํ วา จิตฺตํ.\csromanspacer{} อสมาหิตํ วา จิตฺตํ.\csromanspacer{} วิมุตฺตํ วา จิตฺตํ.\csromanspacer{} อวิมุตฺตํ วา จิตฺตํ ‘อวิมุตฺตํ จิตฺตนฺ’ติ ปชาเนยฺยนฺ”ติ, ตตฺร ตเตฺรว สกฺขิภพฺพตํ ปาปุณาติ สติ สติ อายตเน.}

\prose{โส สเจ อากงฺขติ “อเนกวิหิตํ ปุพฺเพนิวาสํ อนุสฺสเรยฺยํ, เสยฺยถิทํ? เอกมฺปิ ชาติํ เทฺวปิ ชาติโย ฯเปฯ อิติ สาการํ ส-อุทฺเทสํ อเนกวิหิตํ ปุพฺเพนิวาสํ อนุสฺสเรยฺยนฺ”ติ, ตตฺร ตเตฺรว สกฺขิภพฺพตํ ปาปุณาติ สติ สติ อายตเน.}

\prose{โส สเจ อากงฺขติ “ทิพฺเพน จกฺขุนา วิสุทฺเธน อติกฺกนฺตมานุสเกน ฯเปฯ ยาถากมฺมูปเค สตฺเต ปชาเนยฺยนฺ”ติ, ตตฺร ตเตฺรว สกฺขิภพฺพตํ ปาปุณาติ สติ สติ อายตเน.}

\prose{โส สเจ อากงฺขติ “อาสวานํ ขยา อนาสวํ เจโตวิมุตฺติํ ปญฺญาวิมุตฺติํ ทิฏฺเฐว ธมฺเม สยํ อภิญฺญา สจฺฉิกตฺวา อุปสมฺปชฺช วิหเรยฺยนฺ”ติ, ตตฺร ตเตฺรว สกฺขิภพฺพตํ ปาปุณาติ สติ สติ อายตเนติ.\csromanspacer{}\csromanspacer{}อฏฺฐมํ.}

\csromanheader{2}{9. จงฺกมสุตฺต}
\proseitem{29}{ปญฺจิเม ภิกฺขเว จงฺกเม อานิสํสา.\csromanspacer{} กตเม ปญฺจ? อทฺธานกฺขโม โหติ, ปธานกฺขโม โหติ, อปฺปาพาโธ โหติ, อสิตํ ปีตํ ขายิตํ สายิตํ สมฺมา ปริณามํ คจฺฉติ, จงฺกมาธิคโต สมาธิ จิรฏฺฐิติโก โหติ.\csromanspacer{} อิเม โข ภิกฺขเว ปญฺจ จงฺกเม อานิสํสาติ.\csromanspacer{}\csromanspacer{}นวมํ.}

\chapterheadpageread{3. ปญฺจงฺคิกวคฺค \textbf{10. นีคิตสุตฺต}}
\proseitem{30}{\csromanfolio{25}เอวํ เม สุตํ\csromandash{}เอกํ สมยํ ภควา โกสเลสุ จาริกํ จรมาโน มหตา ภิกฺขุสํเฆน สทฺธิํ เยน อิจฺฉานงฺคลํ นาม โกสลานํ พฺราหฺมณคาโม ตทวสริ.\csromanspacer{} ตตฺร สุทํ ภควา อิจฺฉานงฺคเล วิหรติ อิจฺฉานงฺคลวนสณฺเฑ.\csromanspacer{} อสฺโสสุํ โข อิจฺฉานงฺคลกา\footnote{อิจฺฉานงฺคลิกา (สี), อุปริ 300; อํ 3. 154 ปิฏฺเฐสุปิ.} พฺราหฺมณคหปติกา “สมโณ ขลุ โภ โคตโม สกฺยปุตฺโต สกฺยกุลา ปพฺพชิโต อิจฺฉานงฺคลํ อนุปฺปตฺโต อิจฺฉานงฺคเล วิหรติ อิจฺฉานงฺคลวนสณฺเฑ, ตํ โข ปน ภวนฺตํ โคตมํ เอวํ กลฺยาโณ กิตฺติสทฺโท อพฺภุคฺคโต ‘อิติปิ โส ภควา อรหํ สมฺมาสมฺพุทฺโธ วิชฺชาจรณสมฺปนฺโน สุคโต โลกวิทู อนุตฺตโร ปุริสทมฺมสารถิ สตฺถา เทวมนุสฺสานํ พุทฺโธ ภควา’ติ, โส อิมํ โลกํ สเทวกํ สมารกํ สพฺรหฺมกํ สสฺสมณพฺราหฺมณิํ ปชํ สเทวมนุสฺสํ สยํ อภิญฺญา สจฺฉิกตฺวา ปเวเทติ, โส ธมฺมํ เทเสติ อาทิกลฺยาณํ มชฺเฌกลฺยาณํ ปริโยสานกลฺยาณํ สาตฺถํ สพฺยญฺชนํ เกวลปริปุณฺณํ ปริสุทฺธํ พฺรหฺมจริยํ ปกาเสติ.\csromanspacer{} สาธุ โข ปน ตถารูปานํ อรหตํ ทสฺสนํ โหตี”ติ.\csromanspacer{} อถ โข อิจฺฉานงฺคลกา พฺราหฺมณคหปติกา ตสฺสา รตฺติยา อจฺจเยน ปหูตํ ขาทนียํ โภชนียํ อาทาย เยน อิจฺฉานงฺคลวนสณฺโฑ เตนุปสงฺกมิํสุ, อุปสงฺกมิตฺวา พหิทฺวารโกฏฺฐเก อฏฺฐํสุ อุจฺจาสทฺทมหาสทฺทา.}

\prose{เตน โข ปน สมเยน อายสฺมา นาคิโต ภควโต อุปฏฺฐาโก โหติ.\csromanspacer{} อถ โข ภควา อายสฺมนฺตํ นาคิตํ อามนฺเตสิ “เก ปน โข นาคิต อุจฺจาสทฺทมหาสทฺทา เกวฏฺฏา มญฺเญ มจฺฉวิโลเป”ติ.\csromanspacer{} เอเต ภนฺเต อิจฺฉานงฺคลกา พฺราหฺมณคหปติกา ปหูตํ ขาทนียํ โภชนียํ อาทาย พหิทฺวารโกฏฺฐเก ฐิตา ภควนฺตญฺเญว อุทฺทิสฺส ภิกฺขุสํฆญฺจาติ.\csromanspacer{} มาหํ นาคิต ยเสน สมาคมํ, มา จ มยา ยโส.\csromanspacer{} โย โข นาคิต นยิมสฺส เนกฺขมฺมสุขสฺส ปวิเวกสุขสฺส อุปสมสุขสฺส สมฺโพธสุขสฺส นิกามลาภี อสฺส อกิจฺฉลาภี อกสิรลาภี, ยสฺสาหํ เนกฺขมฺม สุขสฺส ปวิเวกสุกฺขสฺส อุปสมสุขสฺส สมฺโพธสุขสฺส นิกามลาภี \csromanfolio{26}อกิจฺฉลาภี อกสิรลาภี, โส ตํ\footnote{โสหํ (ก), โส (สฺยา, กํ)} มีฬฺหสุขํ มิทฺธสุขํ ลาภสกฺการสิโลกสุขํ สาทิเยยฺยาติ.}

\prose{อธิวาเสตุ ทานิ ภนฺเต ภควา, อธิวาเสตุ สุคโต, อธิวาสนกาโล ทานิ ภนฺเต ภควโต, เยน เยเนว ทานิ ภควา คมิสฺสติ, ตํนินฺนาว คมิสฺสนฺติ พฺราหฺมณคหปติกา เนคมา เจว ชานปทา จ.\csromanspacer{} เสยฺยถาปิ ภนฺเต ถุลฺลผุสิตเก เทเว วสฺสนฺเต ยถานินฺนํ อุทกานิ ปวตฺตนฺติ.\csromanspacer{} เอวเมวํ โข ภนฺเต เยน เยเนว ทานิ ภควา คมิสฺสติ, ตํนินฺนาว คมิสฺสนฺติ พฺราหฺมณคหปติกา เนคมา เจว ชานปทา จ.\csromanspacer{} ตํ กิสฺส เหตุ? ตถา หิ ภนฺเต ภควโต สีลปญฺญาณนฺติ.}

\prose{มาหํ นาคิต ยเสน สมาคมํ, มา จ มยา ยโส.\csromanspacer{} โย โข นาคิต นยิมสฺส เนกฺขมฺมสุขสฺส ปวิเวกสุขสฺส อุปสมสุขสฺส สมฺโพธสุขสฺส นิกามลาภี อสฺส อกิจฺฉลาภี อกสิรลาภี, ยสฺสาหํ เนกฺขมฺมสุขสฺส ปวิเวกสุขสฺส อุปสมสุขสฺส สมฺโพธสุขสฺส นิกามลาภี อกิจฺฉลาภี อกสิรลาภี, โส ตํ มีฬฺหสุขํ มิทฺธสุขํ ลาภสกฺการสิโลกสุขํ สาทิเยยฺย.\csromanspacer{} อสิตปีตขายิตสายิตสฺส โข นาคิต อุจฺจารปสฺสาโว, เอโส ตสฺส นิสฺสนฺโท.\csromanspacer{} ปิยานํ โข นาคิต วิปริณามญฺญถาภาวา อุปฺปชฺชนฺติ โสกปริเทวทุกฺขโทมนสฺสุปายาสา, เอโส ตสฺส นิสฺสนฺโท.\csromanspacer{} อสุภนิมิตฺตานุโยคํ อนุยุตฺตสฺส โข นาคิต สุภนิมิตฺเต ปาฏิกุลฺยตา\footnote{ปฏิกฺกูลตา (สี), ปาฏิกฺกุลฺยตา (สฺยา, กํ)} สณฺฐาติ, เอโส ตสฺส นิสฺสนฺโท.\csromanspacer{} ฉสุ โข นาคิต ผสฺสายตเนสุ อนิจฺจานุปสฺสิโน วิหรโต ผสฺเส ปาฏิกุลฺยตา สณฺฐาติ, เอโส ตสฺส นิสฺสนฺโท.\csromanspacer{} ปญฺจสุ โข นาคิต อุปาทานกฺขนฺเธสุ อุทยพฺพยานุปสฺสิโน วิหรโต อุปาทาเน ปาฏิกุลฺยตา สณฺฐาติ, เอโส ตสฺส นิสฺสนฺโทติ.\csromanspacer{}\csromanspacer{}ทสมํ.\csromanspacer{} ปญฺจงฺคิกวคฺโค ตติโย.}

\titlehead{\textbf{ตสฺสุทฺทานํ}}
\csromangathameasure{เทฺว อคารวุปกฺกิเลสา, ทุสฺสีลานุคฺคหิเตน จ. \\ วิมุตฺติสมาธิปญฺจงฺคิกา, จงฺกมํ นาคิเตน จาติ.}
\csromangathagroup{\csromangathastanza{เทฺว อคารวุปกฺกิเลสา, ทุสฺสีลานุคฺคหิเตน จ. \\ วิมุตฺติสมาธิปญฺจงฺคิกา, จงฺกมํ นาคิเตน จาติ.}}

\chapterheadpageread{4. สุมนวคฺค \textbf{4. สุมนวคฺค}}
\csromanheader{2}{1. สุมนสุตฺต}
\proseitem{31}{\csromanfolio{27}เอกํ สมยํ ฯเปฯ อนาถปิณฺฑิกสฺส อาราเม.\csromanspacer{} อถ โข สุมนา ราชกุมารี ปญฺจหิ รถสเตหิ ปญฺจหิ ราชกุมาริสเตหิ ปริวุตา เยน ภควา เตนุปสงฺกมิ, อุปสงฺกมิตฺวา ภควนฺตํ อภิวาเทตฺวา เอกมนฺตํ นิสีทิ, เอกมนฺตํ นิสินฺนา โข สุมนา ราชกุมารี ภควนฺตํ เอตทโวจ\csromandash{}}

\prose{อิธสฺสุ ภนฺเต ภควโต เทฺว สาวกา สมสทฺธา สมสีลา สมปญฺญา เอโก ทายโก เอโก อทายโก.\csromanspacer{} เต กายสฺส เภทา ปรํ มรณา สุคติํ สคฺคํ โลกํ อุปปชฺเชยฺยุํ.\csromanspacer{} เทวภูตานํ ปน เนสํ\footnote{เตสํ (สี)} ภนฺเต สิยา วิเสโส สิยา นานากรณนฺติ.}

\prose{“สิยา สุมเน”ติ ภควา อโวจ.\csromanspacer{} โย โส สุมเน ทายโก, โส อมุํ อทายกํ เทวภูโต สมาโน ปญฺจหิ ฐาเนหิ อธิคณฺหาติ, ทิพฺเพน อายุนา ทิพฺเพน วณฺเณน ทิพฺเพน สุเขน ทิพฺเพน ยเสน ทิพฺเพน อาธิปเตยฺเยน.\csromanspacer{} โย โส สุมเน ทายโก, โส อมุํ อทายกํ เทวภูโต สมาโน อิเมหิ ปญฺจหิ ฐาเนหิ อธิคณฺหาติ.}

\prose{สเจ ปน เต ภนฺเต ตโต จุตา อิตฺถตฺตํ อาคจฺฉนฺติ, มนุสฺสภูตานํ ปน เนสํ ภนฺเต สิยา วิเสโส สิยา นานากรณนฺติ. “สิยา สุมเน”ติ ภควา อโวจ.\csromanspacer{} โย โส สุมเน ทายโก, โส อมุํ อทายกํ มนุสฺสภูโต สมาโน ปญฺจหิ ฐาเนหิ อธิคณฺหาติ มานุสเกน อายุนา มานุสเกน วณฺเณน มานุสเกน สุเขน มานุสเกน ยเสน มานุสเกน อาธิปเตยฺเยน.\csromanspacer{} โย โส สุมเน ทายโก, โส อมุํ อทายกํ มนุสฺสภูโต สมาโน อิเมหิ ปญฺจหิ ฐาเนหิ อธิคณฺหาติ.}

\prose{สเจ ปน เต ภนฺเต อุโภ อคารสฺมา อนคาริยํ ปพฺพชนฺติ, ปพฺพชิตานํ ปน เนสํ ภนฺเต สิยา วิเสโส สิยา นานากรณนฺติ. “สิยา สุมเน”ติ ภควา อโวจ.\csromanspacer{} โย โส สุมเน ทายโก, โส อมุํ อทายกํ ปพฺพชิโต สมาโน ปญฺจหิ ฐาเนหิ อธิคณฺหาติ ยาจิโตว พหุลํ จีวรํ ปริภุญฺชติ อปฺปํ อยาจิโต, ยาจิโตว พหุลํ ปิณฺฑปาตํ \csromanfolio{28}ปริภุญฺชติ อปฺปํ อยาจิโต, ยาจิโตว พหุลํ เสนาสนํ ปริภุญฺชติ อปฺปํ อยาจิโต, ยาจิโตว พหุลํ คิลานปจฺจยเภสชฺชปริกฺขารํ ปริภุญฺชติ อปฺปํ อยาจิโต.\csromanspacer{} เยหิ โข ปน สพฺรหฺมจารีหิ สทฺธิํ วิหรติ, ตฺยสฺส มนาเปเนว พหุลํ กายกมฺเมน สมุทาจรนฺติ อปฺปํ อมนาเปน, มนาเปเนว พหุลํ วจีกมฺเมน สมุทาจรนฺติ อปฺปํ อมนาเปน, มนาเปเนว พหุลํ มโนกมฺเมน สมุทาจรนฺติ อปฺปํ อมนาเปน, มนาปํเยว พหุลํ อุปหารํ อุปหรนฺติ อปฺปํ อมนาปํ.\csromanspacer{} โย โส สุมเน ทายโก, โส อมุํ อทายกํ ปพฺพชิโต สมาโน อิเมหิ ปญฺจหิ ฐาเนหิ อธิคณฺหาตีติ.}

\prose{สเจ ปน เต ภนฺเต อุโภ อรหตฺตํ ปาปุณนฺติ, อรหตฺตปฺปตฺตานํ ปน เนสํ ภนฺเต สิยา วิเสโส, สิยา นานากรณนฺติ.\csromanspacer{} เอตฺถ โข ปเนสาหํ สุมเน น กิญฺจิ นานากรณํ วทามิ ยทิทํ วิมุตฺติยา วิมุตฺตินฺติ\footnote{วิมุตฺตนฺติ (สฺยา, กํ)}.}

\prose{อจฺฉริยํ ภนฺเต, อพฺภุตํ ภนฺเต, ยาวญฺจิทํ ภนฺเต อลเมว ทานานิ ทาตุํ, อลํ ปุญฺญานิ กาตุํ, ยตฺร หิ นาม เทวภูตสฺสาปิ อุปการานิ ปุญฺญานิ, มนุสฺสภูตสฺสาปิ อุปการานิ ปุญฺญานิ, ปพฺพชิตสฺสาปิ อุปการานิ ปุญฺญานีติ.\csromanspacer{} เอวเมตํ สุมเน, อลํ หิ สุมเน ทานานิ ทาตุํ, อลํ ปุญฺญานิ กาตุํ.\csromanspacer{} เทวภูตสฺสาปิ อุปการานิ ปุญฺญานิ, มนุสฺสภูตสฺสาปิ อุปการานิ ปุญฺญานิ, ปพฺพชิตสฺสาปิ อุปการานิ ปุญฺญานีติ.}

\prose{อิทมโวจ ภควา, อิทํ วตฺวาน\footnote{อิทํ วตฺวา (สี, อิ) เอวมุปริปิ.} สุคโต อถาปรํ เอตทโวจ สตฺถา\csromandash{}}

\csromangathameasure{“ยถาปิ จนฺโท วิมโล, คจฺฉํ อากาสธาตุยา. \\ สพฺเพ ตาราคเณ โลเก, อาภาย อติโรจติ. \\ ตเถว สีลสมฺปนฺโน, สทฺโธ ปุริสปุคฺคโล. \\ สพฺเพ มจฺฉริโน โลเก, จาเคน อติโรจติ. \\ ยถาปิ เมโฆ ถนยํ, วิชฺชุมาลี สตกฺกกุ. \\ ถลํ นินฺนญฺจ ปูเรติ, อภิวสฺสํ วสุนฺธรํ. \\ เอวํ ทสฺสนสมฺปนฺโน, สมฺมาสมฺพุทฺธสาวโก. \\ มจฺฉริํ อธิคณฺหาติ, ปญฺจฐาเนหิ ปณฺฑิโต.}
\csromangathagroup{\csromangathastanza{“ยถาปิ จนฺโท วิมโล, คจฺฉํ อากาสธาตุยา. \\ สพฺเพ ตาราคเณ โลเก, อาภาย อติโรจติ.}\csromangathastanza{ตเถว สีลสมฺปนฺโน, สทฺโธ ปุริสปุคฺคโล. \\ สพฺเพ มจฺฉริโน โลเก, จาเคน อติโรจติ.}\csromangathastanza{ยถาปิ เมโฆ ถนยํ, วิชฺชุมาลี สตกฺกกุ. \\ ถลํ นินฺนญฺจ ปูเรติ, อภิวสฺสํ วสุนฺธรํ.}\csromangathastanza{เอวํ ทสฺสนสมฺปนฺโน, สมฺมาสมฺพุทฺธสาวโก. \\ มจฺฉริํ อธิคณฺหาติ, ปญฺจฐาเนหิ ปณฺฑิโต.}}

\chapterheadpageread{4. สุมนวคฺค}
\csromangathameasure{อายุนา ยสสา เจว, วณฺเณน จ สุเขน จ. \\ ส เว โภคปริพฺยูโฬฺห, เปจฺจ สคฺเค ปโมทตี”ติ.ปฐมํ.}
\csromangathagroup{\csromangathastanza{\csromanfolio{29}อายุนา ยสสา เจว\footnote{อายุนา จ ยเสน จ (ก)}, วณฺเณน จ สุเขน จ. \\ ส เว โภคปริพฺยูโฬฺห\footnote{อายุนา จ ยเสน จ (ก)}, เปจฺจ สคฺเค ปโมทตี”ติ.\csromanspacer{}\csromanspacer{}ปฐมํ.}}

\csromanheader{2}{2. จุนฺทีสุตฺต}
\proseitem{32}{เอกํ สมยํ ภควา ราชคเห วิหรติ เวฬุวเน กลนฺทกนิวาเป อถ โข จุนฺที ราชกุมารี ปญฺจหิ รถสเตหิ ปญฺจหิ จ กุมาริสเตหิ ปริวุตา เยน ภควา เตนุปสงฺกมิ, อุปสงฺกมิตฺวา ภควนฺตํ อภิวาเทตฺวา เอกมนฺตํ นิสีทิ, เอกมนฺตํ นิสินฺนา โข จุนฺที ราชกุมารี ภควนฺตํ เอตทโวจ\csromandash{}}

\prose{อมฺหากํ ภนฺเต ภาตา จุนฺโท นาม ราชกุมาโร โส เอวมาห “ยเทว โส โหติ อิตฺถี วา ปุริโส วา พุทฺธํ สรณํ คโต, ธมฺมํ สรณํ คโต, สํฆํ สรณํ คโต, ปาณาติปาตา ปฏิวิรโต, อทินฺนาทานา ปฏิวิรโต, กาเมสุมิจฺฉาจารา ปฏิวิรโต, มุสาวาทา ปฏิวิรโต, สุราเมรยมชฺชปมาทฏฺฐานา ปฏิวิรโต, โส กายสฺส เภทา ปรํ มรณา สุคติํเยว อุปปชฺชติ โน ทุคฺคตินฺ”ติ.\csromanspacer{} สาหํ ภนฺเต ภควนฺตํ ปุจฺฉามิ กถํรูเป โข ภนฺเต สตฺถริ ปสนฺโน กายสฺส เภทา ปรํ มรณา สุคติํเยว อุปปชฺชติ โน ทุคฺคติํ, กถํรูเป ธมฺเม ปสนฺโน กายสฺส เภทา ปรํ มรณา สุคติํเยว อุปปชฺชติ โน ทุคฺคติํ, กถํรูเป สํเฆ ปสนฺโน กายสฺส เภทา ปรํ มรณา สุคติํเยว อุปปชฺชติ โน ทุคฺคติํ, กถํรูเปสุ สีเลสุ ปริปูรการี กายสฺส เภทา ปรํ มรณา สุคติํเยว อุปปชฺชติ โน ทุคฺคตินฺติ.}

\prose{ยาวตา จุนฺทิ สตฺตา อปทา วา ทฺวิปทา วา จตุปฺปทา วา พหุปฺปทา\footnote{อปาทา วา ทฺวิปาทา วา จตุปฺปาทา วา พหุปฺปาทา วา (สี) อํ 1. 343; ขุ 1. 255 ปิฏฺเฐสุปิ.} วา รูปิโน วา อรูปิโน วา สญฺญิโน วา อสญฺญิโน วา เนวสญฺญินาสญฺญิโน วา, ตถาคโต เตสํ อคฺคมกฺขายติ อรหํ สมฺมาสมฺพุทฺโธ.\csromanspacer{} เย โข จุนฺทิ พุทฺเธ ปสนฺนา, อคฺเค เต ปสนฺนา, อคฺเค โข ปน ปสนฺนานํ อคฺโค วิปาโก}

\prose{ยาวโต จุนฺทิ ธมฺมา สงฺขตา, อริโย อฏฺฐงฺคิโก มคฺโค เตสํ อคฺคมกฺขา ยติ.\csromanspacer{} เย จุนฺทิ อริเย อฏฺฐงฺคิเก มคฺเค ปสนฺนา, อคฺเค เต ปสนฺนา, อคฺเค โข ปน ปสนฺนานํ อคฺโค วิปาโก โหติ.}

\prose{\csromanfolio{30}ยาวตา จุนฺทิ ธมฺมา สงฺขตา วา อสงฺขตา วา, วิราโค เตสํ\footnote{เตสํ ธมฺมานํ (สี, อิ, ก)} อคฺคมกฺขายติ, ยทิทํ มทนิมฺมทโน ปิปาสวินโย อาลยสมุคฺฆาโต วฏฺฏุปจฺเฉโท ตณฺหากฺขโย วิราโค นิโรโธ นิพฺพานํ.\csromanspacer{} เย โข จุนฺทิ วิราเค ธมฺเม ปสนฺนา, อคฺเค เต ปสนฺนา, อคฺเค โข ปน ปสนฺนานํ อคฺโค วิปาโก}

\prose{ยาวตา จุนฺทิ สํฆา วา คณา วา, ตถาคตสาวกสํโฆ เตสํ อคฺคมกฺขายติ, ยทิทํ จตฺตาริ ปุริสยุคานิ อฏฺฐ ปุริสปุคฺคลา, เอส ภควโต สาวกสํโฆ อาหุเนยฺโย ปาหุเนยฺโย ทกฺขิเณยฺโย อญฺชลิกรณีโย อนุตฺตรํ ปุญฺญกฺเขตฺตํ โลกสฺส.\csromanspacer{} เย โข จุนฺทิ สํเฆ ปสนฺนา, อคฺเค เต ปสนฺนา, อคฺเค โข ปน ปสนฺนานํ อคฺโค วิปาโก โหติ.}

\prose{ยาวตา จุนฺทิ สีลานิ, อริยกนฺตานิ สีลานิ เตสํ\footnote{อริยกนฺตานิ เตสํ (สี, สฺยา, กํ)} อคฺคมกฺขายติ ยทิทํ อขณฺฑานิ อจฺฉิทฺทานิ อสพลานิ อกมฺมาสานิ ภุชิสฺสานิ วิญฺญุปฺปสตฺถานิ อปรามฏฺฐานิ สมาธิสํวตฺตนิกานิ.\csromanspacer{} เย โข จุนฺทิ อริยกนฺเตสุ สีเลสุ ปริปูรการิโน, อคฺเค เต ปริปูรการิโน, อคฺเค โข ปน ปริปูรการีนํ อคฺโค วิปาโก โหตีติ.}

\csromangathameasure{อคฺคโต เว ปสนฺนานํ, อคฺคํ ธมฺมํ วิชานตํ. \\ อคฺเค พุทฺเธ ปสนฺนานํ, ทกฺขิเณยฺเย อนุตฺตเร. \\ อคฺเค ธมฺเม ปสนฺนานํ, วิราคูปสเม สุเข. \\ อคฺเค สํเฆ ปสนฺนานํ, ปุญฺญกฺเขตฺเต อนุตฺตเร. \\ อคฺคสฺมิํ ทานํ ททตํ, อคฺคํ ปุญฺญํ ปวฑฺฒติ. \\ อคฺคํ อายุ จ วณฺโณ จ, ยโส กิตฺติ สุขํ พลํ. \\ อคฺคสฺส ทาตา เมธาวี, อคฺคธมฺมสมาหิโต. \\ เทวภูโต มนุสฺโส วา, อคฺคปฺปตฺโต ปโมทตีติ.ตติยํ.}
\csromangathagroup{\csromangathastanza{อคฺคโต เว ปสนฺนานํ, อคฺคํ ธมฺมํ วิชานตํ. \\ อคฺเค พุทฺเธ ปสนฺนานํ, ทกฺขิเณยฺเย อนุตฺตเร.}\csromangathastanza{อคฺเค ธมฺเม ปสนฺนานํ, วิราคูปสเม สุเข. \\ อคฺเค สํเฆ ปสนฺนานํ, ปุญฺญกฺเขตฺเต อนุตฺตเร.}\csromangathastanza{อคฺคสฺมิํ ทานํ ททตํ, อคฺคํ ปุญฺญํ ปวฑฺฒติ. \\ อคฺคํ อายุ จ วณฺโณ จ, ยโส กิตฺติ สุขํ พลํ.}\csromangathastanza{อคฺคสฺส ทาตา เมธาวี, อคฺคธมฺมสมาหิโต. \\ เทวภูโต มนุสฺโส วา, อคฺคปฺปตฺโต ปโมทตีติ.\csromanspacer{}\csromanspacer{}ตติยํ.}}

\csromanheader{2}{3. อุคฺคหสุตฺต}
\proseitem{33}{เอกํ สมยํ ภควา ภทฺทิเย วิหรติ ชาติยา วเน.\csromanspacer{} อถ โข อุคฺคโห เมณฺฑกนตฺตา เยน ภควา เตนุปสงฺกมิ, อุปสงฺกมิตฺวา ภควนฺตํ อภิวาเทตฺวา เอกมนฺตํ นิสีทิ, เอกมนฺตํ นิสินฺโน โข อุคฺคโห เมณฺฑกนตฺตา ภควนฺตํ เอตทโวจ\csromandash{}}

\chapterheadpageread{4. สุมนวคฺค}
\prose{\csromanfolio{31}อธิวาเสตุ เม ภนฺเต ภควา สฺวาตนาย อตฺตจตุตฺโถ ภตฺตนฺติ.\csromanspacer{} อธิวาเสสิ ภควา ตุณฺหีภาเวน.\csromanspacer{} อถ โข อุคฺคโห เมณฺฑกนตฺตา ภควโต อธิวาสนํ วิทิตฺวา อุฏฺฐายาสนา ภควนฺตํ อภิวาเทตฺวา ปทกฺขิณํ กตฺวา ปกฺกามิ.}

\prose{อถ โข ภควา ตสฺสา รตฺติยา อจฺจเยน ปุพฺพณฺหสมยํ นิวาเสตฺวา ปตฺตจีวรมาทาย เยน อุคฺคหสฺส เมณฺฑกนตฺตุโน นิเวสนํ เตนุปสงฺกมิ, อุปสงฺกมิตฺวา ปญฺญตฺเต อาสเน นิสีทิ.\csromanspacer{} อถ โข อุคฺคโห เมณฺฑกนตฺตา ภควนฺตํ ปณีเตน ขาทนีเยน โภชนีเยน สหตฺถา สนฺตปฺเปสิ สมฺปวาเรสิ.\csromanspacer{} อถ โข อุคฺคโห เมณฺฑกนตฺตา ภควนฺตํ ภุตฺตาวิํ โอนีตปตฺตปาณิํ เอกมนฺตํ นิสีทิ, เอกมนฺตํ นิสินฺโน โข อุคฺคโห เมณฺฑกนตฺตา ภควนฺตํ เอตทโวจ “อิมา เม ภนฺเต กุมาริโย ปติกุลานิ คมิสฺสนฺติ, โอวทตุ ตาสํ ภนฺเต ภควา, อนุสาสตุ ตาสํ ภนฺเต ภควา, ยํ ตาสํ อสฺส ทีฆรตฺตํ หิตาย สุขายา”ติ.}

\prose{อถ โข ภควา ตา กุมาริโย เอตทโวจ\csromandash{}ตสฺมาติห กุมาริโย เอวํ สิกฺขิตพฺพํ “ยสฺส โว\footnote{ยสฺส โข (สี, สฺยา, กํ)} มาตาปิตโร ภตฺตุโน ทสฺสนฺติ อตฺถกามา หิเตสิโน อนุกมฺปกา อนุกมฺปํ อุปาทาย, ตสฺส ภวิสฺสาม ปุพฺพุฏฺฐายินิโย ปจฺฉานิปาตินิโย กิํการปฏิสฺสาวินิโย มนาปจารินิโย ปิยวาทินิโย”ติ.\csromanspacer{} เอวํ หิ โว กุมาริโย สิกฺขิตพฺพํ.}

\prose{ตสฺมาติห กุมาริโย เอวํ สิกฺขิตพฺพํ “เย เต ภตฺตุ ครุโน\footnote{คุรุโน (ก)} ภวิสฺสนฺติ ‘มาตา’ติ วา ‘ปิตา’ติ วา ‘สมณพฺราหฺมณา’ติ วา, เต สกฺกริสฺสาม ครุํ กริสฺสาม\footnote{ครุกริสฺสาม (สี, สฺยา, กํ, อิ)} มาเนสฺสาม ปูเชสฺสาม, อพฺภาคเต จ อาสโนทเกน ปฏิปูเชสฺสามา”ติ\footnote{ปูเชสฺสามาติ (สี)}.\csromanspacer{} เอวํ หิ โว กุมาริโย สิกฺขิตพฺพํ.}

\prose{ตสฺมาติห กุมาริโย เอวํ สิกฺขิตพฺพํ “เย เต ภตฺตุ อพฺภนฺตรา กมฺมนฺตา ‘อุณฺณา’ติ วา ‘กปฺปาสา’ติ วา, ตตฺถ ทกฺขา ภวิสฺสาม}

\prose{\csromanfolio{32}อนลสา ตตฺรุปายาย วีมํสาย สมนฺนาคตา อลํ กาตุํ อลํ สํวิธาตุนฺ”ติ.\csromanspacer{} เอวํ หิ โว กุมาริโย สิกฺขิตพฺพํ.}

\prose{ตสฺมาติห กุมาริโย เอวํ สิกฺขิตพฺพํ “โย โส ภตฺตุ อพฺภนฺตโร\footnote{อพฺภนฺตเร (ก) 2. ปจฺจยํ เตน (ก-สี), ปจฺจยํเสน (สฺยา, กํ), ปจฺจยํ เสนาสนํ ปจฺจตฺตํเสน (ก) อํ 3. 93 ปิฏฺเฐ. 3. วิภชิสฺสามาติ (สี, สฺยา, กํ)} อนฺโตชโน ทาสาติ วา เปสฺสาติ วา กมฺมกราติ วา, เตสํ กตญฺจ กตโต ชานิสฺสาม, อกตญฺจ อกตโต ชานิสฺสาม, คิลานกานญฺจ พลาพลํ ชานิสฺสาม, ขาทนียํ โภชนียญฺจสฺส ปจฺจํเสน2 สํวิภชิสฺสามา”ติ3.\csromanspacer{} เอวํ หิ โว กุมาริโย สิกฺขิตพฺพํ.}

\prose{ตสฺมา ติห กุมาริโย เอวํ สิกฺขิตพฺพํ “ยํ ภตฺตา อาหริสฺสติ ธนํ วา ธญฺญํ วา รชตํ วา ชาตรูปํ วา, ตํ อารกฺเขน\footnote{ตํ อารกฺขาย (สี) 5. อิจฺฉาจาเรน (สี), อิสฺสาวาเทน (อิ)} คุตฺติยา สมฺปาเทสฺสาม, ตตฺถ จ ภวิสฺสาม อธุตฺตี อเถนี อโสณฺฑี อวินาสิกาโย”ติ.\csromanspacer{} เอวํ หิ โว กุมาริโย สิกฺขิตพฺพํ.\csromanspacer{} อิเมหิ โข กุมาริโย ปญฺจหิ ธมฺเมหิ สมนฺนาคโต มาตุคาโม กายสฺส เภทา ปรํ มรณา มนาปกายิกานํ เทวานํ สหพฺยตํ อุปปชฺชตีติ.}

\csromangathameasure{โย นํ ภรติ สพฺพทา, นิจฺจํ อาตาปิ อุสฺสุโก. \\ สพฺพกามหรํ โปสํ, ภตฺตารํ นาติมญฺญติ. \\ น จาปิ โสตฺถิ ภตฺตารํ, อิสฺสาจาเรน5 โรสเย. \\ ภตฺตุ จ ครุโน สพฺเพ, ปฏิปูเชติ ปณฺฑิตา. \\ อุฏฺฐาหิกา อนลสา, สงฺคหีตปริชฺชนา. \\ ภตฺตุ มนาปํ จรติ, สมฺภตํ อนุรกฺขติ. \\ ยา เอวํ วตฺตตี นารี, ภตฺตุ ฉนฺทวสานุคา. \\ มนาปา นาม เต เทวา, ยตฺถ สา อุปปชฺชตีติ.ตติยํ.}
\csromangathagroup{\csromangathastanza{โย นํ ภรติ สพฺพทา, นิจฺจํ อาตาปิ อุสฺสุโก. \\ สพฺพกามหรํ โปสํ, ภตฺตารํ นาติมญฺญติ.}\csromangathastanza{น จาปิ โสตฺถิ ภตฺตารํ, อิสฺสาจาเรน5 โรสเย. \\ ภตฺตุ จ ครุโน สพฺเพ, ปฏิปูเชติ ปณฺฑิตา.}\csromangathastanza{อุฏฺฐาหิกา\footnote{อุฏฺฐายิกา (สฺยา, กํ, ก)} อนลสา, สงฺคหีตปริชฺชนา. \\ ภตฺตุ มนาปํ\footnote{อุฏฺฐายิกา (สฺยา, กํ, ก)} จรติ, สมฺภตํ อนุรกฺขติ.}\csromangathastanza{ยา เอวํ วตฺตตี นารี, ภตฺตุ ฉนฺทวสานุคา. \\ มนาปา นาม เต เทวา, ยตฺถ สา อุปปชฺชตีติ.\csromanspacer{}\csromanspacer{}ตติยํ.}}

\csromanheader{2}{4. สีหเสนาปติสุตฺต}
\proseitem{34}{เอกํ สมยํ ภควา เวสาลิยํ วิหรติ มหาวเน กูฏาคารสาลายํ.\csromanspacer{} อถ โข สีโห เสนาปติ เยน ภควา เตนุปสงฺกมิ,}

\vspace{\tipitakaparskip}
\csromancenter{สุมนวคฺค อุปสงฺกมิตฺวา ภควนฺตํ อภิวาเทตฺวา เอกมนฺตํ นิสีทิ, เอกมนฺตํ นิสินฺโน โข สีโห เสนาปติ ภควนฺตํ เอตทโวจ\csromandash{}สกฺกา นุ โข ภนฺเต ภควา สนฺทิฏฺฐิกํ ทานผลํ ปญฺญาเปตุนฺติ.}
\prose{\csromanfolio{33}“สกฺกา สีหา”ติ ภควา อโวจ, ทายโก สีห ทานปติ พหุโน ชนสฺส ปิโย โหติ มนาโป, ยมฺปิ สีห ทายโก ทานปติ พหุโน ชนสฺส ปิโย โหติ มนาโป, อิทมฺปิ สนฺทิฏฺฐิกํ ทานผลํ.}

\prose{ปุน จปรํ สีห ทายกํ ทานปติํ สนฺโต สปฺปุริสา ภชนฺติ, ยมฺปิ สีห ทายกํ ทานปติํ สนฺโต สปฺปุริสา ภชนฺติ, อิทมฺปิ สนฺทิฏฺฐิกํ ทานผลํ.}

\prose{ปุน จปรํ สีห ทายกสฺส ทานปติโน กลฺยาโณ กิตฺติสทฺโท อพฺภุคฺคจฺฉติ, ยมฺปิ สีห ทายกสฺส ทานปติโน กลฺยาโณ กิตฺติสทฺโท อพฺภุคฺคจฺฉติ, อิทมฺปิ สนฺทิฏฺฐิกํ ทานผลํ.}

\prose{ปุน จปรํ สีห ทายโก ทานปติ ยํ ยเทว ปริสํ อุปสงฺกมติ ยทิ ขตฺติยปริสํ ยทิ พฺราหฺมณปริสํ ยทิ คหปติปริสํ ยทิ สมณปริสํ, วิสารโท\footnote{วิสารโทว (สี), อุปริ 458 ปิฏฺเฐปิ ปสฺสิตพฺพํ.} อุปสงฺกมติ อมงฺกุภูโต.\csromanspacer{} ยมฺปิ สีห ทายโก ทานปติ ยํ ยเทว ปริสํ อุปสงฺกมติ ยทิ ขตฺติยปริสํ ยทิ พฺราหฺมณปริสํ ยทิ คหปติปริสํ ยทิ สมณปริสํ, วิสารโท อุปสงฺกมติ อมงฺกุภูโต.\csromanspacer{} อิทมฺปิ สนฺทิฏฺฐิกํ ทานผลํ.}

\prose{ปุน จปรํ สีห ทายโก ทานปติ กายสฺส เภทา ปรํ มรณา สุคติํ สคฺคํ โลกํ อุปปชฺชติ, ยมฺปิ สีห ทายโก ทานปติ กายสฺส เภทา ปรํ มรณา สุคติํ สคฺคํ โลกํ อุปปชฺชติ.\csromanspacer{} อิทํ\footnote{อิทมฺปิ สีห (ก)} สมฺปรายิกํ ทานผลนฺติ.}

\prose{เอวํ วุตฺเต สีโห เสนาปติ ภควนฺตํ เอตทโวจ “ยานิมานิ ภนฺเต ภควตา จตฺตาริ สนฺทิฏฺฐิกานิ ทานผลานิ อกฺขาตานิ, นาหํ เอตฺถ ภควโต สทฺธาย คจฺฉามิ, อหํเปตานิ ชานามิ.\csromanspacer{} อหํ ภนฺเต ทายโก ทานปติ, พหุโน ชนสฺส ปิโย มนาโป, อหํ ภนฺเต ทายโก ทานปติ, มํ สนฺโต สปฺปุริสา ภชนฺติ, อหํ ภนฺเต ทายโก ทานปติ, มยฺหํ กลฺยาโณ กิตฺติสทฺโท อพฺภุคฺคโต “สีโห เสนาปติ ทายโก \csromanfolio{34}การโก สํฆุปฏฺฐาโก”ติ, อหํ ภนฺเต ทายโก ทานปติ, ยํ ยเทว ปริสํ อุปสงฺกมามิ ยทิ ขตฺติยปริสํ ยทิ พฺราหฺมณปริสํ ยทิ คหปติปริสํ ยทิ สมณปริสํ, วิสารโท อุปสงฺกมามิ อมงฺกุภูโต.\csromanspacer{} ยานิมานิ ภนฺเต ภควตา จตฺตาริ สนฺทิฏฺฐิกานิ ทานผลานิ อกฺขาตานิ, นาหํ เอตฺถ ภควโต สทฺธาย คจฺฉามิ, อหํเปตานิ ชานามิ.\csromanspacer{} ยญฺจ โข มํ ภนฺเต ภควา เอวมาห “ทายโก สีห ทานปติ กายสฺส เภทา ปรํ มรณา สุคติํ สคฺคํ โลกํ อุปปชฺชตี”ติ.\csromanspacer{} เอตาหํ น ชานามิ, เอตฺถ จ ปนาหํ ภควโต สทฺธาย คจฺฉามีติ.\csromanspacer{} เอวเมตํ สีห เอวเมตํ สีห, ทายโก ทานปติ กายสฺส เภทา ปรํ มรณา สุคติํ สคฺคํ โลกํ อุปปชฺชตีติ.}

\csromangathameasure{ททํ ปิโย โหติ ภชนฺติ นํ พหู, \\ กิตฺติญฺจ ปปฺโปติ ยโส จ วฑฺฒติ. \\ อมงฺกุภูโต ปริสํ วิคาหติ, \\ วิสารโท โหติ นโร อมจฺฉรี. \\ ตสฺมา หิ ทานานิ ททนฺติ ปณฺฑิตา, \\ วิเนยฺย มจฺเฉรมลํ สุเขสิโน. \\ เต ทีฆรตฺตํ ติทิเว ปติฏฺฐิตา, \\ เทวานํ สหพฺยคตา รมนฺติ เต. \\ กตาวกาสา กตกุสลา อิโต จุตา, \\ สยํปภา อนุวิจรนฺติ นนฺทนํ. \\ เต ตตฺถ นนฺทนฺติ รมนฺติ โมทเร, \\ สมปฺปิตา กามคุเณหิ ปญฺจหิ. \\ กตฺวาน วากฺยํ อสิตสฺส ตาทิโน, \\ รมนฺติ สคฺเค สุคตสฺส สาวกาติ.จตุตฺถํ.}
\csromangathagroup{\csromangathastanza{ททํ ปิโย โหติ ภชนฺติ นํ พหู, \\ กิตฺติญฺจ ปปฺโปติ ยโส จ วฑฺฒติ\footnote{ยสสฺส วฑฺฒติ (สฺยา, กํ), ยสํ ปวฑฺฒติ (ก)}. \\ อมงฺกุภูโต ปริสํ วิคาหติ, \\ วิสารโท โหติ นโร อมจฺฉรี.}\csromangathastanza{ตสฺมา หิ ทานานิ ททนฺติ ปณฺฑิตา, \\ วิเนยฺย มจฺเฉรมลํ สุเขสิโน. \\ เต ทีฆรตฺตํ ติทิเว ปติฏฺฐิตา, \\ เทวานํ สหพฺยคตา รมนฺติ เต\footnote{สหพฺยตํ คตา รมนฺติ (สี), สหพฺยตา รมนฺติ เต (ก)}.}\csromangathastanza{กตาวกาสา กตกุสลา อิโต จุตา\footnote{ตโต จุตา (สี)}, \\ สยํปภา อนุวิจรนฺติ นนฺทนํ\footnote{ตโต จุตา (สี)}. \\ เต ตตฺถ นนฺทนฺติ รมนฺติ โมทเร, \\ สมปฺปิตา กามคุเณหิ ปญฺจหิ. \\ กตฺวาน วากฺยํ อสิตสฺส ตาทิโน, \\ รมนฺติ สคฺเค\footnote{ตโต จุตา (สี)} สุคตสฺส สาวกาติ.\csromanspacer{}\csromanspacer{}จตุตฺถํ.}}

\csromanheader{2}{5. ทานานิสํสสุตฺต}
\proseitem{35}{ปญฺจิเม ภิกฺขเว ทาเน อานิสํสา.\csromanspacer{} กตเม ปญฺจ? พหุโน ชนสฺส ปิโย โหติ มนาโป, สนฺโต สปฺปุริสา ภชนฺติ, กลฺยาโณ}

\vspace{\tipitakaparskip}
\csromancenter{สุมนวคฺค กิตฺติสทฺโท อพฺภุคฺคจฺฉติ, คิหิธมฺมา อนปคโต\footnote{คิหิธมฺมา อนเปโต (สี, อิ), คิหิธมฺมมนุปคโต (ก)} โหติ, กายสฺส เภทา ปรํ มรณา สุคติํ สคฺคํ โลกํ อุปปชฺชติ.\csromanspacer{} อิเม โข ภิกฺขเว ปญฺจ ทาเน อานิสํสาติ.}
\vspace{0.5\baselineskip}
\csromangathameasure{ททมาโน ปิโย โหติ, สตํ ธมฺมํ อนุกฺกมํ. \\ สนฺโต นํ สทา ภชนฺติ, สญฺญตา พฺรหฺมจารโย. \\ เต ตสฺส ธมฺมํ เทเสนฺติ, สพฺพทุกฺขาปนูทนํ. \\ ยํ โส ธมฺมํ อิธญฺญาย, ปรินิพฺพาติ อนาสโวติ.ปญฺจมํ.}
\csromangathagroup{\csromangathastanza{\csromanfolio{35}ททมาโน ปิโย โหติ, สตํ ธมฺมํ อนุกฺกมํ. \\ สนฺโต นํ สทา ภชนฺติ\footnote{สนฺโต ภชนฺติ สปฺปุริสา (สี)}, สญฺญตา พฺรหฺมจารโย.}\csromangathastanza{เต ตสฺส ธมฺมํ เทเสนฺติ, สพฺพทุกฺขาปนูทนํ. \\ ยํ โส ธมฺมํ อิธญฺญาย, ปรินิพฺพาติ อนาสโวติ.\csromanspacer{}\csromanspacer{}ปญฺจมํ.}}

\csromanheader{2}{6. กาลทานสุตฺต}
\proseitem{36}{ปญฺจิมานิ ภิกฺขเว กาลทานานิ.\csromanspacer{} กตมานิ ปญฺจ? อาคนฺตุกสฺส ทานํ เทติ, คมิกสฺส ทานํ เทติ, คิลานสฺส ทานํ เทติ, ทุพฺภิกฺเข ทานํ เทติ, ยานิ ตานิ นวสสฺสานิ นวผลานิ, ตานิ ปฐมํ สีลวนฺเตสุ ปติฏฺฐาเปติ.\csromanspacer{} อิมานิ โข ภิกฺขเว ปญฺจ กาลทานานีติ.}

\csromangathameasure{กาเล ททนฺติ สปฺปญฺญา, วทญฺญู วีตมจฺฉรา. \\ กาเลน ทินฺนํ อริเยสุ, อุชุภูเตสุ ตาทิสุ. \\ วิปฺปสนฺนมนา ตสฺส, วิปุลา โหติ ทกฺขิณา. \\ เย ตตฺถ อนุโมทนฺติ, เวยฺยาวจฺจํ กโรนฺติ วา. \\ น เตน ทกฺขินา อูนา, เตปิ ปุญฺญสฺส ภาคิโน. \\ ตสฺมา ทเท อปฺปฏิวานจิตฺโต, ยตฺถ ทินฺนํ มหปฺผลํ. \\ ปุญฺญานิ ปรโลกสฺมิํ, ปติฏฺฐา โหนฺติ ปาณินนฺติ.ฉฏฺฐํ.}
\csromangathagroup{\csromangathastanza{กาเล ททนฺติ สปฺปญฺญา, วทญฺญู วีตมจฺฉรา. \\ กาเลน ทินฺนํ อริเยสุ, อุชุภูเตสุ ตาทิสุ.}\csromangathastanza{วิปฺปสนฺนมนา ตสฺส, วิปุลา โหติ ทกฺขิณา. \\ เย ตตฺถ อนุโมทนฺติ, เวยฺยาวจฺจํ กโรนฺติ วา.}\csromangathastanza{น เตน\footnote{น เตสํ (อิ, ก)} ทกฺขินา อูนา, เตปิ ปุญฺญสฺส ภาคิโน. \\ ตสฺมา ทเท อปฺปฏิวานจิตฺโต, ยตฺถ ทินฺนํ มหปฺผลํ. \\ ปุญฺญานิ ปรโลกสฺมิํ, ปติฏฺฐา โหนฺติ ปาณินนฺติ.\csromanspacer{}\csromanspacer{}ฉฏฺฐํ.}}

\csromanheader{2}{7. โภชนสุตฺต}
\proseitem{37}{โภชนํ ภิกฺขเว ททมาโน ทายโก ปฏิคฺคาหกานํ ปญฺจ ฐานานิ เทติ.\csromanspacer{} กตมานิ ปญฺจ? อายุํ เทติ, วณฺณํ เทติ, สุขํ เทติ, พลํ เทติ, ปฏิภานํ\footnote{ปฏิภาณํ (สี)} เทติ.\csromanspacer{} อายุํ โข ปน ทตฺวา อายุสฺส ภาคี โหติ ทิพฺพสฺส วา มานุสสฺส วา, วณฺณํ ทตฺวา วณฺณสฺส ภาคี โหติ ทิพฺพสฺส วา มานุสสฺส วา, สุขํ ทตฺวา สุขสฺส ภาคี โหติ ทิพฺพสฺส วา \csromanfolio{36}มานุสสฺส วา, พลํ ทตฺวา พลสฺส ภาคี โหติ ทิพฺพสฺส วา มานุสสฺส วา, ปฏิภานํ ทตฺวา ปฏิภานสฺส ภาคี โหติ ทิพฺพสฺส วา มานุสสฺส วา.\csromanspacer{} โภชนํ ภิกฺขเว ททมาโน ทายโก ปฏิคฺคาหกานํ อิมานิ ปญฺจ ฐานานิ เทตีติ.}

\csromangathameasure{อายุโท พลโท ธีโร, วณฺณโท ปฏิภานโท. \\ สุขสฺส ทาตา เมธาวี, สุขํ โส อธิคจฺฉติ. \\ อายุํ ทตฺวา พลํ วณฺณํ, สุขญฺจ ปฏิภานกํ. \\ ทีฆายุ ยสวา โหติ, ยตฺถ ยตฺถูปปชฺชตีติ.สตฺตมํ.}
\csromangathagroup{\csromangathastanza{อายุโท พลโท ธีโร, วณฺณโท ปฏิภานโท. \\ สุขสฺส ทาตา เมธาวี, สุขํ โส อธิคจฺฉติ.}\csromangathastanza{อายุํ ทตฺวา พลํ วณฺณํ, สุขญฺจ ปฏิภานกํ\footnote{ปฏิภาณกํ (สี), ปฏิภานโท (สฺยา, กํ, อิ, ก)}. \\ ทีฆายุ ยสวา โหติ, ยตฺถ ยตฺถูปปชฺชตีติ.\csromanspacer{}\csromanspacer{}สตฺตมํ.}}

\csromanheader{2}{8. สทฺธสุตฺต}
\proseitem{38}{ปญฺจิเม ภิกฺขเว สทฺเธ กุลปุตฺเต อานิสํสา.\csromanspacer{} กตเม ปญฺจ? เย เต ภิกฺขเว โลเก สนฺโต สปฺปุริสา, เต สทฺธญฺเญว ปฐมํ อนุกมฺปนฺตา อนุกมฺปนฺติ โน ตถา อสฺสทฺธํ, สทฺธญฺเญว ปฐมํ อุปสงฺกมนฺตา อุปสงฺกมนฺติ โน ตถา อสฺสทฺธํ, สทฺธญฺเญว ปฐมํ ปฏิคฺคณฺหนฺตา ปฏิคฺคณฺหนฺติ โน ตถา อสฺสทฺธํ, สทฺธญฺเญว ปฐมํ ธมฺมํ เทเสนฺตา เทเสนฺติ โน ตถา อสฺสทฺธํ, สทฺโธ กายสฺส เภทา ปรํ มรณา สุคติํ สคฺคํ โลกํ อุปปชฺชติ.\csromanspacer{} อิเม โข ภิกฺขเว ปญฺจ สทฺเธ กุลปุตฺเต อานิสํสา.}

\prose{เสยฺยถาปิ ภิกฺขเว สุภูมิยํ จตุมหาปเถ มหานิโคฺรโธ สมนฺตา ปกฺขีนํ ปฏิสรณํ โหติ.\csromanspacer{} เอวเมวํ โข ภิกฺขเว สทฺโธ กุลปุตฺโต พหุโน ชนสฺส ปฏิสรณํ โหติ ภิกฺขูนํ ภิกฺขูนีนํ อุปาสกานํ อุปาสิกานนฺติ.}

\csromangathameasure{สาขาปตฺตผลูเปโต, ขนฺธิมาว มหาทุโม. \\ มูลวา ผลสมฺปนฺโน, ปติฏฺฐา โหติ ปกฺขินํ. \\ มโนรเม อายตเน, เสวนฺติ นํ วิหงฺคมา. \\ ฉายํ ฉายตฺถิกา ยนฺติ, ผลตฺถา ผลโภชิโน. \\ ตเถว สีลสมฺปนฺนํ, สทฺธํ ปุริสปุคฺคลํ. \\ นิวาตวุตฺติํ อตฺถทฺธํ, โสรตํ สขิลํ มุทุํ.}
\csromangathagroup{\csromangathastanza{สาขาปตฺตผลูเปโต\footnote{สาขาปตฺตพหุเปโต (กตฺถจิ), สาขาปตฺตปลาสูเปโต (?)}, ขนฺธิมาว\footnote{ขนฺธิมา จ (สี)} มหาทุโม. \\ มูลวา ผลสมฺปนฺโน, ปติฏฺฐา โหติ ปกฺขินํ.}\csromangathastanza{มโนรเม อายตเน, เสวนฺติ นํ วิหงฺคมา. \\ ฉายํ ฉายตฺถิกา\footnote{ฉายตฺถิโน (สี)} ยนฺติ, ผลตฺถา ผลโภชิโน.}\csromangathastanza{ตเถว สีลสมฺปนฺนํ, สทฺธํ ปุริสปุคฺคลํ. \\ นิวาตวุตฺติํ อตฺถทฺธํ, โสรตํ สขิลํ มุทุํ.}}

\chapterheadpageread{4. สุมนวคฺค}
\csromangathameasure{วีตราคา วีตโทสา, วีตโมหา อนาสวา. \\ ปุญฺญกฺเขตฺตานิ โลกสฺมิํ, เสวนฺติ ตาทิสํ นรํ. \\ เต ตสฺส ธมฺมํ เทเสนฺติ, สพฺพทุกฺขาปนูทนํ. \\ ยํ โส ธมฺมํ อิธญฺญาย, ปรินิพฺพาติ อนาสโวติ.อฏฺฐมํ.}
\csromangathagroup{\csromangathastanza{\csromanfolio{37}วีตราคา วีตโทสา, วีตโมหา อนาสวา. \\ ปุญฺญกฺเขตฺตานิ โลกสฺมิํ, เสวนฺติ ตาทิสํ นรํ.}\csromangathastanza{เต ตสฺส ธมฺมํ เทเสนฺติ, สพฺพทุกฺขาปนูทนํ. \\ ยํ โส ธมฺมํ อิธญฺญาย, ปรินิพฺพาติ อนาสโวติ.\csromanspacer{}\csromanspacer{}อฏฺฐมํ.}}

\csromanheader{2}{9. ปุตฺตสุตฺต}
\proseitem{39}{ปญฺจิมานิ ภิกฺขเว ฐานานิ สมฺปสฺสนฺตา มาตาปิตโร ปุตฺตํ อิจฺฉนฺติ กุเล ชายมานํ.\csromanspacer{} กตมานิ ปญฺจ? ภโต วา โน ภริสฺสติ, กิจฺจํ วา โน กริสฺสติ, กุลวํโส จิรํ ฐสฺสติ, ทายชฺชํ ปฏิปชฺชิสฺสติ, อถ วา ปน เปตานํ กาลงฺกตานํ ทกฺขิณํ อนุปฺปทสฺสตีติ.\csromanspacer{} อิมานิ โข ภิกฺขเว ปญฺจ ฐานานิ สมฺปสฺสนฺตา มาตาปิตโร ปุตฺตํ อิจฺฉนฺติ กุเล ชายมานนฺติ.}

\prose{\csromansymbolfootnote{*}{อภิ 4. 261 ปิฏฺเฐปิ.}ปญฺจ ฐานานิ สมฺปสฺสํ, ปุตฺตํ อิจฺฉนฺติ ปณฺฑิตา.}

\csromangathameasure{ภโต วา โน ภริสฺสติ, กิจฺจํ วา โน กริสฺสติ. \\ กุลวํโส จิรํ ติฏฺเฐ, ทายชฺชํ ปฏิปชฺชติ. \\ อถ วา ปน เปตานํ, ทกฺขิณํ อนุปฺปทสฺสติ. \\ ฐานาเนตานิ สมฺปสฺสํ, ปุตฺตํ อิจฺฉนฺติ ปณฺฑิตา. \\ ตสฺมา สนฺโต สปฺปุริสา, กตญฺญู กตเวทิโน. \\ ภรนฺติ มาตาปิตโร, ปุพฺเพ กตมนุสฺสรํ. \\ กโรนฺติ เนสํ กิจฺจานิ, ยถา ตํ ปุพฺพการินํ. \\ โอวาทการี ภตโปสี, กุลวํสํ อหาปยํ. \\ สทฺโธ สีเลน สมฺปนฺโน, ปุตฺโต โหติ ปสํสิโยติ. นวมํ.}
\csromangathagroup{\csromangathastanza{ภโต วา โน ภริสฺสติ, กิจฺจํ วา โน กริสฺสติ. \\ กุลวํโส จิรํ ติฏฺเฐ, ทายชฺชํ ปฏิปชฺชติ.}\csromangathastanza{อถ วา ปน เปตานํ, ทกฺขิณํ อนุปฺปทสฺสติ. \\ ฐานาเนตานิ สมฺปสฺสํ, ปุตฺตํ อิจฺฉนฺติ ปณฺฑิตา.}\csromangathastanza{ตสฺมา สนฺโต สปฺปุริสา, กตญฺญู กตเวทิโน. \\ ภรนฺติ มาตาปิตโร, ปุพฺเพ กตมนุสฺสรํ.}\csromangathastanza{กโรนฺติ เนสํ กิจฺจานิ, ยถา ตํ ปุพฺพการินํ. \\ โอวาทการี ภตโปสี, กุลวํสํ อหาปยํ. \\ สทฺโธ สีเลน สมฺปนฺโน, ปุตฺโต โหติ ปสํสิโยติ.\csromanspacer{} นวมํ.}}

\csromanheader{2}{10. มหาสาลปุตฺตสุตฺต}
\proseitem{40}{หิมวนฺตํ ภิกฺขเว ปพฺพตราชํ นิสฺสาย มหาสาลา ปญฺจหิ วฑฺฒีหิ วฑฺฒนฺติ.\csromanspacer{} กตมาหิ ปญฺจหิ? สาขาปตฺตปลาเสน วฑฺฒนฺติ, ตเจน วฑฺฒนฺติ, ปปฏิกาย วฑฺฒนฺติ, เผคฺคุนา วฑฺฒนฺติ, สาเรน วฑฺฒนฺติ.\csromanspacer{} หิมวนฺตํ ภิกฺขเว ปพฺพตราชํ นิสฺสาย มหาสาลา อิมาหิ ปญฺจหิ วฑฺฒีหิ วฑฺฒนฺติ.\csromanspacer{} เอวเมวํ โข ภิกฺขเว สทฺธํ กุลปุตฺตํ นิสฺสาย อนฺโตชโน ปญฺจหิ วฑฺฒีหิ วฑฺฒติ. \csromanfolio{38}กตมาหิ ปญฺจหิ, สทฺธาย วฑฺฒติ, สีเลน วฑฺฒติ, สุเตน วฑฺฒติ, จาเคน วฑฺฒติ, ปญฺญาย วฑฺฒติ.\csromanspacer{} สทฺธํ ภิกฺขเว กุลปุตฺตํ นิสฺสาย อนฺโตชโน อิมาหิ ปญฺจหิ วฑฺฒีหิ วฑฺฒตีติ.}

\csromangathameasure{ยถา หิ ปพฺพโต เสโล, อรญฺญสฺมิํ พฺรหาวเน. \\ ตํ รุกฺขา อุปนิสฺสาย, วฑฺฒนฺเต เต วนปฺปตี. \\ ตเถว สีลสมฺปนฺนํ, สทฺธํ กุลปุตฺตํ อิมํ. \\ อุปนิสฺสาย วฑฺฒนฺติ, ปุตฺตทารา จ พนฺธวา. \\ อมจฺจา ญาติสํฆา จ, เย จสฺส อนุชีวิโน. \\ ตฺยสฺส สีลวโต สีลํ, จาคํ สุจริตานิ จ. \\ ปสฺสมานานุกุพฺพนฺติ, เย ภวนฺติ วิจกฺขณา. \\ อิธ ธมฺมํ จริตฺวาน, มคฺคํ สุคติคามินํ. \\ นนฺทิโน เทวโลกสฺมิํ, โมทนฺติ กามกามิโนติ.ทสมํ. สุมนวคฺโค จตุตฺโถ.}
\csromangathagroup{\csromangathastanza{ยถา หิ ปพฺพโต เสโล, อรญฺญสฺมิํ พฺรหาวเน. \\ ตํ รุกฺขา อุปนิสฺสาย, วฑฺฒนฺเต เต วนปฺปตี.}\csromangathastanza{ตเถว สีลสมฺปนฺนํ, สทฺธํ กุลปุตฺตํ อิมํ\footnote{กุลปติํ อิธ (สี), กุลปุตฺตํ อิธ (สฺยา)}. \\ อุปนิสฺสาย วฑฺฒนฺติ, ปุตฺตทารา จ พนฺธวา.}\csromangathastanza{อมจฺจา ญาติสํฆา จ, เย จสฺส อนุชีวิโน. \\ ตฺยสฺส สีลวโต สีลํ, จาคํ สุจริตานิ จ.}\csromangathastanza{ปสฺสมานานุกุพฺพนฺติ, เย ภวนฺติ วิจกฺขณา. \\ อิธ ธมฺมํ จริตฺวาน, มคฺคํ\footnote{สคฺคํ (สฺยา, ก)} สุคติคามินํ. \\ นนฺทิโน เทวโลกสฺมิํ, โมทนฺติ กามกามิโนติ.\csromanspacer{}\csromanspacer{}ทสมํ.\csromanspacer{} สุมนวคฺโค จตุตฺโถ.}}

\titlehead{\textbf{ตสฺสุทฺทานํ}}
\csromangathameasure{สุมนา จุนฺที อุคฺคโห, สีโห ทานานิสํสโก. \\ กาลโภชนสทฺธา จ, ปุตฺตสาเลหิ เต ทสาติ.}
\csromangathagroup{\csromangathastanza{สุมนา จุนฺที อุคฺคโห, สีโห ทานานิสํสโก. \\ กาลโภชนสทฺธา จ, ปุตฺตสาเลหิ เต ทสาติ.}}

\chapterhead{5. มุณฺฑราชวคฺค}
\csromanheader{2}{1. อาทิยสุตฺต}
\proseitem{41}{เอกํ สมยํ ภควา สาวตฺถิยํ วิหรติ เชตวเน อนาถปิณฺฑิกสฺส อาราเม.\csromanspacer{} อถ โข อนาถปิณฺฑิโก คหปติ เยน ภควา เตนุปสงฺกมิ, อุปสงฺกมิตฺวา ภควนฺตํ อภิวาเทตฺวา เอกมนฺตํ นิสีทิ, เอกมนฺตํ นิสินฺนํ โข อนาถปิณฺฑิกํ คหปติํ ภควา เอตทโวจ\csromandash{} ปญฺจิเม คหปติ โภคานํ อาทิยา.\csromanspacer{} กตเม ปญฺจ? อิธ คหปติ อริยสาวโก อุฏฺฐานวีริยาธิคเตหิ โภเคหิ พาหาพลปริจิเตหิ เสทาวกฺขิตฺเตหิ ธมฺมิเกหิ ธมฺมลทฺเธหิ อตฺตานํ สุเขติ ปีเณติ, สมฺมา}

\vspace{\tipitakaparskip}
\csromancenter{มุณฺฑราชวคฺค สุขํ ปริหรติ, มาตาปิตโร สุเขติ ปีเณติ, สมฺมา สุขํ ปริหรติ, ปุตฺตทารทาสกมฺมกรโปริเส สุเขติ ปีเณติ, สมฺมา สุขํ ปริหรติ.\csromanspacer{} อยํ ปฐโม โภคานํ อาทิโย.}
\prose{\csromanfolio{39}ปุน จปรํ คหปติ อริยสาวโก อุฏฺฐานวีริยาธิคเตหิ โภเคหิ พาหาพลปริจิเตหิ เสทาวกฺขิตฺเตหิ ธมฺมิเกหิ ธมฺมลทฺเธหิ มิตฺตามจฺเจ สุเขติ ปีเณติ, สมฺมา สุขํ ปริหรติ.\csromanspacer{} อยํ ทุติโย โภคานํ อาทิโย.}

\prose{ปุน จปรํ คหปติ อริยสาวโก อุฏฺฐานวีริยาธิคเตหิ โภเคหิ พาหาพลปริจิเตหิ เสทาวกฺขิตฺเตหิ ธมฺมิเกหิ ธมฺมลทฺเธหิ, ยา ตา โหนฺติ อาปทา อคฺคิโต วา อุทกโต วา ราชโต วา โจรโต วา อปฺปิยโต วา ทายาทโต\footnote{อปฺปิยโต วา ทายาทโต วา (พหูสุ) อํ 1. 379; อุปริ 152 ปิฏฺเฐสุปิ.}, ตถารูปาสุ อาปทาสุ โภเคหิ ปริโยธาย วตฺตติ, โสตฺถิํ อตฺตานํ กโรติ.\csromanspacer{} อยํ ตติโย โภคานํ อาทิโย.}

\prose{ปุน จปรํ คหปติ อริยสาวโก อุฏฺฐานวีริยาธิคเตหิ โภเคหิ พาหาพลปริจิเตหิ เสทาวกฺขิตฺเตหิ ธมฺมิเกหิ ธมฺมลทฺเธหิ ปญฺจพลิํ กตฺตา โหติ ญาติพลิํ อติถิพลิํ ปุพฺพเปตพลิํ ราชพลิํ เทวตาพลิํ.\csromanspacer{} อยํ จตุตฺโถ โภคานํ อาทิโย.}

\prose{ปุน จปรํ คหปติ อริยสาวโก อุฏฺฐานวีริยาธิคเตหิ โภเคหิ พาหาพลปริจิเตหิ เสทาวกฺขิตฺเตหิ ธมฺมิเกหิ ธมฺมลทฺเธหิ, เย เต สมณพฺราหฺมณา มทปฺปมาทา ปฏิวิรตา ขนฺติโสรจฺเจ นิวิฏฺฐา เอกมตฺตานํ ทเมนฺติ, เอกมตฺตานํ สเมนฺติ, เอกมตฺตานํ ปรินิพฺพาเปนฺติ, ตถารูเปสุ สมณพฺราหฺมเณสุ อุทฺธคฺคิกํ ทกฺขิณํ ปติฏฺฐาเปติ โสวคฺคิกํ สุขวิปากํ สคฺคสํวตฺตนิกํ.\csromanspacer{} อยํ ปญฺจโม โภคานํ อาทิโย.\csromanspacer{} อิเม โข คหปติ ปญฺจ โภคานํ อาทิยา.}

\prose{ตสฺส เจ คหปติ อริยสาวกสฺส อิเม ปญฺจ โภคานํ อาทิเย อาทิยโต โภคา ปริกฺขยํ คจฺฉนฺติ.\csromanspacer{} ตสฺส เอวํ โหติ “เย วต โภคานํ อาทิยา, เต จาหํ อาทิยามิ, โภคา จ เม ปริกฺขยํ คจฺฉนฺตี”ติ.\csromanspacer{} อิติสฺส โหติ อวิปฺปฏิสาโร, ตสฺส เจ คหปติ อริยสาวกสฺส อิเม ปญฺจ โภคานํ อาทิเย อาทิยโต โภคา \csromanfolio{40}อภิวฑฺฒนฺติ.\csromanspacer{} ตสฺส เอวํ โหติ “เย วต โภคานํ อาทิยา, เต จาหํ อาทิยามิ, โภคา จ เม อภิวฑฺฒนฺตี”ติ.\csromanspacer{} อิติสฺส โหติ\footnote{อิติสฺส โหติ อวิปฺปฏิสาโร, (สี, สฺยา)} อุภเยเนว อวิปฺปฏิสาโรติ.}

\prose{ภุตฺตา โภคา ภตา ภจฺจา\footnote{คตา ตจฺฉา (ก)}, วิติณฺณา อาปทาสุ เม.}

\prose{อุทฺธคฺคา ทกฺขิณา ทินฺนา, อโถ ปญฺจพลีกตา.}

\prose{อุปฏฺฐิตา สีลวนฺโต, สญฺญตา พฺรหฺมจารโย.}

\prose{ยทตฺถํ โภคํ อิจฺเฉยฺย, ปณฺฑิโต ฆรมาวสํ.}

\prose{โส เม อตฺโถ อนุปฺปตฺโต, กตํ อนนุตาปิยํ.}

\prose{เอตํ\footnote{เอวํ (ก)} อนุสฺสรํ มจฺโจ, อริยธมฺเม ฐิโต นโร.}

\prose{อิเธว นํ ปสํสนฺติ, เปจฺจ สคฺเค ปโมทตีติ\footnote{เปจฺจ สคฺเค จ โมทตีติ (สี, สฺยา)}.\csromanspacer{}\csromanspacer{}ปฐมํ.}

\csromanheader{2}{2. สปฺปุริสสุตฺต}
\proseitem{42}{สปฺปุริโส ภิกฺขเว กุเล ชายมาโน พหุโน ชนสฺส อตฺถาย หิตาย สุขาย โหติ, มาตาปิตูนํ\footnote{มาตาปิตุนฺนํ (สี, อิ)} อตฺถาย หิตาย สุขาย โหติ, ปุตฺตทารสฺส อตฺถาย หิตาย สุขาย โหติ, ทาสกมฺมกรโปริสสฺส อตฺถาย หิตาย สุขาย โหติ, มิตฺตามจฺจานํ อตฺถาย หิตาย สุขาย โหติ, สมณพฺราหฺมณานํ อตฺถาย หิตาย สุขาย โหติ.}

\prose{เสยฺยถาปิ ภิกฺขเว มหาเมโฆ สพฺพสสฺสานิ สมฺปาเทนฺโต พหุโน ชนสฺส อตฺถาย หิตาย สุขาย โหติ.\csromanspacer{} เอวเมวํ โข ภิกฺขเว สปฺปุริโส กุเล ชายมาโน พหุโน ชนสฺส อตฺถาย หิตาย สุขาย โหติ, มาตาปิตูนํ อตฺถาย หิตาย สุขาย โหติ, ปุตฺตทารสฺส อตฺถาย หิตาย สุขาย โหติ, ทาสกมฺมกรโปริสสฺส อตฺถาย หิตาย สุขาย โหติ, มิตฺตามจฺจานํ อตฺถาย หิตาย สุขาย โหติ, สมณพฺราหฺมณานํ อตฺถาย หิตาย สุขาย โหตีติ.}

\prose{หิโต พหุนฺนํ ปฏิปชฺช โภเค, ตํ เทวตา รกฺขติ ธมฺมคุตฺตํ.}

\prose{พหุสฺสุตํ สีลวตูปปนฺนํ, ธมฺเม ฐิตํ น วิชหติ\footnote{วิชหาติ (สี, สฺยา, กํ, อิ)} กิตฺติ.}

\chapterheadpageread{5. มุณฺฑราชวคฺค}
\prose{\csromanfolio{41}ธมฺมฏฺฐํ สีลสมฺปนฺนํ, สจฺจวาทิํ หิรีมนํ.}

\prose{เนกฺขํ ชมฺโพนทสฺเสว, โก ตํ นินฺทิตุมรหติ.}

\prose{เทวาปิ นํ ปสํสนฺติ, พฺรหฺมุนาปิ ปสํสิโตติ.\csromanspacer{}\csromanspacer{}ทุติยํ.}

\csromanheader{2}{3. อิฏฺฐสุตฺต}
\proseitem{43}{อถ โข อนาถปิณฺฑิโก คหปติ เยน ภควา เตนุปสงฺกมิ, อุปสงฺกมิตฺวา ภควนฺตํ อภิวาเทตฺวา เอกมนฺตํ นิสีทิ, เอกมนฺตํ นิสินฺนํ โข อนาถปิณฺฑิกํ คหปติํ ภควา เอตทโวจ\csromandash{}}

\prose{ปญฺจิเม คหปติ ธมฺมา อิฏฺฐา กนฺตา มนาปา ทุลฺลภา โลกสฺมิํ.\csromanspacer{} กตเม ปญฺจ? อายุ คหปติ อิฏฺโฐ กนฺโต มนาโป ทุลฺลโภ โลกสฺมิํ, วณฺโณ อิฏฺโฐ กนฺโต มนาโป ทุลฺลโภ โลกสฺมิํ, สุขํ อิฏฺฐํ กนฺตํ มนาปํ ทุลฺลภํ โลกสฺมิํ, ยโส อิฏฺโฐ กนฺโต มนาโป ทุลฺลโภ โลกสฺมิํ, สคฺคา อิฏฺฐา กนฺตา มนาปา ทุลฺลภา โลกสฺมิํ.\csromanspacer{} อิเม โข คหปติ ปญฺจ ธมฺมา กนฺตา มนาปา ทุลฺลภา โลกสฺมิํ.}

\prose{อิเมสํ โข คหปติ ปญฺจนฺนํ ธมฺมานํ อิฏฺฐานํ กนฺตานํ มนาปานํ ทุลฺลภานํ โลกสฺมิํ น อายาจนเหตุ วา ปตฺถนาเหตุ วา\footnote{น ปตฺถนาเหตุ วา (สฺยา, กํ, อิ)} ปฏิลาภํ วทามิ.\csromanspacer{} อิเมสํ โข คหปติ ปญฺจนฺนํ ธมฺมานํ อิฏฺฐานํ กนฺตานํ มนาปานํ ทุลฺลภานํ โลกสฺมิํ อายาจนเหตุ วา ปตฺถนาเหตุ วา ปฏิลาโภ อภวิสฺส, โก อิธ เกน หาเยถ.}

\prose{น โข คหปติ อรหติ อริยสาวโก อายุกาโม อายุํ อายาจิตุํ วา อภินนฺทิตุํ วา, อายุสฺส วาปิ เหตุ อายุกาเมน คหปติ อริยสาวเกน อายุสํวตฺตนิกา ปฏิปทา ปฏิปชฺชิตพฺพา, อายุสํวตฺตนิกา หิสฺส ปฏิปทา ปฏิปนฺนา อายุปฏิลาภาย สํวตฺตติ, โส ลาภี โหติ อายุสฺส ทิพฺพสฺส วา มานุสสฺส วา.}

\prose{น โข คหปติ อรหติ อริยสาวโก วณฺณกาโม วณฺณํ อายาจิตุํ วา อภินนฺทิตุํ วา, วณฺณสฺส วาปิ เหตุ วณฺณกาเมน คหปติ อริยสาวเกน วณฺณสํวตฺตนิกา ปฏิปทา ปฏิปชฺชิตพฺพา, วณฺณสํวตฺตนิกา หิสฺส ปฏิปทา ปฏิปนฺนา วณฺณปฏิลาภาย สํวตฺตติ, โส ลาภี โหติ วณฺณสฺส ทิพฺพสฺส วา มานุสสฺส วา.}

\prose{\csromanfolio{42}น โข คหปติ อรหติ อริยสาวโก สุขกาโม สุขํ อายาจิตุํ วา อภินนฺทิตุํ วา, สุขสฺส วาปิ เหตุ สุขกาเมน คหปติ อริยสาวเกน สุขสํวตฺตนิกา ปฏิปทา ปฏิปชฺชิตพฺพา, สุขสํวตฺตนิกา หิสฺส ปฏิปทา ปฏิปนฺนา สุขปฏิลาภาย สํวตฺตติ, โส ลาภี โหติ สุขสฺส ทิพฺพสฺส วา มานุสสฺส วา.}

\prose{น โข คหปติ อรหติ อริยสาวโก ยสกาโม ยสํ อายาจิตุํ วา อภินนฺทิตุํ วา, ยสสฺส วาปิ เหตุ ยสกาเมน คหปติ อริยสาวเกน ยสสํวตฺตนิกา ปฏิปทา ปฏิปชฺชิตพฺพา, ยสสํวตฺตนิกา หิสฺส ปฏิปทา ปฏิปนฺนา ยสปฏิลาภาย สํวตฺตติ, โส ลาภี โหติ ยสสฺส ทิพฺพสฺส วา มานุสสฺส วา.}

\prose{น โข คหปติ อรหติ อริยสาวโก สคฺคกาโม สคฺคํ อายาจิตุํ วา อภินนฺทิตุํ วา, สคฺคานํ วาปิ เหตุ สคฺคกาเมน คหปติ อริยสาวเกน สคฺคสํวตฺตนิกา ปฏิปทา ปฏิปชฺชิตพฺพา, สคฺคสํวตฺตนิกา หิสฺส ปฏิปทา ปฏิปนฺนา สคฺคปฏิลาภาย สํวตฺตติ, โส ลาภี โหติ สคฺคานนฺติ.}

\csromangathameasure{อายุํ วณฺณํ ยสํ กิตฺติํ, สคฺคํ อุจฺจากุลีนตํ. \\ รติโย ปตฺถยาเนน, อุฬารา อปราปรา. \\ อปฺปมาทํ ปสํสนฺติ, ปุญฺญกิริยาสุ ปณฺฑิตา. \\ อปฺปมตฺโต อุโภ อตฺเถ, อธิคณฺหาติ ปณฺฑิโต. \\ ทิฏฺเฐ ธมฺเม จ โย อตฺโถ, โย จตฺโถ สมฺปรายิโก. \\ อตฺถาภิสมยา ธีโร, “ปณฺฑิโต”ติ ปวุจฺจตีติ.ตติยํ.}
\csromangathagroup{\csromangathastanza{อายุํ วณฺณํ ยสํ กิตฺติํ, สคฺคํ อุจฺจากุลีนตํ. \\ รติโย ปตฺถยาเนน\footnote{ปตฺถยมาเนน (ก)}, อุฬารา อปราปรา.}\csromangathastanza{อปฺปมาทํ ปสํสนฺติ, ปุญฺญกิริยาสุ ปณฺฑิตา. \\ อปฺปมตฺโต อุโภ อตฺเถ, อธิคณฺหาติ ปณฺฑิโต.}\csromangathastanza{ทิฏฺเฐ ธมฺเม จ\footnote{ทิฏฺเฐว ธมฺเม (สี)} โย อตฺโถ, โย จตฺโถ สมฺปรายิโก. \\ อตฺถาภิสมยา ธีโร, “ปณฺฑิโต”ติ ปวุจฺจตีติ.\csromanspacer{}\csromanspacer{}ตติยํ.}}

\csromanheader{2}{4. มนาปทายีสุตฺต}
\proseitem{44}{เอกํ สมยํ ภควา เวสาลิยํ วิหรติ มหาวเน กูฏาคารสาลายํ.\csromanspacer{} อถ โข ภควา ปุพฺพณฺหสมยํ นิวาเสตฺวา ปตฺตจีวรมาทาย เยน อุคฺคสฺส คหปติโน เวสาลิกสฺส นิเวสนํ เตนุปสงฺกมิ, อุปสงฺกมิตฺวา ปญฺญตฺเต อาสเน นิสีทิ.\csromanspacer{} อถ โข อุคฺโค คหปติ เวสาลิโก เยน ภควา เตนุปสงฺกมิ, อุปสงฺกมิตฺวา ภควนฺตํ}

\vspace{\tipitakaparskip}
\csromancenter{มุณฺฑราชวคฺค อภิวาเทตฺวา เอกมนฺตํ นิสีทิ, เอกมนฺตํ นิสินฺโน โข อุคฺโค คหปติ เวสาลิโก ภควนฺตํ เอตทโวจ\csromandash{}}
\prose{\csromanfolio{43}สมฺมุขา เมตํ ภนฺเต ภควโต สุตํ สมฺมุขา ปฏิคฺคหิตํ “มนาปทายี ลภเต มนาปนฺ”ติ.\csromanspacer{} มนาปํ เม ภนฺเต สาลปุปฺผกํ\footnote{สาลิปุปฺผกํ (ก)} ขาทนียํ, ตํ เม ภควา ปฏิคฺคณฺหาตุ อนุกมฺปํ อุปาทายาติ.\csromanspacer{} ปฏิคฺคเหสิ ภควา อนุกมฺปํ อุปาทาย.}

\prose{สมฺมุขา เมตํ ภนฺเต ภควโต สุตํ สมฺมุขา ปฏิคฺคหิตํ “มนาปทายี ลภเต มนาปนฺ”ติ.\csromanspacer{} มนาปํ เม ภนฺเต สมฺปนฺนโกลกํ สูกรมํสํ\footnote{สมฺปนฺนสูกรมํสํ (สี), สมฺปนฺนวรสูกรมํสํ (สฺยา)}, ตํ เม ภควา ปฏิคฺคณฺหาตุ อนุกมฺปํ อุปาทายาติ.\csromanspacer{} ปฏิคฺคเหสิ ภควา อนุกมฺปํ อุปาทาย.}

\prose{สมฺมุขา เมตํ ภนฺเต ภควโต สุตํ สมฺมุขา ปฏิคฺคหิตํ “มนาปทายี ลภเต มนาปนฺ”ติ.\csromanspacer{} มนาปํ เม ภนฺเต นิพฺพตฺตเตลกํ\footnote{นิพทฺธเตลกํ (อิ, สี-ฏฺฐ), นิพฺพฏฺฏเตลกํ (สฺยา-ฏฺฐ)} นาลิยสากํ, ตํ เม ภควา ปฏิคฺคณฺหาตุ อนุกมฺปํ อุปาทายาติ.\csromanspacer{} ปฏิคฺคเหสิ ภควา อนุกมฺปํ อุปาทาย.}

\prose{สมฺมุขา เมตํ ภนฺเต ภควโต สุตํ สมฺมุขา ปฏิคฺคหิตํ “มนาปทายี ลภเต มนาปนฺ”ติ.\csromanspacer{} มนาโป เม ภนฺเต สาลีนํ โอทโน วิจิตกาฬโก\footnote{วิคตกาลโก (สฺยา, กํ, อิ, ก)} อเนกสูโป อเนกพฺยญฺชโน, ตํ เม ภควา ปฏิคฺคณฺหาตุ อนุกมฺปํ อุปาทายาติ.\csromanspacer{} ปฏิคฺคเหสิ ภควา อนุกมฺปํ อุปาทาย.}

\prose{สมฺมุขา เมตํ ภนฺเต ภควโต สุตํ สมฺมุขา ปฏิคฺคหิตํ “มนาปทายี ลภเต มนาปนฺ”ติ.\csromanspacer{} มนาปานิ เม ภนฺเต กาสิกานิ วตฺถานิ, ตานิ เม ภควา ปฏิคฺคณฺหาตุ อนุกมฺปํ อุปาทายาติ.\csromanspacer{} ปฏิคฺคเหสิ ภควา อนุกมฺปํ อุปาทาย.}

\prose{สมฺมุขา เมตํ ภนฺเต ภควโต สุตํ สมฺมุขา ปฏิคฺคหิตํ “มนาปทายี ลภเต มนาปนฺ”ติ.\csromanspacer{} มนาโป เม ภนฺเต ปลฺลงฺโก โคนกตฺถโต ปฏลิกตฺถโต กทลิมิคปวรปจฺจตฺถรโณ\footnote{กาทลิมิคปวรปจฺจตฺถรโณ (สี)} ส-อุตฺตรจฺฉโท อุภโตโลหิตกูปธาโน, อปิ จ ภนฺเต มยมฺเปตํ ชานาม “เนตํ ภควโต กปฺปตี”ติ.\csromanspacer{} อิทํ เม ภนฺเต จนฺทนผลกํ อคฺฆติ อธิกสตสหสฺสํ, ตํ เม \csromanfolio{44}ภควา ปฏิคฺคณฺหาตุ อนุกมฺปํ อุปาทายาติ.\csromanspacer{} ปฏิคฺคเหสิ ภควา อนุกมฺปํ อุปาทาย.\csromanspacer{} อถ โข ภควา อุคฺคํ คหปติํ เวสาลิกํ อิมินา อนุโมทนีเยน อนุโมทิ\csromandash{}}

\csromangathameasure{“มนาปทายี ลภเต มนาปํ, \\ โย อุชฺชุภูเตสุ ททาติ ฉนฺทสา. \\ อจฺฉาทนํ สยนมนฺนปานํ, \\ นานาปฺปการานิ จ ปจฺจยานิ. \\ จตฺตญฺจ มุตฺตญฺจ อนุคฺคหีตํ, \\ เขตฺตูปเม อรหนฺเต วิทิตฺวา. \\ โส ทุจฺจชํ สปฺปุริโส จชิตฺวา, \\ มนาปทายี ลภเต มนาปนฺ”ติ.}
\csromangathagroup{\csromangathastanza{“มนาปทายี ลภเต มนาปํ, \\ โย อุชฺชุภูเตสุ\footnote{อุชุภูเตสุ (สฺยา, กํ, อิ)} ททาติ ฉนฺทสา. \\ อจฺฉาทนํ สยนมนฺนปานํ\footnote{อุชุภูเตสุ (สฺยา, กํ, อิ)}, \\ นานาปฺปการานิ จ ปจฺจยานิ.}\csromangathastanza{จตฺตญฺจ มุตฺตญฺจ อนุคฺคหีตํ\footnote{อนคฺคหีตํ (สี, สฺยา, กํ, อิ)}, \\ เขตฺตูปเม อรหนฺเต วิทิตฺวา. \\ โส ทุจฺจชํ สปฺปุริโส จชิตฺวา, \\ มนาปทายี ลภเต มนาปนฺ”ติ.}}

\prosecont{อถ โข ภควา อุคฺคํ คหปติํ เวสาลิกํ อิมินา อนุโมทนีเยน อนุโมทิตฺวา อุฏฺฐายาสนา ปกฺกามิ.}

\prose{อถ โข อุคฺโค คหปติ เวสาลิโก อปเรน สมเยน กาลมกาสิ, กาลงฺกโต\footnote{กาลกโต (สี, สฺยา, กํ, อิ)} จ อุคฺโค คหปติ เวสาลิโก อญฺญตรํ มโนมยํ กายํ อุปปชฺชิ.\csromanspacer{} เตน โข ปน สมเยน ภควา สาวตฺถิยํ วิหรติ เชตวเน อนาถปิณฺฑิกสฺส อาราเม.\csromanspacer{} อถ โข อุคฺโค เทวปุตฺโต อภิกฺกนฺตาย รตฺติยา อภิกฺกนฺตวณฺโณ เกวลกปฺปํ เชตวนํ โอภาเสตฺวา เยน ภควา เตนุปสงฺกมิ, อุปสงฺกมิตฺวา ภควนฺตํ อภิวาเทตฺวา เอกมนฺตํ อฏฺฐาสิ, เอกมนฺตํ ฐิตํ โข อุคฺคํ เทวปุตฺตํ ภควา เอตทโวจ “กจฺจิ เต อุคฺค ยถาธิปฺปาโย”ติ.\csromanspacer{} ตคฺฆ เม ภควา ยถาธิปฺปาโยติ.\csromanspacer{} อถ โข ภควา อุคฺคํ เทวปุตฺตํ คาถาหิ อชฺฌภาสิ\csromandash{}}

\csromangathameasure{“มนาปทายี ลภเต มนาปํ, \\ อคฺคสฺส ทาตา ลภเต ปุนคฺคํ. \\ วรสฺส ทาตา วรลาภิ โหติ, \\ เสฏฺฐํ ทโท เสฏฺฐมุเปติ ฐานํ.}
\csromangathagroup{\csromangathastanza{“มนาปทายี ลภเต มนาปํ, \\ อคฺคสฺส ทาตา ลภเต ปุนคฺคํ. \\ วรสฺส ทาตา วรลาภิ โหติ, \\ เสฏฺฐํ ทโท เสฏฺฐมุเปติ ฐานํ.}}

\chapterheadpageread{5. มุณฺฑราชวคฺค}
\csromangathameasure{โย อคฺคทายี วรทายี, เสฏฺฐทายี จ โย นโร. \\ ทีฆายุ ยสวา โหติ, ยตฺถ ยตฺถูปปชฺชตี”ติ.จตุตฺถํ.}
\csromangathagroup{\csromangathastanza{\csromanfolio{45}โย อคฺคทายี วรทายี, เสฏฺฐทายี จ โย นโร. \\ ทีฆายุ ยสวา โหติ, ยตฺถ ยตฺถูปปชฺชตี”ติ.\csromanspacer{}\csromanspacer{}จตุตฺถํ.}}

\csromanheader{2}{5. ปุญฺญาภิสนฺทสุตฺต}
\proseitem{45}{ปญฺจิเม ภิกฺขเว ปุญฺญาภิสนฺทา กุสลาภิสนฺทา สุขสฺสาหารา โสวคฺคิกา สุขวิปากา สคฺคสํวตฺตนิกา อิฏฺฐาย กนฺตาย มนาปาย หิตาย สุขาย สํวตฺตนฺติ.}

\prose{กตเม ปญฺจ? ยสฺส ภิกฺขเว ภิกฺขุ จีวรํ ปริภุญฺชมาโน อปฺปมาณํ เจโตสมาธิํ อุปสมฺปชฺช วิหรติ, อปฺปมาโณ ตสฺส ปุญฺญาภิสนฺโท กุสลาภิสนฺโท สุขสฺสาหาโร โสวคฺคิโก สุขวิปาโก สคฺคสํวตฺตนิโก อิฏฺฐาย กนฺตาย มนาปาย หิตาย สุขาย สํวตฺตติ.}

\prose{ยสฺส ภิกฺขเว ภิกฺขุ ปิณฺฑปาตํ ปริภุญฺชมาโน ฯเปฯ.\csromanspacer{} ยสฺส ภิกฺขเว ภิกฺขุ วิหารํ ปริภุญฺชมาโน ฯเปฯ.\csromanspacer{} ยสฺส ภิกฺขเว ภิกฺขุ มญฺจปีฐํ ปริภุญฺชมาโน ฯเปฯ.}

\prose{ยสฺส ภิกฺขเว ภิกฺขุ คิลานปจฺจยเภสชฺชปริกฺขารํ ปริภุญฺชมาโน อปฺปมาณํ เจโตสมาธิํ อุปสมฺปชฺช วิหรติ, อปฺปมาโณ ตสฺส ปุญฺญาภิสนฺโท กุสลาภิสนฺโท สุขสฺสาหาโร โสวคฺคิโก สุขวิปาโก สคฺคสํวตฺตนิโก อิฏฺฐาย กนฺตาย มนาปาย หิตาย สุขาย สํวตฺตติ.\csromanspacer{} อิเม โข ภิกฺขเว ปญฺจ ปุญฺญาภิสนฺทา กุสลาภิสนฺทา สุขสฺสาหารา โสวคฺคิกา สุขวิปากา สคฺคสํวตฺตนิกา อิฏฺฐาย กนฺตาย มนาปาย หิตาย สุขาย สํวตฺตนฺติ.}

\prose{อิเมหิ จ ปน ภิกฺขเว ปญฺจหิ ปุญฺญาภิสนฺเทหิ กุสลาภิสนฺเทหิ สมนฺนาคตสฺส อริยสาวกสฺส น สุกรํ ปุญฺญสฺส ปธาณํ คเหตุํ “เอตฺตโก ปุญฺญาภิสนฺโท กุสลาภิสนฺโท สุขสฺสาหาโร โสวคฺคิโก สุขวิปาโก สคฺคสํวตฺตนิโก อิฏฺฐาย กนฺตาย มนาปาย หิตาย สุขาย สํวตฺตตี”ติ, อถ โข อสงฺเขยฺโย อปฺปเมยฺโย มหาปุญฺญกฺขนฺโธเตฺวว สงฺขํ คจฺฉติ.}

\prose{เสยฺยถาปิ ภิกฺขเว มหาสมุทฺเท น สุกรํ อุทกสฺส ปมาณํ คเหตุํ “เอตฺตกานิ อุทกาฬฺหกานี”ติ วา “เอตฺตกานิ อุทกาฬฺหกสตานี”ติ \csromanfolio{46}วา “เอตฺตกานิ อุทกาฬฺหกสหสฺสานี”ติ วา “เอตฺตกานิ อุทกาฬฺหกสตสหสฺสานี”ติ วา, อถ โข อสงฺเขยฺโย อปฺปเมยฺโย มหา-อุทกกฺขนฺโธเตฺวว สงฺขํ คจฺฉติ.\csromanspacer{} เอวเมวํ โข ภิกฺขเว อิเมหิ ปญฺจหิ ปุญฺญาภิสนฺเทหิ กุสลาภิสนฺเทหิ สมนฺนาคตสฺส อริยสาวกสฺส น สุกรํ ปุญฺญสฺส ปมาณํ คเหตุํ “เอตฺตโก ปุญฺญาภิสนฺโท กุสลาภิสนฺโท สุขสฺสาหาโร โสวคฺคิโก สุขวิปาโก สคฺคสํวตฺตนิโก อิฏฺฐาย กนฺตาย มนาปาย หิตาย สุขาย สํวตฺตตี”ติ, อถ โข อสงฺเขยฺโย อปฺปเมยฺโย มหาปุญฺญกฺขนฺโธเตฺวว สงฺขํ คจฺฉตีติ.}

\csromangathameasure{มโหทธิํ อปริมิตํ มหาสรํ, \\ พหุเภรวํ รตฺนคณานมาลยํ. \\ นชฺโช ยถา นรคณสํฆเสวิตา, \\ ปุถู สวนฺตี อุปยนฺติ สาครํ. \\ เอวํ นรํ อนฺนทปานวตฺถทํ, \\ เสยฺยานิสชฺชตฺถรณสฺส ทายกํ. \\ ปุญฺญสฺส ธารา อุปยนฺติ ปณฺฑิตํ, \\ นชฺโช ยถา วาริวหาว สาครนฺติ.ปญฺจมํ.}
\csromangathagroup{\csromangathastanza{มโหทธิํ อปริมิตํ มหาสรํ, \\ พหุเภรวํ รตฺนคณานมาลยํ. \\ นชฺโช ยถา นรคณสํฆเสวิตา\footnote{มจฺฉคณสํฆเสวิตา (สฺยา, กํ, ก) สํ 3. 349 ปิฏฺเฐปิ ปสฺสิตพฺพํ.}, \\ ปุถู สวนฺตี อุปยนฺติ สาครํ.}\csromangathastanza{เอวํ นรํ อนฺนทปานวตฺถทํ, \\ เสยฺยานิสชฺชตฺถรณสฺส ทายกํ. \\ ปุญฺญสฺส ธารา อุปยนฺติ ปณฺฑิตํ, \\ นชฺโช ยถา วาริวหาว สาครนฺติ.\csromanspacer{}\csromanspacer{}ปญฺจมํ.}}

\csromanheader{2}{6. สมฺปทาสุตฺต}
\proseitem{46}{ปญฺจิมา ภิกฺขเว สมฺปทา.\csromanspacer{} กตมา ปญฺจ? สทฺธาสมฺปทา สีลสมฺปทา สุตสมฺปทา จาคสมฺปทา ปญฺญาสมฺปทา.\csromanspacer{} อิมา โข ภิกฺขเว ปญฺจ สมฺปทาติ.\csromanspacer{}\csromanspacer{}ฉฏฺฐํ.}

\csromanheader{2}{7. ธนสุตฺต}
\proseitem{47}{ปญฺจิมานิ ภิกฺขเว ธนานิ.\csromanspacer{} กตมานิ ปญฺจ? สทฺธาธนํ สีลธนํ สุตธนํ จาคธนํ ปญฺญาธนํ.}

\prose{กตมญฺจ ภิกฺขเว สทฺธาธนํ? อิธ ภิกฺขเว อริยสาวโก สทฺโธ โหติ, สทฺทหติ ตถาคตสฺส โพธิํ “อิติปิ โส ภควา ฯเปฯ สตฺถา เทวมนุสฺสานํ พุทฺโธ ภควา”ติ.\csromanspacer{} อิทํ วุจฺจติ ภิกฺขเว สทฺธาธนํ.}

\chapterheadpageread{5. มุณฺฑราชวคฺค}
\prose{\csromanfolio{47}กตมญฺจ ภิกฺขเว สีลธนํ? อิธ ภิกฺขเว อริยสาวโก ปาณาติปาตา ปฏิวิรโต โหติ ฯเปฯ สุราเมรยมชฺชปมาทฏฺฐานา ปฏิวิรโต โหติ.\csromanspacer{} อิทํ วุจฺจติ ภิกฺขเว สีลธนํ.}

\prose{กตมญฺจ ภิกฺขเว สุตธนํ? อิธ ภิกฺขเว อริยสาวโก พหุสฺสุโต โหติ ฯเปฯ ทิฏฺฐิยา สุปฺปฏิวิทฺโธ.\csromanspacer{} อิทํ วุจฺจติ ภิกฺขเว สุตธนํ.}

\prose{กตมญฺจ ภิกฺขเว จาคธนํ? อิธ ภิกฺขเว อริยสาวโก วิคตมลมจฺเฉเรน เจตสา อคารํ อชฺฌาวสติ มุตฺตจาโค ปยตปาณิ โวสคฺครโต ยาจโยโค ทานสํวิภาครโต.\csromanspacer{} อิทํ วุจฺจติ ภิกฺขเว จาคธนํ.}

\prose{กตมญฺจ ภิกฺขเว ปญฺญาธนํ? อิธ ภิกฺขเว อริยสาวโก ปญฺญาวา โหติ อุทยตฺถคามินิยา ปญฺญาย สมนฺนาคโต อริยาย นิพฺเพธิกาย สมฺมา ทุกฺขกฺขยคามินิยา.\csromanspacer{} อิทํ วุจฺจติ ภิกฺขเว ปญฺญาธนํ.\csromanspacer{} อิมานิ โข ภิกฺขเว ปญฺจ ธนานีติ.}

\prose{\csromansymbolfootnote{*}{อํ 1. 368 ปิฏฺเฐปิ.}ยสฺส สทฺธา ตถาคเต, อจลา สุปฺปติฏฺฐิตา.}

\csromangathameasure{สีลญฺจ ยสฺส กลฺยาณํ, อริยกนฺตํ ปสํสิตํ. \\ สํเฆ ปสาโท ยสฺสตฺถิ, อุชุภูตญฺจ ทสฺสนํ. \\ “อทลิทฺโท”ติ ตํ อาหุ, อโมฆํ ตสฺส ชีวิตํ. \\ ตสฺมา สทฺธญฺจ สีลญฺจ, ปสาทํ ธมฺมทสฺสนํ. \\ อนุยุญฺเชถ เมธาวี, สรํ พุทฺธาน สาสนนฺติ.สตฺตมํ.}
\csromangathagroup{\csromangathastanza{สีลญฺจ ยสฺส กลฺยาณํ, อริยกนฺตํ ปสํสิตํ. \\ สํเฆ ปสาโท ยสฺสตฺถิ, อุชุภูตญฺจ ทสฺสนํ.}\csromangathastanza{“อทลิทฺโท”ติ ตํ อาหุ, อโมฆํ ตสฺส ชีวิตํ. \\ ตสฺมา สทฺธญฺจ สีลญฺจ, ปสาทํ ธมฺมทสฺสนํ. \\ อนุยุญฺเชถ เมธาวี, สรํ พุทฺธาน สาสนนฺติ.\csromanspacer{}\csromanspacer{}สตฺตมํ.}}

\csromanheader{2}{8. อลพฺภนียฐานสุตฺต}
\proseitem{48}{ปญฺจิมานิ ภิกฺขเว อลพฺภนียานิ ฐานานิ สมเณน วา พฺราหฺมเณน วา เทเวน วา มาเรน วา พฺรหฺมุนา วา เกนจิ วา โลกสฺมิํ.\csromanspacer{} กตมานิ ปญฺจ? ชราธมฺมํ “มา ชีรี”ติ อลพฺภนียํ ฐานํ สมเณน วา พฺราหฺมเณน วา เทเวน วา มาเรน วา พฺรหฺมุนา วา เกนจิ วา โลกสฺมิํ.\csromanspacer{} พฺยาธิธมฺมํ “มา พฺยาธียี”ติ\footnote{ฌาธิธมฺมํ “มา ฌามียี”ติ (สี, อิ)} ฯเปฯ.\csromanspacer{} มรณธมฺมํ “มา มียี”ติ.\csromanspacer{} ขยธมฺมํ “มา ขียี”ติ.\csromanspacer{} นสฺสนธมฺมํ “มา นสฺสี”ติ อลพฺภนียํ ฐานํ สมเณน วา พฺราหฺมเณน วา เทเวน วา มาเรน วา พฺรหฺมุนา วา เกนจิ วา โลกสฺมิํ.}

\prose{\csromanfolio{48}อสฺสุตวโต ภิกฺขเว ปุถุชฺชนสฺส ชราธมฺมํ ชีรติ, โส ชราธมฺเม ชิณฺเณ น อิติ ปฏิสญฺจิกฺขติ “น โข มเยฺหเวกสฺส\footnote{มยฺหเมเวกสฺส (สี)} ชราธมฺมํ ชีรติ, อถ โข ยาวตา สตฺตานํ อาคติ คติ จุติ อุปปตฺติ, สพฺเพสํ สตฺตานํ ชราธมฺมํ ชีรติ.\csromanspacer{} อหญฺเจว\footnote{อหญฺเจ (?)} โข ปน ชราธมฺเม ชิณฺเณ โสเจยฺยํ กิลเมยฺยํ ปริเทเวยฺยํ, อุรตฺตาฬิํ กนฺเทยฺยํ, สมฺโมหํ อาปชฺเชยฺยํ, ภตฺตมฺปิ เม นจฺฉาเทยฺย, กาเยปิ ทุพฺพณฺณิยํ โอกฺกเมยฺย, กมฺมนฺตาปิ นปฺปวตฺเตยฺยุํ\footnote{กมฺมนฺโตปิ นปฺปวตฺเตยฺย (ก)}, อมิตฺตาปิ อตฺตมนา อสฺสุ, มิตฺตาปิ ทุมฺมนา อสฺสู”ติ.\csromanspacer{} โส ชราธมฺเม ชิณฺเณ โสจติ กิลมติ ปริเทวติ, อุรตฺตาฬิํ กนฺทติ, สมฺโมหํ อาปชฺชติ, อยํ วุจฺจติ ภิกฺขเว อสฺสุตวา ปุถุชฺชโน วิทฺโธ สวิเสน โสกสลฺเลน อตฺตานํเยว ปริตาเปติ.}

\prose{ปุน จปรํ ภิกฺขเว อสฺสุตวโต ปุถุชฺชนสฺส พฺยาธิธมฺมํ พฺยาธียติ ฯเปฯ.\csromanspacer{} มรณธมฺมํ มียติ.\csromanspacer{} ขยธมฺมํ ขียติ.\csromanspacer{} นสฺสนธมฺมํ นสฺสติ, โส นสฺสนธมฺเม นฏฺเฐ น อิติ ปฏิสญฺจิกฺขติ “น โข มเยฺหเวกสฺส นสฺสนธมฺมํ นสฺสติ, อถ โข ยาวตา สตฺตานํ อาคติ คติ จุติ อุปปตฺติ, สพฺเพสํ สตฺตานํ นสฺสนธมฺมํ นสฺสติ, อหญฺเจว โข ปน นสฺสนธมฺเม นฏฺเฐ โสเจยฺยํ กิลเมยฺยํ ปริเทเวยฺยํ, อุรตฺตาฬิํ กนฺเทยฺยํ, สมฺโมหํ อาปชฺเชยฺยํ, ภตฺตมฺปิ เม นจฺฉาเทยฺย, กาเยปิ ทุพฺพณฺณิยํ โอกฺกเมยฺย, กมฺมนฺตาปิ นปฺปวตฺเตยฺยุํ, อมิตฺตาปิ อตฺตมนา อสฺสุ, มิตฺตาปิ ทุมฺมนา อสฺสู”ติ.\csromanspacer{} โส นสฺสนธมฺเม นฏฺเฐ โสจติ กิลมติ ปริเทวติ, อุรตฺตาฬิํ กนฺทติ, สมฺโมหํ อาปชฺชติ.\csromanspacer{} อยํ วุจฺจติ ภิกฺขเว อสฺสุตวา ปุถุชฺชโน วิทฺโธ สวิเสน โสกสลฺเลน อตฺตานํเยว ปริตาเปติ.}

\prose{สุตวโต จ โข ภิกฺขเว อริยสาวกสฺส ชราธมฺมํ ชีรติ, โส ชราธมฺเม ชิณฺเณ อิติ ปฏิสญฺจิกฺขติ “น โข มเยฺหเวกสฺส ชราธมฺมํ ชีรติ, อถ โข ยาวตา สตฺตานํ อาคติ คติ จุติ อุปปตฺติ, สพฺเพสํ สตฺตานํ ชราธมฺมํ ชีรติ.\csromanspacer{} อหญฺเจว โข ปน ชราธมฺเม ชิณฺเณ โสเจยฺยํ กิลเมยฺยํ ปริเทเวยฺยํ, อุรตฺตาฬิํ กนฺเทยฺยํ, สมฺโมหํ อาปชฺเชยฺยํ, ภตฺตมฺปิ เม นจฺฉาเทยฺย, กาเยปิ ทุพฺพณฺณิยํ โอกฺกเมยฺย, กมฺมนฺตาปิ นปฺปวตฺเตยฺยุํ, อมิตฺตาปิ อตฺตมนา อสฺสุ, มิตฺตาปิ ทุมฺมนา อสฺสู”ติ.\csromanspacer{} โส ชราธมฺเม ชิณฺเณ น โสจติ น กิลมติ น ปริเทวติ, น อุรตฺตาฬิํ กนฺทติ, น สมฺโมหํ}

\vspace{\tipitakaparskip}
\csromancenter{มุณฺฑราชวคฺค อาปชฺชติ.\csromanspacer{} อยํ วุจฺจติ ภิกฺขเว สุตวา อริยสาวโก อพฺพุหิ\footnote{อพฺพหิ (สี)} สวิสํ โสกสลฺลํ.\csromanspacer{} เยน วิทฺโธ อสฺสุตวา ปุถุชฺชโน อตฺตานํเยว ปริตาเปติ.\csromanspacer{} อโสโก วิสลฺโล อริยสาวโก อตฺตานํเยว ปรินิพฺพาเปติ.}
\prose{\csromanfolio{49}ปุน จปรํ ภิกฺขเว สุตวโต อริยสาวกสฺส พฺยาธิธมฺมํ พฺยาธียติ ฯเปฯ.\csromanspacer{} มรณธมฺมํ มียติ.\csromanspacer{} ขยธมฺมํ ขียติ.\csromanspacer{} นสฺสนธมฺมํ นสฺสติ, โส นสฺสนธมฺเม นฏฺเฐ อิติ ปฏิสญฺจิกฺขติ “น โข มเยฺหเวกสฺส นสฺสนธมฺมํ นสฺสติ, อถ โข ยาวตา สตฺตานํ อาคติ คติ จุติ อุปปตฺติ, สพฺเพสํ สตฺตานํ นสฺสนธมฺมํ นสฺสติ.\csromanspacer{} อหญฺเจว โข ปน นสฺสนธมฺเม นฏฺเฐ โสเจยฺยํ กิลเมยฺยํ ปริเทเวยฺยํ, อุรตฺตาฬิํ กนฺเทยฺยํ, สมฺโมหํ อาปชฺเชยฺยํ, ภตฺตมฺปิ เม นจฺฉาเทยฺย, กาเยปิ ทุพฺพณฺณิยํ โอกฺกเมยฺย, กมฺมนฺตาปิ นปฺปวตฺเตยฺยุํ, อมิตฺตาปิ อตฺตมนา อสฺสุ, มิตฺตาปิ ทุมฺมนา อสฺสู”ติ.\csromanspacer{} โส นสฺสนธมฺเม นฏฺเฐ น โสจติ น กิลมติ น ปริเทวติ, น อุรตฺตาฬิํ กนฺทติ, น สมฺโมหํ อาปชฺชติ.\csromanspacer{} อยํ วุจฺจติ ภิกฺขเว สุตวา อริยสาวโก อพฺพุหิ สวิสํ โสกสลฺลํ.\csromanspacer{} เยน วิทฺโธ อสฺสุตวา ปุถุชฺชโน อตฺตานํเยว ปริตาเปติ.\csromanspacer{} อโสโก วิสลฺโล อริยสาวโก อตฺตานํเยว ปรินิพฺพาเปตีติ.}

\prose{อิมานิ โข ภิกฺขเว ปญฺจ อลพฺภนียานิ ฐานานิ สมเณน วา พฺราหฺมเณน วา เทเวน วา มาเรน วา พฺรหฺมุนา วา เกนจิ วา โลกสฺมินฺติ.}

\csromangathameasure{*น โสจนาย ปริเทวนาย, \\ อตฺโถธ ลพฺภา อปิ อปฺปโกปิ. \\ โสจนฺตเมนํ ทุขิตํ วิทิตฺวา, \\ ปจฺจตฺถิกา อตฺตมนา ภวนฺติ. \\ ยโต จ โข ปณฺฑิโต อาปทาสุ, \\ น เวธตี อตฺถวินิจฺฉยญฺญู. \\ ปจฺจตฺถิกาสฺส ทุขิตา ภวนฺติ, \\ ทิสฺวา มุขํ อวิการํ ปุราณํ. \\ ชปฺเปน มนฺเตน สุภาสิเตน, \\ อนุปฺปทาเนน ปเวณิยา วา. \\ ยถา ยถา ยตฺถ ลเภถ อตฺถํ, \\ ตถา ตถา ตตฺถ ปรกฺกเมยฺย. \\ สเจ ปชาเนยฺย อลพฺภเนยฺโย, \\ มยา ว อญฺเญน วา เอส อตฺโถ. \\ อโสจมาโน อธิวาสเยยฺย, \\ กมฺมํ ทฬฺหํ กินฺติ กโรมิ ทานีติ.อฏฺฐมํ.}
\csromangathagroup{\csromangathastanza{\csromansymbolfootnote{*}{ขุ 5. 131 ชาตเกปิ.}น โสจนาย ปริเทวนาย, \\ อตฺโถธ ลพฺภา\footnote{อตฺโถ อิธ ลพฺภติ (สฺยา), อตฺโถ อิธ ลพฺภา (อิ)} อปิ อปฺปโกปิ. \\ โสจนฺตเมนํ ทุขิตํ วิทิตฺวา, \\ ปจฺจตฺถิกา อตฺตมนา ภวนฺติ.}\csromangathastanza{ยโต จ โข ปณฺฑิโต อาปทาสุ, \\ น เวธตี อตฺถวินิจฺฉยญฺญู. \\ ปจฺจตฺถิกาสฺส ทุขิตา ภวนฺติ, \\ ทิสฺวา มุขํ อวิการํ ปุราณํ.}\csromangathastanza{\csromanfolio{50}ชปฺเปน มนฺเตน สุภาสิเตน, \\ อนุปฺปทาเนน ปเวณิยา วา. \\ ยถา ยถา ยตฺถ\footnote{ยถา ยถา ยตฺถ ยตฺถ (ก)} ลเภถ อตฺถํ, \\ ตถา ตถา ตตฺถ ปรกฺกเมยฺย.}\csromangathastanza{สเจ ปชาเนยฺย อลพฺภเนยฺโย, \\ มยา ว\footnote{มยา วา (สฺยา, กํ, อิ)} อญฺเญน วา เอส อตฺโถ. \\ อโสจมาโน อธิวาสเยยฺย, \\ กมฺมํ ทฬฺหํ กินฺติ กโรมิ ทานีติ.\csromanspacer{}\csromanspacer{}อฏฺฐมํ.}}

\csromanheader{2}{9. โกสลสุตฺต}
\proseitem{49}{เอกํ สมยํ ภควา สาวตฺถิยํ วิหรติ เชตวเน อนาถปิณฺฑิกสฺส อาราเม.\csromanspacer{} อถ โข ราชา ปเสนทิ โกสโล เยน ภควา เตนุปสงฺกมิ, อุปสงฺกมิตฺวา ภควนฺตํ อภิวาเทตฺวา เอกมนฺตํ นิสีทิ.}

\prose{(เตน โข ปน สมเยน มลฺลิกา เทวี กาลงฺกตา โหติ.)\footnote{( ) นตฺถิ (สี)} อถ โข อญฺญตโร ปุริโส เยน ราชา ปเสนทิ โกสโล เตนุปสงฺกมิ, อุปสงฺกมิตฺวา รญฺโญ ปเสนทิสฺส โกสลสฺส อุปกณฺณเก อาโรเจสิ “มลฺลิกา เทวี เทว\footnote{เทว เทวี (สฺยา, กํ, อิ)} กาลงฺกตา”ติ.\csromanspacer{} เอวํ วุตฺเต ราชา ปเสนทิ โกสโล ทุกฺขี ทุมฺมโน ปตฺตกฺขนฺโธ อโธมุโข ปชฺฌายนฺโต อปฺปฏิภาโน นิสีทิ.}

\prose{อถ โข ภควา ราชานํ ปเสนทิํ โกสลํ ทุกฺขิํ ทุมฺมนํ ปตฺตกฺขนฺธํ อโธมุขํ ปชฺฌายนฺตํ อปฺปฏิภานํ วิทิตฺวา ราชานํ ปเสนทิํ โกสลํ เอตทโวจ\csromandash{}ปญฺจิมานิ มหาราช อลพฺภนียานิ ฐานานิ สมเณน วา พฺราหฺมเณน วา เทเวน วา มาเรน วา พฺรหฺมุนา วา เกนจิ วา โลกสฺมิํ.\csromanspacer{} กตมานิ ปญฺจ? ชราธมฺมํ “มา ชีรี”ติ อลพฺภนียํ ฐานํ ฯเปฯ น โสจนาย ปริเทวนาย ฯเปฯ กมฺมํ ทฬฺหํ กินฺติ กโรมิ ทานีติ.\csromanspacer{}\csromanspacer{}นวมํ.}

\csromanheader{2}{10. นารทสุตฺต}
\proseitem{15}{เอกํ สมยํ อายสฺมา นารโท ปาฏลิปุตฺเต วิหรติ กุกฺกุฏาราเม.\csromanspacer{} เตน โข ปน สมเยน มุณฺฑสฺส รญฺโญ ภทฺทา เทวี}

\vspace{\tipitakaparskip}
\csromancenter{มุณฺฑราชวคฺค กาลงฺกตา โหติ ปิยา มนาปา.\csromanspacer{} โส ภทฺทาย เทวิยา กาลงฺกตาย ปิยาย มนาปาย เนว นฺหายติ\footnote{นหายติ (สี, สฺยา, กํ, อิ)} น วิลิมฺปติ น ภตฺตํ ภุญฺชติ, น กมฺมนฺตํ ปโยเชติ, รตฺตินฺทิวํ\footnote{รตฺติทิวํ (ก)} ภทฺทาย เทวิยา สรีเร อชฺโฌมุจฺฉิโต.\csromanspacer{} อถ โข มุณฺโฑ ราชา ปิยกํ โกสารกฺขํ\footnote{โสการกฺขํ (สฺยา)} อามนฺเตสิ “เตน หิ สมฺม ปิยก ภทฺทาย เทวิยา สรีรํ อายสาย เตลโทณิยา ปกฺขิปิตฺวา อญฺญิสฺสา อายสาย โทณิยา ปฏิกุชฺชถ, ยถา มยํ ภทฺทาย เทวิยา สรีรํ จิรตรํ ปสฺเสยฺยามา”ติ. “เอวํ เทวา”ติ โข ปิยโก โกสารกฺโข มุณฺฑสฺส รญฺโญ ปฏิสฺสุตฺวา ภทฺทาย เทวิยา สรีรํ อายสาย เตลโทณิยา ปกฺขิปิตฺวา อญฺญิสฺสา อายสาย โทณิยา ปฏิกุชฺชิ.}
\prose{\csromanfolio{51}อถ โข ปิยกสฺส โกสารกฺขสฺส เอตทโหสิ “อิมสฺส โข มุณฺฑสฺส รญฺโญ ภทฺทา เทวี กาลงฺกตา ปิยา มนาปา, โส ภทฺทาย เทวิยา กาลงฺกตาย ปิยาย มนาปาย เนว นฺหายติ, น วิลิมฺปติ, น ภตฺตํ ภุญฺชติ, น กมฺมนฺตํ ปโยเชติ, รตฺตินฺทิวํ ภทฺทาย เทวิยา สรีเร อชฺโฌมุจฺฉิโต.\csromanspacer{} กํ\footnote{กิํ (ก)} นุ โข มุณฺโฑ ราชา สมณํ วา พฺราหฺมณํ วา ปยิรุปาเสยฺย, ยสฺส ธมฺมํ สุตฺวา โสกสลฺลํ ปชเหยฺยา”ติ.}

\prose{อถ โข ปิยกสฺส โกสารกฺขสฺส เอตทโหสิ “อยํ โข อายสฺมา นารโท ปาฏลิปุตฺเต วิหรติ กุกฺกุฏาราเม, ตํ โข ปนายสฺมนฺตํ นารทํ เอวํ กลฺยาโณ กิตฺติสทฺโท อพฺภุคฺคโต ‘ปณฺฑิโต วิยตฺโต\footnote{พฺยตฺโต (สี, อิ), พฺยตฺโต (สฺยา, กํ, ที 2. 253)} เมธาวี พหุสฺสุโต จิตฺตกถี กลฺยาณปฏิภาโน วุทฺโธ เจว\footnote{พุทฺโธ เจว (ก)} อรหา จ’\footnote{อรหา จา’ติ (?)}.\csromanspacer{} ยํนูน มุณฺโฑ ราชา อายสฺมนฺตํ นารทํ ปยิรุปาเสยฺย, อปฺเปว นาม มุณฺโฑ ราชา อายสฺมโต นารทสฺส ธมฺมํ สุตฺวา โสกสลฺลํ ปชเหยฺยา”ติ.}

\prose{อถ โข ปิยโก โกสารกฺโข เยน มุณฺโฑ ราชา เตนุปสงฺกมิ, อุปสงฺกมิตฺวา มุณฺฑํ ราชานํ เอตทโวจ “อยํ โข เทว อายสฺมา นารโท ปาฏลิปุตฺเต วิหรติ กุกฺกุฏาราเม, ตํ โข ปนายสฺมนฺตํ นารทํ เอวํ กลฺยาโณ กิตฺติสทฺโท อพฺภุคฺคโต ‘ปณฺฑิโต วิยตฺโต เมธาวี พหุสฺสุโต จิตฺตกถี กลฺยาณปฏิภาโน วุทฺโธ}

\prose{\csromanfolio{52}เจว อรหา จ’\footnote{อรหา จา’ติ (?)}.\csromanspacer{} ยทิ ปน เทโว อายสฺมนฺตํ นารทํ ปยิรุปาเสยฺย, อปฺเปว นาม เทโว อายสฺมโต นารทสฺส ธมฺมํ สุตฺวา โสกสลฺลํ ปชเหยฺยา”ติ.\csromanspacer{} เตน หิ สมฺม ปิยก อายสฺมนฺตํ นารทํ ปฏิเวเทหิ “กถํ หิ นาม มาทิโส สมณํ วา พฺราหฺมณํ วา วิชิเต วสนฺตํ ปุพฺเพ อปฺปฏิสํวิทิโต อุปสงฺกมิตพฺพํ มญฺเญยฺยา”ติ. “เอวํ เทวา”ติ โข ปิยโก โกสารกฺโข มุณฺฑสฺส รญฺโญ ปฏิสฺสุตฺวา เยนายสฺมา นารโท เตนุปสงฺกมิ, อุปสงฺกมิตฺวา อายสฺมนฺตํ นารทํ อภิวาเทตฺวา เอกมนฺตํ นิสีทิ, เอกมนฺตํ นิสินฺโน โข ปิยโก โกสารกฺโข อายสฺมนฺตํ นารทํ เอตทโวจ\csromandash{}}

\prose{“อิมสฺส ภนฺเต มุณฺฑสฺส รญฺโญ ภทฺทา เทวี กาลงฺกตา ปิยา มนาปา, โส ภทฺทาย เทวิยา กาลงฺกตาย ปิยาย มนาปาย เนว นฺหายติ, น วิลิมฺปติ, น ภตฺตํ ภุญฺชติ, น กมฺมนฺตํ ปโยเชติ, รตฺตินฺทิวํ ภทฺทาย เทวิยา สรีเร อชฺโฌมุจฺฉิโต.\csromanspacer{} สาธุ ภนฺเต อายสฺมา นารโท มุณฺฑสฺส รญฺโญ ตถา ธมฺมํ เทเสตุ, ยถา มุณฺโฑ ราชา อายสฺมโต นารทสฺส ธมฺมํ สุตฺวา โสกสลฺลํ ปชเหยฺยา”ติ.\csromanspacer{} ยสฺสทานิ ปิยก มุณฺโฑ ราชา กาลํ มญฺญตีติ.}

\prose{อถ โข ปิยโก โกสารกฺโข อุฏฺฐายาสนา อายสฺมนฺตํ นารทํ อภิวาเทตฺวา ปทกฺขิณํ กตฺวา เยน มุณฺโฑ ราชา เตนุปสงฺกมิ, อุปสงฺกมิตฺวา มุณฺฑํ ราชานํ เอตทโวจ “กตาวกาโส โข เทว อายสฺมตา นารเทน, ยสฺสทานิ เทโว กาลํ มญฺญตี”ติ.\csromanspacer{} เตน หิ สมฺม ปิยก ภทฺรานิ ภทฺรานิ ยานานิ โยชาเปหีติ. “เอวํ เทวา”ติ โข ปิยโก โกสารกฺโข มุณฺฑสฺส รญฺโญ ปฏิสฺสุตฺวา ภทฺรานิ ภทฺรานิ ยานานิ โยชาเปตฺวา มุณฺฑํ ราชานํ เอตทโวจ “ยุตฺตานิ โข เต เทว ภทฺรานิ ภทฺรานิ ยานานิ, ยสฺสทานิ เทโว กาลํ มญฺญตี”ติ.}

\prose{อถ โข มุณฺโฑ ราชา ภทฺรํ ยานํ\footnote{ภทฺรํ ภทฺรํ ยานํ (สฺยา, กํ, ก), ภทฺทํ ยานํ (อิ)} อภิรุหิตฺวา ภเทฺรหิ ภเทฺรหิ ยาเนหิ เยน กุกฺกุฏาราโม เตน ปายาสิ มหจฺจา\footnote{มหจฺจ (พหูสุ)} ราชานุภาเวน อายสฺมนฺตํ นารทํ ทสฺสนาย.\csromanspacer{} ยาวติกา ยานสฺส ภูมิ, ยาเนน คนฺตฺวา ยานา ปจฺโจโรหิตฺวา ปตฺติโกว อารามํ ปาวิสิ.\csromanspacer{} อถ โข มุณฺโฑ}

\vspace{\tipitakaparskip}
\csromancenter{มุณฺฑราชวคฺค ราชา เยน อายสฺมา นารโท เตนุปสงฺกมิ, อุปสงฺกมิตฺวา อายสฺมนฺตํ นารทํ อภิวาเทตฺวา เอกมนฺตํ นิสีทิ, เอกมนฺตํ นิสินฺนํ โข มุณฺฑํ ราชานํ อายสฺมา นารโท เอตทโวจ\csromandash{}}
\prose{\csromanfolio{53}ปญฺจิมานิ มหาราช อลพฺภนียานิ ฐานานิ สมเณน วา พฺราหฺมเณน วา เทเวน วา มาเรน วา พฺรหฺมุนา วา เกนจิ วา โลกสฺมิํ.\csromanspacer{} กตมานิ ปญฺจ? ชราธมฺมํ “มา ชีรี”ติ อลพฺภนียํ ฐานํ สมเณน วา พฺราหฺมเณน วา เทเวน วา มาเรน วา พฺรหฺมุนา วา เกนจิ วา โลกสฺมิํ, พฺยาธิธมฺมํ “มา พฺยาธียี”ติ.\csromanspacer{} มรณธมฺมํ “มา มียี”ติ.\csromanspacer{} ขยธมฺมํ “มา ขียี”ติ.\csromanspacer{} นสฺสนธมฺมํ “มา นสฺสี”ติ อลพฺภนียํ ฐานํ สมเณน วา พฺราหฺมเณน วา เทเวน วา มาเรน วา พฺรหฺมุนา วา เกนจิ วา โลกสฺมิํ.}

\prose{อสฺสุตวโต มหาราช ปุถุชฺชนสฺส ชราธมฺมํ ชีรติ, โส ชราธมฺเม ชิณฺเณ น อิติ ปฏิสญฺจิกฺขติ. “น โข มเยฺหเวกสฺส ชราธมฺมํ ชีรติ, อถ โข ยาวตา สตฺตานํ อาคติ คติ จุติ อุปปตฺติ, สพฺเพสํ สตฺตานํ ชราธมฺมํ ชีรติ.\csromanspacer{} อหญฺเจว โข ปน ชราธมฺเม ชิณฺเณ โสเจยฺยํ กิลเมยฺยํ ปริเทเวยฺยํ, อุรตฺตาฬิํ กนฺเทยฺยํ, สมฺโมหํ อาปชฺเชยฺยํ, ภตฺตมฺปิ เม นจฺฉาเทยฺย, กาเยปิ ทุพฺพณฺณิยํ โอกฺกเมยฺย, กมฺมนฺตาปิ นปฺปวตฺเตยฺยุํ, อมิตฺตาปิ อตฺตมนา อสฺสุ, มิตฺตาปิ ทุมฺมนา อสฺสู”ติ.\csromanspacer{} โส ชราธมฺเม ชิณฺเณ โสจติ กิลมติ ปริเทวติ, อุรตฺตาฬิํ กนฺทติ, สมฺโมหํ อาปชฺชติ.\csromanspacer{} อยํ วุจฺจติ มหาราช อสฺสุตวา ปุถุชฺชโน วิทฺโธ สวิเสน โสกสลฺเลน อตฺตานํเยว ปริตาเปติ.}

\prose{ปุน จปรํ มหาราช อสฺสุตวโต ปุถุชฺชนสฺส พฺยาธิธมฺมํ พฺยาธียติ ฯเปฯ มรณธมฺมํ มียติ.\csromanspacer{} ขยธมฺมํ ขียติ.\csromanspacer{} นสฺสนธมฺมํ นสฺสติ, โส นสฺสนธมฺเม นฏฺเฐ น อิติ ปฏิสญฺจิกฺขติ “น โข มเยฺหเวกสฺส นสฺสนธมฺมํ นสฺสติ, อถ โข ยาวตา สตฺตานํ อาคติ คติ จุติ อุปปตฺติ, สพฺเพสํ สตฺตานํ นสฺสนธมฺมํ นสฺสติ.\csromanspacer{} อหญฺเจว โข ปน นสฺสนธมฺเม นฏฺเฐ โสเจยฺยํ กิลเมยฺยํ ปริเทเวยฺยํ, อุรตฺตาฬิํ กนฺเทยฺยํ, สมฺโมหํ อาปชฺเชยฺยํ, ภตฺตมฺปิ เม นจฺฉาเทยฺย, กาเยปิ ทุพฺพณฺณิยํ โอกฺกเมยฺย, กมฺมนฺตาปิ นปฺปวตฺเตยฺยุํ, อมิตฺตาปิ อตฺตมนา อสฺสุ, มิตฺตาปิ ทุมฺมนา อสฺสู”ติ.\csromanspacer{} โส นสฺสนธมฺเม นฏฺเฐ โสจติ กิลมติ ปริเทวติ, อุรตฺตาฬิํ กนฺทติ, สมฺโมหํ อาปชฺชติ.\csromanspacer{} อยํ วุจฺจติ มหาราช อสฺสุตวา ปุถุชฺชโน วิทฺโธ สวิเสน โสกสลฺเลน อตฺตานํเยว ปริตาเปติ.}

\prose{\csromanfolio{54}สุตวโต จ โข มหาราช อริยสาวกสฺส ชราธมฺมํ ชีรติ, โส ชราธมฺเม ชิณฺเณ อิติ ปฏิสญฺจิกฺขติ, “น โข มเยฺหเวกสฺส ชราธมฺมํ ชีรติ, อถ โข ยาวตา สตฺตานํ อาคติ คติ จุติ อุปปตฺติ, สพฺเพสํ สตฺตานํ ชราธมฺมํ ชีรติ.\csromanspacer{} อหญฺเจว โข ปน ชราธมฺเม ชิณฺเณ โสเจยฺยํ กิลเมยฺยํ ปริเทเวยฺยํ, อุรตฺตาฬิํ กนฺเทยฺยํ, สมฺโมหํ อาปชฺเชยฺยํ, ภตฺตมฺปิ เม นจฺฉาเทยฺย, กาเยปิ ทุพฺพณฺณิยํ โอกฺกเมยฺย, กมฺมนฺตาปิ นปฺปวตฺเตยฺยุํ, อมิตฺตาปิ อตฺตมนา อสฺสุ, มิตฺตาปิ ทุมฺมนา อสฺสู”ติ.\csromanspacer{} โส ชราธมฺเม ชิณฺเณ น โสจติ น กิลมติ น ปริเทวติ, น อุรตฺตาฬิํ กนฺทติ, น สมฺโมหํ อาปชฺชติ.\csromanspacer{} อยํ วุจฺจติ มหาราช สุตวา อริยสาวโก อพฺพุหิ สวิสํ โสกสลฺลํ, เยน วิทฺโธ อสฺสุตวา ปุถุชฺชโน อตฺตานํเยว ปริตาเปติ.\csromanspacer{} อโสโก วิสลฺโล อริยสาวโก อตฺตานํเยว ปรินิพฺพาเปติ.}

\prose{ปุน จปรํ มหาราช สุตวโต อริยสาวกสฺส พฺยาธิธมฺมํ พฺยาธียติ ฯเปฯ มรณธมฺมํ มียติ.\csromanspacer{} ขยธมฺมํ ขียติ.\csromanspacer{} นสฺสนธมฺมํ นสฺสติ, โส นสฺสนธมฺเม นฏฺเฐ อิติ ปฏิสญฺจิกฺขติ, “น โข มเยฺหเวกสฺส นสฺสนธมฺมํ นสฺสติ, อถ โข ยาวตา สตฺตานํ อาคติ คติ จุติ อุปปตฺติ, สพฺเพสํ สตฺตานํ นสฺสนธมฺมํ นสฺสติ.\csromanspacer{} อหญฺเจว โข ปน นสฺสนธมฺเม นฏฺเฐ โสเจยฺยํ กิลเมยฺยํ ปริเทเวยฺยํ, อุรตฺตาฬิํ กนฺเทยฺยํ, สมฺโมหํ อาปชฺเชยฺยํ, ภตฺตมฺปิ เม นจฺฉาเทยฺย, กาเยปิ ทุพฺพณฺณิยํ โอกฺกเมยฺย, กมฺมนฺตาปิ นปฺปวตฺเตยฺยุํ, อมิตฺตาปิ อตฺตมนา อสฺสุ, มิตฺตาปิ ทุมฺมนา อสฺสู”ติ.\csromanspacer{} โส นสฺสนธมฺเม นฏฺเฐ น โสจติ น กิลมติ น ปริเทวติ, น อุรตฺตาฬิํ กนฺทติ, น สมฺโมหํ อาปชฺชติ.\csromanspacer{} อยํ วุจฺจติ มหาราช สุตวา อริยสาวโก อพฺพุหิ สวิสํ โสกสลฺลํ, เยน วิทฺโธ อสฺสุตวา ปุถุชฺชโน อตฺตานํเยว ปริตาเปติ.\csromanspacer{} อโสโก วิสลฺโล อริยสาวโก อตฺตานํเยว ปรินิพฺพาเปติ.}

\prose{อิมานิ โข มหาราช ปญฺจ อลพฺภนียานิ ฐานานิ สมเณน วา พฺราหฺมเณน วา เทเวน วา มาเรน วา พฺรหฺมุนา วา เกนจิ วา โลกสฺมินฺติ.}

\csromangathameasure{น โสจนาย ปริเทวนาย, \\ อตฺโถธ ลพฺภา อปิ อปฺปโกปิ. \\ โสจนฺตเมนํ ทุขิตํ วิทิตฺวา, \\ ปจฺจตฺถิกา อตฺตมนา ภวนฺติ.}
\csromangathagroup{\csromangathastanza{น โสจนาย ปริเทวนาย, \\ อตฺโถธ ลพฺภา อปิ อปฺปโกปิ. \\ โสจนฺตเมนํ ทุขิตํ วิทิตฺวา, \\ ปจฺจตฺถิกา อตฺตมนา ภวนฺติ.}}

\chapterheadpageread{5. มุณฺฑราชวคฺค}
\csromangathameasure{ยโต จ โข ปณฺฑิโต อาปทาสุ, \\ น เวธตี อตฺถวินิจฺฉยญฺญู. \\ ปจฺจตฺถิกาสฺส ทุขิตา ภวนฺติ, \\ ทิสฺวา มุขํ อวิการํ ปุราณํ. \\ ชปฺเปน มนฺเตน สุภาสิเตน, \\ อนุปฺปทาเนน ปเวณิยา วา. \\ ยถา ยถา ยตฺถ ลเภถ อตฺถํ, \\ ตถา ตถา ตตฺถ ปรกฺกเมยฺย. \\ สเจ ปชาเนยฺย อลพฺภเนยฺโย, \\ มยา ว อญฺเญน วา เอส อตฺโถ. \\ อโสจมาโน อธิวาสเยยฺย,}
\csromangathagroup{\csromangathastanza{\csromanfolio{55}ยโต จ โข ปณฺฑิโต อาปทาสุ, \\ น เวธตี อตฺถวินิจฺฉยญฺญู. \\ ปจฺจตฺถิกาสฺส ทุขิตา ภวนฺติ, \\ ทิสฺวา มุขํ อวิการํ ปุราณํ.}\csromangathastanza{ชปฺเปน มนฺเตน สุภาสิเตน, \\ อนุปฺปทาเนน ปเวณิยา วา. \\ ยถา ยถา ยตฺถ ลเภถ อตฺถํ, \\ ตถา ตถา ตตฺถ ปรกฺกเมยฺย.}\csromangathastanza{สเจ ปชาเนยฺย อลพฺภเนยฺโย, \\ มยา ว อญฺเญน วา เอส อตฺโถ. \\ อโสจมาโน อธิวาสเยยฺย,}}

\csromanheader{2}{กมฺมํ ทฬฺหํ กโรมิ ทานีติ.*}
\prose{เอวํ วุตฺเต มุณฺโฑ ราชา อายสฺมนฺตํ นารทํ เอตทโวจ “โก นาโม\footnote{โก นุ โข (สี, อิ)} อยํ ภนฺเต ธมฺมปริยาโย”ติ.\csromanspacer{} โสกสลฺลหรโณ นาม อยํ มหาราช ธมฺมปริยาโยติ.\csromanspacer{} ตคฺฆ ภนฺเต โสกสลฺลหรโณ\footnote{ตคฺฆ ภนฺเต โสกสลฺลหรโณ, ตคฺฆ ภนฺเต โสกสลฺลหรโณ (สี, สฺยา, กํ, อิ)}, อิมํ หิ เม ภนฺเต ธมฺมปริยายํ สุตฺวา โสกสลฺลํ ปหีนนฺติ.}

\prose{อถ โข มุณฺโฑ ราชา ปิยกํ โกสารกฺขํ อามนฺเตสิ “เตน หิ สมฺม ปิยก ภทฺทาย เทวิยา สรีรํ ฌาเปถ, ถูปญฺจสฺสา กโรถ, อชฺชตคฺเค ทานิ มยํ นฺหายิสฺสาม เจว วิลิมฺปิสฺสาม, ภตฺตญฺจ ภุญฺชิสฺสาม, กมฺมนฺเต จ ปโยเชสฺสามา”ติ.\csromanspacer{}\csromanspacer{}ทสมํ.\csromanspacer{} มุณฺฑราชวคฺโค ปญฺจโม.}

\titlehead{\textbf{ตสฺสุทฺทานํ}}
\csromangathameasure{อาทิโย สปฺปุริโส อิฏฺฐา, มนาปทายีภิสนฺทํ. \\ สมฺปทา จ ธนํ ฐานํ, โกสโล นารเทน จาติ. \\ \textbf{ปฐโม ปณฺณาสโก สมตฺโต.}}
\csromangathagroup{\csromangathastanza{อาทิโย สปฺปุริโส อิฏฺฐา, มนาปทายีภิสนฺทํ. \\ สมฺปทา จ ธนํ ฐานํ, โกสโล นารเทน จาติ. \\ \textbf{ปฐโม ปณฺณาสโก สมตฺโต.}}}

\csromanmatikaanchor{mtk.2}
\csromanmatikaanchor{mtk.175}
\csromantocmarkref{subsection}{1-2. สมฺปทาสุตฺตทฺวย}{mtk.2}
\bookmark[dest={mtk.2},level=2]{1-2. สมฺปทาสุตฺตทฺวย}
\chapterheadpageread{2. ทุติยปณฺณาสก}
\chapterhead{\textbf{(6) 1. นีวรณวคฺค}}
\vspace{\tipitakaparskip}
\csromancenter{อาวรณสุตฺต 51.\csromanspacer{} เอวํ เม สุตํ\csromandash{}เอกํ สมยํ ภควา สาวตฺถิยํ วิหรติ เชตวเน อนาถปิณฺฑิกสฺส อาราเม.\csromanspacer{} ตตฺร โข ภควา ภิกฺขู อามนฺเตสิ “ภิกฺขโว”ติ. “ภทนฺเต”ติ เต ภิกฺขู ภควโต ปจฺจสฺเสสุํ.\csromanspacer{} ภควา เอตทโวจ\csromandash{}}
\prose{\csromanfolio{56}ปญฺจิเม ภิกฺขเว อาวรณา นีวรณา เจตโส อชฺฌารุตา ปญฺญาย ทุพฺพลีกรณา.\csromanspacer{} กตเม ปญฺจ? กามจฺฉนฺโท ภิกฺขเว อาวรโณ นีวรโณ เจตโส อชฺฌารุโห ปญฺญาย ทุพฺพลีกรโณ.\csromanspacer{} พฺยาปาโท ภิกฺขเว อาวรโณ นีวรโณ เจตโส อชฺฌารุโห ปญฺญาย ทุพฺพลีกรโณ.\csromanspacer{} ถินมิทฺธํ ภิกฺขเว อาวรณํ นีวรณํ เจตโส อชฺฌารุหํ ปญฺญาย ทุพฺพลีกรณํ.\csromanspacer{} อุทฺธจฺจกุกฺกุจฺจํ ภิกฺขเว อาวรณํ นีวรณํ เจตโส อชฺฌารุหํ ปญฺญาย ทุพฺพลีกรณํ.\csromanspacer{} วิจิกิจฺฉา ภิกฺขเว อาวรณา นีวรณา เจตโส อชฺฌารุหา ปญฺญาย ทุพฺพลีกรณา.\csromanspacer{} อิเม โข ภิกฺขเว ปญฺจ อาวรณา นีวรณา เจตโส อชฺฌารุหา ปญฺญาย ทุพฺพลีกรณา.}

\prose{โส วต ภิกฺขเว ภิกฺขุ อิเม ปญฺจ อาวรเณ นีวรเณ เจตโส อชฺฌารุเห ปญฺญาย ทุพฺพลีกรเณ อปฺปหาย อพลาย ปญฺญาย ทุพฺพลาย อตฺตตฺถํ วา ญสฺสติ, ปรตฺถํ วา ญสฺสติ, อุภยตฺถํ วา ญสฺสติ, อุตฺตริ\footnote{อุตฺตริํ (สี, สฺยา, กํ, อิ)} วา มนุสฺสธมฺมา อลมริยญาณทสฺสนวิเสสํ สจฺฉิกริสฺสตีติ เนตํ ฐานํ วิชฺชติ.\csromanspacer{} เสยฺยถาปิ ภิกฺขเว นที ปพฺพเตยฺยา ทูรงฺคมา\footnote{ทูรคมา (สี)} สีฆโสตา หารหารินี, ตสฺสา ปุริโส อุภโต นงฺคลมุขานิ วิวเรยฺย.\csromanspacer{} เอวํ หิ โส ภิกฺขเว มชฺเฌ นทิยา โสโต วิกฺขิตฺโต วิสโฏ พฺยาทิณฺโณ เนว\footnote{น เจว (ก)} ทูรงฺคโม อสฺส น\footnote{น จ (ก)} สีฆโสโต น4 หารหารี\footnote{หารหาริณี (สี)}.\csromanspacer{} เอวเมวํ โข ภิกฺขเว โส วต ภิกฺขุ อิเม ปญฺจ อาวรเณ นีวรเณ เจตโส อชฺฌารุเห ปญฺญาย ทุพฺพลีกรเณ อปฺปหาย อพลาย ปญฺญาย ทุพฺพลาย}

\vspace{\tipitakaparskip}
\csromancenter{(6) 1.\csromanspacer{} นีวรณวคฺค อตฺตตฺถํ วา ญสฺสติ, ปรตฺถํ วา ญสฺสติ, อุภยตฺถํ วา ญสฺสติ, อุตฺตริ วา มนุสฺสธมฺมา อลมริยญาณทสฺสนวิเสสํ สจฺฉิกริสฺสตีติ เนตํ ฐานํ วิชฺชติ.}
\prose{\csromanfolio{57}โส วต ภิกฺขเว ภิกฺขุ อิเม ปญฺจ อาวรเณ นีวรเณ เจตโส อชฺฌารุเห ปญฺญาย ทุพฺพลีกรเณ ปหาย พลวติยา ปญฺญาย อตฺตตฺถํ วา ญสฺสติ, ปรตฺถํ วา ญสฺสติ, อุภยตฺถํ วา ญสฺสติ, อุตฺตริ วา มนุสฺสธมฺมา อลมริยญาณทสฺสนวิเสสํ สจฺฉิกริสฺสตีติ ฐานเมตํ วิชฺชติ.\csromanspacer{} เสยฺยถาปิ ภิกฺขเว นที ปพฺพเตยฺยา ทูรงฺคมา สีฆโสตา หารหารินี.\csromanspacer{} ตสฺสา ปุริโส อุภโต นงฺคลมุขานิ ปิทเหยฺย.\csromanspacer{} เอวํ หิ โส ภิกฺขเว มชฺเฌ นทิยา โสโต อวิกฺขิตฺโต อวิสโฏ อพฺยาทิณฺโณ ทูรงฺคโม เจว อสฺส สีฆโสโต จ หารหารี จ.\csromanspacer{} เอวเมวํ โข ภิกฺขเว โส วต ภิกฺขุ อิเม ปญฺจ อาวรเณ นีวรเณ เจตโส อชฺฌารุเห ปญฺญาย ทุพฺพลีกรเณ ปหาย พลวติยา ปญฺญาย อตฺตตฺถํ วา ญสฺสติ, ปรตฺถํ วา ญสฺสติ, อุภยตฺถํ วา ญสฺสติ, อุตฺตริ วา มนุสฺสธมฺมา อลมริยญาณทสฺสนวิเสสํ สจฺฉิกริสฺสตีติ ฐานเมตํ วิชฺชตีติ.\csromanspacer{}\csromanspacer{}ปฐมํ.}

\csromanheader{2}{2. อกุสลราสิสุตฺต}
\proseitem{52}{“อกุสลราสี”ติ ภิกฺขเว วทมาโน ปญฺจ นีวรเณ\footnote{อิเม ปญฺจ นีวรเณ (สี)} สมฺมา วทมาโน วเทยฺย, เกวโล หายํ\footnote{หยํ (สี), จายํ (สฺยา, กํ), สายํ (ก)} ภิกฺขเว อกุสลราสิ ยทิทํ ปญฺจ นีวรณา.\csromanspacer{} กตเม ปญฺจ? กามจฺฉนฺทนีวรณํ พฺยาปาทนีวรณํ ถินมิทฺธนีวรณํ อุทฺธจฺจกุกฺกุจฺจนีวรณํ วิจิกิจฺฉานีวรณํ. “อกุสลราสี”ติ ภิกฺขเว วทมาโน อิเม ปญฺจ นีวรเณ สมฺมา วทมาโน วเทยฺย, เกวโล หายํ ภิกฺขเว อกุสลราสิ ยทิทํ ปญฺจ นีวรณาติ.\csromanspacer{}\csromanspacer{}ทุติยํ.}

\csromanheader{2}{3. ปธานิยงฺคสุตฺต}
\proseitem{53}{ปญฺจิมานิ ภิกฺขเว ปธานิยงฺคานิ.\csromanspacer{} กตมานิ ปญฺจ? อิธ ภิกฺขเว ภิกฺขุ สทฺโธ โหติ, สทฺทหติ ตถาคตสฺส โพธิํ “อิติปิ โส ภควา อรหํ สมฺมาสมฺพุทฺโธ วิชฺชาจรณสมฺปนฺโน สุคโต โลกวิทู อนุตฺตโร ปุริสทมฺมสารถิ สตฺถา เทวมนุสฺสานํ พุทฺโธ ภควา”ติ.\csromanspacer{} อปฺปาพาโธ \csromanfolio{58}โหติ อปฺปาตงฺโก สมเวปากินิยา คหณิยา สมนฺนาคโต นาติสีตาย นาจฺจุณฺหาย มชฺฌิมาย ปธานกฺขมาย, อสโฐ โหติ อมายาวี ยถาภูตํ อตฺตานํ อาวิกตฺตา สตฺถริ วา วิญฺญูสุ วา สพฺรหฺมจารีสุ, อารทฺธวีริโย วิหรติ อกุสลานํ ธมฺมานํ ปหานาย กุสลานํ ธมฺมานํ อุปสมฺปทาย, ถามวา ทฬฺหปรกฺกโม อนิกฺขิตฺตธุโร กุสเลสุ ธมฺเมสุ, ปญฺญวา โหติ อุทยตฺถคามินิยา ปญฺญาย สมนฺนาคโต อริยาย นิพฺเพธิกาย สมฺมา ทุกฺขกฺขยคามินิยา.\csromanspacer{} อิมานิ โข ภิกฺขเว ปญฺจ ปธานิยงฺคานีติ.\csromanspacer{}\csromanspacer{}ตติยํ.}

\csromanheader{2}{4. สมยสุตฺต}
\proseitem{54}{ปญฺจิเม ภิกฺขเว อสมยา ปธานาย.\csromanspacer{} กตเม ปญฺจ? อิธ ภิกฺขเว ภิกฺขุ ชิณฺโณ โหติ ชรายาภิภูโต.\csromanspacer{} อยํ ภิกฺขเว ปฐโม อสมโย ปธานาย.}

\prose{ปุน จปรํ ภิกฺขเว ภิกฺขุ พฺยาธิโต โหติ พฺยาธินาภิภูโต.\csromanspacer{} อยํ ภิกฺขเว ทุติโย อสมโย ปธานาย.}

\prose{ปุน จปรํ ภิกฺขเว ทุพฺภิกฺขํ โหติ ทุสฺสสฺสํ ทุลฺลภปิณฺฑํ, น สุกรํ อุญฺเฉน ปคฺคเหน ยาเปตุํ.\csromanspacer{} อยํ ภิกฺขเว ตติโย อสมโย ปธานาย.}

\prose{ปุน จปรํ ภิกฺขเว ภยํ โหติ อฏวิสงฺโกโป, จกฺกสมารูฬฺหา ชานปทา ปริยายนฺติ.\csromanspacer{} อยํ ภิกฺขเว จตุตฺโถ อสมโย ปธานาย.}

\prose{ปุน จปรํ ภิกฺขเว สํโฆ ภินฺโน โหติ, สํเฆ โข ปน ภิกฺขเว ภินฺเน อญฺญมญฺญํ อกฺโกสา จ โหนฺติ, อญฺญมญฺญํ ปริภาสา จ โหนฺติ, อญฺญมญฺญํ ปริกฺเขปา จ โหนฺติ, อญฺญมญฺญํ ปริจฺจชา จ โหนฺติ.\csromanspacer{} ตตฺถ อปฺปสนฺนา เจว นปฺปสีทนฺติ, ปสนฺนานญฺจ เอกจฺจานํ อญฺญถตฺตํ โหติ.\csromanspacer{} อยํ ภิกฺขเว ปญฺจโม อสมโย ปธานาย.\csromanspacer{} อิเม โข ภิกฺขเว ปญฺจ อสมยา ปธานายาติ.}

\prose{ปญฺจิเม ภิกฺขเว สมยา ปธานาย.\csromanspacer{} กตเม ปญฺจ? อิธ ภิกฺขเว ภิกฺขุ ทหโร โหติ ยุวา สุสุ กาฬเกโส ภเทฺรน โยพฺพเนน สมนฺนาคโต ปฐเมน วยสา.\csromanspacer{} อยํ ภิกฺขเว ปฐโม สมโย ปธานาย.}

\chapterheadpageread{(6) 1. นีวรณวคฺค}
\prose{\csromanfolio{59}ปุน จปรํ ภิกฺขเว ภิกฺขุ อปฺปาพาโธ โหติ อปฺปาตงฺโก สมเวปากินิยา คหณิยา สมนฺนาคโต นาติสีตาย นาจฺจุณาย มชฺฌิมาย ปธานกฺขมาย.\csromanspacer{} อยํ ภิกฺขเว ทุติโย สมโย ปธานาย.}

\prose{ปุน จปรํ ภิกฺขเว สุภิกฺขํ โหติ สุสสฺสํ สุลภปิณฺฑํ, สุกรํ อุญฺเฉน ปคฺคเหน ยาเปตุํ.\csromanspacer{} อยํ ภิกฺขเว ตติโย สมโย ปธานาย.}

\prose{ปุน จปรํ ภิกฺขเว มนุสฺสา สมคฺคา สมฺโมทมานา อวิวทมานา ขีโรทกีภูตา อญฺญมญฺญํ ปิยจกฺขูหิ สมฺปสฺสนฺตา วิหรนฺติ.\csromanspacer{} อยํ ภิกฺขเว จตุตฺโถ สมโย ปธานาย.}

\prose{ปุน จปรํ ภิกฺขเว สํโฆ สมคฺโค สมฺโมทมาโน อวิวทมาโน เอกุทฺเทโส ผาสุ วิหรติ, สํเฆ โข ปน ภิกฺขเว สมคฺเค น เจว อญฺญมญฺญํ อกฺโกสา โหนฺติ, น จ อญฺญมญฺญํ ปริภาสา โหนฺติ, น จ อญฺญมญฺญํ ปริกฺเขปา โหนฺติ, น จ อญฺญมญฺญํ ปริจฺจชา โหนฺติ.\csromanspacer{} ตตฺถ อปฺปสนฺนา เจว ปสีทนฺติ, ปสนฺนานญฺจ ภิยฺโยภาโว\footnote{ภิยฺโยภาวาย (ก)} โหติ.\csromanspacer{} อยํ ภิกฺขเว ปญฺจโม สมโย ปธานาย.\csromanspacer{} อิเม โข ภิกฺขเว ปญฺจ สมยา ปธานายาติ.\csromanspacer{}\csromanspacer{}จตุตฺถํ.}

\csromanheader{2}{5. มาตาปุตฺตสุตฺต}
\proseitem{55}{เอกํ สมยํ ภควา สาวตฺถิยํ วิหรติ เชตวเน อนาถปิณฺฑิกสฺส อาราเม.\csromanspacer{} เตน โข ปน สมเยน สาวตฺถิยํ อุโภ มาตาปุตฺตา วสฺสาวาสํ อุปคมิํสุ\footnote{อุปสงฺกมิํสุ (ก)} ภิกฺขุ จ ภิกฺขุนี จ, เต อญฺญมญฺญสฺส อภิณฺหํ ทสฺสนกามา อเหสุํ.\csromanspacer{} มาตาปิ ปุตฺตสฺส อภิณฺหํ ทสฺสนกามา อโหสิ, ปุตฺโตปิ มาตรํ อภิณฺหํ ทสฺสนกาโม อโหสิ.\csromanspacer{} เตสํ อภิณฺหํ ทสฺสนา สํสคฺโค อโหสิ, สํสคฺเค สติ วิสฺสาโส อโหสิ, วิสฺสาเส สติ โอตาโร อโหสิ.\csromanspacer{} เต โอติณฺณจิตฺตา สิกฺขํ อปจฺจกฺขาย ทุพฺพลฺยํ อนาวิกตฺวา เมถุนํ ธมฺมํ ปฏิเสวิํสุ.}

\prose{อถ โข สมฺพหุลา ภิกฺขู เยน ภควา เตนุปสงฺกมิํสุ, อุปสงฺกมิตฺวา ภควนฺตํ อภิวาเทตฺวา เอกมนฺตํ นิสีทิํสุ, เอกมนฺตํ นิสินฺนา โข เต ภิกฺขู ภควนฺตํ เอตทโวจุํ “อิธ ภนฺเต สาวตฺถิยํ อุโภ \csromanfolio{60}มาตาปุตฺตา วสฺสาวาสํ อุปคมิํสุ ภิกฺขุ จ ภิกฺขูนี จ, เต อญฺญมญฺญสฺส อภิณฺหํ ทสฺสนกามา อเหสุํ, มาตาปิ ปุตฺตสฺส อภิณฺหํ ทสฺสนกามา อโหสิ, ปุตฺโตปิ มาตรํ อภิณฺหํ ทสฺสนกาโม อโหสิ.\csromanspacer{} เตสํ อภิณฺหํ ทสฺสนา สํสคฺโค อโหสิ, สํสคฺเค สติ วิสฺสาโส อโหสิ, วิสฺสาเส สติ โอตาโร อโหสิ, เต โอติณฺณจิตฺตา สิกฺขํ อปจฺจกฺขาย ทุพฺพลฺยํ อนาวิกตฺวา เมถุนํ ธมฺมํ ปฏิเสวิํสู”ติ.}

\prose{กิํ นุ โส ภิกฺขเว โมฆปุริโส มญฺญติ “น มาตา ปุตฺเต สารชฺชติ, ปุตฺโต วา ปน มาตรี”ติ.\csromanspacer{} นาหํ ภิกฺขเว อญฺญํ เอกรูปมฺปิ สมนุปสฺสามิ, เอวํ\footnote{ยํ เอวํ (สี)} รชนียํ เอวํ กมนียํ เอวํ มทนียํ เอวํ พนฺธนียํ เอวํ มุจฺฉนียํ เอวํ อนฺตรายกรํ อนุตฺตรสฺส โยคกฺเขมสฺส อธิคมาย.\csromanspacer{} ยถยิทํ ภิกฺขเว อิตฺถิรูปํ, อิตฺถิรูเป ภิกฺขเว สตฺตา รตฺตา คิทฺธา คถิตา\footnote{คธิตา (สฺยา, อิ, ก)} มุจฺฉิตา อชฺโฌสนฺนา\footnote{อชฺโฌปนฺนา (พหูสุ)}, เต ทีฆรตฺตํ โสจนฺติ อิตฺถิรูปวสานุคา.}

\prose{นาหํ ภิกฺขเว อญฺญํ เอกสทฺทมฺปิ.\csromanspacer{} เอกคนฺธมฺปิ.\csromanspacer{} เอกรกมฺปิ.\csromanspacer{} เอกโผฏฺฐพฺพมฺปิ สมนุปสฺสามิ เอวํ รชนียํ เอวํ กมนียํ เอวํ มทนียํ เอวํ พนฺธนียํ เอวํ มุจฺฉนียํ เอวํ อนฺตรายกรํ อนุตฺตรสฺส โยคกฺเขมสฺส อธิคมาย.\csromanspacer{} ยถยิทํ ภิกฺขเว อิตฺถิโผฏฺฐพฺพํ, อิตฺถิโผฏฺฐพฺเพ ภิกฺขเว สตฺตา รตฺตา คิทฺธา คถิตา มุจฺฉิตา อชฺโฌสนฺนา, เต ทีฆรตฺตํ โสจนฺติ อิตฺถิโผฏฺฐพฺพวสานุคา.}

\prose{อิตฺถี ภิกฺขเว คจฺฉนฺตีปิ ปุริสสฺส จิตฺตํ ปริยาทาย ติฏฺฐติ, ฐิตาปิ ฯเปฯ นิสินฺนาปิ.\csromanspacer{} สยานาปิ.\csromanspacer{} หสนฺตีปิ.\csromanspacer{} ภณนฺตีปิ.\csromanspacer{} คายนฺตีปิ.\csromanspacer{} โรทนฺตีปิ.\csromanspacer{} อุคฺฆาติตาปิ\footnote{อุคฺฆานิตาปิ (สี)}.\csromanspacer{} มตาปิ ปุริสสฺส จิตฺตํ ปริยาทาย ติฏฺฐติ.\csromanspacer{} ยํ หิ ตํ ภิกฺขเว สมฺมา วทมาโน วเทยฺย “สมนฺตปาโส มารสฺสา”ติ.\csromanspacer{} มาตุคามํเยว สมฺมา วทมาโน วเทยฺย “สมนฺตปาโส มารสฺสา”ติ.}

\csromangathameasure{สลฺลเป อสิหตฺเถน, ปิสาเจนาปิ สลฺลเป. \\ อาสีวิสมฺปิ อาสีเท, เยน ทฏฺโฐ น ชีวติ. \\ นเตฺวว เอโก เอกาย, มาตุคาเมน สลฺลเป.}
\csromangathagroup{\csromangathastanza{สลฺลเป อสิหตฺเถน, ปิสาเจนาปิ สลฺลเป. \\ อาสีวิสมฺปิ อาสีเท\footnote{อาสทฺเท (สฺยา, กํ)}, เยน ทฏฺโฐ น ชีวติ. \\ นเตฺวว เอโก เอกาย, มาตุคาเมน สลฺลเป.}}

\chapterheadpageread{(6) 1. นีวรณวคฺค}
\csromangathameasure{มุฏฺฐสฺสติํ ตา พนฺธนฺติ, เปกฺขิเตน สิเตน จ. \\ อโถปิ ทุนฺนิวตฺเถน, มญฺชุนา ภณิเตน จ. \\ เนโส ชโน สฺวาสีสโท, อปิ อุคฺฆาติโต มโต. \\ ปญฺจ กามคุณา เอเต, อิตฺถิรูปสฺมิํ ทิสฺสเร. \\ รูปา สทฺทา รสา คนฺธา, โผฏฺฐพฺพา จ มโนรมา. \\ เตสํ กาโมฆวูฬฺหานํ, กาเม อปริชานตํ. \\ กาลํ คติ ภวาภวํ, สํสารสฺมิํ ปุรกฺขตา. \\ เย จ กาเม ปริญฺญาย, จรนฺติ อกุโตภยา. \\ เต เว ปารงฺคตา โลเก, เย ปตฺตา อาสวกฺขยนฺติ. ปญฺจมํ.}
\csromangathagroup{\csromangathastanza{\csromanfolio{61}มุฏฺฐสฺสติํ ตา พนฺธนฺติ, เปกฺขิเตน สิเตน จ\footnote{มฺหิเตน จ (สฺยา, กํ)}. \\ อโถปิ ทุนฺนิวตฺเถน, มญฺชุนา ภณิเตน จ.}\csromangathastanza{เนโส ชโน สฺวาสีสโท, อปิ อุคฺฆาติโต มโต. \\ ปญฺจ กามคุณา เอเต, อิตฺถิรูปสฺมิํ ทิสฺสเร.}\csromangathastanza{รูปา สทฺทา รสา คนฺธา, โผฏฺฐพฺพา จ มโนรมา. \\ เตสํ กาโมฆวูฬฺหานํ, กาเม อปริชานตํ.}\csromangathastanza{กาลํ คติ\footnote{คติํ (สี, สฺยา, กํ, อิ)} ภวาภวํ, สํสารสฺมิํ ปุรกฺขตา. \\ เย จ กาเม ปริญฺญาย, จรนฺติ อกุโตภยา. \\ เต เว ปารงฺคตา โลเก, เย ปตฺตา อาสวกฺขยนฺติ.\csromanspacer{} ปญฺจมํ.}}

\csromanheader{2}{6. อุปชฺฌายสุตฺต}
\proseitem{56}{อถ โข อญฺญตโร ภิกฺขุ เยน สโก อุปชฺฌาโย เตนุปสงฺกมิ, อุปสงฺกมิตฺวา สกํ อุปชฺฌายํ เอตทโวจ “เอตรหิ เม ภนฺเต มธุรกชาโต เจว กาโย, ทิสา จ เม น ปกฺขายนฺติ, ธมฺมา จ มํ นปฺปฏิภนฺติ, ถินมิทฺธญฺจ เม จิตฺตํ ปริยาทาย ติฏฺฐติ, อนภิรโต จ พฺรหฺมจริยํ จรามิ, อตฺถิ จ เม ธมฺเมสุ วิจิกิจฺฉา”ติ.}

\prose{อถ โข โส ภิกฺขุ ตํ สทฺธิวิหาริกํ ภิกฺขุํ อาทาย เยน ภควา เตนุปสงฺกมิ, อุปสงฺกมิตฺวา ภควนฺตํ อภิวาเทตฺวา เอกมนฺตํ นิสีทิ, เอกมนฺตํ นิสินฺโน โข โส ภิกฺขุ ภควนฺตํ เอตทโวจ “อยํ ภนฺเต ภิกฺขุ เอวมาห ‘เอตรหิ เม ภนฺเต มธุรกชาโต เจว กาโย, ทิสา จ มํ น ปกฺขายนฺติ, ธมฺมา จ เม นปฺปฏิภนฺติ, ถินมิทฺธญฺจ เม จิตฺตํ ปริยาทาย ติฏฺฐติ, อนภิรโต จ พฺรหฺมจริยํ จรามิ, อตฺถิ จ เม ธมฺเมสุ วิจิกิจฺฉา’ติ”.}

\prose{เอวํ เหตํ ภิกฺขุ โหติ อินฺทฺริเยสุ อคุตฺตทฺวารสฺส โภชเน อมตฺตญฺญุโน ชาคริยํ อนนุยุตฺตสฺส อวิปสฺสกสฺส กุสลานํ ธมฺมานํ ปุพฺพรตฺตาปรรตฺตํ โพธิปกฺขิยานํ ธมฺมานํ ภาวนานุโยคํ อนนุยุตฺตสฺส วิหรโต ยํ มธุรกชาโต เจว กาโย โหติ, ทิสา จสฺส \csromanfolio{62}น ปกฺขายนฺติ, ธมฺมา จ ตํ นปฺปฏิภนฺติ, ถินมิทฺธญฺจสฺส จิตฺตํ ปริยาทาย ติฏฺฐติ, อนภิรโต จ พฺรหฺมจริยํ จรติ, โหติ จสฺส ธมฺเมสุ วิจิกิจฺฉา.\csromanspacer{} ตสฺมา ติห เต ภิกฺขุ เอวํ สิกฺขิตพฺพํ “อินฺทฺริเยสุ คุตฺตทฺวาโร ภวิสฺสามิ โภชเน มตฺตญฺญู ชาคริยํ อนุยุตฺโต วิปสฺสโก กุสลานํ ธมฺมานํ ปุพฺพรตฺตาปรรตฺตํ โพธิปกฺขิยานํ ธมฺมานํ ภาวนานุโยคํ อนุยุตฺโต วิหริสฺสามี”ติ.\csromanspacer{} เอวํ หิ เต ภิกฺขุ สิกฺขิตพฺพนฺติ.}

\prose{อถ โข โส ภิกฺขุ ภควตา อิมินา โอวาเทน โอวทิโต อุฏฺฐายาสนา ภควนฺตํ อภิวาเทตฺวา ปทกฺขิณํ กตฺวา ปกฺกามิ.\csromanspacer{} อถ โข โส ภิกฺขุ เอโก วูปกฏฺโฐ อปฺปมตฺโต อาตาปี ปหิตตฺโต วิหรนฺโต นจิรสฺเสว ยสฺสตฺถาย กุลปุตฺตา สมฺมเทว อคารสฺมา อนคาริยํ ปพฺพชนฺติ, ตทนุตฺตรํ พฺรหฺมจริยปริโยสานํ ทิฏฺเฐว ธมฺเม สยํ อภิญฺญา สจฺฉิกตฺวา อุปสมฺปชฺช วิหาสิ, “ขีณา ชาติ, วุสิตํ พฺรหฺมจริยํ, กตํ กรณียํ, นาปรํ อิตฺถตฺตายา”ติ อพฺภญฺญาสิ.\csromanspacer{} อญฺญตโร ปน โส ภิกฺขุ อรหตํ อโหสิ.}

\prose{อถ โข โส ภิกฺขุ อรหตฺตํ ปตฺโต เยน สโก อุปชฺฌาโย เตนุปสงฺกมิ, อุปสงฺกมิตฺวา สกํ อุปชฺฌายํ เอตทโวจ “เอตรหิ เม ภนฺเต น เจว\footnote{น เตฺวว (สี)} มธุรกชาโต กาโย, ทิสา จ เม ปกฺขายนฺติ, ธมฺมา จ มํ ปฏิภนฺติ, ถินมิทฺธญฺจ เม จิตฺตํ น ปริยาทาย ติฏฺฐติ, อภิรโต จ พฺรหฺมจริยํ จรามิ, นตฺถิ จ เม ธมฺเมสุ วิจิกิจฺฉา”ติ.\csromanspacer{} อถ โข โส ภิกฺขุ ตํ สทฺธิวิหาริกํ ภิกฺขุํ อาทาย เยน ภควา เตนุปสงฺกมิ, อุปสงฺกมิตฺวา ภควนฺตํ อภิวาเทตฺวา เอกมนฺตํ นิสีทิ, เอกมนฺตํ นิสินฺโน โข โส ภิกฺขุ ภควนฺตํ เอตทโวจ “อยํ ภนฺเต ภิกฺขุ เอวมาห ‘เอตรหิ เม ภนฺเต น เจว มธุรกชาโต กาโย, ทิสา จ เม ปกฺขายนฺติ, ธมฺมา จ มํ ปฏิภนฺติ, ถินมิทฺธญฺจ เม จิตฺตํ น ปริยาทาย ติฏฺฐติ, อภิรโต จ พฺรหฺมจริยํ จรามิ, นตฺถิ จ เม ธมฺเมสุ วิจิกิจฺฉา’ติ”.}

\prose{เอวํ เหตํ ภิกฺขุ โหติ อินฺทฺริเยสุ คุตฺตทฺวารสฺส โภชเน มตฺตญฺญุโน ชาคริยํ อนุยุตฺตส วิปสฺสกสฺส กุสลานํ ธมฺมานํ ปุพฺพรตฺตาปรรตฺตํ โพธิปกฺขิยานํ ธมฺมานํ ภาวนานุโยคํ อนุยุตฺตสฺส วิหรโต ยํ น เจว มธุรกชาโต กาโย โหติ, ทิสา จสฺส ปกฺขายนฺติ, ธมฺมา}

\vspace{\tipitakaparskip}
\csromancenter{(6) 1.\csromanspacer{} นีวรณวคฺค จ ตํ ปฏิภนฺติ, ถินมิทฺธญฺจสฺส จิตฺตํ น ปริยาทาย ติฏฺฐติ, อภิรโต จ พฺรหฺมจริยํ จรติ, น จสฺส โหติ ธมฺเมสุ วิจิกิจฺฉา.\csromanspacer{} ตสฺมา ติห โว ภิกฺขเว เอวํ สิกฺขิตพฺพํ “อินฺทฺริเยสุ คุตฺตทฺวารา ภวิสฺสาม, โภชเน มตฺตญฺญุโน ชาคริยํ อนุยุตฺตา วิปสฺสกา กุสลานํ ธมฺมานํ ปุพฺพรตฺตาปรรตฺตํ โพธิปกฺขิยานํ ธมฺมานํ ภาวนานุโยคํ อนุยุตฺตา วิหริสฺสามา”ติ.\csromanspacer{} เอวํ หิ โว ภิกฺขเว สิกฺขิตพฺพนฺติ.\csromanspacer{}\csromanspacer{}ฉฏฺฐํ.}
\csromanheader{2}{7. อภิณฺหปจฺจเวกฺขิตพฺพ ฐานสุตฺต}
\proseitem{57}{\csromanfolio{63}ปญฺจิมานิ ภิกฺขเว ฐานานิ อภิณฺหํ ปจฺจเวกฺขิตพฺพานิ อิตฺถิยา วา ปุริเสน วา คหฏฺเฐน วา ปพฺพชิเตน วา.\csromanspacer{} กตมานิ ปญฺจ? “ชราธมฺโมมฺหิ, ชรํ อนตีโต”ติ อภิณฺหํ ปจฺจเวกฺขิตพฺพํ อิตฺถิยา วา ปุริเสน วา คหฏฺเฐน วา ปพฺพชิเตน วา, “พฺยาธิธมฺโมมฺหิ, พฺยาธิํ อนตีโต”ติ อภิณฺหํ ปจฺจเวกฺขิตพฺพํ อิตฺถิยา วา ปุริเสน วา คหฏฺเฐน วา ปพฺพชิเตน วา, “มรณธมฺโมมฺหิ, มรณํ อนตีโต”ติ อภิณฺหํ ปจฺจเวกฺขิตพฺพํ อิตฺถิยา วา ปุริเสน วา คหฏฺเฐน วา ปพฺพชิเตน วา, “สพฺเพหิ เม ปิเยหิ มนาเปหิ นานาภาโว วินาภาโว”ติ อภิณฺหํ ปจฺจเวกฺขิตพฺพํ อิตฺถิยา วา ปุริเสน วา คหฏฺเฐน วา ปพฺพชิเตน วา, “กมฺมสฺสโกมฺหิ, กมฺมทายาโท กมฺมโยนิ กมฺมพนฺธุ กมฺมปฏิสรโณ, ยํ กมฺมํ กริสฺสามิ กลฺยาณํ วา ปาปกํ วา, ตสฺส ทายาโท ภวิสฺสามี”ติ อภิณฺหํ ปจฺจเวกฺขิตพฺพํ อิตฺถิยา วา ปุริเสน วา คหฏฺเฐน วา ปพฺพชิเตน วา.}

\prose{กิญฺจ ภิกฺขเว อตฺถวสํ ปฏิจฺจ “ชราธมฺโมมฺหิ, ชรํ อนตีโต”ติ อภิณฺหํ ปจฺจเวกฺขิตพฺพํ อิตฺถิยา วา ปุริเสน วา คหฏฺเฐน วา ปพฺพชิเตน วา.\csromanspacer{} อตฺถิ ภิกฺขเว สตฺตานํ โยพฺพเน โยพฺพนมโท, เยน มเทน มตฺตา กาเยน ทุจฺจริตํ จรนฺติ, วาจาย ทุจฺจริตํ จรนฺติ, มนสา ทุจฺจริตํ จรนฺติ.\csromanspacer{} ตสฺส ตํ ฐานํ อภิณฺหํ ปจฺจเวกฺขโต โย โยพฺพเน โยพฺพนมโท, โส สพฺพโส ปา ปหียติ, ตนุ วา ปน โหติ.\csromanspacer{} อิทํ โข ภิกฺขเว อตฺถวสํ ปฏิจฺจ “ชราธมฺโมมฺหิ, ชรํ อนตีโต”ติ อภิณฺหํ ปจฺจเวกฺขิตพฺพํ อิตฺถิยา วา ปุริเสน วา คหฏฺเฐน วา ปพฺพชิเตน วา.}

\prose{กิญฺจิ ภิกฺขเว อตฺถวสํ ปฏิจฺจ “พฺยาธิธมฺโมมฺหิ, พฺยาธิํ อนตีโต”ติ อภิณฺหํ ปจฺจเวกฺขิตพฺพํ อิตฺถิยา วา ปุริเสน วา คหฏฺเฐน วา ปพฺพชิเตน วา.\csromanspacer{} อตฺถิ ภิกฺขเว สตฺตานํ อาโรเคฺย อาโรคฺยมโท, เยน มเทน มตฺตา \csromanfolio{64}กาเยน ทุจฺจริตํ จรนฺติ, วาจาย ทุจฺจริตํ จรนฺติ, มนสา ทุจฺจริตํ จรนฺติ.\csromanspacer{} ตสฺส ตํ ฐานํ อภิณฺหํ ปจฺจเวกฺขโต โย อาโรเคฺย อาโรคฺยมโท, โส สพฺพโส วา ปหียติ, ตนุ วา ปน โหติ.\csromanspacer{} อิทํ โข ภิกฺขเว อตฺถวสํ ปฏิจฺจ “พฺยาธิธมฺโมมฺหิ, พฺยาธิํ อนตีโต”ติ อภิณฺหํ ปจฺจเวกฺขิตพฺพํ อิตฺถิยา วา ปุริเสน วา คหฏฺเฐน วา ปพฺพชิเตน วา.}

\prose{กิญฺจ ภิกฺขเว อตฺถวสํ ปฏิจฺจ “มรณธมฺโมมฺหิ, มรณํ อนตีโต”ติ อภิณฺหํ ปจฺจเวกฺขิตพฺพํ อิตฺถิยา วา ปุริเสน วา คหฏฺเฐน วา ปพฺพชิเตน วา.\csromanspacer{} อตฺถิ ภิกฺขเว สตฺตานํ ชีวิเต ชีวิตมโท, เยน มเทน มตฺตา กาเยน ทุจฺจริตํ จรนฺติ, วาจาย ทุจฺจริตํ จรนฺติ, มนสา ทุจฺจริตํ จรนฺติ.\csromanspacer{} ตสฺส ตํ ฐานํ อภิณฺหํ ปจฺจเวกฺขโต โย ชีวิเต ชีวิตมโท, โส สพฺพโส วา ปหียติ, ตนุ วา ปน โหติ.\csromanspacer{} อิทํ โข ภิกฺขเว อตฺถวสํ ปฏิจฺจ “มรณธมฺโมมฺหิ, มรณํ อนตีโต”ติ อภิณฺหํ ปจฺจเวกฺขิตพฺพํ อิตฺถิยา วา ปุริเสน วา คหฏฺเฐน วา ปพฺพชิเตน วา.}

\prose{กิญฺจ ภิกฺขเว อตฺถวสํ ปฏิจฺจ “สพฺเพหิ เม ปิเยหิ มนาเปหิ นานาภาโว วินาภาโว’ติ อภิณฺหํ ปจฺจเวกฺขิตพฺพํ อิตฺถิยา วา ปุริเสน วา คหฏฺเฐน วา ปพฺพชิเตน วา.\csromanspacer{} อตฺถิ ภิกฺขเว สตฺตานํ ปิเยสุ มนาเปสุ โย ฉนฺทราโค, เยน ราเคน รตฺตา กาเยน ทุจฺจริตํ จรนฺติ, วาจาย ทุจฺจริตํ จรนฺติ, มนสา ทุจฺจริตํ จรนฺติ.\csromanspacer{} ตสฺส ตํ ฐานํ อภิณฺหํ ปจฺจเวกฺขโต โย ปิเยสุ มนาเปสุ ฉนฺทราโค, โส สพฺพโส วา ปหียติ ตนุ วา ปน โหติ.\csromanspacer{} อิทํ โข ภิกฺขเว อตฺถวสํ ปฏิจฺจ “สพฺเพหิ เม ปิเยหิ มนาเปหิ นนาภาโว วินาภาโว”ติ อภิณฺหํ ปจฺจเวกฺขิตพฺพํ อิตฺถิยา วา ปุริเสน วา คหฏฺเฐน วา ปพฺพชิเตน วา.}

\prose{กิญฺจ ภิกฺขเว อตฺถวสํ ปฏิจฺจ “กมฺมสฺสโกมฺหิ, กมฺมทายาโท กมฺมโยนิ กมฺมพนฺธุ กมฺมปฏิสรโณ.\csromanspacer{} ยํ กมฺมํ กริสฺสามิ กลฺยาณํ วา ปาปกํ วา, ตสฺส ทายาโท ภวิสฺสามี”ติ อภิณฺหํ ปจฺจเวกฺขิตพฺพํ อิตฺถิยา วา ปุริเสน วา คหฏฺเฐน วา ปพฺพชิเตน วา.\csromanspacer{} อตฺถิ ภิกฺขเว สตฺตานํ กายทุจฺจริตํ วจีทุจฺจริตํ มโนทุจฺจริตํ.\csromanspacer{} ตสฺส ตํ ฐานํ อภิณฺหํ ปจฺจเวกฺขโต สพฺพโส วา ทุจฺจริตํ ปหียติ, ตนุ วา ปน โหติ.\csromanspacer{} อิทํ โข ภิกฺขเว อตฺถวสํ ปฏิจฺจ “กมฺมสฺสโกมฺหิ, กมฺมทายาโท กมฺมโยนิ กมฺมพนฺธุ กมฺมปฏิสรโณ, ยํ กมฺมํ กริสฺสามิ กลฺยาณํ วา ปาปกํ วา, ตสฺส ทายาโท}

\vspace{\tipitakaparskip}
\csromancenter{(6) 1.\csromanspacer{} นีวรณวคฺค ภวิสฺสามี”ติ อภิณฺหํ ปจฺจเวกฺขิตพฺพํ อิตฺถิยา วา ปุริเสน วา คหฏฺเฐน วา ปพฺพชิเตน วา.}
\prose{\csromanfolio{65}ส โข\footnote{สเจ (อิ, ก)} โส ภิกฺขเว อริยสาวโก อิติ ปฏิสญฺจิกฺขติ “น โข อหญฺเญเวโก ชราธมฺโม\footnote{อหญฺเจเวโก ชราธมฺโมมฺหิ (ก)} ชรํ อนตีโต.\csromanspacer{} อถ โข ยาวตา สตฺตานํ อาคติ คติ จุติ อุปปตฺติ, สพฺเพ สตฺตา ชราธมฺมา ชรํ อนตีตา”ติ, ตสฺส ตํ ฐานํ อภิณฺหํ ปจฺจเวกฺขโต มคฺโค สญฺชายติ, โส ตํ มคฺคํ อาเสวติ ภาเวติ พหุลีกโรติ.\csromanspacer{} ตสฺส ตํ มคฺคํ อาเสวโต ภาวยโต พหุลีกโรโต สํโยชนานิ สพฺพโส ปหียนฺติ, อนุสยา พฺยนฺตีโหนฺติ.}

\prose{“น โข อหญฺเญเวโก พฺยาธิธมฺโม พฺยาธิํ อนตีโต.\csromanspacer{} อถ โข ยาวตา สตฺตานํ อาคติ คติ จุติ อุปปตฺติ, สพฺเพ สตฺตา พฺยาธิธมฺมา พฺยาธิํ อนตีตา”ติ, ตสฺส ตํ ฐานํ อภิณฺหํ ปจฺจเวกฺขโต มคฺโค สญฺชายติ, โส ตํ มคฺคํ อาเสวติ ภาเวติ พหุลีกโรติ.\csromanspacer{} ตสฺส ตํ มคฺคํ อาเสวโต ภาวยโต พหุลีกโรโต สํโยชนานิ สพฺพโส ปหียนฺติ, อนุสยา พฺยนฺตีโหนฺติ.}

\prose{“น โข อหญฺเญเวโก มรณธมฺโม มรณํ อนตีโต.\csromanspacer{} อถ โข ยาวตา สตฺตานํ อาคติ คติ จุติ อุปปตฺติ, สพฺเพ สตฺตา มรณธมฺมา มรณํ อนตีตา”ติ, ตสฺส ตํ ฐานํ อภิณฺหํ ปจฺจเวกฺขโต มคฺโค สญฺชายติ, โส ตํ มคฺคํ อาเสวติ ภาเวติ พหุลีกโรติ.\csromanspacer{} ตสฺส ตํ มคฺคํ อาเสวโต ภาวยโต พหุลีกโรโต สํโยชนานิ สพฺพโส ปหียนฺติ, อนุสยา พฺยนฺตีโหนฺติ.}

\prose{“น โข มเยฺหเวกสฺส สพฺเพหิ ปิเยหิ มนาเปหิ นานาภาเว วินาภาโว.\csromanspacer{} อถ โข ยาวตา สตฺตานํ อาคติ คติ จุติ อุปปตฺติ, สพฺเพสํ สตฺตานํ ปิเยหิ มนาเปหิ นานาภาโว วินาภาโว”ติ, ตสฺส ตํ ฐานํ อภิณฺหํ ปจฺจเวกฺขโต มคฺโค สญฺชายติ, โส ตํ มคฺคํ อาเสวติ ภาเวติ พหุลีกโรติ.\csromanspacer{} ตสฺส ตํ มคฺคํ อาเสวโต ภาวยโต พหุลีกโรโต สํโยชนานิ สพฺพโส ปหียนฺติ, อนุสยา พฺยนฺตีโหนฺติ.}

\prose{\csromanfolio{66}“น โข อหญฺเญเวโก กมฺมสฺสโก กมฺมทายาโท กมฺมโยนิ กมฺมพนฺธุ กมฺมปฺปฏิสรโณ, ยํ กมฺมํ กริสฺสามิ กลฺยาณํ วา ปาปกํ วา, ตสฺส ทายาโท ภวิสฺสามิ.\csromanspacer{} อถ โข ยาวตา สตฺตานํ อาคติ คติ จุติ อุปปตฺติ, สพฺเพ สตฺตา กมฺมสฺสกา กมฺมทายาทา กมฺมโยนิ กมฺมพนฺธุ กมฺมปฺปฏิสรณา, ยํ กมฺมํ กริสฺสนฺติ กลฺยาณํ วา ปาปกํ วา, ตสฺส ทายาทา ภวิสฺสนฺตี”ติ, ตสฺส ตํ ฐานํ อภิณฺหํ ปจฺจเวกฺขโต มคฺโค สญฺชายติ, โส ตํ มคฺคํ อาเสวติ ภาเวติ พหุลีกโรติ.\csromanspacer{} ตสฺส ตํ มคฺคํ อาเสวโต ภาวยโต พหุลีกโรโต สํโยชนานิ สพฺพโส ปหียนฺติ, อนุสยา พฺยนฺตีโหนฺตีติ.}

\csromangathameasure{พฺยาธิธมฺมา ชราธมฺมา, อโถ มรณธมฺมิโน. \\ ยถา ธมฺมา ตถา สตฺตา, ชิคุจฺฉนฺติ ปุถุชฺชนา. \\ อหญฺเจ ตํ ชิคุจฺเฉยฺยํ, เอวํธมฺเมสุ ปาณิสุ. \\ น เมตํ ปติรูป’สฺส, มม เอวํ วิหาริโน. \\ โสหํ เอวํ วิหรนฺโต, ญตฺวา ธมฺมํ นิรูปธิํ. \\ อาโรเคฺย โยพฺพนสฺมิญฺจ, ชีวิตสฺมิญฺจ เย มทา. \\ สพฺเพ มเท อภิโภสฺมิ, เนกฺขมฺมํ ทฏฺฐุ เขมโต2. \\ ตสฺส เม อหุ อุสฺสาโห, นิพฺพานํ อภิปสฺสโต. \\ นาหํ ภพฺโพ เอตรหิ, กามานิ ปฏิเสวิตุํ. \\ อนิวตฺติ3 ภวิสฺสามิ, พฺรหฺมจริยปรายโณติ.สตฺตมํ.}
\csromangathagroup{\csromangathastanza{พฺยาธิธมฺมา ชราธมฺมา, อโถ มรณธมฺมิโน. \\ ยถา ธมฺมา ตถา สตฺตา\footnote{สนฺตา (สฺยา, กํ) 2. เนกฺขมฺเม ทฏฺฐุ เขมตํ (อํ 1. 145 ปิฏฺเฐ) อุภยตฺถปิ อฏฺฐกถาย สเมติ. 3. อนิวตฺตี (?)}, ชิคุจฺฉนฺติ ปุถุชฺชนา.}\csromangathastanza{อหญฺเจ ตํ ชิคุจฺเฉยฺยํ, เอวํธมฺเมสุ ปาณิสุ. \\ น เมตํ ปติรูป’สฺส, มม เอวํ วิหาริโน.}\csromangathastanza{โสหํ เอวํ วิหรนฺโต, ญตฺวา ธมฺมํ นิรูปธิํ. \\ อาโรเคฺย โยพฺพนสฺมิญฺจ, ชีวิตสฺมิญฺจ เย มทา.}\csromangathastanza{สพฺเพ มเท อภิโภสฺมิ, เนกฺขมฺมํ ทฏฺฐุ เขมโต2. \\ ตสฺส เม อหุ อุสฺสาโห, นิพฺพานํ อภิปสฺสโต.}\csromangathastanza{นาหํ ภพฺโพ เอตรหิ, กามานิ ปฏิเสวิตุํ. \\ อนิวตฺติ3 ภวิสฺสามิ, พฺรหฺมจริยปรายโณติ.\csromanspacer{}\csromanspacer{}สตฺตมํ.}}

\csromanheader{2}{8. ลิจฺฉวิกุมารกสุตฺต}
\proseitem{58}{เอกํ สมยํ ภควา เวสาลิยํ วิหรติ มหาวเน กูฏาคารสาลายํ.\csromanspacer{} อถ โข ภควา ปุพฺพณฺหสมยํ นิวาเสตฺวา ปตฺตจีวรมาทาย เวสาลิํ ปิณฺฑาย ปาวิสิ, เวสาลิยํ ปิณฺฑาย จริตฺวา ปจฺฉาภตฺตํ ปิณฺฑปาตปฏิกฺกนฺโต มหาวนํ อชฺโฌคาเหตฺวา อญฺญตรสฺมิํ รุกฺขมูเล ทิวาวิหารํ นิสีทิ.}

\prose{เตน โข ปน สมเยน สมฺพุหุลา ลิจฺฉวิกุมารกา สชฺชานิ ธนูนิ อาทาย กุกฺกุรสํฆปริวุตา มหาวเน อนุจงฺกมมานา อนุวิจรมานา}

\vspace{\tipitakaparskip}
\csromancenter{(6) 1.\csromanspacer{} นีวรณวคฺค อทฺทสุ ภควนฺตํ อญฺญตรสฺมิํ รุกฺขมูเล นิสินฺนํ, ทิสฺวาน สชฺชานิ ธนูนิ นิกฺขิปิตฺวา กุกฺกุรสํฆํ เอกมนฺตํ อุยฺโยเชตฺวา เยน ภควา เตนุปสงฺกมิํสุ, อุปสงฺกมิตฺวา ภควนฺตํ อภิวาเทตฺวา ตุณฺหีภูตา ตุณฺหีภูตา ปญฺชลิกา ภควนฺตํ ปยิรุปาสนฺติ.}
\prose{\csromanfolio{67}เตน โข ปน สมเยน มหานาโม ลิจฺฉวิ มหาวเน ชงฺฆาวิหารํ อนุจงฺกมมาโน อทฺทส เต ลิจฺฉวิกุมารเก ตุณฺหีภูเต ตุณฺหีภูเต ปญฺชลิเก ภควนฺตํ ปยิรุปาสนฺเต, ทิสฺวา เยน ภควา เตนุปสงฺกมิ, อุปสงฺกมิตฺวา ภควนฺตํ อภิวาเทตฺวา เอกมนฺตํ นิสีทิ, เอกมนฺตํ นิสินฺโน โข มหานาโม ลิจฺฉวิ อุทานํ อุทาเนสิ “ภวิสฺสนฺติ วชฺชี, ภวิสฺสนฺติ วชฺชี”ติ.}

\prose{กิํ ปน ตฺวํ มหานาม เอวํ วเทสิ “ภวิสฺสนฺติ วชฺชี, ภวิสฺสนฺติ วชฺชี”ติ.\csromanspacer{} อิเม ภนฺเต ลิจฺฉวิกุมารกา จณฺฑา ผรุสา อปานุภา\footnote{อปชหาติ (สี), อปาฏุภา (สฺยา, กํ), อปชหา (อิ), อปานุตา (กตฺถจิ)}.\csromanspacer{} ยานิปิ ตานิ กุเลสุ ปเหณกานิ\footnote{ปหีนกานิ (สี), ปหีณกานิ (สฺยา, กํ, อิ)} ปหียนฺติ อุจฺฉูติ วา พทราติ วา ปูวาติ วา โมทกาติ วา สํกุลิกาติ วา\footnote{สกฺขลิกาติ วา (สี, อิ)}, ตานิ วิลุมฺปิตฺวา วิลุมฺปิตฺวา ขาทนฺติ.\csromanspacer{} กุลิตฺถีนมฺปิ กุลกุมารีนมฺปิ ปจฺฉาลิยํ ขิปนฺติ, เต ทานิเม ตุณฺหีภูตา ตุณฺหีภูตา ปญฺชิลิกา ภควนฺตํ ปยิรุปาสนฺตีติ.}

\prose{ยสฺส กสฺสจิ มหานาม กุลปุตฺตสฺส ปญฺจ ธมฺมา สํวิชฺชนฺติ, ยทิ วา รญฺโญ ขตฺติยสฺส มุทฺธาวสิตฺตสฺส\footnote{มุทฺธาภิสิตฺตสฺส (ก), อุปริ 134, 136 ปิฏฺเฐสุ ปสฺสิตพฺพํ.}, ยทิ วา รฏฺฐิกสฺส เปตฺตนิกสฺส, ยทิ วา เสนาย เสนาปติกสฺส, ยทิ วา คามคามณิกสฺส, ยทิ วา ปูคคามณิกสฺส, เย วา ปน กุเลสุ ปจฺเจกาธิปจฺจํ กาเรนฺติ, วุทฺธิเยว ปาฏิกงฺขา โน ปริหานิ.}

\prose{กตเม ปญฺจ? อิธ มหานาม กุลปุตฺโต อุฏฺฐานวีริยาธิคเตหิ โภเคหิ พาหาพลปริจิเตหิ เสทาวกฺขิตฺเตหิ ธมฺมิเกหิ ธมฺมลทฺเธหิ มาตาปิตโร สกฺกโรติ ครุํ กโรติ มาเนติ ปูเชติ, ตเมนํ มาตาปิตโร สกฺกตา ครุกตา มานิตา ปูชิตา กลฺยาเณน มนสา อนุกมฺปนฺติ “จิรํ ชีว, ทีฆมายุํ ปาเลหี”ติ.\csromanspacer{} มาตาปิตานุกมฺปิตสฺส มหานาม กุลปุตฺตสฺส วุทฺธิเยว ปาฏิกงฺขา โน ปริหานิ.}

\prose{\csromanfolio{68}ปุน จปรํ มหานาม กุลปุตฺโต อุฏฺฐานวีริยาธิคเตหิ โภเคหิ พาหาพลปริจิเตหิ เสทาวกฺขิตฺเตหิ ธมฺมิเกหิ ธมฺมลทฺเธหิ ปุตฺตทารทาสกมฺมกรโปริเส สกฺกโรติ ครุํ กโรติ มาเนติ ปูเชติ, ตเมนํ ปุตฺตทารทาสกมฺมกรโปริสา สกฺกตา ครุกตา มานิตา ปูชิตา กลฺยาเณน มนสา อนุกมฺปนฺติ “จิรํ ชีว, ทีฆมายุํ ปาเลหี”ติ.\csromanspacer{} ปุตฺตทารทาสกมฺมกรโปริสานุกมฺปิตสฺส มหานาม กุลปุตฺตสฺส วุทฺธิเยว ปาฏิกงฺขา โน ปริหานิ.}

\prose{ปุน จปรํ มหานาม กุลปุตฺโต อุฏฺฐานวีริยาธิคเตหิ โภเคหิ พาหาพลปริจิเตหิ เสทาวกฺขิตฺเตหิ ธมฺมิเกหิ ธมฺมลทฺเธหิ เขตฺตกมฺมนฺตสามนฺตสโพฺยหาเร\footnote{...สามนฺตสํโวหาเร (สี, อิ)} สกฺกโรติ ครุํ กโรติ มาเนติ ปูเชติ, ตเมนํ เขตฺตกมฺมนฺตสามนฺตสโพฺยหารา สกฺกตา ครุกตา มานิตา ปูชิตา กลฺยาเณน มนสา อนุกมฺปนฺติ “จิรํ ชีว, ทีฆมายุํ ปาเลหี”ติ.\csromanspacer{} เขตฺตกมฺมนฺตสามนฺต สโพฺยหารานุกมฺปิตสฺส มหานาม กุลปุตฺตสฺส วุทฺธิเยว ปาฏิกงฺขา โน ปริหานิ.}

\prose{ปุน จปรํ มหานาม กุลปุตฺโต อุฏฺฐานวีริยาธิคเตหิ โภเคหิ พาหาพลปริจิเตหิ เสทาวกฺขิตฺเตหิ ธมฺมิเกหิ ธมฺมลทฺเธหิ ยาวตา พลิปฏิคฺคาหิกา เทวตา สกฺกโรติ ครุํ กโรติ มาเนติ ปูเชติ ตเมนํ พลิปฏิคฺคาหิกา เทวตา สกฺกตา ครุกตา มานิตา ปูชิตา กลฺยาเณน มนสา อนุกมฺปนฺติ “จิรํ ชีว, ทีฆมายุํ ปาเลหี”ติ.\csromanspacer{} เทวตานุกมฺปิตสฺส มหานาม กุลปุตฺตสฺส วุทฺธิเยว ปาฏิกงฺขา โน ปริหานิ.}

\prose{ปุน จปรํ มหานาม กุลปุตฺโต อุฏฺฐานวีริยาธิคเตหิ โภเคหิ พาหาพลปริจิเตหิ เสทาวกฺขิตฺเตหิ ธมฺมิเกหิ ธมฺมลทฺเธหิ สมณพฺราหฺมเณ สกฺกโรติ ครุํ กโรติ มาเนติ ปูเชติ, ตเมนํ สมณพฺราหฺมณา สกฺกตา ครุกตา มานิตา ปูชิตา กลฺยาเณน มนสา อนุกมฺปนฺติ “จิรํ ชีว, ทีฆมายุํ ปาเลหี”ติ.\csromanspacer{} สมณพฺราหฺมณานุกมฺปิตสฺส มหานาม กุลปุตฺตสฺส วุทฺธิเยว ปาฏิกงฺขา โน ปริหานิ.}

\prose{ยสฺส กสฺสจิ มหานาม กุลปุตฺตสฺส อิเม ปญฺจ ธมฺมา สํวิชฺชนฺติ, ยทิ วา รญฺโญ ขตฺติยสฺส มุทฺธาภิสิตฺตสฺส, ยทิ วา รฏฺฐิกสฺส}

\vspace{\tipitakaparskip}
\csromancenter{(6) 1.\csromanspacer{} นีวรณวคฺค เปตฺตนิกสฺส, ยทิ วา เสนาย เสนาปติกสฺส, ยทิ วา คามคามณิกสฺส, ยทิ วา ปูคคามณิกสฺส, เย วา ปน กุเลสุ ปจฺเจกาธิปจฺจํ กาเรนฺติ, วุทฺธิเยว ปาฏิกงฺขา โน ปริหานีติ.}
\vspace{0.5\baselineskip}
\csromangathameasure{มาตาปิตุกิจฺจกโร, ปุตฺตทารหิโต สทา. \\ อนฺโตชนสฺส อตฺถาย, เย จสฺส อนุชีวิโน. \\ อุภินฺนญฺเจว อตฺถาย, วทญฺญู โหติ สีลวา. \\ ญาตีนํ ปุพฺพเปตานํ, ทิฏฺเฐ ธมฺเม จ ชีวตํ. \\ สมณานํ พฺราหฺมณานํ, เทวตานญฺจ ปณฺฑิโต. \\ วิตฺติสญฺชนโน โหติ, ธมฺเมน ฆรมาวสํ. \\ โส กริตฺวาน กลฺยาณํ, ปุชฺโช โหติ ปสํสิโย. \\ อิเธว นํ ปสํสนฺติ, เปจฺจ สคฺเค ปโมทตีติ.อฏฺฐมํ.}
\csromangathagroup{\csromangathastanza{\csromanfolio{69}มาตาปิตุกิจฺจกโร, ปุตฺตทารหิโต สทา. \\ อนฺโตชนสฺส อตฺถาย, เย จสฺส อนุชีวิโน.}\csromangathastanza{อุภินฺนญฺเจว อตฺถาย, วทญฺญู โหติ สีลวา. \\ ญาตีนํ ปุพฺพเปตานํ, ทิฏฺเฐ ธมฺเม จ ชีวตํ\footnote{ชีวินํ (สี), ชีวิตํ (สฺยา, กํ, อิ, ก)}.}\csromangathastanza{สมณานํ พฺราหฺมณานํ, เทวตานญฺจ ปณฺฑิโต. \\ วิตฺติสญฺชนโน โหติ, ธมฺเมน ฆรมาวสํ.}\csromangathastanza{โส กริตฺวาน กลฺยาณํ, ปุชฺโช โหติ ปสํสิโย. \\ อิเธว นํ ปสํสนฺติ, เปจฺจ สคฺเค ปโมทตีติ.\csromanspacer{}\csromanspacer{}อฏฺฐมํ.}}

\csromanheader{2}{9. ปฐมวุฑฺฒปพฺพชิตสุตฺต}
\proseitem{59}{ปญฺจหิ ภิกฺขเว ธมฺเมหิ สมนฺนาคโต ทุลฺลโภ วุฒุปพฺพชิโต.\csromanspacer{} กตเมหิ ปญฺจหิ? ทุลฺลโภ ภิกฺขเว วุฒุปพฺพชิโต นิปุโณ, ทุลฺลโภ อากปฺปสมฺปนฺโน, ทุลฺลโภ พหุสฺสุโต, ทุลฺลโภ ธมฺมกถิโก, ทุลฺลโภ วินยธโร.\csromanspacer{} อิเมหิ โข ภิกฺขเว ปญฺจหิ ธมฺเมหิ สมนฺนาคโต ทุลฺลโภ วุฒุปพฺพชิโตติ.\csromanspacer{}\csromanspacer{}นวมํ.}

\csromanheader{2}{10. ทุติยวุฑฺฒปพฺพชิตสุตฺต}
\proseitem{60}{ปญฺจหิ ภิกฺขเว ธมฺเมหิ สมนฺนาคโต ทุลฺลโภ วุฒุปพฺพชิโต, กตเมหิ ปญฺจหิ? ทุลฺลโภ ภิกฺขเว วุฒุปพฺพชิโต สุวโจ, ทุลฺลโภ สุคฺคหิตคฺคาหิ, ทุลฺลโภ ปทกฺขิณคฺคาหี, ทุลฺลโภ ธมฺมกถิโก, ทุลฺลโภ วินยธโร.\csromanspacer{} อิเมหิ โข ภิกฺขเว ปญฺจหิ ธมฺเมหิ สมนฺนาคโต ทุลฺลโภ วุฒุปพฺพชิโตติ.\csromanspacer{}\csromanspacer{}ทสมํ.\csromanspacer{} นีวรณวคฺโค ปฐโม.}

\titlehead{\textbf{ตสฺสุทฺทานํ}}
\csromangathameasure{อาวรณํ ราสิ องฺคานิ, สมยํ มาตุปุตฺติกา. \\ อุปชฺฌา ฐานา ลิจฺฉวิ, กุมารา อปรา ทุเวติ.}
\csromangathagroup{\csromangathastanza{อาวรณํ ราสิ องฺคานิ, สมยํ มาตุปุตฺติกา. \\ อุปชฺฌา ฐานา ลิจฺฉวิ, กุมารา อปรา ทุเวติ.}}

\chapterheadpageread{\textbf{(7) 2. สญฺญาวคฺค}}
\csromanheader{2}{1. ปฐมสญฺญาสุตฺต}
\proseitem{61}{\csromanfolio{70}ปญฺจิมา ภิกฺขเว สญฺญา ภาวิตา พหุลีกตา มหปฺผลา โหนฺติ มหานิสํสา อมโตคธา อมตปริโยสานา.\csromanspacer{} กตมา ปญฺจ? อสุภสญฺญา มรณสญฺญา อาทีนวสญฺญา อาหาเร ปฏิกูลสญฺญา สพฺพโลเก อนภิรตสญฺญา\footnote{อนภิรติสญฺญา (ก), อุปริ 126, 244 ปิฏฺเฐสุ ปสฺสิตพฺพํ.}.\csromanspacer{} อิมา โข ภิกฺขเว ปญฺจ สญฺญา ภาวิตา พหุลีกตา มหปฺผลา โหนฺติ มหานิสํสา อมโตคธา อมตปริโยสานาติ.\csromanspacer{}\csromanspacer{}ปฐมํ.}

\csromanheader{2}{2. ทุติยสญฺญาสุตฺต}
\proseitem{62}{ปญฺจิมา ภิกฺขเว สญฺญา ภาวิตา พหุลีกตา มหปฺผลา โหนฺติ มหานิสํสา อมโตคธา อมตปริโยสานา.\csromanspacer{} กตมา ปญฺจ? อนิจฺจสญฺญา อนตฺตสญฺญา มรณสญฺญา อาหาเร ปฏิกูลสญฺญา สพฺพโลเก อนภิรตสญฺญา.\csromanspacer{} อิมา โข ภิกฺขเว ปญฺจ สญฺญา ภาวิตา พหุลีกตา มหปฺผลา โหนฺติ มหานิสํสา อมโตคธา อมตปริโยสานาติ.\csromanspacer{}\csromanspacer{}ทุติยํ.}

\csromanheader{2}{3. ปฐมวฑฺฒิสุตฺต}
\proseitem{63}{ปญฺจหิ ภิกฺขเว วฑฺฒีหิ วฑฺฒมาโน อริยสาวโก อริยาย วฑฺฒิยา วฑฺฒติ, สาราทายี จ โหติ วราทายี จ กายสฺส.\csromanspacer{} กตมาหิ ปญฺจหิ? สทฺธาย วฑฺฒติ, สีเลน วฑฺฒติ, สุเตน วฑฺฒติ, จาเคน วฑฺฒติ, ปญฺญาย วฑฺฒติ.\csromanspacer{} อิมาหิ โข ภิกฺขเว ปญฺจหิ วฑฺฒีหิ วฑฺฒมาโน อริยสาวโก อริยาย วฑฺฒิยา วฑฺฒติ, สาราทายี จ โหติ วราทายี จ กายสฺสาติ.}

\csromangathameasure{สทฺธาย สีเลน จ โย ปวฑฺฒติ, \\ ปญฺญาย จาเคน สุเตน จูภยํ. \\ โส ตาทิโส สปฺปุริโส วิจกฺขโณ, \\ อาทียตี สารมิเธว อตฺตโนติ.ตติยํ.}
\csromangathagroup{\csromangathastanza{สทฺธาย สีเลน จ โย ปวฑฺฒติ\footnote{โยธ วฑฺฒติ (สี)}, \\ ปญฺญาย จาเคน สุเตน จูภยํ. \\ โส ตาทิโส สปฺปุริโส วิจกฺขโณ, \\ อาทียตี สารมิเธว อตฺตโนติ.\csromanspacer{}\csromanspacer{}ตติยํ.}}

\chapterheadpageread{(7) 2. สญฺญาวคฺค \textbf{4. ทุติยวฑฺฒิสุตฺต}}
\proseitem{64}{\csromanfolio{71}ปญฺจหิ ภิกฺขเว วฑฺฒีหิ วฑฺฒมานา อริยสาวิกา อริยาย วฑฺฒิยา วฑฺฒติ, สาราทายินี จ โหติ วราทายินี จ กายสฺส.\csromanspacer{} กตมาหิ ปญฺจหิ? สทฺธาย วฑฺฒติ, สีเลน วฑฺฒติ, สุเตน วฑฺฒติ, จาเคน วฑฺฒติ, ปญฺญาย วฑฺฒติ.\csromanspacer{} อิมาหิ โข ภิกฺขเว ปญฺจหิ วฑฺฒีหิ วฑฺฒมานา อริยสาวิกา อริยาย วฑฺฒิยา วฑฺฒติ, สาราทายินี จ โหติ วราทายินี จ กายสฺสาติ.}

\csromangathameasure{สทฺธาย สีเลน จ ยา ปวฑฺฒติ, \\ ปญฺญาย จาเคน สุเตน จูภยํ. \\ สา ตาทิสี สีลวตี อุปาสิกา, \\ อาทียตี สารมิเธว อตฺตโนติ.จตุตฺถํ.}
\csromangathagroup{\csromangathastanza{สทฺธาย สีเลน จ ยา ปวฑฺฒติ\footnote{ยาธุ วฑฺฒติ (สี)}, \\ ปญฺญาย จาเคน สุเตน จูภยํ. \\ สา ตาทิสี สีลวตี อุปาสิกา, \\ อาทียตี สารมิเธว อตฺตโนติ.\csromanspacer{}\csromanspacer{}จตุตฺถํ.}}

\csromanheader{2}{5. สากจฺฉสุตฺต}
\proseitem{65}{\csromansymbolfootnote{*}{อุปริ 167 ปิฏฺเฐปิ.}ปญฺจหิ ภิกฺขเว ธมฺเมหิ สมนฺนาคโต ภิกฺขุ อลํสากจฺโฉ สพฺรหฺมจารีนํ.\csromanspacer{} กตเมหิ ปญฺจหิ? อิธ ภิกฺขเว ภิกฺขุ อตฺตนา จ สีลสมฺปนฺโน โหติ, สีลสมฺปทาย กถาย จ อาคตํ ปญฺหํ พฺยากตฺตา โหติ, อตฺตนา จ สมาธิสมฺปนฺโน โหติ, สมาธิสมฺปทาย กถาย จ อาคตํ ปญฺหํ พฺยากตฺตา โหติ, อตฺตนา จ ปญฺญาสมฺปนฺโน โหติ, ปญฺญาสมฺปทาย กถาย จ อาคตํ ปญฺหํ พฺยากตฺตา โหติ, อตฺตนา จ วิมุตฺติสมฺปนฺโน โหติ, วิมุตฺติสมฺปทาย กถาย จ อาคตํ ปญฺหํ พฺยากตฺตา โหติ, อตฺตนา จ วิมุตฺติญาณทสฺสนสมฺปนฺโน โหติ, วิมุตฺติญาณทสฺสน สมฺปทาย กถาย จ อาคตํ ปญฺหํ พฺยากตฺตา โหติ.\csromanspacer{} อิเมหิ โข ภิกฺขเว ปญฺจหิ ธมฺเมหิ สมนฺนาคโต ภิกฺขุ อลํสากจฺโฉ สพฺรหฺมจารีนนฺติ.\csromanspacer{}\csromanspacer{}ปญฺจมํ.}

\csromanheader{2}{6. สาชีวสุตฺต}
\proseitem{66}{\csromansymbolfootnote{+}{อุปริ 168 ปิฏฺเฐปิ.}ปญฺจหิ ภิกฺขเว ธมฺเมหิ สมนฺนาคโต ภิกฺขุ อลํสาชีโว สพฺรหฺมจารีนํ.\csromanspacer{} กตเมหิ ปญฺจหิ? อิธ ภิกฺขเว ภิกฺขุ อตฺตนา จ สีลสมฺปนฺโน โหติ, สีลสมฺปทาย กถาย จ กตํ ปญฺหํ พฺยากตฺตา โหติ, อตฺตนา จ สมาธิสมฺปนฺโน โหติ, สมาธิสมฺปทาย กถาย \csromanfolio{72}จ กตํ ปญฺหํ พฺยากตฺตา โหติ, อตฺตนา จ ปญฺญาสมฺปนฺโน โหติ, ปญฺญาสมฺปทาย กถาย จ กตํ ปญฺหํ พฺยากตฺตา โหติ, อตฺตนา จ วิมุตฺติสมฺปนฺโน โหติ, วิมุตฺติสมฺปทาย กถาย จ กตํ ปญฺหํ พฺยากตฺตา โหติ, อตฺตนา จ วิมุตฺติญาณทสฺสนสมฺปนฺโน โหติ, วิมุตฺติญาณทสฺสนสมฺปทาย กถาย จ กตํ ปญฺหํ พฺยากตฺตา โหติ.\csromanspacer{} อิเมหิ โข ภิกฺขเว ปญฺจหิ ธมฺเมหิ สมนฺนาคโต ภิกฺขุ อลํสาชีโว สพฺรหฺมจารีนนฺติ.\csromanspacer{}\csromanspacer{}ฉฏฺฐํ.}

\csromanheader{2}{7. ปฐม อิทฺธิปาทสุตฺต}
\proseitem{67}{โย หิ โกจิ ภิกฺขเว ภิกฺขุ วา ภิกฺขูนี วา ปญฺจ ธมฺเม\footnote{อิเม ปญฺจ ธมฺเม (ก)} ภาเวติ, ปญฺจ ธมฺเม1 พหุลีกโรติ, ตสฺส ทฺวินฺนํ ผลานํ อญฺญตรํ ผลํ ปาฏิกงฺขํ ทิฏฺเฐว ธมฺเม อญฺญา, สติ วา อุปาทิเสเส อนาคามิตา.}

\prose{กตเม ปญฺจ? อิธ ภิกฺขเว ภิกฺขุ ฉนฺทสมาธิปธานสงฺขารสมนฺนาคตํ อิทฺธิปาทํ ภาเวติ.\csromanspacer{} วีริยสมาธิ.\csromanspacer{} จิตฺตสมาธิ.\csromanspacer{} วีมํสาสมาธิปธานสงฺขารสมนฺนาคตํ อิทฺธิปาทํ ภาเวติ อุสฺโสฬฺหิญฺเญว ปญฺจมิํ\footnote{อุสฺโสฬฺหีเยว ปญฺจมี (สี)}.\csromanspacer{} โย หิ โกจิ ภิกฺขเว ภิกฺขุ วา ภิกฺขุนี วา อิเม ปญฺจ ธมฺเม ภาเวติ, อิเม ปญฺจ ธมฺเม พหุลีกโรติ, ตสฺส ทฺวินฺนํ ผลานํ อญฺญตรํ ผลํ ปาฏิกงฺขํ ทิฏฺเฐว ธมฺเม อญฺญา, สติ วา อุปาทิเสเส อนาคามิตาติ.\csromanspacer{}\csromanspacer{}สตฺตมํ.}

\csromanheader{2}{8. ทุติย อิทฺธิปาทสุตฺต}
\proseitem{68}{ปุพฺเพวาหํ ภิกฺขเว สมฺโพธา อนภิสมฺพุทฺโธ โพธิสตฺโตว สมาโน ปญฺจ ธมฺเม ภาเวสิํ, ปญฺจ ธมฺเม พหุลีกาสิํ\footnote{พหุลิมกาสิํ (ก), พหุลมกาสิํ (ก)}.\csromanspacer{} กตเม ปญฺจ? ฉนฺทสมา ธิปธานสงฺขารสมนฺนาคตํ อิทฺธิปาทํ ภาเวสิํ.\csromanspacer{} วีริยสมาธิ.\csromanspacer{} จิตฺต สมาธิ.\csromanspacer{} วีมํสาสมาธิปธานสงฺขารสมนฺนาคตํ อิทฺธิปาทํ ภาเวสิํ อุสฺโสฬฺหิญฺเญว ปญฺจมิํ.\csromanspacer{} โส โข อหํ ภิกฺขเว อิเมสํ อุสฺโสฬฺหิปญฺจมานํ ธมฺมานํ ภาวิตตฺตา พหุลีกตตฺตา ยสฺส ยสฺส อภิญฺญาสจฺฉิกรณียสฺส ธมฺมสฺส จิตฺตํ อภินินฺนาเมสิํ อภิญฺญาสจฺฉิกิริยาย, ตตฺร ตเตฺรว สกฺขิภพฺพตํ ปาปุณิํ สติ สติ อายตเน.}

\chapterheadpageread{(7) 2. สญฺญาวคฺค}
\prose{\csromanfolio{73}โส สเจ อากงฺขิํ “อเนกวิหิตํ อิทฺธิวิธํ ปจฺจนุภเวยฺยํ ฯเปฯ ยาว พฺรหฺมโลกาปิ กาเยน วสํ วตฺเตยฺยนฺ”ติ, ตตฺร ตเตฺรว สกฺขิภพฺพตํ ปาปุณิํ สติ สติ อายตเน.}

\prose{โส สเจ อากงฺขิํ ฯเปฯ “อาสวานํ ขยา อนาสวํ เจโตวิมุตฺติํ ปญฺญาวิมุตฺติํ ทิฏฺเฐว ธมฺเม สยํ อภิญฺญา สจฺฉิกตฺวา อุปสมฺปชฺช วิหเรยฺยนฺ”ติ, ตตฺร ตเตฺรว สกฺขิภพฺพตํ ปาปุณิํ สติ สติ อายตเนติ.\csromanspacer{}\csromanspacer{}อฏฺฐมํ.}

\csromanheader{2}{9. นิพฺพิทาสุตฺต}
\proseitem{69}{ปญฺจิเม ภิกฺขเว ธมฺมา ภาวิตา พหุลีกตา เอกนฺตนิพฺพิทาย วิราคาย นิโรธาย อุปสมาย อภิญฺญาย สมฺโพธาย นิพฺพานาย สํวตฺตนฺติ.}

\prose{กตเม ปญฺจ? อิธ ภิกฺขเว ภิกฺขุ อสุภานุปสฺสี กาเย วิหรติ, อาหาเร ปฏิกูลสญฺญี, สพฺพโลเก อนภิรตสญฺญี\footnote{อนภิรติสญฺญี (ก), อุปริ 126, 244 ปิฏฺเฐสุ ปสฺสิตพฺพํ.}, สพฺพสงฺขาเรสุ อนิจฺจานุปสฺสี, มรณสญฺญา โข ปนสฺส อชฺฌตฺตํ สูปฏฺฐิตา โหติ.\csromanspacer{} อิเม โข ภิกฺขเว ปญฺจ ธมฺมา ภาวิตา พหุลีกตา เอกนฺตนิพฺพิทาย วิราคาย นิโรธาย อุปสมาย อภิญฺญาย สมฺโพธาย นิพฺพานาย สํวตฺตนฺตีติ.\csromanspacer{}\csromanspacer{}นวมํ.}

\csromanheader{2}{10. อาสวกฺขยสุตฺต}
\proseitem{70}{ปญฺจิเม ภิกฺขเว ธมฺมา ภาวิตา พหุลีกตา อาสวานํ ขยาย สํวตฺตนฺติ.\csromanspacer{} กตเม ปญฺจ? อิธ ภิกฺขเว ภิกฺขุ อสุภานุปสฺสี กาเย วิหรติ, อาหาเร ปฏิกูลสญฺญี, สพฺพโลเก อนภิรตสญฺญี, สพฺพสงฺขาเรสุ อนิจฺจานุปสฺสี, มรณสญฺญา โข ปนสฺส อชฺฌตฺตํ สูปฏฺฐิตา โหติ.\csromanspacer{} อิเม โข ภิกฺขเว ปญฺจ ธมฺมา ภาวิตา พหุลีกตา อาสวานํ ขยาย สํวตฺตนฺตีติ.\csromanspacer{}\csromanspacer{}ทสมํ.\csromanspacer{} สญฺญาวคฺโค ทุติโย.}

\titlehead{\textbf{ตสฺสุทฺทานํ}}
\csromangathameasure{เทฺว จ สญฺญา เทฺว วฑฺฒี จ, สากจฺเฉน จ สาชีวํ. \\ อิทฺธิปาทา จ เทฺว วุตฺตา, นิพฺพิทา จาสวกฺขยาติ.}
\csromangathagroup{\csromangathastanza{เทฺว จ สญฺญา เทฺว วฑฺฒี จ, สากจฺเฉน จ สาชีวํ. \\ อิทฺธิปาทา จ เทฺว วุตฺตา, นิพฺพิทา จาสวกฺขยาติ.}}

\chapterheadpageread{\textbf{(8) 3. โยธาชีววคฺค}}
\csromanheader{2}{1. ปฐมเจโตวิมุตฺติผลสุตฺต}
\proseitem{71}{\csromanfolio{74}ปญฺจิเม ภิกฺขเว ธมฺมา ภาวิตา พหุลีกตา เจโตวิมุตฺติผลา จ โหนฺติ เจโตวิมุตฺติผลานิสํสา จ, ปญฺญาวิมุตฺติผลา จ โหนฺติ ปญฺญาวิมุตฺติผลานิสํสา จ.}

\prose{กตเม ปญฺจ? อิธ ภิกฺขเว ภิกฺขุ อสุภานุปสฺสี กาเย วิหรติ, อาหาเร ปฏิกูลสญฺญี\footnote{ปฏิกฺกูลสญฺญี (สี, สฺยา, กํ, อิ)}, สพฺพโลเก อนภิรตสญฺญี, สพฺพสงฺขาเรสุ อนิจฺจานุปสฺสี, มรณ สญฺญา โข ปนสฺส อชฺฌตฺตํ สูปฏฺฐิตา โหติ.\csromanspacer{} อิเม โข ภิกฺขเว ปญฺจ ธมฺมา ภาวิตา พหุลีกตา เจโตวิมุตฺติผลา จ โหนฺติ เจโตวิมุตฺติผลานิสํสา จ, ปญฺญาวิมุตฺติ ผลา จ โหนฺติ ปญฺญาวิมุตฺติผลานิสํสา จ.\csromanspacer{} ยโต โข ภิกฺขเว ภิกฺขุ เจโตวิมุตฺโต จ โหติ, ปญฺญาวิมุตฺโต จ โหติ.\csromanspacer{} อยํ วุจฺจติ ภิกฺขเว ภิกฺขุ “อุกฺขิตฺตปลิโฆ” อิติปิ “สํกิณฺณปริโข”\footnote{สํกิณฺณปริกฺโข (สฺยา, กํ)} อิติปิ “อพฺพูเฬฺหสิโก” อิติปิ “นิรคฺคโฬ” อิติปิ “อริโย ปนฺนทฺธโช ปนฺนภาโร วิสํยุตฺโต” อิติปิ.}

\prose{กถญฺจ ภิกฺขเว ภิกฺขุ อุกฺขิตฺตปลิโฆ โหติ? อิธ ภิกฺขเว ภิกฺขุโน อวิชฺชา ปหีนา โหติ อุจฺฉินฺนมูลา ตาลาวตฺถุกตา อนภาวํกตา อายติํ อนุปฺปาท ธมฺมา.\csromanspacer{} เอวํ โข ภิกฺขเว ภิกฺขุ อุกฺขิตฺตปลิโฆ}

\prose{กถญฺจ ภิกฺขเว ภิกฺขุ สํกิณฺณปริโข โหติ? อิธ ภิกฺขเว ภิกฺขุโน โปโนภวิโก\footnote{โปโนพฺภวิโก (สฺยา, ก)} ชาติสํสาโร ปหีโน โหติ อุจฺฉินฺนมูโล ตาลาวตฺถุกโต อนภาวํกโต อายติํ อนุปฺปาทธมฺโม.\csromanspacer{} เอวํ โข ภิกฺขเว ภิกฺขุ สํกิณฺณปริโข โหติ.}

\prose{กถญฺจ ภิกฺขเว ภิกฺขุ อพฺพูเฬฺหสิโก โหติ? อิธ ภิกฺขเว ภิกฺขุโน ตณฺหา ปหีนา โหติ อุจฺฉินฺนมูลา ตาลาวตฺถุกตา อนภาวํกตา อายติํ อนุปฺปาท ธมฺมา.\csromanspacer{} เอวํ โข ภิกฺขเว ภิกฺขุ อพฺพูเฬฺหสิโก โหติ.}

\chapterheadpageread{(8) 3. โยธาชีววคฺค}
\prose{\csromanfolio{75}กถญฺจ ภิกฺขเว ภิกฺขุ นิรคฺคโฬ โหติ? อิธ ภิกฺขเว ภิกฺขุโน ปญฺโจรมฺภาคิยานิ สํโยชนานิ ปหีนานิ โหนฺติ อุจฺฉินฺนมูลานิ ตาลาวตฺถุกตานิ อนภาวํกตานิ อายติํ อนุปฺปาทธมฺมานิ.\csromanspacer{} เอวํ โข ภิกฺขเว ภิกฺขุ นิรคฺคโฬ โหติ.}

\prose{กถญฺจ ภิกฺขเว ภิกฺขุ อริโย ปนฺนทฺธโช ปนฺนภาโร วิสํยุตฺโต โหติ? อิธ ภิกฺขเว ภิกฺขุโน อสฺมิมาโน ปหีโน โหติ อุจฺฉินฺนมูโล ตาลาวตฺถุกโต อนภาวํกโต อายติํ อนุปฺปาทธมฺโม.\csromanspacer{} เอวํ โข ภิกฺขเว ภิกฺขุ อริโย ปนฺนทฺธโช ปนฺนภาโร วิสํยุตฺโต โหตีติ.\csromanspacer{} ปฐมํ.}

\csromanheader{2}{2. ทุติยเจโตวิมุตฺติผลสุตฺต}
\proseitem{72}{ปญฺจิเม ภิกฺขเว ธมฺมา ภาวิตา พหุลีกตา เจโตวิมุตฺติผลา จ โหนฺติ เจโตวิมุตฺติผลานิสํสา จ, ปญฺญาวิมุตฺติผลา จ โหนฺติ ปญฺญาวิมุตฺติผลานิสํสา จ.\csromanspacer{} กตเม ปญฺจ? อนิจฺจสญฺญา, อนิจฺเจ ทุกฺขสญฺญา, ทุกฺเข อนตฺตสญฺญา, ปหานสญฺญา, วิราคสญฺญา.\csromanspacer{} อิเม โข ภิกฺขเว ปญฺจ ธมฺมา ภาวิตา พหุลีกตา เจโตวิมุตฺติผลา จ โหนฺติ เจโตวิมุตฺติผลานิสํสา จ, ปญฺญาวิมุตฺติผลา จ โหนฺติ ปญฺญาวิมุตฺติ ผลานิสํสา จ.\csromanspacer{} ยโต โข ภิกฺขเว ภิกฺขุ เจโตวิมุตฺโต จ โหติ ปญฺญาวิมุตฺโต จ.\csromanspacer{} อยํ วุจฺจติ ภิกฺขเว ภิกฺขุ “อุกฺขิตฺตปลิโฆ” อิติปิ “สํกิณฺณปริโข” อิติปิ “อพฺพูเฬฺหสิโก” อิติปิ “นิรคฺคโฬ” อิติปิ “อริโย ปนฺนทฺธโช ปนฺนภาโร วิสํยุตฺโต” อิติปิ.}

\prose{กถญฺจ ภิกฺขเว ภิกฺขุ อุกฺขิตฺตปลิโฆ โหติ? อิธ ภิกฺขเว ภิกฺขุโน อวิชฺชา ปหีนา โหติ อุจฺฉินฺนมูลา ตาลาวตฺถุกตา อนภาวํกตา อายติํ อนุปฺปาท ธมฺมา.\csromanspacer{} เอวํ โข ภิกฺขเว ภิกฺขุ อุกฺขิตฺตปลิโฆ}

\prose{กถญฺจ ภิกฺขเว ภิกฺขุ สํกิณฺณปริโข โหติ? อิธ ภิกฺขเว ภิกฺขุโน โปโนภวิโก ชาติสํสาโร ปหีโน โหติ อุจฺฉินฺนมูโล ตาลาวตฺถุกโต อนภาวํกโต อายติํ อนุปฺปาทธมฺโม.\csromanspacer{} เอวํ โข ภิกฺขเว ภิกฺขุ สํกิณฺณปริโข โหติ.}

\prose{กถญฺจ ภิกฺขเว ภิกฺขุ อพฺพูเฬฺหสิโก โหติ? อิธ ภิกฺขเว ภิกฺขุโน ตณฺหา ปหีนา โหติ อุจฺฉินฺนมูลา ตาลาวตฺถุกตา \csromanfolio{76}อนภาวํกตา อายติํ อนุปฺปาทธมฺมา.\csromanspacer{} เอวํ โข ภิกฺขเว ภิกฺขุ อพฺพูเฬฺหสิโก โหติ.}

\prose{กถญฺจ ภิกฺขเว ภิกฺขุ นิรคฺคโฬ โหติ? อิธ ภิกฺขเว ภิกฺขุโน ปญฺโจรมฺภาคิยานิ สํโยชนานิ ปหีนานิ โหนฺติ อุจฺฉินฺนมูลานิ ตาลาวตฺถุกตานิ อนภาวํกตานิ อายติํ อนุปฺปาทธมฺมานิ.\csromanspacer{} เอวํ โข ภิกฺขเว ภิกฺขุ นิรคฺคโฬ โหติ.}

\prose{กถญฺจ ภิกฺขเว ภิกฺขุ อริโย ปนฺนทฺธโช ปนฺนภาโร วิสํยุตฺโต โหติ? อิธ ภิกฺขเว ภิกฺขุโน อสฺมิมาโน ปหีโน โหติ อุจฺฉินฺนมูโล ตาลาวตฺถุกโต อนภาวํกโต อายติํ อนุปฺปาทธมฺโม.\csromanspacer{} เอวํ โข ภิกฺขเว ภิกฺขุ อริโย ปนฺนทฺธโช ปนฺนภาโร วิสํยุตฺโต โหตีติ.\csromanspacer{}\csromanspacer{}ทุติยํ.}

\csromanheader{2}{3. ปฐมธมฺมวิหารีสุตฺต}
\proseitem{73}{อถ โข อญฺญตโร ภิกฺขุ เยน ภควา เตนุปสงฺกมิ, อุปสงฺกมิตฺวา ภควนฺตํ อภิวาเทตฺวา เอกมนฺตํ นิสีทิ, เอกมนฺตํ นิสินฺโน โข โส ภิกฺขุ ภควนฺตํ เอตทโวจ “ธมฺมวิหารี ธมฺมวิหารีติ ภนฺเต วุจฺจติ, กิตฺตาวตา นุ โข ภนฺเต ภิกฺขุ ธมฺมวิหารี โหตี”ติ.}

\prose{อิธ ภิกฺขุ ภิกฺขุ ธมฺมํ ปริยาปุณาติ สุตฺตํ เคยฺยํ เวยฺยากรณํ คาถํ อุทานํ อิติวุตฺตกํ ชาตกํ อพฺภุตธมฺมํ เวทลฺลํ, โส ตาย ธมฺมปริยตฺติยา ทิวสํ อตินาเมติ, ริญฺจติ ปฏิสลฺลานํ, นานุยุญฺชติ อชฺฌตฺตํ เจโตสมถํ.\csromanspacer{} อยํ วุจฺจติ ภิกฺขุ ภิกฺขุ ปริยตฺติพหุโล โน ธมฺมวิหารี.}

\prose{ปุน จปรํ ภิกฺขุ ภิกฺขุ ยถาสุตํ ยถาปริยตฺตํ ธมฺมํ วิตฺถาเรน ปเรสํ เทเสติ, โส ตาย ธมฺมปญฺญตฺติยา ทิวสํ อตินาเมติ, ริญฺจติ ปฏิสลฺลานํ, นานุยุญฺชติ อชฺฌตฺตํ เจโตสมถํ.\csromanspacer{} อยํ วุจฺจติ ภิกฺขุ ภิกฺขุ ปญฺญตฺติพหุโล โน ธมฺมวิหารี.}

\prose{ปุน จปรํ ภิกฺขุ ภิกฺขุ ยถาสุตํ ยถาปริยตฺตํ ธมฺมํ วิตฺถาเรน สชฺฌายํ กโรติ, โส เตน สชฺฌาเยน ทิวสํ อตินาเมติ, ริญฺจติ ปฏิสลฺลานํ, นานุยุญฺชติ อชฺฌตฺตํ เจโตสมถํ.\csromanspacer{} อยํ วุจฺจติ ภิกฺขุ ภิกฺขุ สชฺฌายพหุโล โน ธมฺมวิหารี.}

\chapterheadpageread{(8) 3. โยธาชีววคฺค}
\prose{\csromanfolio{77}ปุน จปรํ ภิกฺขุ ภิกฺขุ ยถาสุตํ ยถาปริยตฺตํ ธมฺมํ เจตสา อนุวิตกฺเกติ อนุจิจาเรติ มนสานุเปกฺขติ, โส เตหิ ธมฺมวิตกฺเกหิ ทิวสํ อตินาเมติ, ริญฺจติ ปฏิสลฺลานํ, นานุยุญฺชติ อชฺฌตฺตํ เจโตสมถํ.\csromanspacer{} อยํ วุจฺจติ ภิกฺขุ ภิกฺขุ วิตกฺกพหุโล โน ธมฺมวิหารี.}

\prose{อิธ ภิกฺขุ ภิกฺขุ ธมฺมํ ปริยาปุณาติ สุตฺตํ เคยฺยํ เวยฺยากรณํ คาถํ อุทานํ อิติวุตฺตกํ ชาตกํ อพฺภุตธมฺมํ เวทลฺลํ, โส ตาย ธมฺมปริยตฺติยา น ทิวสํ อตินาเมติ, นาปิ ริญฺจติ ปฏิสลฺลานํ, อนุยุญฺชติ อชฺฌตฺตํ เจโตสมถํ.\csromanspacer{} เอวํ โข ภิกฺขุ ภิกฺขุ ธมฺมวิหารี โหติ.}

\prose{อิติ โข ภิกฺขุ เทสิโต มยา ปริยตฺติพหุโล, เทสิโต ปญฺญตฺติพหุโล, เทสิโต สชฺฌายพหุโล, เทสิโต วิตกฺกพหุโล, เทสิโต ธมฺมวิหารี.\csromanspacer{} ยํ โข ภิกฺขุ\footnote{ยํ ภิกฺขุ (สฺยา, กํ, อิ)} สตฺถารา กรณียํ สาวกานํ หิเตสินา อนุกมฺปเกน อนุกมฺปํ อุปาทาย, กตํ โว ตํ มยา.\csromanspacer{} เอตานิ ภิกฺขุ รุกฺขมูลานิ เอตานิ สุญฺญาคารานิ, ฌายถ ภิกฺขุ มา ปมาทตฺถ, มา ปจฺฉา วิปฺปฏิสาริโน อหุวตฺถ, อยํ โว อมฺหากํ อนุสาสนีติ.\csromanspacer{}\csromanspacer{}ตติยํ.}

\csromanheader{2}{4. ทุติยธมฺมวิหารีสุตฺต}
\proseitem{74}{อถ โข อญฺญตโร ภิกฺขุ เยน ภควา เตนุปสงฺกมิ, อุปสงฺกมิตฺวา ภควนฺตํ อภิวาเทตฺวา เอกมนฺตํ นิสีทิ, เอกมนฺตํ นิสินฺโน โข โส ภิกฺขุ ภควนฺตํ เอตทโวจ “ธมฺมวิหารี ธมฺมวิหารีติ ภนฺเต วุจฺจติ, กิตฺตาวตา นุ โข ภนฺเต ภิกฺขุ ธมฺมวิหารี โหตี”ติ.}

\prose{อิธ ภิกฺขุ ภิกฺขุ ธมฺมํ ปริยาปุณาติ สุตฺตํ เคยฺยํ เวยฺยากรณํ คาถํ อุทานํ อิติวุตฺตกํ ชาตกํ อพฺภุตธมฺมํ เวทลฺลํ, อุตฺตริ\footnote{อุตฺตริํ (สี, สฺยา, กํ, อิ)} จสฺส ปญฺญาย อตฺถํ นปฺปชานาติ.\csromanspacer{} อยํ วุจฺจติ ภิกฺขุ ภิกฺขุ ปริยตฺติพหุโล โน ธมฺมวิหารี.}

\prose{ปุน จปรํ ภิกฺขุ ภิกฺขุ ยถาสุตํ ยถาปริยตฺตํ ธมฺมํ วิตฺถาเรน ปเรสํ เทเสติ, อุตฺตริ จสฺส ปญฺญาย อตฺถํ นปฺปชานาติ.\csromanspacer{} อยํ วุจฺจติ ภิกฺขุ ภิกฺขุ ปญฺญตฺติพหุโล โน ธมฺมวิหารี.}

\prose{\csromanfolio{78}ปุน จปรํ ภิกฺขุ ภิกฺขุ ยถาสุตํ ยถาปริยตฺตํ ธมฺมํ วิตฺถาเรน สชฺฌายํ กโรติ, อุตฺตริ จสฺส ปญฺญาย อตฺถํ นปฺปชานาติ.\csromanspacer{} อยํ วุจฺจติ ภิกฺขุ ภิกฺขุ สชฺฌายพหุโล โน ธมฺมวิหารี.}

\prose{ปุน จปรํ ภิกฺขุ ภิกฺขุ ยถาสุตํ ยถาปริยตฺตํ ธมฺมํ เจตสา อนุวิตกฺเกติ อนุวิจาเรติ มนสานุเปกฺขติ, อุตฺตริ จสฺส ปญฺญาย อตฺถํ นปฺปชานาติ.\csromanspacer{} อยํ วุจฺจติ ภิกฺขุ ภิกฺขุ วิตกฺกพหุโล โน ธมฺมวิหารี.}

\prose{อิธ ภิกฺขุ ภิกฺขุ ธมฺมํ ปริยาปุณาติ สุตฺตํ เคยฺยํ เวยฺยากรณํ คาถํ อุทานํ อิติวุตฺตกํ ชาตกํ อพฺภุตธมฺมํ เวทลฺลํ, อุตฺตริ จสฺส ปญฺญาย อตฺถํ ปชานาติ.\csromanspacer{} เอวํ โข ภิกฺขุ ภิกฺขุ ธมฺมวิหารี โหติ.}

\prose{อิติ โข ภิกฺขุ เทสิโต มยา ปริยตฺติพหุโล, เทสิโต ปญฺญตฺติพหุโล, เทสิโต สชฺฌายพหุโล, เทสิโต วิตกฺกพหุโล, เทสิโต ธมฺมวิหารี.\csromanspacer{} ยํ โข ภิกฺขุ สตฺถารา กรณียํ สาวกานํ หิเตสินา อนุกมฺปเกน อนุกมฺปํ อุปาทาย, กตํ โว ตํ มยา.\csromanspacer{} เอตานิ ภิกฺขุ รุกฺขมูลานิ เอตานิ สุญฺญาคารานิ, ฌายถ ภิกฺขุ มา ปมาทตฺถ, มา ปจฺฉา วิปฺปฏิสาริโน อหุวตฺถ, อยํ โว อมฺหากํ อนุสาสนีติ.\csromanspacer{}\csromanspacer{}จตุตฺถํ.}

\csromanheader{2}{5. ปฐมโยธาชีวสุตฺต}
\proseitem{75}{ปญฺจิเม ภิกฺขเว โยธาชีวา สนฺโต สํวิชฺชมานา โลกสฺมิํ.\csromanspacer{} กตเม ปญฺจ? อิธ ภิกฺขเว เอกจฺโจ โยธาชีโว รชคฺคญฺเญว ทิสฺวา สํสีทติ วิสีทติ น สนฺถมฺภติ น สกฺโกติ สงฺคามํ โอตริตุํ, เอวรูโปปิ\footnote{เอวรูโป (สี) อภิ 3. 175 ปิฏฺเฐปิ.} ภิกฺขเว อิเธกจฺโจ\footnote{เอกจฺโจ (สี)} โยธาชีโว โหติ.\csromanspacer{} อยํ ภิกฺขเว ปฐโม โยธาชีโว สนฺโต สํวิชฺชมาโน โลกสฺมิํ.}

\prose{ปุน จปรํ ภิกฺขเว อิเธกจฺโจ โยธาชีโว สหติ รชคฺคํ, อปิ จ โข ธชคฺคญฺเญว ทิสฺวา สํสีทติ วิสีทติ น สนฺถมฺภติ น สกฺโกติ สงฺคามํ โอตริตุํ, เอวรูโปปิ ภิกฺขเว อิเธกจฺโจ โยธาชีโว โหติ.\csromanspacer{} อยํ ภิกฺขเว ทุติโย โยธาชีโว สนฺโต สํวิชฺชมาโน โลกสฺมิํ.}

\prose{ปุน จปรํ ภิกฺขเว อิเธกจฺโจ โยธาชีโว สหติ รชคฺคํ, สหติ ธชคฺคํ, อปิ จ โข อุสฺสารณญฺเญว\footnote{อุสฺสาทนํเยว (สี, อิ)} สุตฺวา สํสีทติ วิสีทติ น สนฺถมฺภติ}

\vspace{\tipitakaparskip}
\csromancenter{(8) 3.\csromanspacer{} โยธาชีววคฺค น สกฺโกติ สงฺคามํ โอตริตุํ, เอวรูโปปิ ภิกฺขเว อิเธกจฺโจ โยธาชีโว โหติ.\csromanspacer{} อยํ ภิกฺขเว ตติโย โยธาชีโว สนฺโต สํวิชฺชมาโน โลกสฺมิํ.}
\prose{\csromanfolio{79}ปุน จปรํ ภิกฺขเว อิเธกจฺโจ โยธาชีโว สหติ รชคฺคํ, สหติ ธชคฺคํ, สหติ อุสฺสารณํ, อปิ จ โข สมฺปหาเร หญฺญติ\footnote{อาหญฺญติ (สี)} พฺยาปชฺชติ, เอวรูโปปิ ภิกฺขเว อิเธกจฺโจ โยธาชีโว โหติ.\csromanspacer{} อยํ ภิกฺขเว จตุตฺโถ โยธาชีโว สนฺโต สํวิชฺชมาโน โลกสฺมิํ.}

\prose{ปุน จปรํ ภิกฺขเว อิเธกจฺโจ โยธาชีโว สหติ รชคฺคํ, สหติ ธชคฺคํ, สหติ อุสฺสารณํ, สหติ สมฺปหารํ, โส ตํ สงฺคามํ อภิวิชินิตฺวา วิชิตสงฺคาโม ตเมว สงฺคามสีสํ อชฺฌาวสติ, เอวรูโปปิ ภิกฺขเว อิเธกจฺโจ โยธาชีโว โหติ.\csromanspacer{} อยํ ภิกฺขเว ปญฺจโม โยธาชีโว สนฺโต สํวิชฺชมาโน โลกสฺมิํ.\csromanspacer{} อิเม โข ภิกฺขเว ปญฺจ โยธาชีวา สนฺโต สํวิชฺชมานา โลกสฺมิํ.}

\prose{เอวเมวํ โข ภิกฺขเว ปญฺจิเม โยธาชีวูปมา ปุคฺคลา สนฺโต สํวิชฺชมานา ภิกฺขูสุ.\csromanspacer{} กตเม ปญฺจ? อิธ ภิกฺขเว ภิกฺขุ รชคฺคญฺเญว ทิสฺวา สํสีทติ วิสีทติ น สนฺถมฺภติ น สกฺโกติ พฺรหฺมจริยํ สนฺธาเรตุํ\footnote{สนฺถาเนตุํ (สี, สฺยา, กํ, อิ)}.\csromanspacer{} สิกฺขาทุพฺพลฺยํ อาวิกตฺวา สิกฺขํ ปจฺจกฺขาย หีนายาวตฺตติ.\csromanspacer{} กิมสฺส รชคฺคสฺมิํ, อิธ ภิกฺขเว ภิกฺขุ สุณาติ “อมุกสฺมิํ นาม คาเม วา นิคเม วา อิตฺถี วา กุมารี วา อภิรูปา ทสฺสนียา ปาสาทิกา ปรมาย วณฺณโปกฺขรตาย สมนฺนาคตา”ติ, โส ตํ สุตฺวา สํสีทติ วิสีทติ น สนฺถมฺภติ น สกฺโกติ พฺรหฺมจริยํ สนฺธาเรตุํ, สิกฺขาทุพฺพลฺยํ อาวิกตฺวา สิกฺขํ ปจฺจกฺขาย หีนายาวตฺตติ, อิทมสฺส รชคฺคสฺมิํ.}

\prose{เสยฺยถาปิ โส ภิกฺขเว โยธาชีโว รชคฺคญฺเญว ทิสฺวา สํสีทติ วิสีทติ น สนฺถมฺภติ น สกฺโกติ สงฺคามํ โอตริตุํ, ตถูปมาหํ ภิกฺขเว อิมํ ปุคฺคลํ วทามิ, เอวรูโปปิ ภิกฺขเว อิเธกจฺโจ ปุคฺคโล โหติ.\csromanspacer{} อยํ ภิกฺขเว ปฐโม โยธาชีวูปโม ปุคฺคโล สนฺโต สํวิชฺชมาโน ภิกฺขูสุ.}

\prose{ปุน จปรํ ภิกฺขเว ภิกฺขุ สหติ รชคฺคํ, อปิ จ โข ธชคฺคญฺเญว ทิสฺวา สํสีทติ วิสีทติ น สนฺถมฺภติ น สกฺโกติ พฺรหฺมจริยํ สนฺธาเรตุํ, \csromanfolio{80}สิกฺขาทุพฺพลฺยํ อาวิกตฺวา สิกฺขํ ปจฺจกฺขาย หีนายาวตฺตติ.\csromanspacer{} กิมสฺส ธชคฺคสฺมิํ, อิธ ภิกฺขเว ภิกฺขุ น เหว โข สุณาติ “อมุกสฺมิํ นาม คาเม วา นิคเม วา อิตฺถี วา กุมารี วา อภิรูปา ทสฺสนียา ปาสาทิกา ปรมาย วณฺณโปกฺขรตาย สมนฺนาคตา”ติ, อปิ จ โข สามํ ปสฺสติ อิตฺถิํ วา กุมาริํ วา อภิรูปํ ทสฺสนียํ ปาสาทิกํ ปรมาย วณฺณโปกฺขรตาย สมนฺนาคตํ.\csromanspacer{} โส ตํ ทิสฺวา สํสีทติ วิสีทติ น สนฺถมฺภติ น สกฺโกติ พฺรหฺมจริยํ สนฺธาเรตุํ, สิกฺขาทุพฺพลฺยํ อาวิกตฺวา สิกฺขํ ปจฺจกฺขาย หีนายาวตฺตติ, อิทมสฺส ธชคฺคสฺมิํ.}

\prose{เสยฺยถาปิ โส ภิกฺขเว โยธาชีโว สหติ รชคฺคํ, อปิ จ โข ธชคฺคญฺเญว ทิสฺวา สํสีทติ วิสีทติ น สนฺถมฺภติ น สกฺโกติ สงฺคามํ โอตริตุํ.\csromanspacer{} ตถูปมาหํ ภิกฺขเว อิมํ ปุคฺคลํ วทามิ, เอวรูโปปิ ภิกฺขเว อิเธกจฺโจ ปุคฺคโล โหติ.\csromanspacer{} อยํ ภิกฺขเว ทุติโย โยธาชีวูปโม ปุคฺคโล สนฺโต สํวิชฺชมาโน ภิกฺขูสุ.}

\prose{ปุน จปรํ ภิกฺขเว ภิกฺขุ สหติ รชคฺคํ, สหติ ธชคฺคํ, อปิ จ โข อุสฺสารณญฺเญว สุตฺวา สํสีทติ วิสีทติ น สนฺถมฺภติ น สกฺโกติ พฺรหฺมจริยํ สนฺธาเรตุํ, สิกฺขาทุพฺพลฺยํ อาวิกตฺวา สิกฺขํ ปจฺจกฺขาย หีนายาวตฺตติ.\csromanspacer{} กิมสฺส อุสฺสารณาย? อิธ ภิกฺขเว ภิกฺขุํ อรญฺญคตํ วา รุกฺขมูลคตํ วา สุญฺญาคารคตํ วา มาตุคาโม อุปสงฺกมิตฺวา อูหสติ\footnote{อุหสติ (ก), โอหสติ (สฺยา, กํ) อภิ 3. 177 ปิฏฺเฐปิ.} อุลฺลปติ อุชฺชคฺฆติ อุปฺปณฺเฑติ, โส มาตุคาเมน อูหสิยมาโน อุลฺลปิยมาโน อุชฺชคฺฆิยมาโน อุปฺปณฺฑิยมาโน สํสีทติ วิสีทติ น สนฺถมฺภติ น สกฺโกติ พฺรหฺมจริยํ สนฺธาเรตุํ, สิกฺขาทุพฺพลฺยํ อาวิกตฺวา สิกฺขํ ปจฺจกฺขาย หีนายาวตฺตติ, อิทมสฺส อุสฺสารณาย.}

\prose{เสยฺยถาปิ โส ภิกฺขเว โยธาชีโว สหติ รชคฺคํ, สหติ ธชคฺคํ, อปิ จ โข อุสฺสารณญฺเญว สุตฺวา สํสีทติ วิสีทติ น สนฺถมฺภติ น สกฺโกติ สงฺคามํ โอตริตุํ.\csromanspacer{} ตถูปมาหํ ภิกฺขเว อิมํ ปุคฺคลํ วทามิ, เอวรูโปปิ ภิกฺขเว อิเธกจฺโจ ปุคฺคโล โหติ.\csromanspacer{} อยํ ภิกฺขเว ตติโย โยธาชีวูปโม ปุคฺคโล สนฺโต สํวิชฺชมาโน ภิกฺขูสุ.}

\prose{ปุน จปรํ ภิกฺขเว ภิกฺขุ สหติ รชคฺคํ, สหติ ธชคฺคํ, สหติ อุสฺสารณํ, อปิ จ โข สมฺปหาเร หญฺญติ พฺยาปชฺชติ.\csromanspacer{} กิมสฺส}

\vspace{\tipitakaparskip}
\csromancenter{(8) 3.\csromanspacer{} โยธาชีววคฺค สมฺปหารสฺมิํ? อิธ ภิกฺขเว ภิกฺขุํ อรญฺญคตํ วา รุกฺขมูลคตํ วา สุญฺญาคารคตํ วา มาตุคาโม อุปสงฺกมิตฺวา อภินิสีทติ อภินิปชฺชติ อชฺโฌตฺถรติ.\csromanspacer{} โส มาตุคาเมน อภินิสีทิยมาโน อภินิปชฺชิยมาโน อชฺโฌตฺถริยมาโน สิกฺขํ อปจฺจกฺขาย ทุพฺพลฺยํ อนาวิกตฺวา เมถุนํ ธมฺมํ ปฏิเสวติ.\csromanspacer{} อิทมสฺส สมฺปหารสฺมิํ.}
\prose{\csromanfolio{81}เสยฺยถาปิ โส ภิกฺขเว โยธาชีโว สหติ รชคฺคํ, สหติ ธชคฺคํ, สหติ อุสฺสารณํ, อปิ จ โข สมฺปหาเร หญฺญติ พฺยาปชฺชติ, ตถูปมาหํ ภิกฺขเว อิมํ ปุคฺคลํ วทามิ, เอวรูโปปิ ภิกฺขเว อิเธกจฺโจ ปุคฺคโล โหติ.\csromanspacer{} อยํ ภิกฺขเว จตุตฺโถ โยธาชีวูปโม ปุคฺคโล สนฺโต สํวิชฺชมาโน ภิกฺขูสุ.}

\prose{ปุน จปรํ ภิกฺขเว ภิกฺขุ สหติ รชคฺคํ, สหติ ธชคฺคํ, สหติ อุสฺสารณํ, สหติ สมฺปหารํ.\csromanspacer{} โส ตํ สงฺคามํ อภิวิชินิตฺวา วิชิตสงฺคาโม ตเมว สงฺคามสีสํ อชฺฌาวสติ.\csromanspacer{} กิมสฺส สงฺคามวิชยสฺมิํ? อิธ ภิกฺขเว ภิกฺขุ อรญฺญคตํ วา รุกฺขมูลคตํ วา สุญฺญาคารคตํ วา มาตุคาโม อุปสงฺกมิตฺวา อภินิสีทติ อภินิปชฺชติ อชฺโฌตฺถรติ, โส มาตุคาเมน อภินิสีทิยมาโน อภินิปชฺชิยมาโน อชฺโฌตฺถริยมาโน วินิเวเฐตฺวา วินิโมเจตฺวา เยน กามํ ปกฺกมติ.\csromanspacer{} โส วิวิตฺตํ เสนาสนํ ภชติ อรญฺญํ รุกฺขมูลํ ปพฺพตํ กนฺทรํ คิริคุหํ สุสานํ วนปตฺถํ อพฺโภกาสํ ปลาลปุญฺชํ.}

\prose{โส อรญฺญคโต วา รุกฺขมูลคโต วา สุญฺญาคารคโต วา นิสีทติ ปลฺลงฺกํ อาภุชิตฺวา อุชุํ กายํ ปณิธาย ปริมุขํ สติํ อุปฏฺฐเปตฺวา.\csromanspacer{} โส อภิชฺฌํ โลเก ปหาย วิคตาภิชฺเฌน เจตสา วิหรติ, อภิชฺฌาย จิตฺตํ ปริโสเธติ.\csromanspacer{} พฺยาปาท ปโทสํ ปหาย อพฺยาปนฺนจิตฺโต วิหรติ, สพฺพปาณภูตหิตานุกมฺปี พฺยาปาทปโทสา จิตฺตํ ปริโสเธติ.\csromanspacer{} ถินมิทฺธํ ปหาย วิคตถินมิทฺโธ วิหรติ, อาโลกสญฺญี สโต สมฺปชาโน ถินมิทฺธา จิตฺตํ ปริโสเธติ.\csromanspacer{} อุทฺธจฺจกุกฺกุจฺจํ ปหาย อนุทฺธโต วิหรติ, อชฺฌตฺตํ วูปสนฺตจิตฺโต อุทฺธจฺจกุกฺกุจฺจา จิตฺตํ ปริโสเธติ.\csromanspacer{} วิจิกิจฺฉํ ปหาย ติณฺณวิจิกิจฺโฉ วิหรติ, อกถํกถี กุสเลสุ ธมฺเมสุ วิจิกิจฺฉาย จิตฺตํ ปริโสเธติ.\csromanspacer{} โส อิเม ปญฺจ นีวรเณ ปหาย เจตโส อุปกฺกิเลเส ปญฺญาย ทุพฺพลีกรเณ วิวิจฺเจว กาเมหิ ฯเปฯ จตุตฺถํ ฌานํ อุปสมฺปชฺช วิหรติ.}

\prose{\csromanfolio{82}โส เอวํ สมาหิเต จิตฺเต ปริสุทฺเธ ปริโยทาเต อนงฺคเณ วิคตูปกฺกิเลเส มุทุภูเต กมฺมนิเย ฐิเต อาเนญฺชปฺปตฺเต อาสวานํ ขยญาณาย จิตฺตํ อภินินฺนาเมติ.\csromanspacer{} โส “อิทํ ทุกฺขนฺ”ติ ยถาภูตํ ปชานาติ, “อยํ ทุกฺขสมุทโย”ติ ยถาภูตํ ปชานาติ, “อยํ ทุกฺขนิโรโธ”ติ ยถาภูตํ ปชานาติ, “อยํ ทุกฺขนิโรธคามินี ปฏิปทา”ติ ยถาภูตํ ปชานาติ, “อิเม อาสวา”ติ ยถาภูตํ ปชานาติ, “อยํ อาสวสมุทโย”ติ ยถาภูตํ ปชานาติ, “อยํ อาสวนิโรโธ”ติ ยถาภูตํ ปชานาติ, “อยํ อาสวนิโรธคามินี ปฏิปทา”ติ ยถาภูตํ ปชานาติ.\csromanspacer{} ตสฺส เอวํ ชานโต เอวํ ปสฺสโต กามาสวาปิ จิตฺตํ วิมุจฺจติ ภวาสวาปิ จิตฺตํ วิมุจฺจติ, อวิชฺชาสวาปิ จิตฺตํ วิมุจฺจติ, วิมุตฺตสฺมิํ “วิมุตฺตมฺ”อิติ ญาณํ โหติ, “ขีณา ชาติ วุสิตํ พฺรหฺมจริยํ, กตํ กรณียํ, นาปรํ อิตฺถตฺตายา”ติ ปชานาติ.\csromanspacer{} อิทมสฺส สงฺคามวิชยสฺมิํ.}

\prose{เสยฺยาถาปิ โส ภิกฺขเว โยธาชีโว สหติ รชคฺคํ, สหติ ธชคฺคํ, สหติ อุสฺสารณํ, สหติ สมฺปหารํ.\csromanspacer{} โส ตํ สงฺคามํ อภิวิชินิตฺวา วิชิตสงฺคาโม ตเมว สงฺคามสีสํ อชฺฌาวสติ, ตถูปมาหํ ภิกฺขเว อิมํ ปุคฺคลํ วทามิ, เอวรูโปปิ ภิกฺขเว อิเธกจฺโจ ปุคฺคโล โหติ.\csromanspacer{} อยํ ภิกฺขเว ปญฺจโม โยธาชีวูปโม ปุคฺคโล สนฺโต สํวิชฺชมาโน ภิกฺขูสุ.\csromanspacer{} อิเม โข ภิกฺขเว ปญฺจ โยธาชีวูปมา ปุคฺคลา สนฺโต สํวิชฺชมานา ภิกฺขูสูติ.\csromanspacer{}\csromanspacer{}ปญฺจมํ.}

\csromanheader{2}{6. ทุติยโยธาชีวสุตฺต}
\proseitem{76}{ปญฺจิเม ภิกฺขเว โยธาชีวา สนฺโต สํวิชฺชมานา โลกสฺมิํ.\csromanspacer{} กตเม ปญฺจ? อิธ ภิกฺขเว เอกจฺโจ โยธาชีโว อสิจมฺมํ คเหตฺวา ธนุกลาปํ สนฺนยฺหิตฺวา วิยูฬฺหํ สงฺคามํ โอตรติ, โส ตสฺมิํ สงฺคาเม อุสฺสหติ วายมติ, ตเมนํ อุสฺสหนฺตํ วายมนฺตํ ปเร หนนฺติ ปริยาปาเทนฺติ, เอวรูโปปิ ภิกฺขเว อิเธกจฺโจ โยธาชีโว โหติ.\csromanspacer{} อยํ ภิกฺขเว ปฐโม โยธาชีโว สนฺโต สํวิชฺชมาโน โลกสฺมิํ.}

\prose{ปุน จปรํ ภิกฺขเว อิเธกจฺโจ โยธาชีโว อสิจมฺมํ คเหตฺวา ธนุกลาปํ สนฺนยฺหิตฺวา วิยูฬฺหํ สงฺคามํ โอตรติ, โส ตสฺมิํ สงฺคาเม}

\vspace{\tipitakaparskip}
\csromancenter{(8) 3.\csromanspacer{} โยธาชีววคฺค อุสฺสหติ วายมติ, ตเมนํ อุสฺสหนฺตํ วายมนฺตํ ปเร อุปลิกฺขนฺติ\footnote{อุปลิขนฺติ (ก)}, ตเมนํ อปเนนฺติ, อปเนตฺวา ญาตกานํ เนนฺติ.\csromanspacer{} โส ญาตเกหิ นียมาโน อปฺปตฺวาว ญาตเก อนฺตรามคฺเค กาลํ กโรติ, เอวรูโปปิ ภิกฺขเว อิเธกจฺโจ โยธาชีโว โหติ.\csromanspacer{} อยํ ภิกฺขเว ทุติโย โยธาชีโว สนฺโต สํวิชฺชมาโน โลกสฺมิํ.}
\prose{\csromanfolio{83}ปุน จปรํ ภิกฺขเว อิเธกจฺโจ โยธาชีโว อสิจมฺมํ คเหตฺวา ธนุกลาปํ สนฺนยฺหิตฺวา วิยูฬฺหํ สงฺคามํ โอตรติ, โส ตสฺมิํ สงฺคาเม อุสฺสหติ วายมติ, ตเมนํ อุสฺสหนฺตํ วายมนฺตํ ปเร อุปลิกฺขนฺติ, ตเมนํ อปเนนฺติ, อปเนตฺวา ญาตกานํ เนนฺติ, ตเมนํ ญาตกา อุปฏฺฐหนฺติ ปริจรนฺติ.\csromanspacer{} โส ญาตเกหิ อุปฏฺฐหิยมาโน ปริจริยมาโน เตเนว อาพาเธน กาลํ กโรติ, เอวรูโปปิ ภิกฺขเว อิเธกจฺโจ โยธาชีโว โหติ.\csromanspacer{} อยํ ภิกฺขเว ตติโย โยธาชีโว สนฺโต สํวิชฺชมาโน โลกสฺมิํ.}

\prose{ปุน จปรํ ภิกฺขเว อิเธกจฺโจ โยธาชีโว อสิจมฺมํ คเหตฺวา ธนุกลาปํ สนฺนยฺหิตฺวา วิยูฬฺหํ สงฺคามํ โอตรติ, โส ตสฺมิํ สงฺคาเม อุสฺสหติ วายมติ, ตเมนํ อุสฺสหนฺตํ วายมนฺตํ ปเร อุปลิกฺขนฺติ, ตเมนํ อปเนนฺติ, อปเนตฺวา ญาตกานํ เนนฺติ, ตเมนํ ญาตกา อุปฏฺฐหนฺติ ปริจรนฺติ.\csromanspacer{} โส ญาตเกหิ อุปฏฺฐหิยมาโน ปริจริยมาโน วุฏฺฐาติ ตมฺหา อาพาธา, เอวรูโปปิ ภิกฺขเว อิเธกจฺโจ โยธาชีโว โหติ.\csromanspacer{} อยํ ภิกฺขเว จตุตฺโถ โยธาชีโว สนฺโต สํวิชฺชมาโน โลกสฺมิํ.}

\prose{ปุน จปรํ ภิกฺขเว อิเธกจฺโจ โยธาชีโว อสิจมฺมํ คเหตฺวา ธนุกลาปํ สนฺนยฺหิตฺวา วิยูฬฺหํ สงฺคามํ โอตรติ, โส ตํ สงฺคามํ อภิวิชินิตฺวา วิชิตสงฺคาโม ตเมว สงฺคามสีสํ อชฺฌาวสติ, เอวรูโปปิ ภิกฺขเว อิเธกจฺโจ โยธาชีโว โหติ.\csromanspacer{} อยํ ภิกฺขเว ปญฺจโม โยธาชีโว สนฺโต สํวิชฺชมาโน โลกสฺมิํ.\csromanspacer{} อิเม โข ภิกฺขเว ปญฺจ โยธาชีวา สนฺโต สํวิชฺชมานา โลกสฺมิํ.}

\prose{เอวเมวํ โข ภิกฺขเว ปญฺจิเม โยธาชีวูปมา ปุคฺคลา สนฺโต สํวิชฺชมานา ภิกฺขูสุ.\csromanspacer{} กตเม ปญฺจ? อิธ ภิกฺขเว ภิกฺขุ อญฺญตรํ คามํ วา \csromanfolio{84}นิคมํ วา อุปนิสฺสาย วิหรติ, โส ปุพฺพณฺหสมยํ นิวาเสตฺวา ปตฺตจีวรมาทาย ตเมว คามํ วา นิคมํ วา ปิณฺฑาย ปวิสติ อรกฺขิเตเนว กาเยน อรกฺขิตาย วาจาย อรกฺขิเตน จิตฺเตน อนุปฏฺฐิตาย สติยา อสํวุเตหิ อินฺทฺริเยหิ.\csromanspacer{} โส ตตฺถ ปสฺสติ มาตุคามํ ทุนฺนิวตฺถํ วา ทุปฺปารุตํ วา, ตสฺส ตํ มาตุคามํ ทิสฺวา ทุนฺนิวตฺถํ วา ทุปฺปารุตํ วา ราโค จิตฺตํ อนุทฺธํเสติ.\csromanspacer{} โส ราคานุทฺธํสิเตน จิตฺเตน สิกฺขํ อปจฺจกฺขาย ทุพฺพลฺยํ อนาวิกตฺวา เมถุนํ ธมฺมํ ปฏิเสวติ.}

\prose{เสยฺยถาปิ โส ภิกฺขเว โยธาชีโว อสิจมฺมํ คเหตฺวา ธนุกลาปํ สนฺนยฺหิตฺวา วิยูฬฺหํ สงฺคามํ โอตรติ, โส ตสฺมิํ สงฺคาเม อุสฺสหติ วายมติ, ตเมนํ อุสฺสหนฺตํ วายมนฺตํ ปเร หนนฺติ ปริยาปาเทนฺติ, ตถูปมาหํ ภิกฺขเว อิมํ ปุคฺคลํ วทามิ, เอวรูโปปิ ภิกฺขเว อิเธกจฺโจ ปุคฺคโล โหติ.\csromanspacer{} อยํ ภิกฺขเว ปฐโม โยธาชีวูปโม ปุคฺคโล สนฺโต สํวิชฺชมาโน ภิกฺขูสุ.}

\prose{ปุน จปรํ ภิกฺขเว ภิกฺขุ อญฺญตรํ คามํ วา นิคมํ วา อุปนิสฺสาย วิหรติ, โส ปุพฺพณฺหสมยํ นิวาเสตฺวา ปตฺตจีวรมาทาย ตเมว คามํ วา นิคมํ วา ปิณฺฑาย ปวิสติ อรกฺขิเตเนว กาเยน อรกฺขิตาย วาจาย อรกฺขิเตน จิตฺเตน อนุปฏฺฐิตาย สติยา อสํวุเตหิ อินฺทฺริเยหิ.\csromanspacer{} โส ตตฺถ ปสฺสติ มาตุคามํ ทุนฺนิวตฺถํ วา ทุปฺปารุตํ วา, ตสฺส ตํ มาตุคามํ ทิสฺวา ทุนฺนิวตฺถํ วา ทุปฺปารุตํ วา ราโค จิตฺตํ อนุทฺธํเสติ.\csromanspacer{} โส ราคานุทฺธํสิเตน จิตฺเตน ปริฑยฺหเตว กาเยน, ปริฑยฺหติ เจตสา.\csromanspacer{} ตสฺส เอวํ โหติ “ยํนูนาหํ อารามํ คนฺตฺวา ภิกฺขูนํ อาโรเจยฺยํ ‘ราคปริยุฏฺฐิโตมฺหิ\footnote{ราคายิโตมฺหิ (สี, สฺยา, กํ)} อาวุโส ราคปเรโต, น สกฺโกมิ พฺรหฺมจริยํ สนฺธาเรตุํ, สิกฺขาทุพฺพลฺยํ อาวิกตฺวา สิกฺขํ ปจฺจกฺขาย หีนายาวตฺติสฺสามี’ติ”.\csromanspacer{} โส อารามํ คจฺฉนฺโต อปฺปตฺวาว อารามํ อนฺตรามคฺเค สิกฺขาทุพฺพลฺยํ อาวิกตฺวา สิกฺขํ ปจฺจกฺขาย หีนายาวตฺตติ.}

\prose{เสยฺยถาปิ โส ภิกฺขเว โยธาชีโว อสิจมฺมํ คเหตฺวา ธนุกลาปํ สนฺนยฺหิตฺวา วิยูฬฺหํ สงฺคามํ โอตรติ, โส ตสฺมิํ สงฺคาเม อุสฺสหติ วายมติ, ตเมนํ อุสฺสหนฺตํ วายมนฺตํ ปเร อุปลิกฺขนฺติ, ตเมนํ อปเนนฺติ, อปเนตฺวา ญาตกานํ เนนฺติ.\csromanspacer{} โส ญาตเกหิ นียมาโน}

\vspace{\tipitakaparskip}
\csromancenter{(8) 3.\csromanspacer{} โยธาชีววคฺค อปฺปตฺวาว ญาตเก อนฺตรามคฺเค กาลํ กโรติ, ตถูปมาหํ ภิกฺขเว อิมํ ปุคฺคลํ วทามิ, เอวรูโปปิ ภิกฺขเว อิเธกจฺโจ ปุคฺคโล โหติ.\csromanspacer{} อยํ ภิกฺขเว ทุติโย โยธาชีวูปโม ปุคฺคโล สนฺโต สํวิชฺชมาโน ภิกฺขูสุ.}
\prose{\csromanfolio{85}ปุน จปรํ ภิกฺขเว ภิกฺขุ อญฺญตรํ คามํ วา นิคมํ วา อุปนิสฺสาย วิหรติ, โส ปุพฺพณฺหสมยํ นิวาเสตฺวา ปตฺตจีวรมาทาย ตเมว คามํ วา นิคมํ วา ปิณฺฑาย ปวิสติ อรกฺขิเตเนว กาเยน อรกฺขิตาย วาจาย อรกฺขิเตน จิตฺเตน อนุปฏฺฐิตาย สติยา อสํวุเตหิ อินฺทฺริเยหิ.\csromanspacer{} โส ตตฺถ ปสฺสติ มาตุคามํ ทุนฺนิวตฺถํ วา ทุปฺปารุตํ วา, ตสฺส ตํ มาตุคามํ ทิสฺวา ทุนฺนิวตฺถํ วา ทุปฺปารุตํ วา ราโค จิตฺตํ อนุทฺธํเสติ.\csromanspacer{} โส ราคานุทฺธํสิเตน จิตฺเตน ปริฑยฺหเตว กาเยน, ปริฑยฺหติ เจตสฺสา.\csromanspacer{} ตสฺส เอวํ โหติ “ยํนูนาหํ อารามํ คนฺตฺวา ภิกฺขูนํ อาโรเจยฺยํ ‘ราคปริยุฏฺฐิโตมฺหิ อาวุโส ราคปเรโต, น สกฺโกมิ พฺรหฺมจริยํ สนฺธาเรตุํ, สิกฺขาทุพฺพลฺยํ อาวิกตฺวา สิกฺขํ ปจฺจกฺขาย หีนายาวตฺติสฺสามี’ติ”.\csromanspacer{} โส อารามํ คนฺตฺวา ภิกฺขูนํ อาโรเจติ “ราคปริยุฏฺฐิโตมฺหิ อาวุโส ราคปเรโต น สกฺโกมิ พฺรหฺมจริยํ สนฺธาเรตุํ, สิกฺขาทุพฺพลฺยํ อาวิกตฺวา สิกฺขํ ปจฺจกฺขาย หีนายาวตฺติสฺสามี”ติ.}

\prose{ตเมนํ สพฺรหฺมจารี โอวทนฺติ อนุสาสนฺติ “อปฺปสฺสาทา อาวุโส กามา วุตฺตา ภควตา พหุทุกฺขา พหุปายาสา\footnote{พหูปายาสา (สี, สฺยา, กํ, อิ) วิ 2. 175; วิ 4. 68; ม 1. 182 ปิฏฺเฐสุปิ.}, อาทีนโว เอตฺถ ภิยฺโย.\csromanspacer{} อฏฺฐิกงฺกลูปมา กามา วุตฺตา ภควตา พหุทุกฺขา พหุปายาสา, อาทีนโว เอตฺถ ภิยฺโย.\csromanspacer{} มํสเปสูปมา กามา วุตฺตา ภควตา พหุทุกฺขา พหุปายาสา, อาทีนโว เอตฺถ ภิยฺโย.\csromanspacer{} ติณุกฺกูปมา กามา วุตฺตา ภควตา พหุทุกฺขา พหุปายาสา, อาทีนโว เอตฺถ ภิยฺโย.\csromanspacer{} องฺคารกาสูปมา กามา วุตฺตา ภควตา พหุทุกฺขา พหุปายาสา, อาทีนโว เอตฺถ ภิยฺโย.\csromanspacer{} สุปินกูปมา กามา วุตฺตา ภควตา พหุทุกฺขา พหุปายาสา, อาทีนโว เอตฺถ ภิยฺโย.\csromanspacer{} ยาจิตกูปมา กามา วุตฺตา ภควตา พหุทุกฺขา พหุปายาสา, อาทีนโว เอตฺถ ภิยฺโย.\csromanspacer{} รุกฺขผลูปมา กามา วุตฺตา ภควตา พหุทุกฺขา พหุปายาสา, อาทีนโว เอตฺถ ภิยฺโย.\csromanspacer{} อสิสูนูปมา กามา วุตฺตา ภควตา พหุทุกฺขา พหุปายาสา, อาทีนโว เอตฺถ ภิยฺโย.\csromanspacer{} สตฺติสูลูปมา กามา วุตฺตา \csromanfolio{86}ภควตา พหุทุกฺขา พหุปายาสา, อาทีนโว เอตฺถ ภิยฺโย.\csromanspacer{} สปฺปสิรูปมา กามา วุตฺตา ภควตา พหุทุกฺขา พหุปายาสา, อาทีนโว เอตฺถ ภิยฺโย.\csromanspacer{} อภิรมตายสฺมา พฺรหฺมจริเย, มายสฺมา สิกฺขาทุพฺพลฺยํ อาวิกตฺวา สิกฺขํ ปจฺจกฺขาย หีนายาวตฺตี”ติ.}

\prose{โส สพฺรหฺมจารีหิ เอวํ โอวทิยมาโน เอวํ อนุสาสิยมาโน เอวมาห “กิญฺจาปิ อาวุโส อปฺปสฺสาทา กามา วุตฺตา ภควตา พหุทุกฺขา พหุปายาสา, อาทีนโว เอตฺถ ภิยฺโย.\csromanspacer{} อถ โข เนวาหํ สกฺโกมิ พฺรหฺมจริยํ สนฺธาเรตุํ, สิกฺขาทุพฺพลฺยํ อาวิกตฺวา สิกฺขํ ปจฺจกฺขาย หีนายาวตฺติสฺสามี”ติ.\csromanspacer{} โส สิกฺขาทุพฺพลฺยํ อาวิกตฺวา สิกฺขํ ปจฺจกฺขาย หีนายาวตฺตติ.}

\prose{เสยฺยถาปิ โส ภิกฺขเว โยธาชีโว อสิจมฺมํ คเหตฺวา ธนุกลาปํ สนฺนยฺหิตฺวา วิยูฬฺหํ สงฺคามํ โอตรติ, โส ตสฺมิํ สงฺคาเม อุสฺสหติ วายมติ, ตเมนํ อุสฺสหนฺตํ วายมนฺตํ ปเร อุปลิกฺขนฺติ, ตเมนํ อปเนนฺติ, อปเนตฺวา ญาตกานํ เนนฺติ, ตเมนํ ญาตกา อุปฏฺฐหนฺติ ปริจรนฺติ.\csromanspacer{} โส ญาตเกหิ อุปฏฺฐหิยมาโน ปริจริยมาโน เตเนว อาพาเธน กาลํ กโรติ, ตถูปมาหํ ภิกฺขเว อิมํ ปุคฺคลํ วทามิ, เอวรูโปปิ ภิกฺขเว อิเธกจฺโจ ปุคฺคโล โหติ.\csromanspacer{} อยํ ภิกฺขเว ตติโย โยธาชีวูปโม ปุคฺคโล สนฺโต สํวิชฺชมาโน ภิกฺขูสุ.}

\prose{ปุน จปรํ ภิกฺขเว ภิกฺขุ อญฺญตรํ คามํ วา นิคมํ วา อุปนิสฺสาย วิหรติ, โส ปุพฺพณฺหสมยํ นิวาเสตฺวา ปตฺตจีวรมาทาย ตเมว คามํ วา นิคมํ วา ปิณฺฑาย ปวิสติ อรกฺขิเตเนว กาเยน อรกฺขิตาย วาจาย อรกฺขิเตน จิตฺเตน อนุปฏฺฐิตาย สติยา อสํวุเตหิ อินฺทฺริเยหิ.\csromanspacer{} โส ตตฺถ ปสฺสติ มาตุคามํ ทุนฺนิวตฺถํ วา ทุปฺปารุตํ วา, ตสฺส ตํ มาตุคามํ ทิสฺวา ทุนฺนิวตฺถํ วา ทุปฺปารุตํ วา ราโค จิตฺตํ อนุทฺธํเสติ.\csromanspacer{} โส ราคานุทฺธํสิเตน จิตฺเตน ปริฑยฺหเตว กาเยน, ปริฑยฺหติ เจตสา.\csromanspacer{} ตสฺส เอวํ โหติ “ยํนูนาหํ อารามํ คนฺตฺวา ภิกฺขูนํ อาโรเจยฺยํ ‘ราคปริยุฏฺฐิโตมฺหิ อาวุโส ราคปเรโต, น สกฺโกมิ พฺรหฺมจริยํ สนฺธาเรตุํ, สิกฺขาทุพฺพลฺยํ อาวิกตฺวา สิกฺขํ ปจฺจกฺขาย หีนายาวตฺติสฺสามี’ติ”.\csromanspacer{} โส อารามํ คนฺตฺวา ภิกฺขูนํ อาโรเจติ “ราคปริยุฏฺฐิโตมฺหิ อาวุโส ราคปเรโต, น สกฺโกมิ}

\vspace{\tipitakaparskip}
\csromancenter{(8) 3.\csromanspacer{} โยธาชีววคฺค พฺรหฺมจริยํ สนฺธาเรตุํ, สิกฺขาทุพฺพลฺยํ อาวิกตฺวา สิกฺขํ ปจฺจกฺขาย หีนายาวตฺติสฺสามี”ติ.}
\prose{\csromanfolio{87}ตเมนํ สพฺรหฺมจารี โอวทนฺติ อนุสาสนฺติ “อปฺปสฺสาทา อาวุโส กามา วุตฺตา ภควตา พหุทุกฺขา พหุปายาสา, อาทีนโว เอตฺถ ภิยฺโย.\csromanspacer{} อฏฺฐิกงฺกลูปมา กามา วุตฺตา ภควตา พหุทุกฺขา พหุปายาสา, อาทีนโว เอตฺถ ภิยฺโย.\csromanspacer{} มํสเปสูปมา กามา วุตฺตา ภควตา ฯเปฯ.\csromanspacer{} ติณุกฺกูปมา กามา วุตฺตา ภควตา.\csromanspacer{} องฺคารกาสูปมา กามา วุตฺตา ภควตา.\csromanspacer{} สุปินกูปมา กามา วุตฺตา ภควตา.\csromanspacer{} ยาจิตกูปมา กามา วุตฺตา ภควตา.\csromanspacer{} รุกฺขผลูปมา กามา วุตฺตา ภควตา.\csromanspacer{} อสิสูนูปมา กามา วุตฺตา ภควตา.\csromanspacer{} สตฺติสูลูปมา กามา วุตฺตา ภควตา.\csromanspacer{} สปฺปสิรูปมา กามา วุตฺตา ภควตา พหุทุกฺขา พหุปายาสา, อาทีนโว เอตฺถ ภิยฺโย.\csromanspacer{} อภิรมตายสฺมา พฺรหฺมจริเย, มายสฺมา สิกฺขาทุพฺพลฺยํ อาวิกตฺวา สิกฺขํ ปจฺจกฺขาย หีนายาวตฺตี”ติ.}

\prose{โส สพฺรหฺมจารีหิ เอวํ โอวทิยมาโน เอวํ อนุสาสิยมาโน เอวมาห “อุสฺสหิสฺสามิ อาวุโส วายมิสฺสามิ อาวุโส อภิรมิสฺสามิ อาวุโส, น ทานาหํ อาวุโส สิกฺขาทุพฺพลฺยํ อาวิกตฺวา สิกฺขํ ปจฺจกฺขาย หีนายาวตฺติสฺสามี”ติ.}

\prose{เสยฺยถาปิ โส ภิกฺขเว โยธาชีโว อสิจมฺมํ คเหตฺวา ธนุกลาปํ สนฺนยฺหิตฺวา วิยูฬฺหํ สงฺคามํ โอตรติ, โส ตสฺมิํ สงฺคาเม อุสฺสหติ วายมติ, ตเมนํ อุสฺสหนฺตํ วายมนฺตํ ปเร อุปลิกฺขนฺติ, ตเมนํ อปเนนฺติ, อปเนตฺวา ญาตกานํ เนนฺติ, ตเมนํ ญาตกา อุปฏฺฐหนฺติ ปริจรนฺติ.\csromanspacer{} โส ญาตเกหิ อุปฏฺฐหิยมาโน ปริจริยมาโน วุฏฺฐาติ ตมฺหา อาพาธา, ตถูปมาหํ ภิกฺขเว อิมํ ปุคฺคลํ วทามิ, เอวรูโปปิ ภิกฺขเว อิเธกจฺโจ ปุคฺคโล โหติ.\csromanspacer{} อยํ ภิกฺขเว จตุตฺโถ โยธาชีวูปโม ปุคฺคโล สนฺโต สํวิชฺชมาโน ภิกฺขูสุ.}

\prose{ปุน จปรํ ภิกฺขเว ภิกฺขุ อญฺญตรํ คามํ วา นิคมํ วา อุปนิสฺสาย วิหรติ, โส ปุพฺพณฺหสมยํ นิวาเสตฺวา ปตฺตจีวรมาทาย ตเมว คามํ วา นิคมํ วา ปิณฺฑาย ปวิสติ รกฺขิเตเนว กาเยน รกฺขิตาย วาจาย รกฺขิเตน จิตฺเตน อุปฏฺฐิตาย สติยา สํวุเตหิ อินฺทฺริเยหิ. \csromanfolio{88}โส จกฺขุนา รูปํ ทิสฺวา น นิมิตฺตคฺคาหี โหติ นานุพฺยญฺชนคฺคาหี, ยตฺวาธิกรณเมนํ จกฺขุนฺทฺริยํ อสํวุตํ วิหรนฺตํ อภิชฺฌาโทมนสฺสา ปาปกา อกุสลา ธมฺมา อนฺวาสฺสเวยฺยุํ.\csromanspacer{} ตสฺส สํวราย ปฏิปชฺชติ, รกฺขติ จกฺขุนฺทฺริยํ, จกฺขุนฺทฺริเย สํวรํ อาปชฺชติ.\csromanspacer{} โสเตน สทฺทํ สุตฺวา.\csromanspacer{} ฆาเนน คนฺธํ ฆายิตฺวา.\csromanspacer{} ชิวฺหาย รสํ สายิตฺวา.\csromanspacer{} กาเยน โผฏฺฐพฺพํ ผุสิตฺวา.\csromanspacer{} มนสา ธมฺมํ วิญฺญาย น นิมิตฺตคฺคาหี โหติ นานุพฺยญฺชนคฺคาหี, ยตฺวาธิกรณเมนํ มนินฺทฺริยํ อสํวุตํ วิหรนฺตํ อภิชฺฌาโทมนสฺสา ปาปกา อกุสลา ธมฺมา อนฺวาสฺสเวยฺยุํ.\csromanspacer{} ตสฺส สํวราย ปฏิปชฺชติ, รกฺขติ มนินฺทฺริยํ, มนินฺทฺริเย สํวรํ อาปชฺชติ.\csromanspacer{} โส ปจฺฉาภตฺตํ ปิณฺฑปาตปฏิกฺกนฺโต วิวิตฺตํ เสนาสนํ ภชติ อรญฺญํ รุกฺขมูลํ ปพฺพตํ กนฺทรํ คิริคุหํ สุสานํ วนปตฺถํ อพฺโภกาสํ ปลาลปุญฺชํ, โส อรญฺญคโต วา รุกฺขมูลคโต วา สุญฺญาคารคโต วา นิสีทติ ปลฺลงฺกํ อาภุชิตฺวา อุชุํ กายํ ปณิธาย ปริมุขํ สติํ อุปฏฺฐเปตฺวา.\csromanspacer{} โส อภิชฺฌํ โลเก ปหาย ฯเปฯ.\csromanspacer{} โส อิเม ปญฺจ นีวรเณ ปหาย เจตโส อุปกฺกิเลเส ปญฺญาย ทุพฺพลีกรเณ วิวิจฺเจว กาเมหิ ฯเปฯ จตุตฺถํ ฌานํ อุปสมฺปชฺช วิหรติ.}

\prose{โส เอวํ สมาหิเต จิตฺเต ปริสุทฺเธ ปริโยทาเต อนงฺคเณ วิคตูปกฺกิเลเส มุทุภูเต กมฺมนิเย ฐิเต อาเนญฺชปฺปตฺเต อาสวานํ ขยญาณาย จิตฺตํ อภินินฺนาเมติ, โส “อิทํ ทุกฺขนฺ”ติ ยถาภูตํ ปชานาติ ฯเปฯ นาปรํ อิตฺถตฺตายาติ ปชานาติ.}

\prose{เสยฺยถาปิ โส ภิกฺขเว โยธาชีโว อสิจมฺมํ คเหตฺวา ธนุกลาปํ สนฺนยฺหิตฺวา วิยูฬฺหํ สงฺคามํ โอตรติ, โส ตํ สงฺคามํ อภิวิชินิตฺวา วิชิตสงฺคาโม ตเมว สงฺคามสีสํ อชฺฌาวสติ, ตถูปมาหํ ภิกฺขเว อิมํ ปุคฺคลํ วทามิ, เอวรูโปปิ ภิกฺขเว อิเธกจฺโจ ปุคฺคโล โหติ.\csromanspacer{} อยํ ภิกฺขเว ปญฺจโม โยธาชีวูปโม ปุคฺคโล สนฺโต สํวิชฺชมาโน ภิกฺขูสุ.\csromanspacer{} อิเม โข ภิกฺขเว ปญฺจ โยธาชีวูปมา ปุคฺคลา สนฺโต สํวิชฺชมานา ภิกฺขูสูติ.\csromanspacer{}\csromanspacer{}ฉฏฺฐํ.}

\csromanheader{2}{7. ปฐม อนาคตภยสุตฺต}
\proseitem{77}{ปญฺจิมานิ ภิกฺขเว อนาคตภยานิ สมฺปสฺสมาเนน อลเมว อารญฺญิเกน\footnote{อารญฺญเกน (สพฺพตฺถ, อุปริ 193; วิ 5. 335 ปิฏฺเฐสุปิ.)} ภิกฺขุนา อปฺปมตฺเตน อาตาปินา ปหิตตฺเตน วิหริตุํ}

\vspace{\tipitakaparskip}
\csromancenter{(8) 3.\csromanspacer{} โยธาชีววคฺค อปฺปตฺตสฺส ปตฺติยา อนธิคตสฺส อธิคมาย อสจฺฉิกตสฺส สจฺฉิกิริยาย.}
\prose{\csromanfolio{89}กตมานิ ปญฺจ? อิธ ภิกฺขเว อารญฺญิโก ภิกฺขุ อิติ ปฏิสญฺจิกฺขติ “อหํ โข เอตรหิ เอกโก อรญฺเญ วิหรามิ, เอกกํ โข ปน มํ\footnote{เอกกํ โข ปน (สฺยา, กํ)} อรญฺเญ วิหรนฺตํ อหิ วา มํ ฑํเสยฺย, วิจฺฉิโก\footnote{วิจฺฉิกา (สฺยา)} วา มํ ฑํเสยฺย, สตปที วา มํ ฑํเสยฺย, เตน เม อสฺส กาลํกิริยา, โส มมสฺส อนฺตราโย.\csromanspacer{} หนฺทาหํ วีริยํ อารภามิ อปฺปตฺตสฺส ปตฺติยา อนธิคตสฺส อธิคมาย อสจฺฉิกตสฺส สจฺฉิกิริยายา”ติ.\csromanspacer{} อิทํ ภิกฺขเว ปฐมํ อนาคตภยํ สมฺปสฺสมาเนน อลเมว อารญฺญิเกน ภิกฺขุนา อปฺปมตฺเตน อาตาปินา ปหิตตฺเตน วิหริตุํ อปฺปตฺตสฺส ปตฺติยา อนธิคตสฺส อธิคมาย อสจฺฉิกตสฺส สจฺฉิกิริยาย.}

\prose{ปุน จปรํ ภิกฺขเว อารญฺญิโก ภิกฺขุ อิติ ปฏิสญฺจิกฺขติ “อหํ โข เอตรหิ เอกโก อรญฺเญ วิหรามิ, เอกโก โข ปนาหํ อรญฺเญ วิหรนฺโต อุปกฺขลิตฺวา วา ปปเตยฺยํ, ภตฺตํ วา ภุตฺตํ เม พฺยาปชฺเชยฺย, ปิตฺตํ วา เม กุปฺเปยฺย, เสมฺหํ วา เม กุปฺเปยฺย, สตฺถกา วา เม วาตา กุปฺเปยฺยุํ, เตน เม อสฺส กาลํกิริยา, โส มมสฺส อนฺตราโย.\csromanspacer{} หนฺทาหํวีริยํ อารภามิ อปฺปตฺตสฺส ปตฺติยา อนธิคตสฺส อธิคมาย อสจฺฉิกตสฺส สจฺฉิกิริยายา”ติ.\csromanspacer{} อิทํ ภิกฺขเว ทุติยํ อนาคตภยํ สมฺปสฺสมาเนน อลเมว อารญฺญิเกน ภิกฺขุนา อปฺปมตฺเตน อาตาปินา ปหิตตฺเตน วิหริตุํ อปฺปตฺตสฺส ปตฺติยา อนธิคตสฺส อธิคมาย อสจฺฉิกตสฺส สจฺฉิกิริยาย.}

\prose{ปุน จปรํ ภิกฺขเว อารญฺญิโก ภิกฺขุ อิติ ปฏิสญฺจิกฺขติ “อหํ โข เอตรหิ เอกโก อรญฺเญ วิหรามิ, เอกโก โข ปนาหํ อรญฺเญ วิหรนฺโต วาเฬหิ สมาคจฺเฉยฺยํ สีเหน วา พฺยคฺเฆน วา ทีปินา วา อจฺเฉน วา ตรจฺเฉน วา, เต มํ ชีวิตา โวโรเปยฺยุํ, เตน เม อสฺส กาลํกิริยา, โส มมสฺส อนฺตราโย.\csromanspacer{} หนฺทาหํ วีริยํ อารภามิ อปฺปตฺตสฺส ปตฺติยา อนธิคตสฺส อธิคมาย อสจฺฉิกตสฺส สจฺฉิกิริยายา”ติ.\csromanspacer{} อิทํ ภิกฺขเว ตติยํ อนาคตภยํ สมฺปสฺสมาเนน อลเมว อารญฺญิเกน ภิกฺขุนา อปฺปมตฺเตน อาตาปินา ปหิตตฺเตน \csromanfolio{90}วิหริตุํ อปฺปตฺตสฺส ปตฺติยา อนธิคตสฺส อธิคมาย อสจฺฉิกตสฺส สจฺฉิกิริยาย.}

\prose{ปุน จปรํ ภิกฺขเว อารญฺญิโก ภิกฺขุ อิติ ปฏิสญฺจิกฺขติ “อหํ โข เอตรหิ เอกโก อรญฺเญ วิหรามิ, เอกโก โข ปนาหํ อรญฺเญ วิหรนฺโต มาณเวหิ สมาคจฺเฉยฺยํ กตกมฺเมหิ วา อกตกมฺเมหิ วา, เต มํ ชีวิตา โวโรเปยฺยุํ, เตน เม อสฺส กาลํกิริยา, โส มมสฺส อนฺตราโย.\csromanspacer{} หนฺทาหํ วีริยํ อารภามิ อปฺปตฺตสฺส ปตฺติยา อนธิคตสฺส อธิคมาย อสจฺฉิกตสฺส สจฺฉิกิริยายา”ติ.\csromanspacer{} อิทํ ภิกฺขเว จตุตฺถํ อนาคตภยํ สมฺปสฺสมาเนน อลเมว อารญฺญิเกน ภิกฺขุนา อปฺปมตฺเตน อาตาปินา ปหิตตฺเตน วิหริตุํ อปฺปตฺตสฺส ปตฺติยา อนธิคตสฺส อธิคมาย อสจฺฉิกตสฺส สจฺฉิกิริยาย.}

\prose{ปุน จปรํ ภิกฺขเว อารญฺญิโก ภิกฺขุ อิติ ปฏิสญฺจิกฺขติ “อหํ โข เอตรหิ เอกโก อรญฺเญ วิหรามิ, สนฺติ โข ปนารญฺเญ วาฬา อมนุสฺสา, เต มํ ชีวิตา โวโรเปยฺยุํ, เตน เม อสฺส กาลํกิริยา, โส มมสฺส อนฺตราโย.\csromanspacer{} หนฺทาหํ วีริยํ อารภามิ อปฺปตฺตสฺส ปตฺติยา อนธิคตสฺส อธิคมาย อสจฺฉิกตสฺส สจฺฉิกิริยายา”ติ.\csromanspacer{} อิทํ ภิกฺขเว ปญฺจมํ อนาคตภยํ สมฺปสฺสมาเนน อลเมว อารญฺญิเกน ภิกฺขุนา อปฺปมตฺเตน อาตาปินา ปหิตตฺเตน วิหริตุํ อปฺปตฺตสฺส ปตฺติยา อนธิคตสฺส อธิคมาย อสจฺฉิกตสฺส สจฺฉิกิริยาย.}

\prose{อิมานิ โข ภิกฺขเว ปญฺจ อนาคตภยานิ สมฺปสฺสมาเนน อลเมว อารญฺญิเกน ภิกฺขุนา อปฺปมตฺเตน อาตาปินา ปหิตตฺเตน วิหริตุํ อปฺปตฺตสฺส ปตฺติยา อนธิคตสฺส อธิคมาย อสจฺฉิกตสฺส สจฺฉิกิริยายาติ.\csromanspacer{}\csromanspacer{}สตฺตมํ.}

\csromanheader{2}{8. ทุติย อนาคตภยสุตฺต}
\proseitem{78}{ปญฺจิมานิ ภิกฺขเว อนาคตภยานิ สมฺปสฺสมาเนน อลเมว ภิกฺขุนา อปฺปมตฺเตน อาตาปินา ปหิตตฺเตน วิหริตุํ อปฺปตฺตสฺส ปตฺติยา อนธิคตสฺส อธิคมาย อสจฺฉิกตสฺส สจฺฉิกิริยาย.\csromanspacer{} กตมานิ ปญฺจ? อิธ ภิกฺขเว ภิกฺขุ อิติ ปฏิสญฺจิกฺขติ “อหํ โข เอตรหิ ทหโร ยุวา สุสุกาฬเกโส ภเทฺรน โยพฺพเนน สมนฺนาคโต ปฐเมน}

\vspace{\tipitakaparskip}
\csromancenter{(8) 3.\csromanspacer{} โยธาชีววคฺค วยสา.\csromanspacer{} โหติ โข ปน โส สมโย, ยํ อิมํ กายํ ชรา ผุสติ, ชิณฺเณน โข ปน ชราย อภิภูเตน น สุกรํ พุทฺธานํ สาสนํ มนสิ กาตุํ, น สุกรานิ อรญฺญวนปตฺถานิ ปนฺตานิ เสนาสนานิ ปฏิเสวิตุํ.\csromanspacer{} ปุรา มํ โส ธมฺโม อาคจฺฉติ อนิฏฺโฐ อกนฺโต อมนาโป, หนฺทาหํ ปฏิกจฺเจว\footnote{ปฏิคจฺเจว (สี)} วีริยํ อารภามิ อปฺปตฺตสฺส ปตฺติยา อนธิคตสฺส อธิคมาย อสจฺฉิกตสฺส สจฺฉิกิริยาย, เยนาหํ ธมฺเมน สมนฺนาคโต ชิณฺณโกปิ ผาสุํ\footnote{ผาสุ (อิ, ก)} วิหริสฺสามี”ติ.\csromanspacer{} อิทํ ภิกฺขเว ปฐมํ อนาคตภยํ สมฺปสฺสมาเนน อลเมว ภิกฺขุนา อปฺปมตฺเตน อาตาปินา ปหิตตฺเตน วิหริตุํ อปฺปตฺตสฺส ปตฺติยา อนธิคตสฺส อธิคมาย อสจฺฉิกตสฺส สจฺฉิกิริยาย.}
\prose{\csromanfolio{91}ปุน จปรํ ภิกฺขเว ภิกฺขุ อิติ ปฏิสญฺจิกฺขติ “อหํ โข เอตรหิ อปฺปาพาโธ อปฺปาตงฺโก สมเวปากินิยา คหณิยา สมนฺนาคโต นาติสีตาย นาจฺจุณฺหาย มชฺฌิมาย ปธานกฺขมาย.\csromanspacer{} โหติ โข ปน โส สมโย, ยํ อิมํ กายํ พฺยาธิ ผุสติ, พฺยาธิเตน โข ปน พฺยาธินา อภิภูเตน\footnote{พฺยาธาภิภูเตน (สี, อิ, ก)} น สุกรํ พุทฺธานํ สาสนํ มนสิ กาตุํ, น สุกรานิ อรญฺญวนปตฺถานิ ปนฺตานิ เสนาสนานิ ปฏิเสวิตุํ.\csromanspacer{} ปุรา มํ โส ธมฺโม อาคจฺฉติ อนิฏฺโฐ อกนฺโต อมนาโป, หนฺทาหํ ปฏิกจฺเจว วีริยํ อารภามิ อปฺปตฺตสฺส ปตฺติยา อนธิคตสฺส อธิคมาย อสจฺฉิกตสฺส สจฺฉิกิริยาย, เยนาหํ ธมฺเมน สมนฺนาคโต พฺยาธิโตปิ ผาสุํ วิหริสฺสามี”ติ.\csromanspacer{} อิทํ ภิกฺขเว ทุติยํ อนาคตภยํ สมฺปสฺสมาเนน อลเมว ภิกฺขุนา อปฺปมตฺเตน อาตาปินา ปหิตตฺเตน วิหริตุํ อปฺปตฺตสฺส ปตฺติยา อนธิคตสฺส อธิคมาย อสจฺฉิกตสฺส สจฺฉิกิริยาย.}

\prose{ปุน จปรํ ภิกฺขเว ภิกฺขุ อิติ ปฏิสญฺจิกฺขติ “เอตรหิ โข สุภิกฺขํ สุสสฺสํ สุลภปิณฺฑํ สุกรํ อุญฺเฉน ปคฺคเหน ยาเปตุํ.\csromanspacer{} โหติ โข ปน โส สมโย, ยํ ทุพฺภิกฺขํ โหติ ทุสฺสสฺสํ ทุลฺลภปิณฺฑํ น สุกรํ อุญฺเฉน ปคฺคเหน ยาเปตุํ, ทุพฺภิกฺเข โข ปน มนุสฺสา เยน สุภิกฺขํ เตน สงฺกมนฺติ\footnote{เตนุปสงฺกมนฺติ (ก)}, ตตฺถ สงฺคณิกวิหาโร โหติ อากิณฺณวิหาโร.\csromanspacer{} สงฺคณิกวิหาเร โข ปน สติ อากิณฺณวิหาเร น สุกรํ พุทฺธานํ สาสนํ \csromanfolio{92}มนสิ กาตุํ, น สุกรานิ อรญฺญวนปตฺถานิ ปนฺตานิ เสนาสนานิ ปฏิเสวิตุํ.\csromanspacer{} ปุรา มํ โส ธมฺโม อาคจฺฉติ อนิฏฺโฐ อกนฺโต อมนาโป, หนฺทาหํ ปฏิกจฺเจว วีริยํ อารภามิ อปฺปตฺตสฺส ปตฺติยา อนธิคตสฺส อธิคมาย อสจฺฉิกตสฺส สจฺฉิกิริยาย, เยนาหํ ธมฺเมน สมนฺนาคโต ทุพฺภิกฺเขปิ ผาสุํ วิหริสฺสามี”ติ.\csromanspacer{} อิทํ ภิกฺขเว ตติยํ อนาคตภยํ สมฺปสฺสมาเนน อลเมว ภิกฺขุนา อปฺปมตฺเตน อาตาปินา ปหิตตฺเตน วิหริตุํ อปฺปตฺตสฺส ปตฺติยา อนธิคตสฺส อธิคมาย อสจฺฉิกตสฺส สจฺฉิกิริยาย.}

\prose{ปุน จปรํ ภิกฺขเว ภิกฺขุ อิติ ปฏิสญฺจิกฺขติ “เอตรหิ โข มนุสฺสา สมคฺคา สมฺโมทมานา อวิวทมานา ขีโรทกีภูตา อญฺญมญฺญํ ปิยจกฺขูหิ สมฺปสฺสนฺตา วิหรนฺติ.\csromanspacer{} โหติ โข ปน โส สมโย, ยํ ภยํ โหติ อฏวิสงฺโกโป, จกฺกสมารูฬฺหา ชานปทา ปริยายนฺติ, ภเย โข ปน สติ มนุสฺสา เยน เขมํ เตน สงฺกมนฺติ.\csromanspacer{} ตตฺถ สงฺคณิกวิหาโร โหติ อากิณฺณวิหาโร.\csromanspacer{} สงฺคณิกวิหาเร โข ปน สติ อากิณฺณวิหาเร น สุกรํ พุทฺธานํ สาสนํ มนสิ กาตุํ, น สุกรานิ อรญฺญวนปตฺถานิ ปนฺตานิ เสนาสนานิ ปฏิเสวิตุํ.\csromanspacer{} ปุรา มํ โส ธมฺโม อาคจฺฉติ อนิฏฺโฐ อกนฺโต อมนาโป, หนฺทาหํ ปฏิกจฺเจว วีริยํ อารภามิ อปฺปตฺตสฺส ปตฺติยา อนธิคตสฺส อธิคมาย อสจฺฉิกตสฺส สจฺฉิกิริยาย, เยนาหํ ธมฺเมน สมนฺนาคโต ภเยปิ ผาสุํ วิหริสฺสามี”ติ.\csromanspacer{} อิทํ ภิกฺขเว จตุตฺถํ อนาคตภยํ สมฺปสฺสมาเนน อลเมว ภิกฺขุนา อปฺปมตฺเตน อาตาปินา ปหิตตฺเตน วิหริตุํ อปฺปตฺตสฺส ปตฺติยา อนธิคตสฺส อธิคมาย อสจฺฉิกตสฺส สจฺฉิกิริยาย.}

\prose{ปุน จปรํ ภิกฺขเว ภิกฺขุ อิติ ปฏิสญฺจิกฺขติ “เอตรหิ โข สํโฆ สมคฺโค สมฺโมทมาโน อวิวทมาโน เอกุทฺเทโส ผาสุ วิหรติ.\csromanspacer{} โหติ โข ปน โส สมโย, ยํ สํโฆ ภิชฺชติ.\csromanspacer{} สํเฆ โข ปน ภินฺเน น สุกรํ พุทฺธานํ สาสนํ มนสิ กาตุํ, น สุกรานิ อรญฺญวนปตฺถานิ ปนฺตานิ เสนาสนานิ ปฏิเสวิตุํ.\csromanspacer{} ปุรา มํ โส ธมฺโม อาคจฺฉติ อนิฏฺโฐ อกนฺโต อมนาโป, หนฺทาหํ ปฏิกจฺเจว วีริยํ อารภามิ อปฺปตฺตสฺส ปตฺติยา อนธิคตสฺส อธิคมาย อสจฺฉิกตสฺส สจฺฉิกิริยาย, เยนาหํ ธมฺเมน สมนฺนาคโต ภินฺเนปิ สํเฆ ผาสุํ วิหริสฺสามี”ติ.\csromanspacer{} อิทํ ภิกฺขเว ปญฺจมํ อนาคตภยํ สมฺปสฺสมาเนน อลเมว ภิกฺขุนา}

\vspace{\tipitakaparskip}
\csromancenter{(8) 3.\csromanspacer{} โยธาชีววคฺค อปฺปมตฺเตน อาตาปินา ปหิตตฺเตน วิหริตุํ อปฺปตฺตสฺส ปตฺติยา อนธิคตสฺส อธิคมาย อสจฺฉิกตสฺส สจฺฉิกิริยาย.}
\prose{\csromanfolio{93}อิมานิ โข ภิกฺขเว ปญฺจ อนาคตภยานิ สมฺปสฺสมาเนน อลเมว ภิกฺขุนา อปฺปมตฺเตน อาตาปินา ปหิตตฺเตน วิหริตุํ อปฺปตฺตสฺส ปตฺติยา อนธิคตสฺส อธิคมาย อสจฺฉิกตสฺส สจฺฉิกิริยายาติ.\csromanspacer{}\csromanspacer{}อฏฺฐมํ.}

\csromanheader{2}{9. ตติย อนาคตภยสุตฺต}
\proseitem{79}{ปญฺจิมานิ ภิกฺขเว อนาคตภยานิ เอตรหิ อสมุปฺปนฺนานิ อายติํ สมุปฺปชฺชิสฺสนฺติ, ตานิ โว\footnote{โข (กตฺถจิ)} ปฏิพุชฺฌิตพฺพานิ, ปฏิพุชฺฌิตฺวา จ เตสํ ปหานาย วายมิตพฺพํ.}

\prose{กตมานิ ปญฺจ? ภวิสฺสนฺติ ภิกฺขเว ภิกฺขู อนาคตมทฺธานํ อภาวิตกายา อภาวิตสีลา อภาวิตจิตฺตา อภาวิตปญฺญา, เต อภาวิตกายา สมานา อภาวิตสีลา อภาวิตจิตฺตา อภาวิตปญฺญา อญฺเญ อุปสมฺปาเทสฺสนฺติ, เตปิ\footnote{เต (สี)} น สกฺขิสฺสนฺติ วิเนตุํ อธิสีเล อธิจิตฺเต อธิปญฺญาย.\csromanspacer{} เตปิ ภวิสฺสนฺติ อภาวิตกายา อภาวิตสีลา อภาวิตจิตฺตา อภาวิตปญฺญา, เต อภาวิตกายา สมานา อภาวิตสีลา อภาวิตจิตฺตา อภาวิตปญฺญา อญฺเญ อุปสมฺปาเทสฺสนฺติ, เตปิ\footnote{เต (?)} น สกฺขิสฺสนฺติ วิเนตุํ อธิสีเล อธิจิตฺเต อธิปญฺญาย.\csromanspacer{} เตปิ ภวิสฺสนฺติ อภาวิตกายา อภาวิตสีลา อภาวิตจิตฺตา อภาวิตปญฺญา.\csromanspacer{} อิติ โข ภิกฺขเว ธมฺมสนฺโทสา วินยสนฺโทโส, วินยสนฺโทสา ธมฺมสนฺโทโส.\csromanspacer{} อิทํ ภิกฺขเว ปฐมํ อนาคตภยํ เอตรหิ อสมุปฺปนฺนํ อายติํ สมุปฺปชฺชิสฺสติ, ตํ โว ปฏิพุชฺฌิตพฺพํ, ปฏิพุชฺฌิตฺวา จ ตสฺส ปหานาย วายมิตพฺพํ.}

\prose{ปุน จปรํ ภิกฺขเว ภวิสฺสนฺติ ภิกฺขู อนาคตมทฺธานํ อภาวิตกายา อภาวิตสีลา อภาวิตจิตฺตา อภาวิตปญฺญา, เต อภาวิตกายา สมานา อภาวิตสีลา อภาวิตจิตฺตา อภาวิตปญฺญา อญฺเญสํ นิสฺสยํ ทสฺสนฺติ, เตปิ น สกฺขิสฺสนฺติ วิเนตุํ อธิสีเล อธิจิตฺเต อธิปญฺญาย.\csromanspacer{} เตปิ ภวิสฺสนฺติ อภาวิตกายา อภาวิตสีลา \csromanfolio{94}อภาวิตจิตฺตา อภาวิตปญฺญา, เต อภาวิตกายา สมานา อภาวิตสีลา อภาวิตจิตฺตา อภาวิตปญฺญา อญฺเญสํ นิสฺสยํ ทสฺสนฺติ, เตปิ น สกฺขิสฺสนฺติ วิเนตุํ อธิสีเล อธิจิตฺเต อธิปญฺญาย.\csromanspacer{} เตปิ ภวิสฺสนฺติ อภาวิตกายา อภาวิตสีลา อภาวิตจิตฺตา อภาวิตปญฺญา.\csromanspacer{} อิติ โข ภิกฺขเว ธมฺมสนฺโทสา วินยสนฺโทโส, วินยสนฺโทสา ธมฺมสนฺโทโส.\csromanspacer{} อิทํ ภิกฺขเว ทุติยํ อนาคตภยํ เอตรหิ อสมุปฺปนฺนํ อายติํ สมุปฺปชฺชิสฺสติ, ตํ โว ปฏิพุชฺฌิตพฺพํ, ปฏิพุชฺฌิตฺวา จ ตสฺส ปหานาย วายมิตพฺพํ.}

\prose{ปุน จปรํ ภิกฺขเว ภวิสฺสนฺติ ภิกฺขู อนาคตมทฺธานํ อภาวิตกายา อภาวิตสีลา อภาวิตจิตฺตา อภาวิตปญฺญา, เต อภาวิตกายา สมานา อภาวิตสีลา อภาวิตจิตฺตา อภาวิตปญฺญา อภิธมฺมกถํ เวทลฺลกถํ กเถนฺตา กณฺหธมฺมํ โอกฺกมมานา น พุชฺฌิสฺสนฺติ.\csromanspacer{} อิติ โข ภิกฺขเว ธมฺมสนฺโทสา วินยสนฺโทโส, วินยสนฺโทสา ธมฺมสนฺโทโส.\csromanspacer{} อิทํ ภิกฺขเว ตติยํ อนาคตภยํ เอตรหิ อสมุปฺปนฺนํ อายติํ สมุปฺปชฺชิสฺสติ, ตํ โว ปฏิพุชฺฌิตพฺพํ, ปฏิพุชฺฌิตฺวา จ ตสฺส ปหานาย วายมิตพฺพํ.}

\prose{ปุน จปรํ ภิกฺขเว ภวิสฺสนฺติ ภิกฺขู อนาคตมทฺธานํ อภาวิตกายา อภาวิตสีลา อภาวิตจิตฺตา อภาวิตปญฺญา, เต อภาวิตกายา สมานา อภาวิตสีลา อภาวิตจิตฺตา อภาวิตปญฺญา เย เต สุตฺตนฺตา ตถาคตภาสิตา คมฺภีรา คมฺภีรตฺถา โลกุตฺตรา สุญฺญตาปฺปฏิสํยุตฺตา, เตสุ ภญฺญมาเนสุ น สุสฺสูสิสฺสนฺติ, น โสตํ โอทหิสฺสนฺติ, น อญฺญา จิตฺตํ อุปฏฺฐเปสฺสนฺติ, น จ เต ธมฺเม อุคฺคเหตพฺพํ ปริยาปุณิตพฺพํ มญฺญิสฺสนฺติ.\csromanspacer{} เย ปน เต สุตฺตนฺตา กวิตา\footnote{กวิกตา (สี, สฺยา, กํ อิ สํ 1. 457 ปิฏฺเฐ จ) ฏีกา โอโลเกตพฺพา.} กาเวยฺยา จิตฺตกฺขรา จิตฺตพฺยญฺชนา พาหิรกา สาวกภาสิตา, เตสุ ภญฺญมาเนสุ สุสฺสูสิสฺสนฺติ, โสตํ โอทหิสฺสนฺติ, อญฺญา จิตฺตํ อุปฏฺฐเปสฺสนฺติ, เต จ ธมฺเม อุคฺคเหตพฺพํ ปริยาปุณิตพฺพํ มญฺญิสฺสนฺติ.\csromanspacer{} อิติ โข ภิกฺขเว ธมฺมสนฺโทสา วินยสนฺโทโส, วินยสนฺโทสา ธมฺมสนฺโทโส.\csromanspacer{} อิทํ ภิกฺขเว จตุตฺถํ อนาคตภยํ เอตรหิ อสมุปฺปนฺนํ อายติํ สมุปฺปชฺชิสฺสติ, ตํ โว ปฏิพุชฺฌิตพฺพํ, ปฏิพุชฺฌิตฺวา จ ตสฺส ปหานาย วายมิตพฺพํ.}

\chapterheadpageread{(8) 3. โยธาชีววคฺค}
\prose{\csromanfolio{95}ปุน จปรํ ภิกฺขเว ภวิสฺสนฺติ ภิกฺขู อนาคตมทฺธานํ อภาวิตกายา อภาวิตสีลา อภาวิตจิตฺตา อภาวิตปญฺญา, เต อภาวิตกายา สมานา อภาวิตสีลา อภาวิตจิตฺตา อภาวิตปญฺญา เถรา ภิกฺขู พาหุลิกา\footnote{พาหุลฺลิกา (สฺยา, กํ) อํ 1. 71, 245 ปิฏฺเฐสุปิ.} ภวิสฺสนฺติ สาถลิกา โอกฺกมเน ปุพฺพงฺคมา ปวิเวเก นิกฺขิตฺตธุรา, น วีริยํ อารภิสฺสนฺติ อปฺปตฺตสฺส ปตฺติยา อนธิคตสฺส อธิคมาย อสจฺฉิกตสฺส สจฺฉิกิริยาย.\csromanspacer{} เตสํ ปจฺฉิมา ชนตา ทิฏฺฐานุคติํ อาปชฺชิสฺสติ.\csromanspacer{} สาปิ ภวิสฺสติ พาหุลิกา สาถลิกา โอกฺกมเน ปุพฺพงฺคมา ปวิเวเก นิกฺขิตฺตธุรา, น วีริยํ อารภิสฺสติ อปฺปตฺตสฺส ปตฺติยา อนธิคตสฺส อธิคมาย อสจฺฉิกตสฺส สจฺฉิกิริยาย.\csromanspacer{} อิติ โข ภิกฺขเว ธมฺมสนฺโทสา วินยสนฺโทโส, วินยสนฺโทสา ธมฺมสนฺโทโส.\csromanspacer{} อิทํ ภิกฺขเว ปญฺจมํ อนาคตภยํ เอตรหิ อสมุปฺปนฺนํ อายติํ สมุปฺปชฺชิสฺสติ, ตํ โว ปฏิพุชฺฌิตพฺพํ, ปฏิพุชฺฌิตฺวา จ ตสฺส ปหานาย วายมิตพฺพํ.\csromanspacer{} อิมานิ โข ภิกฺขเว ปญฺจ อนาคตภยานิ เอตรหิ อสมุปฺปนฺนานิ อายติํ สมุปฺปชฺชิสฺสนฺติ, ตานิ โว ปฏิพุชฺฌิตพฺพานิ, ปฏิพุชฺฌิตฺวา จ เตสํ ปหานาย วายมิตพฺพนฺติ.\csromanspacer{}\csromanspacer{}นวมํ.}

\csromanheader{2}{10. จตุตฺถ อนาคตภยสุตฺต}
\proseitem{80}{ปญฺจิมานิ ภิกฺขเว อนาคตภยานิ เอตรหิ อสมุปฺปนฺนานิ อายติํ สมุปฺปชฺชิสฺสนฺติ, ตานิ โว ปฏิพุชฺฌิตพฺพานิ, ปฏิพุชฺฌิตฺวา จ เตสํ ปหานาย วายมิตพฺพํ.}

\prose{กตมานิ ปญฺจ? ภวิสฺสนฺติ ภิกฺขเว ภิกฺขู อนาคตมทฺธานํ จีวเร กลฺยาณกามา, เต จีวเร กลฺยาณกามา สมานา ริญฺจิสฺสนฺติ ปํสุกูลิกตฺตํ, ริญฺจิสฺสนฺติ อรญฺญวนปตฺถานิ ปนฺตานิ เสนาสนานิ, คามนิคมราชธานีสุ โอสริตฺวา วาสํ กปฺเปสฺสนฺติ, จีวรเหตุ จ อเนกวิหิตํ อเนสนํ อปฺปติรูปํ อาปชฺชิสฺสนฺติ.\csromanspacer{} อิทํ ภิกฺขเว ปฐมํ อนาคตภยํ เอตรหิ อสมุปฺปนฺนํ อายติํ สมุปฺปชฺชิสฺสติ, ตํ โว ปฏิพุชฺฌิตพฺพํ, ปฏิพุชฺฌิตฺวา จ ตสฺส ปหานาย วายมิตพฺพํ.}

\prose{ปุน จปรํ ภิกฺขเว ภวิสฺสนฺติ ภิกฺขู อนาคตมทฺธานํ ปิณฺฑปาเต กลฺยาณกามา, เต ปิณฺฑปาเต กลฺยาณกามา สมานา ริญฺจิสฺสนฺติ \csromanfolio{96}ปิณฺฑปาติกตฺตํ, ริญฺจิสฺสนฺติ อรญฺญวนปตฺถานิ ปนฺตานิ เสนาสนานิ, คามนิคมราชธานีสุ โอสริตฺวา วาสํ กปฺเปสฺสนฺติ ชิวฺหคฺเคน รสคฺคานิ ปริเยสมานา, ปิณฺฑปาตเหตุ จ อเนกวิหิตํ อเนสนํ อปฺปติรูปํ อาปชฺชิสฺสนฺติ.\csromanspacer{} อิทํ ภิกฺขเว ทุติยํ อนาคตภยํ เอตรหิ อสมุปฺปนฺนํ อายติํ สมุปฺปชฺชิสฺสติ, ตํ โว ปฏิพุชฺฌิตพฺพํ, ปฏิพุชฺฌิตฺวา จ ตสฺส ปหานาย วายมิตพฺพํ.}

\prose{ปุน จปรํ ภิกฺขเว ภวิสฺสนฺติ ภิกฺขู อนาคตมทฺธานํ เสนาสเน กลฺยาณกามา, เต เสนาสเน กลฺยาณกามา สมานา ริญฺจิสฺสนฺติ รุกฺขมูลิกตฺตํ\footnote{อารญฺญกตฺตํ (สฺยา, กํ)}, ริญฺจิสฺสนฺติ อรญฺญวนปตฺถานิ ปนฺตานิ เสนาสนานิ, คามนิคมราชธานีสุ โอสริตฺวา วาสํ กปฺเปสฺสนฺติ, เสนาสนเหตุ จ อเนกวิหิตํ อเนสนํ อปฺปติรูปํ อาปชฺชิสฺสนฺติ.\csromanspacer{} อิทํ ภิกฺขเว ตติยํ อนาคตภยํ เอตรหิ อสมุปฺปนฺนํ อายติํ สมุปฺปชฺชิสฺสติ, ตํ โว ปฏิพุชฺฌิตพฺพํ, ปฏิพุชฺฌิตฺวา จ ตสฺส ปหานาย วายมิตพฺพํ.}

\prose{ปุน จปรํ ภิกฺขเว ภวิสฺสนฺติ ภิกฺขู อนาคตมทฺธานํ ภิกฺขุนี สิกฺขมานาสมณุทฺเทเสหิ สํสฏฺฐา วิหริสฺสนฺติ, ภิกฺขุนีสิกฺขมานาสมณุทฺเทเสหิ สํสคฺเค โข ปน ภิกฺขเว สติ เอตํ ปาฏิกงฺขํ “อนภิรตา วา พฺรหฺมจริยํ จริสฺสนฺติ, อญฺญตรํ วา สํกิลิฏฺฐํ อาปตฺติํ อาปชฺชิสฺสนฺติ, สิกฺขํ วา ปจฺจกฺขาย หีนายาวตฺติสฺสนฺติ”.\csromanspacer{} อิทํ ภิกฺขเว จตุตฺถํ อนาคตภยํ เอตรหิ อสมุปฺปนฺนํ อายติํ สมุปฺปชฺชิสฺสติ, ตํ โว ปฏิพุชฺฌิตพฺพํ, ปฏิพุชฺฌิตฺวา จ ตสฺส ปหานาย วายมิตพฺพํ.}

\prose{ปุน จปรํ ภิกฺขเว ภวิสฺสนฺติ ภิกฺขู อนาคตมทฺธานํ อารามิกสมณุทฺเทเสหิ สํสฏฺฐา วิหริสฺสนฺติ, อารามิกสมณุทฺเทเสหิ สํสคฺเค โข ปน ภิกฺขเว สติ เอตํ ปาฏิกงฺขํ “อเนกวิหิตํ สนฺนิธิการปริโภคํ อนุยุตฺตา วิหริสฺสนฺติ, โอฬาริกมฺปิ นิมิตฺตํ กริสฺสนฺติ ปถวิยาปิ หริตคฺเคปิ”.\csromanspacer{} อิทํ ภิกฺขเว ปญฺจมํ อนาคตภยํ เอตรหิ อสมุปฺปนฺนํ อายติํ สมุปฺปชฺชิสฺสติ, ตํ โว ปฏิพุชฺฌิตพฺพํ, ปฏิพุชฺฌิตฺวา จ ตสฺส ปหานาย วายมิตพฺพํ.}

\chapterheadpageread{\textbf{(9) 4. เถรวคฺค}}
\prose{\csromanfolio{97}อิมานิ โข ภิกฺขเว ปญฺจ อนาคตภยานิ เอตรหิ อสมุปฺปนฺนานิ อายติํ สมุปฺปชฺชิสฺสนฺติ, ตานิ โว ปฏิพุชฺฌิตพฺพานิ, ปฏิพุชฺฌิตฺวา จ เตสํ ปหานาย วายมิตพฺพนฺติ.\csromanspacer{}\csromanspacer{}ทสมํ.\csromanspacer{} โยธาชีววคฺโค ตติโย.}

\titlehead{\textbf{ตสฺสุทฺทานํ}}
\csromangathameasure{เทฺว เจโตวิมุตฺติผลา, เทฺว จ ธมฺมวิหาริโน. \\ โยธาชีวา จ เทฺว วุตฺตา, จตฺตาโร จ อนาคตาติ.}
\csromangathagroup{\csromangathastanza{เทฺว เจโตวิมุตฺติผลา, เทฺว จ ธมฺมวิหาริโน. \\ โยธาชีวา จ เทฺว วุตฺตา, จตฺตาโร จ อนาคตาติ.}}

\chapterhead{\textbf{(9) 4. เถรวคฺค}}
\csromanheader{2}{1. รชนียสุตฺต}
\proseitem{81}{ปญฺจหิ ภิกฺขเว ธมฺเมหิ สมนฺนาคโต เถโร ภิกฺขุ สพฺรหฺมจารีนํ อปฺปิโย จ โหติ อมนาโป จ อครุ จ อภาวนีโย จ.\csromanspacer{} กตเมหิ ปญฺจหิ? รชนีเย รชฺชติ, ทุสฺสนีเย\footnote{ทุสนีเย (สี, สฺยา, กํ, อิ)} ทุสฺสติ, โมหนีเย มุยฺหติ, กุปฺปนีเย\footnote{กุปนีเย (สี, สฺยา, กํ), โกปนีเย (อิ)} กุปฺปติ, มทนีเย มชฺชติ.\csromanspacer{} อิเมหิ โข ภิกฺขเว ปญฺจหิ ธมฺเมหิ สมนฺนาคโต เถโร ภิกฺขุ สพฺรหฺมจารีนํ อปฺปิโย จ โหติ อมนาโป จ อครุ จ อภาวนีโย จ.}

\prose{ปญฺจหิ ภิกฺขเว ธมฺเมหิ สมนฺนาคโต เถโร ภิกฺขุ สพฺรหฺมจารีนํ ปิโย จ โหติ มนาโป จ ครุ จ ภาวนีโย จ.\csromanspacer{} กตเมหิ ปญฺจหิ? รชนีเย น รชฺชติ, ทุสฺสนีเย น ทุสฺสติ, โมหนีเย น มุยฺหติ, กุปฺปนีเย น กุปฺปติ, มทนีเย น มชฺชติ.\csromanspacer{} อิเมหิ โข ภิกฺขเว ปญฺจหิ ธมฺเมหิ สมนฺนาคโต เถโร ภิกฺขุ สพฺรหฺมจารีนํ ปิโย จ โหติ มนาโป จ ครุ จ ภาวนีโย จาติ.\csromanspacer{}\csromanspacer{}ปฐมํ.}

\csromanheader{2}{2. วีตราคสุตฺต}
\proseitem{82}{ปญฺจหิ ภิกฺขเว ธมฺเมหิ สมนฺนาคโต เถโร ภิกฺขุ สพฺรหฺมจารีนํ อปฺปิโย จ โหติ อมนาโป จ อครุ จ อภาวนีโย จ. \csromanfolio{98}กตเมหิ ปญฺจหิ? อวีตราโค โหติ อวีตโทโส โหติ, อวีตโมโห โหติ, มกฺขี จ ปฬาสี จ.\csromanspacer{} อิเมหิ โข ภิกฺขเว ปญฺจหิ ธมฺเมหิ สมนฺนาคโต เถโร ภิกฺขุ สพฺรหฺมจารีนํ อปฺปิโย จ โหติ อมนาโป จ อครุ จ อภาวนีโย จ.}

\prose{ปญฺจหิ ภิกฺขเว ธมฺเมหิ สมนฺนาคโต เถโร ภิกฺขุ สพฺรหฺมจารีนํ ปิโย จ โหติ มนาโป จ ครุ จ ภาวนีโย จ.\csromanspacer{} กตเมหิ ปญฺจหิ? วีตราโค โหติ, วีตโทโส โหติ, วีตโมโห โหติ, อมกฺขี จ อปฬาสี จ.\csromanspacer{} อิเมหิ โข ภิกฺขเว ปญฺจหิ ธมฺเมหิ สมนฺนาคโต เถโร ภิกฺขุ สพฺรหฺมจารีนํ ปิโย จ โหติ มนาโป จ ครุ จ ภาวนีโย จาติ.\csromanspacer{}\csromanspacer{}ทุติยํ.}

\csromanheader{2}{3. กุหกสุตฺต}
\proseitem{83}{ปญฺจหิ ภิกฺขเว ธมฺเมหิ สมนฺนาคโต เถโร ภิกฺขุ สพฺรหฺมจารีนํ อปฺปิโย จ โหติ อมนาโป จ อครุ จ อภาวนีโย จ.\csromanspacer{} กตเมหิ ปญฺจหิ? กุหโก จ โหติ ลปโก จ เนมิตฺติโก\footnote{นิมิตฺติโก (สฺยา, กํ), นิมิตฺตโก (ก)} จ นิปฺเปสิโก จ ลาเภน จ ลาภํ นิชิคีสิตา\footnote{นิชิคิํสิตา (สี, สฺยา, กํ, อิ)}.\csromanspacer{} อิเมหิ โข ภิกฺขเว ปญฺจหิ ธมฺเมหิ สมนฺนาคโต เถโร ภิกฺขุ สพฺรหฺมจารีนํ อปฺปิโย จ โหติ อมนาโป จ อครุ จ อภาวนีโย จ.}

\prose{ปญฺจหิ ภิกฺขเว ธมฺเมหิ สมนฺนาคโต เถโร ภิกฺขุ สพฺรหฺมจารีนํ ปิโย จ โหติ มนาโป จ ครุ จ ภาวนีโย จ.\csromanspacer{} กตเมหิ ปญฺจหิ? น จ กุหโก โหติ น จ ลปโก น จ เนมิตฺติโก น จ นิปฺเปสิโก น จ ลาเภน ลาภํ นิชิคีสิตา.\csromanspacer{} อิเมหิ โข ภิกฺขเว ปญฺจหิ ธมฺเมหิ สมนฺนาคโต เถโร ภิกฺขุ สพฺรหฺมจารีนํ ปิโย จ โหติ มนาโป จ ครุ จ ภาวนีโย จาติ.\csromanspacer{}\csromanspacer{}ตติยํ.}

\csromanheader{2}{4. อสฺสทฺธสุตฺต}
\proseitem{84}{ปญฺจหิ ภิกฺขเว ธมฺเมหิ สมนฺนาคโต เถโร ภิกฺขุ สพฺรหฺมจารีนํ อปฺปิโย จ โหติ อมนาโป จ อครุ จ อภาวนีโย จ.\csromanspacer{} กตเมหิ ปญฺจหิ? อสฺสทฺโธ โหติ, อหิริโก โหติ, อโนตฺตปฺปี}

\vspace{\tipitakaparskip}
\csromancenter{(9) 4.\csromanspacer{} เถรวคฺค โหติ, กุสีโต โหติ, ทุปฺปญฺโญ โหติ.\csromanspacer{} อิเมหิ โข ภิกฺขเว ปญฺจหิ ธมฺเมหิ สมนฺนาคโต เถโร ภิกฺขุ สพฺรหฺมจารีนํ อปฺปิโย จ โหติ อมนาโป จ อครุ จ อภาวนีโย จ.}
\prose{\csromanfolio{99}ปญฺจหิ ภิกฺขเว ธมฺเมหิ สมนฺนาคโต เถโร ภิกฺขุ สพฺรหฺมจารีนํ ปิโย จ โหติ มนาโป จ ครุ จ ภาวนีโย จ.\csromanspacer{} กตเมหิ ปญฺจหิ? สทฺโธ โหติ, หิรีมา โหติ, โอตฺตปฺปี โหติ, อารทฺธวีริโย โหติ, ปญฺญวา โหติ.\csromanspacer{} อิเมหิ โข ภิกฺขเว ปญฺจหิ ธมฺเมหิ สมนฺนาคโต เถโร ภิกฺขุ สพฺรหฺมจารีนํ ปิโย จ โหติ มนาโป จ ครุ จ ภาวนีโย จาติ.\csromanspacer{}\csromanspacer{}จตุตฺถํ.}

\csromanheader{2}{5. อกฺขมสุตฺต}
\proseitem{85}{ปญฺจหิ ภิกฺขเว ธมฺเมหิ สมนฺนาคโต เถโร ภิกฺขุ สพฺรหฺมจารีนํ อปฺปิโย จ โหติ อมนาโป จ อครุ จ อภาวนีโย จ.\csromanspacer{} กตเมหิ ปญฺจหิ? อกฺขโม โหติ รูปานํ, อกฺขโม สทฺทานํ, อกฺขโม คนฺธานํ, อกฺขโม รสานํ, อกฺขโม โผฏฺฐพฺพานํ.\csromanspacer{} อิเมหิ โข ภิกฺขเว ปญฺจหิ ธมฺเมหิ สมนฺนาคโต เถโร ภิกฺขุ สพฺรหฺมจารีนํ อปฺปิโย จ โหติ อมนาโป จ อครุ จ อภาวนีโย จ.}

\prose{ปญฺจหิ ภิกฺขเว ธมฺเมหิ สมนฺนาคโต เถโร ภิกฺขุ สพฺรหฺมจารีนํ ปิโย จ โหติ มนาโป จ ครุ จ ภาวนีโย จ.\csromanspacer{} กตเมหิ ปญฺจหิ? ขโม โหติ รูปานํ, ขโม สทฺทานํ, ขโม คนฺธานํ, ขโม รสานํ, ขโม โผฏฺฐพฺพานํ.\csromanspacer{} อิเมหิ โข ภิกฺขเว ปญฺจหิ ธมฺเมหิ สมนฺนาคโต เถโร ภิกฺขุ สพฺรหฺมจารีนํ ปิโย จ โหติ มนาโป จ ครุ จ ภาวนีโย จาติ.\csromanspacer{}\csromanspacer{}ปญฺจมํ.}

\csromanheader{2}{6. ปฏิสมฺภิทาปฺปตฺตสุตฺต}
\proseitem{86}{ปญฺจหิ ภิกฺขเว ธมฺเมหิ สมนฺนาคโต เถโร ภิกฺขุ สพฺรหฺมจารีนํ ปิโย จ โหติ มนาโป จ ครุ จ ภาวนีโย จ.\csromanspacer{} กตเมหิ ปญฺจหิ? อตฺถปฏิ สมฺภิทาปตฺโต โหติ, ธมฺมปฏิสมฺภิทาปตฺโต โหติ, นิรุตฺติปฏิสมฺภิทาปตฺโต โหติ, ปฏิภานปฏิสมฺภิทาปตฺโต โหติ, ยานิ ตานิ สพฺรหฺมจารีนํ อุจฺจาวจานิ กิํกรณียานิ, ตตฺถ ทกฺโข โหติ อนลโส, ตตฺรุปายาย วีมํสาย สมนฺนาคโต อลํ \csromanfolio{100}กาตุํ อลํ สํวิธาตุํ.\csromanspacer{} อิเมหิ โข ภิกฺขเว ปญฺจหิ ธมฺเมหิ สมนฺนาคโต เถโร ภิกฺขุ สพฺรหฺมจารีนํ ปิโย จ โหติ มนาโป จ ครุ จ ภาวนีโย จาติ.\csromanspacer{}\csromanspacer{}ฉฏฺฐํ.}

\csromanheader{2}{7. สีลวนฺตสุตฺต}
\proseitem{87}{ปญฺจหิ ภิกฺขเว ธมฺเมหิ สมนฺนาคโต เถโร ภิกฺขุ สพฺรหฺมจารีนํ ปิโย จ โหติ มนาโป จ ครุ จ ภาวนีโย จ.\csromanspacer{} กตเมหิ ปญฺจหิ? สีลวา โหติ, ปาติโมกฺขสํวรสํวุโต วิหรติ อาจารโคจรสมฺปนฺโน, อณุมตฺเตสุ วชฺเชสุ ภยทสฺสาวี สมาทาย สิกฺขติ สิกฺขาปเทสุ.\csromanspacer{} พหุสฺสุโต โหติ สุตธโร สุตสนฺนิจโย, เย เต ธมฺมา อาทิกลฺยาณา มชฺเฌกลฺยาณา ปริโยสานกลฺยาณา สาตฺถํ สพฺยญฺชนํ\footnote{สาตฺถา สพฺยญฺชนา (สี)} เกวลปริปุณฺณํ ปริสุทฺธํ พฺรหฺมจริยํ อภิวทนฺติ, ตถารูปาสฺส ธมฺมา พหุสฺสุตา โหนฺติ ธาตา\footnote{ธตา (สี, สฺยา, กํ, อิ)} วจสา ปริจิตา มนสานุเปกฺขิตา ทิฏฺฐิยา สุปฺปฏิวิทฺธา.\csromanspacer{} กลฺยาณวาโจ โหติ กลฺยาณวากฺกรโณ โปริยา วาจาย สมนฺนาคโต วิสฺสฏฺฐาย อเนลคฬาย อตฺถสฺส วิญฺญาปนิยา, จตุนฺนํ ฌานานํ อาภิเจตสิกานํ ทิฏฺฐธมฺมสุขวิหารานํ นิกามลาภี โหติ อกิจฺฉลาภี อกสิรลาภี.\csromanspacer{} อาสวานํ ขยา อนาสวํ เจโตวิมุตฺติํ ปญฺญาวิมุตฺติํ ทิฏฺเฐว ธมฺเม สยํ อภิญฺญา สจฺฉิกตฺวา อุปสมฺปชฺช วิหรติ.\csromanspacer{} อิเมหิ โข ภิกฺขเว ปญฺจหิ ธมฺเมหิ สมนฺนาคโต เถโร ภิกฺขุ สพฺรหฺมจารีนํ ปิโย จ โหติ มนาโป จ ครุ จ ภาวนีโย จาติ.\csromanspacer{}\csromanspacer{}สตฺตมํ.}

\csromanheader{2}{8. เถรสุตฺต}
\proseitem{88}{ปญฺจหิ ภิกฺขเว ธมฺเมหิ สมนฺนาคโต เถโร ภิกฺขุ พหุชน-อหิตาย ปฏิปนฺโน โหติ พหุชน-อสุขาย พหุโน ชนสฺส อนตฺถาย อหิตาย ทุกฺขาย เทวมนุสฺสานํ.}

\prose{กตเมหิ ปญฺจหิ? เถโร โหติ รตฺตญฺญู จิรปพฺพชิโต, ญาโต โหติ ยสสฺสี สคหฏฺฐปพฺพชิตานํ\footnote{คหฏฺฐปพฺพชิตานํ (สี)} พหุชนปริวาโร, ลาภี โหติ จีวรปิณฺฑปาต- เสนาสนคิลานปฺปจฺจยเภสชฺชปริกฺขารานํ, พหุสฺสุโต โหติ สุตธโร สุตสนฺนิจโย, เย เต ธมฺมา อาทิกลฺยาณา}

\vspace{\tipitakaparskip}
\csromancenter{(9) 4.\csromanspacer{} เถรวคฺค มชฺเฌกลฺยาณา ปริโยสานกลฺยาณา สาตฺถํ สพฺยญฺชนํ เกวลปริปุณฺณํ ปริสุทฺธํ พฺรหฺมจริยํ อภิวทนฺติ, ตถารูปาสฺส ธมฺมา พหุสฺสุตา โหนฺติ ธาตา วจสา ปริจิตา มนสานุเปกฺขิตา ทิฏฺฐิยา อปฺปฏิวิทฺธา.\csromanspacer{} มิจฺฉาทิฏฺฐิโก โหติ วิปรีตทสฺสโน, โส พหุชนํ สทฺธมฺมา วุฏฺฐาเปตฺวา อสทฺธมฺเม ปติฏฺฐาเปติ. “เถโร ภิกฺขุ รตฺตญฺญู จิรปพฺพชิโต” อิติปิสฺส ทิฏฺฐานุคติํ อาปชฺชนฺติ, “ญาโต เถโร ภิกฺขุ ยสสฺสี สคหฏฺฐปพฺพชิตานํ พหุชนปริวาโร” อิติปิสฺส ทิฏฺฐานุคติํ อาปชฺชนฺติ, “ลาภี เถโร ภิกฺขุ จีวรปิณฺฑปาตเสนาสนคิลานปฺปจฺจยเภสชฺชปริกฺขารานํ” อิติปิสฺส ทิฏฺฐานุคติํ อาปชฺชนฺติ, “พหุสฺสุโต เถโร ภิกฺขุ สุตธโร สุตสนฺนิจโย” อิติปิสฺส ทิฏฺฐานุคติํ อาปชฺชนฺติ.\csromanspacer{} อิเมหิ โข ภิกฺขเว ปญฺจหิ ธมฺเมหิ สมนฺนาคโต เถโร ภิกฺขุ พหุชน-อหิตาย ปฏิปนฺโน โหติ พหุชน-อสุขาย พหุโน ชนสฺส อนตฺถาย อหิตาย ทุกฺขาย เทวมนุสฺสานํ.}
\prose{\csromanfolio{101}ปญฺจหิ ภิกฺขเว ธมฺเมหิ สมนฺนาคโต เถโร ภิกฺขุ พหุชนหิตาย ปฏิปนฺโน โหติ พหุชนสุขาย พหุโน ชนสฺส อตฺถาย หิตาย สุขาย เทวมนุสฺสานํ.}

\prose{กตเมหิ ปญฺจหิ? เถโร โหติ รตฺตญฺญู จิรปพฺพชิโต, ญาโต โหติ ยสสฺสี สคหฏฺฐปพฺพชิตานํ พหุชนปริวาโร, ลาภี โหติ จีวร ปิณฺฑปาต เสนาสน คิลานปฺปจฺจย เภสชฺช ปริกฺขารานํ, พหุสฺสุโต โหติ สุตธโร สุตสนฺนิจโย, เย เต ธมฺมา อาทิกลฺยาณา มชฺเฌกลฺยาณา ปริโยสานกลฺยาณา สาตฺถํ สพฺยญฺชนํ เกวลปริปุณฺณํ ปริสุทฺธํ พฺรหฺมจริยํ อภิวทนฺติ, ตถารูปาสฺส ธมฺมา พหุสฺสุตา โหนฺติ ธาตา วจสา ปริจิตา มนสานุเปกฺขิตา ทิฏฺฐิยา สุปฺปฏิวิทฺธา.\csromanspacer{} สมฺมาทิฏฺฐิโก โหติ อวิปรีตทสฺสโน, โส พหุชนํ อสทฺธมฺมา วุฏฺฐาเปตฺวา สทฺธมฺเม ปติฏฺฐาเปติ. “เถโร ภิกฺขุ รตฺตญฺญู จิรปพฺพชิโต” อิติปิสฺส ทิฏฺฐานุคติํ อาปชฺชนฺติ, “ญาโต เถโร ภิกฺขุ ยสสฺสี สคหฏฺฐปพฺพชิตานํ พหุชนปริวาโร” อิติปิสฺส ทิฏฺฐานุคติํ อาปชฺชนฺติ, “ลาภี เถโร ภิกฺขุ จีวรปิณฺฑปาตเสนาสนคิลานปฺปจฺจยเภสชฺชปริกฺขารานํ” อิติปิสฺส ทิฏฺฐานุคติํ อาปชฺชนฺติ, “พหุสฺสุโต เถโร ภิกฺขุ สุตธโร สุตสนฺนิจโย” อิติปิสฺส ทิฏฺฐานุคติํ อาปชฺชนฺติ.\csromanspacer{} อิเมหิ โข ภิกฺขเว ปญฺจหิ ธมฺเมหิ สมนฺนาคโต เถโร ภิกฺขุ พหุชนหิตาย ปฏิปนฺโน โหติ \csromanfolio{102}พหุชนสุขาย พหุโน ชนสฺส อตฺถาย หิตาย สุขาย เทวมนุสฺสานนฺติ.\csromanspacer{}\csromanspacer{}อฏฺฐมํ.}

\csromanheader{2}{9. ปฐมเสขสุตฺต}
\proseitem{89}{ปญฺจิเม ภิกฺขเว ธมฺมา เสขสฺส ภิกฺขุโน ปริหานาย สํวตฺตนฺติ.\csromanspacer{} กตเม ปญฺจ? กมฺมารามตา ภสฺสารามตา นิทฺทารามตา สงฺคณิการามตา ยถา วิมุตฺตํ จิตฺตํ น ปจฺจเวกฺขติ.\csromanspacer{} อิเม โข ภิกฺขเว ปญฺจ ธมฺมา เสขสฺส ภิกฺขุโน ปริหานาย สํวตฺตนฺติ.}

\prose{ปญฺจิเม ภิกฺขเว ธมฺมา เสขสฺส ภิกฺขุโน อปริหานาย สํวตฺตนฺติ.\csromanspacer{} กตเม ปญฺจ? น กมฺมารามตา น ภสฺสารามตา น นิทฺทารามตา น สงฺคณิ การามตา ยถาวิมุตฺตํ จิตฺตํ ปจฺจเวกฺขติ.\csromanspacer{} อิเม โข ภิกฺขเว ปญฺจ ธมฺมา เสขสฺส ภิกฺขุโน อปริหานาย สํวตฺตนฺตีติ.\csromanspacer{}\csromanspacer{}นวมํ.}

\csromanheader{2}{10. ทุติยเสขสุตฺต}
\proseitem{90}{ปญฺจิเม ภิกฺขเว ธมฺมา เสขสฺส ภิกฺขุโน ปริหานาย สํวตฺตนฺติ.\csromanspacer{} กตเม ปญฺจ? อิธ ภิกฺขเว เสโข ภิกฺขุ พหุกิจฺโจ โหติ พหุกรณีโย, วิยตฺโต กิํกรณีเยสุ, ริญฺจติ ปฏิสลฺลานํ, นานุยุญฺชติ อชฺฌตฺตํ เจโตสมถํ.\csromanspacer{} อยํ ภิกฺขเว ปฐโม ธมฺโม เสขสฺส ภิกฺขุโน ปริหานาย สํวตฺตติ.}

\prose{ปุน จปรํ ภิกฺขเว เสโข ภิกฺขุ อปฺปมตฺตเกน กมฺเมน ทิวสํ อตินาเมติ, ริญฺจติ ปฏิสลฺลานํ, นานุยุญฺชติ อชฺฌตฺตํ เจโตสมถํ.\csromanspacer{} อยํ ภิกฺขเว ทุติโย ธมฺโม เสขสฺส ภิกฺขุโน ปริหานาย สํวตฺตติ.}

\prose{ปุน จปรํ ภิกฺขเว เสโข ภิกฺขุ สํสฏฺโฐ วิหรติ คหฏฺฐปพฺพชิเตหิ อนนุโลมิเกน คิหิสํสคฺเคน, ริญฺจติ ปฏิสลฺลานํ, นานุยุญฺชติ อชฺฌตฺตํ เจโตสมถํ.\csromanspacer{} อยํ ภิกฺขเว ตติโย ธมฺโม เสขสฺส ภิกฺขุโน ปริหานาย สํวตฺตติ.}

\prose{ปุน จปรํ ภิกฺขเว เสโข ภิกฺขุ อกาเลน คามํ ปวิสติ, อติทิวา ปฏิกฺกมติ, ริญฺจติ ปฏิสลฺลานํ, นานุยุญฺชติ อชฺฌตฺตํ เจโตสมถํ.\csromanspacer{} อยํ ภิกฺขเว จตุตฺโถ ธมฺโม เสขสฺส ภิกฺขุโน ปริหานาย สํวตฺตติ.}

\chapterheadpageread{(9) 4. เถรวคฺค}
\prose{\csromanfolio{103}ปุน จปรํ ภิกฺขเว เสโข ภิกฺขุ ยายํ กถา อาภิสลฺเลขิกา เจโตวิวรณสปฺปายา.\csromanspacer{} เสยฺยถิทํ? อปฺปิจฺฉกถา สนฺตุฏฺฐิกถา ปวิเวกกถา อสํสคฺคกถา วีริยารมฺภกถา สีลกถา สมาธิกถา ปญฺญากถา วิมุตฺติกถา วิมุตฺติญาณทสฺสนกถา.\csromanspacer{} เอวรูปิยา กถาย น นิกามลาภี โหติ น อกิจฺฉลาภี น อกสิรลาภี\footnote{กิจฺฉลาภี กสิรลาภี (สี, สฺยา, กํ, อิ)}, ริญฺจติ ปฏิสลฺลานํ, นานุยุญฺชติ อชฺฌตฺตํ เจโตสมถํ.\csromanspacer{} อยํ ภิกฺขเว ปญฺจโม ธมฺโม เสขสฺส ภิกฺขุโน ปริหานาย สํวตฺตติ.\csromanspacer{} อิเม โข ภิกฺขเว ปญฺจ ธมฺมา เสขสฺส ภิกฺขุโน ปริหานาย สํวตฺตนฺติ.}

\prose{ปญฺจิเม ภิกฺขเว ธมฺมา เสขสฺส ภิกฺขุโน อปริหานาย สํวตฺตนฺติ.\csromanspacer{} กตเม ปญฺจ? อิธ ภิกฺขเว เสโข ภิกฺขุ น พหุกิจฺโจ โหติ น พหุกรณีโย, วิยตฺโต กิํกรณีเยสุ, น ริญฺจติ ปฏิสลฺลานํ, อนุยุญฺชติ อชฺฌตฺตํ เจโตสมถํ.\csromanspacer{} อยํ ภิกฺขเว ปฐโม ธมฺโม เสขสฺส ภิกฺขุโน อปริหานาย สํวตฺตติ.}

\prose{ปุน จปรํ ภิกฺขเว เสโข ภิกฺขุ น อปฺปมตฺตเกน กมฺเมน ทิวสํ อตินาเมติ, น ริญฺจติ ปฏิสลฺลานํ, อนุยุญฺชติ อชฺฌตฺตํ เจโตสมถํ.\csromanspacer{} อยํ ภิกฺขเว ทุติโย ธมฺโม เสขสฺส ภิกฺขุโน อปริหานาย สํวตฺตติ.}

\prose{ปุน จปรํ ภิกฺขเว เสโข ภิกฺขุ อสํสฏฺโฐ วิหรติ คหฏฺฐ ปพฺพชิเตหิ อนนุโลมิเกน คิหิสํสคฺเคน, น ริญฺจติ ปฏิสลฺลานํ, อนุยุญฺชติ อชฺฌตฺตํ เจโตสมถํ.\csromanspacer{} อยํ ภิกฺขเว ตติโย ธมฺโม เสขสฺส ภิกฺขุโน อปริหานาย สํวตฺตติ.}

\prose{ปุน จปรํ ภิกฺขเว เสโข ภิกฺขุ น อติกาเลน คามํ ปวิสติ, นาติทิวา ปฏิกฺกมติ, น ริญฺจติ ปฏิสลฺลานํ, อนุยุญฺชติ อชฺฌตฺตํ เจโตสมถํ.\csromanspacer{} อยํ ภิกฺขเว จตุตฺโถ ธมฺโม เสขสฺส ภิกฺขุโน อปริหานาย สํวตฺตติ.}

\prose{ปุน จปรํ ภิกฺขเว เสโข ภิกฺขุ ยายํ กถา อาภิสลฺเลขิกา เจโตวิวรณสปฺปายา.\csromanspacer{} เสยฺยถิทํ? อปฺปิจฺฉกถา สนฺตุฏฺฐิกถา ปวิเวกกถา อสํสคฺคกถา วีริยารมฺภกถา สีลกถา สมาธิกถา ปญฺญากถา วิมุตฺติกถา วิมุตฺติญาณทสฺสนกถา.\csromanspacer{} เอวรูปิยา}

\prose{\csromanfolio{104}กถาย นิกามลาภี โหติ อกิจฺฉลาภี อกสิรลาภี, น ริญฺจติ ปฏิสลฺลานํ, อนุยุญฺชติ อชฺฌตฺตํ เจโตสมถํ.\csromanspacer{} อยํ ภิกฺขเว ปญฺจโม ธมฺโม เสขสฺส ภิกฺขุโน อปริหานาย สํวตฺตติ.\csromanspacer{} อิเม โข ภิกฺขเว ปญฺจ ธมฺมา เสขสฺส ภิกฺขุโน อปริหานาย สํวตฺตนฺตีติ.\csromanspacer{}\csromanspacer{}ทสมํ.\csromanspacer{} เถรวคฺโค จตุตฺโถ.}

\titlehead{\textbf{ตสฺสุทฺทานํ}}
\csromangathameasure{รชนีโย วีตราโค, กุหกาสฺสทฺธ-อกฺขมา. \\ ปฏิสมฺภิทา จ สีเลน, เถโร เสขา ปเร ทุเวติ.}
\csromangathagroup{\csromangathastanza{รชนีโย วีตราโค, กุหกาสฺสทฺธ-อกฺขมา. \\ ปฏิสมฺภิทา จ สีเลน, เถโร เสขา ปเร ทุเวติ.}}

\csromanheader{1}{\textbf{(10) 5. กกุธวคฺค}}
\csromanheader{2}{1. ปฐมสมฺปทาสุตฺต}
\proseitem{91}{ปญฺจิมา ภิกฺขเว สมฺปทา.\csromanspacer{} กตมา ปญฺจ? สทฺธาสมฺปทา สีลสมฺปทา สุตสมฺปทา จาคสมฺปทา ปญฺญาสมฺปทา.\csromanspacer{} อิมา โข ภิกฺขเว ปญฺจ สมฺปทาติ.\csromanspacer{}\csromanspacer{}ปฐมํ.}

\csromanheader{2}{2. ทุติยสมฺปทาสุตฺต}
\proseitem{92}{ปญฺจิมา ภิกฺขเว สมฺปทา.\csromanspacer{} กตมา ปญฺจ? สีลสมฺปทา สมาธิสมฺปทา ปญฺญาสมฺปทา วิมุตฺติสมฺปทา วิมุตฺติญาณทสฺสนสมฺปทา.\csromanspacer{} อิมา โข ภิกฺขเว ปญฺจ สมฺปทาติ.\csromanspacer{}\csromanspacer{}ทุติยํ.}

\csromanmatikaanchor{mtk.3}
\csromantocmarkref{subsection}{3. (อญฺญา) พฺยากรณสุตฺต}{mtk.3}
\bookmark[dest={mtk.3},level=2]{3. (อญฺญา) พฺยากรณสุตฺต}
\csromanheader{2}{3. พฺยากรณสุตฺต}
\proseitem{93}{ปญฺจิมานิ ภิกฺขเว อญฺญาพฺยากรณานิ.\csromanspacer{} กตมานิ ปญฺจ? มนฺทตฺตา โมมูหตฺตา อญฺญํ พฺยากโรติ, ปาปิจฺโฉ อิจฺฉาปกโต อญฺญํ พฺยากโรติ, อุมฺมาทา จิตฺตกฺเขปา อญฺญํ พฺยากโรกิ, อธิมาเนน อญฺญํ พฺยากโรติ, สมฺมเทว อญฺญํ พฺยากโรหิ.\csromanspacer{} อิมานิ โข ภิกฺขเว ปญฺจ อญฺญาพฺยากรณานีติ.\csromanspacer{}\csromanspacer{}ตติยํ.}

\chapterheadpageread{(10) 5. กกุธวคฺค \textbf{4. ผาสุวิหารสุตฺต}}
\csromanmatikaanchor{mtk.4}
\csromantocmarkref{subsection}{4-8. ผาสุวิหารสุตฺตาทิ}{mtk.4}
\bookmark[dest={mtk.4},level=2]{4-8. ผาสุวิหารสุตฺตาทิ}
\proseitem{94}{\csromanfolio{105}ปญฺจิเม ภิกฺขเว ผาสุวิหารา.\csromanspacer{} กตเม ปญฺจ? อิธ ภิกฺขเว ภิกฺขุ วิวิจฺเจว กาเมหิ วิวิจฺจ อกุสเลหิ ธมฺเมหิ สวิตกฺกํ สวิจารํ วิเวกชํ ปีติสุขํ ปฐมํ ฌานํ อุปสมฺปชฺช วิหรติ.\csromanspacer{} วิตกฺกวิจารานํ วูปสมา ฯเปฯ ทุติยํ ฌานํ.\csromanspacer{} ตติยํ ฌานํ.\csromanspacer{} จตุตฺถํ ฌานํ อุปสมฺปชฺช วิหรติ.\csromanspacer{} อาสวานํ ขยา อนาสวํ เจโตวิมุตฺติํ ปญฺญาวิมุตฺติํ ทิฏฺเฐว ธมฺเม สยํ อภิญฺญา สจฺฉิกตฺวา อุปสมฺปชฺช วิหรติ.\csromanspacer{} อิเม โข ภิกฺขเว ปญฺจ ผาสุวิหาราติ.\csromanspacer{}\csromanspacer{}จตุตฺถํ.}

\csromanheader{2}{5. อกุปฺปสุตฺต}
\proseitem{95}{ปญฺจหิ ภิกฺขเว ธมฺเมหิ สมนฺนาคโต ภิกฺขุ นจิรสฺเสว อกุปฺปํ ปฏิวิชฺฌติ.\csromanspacer{} กตเมหิ ปญฺจหิ? อิธ ภิกฺขเว ภิกฺขุ อตฺถปฏิสมฺภิทาปตฺโต โหติ, ธมฺมปฏิสมฺภิทาปตฺโต โหติ, นิรุตฺติปฏิสมฺภิทาปตฺโต โหติ, ปฏิภานปฏิสมฺภิ ทาปตฺโต โหติ, ยถาวิมุตฺตํ จิตฺตํ ปจฺจเวกฺขติ.\csromanspacer{} อิเมหิ โข ภิกฺขเว ปญฺจหิ ธมฺเมหิ สมนฺนาคโต ภิกฺขุ นจิรสฺเสว อกุปฺปํ ปฏิวิชฺฌตีติ.\csromanspacer{}\csromanspacer{}ปญฺจมํ.}

\csromanheader{2}{6. สุตธรสุตฺต}
\proseitem{96}{ปญฺจหิ ภิกฺขเว ธมฺเมหิ สมนฺนาคโต ภิกฺขุ อานาปานสฺสติํ อาเสวนฺโต นจิรสฺเสว อกุปฺปํ ปฏิวิชฺฌติ.\csromanspacer{} กตเมหิ ปญฺจหิ? อิธ ภิกฺขเว ภิกฺขุ อปฺปฏฺโฐ โหติ อปฺปกิจฺโจ สุภโร สุสนฺโตโส ชีวิตปริกฺขาเรสุ.\csromanspacer{} อปฺปาหาโร โหติ อโนทริกตฺตํ อนุยุตฺโต.\csromanspacer{} อปฺปมิทฺโธ โหติ ชาคริยํ อนุยุตฺโต.\csromanspacer{} พหุสฺสุโต โหติ สุตธโร สุตสนฺนิจโย, เย เต ธมฺมา อาทิกลฺยาณา มชฺเฌกลฺยาณา ปริโยสานกลฺยาณา สาตฺถํ สพฺยญฺชนํ เกวลปริปุณฺณํ ปริสุทฺธํ พฺรหฺมจริยํ อภิวทนฺติ, ตถารูปาสฺส ธมฺมา พหุสฺสุตา โหนฺติ ธาตา วจสา ปริจิตา มนสานุเปกฺขิตา ทิฏฺฐิยา สุปฺปฏิวิทฺธา, ยถาวิมุตฺตํ จิตฺตํ ปจฺจเวกฺขติ.\csromanspacer{} อิเมหิ โข ภิกฺขเว ปญฺจหิ ธมฺเมหิ สมนฺนาคโต ภิกฺขุ อานาปานสฺสติํ อาเสวนฺโต นจิรสฺเสว อกุปฺปํ ปฏิวิชฺฌตีติ.\csromanspacer{}\csromanspacer{}ฉฏฺฐํ.}

\csromanheader{2}{7. กถาสุตฺต}
\proseitem{97}{\csromanfolio{106}ปญฺจหิ ภิกฺขเว ธมฺเมหิ สมนฺนาคโต ภิกฺขุ อานาปานสฺสติํ ภาเวนฺโต นจิรสฺเสว อกุปฺปํ ปฏิวิชฺฌติ.\csromanspacer{} กตเมหิ ปญฺจหิ? อิธ ภิกฺขเว ภิกฺขุ อปฺปฏฺโฐ โหติ อปฺปกิจฺโจ สุภโร สุสนฺโตโส ชีวิตปริกฺขาเรสุ.\csromanspacer{} อปฺปาหาโร โหติ อโนทริกตฺตํ อนุยุตฺโต.\csromanspacer{} อปฺปมิทฺโธ โหติ ชาคริยํ อนุยุตฺโต.\csromanspacer{} ยายํ กถา อาภิสลฺเลขิกา เจโตวิวรณสปฺปายา.\csromanspacer{} เสยฺยถิทํ? อปฺปิจฺฉกถา ฯเปฯ วิมุตฺติญาณทสฺสนกถา, เอวรูปิยา กถาย นิกามลาภี โหติ อกิจฺฉลาภี อกสิรลาภี.\csromanspacer{} ยถาวิมุตฺตํ จิตฺตํ ปจฺจเวกฺขติ.\csromanspacer{} อิเมหิ โข ภิกฺขเว ปญฺจหิ ธมฺเมหิ สมนฺนาคโต ภิกฺขุ อานาปานสฺสติํ ภาเวนฺโต นจิรสฺเสว อกุปฺปํ ปฏิวิชฺฌตีติ.\csromanspacer{}\csromanspacer{}สตฺตมํ.}

\csromanheader{2}{8. อารญฺญกสุตฺต}
\proseitem{98}{ปญฺจหิ ภิกฺขเว ธมฺเมหิ สมนฺนาคโต ภิกฺขุ อานาปานสฺสติํ พหุลีกโรนฺโต นจิรสฺเสว อกุปฺปํ ปฏิวิชฺฌติ.\csromanspacer{} กตเมหิ ปญฺจหิ? อิธ ภิกฺขเว ภิกฺขุ อปฺปฏฺโฐ โหติ อปฺปกิจฺโจ สุภโร สุสนฺโตโส ชีวิตปริกฺขาเรสุ.\csromanspacer{} อปฺปาหาโร โหติ อโนทริกตฺตํ อนุยุตฺโต.\csromanspacer{} อปฺปมิทฺโธ โหติ ชาคริยํ อนุยุตฺโต.\csromanspacer{} อารญฺญโก โหติ ปนฺตเสนาสโน, ยถาวิมุตฺตํ จิตฺตํ ปจฺจเวกฺขติ.\csromanspacer{} อิเมหิ โข ภิกฺขเว ปญฺจหิ ธมฺเมหิ สมนฺนาคโต ภิกฺขุ อานาปานสฺสติํ พหุลีกโรนฺโต นจิรสฺเสว อกุปฺปํ ปฏิวิชฺฌตีติ.\csromanspacer{}\csromanspacer{}อฏฺฐมํ.}

\csromanmatikaanchor{mtk.5}
\csromantocmarkref{subsection}{9. สีหสุตฺต}{mtk.5}
\bookmark[dest={mtk.5},level=2]{9. สีหสุตฺต}
\csromanheader{2}{9. สีหสุตฺต}
\proseitem{99}{สีโห ภิกฺขเว มิคราชา สายนฺหสมยํ อาสยา นิกฺขมติ อาสยา นิกฺขมิตฺวา วิชมฺภติ, วิชมฺภิตฺวา สมนฺตา จตุทฺทิสํ\footnote{จตุทฺทิสา (สฺยา, กํ, อิ, ก); อํ 1. 342; สํ 2. 70 ปิฏฺเฐสุปิ ปสฺสิตพฺพํ.} อนุวิโลเกติ, สมนฺตา จตุทฺทิสํ1 อนุวิโลเกตฺวา ติกฺขตฺตุํ สีหนาทํ นทติ, ติกฺขตฺตุํ สีหนาทํ นทิตฺวา โคจราย ปกฺกมติ.\csromanspacer{} โส หตฺถิสฺส เจปิ ปหารํ เทติ, สกฺกจฺจญฺเญว ปหารํ เทติ โน อสกฺกจฺจํ.\csromanspacer{} มหิํสสฺส\footnote{มหิสสฺส (สี, สฺยา, กํ, อิ)} เจปิ ปหารํ เทติ, สกฺกจฺจญฺเญว ปหารํ เทติ โน อสกฺกจฺจํ.\csromanspacer{} ควสฺส เจปิ ปหารํ เทติ, สกฺกจฺจญฺเญว ปหารํ เทติ โน อสกฺกจฺจํ.\csromanspacer{} ทีปิสฺส}

\vspace{\tipitakaparskip}
\csromancenter{(10) 5.\csromanspacer{} กกุธวคฺค เจปิ ปหารํ เทติ, สกฺกจฺจญฺเญว ปหารํ เทติ โน อสกฺกจฺจํ.\csromanspacer{} ขุทฺทกานํ เจปิ ปาณานํ ปหารํ เทติ อนฺตมโส สสพิฬารานมฺปิ\footnote{สสพิฬารานํ (ก)}, สกฺกจฺจญฺเญว ปหารํ เทติ โน อสกฺกจฺจํ.\csromanspacer{} ตํ กิสฺส เหตุ? “มา เม โยคฺคปโถ นสฺสา”ติ.}
\prose{\csromanfolio{107}“สีโห”ติ โข ภิกฺขเว ตถาคตสฺเสตํ อธิวจนํ อรหโต สมฺมาสมฺพุทฺธสฺส.\csromanspacer{} ยํ โข ภิกฺขเว ตถาคโต ปริสาย ธมฺมํ เทเสติ, อิทมสฺส โหติ สีหนาทสฺมิํ, ภิกฺขูนญฺเจปิ ภิกฺขเว ตถาคโต ธมฺมํ เทเสติ, สกฺกจฺจญฺเญว ตถาคโต ธมฺมํ เทเสติ โน อสกฺกจฺจํ.\csromanspacer{} ภิกฺขุนีนญฺเจปิ ภิกฺขเว ตถาคโต ธมฺมํ เทเสติ, สกฺกจฺจญฺเญว ตถาคโต ธมฺมํ เทเสติ โน อสกฺกจฺจํ.\csromanspacer{} อุปาสกานญฺเจปิ ภิกฺขเว ตถาคโต ธมฺมํ เทเสติ, สกฺกจฺจญฺเญว ตถาคโต ธมฺมํ เทเสติ โน อสกฺกจฺจํ.\csromanspacer{} อุปาสิกานญฺเจปิ ภิกฺขเว ตถาคโต ธมฺมํ เทเสติ, สกฺกจฺจญฺเญว ตถาคโต ธมฺมํ เทเสติ โน อสกฺกจฺจํ.\csromanspacer{} ปุถุชฺชนานญฺเจปิ ภิกฺขเว ตถาคโต ธมฺมํ เทเสติ อนฺตมโส อนฺนภารเนสาทานมฺปิ\footnote{อนฺนภารเนสาทานํ (ก)}, สกฺกจฺจญฺเญว ตถาคโต ธมฺมํ เทเสติ โน อสกฺกจฺจํ.\csromanspacer{} ตํ กิสฺส เหตุ? ธมฺมครุ ภิกฺขเว ตถาคโต ธมฺมคารโวติ.\csromanspacer{}\csromanspacer{}นวมํ.}

\csromanmatikaanchor{mtk.6}
\csromantocmarkref{subsection}{10. กกุธเถรสุตฺต}{mtk.6}
\bookmark[dest={mtk.6},level=2]{10. กกุธเถรสุตฺต}
\csromanheader{2}{10. กกุธเถรสุตฺต}
\proseitem{100}{เอวํ เม สุตํ\csromandash{}เอกํ สมยํ ภควา โกสมฺพิยํ วิหรติ โฆสิตาราเม. * เตน โข ปน สมเยน กกุโธ นาม โกลิยปุตฺโต\footnote{โกฬิยปุตฺโต (สี, สฺยา, ก)} อายสฺมโต มหาโมคฺคลฺลานสฺส อุปฏฺฐาโก อธุนากาลงฺกโต อญฺญตรํ มโนมยํ กายํ อุปปนฺโน, ตสฺส เอวรูโป อตฺตภาวปฏิลาโภ โหติ, เสยฺยถาปิ นาม เทฺว วา ตีณิ วา มาคธกานิ\footnote{มาคธิกานิ (สี, อิ, ก)} คามกฺเขตฺตานิ.\csromanspacer{} โส เตน อตฺตภาวปฏิลาเภน เนว อตฺตานํ\footnote{เนวตฺตานํ พฺยาพาเธติ (สี)} โน ปรํ พฺยาพาเธติ.}

\prose{อถ โข กกุโธ เทวปุตฺโต เยนายสฺมา มหาโมคฺคลฺลาโน เตนุปสงฺกมิ, อุปสงฺกมิตฺวา อายสฺมนฺตํ มหาโมคฺคลฺลานํ อภิวาเทตฺวา เอกมนฺตํ อฏฺฐาสิ, เอกมนฺตํ ฐิโต โข กกุโธ เทวปุตฺโต อายสฺมนฺตํ \csromanfolio{108}มหาโมคฺคลฺลานํ เอตทโวจ “เทวทตฺตสฺส ภนฺเต เอวรูปํ อิจฺฉาคตํ อุปฺปชฺชิ ‘อหํ ภิกฺขุสํฆํ ปริหริสฺสามี’ติ.\csromanspacer{} สหจิตฺตุปฺปาทา จ ภนฺเต เทวทตฺโต ตสฺสา อิทฺธิยา ปริหีโน”ติ.\csromanspacer{} อิทมโวจ กกุโธ เทวปุตฺโต, อิทํ วตฺวา อายสฺมนฺตํ มหาโมคฺคลฺลานํ อภิวาเทตฺวา ปทกฺขิณํ กตฺวา ตตฺเถวนฺตรธายิ.}

\prose{อถ โข อายสฺมา มหาโมคฺคลฺลาโน เยน ภควา เตนุปสงฺกมิ, อุปสงฺกมิตฺวา ภควนฺตํ อภิวาเทตฺวา เอกมนฺตํ นิสีทิ, เอกมนฺตํ นิสินฺโน โข อายสฺมา มหาโมคฺคลฺลาโน ภควนฺตํ เอตทโวจ\csromandash{}}

\prose{กกุโธ นาม ภนฺเต โกลิยปุตฺโต มมํ อุปฏฺฐาโก อธุนากาลงฺกโต อญฺญตรํ มโนมยํ กายํ อุปปนฺโน โหติ, ตสฺส เอวรูโป อตฺตภาวปฏิลาโภ.\csromanspacer{} เสยฺยถาปิ นาม เทฺว วา ตีณิ วา มาคธกานิ คามกฺเขตฺตานิ, โส เตน อตฺตภาวปฏิลาเภน เนว อตฺตานํ โน ปรํ พฺยาพาเธติ.\csromanspacer{} อถ โข ภนฺเต กกุโธ เทวปุตฺโต เยนาหํ เตนุปสงฺกมิ, อุปสงฺกมิตฺวา มํ อภิวาเทตฺวา เอกมนฺตํ อฏฺฐาสิ, เอกมนฺตํ ฐิโต โข ภนฺเต กกุโธ เทวปุตฺโต มํ เอตทโวจ “เทวทตฺตสฺส ภนฺเต เอวรูปํ อิจฺฉาคตํ อุปฺปชฺชิ ‘อหํ ภิกฺขุสํฆํ ปริหริสฺสามี’ติ.\csromanspacer{} สหจิตฺตุปฺปาทา จ ภนฺเต เทวทตฺโต ตสฺสา อิทฺธิยา ปริหีโน”ติ.\csromanspacer{} อิทมโวจ ภนฺเต กกุโธ เทวปุตฺโต, อิทํ วตฺวา มํ อภิวาเทตฺวา ปทกฺขิณํ กตฺวา ตตฺเถวนฺตรธายีติ.}

\prose{กิํ ปน เต โมคฺคลฺลาน กกุโธ เทวปุตฺโต เจตสา เจโต ปริจฺจ วิทิโต “ยํ กิญฺจิ กกุโธ เทวปุตฺโต ภาสติ, สพฺพํ ตํ ตเถว โหติ โน อญฺญถา”ติ.\csromanspacer{} เจตสา เจโต ปริจฺจ วิทิโต เม ภนฺเต กกุโธ เทวปุตฺโต “ยํ กิญฺจิ กกุโธ เทวปุตฺโต ภาสติ, สพฺพํ ตํ ตเถว โหติ โน อญฺญถา”ติ.\csromanspacer{} รกฺขสฺเสตํ โมคฺคลฺลานํ วาจํ (รกฺขสฺเสตํ โมคฺคลฺลาน วาจํ)\footnote{( ) สี-สฺยา-กํ-อิ-โปตฺถเกสุ นตฺถิ; วิ 4. 343 ปิฏฺเฐ ปน สพฺพตฺถปิ ทิสฺสติเยว.}.\csromanspacer{} อิทานิ โส โมฆปุริโส อตฺตนาว อตฺตานํ ปาตุกริสฺสติ.}

\prose{ปญฺจิเม โมคฺคลฺลาน สตฺถาโร สนฺโต สํวิชฺชมานา โลกสฺมิํ.\csromanspacer{} กตเม ปญฺจ? อิธ โมคฺคลฺลาน เอกจฺโจ สตฺถา อปริสุทฺธสีโล สมาโน “ปริสุทฺธสีโลมฺหี”ติ ปฏิชานาติ “ปริสุทฺธํ เม สีลํ}

\vspace{\tipitakaparskip}
\csromancenter{(10) 5.\csromanspacer{} กกุธวคฺค ปริโยทาตํ อสํกิลิฏฺฐนฺ”ติ.\csromanspacer{} ตเมนํ สาวกา เอวํ ชานนฺติ “อยํ โข ภวํ สตฺถา อปริสุทฺธสีโล สมาโน ‘ปริสุทฺธสีโลมฺหี’ติ ปฏิชานาติ ‘ปริสุทฺธํ เม สีลํ ปริโยทาตํ อสํกิลิฏฺฐนฺ’ติ.\csromanspacer{} มยญฺเจว โข ปน คิหีนํ อาโรเจยฺยาม, นาสฺสสฺส มนาปํ.\csromanspacer{} ยํ โข ปนสฺส อมนาปํ, กถํ นํ\footnote{กถํ นุ ตํ (สี), กถํ นุ (สฺยา, กํ, อิ, ก), กถํ ตํ (กตฺถจิ)} มยํ เตน สมุทาจเรยฺยาม.\csromanspacer{} สมฺมนฺนติ โข ปน จีวรปิณฺฑปาตเสนาสนคิลานปฺปจฺจย เภสชฺชปริกฺขาเรน.\csromanspacer{} ยํ ตุโม กริสฺสติ, ตุโมว เตน ปญฺญายิสฺสตี”ติ.\csromanspacer{} เอวรูปํ โข โมคฺคลฺลาน สตฺถารํ สาวกา สีลโต รกฺขนฺติ, เอวรูโป จ ปน สตฺถา สาวเกหิ สีลโต รกฺขํ ปจฺจาสีสติ\footnote{ปจฺจาสิํสติ (สี, สฺยา, กํ, อิ)}.}
\prose{\csromanfolio{109}ปุน จปรํ โมคฺคลฺลาน อิเธกจฺโจ สตฺถา อปริสุทฺธาชีโว สมาโน “ปริสุทฺธาชีโวมฺหี”ติ ปฏิชานาติ “ปริสุทฺโธ เม อาชีโว ปริโยทาโต อสํกิลิฏฺโฐ”ติ.\csromanspacer{} ตเมนํ สาวกา เอวํ ชานนฺติ “อยํ โข ภวํ สตฺถา อปริสุทฺธาชีโว สมาโน ‘ปริสุทฺธาชีโวมฺหี’ติ ปฏิชานาติ ‘ปริสุทฺโธ เม อาชีโว ปริโยทาโต อสํกิลิฏฺโฐ’ติ.\csromanspacer{} มยญฺเจว โข ปน คิหีนํ อาโรเจยฺยาม, นาสฺสสฺส มนาปํ.\csromanspacer{} ยํ โข ปนสฺส อมนาปํ, กถํ นํ มยํ เตน สมุทาจเรยฺยาม, สมฺมนฺนติ โข ปน จีวรปิณฺฑปาต เสนาสน คิลานปฺปจฺจย เภสชฺช ปริกฺขาเรน.\csromanspacer{} ยํ ตุโม กริสฺสติ, ตุโมว เตน ปญฺญายิสฺสตี”ติ.\csromanspacer{} เอวรูปํ โข โมคฺคลฺลาน สตฺถารํ สาวกา อาชีวโต รกฺขนฺติ, เอวรูโป จ ปน สตฺถา สาวเกหิ อาชีวโต รกฺขํ ปจฺจาสีสติ.}

\prose{ปุน จปรํ โมคฺคลฺลาน อิเธกจฺโจ สตฺถา อปริสุทฺธธมฺมเทสโน สมาโน “ปริสุทฺธธมฺมเทสโนมฺหี”ติ ปฏิชานาติ “ปริสุทฺธา เม ธมฺมเทสนา ปริโยทาตา อสํกิลิฏฺฐา”ติ.\csromanspacer{} ตเมนํ สาวกา เอวํ ชานนฺติ “อยํ โข ภวํ สตฺถา อปริสุทฺธ ธมฺมเทสโน สมาโน ‘ปริสุทฺธธมฺมเทสโนมฺหี’ติ ปฏิชานาติ ‘ปริสุทฺธา เม ธมฺมเทสนา ปริโยทาตา อสํกิลิฏฺฐา’ติ.\csromanspacer{} มยญฺเจว โข ปน คิหีนํ อาโรเจยฺยาม, นาสฺสสฺส มนาปํ.\csromanspacer{} ยํ โข ปนสฺส อมนาปํ, กถํ นํ มยํ เตน สมุทาจเรยฺยาม.\csromanspacer{} สมฺมนฺนติ โข ปน จีวร ปิณฺฑปาต เสนาสน คิลานปฺปจฺจยเภสชฺช ปริกฺขาเรน.\csromanspacer{} ยํ ตุโม กริสฺสติ, ตุโมว เตน ปญฺญายิสฺสตี”ติ.\csromanspacer{} เอวรูปํ โข โมคฺคลฺลาน สตฺถารํ สาวกา ธมฺมเทสนโต \csromanfolio{110}รกฺขนฺติ, เอวรูโป จ ปน สตฺถา สาวเกหิ ธมฺมเทสนโต รกฺขํ ปจฺจาสีสติ.}

\prose{ปุน จปรํ โมคฺคลฺลาน อิเธกจฺโจ สตฺถา อปริสุทฺธเวยฺยากรโณ สมาโน “ปริสุทฺธเวยฺยากรโณมฺหี”ติ ปฏิชานาติ “ปริสุทฺธํ เม เวยฺยากรณํ ปริโยทาตํ อสํกิลิฏฺฐนฺ”ติ.\csromanspacer{} ตเมนํ สาวกา เอวํ ชานนฺติ “อยํ โข ภวํ สตฺถา อปริสุทฺธเวยฺยากรโณ สมาโน ‘ปริสุทฺธเวยฺยากรโณมฺหี’ติ ปฏิชานาติ ‘ปริสุทฺธํ เม เวยฺยากรณํ ปริโยทาตํ อสํกิลิฏฺฐนฺ’ติ.\csromanspacer{} มยญฺเจว โข ปน คิหีนํ อาโรเจยฺยาม, นาสฺสสฺส มนาปํ.\csromanspacer{} ยํ โข ปนสฺส อมนาปํ, กถํ นํ มยํ เตน สมุทาจเรยฺยาม.\csromanspacer{} สมฺมนฺนติ โข ปน จีวร ปิณฺฑปาต เสนาสน คิลานปฺปจฺจย เภสชฺช ปริกฺขาเรน.\csromanspacer{} ยํ ตุโม กริสฺสติ, ตุโมว เตน ปญฺญายิสฺสตี”ติ.\csromanspacer{} เอวรูปํ โข โมคฺคลฺลานํ สตฺถารํ สาวกา เวยฺยากรณโต รกฺขนฺติ, เอวรูโป จ ปน สตฺถา สาวเกหิ เวยฺยากรณโต รกฺขํ ปจฺจาสีสติ.}

\prose{ปุน จปรํ โมคฺคลฺลาน อิเธกจฺโจ สตฺถา อปริสุทฺธญาณทสฺสโน สมาโน “ปริสุทฺธญาณทสฺสโนมฺหี”ติ ปฏิชานาติ “ปริสุทฺธํ เม ญาณทสฺสนํ ปริโยทาตํ อสํกิลิฏฺฐนฺ”ติ.\csromanspacer{} ตเมนํ สาวกา เอวํ ชานนฺติ “อยํ โข ภวํ สตฺถา อปริ สุทฺธญาณทสฺสโน สมาโน ‘ปริสุทฺธญาณทสฺสโนมฺหี’ติ ปฏิชานาติ ‘ปริสุทฺธํ เม ญาณทสฺสนํ ปริโยทาตํ อสํกิลิฏฺฐนฺ’ติ.\csromanspacer{} มยญฺเจว โข ปน คิหีนํ อาโรเจยฺยาม, นาสฺสสฺส มนาปํ.\csromanspacer{} ยํ โข ปนสฺส อมนาปํ, กถํ นํ มยํ เตน สมุทาจเรยฺยาม.\csromanspacer{} สมฺมนฺนติ โข ปน จีวร ปิณฺฑปาต เสนาสนคิลาน ปฺปจฺจย เภสชฺช ปริกฺขาเรน.\csromanspacer{} ยํ ตุโม กริสฺสติ, ตุโมว เตน ปญฺญายิสฺสตี”ติ.\csromanspacer{} เอวรูปํ โข โมคฺคลฺลาน สตฺถารํ สาวกา ญาณทสฺสนโต รกฺขนฺติ, เอวรูโป จ ปน สตฺถา สาวเกหิ ญาณทสฺสนโต รกฺขํ ปจฺจาสีสติ.\csromanspacer{} อิเม โข โมคฺคลฺลาน ปญฺจ สตฺถาโร สนฺโต สํวิชฺชมานา โลกสฺมิํ.}

\prose{อหํ โข ปน โมคฺคลฺลาน ปริสุทฺธสีโล สมาโน “ปริสุทฺธสีโลมฺหี”ติ ปฏิชานามิ “ปริสุทฺธํ เม สีลํ ปริโยทาตํ อสํกิลิฏฺฐนฺ”ติ.\csromanspacer{} น จ มํ สาวกา สีลโต รกฺขนฺติ.\csromanspacer{} น จาหํ สาวเกหิ สีลโต รกฺขํ ปจฺจาสีสามิ.\csromanspacer{} ปริสุทฺธาชีโว สมาโน “ปริสุทฺธาชีโวมฺหี”ติ ปฏิชานามิ “ปริสุทฺโธ เม อาชีโว ปริโยทาโต อสํกิลิฏฺโฐ”ติ.\csromanspacer{} น จ มํ}

\vspace{\tipitakaparskip}
\csromancenter{(10) 5.\csromanspacer{} กกุธวคฺค สาวกา อาชีวโต รกฺขนฺติ.\csromanspacer{} น จาหํ สาวเกหิ อาชีวโต รกฺขํ ปจฺจาสีสามิ.\csromanspacer{} ปริสุทฺธธมฺมเทสโน สมาโน “ปริสุทฺธธมฺมเทสโนมฺหี”ติ ปฏิชานามิ “ปริสุทฺธา เม ธมฺมเทสนา ปริโยทาตา อสํกิลิฏฺฐา”ติ.\csromanspacer{} น จ มํ สาวกา ธมฺมเทสนโต รกฺขนฺติ.\csromanspacer{} น จาหํ สาวเกหิ ธมฺมเทสนโต รกฺขํ ปจฺจาสีสามิ.\csromanspacer{} ปริสุทฺธเวยฺยากรโณ สมาโน “ปริสุทฺธเวยฺยากรโณมฺหี”ติ ปฏิชานามิ “ปริสุทฺธํ เม เวยฺยากรณํ ปริโยทาตํ อสํกิลิฏฺฐนฺ”ติ.\csromanspacer{} น จ มํ สาวกา เวยฺยากรณโต รกฺขนฺติ.\csromanspacer{} น จาหํ สาวเกหิ เวยฺยากรณโต รกฺขํ ปจฺจาสีสามิ.\csromanspacer{} ปริสุทฺธญาณทสฺสโน สมาโน “ปริสุทฺธญาณทสฺสโนมฺหี”ติ ปฏิชานามิ “ปริสุทฺธํ เม ญาณทสฺสนํ ปริโยทาตํ อสํกิลิฏฺฐนฺ”ติ.\csromanspacer{} น จ มํ สาวกา ญาณทสฺสนโต รกฺขนฺติ.\csromanspacer{} น จาหํ สาวเกหิ ญาณทสฺสนโต รกฺขํ ปจฺจาสีสามีติ.\csromanspacer{}\csromanspacer{}ทสมํ.\csromanspacer{} กกุธวคฺโค ปญฺจโม.}
\titlehead{\textbf{ตสฺสุทฺทานํ}}
\csromangathameasure{เทฺว สมฺปทา พฺยากรณํ, ผาสุ อกุปฺปปญฺจมํ. \\ สุตํ กถา อารญฺญโก, สีโห จ กกุโธ ทสาติ. \\ \textbf{ทุติโย ปณฺณาสโก สมตฺโต.}}
\csromangathagroup{\csromangathastanza{\csromanfolio{111}เทฺว สมฺปทา พฺยากรณํ, ผาสุ อกุปฺปปญฺจมํ. \\ สุตํ กถา อารญฺญโก, สีโห จ กกุโธ ทสาติ. \\ \textbf{ทุติโย ปณฺณาสโก สมตฺโต.}}}

\csromanmatikaanchor{mtk.7}
\csromantocmarknopage{section}{(3) ตติยปณฺณาสก}{mtk.7}
\bookmark[dest={mtk.7},level=1]{(3) ตติยปณฺณาสก}
\csromanheader{1}{3. ตติยปณฺณาสก}
\csromanmatikaanchor{mtk.8}
\csromantocmarknopage{section}{(11) 1. ผาสุวิหารวคฺค}{mtk.8}
\bookmark[dest={mtk.8},level=1]{(11) 1. ผาสุวิหารวคฺค}
\csromanheader{1}{\textbf{(11) 1. ผาสุวิหารวคฺค}}
\csromanmatikaanchor{mtk.9}
\csromantocmarkref{subsection}{1. สารชฺชสุตฺต}{mtk.9}
\bookmark[dest={mtk.9},level=2]{1. สารชฺชสุตฺต}
\csromanheader{2}{1. สารชฺชสุตฺต}
\proseitem{101}{\csromanfolio{112}ปญฺจิเม ภิกฺขเว เสขเวสารชฺชกรณา ธมฺมา.\csromanspacer{} กตเม ปญฺจ? อิธ ภิกฺขเว ภิกฺขุ สทฺโธ โหติ, สีลวา โหติ, พหุสฺสุโต โหติ, อารทฺธวีริโย โหติ, ปญฺญวา โหติ.}

\prose{ยํ ภิกฺขเว อสฺสทฺธสฺส สารชฺชํ โหติ, สทฺธสฺส ตํ สารชฺชํ น โหติ, ตสฺมายํ ธมฺโม เสขเวสารชฺชกรโณ.}

\prose{ยํ ภิกฺขเว ทุสฺสีลสฺส สารชฺชํ โหติ, สีลวโต ตํ สารชฺชํ น โหติ, ตสฺมายํ ธมฺโม เสขเวสารชฺชกรโณ.}

\prose{ยํ ภิกฺขเว อปฺปสฺสุตสฺส สารชฺชํ โหติ, พหุสฺสุตสฺส ตํ สารชฺชํ น โหติ, ตสฺมายํ ธมฺโม เสขเวสารชฺชกรโณ.}

\prose{ยํ ภิกฺขเว กุสีตสฺส สารชฺชํ โหติ, อารทฺธวีริยสฺส ตํ สารชฺชํ น โหติ, ตสฺมายํ ธมฺโม เสขเวสารชฺชกรโณ.}

\prose{ยํ ภิกฺขเว ทุปฺปญฺญสฺส สารชฺชํ โหติ, ปญฺญวโต ตํ สารชฺชํ น โหติ, ตสฺมายํ ธมฺโม เสขเวสารชฺชกรโณ.\csromanspacer{} อิเม โข ภิกฺขเว ปญฺจ เสขเวสารชฺชกรณา ธมฺมาติ.\csromanspacer{}\csromanspacer{}ปฐมํ.}

\csromanmatikaanchor{mtk.10}
\csromantocmarkref{subsection}{2. อุสฺสงฺกิตสุตฺต}{mtk.10}
\bookmark[dest={mtk.10},level=2]{2. อุสฺสงฺกิตสุตฺต}
\csromanheader{2}{2. อุสฺสงฺกิตสุตฺต}
\proseitem{102}{ปญฺจหิ ภิกฺขเว ธมฺเมหิ สมนฺนาคโต ภิกฺขุ อุสฺสงฺกิตปริสงฺกิโต โหติ “ปาปภิกฺขู”ติ อปิ อกุปฺปธมฺโมปิ\footnote{อปิ อกุปฺปธมฺโม (สี, สฺยา, กํ)}.}

\prose{กตเมหิ ปญฺจหิ? อิธ ภิกฺขเว ภิกฺขุ เวสิยาโคจโร วา โหติ, วิธวาโคจโร วา โหติ, ถุลฺลกุมาริกาโคจโร วา โหติ, ปณฺฑกโคจโร วา โหติ, ภิกฺขุนีโคจโร วา โหติ.}

\chapterheadpageread{(11) 1. ผาสุวิหารวคฺค}
\prose{\csromanfolio{113}อิเมหิ โข ภิกฺขเว ปญฺจหิ ธมฺเมหิ สมนฺนาคโต ภิกฺขุ อุสฺสงฺกิตปริสงฺกิโต โหติ “ปาปภิกฺขู”ติ อปิ อกุปฺปธมฺโมปีติ.\csromanspacer{}\csromanspacer{}ทุติยํ.}

\csromanmatikaanchor{mtk.11}
\csromantocmarkref{subsection}{3. มหาโจรสุตฺต}{mtk.11}
\bookmark[dest={mtk.11},level=2]{3. มหาโจรสุตฺต}
\csromanheader{2}{3. มหาโจรสุตฺต}
\proseitem{103}{ปญฺจหิ ภิกฺขเว องฺเคหิ สมนฺนาคโต มหาโจโร สนฺธิมฺปิ ฉินฺทติ, นิลฺโลปมฺปิ หรติ, เอกาคาริกมฺปิ กโรติ, ปริปนฺเถปิ ติฏฺฐติ.\csromanspacer{} กตเมหิ ปญฺจหิ? อิธ ภิกฺขเว มหาโจโร วิสมนิสฺสิโต จ โหติ คหนนิสฺสิโต จ พลวนิสฺสิโต จ โภคจาคี จ เอกจารี จ.}

\prose{กถญฺจ ภิกฺขเว มหาโจโร วิสมนิสฺสิโต โหติ? อิธ ภิกฺขเว มหาโจโร นทีวิทุคฺคํ วา นิสฺสิโต โหติ ปพฺพตวิสมํ วา.\csromanspacer{} เอวํ โข ภิกฺขเว มหาโจโร วิสมนิสฺสิโต โหติ.}

\prose{กถญฺจ ภิกฺขเว มหาโจโร คหนนิสฺสิโต โหติ? อิธ ภิกฺขเว มหาโจโร ติณคหนํ วา นิสฺสิโต โหติ รุกฺขคหนํ วา โรธํ\footnote{เคธํ (สี) อํ 1. 152 ปิฏฺเฐปิ.} วา มหาวนสณฺฑํ วา.\csromanspacer{} เอวํ โข ภิกฺขเว มหาโจโร คหนนิสฺสิโต โหติ.}

\prose{กถญฺจ ภิกฺขเว มหาโจโร พลวนิสฺสิโต โหติ? อิธ ภิกฺขเว มหาโจโร ราชานํ วา ราชมหามตฺตานํ วา นิสฺสิโต โหติ, ตสฺส เอวํ โหติ “สเจ มํ โกจิ กิญฺจิ วกฺขติ, อิเม เม ราชาโน วา ราชมหามตฺตา วา ปริโยธาย อตฺถํ ภณิสฺสนฺตี”ติ.\csromanspacer{} สเจ นํ โกจิ กิญฺจิ อาห, ตฺยสฺส ราชาโน วา ราชมหามตฺตา วา ปริโยธาย อตฺถํ ภณนฺติ.\csromanspacer{} เอวํ โข ภิกฺขเว มหาโจโร พลวนิสฺสิโต โหติ.}

\prose{กถญฺจ ภิกฺขเว มหาโจโร โภคจาคี โหติ? อิธ ภิกฺขเว มหาโจโร อฑฺโฒ โหติ มหทฺธโน มหาโภโค, ตสฺส เอวํ โหติ “สเจ มํ โกจิ กิญฺจิ วกฺขติ, อิโต โภเคน ปฏิสนฺถริสฺสามี”ติ.\csromanspacer{} สเจ นํ โกจิ กิญฺจิ อาห, ตโต โภเคน ปฏิสนฺถรติ.\csromanspacer{} เอวํ โข ภิกฺขเว มหาโจโร โภคจาคี}

\prose{กถญฺจ ภิกฺขเว มหาโจโร เอกจารี โหติ? อิธ ภิกฺขเว มหาโจโร เอกโกว คหณานิ\footnote{นิคฺคหณานิ (สี, สฺยา, กํ, อิ)} กตฺตา โหติ, ตํ กิสฺส เหตุ? “มา \csromanfolio{114}เม คุยฺหมนฺตา พหิทฺธา สมฺเภทํ อคมํสู”ติ.\csromanspacer{} เอวํ โข ภิกฺขเว มหาโจโร เอกจารี โหติ.}

\prose{อิเมหิ โข ภิกฺขเว ปญฺจหงฺเคหิ สมนฺนาคโต มหาโจโร สนฺธิมฺปิ ฉินฺทติ, นิลฺโลปมฺปิ หรติ, เอกาคาริกมฺปิ กโรติ, ปริปนฺเถปิ ติฏฺฐติ.}

\prose{เอวเมวํ โข ภิกฺขเว ปญฺจหิ ธมฺเมหิ สมนฺนาคโต ปาปภิกฺขุ ขตํ อุปหตํ อตฺตานํ ปริหรติ, สาวชฺโช จ โหติ สานุวชฺโช วิญฺญูนํ, พหุญฺจ อปุญฺญํ ปสวติ.\csromanspacer{} กตเมหิ ปญฺจหิ, อิธ ภิกฺขเว ปาปภิกฺขุ วิสมนิสฺสิโต จ โหติ คหนนิสฺสิโต จ พลวนิสฺสิโต จ โภคจาคี จ เอกจารี จ.}

\prose{กถญฺจ ภิกฺขเว ปาปภิกฺขุ วิสมนิสฺสิโต โหติ? อิธ ภิกฺขเว ปาปภิกฺขุ วิสเมน กายกมฺเมน สมนฺนาคโต โหติ, วิสเมน วจีกมฺเมน สมนฺนาคโต โหติ, วิสเมน มโนกมฺเมน สมนฺนาคโต โหติ.\csromanspacer{} เอวํ โข ภิกฺขเว ปาปภิกฺขุ วิสมนิสฺสิโต โหติ.}

\prose{กถญฺจ ภิกฺขเว ปาปภิกฺขุ คหนนิสฺสิโต โหติ? อิธ ภิกฺขเว ปาปภิกฺขุ มิจฺฉาทิฏฺฐิโก โหติ อนฺตคฺคาหิกาย ทิฏฺฐิยา สมนฺนาคโต.\csromanspacer{} เอวํ โข ภิกฺขเว ปาปภิกฺขุ คหนนิสฺสิโต โหติ.}

\prose{กถญฺจ ภิกฺขเว ปาปภิกฺขุ พลวนิสฺสิโต โหติ? อิธ ภิกฺขเว ปาปภิกฺขุ ราชานํ วา ราชมหามตฺตานํ วา นิสฺสิโต โหติ, ตสฺส เอวํ โหติ “สเจ มํ โกจิ กิญฺจิ วกฺขติ, อิเม เม ราชาโน วา ราชมหามตฺตา วา ปริโยธาย อตฺถํ ภณิสฺสนฺตี”ติ.\csromanspacer{} สเจ นํ โกจิ กิญฺจิ อาห, ตฺยสฺส ราชาโน วา ราชมหามตฺตา วา ปริโยธาย อตฺถํ ภณนฺติ.\csromanspacer{} เอวํ โข ภิกฺขเว ปาปภิกฺขุ พลวนิสฺสิโต โหติ.}

\prose{กถญฺจ ภิกฺขเว ปาปภิกฺขุ โภคจาคี โหติ? อิธ ภิกฺขเว ปาปภิกฺขุ ลาภี โหติ จีวรปิณฺฑปาต เสนาสนคิลานปฺปจฺจย เภสชฺช ปริกฺขารานํ.\csromanspacer{} ตสฺส เอวํ โหติ “สเจ มํ โกจิ กิญฺจิ วกฺขติ, อิโต ลาเภน ปฏิสนฺถริสฺสามี”ติ.\csromanspacer{} สเจ นํ โกจิ กิญฺจิ อาห, ตโต ลาเภน ปฏิสนฺถรติ.\csromanspacer{} เอวํ โข ภิกฺขเว ปาปภิกฺขุ โภคจาคี โหติ.}

\chapterheadpageread{(11) 1. ผาสุวิหารวคฺค}
\prose{\csromanfolio{115}กถญฺจ ภิกฺขเว ปาปภิกฺขุ เอกจารี โหติ? อิธ ภิกฺขเว ปาปภิกฺขุ เอกโกว ปจฺจนฺติเมสุ ชนปเทสุ นิวาสํ กปฺเปติ, โส ตตฺถ กุลานิ อุปสงฺกมนฺโต ลาภํ ลภติ.\csromanspacer{} เอวํ โข ภิกฺขเว ปาปภิกฺขุ เอกจารี โหติ.}

\prose{อิเมหิ โข ภิกฺขเว ปญฺจหิ ธมฺเมหิ สมนฺนาคโต ปาปภิกฺขุ ขตํ อุปหตํ อตฺตานํ ปริหรติ, สาวชฺโช จ โหติ สานุวชฺโช วิญฺญูนํ, พหุญฺจ อปุญฺญํ ปสวตีติ.\csromanspacer{}\csromanspacer{}ตติยํ.}

\csromanmatikaanchor{mtk.12}
\csromantocmarkref{subsection}{4. สมณสุขุมาลสุตฺต}{mtk.12}
\bookmark[dest={mtk.12},level=2]{4. สมณสุขุมาลสุตฺต}
\csromanheader{2}{4. สมณสุขุมาลสุตฺต}
\proseitem{104}{ปญฺจหิ ภิกฺขเว ธมฺเมหิ สมนฺนาคโต ภิกฺขุ สมเณสุ สมณสุขุมาโล โหติ.}

\prose{กตเมหิ ปญฺจหิ? อิธ ภิกฺขเว ภิกฺขุ ยาจิโตว พหุลํ จีวรํ ปริภุญฺชติ อปฺปํ อยาจิโต, ยาจิโตว พหุลํ ปิณฺฑปาตํ ปริภุญฺชติ อปฺปํ อยาจิโต, ยาจิโตว พหุลํ เสนาสนํ ปริภุญฺชติ อปฺปํ อยาจิโต, ยาจิโตว พหุลํ คิลานปฺปจฺจยเภสชฺชปริกฺขารํ ปริภุญฺชติ อปฺปํ อยาจิโต.\csromanspacer{} เยหิ โข ปน สพฺรหฺมจารีหิ สทฺธิํ วิหรติ, ตฺยสฺส\footnote{ตฺยาสฺส (ก) อํ 1. 399 ปิฏฺเฐปิ.} มนาเปเนว พหุลํ กายกมฺเมน สมุทาจรนฺติ อปฺปํ อมนาเปน, มนาเปเนว พหุลํ วจีกมฺเมน สมุทาจรนฺติ อปฺปํ อมนาเปน, มนาเปเนว พหุลํ มโนกมฺเมน สมุทาจรนฺติ อปฺปํ อมนาเปน, มนาปํเยว อุปหารํ อุปหรนฺติ อปฺปํ อมนาปํ.\csromanspacer{} ยานิ โข ปน ตานิ เวทยิตานิ ปิตฺตสมุฏฺฐานานิ วา เสมฺหสมุฏฺฐานานิ วา วาตสมุฏฺฐานานิ วา สนฺนิปาติกานิ วา อุตุปริณามชานิ วา วิสมปริหารชานิ วา โอปกฺกมิกานิ วา กมฺมวิปากชานิ วา, ตานิสฺส น พหุเทว อุปฺปชฺชนฺติ, อปฺปาพาโธ โหติ.\csromanspacer{} จตุนฺนํ ฌานานํ อาภิเจตสิกานํ ทิฏฺฐธมฺมสุขวิหารานํ นิกามลาภี โหติ อกิจฺฉลาภี อกสิรลาภี.\csromanspacer{} อาสวานํ ขยา อนาสวํ เจโตวิมุตฺติํ ปญฺญาวิมุตฺติํ ทิฏฺเฐว ธมฺเม สยํ อภิญฺญา สจฺฉิกตฺวา อุปสมฺปชฺช วิหรติ.\csromanspacer{} อิเมหิ โข ภิกฺขเว ปญฺจหิ ธมฺเมหิ สมนฺนาคโต ภิกฺขุ สมเณสุ สมณสุขุมาโล โหติ.}

\prose{ยํ หิ ตํ ภิกฺขเว สมฺมา วทมาโน วเทยฺย “สมเณสุ สมณ สุขุมาโล”ติ, มเมว ตํ ภิกฺขเว สมฺมา\footnote{มเมว ตํ สมฺมา (?)} วทมาโน วเทยฺย “สมเณสุ \csromanfolio{116}สมณสุขุมาโล”ติ.\csromanspacer{} อหํ หิ ภิกฺขเว ยาจิโตว พหุลํ จีวรํ ปริภุญฺชามิ อปฺปํ อยาจิโต, ยาจิโตว พหุลํ ปิณฺฑปาตํ ปริภุญฺชามิ อปฺปํ อยาจิโต, ยาจิโตว พหุลํ เสนาสนํ ปริภุญฺชามิ อปฺปํ อยาจิโต, ยาจิโตว พหุลํ คิลานปฺปจฺจย เภสชฺช ปริกฺขารํ ปริภุญฺชามิ อปฺปํ อยาจิโต.\csromanspacer{} เยหิ โข ปน ภิกฺขูหิ สทฺธิํ วิหรามิ, เต มํ มนาเปเนว พหุลํ กายกมฺเมน สมุทาจรนฺติ อปฺปํ อมนาเปน, มนาเปเนว พหุลํ วจีกมฺเมน สมุทาจรนฺติ อปฺปํ อมนาเปน, มนาเปเนว พหุลํ มโนกมฺเมน สมุทาจรนฺติ อปฺปํ อมนาเปน, มนาปํเยว อุปหารํ อุปหรนฺติ อปฺปํ อมนาปํ.\csromanspacer{} ยานิ โข ปน ตานิ เวทยิตานิ ปิตฺตสมุฏฺฐานานิ วา เสมฺหสมุฏฺฐานานิ วา วาตสมุฏฺฐานานิ วา สนฺนิปาติกานิ วา อุตุปริณามชานิ วา วิสมปริหารชานิ วา โอปกฺกมิกานิ วา กมฺมวิปากชานิ วา, ตานิ เม น พหุเทว อุปฺปชฺชนฺติ, อปฺปาพาโธหมสฺมิ.\csromanspacer{} จตุนฺนํ โข ปนสฺมิ\footnote{จตุนฺนํ โข ปน (สี), จตุนฺนํ (สฺยา, อิ)} ฌานานํ อาภิเจตสิกานํ ทิฏฺฐธมฺมสุขวิหารานํ นิกามลาภี\footnote{นิกามลาภี โหมิ (สี, ก)} อกิจฺฉลาภี อกสิรลาภี.\csromanspacer{} อาสวานํ ขยา ฯเปฯ สจฺฉิกตฺวา อุปสมฺปชฺช วิหรามิ.}

\prose{ยํ หิ ตํ ภิกฺขเว สมฺมา วทมาโน วเทยฺย “สมเณสุ สมณ สุขุมาโล”ติ, มเมว ตํ ภิกฺขเว สมฺมา วทมาโน วเทยฺย “สมเณสุ สมณสุขุมาโล”ติ.\csromanspacer{}\csromanspacer{}จตุตฺถํ.}

\csromanmatikaanchor{mtk.13}
\csromantocmarkref{subsection}{5. ผาสุวิหารสุตฺต}{mtk.13}
\bookmark[dest={mtk.13},level=2]{5. ผาสุวิหารสุตฺต}
\csromanheader{2}{5. ผาสุวิหารสุตฺต}
\proseitem{105}{ปญฺจิเม ภิกฺขเว ผาสุวิหารา.\csromanspacer{} กตเม ปญฺจ? อิธ ภิกฺขเว ภิกฺขุโน เมตฺตํ กายกมฺมํ ปจฺจุปฏฺฐิตํ โหติ สพฺรหฺมจารีสุ อาวิ เจว รโห จ.\csromanspacer{} เมตฺตํ วจีกมฺมํ.\csromanspacer{} เมตฺตํ มโนกมฺมํ ปจฺจุปฏฺฐิตํ โหติ สพฺรหฺมจารีสุ อาวิ เจว รโห จ.\csromanspacer{} ยานิ ตานิ สีลานิ อขณฺฑานิ อจฺฉิทฺทานิ อสพลานิ อกมฺมาสานิ ภุชิสฺสานิ วิญฺญุปฺปสตฺถานิ อปรามฏฺฐานิ สมาธิสํวตฺตนิกานิ, ตถารูเปหิ สีเลหิ สีลสามญฺญคโต วิหรติ สพฺรหฺมจารีหิ อาวิ เจว รโห จ.\csromanspacer{} ยายํ ทิฏฺฐิ อริยา นิยฺยานิกา นิยฺยาติ ตกฺกรสฺส สมฺมา ทุกฺขกฺขยาย, ตถารูปาย ทิฏฺฐิยา ทิฏฺฐิสามญฺญคโต วิหรติ สพฺรหฺมจารีหิ อาวิ เจว รโห จ.\csromanspacer{} อิเม โข ภิกฺขเว ปญฺจ ผาสุวิหาราติ.\csromanspacer{}\csromanspacer{}ปญฺจมํ.}

\csromanmatikaanchor{mtk.14}
\csromantocmarkref{subsection}{6. อานนฺทสุตฺต}{mtk.14}
\bookmark[dest={mtk.14},level=2]{6. อานนฺทสุตฺต}
\csromanheader{2}{(11) 1. ผาสุวิหารวคฺค \textbf{6. อานนฺทสุตฺต}}
\proseitem{106}{\csromanfolio{117}เอกํ สมยํ ภควา โกสมฺพิยํ วิหรติ โฆสิตาราเม.\csromanspacer{} อถ โข อายสฺมา อานนฺโท เยน ภควา เตนุปสงฺกมิ, อุปสงฺกมิตฺวา ภควนฺตํ อภิวาเทตฺวา เอกมนฺตํ นิสีทิ, เอกมนฺตํ นิสินฺโน โข อายสฺมา อานนฺโท ภควนฺตํ เอตทโวจ\csromandash{}}

\prose{กิตฺตาวตา นุ โข ภนฺเต ภิกฺขุ สํเฆ\footnote{ภิกฺขุสํโฆ (สฺยา, อิ)} วิหรนฺโต ผาสุํ วิหเรยฺยาติ.\csromanspacer{} ยโต โข อานนฺท ภิกฺขุ อตฺตนา\footnote{อตฺตนา จ (อิ, ก)} สีลสมฺปนฺโน โหติ โน\footnote{โน จ (ก)} ปรํ อธิสีเล สมฺปวตฺตา\footnote{สมฺปวตฺตา โหติ (ก)}.\csromanspacer{} เอตฺตาวตาปิ โข อานนฺท ภิกฺขุ สํเฆ วิหรนฺโต ผาสุํ วิหเรยฺยาติ.}

\prose{สิยา ปน ภนฺเต อญฺโญปิ ปริยาโย ยถา ภิกฺขุ สํเฆ วิหรนฺโต ผาสุํ วิหเรยฺยาติ.\csromanspacer{} สิยา อานนฺท\footnote{อานนฺทาติ ภควา อโวจ (สฺยา, อิ)}, ยโต โข อานนฺท ภิกฺขุ อตฺตนา สีลสมฺปนฺโน โหติ โน ปรํ อธิสีเล สมฺปวตฺตา, อตฺตานุเปกฺขี จ โหติ โน ปรานุเปกฺขี.\csromanspacer{} เอตฺตาวตาปิ โข อานนฺท ภิกฺขุ สํเฆ วิหรนฺโต ผาสุํ วิหเรยฺยาติ.}

\prose{สิยา ปน ภนฺเต อญฺโญปิ ปริยาโย ยถา ภิกฺขุ สํเฆ วิหรนฺโต ผาสุํ วิหเรยฺยาติ.\csromanspacer{} สิยา อานนฺท, ยโต โข อานนฺท ภิกฺขุ อตฺตนา สีลสมฺปนฺโน โหติ โน ปรํ อธิสีเล สมฺปวตฺตา, อตฺตานุเปกฺขี จ โหติ โน ปรานุเปกฺขี, อปญฺญาโต จ โหติ, เตน จ อปญฺญาตเกน โน ปริตสฺสติ.\csromanspacer{} เอตฺตาวตาปิ โข อานนฺท ภิกฺขุ สํเฆ วิหรนฺโต ผาสุํ วิหเรยฺยาติ.}

\prose{สิยา ปน ภนฺเต อญฺโญปิ ปริยาโย ยถา ภิกฺขุ สํเฆ วิหรนฺโต ผาสุํ วิหเรยฺยาติ.\csromanspacer{} สิยา อานนฺท, ยโต โข อานนฺท ภิกฺขุ อตฺตนา สีลสมฺปนฺโน โหติ โน ปรํ อธิสีเล สมฺปวตฺตา, อตฺตานุเปกฺขี จ โหติ โน ปรานุเปกฺขี, อปญฺญาโต จ โหติ, เตน จ อปญฺญาตเกน โน ปริตสฺสติ, จตุนฺนญฺจ\footnote{จตุนฺนํ (อิ, ก)} ฌานานํ อาภิเจตสิกานํ ทิฏฺฐธมฺม สุขวิหารานํ นิกามลาภี โหติ อกิจฺฉลาภี อกสิรลาภี.\csromanspacer{} เอตฺตาวตาปิ โข อานนฺท ภิกฺขุ สํเฆ วิหรนฺโต ผาสุํ วิหเรยฺยาติ.}

\csromanmatikaanchor{mtk.15}
\csromantocmarkref{subsection}{7-10. สีลสุตฺตาทิ}{mtk.15}
\bookmark[dest={mtk.15},level=2]{7-10. สีลสุตฺตาทิ}
\prose{\csromanfolio{118}สิยา ปน ภนฺเต อญฺโญปิ ปริยาโย ยถา ภิกฺขุ สํเฆ วิหรนฺโต ผาสุํ วิหเรยฺยาติ.\csromanspacer{} สิยา อานนฺท, ยโต โข อานนฺท ภิกฺขุ อตฺตนา สีลสมฺปนฺโน โหติ โน ปรํ อธิสีเล สมฺปวตฺตา, อตฺตานุเปกฺขี จ โหติ โน ปรานุเปกฺขี, อปญฺญาโต จ โหติ, เตน จ อปญฺญาตเกน โน ปริตสฺสติ, จตุนฺนญฺจ ฌานานํ อาภิเจตสิกานํ ทิฏฺฐธมฺม สุขวิหารานํ นิกามลาภี โหติ อกิจฺฉลาภี อกสิรลาภี, อาสวานญฺจ\footnote{อาสวานํ (สี, อิ, ก)} ขยา อนาสวํ เจโตวิมุตฺติํ ปญฺญาวิมุตฺติํ ทิฏฺเฐว ธมฺเม สยํ อภิญฺญา สจฺฉิกตฺวา อุปสมฺปชฺช วิหรติ.\csromanspacer{} เอตฺตาวตาปิ โข อานนฺท ภิกฺขุ สํเฆ วิหรนฺโต ผาสุํ วิหเรยฺย.}

\prose{อิมมฺหา จาหํ อานนฺท ผาสุวิหารา อญฺโญ ผาสุวิหาโร อุตฺตริตโร วา ปณีตตโร วา นตฺถีติ วทามีติ.\csromanspacer{}\csromanspacer{}ฉฏฺฐํ.}

\csromanheader{2}{7. สีลสุตฺต}
\proseitem{107}{ปญฺจหิ ภิกฺขเว ธมฺเมหิ สมนฺนาคโต ภิกฺขุ อาหุเนยฺโย โหติ ปาหุเนยฺโย ทกฺขิเณยฺโย อญฺชลิกรณีโย อนุตฺตรํ ปุญฺญกฺเขตฺตํ โลกสฺส.}

\prose{กตเมหิ ปญฺจหิ? อิธ ภิกฺขเว ภิกฺขุ สีลสมฺปนฺโน โหติ, สมาธิ สมฺปนฺโน โหติ, ปญฺญาสมฺปนฺโน โหติ, วิมุตฺติสมฺปนฺโน โหติ, วิมุตฺติญาณทสฺสน สมฺปนฺโน โหติ.}

\prose{อิเมหิ โข ภิกฺขเว ปญฺจหิ ธมฺเมหิ สมนฺนาคโต ภิกฺขุ อาหุเนยฺโย โหติ ปาหุเนยฺโย ทกฺขิเณยฺโย อญฺชลิกรณีโย อนุตฺตรํ ปุญฺญกฺเขตฺตํ โลกสฺสาติ.\csromanspacer{}\csromanspacer{}สตฺตมํ.}

\csromanheader{2}{8. อเสขสุตฺต}
\proseitem{108}{ปญฺจหิ ภิกฺขเว ธมฺเมหิ สมนฺนาคโต ภิกฺขุ อาหุเนยฺโย โหติ ปาหุเนยฺโย ทกฺขิเณยฺโย ฯเปฯ อนุตฺตรํ ปุญฺญกฺเขตฺตํ โลกสฺส.}

\prose{กตเมหิ ปญฺจหิ? อิธ ภิกฺขเว ภิกฺขุ อเสเขน สีลกฺขนฺเธน สมนฺนาคโต โหติ, อเสเขน สมาธิกฺขนฺเธน สมนฺนาคโต โหติ,}

\vspace{\tipitakaparskip}
\csromancenter{(11) 1.\csromanspacer{} ผาสุวิหารวคฺค อเสเขน ปญฺญากฺขนฺเธน สมนฺนาคโต โหติ, อเสเขน วิมุตฺติกฺขนฺเธน สมนฺนาคโต โหติ, อเสเขน วิมุตฺติญาณทสฺสนกฺขนฺเธน สมนฺนาคโต โหติ.\csromanspacer{} อิเมหิ โข ภิกฺขเว ปญฺจหิ ธมฺเมหิ สมนฺนาคโต ภิกฺขุ อาหุเนยฺโย โหติ ฯเปฯ อนุตฺตรํ ปุญฺญกฺเขตฺตํ โลกสฺสาติ.\csromanspacer{}\csromanspacer{}อฏฺฐมํ.}
\csromanheader{2}{9. จาตุทฺทิสสุตฺต}
\proseitem{109}{\csromanfolio{119}ปญฺจหิ ภิกฺขเว ธมฺเมหิ สมนฺนาคโต ภิกฺขุ จาตุทฺทิโส โหติ.\csromanspacer{} กตเมหิ ปญฺจหิ? อิธ ภิกฺขเว ภิกฺขุ สีลวา โหติ, ปาติโมกฺขสํวรสํวุโต วิหรติ อาจารโคจรสมฺปนฺโน, อณุมตฺเตสุ วชฺเชสุ ภยทสฺสาวี สมาทาย สิกฺขติ สิกฺขาปเทสุ.\csromanspacer{} พหุสฺสุโต โหติ สุตธโร สุตสนฺนิจโย, เย เต ธมฺมา อาทิกลฺยาณา มชฺเฌกลฺยาณา ปริโยสานกลฺยาณา สาตฺถํ สพฺยญฺชนํ เกวลปริปุณฺณํ ปริสุทฺธํ พฺรหฺมจริยํ อภิวทนฺติ, ตถารูปาสฺส ธมฺมา พหุสฺสุตา โหนฺติ ธาตา วจสา ปริจิตา มนสานุเปกฺขิตา ทิฏฺฐิยา สุปฺปฏิวิทฺธา.\csromanspacer{} สนฺตุฏฺโฐ โหติ อิตรีตร จีวร ปิณฺฑปาต เสนาสน คิลานปฺปจฺจย เภสชฺช ปริกฺขาเรน.\csromanspacer{} จตุนฺนํ ฌานานํ อาภิเจตสิกานํ ทิฏฺฐธมฺมสุขวิหารานํ นิกามลาภี โหติ อกิจฺฉลาภี อกสิรลาภี.\csromanspacer{} อาสวานํ ขยา อนาสวํ เจโตวิมุตฺติํ ปญฺญาวิมุตฺติํ ทิฏฺเฐว ธมฺเม สยํ อภิญฺญา สจฺฉิกตฺวา อุปสมฺปชฺช วิหรติ.\csromanspacer{} อิเมหิ โข ภิกฺขเว ปญฺจหิ ธมฺเมหิ สมนฺนาคโต ภิกฺขุ จาตุทฺทิโส โหตีติ.\csromanspacer{}\csromanspacer{}นวมํ.}

\csromanheader{2}{10. อรญฺญสุตฺต}
\proseitem{110}{ปญฺจหิ ภิกฺขเว ธมฺเมหิ สมนฺนาคโต ภิกฺขุ อลํ อรญฺญวนปตฺถานิ ปนฺตานิ เสนาสนานิ ปฏิเสวิตุํ.\csromanspacer{} กตเมหิ ปญฺจหิ? อิธ ภิกฺขเว ภิกฺขุ สีลวา โหติ ฯเปฯ สมาทาย สิกฺขติ สิกฺขาปเทสุ.\csromanspacer{} พหุสฺสุโต โหติ ฯเปฯ ทิฏฺฐิยา สุปฺปฏิวิทฺธา.\csromanspacer{} อารทฺธวีริโย วิหรติ ถามวา ทฬฺหปรกฺกโม อนิกฺขิตฺตธุโร กุสเลสุ ธมฺเมสุ.\csromanspacer{} จตุนฺนํ ฌานานํ อาภิเจตสิกานํ ทิฏฺฐธมฺมสุข วิหารานํ นิกามลาภี โหติ อกิจฺฉลาภี อกสิรลาภี.\csromanspacer{} อาสวานํ ขยา อนาสวํ เจโตวิมุตฺติํ ปญฺญาวิมุตฺติํ ทิฏฺเฐว ธมฺเม สยํ อภิญฺญา สจฺฉิกตฺวา อุปสมฺปชฺช วิหรติ. \csromanfolio{120}อิเมหิ โข ภิกฺขเว ปญฺจหิ ธมฺเมหิ สมนฺนาคโต ภิกฺขุ อลํ อรญฺญวนปตฺถานิ ปนฺตานิ เสนาสนานิ ปฏิเสวิตุนฺติ.\csromanspacer{}\csromanspacer{}ทสมํ.\csromanspacer{} ผาสุวิหารวคฺโค ปฐโม.}

\titlehead{\textbf{ตสฺสุทฺทานํ}}
\csromangathameasure{สารชฺชํ สงฺกิโต โจโร, สุขุมาลํ ผาสุ ปญฺจมํ. \\ อานนฺท สีลาเสขา จ, จาตุทฺทิโส อรญฺเญน จาติ.}
\csromangathagroup{\csromangathastanza{สารชฺชํ สงฺกิโต โจโร, สุขุมาลํ ผาสุ ปญฺจมํ. \\ อานนฺท สีลาเสขา จ, จาตุทฺทิโส อรญฺเญน จาติ.}}

\csromanmatikaanchor{mtk.16}
\csromantocmarknopage{section}{(12) 2. อนฺธกวินฺทวคฺค}{mtk.16}
\bookmark[dest={mtk.16},level=1]{(12) 2. อนฺธกวินฺทวคฺค}
\csromanheader{1}{\textbf{(12) 2. อนฺธกวินฺทวคฺค}}
\csromanmatikaanchor{mtk.17}
\csromantocmarkref{subsection}{1. กุลูปกสุตฺต}{mtk.17}
\bookmark[dest={mtk.17},level=2]{1. กุลูปกสุตฺต}
\csromanheader{2}{1. กุลูปกสุตฺต}
\proseitem{111}{ปญฺจหิ ภิกฺขเว ธมฺเมหิ สมนฺนาคโต กุลูปโก ภิกฺขุ กุเลสุ อปฺปิโย จ โหติ อมนาโป จ อครุ จ อภาวนีโย จ.\csromanspacer{} กตเมหิ ปญฺจหิ? อสนฺถววิสฺสาสี\footnote{อสนฺถุตวิสฺสาสี (สี), อสนฺธววิสฺสาสี (ก)} จ โหติ อนิสฺสรวิกปฺปี จ วิสฺสฏฺฐุปเสวี\footnote{วิยตฺถุปเสวี (สี), พฺยตฺถุปเสวี (สฺยา, กํ), วฺยตฺตูปเสวี (อิ)} จ อุปกณฺณกชปฺปี จ อติยาจนโก จ.\csromanspacer{} อิเมหิ โข ภิกฺขเว ปญฺจหิ ธมฺเมหิ สมนฺนาคโต กุลูปโก ภิกฺขุ กุเลสุ อปฺปิโย จ โหติ อมนาโป จ อครุ จ อภาวนีโย จ.}

\prose{ปญฺจหิ ภิกฺขเว ธมฺเมหิ สมนฺนาคโต กุลูปโก ภิกฺขุ กุเลสุ ปิโย จ โหติ มนาโป จ ครุ จ ภาวนีโย จ.\csromanspacer{} กตเมหิ ปญฺจหิ? น อสนฺถววิสฺสาสี จ โหติ น อนิสฺสรวิกปฺปี จ น วิสฺสฏฺฐุปเสวี จ น อุปกณฺณกชปฺปี จ น อติยาจนโก จ.\csromanspacer{} อิเมหิ โข ภิกฺขเว ปญฺจหิ ธมฺเมหิ สมนฺนาคโต กุลูปโก ภิกฺขุ กุเลสุ ปิโย จ โหติ มนาโป จ ครุ จ ภาวนีโย จาติ.\csromanspacer{}\csromanspacer{}ปฐมํ.}

\csromanmatikaanchor{mtk.18}
\csromantocmarkref{subsection}{2. ปจฺฉาสมณสุตฺต}{mtk.18}
\bookmark[dest={mtk.18},level=2]{2. ปจฺฉาสมณสุตฺต}
\csromanheader{2}{2. ปจฺฉาสมณสุตฺต}
\proseitem{112}{ปญฺจหิ ภิกฺขเว ธมฺเมหิ สมนฺนาคโต ปจฺฉาสมโณ น อาทาตพฺโพ.\csromanspacer{} กตเมหิ ปญฺจหิ? อติทูเร วา คจฺฉติ อจฺจาสนฺเน วา, น}

\vspace{\tipitakaparskip}
\csromancenter{(12) 2.\csromanspacer{} อนฺธกวินฺทวคฺค ปตฺตปริยาปนฺนํ คณฺหติ, อาปตฺติสามนฺตา ภณมานํ น นิวาเรติ, ภณมานสฺส อนฺตรนฺตรา กถํ โอปาเตติ, ทุปฺปญฺโญ โหติ ชโฬ เอฬมูโค.\csromanspacer{} อิเมหิ โข ภิกฺขเว ปญฺจหิ ธมฺเมหิ สมนฺนาคโต ปจฺฉาสมโณ น อาทาตพฺโพ.}
\prose{\csromanfolio{121}ปญฺจหิ ภิกฺขเว ธมฺเมหิ สมนฺนาคโต ปจฺฉาสมโณ อาทาตพฺโพ.\csromanspacer{} กตเมหิ ปญฺจหิ? นาติทูเร คจฺฉติ น อจฺจาสนฺเน, ปตฺตปริยาปนฺนํ คณฺหติ, อาปตฺติสามนฺตา ภณมานํ นิวาเรติ, ภณมานสฺส น อนฺตรนฺตรา กถํ โอปาเตติ, ปญฺญวา โหติ อชโฬ อเนฬมูโค.\csromanspacer{} อิเมหิ โข ภิกฺขเว ปญฺจหิ ธมฺเมหิ สมนฺนาคโต ปจฺฉาสมโณ อาทาตพฺโพติ.\csromanspacer{}\csromanspacer{}ทุติยํ.}

\csromanmatikaanchor{mtk.19}
\csromantocmarkref{subsection}{3. สมฺมาสมาธิสุตฺต}{mtk.19}
\bookmark[dest={mtk.19},level=2]{3. สมฺมาสมาธิสุตฺต}
\csromanheader{2}{3. สมฺมาสมาธิสุตฺต}
\proseitem{113}{ปญฺจหิ ภิกฺขเว ธมฺเมหิ สมนฺนาคโต ภิกฺขุ อภพฺโพ สมฺมาสมาธิํ อุปสมฺปชฺช วิหริตุํ.\csromanspacer{} กตเมหิ ปญฺจหิ? อิธ ภิกฺขเว ภิกฺขุ อกฺขโม โหติ รูปานํ, อกฺขโม สทฺทานํ, อกฺขโม คนฺธานํ, อกฺขโม รสานํ, อกฺขโม โผฏฺฐพฺพานํ.\csromanspacer{} อิเมหิ โข ภิกฺขเว ปญฺจหิ ธมฺเมหิ สมนฺนาคโต ภิกฺขุ อภพฺโพ สมฺมาสมาธิํ อุปสมฺปชฺช วิหริตุํ.}

\prose{ปญฺจหิ ภิกฺขเว ธมฺเมหิ สมนฺนาคโต ภิกฺขุ ภพฺโพ สมฺมาสมาธิํ อุปสมฺปชฺช วิหริตุํ.\csromanspacer{} กตเมหิ ปญฺจหิ? อิธ ภิกฺขเว ภิกฺขุ ขโม โหติ รูปานํ, ขโม สทฺทานํ, ขโม คนฺธานํ, ขโม รสานํ, ขโม โผฏฺฐพฺพานํ.\csromanspacer{} อิเมหิ โข ภิกฺขเว ปญฺจหิ ธมฺเมหิ สมนฺนาคโต ภิกฺขุ ภพฺโพ สมฺมสมาธิํ อุปสมฺปชฺช วิหริตุนฺติ.\csromanspacer{}\csromanspacer{}ตติยํ.}

\csromanmatikaanchor{mtk.20}
\csromantocmarkref{subsection}{4. อนฺธกวินฺทสุตฺต}{mtk.20}
\bookmark[dest={mtk.20},level=2]{4. อนฺธกวินฺทสุตฺต}
\csromanheader{2}{4. อนฺธกวินฺทสุตฺต}
\proseitem{114}{เอกํ สมยํ ภควา มคเธสุ วิหรติ อนฺธกวินฺเท.\csromanspacer{} อถ โข อายสฺมา อานนฺโท เยน ภควา เตนุปสงฺกมิ, อุปสงฺกมิตฺวา ภควนฺตํ อภิวาเทตฺวา เอกมนฺตํ นิสีทิ, เอกมนฺตํ นิสินฺนํ โข อายสฺมนฺตํ อานนฺทํ ภควา เอตทโวจ\csromandash{}}

\csromanmatikaanchor{mtk.21}
\csromantocmarkref{subsection}{5-10. มจฺฉรินีสุตฺตาทิ}{mtk.21}
\bookmark[dest={mtk.21},level=2]{5-10. มจฺฉรินีสุตฺตาทิ}
\prose{เย เต อานนฺท ภิกฺขู นวา อจิรปพฺพชิตา อธุนาคตา อิมํ ธมฺมวินยํ, เต โว อานนฺท ภิกฺขู ปญฺจสุ ธมฺเมสุ สมาทเปตพฺพา\footnote{สมาทาเปตพฺพา (?)} นิเวเสตพฺพา \csromanfolio{122}ปติฏฺฐาเปตพฺพา.\csromanspacer{} กตเมสุ ปญฺจสุ? เอถ ตุเมฺห อาวุโส สีลวา โหถ, ปาติโมกฺข สํวรสํวุตา วิหรถ อาจารโคจรสมฺปนฺนา, อณุมตฺเตสุ วชฺเชสุ ภยทสฺสาวิโน สมาทาย สิกฺขถ สิกฺขาปเทสูติ.\csromanspacer{} อิติ ปาติโมกฺขสํวเร สมาทเปตพฺพา นิเวเสตพฺพา ปติฏฺฐาเปตพฺพา.}

\prose{เอถ ตุเมฺห อาวุโส อินฺทฺริเยสุ คุตฺตทฺวารา วิหรถ อารกฺขสติโน นิปกฺกสติโน\footnote{นิปกสติโน (สี, สฺยา), เนปกฺกสติโน (?)} สารกฺขิตมานสา สตารกฺเขน เจตสา สมนฺนาคตาติ.\csromanspacer{} อิติ อินฺทฺริยสํวเร สมาท เปตพฺพา นิเวเสตพฺพา ปติฏฺฐาเปตพฺพา.}

\prose{เอถ ตุเมฺห อาวุโส อปฺปภสฺสา โหถ ภสฺเส ปริยนฺตการิโนติ.\csromanspacer{} อิติ ภสฺสปริยนฺเต สมาทเปตพฺพา นิเวเสตพฺพา ปติฏฺฐาเปตพฺพา.}

\prose{เอถ ตุเมฺห อาวุโส อารญฺญิกา โหถ, อรญฺญวนปตฺถานิ ปนฺตานิ เสนาสนานิ ปฏิเสวถาติ.\csromanspacer{} อิติ กายวูปกาเส สมาทเปตพฺพา นิเวเสตพฺพา ปติฏฺฐาเปตพฺพา.}

\prose{เอถ ตุเมฺห อาวุโส สมฺมาทิฏฺฐิกา โหถ สมฺมาทสฺสเนน สมนฺนาคตาติ.\csromanspacer{} อิติ สมฺมาทสฺสเน สมาทเปตพฺพา นิเวเสตพฺพา ปติฏฺฐาเปตพฺพา.\csromanspacer{} เย เต อานนฺท ภิกฺขู นวา อจิรปพฺพชิตา อธุนาคตา อิมํ ธมฺมวินยํ, เต โว อานนฺท ภิกฺขู อิเมสุ ปญฺจสุ ธมฺเมสุ สมาทเปตพฺพา นิเวเสตพฺพา ปติฏฺฐาเปตพฺพาติ.\csromanspacer{}\csromanspacer{}จตุตฺถํ.}

\csromanheader{2}{5. มจฺฉรินีสุตฺต}
\proseitem{115}{ปญฺจหิ ภิกฺขเว ธมฺเมหิ สมนฺนาคตา ภิกฺขูนี ยถาภตํ นิกฺขิตฺตา เอวํ นิรเย.\csromanspacer{} กตเมหิ ปญฺจหิ? อาวาสมจฺฉรินี โหติ, กุลมจฺฉรินี โหติ, ลาภมจฺฉรินี โหติ, วณฺณมจฺฉรินี โหติ, ธมฺมมจฺฉรินี โหติ.\csromanspacer{} อิเมหิ โข ภิกฺขเว ปญฺจหิ ธมฺเมหิ สมนฺนาคตา ภิกฺขูนี ยถาภตํ นิกฺขิตฺตา เอวํ นิรเย.}

\prose{ปญฺจหิ ภิกฺขเว ธมฺเมหิ สมนฺนาคตา ภิกฺขูนี ยถาภตํ นิกฺขิตฺตา เอวํ สคฺเค.\csromanspacer{} กตเมหิ ปญฺจหิ? น อาวาสมจฺฉรินี โหติ, น กุลมจฺฉรินี โหติ, น ลาภมจฺฉรินี โหติ, น วณฺณมจฺฉรินี โหติ, น ธมฺมมจฺฉรินี}

\vspace{\tipitakaparskip}
\csromancenter{(12) 2.\csromanspacer{} อนฺธกวินฺทวคฺค โหติ.\csromanspacer{} อิเมหิ โข ภิกฺขเว ปญฺจหิ ธมฺเมหิ สมนฺนาคตา ภิกฺขูนี ยถาภตํ นิกฺขิตฺตา เอวํ สคฺเคติ.\csromanspacer{}\csromanspacer{}ปญฺจมํ.}
\csromanheader{2}{6. วณฺณนาสุตฺต}
\proseitem{116}{\csromanfolio{123}ปญฺจหิ ภิกฺขเว ธมฺเมหิ สมนฺนาคตา ภิกฺขูนี ยถาภตํ นิกฺขิตฺตา เอวํ นิรเย.\csromanspacer{} กตเมหิ ปญฺจหิ? อนนุวิจฺจ อปริโยคาเหตฺวา อวณฺณารหสฺส วณฺณํ ภาสติ, อนนุวิจฺจ อปริโยคาเหตฺวา วณฺณารหสฺส อวณฺณํ ภาสติ, อนนุวิจฺจ อปริโยคาเหตฺวา อปฺปสาทนีเย ฐาเน ปสาทํ อุปทํเสติ, อนนุวิจฺจ อปริโยคาเหตฺวา ปสาทนีเย ฐาเน อปฺปสาทํ อุปทํเสติ, สทฺธาเทยฺยํ วินิปาเตติ.\csromanspacer{} อิเมหิ โข ภิกฺขเว ปญฺจหิ ธมฺเมหิ สมนฺนาคตา ภิกฺขูนี ยถาภตํ นิกฺขิตฺตา เอวํ นิรเย.}

\prose{ปญฺจหิ ภิกฺขเว ธมฺเมหิ สมนฺนาคตา ภิกฺขุนี ยถาภตํ นิกฺขิตฺตา เอวํ สคฺเค.\csromanspacer{} กตเมหิ ปญฺจหิ? อนุวิจฺจ ปริโยคาเหตฺวา อวณฺณารหสฺส อวณฺณํ ภาสติ, อนุวิจฺจ ปริโยคาเหตฺวา วณฺณารหสฺส วณฺณํ ภาสติ, อนุวิจฺจ ปริโยคาเหตฺวา อปฺปสาทนีเย ฐาเน อปฺปสาทํ อุปทํเสติ, อนุวิจฺจ ปริโยคาเหตฺวา ปสาทนีเย ฐาเน ปสาทํ อุปทํเสติ, สทฺธาเทยฺยํ น วินิปาเตติ.\csromanspacer{} อิเมหิ โข ภิกฺขเว ปญฺจหิ ธมฺเมหิ สมนฺนาคตา ภิกฺขุนี ยถาภตํ นิกฺขิตฺตา เอวํ สคฺเคติ.\csromanspacer{}\csromanspacer{}ฉฏฺฐํ.}

\csromanheader{2}{7. อิสฺสุกินีสุตฺต}
\proseitem{117}{ปญฺจหิ ภิกฺขเว ธมฺเมหิ สมนฺนาคตา ภิกฺขูนี ยถาภตํ นิกฺขิตฺตา เอวํ นิรเย.\csromanspacer{} กตเมหิ ปญฺจหิ? อนนุวิจฺจ อปริโยคาเหตฺวา อวณฺณารหสฺส วณฺณํ ภาสติ, อนนุวิจฺจ อปริโยคาเหตฺวา วณฺณารหสฺส อวณฺณํ ภาสติ, อิสฺสุกินี จ โหติ มจฺฉรินี จ, สทฺธาเทยฺยํ\footnote{สทฺธาเทยฺยญฺจ (สฺยา)} วินิปาเตติ.\csromanspacer{} อิเมหิ โข ภิกฺขเว ปญฺจหิ ธมฺเมหิ สมนฺนาคตา ภิกฺขุนี ยถาภตํ นิกฺขิตฺตา เอวํ นิรเย.}

\prose{ปญฺจหิ ภิกฺขเว ธมฺเมหิ สมนฺนาคตา ภิกฺขุนี ยถาภตํ นิกฺขิตฺตา เอวํ สคฺเค.\csromanspacer{} กตเมหิ ปญฺจหิ? อนุวิจฺจ ปริโยคาเหตฺวา อวณฺณารหสฺส อวณฺณํ ภาสติ, อนุวิจฺจ ปริโยคาเหตฺวา วณฺณารหสฺส วณฺณํ \csromanfolio{124}ภาสติ, อนิสฺสุกินี จ โหติ อมจฺฉรินี จ, สทฺธาเทยฺยํ น วินิปาเตติ.\csromanspacer{} อิเมหิ โข ภิกฺขเว ปญฺจหิ ธมฺเมหิ สมนฺนาคตา ภิกฺขุนี ยถาภตํ นิกฺขิตฺตา เอวํ สคฺเคติ.\csromanspacer{}\csromanspacer{}สตฺตมํ.}

\csromanheader{2}{8. มิจฺฉาทิฏฺฐิกสุตฺต}
\proseitem{118}{ปญฺจหิ ภิกฺขเว ธมฺเมหิ สมนฺนาคตา ภิกฺขุนี ยถาภตํ นิกฺขิตฺตา เอวํ นิรเย.\csromanspacer{} กตเมหิ ปญฺจหิ? อนนุวิจฺจ อปริโยคาเหตฺวา อวณฺณารหสฺส วณฺณํ ภาสติ, อนนุวิจฺจ อปริโยคาเหตฺวา วณฺณารหสฺส อวณฺณํ ภาสติ, มิจฺฉาทิฏฺฐิกา จ โหติ มิจฺฉาสงฺกปฺปา จ, สทฺธาเทยฺยํ วินิปาเตติ.\csromanspacer{} อิเมหิ โข ภิกฺขเว ปญฺจหิ ธมฺเมหิ สมนฺนาคตา ภิกฺขุนี ยถาภตํ นิกฺขิตฺตา เอวํ นิรเย.}

\prose{ปญฺจหิ ภิกฺขเว ธมฺเมหิ สมนฺนาคตา ภิกฺขุนี ยถาภตํ นิกฺขิตฺตา เอวํ สคฺเค.\csromanspacer{} กตเมหิ ปญฺจหิ? อนุวิจฺจ ปริโยคาเหตฺวา อวณฺณารหสฺส อวณฺณํ ภาสติ, อนุวิจฺจ ปริโยคาเหตฺวา วณฺณารหสฺส วณฺณํ ภาสติ, สมฺมาทิฏฺฐิกา จ โหติ สมฺมาสงฺกปฺปา จ, สทฺธาเทยฺยํ น วินิปาเตติ.\csromanspacer{} อิเมหิ โข ภิกฺขเว ปญฺจหิ ธมฺเมหิ สมนฺนาคตา ภิกฺขุนี ยถาภตํ นิกฺขิตฺตา เอวํ สคฺเคติ.\csromanspacer{}\csromanspacer{}อฏฺฐมํ.}

\csromanheader{2}{9. มิจฺฉาวาจาสุตฺต}
\proseitem{119}{ปญฺจหิ ภิกฺขเว ธมฺเมหิ สมนฺนาคตา ภิกฺขุนี ยถาภตํ นิกฺขิตฺตา เอวํ นิรเย.\csromanspacer{} กตเมหิ ปญฺจหิ? อนนุวิจฺจ อปริโยคาเหตฺวา อวณฺณารหสฺส วณฺณํ ภาสติ, อนนุวิจฺจ อปริโยคาเหตฺวา วณฺณารหสฺส อวณฺณํ ภาสติ, มิจฺฉาวาจา จ โหติ มิจฺฉากมฺมนฺตา จ, สทฺธาเทยฺยํ วินิปาเตติ.\csromanspacer{} อิเมหิ โข ภิกฺขเว ปญฺจหิ ธมฺเมหิ สมนฺนาคตา ภิกฺขุนี ยถาภตํ นิกฺขิตฺตา เอวํ นิรเย.}

\prose{ปญฺจหิ ภิกฺขเว ธมฺเมหิ สมนฺนาคตา ภิกฺขุนี ยถาภตํ นิกฺขิตฺตา เอวํ สคฺเค.\csromanspacer{} กตเมหิ ปญฺจหิ? อนุวิจฺจ ปริโยคาเหตฺวา อวณฺณารหสฺส อวณฺณํ ภาสติ, อนุวิจฺจ ปริโยคาเหตฺวา วณฺณารหสฺส วณฺณํ ภาสติ, สมฺมาวาจา จ โหติ สมฺมากมฺมนฺตา จ, สทฺธาเทยฺยํ น วินิปาเตติ.\csromanspacer{} อิเมหิ โข ภิกฺขเว ปญฺจหิ ธมฺเมหิ สมนฺนาคตา ภิกฺขุนี ยถาภตํ นิกฺขิตฺตา เอวํ สคฺเคติ.\csromanspacer{}\csromanspacer{}นวมํ.}

\chapterheadpageread{\textbf{(13) 3. คิลานวคฺค} \textbf{10. มิจฺฉาวายามสุตฺต}}
\proseitem{120}{\csromanfolio{125}ปญฺจหิ ภิกฺขเว ธมฺเมหิ สมนฺนาคตา ภิกฺขุนี ยถาภตํ นิกฺขิตฺตา เอวํ นิรเย.\csromanspacer{} กตเมหิ ปญฺจหิ? อนนุวิจฺจ อปริโยคาเหตฺวา อวณฺณารหสฺส วณฺณํ ภาสติ, อนนุวิจฺจ อปริโยคาเหตฺวา วณฺณารหสฺส อวณฺณํ ภาสติ, มิจฺฉาวายามา จ โหติ มิจฺฉาสตินี จ\footnote{มิจฺฉาสติ จ (สฺยา)}, สทฺธาเทยฺยํ วินิปาเตติ.\csromanspacer{} อิเมหิ โข ภิกฺขเว ปญฺจหิ ธมฺเมหิ สมนฺนาคตา ภิกฺขุนี ยถาภตํ นิกฺขิตฺตา เอวํ นิรเย.}

\prose{ปญฺจหิ ภิกฺขเว ธมฺเมหิ สมนฺนาคตา ภิกฺขุนี ยถาภตํ นิกฺขิตฺตา เอวํ สคฺเค.\csromanspacer{} กตเมหิ ปญฺจหิ? อนุวิจฺจ ปริโยคาเหตฺวา อวณฺณารหสฺส อวณฺณํ ภาสติ, อนุวิจฺจ ปริโยคาเหตฺวา วณฺณารหสฺส วณฺณํ ภาสติ, สมฺมาวายามา จ โหติ สมฺมาสตินี จ, สทฺธาเทยฺยํ น วินิปาเตติ.\csromanspacer{} อิเมหิ โข ภิกฺขเว ปญฺจหิ ธมฺเมหิ สมนฺนาคตา ภิกฺขุนี ยถาภตํ นิกฺขิตฺตา เอวํ สคฺเคติ.\csromanspacer{}\csromanspacer{}ทสมํ.\csromanspacer{} อนฺธกวินฺทวคฺโค ทุติโย.}

\titlehead{\textbf{ตสฺสุทฺทานํ}}
\csromangathameasure{กุลูปโก ปจฺฉาสมโณ, สมาธิ-อนฺธกวินฺทํ. \\ มจฺฉรี วณฺณนา อิสฺสา, ทิฏฺฐิวาจาย วายมาติ.}
\csromangathagroup{\csromangathastanza{กุลูปโก ปจฺฉาสมโณ, สมาธิ-อนฺธกวินฺทํ. \\ มจฺฉรี วณฺณนา อิสฺสา, ทิฏฺฐิวาจาย วายมาติ.}}

\csromanmatikaanchor{mtk.22}
\csromantocmarknopage{section}{(13) 3. คิลานวคฺค}{mtk.22}
\bookmark[dest={mtk.22},level=1]{(13) 3. คิลานวคฺค}
\csromanheader{1}{\textbf{(13) 3. คิลานวคฺค}}
\csromanmatikaanchor{mtk.23}
\csromantocmarkref{subsection}{1. คิลานสุตฺต}{mtk.23}
\bookmark[dest={mtk.23},level=2]{1. คิลานสุตฺต}
\csromanheader{2}{1. คิลานสุตฺต}
\proseitem{121}{เอกํ สมยํ ภควา เวสาลิยํ วิหรติ มหาวเน กูฏาคารสาลายํ.\csromanspacer{} อถ โข ภควา สายนฺหสมยํ ปฏิสลฺลานา วุฏฺฐิโต เยน คิลานสาลา เตนุปสงฺกมิ.\csromanspacer{} อทฺทสา โข ภควา อญฺญตรํ ภิกฺขุํ ทุพฺพลํ คิลานกํ, ทิสฺวา ปญฺญตฺเต อาสเน นิสีทิ, นิสชฺช โข ภควา ภิกฺขู อามนฺเตสิ\csromandash{}}

\csromanmatikaanchor{mtk.24}
\csromanmatikaanchor{mtk.25}
\csromantocmarkref{subsection}{2. สติสูปฏฺฐิตสุตฺต}{mtk.24}
\bookmark[dest={mtk.24},level=2]{2. สติสูปฏฺฐิตสุตฺต}
\csromantocmarkref{subsection}{3-4. อุปฏฺฐากสุตฺตทฺวย}{mtk.25}
\bookmark[dest={mtk.25},level=2]{3-4. อุปฏฺฐากสุตฺตทฺวย}
\prose{\csromanfolio{126}ยํ กิญฺจิ\footnote{ยํ กญฺจิ (สฺยา, กํ)} ภิกฺขเว ภิกฺขุํ ทุพฺพลํ\footnote{ภิกฺขเว ทุพฺพลํ (สี, สฺยา, กํ, อิ)} คิลานกํ ปญฺจ ธมฺมา น วิชหนฺติ, หสฺเสตํ ปาฏิกงฺขํ “นจิรสฺเสว อาสวานํ ขยา อนาสวํ เจโตวิมุตฺติํ ปญฺญาวิมุตฺติํ ทิฏฺเฐว ธมฺเม สยํ อภิญฺญา สจฺฉิกตฺวา อุปสมฺปชฺช วิหริสฺสตี”ติ.}

\prose{กตเม ปญฺจ? อิธ ภิกฺขเว ภิกฺขุ อสุภานุปสฺสี กาเย วิหรติ, อาหาเร ปฏิกูลสญฺญี, สพฺพโลเก อนภิรตสญฺญี\footnote{สพฺพตฺถปิ เอวเมว ทิสฺสติ.}, สพฺพสงฺขาเรสุ อนิจฺจานุปสฺสี, มรณสญฺญา โข ปนสฺส อชฺฌตฺตํ สูปฏฺฐิตา โหติ.\csromanspacer{} ยํ กิญฺจิ ภิกฺขเว ภิกฺขุํ ทุพฺพลํ คิลานกํ อิเม ปญฺจ ธมฺมา น วิชหนฺติ, ตสฺเสตํ ปาฏิกงฺขํ “นจิรสฺเสว อาสวานํ ขยา ฯเปฯ สจฺฉิกตฺวา อุปสมฺปชฺช วิหริสฺสตี”ติ.\csromanspacer{}\csromanspacer{}ปฐมํ.}

\csromanheader{2}{2. สติสูปฏฺฐิตสุตฺต}
\proseitem{122}{โย หิ โกจิ ภิกฺขเว ภิกฺขุ วา ภิกฺขุนี วา ปญฺจ ธมฺเม ภาเวติ, ปญฺจ ธมฺเม พหุลีกโรติ, ตสฺส ทฺวินฺนํ ผลานํ อญฺญตรํ ผลํ ปาฏิกงฺขํ “ทิฏฺเฐว ธมฺเม อญฺญา, สติ วา อุปาทิเสเส อนาคามิตา”.}

\prose{กตเม ปญฺจ? อิธ ภิกฺขเว ภิกฺขุโน อชฺฌตฺตญฺเญว สติ สูปฏฺฐิตา โหติ, ธมฺมานํ อุทยตฺถคามินิยา ปญฺญาย อสุภานุปสฺสี กาเย วิหรติ, อาหาเร ปฏิกูลสญฺญี, สพฺพโลเก อนภิรตสญฺญี, สพฺพสงฺขาเรสุ อนิจฺจานุปสฺสี.\csromanspacer{} โย หิ โกจิ ภิกฺขเว ภิกฺขุ วา ภิกฺขุนี วา อิเม ปญฺจ ธมฺเม ภาเวติ.\csromanspacer{} อิเม ปญฺจ ธมฺเม พหุลีกโรติ, ตสฺส ทฺวินฺนํ ผลานํ อญฺญตรํ ผลํ ปาฏิกงฺขํ “ทิฏฺเฐว ธมฺเม อญฺญา, สติ วา อุปาทิเสเส อนาคามิตา”ติ.\csromanspacer{}\csromanspacer{}ทุติยํ.}

\csromanheader{2}{3. ปฐม อุปฏฺฐากสุตฺต}
\proseitem{123}{ปญฺจหิ ภิกฺขเว ธมฺเมหิ สมนฺนาคโต คิลาโน ทูปฏฺฐาโก\footnote{ทุปฏฺฐาโก (สฺยา, กํ, อิ, ก) วิ 3. 418 ปิฏฺเฐปิ.} โหติ.\csromanspacer{} กตเมหิ ปญฺจหิ? อสปฺปายการี โหติ, สปฺปาเย มตฺตํ น ชานาติ, เภสชฺชํ นปฺปฏิเสวิตา โหติ, อตฺถกามสฺส คิลานุปฏฺฐากสฺส น ยถาภูตํ อาพาธํ อาวิกตฺตา โหติ “อภิกฺกมนฺตํ วา ‘อภิกฺกมตี’ติ, ปฏิกฺกมนฺตํ วา ‘ปฏิกฺกมตี’ติ, ฐิตํ วา ‘ฐิโต’ติ”, อุปฺปนฺนานํ}

\vspace{\tipitakaparskip}
\csromancenter{(13) 3.\csromanspacer{} คิลานวคฺค สารีริกานํ เวทนานํ ทุกฺขานํ ติพฺพานํ\footnote{ติปฺปานํ (สี)} ขรานํ กฏุกานํ อสาตานํ อมนาปานํ ปาณหรานํ อนธิวาสกชาติโก โหติ.\csromanspacer{} อิเมหิ โข ภิกฺขเว ปญฺจหิ ธมฺเมหิ สมนฺนาคโต คิลาโน ทูปฏฺฐาโก โหติ.}
\prose{\csromanfolio{127}ปญฺจหิ ภิกฺขเว ธมฺเมหิ สมนฺนาคโต คิลาโน สูปฏฺฐาโก โหติ.\csromanspacer{} กตเมหิ ปญฺจหิ, สปฺปายการี โหติ, สปฺปาเย มตฺตํ ชานาติ, เภสชฺชํ ปฏิเสวิตา โหติ, อตฺถกามสฺส คิลานุปฏฺฐากสฺส ยถาภูตํ อาพาธํ อาวิกตฺตา โหติ อภิกฺกมนฺตํ วา “อภิกฺกมตี”ติ ปฏิกฺกมนฺตํ วา “ปฏิกฺกมตี”ติ ฐิตํ วา “ฐิโต”ติ, อุปฺปนฺนานํ สารีริกานํ เวทนานํ ทุกฺขานํ ติพฺพานํ ขรานํ กฏุกานํ อสาตานํ อมนาปานํ ปาณหรานํ อธิวาสกชาติโก โหติ.\csromanspacer{} อิเมหิ โข ภิกฺขเว ปญฺจหิ ธมฺเมหิ สมนฺนาคโต คิลาโน สูปฏฺฐาโก โหตีติ.\csromanspacer{}\csromanspacer{}ตติยํ.}

\csromanheader{2}{4. ทุติย อุปฏฺฐากสุตฺต}
\proseitem{124}{ปญฺจหิ ภิกฺขเว ธมฺเมหิ สมนฺนาคโต คิลานุปฏฺฐาโก นาลํ คิลานํ อุปฏฺฐาตุํ.\csromanspacer{} กตเมหิ ปญฺจหิ, นปฺปฏิพโล โหติ เภสชฺชํ สํวิธาตุํ, สปฺปายาสปฺปายํ น ชานาติ, อสปฺปายํ อุปนาเมติ, สปฺปายํ อปนาเมติ, อามิสนฺตโร คิลานํ อุปฏฺฐาติ โน เมตฺตจิตฺโต, เชคุจฺฉี โหติ อุจฺจารํ วา ปสฺสาวํ วา วนฺตํ วา เขฬํ วา นีหริตุํ, นปฺปฏิพโล โหติ คิลานํ กาเลน กาลํ ธมฺมิยา กถาย สนฺทสฺเสตุํ สมาทเปตุํ\footnote{สมาทาเปตุํ (?)} สมุตฺเตเชตุํ สมฺปหํเสตุํ.\csromanspacer{} อิเมหิ โข ภิกฺขเว ปญฺจหิ ธมฺเมหิ สมนฺนาคโต คิลานุปฏฺฐาโก นาลํ คิลานํ อุปฏฺฐาตุํ.}

\prose{ปญฺจหิ ภิกฺขเว ธมฺเมหิ สมนฺนาคโต คิลานุปฏฺฐาโก อลํ คิลานํ อุปฏฺฐาตุํ.\csromanspacer{} กตเมหิ ปญฺจหิ, ปฏิพโล โหติ เภสชฺชํ สํวิธาตุํ, สปฺปายา สปฺปายํ ชานาติ, อสปฺปายํ อปนาเมติ, สปฺปายํ อุปนาเมติ, เมตฺตจิตฺโต คิลานํ อุปฏฺฐาติ โน อามิสนฺตโร, อเชคุจฺฉี โหติ อุจฺจารํ วา ปสฺสาวํ วา วนฺตํ วา เขฬํ วา นีหริตุํ, ปฏิพโล โหติ คิลานํ กาเลน กาลํ ธมฺมิยา กถาย สนฺทสฺเสตุํ สมาทเปตุํ สมุตฺเตเชตุํ สมฺปหํเสตุํ.\csromanspacer{} อิเมหิ โข ภิกฺขเว ปญฺจหิ ธมฺเมหิ สมนฺนาคโต คิลานุปฏฺฐาโก อลํ คิลานํ อุปฏฺฐาตุนฺติ.\csromanspacer{}\csromanspacer{}จตุตฺถํ.}

\csromanheader{2}{5. ปฐม อนายุสฺสาสุตฺต}
\csromanmatikaanchor{mtk.26}
\csromantocmarkref{subsection}{5-6. อนายุสฺสาสุตฺตทฺวย}{mtk.26}
\bookmark[dest={mtk.26},level=2]{5-6. อนายุสฺสาสุตฺตทฺวย}
\proseitem{125}{\csromanfolio{128}ปญฺจิเม ภิกฺขเว ธมฺมา อนายุสฺสา.\csromanspacer{} กตเม ปญฺจ, อสปฺปายการี โหติ, สปฺปาเย มตฺตํ น ชานาติ, อปริณตโภชี จ โหติ, อกาลจารี จ โหติ อพฺรหฺมจารี จ.\csromanspacer{} อิเม โข ภิกฺขเว ปญฺจ ธมฺมา อนายุสฺสา.}

\prose{ปญฺจิเม ภิกฺขเว ธมฺมา อายุสฺสา.\csromanspacer{} กตเม ปญฺจ, สปฺปายการี โหติ, สปฺปาเย มตฺตํ ชานาติ, ปริณตโภชี จ โหติ, กาลจารี จ โหติ พฺรหฺมจารี จ.\csromanspacer{} อิเม โข ภิกฺขเว ปญฺจ ธมฺมา อายุสฺสาติ.\csromanspacer{}\csromanspacer{}ปญฺจมํ.}

\csromanheader{2}{6. ทุติย อนายุสฺสาสุตฺต}
\proseitem{126}{ปญฺจิเม ภิกฺขเว ธมฺมา อนายุสฺสา.\csromanspacer{} กตเม ปญฺจ, อสปฺปายการี โหติ, สปฺปาเย มตฺตํ น ชานาติ, อปริณตโภชี จ โหติ, ทุสฺสีโล จ ปาปมิตฺโต จ.\csromanspacer{} อิเม โข ภิกฺขเว ปญฺจ ธมฺมา อนายุสฺสา.}

\prose{ปญฺจิเม ภิกฺขเว ธมฺมา อายุสฺสา.\csromanspacer{} กตเม ปญฺจ, สปฺปายการี โหติ, สปฺปาเย มตฺตํ ชานาติ, ปริณตโภชี จ โหติ, สีลวา จ กลฺยาณมิตฺโต จ.\csromanspacer{} อิเม โข ภิกฺขเว ปญฺจ ธมฺมา อายุสฺสาติ.\csromanspacer{}\csromanspacer{}ฉฏฺฐํ.}

\csromanmatikaanchor{mtk.27}
\csromantocmarkref{subsection}{7. วปกาสสุตฺต}{mtk.27}
\bookmark[dest={mtk.27},level=2]{7. วปกาสสุตฺต}
\csromanheader{2}{7. วปกาสสุตฺต}
\proseitem{127}{ปญฺจหิ ภิกฺขเว ธมฺเมหิ สมนฺนาคโต ภิกฺขุ นาลํ สํฆมฺหา วปกาสิตุํ\footnote{วิ + อป + กาสิตุํ = วปกาสิตุํ.}.\csromanspacer{} กตเมหิ ปญฺจหิ, อิธ ภิกฺขเว ภิกฺขุ อสนฺตุฏฺโฐ โหติ อิตรีตเรน จีวเรน, อสนฺตุฏฺโฐ โหติ อิตรีตเรน ปิณฺฑปาเตน, อสนฺตุฏฺโฐ โหติ อิตรีตเรน เสนาสเนน, อสนฺตุฏฺโฐ โหติ อิตรีตเรน คิลานปฺปจฺจยเภสชฺชปริกฺขาเรน, กามสงฺกปฺปพหุโล จ วิหรติ.\csromanspacer{} อิเมหิ โข ภิกฺขเว ปญฺจหิ ธมฺเมหิ สมนฺนาคโต ภิกฺขุ นาลํ สํฆมฺหา วปกาสิตุํ.}

\prose{ปญฺจหิ ภิกฺขเว ธมฺเมหิ สมนฺนาคโต ภิกฺขุ อลํ สํฆมฺหา วปกาสิตุํ.\csromanspacer{} กตเมหิ ปญฺจหิ, อิธ ภิกฺขเว ภิกฺขุ สนฺตุฏฺโฐ โหติ อิตรีตเรน จีวเรน, สนฺตุฏฺโฐ โหติ อิตรีตเรน ปิณฺฑปาเตน, สนฺตุฏฺโฐ โหติ อิตรีตเรน เสนาสเนน, สนฺตุฏฺโฐ โหติ อิตรีตเรน คิลานปฺปจฺจย}

\vspace{\tipitakaparskip}
\csromancenter{(13) 3.\csromanspacer{} คิลานวคฺค เภสชฺชปริกฺขาเรน, เนกฺขมฺมสงฺกปฺปพหุโล\footnote{น กามสงฺกปฺปพหุโล (ก)} จ วิหรติ.\csromanspacer{} อิเมหิ โข ภิกฺขเว ปญฺจหิ ธมฺเมหิ สมนฺนาคโต ภิกฺขุ อลํ สํฆมฺหา วปกาสิตุนฺติ.\csromanspacer{}\csromanspacer{}สตฺตมํ.}
\csromanmatikaanchor{mtk.28}
\csromantocmarkref{subsection}{8. สมณสุขสุตฺต}{mtk.28}
\bookmark[dest={mtk.28},level=2]{8. สมณสุขสุตฺต}
\csromanheader{2}{8. สมณสุขสุตฺต}
\proseitem{128}{\csromanfolio{129}ปญฺจิมานิ ภิกฺขเว สมณทุกฺขานิ.\csromanspacer{} กตมานิ ปญฺจ? อิธ ภิกฺขเว ภิกฺขุ อสนฺตุฏฺโฐ โหติ อิตรีตเรน จีวเรน, อสนฺตุฏฺโฐ โหติ อิตรีตเรน ปิณฺฑปาเตน, อสนฺตุฏฺโฐ โหติ อิตรีตเรน เสนาสเนน, อสนฺตุฏฺโฐ โหติ อิตรีตเรน คิลานปฺปจฺจยเภสชฺชปริกฺขาเรน, อนภิรโต จ พฺรหฺมจริยํ จรติ.\csromanspacer{} อิมานิ โข ภิกฺขเว ปญฺจ สมณทุกฺขานิ.}

\prose{ปญฺจิมานิ ภิกฺขเว สมณสุขานิ.\csromanspacer{} กตมานิ ปญฺจ? อิธ ภิกฺขเว ภิกฺขุ สนฺตุฏฺโฐ โหติ อิตรีตเรน จีวเรน, สนฺตุฏฺโฐ โหติ อิตรีตเรน ปิณฺฑปาเตน, สนฺตุฏฺโฐ โหติ อิตรีตเรน เสนาสเนน, สนฺตุฏฺโฐ โหติ อิตรีตเรน คิลานปฺปจฺจยเภสชฺช ปริกฺขาเรน, อภิรโต จ พฺรหฺมจริยํ จรติ.\csromanspacer{} อิมานิ โข ภิกฺขเว ปญฺจ สมณสุขานีติ.\csromanspacer{}\csromanspacer{}อฏฺฐมํ.}

\csromanmatikaanchor{mtk.29}
\csromantocmarkref{subsection}{9. ปริกุปฺปสุตฺต}{mtk.29}
\bookmark[dest={mtk.29},level=2]{9. ปริกุปฺปสุตฺต}
\csromanheader{2}{9. ปริกุปฺปสุตฺต}
\proseitem{129}{ปญฺจิเม ภิกฺขเว อาปายิกา เนรยิกา ปริกุปฺปา อเตกิจฺฉา.\csromanspacer{} กตเม ปญฺจ? มาตา\footnote{มาตรํ (ก)} ชีวิตา โวโรปิตา โหติ, ปิตา\footnote{ปิตรํ (ก)} ชีวิตา โวโรปิโต\footnote{โวโรปิตา (ก)} โหติ, อรหํ\footnote{อรหนฺตํ (ก), อรหา (สฺยา)} ชีวิตา โวโรปิโต โหติ, ตถาคตสฺส ทุฏฺเฐน จิตฺเตน โลหิตํ อุปฺปาทิตํ โหติ, สํโฆ ภินฺโน โหติ.\csromanspacer{} อิเม โข ภิกฺขเว ปญฺจ อาปายิกา เนรยิกา ปริกุปฺปา อเตกิจฺฉาติ.\csromanspacer{}\csromanspacer{}นวมํ.}

\csromanmatikaanchor{mtk.30}
\csromantocmarkref{subsection}{10. พฺยสนสุตฺต}{mtk.30}
\bookmark[dest={mtk.30},level=2]{10. พฺยสนสุตฺต}
\csromanheader{2}{10. พฺยสนสุตฺต}
\csromanmatikaanchor{mtk.31}
\csromanmatikaanchor{mtk.32}
\csromantocmarknopage{section}{(14) 4. ราชวคฺค}{mtk.31}
\bookmark[dest={mtk.31},level=1]{(14) 4. ราชวคฺค}
\csromantocmarkref{subsection}{1-2. จกฺกานุวตฺตนสุตฺตทฺวย}{mtk.32}
\bookmark[dest={mtk.32},level=2]{1-2. จกฺกานุวตฺตนสุตฺตทฺวย}
\proseitem{130}{ปญฺจิมานิ ภิกฺขเว พฺยสนานิ.\csromanspacer{} กตมานิ ปญฺจ? ญาติพฺยสนํ โภคพฺยสนํ โรคพฺยสนํ สีลพฺยสนํ ทิฏฺฐิพฺยสนํ.\csromanspacer{} น ภิกฺขเว สตฺตา ญาติพฺยสนเหตุ \csromanfolio{130}วา โภคพฺยสนเหตุ วา โรคพฺยสนเหตุ วา กายสฺส เภทา ปรํ มรณา อปายํ ทุคฺคติํ วินิปาตํ นิรยํ อุปปชฺชนฺติ.\csromanspacer{} สีลพฺยสนเหตุ วา ภิกฺขเว สตฺตา ทิฏฺฐิพฺยสนเหตุ วา กายสฺส เภทา ปรํ มรณา อปายํ ทุคฺคติํ วินิปาตํ นิรยํ อุปปชฺชนฺติ.\csromanspacer{} อิมานิ โข ภิกฺขเว ปญฺจ พฺยสนานิ.}

\prose{ปญฺจิมา ภิกฺขเว สมฺปทา.\csromanspacer{} กตมา ปญฺจ, ญาติสมฺปทา โภคสมฺปทา อาโรคฺยสมฺปทา สีลสมฺปทา ทิฏฺฐิสมฺปทา.\csromanspacer{} น ภิกฺขเว สตฺตา ญาติสมฺปทาเหตุ วา โภคสมฺปทาเหตุ วา อาโรคฺยสมฺปทาเหตุ วา กายสฺส เภทา ปรํ มรณา สุคติํ สคฺคํ โลกํ อุปปชฺชนฺติ.\csromanspacer{} สีลสมฺปทาเหตุ วา ภิกฺขเว สตฺตา ทิฏฺฐิสมฺปทาเหตุ วา กายสฺส เภทา ปรํ มรณา สุคติํ สคฺคํ โลกํ อุปปชฺชนฺติ.\csromanspacer{} อิมา โข ภิกฺขเว ปญฺจ สมฺปทาติ.\csromanspacer{}\csromanspacer{}ทสมํ.\csromanspacer{} คิลานวคฺโค ตติโย.}

\titlehead{\textbf{ตสฺสุทฺทานํ}}
\csromangathameasure{คิลาโน สติสูปฏฺฐิ, เทฺว อุปฏฺฐากา ทุวายุสา. \\ วปกาสสมณสุขา, ปริกุปฺปํ พฺยสเนน จาติ.}
\csromangathagroup{\csromangathastanza{คิลาโน สติสูปฏฺฐิ, เทฺว อุปฏฺฐากา ทุวายุสา. \\ วปกาสสมณสุขา, ปริกุปฺปํ พฺยสเนน จาติ.}}

\csromanheader{1}{\textbf{(14) 4. ราชวคฺค}}
\csromanheader{2}{1. ปฐมจกฺกานุวตฺตนสุตฺต}
\proseitem{131}{ปญฺจหิ ภิกฺขเว องฺเคหิ สมนฺนาคโต ราชา จกฺกวตฺตี ธมฺเมเนว จกฺกํ วตฺเตติ\footnote{ปวตฺเตติ (สฺยา, อิ, ก)}, ตํ โหติ จกฺกํ อปฺปฏิวตฺติยํ\footnote{อปฺปติวตฺติยํ (สี)} เกนจิ มนุสฺสภูเตน ปจฺจตฺถิเกน ปาณินา.}

\prose{กตเมหิ ปญฺจหิ, อิธ ภิกฺขเว ราชา จกฺกวตฺตี อตฺถญฺญู จ โหติ ธมฺมญฺญู จ มตฺตญฺญู จ กาลญฺญู จ ปริสญฺญู จ.\csromanspacer{} อิเมหิ โข ภิกฺขเว ปญฺจหิ องฺเคหิ สมนฺนาคโต ราชา จกฺกวตฺตี ธมฺเมเนว จกฺกํ ปวตฺเตติ, ตํ โหติ จกฺกํ อปฺปฏิวตฺติยํ เกนจิ มนุสฺสภูเตน ปจฺจตฺถิเกน ปาณินา.}

\chapterheadpageread{(14) 4. ราชวคฺค}
\prose{\csromanfolio{131}เอวเมวํ โข ภิกฺขเว ปญฺจหิ ธมฺเมหิ สมนฺนาคโต ตถาคโต อรหํ สมฺมาสมฺพุทฺโธ ธมฺเมเนว อนุตฺตรํ ธมฺมจกฺกํ ปวตฺเตติ, ตํ โหติ จกฺกํ อปฺปฏิวตฺติยํ สมเณน วา พฺราหฺมเณน วา เทเวน วา มาเรน วา พฺรหฺมุนา วา เกนจิ วา โลกสฺมิํ.}

\prose{กตเมหิ ปญฺจหิ? อิธ ภิกฺขเว ตถาคโต อรหํ สมฺมาสมฺพุทฺโธ อตฺถญฺญู ธมฺมญฺญู มตฺตญฺญู กาลญฺญู ปริสญฺญู.\csromanspacer{} อิเมหิ โข ภิกฺขเว ปญฺจหิ ธมฺเมหิ สมนฺนาคโต ตถาคโต อรหํ สมฺมาสมฺพุทฺโธ ธมฺเมเนว อนุตฺตรํ ธมฺมจกฺกํ ปวตฺเตติ, ตํ โหติ ธมฺมจกฺกํ อปฺปฏิวตฺติยํ สมเณน วา พฺราหฺมเณน วา เทเวน วา มาเรน วา พฺรหฺมุนา วา เกนจิ วา โลกสฺมินฺติ.\csromanspacer{}\csromanspacer{}ปฐมํ.}

\csromanheader{2}{2. ทุติยจกฺกานุวตฺตนสุตฺต}
\proseitem{132}{ปญฺจหิ ภิกฺขเว องฺเคหิ สมนฺนาคโต รญฺโญ จกฺกวตฺติสฺส เชฏฺโฐ ปุตฺโต ปิตรา ปวตฺติตํ จกฺกํ ธมฺเมเนว อนุปฺปวตฺเตติ, ตํ โหติ จกฺกํ อปฺปฏิวตฺติยํ เกนจิ มนุสฺสภูเตน ปจฺจตฺถิเกน ปาณินา.}

\prose{กตเมหิ ปญฺจหิ? อิธ ภิกฺขเว รญฺโญ จกฺกวตฺติสฺส เชฏฺโฐ ปุตฺโต อตฺถญฺญู จ โหติ ธมฺมญฺญู จ มตฺตญฺญู จ กาลญฺญู จ ปริสญฺญู จ.\csromanspacer{} อิเมหิ โข ภิกฺขเว ปญฺจหิ องฺเคหิ สมนฺนาคโต รญฺโญ จกฺกวตฺติสฺส เชฏฺโฐ ปุตฺโต ปิตรา ปวตฺติตํ จกฺกํ ธมฺเมเนว อนุปฺปวตฺเตติ, ตํ โหติ จกฺกํ อปฺปฏิวตฺติยํ เกนจิ มนุสฺสภูเตน ปจฺจตฺถิเกน ปาณินา.}

\prose{เอวเมวํ โข ภิกฺขเว ปญฺจหิ ธมฺเมหิ สมนฺนาคโต สาริปุตฺโต ตถาคเตน อนุตฺตรํ ธมฺมจกฺกํ ปวตฺติตํ สมฺมเทว อนุปฺปวตฺเตติ, ตํ โหติ จกฺกํ อปฺปฏิวตฺติยํ สมเณน วา พฺราหฺมเณน วา เทเวน วา มาเรน วา พฺรหฺมุนา วา เกนจิ วา โลกสฺมิํ.}

\prose{กตเมหิ ปญฺจหิ? อิธ ภิกฺขเว สาริปุตฺโต อตฺถญฺญู ธมฺมญฺญู มตฺตญฺญู กาลญฺญู ปริสญฺญู.\csromanspacer{} อิเมหิ โข ภิกฺขเว ปญฺจหิ ธมฺเมหิ สมนฺนาคโต สาริปุตฺโต ตถาคเตน อนุตฺตรํ ธมฺมจกฺกํ ปวตฺติตํ สมฺมเทว อนุปฺปวตฺเตติ, ตํ โหติ จกฺกํ อปฺปฏิวตฺติยํ สมเณน วา พฺราหฺมเณน วา เทเวน วา มาเรน วา พฺรหฺมุนา วา เกนจิ วา โลกสฺมินฺติ.\csromanspacer{}\csromanspacer{}ทุติยํ.}

\csromanmatikaanchor{mtk.33}
\csromantocmarkref{subsection}{3. ธมฺมราชาสุตฺต}{mtk.33}
\bookmark[dest={mtk.33},level=2]{3. ธมฺมราชาสุตฺต}
\csromanheader{2}{3. ธมฺมราชาสุตฺต}
\proseitem{133}{\csromanfolio{132}โยปิ โส\footnote{โยปิ โข (สี, สฺยา, อิ)} ภิกฺขเว ราชา จกฺกวตฺตี ธมฺมิโก ธมฺมราชา, โสปิ น อราชกํ จกฺกํ วตฺเตตีติ.\csromanspacer{} เอวํ วุตฺเต อญฺญตโร ภิกฺขุ ภควนฺตํ เอตทโวจ “โก ปน ภนฺเต รญฺโญ จกฺกวตฺติสฺส ธมฺมิกสฺส ธมฺมรญฺโญ ราชา”ติ. “ธมฺโม ภิกฺขู”ติ ภควา อโวจ\csromandash{}}

\prose{อิธ ภิกฺขุ ราชา จกฺกวตฺตี ธมฺมิโก ธมฺมราชา ธมฺมญฺเญว นิสฺสาย ธมฺมํ สกฺกโรนฺโต ธมฺมํ ครุํ กโรนฺโต ธมฺมํ อปจายมาโน ธมฺมทฺธโช ธมฺมเกตุ ธมฺมาธิปเตยฺโย ธมฺมิกํ รกฺขาวรณคุตฺติํ สํวิทหติ อนฺโตชนสฺมิํ.}

\prose{ปุน จปรํ ภิกฺขุ ราชา จกฺกวตฺตี ธมฺมิโก ธมฺมราชา ธมฺมญฺเญว นิสฺสาย ธมฺมํ สกฺกโรนฺโต ธมฺมํ ครุํ กโรนฺโต ธมฺมํ อปจายมาโน ธมฺมทฺธโช ธมฺมเกตุ ธมฺมาธิปเตยฺโย ธมฺมิกํ รกฺขาวรณคุตฺติํ สํวิทหติ ขตฺติเยสุ อนุยนฺเตสุ\footnote{อนุยุตฺเตสุ (สี)} ฯเปฯ พลกายสฺมิํ.\csromanspacer{} พฺราหฺมณคหปติเกสุ เนคมชานปเทสุ.\csromanspacer{} สมณพฺราหฺมเณสุ มิคปกฺขีสุ.\csromanspacer{} ส โข โส ภิกฺขุ ราชา จกฺกวตฺตี ธมฺมิโก ธมฺมราชา ธมฺมญฺเญว นิสฺสาย ธมฺมํ สกฺกโรนฺโต ธมฺมํ ครุํ กโรนฺโต ธมฺมํ อปจายมาโน ธมฺมทฺธโช ธมฺมเกตุ ธมฺมาธิปเตยฺโย ธมฺมิกํ รกฺขาวรณคุตฺติํ สํวิทหิตฺวา อนฺโตชนสฺมิํ, ธมฺมิกํ รกฺขาวรณคุตฺติํ สํวิทหิตฺวา ขตฺติเยสุ, อนุยนฺเตสุ, พลกายสฺมิํ, พฺราหฺมณคหปติเกสุ เนคมชานปเทสุ, สมณพฺราหฺมเณสุ มิคปกฺขีสุ ธมฺเมเนว จกฺกํ ปวตฺเตติ, ตํ โหติ จกฺกํ อปฺปฏิวตฺติยํ เกนจิ มนุสฺส ภูเตน ปจฺจตฺถิเกน ปาณินา.}

\prose{เอวเมว โข ภิกฺขุ ตถาคโต อรหํ สมฺมาสมฺพุทฺโธ ธมฺมิโก ธมฺมราชา ธมฺมญฺเญว นิสฺสาย ธมฺมํ สกฺกโรนฺโต ธมฺมํ ครุํ กโรนฺโต ธมฺมํ อปจายมาโน ธมฺมทฺธโช ธมฺมเกตุ ธมฺมาธิปเตยฺโย ธมฺมิกํ รกฺขาวรณคุตฺติํ สํวิทหติ ภิกฺขูสุ “เอวรูปํ กายกมฺมํ เสวิตพฺพํ, เอวรูปํ กายกมฺมํ น เสวิตพฺพํ, เอวรูปํ วจีกมฺมํ เสวิตพฺพํ, เอวรูปํ วจีกมฺมํ น เสวิตพฺพํ, เอวรูปํ มโนกมฺมํ เสวิตพฺพํ, เอวรูปํ มโนกมฺมํ น เสวิตพฺพํ, เอวรูโป อาชีโว เสวิตพฺโพ, เอวรูโป อาชีโว น เสวิตพฺโพ, เอวรูโป คามนิคโม เสวิตพฺโพ, เอวรูโป คามนิคโม น เสวิตพฺโพ”ติ.}

\csromanmatikaanchor{mtk.34}
\csromantocmarkref{subsection}{4. ยสฺสํทิสสุตฺต}{mtk.34}
\bookmark[dest={mtk.34},level=2]{4. ยสฺสํทิสสุตฺต}
\csromanheader{2}{(14) 4. ราชวคฺค}
\prose{\csromanfolio{133}ปุน จปรํ ภิกฺขุ ตถาคโต อรหํ สมฺมาสมฺพุทฺโธ ธมฺมิโก ธมฺมราชา ธมฺมญฺเญว นิสฺสาย ธมฺมํ สกฺกโรนฺโต ธมฺมํ ครุํ กโรนฺโต ธมฺมํ อปจายมาโน ธมฺมทฺธโช ธมฺมเกตุ ธมฺมาธิปเตยฺโย ธมฺมิกํ รกฺขาวรณคุตฺติํ สํวิทหติ ภิกฺขุนีสุ\footnote{ภิกฺขูสุ ภิกฺขุนีสุ (สี, อิ)} ฯเปฯ อุปาสเกสุ ฯเปฯ อุปาสิกาสุ “เอวรูปํ กายกมฺมํ เสวิตพฺพํ, เอวรูปํ กายกมฺมํ น เสวิตพฺพํ.\csromanspacer{} เอวรูปํ วจีกมฺมํ เสวิตพฺพํ, เอวรูปํ วจีกมฺมํ น เสวิตพฺพํ.\csromanspacer{} เอวรูปํ มโนกมฺมํ เสวิตพฺพํ, เอวรูปํ มโนกมฺมํ น เสวิตพฺพํ.\csromanspacer{} เอวรูโป อาชีโว เสวิตพฺโพ, เอวรูโป อาชีโว น เสวิตพฺโพ.\csromanspacer{} เอวรูโป คามนิคโม เสวิตพฺโพ, เอวรูโป คามนิคโม น เสวิตพฺโพ”ติ.}

\prose{ส โข โส ภิกฺขุ ตถาคโต อรหํ สมฺมาสมฺพุทฺโธ ธมฺมิโก ธมฺมราชา ธมฺมญฺเญว นิสฺสาย ธมฺมํ สกฺกโรนฺโต ธมฺมํ ครุํ กโรนฺโต ธมฺมํ อปจายมาโน ธมฺมทฺธโช ธมฺมเกตุ ธมฺมาธิปเตยฺโย ธมฺมิกํ รกฺขาวรณคุตฺติํ สํวิทหิตฺวา ภิกฺขูสุ, ธมฺมิกํ รกฺขาวรณคุตฺติํ สํวิทหิตฺวา ภิกฺขุนีสุ, ธมฺมิกํ รกฺขาวรณคุตฺติํ สํวิทหิตฺวา อุปาสเกสุ, ธมฺมิกํ รกฺขาวรณคุตฺติํ สํวิทหิตฺวา อุปาสิกาสุ ธมฺเมเนว อนุตฺตรํ ธมฺมจกฺกํ ปวตฺเตติ, ตํ โหติ จกฺกํ อปฺปฏิวตฺติยํ สมเณน วา พฺราหฺมเณน วา เทเวน วา มาเรน วา พฺรหฺมุนา วา เกนจิ วา โลกสฺมินฺติ.\csromanspacer{}\csromanspacer{}ตติยํ.}

\csromanheader{2}{4. ยสฺสํทิสํสุตฺต}
\proseitem{134}{ปญฺจหิ ภิกฺขเว องฺเคหิ สมนฺนาคโต ราชา ขตฺติโย มุทฺธาวสิตฺโต ยสฺสํ ยสฺสํ ทิสายํ วิหรติ, สกฺกสฺมิํเยว วิชิเต วิหรติ.}

\csromanmatikaanchor{mtk.35}
\csromantocmarkref{subsection}{5-6. ปตฺถนาสุตฺต}{mtk.35}
\bookmark[dest={mtk.35},level=2]{5-6. ปตฺถนาสุตฺต}
\prose{กตเมหิ ปญฺจหิ? อิธ ภิกฺขเว ราชา ขตฺติโย มุทฺธาวสิตฺโต อุภโต สุชาโต โหติ มาติโต จ ปิติโต จ สํสุทฺธคหณิโก ยาว สตฺตมา ปิตามหยุคา อกฺขิตฺโต อนุปกฺกุฏฺโฐ ชาติวาเทน, อฑฺโฒ โหติ มหทฺธโน มหาโภโค ปริปุณฺณโกสโก ฏฺฐาคาโร, พลวา โข ปน โหติ จตุรงฺคินิยา เสนาย สมนฺนาคโต อสฺสวาย โอวาทปฏิกราย, ปริณายโก โข ปนสฺส โหติ ปณฺฑิโต วิยตฺโต เมธาวี ปฏิพโล อตีตานาคตปจฺจุปฺปนฺเน อตฺเถ จินฺเตตุํ.\csromanspacer{} ตสฺสิเม จตฺตาโร ธมฺมา ยสํ ปริปาเจนฺติ, โส อิมินา \csromanfolio{134}ยสปญฺจเมน\footnote{ยเสน ปญฺจเมน (ก), ปญฺจเมน (สี)} ธมฺเมน สมนฺนาคโต ยสฺสํ ยสฺสํ ทิสายํ วิหรติ, สกสฺมิํ เยว วิชิเต วิหรติ.\csromanspacer{} ตํ กิสฺส เหตุ, เอวํ เหตํ ภิกฺขเว โหติ วิชิตาวีนํ.}

\prose{เอวเมว โข ภิกฺขเว ปญฺจหิ ธมฺเมหิ สมนฺนาคโต ภิกฺขุ ยสฺสํ ยสฺสํ ทิสายํ วิหรติ, วิมุตฺตจิตฺโตว\footnote{วิมุตฺตจิตฺโต (สี, อิ), วิมุตฺตจิตฺโต จ (ก)} วิหรติ.\csromanspacer{} กตเมหิ ปญฺจหิ, อิธ ภิกฺขเว ภิกฺขุ สีลวา โหติ, ปาติโมกฺขสํวรสํวุโต วิหรติ อาจารโคจรสมฺปนฺโน, อณุมตฺเตสุ วชฺเชสุ ภยทสฺสาวี สมาทาย สิกฺขติ สิกฺขาปเทสุ, ราชาว ขตฺติโย มุทฺธาวสิตฺโต ชาติสมฺปนฺโน.\csromanspacer{} พหุสฺสุโต โหติ สุตธโร สุตสนฺนิจโย, เย เต ธมฺมา อาทิกลฺยาณา มชฺเฌกลฺยาณา ปริโยสานกลฺยาณา สาตฺถํ สพฺยญฺชนํ เกวลปริปุณฺณํ ปริสุทฺธํ พฺรหฺมจริยํ อภิวทนฺติ, ตถารูปาสฺส ธมฺมา พหุสฺสุตา โหนฺติ ธาตา วจสา ปริจิตา มนสานุเปกฺขิตา ทิฏฺฐิยา สุปฺปฏิวิทฺธา, ราชาว ขตฺติโย มุทฺธาว สิตฺโต อฑฺโฒ มหทฺธโน มหาโภโค ปริปุณฺณโกสโกฏฺฐาคาโร.\csromanspacer{} อารทฺธวีริโย วิหรติ อกุสลานํ ธมฺมานํ ปหานาย กุสลานํ ธมฺมานํ อุปสมฺปทาย, ถามวา ทฬฺหปรกฺกโม อนิกฺขิตฺตธุโร กุสเลสุ ธมฺเมสุ, ราชาว ขตฺติโย มุทฺธาวสิตฺโต พลสมฺปนฺโน.\csromanspacer{} ปญฺญวา โหติ อุทยตฺถคามินิยา ปญฺญาย สมนฺนาคโต อริยาย นิพฺเพธิกาย สมฺมา ทุกฺขกฺขยคามินิยา, ราชาว ขตฺติโย มุทฺธาวสิตฺโต ปริณายกสมฺปนฺโน.\csromanspacer{} ตสฺสิเม จตฺตาโร ธมฺมา วิมุตฺติํ ปริปาเจนฺติ, โส อิมินา วิมุตฺติปญฺจเมน ธมฺเมน สมนฺนาคโต ยสฺสํ ยสฺสํ ทิสายํ วิหรติ, วิมุตฺตจิตฺโตว วิหรติ.\csromanspacer{} ตํ กิสฺส เหตุ, เอวํ เหตํ ภิกฺขเว โหติ วิมุตฺตจิตฺตานนฺติ.\csromanspacer{}\csromanspacer{}จตุตฺถํ.}

\csromanheader{2}{5. ปฐมปตฺถนาสุตฺต}
\proseitem{135}{ปญฺจหิ ภิกฺขเว องฺเคหิ สมนฺนาคโต รญฺโญ ขตฺติยสฺส มุทฺธาวสิตฺตสฺส เชฏฺโฐ ปุตฺโต รชฺชํ ปตฺเถติ.\csromanspacer{} กตเมหิ ปญฺจหิ, อิธ ภิกฺขเว รญฺโญ ขตฺติยสฺส มุทฺธาวสิตฺตสฺส เชฏฺโฐ ปุตฺโต อุภโต สุชาโต โหติ มาติโต จ ปิติโต จ สํสุทฺธคหณิโก ยาว สตฺตมา ปิตามหยุคา อกฺขิตฺโต อนุปกฺกุฏฺโฐ ชาติวาเทน, อภิรูโป}

\vspace{\tipitakaparskip}
\csromancenter{(14) 4.\csromanspacer{} ราชวคฺค โหติ ทสฺสนีโย ปาสาทิโก ปรมาย วณฺณโปกฺขรตาย สมนฺนาคโต, มาตาปิตูนํ ปิโย โหติ มนาโป, เนคมชานปทสฺส ปิโย โหติ มนาโป, ยานิ ตานิ รญฺญํ ขตฺติยานํ มุทฺธาวสิตฺตานํ สิปฺปฏฺฐานานิ หตฺถิสฺมิํ วา อสฺสสฺมิํ วา รถสฺมิํ วา ธนุสฺมิํ วา ถรุสฺมิํ วา, ตตฺถ สิกฺขิโต โหติ อนวโย.}
\prose{\csromanfolio{135}ตสฺส เอวํ โหติ “อหํ โขมฺหิ อุภโต สุชาโต มาติโต จ ปิติโต จ สํสุทฺธ คหณิโก ยาว สตฺตมา ปิตามหยุคา อกฺขิตฺโต อนุปกฺกุฏฺโฐ ชาติวาเทน, กสฺมาหํ รชฺชํ น ปตฺเถยฺยํ.\csromanspacer{} อหํ โขมฺหิ อภิรูโป ทสฺสนีโย ปาสาทิโก ปรมาย วณฺณโปกฺขรตาย สมนฺนาคโต, กสฺมาหํ รชฺชํ น ปตฺเถยฺยํ.\csromanspacer{} อหํ โขมฺหิ มาตาปิตูนํ ปิโย มนาโป, กสฺมาหํ รชฺชํ น ปตฺเถยฺยํ.\csromanspacer{} อหํ โขมฺหิ เนคมชานปทสฺส ปิโย มนาโป, กสฺมาหํ รชฺชํ น ปตฺเถยฺยํ.\csromanspacer{} อหํ โขมฺหิ ยานิ ตานิ รญฺญํ ขตฺติยานํ มุทฺธาวสิตฺตานํ สิปฺปฏฺฐานานิ หตฺถิสฺมิํ วา อสฺสสฺมิํ วา รถสฺมิํ วา ธนุสฺมิํ วา ถรุสฺมิํ วา, ตตฺถ\footnote{ตตฺถมฺหิ (สี), ตตฺถปิ (ก)} สิกฺขิโต อนวโย, กสฺมาหํ รชฺชํ น ปตฺเถยฺยนฺติ.\csromanspacer{} อิเมหิ โข ภิกฺขเว ปญฺจหิ องฺเคหิ สมนฺนาคโต รญฺโญ ขตฺติยสฺส มุทฺธาวสิตฺตสฺส เชฏฺโฐ ปุตฺโต รชฺชํ ปตฺเถติ.}

\prose{เอวเมว โข ภิกฺขเว ปญฺจหิ ธมฺเมหิ สมนฺนาคโต ภิกฺขุ อาสวานํ ขยํ ปตฺเถติ.\csromanspacer{} กตเมหิ ปญฺจหิ, อิธ ภิกฺขเว ภิกฺขุ สทฺโธ โหติ, สทฺทหติ ตถาคตสฺส โพธิํ “อิติปิ โส ภควา อรหํ สมฺมาสมฺพุทฺโธ วิชฺชาจรณสมฺปนฺโน สุคโต โลกวิทู อนุตฺตโร ปุริสทมฺมสารถิ สตฺถา เทวมนุสฺสานํ พุทฺโธ ภควา”ติ.\csromanspacer{} อปฺปาพาโธ โหติ อปฺปาตงฺโก สมเวปากินิยา คหณิยา สมนฺนาคโต นาติสีตาย นาจฺจุณฺหาย มชฺฌิมาย ปธานกฺขมาย.\csromanspacer{} อสโฐ โหติ อมายาวี, ยถาภูตํ อตฺตานํ อาวิกตฺตา สตฺถริ วา วิญฺญูสุ วา สพฺรหฺมจารีสุ.\csromanspacer{} อารทฺธวีริโย วิหรติ อกุสลานํ ธมฺมานํ ปหานาย กุสลานํ ธมฺมานํ อุปสมฺปทาย, ถามวา ทฬฺหปรกฺกโม อนิกฺขิตฺตธุโร กุสเลสุ ธมฺเมสุ.\csromanspacer{} ปญฺญวา โหติ อุทยตฺถคามินิยา ปญฺญาย สมนฺนาคโต อริยาย นิพฺเพธิกาย สมฺมา ทุกฺขกฺขยคามินิยา.}

\prose{\csromanfolio{136}ตสฺส เอวํ โหติ “อหํ โขมฺหิ สทฺโธ, สทฺทหามิ ตถาคตสฺส โพธิํ ‘อิติปิ โส ภควา อรหํ สมฺมาสมฺพุทฺโธ ฯเปฯ สตฺถา เทวมนุสฺสานํ พุทฺโธ ภควา’ติ, กสฺมาหํ อาสวานํ ขยํ น ปตฺเถยฺยํ.\csromanspacer{} อหํ โขมฺหิ อปฺปาพาโธ อปฺปาตงฺโก สมเวปากินิยา คหณิยา สมนฺนาคโต นาติสีตาย นาจฺจุณฺหาย มชฺฌิมาย ปธานกฺขมาย, กสฺมาหํ อาสวานํ ขยํ น ปตฺเถยฺยํ.\csromanspacer{} อหํ โขมฺหิ อสโฐ อมายาวี, ยถาภูตํ อตฺตานํ อาวิกตฺตา สตฺถริ วา วิญฺญูสุ วา สพฺรหฺมจารีสุ, กสฺมาหํ อาสวานํ ขยํ น ปตฺเถยฺยํ.\csromanspacer{} อหํ โขมฺหิ อารทฺธวีริโย วิหรามิ อกุสลานํ ธมฺมานํ ปหานาย กุสลานํ ธมฺมานํ อุปสมฺปทาย, ถามวา ทฬฺหปรกฺกโม อนิกฺขิตฺตธุโร กุสเลสุ ธมฺเมสุ, กสฺมาหํ อาสวานํ ขยํ น ปตฺเถยฺยํ.\csromanspacer{} อหํ โขมฺหิ ปญฺญวา อุทยตฺถคามินิยา ปญฺญาย สมนฺนาคโต อริยาย นิพฺเพธิกาย สมฺมา ทุกฺขกฺขยคามินิยา, กสฺมาหํ อาสวานํ ขยํ น ปตฺเถยฺยนฺ”ติ.\csromanspacer{} อิเมหิ โข ภิกฺขเว ปญฺจหิ ธมฺเมหิ สมนฺนาคโต ภิกฺขุ อาสวานํ ขยํ ปตฺเถตีติ.\csromanspacer{}\csromanspacer{}ปญฺจมํ.}

\csromanheader{2}{6. ทุติยปตฺถนาสุตฺต}
\proseitem{136}{ปญฺจหิ ภิกฺขเว องฺเคหิ สมนฺนาคโต รญฺโญ ขตฺติยสฺส มุทฺธาว สิตฺตสฺส เชฏฺโฐ ปุตฺโต โอปรชฺชํ\footnote{อุปรชฺชํ (สฺยา, อิ, ก)} ปตฺเถติ.\csromanspacer{} กตเมหิ ปญฺจหิ? อิธ ภิกฺขเว รญฺโญ กตฺติยสฺส มุทฺธาวสิตฺตสฺส เชฏฺโฐ ปุตฺโต อุภโต สุชาโต โหติ มาติโต จ ปิติโต จ สํสุทฺธคหณิโก ยาว สตฺตมา ปิตามหยุคา อกฺขิตฺโต อนุปกฺกุฏฺโฐ ชาติวาเทน, อภิรูโป โหติ ทสฺสนีโย ปาสาทิโก ปรมาย วณฺณโปกฺขรตาย สมนฺนาคโต, มาตาปิตูนํ ปิโย โหติ มนาโป, พลกายสฺส ปิโย โหติ มนาโป, ปณฺฑิโต โหติ วิยตฺโต เมธาวี, ปฏิพโล อตีตานาคตปจฺจุปฺปนฺเน อตฺเถ จินฺเตตุํ.}

\prose{ตสฺส เอวํ โหติ “อหํ โขมฺหิ อุภโต สุชาโต มาติโต จ ปิติโต จ สํสุทฺธ คหณิโก ยาว สตฺตมา ปิตามหยุคา อกฺขิตฺโต อนุปกฺกุฏฺโฐ ชาติวาเทน, กสฺมาหํ โอปรชฺชํ น ปตฺเถยฺยํ.\csromanspacer{} อหํ โขมฺหิ อภิรูโป ทสฺสนีโย ปาสาทิโก ปรมาย วณฺณโปกฺขรตาย สมนฺนาคโต, กสฺมาหํ โอปรชฺชํ น ปตฺเถยฺยํ.\csromanspacer{} อหํ โขมฺหิ มาตาปิตูนํ}

\vspace{\tipitakaparskip}
\csromancenter{(14) 4.\csromanspacer{} ราชวคฺค ปิโย มนาโป, กสฺมาหํ โอปรชฺชํ น ปตฺเถยฺยํ.\csromanspacer{} อหํ โขมฺหิ พลกายสฺส ปิโย มนาโป, กสฺมาหํ โอปรชฺชํ น ปตฺเถยฺยํ.\csromanspacer{} อหํ โขมฺหิ ปณฺฑิโต วิยตฺโต เมธาวี ปฏิพโล อตีตานาคตปจฺจุปฺปนฺเน อตฺเถ จินฺเตตุํ, กสฺมาหํ โอปรชฺชํ น ปตฺเถยฺยนฺ”ติ.\csromanspacer{} อิเมหิ โข ภิกฺขเว ปญฺจหิ องฺเคหิ สมนฺนาคโต รญฺโญ กตฺติยสฺส มุทฺธาวสิตฺตสฺส เชฏฺโฐ ปุตฺโต โอปรชฺชํ ปตฺเถติ.}
\prose{\csromanfolio{137}เอวเมวํ โข ภิกฺขเว ปญฺจหิ ธมฺเมหิ สมนฺนาคโต ภิกฺขุ อาสวานํ ขยํ ปตฺเถติ.\csromanspacer{} กตเมหิ ปญฺจหิ? อิธ ภิกฺขเว ภิกฺขุ สีลวา โหติ ฯเปฯ สมาทาย สิกฺขติ สิกฺขาปเทสุ.\csromanspacer{} พหุสฺสุโต โหติ ฯเปฯ ทิฏฺฐิยา สุปฺปฏิวิทฺธา.\csromanspacer{} จตูสุ สติปฏฺฐาเนสุ สุปฺปติฏฺฐิตจิตฺโต\footnote{สุปฏฺฐิตจิตฺโต (สี, สฺยา), สูปฏฺฐิตจิตฺโต (ก)} โหติ.\csromanspacer{} อารทฺธวีริโย วิหรติ อกุสลานํ ธมฺมานํ ปหานาย กุสลานํ ธมฺมานํ อุปสมฺปทาย ถามวา ทฬฺหปรกฺกโม อนิกฺขิตฺต ธุโร กุสเลสุ ธมฺเมสุ.\csromanspacer{} ปญฺญวา โหติ อุทยตฺถคามินิยา ปญฺญาย สมนฺนาคโต อริยาย นิพฺเพธิกาย สมฺมา ทุกฺขกฺขยคามินิยา.}

\prose{ตสฺส เอวํ โหติ “อหํ โขมฺหิ สีลวา ปาติโมกฺขสํวรสํวุโต วิหรามิ อาจารโคจรสมฺปนฺโน, อณุมตฺเตสุ วชฺเชสุ ภยทสฺสาวี สมาทาย สิกฺขามิ สิกฺขาปเทสุ, กสฺมาหํ อาสวานํ ขยํ น ปตฺเถยฺยํ.\csromanspacer{} อหํ โขมฺหิ พหุสฺสุโต สุตธโร สุตสนฺนิจโย, เย เต ธมฺมา อาทิกลฺยาณา มชฺเฌกลฺยาณา ปริโยสานกลฺยาณา สาตฺถํ สพฺยญฺชนํ\footnote{สตฺถา พฺยญฺชนา (สี)} เกวลปริปุณฺณํ ปริสุทฺธํ พฺรหฺมจริยํ อภิวทนฺติ, ตถารูปา เม ธมฺมา พหุสฺสุตา โหนฺติ ธาตา วจสา ปริจิตา มนสานุเปกฺขิตา ทิฏฺฐิยา สุปฺปฏิวิทฺธา, กสฺมาหํ อาสวานํ ขยํ น ปตฺเถยฺยํ.\csromanspacer{} อหํ โขมฺหิ จตูสุ สติปฏฺฐาเนสุ สุปฺปติฏฺฐิตจิตฺโต, กสฺมาหํ อาสวานํ ขยํ น ปตฺเถยฺยํ.\csromanspacer{} อหํ โขมฺหิ อารทฺธวีริโย วิหรามิ อกุสลานํ ธมฺมานํ ปหานาย กุสลานํ ธมฺมานํ อุปสมฺปทาย ถามวา ทฬฺหปรกฺกโม อนิกฺขิตฺตธุโร กุสเลสุ ธมฺเมสุ, กสฺมาหํ อาสวานํ ขยํ น ปตฺเถยฺยํ.\csromanspacer{} อหํ โขมฺหิ ปญฺญวา อุทยตฺถคามินิยา ปญฺญาย สมนฺนาคโต อริยาย นิพฺเพธิกาย สมฺมา ทุกฺขกฺขยคามินิยา, กสฺมาหํ อาสวานํ ขยํ น ปตฺเถยฺยนฺ”ติ.\csromanspacer{} อิเมหิ โข ภิกฺขเว ปญฺจหิ ธมฺเมหิ สมนฺนาคโต ภิกฺขุ อาสวานํ ขยํ ปตฺเถตีติ.\csromanspacer{}\csromanspacer{}ฉฏฺฐํ.}

\csromanmatikaanchor{mtk.36}
\csromantocmarkref{subsection}{7. อปฺปํสุปติสุตฺต}{mtk.36}
\bookmark[dest={mtk.36},level=2]{7. อปฺปํสุปติสุตฺต}
\csromanheader{2}{7. อปฺปํสุปติสุตฺต}
\proseitem{137}{\csromanfolio{138}ปญฺจิเม ภิกฺขเว อปฺปํ รตฺติยา สุปนฺติ, พหุํ ชคฺคนฺติ.\csromanspacer{} กตเม ปญฺจ? อิตฺถี ภิกฺขเว ปุริสาธิปฺปายา อปฺปํ รตฺติยา สุปติ, พหุํ ชคฺคติ.\csromanspacer{} ปุริโส ภิกฺขเว อิตฺถาธิปฺปาโย อปฺปํ รตฺติยา สุปติ, พหุํ ชคฺคติ.\csromanspacer{} โจโร ภิกฺขเว อาทานาธิปฺปาโย อปฺปํ รตฺติยา สุปติ, พหุํ ชคฺคติ.\csromanspacer{} ราชา\footnote{ราชยุตฺโต (อิ, ก)} ภิกฺขเว ราชกรณีเยสุ ยุตฺโต อปฺปํ รตฺติยา สุปติ, พหุํ ชคฺคติ.\csromanspacer{} ภิกฺขุ ภิกฺขเว วิสํโยคาธิปฺปาโย อปฺปํ รตฺติยา สุปติ, พหุํ ชคฺคติ.\csromanspacer{} อิเม โข ภิกฺขเว ปญฺจ อปฺปํ รตฺติยา สุปนฺติ, พหุํ ชคฺคนฺตีติ.\csromanspacer{}\csromanspacer{}สตฺตมํ.}

\csromanmatikaanchor{mtk.37}
\csromantocmarkref{subsection}{8. ภตฺตาทกสุตฺต}{mtk.37}
\bookmark[dest={mtk.37},level=2]{8. ภตฺตาทกสุตฺต}
\csromanheader{2}{8. ภตฺตาทกสุตฺต}
\proseitem{138}{ปญฺจหิ ภิกฺขเว องฺเคหิ สมนฺนาคโต รญฺโญ นาโค ภตฺตาทโก จ โหติ โอกาสผรโณ จ ลณฺฑสารโณ จ สลากคฺคาหี จ, รญฺโญ นาโคเตฺวว สงฺขํ คจฺฉติ.\csromanspacer{} กตเมหิ ปญฺจหิ? อิธ ภิกฺขเว รญฺโญ นาโค อกฺขโม โหติ รูปานํ, อกฺขโม สทฺทานํ, อกฺขโม คนฺธานํ, อกฺขโม รสานํ, อกฺขโม โผฏฺฐพฺพานํ.\csromanspacer{} อิเมหิ โข ภิกฺขเว ปญฺจหิ องฺเคหิ สมนฺนาคโต รญฺโญ นาโค ภตฺตาทโก จ โอกาสผรโณ จ ลณฺฑสารโณ จ สลากคฺคาหี จ, รญฺโญ นาโคเตฺวว สงฺขํ คจฺฉติ.}

\prose{เอวเมวํ โข ภิกฺขเว ปญฺจหิ ธมฺเมหิ สมนฺนาคโต ภิกฺขุ ภตฺตาทโก จ โหติ โอกาสผรโณ จ มญฺจปีฐมทฺทโน\footnote{ปีฐมทฺทโน (สี, สฺยา, กํ, อิ)} จ สลากคฺคาหี จ, ภิกฺขุเตฺวว สงฺขํ คจฺฉติ.\csromanspacer{} กตเมหิ ปญฺจหิ? อิธ ภิกฺขเว ภิกฺขุ อกฺขโม โหติ รูปานํ, อกฺขโม สทฺทานํ, อกฺขโม คนฺธานํ, อกฺขโม รสานํ, อกฺขโม โผฏฺฐพฺพานํ.\csromanspacer{} อิเมหิ โข ภิกฺขเว ปญฺจหิ ธมฺเมหิ สมนฺนาคโต ภิกฺขุ ภตฺตาทโก จ โหติ โอกาสผรโณ จ มญฺจปีฐมทฺทโน จ สลากคฺคาหี จ, ภิกฺขุเตฺวว สงฺขํ คจฺฉตีติ.\csromanspacer{}\csromanspacer{}อฏฺฐมํ.}

\csromanmatikaanchor{mtk.38}
\csromantocmarkref{subsection}{9. อกฺขมสุตฺต}{mtk.38}
\bookmark[dest={mtk.38},level=2]{9. อกฺขมสุตฺต}
\csromanheader{2}{9. อกฺขมสุตฺต}
\proseitem{139}{ปญฺจหิ ภิกฺขเว องฺเคหิ สมนฺนาคโต รญฺโญ นาโค น ราชารโห โหติ น ราชโภคฺโค, น รญฺโญ องฺคํเตฺวว สงฺขํ}

\vspace{\tipitakaparskip}
\csromancenter{(14) 4.\csromanspacer{} ราชวคฺค คจฺฉติ.\csromanspacer{} กตเมหิ ปญฺจหิ? อิธ ภิกฺขเว รญฺโญ นาโค อกฺขโม โหติ รูปานํ, อกฺขโม สทฺทานํ, อกฺขโม คนฺธานํ, อกฺขโม รสานํ, อกฺขโม โผฏฺฐพฺพานํ.}
\prose{\csromanfolio{139}กถญฺจ ภิกฺขเว รญฺโญ นาโค อกฺขโม โหติ รูปานํ? อิธ ภิกฺขเว รญฺโญ นาโค สงฺคามคโต หตฺถิกายํ วา ทิสฺวา อสฺสกายํ วา ทิสฺวา รถกายํ วา ทิสฺวา ปตฺติกายํ วา ทิสฺวา สํสีทติ วิสีทติ น สนฺถมฺภติ น สกฺโกติ สงฺคามํ โอตริตุํ.\csromanspacer{} เอวํ โข ภิกฺขเว รญฺโญ นาโค อกฺขโม โหติ รูปานํ. (1)}

\prose{กถญฺจ ภิกฺขเว รญฺโญ นาโค อกฺขโม โหติ สทฺทานํ? อิธ ภิกฺขเว รญฺโญ นาโค สงฺคามคโต หตฺถิสทฺทํ วา สุตฺวา อสฺสสทฺทํ วา สุตฺวา รถสทฺทํ วา สุตฺวา ปตฺติสทฺทํ วา สุตฺวา เภริปณวสงฺขติณวนินฺนาทสทฺทํ วา สุตฺวา สํสีทติ วิสีทติ น สนฺถมฺภติ น สกฺโกติ สงฺคามํ โอตริตุํ.\csromanspacer{} เอวํ โข ภิกฺขเว รญฺโญ นาโค อกฺขโม โหติ สทฺทานํ. (2)}

\prose{กถญฺจ ภิกฺขเว รญฺโญ นาโค อกฺขโม โหติ คนฺธานํ? อิธ ภิกฺขเว รญฺโญ นาโค สงฺคามคโต เย เต รญฺโญ นาคา อภิชาตา สงฺคามาวจรา เตสํ มุตฺตกรีสสฺส คนฺธํ ฆายิตฺวา สํสีทติ วิสีทติ น สนฺถมฺภติ น สกฺโกติ สงฺคามํ โอตริตุํ.\csromanspacer{} เอวํ โข ภิกฺขเว รญฺโญ นาโค อกฺขโม โหติ คนฺธานํ. (3)}

\prose{กถญฺจ ภิกฺขเว รญฺโญ นาโค อกฺขโม โหติ รสานํ? อิธ ภิกฺขเว รญฺโญ นาโค สงฺคามคโต เอกิสฺสา วา ติโณทกทตฺติยา วิมานิโต\footnote{วิหนีโต (สฺยา), วิหานิโต (กตฺถจิ)} ทฺวีหิ วา ตีหิ วา จตูหิ วา ปญฺจหิ วา ติโณทกทตฺตีหิ วิมานิโต สํสีทติ วิสีทติ น สนฺถมฺภติ น สกฺโกติ สงฺคามํ โอตริตุํ.\csromanspacer{} เอวํ โข ภิกฺขเว รญฺโญ นาโค อกฺขโม โหติ รสานํ. (4)}

\prose{กถญฺจ ภิกฺขเว รญฺโญ นาโค อกฺขโม โหติ โผฏฺฐพฺพานํ? อิธ ภิกฺขเว รญฺโญ นาโค สงฺคามคโต เอเกน วา สรเวเคน วิทฺโธ ทฺวีหิ วา ตีหิ วา จตูหิ วา ปญฺจหิ วา สรเวเคหิ วิทฺโธ สํสีทติ วิสีทติ \csromanfolio{140}น สนฺถมฺภติ น สกฺโกติ สงฺคามํ โอตริตุํ.\csromanspacer{} เอวํ โข ภิกฺขเว รญฺโญ นาโค อกฺขโม โหติ โผฏฺฐพฺพานํ. (5)}

\prose{อิเมหิ โข ภิกฺขเว ปญฺจหิ องฺเคหิ สมนฺนาคโต รญฺโญ นาโค น ราชารโห โหติ น ราชโภคฺโค, น รญฺโญ องฺคํเตฺวว สงฺขํ คจฺฉติ.}

\prose{เอวเมวํ โข ภิกฺขเว ปญฺจหิ องฺเคหิ สมนฺนาคโต ภิกฺขุ น อาหุเนยฺโย โหติ น ปาหุเนยฺโย น ทกฺขิเณยฺโย น อญฺชลิกรณีโย น อนุตฺตรํ ปุญฺญกฺเขตฺตํ โลกสฺส.\csromanspacer{} กตเมหิ ปญฺจหิ? อิธ ภิกฺขเว ภิกฺขุ อกฺขโม โหติ รูปานํ, อกฺขโม สทฺทานํ, อกฺขโม คนฺธานํ, อกฺขโม รสานํ, อกฺขโม โผฏฺฐพฺพานํ.}

\prose{กถญฺจ ภิกฺขเว ภิกฺขุ อกฺขโม โหติ รูปานํ? อิธ ภิกฺขเว ภิกฺขุ จกฺขุนา รูปํ ทิสฺวา รชนีเย รูเป สารชฺชติ, น สกฺโกติ จิตฺตํ สมาทหิตุํ.\csromanspacer{} เอวํ โข ภิกฺขเว ภิกฺขุ อกฺขโม โหติ รูปานํ. (1)}

\prose{กถญฺจ ภิกฺขเว ภิกฺขุ อกฺขโม โหติ สทฺทานํ? อิธ ภิกฺขเว ภิกฺขุ โสเตน สทฺทํ สุตฺวา รชนีเย สทฺเท สารชฺชติ, น สกฺโกติ จิตฺตํ สมาทหิตุํ.\csromanspacer{} เอวํ โข ภิกฺขเว ภิกฺขุ อกฺขโม โหติ สทฺทานํ. (2)}

\prose{กถญฺจ ภิกฺขเว ภิกฺขุ อกฺขโม โหติ คนฺธานํ? อิธ ภิกฺขเว ภิกฺขุ ฆาเนน คนฺธํ ฆายิตฺวา รชนีเย คนฺเธ สารชฺชติ, น สกฺโกติ จิตฺตํ สมาทหิตุํ.\csromanspacer{} เอวํ โข ภิกฺขเว ภิกฺขุ อกฺขโม โหติ คนฺธานํ. (3)}

\prose{กถญฺจ ภิกฺขเว ภิกฺขุ อกฺขโม โหติ รสานํ? อิธ ภิกฺขเว ภิกฺขุ ชิวฺหาย รสํ สายิตฺวา รชนีเย รเส สารชฺชติ, น สกฺโกติ จิตฺตํ สมาทหิตุํ.\csromanspacer{} เอวํ โข ภิกฺขเว ภิกฺขุ อกฺขโม โหติ รสานํ. (4)}

\prose{กถญฺจ ภิกฺขเว ภิกฺขุ อกฺขโม โหติ โผฏฺฐพฺพานํ? อิธ ภิกฺขเว ภิกฺขุ กาเยน โผฏฺฐพฺพํ ผุสิตฺวา รชนีเย โผฏฺฐพฺเพ สารชฺชติ, น สกฺโกติ จิตฺตํ สมาทหิตุํ.\csromanspacer{} เอวํ โข ภิกฺขเว ภิกฺขุ อกฺขโม โหติ โผฏฺฐพฺพานํ. (5)}

\prose{อิเมหิ โข ภิกฺขเว ปญฺจหิ ธมฺเมหิ สมนฺนาคโต ภิกฺขุ น อาหุเนยฺโย โหติ น ปาหุเนยฺโย น ทกฺขิเณยฺโย น อญฺชลิกรณีโย น อนุตฺตรํ ปุญฺญกฺเขตฺตํ โลกสฺส.}

\chapterheadpageread{(14) 4. ราชวคฺค}
\prose{\csromanfolio{141}ปญฺจหิ ภิกฺขเว องฺเคหิ สมนฺนาคโต รญฺโญ นาโค ราชารโห โหติ ราชโภคฺโค, รญฺโญ องฺคํเตฺวว สงฺขํ คจฺฉติ.\csromanspacer{} กตเมหิ ปญฺจหิ? อิธ ภิกฺขเว รญฺโญ นาโค ขโม โหติ รูปานํ, ขโม สทฺทานํ, ขโม คนฺธานํ, ขโม รสานํ, ขโม โผฏฺฐพฺพานํ.}

\prose{กถญฺจ ภิกฺขเว รญฺโญ นาโค ขโม โหติ รูปานํ? อิธ ภิกฺขเว รญฺโญ นาโค สงฺคามคโต หตฺถิกายํ วา ทิสฺวา อสฺสกายํ วา ทิสฺวา รถกายํ วา ทิสฺวา ปตฺติกายํ วา ทิสฺวา น สํสีทติ น วิสีทติ สนฺถมฺภติ สกฺโกติ สงฺคามํ โอตริตุํ.\csromanspacer{} เอวํ โข ภิกฺขเว รญฺโญ นาโค ขโม โหติ รูปานํ. (1)}

\prose{กถญฺจ ภิกฺขเว รญฺโญ นาโค ขโม โหติ สทฺทานํ? อิธ ภิกฺขเว รญฺโญ นาโค สงฺคามคโต หตฺถิสทฺทํ วา สุตฺวา อสฺสสทฺทํ วา สุตฺวา รถสทฺทํ วา สุตฺวา ปตฺติสทฺทํ วา สุตฺวา เภริปณวสงฺขติณวนินฺนาทสทฺทํ วา สุตฺวา น สํสีทติ น วิสีทติ สนฺถมฺภติ สกฺโกติ สงฺคามํ โอตริตุํ.\csromanspacer{} เอวํ โข ภิกฺขเว รญฺโญ นาโค ขโม โหติ สทฺทานํ. (2)}

\prose{กถญฺจ ภิกฺขเว รญฺโญ นาโค ขโม โหติ คนฺธานํ? อิธ ภิกฺขเว รญฺโญ นาโค สงฺคามคโต เย เต รญฺโญ นาคา อภิชาตา สงฺคามาวจรา, เตสํ มุตฺตกรีสสฺส คนฺธํ ฆายิตฺวา น สํสีทติ น วิสีทติ สนฺถมฺภติ สกฺโกติ สงฺคามํ โอตริตุํ.\csromanspacer{} เอวํ โข ภิกฺขเว รญฺโญ นาโค ขโม โหติ คนฺธานํ. (3)}

\prose{กถญฺจ ภิกฺขเว รญฺโญ นาโค ขโม โหติ รสานํ? อิธ ภิกฺขเว รญฺโญ นาโค สงฺคามคโต เอกิสฺสา วา ติโณทกทตฺติยา วิมานิโต ทฺวีหิ วา ตีหิ วา จตูหิ วา ปญฺจหิ วา ติโณทกทตฺตีหิ วิมานิโต น สํสีทติ น วิสีทติ สนฺถมฺภติ สกฺโกติ สงฺคามํ โอตริตุํ.\csromanspacer{} เอวํ โข ภิกฺขเว รญฺโญ นาโค ขโม โหติ รสานํ. (4)}

\prose{กถญฺจ ภิกฺขเว รญฺโญ นาโค ขโม โหติ โผฏฺฐาพฺพานํ? อิธ ภิกฺขเว รญฺโญ นาโค สงฺคามคโต เอเกน วา สรเวเคน วิทฺโธ ทฺวีหิ วา ตีหิ วา จตูหิ วา ปญฺจหิ วา สรเวเคหิ วิทฺโธ น สํสีทติ น วิสีทติ สนฺถมฺภติ สกฺโกติ สงฺคามํ โอตริตุํ.\csromanspacer{} เอวํ โข ภิกฺขเว รญฺโญ นาโค ขโม โหติ โผฏฺฐาพฺพานํ. (5)}

\csromanmatikaanchor{mtk.39}
\csromantocmarkref{subsection}{10. โสตาสุตฺต}{mtk.39}
\bookmark[dest={mtk.39},level=2]{10. โสตาสุตฺต}
\prose{\csromanfolio{142}อิเมหิ โข ภิกฺขเว ปญฺจหิ องฺเคหิ สมนฺนาคโต รญฺโญ นาโค ราชารโห โหติ ราชโภคฺโค, รญฺโญ องฺคํเตฺวว สงฺขํ คจฺฉติ.}

\prose{เอวเมวํ โข ภิกฺขเว ปญฺจหิ ธมฺเมหิ สมนฺนาคโต ภิกฺขุ อาหุเนยฺโย โหติ ปาหุเนยฺโย ทกฺขิเณยฺโย อญฺชลิกรณีโย อนุตฺตรํ ปุญฺญกฺเขตฺตํ โลกสฺส.\csromanspacer{} กตเมหิ ปญฺจหิ? อิธ ภิกฺขเว ภิกฺขุ ขโม โหติ รูปานํ, ขโม สทฺทานํ, ขโม คนฺธานํ, ขโม รสานํ, ขโม โผฏฺฐพฺพานํ.}

\prose{กถญฺจ ภิกฺขเว ภิกฺขุ ขโม โหติ รูปานํ? อิธ ภิกฺขเว ภิกฺขุ จกฺขุนา รูปํ ทิสฺวา รชนีเย รูเป น สารชฺชติ, สกฺโกติ จิตฺตํ สมาทหิตุํ.\csromanspacer{} เอวํ โข ภิกฺขเว ภิกฺขุ ขโม โหติ รูปานํ. (1)}

\prose{กถญฺจ ภิกฺขเว ภิกฺขุ ขโม โหติ สทฺทานํ? อิธ ภิกฺขเว ภิกฺขุ โสเตน สทฺทํ สุตฺวา รชนีเย สทฺเท น สารชฺชติ, สกฺโกติ จิตฺตํ สมาทหิตุํ.\csromanspacer{} เอวํ โข ภิกฺขเว ภิกฺขุ ขโม โหติ สทฺทานํ. (2)}

\prose{กถญฺจ ภิกฺขเว ภิกฺขุ ขโม โหติ คนฺธานํ? อิธ ภิกฺขเว ภิกฺขุ ฆาเนน คนฺธํ ฆายิตฺวา รชนีเย คนฺเธ น สารชฺชติ, สกฺโกติ จิตฺตํ สมาทหิตุํ.\csromanspacer{} เอวํ โข ภิกฺขเว ภิกฺขุ ขโม โหติ คนฺธานํ. (3)}

\prose{กถญฺจ ภิกฺขเว ภิกฺขุ ขโม โหติ รสานํ? อิธ ภิกฺขเว ภิกฺขุ ชิวฺหาย รสํ สายิตฺวา รชนีเย รเส น สารชฺชติ, สกฺโกติ จิตฺตํ สมาทหิตุํ.\csromanspacer{} เอวํ โข ภิกฺขเว ภิกฺขุ ขโม โหติ รสานํ. (4)}

\prose{กถญฺจ ภิกฺขเว ภิกฺขุ ขโม โหติ โผฏฺฐพฺพานํ? อิธ ภิกฺขเว ภิกฺขุ กาเยน โผฏฺฐพฺพํ ผุสิตฺวา รชนีเย โผฏฺฐพฺเพ น สารชฺชติ, สกฺโกติ จิตฺตํ สมาทหิตุํ.\csromanspacer{} เอวํ โข ภิกฺขเว ภิกฺขุ ขโม โหติ โผฏฺฐพฺพานํ. (5)}

\prose{อิเมหิ โข ภิกฺขเว ปญฺจหิ ธมฺเมหิ สมนฺนาคโต ภิกฺขุ อาหุเนยฺโย โหติ ปาหุเนยฺโย ทกฺขิเณยฺโย อญฺชลิกรณีโย อนุตฺตรํ ปุญฺญกฺเขตฺตํ โลกสฺสาติ.\csromanspacer{}\csromanspacer{}นวมํ.}

\csromanheader{2}{10. โสตสุตฺต}
\proseitem{140}{ปญฺจหิ ภิกฺขเว องฺเคหิ สมนฺนาคโต รญฺโญ นาโค ราชารโห โหติ ราชโภคฺโค, รญฺโญ องฺคํเตฺวว สงฺขํ คจฺฉติ.}

\vspace{\tipitakaparskip}
\csromancenter{(14) 4.\csromanspacer{} ราชวคฺค กตเมหิ ปญฺจหิ? อิธ ภิกฺขเว รญฺโญ นาโค โสตา จ โหติ หนฺตา จ รกฺขิตา จ ขนฺตา จ คนฺตา จ.}
\prose{\csromanfolio{143}กถญฺจ ภิกฺขเว รญฺโญ นาโค โสตา โหติ? อิธ ภิกฺขเว รญฺโญ นาโค ยเมนํ หตฺถิทมฺมสารถิ\footnote{หตฺถิทมฺมสารถี (สี)} การณํ กาเรติ ยทิ วา กตปุพฺพํ ยทิ วา อกตปุพฺพํ, ตํ อฏฺฐิํ กตฺวา\footnote{อฏฺฐิกตฺวา (สี, สฺยา, กํ, อิ) อํ 1. 431 ปิฏฺเฐปิ.} มนสิ กตฺวา สพฺพํ เจตสา\footnote{สพฺพเจตสา (?)} สมนฺนาหริตฺวา โอหิตโสโต สุณาติ.\csromanspacer{} เอวํ โข ภิกฺขเว รญฺโญ นาโค โสตา โหติ. (1)}

\prose{กถญฺจ ภิกฺขเว รญฺโญ นาโค หนฺตา โหติ? อิธ ภิกฺขเว รญฺโญ นาโค สงฺคามคโต หตฺถิมฺปิ หนติ\footnote{หนฺติ (สี, อิ)}, หตฺถารุหมฺปิ หนติ, อสฺสมฺปิ หนติ, อสฺสารุหมฺปิ หนติ, รถมฺปิ หนติ, รถิกมฺปิ\footnote{รถารุหมฺปิ (อิ)} หนติ, ปตฺติกมฺปิ หนติ.\csromanspacer{} เอวํ โข ภิกฺขเว รญฺโญ นาโค หนฺตา โหติ. (2)}

\prose{กถญฺจ ภิกฺขเว รญฺโญ นาโค รกฺขิตา โหติ? อิธ ภิกฺขเว รญฺโญ นาโค สงฺคามคโต รกฺขติ ปุริมํ กายํ, รกฺขติ ปจฺฉิมํ กายํ, รกฺขติ ปุริเม ปาเท, รกฺขติ ปจฺฉิเม ปาเท, รกฺขติ สีสํ, รกฺขติ กณฺเณ, รกฺขติ ทนฺเต, รกฺขติ โสณฺฑํ, รกฺขติ วาลธิํ, รกฺขติ หตฺถารุหํ.\csromanspacer{} เอวํ โข ภิกฺขเว รญฺโญ นาโค รกฺขิตา โหติ. (3)}

\prose{กถญฺจ ภิกฺขเว รญฺโญ นาโค ขนฺตา โหติ? อิธ ภิกฺขเว รญฺโญ นาโค สงฺคามคโต ขโม โหติ สตฺติปฺปหารานํ อสิปฺปหารานํ อุสุปฺปหารานํ ผรสุปฺปหารานํ เภริปณวสงฺขติณวนินฺนาทสทฺทานํ.\csromanspacer{} เอวํ โข ภิกฺขเว รญฺโญ นาโค ขนฺตา โหติ. (4)}

\prose{กถญฺจ ภิกฺขเว รญฺโญ นาโค คนฺตา โหติ? อิธ ภิกฺขเว รญฺโญ นาโค ยเมนํ หตฺถิทมฺมสารถิ ทิสํ เปเสติ ยทิ วา คตปุพฺพํ, ยทิ วา อคตปุพฺพํ, ตํ ขิปฺปเมว คนฺตา โหติ.\csromanspacer{} เอวํ โข ภิกฺขเว รญฺโญ นาโค คนฺตา โหติ. (5)}

\prose{อิเมหิ โข ภิกฺขเว ปญฺจหิ องฺเคหิ สมนฺนาคโต รญฺโญ นาโค ราชารโห โหติ ราชโภคฺโค, รญฺโญ องฺคํเตฺวว สงฺขํ คจฺฉติ.}

\prose{\csromanfolio{144}เอวเมวํ โข ภิกฺขเว ปญฺจหิ ธมฺเมหิ สมนฺนาคโต ภิกฺขุ อาหุเนยฺโย โหติ ปาหุเนยฺโย ทกฺขิเณยฺโย อญฺชลิกรณีโย อนุตฺตรํ ปุญฺญกฺเขตฺตํ โลกสฺส.\csromanspacer{} กตเมหิ ปญฺจหิ? อิธ ภิกฺขเว ภิกฺขุ โสตา จ โหติ หนฺตา จ รกฺขิตา จ ขนฺตา จ คนฺตา จ.}

\prose{กถญฺจ ภิกฺขเว ภิกฺขุ โสตา โหติ? อิธ ภิกฺขเว ภิกฺขุ ตถาคตปฺปเวทิเต ธมฺมวินเย เทสิยมาเน อฏฺฐิํกตฺวา มนสิ กตฺวา สพฺพํ เจตสา สมนฺนา หริตฺวา โอหิตโสโต ธมฺมํ สุณาติ.\csromanspacer{} เอวํ โข ภิกฺขเว ภิกฺขุ โสตา โหติ. (1)}

\prose{กถญฺจ ภิกฺขเว ภิกฺขุ หนฺตา โหติ? อิธ ภิกฺขเว ภิกฺขุ อุปฺปนฺนํ กามวิตกฺกํ นาธิวาเสติ ปชหติ วิโนเทติ (หนติ)\footnote{( ) นตฺถิ สี-อิ-โปตฺถเกสุ; อํ 1. 432 ปิฏฺเฐปิ.} พฺยนฺตีกโรติ อนภาวํ คเมติ, อุปฺปนฺนํ พฺยาปาทวิตกฺกํ ฯเปฯ อุปฺปนฺนํ วิหิํสาวิตกฺกํ ฯเปฯ อุปฺปนฺนุปฺปนฺเน ปาปเก อกุสเล ธมฺเม นาธิวาเสติ ปชหติ วิโนเทติ (หนติ)1 พฺยนฺตีกโรติ อนภาวํ คเมติ.\csromanspacer{} เอวํ โข ภิกฺขเว ภิกฺขุ หนฺตา โหติ. (2)}

\prose{กถญฺจ ภิกฺขเว ภิกฺขุ รกฺขิตา โหติ? อิธ ภิกฺขเว ภิกฺขุ จกฺขุนา รูปํ ทิสฺวา น นิมิตฺตคฺคาหี โหติ นานุพฺยญฺชนคฺคาหี, ยตฺวาธิกรณเมนํ จกฺขุนฺทฺริยํ อสํวุตํ วิหรนฺตํ อภิชฺฌาโทมนสฺสา ปาปกา อกุสลา ธมฺมา อนฺวาสฺสเวยฺยุํ.\csromanspacer{} ตสฺส สํวราย ปฏิปชฺชติ, รกฺขติ จกฺขุนฺทฺริยํ, จกฺขุนฺทฺริเย สํวรํ อาปชฺชติ.\csromanspacer{} โสเตน สทฺทํ สุตฺวา.\csromanspacer{} ฆาเนน คนฺธํ ฆายิตฺวา.\csromanspacer{} ชิวฺหาย รสํ สายิตฺวา.\csromanspacer{} กาเยน โผฏฺฐพฺพํ ผุสิตฺวา.\csromanspacer{} มนสา ธมฺมํ วิญฺญาย น นิมิตฺตคฺคาหี โหติ นานุพฺยญฺชนคฺคาหี, ยตฺวาธิกรณเมนํ มนินฺทฺริยํ อสํวุตํ วิหรนฺตํ อภิชฺฌาโทมนสฺสา ปาปกา อกุสลา ธมฺมา อนฺวาสฺสเวยฺยุํ.\csromanspacer{} ตสฺส สํวราย ปฏิปชฺชติ, รกฺขติ มนินฺทฺริยํ, มนินฺทฺริเย สํวรํ อาปชฺชติ.\csromanspacer{} เอวํ โข ภิกฺขเว ภิกฺขุ รกฺขิตา โหติ. (3)}

\prose{กถญฺจ ภิกฺขเว ภิกฺขุ ขนฺตา โหติ? อิธ ภิกฺขเว ภิกฺขุ ขโม โหติ สีตสฺส อุณฺหสฺส ชิฆจฺฉาย ปิปาสาย ฑํสมกสวาตาตปสรีสป2สมฺผสฺสานํ ทุรุตฺตานํ ทุราคตานํ วจนปถานํ, อุปฺปนฺนานํ สารีริกานํ เวทนานํ ทุกฺขานํ ติพฺพานํ ขรานํ กฏุกานํ อสาตานํ อมนาปานํ}

\vspace{\tipitakaparskip}
\csromancenter{\textbf{(15) 5.}\csromanspacer{}\textbf{ ติกณฺฑกีวคฺค} ปาณหรานํ อธิวาสกชาติโก โหติ.\csromanspacer{} เอวํ โข ภิกฺขเว ภิกฺขุ ขนฺตา โหติ. (4)}
\prose{\csromanfolio{145}กถญฺจ ภิกฺขเว ภิกฺขุ คนฺตา โหติ? อิธ ภิกฺขเว ภิกฺขุ ยา สา ทิสา อคตปุพฺพา อิมินา ทีเฆน อทฺธุนา ยทิทํ สพฺพสงฺขารสมโถ สพฺพูปธิ ปฏินิสฺสคฺโค ตณฺหากฺขโย วิราโค นิโรโธ นิพฺพานํ, ตํ ขิปฺปญฺเญว คนฺตา โหติ.\csromanspacer{} เอวํ โข ภิกฺขเว ภิกฺขุ คนฺตา โหติ. (5)}

\prose{อิเมหิ โข ภิกฺขเว ปญฺจหิ ธมฺเมหิ สมนฺนาคโต ภิกฺขุ อาหุเนยฺโย โหติ ฯเปฯ อนุตฺตรํ ปุญฺญกฺเขตฺตํ โลกสฺสาติ.\csromanspacer{}\csromanspacer{}ทสมํ.\csromanspacer{} ราชวคฺโค จตุตฺโถ.}

\titlehead{\textbf{ตสฺสุทฺทานํ}}
\csromangathameasure{จกฺกานุวตฺตนา ราชา, ยสฺสํทิสํ เทฺว เจว ปตฺถนา. \\ อปฺปํสุปติ ภตฺตาโท, อกฺขโม จ โสเตน จาติ.}
\csromangathagroup{\csromangathastanza{จกฺกานุวตฺตนา ราชา, ยสฺสํทิสํ เทฺว เจว ปตฺถนา. \\ อปฺปํสุปติ ภตฺตาโท, อกฺขโม จ โสเตน จาติ.}}

\csromanmatikaanchor{mtk.40}
\csromantocmarknopage{section}{(15) 5. ติกณฺฑกีวคฺค}{mtk.40}
\bookmark[dest={mtk.40},level=1]{(15) 5. ติกณฺฑกีวคฺค}
\csromanheader{1}{\textbf{(15) 5. ติกณฺฑกีวคฺค}}
\csromanmatikaanchor{mtk.41}
\csromantocmarkref{subsection}{1. อวชานาติสุตฺต}{mtk.41}
\bookmark[dest={mtk.41},level=2]{1. อวชานาติสุตฺต}
\csromanheader{2}{1. อวชานาติสุตฺต}
\proseitem{141}{ปญฺจิเม ภิกฺขเว ปุคฺคลา สนฺโต สํวิชฺชมานา โลกสฺมิํ.\csromanspacer{} กตเม ปญฺจ? ทตฺวา อวชานาติ, สํวาเสน อวชานาติ, อาเธยฺยมุโข\footnote{อาทิยฺยมุโข (สี), อาเทยฺยมุโข (สฺยา, กํ), อาทิยมุโข (อิ) อฏฺฐกถาย ปฐมสํวณฺณนานุรูปํ. อภิ 3. 175 ปิฏฺเฐปิ ปสฺสิตพฺพํ.} โหติ, โลโล โหติ, มนฺโท โมมูโห โหติ\footnote{มนฺโท โหติ โมมูโห (สี)}.}

\prose{กถญฺจ ภิกฺขเว ปุคฺคโล ทตฺวา อวชานาติ? อิธ ภิกฺขเว ปุคฺคโล ปุคฺคลสฺส เทติ จีวรปิณฺฑปาตเสนาสนคิลานปฺปจฺจยเภสชฺชปริกฺขารํ, ตสฺส เอวํ โหติ “อหํ เทมิ, อยํ ปฏิคฺคณฺหาตี”ติ, ตเมนํ ทตฺวา อวชานาติ.\csromanspacer{} เอวํ โข ภิกฺขเว ปุคฺคโล ทตฺวา อวชานาติ.}

\prose{\csromanfolio{146}กถญฺจ ภิกฺขเว ปุคฺคโล สํวาเสน อวชานาติ? อิธ ภิกฺขเว ปุคฺคโล ปุคฺคเลน สทฺธิํ สํวสติ เทฺว วา ตีณิ วา วสฺสานิ, ตเมนํ สํวาเสน อวชานาติ.\csromanspacer{} เอวํ โข ภิกฺขเว ปุคฺคโล สํวาเสน อวชานาติ.}

\prose{กถญฺจ ภิกฺขเว ปุคฺคโล อาเธยฺยมุโข โหติ? อิธ ภิกฺขเว เอกจฺโจ ปุคฺคโล ปรสฺส วณฺเณ วา อวณฺเณ วา ภาสิยมาเน ตํ ขิปฺปญฺเญว อธิมุจฺจิตา\footnote{อธิมุจฺจิโต (สฺยา)} โหติ.\csromanspacer{} เอวํ โข ภิกฺขเว ปุคฺคโล อาเธยฺยมุโข โหติ.}

\prose{กถญฺจ ภิกฺขเว ปุคฺคโล โลโล โหติ? อิธ ภิกฺขเว เอกจฺโจ ปุคฺคโล อิตฺตรสทฺโธ โหติ อิตฺตรภตฺตี อิตฺตรเปโม อิตฺตรปฺปสาโท.\csromanspacer{} เอวํ โข ภิกฺขเว ปุคฺคโล โลโล โหติ.}

\prose{กถญฺจ ภิกฺขเว ปุคฺคโล มนฺโท โมมูโห โหติ? อิธ ภิกฺขเว เอกจฺโจ ปุคฺคโล กุสลากุสเล ธมฺเม น ชานาติ, สาวชฺชานวชฺเช ธมฺเม น ชานาติ, หีนปฺปณีเต ธมฺเม น ชานาติ, กณฺหสุกฺกสปฺปฏิภาเค ธมฺเม น ชานาติ.\csromanspacer{} เอวํ โข ภิกฺขเว ปุคฺคโล มนฺโท โมมูโห โหติ.\csromanspacer{} อิเม โข ภิกฺขเว ปญฺจ ปุคฺคลา สนฺโต สํวิชฺชมานา โลกสฺมินฺติ.\csromanspacer{}\csromanspacer{}ปฐมํ.}

\csromanmatikaanchor{mtk.42}
\csromantocmarkref{subsection}{2. อารภติสุตฺต}{mtk.42}
\bookmark[dest={mtk.42},level=2]{2. อารภติสุตฺต}
\csromanheader{2}{2. อารภติสุตฺต}
\proseitem{142}{ปญฺจิเม ภิกฺขเว ปุคฺคลา สนฺโต สํวิชฺชมานา โลกสฺมิํ.\csromanspacer{} กตเม ปญฺจ? อิธ ภิกฺขเว เอกจฺโจ ปุคฺคโล อารภติ จ วิปฺปฏิสารี จ โหติ, ตญฺจ เจโตวิมุตฺติํ ปญฺญาวิมุตฺติํ ยถาภูตํ นปฺปชานาติ, ยตฺถสฺส เต อุปฺปนฺนา ปาปกา อกุสลา ธมฺมา อปริเสสา นิรุชฺฌนฺติ. (1)}

\prose{\csromansymbolfootnote{*}{อภิ 3. 173 ปิฏฺเฐปิ.}อิธ ปน ภิกฺขเว เอกจฺโจ ปุคฺคโล อารภติ น วิปฺปฏิสารี โหติ, ตญฺจ เจโตวิมุตฺติํ ปญฺญาวิมุตฺติํ ยถาภูตํ นปฺปชานาติ, ยตฺถสฺส เต อุปฺปนฺนา ปาปกา อกุสลา ธมฺมา อปริเสสา นิรุชฺฌนฺติ. (2)}

\prose{อิธ ปน ภิกฺขเว เอกจฺโจ ปุคฺคโล น อารภติ วิปฺปฏิสารี โหติ, ตญฺจ เจโตวิมุตฺติํ ปญฺญาวิมุตฺติํ ยถาภูตํ นปฺปชานาติ, หตฺถสฺส เต อุปฺปนฺนา ปาปกา อกุสลา ธมฺมา อปริเสสา นิรุชฺฌนฺติ. (3)}

\chapterheadpageread{(15) 5. ติกณฺฑกีวคฺค}
\prose{\csromanfolio{147}อิธ ปน ภิกฺขเว เอกจฺโจ ปุคฺคโล น อารภติ น วิปฺปฏิสารี โหติ, ตญฺจ เจโตวิมุตฺติํ ปญฺญาวิมุตฺติํ ยถาภูตํ นปฺปชานาติ, ยตฺถสฺส เต อุปฺปนฺนา ปาปกา อกุสลา ธมฺมา อปริเสสา นิรุชฺฌนฺติ. (4)}

\prose{อิธ ปน ภิกฺขเว เอกจฺโจ ปุคฺคโล น อารภติ น วิปฺปฏิสารี โหติ, ตญฺจ เจโตวิมุตฺติํ ปญฺญาวิมุตฺติํ ยถาภูตํ ปชานาติ, ยตฺถสฺส เต อุปฺปนฺนา ปาปกา อกุสลา ธมฺมา อปริเสสา นิรุชฺฌนฺติ. (5)}

\prose{ตตฺร ภิกฺขเว ยฺวายํ ปุคฺคโล อารภติ จ วิปฺปฏิสารี จ โหติ, ตญฺจ เจโตวิมุตฺติํ ปญฺญาวิมุตฺติํ ยถาภูตํ นปฺปชานาติ, ยตฺถสฺส เต อุปฺปนฺนา ปาปกา อกุสลา ธมฺมา อปริเสสา นิรุชฺฌนฺติ.\csromanspacer{} โส เอวมสฺส วจนีโย “อายสฺมโต โข อารมฺภชา\footnote{อารพฺภชา (อิ, ก), อารภชา (สฺยา, กํ)} อาสวา สํวิชฺชนฺติ, วิปฺปฏิสารชา อาสวา ปวฑฺฒนฺติ\footnote{สํวฑฺฒนฺติ (ก)}, สาธุ วตายสฺมา อารมฺภเช อาสเว ปหาย วิปฺปฏิสารเช อาสเว ปฏิวิโนเทตฺวา จิตฺตํ ปญฺญญฺจ ภาเวตุ\footnote{ภาเวตุํ (สี, อิ)}, เอวมายสฺมา อมุนา ปญฺจเมน ปุคฺคเลน สมสโม ภวิสฺสตี”ติ. (1)}

\prose{ตตฺร ภิกฺขเว ยฺวายํ ปุคฺคโล อารภติ น วิปฺปฏิสารี โหติ, ตญฺจ เจโตวิมุตฺติํ ปญฺญาวิมุตฺติํ ยถาภูตํ นปฺปชานาติ, ยตฺถสฺส เต อุปฺปนฺนา ปาปกา อกุสลา ธมฺมา อปริเสสา นิรุชฺฌนฺติ.\csromanspacer{} โส เอวมสฺส วจนีโย “อายสฺมโต โข อารมฺภชา อาสวา สํวิชฺชนฺติ, วิปฺปฏิสารชา อาสวา น ปวฑฺฒนฺติ, สาธุ วตายสฺมา อารมฺภเช อาสเว ปหาย จิตฺตํ ปญฺญญฺจ ภาเวตุ, เอวมายสฺมา อมุนา ปญฺจเมน ปุคฺคเลน สมสโม ภวิสฺสตี”ติ. (2)}

\prose{ตตฺร ภิกฺขเว ยฺวายํ ปุคฺคโล น อารภติ วิปฺปฏิสารี โหติ, ตญฺจ เจโตวิมุตฺติํ ปญฺญาวิมุตฺติํ ยถาภูตํ นปฺปชานาติ, ยตฺถสฺส เต อุปฺปนฺนา ปาปกา อกุสลา ธมฺมา อปริเสสา นิรุชฺฌนฺติ.\csromanspacer{} โส เอวมสฺส วจนีโย “อายสฺมโต โข อารมฺภชา อาสวา น สํวิชฺชนฺติ, วิปฺปฏิสารชา อาสวา ปวฑฺฒนฺติ, สาธุ วตายสฺมา วิปฺปฏิสารเช อาสเว ปฏิวิโนเทตฺวา จิตฺตํ ปญฺญญฺจ ภาเวตุ, เอวมายสฺมา อมุนา ปญฺจเมน ปุคฺคเลน สมสโม ภวิสฺสตี”ติ. (3)}

\prose{ตตฺร ภิกฺขเว ยฺวายํ ปุคฺคโล น อารภติ น วิปฺปฏิสารี โหติ, ตญฺจ เจโตวิมุตฺติํ ปญฺญาวิมุตฺติํ ยถาภูตํ นปฺปชานาติ, ยตฺถสฺส เต อุปฺปนฺนา}

\prose{\csromanfolio{148}ปาปกา อกุสลา ธมฺมา อปริเสสา นิรุชฺฌนฺติ.\csromanspacer{} โส เอวมสฺส วจนีโย “อายสฺมโต โข อารมฺภชา อาสวา น สํวิชฺชนฺติ, วิปฺปฏิสารชา อาสวา น ปวฑฺฒนฺติ, สาธุ วตายสฺมา จิตฺตํ ปญฺญญฺจ ภาเวตุ, เอวมายสฺมา อมุนา ปญฺจเมน ปุคฺคเลน สมสโม ภวิสฺสตี”ติ. (4)}

\prose{อิติ โข ภิกฺขเว อิเม จตฺตาโร ปุคฺคลา อมุนา ปญฺจเมน ปุคฺคเลน เอวํ โอวทิยมานา เอวํ อนุสาสิยมานา อนุปุพฺเพน อาสวานํ ขยํ ปาปุณนฺตีติ*.\csromanspacer{}\csromanspacer{}ทุติยํ.}

\csromanmatikaanchor{mtk.43}
\csromantocmarkref{subsection}{3. สารนฺททสุตฺต}{mtk.43}
\bookmark[dest={mtk.43},level=2]{3. สารนฺททสุตฺต}
\csromanheader{2}{3. สารนฺททสุตฺต}
\proseitem{143}{เอกํ สมยํ ภควา เวสาลิยํ วิหรติ มหาวเน กูฏาคารสาลายํ.\csromanspacer{} อถ โข ภควา ปุพฺพณฺหสมยํ นิวาเสตฺวา ปตฺตจีวรมาทาย เวสาลิํ ปิณฺฑาย ปาวิสิ.\csromanspacer{} เตน โข ปน สมเยน ปญฺจมตฺตานํ ลิจฺฉวิสตานํ สารนฺทเท เจติเย สนฺนิสินฺนานํ สนฺนิปติตานํ อยมนฺตรากถา อุทปาทิ\csromandash{}ปญฺจนฺนํ รตนานํ ปาตุภาโว ทุลฺลโภ โลกสฺมิํ.\csromanspacer{} กตเมสํ ปญฺจนฺนํ? หตฺถิรตนสฺส ปาตุภาโว ทุลฺลโภ โลกสฺมิํ, อสฺสรตนสฺส ปาตุภาโว ทุลฺลโภ โลกสฺมิํ, มณิรตนสฺส ปาตุภาโว ทุลฺลโภ โลกสฺมิํ, อิตฺถิรตนสฺส ปาตุภาโว ทุลฺลโภ โลกสฺมิํ, คหปติรตนสฺส ปาตุภาโว ทุลฺลโภ โลกสฺมิํ.\csromanspacer{} อิเมสํ ปญฺจนฺนํ รตนานํ ปาตุภาโว ทุลฺลโภ โลกสฺมินฺติ.}

\prose{อถ โข เต ลิจฺฉวี มคฺเค ปุริสํ ฐเปสุํ\footnote{เปเสสุํ (สฺยา, ก)} “ยทา ตฺวํ\footnote{ยถา ตฺวํ (สี, อิ)} อมฺโภ ปุริส ปสฺเสยฺยาสิ ภควนฺตํ, อถ อมฺหากํ อาโรเจยฺยาสี”ติ.\csromanspacer{} อทฺทสา โข โส ปุริโส ภควนฺตํ ทูรโตว อาคจฺฉนฺตํ, ทิสฺวาน เยน เต ลิจฺฉวี เตนุปสงฺกมิ, อุปสงฺกมิตฺวา เต ลิจฺฉวี เอตทโวจ “อยํ โส ภนฺเต ภควา คจฺฉติ อรหํ สมฺมาสมฺพุทฺโธ, ยสฺสทานิ กาลํ มญฺญถา”ติ.}

\prose{อถ โข เต ลิจฺฉวี เยน ภควา เตนุปสงฺกมิํสุ, อุปสงฺกมิตฺวา ภควนฺตํ อภิวาเทตฺวา เอกมนฺตํ อฏฺฐํสุ, เอกมนฺตํ ฐิตา โข เต ลิจฺฉวี ภควนฺตํ เอตทโวจุํ\csromandash{}}

\prose{“สาธุ ภนฺเต เยน สารนฺททํ เจติยํ เตนุปสงฺกมตุ อนุกมฺปํ อุปาทายา”ติ.\csromanspacer{} อธิวาเสสิ ภควา ตุณฺหีภาเวน.\csromanspacer{} อถ โข ภควา}

\vspace{\tipitakaparskip}
\csromancenter{(15) 5.\csromanspacer{} ติกณฺฑกีวคฺค เยน สารนฺททํ เจติยํ เตนุปสงฺกมิ, อุปสงฺกมิตฺวา ปญฺญตฺเต อาสเน นิสีทิ, นิสชฺช โข ภควา เต ลิจฺฉวี เอตทโวจ “กาย นุตฺถ ลิจฺฉวี เอตรหิ กถาย สนฺนิสินฺนา, กา จ ปน โว อนฺตรากถา วิปฺปกตา”ติ.\csromanspacer{} อิธ ภนฺเต อมฺหากํ สนฺนิสินฺนานํ สนฺนิปติตานํ อยมนฺตรากถา อุทปาทิ\csromandash{} ปญฺจนฺนํ รตนานํ ปาตุภาโว ทุลฺลโภ โลกสฺมิํ.\csromanspacer{} กตเมสํ ปญฺจนฺนํ? หตฺถิรตนสฺส ปาตุภาโว ทุลฺลโภ โลกสฺมิํ, อสฺสรตนสฺส ปาตุภาโว ทุลฺลโภ โลกสฺมิํ, มณิรตนสฺส ปาตุภาโว ทุลฺลโภ โลกสฺมิํ, อิตฺถิรตนสฺส ปาตุภาโว ทุลฺลโภ โลกสฺมิํ.\csromanspacer{} คหปติรตนสฺส ปาตุภาโว ทุลฺลโภ โลกสฺมิํ.\csromanspacer{} อิเมสํ ปญฺจนฺนํ รตนานํ ปาตุภาโว ทุลฺลโภ โลกสฺมินฺติ.}
\prose{\csromanfolio{149}กามาธิมุตฺตานํ วต โภ ลิจฺฉวีนํ\footnote{กามาธิมุตฺตานํ วต โว ลิจฺฉวีนํ (สี), กามาธิมุตฺตานํ วต โว ลิจฺฉวี (สฺยา), กามาธิมุตฺตานํว โว ลิจฺฉวี (?)} กามํเยว อารพฺภ อนฺตรากถา อุทปาทิ.\csromanspacer{} ปญฺจนฺนํ ลิจฺฉวี รตนานํ ปาตุภาโว ทุลฺลโภ โลกสฺมิํ.\csromanspacer{} กตเมสํ ปญฺจนฺนํ? ตถาคตสฺส อรหโต สมฺมาสมฺพุทฺธสฺส ปาตุภาโว ทุลฺลโภ โลกสฺมิํ, ตถาคตปฺปเวทิตสฺส ธมฺมวินยสฺส เทเสตา ปุคฺคโล ทุลฺลโภ โลกสฺมิํ, ตถาคตปฺปเวทิตสฺส ธมฺมวินยสฺส เทสิตสฺส วิญฺญาตา ปุคฺคโล ทุลฺลโภ โลกสฺมิํ, ตถาคตปฺปเวทิตสฺส ธมฺมวินยสฺส เทสิตสฺส วิญฺญาตา\footnote{วิญฺญาตสฺส (สี, อิ), อุปริ 209 ปิฏฺเฐปิ.} ธมฺมานุ ธมฺมปฺปฏิปนฺโน ปุคฺคโล ทุลฺลโภ โลกสฺมิํ, กตญฺญู กตเวที ปุคฺคโล ทุลฺลโภ โลกสฺมิํ.\csromanspacer{} อิเมสํ โข ลิจฺฉวี ปญฺจนฺนํ รตนานํ ปาตุภาโว ทุลฺลโภ โลกสฺมินฺติ.\csromanspacer{}\csromanspacer{}ตติยํ.}

\csromanmatikaanchor{mtk.44}
\csromantocmarkref{subsection}{4. ติกณฺฑกีสุตฺต}{mtk.44}
\bookmark[dest={mtk.44},level=2]{4. ติกณฺฑกีสุตฺต}
\csromanheader{2}{4. ติกณฺฑกีสุตฺต}
\proseitem{144}{เอกํ สมยํ ภควา สาเกเต วิหรติ ติกณฺฑกีวเน\footnote{กณฺฑกีวเน (สํ 3. 260 ปิฏฺเฐปิ)}.\csromanspacer{} ตตฺร โข ภควา ภิกฺขู อามนฺเตสิ “ภิกฺขโว”ติ. “ภทนฺเต”ติ เต ภิกฺขู ภควโต ปจฺจสฺโสสุํ.\csromanspacer{} ภควา เอตทโวจ\csromandash{}}

\csromanmatikaanchor{mtk.45}
\csromantocmarkref{subsection}{5-6. นิรยสุตฺตาทิ}{mtk.45}
\bookmark[dest={mtk.45},level=2]{5-6. นิรยสุตฺตาทิ}
\prose{สาธุ ภิกฺขเว ภิกฺขุ กาเลน กาลํ อปฺปฏิกูเล ปฏิกูลสญฺญี\footnote{อปฺปฏิกฺกูเล ปฏิกฺกูลสญฺญี (สี, สฺยา, กํ, อิ)} วิหเรยฺย, สาธุ ภิกฺขเว ภิกฺขุ กาเลน กาลํ ปฏิกูเล อปฺปฏิกูลสญฺญี วิหเรยฺย, สาธุ ภิกฺขเว ภิกฺขุ กาเลน กาลํ อปฺปฏิกูเล จ \csromanfolio{150}ปฏิกูเล จ ปฏิกูลสญฺญี วิหเรยฺย, สาธุ ภิกฺขเว ภิกฺขุ กาเลน กาลํ ปฏิกูเล จ อปฺปฏิกูเล จ อปฺปฏิกูลสญฺญี วิหเรยฺย, สาธุ ภิกฺขเว ภิกฺขุ กาเลน กาลํ ปฏิกูลญฺจ อปฺปฏิกูลญฺจ ตทุภยํ อภินิวชฺเชตฺวา อุเปกฺขโก วิหเรยฺย สโต สมฺปชาโน.}

\prose{กิญฺจ\footnote{กถญฺจ (สี, อิ, ก)} ภิกฺขเว ภิกฺขุ อตฺถวสํ ปฏิจฺจ อปฺปฏิกูเล ปฏิกูลสญฺญี วิหเรยฺย.\csromanspacer{} มา เม รชนีเยสุ ธมฺเมสุ ราโค อุทปาทีติ, อิทํ โข ภิกฺขเว ภิกฺขุ อตฺถวสํ ปฏิจฺจ อปฺปฏิกูเล ปฏิกูลสญฺญี วิหเรยฺย.}

\prose{กิญฺจ ภิกฺขเว ภิกฺขุ อตฺถวสํ ปฏิจฺจ ปฏิกูเล อปฺปฏิกูลสญฺญี วิหเรยฺย.\csromanspacer{} มา เม โทสนีเยสุ ธมฺเมสุ โทโส อุทปาทีติ, อิทํ โข ภิกฺขเว ภิกฺขุ อตฺถวสํ ปฏิจฺจ ปฏิกูเล อปฺปฏิกูลสญฺญี วิหเรยฺย.}

\prose{กิญฺจ ภิกฺขเว ภิกฺขุ อตฺถวสํ ปฏิจฺจ อปฺปฏิกูเล จ ปฏิกูเล จ ปฏิกูลสญฺญี วิหเรยฺย.\csromanspacer{} มา เม รชนีเยสุ ธมฺเมสุ ราโค อุทปาทิ, มา เม โทสนีเยสุ ธมฺเมสุ โทโส อุทปาทีติ, อิทํ โข ภิกฺขเว ภิกฺขุ อตฺถวสํ ปฏิจฺจ อปฺปฏิกูเล จ ปฏิกูเล จ ปฏิกูลสญฺญี วิหเรยฺย.}

\prose{กิญฺจ ภิกฺขเว ภิกฺขุ อตฺถวสํ ปฏิจฺจ ปฏิกูเล จ อปฺปฏิกูเล จ อปฺปฏิกูลสญฺญี วิหเรยฺย.\csromanspacer{} มา เม โทสนีเยสุ ธมฺเมสุ โทโส อุทปาทิ, มา เม รชนีเยสุ ธมฺเมสุ ราโค อุทปาทีติ, อิทํ โข ภิกฺขเว ภิกฺขุ อตฺถวสํ ปฏิจฺจ ปฏิกูเล จ อปฺปฏิกูเล จ อปฺปฏิกูลสญฺญี วิหเรยฺย.}

\prose{กิญฺจ ภิกฺขเว ภิกฺขุ อตฺถวสํ ปฏิจฺจ ปฏิกูลญฺจ อปฺปฏิกูลญฺจ ตทุภยํ อภินิวชฺเชตฺวา อุเปกฺขโก วิหเรยฺย สโต สมฺปชาโน.\csromanspacer{} มา เม กฺวจนิ\footnote{กฺวจินิ (สี, สฺยา, อิ)} กตฺถจิ กิญฺจนํ\footnote{กิญฺจน (สิ, อิ)} รชนีเยสุ ธมฺเมสุ ราโค อุทปาทิ, มา เม กฺวจนิ กตฺถจิ กิญฺจนํ โทสนีเยสุ ธมฺเมสุ โทโส อุทปาทิ, มา เม กฺวจนิ กตฺถจิ กิญฺจนํ โมหนีเยสุ ธมฺเมสุ โมโห อุทปาทีติ, อิทํ โข ภิกฺขเว ภิกฺขุ อตฺถวสํ ปฏิจฺจ ปฏิกูลญฺจ อปฺปฏิกูลญฺจ ตทุภยํ อภินิวชฺเชตฺวา อุเปกฺขโก วิหเรยฺย สโต สมฺปชาโนติ.\csromanspacer{}\csromanspacer{}จตุตฺถํ.}

\csromanheader{2}{5. นิรยสุตฺต}
\proseitem{145}{ปญฺจหิ ภิกฺขเว ธมฺเมหิ สมนฺนาคโต ยถาภตํ นิกฺขิตฺโต เอวํ นิรเย.\csromanspacer{} กตเมหิ ปญฺจหิ? ปาณาติปาตี โหติ, อทินฺนาทายี โหติ,}

\vspace{\tipitakaparskip}
\csromancenter{(15) 5.\csromanspacer{} ติกณฺฑกีวคฺค กาเมสุมิจฺฉาจารี โหติ, มุสาวาที โหติ, สุราเมรยมชฺชปมาทฏฺฐายี โหติ.\csromanspacer{} อิเมหิ โข ภิกฺขเว ปญฺจหิ ธมฺเมหิ สมนฺนาคโต ยถาภตํ นิกฺขิตฺโต เอวํ นิรเย.}
\csromanmatikaanchor{mtk.46}
\csromantocmarkref{subsection}{7-10. อสปฺปุริสทานสุตฺตาทิ}{mtk.46}
\bookmark[dest={mtk.46},level=2]{7-10. อสปฺปุริสทานสุตฺตาทิ}
\prose{\csromanfolio{151}ปญฺจหิ ภิกฺขเว ธมฺเมหิ สมนฺนาคโต ยถาภตํ นิกฺขิตฺโต เอวํ สคฺเค.\csromanspacer{} กตเมหิ ปญฺจหิ? ปาณติปาตา ปฏิวิรโต โหติ, อทินฺนาทานา ปฏิวิรโต โหติ, กาเมสุ มิจฺฉาจารา ปฏิวิรโต โหติ, มุสาวาทา ปฏิวิรโต โหติ, สุราเมรยมชฺชปมาทฏฺฐานา ปฏิวิรโต โหติ.\csromanspacer{} อิเมหิ โข ภิกฺขเว ปญฺจหิ ธมฺเมหิ สมนฺนาคโต ยถาภตํ นิกฺขิตฺโต เอวํ สคฺเคติ.\csromanspacer{}\csromanspacer{}ปญฺจมํ.}

\csromanheader{2}{6. มิตฺตสุตฺต}
\proseitem{146}{ปญฺจหิ ภิกฺขเว ธมฺเมหิ สมนฺนาคโต ภิกฺขุ มิตฺโต น เสวิตพฺโพ.\csromanspacer{} กตเมหิ ปญฺจหิ? กมฺมนฺตํ กาเรติ, อธิกรณํ อาทิยติ, ปาโมกฺเขสุ ภิกฺขูสุ ปฏิวิรุทฺโธ โหติ, ทีฆจาริกํ อนวตฺถจาริกํ\footnote{อวตฺถานจาริกํ (สฺยา)} อนุยุตฺโต วิหรติ, นปฺปฏิพโล โหติ กาเลน กาลํ ธมฺมิยา กถาย สนฺทสฺเสตุํ สมาทเปตุํ สมุตฺเตเชตุํ สมฺปหํเสตุํ.\csromanspacer{} อิเมหิ โข ภิกฺขเว ปญฺจหิ ธมฺเมหิ สมนฺนาคโต ภิกฺขุ มิตฺโต น เสวิตพฺโพ.}

\prose{ปญฺจหิ ภิกฺขเว ธมฺเมหิ สมนฺนาคโต ภิกฺขุ มิตฺโต เสวิตพฺโพ.\csromanspacer{} กตเมหิ ปญฺจหิ? น กมฺมนฺตํ กาเรติ, น อธิกรณํ อาทิยติ, น ปาโมกฺเขสุ ภิกฺขูสุ ปฏิวิรุทฺโธ โหติ, น ทีฆจาริกํ อนวตฺถจาริกํ อนุยุตฺโต วิหรติ, ปฏิพโล โหติ กาเลน กาลํ ธมฺมิยา กถาย สนฺทสฺเสตุํ สมาทเปตุํ สมุตฺเตเชตุํ สมฺปหํเสตุํ.\csromanspacer{} อิเมหิ โข ภิกฺขเว ปญฺจหิ ธมฺเมหิ สมนฺนาคโต ภิกฺขุ มิตฺโต เสวิตพฺโพติ.\csromanspacer{}\csromanspacer{}ฉฏฺฐํ.}

\csromanheader{2}{7. อสปฺปุริสทานสุตฺต}
\proseitem{147}{ปญฺจิมานิ ภิกฺขเว อสปฺปุริสทานานิ.\csromanspacer{} กตมานิ ปญฺจ? อสกฺกจฺจํ เทติ, อจิตฺตีกตฺวา\footnote{อจิตฺติกตฺวา (อิ), อจิติํ กตฺวา (สฺยา), อจิตฺติํ กตฺวา (ก)} เทติ, อสหตฺถา เทติ, อปวิทฺธํ\footnote{อปวิฏฺฐํ (สฺยา, กํ)} เทติ, อนาคมนทิฏฺฐิโก เทติ.\csromanspacer{} อิมานิ โข ภิกฺขเว ปญฺจ อสปฺปุริสทานานิ.}

\prose{\csromanfolio{152}ปญฺจิมานิ ภิกฺขเว สปฺปุริสทานานิ.\csromanspacer{} กตมานิ ปญฺจ? สกฺกจฺจํ เทติ, จิตฺตีกตฺวา เทติ, สหตฺถา เทติ, อนปวิทฺธํ เทติ, อาคมนทิฏฺฐิโก เทติ.\csromanspacer{} อิมานิ โข ภิกฺขเว ปญฺจ สปฺปุริสทานานีติ.\csromanspacer{}\csromanspacer{}สตฺตมํ.}

\csromanheader{2}{8. สปฺปุริสทานสุตฺต}
\proseitem{148}{ปญฺจิมานิ ภิกฺขเว สปฺปุริสทานานิ.\csromanspacer{} กตมานิ ปญฺจ? สทฺธาย ทานํ เทติ, สกฺกจฺจํ ทานํ เทติ, กาเลน ทานํ เทติ, อนุคฺคหิตจิตฺโต\footnote{อนคฺคหิตจิตฺโต (สี)} ทานํ เทติ, อตฺตานญฺจ ปรญฺจ อนุปหจฺจ ทานํ เทติ.}

\prose{สทฺธาย โข ปน ภิกฺขเว ทานํ ทตฺวา ยตฺถ ยตฺถ ตสฺส ทานสฺส วิปาโก นิพฺพตฺตติ, อฑฺโฒ จ โหติ มหทฺธโน มหาโภโค, อภิรูโป จ โหติ ทสฺสนีโย ปาสาทิโก ปรมาย วณฺณโปกฺขรตาย สมนฺนาคโต.}

\prose{สกฺกจฺจํ โข ปน ภิกฺขเว ทานํ ทตฺวา ยตฺถ ยตฺถ ตสฺส ทานสฺส วิปาโก นิพฺพตฺตติ, อฑฺโฒ จ โหติ มหทฺธโน มหาโภโค, เยปิสฺส เต โหนฺติ “ปุตฺตา”ติ วา “ทารา”ติ วา “ทาสา”ติ วา “เปสฺสา”ติ วา “กมฺมกรา”ติ\footnote{กมฺมการาติ (ก)} วา, เตปิ สุสฺสูสนฺติ, โสตํ โอทหนฺติ, อญฺญา จิตฺตํ อุปฏฺฐเปนฺติ.}

\prose{กาเลน โข ปน ภิกฺขเว ทานํ ทตฺวา ยตฺถ ยตฺถ ตสฺส ทานสฺส วิปาโก นิพฺพตฺตติ, อฑฺโฒ จ โหติ มหทฺธโน มหโภโค, กาลาคตา จสฺส อตฺถา ปจุรา โหนฺติ.}

\prose{อนุคฺคหิตจิตฺโต โข ปน ภิกฺขเว ทานํ ทตฺวา ยตฺถ ยตฺถ ตสฺส ทานสฺส วิปาโก นิพฺพตฺตติ, อฑฺโฒ จ โหติ มหทฺธโน มหาโภโค, อุฬาเรสุ จ ปญฺจสุ กามคุเณสุ โภคาย จิตฺตํ นมติ.}

\prose{อตฺตานญฺจ ปรญฺจ อนุปหจฺจ โข ปน ภิกฺขเว ทานํ ทตฺวา ยตฺถ ยตฺถ ตสฺส ทานสฺส วิปาโก นิพฺพตฺตติ, อฑฺโฒ จ โหติ มหทฺธโน มหาโภโค, น จสฺส กุโตจิ โภคานํ อุปฆาโต อาคจฺฉติ อคฺคิโต วา อุทกโต วา ราชโต วา โจรโต วา อปฺปิยโต วา ทายาทโต\footnote{อปฺปิยโต วา ทายาทโต วา (สี, สฺยา, กํ, อิ), อปฺปิยทายาทโต วา (ก)}.\csromanspacer{} อิมานิ โข ภิกฺขเว ปญฺจ สปฺปุริสทานานีติ.\csromanspacer{}\csromanspacer{}อฏฺฐมํ.}

\chapterheadpageread{(15) 5. ติกณฺฑกีวคฺค \textbf{9. ปฐมสมยวิมุตฺตสุตฺต}}
\proseitem{149}{\csromanfolio{153}ปญฺจิเม ภิกฺขเว ธมฺมา สมยวิมุตฺตสฺส ภิกฺขุโน ปริหานาย สํวตฺตนฺติ.\csromanspacer{} กตเม ปญฺจ? กมฺมารามตา ภสฺสารามตา นิทฺทารามตา สงฺคณิการามตา ยถาวิมุตฺตํ จิตฺตํ น ปจฺจเวกฺขติ.\csromanspacer{} อิเม โข ภิกฺขเว ปญฺจ ธมฺมา สมยวิมุตฺตสฺส ภิกฺขุโน ปริหานาย สํวตฺตนฺติ.}

\prose{ปญฺจิเม ภิกฺขเว ธมฺมา สมยวิมุตฺตสฺส ภิกฺขุโน อปริหานาย สํวตฺตนฺติ.\csromanspacer{} กตเม ปญฺจ? น กมฺมารามตา น ภสฺสารามตา น นิทฺทารามตา น สงฺคณิการามตา ยถาวิมุตฺตํ จิตฺตํ ปจฺจเวกฺขติ.\csromanspacer{} อิเม โข ภิกฺขเว ปญฺจ ธมฺมา สมยวิมุตฺตสฺส ภิกฺขุโน อปริหานาย สํวตฺตนฺตีติ.\csromanspacer{}\csromanspacer{}นวมํ.}

\csromanheader{2}{10. ทุติยสมยวิมุตฺตสุตฺต}
\proseitem{150}{\csromansymbolfootnote{*}{อภิ 4. 78 ปิฏฺเฐสุ.}ปญฺจิเม ภิกฺขเว ธมฺมา สมยวิมุตฺตสฺส ภิกฺขุโน ปริหานาย สํวตฺตนฺติ.\csromanspacer{} กตเม ปญฺจ? กมฺมารามตา ภสฺสารามตา นิทฺทารามตา อินฺทฺริเยสุ อคุตฺตทฺวารตา โภชเน อมตฺตญฺญุตา.\csromanspacer{} อิเม โข ภิกฺขเว ปญฺจ ธมฺมา สมย วิมุตฺตสฺส ภิกฺขุโน ปริหานาย สํวตฺตนฺติ.}

\prose{ปญฺจิเม ภิกฺขเว ธมฺมา สมยวิมุตฺตสฺส ภิกฺขุโน อปริหานาย สํวตฺตนฺติ.\csromanspacer{} กตเม ปญฺจ? น กมฺมารามตา น ภสฺสารามตา น นิทฺทารามตา อินฺทฺริเยสุ คุตฺตทฺวารตา โภชเน มตฺตญฺญุตา.\csromanspacer{} อิเม โข ภิกฺขเว ปญฺจ ธมฺมา สมยวิมุตฺตสฺส ภิกฺขุโน อปริหานาย สํวตฺตนฺตีติ.\csromanspacer{}\csromanspacer{}ทสมํ.\csromanspacer{} ติกณฺฑกีวคฺโค ปญฺจโม.}

\vspace{\tipitakaparskip}
\csromancenter{\textbf{ตสฺสุทฺทานํ.}}
\chapterhead{ทตฺวา อวชานาติ อารภติ จ, สารนฺทท ติกณฺฑ นิรเยน จ.}
\csromangathameasure{มิตฺโต อสปฺปุริสสปฺปุริเสน, สมยวิมุตฺตํ อปเร เทฺวติ.}
\csromangathagroup{\csromangathastanza{มิตฺโต อสปฺปุริสสปฺปุริเสน, สมยวิมุตฺตํ อปเร เทฺวติ.}}

\vspace{\tipitakaparskip}
\csromancenter{\textbf{ตติโย ปณฺณาสโก สมตฺโต.}}
\csromanmatikaanchor{mtk.47}
\csromantocmarknopage{chapter}{4. จตุตฺถปณฺณาสก}{mtk.47}
\bookmark[dest={mtk.47},level=0]{4. จตุตฺถปณฺณาสก}
\chapterheadpageread{4. จตุตฺถปณฺณาสก}
\csromanmatikaanchor{mtk.48}
\csromantocmarknopage{section}{(16) 1. สทฺธมฺมวคฺค}{mtk.48}
\bookmark[dest={mtk.48},level=1]{(16) 1. สทฺธมฺมวคฺค}
\csromanheader{1}{\textbf{(16) 1. สทฺธมฺมวคฺค}}
\csromanheader{2}{1. ปฐมสมฺมตฺตนิยามสุตฺต}
\csromanmatikaanchor{mtk.49}
\csromantocmarkref{subsection}{1-3. สมฺมตฺตนิยามสุตฺตานิ}{mtk.49}
\bookmark[dest={mtk.49},level=2]{1-3. สมฺมตฺตนิยามสุตฺตานิ}
\proseitem{151}{\csromanfolio{154}ปญฺจหิ ภิกฺขเว ธมฺเมหิ สมนฺนาคโต สุณนฺโตปิ สทฺธมฺมํ อภพฺโพ นิยามํ โอกฺกมิตุํ กุสเลสุ ธมฺเมสุ สมฺมตฺตํ.\csromanspacer{} กตเมหิ ปญฺจหิ? กถํ ปริโภติ, กถิกํ\footnote{กถิตํ (ก)} ปริโภติ, อตฺตานํ ปริโภติ, วิกฺขิตฺตจิตฺโต ธมฺมํ สุณาติ อเนกคฺคจิตฺโต, อโยนิโส จ\footnote{อโยนิโส (สฺยา, กํ)} มนสิ กโรติ.\csromanspacer{} อิเม โข ภิกฺขเว ปญฺจหิ ธมฺเมหิ สมนฺนาคโต สุณนฺโตปิ สทฺธมฺมํ อภพฺโพ นิยาเม โอกฺกมิตุํ กุสเลสุ ธมฺเมสุ สมฺมตฺตํ.}

\prose{ปญฺจหิ ภิกฺขเว ธมฺเมหิ สมนฺนาคโต สุณนฺโต สทฺธมฺมํ ภพฺโพ นิยามํ โอกฺกมิตุํ กุสเลสุ ธมฺเมสุ สมฺมตฺตํ.\csromanspacer{} กตเมหิ ปญฺจหิ? น กถํ ปริโภติ, น กถิกํ ปริโภติ, น อตฺตานํ ปริโภติ, อวิกฺขิตฺตจิตฺโต ธมฺมํ สุณาติ เอกคฺคจิตฺโต, โยนิโส จ มนสิ กโรติ.\csromanspacer{} อิเมหิ โข ภิกฺขเว ปญฺจหิ ธมฺเมหิ สมนฺนาคโต สุณนฺโต สทฺธมฺมํ ภพฺโพ นิยามํ โอกฺกมิตุํ กุสเลสุ ธมฺเมสุ สมฺมตฺตนฺติ.\csromanspacer{}\csromanspacer{}ปฐมํ.}

\csromanheader{2}{2. ทุติยสมฺมตฺตนิยามสุตฺต}
\proseitem{152}{ปญฺจหิ ภิกฺขเว ธมฺเมหิ สมนฺนาคโต สุณนฺโตปิ สทฺธมฺมํ อภพฺโพ นิยามํ โอกฺกมิตุํ กุสเลสุ ธมฺเมสุ สมฺมตฺตํ.\csromanspacer{} กตเมหิ ปญฺจหิ? กถํ ปริโภติ, กถิกํ ปริโภติ, อตฺตานํ ปริโภติ, ทุปฺปญฺโญ โหติ ชโฬ เอฬมูโค, อนญฺญาเต อญฺญาตมานี โหติ.\csromanspacer{} อิเมหิ โข ภิกฺขเว ปญฺจหิ ธมฺเมหิ สมนฺนาคโต สุณนฺโตปิ สทฺธมฺมํ อภพฺโพ นิยามํ โอกฺกมิตุํ กุสเลสุ ธมฺเมสุ สมฺมตฺตํ.}

\prose{ปญฺจหิ ภิกฺขเว ธมฺเมหิ สมนฺนาคโต สุณนฺโต สทฺธมฺมํ ภพฺโพ นิยามํ โอกฺกมิตุํ กุสเลสุ ธมฺเมสุ สมฺมตฺตํ.\csromanspacer{} กตเมหิ ปญฺจหิ? น กถํ ปริโภติ, น กถิกํ ปริโภติ, น อตฺตานํ ปริโภติ, ปญฺญวา}

\vspace{\tipitakaparskip}
\csromancenter{(16) 1.\csromanspacer{} สทฺธมฺมวคฺค โหติ อชโฬ อเนฬมูโค น อนญฺญาเต อญฺญาตมานี โหติ.\csromanspacer{} อิเมหิ โข ภิกฺขเว ปญฺจหิ ธมฺเมหิ สมนฺนาคโต สุณนฺโต สทฺธมฺมํ ภพฺโพ นิยามํ โอกฺกมิตุํ กุสเลสุ ธมฺเมสุ สมฺมตฺตนฺติ.\csromanspacer{}\csromanspacer{}ทุติยํ.}
\csromanheader{2}{3. ตติยสมฺมตฺตนิยามสุตฺต}
\csromanmatikaanchor{mtk.50}
\csromantocmarkref{subsection}{4-6. สทฺธมฺมสมฺโมสสุตฺตานิ}{mtk.50}
\bookmark[dest={mtk.50},level=2]{4-6. สทฺธมฺมสมฺโมสสุตฺตานิ}
\proseitem{153}{\csromanfolio{155}ปญฺจหิ ภิกฺขเว ธมฺเมหิ สมนฺนาคโต สุณนฺโตปิ สทฺธมฺมํ อภพฺโพ นิยามํ โอกฺกมิตุํ กุสเลสุ ธมฺเมสุ สมฺมตฺตํ.\csromanspacer{} กตเมหิ ปญฺจหิ? มกฺขี ธมฺมํ สุณาติ มกฺขปริยุฏฺฐิโต, อุปารมฺภจิตฺโต\footnote{ส-อุปารมฺภจิตฺโต (สฺยา, กํ)} ธมฺมํ สุณาติ รนฺธคเวสี, ธมฺมเทสเก อาหตจิตฺโต โหติ ขีลชาโต\footnote{ขิลชาโต (สฺยา, อิ)}, ทุปฺปญฺโญ โหติ ชโฬ เอฬมูโค, อนญฺญาเต อญฺญาตมานี โหติ.\csromanspacer{} อิเมหิ โข ภิกฺขเว ปญฺจหิ ธมฺเมหิ สมนฺนาคโต สุณนฺโตปิ สทฺธมฺมํ อภพฺโพ นิยามํ โอกฺกมิตุํ กุสเลสุ ธมฺเมสุ สมฺมตฺตํ.}

\prose{ปญฺจหิ ภิกฺขเว ธมฺเมหิ สมนฺนาคโต สุณนฺโต สทฺธมฺมํ ภพฺโพ นิยามํ โอกฺกมิตุํ กุสเลสุ ธมฺเมสุ สมฺมตฺตํ.\csromanspacer{} กตเมหิ ปญฺจหิ? อมกฺขี ธมฺมํ สุณาติ น มกฺขปริยุฏฺฐิโต, อนุปารมฺภจิตฺโต ธมฺมํ สุณาติ น รนฺธคเวสี, ธมฺมเทสเก อนาหตจิตฺโต โหติ อขีลชาโต, ปญฺญวา โหติ อชโฬ อเนฬมูโค, น อนญฺญาเต อญฺญาตมานี โหติ.\csromanspacer{} อิเมหิ โข ภิกฺขเว ปญฺจหิ ธมฺเมหิ สมนฺนาคโต สุณนฺโต สทฺธมฺมํ ภพฺโพ นิยามํ โอกฺกมิตุํ กุสเลสุ ธมฺเมสุ สมฺมตฺตนฺติ.\csromanspacer{}\csromanspacer{}ตติยํ.}

\csromanheader{2}{4. ปฐมสทฺธมฺมสมฺโมสสุตฺต}
\proseitem{154}{ปญฺจิเม ภิกฺขเว ธมฺมา สทฺธมฺมสฺส สมฺโมสาย อนฺตรธานาย สํวตฺตนฺติ.\csromanspacer{} กตเม ปญฺจ? อิธ ภิกฺขเว ภิกฺขู น สกฺกจฺจํ ธมฺมํ สุณนฺติ, น สกฺกจฺจํ ธมฺมํ ปริยาปุณนฺติ, น สกฺกจฺจํ ธมฺมํ ธาเรนฺติ, น สกฺกจฺจํ ธาตานํ\footnote{ธตานํ (สี, สฺยา, กํ, อิ)} ธมฺมานํ อตฺถํ อุปปริกฺขนฺติ, น สกฺกจฺจํ อตฺถมญฺญาย ธมฺมมญฺญาย ธมฺมานุธมฺมํ ปฏิปชฺชนฺติ.\csromanspacer{} อิเม โข ภิกฺขเว ปญฺจ ธมฺมา สทฺธมฺมสฺส สมฺโมสาย อนฺตรธานาย สํวตฺตนฺติ.}

\prose{ปญฺจิเม ภิกฺขเว ธมฺมา สทฺธมฺมสฺส ฐิติยา อสมฺโมสาย อนนฺตรธานาย สํวตฺตนฺติ.\csromanspacer{} กตเม ปญฺจ? อิธ ภิกฺขเว ภิกฺขู สกฺกจฺจํ ธมฺมํ สุณนฺติ, \csromanfolio{156}สกฺกจฺจํ ธมฺมํ ปริยาปุณนฺติ, สกฺกจฺจํ ธมฺมํ ธาเรนฺติ, สกฺกจฺจํ ธาตานํ ธมฺมานํ อตฺถํ อุปปริกฺขนฺติ, สกฺกจฺจํ อตฺถมญฺญาย ธมฺมมญฺญาย ธมฺมานุธมฺมํ ปฏิปชฺชนฺติ.\csromanspacer{} อิเม โข ภิกฺขเว ปญฺจ ธมฺมา สทฺธมฺมสฺส ฐิติยา อสมฺโมสาย อนนฺตรธานาย สํวตฺตนฺตีติ.\csromanspacer{}\csromanspacer{}จตุตฺถํ.}

\csromanheader{2}{5. ทุติยสทฺธมฺมสมฺโมสสุตฺต}
\proseitem{155}{ปญฺจิเม ภิกฺขเว ธมฺมา สทฺธมฺมสฺส สมฺโมสาย อนฺตรธานาย สํวตฺตนฺติ.\csromanspacer{} กตเม ปญฺจ? อิธ ภิกฺขเว ภิกฺขู ธมฺมํ น ปริยาปุณนฺติ สุตฺตํ เคยฺยํ เวยฺยากรณํ คาถํ อุทานํ อิติวุตฺตกํ ชาตกํ อพฺภุตธมฺมํ เวทลฺลํ.\csromanspacer{} อยํ ภิกฺขเว ปฐโม ธมฺโม สทฺธมฺมสฺส สมฺโมสาย อนฺตรธานาย สํวตฺตติ.}

\prose{ปุน จปรํ ภิกฺขเว ภิกฺขู ยถาสุตํ ยถาปริยตฺตํ ธมฺมํ น วิตฺถาเรน ปเรสํ เทเสนฺติ.\csromanspacer{} อยํ ภิกฺขเว ทุติโย ธมฺโม สทฺธมฺมสฺส สมฺโมสาย อนฺตรธานาย สํวตฺตติ.}

\prose{ปุน จปรํ ภิกฺขเว ภิกฺขู ยถาสุตํ ยถาปริยตฺตํ ธมฺมํ น วิตฺถาเรน ปรํ\footnote{ปเรสํ (สี, สฺยา, กํ, อิ), ปเร (?)} วาเจนฺติ.\csromanspacer{} อยํ ภิกฺขเว ตติโย ธมฺโม สทฺธมฺมสฺส สมฺโมสาย อนฺตรธานาย สํวตฺตติ.}

\prose{ปุน จปรํ ภิกฺขเว ภิกฺขู ยถาสุตํ ยถาปริยตฺตํ ธมฺมํ น วิตฺถาเรน สชฺฌายํ กโรนฺติ.\csromanspacer{} อยํ ภิกฺขเว จตุตฺโถ ธมฺโม สทฺธมฺมสฺส สมฺโมสาย อนฺตรธานาย สํวตฺตติ.}

\prose{ปุน จปรํ ภิกฺขเว ภิกฺขู ยถาสุตํ ยถาปริยตฺตํ ธมฺมํ น เจตสา อนุวิตกฺเกนฺติ อนุวิจาเรนฺติ มนสานุเปกฺขนฺติ.\csromanspacer{} อยํ ภิกฺขเว ปญฺจโม ธมฺโม สทฺธมฺมสฺส สมฺโมสาย อนฺตรธานาย สํวตฺตติ.\csromanspacer{} อิเม โข ภิกฺขเว ปญฺจ ธมฺมา สทฺธมฺมสฺส สมฺโมสาย อนฺตรธานาย สํวตฺตนฺติ.}

\prose{ปญฺจิเม ภิกฺขเว ธมฺมา สทฺธมฺมสฺส ฐิติยา อสมฺโมสาย อนนฺตร ธานาย สํวตฺตนฺติ.\csromanspacer{} กตเม ปญฺจ? อิธ ภิกฺขเว ภิกฺขู ธมฺมํ ปริยาปุณนฺติ สุตฺตํ เคยฺยํ เวยฺยากรณํ คาถํ อุทานํ อิติวุตฺตกํ ชาตกํ อพฺภุตธมฺมํ เวทลฺลํ.\csromanspacer{} อยํ ภิกฺขเว ปฐโม ธมฺโม สทฺธมฺมสฺส ฐิติยา อสมฺโมสาย อนนฺตรธานาย สํวตฺตติ.}

\chapterheadpageread{(16) 1. สทฺธมฺมวคฺค}
\prose{\csromanfolio{157}ปุน จปรํ ภิกฺขเว ภิกฺขู ยถาสุตํ ยถาปริยตฺตํ ธมฺมํ วิตฺถาเรน ปเรสํ เทเสนฺติ.\csromanspacer{} อยํ ภิกฺขเว ทุติโย ธมฺโม สทฺธมฺมสฺส ฐิติยา อสมฺโมสาย อนนฺตรธานาย สํวตฺตติ.}

\prose{ปุน จปรํ ภิกฺขเว ภิกฺขู ยถาสุตํ ยถาปริยตฺตํ ธมฺมํ วิตฺถาเรน ปรํ วาเจนฺติ.\csromanspacer{} อยํ ภิกฺขเว ตติโย ธมฺโม สทฺธมฺมสฺส ฐิติยา อสมฺโมสาย อนนฺตรธานาย สํวตฺตติ.}

\prose{ปุน จปรํ ภิกฺขเว ภิกฺขู ยถาสุตํ ยถาปริยตฺตํ ธมฺมํ วิตฺถาเรน สชฺฌายํ กโรนฺติ.\csromanspacer{} อยํ ภิกฺขเว จตุตฺโถ ธมฺโม สทฺธมฺมสฺส ฐิติยา อสมฺโมสาย อนนฺตรธานาย สํวตฺตติ.}

\prose{ปุน จปรํ ภิกฺขเว ภิกฺขู ยถาสุตํ ยถาปริยตฺตํ ธมฺมํ เจตสา อนุวิตกฺเกนฺติ อนุวิจาเรนฺติ มนสานุเปกฺขนฺติ.\csromanspacer{} อยํ ภิกฺขเว ปญฺจโม ธมฺโม สทฺธมฺมสฺส ฐิติยา อสมฺโมสาย อนนฺตรธานาย สํวตฺตติ.\csromanspacer{} อิเม โข ภิกฺขเว ปญฺจ ธมฺมา สทฺธมฺมสฺส ฐิติยา อสมฺโมสาย อนนฺตรธานาย สํวตฺตนฺตีติ.\csromanspacer{}\csromanspacer{}ปญฺจมํ.}

\csromanheader{2}{6. ตติยสทฺธมฺมสมฺโมสสุตฺต}
\proseitem{156}{\csromansymbolfootnote{*}{อํ 1. 465 ปิฏฺเฐปิ.}ปญฺจิเม ภิกฺขเว ธมฺมา สทฺธมฺมสฺส สมฺโมสาย อนฺตรธานาย สํวตฺตนฺติ.\csromanspacer{} กตเม ปญฺจ? อิธ ภิกฺขเว ภิกฺขู ทุคฺคหิตํ สุตฺตนฺตํ ปริยาปุณนฺติ ทุนฺนิกฺขิตฺเตหิ ปทพฺยญฺชเนหิ.\csromanspacer{} ทุนฺนิกฺขิตฺตสฺส ภิกฺขเว ปทพฺยญฺชนสฺส อตฺโถปิ ทุนฺนโย โหติ.\csromanspacer{} อยํ ภิกฺขเว ปฐโม ธมฺโม สทฺธมฺมสฺส สมฺโมสาย อนฺตรธานาย สํวตฺตติ.}

\prose{ปุน จปรํ ภิกฺขเว ภิกฺขู ทุพฺพจา โหนฺติ โทวจสฺสกรเณหิ ธมฺเมหิ สมนฺนาคตา อกฺขมา อปฺปทกฺขิณคฺคาหิโน อนุสาสนิํ.\csromanspacer{} อยํ ภิกฺขเว ทุติโย ธมฺโม สทฺธมฺมสฺส สมฺโมสาย อนฺตรธานาย สํวตฺตติ.}

\prose{ปุน จปรํ ภิกฺขเว เย เต ภิกฺขู พหุสฺสุตา อาคตาคมา ธมฺมธรา วินยธรา มาติกาธรา, เต น สกฺกจฺจํ สุตฺตนฺตํ ปรํ วาเจนฺติ, เตสํ อจฺจเยน ฉินฺนมูลโก สุตฺตนฺโต โหติ อปฺปฏิสรโณ.\csromanspacer{} อยํ ภิกฺขเว ตติโย ธมฺโม สทฺธมฺมสฺส สมฺโมสาย อนฺตรธานาย สํวตฺตติ.}

\prose{ปุน จปรํ ภิกฺขเว เถรา ภิกฺขู พาหุลิกา โหนฺติ สาถลิกา โอกฺกมเน ปุพฺพงฺคมา ปวิเวเก นิกฺขิตฺตธุรา, น วีริยํ อารภนฺติ อปฺปตฺตสฺส \csromanfolio{158}ปตฺติยา อนธิคตสฺส อธิคมาย อสจฺฉิกตสฺส สจฺฉิกิริยาย, เตสํ ปจฺฉิมา ชนตา ทิฏฺฐานุคติํ อาปชฺชติ.\csromanspacer{} สาปิ โหติ พาหุลิกา สาถลิกา โอกฺกมเน ปุพฺพงฺคมา ปวิเวเก นิกฺขิตฺตธุรา, น วีริยํ อารภติ อปฺปตฺตสฺส ปตฺติยา อนธิคตสฺส อธิคมาย อสจฺฉิกตสฺส สจฺฉิกิริยาย.\csromanspacer{} อยํ ภิกฺขเว จตุตฺโถ ธมฺโม สทฺธมฺมสฺส สมฺโมสาย อนฺตรธานาย สํวตฺตติ.}

\prose{ปุน จปรํ ภิกฺขเว สํโฆ ภินฺโน โหติ, สํเฆ โข ปน ภิกฺขเว ภินฺเน อญฺญมญฺญํ อกฺโกสา จ โหนฺติ, อญฺญมญฺญํ ปริภาสา จ โหนฺติ, อญฺญมญฺญํ ปริกฺเขปา จ โหนฺติ, อญฺญมญฺญํ ปริจฺจชนา\footnote{ปริจฺจชา (สฺยา, กํ)} จ โหนฺติ.\csromanspacer{} ตตฺถ อปฺปสนฺนา เจว นปฺปสีทนฺติ, ปสนฺนานญฺจ เอกจฺจานํ อญฺญถตฺตํ โหติ.\csromanspacer{} อยํ ภิกฺขเว ปญฺจโม ธมฺโม สทฺธมฺมสฺส สมฺโมสาย อนฺตรธานาย สํวตฺตติ.\csromanspacer{} อิเม โข ภิกฺขเว ปญฺจ ธมฺมา สทฺธมฺมสฺส สมฺโมสาย อนฺตรธานาย สํวตฺตนฺติ.}

\prose{ปญฺจิเม ภิกฺขเว ธมฺมา สทฺธมฺมสฺส ฐิติยา อสมฺโมสาย อนนฺตร ธานาย สํวตฺตนฺติ.\csromanspacer{} กตเม ปญฺจ? อิธ ภิกฺขเว ภิกฺขู สุคฺคหิตํ สุตฺตนฺตํ ปริยาปุณนฺติ สุนิกฺขิตฺเตหิ ปทพฺยญฺชเนหิ.\csromanspacer{} สุนิกฺขิตฺตสฺส ภิกฺขเว ปทพฺยญฺชนสฺส อตฺโถปิ สุนโย โหติ.\csromanspacer{} อยํ ภิกฺขเว ปฐโม ธมฺโม สทฺธมฺมสฺส ฐิติยา อสมฺโมสาย อนนฺตรธานาย สํวตฺตติ.}

\prose{ปุน จปรํ ภิกฺขเว ภิกฺขู สุวจา โหนฺติ, โสวจสฺสกรเณหิ ธมฺเมหิ สมนฺนาคตา ขมา ปทกฺขิณคฺคาหิโน อนุสาสนิํ.\csromanspacer{} อยํ ภิกฺขเว ทุติโย ธมฺโม สทฺธมฺมสฺส ฐิติยา อสมฺโมสาย อนนฺตรธานาย สํวตฺตติ.}

\prose{ปุน จปรํ ภิกฺขเว เย เต ภิกฺขู พหุสฺสุตา อาคตาคมา ธมฺมธรา วินยธรา มาติกาธรา, เต สกฺกจฺจํ สุตฺตนฺตํ ปรํ วาเจนฺติ, เตสํ อจฺจเยน นจฺฉินฺนมูลโก\footnote{อจฺฉินฺนมูลโก (ก) อํ 1. 466 ปิฏฺเฐปิ.} สุตฺตนฺโต โหติ สปฺปฏิสรโณ.\csromanspacer{} อยํ ภิกฺขเว ตติโย ธมฺโม สทฺธมฺมสฺส ฐิติยา อสมฺโมสาย อนนฺตรธานาย สํวตฺตติ.}

\prose{ปุน จปรํ ภิกฺขเว เถรา ภิกฺขู น พาหุลิกา โหนฺติ น สาถลิกา โอกฺกมเน นิกฺขิตฺตธุรา ปวิเวเก ปุพฺพงฺคมา, วีริยํ อารภนฺติ อปฺปตฺตสฺส}

\vspace{\tipitakaparskip}
\csromancenter{(16) 1.\csromanspacer{} สทฺธมฺมวคฺค ปตฺติยา อนธิคตสฺส อธิคมาย อสจฺฉิกตสฺส สจฺฉิกิริยาย, เตสํ ปจฺฉิมา ชนตา ทิฏฺฐานุคติํ อาปชฺชติ, สาปิ โหติ น พาหุลิกา น สาถลิกา โอกฺกมเน นิกฺขิตฺตธุรา ปวิเวเก ปุพฺพงฺคมา, วีริยํ อารภติ อปฺปตฺตสฺส ปตฺติยา อนธิคตสฺส อธิคมาย อสจฺฉิกตสฺส สจฺฉิกิริยาย.\csromanspacer{} อยํ ภิกฺขเว จตุตฺโถ ธมฺโม สทฺธมฺมสฺส ฐิติยา อสมฺโมสาย อนนฺตรธานาย สํวตฺตติ.}
\prose{\csromanfolio{159}ปุน จปรํ ภิกฺขเว สํโฆ สมคฺโค สมฺโมทมาโน อวิวทมาโน เอกุทฺเทโส ผาสุํ วิหรติ, สํเฆ โข ปน ภิกฺขเว สมคฺเค น เจว อญฺญมญฺญํ อกฺโกสา โหนฺติ, น จ อญฺญมญฺญํ ปริภาสา โหนฺติ, น จ อญฺญมญฺญํ ปริกฺเขปา โหนฺติ, น จ อญฺญมญฺญํ ปริจฺจชนา โหนฺติ, ตตฺถ อปฺปสนฺนา เจว ปสีทนฺติ, ปสนฺนานญฺจ ภิยฺโยภาโว โหติ.\csromanspacer{} อยํ ภิกฺขเว ปญฺจโม ธมฺโม สทฺธมฺมสฺส ฐิติยา อสมฺโมสาย อนนฺตรธานาย สํวตฺตติ.\csromanspacer{} อิเม โข ภิกฺขเว ปญฺจ ธมฺมา สทฺธมฺมสฺส ฐิติยา อสมฺโมสาย อนนฺตรธานาย สํวตฺตนฺตีติ.\csromanspacer{}\csromanspacer{}ฉฏฺฐํ.}

\csromanmatikaanchor{mtk.51}
\csromantocmarkref{subsection}{7. ทุกฺกถาสุตฺต}{mtk.51}
\bookmark[dest={mtk.51},level=2]{7. ทุกฺกถาสุตฺต}
\csromanheader{2}{7. ทุกฺกถาสุตฺต}
\proseitem{157}{ปญฺจนฺนํ ภิกฺขเว ปุคฺคลานํ กถา ทุกฺกถา ปุคฺคเล ปุคฺคลํ\footnote{ปุคฺคลํ ปุคฺคลํ (สี, อิ)} อุปนิธาย.\csromanspacer{} กตเมสํ ปญฺจนฺนํ? อสฺสทฺธสฺส ภิกฺขเว สทฺธากถา ทุกฺกถา, ทุสฺสีลสฺส สีลกถา ทุกฺกถา, อปฺปสฺสุตสฺส พาหุสจฺจกถา ทุกฺกถา, มจฺฉริสฺส\footnote{มจฺฉริยสฺส (สี, อิ, ก)} จาคกถา ทุกฺกถา, ทุปฺปญฺญสฺส ปญฺญากถา ทุกฺกถา.}

\prose{กสฺมา จ ภิกฺขเว อสฺสทฺธสฺส สทฺธากถา ทุกฺกถา? อสฺสทฺโธ ภิกฺขเว สทฺธากถาย กจฺฉมานาย อภิสชฺชติ กุปฺปติ พฺยาปชฺชติ ปติตฺถียติ, โกปญฺจ โทสญฺจ อปฺปจฺจยญฺจ ปาตุกโรติ.\csromanspacer{} ตํ กิสฺส เหตุ? ตํ หิ โส ภิกฺขเว สทฺธาสมฺปทํ อตฺตนิ น สมนุปสฺสติ\footnote{น สมฺปสฺสติ (สี)}, น จ ลภติ ตโตนิทานํ ปีติปาโมชฺชํ.\csromanspacer{} ตสฺมา อสฺสทฺธสฺส สทฺธากถา ทุกฺกถา.}

\prose{กสฺมา จ ภิกฺขเว ทุสฺสีลสฺส สีลกถา ทุกฺกถา? ทุสฺสีโล ภิกฺขเว สีลกถาย กจฺฉมานาย อภิสชฺชติ กุปฺปติ พฺยาปชฺชติ ปติตฺถียติ, โกปญฺจ โทสญฺจ อปฺปจฺจยญฺจ ปาตุกโรติ.\csromanspacer{} ตํ กิสฺส เหตุ? ตํ หิ \csromanfolio{160}โส ภิกฺขเว สีลสมฺปทํ อตฺตนิ น สมนุปสฺสติ, น จ ลภติ ตโตนิทานํ ปีติปาโมชฺชํ.\csromanspacer{} ตสฺมา ทุสฺสีลสฺส สีลกถา ทุกฺกถา.}

\prose{กสฺมา จ ภิกฺขเว อปฺปสฺสุตสฺส พาหุสจฺจกถา ทุกฺกถา? อปฺปสฺสุโต ภิกฺขเว พาหุสจฺจกถาย กจฺฉมานาย อภิสชฺชติ กุปฺปติ พฺยาปชฺชติ ปติตฺถียติ, โกปญฺจ โทสญฺจ อปฺปจฺจยญฺจ ปาตุกโรติ.\csromanspacer{} ตํ กิสฺส เหตุ? ตํ หิ โส ภิกฺขเว สุตสมฺปทํ อตฺตนิ น สมนุปสฺสติ, น จ ลภติ ตโตนิทานํ ปีติปาโมชฺชํ.\csromanspacer{} ตสฺมา อปฺปสฺสุตสฺส พาหุสจฺจกถา ทุกฺกถา.}

\prose{กสฺมา จ ภิกฺขเว มจฺฉริสฺส จาคกถา ทุกฺกถา? มจฺฉรี ภิกฺขเว จาคกถาย กจฺฉมานาย อภิสชฺชติ กุปฺปติ พฺยาปชฺชติ ปติตฺถียติ, โกปญฺจ โทสญฺจ อปฺปจฺจยญฺจ ปาตุกโรติ.\csromanspacer{} ตํ กิสฺส เหตุ? ตํ หิ โส ภิกฺขเว จาคสมฺปทํ อตฺตนิ น สมนุปสฺสติ, น จ ลภติ ตโตนิทานํ ปีติปาโมชฺชํ.\csromanspacer{} ตสฺมา มจฺฉริสฺส จาคกถา ทุกฺกถา.}

\prose{กสฺมา จ ภิกฺขเว ทุปฺปญฺญสฺส ปญฺญากถา ทุกฺกถา? ทุปฺปญฺโญ ภิกฺขเว ปญฺญากถาย กจฺฉมานาย อภิสชฺชติ กุปฺปติ พฺยาปชฺชติ ปติตฺถียติ, โกปญฺจ โทสญฺจ อปฺปจฺจยญฺจ ปาตุกโรติ.\csromanspacer{} ตํ กิสฺส เหตุ? ตํ หิ โส ภิกฺขเว ปญฺญาสมฺปทํ อตฺตนิ น สมนุปสฺสติ, น จ ลภติ ตโตนิทานํ ปีติปาโมชฺชํ.\csromanspacer{} ตสฺมา ทุปฺปญฺญสฺส ปญฺญากถา ทุกฺกถา.\csromanspacer{} อิเมสํ โข ภิกฺขเว ปญฺจนฺนํ ปุคฺคลานํ กถา ทุกฺกถา ปุคฺคเล ปุคฺคลํ อุปนิธาย.}

\prose{ปญฺจนฺนํ ภิกฺขเว ปุคฺคลานํ กถา สุกถา ปุคฺคเล ปุคฺคลํ อุปนิธาย.\csromanspacer{} กตเมสํ ปญฺจนฺนํ? สทฺธสฺส ภิกฺขเว สทฺธากถา สุกถา, สีลวโต สีลกถา สุกถา, พหุสฺสุตสฺส พาหุสจฺจกถา สุกถา, จาควโต จาคกถา สุกถา, ปญฺญวโต ปญฺญากถา สุกถา.}

\prose{กสฺมา จ ภิกฺขเว สทฺธสฺส สทฺธากถา สุกถา? สทฺโธ ภิกฺขเว สทฺธากถาย กจฺฉมานาย นาภิสชฺชติ น กุปฺปติ น พฺยาปชฺชติ น ปติตฺถียติ, น โกปญฺจ โทสญฺจ อปฺปจฺจยญฺจ ปาตุกโรติ.\csromanspacer{} ตํ กิสฺส เหตุ? ตํ หิ โส ภิกฺขเว สทฺธาสมฺปทํ อตฺตนิ สมนุปสฺสติ, ลภติ จ ตโตนิทานํ ปีติปาโมชฺชํ.\csromanspacer{} ตสฺมา สทฺธสฺส สทฺธากถา สุกถา.}

\prose{กสฺมา จ ภิกฺขเว สีลวโต สีลกถา สุกถา? สีลวา ภิกฺขเว สีลกถาย กจฺฉมานาย นาภิสชฺชติ น กุปฺปติ น พฺยาปชฺชติ}

\vspace{\tipitakaparskip}
\csromancenter{(16) 1.\csromanspacer{} สทฺธมฺมวคฺค น ปติตฺถียติ, น โกปญฺจ โทสญฺจ อปฺปจฺจยญฺจ ปาตุกโรติ.\csromanspacer{} ตํ กิสฺส เหตุ? ตํ หิ โส ภิกฺขเว สีลสมฺปทํ อตฺตนิ สมนุปสฺสติ, ลภติ จ ตโตนิทานํ ปีติปาโมชฺชํ.\csromanspacer{} ตสฺมา สีลวโต สีลกถา สุกถา.}
\prose{\csromanfolio{161}กสฺมา จ ภิกฺขเว พหุสฺสุตสฺส พาหุสจฺจกถา สุกถา? พหุสฺสุโต ภิกฺขเว พาหุสจฺจกถาย กจฺฉมานาย นาภิสชฺชติ น กุปฺปติ น พฺยาปชฺชติ น ปติตฺถียติ, น โกปญฺจ โทสญฺจ อปฺปจฺจยญฺจ ปาตุกโรติ.\csromanspacer{} ตํ กิสฺส เหตุ? ตํ หิ โส ภิกฺขเว สุตสมฺปทํ อตฺตนิ สมนุปสฺสติ, ลภติ จ ตโตนิทานํ ปีติปาโมชฺชํ.\csromanspacer{} ตสฺมา พหุสฺสุตสฺส พาหุสจฺจกถา สุกถา.}

\prose{กสฺมา จ ภิกฺขเว จาควโต จาคกถา สุกถา? จาควา ภิกฺขเว จาคกถาย กจฺฉมานาย นาภิสชฺชติ น กุปฺปติ น พฺยาปชฺชติ น ปติตฺถียติ, น โกปญฺจ โทสญฺจ อปฺปจฺจยญฺจ ปาตุกโรติ.\csromanspacer{} ตํ กิสฺส เหตุ? ตํ หิ โส ภิกฺขเว จาคสมฺปทํ อตฺตนิ สมนุปสฺสติ, ลภติ จ ตโตนิทานํ ปีติปาโมชฺชํ.\csromanspacer{} ตสฺมา จาควโต จาคกถา สุกถา.}

\prose{กสฺมา จ ภิกฺขเว ปญฺญวโต ปญฺญากถา สุกถา? ปญฺญวา ภิกฺขเว ปญฺญากถาย กจฺฉมานาย นาภิสชฺชติ น กุปฺปติ น พฺยาปชฺชติ น ปติตฺถียติ, น โกปญฺจ โทสญฺจ อปฺปจฺจยญฺจ ปาตุกโรติ.\csromanspacer{} ตํ กิสฺส เหตุ? ตํ หิ โส ภิกฺขเว ปญฺญาสมฺปทํ อตฺตนิ สมนุปสฺสติ, ลภติ จ ตโตนิทานํ ปีติปาโมชฺชํ.\csromanspacer{} ตสฺมา ปญฺญวโต ปญฺญากถา สุกถา.\csromanspacer{} อิเมสํ โข ภิกฺขเว ปญฺจนฺนํ ปุคฺคลานํ กถา สุกถา ปุคฺคเล ปุคฺคลํ อุปนิธายาติ.\csromanspacer{}\csromanspacer{}สตฺตมํ.}

\csromanmatikaanchor{mtk.52}
\csromantocmarkref{subsection}{8. สารชฺชสุตฺต}{mtk.52}
\bookmark[dest={mtk.52},level=2]{8. สารชฺชสุตฺต}
\csromanheader{2}{8. สารชฺชสุตฺต}
\proseitem{158}{ปญฺจหิ ภิกฺขเว ธมฺเมหิ สมนฺนาคโต ภิกฺขุ สารชฺชํ โอกฺกนฺโต\footnote{โอกฺกมนฺโต (ก)} โหติ.\csromanspacer{} กตเมหิ ปญฺจหิ? อิธ ภิกฺขเว ภิกฺขุ อสฺสทฺโธ โหติ, ทุสฺสีโล โหติ, อปฺปสฺสุโต โหติ, กุสีโต โหติ, ทุปฺปญฺโญ โหติ.\csromanspacer{} อิเมหิ โข ภิกฺขเว ปญฺจหิ ธมฺเมหิ สมนฺนาคโต ภิกฺขุ สารชฺชํ โอกฺกนฺโต โหติ.}

\prose{\csromanfolio{162}ปญฺจหิ ภิกฺขเว ธมฺเมหิ สมนฺนาคโต ภิกฺขุ วิสารโท โหติ.\csromanspacer{} กตเมหิ ปญฺจหิ, อิธ ภิกฺขเว ภิกฺขุ สทฺโธ โหติ, สีลวา โหติ, พหุสฺสุโต โหติ, อารทฺธวีริโย โหติ, ปญฺญวา โหติ.\csromanspacer{} อิเมหิ โข ภิกฺขเว ปญฺจหิ ธมฺเมหิ สมนฺนาคโต ภิกฺขุ วิสารโท โหตีติ.\csromanspacer{}\csromanspacer{}อฏฺฐมํ.}

\csromanmatikaanchor{mtk.53}
\csromantocmarkref{subsection}{9. อุทายีสุตฺต}{mtk.53}
\bookmark[dest={mtk.53},level=2]{9. อุทายีสุตฺต}
\csromanheader{2}{9. อุทายีสุตฺต}
\proseitem{159}{เอวํ เม สุตํ\csromandash{}เอกํ สมยํ ภควา โกสมฺพิยํ วิหรติ โฆสิตาราเม.\csromanspacer{} เตน โข ปน สมเยน อายสฺมา อุทายี มหติยา คิหิปริสาย ปริวุโต ธมฺมํ เทเสนฺโต นิสินฺโน โหติ.\csromanspacer{} อทฺทสา โข อายสฺมา อานนฺโท อายสฺมนฺตํ อุทายิํ มหติยา คิหิปริสาย ปริวุตํ ธมฺมํ เทเสนฺตํ นิสินฺนํ.\csromanspacer{} ทิสฺวา เยน ภควา เตนุปสงฺกมิ, อุปสงฺกมิตฺวา ภควนฺตํ อภิวาเทตฺวา เอกมนฺตํ นิสีทิ, เอกมนฺตํ นิสินฺโน โข อายสฺมา อานนฺโท ภควนฺตํ เอตทโวจ “อายสฺมา ภนฺเต อุทายี มหติยา คิหิปริสาย ปริวุโต ธมฺมํ เทเสตี”ติ\footnote{เทเสนฺโต นิสินฺโน”ติ (สฺยา)}.}

\prose{น โข อานนฺท สุกรํ ปเรสํ ธมฺมํ เทเสตุํ, ปเรสํ อานนฺท ธมฺมํ เทเสนฺเตน ปญฺจ ธมฺเม อชฺฌตฺตํ อุปฏฺฐาเปตฺวา ปเรสํ ธมฺโม เทเสตพฺโพ.\csromanspacer{} กตเม ปญฺจ? “อนุปุพฺพิํ กถํ\footnote{อานุปุพฺพิกถํ (สี), อนุปุพฺพิกถํ (สฺยา, อิ, ก)} กเถสฺสามี”ติ ปเรสํ ธมฺโม เทเสตพฺโพ, “ปริยายทสฺสาวี กถํ กเถสฺสามี”ติ ปเรสํ ธมฺโม เทเสตพฺโพ, “อนุทฺทยตํ ปฏิจฺจ กถํ กเถสฺสามี”ติ ปเรสํ ธมฺโม เทเสตพฺโพ, “น อามิสนฺตโร กถํ กเถสฺสามี”ติ ปเรสํ ธมฺโม เทเสตพฺโพ, “อตฺตานญฺจ ปรญฺจ อนุปหจฺจ กถํ กเถสฺสามี”ติ ปเรสํ ธมฺโม เทเสตพฺโพ.\csromanspacer{} น โข อานนฺท สุกรํ ปเรสํ ธมฺมํ เทเสตุํ, ปเรสํ อานนฺท ธมฺมํ เทเสนฺเตน อิเม ปญฺจ ธมฺเม อชฺฌตฺตํ อุปฏฺฐาเปตฺวา ปเรสํ ธมฺโม เทเสตพฺโพติ.\csromanspacer{}\csromanspacer{}นวมํ.}

\csromanmatikaanchor{mtk.54}
\csromantocmarkref{subsection}{10. ทุปฺปฏิวิโนทยสุตฺต}{mtk.54}
\bookmark[dest={mtk.54},level=2]{10. ทุปฺปฏิวิโนทยสุตฺต}
\csromanheader{2}{10. ทุปฺปฏิวิโนทยสุตฺต}
\proseitem{160}{ปญฺจิเม ภิกฺขเว อุปฺปนฺนา ทุปฺปฏิวิโนทยา.\csromanspacer{} กตเม ปญฺจ? อุปฺปนฺโน ราโค ทุปฺปฏิวิโนทโย, อุปฺปนฺโน โทโส ทุปฺปฏิวิโนทโย, อุปฺปนฺโน}

\vspace{\tipitakaparskip}
\csromancenter{(17) 2.\csromanspacer{} อาคาตวคฺค โมโห ทุปฺปฏิวิโนทโย, อุปฺปนฺนํ ปฏิภานํ ทุปฺปฏิวิโนทยํ, อุปฺปนฺนํ คมิกจิตฺตํ ทุปฺปฏิวิโนทยํ.\csromanspacer{} อิเม โข ภิกฺขเว ปญฺจ อุปฺปนฺนา ทุปฺปฏิวิโนทยาติ.\csromanspacer{}\csromanspacer{}ทสมํ.\csromanspacer{} สทฺธมฺมวคฺโค ปฐโม.}
\titlehead{\textbf{ตสฺสุทฺทานํ}}
\csromanmatikaanchor{mtk.55}
\csromanmatikaanchor{mtk.56}
\csromantocmarknopage{section}{(17) 2. อาฆาตวคฺค}{mtk.55}
\bookmark[dest={mtk.55},level=1]{(17) 2. อาฆาตวคฺค}
\csromantocmarkref{subsection}{1-2. อาฆาตปฏิวินยสุตฺตทฺวย}{mtk.56}
\bookmark[dest={mtk.56},level=2]{1-2. อาฆาตปฏิวินยสุตฺตทฺวย}
\csromangathameasure{ตโย สมฺมตฺตนิยามา, ตโย สทฺธมฺมสมฺโมสา. \\ ทุกฺกถา เจว สารชฺชํ, อุทายิทุพฺพิโนทยาติ.}
\csromangathagroup{\csromangathastanza{\csromanfolio{163}ตโย สมฺมตฺตนิยามา, ตโย สทฺธมฺมสมฺโมสา. \\ ทุกฺกถา เจว สารชฺชํ, อุทายิทุพฺพิโนทยาติ.}}

\csromanheader{1}{\textbf{(17) 2. อาฆาตวคฺค}}
\csromanheader{2}{1. ปฐม อาฆาตปฏิวินยสุตฺต}
\proseitem{161}{ปญฺจิเม ภิกฺขเว อาฆาตปฏิวินยา, ยตฺถ ภิกฺขุโน อุปฺปนฺโน อาฆาโต สพฺพโส ปฏิวิเนตพฺโพ.\csromanspacer{} กตเม ปญฺจ? ยสฺมิํ ภิกฺขเว ปุคฺคเล อาฆาโต ชาเยถ, เมตฺตา ตสฺมิํ ปุคฺคเล ภาเวตพฺพา, เอวํ ตสฺมิํ ปุคฺคเล อาฆาโต ปฏิวิเนตพฺโพ.\csromanspacer{} ยสฺมิํ ภิกฺขเว ปุคฺคเล อาฆาโต ชาเยถ, กรุณา ตสฺมิํ ปุคฺคเล ภาเวตพฺพา, เอวํ ตสฺมิํ ปุคฺคเล อาฆาโต ปฏิวิเนตพฺโพ.\csromanspacer{} ยสฺมิํ ภิกฺขเว ปุคฺคเล อาฆาโต ชาเยถ, อุเปกฺขา ตสฺมิํ ปุคฺคเล ภาเวตพฺพา, เอวํ ตสฺมิํ ปุคฺคเล อาฆาโต ปฏิวิเนตพฺโพ.\csromanspacer{} ยสฺมิํ ภิกฺขเว ปุคฺคเล อาฆาโต ชาเยถ, อสติ-อมนสิกาโร ตสฺมิํ ปุคฺคเล อาปชฺชิตพฺโพ, เอวํ ตสฺมิํ ปุคฺคเล อาฆาโต ปฏิวิเนตพฺโพ.\csromanspacer{} ยสฺมิํ ภิกฺขเว ปุคฺคเล อาฆาโต ชาเยถ, กมฺมสฺสกตา ตสฺมิํ ปุคฺคเล อธิฏฺฐาตพฺพา, “กมฺมสฺสโก อยมายสฺมา กมฺมทายาโท กมฺมโยนิ กมฺมพนฺธุ กมฺมปฺปฏิสรโณ, ยํ กมฺมํ กริสฺสติ กลฺยาณํ วา ปาปกํ วา, ตสฺส ทายาโท ภวิสฺสตี”ติ, เอวํ ตสฺมิํ ปุคฺคเล อาฆาโต ปฏิวิเนตพฺโพ.\csromanspacer{} อิเม โข ภิกฺขเว ปญฺจ อาฆาตปฏิวินยา, ยตฺถ ภิกฺขุโน อุปฺปนฺโน อาฆาโต สพฺพโส ปฏิวิเนตพฺโพติ.\csromanspacer{}\csromanspacer{}ปฐมํ.}

\csromanheader{2}{2. ทุติย อาฆาตปฏิวินยสุตฺต}
\proseitem{162}{\csromanfolio{164}ตตฺร โข อายสฺมา สาริปุตฺโต ภิกฺขู อามนฺเตสิ “อาวุโส ภิกฺขเว”ติ. “อาวุโส”ติ โข เต ภิกฺขู อายสฺมโต สาริปุตฺตสฺส ปจฺจสฺโสสุํ.\csromanspacer{} อายสฺมา สาริปุตฺโต เอตทโวจ\csromandash{}}

\prose{ปญฺจิเม อาวุโส อาฆาตปฏิวินยา, ยตฺถ ภิกฺขุโน อุปฺปนฺโน อาฆาโต สพฺพโส ปฏิวิเนตพฺโพ.\csromanspacer{} กตเม ปญฺจ? อิธาวุโส เอกจฺโจ ปุคฺคโล อปริสุทฺธกายสมาจาโร โหติ ปริสุทฺธวจีสมาจาโร, เอวรูเปปิ อาวุโส ปุคฺคเล อาฆาโต ปฏิวิเนตพฺโพ.\csromanspacer{} อิธ ปนาวุโส เอกจฺโจ ปุคฺคโล อปริสุทฺธวจีสมาจาโร โหติ ปริสุทฺธกายสมาจาโร, เอวรูเปปิ อาวุโส ปุคฺคเล อาฆาโต ปฏิวิเนตพฺโพ.\csromanspacer{} อิธ ปนาวุโส เอกจฺโจ ปุคฺคโล อปริสุทฺธกายสมาจาโร โหติ อปริสุทฺธวจีสมาจาโร, ลภติ จ กาเลน กาลํ เจตโส วิวรํ เจตโส ปสาทํ.\csromanspacer{} เอวรูเปปิ อาวุโส ปุคฺคเล อาฆาโต ปฏิวิเนตพฺโพ.\csromanspacer{} อิธ ปนาวุโส เอกจฺโจ ปุคฺคโล อปริสุทฺธกายสมาจาโร โหติ อปริสุทฺธวจีสมาจาโร, น จ ลภติ กาเลน กาลํ เจตโส วิวรํ เจตโส ปสาทํ.\csromanspacer{} เอวรูเปปิ อาวุโส ปุคฺคเล อาฆาโต ปฏิวิเนตพฺโพ.\csromanspacer{} อิธ ปนาวุโส เอกจฺโจ ปุคฺคโล ปริสุทฺธกายสมาจาโร ปริสุทฺธวจีสมาจาโร, ลภติ จ กาเลน กาลํ เจตโส วิวรํ เจตโส ปสาทํ.\csromanspacer{} เอวรูเปปิ อาวุโส ปุคฺคเล อาฆาโต ปฏิวิเนตพฺโพ.}

\prose{ตตฺราวุโส ยฺวายํ ปุคฺคโล อปริสุทฺธกายสมาจาโร ปริสุทฺธวจีสมาจาโร, กถํ ตสฺมิํ ปุคฺคเล อาฆาโต ปฏิวิเนตพฺโพ? เสยฺยถาปิ อาวุโส ภิกฺขุ ปํสุกูลิโก รถิยาย นนฺตกํ ทิสฺวา วาเมน ปาเทน นิคฺคณฺหิตฺวา ทกฺขิเณน ปาเทน ปตฺถริตฺวา\footnote{วิตฺถาเรตฺวา (สี, อิ)} โย ตตฺถ สาโร, ตํ ปริปาเตตฺวา อาทาย ปกฺกเมยฺย.\csromanspacer{} เอวเมวํ ขฺวาวุโส ยฺวายํ ปุคฺคโล อปริสุทฺธกายสมาจาโร ปริสุทฺธวจีสมาจาโร.\csromanspacer{} ยา’สฺส อปริสุทฺธกายสมาจารตา, น สา’สฺส ตสฺมิํ สมเย มนสิ กาตพฺพา.}

\vspace{\tipitakaparskip}
\csromancenter{(17) 2.\csromanspacer{} อาคาตวคฺค ยา จ ขฺวาสฺส ปริสุทฺธวจีสมาจารตา, สา’สฺส ตสฺมิํ สมเย มนสิ กาตพฺพา.\csromanspacer{} เอวํ ตสฺมิํ ปุคฺคเล อาฆาโต ปฏิวิเนตพฺโพ. (1)}
\prose{\csromanfolio{165}ตตฺราวุโส ยฺวายํ ปุคฺคโล อปริสุทฺธวจีสมาจาโร ปริสุทฺธกายสมาจาโร, กถํ ตสฺมิํ ปุคฺคเล อาฆาโต ปฏิวิเนตพฺโพ? เสยฺยถาปิ อาวุโส โปกฺขรณี เสวาลปณกปริโยนทฺธา, อถ ปุริโส อาคจฺเฉยฺย ฆมฺมาภิตตฺโต ฆมฺมปเรโต กิลนฺโต ตสิโต ปิปาสิโต, โส ตํ โปกฺขรณิํ โอคาเหตฺวา อุโภหิ หตฺเถหิ อิติจิติจ เสวาลปณกํ อปวิยูหิตฺวา อญฺชลินา ปิวิตฺวา ปกฺกเมยฺย.\csromanspacer{} เอวเมวํ โข อาวุโส ยฺวายํ ปุคฺคโล อปริสุทฺธวจีสมาจาโร ปริสุทฺธกายสมาจาโร.\csromanspacer{} ยา’สฺส อปริสุทฺธวจีสมาจารตา, น สา’สฺส ตสฺมิํ สมเย มนสิ กาตพฺพา.\csromanspacer{} ยา จ ขฺวาสฺส ปริสุทฺธกายสมาจารตา, สา’สฺส ตสฺมิํ สมเย มนสิ กาตพฺพา.\csromanspacer{} เอวํ ตสฺมิํ ปุคฺคเล อาฆาโต ปฏิวิเนตพฺโพ. (2)}

\prose{ตตฺราวุโส ยฺวายํ ปุคฺคโล อปริสุทฺธกายสมาจาโร อปริสุทฺธ วจีสมาจาโร, ลภติ จ กาเลน กาลํ เจตโส วิวรํ เจตโส ปสาทํ, กถํ ตสฺมิํ ปุคฺคเล อาฆาโต ปฏิวิเนตพฺโพ? เสยฺยถาปิ อาวุโส ปริตฺตํ โคปเท\footnote{โคปทเก (สี, สฺยา)} อุทกํ, อถ ปุริโส อาคจฺเฉยฺย ฆมฺมาภิตตฺโต ฆมฺมปเรโต กิลนฺโต ตสิโต ปิปาสิโต, ตสฺส เอวมสฺส “อิทํ โข ปริตฺตํ โคปเท อุทกํ, สจาหํ อญฺชลินา วา ปิวิสฺสามิ ภาชเนน วา, โขเภสฺสามิปิ ตํ โลเฬสฺสามิปิ ตํ อเปยฺยมฺปิ ตํ กริสฺสามิ, ยํนูนาหํ จตุกฺกุณฺฑิโก\footnote{จตุคุณฺฑิโก (สี), จตุกุณฺฑิโก (สฺยา, กํ, อิ), จตุโกณฺฑิโก (ที 3. 4 ปิฏฺเฐ)} นิปติตฺวา โคปีตกํ ปิวิตฺวา ปกฺกเมยฺยนฺ”ติ.\csromanspacer{} โส จตุกฺกุณฺฑิโก นิปติตฺวา โคปีตกํ ปิวิตฺวา ปกฺกเมยฺย.\csromanspacer{} เอวเมวํ โข อาวุโส ยฺวายํ ปุคฺคโล อปริสุทฺธกายสมาจาโร อปริสุทฺธวจีสมาจาโร, ลภติ จ กาเลน กาลํ เจตโส วิวรํ เจตโส ปสาทํ, ยา’สฺส อปริสุทฺธกายสมาจารตา, น สา’สฺส ตสฺมิํ สมเย มนสิ กาตพฺพา.\csromanspacer{} ยาปิสฺส อปริสุทฺธวจีสมาจารตา, น สาปิสฺส ตสฺมิํ สมเย มนสิ กาตพฺพา.\csromanspacer{} ยญฺจ โข โส ลภติ กาเลน กาลํ เจตโส วิวรํ เจตโส ปสาทํ, ตเมวสฺส\footnote{ตเทวสฺส (สี, สฺยา)} ตสฺมิํ สมเย มนสิ กาตพฺพํ.\csromanspacer{} เอวํ ตสฺมิํ ปุคฺคเล อาฆาโต ปฏิวิเนตพฺโพ. (3)}

\prose{\csromanfolio{166}ตตฺราวุโส ยฺวายํ ปุคฺคโล อปริสุทฺธกายสมาจาโร อปริสุทฺธ วจีสมาจาโร, น จ ลภติ กาเลน กาลํ เจตโส วิวรํ เจตโส ปสาทํ, กถํ ตสฺมิํ ปุคฺคเล อาฆาโต ปฏิวิเนตพฺโพ? เสยฺยถาปิ อาวุโส ปุริโส อาพาธิโก ทุกฺขิโต พาฬฺหคิลาโน อทฺธานมคฺคปฺปฏิปนฺโน, ตสฺส ปุรโตปิสฺส ทูเร คาโม, ปจฺฉโตปิสฺส ทูเร คาโม, โส น ลเภยฺย สปฺปายานิ โภชนานิ, น ลเภยฺย สปฺปายานิ เภสชฺชานิ, น ลเภยฺย ปติรูปํ อุปฏฺฐากํ, น ลเภยฺย คามนฺตนายกํ.\csromanspacer{} ตเมนํ อญฺญตโร ปุริโส ปสฺเสยฺย อทฺธานมคฺคปฺปฏิปนฺโน, โส ตสฺมิํ ปุริเส การุญฺญํเยว อุปฏฺฐาเปยฺย, อนุทฺทยํเยว อุปฏฺฐาเปยฺย, อนุกมฺปํเยว อุปฏฺฐาเปยฺย, อโห วตายํ ปุริโส ลเภยฺย สปฺปายานิ โภชนานิ, ลเภยฺย สปฺปายานิ เภสชฺชานิ, ลเภยฺย ปติรูปํ อุปฏฺฐากํ, ลเภยฺย คามนฺตนายกํ.\csromanspacer{} ตํ กิสฺส เหตุ? มายํ\footnote{อยํ (ก)} ปุริโส อิเธว อนยพฺยสนํ อาปชฺชีติ\footnote{อาปชฺเชยฺย (ก)}.\csromanspacer{} เอวเมวํ โข อาวุโส ยฺวายํ ปุคฺคโล อปริสุทฺธกายสมาจาโร อปริสุทฺธวจีสมาจาโร, น จ ลภติ กาเลน กาลํ เจตโส วิวรํ เจตโส ปสาทํ.\csromanspacer{} เอวรูเปปิ\footnote{เอวรูเป (อิ)} อาวุโส ปุคฺคเล การุญฺญํเยว อุปฏฺฐาเปตพฺพํ, อนุทฺทยาเยว อุปฏฺฐาเปตพฺพา, อนุกมฺปาเยว อุปฏฺฐาเปตพฺพา, อโห วต อยมายสฺมา กายทุจฺจริตํ ปหาย กายสุจริตํ ภาเวยฺย, วจีทุจฺจริตํ ปหาย วจีสุจริตํ ภาเวยฺย, มโนทุจฺจริตํ ปหาย มโนสุจริตํ ภาเวยฺย.\csromanspacer{} ตํ กิสฺส เหตุ? มายํ อายสฺมา\footnote{อยมายสฺมา (ก)} กายสฺส เภทา ปรํ มรณา อปายํ ทุคฺคติํ วินิปาตํ นิรยํ อุปปชฺชีติ\footnote{อุปปชฺชตีติ (ก)}.\csromanspacer{} เอวํ ตสฺมิํ ปุคฺคเล อาฆาโต ปฏิวิเนตพฺโพ. (4)}

\prose{ตตฺราวุโส ยฺวายํ ปุคฺคโล ปริสุทฺธกายสมาจาโร ปริสุทฺธวจีสมาจาโร, ลภติ จ กาเลน กาลํ เจตโส วิวรํ เจตโส ปสาทํ, กถํ ตสฺมิํ ปุคฺคเล อาฆาโต ปฏิวิเนตพฺโพ? เสยฺยถาปิ อาวุโส โปกฺขรณี อจฺโฉทกา สาโตทกา สีโตทกา\footnote{อจฺโฉทิกา สาโตทิกา สีโตทิกา (สี)} เสตกา\footnote{เสโตทกา (ก)} สุปติตฺถา รมณียา นานารุกฺเขหิ สญฺฉนฺนา, อถ ปุริโส อาคจฺเฉยฺย ฆมฺมาภิตตฺโต ฆมฺมปเรโต กิลนฺโต ตสิโต}

\vspace{\tipitakaparskip}
\csromancenter{(17) 2.\csromanspacer{} อาคาตวคฺค ปิปาสิโต, โส ตํ โปกฺขรณิํ โอคาเหตฺวา นฺหาตฺวา จ ปิวิตฺวา จ ปจฺจุตฺตริตฺวา ตตฺเถว รุกฺขจฺฉายาย นิสีเทยฺย วา นิปชฺเชยฺย วา.}
\csromanmatikaanchor{mtk.57}
\csromantocmarkref{subsection}{3-5. สากจฺฉสุตฺตาทิ}{mtk.57}
\bookmark[dest={mtk.57},level=2]{3-5. สากจฺฉสุตฺตาทิ}
\prose{\csromanfolio{167}เอวเมว โข อาวุโส ยฺวายํ ปุคฺคโล ปริสุทฺธกายสมาจาโร ปริสุทฺธ วจีสมาจาโร, ลภติ จ กาเลน กาลํ เจตโส วิวรํ เจตโส ปสาทํ.\csromanspacer{} ยาปิสฺส ปริสุทฺธกายสมาจารตา, สาปิสฺส ตสฺมิํ สมเย มนสิ กาตพฺพา.\csromanspacer{} ยาปิสฺส ปริสุทฺธวจีสมาจารตา, สาปิสฺส ตสฺมิํ สมเย มนสิ กาตพฺพา.\csromanspacer{} ยมฺปิ ลภติ กาเลน กาลํ เจตโส วิวรํ เจตโส ปสาทํ, ตมฺปิสฺส ตสฺมิํ สมเย มนสิ กาตพฺพํ.\csromanspacer{} เอวํ ตสฺมิํ ปุคฺคเล อาฆาโต ปฏิวิเนตพฺโพ.\csromanspacer{} สมนฺตปาสาทิกํ อาวุโส ปุคฺคลํ อาคมฺม จิตฺตํ ปสีทติ. (5)}

\prose{อิเม โข อาวุโส ปญฺจ อาฆาตปฏิวินยา, ยตฺถ ภิกฺขุโน อุปฺปนฺโน อาฆาโต สพฺพโส ปฏิวิเนตพฺโพติ.\csromanspacer{}\csromanspacer{}ทุติยํ.}

\csromanheader{2}{3. สากจฺฉสุตฺต}
\proseitem{163}{ตตฺร โข อายสฺมา สาริปุตฺโต ภิกฺขู อามนฺเตสิ “อาวุโส ภิกฺขเว”ติ. “อาวุโส”ติ โข เต ภิกฺขู อายสฺมโต สาริปุตฺตสฺส ปจฺจสฺโสสุํ.\csromanspacer{} อายสฺมา สาริปุตฺโต เอตทโวจ\csromandash{}}

\prose{ปญฺจหาวุโส ธมฺเมหิ สมนฺนาคโต ภิกฺขุ อลํสากจฺโฉ สพฺรหฺมจารีนํ.\csromanspacer{} กตเมหิ ปญฺจหิ, อิธาวุโส ภิกฺขุ อตฺตนา จ สีลสมฺปนฺโน โหติ, สีลสมฺปทากถาย จ อาคตํ ปญฺหํ พฺยากตฺตา โหติ.\csromanspacer{} อตฺตนา จ สมาธิสมฺปนฺโน โหติ, สมาธิสมฺป ทากถาย จ อาคตํ ปญฺหํ พฺยากตฺตา โหติ.\csromanspacer{} อตฺตนา จ ปญฺญาสมฺปนฺโน โหติ, ปญฺญา สมฺปทากถาย จ อาคตํ ปญฺหํ พฺยากตฺตา โหติ.\csromanspacer{} อตฺตนา จ วิมุตฺติสมฺปนฺโน โหติ, วิมุตฺติสมฺปทากถาย จ อาคตํ ปญฺหํ พฺยากตฺตา โหติ.\csromanspacer{} อตฺตนา จ วิมุตฺติญาณ ทสฺสนสมฺปนฺโน โหติ, วิมุตฺติญาณทสฺสนสมฺปทากถาย จ อาคตํ ปญฺหํ พฺยากตฺตา โหติ.\csromanspacer{} อิเมหิ โข อาวุโส ปญฺจหิ ธมฺเมหิ สมนฺนาคโต ภิกฺขุ อลํสากจฺโฉ สพฺรหฺมจารีนนฺติ.\csromanspacer{}\csromanspacer{}ตติยํ.}

\csromanheader{2}{4. สาชีวสุตฺต}
\proseitem{164}{\csromanfolio{168}\csromansymbolfootnote{*}{เหฏฺฐา 71 ปิฏฺเฐปิ.}ตตฺร โข อายสฺมา สาริปุตฺโต ภิกฺขู อามนฺเตสิ ฯเปฯ.\csromanspacer{} ปญฺจหิ อาวุโส ธมฺเมหิ สมนฺนาคโต ภิกฺขุ อลํสาชีโว สพฺรหฺมจารีนํ.\csromanspacer{} กตเมหิ ปญฺจหิ? อิธาวุโส ภิกฺขุ อตฺตนา จ สีลสมฺปนฺโน โหติ, สีลสมฺปทากถาย จ อาคตํ ปญฺหํ พฺยากตฺตา โหติ.\csromanspacer{} อตฺตนา จ สมาธิสมฺปนฺโน โหติ, สมาธิสมฺปทากถาย จ อาคตํ ปญฺหํ พฺยากตฺตา โหติ.\csromanspacer{} อตฺตนา จ ปญฺญาสมฺปนฺโน โหติ, ปญฺญาสมฺปทากถาย จ อาคตํ ปญฺหํ พฺยากตฺตา โหติ.\csromanspacer{} อตฺตนา จ วิมุตฺติสมฺปนฺโน โหติ, วิมุตฺติสมฺป ทากถาย จ อาคตํ ปญฺหํ พฺยากตฺตา โหติ.\csromanspacer{} อตฺตนา จ วิมุตฺติญาณทสฺสนสมฺปนฺโน โหติ, วิมุตฺติญาณทสฺสนสมฺปทากถาย จ อาคตํ ปญฺหํ พฺยากตฺตา โหติ.\csromanspacer{} อิเมหิ โข อาวุโส ปญฺจหิ ธมฺเมหิ สมนฺนาคโต ภิกฺขุ อลํสาชีโว สพฺรหฺมจารีนนฺติ.\csromanspacer{}\csromanspacer{}จตุตฺถํ.}

\csromanheader{2}{5. ปญฺหปุจฺฉาสุตฺต}
\proseitem{165}{ตตฺร โข อายสฺมา สาริปุตฺโต ภิกฺขู อามนฺเตสิ ฯเปฯ.\csromanspacer{} โย หิ โกจิ อาวุโส ปรํ ปญฺหํ ปุจฺฉติ, สพฺโพ โส ปญฺจหิ ฐาเนหิ เอเตสํ วา อญฺญตเรน.\csromanspacer{} กตเมหิ ปญฺจหิ? มนฺทตฺตา โมมูหตฺตา\footnote{โมมุหตฺตา (สี)} ปรํ ปญฺหํ ปุจฺฉติ.\csromanspacer{} ปาปิจฺโฉ อิจฺฉาปกโต ปรํ ปญฺหํ ปุจฺฉติ, ปริภวํ ปรํ ปญฺหํ ปุจฺฉติ, อญฺญาตุกาโม ปรํ ปญฺหํ ปุจฺฉติ, อถ วา ปเนวํจิตฺโต\footnote{อถ วา ปกุปฺปนฺโต (สี, อิ)} ปรํ ปญฺหํ ปุจฺฉติ “สเจ เม ปญฺหํ ปุฏฺโฐ สมฺมเทว พฺยากริสฺสติ, อิจฺเจตํ กุสลํ.\csromanspacer{} โน เจ\footnote{โน จ (สฺยา)} เม ปญฺหํ ปุฏฺโฐ สมฺมเทว พฺยากริสฺสติ, อหมสฺส สมฺมเทว พฺยากริสฺสามี”ติ.\csromanspacer{} โย หิ โกจิ อาวุโส ปรํ ปญฺหํ ปุจฺฉติ, สพฺโพ โส อิเมหิ ปญฺจหิ ฐาเนหิ เอเตสํ วา อญฺญตเรน.\csromanspacer{} อหํ โข ปนาวุโส เอวํจิตฺโต ปรํ ปญฺหํ ปุจฺฉามิ “สเจ เม ปญฺหํ ปุฏฺโฐ สมฺมเทว พฺยากริสฺสติ, อิจฺเจตํ กุสลํ.\csromanspacer{} โน เจ เม ปญฺหํ ปุฏฺโฐ สมฺมเทว พฺยากริสฺสติ, อหมสฺส สมฺมเทว พฺยากริสฺสามี”ติ.\csromanspacer{}\csromanspacer{}ปญฺจมํ.}

\csromanmatikaanchor{mtk.58}
\csromantocmarkref{subsection}{6. นิโรธสุตฺต}{mtk.58}
\bookmark[dest={mtk.58},level=2]{6. นิโรธสุตฺต}
\csromanheader{2}{6. นิโรธสุตฺต}
\proseitem{166}{ตตฺร โข อายสฺมา สาริปุตฺโต ภิกฺขู อามนฺเตสิ ฯเปฯ.\csromanspacer{} อิธาวุโส ภิกฺขุ สีลสมฺปนฺโน สมาธิสมฺปนฺโน ปญฺญาสมฺปนฺโน สญฺญา}

\vspace{\tipitakaparskip}
\csromancenter{(17) 2.\csromanspacer{} อาคาตวคฺค เวทยิตนิโรธํ สมาปชฺเชยฺยาปิ วุฏฺฐเหยฺยาปิ\footnote{สมาปชฺเชยฺยปิ วุฏฺฐเหยฺยปิ (สี, สฺยา, กํ, อิ) 2. กพฬิํการาหารภกฺขานํ (สี, สฺยา, กํ, อิ)}, อตฺเถตํ ฐานํ.\csromanspacer{} โน เจ ทิฏฺเฐว ธมฺเม อญฺญํ อาราเธยฺย, อติกฺกมฺเมว กพฬีการาหารภกฺขานํ2 เทวานํ สหพฺยตํ อญฺญตรํ มโนมยํ กายํ อุปปนฺโน สญฺญาเวทยิตนิโรธํ สมาปชฺเชยฺยาปิ วุฏฺฐเหยฺยาปิ, อตฺเถตํ ฐานนฺติ.}
\prose{\csromanfolio{169}เอวํ วุตฺเต อายสฺมา อุทายี อายสฺมนฺตํ สาริปุตฺตํ เอตทโวจ “อฏฺฐานํ โข เอตํ อาวุโส สาริปุตฺต อนวกาโส, ยํ โส ภิกฺขุ อติกฺกมฺเมว กพฬีการาหาร ภกฺขานํ เทวานํ สหพฺยตํ อญฺญตรํ มโนมยํ กายํ อุปปนฺโน สญฺญาเวทยิตนิโรธํ สมาปชฺเชยฺยาปิ, วุฏฺฐเหยฺยาปิ, นตฺเถตํ ฐานนฺ”ติ.}

\prose{ทุติยมฺปิ โข ฯเปฯ.\csromanspacer{} ตติยมฺปิ โข อายสฺมา สาริปุตฺโต ภิกฺขู อามนฺเตสิ “อิธาวุโส ภิกฺขุ สีลสมฺปนฺโน สมาธิสมฺปนฺโน ปญฺญาสมฺปนฺโน สญฺญาเวทยิต นิโรธํ สมาปชฺเชยฺยาปิ วุฏฺฐเหยฺยาปิ, อตฺเถตํ ฐานํ.\csromanspacer{} โน เจ ทิฏฺเฐว ธมฺเม อญฺญํ อาราเธยฺย, อติกฺกมฺเมว กพฬีการาหารภกฺขานํ เทวานํ สหพฺยตํ อญฺญตรํ มโนมยํ กายํ อุปปนฺโน สญฺญาเวทยิตนิโรธํ สมาปชฺเชยฺยาปิ วุฏฺฐเหยฺยาปิ, อตฺเถตํ ฐานนฺ”ติ.}

\prose{ตติยมฺปิ โข อายสฺมา อุทายี อายสฺมนฺตํ สาริปุตฺตํ เอตทโวจ “อฏฺฐานํ โข เอตํ อาวุโส สาริปุตฺต อนวกาโส, ยํ โส ภิกฺขุ อติกฺกมฺเมว กพฬีการาหาร ภกฺขานํ เทวานํ สหพฺยตํ อญฺญตรํ มโนมยํ กายํ อุปปนฺโน สญฺญาเวทยิตนิโรธํ สมาปชฺเชยฺยาปิ วุฏฺฐเหยฺยาปิ, นตฺเถตํ ฐานนฺ”ติ}

\prose{อถ โข อายสฺมโต สาริปุตฺตสฺส เอตทโหสิ “ยาวตติยกมฺปิ\footnote{ยาวตติยมฺปิ (สี, สฺยา, อิ)} โข เม อายสฺมา อุทายี ปฏิกฺโกสติ, น จ เม โกจิ ภิกฺขุ อนุโมทติ, ยํนูนาหํ เยน ภควา เตนุปสงฺกเมยฺยนฺ”ติ.\csromanspacer{} อถ โข อายสฺมา สาริปุตฺโต เยน ภควา เตนุปสงฺกมิ, อุปสงฺกมิตฺวา ภควนฺตํ อภิวาเทตฺวา เอกมนฺตํ นิสีทิ, เอกมนฺตํ นิสินฺโน โข อายสฺมา สาริปุตฺโต ภิกฺขู อามนฺเตสิ “อิธาวุโส ภิกฺขุ สีลสมฺปนฺโน สมาธิ สมฺปนฺโน ปญฺญาสมฺปนฺโน สญฺญาเวทยิตนิโรธํ สมาปชฺเชยฺยาปิ วุฏฺฐเหยฺยาปิ, อตฺเถตํ ฐานํ, โน เจ ทิฏฺเฐว ธมฺเม \csromanfolio{170}อญฺญํ อาราเธยฺย, อติกฺกมฺเมว กพฬีการาหารภกฺขานํ เทวานํ สหพฺยตํ อญฺญตรํ มโนมยํ กายํ อุปปนฺโน สญฺญาเวทยิตนิโรธํ สมาปชฺเชยฺยาปิ วุฏฺฐเหยฺยาปิ, อตฺเถตํ ฐานนฺ”ติ.}

\prose{เอวํ วุตฺเต อายสฺมา อุทายี อายสฺมนฺตํ สาริปุตฺตํ เอตทโวจ “อฏฺฐานํ โข เอตํ อาวุโส สาริปุตฺต อนวกาโส, ยํ โส ภิกฺขุ อติกฺกมฺเมว กพฬีการาหาร ภกฺขานํ เทวานํ สหพฺยตํ อญฺญตรํ มโนมยํ กายํ อุปปนฺโน สญฺญาเวทยิตนิโรธํ สมาปชฺเชยฺยาปิ วุฏฺฐเหยฺยาปิ, นตฺเถตํ ฐานนฺ”ติ.}

\prose{ทุติยมฺปิ โข ฯเปฯ.\csromanspacer{} ตติยมฺปิ โข อายสฺมา สาริปุตฺโต ภิกฺขู อามนฺเตสิ “อิธาวุโส ภิกฺขุ สีลสมฺปนฺโน สมาธิสมฺปนฺโน ปญฺญาสมฺปนฺโน สญฺญาเวทยิตนิโรธํ สมาปชฺเชยฺยาปิ วุฏฺฐเหยฺยาปิ, อตฺเถตํ ฐานํ.\csromanspacer{} โน เจ ทิฏฺเฐว ธมฺเม อญฺญํ อาราเธยฺย, อติกฺกมฺเมว กพฬีการาหารภกฺขานํ เทวานํ สหพฺยตํ อญฺญตรํ มโนมยํ กายํ อุปปนฺโน สญฺญาเวทยิต นิโรธํ สมาปชฺเชยฺยาปิ วุฏฺฐเหยฺยาปิ, อตฺเถตํ ฐานนฺ”ติ.}

\prose{ตติยมฺปิ โข อายสฺมา อุทายี อายสฺมนฺตํ สาริปุตฺตํ เอตทโวจ “อฏฺฐานํ โข เอตํ อาวุโส สาริปุตฺต อนวกาโส, ยํ โส ภิกฺขุ อติกฺกมฺเมว กพฬีการาหาร ภกฺขานํ เทวานํ สหพฺยตํ อญฺญตรํ มโนมยํ กายํ อุปปนฺโน สญฺญา เวทยิตนิโรธํ สมาปชฺเชยฺยาปิ วุฏฺฐเหยฺยาปิ, นตฺเถตํ ฐานนฺ”ติ.}

\prose{อถ โข อายสฺมโต สาริปุตฺตสฺส เอตทโหสิ “ภควโตปิ โข เม สมฺมุขา อายสฺมา อุทายี ยาวตติยกํ ปฏิกฺโกสติ, น จ เม โกจิ ภิกฺขุ อนุโมทติ, ยํนูนาหํ ตุณฺหี อสฺสนฺ”ติ.\csromanspacer{} อถ โข อายสฺมา สาริปุตฺโต ตุณฺหี อโหสิ.}

\prose{อถ โข ภควา อายสฺมนฺตํ อุทายิํ อามนฺเตสิ “กํ ปน ตฺวํ อุทายิ มโนมยํ กายํ ปจฺเจสี”ติ.\csromanspacer{} เย เต ภนฺเต เทวา อรูปิโน สญฺญามยาติ.\csromanspacer{} กิํ นุ โข ตุยฺหํ อุทายิ พาลสฺส อพฺยตฺตสฺส ภณิเตน.\csromanspacer{} ตฺวมฺปิ นาม ภณิตพฺพํ มญฺญสีติ.\csromanspacer{} อถ โข ภควา อายสฺมนฺตํ อานนฺทํ อามนฺเตสิ “อตฺถิ นาม อานนฺท เถรํ ภิกฺขุํ วิเหสิยมานํ}

\vspace{\tipitakaparskip}
\csromancenter{(17) 2.\csromanspacer{} อาคาตวคฺค อชฺฌุเปกฺขิสฺสถ.\csromanspacer{} น หิ นาม อานนฺท การุญฺญมฺปิ ภวิสฺสติ เถรมฺหิ\footnote{พฺยตฺตมฺหิ (สฺยา, กํ, ก)} ภิกฺขุมฺหิ วิเหสิยมานมฺหี”ติ.}
\prose{\csromanfolio{171}อถ โข ภควา ภิกฺขู อามนฺเตสิ “อิธ ภิกฺขเว ภิกฺขุ สีลสมฺปนฺโน สมาธิสมฺปนฺโน ปญฺญาสมฺปนฺโน สญฺญาเวทยิตนิโรธํ สมาปชฺเชยฺยาปิ วุฏฺฐเหยฺยาปิ, อตฺเถตํ ฐานํ.\csromanspacer{} โน เจ ทิฏฺเฐว ธมฺเม อญฺญํ อาราเธยฺย, อติกฺกมฺเมว กพฬีการาหารภกฺขานํ เทวานํ สหพฺยตํ อญฺญตรํ มโนมยํ กายํ อุปปนฺโน สญฺญาเวทยิตนิโรธํ สมาปชฺเชยฺยาปิ วุฏฺฐเหยฺยาปิ, อตฺเถตํ ฐานนฺ”ติ.\csromanspacer{} อิทมโวจ ภควา, อิทํ วตฺวาน สุคโต อุฏฺฐายาสนา วิหารํ ปาวิสิ.}

\prose{อถ โข อายสฺมา อานนฺโท อจิรปกฺกนฺตสฺส ภควโต เยนายสฺมา อุปวาโน เตนุปสงฺกมิ, อุปสงฺกมิตฺวา อายสฺมนฺตํ อุปวานํ เอตทโวจ “อิธาวุโส อุปวาน อญฺเญ เถเร ภิกฺขู วิเหเสนฺติ, มยํ เตน น มุจฺจาม.\csromanspacer{} อนจฺฉริยํ โข ปเนตํ อาวุโส อุปวาน, ยํ ภควา สายนฺหสมยํ ปฏิสลฺลานา วุฏฺฐิโต เอตเทว อารพฺภ อุทาหเรยฺย ยถา อายสฺมนฺตํเยเวตฺถ อุปวานํ ปฏิภาเสยฺย อิทาเนว อมฺหากํ สารชฺชํ โอกฺกนฺตนฺ”ติ.\csromanspacer{} อถ โข ภควา สายนฺหสมยํ ปฏิสลฺลานา วุฏฺฐิโต เยน อุปฏฺฐานสาลา เตนุปสงฺกมิ, อุปสงฺกมิตฺวา ปญฺญตฺเต อาสเน นิสีทิ, นิสชฺช โข ภควา อายสฺมนฺตํ อุปวานํ เอตทโวจ\csromandash{}}

\csromanmatikaanchor{mtk.59}
\csromantocmarkref{subsection}{7-9. โจทนาสุตฺตาทิ}{mtk.59}
\bookmark[dest={mtk.59},level=2]{7-9. โจทนาสุตฺตาทิ}
\prose{กติหิ นุ โข อุปวาน ธมฺเมหิ สมนฺนาคโต เถโร ภิกฺขุ สพฺรหฺม จารีนํ ปิโย จ โหติ มนาโป จ ครุ จ ภาวนีโย จาติ.\csromanspacer{} ปญฺจหิ ภนฺเต ธมฺเมหิ สมนฺนาคโต เถโร ภิกฺขุ สพฺรหฺมจารีนํ ปิโย จ โหติ มนาโป จ ครุ จ ภาวนีโย จ.\csromanspacer{} กตเมหิ ปญฺจหิ, อิธ ภนฺเต เถโร ภิกฺขุ สีลวา โหติ ฯเปฯ สมาทาย สิกฺขติ สิกฺขาปเทสุ.\csromanspacer{} พหุสฺสุโต โหติ ฯเปฯ ทิฏฺฐิยา สุปฺปฏิวิทฺธา.\csromanspacer{} กลฺยาณวาโจ โหติ กลฺยาณวากฺกรโณ โปริยา วาจาย สมนฺนาคโต วิสฺสฏฺฐาย อเนลคลาย\footnote{อเนลคฬาย (สฺยา, กํ)} อตฺถสฺส วิญฺญาปนิยา.\csromanspacer{} จตุนฺนํ ฌานานํ อาภิเจตสิกานํ ทิฏฺฐธมฺมสุขวิหารานํ นิกามลาภี โหติ อกิจฺฉลาภี อกสิรลาภี.\csromanspacer{} อาสวานํ ขยา ฯเปฯ สจฺฉิกตฺวา อุปสมฺปชฺช วิหรติ.\csromanspacer{} อิเมหิ โข ภนฺเต ปญฺจหิ ธมฺเมหิ สมนฺนาคโต เถโร \csromanfolio{172}ภิกฺขุ สพฺรหฺมจารีนํ ปิโย จ โหติ มนาโป จ ครุ จ ภาวนีโย จาติ.}

\prose{สาธุ สาธุ อุปวาน, อิเมหิ โข อุปวาน ปญฺจหิ ธมฺเมหิ สมนฺนาคโต เถโร ภิกฺขุ สพฺรหฺมจารีนํ ปิโย จ โหติ มนาโป จ ครุ จ ภาวนีโย จ.\csromanspacer{} อิเม เจ อุปวาน ปญฺจ ธมฺมา เถรสฺส ภิกฺขุโน น สํวิชฺเชยฺยุํ, ตํ สพฺรหฺมจารี น สกฺกเรยฺยุํ น ครุํ กเรยฺยุํ น มาเนยฺยุํ น ปูเชยฺยุํ ขณฺฑิจฺเจน ปาลิจฺเจน วลิตฺตจตาย.\csromanspacer{} ยสฺมา จ โข อุปวาน อิเม ปญฺจ ธมฺมา เถรสฺส ภิกฺขุโน สํวิชฺชนฺติ, ตสฺมา ตํ สพฺรหฺมจารี สกฺกโรนฺติ ครุํ กโรนฺติ มาเนนฺติ ปูเชนฺตีติ.\csromanspacer{}\csromanspacer{}ฉฏฺฐํ.}

\csromanheader{2}{7. โจทนาสุตฺต}
\proseitem{167}{ตตฺร โข อายสฺมา สาริปุตฺโต ภิกฺขู อามนฺเตสิ\csromandash{}โจทเกน อาวุโส ภิกฺขุนา ปรํ โจเทตุกาเมน ปญฺจ ธมฺเม อชฺฌตฺตํ อุปฏฺฐาเปตฺวา ปโร โจเทตพฺโพ.}

\prose{กตเม ปญฺจ? กาเลน วกฺขามิ โน อกาเลน, ภูเตน วกฺขามิ โน อภูเตน, สเณฺหน วกฺขามิ โน ผรุเสน, อตฺถสํหิเตน วกฺขามิ โน อนตฺถสํหิเตน, เมตฺตจิตฺโต\footnote{เมตฺตจิตฺเตน (สี, อิ, ก) วิ 4. 438 ปิฏฺเฐ ปสฺสิตพฺพํ.} วกฺขามิ โน โทสนฺตโร\footnote{โทสนฺตเรน (สี, อิ, ก)}.\csromanspacer{} โจทเกน อาวุโส ภิกฺขุนา ปรํ โจเทตุ กาเมน อิเม ปญฺจ ธมฺเม อชฺฌตฺตํ อุปฏฺฐาเปตฺวา ปโร โจเทตพฺโพ.}

\prose{อิธาหํ อาวุโส เอกจฺจํ ปุคฺคลํ ปสฺสามิ อกาเลน โจทิยมานํ โน กาเลน กุปิตํ, อภูเตน โจทิยมานํ โน ภูเตน กุปิตํ, ผรุเสน โจทิยมานํ โน สเณฺหน กุปิตํ, อนตฺถสํหิเตน โจทิยมานํ โน อตฺถสํหิเตน กุปิตํ, โทสนฺตเรน โจทิยมานํ โน เมตฺตจิตฺเตน กุปิตํ.}

\prose{อธมฺมจุทิตสฺส อาวุโส ภิกฺขุโน ปญฺจหากาเรหิ อวิปฺปฏิสาโร อุปทหา ตพฺโพ\footnote{อุปทหิตพฺโพ (สี, สฺยา, อิ)}.\csromanspacer{} อกาเลนายสฺมา จุทิโต โน กาเลน, อลํ เต อวิปฺปฏิสาราย.\csromanspacer{} อภูเตนายสฺมา จุทิโต โน ภูเตน, อลํ เต อวิปฺปฏิสาราย.\csromanspacer{} ผรุเสนายสฺมา จุทิโต โน สเณฺหน, อลํ เต}

\vspace{\tipitakaparskip}
\csromancenter{(17) 2.\csromanspacer{} อาคาตวคฺค อวิปฺปฏิสาราย.\csromanspacer{} อนตฺถสํหิเตนายสฺมา จุทิโต โน อตฺถสํหิเตน, อลํ เต อวิปฺปฏิสาราย.\csromanspacer{} โทสนฺตเรนายสฺมา จุทิโต โน เมตฺตจิตฺเตน, อลํ เต อวิปฺปฏิสารายาติ.\csromanspacer{} อธมฺมจุทิตสฺส อาวุโส ภิกฺขุโน อิเมหิ ปญฺจหากาเรหิ อวิปฺปฏิสาโร อุปทหาตพฺโพ.}
\prose{\csromanfolio{173}อธมฺมโจทกสฺส อาวุโส ภิกฺขุโน ปญฺจหากาเรหิ วิปฺปฏิสาโร อุปทหา ตพฺโพ.\csromanspacer{} อกาเลน เต อาวุโส โจทิโต\footnote{จุทิโต (สี, สฺยา, อิ)} โน กาเลน, อลํ เต วิปฺปฏิสาราย.\csromanspacer{} อภูเตน เต อาวุโส โจทิโต โน ภูเตน, อลํ เต วิปฺปฏิสาราย.\csromanspacer{} ผรุเสน เต อาวุโส โจทิโต โน สเณฺหน, อลํ เต วิปฺปฏิสาราย.\csromanspacer{} อนตฺถสํหิเตน เต อาวุโส โจทิโต โน อตฺถสํหิเตน, อลํ เต วิปฺปฏิสาราย.\csromanspacer{} โทสนฺตเรน เต อาวุโส โจทิโต โน เมตฺตจิตฺเตน, อลํ เต วิปฺปฏิสารายาติ.\csromanspacer{} อธมฺมโจทกสฺส อาวุโส ภิกฺขุโน อิเมหิ ปญฺจหากาเรหิ วิปฺปฏิสาโร อุปทหาตพฺโพ.\csromanspacer{} ตํ กิสฺส เหตุ? ยถา น อญฺโญปิ ภิกฺขุ อภูเตน โจเทตพฺพํ มญฺเญยฺยาติ.}

\prose{อิธ ปนาหํ อาวุโส เอกจฺจํ ปุคฺคลํ ปสฺสามิ กาเลน โจทิยมานํ โน อกาเลน กุปิตํ.\csromanspacer{} ภูเตน โจทิยมานํ โน อภูเตน กุปิตํ.\csromanspacer{} สเณฺหน โจทิยมานํ โน ผรุเสน กุปิตํ.\csromanspacer{} อตฺถสํหิเตน โจทิยมานํ โน อนตฺถสํหิเตน กุปิตํ.\csromanspacer{} เมตฺตจิตฺเตน โจทิยมานํ โน โทสนฺตเรน กุปิตํ.}

\prose{ธมฺมจุทิตสฺส อาวุโส ภิกฺขุโน ปญฺจหากาเรหิ วิปฺปฏิสาโร อุปทหาตพฺโพ.\csromanspacer{} กาเลนายสฺมา จุทิโต โน อกาเลน, อลํ เต วิปฺปฏิสาราย.\csromanspacer{} ภูเตนายสฺมา จุทิโต โน อภูเตน, อลํ เต วิปฺปฏิสาราย.\csromanspacer{} สเณฺหนายสฺมา จุทิโต โน ผรุเสน, อลํ เต วิปฺปฏิสาราย.\csromanspacer{} อตฺถสํหิเตนายสฺมา จุทิโต โน อนตฺถสํหิเตน, อลํ เต วิปฺปฏิสาราย.\csromanspacer{} เมตฺตจิตฺเตนายสฺมา จุทิโต โน โทสนฺตเรน, อลํ เต วิปฺปฏิสารายาติ.\csromanspacer{} ธมฺม จุทิตสฺส อาวุโส ภิกฺขุโน อิเมหิ ปญฺจหากาเรหิ วิปฺปฏิสาโร อุปทหาตพฺโพ.}

\prose{ธมฺมโจทกสฺส อาวุโส ภิกฺขุโน ปญฺจหากาเรหิ อวิปฺปฏิสาโร อุปทหาตพฺโพ.\csromanspacer{} กาเลน เต อาวุโส โจทิโต โน อกาเลน, \csromanfolio{174}อลํ เต อวิปฺปฏิสาราย.\csromanspacer{} ภูเตน เต อาวุโส โจทิโต โน อภูเตน, อลํ เต อวิปฺปฏิสาราย.\csromanspacer{} สเณฺหน เต อาวุโส โจทิโต โน ผรุเสน, อลํ เต อวิปฺปฏิสาราย.\csromanspacer{} อตฺถสํหิเตน เต อาวุโส โจทิโต โน อนตฺถสํหิเตน, อลํ เต อวิปฺปฏิสาราย.\csromanspacer{} เมตฺตจิตฺเตน เต อาวุโส โจทิโต โน โทสนฺตเรน, อลํ เต อวิปฺปฏิสารายาติ.\csromanspacer{} ธมฺมโจทกสฺส อาวุโส ภิกฺขุโน อิเมหิ ปญฺจหากาเรหิ อวิปฺปฏิสาโร อุปทหาตพฺโพ.\csromanspacer{} ตํ กิสฺส เหตุ? ยถา อญฺโญปิ ภิกฺขุ ภูเตน โจทิตพฺพํ มญฺเญยฺยาติ.}

\prose{จุทิเตน อาวุโส ปุคฺคเลน ทฺวีสุ ธมฺเมสุ ปติฏฺฐาตพฺพํ สจฺเจ จ อกุปฺเป จ.\csromanspacer{} มํ เจวิ\footnote{ปญฺจหิ (สฺยา, ก)} อาวุโส ปเร โจเทยฺยุํ กาเลน วา อกาเลน วา ภูเตน วา อภูเตน วา สเณฺหน วา ผรุเสน วา อตฺถสํหิเตน วา อนตฺถสํหิเตน วา เมตฺตจิตฺตา\footnote{เมตฺตจิตฺเตน (สี, อิ, ก) ม 1. 178 ปิฏฺเฐ ปสฺสิตพฺพํ. 3. โทสนฺตเรน (สี, อิ, ก)} วา โทสนฺตรา3 วา.\csromanspacer{} อหมฺปิ ทฺวีสุเยว ธมฺเมสุ ปติฏฺฐเหยฺยํ สจฺเจ จ อกุปฺเป จ.\csromanspacer{} สเจ ชาเนยฺยํ “อตฺเถโส มยิ ธมฺโม”ติ. “อตฺถี”ติ นํ วเทยฺยํ “สํวิชฺชเตโส มยิ ธมฺโม”ติ.\csromanspacer{} สเจ ชาเนยฺยํ “นตฺเถโส มยิ ธมฺโม”ติ. “นตฺถี”ติ นํ วเทยฺยํ “เนโส ธมฺโม มยิ สํวิชฺชตี”ติ.}

\prose{เอวมฺปิ โข เต\footnote{เอวมฺปิ โข (ก)} สาริปุตฺต วุจฺจมานา อถ จ ปนิเธกจฺเจ โมฆปุริสา น ปทกฺขิณํ คณฺหนฺตีติ.}

\prose{เย เต ภนฺเต ปุคฺคลา อสฺสทฺธา ชีวิกตฺถา น สทฺธา อคารสฺมา อนคาริยํ ปพฺพชิตา สฐา มายาวิโน เกตพิโน\footnote{เกฏุภิโน (สี, สฺยา, กํ, อิ)} อุทฺธตา อุนฺนฬา จปลา มุขรา วิกิณฺณวาจา อินฺทฺริเยสุ อคุตฺตทฺวารา โภชเน อมตฺตญฺญุโน ชาคริยํ อนนุยุตฺตา สามญฺเญ อนเปกฺขวนฺโต สิกฺขาย น ติพฺพคารวา พาหุลิกา สาถลิกา โอกฺกมเน ปุพฺพงฺคมา ปวิเวเก นิกฺขิตฺตธุรา กุสีตา หีนวีริยา มุฏฺฐสฺสติโน อสมฺปชานา อสมาหิตา วิพฺภนฺตจิตฺตา ทุปฺปญฺญา เอฬมูคา, เต มยา เอวํ วุจฺจมานา น ปทกฺขิณํ คณฺหนฺติ.}

\prose{เย ปน เต ภนฺเต กุลปุตฺตา สทฺธา อคารสฺมา อนคาริยํ ปพฺพชิตา อสฐา อมายาวิโน อเกตพิโน อนุทฺธตา อนุนฺนฬา อจปลา อมุขรา อวิกิณฺณวาจา อินฺทฺริเยสุ คุตฺตทฺวารา โภชเน มตฺตญฺญุโน}

\vspace{\tipitakaparskip}
\csromancenter{(17) 2.\csromanspacer{} อาคาตวคฺค ชาคริยํ อนุยุตฺตา สามญฺเญ อเปกฺขวนฺโต สิกฺขาย ติพฺพคารวา น พาหุลิกา น สาถลิกา โอกฺกมเน นิกฺขิตฺตธุรา ปวิเวเก ปุพฺพงฺคมา อารทฺธวีริยา ปหิตตฺตา อุปฏฺฐิตสฺสติโน สมฺปชานา สมาหิตา เอกคฺคจิตฺตา ปญฺญวนฺโต อเนฬมูคา, เต มยา เอวํ วุจฺจมานา ปทกฺขิณํ คณฺหนฺตีติ.}
\prose{\csromanfolio{175}เย เต สาริปุตฺต ปุคฺคลา อสฺสทฺธา ชีวิกตฺถา น สทฺธา อคารสฺมา อนคาริยํ ปพฺพชิตา สฐา มายาวิโน เกตพิโน อุทฺธตา อุนฺนฬา จปลา มุขรา วิกิณฺณวาจา อินฺทฺริเยสุ อคุตฺตทฺวารา โภชเน อมตฺตญฺญุโน ชาคริยํ อนนุยุตฺตา สามญฺเญ อนเปกฺขวนฺโต สิกฺขาย น ติพฺพคารวา พาหุลิกา สาถลิกา โอกฺกมเน ปุพฺพงฺคมา ปวิเวเก นิกฺขิตฺตธุรา กุสีตา หีนวีริยา มุฏฺฐสฺสติโน อสมฺปชานา อสมาหิตา วิพฺภนฺตจิตฺตา ทุปฺปญฺญา เอฬมูคา, ติฏฺฐนฺตุ เต.}

\prose{เย ปน เต สาริปุตฺต กุลปุตฺตา สทฺธา อคารสฺมา อนคาริยํ ปพฺพชิตา อสฐา อมายาวิโน อเกตพิโน อนุทฺธตา อนุนฺนฬา อจปลา อมุขรา อวิกิณฺณวาจา อินฺทฺริเยสุ คุตฺตทฺวารา โภชเน มตฺตญฺญุโน ชาคริยํ อนุยุตฺตา สามญฺเญ อเปกฺขวนฺโต สิกฺขาย ติพฺพคารวา น พาหุลิกา น สาถลิกา โอกฺกมเน นิกฺขิตฺตธุรา ปวิเวเก ปุพฺพงฺคมา อารทฺธวีริยา ปหิตตฺตา อุปฏฺฐิตสฺสติโน สมฺปชานา สมาหิตา เอกคฺคจิตฺตา ปญฺญวนฺโต อเนฬมูคา, เต ตฺวํ สาริปุตฺต วเทยฺยาสิ.\csromanspacer{} โอวท สาริปุตฺต สพฺรหฺมจารี, อนุสาส สาริปุตฺต สพฺรหฺมจารี “อสทฺธมฺมา วุฏฺฐาเปตฺวา สทฺธมฺเม ปติฏฺฐาเปสฺสามิ สพฺราหฺมจารี”ติ.\csromanspacer{} เอวํ หิ เต สาริปุตฺต สิกฺขิตพฺพนฺติ.\csromanspacer{}\csromanspacer{}สตฺตมํ.}

\csromanheader{2}{8. สีลสุตฺต}
\proseitem{168}{ตตฺร โข อายสฺมา สาริปุตฺโต ภิกฺขู อามนฺเตสิ, ทุสฺสีลสฺส อาวุโส สีลวิปนฺนสฺส หตูปนิโส โหติ สมฺมาสมาธิ, สมฺมาสมาธิมฺหิ อสติ สมฺมาสมาธิวิปนฺนสฺส หตูปนิสํ โหติ ยถาภูตญาณทสฺสนํ, ยถาภูตญาณทสฺสเน อสติ ยถาภูตญาณทสฺสนวิปนฺนสฺส หตูปนิโส โหติ นิพฺพิทาวิราโค, นิพฺพิทาวิราเค อสติ นิพฺพิทาวิราคคิปนฺนสฺส หตูปนิสํ โหติ วิมุตฺติญาณทสฺสนํ.\csromanspacer{} เสยฺยถาปิ อาวุโส รุกฺโข สาขาปลาสวิปนฺโน ตสฺส ปปฏิกาปิ น ปาริปูริํ คจฺฉติ, \csromanfolio{176}ตโจปิ เผคฺคุปิ สาโรปิ น ปาริปูริํ คจฺฉติ.\csromanspacer{} เอวเมวํ โข อาวุโส ทุสฺสีลสฺส สีลวิปนฺนสฺส หตูปนิโส โหติ สมฺมาสมาธิ, สมฺมาสมาธิมฺหิ อสติ สมฺมาสมาธิวิปนฺนสฺส หตูปนิสํ โหติ ยถาภูตญาณทสฺสนํ, ยถาภูตญาณทสฺสเน อสติ ยถาภูตญาณทสฺสนวิปนฺนสฺส หตูปนิโส โหติ นิพฺพิทาวิราโค, นิพฺพิทาวิราเค อสติ นิพฺพิทาวิราควิปนฺนสฺส หตูปนิสํ โหติ วิมุตฺติญาณทสฺสนํ.}

\prose{สีลวโต อาวุโส สีลสมฺปนฺนสฺส อุปนิสสมฺปนฺโน โหติ สมฺมาสมาธิ, สมฺมาสมาธิมฺหิ สติ สมฺมาสมาธิสมฺปนฺนสฺส อุปนิสสมฺปนฺนํ โหติ ยถาภูตญาณทสฺสนํ, ยถาภูตญาณทสฺสเน สติ ยถาภูตญาณทสฺสน สมฺปนฺนสฺส อุปนิสสมฺปนฺโน โหติ นิพฺพิทาวิราโค, นิพฺพิทาวิราเค สติ นิพฺพิทาวิราค สมฺปนฺนสฺส อุปนิสสมฺปนฺนํ โหติ วิมุตฺติญาณทสฺสนํ.\csromanspacer{} เสยฺยถาปิ อาวุโส รุกฺโข สาขาปลาสสมฺปนฺโน ตสฺส ปปฏิกาปิ ปาริปูริํ คจฺฉติ, ตโจปิ เผคฺคุปิ สาโรปิ ปาริปูริํ คจฺฉติ.\csromanspacer{} เอวเมวํ โข อาวุโส สีลวโต สีลสมฺปนฺนสฺส อุปนิสสมฺปนฺโน โหติ สมฺมาสมาธิ, สมฺมาสมาธิมฺหิ สติ สมฺมาสมาธิสมฺปนฺนสฺส อุปนิส สมฺปนฺนํ โหติ ยถาภูตญาณทสฺสนํ, ยถาภูตญาณทสฺสเน สติ ยถาภูตญาณ ทสฺสนสมฺปนฺนสฺส อุปนิสสมฺปนฺโน โหติ นิพฺพิทาวิราโค, นิพฺพิทาวิราเค สติ นิพฺพิทา วิราคสมฺปนฺนสฺส อุปนิสสมฺปนฺนํ โหติ วิมุตฺติญาณทสฺสนนฺติ\footnote{เหฏฺฐา 17 ปิฏฺเฐ; อุปริ 316, 472 ปิฏฺเฐปิ.}.\csromanspacer{}\csromanspacer{}อฏฺฐมํ.}

\csromanheader{2}{9. ขิปฺปนิสนฺติสุตฺต}
\proseitem{169}{อถ โข อายสฺมา อานนฺโท เยนายสฺมา สาริปุตฺโต เตนุปสงฺกมิ, อุปสงฺกมิตฺวา อายสฺมตา สาริปุตฺเตน สทฺธิํ สมฺโมทิ, สมฺโมทนียํ กถํ สารณียํ\footnote{สาราณียํ (สี, สฺยา, กํ, อิ)} วีติสาเรตฺวา เอกมนฺตํ นิสีทิ, เอกมนฺตํ นิสินฺโน โข อายสฺมา อานนฺโท อายสฺมนฺตํ สาริปุตฺตํ เอตทโวจ\csromandash{}}

\prose{กิตฺตาวตา นุ โข อาวุโส สาริปุตฺต ภิกฺขุ ขิปฺปนิสนฺติ จ โหติ กุสเลสุ ธมฺเมสุ สุคฺคหิตคฺคาหี จ, พหุญฺจ คณฺหาติ, คหิตํ จสฺส นปฺปมุสฺสตีติ.\csromanspacer{} อายสฺมา โข อานนฺโท พหุสฺสุโต, ปฏิภาตุ อายสฺมนฺตํเยว อานนฺทนฺติ.\csromanspacer{} เตนหาวุโส สาริปุตฺต สุณาหิ สาธุกํ}

\vspace{\tipitakaparskip}
\csromancenter{(17) 2.\csromanspacer{} อาคาตวคฺค มนสิ กโรหิ, ภาสิสฺสามีติ. “เอวมาวุโส”ติ โข อายสฺมา สาริปุตฺโต อายสฺมโต อานนฺทสฺส ปจฺจสฺโสสิ, อายสฺมา อานนฺโท เอตทโวจ\csromandash{}}
\prose{\csromanfolio{177}อิธาวุโส สาริปุตฺต ภิกฺขุ อตฺถกุสโล จ โหติ ธมฺมกุสโล จ พฺยญฺชน กุสโล จ นิรุตฺติกุสโล จ ปุพฺพาปรกุสโล จ.\csromanspacer{} เอตฺตาวตา โข อาวุโส สาริปุตฺต ภิกฺขุ ขิปฺปนิสนฺติ จ โหติ กุสเลสุ ธมฺเมสุ สุคฺคหิตคฺคาหี จ, พหุญฺจ คณฺหาติ, คหิตํ จสฺส นปฺปมุสฺสตีติ.\csromanspacer{} อจฺฉริยํ อาวุโส, อพฺภุตํ อาวุโส, ยาว สุภาสิตํ จิทํ อายสฺมตา อานนฺเทน.\csromanspacer{} อิเมหิ จ มยํ ปญฺจหิ ธมฺเมหิ สมนฺนาคตํ อายสฺมนฺตํ อานนฺทํ ธาเรม, อายสฺมา อานนฺโท อตฺถกุสโล ธมฺมกุสโล พฺยญฺชนกุสโล นิรุตฺติกุสโล ปุพฺพาปรกุสโลติ.\csromanspacer{}\csromanspacer{}นวมํ.}

\csromanmatikaanchor{mtk.60}
\csromantocmarkref{subsection}{10. ภทฺทชิสุตฺต}{mtk.60}
\bookmark[dest={mtk.60},level=2]{10. ภทฺทชิสุตฺต}
\csromanheader{2}{10. ภทฺทชิสุตฺต}
\proseitem{170}{เอกํ สมยํ อายสฺมา อานนฺโท โกสมฺพิยํ วิหรติ โฆสิตาราเม.\csromanspacer{} อถ โข อายสฺมา ภทฺทชิ เยนายสฺมา อานนฺโท เตนุปสงฺกมิ, อุปสงฺกมิตฺวา อายสฺมตา อานนฺเทน สทฺธิํ สมฺโมทิ, สมฺโมทนียํ กถํ สารณียํ วีติสาเรตฺวา เอกมนฺตํ นิสีทิ, เอกมนฺตํ นิสินฺนํ โข อายสฺมนฺตํ ภทฺทชิํ อายสฺมา อานนฺโท เอตทโวจ “กิํ นุ โข อาวุโส ภทฺทชิ ทสฺสนานํ อคฺคํ, กิํ สวนานํ อคฺคํ, กิํ สุขานํ อคฺคํ, กิํ สญฺญานํ อคฺคํ, กิํ ภวานํ อคฺคนฺ”ติ.}

\prose{อตฺถาวุโส พฺรหฺมา อภิภู อนภิภูโต อญฺญทตฺถุทโส วสวตฺตี, โย ตํ พฺรหฺมานํ ปสฺสติ, อิทํ ทสฺสนานํ อคฺคํ.\csromanspacer{} อตฺถาวุโส อาภสฺสรา นาม เทวา สุเขน อภิสนฺนา ปริสนฺนา, เต กทาจิ กรหจิ อุทานํ อุทาเนนฺติ “อโห สุขํ อโห สุขนฺ”ติ.\csromanspacer{} โย ตํ สทฺทํ สุณาติ, อิทํ สวนานํ อคฺคํ.\csromanspacer{} อตฺถาวุโส สุภกิณฺหา นาม เทวา, เต สนฺตํเยว ตุสิตา สุขํ ปฏิเวเทนฺติ, อิทํ สุขานํ อคฺคํ.\csromanspacer{} อตฺถาวุโส อากิญฺจญฺญายตนูปคา เทวา, อิทํ สญฺญานํ อคฺคํ.\csromanspacer{} อตฺถาวุโส เนวสญฺญานาสญฺญายตนูปคา เทวา, อิทํ ภวานํ อคฺคนฺติ.\csromanspacer{} สเมติ โข อิทํ อายสฺมโต ภทฺทชิสฺส, ยทิทํ พหุนา ชเนนาติ.}

\csromanmatikaanchor{mtk.61}
\csromanmatikaanchor{mtk.62}
\csromantocmarknopage{section}{(18) 3. อุปาสกวคฺค}{mtk.61}
\bookmark[dest={mtk.61},level=1]{(18) 3. อุปาสกวคฺค}
\csromantocmarkref{subsection}{1-3. สารชฺชสุตฺตาทิ}{mtk.62}
\bookmark[dest={mtk.62},level=2]{1-3. สารชฺชสุตฺตาทิ}
\prose{\csromanfolio{178}อายสฺมา โข อานนฺโท พหุสฺสุโต, ปฏิภาตุ อายสฺมนฺตํเยว อานนฺทนฺติ.\csromanspacer{} เตนหาวุโส ภทฺทชิ สุณาหิ สาธุกํ มนสิ กโรหิ, ภาสิสฺสามีติ. “เอวมาวุโส”ติ โข อายสฺมา ภทฺทชิ อายสฺมโต อานนฺทสฺส ปจฺจสฺโสสิ.\csromanspacer{} อายสฺมา อานนฺโท เอตทโวจ\csromandash{}}

\prose{ยถา ปสฺสโต โข อาวุโส อนนฺตรา อาสวานํ ขโย โหติ, อิทํ ทสฺสนานํ อคฺคํ.\csromanspacer{} ยถา สุณโต อนนฺตรา อาสวานํ ขโย โหติ, อิทํ สวนานํ อคฺคํ.\csromanspacer{} ยถา สุขิตสฺส อนนฺตรา อาสวานํ ขโย โหติ, อิทํ สุขานํ อคฺคํ.\csromanspacer{} ยถา สญฺญิสฺส อนนฺตรา อาสวานํ ขโย โหติ, อิทํ สญฺญานํ อคฺคํ.\csromanspacer{} ยถา ภูตสฺส อนนฺตรา อาสวานํ ขโย โหติ, อิทํ ภวานํ อคฺคนฺติ.\csromanspacer{}\csromanspacer{}ทสมํ.\csromanspacer{} อาฆาตวคฺโค ทุติโย.}

\titlehead{\textbf{ตสฺสุทฺทานํ}}
\csromangathameasure{เทฺว อาฆาตวินยา, สากจฺฉา สาชีวโต ปญฺหํ. \\ ปุจฺฉา นิโรโธ โจทนา, สีลํ นิสนฺติ ภทฺทชีติ.}
\csromangathagroup{\csromangathastanza{เทฺว อาฆาตวินยา, สากจฺฉา สาชีวโต ปญฺหํ. \\ ปุจฺฉา นิโรโธ โจทนา, สีลํ นิสนฺติ ภทฺทชีติ.}}

\csromanheader{1}{\textbf{(18) 3. อุปาสกวคฺค}}
\csromanheader{2}{1. สารชฺชสุตฺต}
\proseitem{171}{เอวํ เม สุตํ\csromandash{}เอกํ สมยํ ภควา สาวตฺถิยํ วิหรติ เชตวเน อนาถปิณฺฑิกสฺส อาราเม.\csromanspacer{} ตตฺร โข ภควา ภิกฺขู อามนฺเตสิ “ภิกฺขโว”ติ. “ภทนฺเต”ติ เต ภิกฺขู ภควโต ปจฺจสฺโสสุํ.\csromanspacer{} ภควา เอตทโวจ\csromandash{}}

\prose{ปญฺจหิ ภิกฺขเว ธมฺเมหิ สมนฺนาคโต อุปาสโก สารชฺชํ โอกฺกนฺโต โหติ.\csromanspacer{} กตเมหิ ปญฺจหิ? ปาณาติปาตี โหติ, อทินฺนาทายี โหติ, กาเมสุมิจฺฉาจารี โหติ, มุสาวาที โหติ, สุราเมรมชฺชปมาทฏฺฐายี โหติ.\csromanspacer{} อิเมหิ โข ภิกฺขเว ปญฺจหิ ธมฺเมหิ สมนฺนาคโต อุปาสโก สารชฺชํ โอกฺกนฺโต โหติ.}

\chapterheadpageread{(18) 3. อุปาสกวคฺค}
\prose{\csromanfolio{179}ปญฺจหิ ภิกฺขเว ธมฺเมหิ สมนฺนาคโต อุปาสโก วิสารโท โหติ.\csromanspacer{} กตเมหิ ปญฺจหิ? ปาณาติปาตา ปฏิวิรโต โหติ, อทินฺนาทานา ปฏิวิรโต โหติ, กาเมสุมิจฺฉาจารา ปฏิวิรโต โหติ, มุสาวาทา ปฏิวิรโต โหติ, สุราเมรยมชฺชปมาทฏฺฐานา ปฏิวิรโต โหติ.\csromanspacer{} อิเมหิ โข ภิกฺขเว ปญฺจหิ ธมฺเมหิ สมนฺนาคโต อุปาสโก วิสารโท โหตีติ.\csromanspacer{}\csromanspacer{}ปฐมํ.}

\csromanheader{2}{2. วิสารทสุตฺต}
\proseitem{172}{ปญฺจหิ ภิกฺขเว ธมฺเมหิ สมนฺนาคโต อุปาสโก อวิสารโท อคารํ อชฺฌาวสติ.\csromanspacer{} กตเมหิ ปญฺจหิ? ปาณาติปาตี โหติ ฯเปฯ สุราเมรยมชฺชปมาทฏฺฐายี โหติ.\csromanspacer{} อิเมหิ โข ภิกฺขเว ปญฺจหิ ธมฺเมหิ สมนฺนาคโต อุปาสโก อวิสารโท อคารํ อชฺฌาวสติ.}

\prose{ปญฺจหิ ภิกฺขเว ธมฺเมหิ สมนฺนาคโต อุปาสโก วิสารโท อคารํ อชฺฌาวสติ.\csromanspacer{} กตเมหิ ปญฺจหิ? ปาณาติปาตา ปฏิวิรโต โหติ ฯเปฯ สุราเมรยมชฺชปมาทฏฺฐานา ปฏิวิรโต โหติ.\csromanspacer{} อิเมหิ โข ภิกฺขเว ปญฺจหิ ธมฺเมหิ สมนฺนาคโต อุปาสโก วิสารโท อคารํ อชฺฌาวสตีติ.\csromanspacer{}\csromanspacer{}ทุติยํ.}

\csromanheader{2}{3. นิรยสุตฺต}
\proseitem{173}{ปญฺจหิ ภิกฺขเว ธมฺเมหิ สมนฺนาคโต อุปาสโก ยถาภตํ นิกฺขิตฺโต เอวํ นิรเย.\csromanspacer{} กตเมหิ ปญฺจหิ? ปาณาติปาตี โหติ ฯเปฯ สุราเมรยมชฺชปมาทฏฺฐายี โหติ.\csromanspacer{} อิเมหิ โข ภิกฺขเว ปญฺจหิ ธมฺเมหิ สมนฺนาคโต อุปาสโก ยถาภตํ นิกฺขิตฺโต เอวํ นิรเย.}

\prose{ปญฺจหิ ภิกฺขเว ธมฺเมหิ สมนฺนาคโต อุปาสโก ยถาภตํ นิกฺขิตฺโต เอวํ สคฺเค.\csromanspacer{} กตเมหิ ปญฺจหิ? ปาณาติปาตา ปฏิวิรโต โหติ ฯเปฯ สุราเมรยมชฺชปมาทฏฺฐานา ปฏิวิรโต โหติ.\csromanspacer{} อิเมหิ โข ภิกฺขเว ปญฺจหิ ธมฺเมหิ สมนฺนาคโต อุปาสโก ยถาภตํ นิกฺขิตฺโต เอวํ สคฺเคติ.\csromanspacer{}\csromanspacer{}ตติยํ.}

\csromanmatikaanchor{mtk.63}
\csromantocmarkref{subsection}{4. เวรสุตฺต}{mtk.63}
\bookmark[dest={mtk.63},level=2]{4. เวรสุตฺต}
\csromanheader{2}{4. เวรสุตฺต}
\proseitem{174}{\csromanfolio{180}อถ โข อนาถปิณฺฑิโก คหปติ เยน ภควา เตนุปสงฺกมิ, อุปสงฺก มิตฺวา ภควนฺตํ อภิวาเทตฺวา เอกมนฺตํ นิสีทิ, เอกมนฺตํ นิสินฺนํ โข อนาถปิณฺฑิกํ คหปติํ ภควา เอตทโวจ\csromandash{}}

\prose{ปญฺจ คหปติ ภยานิ เวรานิ อปฺปหาย “ทุสฺสีโล” อิติ วุจฺจติ, นิรยญฺจ อุปปชฺชติ.\csromanspacer{} กตมานิ ปญฺจ? ปาณาติปาตํ อทินฺนาทานํ กาเมสุมิจฺฉาจารํ มุสาวาทํ สุราเมรยมชฺชปมาทฏฺฐานํ.\csromanspacer{} อิมานิ โข คหปติ ปญฺจ ภยานิ เวรานิ อปฺปหาย “ทุสฺสีโล” อิติ วุจฺจติ, นิรยญฺจ อุปปชฺชติ.}

\prose{ปญฺจ คหปติ ภยานิ เวรานิ ปหาย “สีลวา” อิติ วุจฺจติ, สุคติญฺจ อุปปชฺชติ.\csromanspacer{} กตมานิ ปญฺจ? ปาณาติปาตํ อทินฺนาทานํ กาเมสุมิจฺฉาจารํ มุสาวาทํ สุราเมรยมชฺชปมาทฏฺฐานํ.\csromanspacer{} อิมานิ โข คหปติ ปญฺจ ภยานิ เวรานิ ปหาย “สีลวา” อิติ วุจฺจติ, สุคติญฺจ อุปปชฺชติ.}

\prose{ยํ คหปติ ปาณาติปาตี ปาณาติปาตปจฺจยา ทิฏฺฐธมฺมิกมฺปิ ภยํ เวรํ ปสวติ, สมฺปรายิกมฺปิ ภยํ เวรํ ปสวติ, เจตสิกมฺปิ ทุกฺขํ โทมนสฺสํ ปฏิสํเวเทติ.\csromanspacer{} ปาณาติปาตา ปฏิวิรโต เนว ทิฏฺฐธมฺมิกํ ภยํ เวรํ ปสวติ, น สมฺปรายิกํ ภยํ เวรํ ปสวติ, น เจตสิกํ ทุกฺขํ โทมนสฺสํ ปฏิสํเวเทติ.\csromanspacer{} ปาณาติปาตา ปฏิวิรตสฺส เอวํ ตํ ภยํ เวรํ วูปสนฺตํ โหติ.}

\prose{ยํ คหปติ อทินฺนาทายี ฯเปฯ.}

\prose{ยํ คหปติ กาเมสุมิจฺฉาจารี ฯเปฯ.}

\prose{ยํ คหปติ มุสาวาที ฯเปฯ.}

\prose{ยํ คหปติ สุราเมรยมชฺชปมาทฏฺฐายี สุราเมรยมชฺชปมาทฏฺฐาน ปจฺจยา ทิฏฺฐธมฺมิกมฺปิ ภยํ เวรํ ปสวติ, สมฺปรายิกมฺปิ ภยํ เวรํ ปสวติ, เจตสิกมฺปิ ทุกฺขํ โทมนสฺสํ ปฏิสํเวเทติ.\csromanspacer{} สุราเมรยมชฺชปมาทฏฺฐานา ปฏิวิรโต เนว ทิฏฺฐธมฺมิกํ ภยํ เวรํ ปสวติ, น สมฺปรายิกํ ภยํ เวรํ ปสวติ, น เจตสิกํ ทุกฺขํ โทมนสฺสํ ปฏิสํเวเทติ.\csromanspacer{} สุราเมรยมชฺชปมาทฏฺฐานา ปฏิวิรตสฺส เอวํ ตํ ภยํ เวรํ วูปสนฺตํ โหตีติ.}

\chapterheadpageread{(18) 3. อุปาสกวคฺค}
\csromangathameasure{โย ปาณมติปาเตติ, มุสาวาทญฺจ ภาสติ. \\ โลเก อทินฺนํ อาทิยติ, ปรทารญฺจ คจฺฉติ. \\ สุราเมรยปานญฺจ, โย นโร อนุยุญฺชติ. \\ อปฺปหาย ปญฺจ เวรานิ, “ทุสฺสีโล” อิติ วุจฺจติ. \\ กายสฺส เภทา ทุปฺปญฺโญ, นิรยํ โสปปชฺชติ. \\ โย ปาณํ นาติปาเตติ, มุสาวาทํ น ภาสติ. \\ โลเก อทินฺนํ นาทิยติ, ปรทารํ น คจฺฉติ. \\ สุราเมรยปานญฺจ, โย นโร นานุยุญฺชติ. \\ ปหาย ปญฺจ เวรานิ, “สีลวา” อิติ วุจฺจติ. \\ กายสฺส เภทา สปฺปญฺโญ, สุคติํ โสปปชฺชตีติ. จตุตฺถํ.}
\csromangathagroup{\csromangathastanza{\csromanfolio{181}โย ปาณมติปาเตติ, มุสาวาทญฺจ ภาสติ. \\ โลเก อทินฺนํ อาทิยติ, ปรทารญฺจ คจฺฉติ.}\csromangathastanza{สุราเมรยปานญฺจ, โย นโร อนุยุญฺชติ. \\ อปฺปหาย ปญฺจ เวรานิ, “ทุสฺสีโล” อิติ วุจฺจติ.}\csromangathastanza{กายสฺส เภทา ทุปฺปญฺโญ, นิรยํ โสปปชฺชติ. \\ โย ปาณํ นาติปาเตติ, มุสาวาทํ น ภาสติ.}\csromangathastanza{โลเก อทินฺนํ นาทิยติ, ปรทารํ น คจฺฉติ. \\ สุราเมรยปานญฺจ, โย นโร นานุยุญฺชติ.}\csromangathastanza{ปหาย ปญฺจ เวรานิ, “สีลวา” อิติ วุจฺจติ. \\ กายสฺส เภทา สปฺปญฺโญ, สุคติํ โสปปชฺชตีติ.\csromanspacer{} จตุตฺถํ.}}

\csromanmatikaanchor{mtk.64}
\csromantocmarkref{subsection}{5. จณฺฑาลสุตฺต}{mtk.64}
\bookmark[dest={mtk.64},level=2]{5. จณฺฑาลสุตฺต}
\csromanheader{2}{5. จณฺฑาลสุตฺต}
\proseitem{175}{ปญฺจหิ ภิกฺขเว ธมฺเมหิ สมนฺนาคโต อุปาสโก อุปาสกจณฺฑาโล จ โหติ อุปาสกมลญฺจ อุปาสกปติกุฏฺโฐ จ\footnote{อุปาสกปติกิฏฺโฐ จ (สี, สฺยา, กํ, อิ)}.\csromanspacer{} กตเมหิ ปญฺจหิ? อสฺสทฺโธ โหติ, ทุสฺสีโล โหติ, โกตูหลมงฺคลิโก โหติ, มงฺคลํ ปจฺเจติ โน กมฺมํ, อิโต จ พหิทฺธา ทกฺขิเณยฺยํ คเวสติ, ตตฺถ จ ปุพฺพการํ กโรติ.\csromanspacer{} อิเมหิ โข ภิกฺขเว ปญฺจหิ ธมฺเมหิ สมนฺนาคโต อุปาสโก อุปาสกจณฺฑาโล จ โหติ อุปาสกมลญฺจ อุปาสก ปติกุฏฺโฐ จ.}

\prose{ปญฺจหิ ภิกฺขเว ธมฺเมหิ สมนฺนาคโต อุปาสโก อุปาสกรตนญฺจ โหติ อุปาสกปทุมญฺจ อุปาสกปุณฺฑรีกญฺจ\footnote{อุปาสกปุณฺฑรีโก จ (อิ, ก)}.\csromanspacer{} กตเมหิ ปญฺจหิ? สทฺโธ โหติ, สีลวา โหติ, อโกตูหลมงฺคลิโก โหติ, กมฺมํ ปจฺเจติ โน มงฺคลํ, น อิโต พหิทฺธา ทกฺขิเณยฺยํ คเวสติ, อิธ จ ปุพฺพการํ กโรติ.\csromanspacer{} อิเมหิ โข ภิกฺขเว ปญฺจหิ ธมฺเมหิ สมนฺนาคโต อุปาสโก อุปาสกรตนญฺจ โหติ อุปาสกปทุมญฺจ อุปาสกปุณฺฑรีกญฺจาติ.\csromanspacer{}\csromanspacer{}ปญฺจมํ.}

\csromanmatikaanchor{mtk.65}
\csromantocmarkref{subsection}{6. ปีติสุตฺต}{mtk.65}
\bookmark[dest={mtk.65},level=2]{6. ปีติสุตฺต}
\csromanheader{2}{6. ปีติสุตฺต}
\proseitem{176}{\csromanfolio{182}อถ โข อนาถปิณฺฑิโก คหปติ ปญฺจมตฺเตหิ อุปาสกสเตหิ ปริวุโต เยน ภควา เตนุปสงฺกมิ, อุปสงฺกมิตฺวา ภควนฺตํ อภิวาเทตฺวา เอกมนฺตํ นิสีทิ.\csromanspacer{} เอกมนฺตํ นิสินฺนํ โข อนาถปิณฺฑิกํ คหปติํ ภควา เอตทโวจ\csromandash{}}

\prose{ตุเมฺห โข คหปติ ภิกฺขุสํฆํ ปจฺจุปฏฺฐิตา จีวรปิณฺฑปาตเสนาสน คิลานปฺปจฺจยเภสชฺชปริกฺขาเรน.\csromanspacer{} น โข คหปติ ตาวตเกเนว ตุฏฺฐิ กรณียา “มยํ ภิกฺขุสํฆํ ปจฺจุปฏฺฐิตา จีวรปิณฺฑปาตเสนาสนคิลานปฺปจฺจย เภสชฺชปริกฺขาเรนา”ติ.\csromanspacer{} ตสฺมาติห คหปติ เอวํ สิกฺขิตพฺพํ “กินฺติ มยํ กาเลน กาลํ ปวิเวกํ ปีติํ อุปสมฺปชฺช วิหเรยฺยามา”ติ.\csromanspacer{} เอวํ หิ โว คหปติ สิกฺขิ ตพฺพนฺติ.}

\prose{เอวํ วุตฺเต อายสฺมา สาริปุตฺโต ภควนฺตํ เอตทโวจ\csromandash{}อจฺฉริยํ ภนฺเต, อพฺภุตํ ภนฺเต, ยาว สุภาสิตํ จิทํ ภนฺเต ภควตา “ตุเมฺห โข คหปติ ภิกฺขุสํฆํ ปจฺจุปฏฺฐิตา จีวรปิณฺฑปาตเสนาสนคิลานปฺปจฺจย เภสชฺชปริกฺขาเรน.\csromanspacer{} น โข คหปติ ตาวตเกเนว ตุฏฺฐิ กรณียา ‘มยํ ภิกฺขุสํฆํ ปจฺจุปฏฺฐิตา จีวรปิณฺฑปาตเสนาสนคิลานปฺปจฺจยเภสชฺช ปริกฺขาเรนา’ติ.\csromanspacer{} ตสฺมาติห คหปติ เอวํ สิกฺขิตพฺพํ ‘กินฺติ มยํ กาเลน กาลํ ปวิเวกํ ปีติํ อุปสมฺปชฺช วิหเรยฺยามา’ติ.\csromanspacer{} เอวํ หิ โว คหปติ สิกฺขิตพฺพนฺ”ติ.\csromanspacer{} ยสฺมิํ ภนฺเต สมเย อริยสาวโก ปวิเวกํ ปีติํ อุปสมฺปชฺช วิหรติ, ปญฺจสฺส ฐานานิ ตสฺมิํ สมเย น โหนฺติ.\csromanspacer{} ยมฺปิสฺส กามูปสํหิตํ ทุกฺขํ โทมนสฺสํ, ตมฺปิสฺส ตสฺมิํ สมเย น โหติ.\csromanspacer{} ยมฺปิสฺส กามูปสํหิตํ สุขํ โสมนสฺสํ, ตมฺปิสฺส ตสฺมิํ สมเย น โหติ.\csromanspacer{} ยมฺปิสฺส อกุสลูปสํหิตํ ทุกฺขํ โทมนสฺสํ, ตมฺปิสฺส ตสฺมิํ สมเย น โหติ.\csromanspacer{} ยมฺปิสฺส อกุสลูปสํหิตํ สุขํ โสมนสฺสํ, ตมฺปิสฺส ตสฺมิํ สมเย น โหติ.\csromanspacer{} ยมฺปิสฺส กุสลูปสํหิตํ ทุกฺขํ โทมนสฺสํ, ตมฺปิสฺส ตสฺมิํ สมเย น โหติ.\csromanspacer{} ยสฺมิํ ภนฺเต สมเย อริยสาวโก ปวิเวกํ ปีติํ อุปสมฺปชฺช วิหรติ, อิมานิสฺส ปญฺจ\footnote{อิมานิ ปญฺจสฺส (สฺยา, กํ)} ฐานานิ ตสฺมิํ สมเย น โหนฺตีติ.}

\chapterheadpageread{(18) 3. อุปาสกวคฺค}
\csromanmatikaanchor{mtk.66}
\csromanmatikaanchor{mtk.67}
\csromantocmarkref{subsection}{7. วณิชฺชาสุตฺต}{mtk.66}
\bookmark[dest={mtk.66},level=2]{7. วณิชฺชาสุตฺต}
\csromantocmarkref{subsection}{8. ราชาสุตฺต}{mtk.67}
\bookmark[dest={mtk.67},level=2]{8. ราชาสุตฺต}
\prose{\csromanfolio{183}สาธุ สาธุ สาริปุตฺต, ยสฺมิํ สาริปุตฺต สมเย อริยสาวโก ปวิเวกํ ปีติํ อุปสมฺปชฺช วิหรติ, ปญฺจสฺส ฐานานิ ตสฺมิํ สมเย น โหนฺติ.\csromanspacer{} ยมฺปิสฺส กามูปสํหิตํ ทุกฺขํ โทมนสฺสํ, ตมฺปิสฺส ตสฺมิํ สมเย น โหติ.\csromanspacer{} ยมฺปิสฺส กามูปสํหิตํ สุขํ โสมนสฺสํ, ตมฺปิสฺส ตสฺมิํ สมเย น โหติ.\csromanspacer{} ยมฺปิสฺส อกุสลูปสํหิตํ ทุกฺขํ โทมนสฺสํ, ตมฺปิสฺส ตสฺมิํ สมเย น โหติ.\csromanspacer{} ยมฺปิสฺส อกุสลูปสํหิตํ สุขํ โสมนสฺสํ, ตมฺปิสฺส ตสฺมิํ สมเย น โหติ.\csromanspacer{} ยมฺปิสฺส กุสลูปสํหิตํ ทุกฺขํ โทมนสฺสํ, ตมฺปิสฺส ตสฺมิํ สมเย น โหติ.\csromanspacer{} ยสฺมิํ สาริปุตฺต สมเย อริยสาวโก ปวิเวกํ ปีติํ อุปสมฺปชฺช วิหรติ, อิมานิสฺส\footnote{อิมาเนตฺถ (สี)} ปญฺจ ฐานานิ ตสฺมิํ สมเย น โหนฺตีติ.\csromanspacer{}\csromanspacer{}ฉฏฺฐํ.}

\csromanheader{2}{7. วณิชฺชาสุตฺต}
\proseitem{177}{ปญฺจิมา ภิกฺขเว วณิชฺชา อุปาสเกน อกรณียา.\csromanspacer{} กตมา ปญฺจ? สตฺถวณิชฺชา สตฺตวณิชฺชา มํสวณิชฺชา มชฺชวณิชฺชา วิสวณิชฺชา.\csromanspacer{} อิมา โข ภิกฺขเว ปญฺจ วณิชฺชา อุปาสเกน อกรณียาติ.\csromanspacer{}\csromanspacer{}สตฺตมํ.}

\csromanheader{2}{8. ราชสุตฺต}
\proseitem{178}{ตํ กิํ มญฺญถ ภิกฺขเว, อปิ นุ ตุเมฺหหิ ทิฏฺฐํ วา สุตํ วา “อยํ ปุริโส ปาณาติปาตํ ปหาย ปาณาติปาตา ปฏิวิรโต”ติ\footnote{ปฏิวิรโต โหตีติ (สี), ปฏิวิรโต โหติ (สฺยา, กํ, อิ) 3. ตเถว ปาปกํ กมฺมํ ปเวทยนฺติ (สี), ตเทว ปาปกมฺมํ ปเวเทติ (สฺยา, กํ) 4. โวโรเปตีติ (สฺยา, กํ)} ตเมนํ ราชาโน คเหตฺวา ปาณาติปาตา เวรมณิเหตุ หนนฺติ วา พนฺธนฺติ วา ปพฺพาเชนฺติ วา ยถาปจฺจยํ วา กโรนฺตีติ.\csromanspacer{} โน เหตํ ภนฺเต.\csromanspacer{} สาธุ ภิกฺขเว, มยาปิ โข เอตํ ภิกฺขเว เนว ทิฏฺฐํ น สุตํ “อยํ ปุริโส ปาณาติปาตํ ปหาย ปาณาติปาตา ปฏิวิรโต”ติ ตเมนํ ราชาโน คเหตฺวา ปาณาติปาตา เวรมณิเหตุ หนนฺติ วา พนฺธนฺติ วา ปพฺพาเชนฺติ วา ยถาปจฺจยํ วา กโรนฺตีติ.\csromanspacer{} อปิ จ ขฺวสฺส ตเถว ปาปกมฺมํ ปเวเทนฺติ3 “อยํ ปุริโส อิตฺถิํ วา ปุริสํ วา ชีวิตา โวโรเปสี”ติ4.\csromanspacer{} ตเมนํ ราชาโน คเหตฺวา ปาณาติปาตเหตุ หนนฺติ วา \csromanfolio{184}พนฺธนฺติ วา ปพฺพาเชนฺติ วา ยถาปจฺจยํ วา กโรนฺติ.\csromanspacer{} อปิ นุ ตุเมฺหหิ เอวรูปํ ทิฏฺฐํ วา สุตํ วาติ.\csromanspacer{} ทิฏฺฐํ จ โน ภนฺเต สุตํ จ สุยฺยิสฺสติ\footnote{สูยิสฺสติ (สี, อิ)} จาติ.}

\prose{ตํ กิํ มญฺญถ ภิกฺขเว, อปิ นุ ตุเมฺหหิ ทิฏฺฐํ วา สุตํ วา “อยํ ปุริโส อทินฺนาทานํ ปหาย อทินฺนาทานา ปฏิวิรโต”ติ ตเมนํ ราชาโน คเหตฺวา อทินฺนาทานา เวรมณิเหตุ หนนฺติ วา พนฺธนฺติ วา ปพฺพาเชนฺติ วา ยถาปจฺจยํ วา กโรนฺตีติ.\csromanspacer{} โน เหตํ ภนฺเต.\csromanspacer{} สาธุ ภิกฺขเว, มยาปิ โข เอตํ ภิกฺขเว เนว ทิฏฺฐํ น สุตํ “อยํ ปุริโส อทินฺนาทานํ ปหาย อทินฺนาทานา ปฏิวิรโต”ติ ตเมนํ ราชาโน คเหตฺวา อทินฺนาทานา เวรมณิเหตุ หนนฺติ วา พนฺธนฺติ วา ปพฺพาเชนฺติ วา ยถาปจฺจยํ วา กโรนฺตีติ.\csromanspacer{} อปิ จ ขฺวสฺส ตเถว ปาปกมฺมํ ปเวเทนฺติ “อยํ ปุริโส คามา วา อรญฺญา วา อทินฺนํ เถยฺยสงฺขาตํ อาทิยี”ติ\footnote{อาทิยติ (สฺยา, กํ)}.\csromanspacer{} ตเมนํ ราชาโน คเหตฺวา อทินฺนาทานเหตุ หนนฺติ วา พนฺธนฺติ วา ปพฺพาเชนฺติ วา ยถาปจฺจยํ วา กโรนฺติ.\csromanspacer{} อปิ นุ ตุเมฺหหิ เอวรูปํ ทิฏฺฐํ วา สุตํ วาติ.\csromanspacer{} ทิฏฺฐํ จ โน ภนฺเต สุตํ จ สุยฺยิสฺสติ จาติ.}

\prose{ตํ กิํ มญฺญถ ภิกฺขเว, อปิ นุ ตุเมฺหหิ ทิฏฺฐํ วา สุตํ วา “อยํ ปุริโส กาเมสุมิจฺฉาจารํ ปหาย กาเมสุมิจฺฉาจารา ปฏิวิรโต”ติ ตเมนํ ราชาโน คเหตฺวา กาเมสุมิจฺฉาจารา เวรมณิเหตุ หนนฺติ วา พนฺธนฺติ วา ปพฺพาเชนฺติ วา ยถาปจฺจยํ วา กโรนฺตีติ.\csromanspacer{} โน เหตํ ภนฺเต.\csromanspacer{} สาธุ ภิกฺขเว, มยาปิ โข เอตํ ภิกฺขเว เนว ทิฏฺฐํ น สุตํ “อยํ ปุริโส กาเมสุมิจฺฉาจารํ ปหาย กาเมสุมิจฺฉาจารา ปฏิวิรโต”ติ ตเมนํ ราชาโน คเหตฺวา กาเมสุมิจฺฉาจารา เวรมณิเหตุ หนนฺติ วา พนฺธนฺติ วา ปพฺพาเชนฺติ วา ยถาปจฺจยํ วา กโรนฺตีติ.\csromanspacer{} อปิ จ ขฺวสฺส ตเถว ปาปกมฺมํ ปเวเทนฺติ “อยํ ปุริโส ปริตฺถีสุ ปรกุมารีสุ จาริตฺตํ อาปชฺชี”ติ\footnote{อาปชฺชติ (สฺยา, กํ)}.\csromanspacer{} ตเมนํ ราชาโน คเหตฺวา กาเมสุมิจฺฉาจารเหตุ หนนฺติ วา พนฺธนฺติ วา ปพฺพาเชนฺติ วา ยถาปจฺจยํ วา กโรนฺติ.\csromanspacer{} อปิ นุ ตุเมฺหหิ เอวรูปํ ทิฏฺฐํ วา สุตํ วาติ.\csromanspacer{} ทิฏฺฐํ จ โน ภนฺเต สุตํ จ สุยฺยิสฺสติ จาติ.}

\prose{ตํ กิํ มญฺญถ ภิกฺขเว, อปิ นุ ตุเมฺหหิ ทิฏฺฐํ วา สุตํ วา “อยํ ปุริโส มุสาวาทํ ปหาย มุสาวาทา ปฏิวิรโต”ติ ตเมนํ ราชาโน}

\vspace{\tipitakaparskip}
\csromancenter{(18) 3.\csromanspacer{} อุปาสกวคฺค คเหตฺวา มุสาวาทา เวรมณิเหตุ หนนฺติ วา พนฺธนฺติ วา ปพฺพาเชนฺติ วา ยถาปจฺจยํ วา กโรนฺตีติ.\csromanspacer{} โน เหตํ ภนฺเต.\csromanspacer{} สาธุ ภิกฺขเว, มยาปิ โข เอตํ ภิกฺขเว เนว ทิฏฺฐํ น สุตํ “อยํ ปุริโส มุสาวาทํ ปหาย มุสาวาทา ปฏิวิรโต”ติ ตเมนํ ราชาโน คเหตฺวา มุสาวาทา เวรมณิเหตุ หนนฺติ วา พนฺธนฺติ วา ปพฺพาเชนฺติ วา ยถาปจฺจยํ วา กโรนฺตีติ.\csromanspacer{} อปิ จ ขฺวสฺส ตเถว ปาปกมฺมํ ปเวเทนฺติ “อยํ ปุริโส คหปติสฺส วา คหปติปุตฺตสฺส วา มุสาวาเทน อตฺถํ ปภญฺชี”ติ\footnote{ภญฺชตีติ (สี), ภญฺชติ (สฺยา, กํ), ภญฺชีติ (อิ)}.\csromanspacer{} ตเมนํ ราชาโน คเหตฺวา มุสาวาทเหตุ หนนฺติ วา พนฺธนฺติ วา ปพฺพาเชนฺติ วา ยถาปจฺจยํ วา กโรนฺติ.\csromanspacer{} อปิ นุ ตุเมฺหหิ เอวรูปํ ทิฏฺฐํ วา สุตํ วาติ.\csromanspacer{} ทิฏฺฐํ จ โน ภนฺเต สุตํ จ สุยฺยิสฺสติ จาติ.}
\prose{\csromanfolio{185}ตํ กิํ มญฺญถ ภิกฺขเว, อปิ นุ ตุเมฺหหิ ทิฏฺฐํ วา สุตํ วา “อยํ ปุริโส สุราเมรยมชฺชปมาทฏฺฐานํ ปหาย สุราเมรยมชฺชปมาทฏฺฐานา ปฏิวิรโต”ติ ตเมนํ ราชาโน คเหตฺวา สุราเมรยมชฺชปมาทฏฺฐานา เวรมณิเหตุ หนนฺติ วา พนฺธนฺติ วา ปพฺพาเชนฺติ วา ยถาปจฺจยํ วา กโรนฺตีติ.\csromanspacer{} โน เหตํ ภนฺเต.\csromanspacer{} สาธุ ภิกฺขเว, มยาปิ โข เอตํ ภิกฺขเว เนว ทิฏฺฐํ น สุตํ “อยํ ปุริโส สุราเมรยมชฺชปมาทฏฺฐานํ ปหาย สุราเมรยมชฺชปมาทฏฺฐานา ปฏิวิรโต”ติ ตเมนํ ราชาโน คเหตฺวา สุราเมรยมชฺชปมาทฏฺฐานา เวรมณิเหตุ หนนฺติ วา พนฺธนฺติ วา ปพฺพาเชนฺติ วา ยถาปจฺจยํ วา กโรนฺตีติ.\csromanspacer{} อปิ จ ขฺวสฺส ตเถว ปาปกมฺมํ ปเวเทนฺติ “อยํ ปุริโส สุราเมรยมชฺชปมาทฏฺฐานํ อนุยุตฺโต อิตฺถิํ วา ปุริสํ วา ชีวิตา โวโรเปสิ\footnote{โวโรเปติ (สฺยา)}.\csromanspacer{} อยํ ปุริโส สุราเมรยมชฺชปมาทฏฺฐานํ อนุยุตฺโต คามา วา อรญฺญา วา อทินฺนํ เถยฺยสงฺขาตํ อาทิยิ\footnote{อาทิยติ (สี, สฺยา)}.\csromanspacer{} อยํ ปุริโส สุราเมรยมชฺชปมาทฏฺฐานํ อนุยุตฺโต ปริตฺถีสุ ปรกุมารีสุ จาริตฺตํ อาปชฺชิ\footnote{อาปชฺชติ (สี, สฺยา)}.\csromanspacer{} อยํ ปุริโส สุราเมรยมชฺชปมาทฏฺฐานํ อนุยุตฺโต คหปติสฺส วา คหปติปุตฺตสฺส วา มุสาวาเทน อตฺถํ ปภญฺชี”ติ.\csromanspacer{} ตเมนํ ราชาโน คเหตฺวา สุราเมรยมชฺชปมาทฏฺฐานเหตุ หนนฺติ วา พนฺธนฺติ วา ปพฺพาเชนฺติ วา ยถาปจฺจยํ วา กโรนฺติ.\csromanspacer{} อปิ นุ ตุเมฺหหิ เอวรูปํ ทิฏฺฐํ วา สุตํ วาติ.\csromanspacer{} ทิฏฺฐํ จ โน ภนฺเต สุตํ จ สุยฺยิสฺสติ จาติ.\csromanspacer{}\csromanspacer{}อฏฺฐมํ.}

\csromanmatikaanchor{mtk.68}
\csromantocmarkref{subsection}{9. คิหิสุตฺต}{mtk.68}
\bookmark[dest={mtk.68},level=2]{9. คิหิสุตฺต}
\csromanheader{2}{9. คิหิสุตฺต}
\proseitem{179}{\csromanfolio{186}อถ โข อนาถปิณฺฑิโก คหปติ ปญฺจมตฺเตหิ อุปาสกสเตหิ ปริวุโต เยน ภควา เตนุปสงฺกมิ, อุปสงฺกมิตฺวา ภควนฺตํ อภิวาเทตฺวา เอกมนฺตํ นิสีทิ.\csromanspacer{} อถ โข ภควา อายสฺมนฺตํ สาริปุตฺตํ อามนฺเตสิ\csromandash{}ยํ กญฺจิ\footnote{ยํ กิญฺจิ (สี, อิ)} สาริปุตฺต ชาเนยฺยาถ คิหิํ โอทาตวสนํ ปญฺจสุ สิกฺขาปเทสุ สํวุตกมฺมนฺตํ จตุนฺนํ อาภิเจตสิกานํ ทิฏฺฐธมฺมสุขวิหารานํ นิกามลาภิํ อกิจฺฉลาภิํ อกสิรลาภิํ.\csromanspacer{} โส อากงฺขมาโน อตฺตนาว อตฺตานํ พฺยากเรยฺย “ขีณนิรโยมฺหิ ขีณติรจฺฉานโยนิ ขีณเปตฺติวิสโย ขีณาปายทุคฺคติวินิปาโต, โสตาปนฺโนหมสฺมิ อวินิปาตธมฺโม นิยโต สมฺโพธิปรายโณ”ติ.}

\prose{กตเมสุ ปญฺจสุ สิกฺขาปเทสุ สํวุตกมฺมนฺโต โหติ? อิธ สาริปุตฺต อริยสาวโก ปาณาติปาตา ปฏิวิรโต โหติ, อทินฺนาทานา ปฏิวิรโต โหติ, กาเมสุมิจฺฉาจารา ปฏิวิรโต โหติ, มุสาวาทา ปฏิวิรโต โหติ, สุราเมรยมชฺชปมาทฏฺฐานา ปฏิวิรโต โหติ.\csromanspacer{} อิเมสุ ปญฺจสุ สิกฺขาปเทสุ สํวุตกมฺมนฺโต โหติ.}

\prose{กตเมสํ จตุนฺนํ อาภิเจตสิกานํ ทิฏฺฐธมฺมสุขวิหารานํ นิกามลาภี โหติ อกิจฺฉลาภี อกสิรลาภี? อิธ สาริปุตฺต อริยสาวโก พุทฺเธ อเวจฺจปฺปสาเทน สมนฺนาคโต โหติ “อิติปิ โส ภควา อรหํ สมฺมาสมฺพุทฺโธ วิชฺชาจรณสมฺปนฺโน สุคโต โลกวิทู อนุตฺตโร ปุริสทมฺมสารถิ สตฺถา เทวมนุสฺสานํ พุทฺโธ ภควา”ติ.\csromanspacer{} อยมสฺส ปฐโม อาภิเจตสิโก ทิฏฺฐธมฺมสุขวิหาโร อธิคโต โหติ อวิสุทฺธสฺส จิตฺตสฺส วิสุทฺธิยา อปริโยทาตสฺส จิตฺตสฺส ปริโยทปนาย.}

\prose{ปุน จปรํ สาริปุตฺต อริยสาวโก ธมฺเม อเวจฺจปฺปสาเทน สมนฺนาคโต โหติ “สฺวากฺขาโต ภควตา ธมฺโม สนฺทิฏฺฐิโก อกาลิโก เอหิปสฺสิโก โอปเนยฺยิโก ปจฺจตฺตํ เวทิตพฺโพ วิญฺญูหี”ติ.\csromanspacer{} อยมสฺส ทุติโย อาภิเจตสิโก ทิฏฺฐธมฺมสุขวิหาโร อธิคโต โหติ อวิสุทฺธสฺส จิตฺตสฺส วิสุทฺธิยา อปริโยทาตสฺส จิตฺตสฺส ปริโยทปนาย.}

\chapterheadpageread{(18) 3. อุปาสกวคฺค}
\prose{\csromanfolio{187}ปุน จปรํ สาริปุตฺต อริยสาวโก สํเฆ อเวจฺจปฺปสาเทน สมนฺนาคโต โหติ “สุปฺปฏิปนฺโน ภควโต สาวกสํโฆ, อุชุปฺปฏิปนฺโน ภควโต สาวกสํโฆ, ญายปฺปฏิปนฺโน ภควโต สาวกสํโฆ, สามีจิปฺปฏิปนฺโน ภควโต สาวกสํโฆ, ยทิทํ จตฺตาริ ปุริสยุคานิ อฏฺฐ ปุริสปุคฺคลา, เอส ภควโต สาวกสํโฆ อาหุเนยฺโย ปาหุเนยฺโย ทกฺขิเณยฺโย อญฺชลิกรณีโย อนุตฺตรํ ปุญฺญกฺเขตฺตํ โลกสฺสา”ติ.\csromanspacer{} อยมสฺส ตติโย อาภิเจตสิโก ทิฏฺฐธมฺมสุขวิหาโร อธิคโต โหติ อวิสุทฺธสฺส จิตฺตสฺส วิสุทฺธิยา อปริโยทาตสฺส จิตฺตสฺส ปริโยทปนาย.}

\prose{ปุน จปรํ สาริปุตฺต อริยสาวโก อริยกนฺเตหิ สีเลหิ สมนฺนาคโต โหติ อขณฺเฑหิ อจฺฉิทฺเทหิ อสพเลหิ อกมฺมาเสหิ ภุชิสฺเสหิ วิญฺญุปฺปสตฺเถหิ อปรามฏฺเฐหิ สมาธิสํวตฺตนิเกหิ.\csromanspacer{} อยมสฺส จตุตฺโถ อาภิเจตสิโก ทิฏฺฐธมฺมสุขวิหาโร อธิคโต โหติ อวิสุทฺธสฺส จิตฺตสฺส วิสุทฺธิยา อปริโยทาตสฺส จิตฺตสฺส ปริโยทปนาย.\csromanspacer{} อิเมสํ จตุนฺนํ อาภิเจตสิกานํ ทิฏฺฐธมฺมสุขวิหารานํ นิกามลาภี โหติ อกิจฺฉลาภี อกสิรลาภี.}

\prose{ยํ กญฺจิ สาริปุตฺต ชาเนยฺยาถ คิหิํ โอทาตวสนํ อิเมสุ ปญฺจสุ สิกฺขาปเทสุ สํวุตกมฺมนฺตํ อิเมสญฺจ จตุนฺนํ อาภิเจตสิกานํ ทิฏฺฐธมฺมสุขวิหารานํ นิกามลาภิํ อกิจฺฉลาภิํ อกสิรลาภิํ.\csromanspacer{} โส อากงฺขมาโน อตฺตนาว อตฺตานํ พฺยากเรยฺย “ขีณนิรโยมฺหิ ขีณติรจฺฉานโยนิ ขีณเปตฺติวิสโย ขีณาปายทุคฺคติวินิปาโต โสตาปนฺโนหมสฺมิ อวินิปาตธมฺโม นิยโต สมฺโพธิปรายโณ”ติ.}

\csromangathameasure{นิรเยสุ ภยํ ทิสฺวา, ปาปานิ ปริวชฺชเย. \\ อริยธมฺมํ สมาทาย, ปณฺฑิโต ปริวชฺชเย. \\ น หิํเส ปาณภูตานิ, วิชฺชมาเน ปรกฺกเม. \\ มุสา จ น ภเณ ชานํ, อทินฺนํ น ปรามเส. \\ เสหิ ทาเรหิ สนฺตุฏฺโฐ, ปรทารญฺจ อารเม. \\ เมรยํ วารุณิํ ชนฺตุ, น ปิเว จิตฺตโมหนิํ. \\ อนุสฺสเรยฺย สมฺพุทฺธํ, ธมฺมญฺจานุวิตกฺกเย. \\ อพฺยาปชฺชํ หิตํ จิตฺตํ, เทวโลกาย ภาวเย. \\ อุปฏฺฐิเต เทยฺยธมฺเม, ปุญฺญตฺถสฺส ชิคีสโต. \\ สนฺเตสุ ปฐมํ ทินฺนา, วิปุลา โหติ ทกฺขิณา. \\ สนฺโต หเว ปวกฺขามิ, สาริปุตฺต สุโณหิ เม. \\ อิติ กณฺหาสุ เสตาสุ, โรหิณีสุ หรีสุ วา. \\ กมฺมาสาสุ สรูปาสุ, โคสุ ปาเรวตาสุ วา. \\ ยาสุ กาสุจิ เอตาสุ, ทนฺโต ชายติ ปุงฺคโว. \\ โธรโยฺห พลสมฺปนฺโน, กลฺยาณชวนิกฺกโม. \\ ตเมว ภาเร ยุญฺชนฺติ, นาสฺส วณฺณํ ปริกฺขเร. \\ เอวเมวํ มนุสฺเสสุ, ยสฺมิํ กิสฺมิญฺจิ ชาติเย. \\ ขตฺติเย พฺราหฺมเณ เวสฺเส, สุทฺเท จณฺฑาลปุกฺกุเส. \\ ยาสุ กาสุจิ เอตาสุ, ทนฺโต ชายติ สุพฺพโต. \\ ธมฺมฏฺโฐ สีลสมฺปนฺโน, สจฺจวาที หิรีมโน. \\ ปหีนชาติมรโณ, พฺรหฺมจริยสฺส เกวลี. \\ ปนฺนภาโร วิสํยุตฺโต, กตกิจฺโจ อนาสโว. \\ ปารคู สพฺพธมฺมานํ, อนุปาทาย นิพฺพุโต. \\ ตสฺมิญฺจ วิรเช เขตฺเต, วิปุลา โหติ ทกฺขิณา. \\ พาลา จ อวิชานนฺตา, ทุมฺเมธา อสฺสุตาวิโน. \\ พหิทฺธา ททนฺติ ทานานิ, น หิ สนฺเต อุปาสเร. \\ เย จ สนฺเต อุปาสนฺติ, สปฺปญฺเญ ธีรสมฺมเต. \\ สทฺธา จ เนสํ สุคเต, มูลชาตา ปติฏฺฐิตา. \\ เทวโลกญฺจ เต ยนฺติ, กุเล วา อิธ ชายเร. \\ อนุปุพฺเพน นิพฺพานํ, อธิคจฺฉนฺติ ปณฺฑิตาติ.นวมํ.}
\csromangathagroup{\csromangathastanza{นิรเยสุ ภยํ ทิสฺวา, ปาปานิ ปริวชฺชเย. \\ อริยธมฺมํ สมาทาย, ปณฺฑิโต ปริวชฺชเย.}\csromangathastanza{น หิํเส ปาณภูตานิ, วิชฺชมาเน ปรกฺกเม. \\ มุสา จ น ภเณ ชานํ, อทินฺนํ น ปรามเส.}\csromangathastanza{เสหิ ทาเรหิ สนฺตุฏฺโฐ, ปรทารญฺจ อารเม\footnote{นารเม (สี, สฺยา)}. \\ เมรยํ วารุณิํ ชนฺตุ, น ปิเว จิตฺตโมหนิํ.}\csromangathastanza{\csromanfolio{188}อนุสฺสเรยฺย สมฺพุทฺธํ, ธมฺมญฺจานุวิตกฺกเย. \\ อพฺยาปชฺชํ\footnote{อพฺยาพชฺฌํ (?) อพฺยาปชฺฌํ (ก)} หิตํ จิตฺตํ, เทวโลกาย ภาวเย.}\csromangathastanza{อุปฏฺฐิเต เทยฺยธมฺเม, ปุญฺญตฺถสฺส ชิคีสโต\footnote{ชิคิํสโต (สี, สฺยา, กํ, อิ)}. \\ สนฺเตสุ ปฐมํ ทินฺนา, วิปุลา โหติ ทกฺขิณา.}\csromangathastanza{สนฺโต หเว ปวกฺขามิ, สาริปุตฺต สุโณหิ เม. \\ อิติ กณฺหาสุ เสตาสุ, โรหิณีสุ หรีสุ วา.}\csromangathastanza{กมฺมาสาสุ สรูปาสุ, โคสุ ปาเรวตาสุ วา. \\ ยาสุ กาสุจิ เอตาสุ, ทนฺโต ชายติ ปุงฺคโว.}\csromangathastanza{โธรโยฺห พลสมฺปนฺโน, กลฺยาณชวนิกฺกโม. \\ ตเมว ภาเร ยุญฺชนฺติ, นาสฺส วณฺณํ ปริกฺขเร.}\csromangathastanza{เอวเมวํ มนุสฺเสสุ, ยสฺมิํ กิสฺมิญฺจิ ชาติเย. \\ ขตฺติเย พฺราหฺมเณ เวสฺเส, สุทฺเท จณฺฑาลปุกฺกุเส.}\csromangathastanza{ยาสุ กาสุจิ เอตาสุ, ทนฺโต ชายติ สุพฺพโต. \\ ธมฺมฏฺโฐ สีลสมฺปนฺโน, สจฺจวาที หิรีมโน.}\csromangathastanza{ปหีนชาติมรโณ, พฺรหฺมจริยสฺส เกวลี. \\ ปนฺนภาโร วิสํยุตฺโต, กตกิจฺโจ อนาสโว.}\csromangathastanza{ปารคู สพฺพธมฺมานํ, อนุปาทาย นิพฺพุโต. \\ ตสฺมิญฺจ วิรเช เขตฺเต, วิปุลา โหติ ทกฺขิณา.}\csromangathastanza{พาลา จ อวิชานนฺตา, ทุมฺเมธา อสฺสุตาวิโน. \\ พหิทฺธา ททนฺติ ทานานิ, น หิ สนฺเต อุปาสเร.}\csromangathastanza{เย จ สนฺเต อุปาสนฺติ, สปฺปญฺเญ ธีรสมฺมเต. \\ สทฺธา จ เนสํ สุคเต, มูลชาตา ปติฏฺฐิตา.}\csromangathastanza{เทวโลกญฺจ เต ยนฺติ, กุเล วา อิธ ชายเร. \\ อนุปุพฺเพน นิพฺพานํ, อธิคจฺฉนฺติ ปณฺฑิตาติ.\csromanspacer{}\csromanspacer{}นวมํ.}}

\csromanmatikaanchor{mtk.69}
\csromantocmarkref{subsection}{10. คเวสีสุตฺต}{mtk.69}
\bookmark[dest={mtk.69},level=2]{10. คเวสีสุตฺต}
\csromanheader{2}{(18) 3. อุปาสกวคฺค \textbf{10. คเวสีสุตฺต}}
\proseitem{180}{\csromanfolio{189}เอกํ สมยํ ภควา โกสเลสุ จาริกํ จรติ มหตา ภิกฺขุ สํเฆน สทฺธิํ.\csromanspacer{} อทฺทสา โข ภควา อทฺธานมคฺคปฺปฏิปนฺโน อญฺญตรสฺมิํ ปเทเส มหนฺตํ สาลวนํ, ทิสฺวาน\footnote{ทิสฺวา (สี, อิ)} มคฺคา โอกฺกมฺม\footnote{อุกฺกมฺม (กตฺถจิ)} เยน ตํ สาลวนํ เตนุปสงฺกมิ, อุปสงฺกมิตฺวา ตํ สาลวนํ อชฺโฌคาเหตฺวา อญฺญตรสฺมิํ ปเทเส สิตํ ปาตฺวากาสิ.}

\prose{อถ โข อายสฺมโต อานนฺทสฺส เอตทโหสิ “โก นุ โข เหตุ โก ปจฺจโย ภควโต สิตสฺส ปาตุกมฺมาย, น อการเณน ตถาคตา สิตํ ปาตุกโรนฺตี”ติ.\csromanspacer{} อถ โข อายสฺมา อานนฺโท ภควนฺตํ เอตทโวจ “โก นุ โข ภนฺเต เหตุ โก ปจฺจโย ภควโต สิตสฺส ปาตุกมฺมาย, น อการเณน ตถาคตา สิตํ ปาตุกโรนฺตี”ติ.}

\prose{ภูตปุพฺพํ อานนฺท อิมสฺมิํ ปเทเส นครํ อโหสิ อิทฺธญฺเจว ผีตญฺจ พหุชนํ อากิณฺณมนุสฺสํ.\csromanspacer{} ตํ โข ปนานนฺท นครํ กสฺสโป ภควา อรหํ สมฺมาสมฺพุทฺโธ อุปนิสฺสาย วิหาสิ.\csromanspacer{} กสฺสปสฺส โข ปนานนฺท ภควโต อรหโต สมฺมาสมฺพุทฺธสฺส คเวสี นาม อุปาสโก อโหสิ สีเลสุ อปริปูรการี.\csromanspacer{} คเวสินา โข อานนฺท อุปาสเกน ปญฺจมตฺตานิ อุปาสกสตานิ ปฏิเทสิตานิ สมาทปิตานิ\footnote{สมาทาปิตานิ (?)} อเหสุํ สีเลสุ อปริปูรการิโน.\csromanspacer{} อถ โข อานนฺท คเวสิสฺส อุปาสกสฺส เอตทโหสิ “อหํ โข อิเมสํ ปญฺจนฺนํ อุปาสกสตานํ พหูปกาโร\footnote{พหุกาโร (กตฺถจิ)} ปุพฺพงฺคโม สมาทเปตา\footnote{สมาทาเปตา (?)}, อหญฺจมฺหิ สีเลสุ อปริปูรการี, อิมานิ จ ปญฺจ อุปาสกสตานิ สีเลสุ อปริปูรการิโน, อิจฺเจตํ สมสมํ นตฺถิ กิญฺจิ อติเรกํ, หนฺทาหํ อติเรกายา”ติ.}

\prose{อถ โข อานนฺท คเวสี อุปาสโก เยน ตานิ ปญฺจ อุปาสกสตานิ เตนุปสงฺกมิ, อุปสงฺกมิตฺวา ตานิ ปญฺจ อุปาสกสตานิ เอตทโวจ “อชฺชตคฺเค มํ อายสฺมนฺโต สีเลสุ ปริปูรการิํ ธาเรถา”ติ.\csromanspacer{} อถ โข อานนฺท เตสํ ปญฺจนฺนํ อุปาสกสตานํ เอตทโหสิ “อยฺโย โข คเวสี อมฺหากํ พหูปกาโร ปุพฺพงฺคโม สมาทเปตา, อยฺโย หิ นาม \csromanfolio{190}คเวสี สีเลสุ ปริปูรการี ภวิสฺสติ, กิมงฺคํ\footnote{กิมงฺค (สี, อิ)} ปน มยนฺ”ติ\footnote{ปน น มยนฺติ (สี) อํ 1. 464; วิ 4. 339; สํ 3. 328 ปิฏฺเฐสุ ปาฬิยา สํสนฺเทตพฺพํ.}.\csromanspacer{} อถ โข อานนฺท ตานิ ปญฺจ อุปาสกสตานิ เยน คเวสี อุปาสโก เตนุปสงฺกมิํสุ, อุปสงฺกมิตฺวา คเวสิํ อุปาสกํ เอตทโวจุํ “อชฺชตคฺเค อยฺโย คเวสี อิมานิปิ ปญฺจ อุปาสกสตานิ สีเลสุ ปริปูรการิโน ธาเรตู”ติ.\csromanspacer{} อถ โข อานนฺท คเวสิสฺส อุปาสกสฺส เอตทโหสิ “อหํ โข อิเมสํ ปญฺจนฺนํ อุปาสกสตานํ พหูปกาโร ปุพฺพงฺคโม สมาทเปตา, อหญฺจมฺหิ สีเลสุ ปริปูรการี, อิมานิปิ ปญฺจ อุปาสกสตานิ สีเลสุ ปริปูรการิโน, อิจฺเจตํ สมสมํ นตฺถิ กิญฺจิ อติเรกํ, หนฺทาหํ อติเรกายา”ติ.}

\prose{อถ โข อานนฺท คเวสี อุปาสโก เยน ตานิ ปญฺจ อุปาสกสตานิ เตนุปสงฺกมิ, อุปสงฺกมิตฺวา ตานิ ปญฺจ อุปาสกสตานิ เอตทโวจ “อชฺชตคฺเค มํ อายสฺมนฺโต พฺรหฺมจาริํ ธาเรถ อาราจาริํ\footnote{อนาจาริํ (อิ)} วิรตํ เมถุนา คามธมฺมา”ติ.\csromanspacer{} อถ โข อานนฺท เตสํ ปญฺจนฺนํ อุปาสกสตานํ เอตทโหสิ “อยฺโย โข คเวสี อมฺหากํ พหูปกาโร ปุพฺพงฺคโม สมาทเปตา, อยฺโย หิ นาม คเวสี พฺรหฺมจารี ภวิสฺสติ อาราจารี วิรโต เมถุนา คามธมฺมา, กิมงฺคํ ปน มยนฺ”ติ.\csromanspacer{} อถ โข อานนฺท ตานิ ปญฺจ อุปาสกสตานิ เยน คเวสี อุปาสโก เตนุปสงฺกมิํสุ, อุปสงฺกมิตฺวา คเวสิํ อุปาสกํ เอตทโวจุํ “อชฺชตคฺเค อยฺโย คเวสี อิมานิปิ ปญฺจ อุปาสกสตานิ พฺรหฺมจาริโน ธาเรตุ อาราจาริโน วิรตา เมถุนา คามธมฺมา”ติ.\csromanspacer{} อถ โข อานนฺท คเวสิสฺส อุปาสกสฺส เอตทโหสิ “อหํ โข อิเมสํ ปญฺจนฺนํ อุปาสกสตานํ พหูปกาโร ปุพฺพงฺคโม สมาทเปตา, อหญฺจมฺหิ สีเลสุ ปริปูรการี, อิมานิปิ ปญฺจ อุปาสกสตานิ สีเลสุ ปริปูรการิโน, อหญฺจมฺหิ พฺรหฺมจารี อาราจารี วิรโต เมถุนา คามธมฺมา, อิมานิปิ ปญฺจ อุปาสกสตานิ พฺรหฺมจาริโน อาราจาริโน วิรตา เมถุนา คามธมฺมา, อิจฺเจตํ สมสมํ นตฺถิ กิญฺจิ อติเรกํ, หนฺทาหํ อติเรกายา”ติ.}

\prose{อถ โข อานนฺท คเวสี อุปาสโก เยน ตานิ ปญฺจ อุปาสกสตานิ เตนุปสงฺกมิ, อุปสงฺกมิตฺวา ตานิ ปญฺจ อุปาสกสตานิ เอตทโวจ “อชฺชตคฺเค มํ อายสฺมนฺโต เอกภตฺติกํ ธาเรถ รตฺตูปรตํ}

\vspace{\tipitakaparskip}
\csromancenter{(18) 3.\csromanspacer{} อุปาสกวคฺค วิรตํ วิกาลโภชนา”ติ.\csromanspacer{} อถ โข อานนฺท เตสํ ปญฺจนฺนํ อุปาสกสตานํ เอตทโหสิ “อยฺโย โข คเวสี พหูปกาโร ปุพฺพงฺคโม สมาทเปตา, อยฺโย หิ นาม คเวสี เอกภตฺติโก ภวิสฺสติ รตฺตูปรโต วิรโต วิกาลโภชนา, กิมงฺคํ ปน มยนฺ”ติ.\csromanspacer{} อถ โข อานนฺท ตานิ ปญฺจ อุปาสกสตานิ เยน คเวสี อุปาสโก เตนุปสงฺกมิํสุ, อุปสงฺกมิตฺวา คเวสิํ อุปาสกํ เอตทโวจุํ “อชฺชตคฺเค อยฺโย คเวสี อิมานิปิ ปญฺจ อุปาสกสตานิ เอกภตฺติเก ธาเรตุ รตฺตูปรเต วิรเต วิกาลโภชนา”ติ.\csromanspacer{} อถ โข อานนฺท คเวสิสฺส อุปาสกสฺส เอตทโหสิ “อหํ โข อิเมสํ ปญฺจนฺนํ อุปาสกสตานํ พหูปกาโร ปุพฺพงฺคโม สมาทเปตา, อหญฺจมฺหิ สีเลสุ ปริปูรการี, อิมานิปิ ปญฺจ อุปาสกสตานิ สีเลสุ ปริปูรการิโน, อหญฺจมฺหิ พฺรหฺมจารี อาราจารี วิรโต เมถุนา คามธมฺมา, อิมานิปิ ปญฺจ อุปาสกสตานิ พฺรหฺมจาริโน อาราจาริโน วิรตา เมถุนา คามธมฺมา, อหญฺจมฺหิ เอกภตฺติโก รตฺตูปรโต วิรโต วิกาลโภชนา, อิมานิปิ ปญฺจ อุปาสกสตานิ เอกภตฺติกา รตฺตูปรตา วิรตา วิกาลโภชนา, อิจฺเจตํ สมสมํ นตฺถิ กิญฺจิ อติเรกํ, หนฺทาหํ อติเรกายา”ติ.}
\prose{\csromanfolio{191}อถ โข อานนฺท คเวสี อุปาสโก เยน กสฺสโป ภควา อรหํ สมฺมาสมฺพุทฺโธ เตนุปสงฺกมิ, อุปสงฺกมิตฺวา กสฺสปํ ภควนฺตํ อรหนฺตํ สมฺมาสมฺพุทฺธํ เอตทโวจ “ลเภยฺยาหํ ภนฺเต ภควโต สนฺติเก ปพฺพชฺชํ, ลเภยฺยํ อุปสมฺปทนฺ”ติ.\csromanspacer{} อลตฺถ โข อานนฺท คเวสี อุปาสโก กสฺสปสฺส ภควโต อรหโต สมฺมาสมฺพุทฺธสฺส สนฺติเก ปพฺพชฺชํ, อลตฺถ อุปสมฺปทํ.\csromanspacer{} อจิรูปสมฺปนฺโน โข ปนานนฺท คเวสี ภิกฺขุ เอโก วูปกฏฺโฐ อปฺปมตฺโต อาตาปี ปหิตตฺโต วิหรนฺโต นจิรสฺเสว ยสฺสตฺถาย กุลปุตฺตา สมฺมเทว อคารสฺมา อนคาริยํ ปพฺพชนฺติ, ตทนุตฺตรํ พฺรหฺมจริยปริโยสานํ ทิฏฺเฐว ธมฺเม สยํ อภิญฺญา สจฺฉิกตฺวา อุปสมฺปชฺช วิหาสิ, “ขีณา ชาติ, วุสิตํ พฺรหฺมจริยํ, กตํ กรณียํ, นาปรํ อิตฺถตฺตายา”ติ อพฺภญฺญาสิ.\csromanspacer{} อญฺญาตโร จ ปนานนฺท คเวสี ภิกฺขุ อรหตํ อโหสิ. อถ โข อานนฺท เตสํ ปญฺจนฺนํ อุปาสกสตานํ เอตทโหสิ “อยฺโย โข คเวสี อมฺหากํ พหูปกาโร ปุพฺพงฺคโม สมาทเปตา.\csromanspacer{} อยฺโย หิ นาม คเวสี เกสมสฺสุํ โอหาเรตฺวา กาสายานิ วตฺถานิ \csromanfolio{192}อจฺฉาเทตฺวา อคารสฺมา อนคาริยํ ปพฺพชิสฺสติ, กิมงฺคํ ปน มยนฺ”ติ.\csromanspacer{} อถ โข อานนฺท ตานิ ปญฺจ อุปาสกสตานิ เยน กสฺสโป ภควา อรหํ สมฺมาสมฺพุทฺโธ เตนุปสงฺกมิํสุ, อุปสงฺกมิตฺวา กสฺสปํ ภควนฺตํ อรหนฺตํ สมฺมาสมฺพุทฺธํ เอตทโวจุํ “ลเภยฺยาม มยํ ภนฺเต ภควโต สนฺติเก ปพฺพชฺชํ, ลเภยฺยาม อุปสมฺปทนฺ”ติ.\csromanspacer{} อลภิํสุ โข อานนฺท ตานิ ปญฺจ อุปาสกสตานิ กสฺสปสฺส ภควโต อรหโต สมฺมาสมฺพุทฺธสฺส สนฺติเก ปพฺพชฺชํ, อลภิํสุ อุปสมฺปทํ.}

\prose{อถ โข อานนฺท คเวสิสฺส ภิกฺขุโน เอตทโหสิ “อหํ โข อิมสฺส อนุตฺตรสฺส วิมุตฺติสุขสฺส นิกามลาภี โหมิ อกิจฺฉลาภี อกสิรลาภี, อโห วติมานิปิ ปญฺจ ภิกฺขุสตานิ อิมสฺส อนุตฺตรสฺส วิมุตฺติสุขสฺส นิกามลาภิโน อสฺสุ อกิจฺฉลาภิโน อกสิรลาภิโน”ติ.\csromanspacer{} อถ โข อานนฺท ตานิ ปญฺจ ภิกฺขุสตานิ วูปกฏฺฐา\footnote{ภิกฺขุสตานิ เอเกกา วูปกฏฺฐา (สฺยา, กํ)} อปฺปมตฺตา อาตาปิโน ปหิตตฺตา วิหรนฺตา นจิรสฺเสว ยสฺสตฺถาย กุลปุตฺตา สมฺมเทว อคารสฺมา อนคาริยํ ปพฺพชนฺติ, ตทนุตฺตรํ พฺรหฺมจริยปริโยสานํ ทิฏฺเฐว ธมฺเม สยํ อภิญฺญา สจฺฉิกตฺวา อุปสมฺปชฺช วิหริํสุ, “ขีณา ชาติ, วุสิตํ พฺรหฺมจริยํ, กตํ กรณียํ, นาปรํ อิตฺถตฺตายา”ติ อพฺภญฺญิํสุ.}

\prose{อิติ โข อานนฺท ตานิ ปญฺจ ภิกฺขุสตานิ คเวสีปมุขานิ อุตฺตรุตฺตริ\footnote{อุตฺตรุตฺตริํ (สี, สฺยา, กํ, อิ)} ปณีตปณีตํ วายมมานา อนุตฺตรํ วิมุตฺติํ สจฺฉากํสุ.\csromanspacer{} ตสฺมาติห อานนฺท เอวํ สิกฺขิตพฺพํ “อุตฺตรุตฺตริ ปณีตปณีตํ วายมมานา อนุตฺตรํ วิมุตฺติํ สจฺฉิกริสฺสามา”ติ.\csromanspacer{} เอวํ หิ โว อานนฺท สิกฺขิตพฺพนฺติ.\csromanspacer{}\csromanspacer{}ทสมํ.\csromanspacer{} อุปาสกวคฺโค ตติโย.}

\vspace{\tipitakaparskip}
\csromancenter{\textbf{ตสฺสุทฺทานํ.}}
\vspace{0.5\baselineskip}
\csromangathameasure{สารชฺชํ วิสารโท นิรยํ, เวรํ จณฺฑาลปญฺจมํ. \\ ปีติ วณิชฺชา ราชาโน, คิหี เจว คเวสินาติ.}
\csromangathagroup{\csromangathastanza{สารชฺชํ วิสารโท นิรยํ, เวรํ จณฺฑาลปญฺจมํ. \\ ปีติ วณิชฺชา ราชาโน, คิหี เจว คเวสินาติ.}}

\csromanmatikaanchor{mtk.70}
\csromantocmarknopage{section}{(19) 4. อรญฺญวคฺค}{mtk.70}
\bookmark[dest={mtk.70},level=1]{(19) 4. อรญฺญวคฺค}
\csromanheader{1}{\textbf{(19) 4. อรญฺญวคฺค} \textbf{(19) 4. อรญฺญวคฺค}}
\csromanheader{2}{1. อารญฺญิกสุตฺต}
\csromanmatikaanchor{mtk.71}
\csromantocmarkref{subsection}{1-10. อารญฺญิกสุตฺตาทิ}{mtk.71}
\bookmark[dest={mtk.71},level=2]{1-10. อารญฺญิกสุตฺตาทิ}
\proseitem{181}{\csromanfolio{193}ปญฺจิเม ภิกฺขเว อารญฺญิกา\footnote{อารญฺญกา (สพฺพตฺถ) วิ 5. 335 ปิฏฺเฐ ปสฺสิตพฺพํ.}.\csromanspacer{} กตเม ปญฺจ? มนฺทตฺตา โมมูหตฺตา อารญฺญิโก โหติ, ปาปิจฺโฉ อิจฺฉาปกโต อารญฺญิโก โหติ, อุมฺมาทา จิตฺตกฺเขปา อารญฺญิโก โหติ, “วณฺณิตํ พุทฺเธหิ พุทฺธสาวเกหี”ติ อารญฺญิโก โหติ, อปฺปิจฺฉตํเยว นิสฺสาย สนฺตุฏฺฐิํเยว นิสฺสาย สลฺเลขํเยว นิสฺสาย ปวิเวกํเยว นิสฺสาย อิทมตฺถิตํเยว\footnote{อิทมฏฺฐิตํเยว (สี, อิ)} นิสฺสาย อารญฺญิโก โหติ.\csromanspacer{} อิเม โข ภิกฺขเว ปญฺจ อารญฺญิกา.\csromanspacer{} อิเมสํ โข ภิกฺขเว ปญฺจนฺนํ อารญฺญิกานํ ยฺวายํ อารญฺญิโก อปฺปิจฺฉตํเยว นิสฺสาย สนฺตุฏฺฐิํเยว นิสฺสาย สลฺเลขํเยว นิสฺสาย ปวิเวกํเยว นิสฺสาย อิทมตฺถิตํเยว นิสฺสาย อารญฺญิโก โหติ.\csromanspacer{} อยํ อิเมสํ ปญฺจนฺนํ อารญฺญิกานํ อคฺโค จ เสฏฺโฐ จ โมกฺโข\footnote{ปาโมกฺโข (อํ 1. 408; อํ 3. 403 ปิฏฺเฐ)} จ อุตฺตโม จ ปวโร จ.}

\prose{เสยฺยถาปิ ภิกฺขเว ควา ขีรํ, ขีรมฺหา ทธิ, ทธิมฺหา นวนีตํ, นวนีตมฺหา สปฺปิ, สปฺปิมฺหา สปฺปิมณฺโฑ, สปฺปิมณฺโฑ\footnote{สปฺปิมฺหา สปฺปิมณฺโฑ (ก) สํ 2. 224 ปิฏฺเฐปิ.} ตตฺถ อคฺคมกฺขายติ.\csromanspacer{} เอวเมวํ โข ภิกฺขเว อิเมสํ ปญฺจนฺนํ อารญฺญิกานํ ยฺวายํ อารญฺญิโก อปฺปิจฺฉตํเยว นิสฺสาย สนฺตุฏฺฐิํเยว นิสฺสาย สลฺเลขํเยว นิสฺสาย ปวิเวกํเยว นิสฺสาย อิทมตฺถิตํเยว นิสฺสาย อารญฺญิโก โหติ, อยํ อิเมสํ ปญฺจนฺนํ อารญฺญิกานํ อคฺโค จ เสฏฺโฐ จ โมกฺโข จ อุตฺตโม จ ปวโร จาติ.\csromanspacer{}\csromanspacer{}ปฐมํ.}

\csromanheader{2}{2. จีวรสุตฺต}
\proseitem{182}{ปญฺจิเม ภิกฺขเว ปํสุกูลิกา.\csromanspacer{} กตเม ปญฺจ? มนฺทตฺตา โมมูหตฺตา ปํสุกูลิโก โหติ ฯเปฯ อิทมตฺถิตํเยว นิสฺสาย ปํสุกูลิโก โหติ.\csromanspacer{} อิเม โข ภิกฺขเว ปญฺจ ปํสุกูลิกาติ.\csromanspacer{}\csromanspacer{}ทุติยํ.}

\csromanheader{2}{3. รุกฺขมูลิกสุตฺต}
\proseitem{183}{ปญฺจิเม ภิกฺขเว รุกฺขมูลิกา.\csromanspacer{} กตเม ปญฺจ? มนฺทตฺตา โมมูหตฺตา รุกฺขมูลิโก โหติ ฯเปฯ อิทมตฺถิตํเยว นิสฺสาย รุกฺขมูลิโก โหติ.\csromanspacer{} อิเม โข ภิกฺขเว ปญฺจ รุกฺขมูลิกาติ.\csromanspacer{}\csromanspacer{}ตติยํ.}

\csromanheader{2}{4. โสสานิกสุตฺต}
\proseitem{184}{\csromanfolio{194}ปญฺจิเม ภิกฺขเว โสสานิกา.\csromanspacer{} กตเม ปญฺจ? มนฺทตฺตา โมมูหตฺตา โสสานิโก โหติ ฯเปฯ อิทมตฺถิตํเยว นิสฺสาย โสสานิโก โหติ.\csromanspacer{} อิเม โข ภิกฺขเว ปญฺจ โสสานิกาติ.\csromanspacer{}\csromanspacer{}จตุตฺถํ.}

\csromanheader{2}{5. อพฺโภกาสิกสุตฺต}
\vspace{\tipitakaparskip}
\csromancenter{ปญฺจิเม ภิกฺขเว อพฺโภกาสิกา ฯเปฯ.\csromanspacer{}\csromanspacer{}ปญฺจมํ.}
\csromanheader{2}{6. เนสชฺชิกสุตฺต}
\vspace{\tipitakaparskip}
\csromancenter{ปญฺจิเม ภิกฺขเว เนสชฺชิกา ฯเปฯ.\csromanspacer{}\csromanspacer{}ฉฏฺฐํ.}
\csromanheader{2}{7. ยถาสนฺถติกสุตฺต}
\vspace{\tipitakaparskip}
\csromancenter{ปญฺจิเม ภิกฺขเว ยถาสนฺถติกา ฯเปฯ.\csromanspacer{}\csromanspacer{}สตฺตมํ.}
\csromanheader{2}{8. เอกาสนิกสุตฺต}
\vspace{\tipitakaparskip}
\csromancenter{ปญฺจิเม ภิกฺขเว เอกาสนิกา ฯเปฯ.\csromanspacer{}\csromanspacer{}อฏฺฐมํ.}
\csromanheader{2}{9. ขลุปจฺฉาภตฺติกสุตฺต}
\vspace{\tipitakaparskip}
\csromancenter{ปญฺจิเม ภิกฺขเว ขลุปจฺฉาภตฺติกา ฯเปฯ.\csromanspacer{}\csromanspacer{}นวมํ.}
\csromanheader{2}{10. ปตฺตปิณฺฑิกสุตฺต}
\proseitem{190}{ปญฺจิเม ภิกฺขเว ปตฺตปิณฺฑิกา.\csromanspacer{} กตเม ปญฺจ? มนฺทตฺตา โมมูหตฺตา ปตฺตปิณฺฑิโก โหติ, ปาปิจฺโฉ อิจฺฉาปกโต ปตฺตปิณฺฑิโก โหติ, อุมฺมาทา จิตฺตกฺเขปา ปตฺตปิณฺฑิโก โหติ, “วณฺณิตํ พุทฺเธหิ พุทฺธสาวเกหี”ติ ปตฺตปิณฺฑิโก โหติ.\csromanspacer{} อปฺปิจฺฉตํเยว นิสฺสาย สนฺตุฏฺฐิํเยว นิสฺสาย สลฺเลขํเยว นิสฺสาย ปวิเวกํเยว นิสฺสาย อิทมตฺถิตํเยว นิสฺสาย ปตฺตปิณฺฑิโก โหติ.\csromanspacer{} อิเม โข ภิกฺขเว ปญฺจ ปตฺตปิณฺฑิกา.\csromanspacer{} อิเมสํ โข ภิกฺขเว ปญฺจนฺนํ ปตฺตปิณฺฑิกานํ ยฺวายํ ปตฺตปิณฺฑิโก อปฺปิจฺฉตํเยว นิสฺสาย สนฺตุฏฺฐิํเยว นิสฺสาย}

\vspace{\tipitakaparskip}
\csromancenter{\textbf{(20) 5.}\csromanspacer{}\textbf{ พฺราหฺมณวคฺค} สลฺเลขํเยว นิสฺสาย ปวิเวกํเยว นิสฺสาย อิทมตฺถิตํเยว นิสฺสาย ปตฺตปิณฺฑิโก โหติ.\csromanspacer{} อยํ อิเมสํ ปญฺจนฺนํ ปตฺตปิณฺฑิกานํ อคฺโค จ เสฏฺโฐ จ โมกฺโข จ อุตฺตโม จ ปวโร จ.}
\prose{\csromanfolio{195}เสยฺยถาปิ ภิกฺขเว ควา ขีรํ, ขีรมฺหา ทธิ, ทธิมฺหา นวนีตํ, นวนีตมฺหา สปฺปิ, สปฺปิมฺหา สปฺปิมณฺโฑ, สปฺปิมณฺโฑ ตตฺถ อคฺคมกฺขายติ.\csromanspacer{} เอวเมวํ โข ภิกฺขเว อิเมสํ ปญฺจนฺนํ ปตฺตปิณฺฑิกานํ ยฺวายํ ปตฺตปิณฺฑิโก อปฺปิจฺฉตํเยว นิสฺสาย สนฺตุฏฺฐิํเยว นิสฺสาย สลฺเลขํเยว นิสฺสาย ปวิเวกํ เยว นิสฺสาย อิทมตฺถิตํเยว นิสฺสาย ปตฺตปิณฺฑิโก โหติ.\csromanspacer{} อยํ อิเมสํ ปญฺจนฺนํ ปตฺตปิณฺฑิกานํ อคฺโค จ เสฏฺโฐ จ โมกฺโข จ อุตฺตโม จ ปวโร จาติ.\csromanspacer{}\csromanspacer{}ทสมํ.\csromanspacer{} อรญฺญวคฺโค จตุตฺโถ.}

\titlehead{\textbf{ตสฺสุทฺทานํ}}
\csromangathameasure{อรญฺญํ จีวรํ รุกฺข, สุสานํ อพฺโภกาสิกํ. \\ เนสชฺชํ สนฺถตํ เอกาสนิกํ, ขลุปจฺฉาปิณฺฑิเกน จาติ.}
\csromangathagroup{\csromangathastanza{อรญฺญํ จีวรํ รุกฺข, สุสานํ อพฺโภกาสิกํ. \\ เนสชฺชํ สนฺถตํ เอกาสนิกํ, ขลุปจฺฉาปิณฺฑิเกน จาติ.}}

\csromanmatikaanchor{mtk.72}
\csromantocmarknopage{section}{(20) 5. พฺราหฺมณวคฺค}{mtk.72}
\bookmark[dest={mtk.72},level=1]{(20) 5. พฺราหฺมณวคฺค}
\csromanheader{1}{\textbf{(20) 5. พฺราหฺมณวคฺค}}
\csromanmatikaanchor{mtk.73}
\csromantocmarkref{subsection}{1. โสณสุตฺต}{mtk.73}
\bookmark[dest={mtk.73},level=2]{1. โสณสุตฺต}
\csromanheader{2}{1. โสณสุตฺต}
\proseitem{191}{ปญฺจิเม ภิกฺขเว โปราณา พฺราหฺมณธมฺมา เอตรหิ สุนเขสุ สนฺทิสฺสนฺติ โน พฺราหฺมเณสุ.\csromanspacer{} กตเม ปญฺจ? ปุพฺเพ สุทํ\footnote{ปุพฺพสฺสุทํ (ก)} ภิกฺขเว พฺราหฺมณา พฺราหฺมณิํเยว คจฺฉนฺติ โน อพฺราหฺมณิํ.\csromanspacer{} เอตรหิ ภิกฺขเว พฺราหฺมณา พฺราหฺมณิมฺปิ คจฺฉนฺติ, อพฺราหฺมณิมฺปิ คจฺฉนฺติ.\csromanspacer{} เอตรหิ ภิกฺขเว สุนขา สุนขิํเยว คจฺฉนฺติ โน อสุนขิํ.\csromanspacer{} อยํ ภิกฺขเว ปฐโม โปราโณ พฺราหฺมณธมฺโม เอตรหิ สุนเขสุ สนฺทิสฺสติ โน พฺราหฺมเณสุ.}

\csromanmatikaanchor{mtk.74}
\csromantocmarkref{subsection}{2. โทณพฺราหฺมณสุตฺต}{mtk.74}
\bookmark[dest={mtk.74},level=2]{2. โทณพฺราหฺมณสุตฺต}
\prose{ปุพฺเพ สุทํ ภิกฺขเว พฺราหฺมณา พฺราหฺมณิํ อุตุนิํเยว คจฺฉนฺติ โน อนุตุนิํ.\csromanspacer{} เอตรหิ ภิกฺขเว พฺราหฺมณา พฺราหฺมณิํ อุตุนิมฺปิ คจฺฉนฺติ, อนุตุนิมฺปิ \csromanfolio{196}คจฺฉนฺติ.\csromanspacer{} เอตรหิ ภิกฺขเว สุนขา สุนขิํ อุตุนิํเยว คจฺฉนฺติ โน อนุตุนิํ.\csromanspacer{} อยํ ภิกฺขเว ทุติโย โปราโณ พฺราหฺมณธมฺโม เอตรหิ สุนเขสุ สนฺทิสฺสติ โน พฺราหฺมเณสุ.}

\prose{ปุพฺเพ สุทํ ภิกฺขเว พฺราหฺมณา พฺราหฺมณิํ เนว กิณนฺติ โน วิกฺกิณนฺติ, สมฺปิเยเนว สํวาสํ สํพนฺธาย\footnote{สํสคฺคตฺถาย (สี, อิ)} สํปวตฺเตนฺติ.\csromanspacer{} เอตรหิ ภิกฺขเว พฺราหฺมณา พฺราหฺมณิํ กิณนฺติปิ วิกฺกิณนฺติปิ, สมฺปิเยนปิ สํวาสํ สํพนฺธาย สํปวตฺเตนฺติ.\csromanspacer{} เอตรหิ ภิกฺขเว สุนขา สุนขิํ เนว กิณนฺติ โน วิกฺกิณนฺติ, สมฺปิเยเนว สํวาสํ สํพนฺธาย สํปวตฺเตนฺติ.\csromanspacer{} อยํ ภิกฺขเว ตติโย โปราโณ พฺราหฺมณธมฺโม เอตรหิ สุนเขสุ สนฺทิสฺสติ โน พฺราหฺมเณสุ.}

\prose{ปุพฺเพ สุทํ ภิกฺขเว พฺราหฺมณา น สนฺนิธิํ กโรนฺติ ธนสฺสปิ ธญฺญสฺสปิ รชตสฺสปิ ชาตรูปสฺสปิ.\csromanspacer{} เอตรหิ ภิกฺขเว พฺราหฺมณา สนฺนิธิํ กโรนฺติ ธนสฺสปิ ธญฺญสฺสปิ รชตสฺสปิ ชาตรูปสฺสปิ.\csromanspacer{} เอตรหิ ภิกฺขเว สุนขา น สนฺนิธิํ กโรนฺติ ธนสฺสปิ ธญฺญสฺสปิ รชตสฺสปิ ชาตรูปสฺสปิ.\csromanspacer{} อยํ ภิกฺขเว จตุตฺโถ โปราโณ พฺราหฺมณธมฺโม เอตรหิ สุนเขสุ สนฺทิสฺสติ โน พฺราหฺมเณสุ.}

\prose{ปุพฺเพ สุทํ ภิกฺขเว พฺราหฺมณา สายํ สายมาสาย, ปาโต ปาตราสาย ภิกฺขํ ปริเยสนฺติ.\csromanspacer{} เอตรหิ ภิกฺขเว พฺราหฺมณา ยาวทตฺถํ อุทราวเทหกํ ภุญฺชิตฺวา อวเสสํ อาทาย ปกฺกมนฺติ.\csromanspacer{} เอตรหิ ภิกฺขเว สุนขา สายํ สายมาสาย, ปาโต ปาตราสาย ภิกฺขํ ปริเยสนฺติ.\csromanspacer{} อยํ ภิกฺขเว ปญฺจโม โปราโณ พฺราหฺมณธมฺโม เอตรหิ สุนเขสุ สนฺทิสฺสติ โน พฺราหฺมเณสุ.\csromanspacer{} อิเม โข ภิกฺขเว ปญฺจ โปราณา พฺราหฺมณธมฺมา เอตรหิ สุนเขสุ สนฺทิสฺสนฺติ โน พฺราหฺมเณสูติ.\csromanspacer{}\csromanspacer{}ปฐมํ.}

\csromanheader{2}{2. โทณ พฺราหฺมณสุตฺต}
\proseitem{192}{อถ โข โทโณ พฺราหฺมโณ เยน ภควา เตนุปสงฺกมิ, อุปสงฺกมิตฺวา ภควตา สทฺธิํ สมฺโมทิ, สมฺโมทนียํ กถํ สารณียํ วีติสาเรตฺวา เอกมนฺตํ นิสีทิ, เอกมนฺตํ นิสินฺโน โข โทโณ พฺราหฺมโณ ภควนฺตํ เอตทโวจ\csromandash{}}

\chapterheadpageread{(20) 5. พฺราหฺมณวคฺค}
\prose{\csromanfolio{197}สุตํ เมตํ โภ โคตม, “น สมโณ โคตโม พฺราหฺมเณ ชิณฺเณ วุทฺเธ มหลฺลเก อทฺธคเต วโย อนุปฺปตฺเต อภิวาเทติ วา ปจฺจุฏฺเฐติ วา อาสเนน วา นิมนฺเตตี”ติ, ตยิทํ โภ โคตม ตเถว.\csromanspacer{} น หิ ภวํ โคตโม พฺราหฺมเณ ชิณฺเณ วุฑฺเฒ มหลฺลเก อทฺธคเต วโย อนุปฺปตฺเต อภิวาเทติ วา ปจฺจุฏฺเฐติ วา อาสเนน วา นิมนฺเตติ.\csromanspacer{} ตยิทํ โภ โคตม น สมฺปนฺนเมวาติ.\csromanspacer{} ตฺวมฺปิ โน โทณ พฺราหฺมโณ ปฏิชานาสีติ.\csromanspacer{} ยํ หิ ตํ โภ โคตม สมฺมา วทมาโน วเทยฺย “พฺราหฺมโณ อุภโต สุชาโต มาติโต จ ปิติโต จ สํสุทฺธคหณิโก ยาว สตฺตมา ปิตามหยุคา อกฺขิตฺโต อนุปกฺกุฏฺโฐ ชาติวาเทน, อชฺฌายโก มนฺตธโร ติณฺณํ เวทานํ ปารคู สนิฆณฺฑุเกฏุภานํ สากฺขรปฺปเภทานํ อิติหาสปญฺจมานํ ปทโก เวยฺยากรโณ โลกายตมหาปุริสลกฺขเณสุ อนวโย”ติ, มเมว ตํ โภ โคตม สมฺมา วทมาโน วเทยฺย “อหํ หิ โภ โคตม พฺราหฺมโณ อุภโต สุชาโต มาติโต จ ปิติโต จ สํสุทฺธคหณิโก ยาว สตฺตมา ปิตามหยุคา อกฺขิตฺโต อนุปกฺกุฏฺโฐ ชาติวาเทน, อชฺฌายโก มนฺตธโร ติณฺณํ เวทานํ ปารคู สนิฆณฺฑุเกฏุภานํ สากฺขรปฺปเภทานํ อิติหาสปญฺจมานํ ปทโก เวยฺยากรโณ โลกายตมหาปุริสลกฺขเณสุ อนวโย”ติ.}

\prose{เย โข เต โทณ พฺราหฺมณานํ ปุพฺพกา อิสโย มนฺตานํ กตฺตาโร มนฺตานํ ปวตฺตาโร, เยสมิทํ เอตรหิ พฺราหฺมณา โปราณํ มนฺตปทํ คีตํ ปวุตฺตํ สมิหิตํ ตทนุคายนฺติ ตทนุภาสนฺติ ภาสิตมนุภาสนฺติ สชฺฌายิตมนุสชฺฌายนฺติ วาจิตมนุวาเจนฺติ.\csromanspacer{} เสยฺยถิทํ, อฏฺฐโก วามโก วามเทโว เวสฺสามิตฺโต ยมทคฺคิ\footnote{ยมตคฺคิ (สี) ที 1. 97, 224, 228; ม 2. 381; วิ 3. 343; อุปริ 202 ปิฏฺเฐสุปิ ปสฺสิตพฺพํ.} องฺคีรโส ภารทฺวาโช วาเสฏฺโฐ กสฺสโป ภคุ, ตฺยสฺสุ’เม ปญฺจ พฺราหฺมเณ ปญฺญาเปนฺติ พฺรหฺมสมํ เทวสมํ มริยาทํ สมฺภินฺนมริยาทํ พฺราหฺมณจณฺฑาลํเยว ปญฺจมํ, เตสํ ตฺวํ โทณ กตโมติ.}

\prose{น โข มยํ โภ โคตม ปญฺจ พฺราหฺมเณ ชานาม, อถ โข มยํ พฺราหฺมณาเตฺวว ชานาม.\csromanspacer{} สาธุ เม ภวํ โคตโม ตถา ธมฺมํ เทเสตุ, ยถา อหํ อิเม ปญฺจ พฺราหฺมเณ ชาเนยฺยนฺติ.\csromanspacer{} เตน หิ พฺราหฺมณ สุโณหิ}

\prose{\csromanfolio{198}สาธุกํ มนสิ กโรหิ, ภาสิสฺสามีติ. “เอวํ โภ”ติ โข โทโณ พฺราหฺมโณ ภควโต ปจฺจสฺโสสิ.\csromanspacer{} ภควา เอตทโวจ\csromandash{}}

\prose{กถญฺจ โทณ พฺราหฺมโณ พฺรหฺมสโม โหติ? อิธ โทณ พฺราหฺมโณ อุภโต สุชาโต โหติ มาติโต จ ปิติโต จ สํสุทฺธคหณิโก ยาว สตฺตมา ปิตามหยุคา อกฺขิตฺโต อนุปกฺกุฏฺโฐ ชาติวาเทน.\csromanspacer{} โส อฏฺฐจตฺตาลีสวสฺสานิ โกมารพฺรหฺมจริยํ\footnote{โกมารํ พฺรหฺมจริยํ (สฺยา, ก)} จรติ มนฺเต อธียมาโน, อฏฺฐจตฺตาลีสวสฺสานิ โกมารพฺรหฺมจริยํ จริตฺวา มนฺเต อธียิตฺวา อาจริยสฺส อาจริยธนํ ปริเยสติ ธมฺเมเนว โน อธมฺเมน.}

\prose{ตตฺถ จ โทณ โก ธมฺโม, เนว กสิยา น วณิชฺชาย น โครกฺเขน น อิสฺสตฺเถน\footnote{น อิสฺสตฺเตน (ก)} น ราชโปริเสน น สิปฺปญฺญตเรน, เกวลํ ภิกฺขาจริยาย กปาลํ อนติมญฺญมาโน.\csromanspacer{} โส อาจริยสฺส อาจริยธนํ นิยฺยาเทตฺวา\footnote{นียฺยาเทตฺวา (สี), นียาเทตฺวา (อิ), นิยฺยาเตตฺวา (กตฺถจิ)} เกสมสฺสุํ โอหาเรตฺวา กาสายานิ วตฺถานิ อจฺฉาเทตฺวา อคารสฺมา อนคาริยํ ปพฺพชติ.\csromanspacer{} โส เอวํ ปพฺพชิโต สมาโน เมตฺตาสหคเตน เจตสา เอกํ ทิสํ ผริตฺวา วิหรติ.\csromanspacer{} ตถา ทุติยํ.\csromanspacer{} ตถา ตติยํ.\csromanspacer{} ตถา จตุตฺถํ\footnote{จตุตฺถิํ (สี)}.\csromanspacer{} อิติ อุทฺธมโธ ติริยํ สพฺพธิ สพฺพตฺตตาย สพฺพาวนฺตํ โลกํ เมตฺตาสหคเตน เจตสา วิปุเลน มหคฺคเตน อปฺปมาเณน อเวเรน อพฺยาปชฺเชน\footnote{อพฺยาปชฺเฌน (ก) อพฺยาพชฺเฌน (?)} ผริตฺวา วิหรติ.\csromanspacer{} กรุณา ฯเปฯ.\csromanspacer{} มุทิตา.\csromanspacer{} อุเปกฺขาสหคเตน เจตสา เอกํ ทิสํ ผริตฺวา วิหรติ.\csromanspacer{} ตถา ทุติยํ.\csromanspacer{} ตถา ตติยํ.\csromanspacer{} ตถา จตุตฺถํ.\csromanspacer{} อิติ อุทฺธมโธ ติริยํ สพฺพธิ สพฺพตฺตตาย สพฺพาวนฺตํ โลกํ อุเปกฺขาสหคเตน เจตสา วิปุเลน มหคฺคเตน อปฺปมาเณน อเวเรน อพฺยาปชฺเชน ผริตฺวา วิหรติ.\csromanspacer{} โส อิเม จตฺตาโร พฺรหฺมวิหาเร ภาเวตฺวา กายสฺส เภทา ปรํ มรณา สุคติํ พฺรหฺม โลกํ อุปปชฺชติ.\csromanspacer{} เอวํ โข โทณ พฺราหฺมโณ พฺรหฺมสโม โหติ. (1)}

\prose{กถญฺจ โทณ พฺราหฺมโณ เทวสโม โหติ? อิธ โทณ พฺราหฺมโณ อุภโต สุชาโต โหติ มาติโต จ ปิติโต จ สํสุทฺธคหณิโก ยาว สตฺตมา ปิตามหยุคา อกฺขิตฺโต อนุปกฺกุฏฺโฐ}

\vspace{\tipitakaparskip}
\csromancenter{(20) 5.\csromanspacer{} พฺราหฺมณวคฺค ชาติวาเทน.\csromanspacer{} โส อฏฺฐจตฺตาลีสวสฺสานิ โกมารพฺรหฺมจริยํ จรติ มนฺเต อธียมาโน, อฏฺฐจตฺตาลีสวสฺสานิ โกมารพฺรหฺมจริยํ จริตฺวา มนฺเต อธียิตฺวา อาจริยสฺส อาจริยธนํ ปริเยสติ ธมฺเมเนว โน อธมฺเมน.\csromanspacer{} ตตฺถ จ โทณ โก ธมฺโม เนว กสิยา น วณิชฺชาย น โครกฺเขน น อิสฺสตฺเถน น ราชโปริเสน น สิปฺปญฺญตเรน, เกวลํ ภิกฺขาจริยาย กปาลํ อนติมญฺญมาโน.\csromanspacer{} โส อาจริยสฺส อาจริยธนํ นิยฺยาเทตฺวา ทารํ ปริเยสติ ธมฺเมเนว โน อธมฺเมน.}
\prose{\csromanfolio{199}ตตฺถ จ โทณ โก ธมฺโม, เนว กเยน น วิกฺกเยน พฺราหฺมณิํเยว อุทกูปสฺสฏฺฐํ.\csromanspacer{} โส พฺราหฺมณิํเยว คจฺฉติ, น กตฺติยิํ น เวสฺสิํ น สุทฺทิํ น จณฺฑาลิํ น เนสาทิํ น เวนิํ\footnote{น เวณิํ (สี, สฺยา, กํ, อิ)} น รถการิํ น ปุกฺกุสิํ คจฺฉติ, น คพฺภินิํ คจฺฉติ, น ปายมานํ คจฺฉติ, น อนุตุนิํ คจฺฉติ.\csromanspacer{} กสฺมา จ โทณ พฺราหฺมโณ น คพฺภินิํ คจฺฉติ.\csromanspacer{} สเจ โทณ พฺราหฺมโณ คพฺภินิํ คจฺฉติ, อติมีฬฺหโช นาม โส โหติ มาณวโก วา มาณวิกา\footnote{มาณวกี (ก)} วา, ตสฺมา โทณ พฺราหฺมโณ น คพฺภินิํ คจฺฉติ.\csromanspacer{} กสฺมา จ โทณ พฺราหฺมโณ น ปายมานํ คจฺฉติ.\csromanspacer{} สเจ โทณ พฺราหฺมโณ ปายมานํ คจฺฉติ, อสุจิปฏิปีฬิโต นาม โส โหติ มาณวโก วา มาณวิกา วา, ตสฺมา โทณ พฺราหฺมโณ น ปายมานํ คจฺฉติ.\csromanspacer{} ตสฺส สา โหติ พฺราหฺมณี เนว กามตฺถา น ทวตฺถา น รตตฺถา, ปชตฺถาว พฺราหฺมณสฺส พฺราหฺมณี โหติ.\csromanspacer{} โส เมถุนํ อุปฺปาเทตฺวา เกสมสฺสุํ โอหาเรตฺวา กาสายานิ วตฺถานิ อจฺฉาเทตฺวา อคารสฺมา อนคาริยํ ปพฺพชติ.\csromanspacer{} โส เอวํ ปพฺพชิโต สมาโน วิวิจฺเจว กาเมหิ ฯเปฯ จตุตฺถํ ฌานํ อุปสมฺปชฺช วิหรติ.\csromanspacer{} โส อิเม จตฺตาโร ฌาเน ภาเวตฺวา กายสฺส เภทา ปรํ มรณา สุคติํ สคฺคํ โลกํ อุปปชฺชติ.\csromanspacer{} เอวํ โข โทณ พฺราหฺมโณ เทวสโม โหติ. (2)}

\prose{กถญฺจ โทณ พฺราหฺมโณ มริยาโท โหติ? อิธ โทณ พฺราหฺมโณ อุภโต สุชาโต โหติ มาติโต จ ปิติโต จ สํสุทฺธคหณิโก ยาว สตฺตมา ปิตามหยุคา อกฺขิตฺโต อนุปกฺกุฏฺโฐ ชาติวาเทน.\csromanspacer{} โส อฏฺฐจตฺตาลีสวสฺสานิ โกมารพฺรหฺมจริยํ จรติ มนฺเต อธียมาโน, อฏฺฐจตฺตาลีสวสฺสานิ โกมารพฺรหฺมจริยํ จริตฺวา มนฺเต อธียิตฺวา อาจริยสฺส อาจริยธนํ ปริเยสติ ธมฺเมเนว โน}

\prose{\csromanfolio{200}อธมฺเมน.\csromanspacer{} ตตฺถ จ โทณ โก ธมฺโม, เนว กสิยา น วณิชฺชาย น โครกฺเขน น อิสฺสตฺเถน น ราชโปริเสน น สิปฺปญฺญตเรน, เกวลํ ภิกฺขาจริยาย กปาลํ อนติมญฺญมาโน.\csromanspacer{} โส อาจริยสฺส อาจริยธนํ นิยฺยาเทตฺวา ทารํ ปริเยสติ ธมฺเมเนว โน อธมฺเมน.}

\prose{ตตฺถ จ โทณ โก ธมฺโม, เนว กเยน น วิกฺกเยน พฺราหฺมณิํเยว อุทกูปสฺสฏฺฐํ.\csromanspacer{} โส พฺราหฺมณิํเยว คจฺฉติ, น ขตฺติยิํ น เวสฺสิํ น สุทฺทิํ น จณฺฑาลิํ น เนสาทิํ น เวนิํ น รถการิํ น ปุกฺกุสิํ คจฺฉติ, น คพฺภินิํ คจฺฉติ, น ปายมานํ คจฺฉติ, น อนุตุนิํ คจฺฉติ.\csromanspacer{} กสฺมา จ โทณ พฺราหฺมโณ น คพฺภินิํ คจฺฉติ.\csromanspacer{} สเจ โทณ พฺราหฺมโณ คพฺภินิํ คจฺฉติ, อติมีฬฺหโช นาม โส โหติ มาณวโก วา มาณวิกา วา, ตสฺมา โทณ พฺราหฺมโณ น คพฺภินิํ คจฺฉติ.\csromanspacer{} กสฺมา จ โทณ พฺราหฺมโณ น ปายมานํ คจฺฉติ.\csromanspacer{} สเจ โทณ พฺราหฺมโณ ปายมานํ คจฺฉติ, อสุจิปฏิปีฬิโต นาม โส โหติ มาณวโก วา มาณวิกา วา, ตสฺมา โทณ พฺราหฺมโณ น ปายมานํ คจฺฉติ.\csromanspacer{} ตสฺส สา โหติ พฺราหฺมณี เนว กามตฺถา น ทวตฺถา น รตตฺถา, ปชตฺถาว พฺราหฺมณสฺส พฺราหฺมณี โหติ.\csromanspacer{} โส เมถุนํ อุปฺปาเทตฺวา ตเมว ปุตฺตสฺสาทํ นิกามยมาโน กุฏุมฺพํ อชฺฌาวสติ, น อคารสฺมา อนคาริยํ ปพฺพชติ.\csromanspacer{} ยาว โปราณานํ พฺราหฺมณานํ มริยาโท, ตตฺถ ติฏฺฐติ, ตํ น วีติกฺกมติ. “ยาว โปราณานํ พฺราหฺมณานํ มริยาโท, ตตฺถ พฺราหฺมโณ ฐิโต ตํ น วีติกฺกมตี”ติ โข โทณ ตสฺมา พฺราหฺมโณ “มริยาโท”ติ วุจฺจติ.\csromanspacer{} เอวํ โข โทณ พฺราหฺมโณ มริยาโท โหติ. (3)}

\prose{กถญฺจ โทณ พฺราหฺมโณ สมฺภินฺนมริยาโท โหติ? อิธ โทณ พฺราหฺมโณ อุภโต สุชาโต โหติ มาติโต จ ปิติโต จ สํสุทฺธคหณิโก ยาว สตฺตมา ปิตามหยุคา อกฺขิตฺโต อนุปกฺกุฏฺโฐ ชาติวาเทน.\csromanspacer{} โส อฏฺฐจตฺตาลีสวสฺสานิ โกมารพฺรหฺมจริยํ จรติ มนฺเต อธียมาโน, อฏฺฐจตฺตาลีสวสฺสานิ โกมารพฺรหฺมจริยํ จริตฺวา มนฺเต อธียิตฺวา อาจริยสฺส อาจริยธนํ ปริเยสติ ธมฺเมเนว โน อธมฺเมน.}

\prose{ตตฺถ จ โทณ โก ธมฺโม, เนว กสิยา น วณิชฺชาย น โครกฺเขน น อิสฺสตฺเถน น ราชโปริเสน น สิปฺปญฺญตเรน, เกวลํ ภิกฺขาจริยาย กปาลํ อนติมญฺญมาโน.\csromanspacer{} โส อาจริยสฺส อาจริยธนํ นิยฺยาเทตฺวา}

\vspace{\tipitakaparskip}
\csromancenter{(20) 5.\csromanspacer{} พฺราหฺมณวคฺค ทารํ ปริเยสติ ธมฺเมนปิ อธมฺเมนปิ กเยนปิ วิกฺกเยนปิ พฺราหฺมณิมฺปิ อุทกูปสฺสฏฺฐํ.\csromanspacer{} โส พฺราหฺมณิมฺปิ คจฺฉติ, ขตฺติยิมฺปิ คจฺฉติ, เวสฺสิมฺปิ คจฺฉติ, สุทฺทิมฺปิ คจฺฉติ, จณฺฑาลิมฺปิ คจฺฉติ, เนสาทิมฺปิ คจฺฉติ, เวนิมฺปิ คจฺฉติ, รถการิมฺปิ คจฺฉติ, ปุกฺกุสิมฺปิ คจฺฉติ, คพฺภินิมฺปิ คจฺฉติ, ปายมานมฺปิ คจฺฉติ, อุตุนิมฺปิ คจฺฉติ, อนุตุนิมฺปิ คจฺฉติ.\csromanspacer{} ตสฺส สา โหติ พฺรหฺมณี กามตฺถาปิ ทวตฺถาปิ รตตฺถาปิ ปชตฺถาปิ พฺราหฺมณสฺส พฺราหฺมณี โหติ.\csromanspacer{} ยาว โปราณานํ พฺราหฺมณานํ มริยาโท, ตตฺถ น ติฏฺฐติ, ตํ วีติกฺกมติ. “ยาว โปราณานํ พฺราหฺมณานํ มริยาโท, ตตฺถ พฺราหฺมโณ น ฐิโต ตํ วีติกฺกมตี”ติ โข โทณ ตสฺมา พฺราหฺมโณ “สมฺภินฺนมริยาโท”ติ วุจฺจติ.\csromanspacer{} เอวํ โข โทณ พฺราหฺมโณ สมฺภินฺนมริยาโท โหติ. (4)}
\prose{\csromanfolio{201}กถญฺจ โทณ พฺราหฺมโณ พฺราหฺมณจณฺฑาโล โหติ? อิธ โทณ พฺราหฺมโณ อุภโต สุชาโต โหติ มาติโต จ ปิติโต จ สํสุทฺธคหณิโก ยาว สตฺตมา ปิตามหยุคา อกฺขิตฺโต อนุปกฺกุฏฺโฐ ชาติวาเทน.\csromanspacer{} โส อฏฺฐจตฺตาลีสวสฺสานิ โกมารพฺรหฺมจริยํ จรติ มนฺเต อธียมาโน, อฏฺฐจตฺตาลีสวสฺสานิ โกมารพฺรหฺมจริยํ จริตฺวา มนฺเต อธียิตฺวา อาจริยสฺส อาจริยธนํ ปริเยสติ ธมฺเมนปิ อธมฺเมนปิ กสิยาปิ วณิชฺชายปิ โครกฺเขนปิ อิสฺสตฺเถนปิ ราชโปริเสนปิ สิปฺปญฺญตเรนปิ, เกวลมฺปิ ภิกฺขาจริยาย กปาลํ อนติมญฺญมาโน.}

\prose{โส อาจริยสฺส อาจริยธนํ นิยฺยาเทตฺวา ทารํ ปริเยสติ ธมฺเมนปิ อธมฺเมนปิ กเยนปิ วิกฺกเยนปิ พฺราหฺมณิมฺปิ อุทกูปสฺสฏฺฐํ.\csromanspacer{} โส พฺราหฺมณิมฺปิ คจฺฉติ, ขตฺติยิมฺปิ คจฺฉติ, เกสฺสิมฺปิ คจฺฉติ, สุทฺทิมฺปิ คจฺฉติ, จณฺฑาลิมฺปิ คจฺฉติ, เนสาทิมฺปิ คจฺฉติ, เวนิมฺปิ คจฺฉติ, รถการิมฺปิ คจฺฉติ, ปุกฺกุสิมฺปิ คจฺฉติ, คพฺภินิมฺปิ คจฺฉติ, ปายมานมฺปิ คจฺฉติ, อุตุนิมฺปิ คจฺฉติ, อนุตุนิมฺปิ คจฺฉติ.\csromanspacer{} ตสฺส สา โหติ พฺราหฺมณี กามตฺถาปิ ทวตฺถาปิ รตตฺถาปิ ปชตฺถาปิ พฺราหฺมณสฺส พฺราหฺมณี โหติ.\csromanspacer{} โส สพฺพกมฺเมหิ ชีวิกํ\footnote{ชีวิตํ (ก)} กปฺเปติ.\csromanspacer{} ตเมนํ พฺราหฺมณา เอวมาหํสุ “กสฺมา ภวํ พฺราหฺมโณ ปฏิชานมาโน สพฺพกมฺเมหิ ชีวิกํ กปฺเปตี”ติ.\csromanspacer{} โส เอวมาห “เสยฺยถาปิ โภ อคฺคิ สุจิมฺปิ ฑหติ, อสุจิมฺปิ ฑหติ, น จ เตน อคฺคิ อุปลิปฺปติ\footnote{อุปลิมฺปติ (ก)}.\csromanspacer{} เอวเมวํ โข โภ สพฺพกมฺเมหิ เจปิ พฺราหฺมโณ ชีวิกํ กปฺเปติ, น จ เตน พฺราหฺมโณ อุปลิปฺปติ. “สพฺพกมฺเมหิ ชีวิกํ กปฺเปตี”ติ โข}

\prose{\csromanfolio{202}โทณ ตสฺมา พฺราหฺมโณ “พฺราหฺมณจณฺฑาโล”ติ วุจฺจติ.\csromanspacer{} เอวํ โข โทณ พฺราหฺมโณ พฺราหฺมณจณฺฑาโล โหติ.}

\prose{เย โข เต โทณ พฺราหฺมณานํ ปุพฺพกา อิสโย มนฺตานํ กตฺตาโร มนฺตานํ ปวตฺตาโร.\csromanspacer{} เยสมิทํ เอตรหิ พฺราหฺมณา โปราณํ มนฺตปทํ คีตํ ปวุตฺตํ สมิหิตํ ตทนุคายนฺติ ตทนุภาสนฺติ ภาสิตมนุภาสนฺติ สชฺฌายิตมนุ สชฺฌายนฺติ วาจิตมนุวาเจนฺติ.\csromanspacer{} เสยฺยถิทํ, อฏฺฐโก วามโก วามเทโว เวสฺสามิตฺโต ยมทคฺคิ องฺคีรโส ภารทฺวาโช วาเสฏฺโฐ กสฺสโป ภคุ, ตฺยสฺสุ’เม ปญฺจ พฺราหฺมเณ ปญฺญาเปนฺติ พฺรหฺมสมํ เทวสมํ มริยาทํ สมฺภินฺนมริยาทํ พฺราหฺมณ จณฺฑาลํเยว ปญฺจมํ, เตสํ ตฺวํ โทณ กตโมติ.}

\prose{เอวํ สนฺเต มยํ โภ โคตม พฺราหฺมณจณฺฑาลมฺปิ น ปูเรม.\csromanspacer{} อภิกฺกนฺตํ โภ โคตม ฯเปฯ อุปาสกํ มํ ภวํ โคตโม ธาเรตุ อชฺชตคฺเค ปาณุเปตํ สรณํ คตนฺติ.\csromanspacer{}\csromanspacer{}ทุติยํ.}

\csromanmatikaanchor{mtk.75}
\csromantocmarkref{subsection}{3. สงฺคารวสุตฺต}{mtk.75}
\bookmark[dest={mtk.75},level=2]{3. สงฺคารวสุตฺต}
\csromanheader{2}{3. สงฺคารวสุตฺต}
\proseitem{193}{อถ โข สงฺคารโว พฺราหฺมโณ เยน ภควา เตนุปสงฺกมิ, อุปสงฺกมิตฺวา ภควตา สทฺธิํ สมฺโมทิ, สมฺโมทนียํ กถํ สารณียํ วีติสาเรตฺวา เอกมนฺตํ นิสีทิ, เอกมนฺตํ นิสินฺโน โข สงฺคารโว พฺราหฺมโณ ภควนฺตํ เอตทโวจ “โก นุ โข โภ โคตม เหตุ โก ปจฺจโย, เยน กทาจิ ทีฆรตฺตํ สชฺฌายกตาปิ มนฺตา นปฺปฏิภนฺติ, ปเคว อสชฺฌายกตา.\csromanspacer{} โก ปน โภ โคตม เหตุ โก ปจฺจโย, เยน กทาจิ ทีฆรตฺตํ อสชฺฌายกตาปิ มนฺตา ปฏิภนฺติ, ปเคว สชฺฌายกตา”ติ.}

\prose{ยสฺมิํ พฺราหฺมณ สมเย กามราคปริยุฏฺฐิเตน เจตสา วิหรติ กามราคปเรเตน, อุปฺปนฺนสฺส จ กามราคสฺส นิสฺสรณํ ยถาภูตํ นปฺปชานาติ.\csromanspacer{} อตฺตตฺถมฺปิ ตสฺมิํ สมเย ยถาภูตํ นปฺปชานาติ น ปสฺสติ, ปรตฺถมฺปิ ตสฺมิํ สมเย ยถาภูตํ นปฺปชานาติ น ปสฺสติ, อุภยตฺถมฺปิ ตสฺมิํ สมเย ยถาภูตํ นปฺปชานาติ น ปสฺสติ, ทีฆรตฺตํ สชฺฌายกตาปิ มนฺตา นปฺปฏิภนฺติ, ปเคว อสชฺฌายกตา.\csromanspacer{} เสยฺยถาปิ พฺราหฺมณ อุทปตฺโต สํสฏฺโฐ ลาขาย วา หลิทฺทิยา วา นีลิยา วา มญฺชิฏฺฐาย วา, ตตฺถ จกฺขุมา ปุริโส สกํ มุขนิมิตฺตํ ปจฺจเวกฺขมาโน ยถาภูตํ}

\vspace{\tipitakaparskip}
\csromancenter{(20) 5.\csromanspacer{} พฺราหฺมณวคฺค นปฺปชาเนยฺย น ปสฺเสยฺย.\csromanspacer{} เอวเมวํ โข พฺราหฺมณ ยสฺมิํ สมเย กามราคปริยุฏฺฐิเตน เจตสา วิหรติ กามราคปเรเตน, อุปฺปนฺนสฺส จ กามราคสฺส นิสฺสรณํ ยถาภูตํ นปฺปชานาติ.\csromanspacer{} อตฺตตฺถมฺปิ ตสฺมิํ สมเย ยถาภูตํ นปฺปชานาติ น ปสฺสติ, ปรตฺถมฺปิ ฯเปฯ.\csromanspacer{} อุภยตฺถมฺปิ ตสฺมิํ สมเย ยถาภูตํ นปฺปชานาติ น ปสฺสติ, ทีฆรตฺตํ สชฺฌายกตาปิ มนฺตา นปฺปฏิภนฺติ, ปเคว อสชฺฌายกตา. (1)}
\prose{\csromanfolio{203}ปุน จปรํ พฺราหฺมณ ยสฺมิํ สมเย พฺยาปาทปริยุฏฺฐิเตน เจตสา วิหรติ พฺยาปาทปเรเตน, อุปฺปนฺนสฺส จ พฺยาปาทสฺส นิสฺสรณํ ยถาภูตํ นปฺปชานาติ.\csromanspacer{} อตฺตตฺถมฺปิ ตสฺมิํ สมเย ยถาภูตํ นปฺปชานาติ น ปสฺสติ, ปรตฺถมฺปิ ฯเปฯ.\csromanspacer{} อุภยตฺถมฺปิ ตสฺมิํ สมเย ยถาภูตํ นปฺปชานาติ น ปสฺสติ, ทีฆรตฺตํ สชฺฌายกตาปิ มนฺตา นปฺปฏิภนฺติ, ปเคว อสชฺฌายกตา.\csromanspacer{} เสยฺยถาปิ พฺราหฺมณ อุทปตฺโต อคฺคินา สนฺตตฺโต อุกฺกุธิโต\footnote{อุกฺกฏฺฐิโต (สี, อิ), อุกฺกุฏฺฐิโต (สฺยา, กํ)} อุสฺสทกชาโต\footnote{อุสุมกชาโต (กตฺถจิ), อุสฺสุรกชาโต (ก), อุสฺมุทกชาโต (ม 3. 107 ปิฏฺเฐ)}, ตตฺถ จกฺขุมา ปุริโส สกํ มุขนิมิตฺตํ ปจฺจเวกฺขมาโน ยถาภูตํ นปฺปชาเนยฺย น ปสฺเสยฺย.\csromanspacer{} เอวเมวํ โข พฺราหฺมณ ยสฺมิํ สมเย พฺยาปาทปริยุฏฺฐิเตน เจตสา วิหรติ พฺยาปาทปเรเตน, อุปฺปนฺนสฺส จ พฺยาปาทสฺส นิสฺสรณํ ยถาภูตํ นปฺปชานาติ.\csromanspacer{} อตฺตตฺถมฺปิ ตสฺมิํ สมเย ยถาภูตํ นปฺปชานาติ น ปสฺสติ, ปรตฺถมฺปิ ฯเปฯ.\csromanspacer{} อุภยตฺถมฺปิ ตสฺมิํ สมเย ยถาภูตํ นปฺปชานาติ น ปสฺสติ, ทีฆรตฺตํ สชฺฌายกตาปิ มนฺตา นปฺปฏิภนฺติ, ปเคว อสชฺฌายกตา. (2)}

\prose{ปุน จปรํ พฺราหฺมณ ยสฺมิํ สมเย ถินมิทฺธปริยุฏฺฐิเตน เจตสา วิหรติ ถินมิทฺธปเรเตน, อุปฺปนฺนสฺส จ ถินมิทฺธสฺส นิสฺสรณํ ยถาภูตํ นปฺปชานาติ.\csromanspacer{} อตฺตตฺถมฺปิ ตสฺมิํ สมเย ยถาภูตํ นปฺปชานาติ น ปสฺสติ, ปรตฺถมฺปิ ฯเปฯ.\csromanspacer{} อุภยตฺถมฺปิ ตสฺมิํ สมเย ยถาภูตํ นปฺปชานาติ น ปสฺสติ, ทีฆรตฺตํ สชฺฌายกตาปิ มนฺตา นปฺปฏิภนฺติ, ปเคว อสชฺฌายกตา.\csromanspacer{} เสยฺยถาปิ พฺราหฺมณ อุทปตฺโต เสวาลปณกปริโยนทฺโธ, ตตฺถ จกฺขุมา ปุริโส สกํ มุขนิมิตฺตํ ปจฺจเวกฺขมาโน ยถาภูตํ นปฺปชาเนยฺย น ปสฺเสยฺย.\csromanspacer{} เอวเมวํ โข พฺราหฺมณ ยสฺมิํ สมเย ถินมิทฺธปริยุฏฺฐิเตน เจตสา วิหรติ ถินมิทฺธปเรเตน, อุปฺปนฺนสฺส จ ถินมิทฺธสฺส นิสฺสรณํ ยถาภูตํ นปฺปชานาติ.\csromanspacer{} อตฺตตฺถมฺปิ \csromanfolio{204}ตสฺมิํ สมเย ยถาภูตํ นปฺปชานาติ น ปสฺสติ, ปรตฺถมฺปิ ฯเปฯ.\csromanspacer{} อุภยตฺถมฺปิ ตสฺมิํ สมเย ยถาภูตํ นปฺปชานาติ น ปสฺสติ, ทีฆรตฺตํ สชฺฌายกตาปิ มนฺตา นปฺปฏิภนฺติ, ปเคว อสชฺฌายกตา. (3)}

\prose{ปุน จปรํ พฺราหฺมณ ยสฺมิํ สมเย อุทฺธจฺจกุกฺกุจฺจปริยุฏฺฐิเตน เจตสา วิหรติ อุทฺธจฺจกุกฺกุจฺจปเรเตน, อุปฺปนฺนสฺส จ อุทฺธจฺจกุกฺกุจฺจสฺส นิสฺสรณํ ยถาภูตํ นปฺปชานาติ.\csromanspacer{} อตฺตตฺถมฺปิ ตสฺมิํ สมเย ยถาภูตํ นปฺปชานาติ น ปสฺสติ, ปรตฺถมฺปิ ฯเปฯ.\csromanspacer{} อุภยตฺถมฺปิ ตสฺมิํ สมเย ยถาภูตํ นปฺปชานาติ น ปสฺสติ, ทีฆรตฺตํ สชฺฌายกตาปิ มนฺตา นปฺปฏิภนฺติ, ปเคว อสชฺฌายกตา.\csromanspacer{} เสยฺยถาปิ พฺราหฺมณ อุทปตฺโต วาเตริโต จลิโต ภนฺโต อูมิชาโต\footnote{อุมฺมิชาโต (อิ)}, ตตฺถ จกฺขุมา ปุริโส สกํ มุขนิมิตฺตํ ปจฺจเวกฺขมาโน ยถาภูตํ นปฺปชาเนยฺย น ปสฺเสยฺย.\csromanspacer{} เอวเมวํ โข พฺราหฺมณ ยสฺมิํ สมเย อุทฺธจฺจกุกฺกุจฺจปริยุฏฺฐิเตน เจตสา วิหรติ อุทฺธจฺจกุกฺกุจฺจปเรเตน, อุปฺปนฺนสฺส จ อุทฺธจฺจกุกฺกุจฺจสฺส นิสฺสรณํ ยถาภูตํ นปฺปชานาติ.\csromanspacer{} อตฺตตฺถมฺปิ ตสฺมิํ สมเย ยถาภูตํ นปฺปชานาติ น ปสฺสติ, ปรตฺถมฺปิ ฯเปฯ.\csromanspacer{} อุภยตฺถมฺปิ ตสฺมิํ สมเย ยถาภูตํ นปฺปชานาติ น ปสฺสติ, ทีฆรตฺตํ สชฺฌายกตาปิ มนฺตา นปฺปฏิภนฺติ, ปเคว อสชฺฌายกตา. (4)}

\prose{ปุน จปรํ พฺราหฺมณ ยสฺมิํ สมเย วิจิกิจฺฉาปริยุฏฺฐิเตน เจตสา วิหรติ วิจิกิจฺฉาปเรเตน, อุปฺปนฺนาย จ วิจิกิจฺฉาย นิสฺสรณํ ยถาภูตํ นปฺปชานาติ.\csromanspacer{} อตฺตตฺถมฺปิ ตสฺมิํ สมเย ยถาภูตํ นปฺปชานาติ น ปสฺสติ, ปรตฺถมฺปิ ฯเปฯ.\csromanspacer{} อุภยตฺถมฺปิ ตสฺมิํ สมเย ยถาภูตํ นปฺปชานาติ น ปสฺสติ, ทีฆรตฺตํ สชฺฌายกตาปิ มนฺตา นปฺปฏิภนฺติ, ปเคว อสชฺฌายกตา.\csromanspacer{} เสยฺยถาปิ พฺราหฺมณ อุทปตฺโต อาวิโล ลุฬิโต กลลีภูโต อนฺธกาเร นิกฺขิตฺโต, ตตฺถ จกฺขุมา ปุริโส สกํ มุขนิมิตฺตํ ปจฺจเวกฺขมาโน ยถาภูตํ นปฺปชาเนยฺย น ปสฺเสยฺย.\csromanspacer{} เอวเมวํ โข พฺราหฺมณ ยสฺมิํ สมเย วิจิกิจฺฉาปริยุฏฺฐิเตน เจตสา วิหรติ วิจิกิจฺฉาปเรเตน, อุปฺปนฺนาย จ วิจิกิจฺฉาย นิสฺสรณํ ยถาภูตํ นปฺปชานาติ.\csromanspacer{} อตฺตตฺถมฺปิ ตสฺมิํ สมเย ยถาภูตํ นปฺปชานาติ น ปสฺสติ, ปรตฺถมฺปิ ฯเปฯ.\csromanspacer{} อุภยตฺถมฺปิ ตสฺมิํ สมเย ยถาภูตํ นปฺปชานาติ น ปสฺสติ, ทีฆรตฺตํ สชฺฌายกตาปิ มนฺตา นปฺปฏิภนฺติ, ปเคว อสชฺฌายกตา. (5)}

\chapterheadpageread{(20) 5. พฺราหฺมณวคฺค}
\prose{\csromanfolio{205}ยสฺมิํ จ โข พฺราหฺมณ สมเย น กามราคปริยุฏฺฐิเตน เจตสา วิหรติ น กามราคปเรเตน, อุปฺปนฺนสฺส จ กามราคสฺส นิสฺสรณํ ยถาภูตํ ปชานาติ.\csromanspacer{} อตฺตตฺถมฺปิ ตสฺมิํ สมเย ยถาภูตํ ปชานาติ ปสฺสติ, ปรตฺถมฺปิ ตสฺมิํ สมเย ยถาภูตํ ปชานาติ ปสฺสติ, อุภยตฺถมฺปิ ตสฺมิํ สมเย ยถาภูตํ ปชานาติ ปสฺสติ, ทีฆรตฺตํ อสชฺฌายกตาปิ มนฺตา ปฏิภนฺติ, ปเคว สชฺฌายกตา.\csromanspacer{} เสยฺยถาปิ พฺราหฺมณ อุทปตฺโต อสํสฏฺโฐ ลาขาย วา หลิทฺทิยา วา นีลิยา วา มญฺชิฏฺฐาย วา, ตตฺถ จกฺขุมา ปุริโส สกํ มุขนิมิตฺตํ ปจฺจเวกฺขมาโน ยถาภูตํ ปชาเนยฺย ปสฺเสยฺย.\csromanspacer{} เอวเมวํ โข พฺราหฺมณ ยสฺมิํ สมเย น กามราคปริยุฏฺฐิเตน เจตสา วิหรติ ฯเปฯ. (1)}

\prose{ปุน จปรํ พฺราหฺมณ ยสฺมิํ สมเย น พฺยาปาทปริยุฏฺฐิเตน เจตสา วิหรติ ฯเปฯ.\csromanspacer{} เสยฺยถาปิ พฺราหฺมณ อุทปตฺโต อคฺคินา อสนฺตตฺโต อนุกฺกุธิโต อนุสฺสทกชาโต, ตตฺถ จกฺขุมา ปุริโส สกํ มุขนิมิตฺตํ ปจฺจเวกฺขมาโน ยถาภูตํ ปชาเนยฺย ปสฺเสยฺย.\csromanspacer{} เอวเมวํ โข พฺราหฺมณ ยสฺมิํ สมเย น พฺยาปาทปริยุฏฺฐิเตน เจตสา วิหรติ ฯเปฯ. (2)}

\prose{ปุน จปรํ พฺราหฺมณ ยสฺมิํ สมเย น ถินมิทฺธปริยุฏฺฐิเตน เจตสา วิหรติ ฯเปฯ.\csromanspacer{} เสยฺยถาปิ พฺราหฺมณ อุทปตฺโต น เสวาลปณกปริโยนทฺโธ, ตตฺถ จกฺขุมา ปุริโส สกํ มุขนิมิตฺตํ ปจฺจเวกฺขมาโน ยถาภูตํ ปชาเนยฺย ปสฺเสยฺย.\csromanspacer{} เอวเมวํ โข พฺราหฺมณ ยสฺมิํ สมเย น ถินมิทฺธปริยุฏฺฐิเตน เจตสา วิหรติ ฯเปฯ. (3)}

\prose{ปุน จปรํ พฺราหฺมณ ยสฺมิํ สมเย น อุทฺธจฺจกุกฺกุจฺจ- ปริยุฏฺฐิเตน เจตสา วิหรติ ฯเปฯ.\csromanspacer{} เสยฺยถาปิ พฺราหฺมณ อุทปตฺโต น วาเตริโต น จลิโต น ภนฺโต น อูมิชาโต, ตตฺถ จกฺขุมา ปุริโส สกํ มุขนิมิตฺตํ ปจฺจเวกฺขมาโน ยถาภูตํ ปชาเนยฺย ปสฺเสยฺย.\csromanspacer{} เอวเมวํ โข พฺราหฺมณ ยสฺมิํ สมเย น อุทฺธจฺจกุกฺกุจฺจ- ปริยุฏฺฐิเตน เจตสา วิหรติ ฯเปฯ. (4)}

\prose{ปุน จปรํ พฺราหฺมณ ยสฺมิํ สมเย น วิจิกิจฺฉาปริยุฏฺฐิเตน เจตสา วิหรติ น วิจิกิจฺฉาปเรเตน, อุปฺปนฺนาย จ วิจิกิจฺฉาย นิสฺสรณํ ยถาภูตํ ปชานาติ.\csromanspacer{} อตฺตตฺถมฺปิ ตสฺมิํ สมเย ยถาภูตํ ปชานาติ ปสฺสติ, ปรตฺถมฺปิ ตสฺมิํ สมเย ยถาภูตํ ปชานาติ ปสฺสติ, อุภยตฺถมฺปิ ตสฺมิํ สมเย ยถาภูตํ ปชานาติ ปสฺสติ, ทีฆรตฺตํ}

\prose{\csromanfolio{206}อสชฺฌายกตาปิ มนฺตา ปฏิภนฺติ, ปเคว สชฺฌายกตา.\csromanspacer{} เสยฺยถาปิ พฺราหฺมณ อุทปตฺโต อจฺโฉ วิปฺปสนฺโน อนาวิโล อาโลเก นิกฺขิตฺโต, ตตฺถ จกฺขุมา ปุริโส สกํ มุข นิมิตฺตํ ปจฺจเวกฺขมาโน ยถาภูตํ ปชาเนยฺย ปสฺเสยฺย.\csromanspacer{} เอวเมวํ โข พฺราหฺมณ ยสฺมิํ สมเย น วิจิกิจฺฉาปริยุฏฺฐิเตน เจตสา วิหรติ น วิจิกิจฺฉาปเรเตน, อุปฺปนฺนาย จ วิจิกิจฺฉาย นิสฺสรณํ ยถาภูตํ ปชานาติ.\csromanspacer{} อตฺตตฺถมฺปิ ตสฺมิํ สมเย ยถาภูตํ ปชานาติ ปสฺสติ, ปรตฺถมฺปิ ฯเปฯ.\csromanspacer{} อุภยตฺถมฺปิ ตสฺมิํ สมเย ยถาภูตํ ปชานาติ ปสฺสติ, ทีฆรตฺตํ อสชฺฌายกตาปิ มนฺตา ปฏิภนฺติ, ปเคว สชฺฌายกตา. (5)}

\prose{อยํ โข พฺราหฺมณ เหตุ อยํ ปจฺจโย, เยน กทาจิ ทีฆรตฺตํ สชฺฌายกตาปิ มนฺตา นปฺปฏิภนฺติ, ปเคว อสชฺฌายกตา.\csromanspacer{} อยํ ปน พฺราหฺมณ เหตุ อยํ ปจฺจโย, เยน กทาจิ ทีฆรตฺตํ อสชฺฌายกตาปิ มนฺตา ปฏิภนฺติ, ปเคว สชฺฌาย กตาติ.}

\prose{อภิกฺกนฺตํ โภ โคตม ฯเปฯ อุปาสกํ มํ ภวํ โคตโม ธาเรตุ อชฺชตคฺเค ปาณุเปตํ สรณํ คตนฺติ.\csromanspacer{}\csromanspacer{}ตติยํ.}

\csromanmatikaanchor{mtk.76}
\csromantocmarkref{subsection}{4. การณปาลีสุตฺต}{mtk.76}
\bookmark[dest={mtk.76},level=2]{4. การณปาลีสุตฺต}
\csromanheader{2}{4. การณปาลีสุตฺต}
\proseitem{194}{เอกํ สมยํ ภควา เวสาลิยํ วิหรติ มหาวเน กูฏาคารสาลายํ.\csromanspacer{} เตน โข ปน สมเยน การณปาลี\footnote{กรณปาลี (ก)} พฺราหฺมโณ ลิจฺฉวีนํ กมฺมนฺตํ กาเรติ.\csromanspacer{} อทฺทสา โข การณปาลี พฺราหฺมโณ ปิงฺคิยานิํ พฺราหฺมณํ ทูรโตว อาคจฺฉนฺตํ, ทิสฺวา ปิงฺคิยานิํ พฺราหฺมณํ เอตทโวจ\csromandash{}}

\prose{หนฺท กุโต นุ ภวํ ปิงฺคิยานี อาคจฺฉติ ทิวา ทิวสฺสาติ.\csromanspacer{} อิโตหํ\footnote{อิธาหํ (สฺยา, กํ), อิโต หิ โข อหํ (ม 1. 232)} โภ อาคจฺฉามิ สมณสฺส โคตมสฺส สนฺติกาติ.\csromanspacer{} ตํ กิํ มญฺญติ ภวํ ปิงฺคิยานี, สมณสฺส โคตมสฺส ปญฺญาเวยฺยตฺติยํ ปณฺฑิโต มญฺเญติ.\csromanspacer{} โก จาหํ โภ, โก จ สมณสฺส โคตมสฺส ปญฺญาเวยฺยตฺติยํ ชานิสฺสามิ, โสปิ นูนสฺส ตาทิโสว, โย สมณสฺส โคตมสฺส ปญฺญาเวยฺยตฺติยํ ชาเนยฺยาติ.\csromanspacer{} อุฬาราย ขลุ ภวํ ปิงฺคิยานี สมณํ โคตมํ ปสํสาย ปสํสตีติ.\csromanspacer{} โก จาหํ โภ, โก จ}

\vspace{\tipitakaparskip}
\csromancenter{(20) 5.\csromanspacer{} พฺราหฺมณวคฺค สมณํ โคตมํ ปสํสิสฺสามิ, ปสตฺถปฺปสตฺโถว\footnote{ปสฏฺฐปสฏฺโฐ จ (สฺยา, กํ, ก)} โส ภวํ โคตโม เสฏฺโฐ เทวมนุสฺสานนฺติ.\csromanspacer{} กิํ ปน ภวํ ปิงฺคิยานี อตฺถวสํ สมฺปสฺสมาโน สมเณ โคตเม เอวํ อภิปฺปสนฺโนติ.}
\prose{\csromanfolio{207}เสยฺยถาปิ โภ ปุริโส อคฺครสปริติตฺโต น อญฺเญสํ หีนานํ รสานํ ปิเหติ.\csromanspacer{} เอวเมวํ โข โภ ยโต ยโต ตสฺส โภโต โคตมสฺส ธมฺมํ สุณาติ ยทิ สุตฺตโส ยทิ เคยฺยโส ยทิ เวยฺยากรณโส ยทิ อพฺภุตธมฺมโส, ตโต ตโต น อญฺเญสํ ปุถุสมณพฺราหฺมณปฺปวาทานํ ปิเหติ.}

\prose{เสยฺยถาปิ โภ ปุริโส ชิฆจฺฉาทุพฺพลฺยปเรโต มธุปิณฺฑิกํ อธิคจฺเฉยฺย, โส ยโต ยโต สาเยถ, ลภเตว\footnote{สาเยยฺย, ลเภเถว (ม 1. 161 ปิฏฺเฐ)} สาทุรสํ อเสจนกํ.\csromanspacer{} เอวเมวํ โข โภ ยโต ยโต ตสฺส โภโต โคตมสฺส ธมฺมํ สุณาติ ยทิ สุตฺตโส ยทิ เคยฺยโส ยทิ เวยฺยากรณโส ยทิ อพฺภุตธมฺมโส, ตโต ตโต ลภเตว อตฺตมนตํ, ลภติ เจตโส ปสาทํ.}

\prose{เสยฺยถาปิ โภ ปุริโส จนฺทนฆฏิกํ อธิคจฺเฉยฺย หริจนฺทนสฺส วา โลหิตจนฺทนสฺส วา, โส ยโต ยโต ฆาเยถ ยทิ มูลโต ยทิ มชฺฌโต ยทิ อคฺคโต, อธิคจฺฉเตว\footnote{อธิคจฺเฉเถว (?)} สุรภิคนฺธํ อเสจนกํ.\csromanspacer{} เอวเมวํ โข โภ ยโต ยโต ตสฺส โภโต โคตมสฺส ธมฺมํ สุณาติ ยทิ สุตฺตโส ยทิ เคยฺยโส ยทิ เวยฺยากรณโส ยทิ อพฺภุตธมฺมโส, ตโต ตโต อธิคจฺฉติ ปาโมชฺชํ, อธิคจฺฉติ โสมนสฺสํ.}

\prose{เสยฺยถาปิ โภ ปุริโส อาพาธิโก ทุกฺขิโต พาฬฺหคิลาโน, ตสฺส กุสโล ภิสกฺโก ฐานโส อาพาธํ นีหเรยฺย.\csromanspacer{} เอวเมวํ โข โภ ยโต ยโต ตสฺส โภโต โคตมสฺส ธมฺมํ สุณาติ ยทิ สุตฺตโส ยทิ เคยฺยโส ยทิ เวยฺยากรณโส ยทิ อพฺภุตธมฺมโส, ตโต ตโต โสกปริเทวทุกฺขโทมนสฺสุปายาสา อพฺภตฺถํ คจฺฉนฺติ.}

\prose{\csromanfolio{208}เสยฺยถาปิ โภ โปกฺขรณี อจฺโฉทกา สาโตทกา สีโตทกา เสตกา สุปติตฺถา รมณียา, อถ ปุริโส อาคจฺเฉยฺย ฆมฺมาภิตตฺโต ฆมฺมปเรโต กิลนฺโต ตสิโต ปิปาสิโต, โส ตํ โปกฺขรณิํ โอคาเหตฺวา นฺหาตฺวา จ ปิวิตฺวา จ สพฺพทรถกิลมถปริฬาหํ ปฏิปฺปสฺสมฺเภยฺย.\csromanspacer{} เอวเมวํ โข โภ ยโต ยโต ตสฺส โภโต โคตมสฺส ธมฺมํ สุณาติ ยทิ สุตฺตโส ยทิ เคยฺยโส ยทิ เวยฺยากรณโส ยทิ อพฺภุตธมฺมโส, ตโต ตโต สพฺพทรถกิลมถปริฬาหา ปฏิปฺปสฺสมฺภนฺตีติ.}

\prose{เอวํ วุตฺเต การณปาลี พฺราหฺมโณ อุฏฺฐายาสนา เอกํสํ อุตฺตราสงฺคํ กริตฺวา ทกฺขิณํ ชาณุมณฺฑลํ ปถวิยํ นิหนฺตฺวา เยน ภควา เตนญฺชลิํ ปณาเมตฺวา ติกฺขตฺตุํ อุทานํ อุทาเนสิ\csromandash{}}

\csromanheader{2}{“นโม ตสฺส ภควโต อรหโต สมฺมาสมฺพุทฺธสฺส.}
\namakkaram{นโม ตสฺส ภควโต อรหโต สมฺมาสมฺพุทฺธสฺส.}
\namakkaram{นโม ตสฺส ภควโต อรหโต สมฺมาสมฺพุทฺธสฺสา”ติ.}
\prose{อภิกฺกนฺตํ โภ ปิงฺคิยานิ อภิกฺกนฺตํ โภ ปิงฺคิยานิ, เสยฺยถาปิ โภ ปิงฺคิยานิ นิกฺกุชฺชิตํ\footnote{นิกุชฺชิตํ (ก)} วา อุกฺกุชฺเชยฺย, ปฏิจฺฉนฺนํ วา วิวเรยฺย, มูฬฺหสฺส วา มคฺคํ อาจิกฺเขยฺย, อนฺธกาเร วา เตลปชฺโชตํ ธาเรยฺย “จกฺขุมนฺโต รูปานิ ทกฺขนฺตี”ติ.\csromanspacer{} เอวเมวํ โภตา ปิงฺคิยานินา อเนกปริยาเยน ธมฺโม ปกาสิโต, เอสาหํ โภ ปิงฺคิยานิ ตํ ภวนฺตํ โคตมํ สรณํ คจฺฉามิ ธมฺมญฺจ ภิกฺขุสํฆญฺจ, อุปาสกํ มํ ภวํ ปิงฺคิยานี ธาเรตุ อชฺชตคฺเค ปาณุเปตํ สรณํ คตนฺติ.\csromanspacer{}\csromanspacer{}จตุตฺถํ.}

\csromanmatikaanchor{mtk.77}
\csromantocmarkref{subsection}{5. ปิงฺคิยานีสุตฺต}{mtk.77}
\bookmark[dest={mtk.77},level=2]{5. ปิงฺคิยานีสุตฺต}
\csromanheader{2}{5. ปิงฺคิยานีสุตฺต}
\proseitem{195}{เอกํ สมยํ ภควา เวสาลิยํ วิหรติ มหาวเน กูฏาคารสาลายํ.\csromanspacer{} เตน โข ปน สมเยน ปญฺจมตฺตานิ ลิจฺฉวิสตานิ ภควนฺตํ ปยิรุปาสนฺติ.\csromanspacer{} อปฺเปกจฺเจ ลิจฺฉวี นีลา โหนฺติ นีลวณฺณา นีลวตฺถา นีลาลงฺการา, อปฺเปกจฺเจ ลิจฺฉวี ปีตา โหนฺติ ปีตวณฺณา ปีตวตฺถา ปีตาลงฺการา, อปฺเปกจฺเจ ลิจฺฉวี โลหิตกา โหนฺติ โลหิตกวณฺณา โลหิตกวตฺถา โลหิตกาลงฺการา, อปฺเปกจฺเจ ลิจฺฉวี}

\vspace{\tipitakaparskip}
\csromancenter{(20) 5.\csromanspacer{} พฺราหฺมณวคฺค โอทาตา โหนฺติ โอทาตวณฺณา โอทาตวตฺถา โอทาตาลงฺการา, ตฺยสฺสุทํ ภควา อติโรจติ วณฺเณน เจว ยสสา จ.}
\prose{\csromanfolio{209}อถ โข ปิงฺคิยานี พฺราหฺมโณ อุฏฺฐายาสนา เอกํสํ อุตฺตราสงฺคํ กริตฺวา เยน ภควา เตนญฺชลิํ ปณาเมตฺวา ภควนฺตํ เอตทโวจ “ปฏิภาติ มํ ภควา, ปฏิภาติ มํ สุคตา”ติ. “ปฏิภาตุ ตํ ปิงฺคิยานี”ติ ภควา อโวจ.\csromanspacer{} อถ ปิงฺคิยานี พฺราหฺมโณ ภควโต สมฺมุขา สารุปฺปาย คาถาย อภิตฺถวิ\csromandash{}}

\csromangathameasure{“ปทฺมํ ยถา โกกนทํ สุคนฺธํ, \\ ปาโต สิยา ผุลฺลมวีตคนฺธํ. \\ องฺคีรสํ ปสฺส วิโรจมานํ, \\ ตปนฺตมาทิจฺจมิวนฺตลิกฺเข”ติ.}
\csromangathagroup{\csromangathastanza{“ปทฺมํ\footnote{ปทุมํ (ก) สํ 1. 81 ปิฏฺเฐปิ.} ยถา โกกนทํ\footnote{โกกนุทํ (สฺยา, กํ)} สุคนฺธํ, \\ ปาโต สิยา ผุลฺลมวีตคนฺธํ. \\ องฺคีรสํ ปสฺส วิโรจมานํ, \\ ตปนฺตมาทิจฺจมิวนฺตลิกฺเข”ติ.}}

\prose{อถ โข เต ลิจฺฉวี ปญฺจหิ อุตฺตราสงฺคสเตหิ ปิงฺคิยานิํ พฺราหฺมณํ อจฺฉาเทสุํ.\csromanspacer{} อถ โข ปิงฺคิยานี พฺราหฺมโณ เตหิ ปญฺจหิ อุตฺตราสงฺคสเตหิ ภควนฺตํ อจฺฉาเทสิ.}

\prose{อถ โข ภควา เต ลิจฺฉวี เอตทโวจ “ปญฺจนฺนํ ลิจฺฉวี รตนานํ ปาตุภาโว ทุลฺลโภ โลกสฺมิํ.\csromanspacer{} กตเมสํ ปญฺจนฺนํ? ตถาคตสฺส อรหโต สมฺมาสมฺพุทฺธสฺส ปาตุภาโว ทุลฺลโภ โลกสฺมิํ, ตถาคตปฺปเวทิตสฺส ธมฺมวินยสฺส เทเสตา ปุคฺคโล ทุลฺลโภ โลกสฺมิํ, ตถาคตปฺปเวทิตสฺส ธมฺมวินยสฺส เทสิตสฺส วิญฺญาตา ปุคฺคโล ทุลฺลโภ โลกสฺมิํ, ตถาคตปฺปเวทิตสฺส ธมฺมวินยสฺส เทสิตสฺส วิญฺญาตา ธมฺมานุธมฺมปฺปฏิปนฺโน ปุคฺคโล ทุลฺลโภ โลกสฺมิํ, กตญฺญู กตเวที ปุคฺคโล ทุลฺลโภ โลกสฺมิํ.\csromanspacer{} อิเมสํ โข ลิจฺฉวี ปญฺจนฺนํ รตนานํ ปาตุภาโว ทุลฺลโภ โลกสฺมินฺ”ติ\footnote{เหฏฺฐา 149 ปิฏฺเฐปิ.}.\csromanspacer{}\csromanspacer{}ปญฺจมํ.}

\csromanmatikaanchor{mtk.78}
\csromantocmarkref{subsection}{6. มหาสุปินสุตฺต}{mtk.78}
\bookmark[dest={mtk.78},level=2]{6. มหาสุปินสุตฺต}
\csromanheader{2}{6. มหาสุปินสุตฺต}
\proseitem{196}{ตถาคตสฺส ภิกฺขเว อรหโต สมฺมาสมฺพุทฺธสฺส ปุพฺเพว สมฺโพธา อนภิสมฺพุทฺธสฺส โพธิสตฺตสฺเสว สโต ปญฺจ มหาสุปินา ปาตุรเหสุํ.\csromanspacer{} กตเม ปญฺจ? ตถาคตสฺส ภิกฺขเว อรหโต สมฺมาสมฺพุทฺธสฺส ปุพฺเพว สมฺโพธา อนภิสมฺพุทฺธสฺส โพธิสตฺตสฺเสว \csromanfolio{210}สโต อยํ มหาปถวี มหาสยนํ อโหสิ, หิมวา ปพฺพตราชา พิพฺโพหนํ\footnote{พิมฺโพหนํ (สี, สฺยา, กํ, อิ), พิมฺพ + โอหนํ อิติ ปทวิภาโค.} อโหสิ, ปุรตฺถิเม สมุทฺเท วาโม หตฺโถ โอหิโต อโหสิ, ปจฺฉิเม สมุทฺเท ทกฺขิโณ หตฺโถ โอหิโต อโหสิ, ทกฺขิเณ สมุทฺเท อุโภ ปาทา โอหิตา อเหสุํ.\csromanspacer{} ตถาคตสฺส ภิกฺขเว อรหโต สมฺมาสมฺพุทฺธสฺส ปุพฺเพว สมฺโพธา อนภิสมฺพุทฺธสฺส โพธิสตฺตสฺเสว สโต อยํ ปฐโม มหาสุปิโน ปาตุรโหสิ.}

\prose{ปุน จปรํ ภิกฺขเว ตถาคตสฺส อรหโต สมฺมาสมฺพุทฺธสฺส ปุพฺเพว สมฺโพธา อนภิสมฺพุทฺธสฺส โพธิสตฺตสฺเสว สโต ติริยา นาม ติณชาติ นาภิยา อุคฺคนฺตฺวา นภํ อาหจฺจ ฐิตา อโหสิ.\csromanspacer{} ตถาคตสฺส ภิกฺขเว อรหโต สมฺมาสมฺพุทฺธสฺส ปุพฺเพว สมฺโพธา อนภิสมฺพุทฺธสฺส โพธิสตฺตสฺเสว สโต อยํ ทุติโย มหาสุปิโน ปาตุรโหสิ.}

\prose{ปุน จปรํ ภิกฺขเว ตถาคตสฺส อรหโต สมฺมาสมฺพุทฺธสฺส ปุพฺเพว สมฺโพธา อนภิสมฺพุทฺธสฺส โพธิสตฺตสฺเสว สโต เสตา กิมี กณฺหสีสา ปาเทหิ อุสฺสกฺกิตฺวา ( )\footnote{(อคฺคนขโต) กตฺถจิ ทิสฺสติ.} ยาว ชาณุมณฺฑลา ปฏิจฺฉาเทสุํ.\csromanspacer{} ตถาคตสฺส ภิกฺขเว อรหโต สมฺมาสมฺพุทฺธสฺส ปุพฺเพว สมฺโพธา อนภิสมฺพุทฺธสฺส โพธิสตฺตสฺเสว สโต อยํ ตติโย มหาสุปิโน ปาตุรโหสิ.}

\prose{ปุน จปรํ ภิกฺขเว ตถาคตสฺส อรหโต สมฺมาสมฺพุทฺธสฺส ปุพฺเพว สมฺโพธา อนภิสมฺพุทฺธสฺส โพธิสตฺตสฺเสว สโต จตฺตาโร สกุณา นานาวณฺณา จตูหิ ทิสาหิ อาคนฺตฺวา ปาทมูเล นิปติตฺวา สพฺพเสตา สมฺปชฺชิํสุ.\csromanspacer{} ตถาคตสฺส ภิกฺขเว อรหโต สมฺมาสมฺพุทฺธสฺส ปุพฺเพว สมฺโพธา อนภิสมฺพุทฺธสฺส โพธิสตฺตสฺเสว สโต อยํ จตุตฺโถ มหาสุปิโน ปาตุรโหสิ.}

\prose{ปุน จปรํ ภิกฺขเว ตถาคโต อรหํ สมฺมาสมฺพุทฺโธ ปุพฺเพว สมฺโพธา อนภิสมฺพุทฺโธ โพธิสตฺโตว สมาโน มหโต มีฬฺหปพฺพตสฺส อุปรูปริ จงฺกมติ อลิปฺปมาโน มีเฬฺหน.\csromanspacer{} ตถาคตสฺส ภิกฺขเว อรหโต สมฺมาสมฺพุทฺธสฺส ปุพฺเพว สมฺโพธา อนภิสมฺพุทฺธสฺส โพธิสตฺตสฺเสว สโต อยํ ปญฺจโม มหาสุปิโน ปาตุรโหสิ.}

\chapterheadpageread{(20) 5. พฺราหฺมณวคฺค}
\prose{\csromanfolio{211}ยมฺปิ ภิกฺขเว ตถาคตสฺส อรหโต สมฺมาสมฺพุทฺธสฺส ปุพฺเพว สมฺโพธา อนภิสมฺพุทฺธสฺส โพธิสตฺตสฺเสว สโต อยํ มหาปถวี มหาสยนํ อโหสิ, หิมวา ปพฺพตราชา พิพฺโพหนํ อโหสิ, ปุรตฺถิเม สมุทฺเท วาโม หตฺโถ โอหิโต อโหสิ, ปจฺฉิเม สมุทฺเท ทกฺขิโณ หตฺโถ โอหิโต อโหสิ, ทกฺขิเณ สมุทฺเท อุโภ ปาทา โอหิตา อเหสุํ.\csromanspacer{} ตถาคเตน ภิกฺขเว อรหตา สมฺมาสมฺพุทฺเธน อนุตฺตรา สมฺมาสมฺโพธิ อภิสมฺพุทฺธา, ตสฺสา อภิสมฺโพธาย อยํ ปฐโม มหาสุปิโน ปาตุรโหสิ.}

\prose{ยมฺปิ ภิกฺขเว ตถาคตสฺส อรหโต สมฺมาสมฺพุทฺธสฺส ปุพฺเพว สมฺโพธา อนภิสมฺพุทฺธสฺส โพธิสตฺตสฺเสว สโต ติริยา นาม ติณชาติ นาภิยา อุคฺคนฺตฺวา นภํ อาหจฺจ ฐิตา อโหสิ.\csromanspacer{} ตถาคเตน ภิกฺขเว อรหตา สมฺมาสมฺพุทฺเธน อริโย อฏฺฐงฺคิโก มคฺโค อภิสมฺพุชฺฌิตฺวา ยาว เทวมนุสฺเสหิ สุปฺปกาสิโต, ตสฺส อภิสมฺโพธาย อยํ ทุติโย มหาสุปิโน ปาตุรโหสิ.}

\prose{ยมฺปิ ภิกฺขเว ตถาคตสฺส อรหโต สมฺมาสมฺพุทฺธสฺส ปุพฺเพว สมฺโพธา อนภิสมฺพุทฺธสฺส โพธิสตฺตสฺเสว สโต เสตา กิมี กณฺหสีสา ปาเทหิ อุสฺสกฺกิตฺวา ยาว ชาณุมณฺฑลา ปฏิจฺฉาเทสุํ.\csromanspacer{} พหู ภิกฺขเว คิหี โอทาตวสนา ตถาคตํ ปาณุเปตา\footnote{ปาณุเปตํ (สี, สฺยา, กํ, อิ)} สรณํ คตา.\csromanspacer{} ตสฺส อภิสมฺโพธาย อยํ ตติโย มหาสุปิโน ปาตรโหสิ.}

\prose{ยมฺปิ ภิกฺขเว ตถาคตสฺส อรหโต สมฺมาสมฺพุทฺธสฺส ปุพฺเพว สมฺโพธา อนภิสมฺพุทฺธสฺส โพธิสตฺตสฺเสว สโต จตฺตาโร สกุณา นานาวณฺณา จตูหิ ทิสาหิ อาคนฺตฺวา ปาทมูเล นิปติตฺวา สพฺพเสตา สมฺปชฺชิํสุ.\csromanspacer{} จตฺตาโรเม ภิกฺขเว วณฺณา ขตฺติยา พฺราหฺมณา เวสฺสา สุทฺทา, เต ตถาคตปฺปเวทิเต ธมฺมวินเย อคารสฺมา อนคาริยํ ปพฺพชิตฺวา อนุตฺตรํ วิมุตฺติํ สจฺฉิกโรนฺติ.\csromanspacer{} ตสฺส อภิสมฺโพธาย อยํ จตุตฺโถ มหาสุปิโน ปาตุรโหสิ.}

\prose{ยมฺปิ ภิกฺขเว ตถาคโต อรหํ สมฺมาสมฺพุทฺโธ ปุพฺเพว สมฺโพธา อนภิสมฺพุทฺโธ โพธิสตฺโตว สมาโน มหโต มีฬฺหปพฺพตสฺส อุปรูปริ}

\prose{\csromanfolio{212}จงฺกมติ อลิปฺปมาโน มีเฬฺหน.\csromanspacer{} ลาภี ภิกฺขเว ตถาคโต จีวรปิณฺฑปาตเสนาสน คิลานปฺปจฺจยเภสชฺชปริกฺขารานํ, ตํ\footnote{ตตฺถ จ (สี, สฺยา, กํ, อิ)} ตถาคโต อคถิโต\footnote{อคธิโต (สฺยา, อิ, ก)} อมุจฺฉิโต อนชฺโฌสนฺโน\footnote{อนชฺฌาปนฺโน (ก) อนชฺโฌปนฺโน (สี, สฺยา)} อาทีนวทสฺสาวี นิสฺสรณปญฺโญ ปริภุญฺชติ.\csromanspacer{} ตสฺส อภิสมฺโพธาย อยํ ปญฺจโม มหาสุปิโน ปาตุรโหสิ.}

\prose{ตถาคตสฺส ภิกฺขเว อรหโต สมฺมาสมฺพุทฺธสฺส ปุพฺเพว สมฺโพธา อนภิสมฺพุทฺธสฺส โพธิสตฺตสฺเสว สโต อิเม ปญฺจ มหาสุปินา ปาตุรเหสุนฺติ.\csromanspacer{}\csromanspacer{}ฉฏฺฐํ.}

\csromanmatikaanchor{mtk.79}
\csromantocmarkref{subsection}{7. วสฺสสุตฺต}{mtk.79}
\bookmark[dest={mtk.79},level=2]{7. วสฺสสุตฺต}
\csromanheader{2}{7. วสฺสสุตฺต}
\proseitem{197}{ปญฺจิเม ภิกฺขเว วสฺสสฺส อนฺตรายา, ยํ เนมิตฺตา\footnote{เนมิตฺตกา (กตฺถจิ)} น ชานนฺติ, ยตฺถ เนมิตฺตานํ จกฺขุ น กมติ.\csromanspacer{} กตเม ปญฺจ? อุปริ ภิกฺขเว อากาเส เตโชธาตุ ปกุปฺปติ, เตน อุปฺปนฺนา เมฆา ปฏิวิคจฺฉนฺติ.\csromanspacer{} อยํ ภิกฺขเว ปฐโม วสฺสสฺส อนฺตราโย.\csromanspacer{} ยํ เนมิตฺตา น ชานนฺติ, ยตฺถ เนมิตฺตานํ จกฺขุ น กมติ.}

\prose{ปุน จปรํ ภิกฺขเว อุปริ อากาเส วาโยธาตุ ปกุปฺปติ, เตน อุปฺปนฺนา เมฆา ปฏิวิคจฺฉนฺติ.\csromanspacer{} อยํ ภิกฺขเว ทุติโย วสฺสสฺส อนฺตราโย.\csromanspacer{} ยํ เนมิตฺตา น ชานนฺติ, ยตฺถ เนมิตฺตานํ จกฺขุ น กมติ.}

\prose{ปุน จปรํ ภิกฺขเว ราหุ อสุรินฺโท ปาณินา อุทกํ สมฺปฏิจฺฉิตฺวา มหาสมุทฺเท ฉฑฺเฑติ.\csromanspacer{} อยํ ภิกฺขเว ตติโย วสฺสสฺส อนฺตราโย.\csromanspacer{} ยํ เนมิตฺตา น ชานนฺติ, ยตฺถ เนมิตฺตานํ จกฺขุ น กมติ.}

\prose{ปุน จปรํ ภิกฺขเว วสฺสวลาหกา เทวา ปมตฺตา โหนฺติ.\csromanspacer{} อยํ ภิกฺขเว จตุตฺโถ วสฺสสฺส อนฺตราโย.\csromanspacer{} ยํ เนมิตฺตา น ชานนฺติ, ยตฺถ เนมิตฺตานํ จกฺขุ น กมติ.}

\prose{ปุน จปรํ ภิกฺขเว มนุสฺสา อธมฺมิกา โหนฺติ.\csromanspacer{} อยํ ภิกฺขเว ปญฺจโม วสฺสสฺส อนฺตราโย.\csromanspacer{} ยํ เนมิตฺตา น ชานนฺติ, ยตฺถ เนมิตฺตานํ จกฺขุ น กมติ.\csromanspacer{} อิเม โข ภิกฺขเว ปญฺจ วสฺสสฺส อนฺตรายา, ยํ เนมิตฺตา น ชานนฺติ, ยตฺถ เนมิตฺตานํ จกฺขุ น กมตีติ.\csromanspacer{}\csromanspacer{}สตฺตมํ.}

\csromanmatikaanchor{mtk.80}
\csromantocmarkref{subsection}{8. วาจาสุตฺต}{mtk.80}
\bookmark[dest={mtk.80},level=2]{8. วาจาสุตฺต}
\csromanheader{2}{(20) 5. พฺราหฺมณวคฺค \textbf{8. วาจาสุตฺต}}
\proseitem{198}{\csromanfolio{213}ปญฺจหิ ภิกฺขเว องฺเคหิ สมนฺนาคตา วาจา สุภาสิตา โหติ โน ทุพฺภาสิตา, อนวชฺชา จ อนนุวชฺชา จ วิญฺญูนํ\footnote{อนนุวชฺชา วิญฺญูนํ (สฺยา, กํ)}.\csromanspacer{} กตเมหิ ปญฺจหิ? กาเลน จ ภาสิตา โหติ, สจฺจา จ ภาสิตา โหติ, สณฺหา จ ภาสิตา โหติ, อตฺถสํหิตา จ ภาสิตา โหติ, เมตฺตจิตฺเตน จ ภาสิตา โหติ.\csromanspacer{} อิเมหิ โข ภิกฺขเว ปญฺจหิ องฺเคหิ สมนฺนาคตา วาจา สุภาสิตา โหติ โน ทุพฺภาสิตา, อนวชฺชา จ อนนุวชฺชา จ วิญฺญูนนฺติ.\csromanspacer{}\csromanspacer{}อฏฺฐมํ.}

\csromanmatikaanchor{mtk.81}
\csromantocmarkref{subsection}{9. กุลสุตฺต}{mtk.81}
\bookmark[dest={mtk.81},level=2]{9. กุลสุตฺต}
\csromanheader{2}{9. กุลสุตฺต}
\proseitem{199}{ยํ ภิกฺขเว สีลวนฺโต ปพฺพชิตา กุลํ อุปสงฺกมนฺติ, ตตฺถ มนุสฺสา ปญฺจหิ ฐาเนหิ พหุํ ปุญฺญํ ปสวนฺติ.\csromanspacer{} กตเมหิ ปญฺจหิ? ยสฺมิํ ภิกฺขเว สมเย สีลวนฺเต ปพฺพชิเต กุลํ อุปสงฺกมนฺเต มนุสฺสา ทิสฺวา จิตฺตานิ ปสาเทนฺติ\footnote{ปสีทนฺติ (สฺยา, กํ, ก)}, สคฺคสํวตฺตนิกํ ภิกฺขเว ตํ กุลํ ตสฺมิํ สมเย ปฏิปทํ ปฏิปนฺนํ โหติ.}

\prose{ยสฺมิํ ภิกฺขเว สมเย สีลวนฺเต ปพฺพชิเต กุลํ อุปสงฺกมนฺเต มนุสฺสา ปจฺจุฏฺเฐนฺติ อภิวาเทนฺติ อาสนํ เทนฺติ, อุจฺจากุลีนสํวตฺตนิกํ ภิกฺขเว ตํ กุลํ ตสฺมิํ สมเย ปฏิปทํ ปฏิปนฺนํ โหติ.}

\prose{ยสฺมิํ ภิกฺขเว สมเย สีลวนฺเต ปพฺพชิเต กุลํ อุปสงฺกมนฺเต มนุสฺสา มจฺเฉรมลํ ปฏิวิเนนฺติ\footnote{ปฏิวิโนเทนฺติ (สี, อิ)}, มเหสกฺขสํวตฺตนิกํ ภิกฺขเว ตํ กุลํ ตสฺมิํ สมเย ปฏิปทํ ปฏิปนฺนํ โหติ.}

\prose{ยสฺมิํ ภิกฺขเว สมเย สีลวนฺเต ปพฺพชิเต กุลํ อุปสงฺกมนฺเต มนุสฺสา ยถาสตฺติ ยถาพลํ สํวิภชนฺติ, มหาโภคสํวตฺตนิกํ ภิกฺขเว ตํ กุลํ ตสฺมิํ สมเย ปฏิปทํ ปฏิปนฺนํ โหติ.}

\prose{ยสฺมิํ ภิกฺขเว สมเย สีลวนฺเต ปพฺพชิเต กุลํ อุปสงฺกมนฺเต มนุสฺสา ปริปุจฺฉนฺติ ปริปญฺหนฺติ ธมฺมํ สุณนฺติ, มหาปญฺญาสํวตฺตนิกํ ภิกฺขเว ตํ กุลํ ตสฺมิํ สมเย ปฏิปทํ ปฏิปนฺนํ โหติ.\csromanspacer{} ยํ ภิกฺขเว สีลวนฺโต ปพฺพชิตา กุลํ อุปสงฺกมนฺติ, ตตฺถ มนุสฺสา อิเมหิ ปญฺจหิ ฐาเนหิ พหุํ ปุญฺญํ ปสวนฺตีติ.\csromanspacer{}\csromanspacer{}นวมํ.}

\csromanmatikaanchor{mtk.82}
\csromantocmarkref{subsection}{10. นิสฺสารณียสุตฺต}{mtk.82}
\bookmark[dest={mtk.82},level=2]{10. นิสฺสารณียสุตฺต}
\csromanheader{2}{10. นิสฺสารณียสุตฺต}
\proseitem{200}{\csromanfolio{214}ปญฺจิมา ภิกฺขเว นิสฺสารณียา\footnote{นิสฺสรณียา (อิ), นิสฺสรณิยา (ที 3. 199)} ธาตุโย.\csromanspacer{} กตมา ปญฺจ? อิธ ภิกฺขเว ภิกฺขุโน กามํ\footnote{กาเม (สฺยา, กํ) ที 3. 199 ปิฏฺเฐปิ.} มนสิกโรโต กาเมสุ จิตฺตํ น ปกฺขนฺทติ นปฺปสีทติ น สนฺติฏฺฐติ น วิมุจฺจติ, เนกฺขมฺมํ โข ปนสฺส มนสิกโรโต เนกฺขมฺเม จิตฺตํ ปกฺขนฺทติ ปสีทติ สนฺติฏฺฐติ วิมุจฺจติ, ตสฺส ตํ จิตฺตํ สุคตํ\footnote{สุกตํ (อิ, ก)} สุภาวิตํ สุวุฏฺฐิตํ สุวิมุตฺตํ สุวิสํยุตฺตํ\footnote{วิสํยุตฺตํ (กตฺถจิ, ที 3. 199 ปิฏฺเฐปิ)} กาเมหิ, เย จ กามปจฺจยา อุปฺปชฺชนฺติ อาสวา วิฆาตปริฬาหา, มุตฺโต โส เตหิ, น โส ตํ เวทนํ เวทิยติ, อิทมกฺขาตํ กามานํ นิสฺสรณํ. (1)}

\prose{ปุน จปรํ ภิกฺขเว ภิกฺขุโน พฺยาปาทํ มนสิกโรโต พฺยาปาเท จิตฺตํ น ปกฺขนฺทติ นปฺปสีทติ น สนฺติฏฺฐติ น วิมุจฺจติ, อพฺยาปาทํ โข ปนสฺส มนสิกโรโต อพฺยาปาเท จิตฺตํ ปกฺขนฺทติ ปสีทติ สนฺติฏฺฐติ วิมุจฺจติ, ตสฺส ตํ จิตฺตํ สุคตํ สุภาวิตํ สุวุฏฺฐิตํ สุวิมุตฺตํ สุวิสํยุตฺตํ พฺยาปาเทน, เย จ พฺยาปาทปจฺจยา อุปฺปชฺชนฺติ อาสวา วิฆาตปริฬาหา, มุตฺโต โส เตหิ, น โส ตํ เวทนํ เวทิยติ, อิทมกฺขาตํ พฺยาปาทสฺส นิสฺสรณํ. (2)}

\prose{ปุน จปรํ ภิกฺขเว ภิกฺขุโน วิเหสํ มนสิกโรโต วิเหสาย จิตฺตํ น ปกฺขนฺทติ นปฺปสีทติ น สนฺติฏฺฐติ น วิมุจฺจติ, อวิเหสํ โข ปนสฺส มนสิกโรโต อวิเหสาย จิตฺตํ ปกฺขนฺทติ ปสีทติ สนฺติฏฺฐติ วิมุจฺจติ, ตสฺส ตํ จิตฺตํ สุคตํ สุภาวิตํ สุวุฏฺฐิตํ สุวิมุตฺตํ สุวิสํยุตฺตํ วิเหสาย, เย จ วิเหสาปจฺจยา อุปฺปชฺชนฺติ อาสวา วิฆาตปริฬาหา, มุตฺโต โส เตหิ, น โส ตํ เวทนํ เวทิยติ, อิทมกฺขาตํ วิเหสาย นิสฺสรณํ. (3)}

\prose{ปุน จปรํ ภิกฺขเว ภิกฺขุโน รูปํ มนสิกโรโต รูเป จิตฺตํ น ปกฺขนฺทติ นปฺปสีทติ น สนฺติฏฺฐติ น วิมุจฺจติ, อรูปํ โข ปนสฺส มนสิกโรโต อรูเป จิตฺตํ ปกฺขนฺทติ ปสีทติ สนฺติฏฺฐติ วิมุจฺจติ, ตสฺส ตํ จิตฺตํ สุคตํ สุภาวิตํ สุวุฏฺฐิตํ สุวิมุตฺตํ สุวิสํยุตฺตํ รูเปหิ, เย จ รูปปจฺจยา อุปฺปชฺชนฺติ อาสวา วิฆาตปริฬาหา, มุตฺโต โส เตหิ, น โส ตํ เวทนํ เวทิยติ, อิทมกฺขาตํ รูปานํ นิสฺสรณํ. (4)}

\chapterheadpageread{(20) 5. พฺราหฺมณวคฺค}
\prose{\csromanfolio{215}ปุน จปรํ ภิกฺขเว ภิกฺขุโน สกฺกายํ มนสิกโรโต สกฺกาเย จิตฺตํ น ปกฺขนฺทติ นปฺปสีทติ น สนฺติฏฺฐติ น วิมุจฺจติ, สกฺกายนิโรธํ โข ปนสฺส มนสิกโรโต สกฺกายนิโรเธ จิตฺตํ ปกฺขนฺทติ ปสีทติ สนฺติฏฺฐติ วิมุจฺจติ.\csromanspacer{} ตสฺส ตํ จิตฺตํ สุคตํ สุภาวิตํ สุวุฏฺฐิตํ สุวิมุตฺตํ สุวิสํยุตฺตํ สกฺกาเยน, เย จ สกฺกายปจฺจยา อุปฺปชฺชนฺติ อาสวา วิฆาตปริฬาหา, มุตฺโต โส เตหิ, น โส ตํ เวทนํ เวทิยติ, อิทมกฺขาตํ สกฺกายสฺส นิสฺสรณํ. (5)}

\prose{ตสฺส กามนนฺทีปิ นานุเสติ, พฺยาปาทนนฺทีปิ นานุเสติ, วิเหสานนฺทีปิ นานุเสติ, รูปนนฺทีปิ นานุเสติ, สกฺกายนนฺทีปิ นานุเสติ. (โส)\footnote{( ) กตฺถจิ นตฺถิ.} กามนนฺทิยาปิ อนนุสยา พฺยาปาทนนฺทิยาปิ อนนุสยา วิเหสานนฺทิยาปิ อนนุสยา รูปนนฺทิยาปิ อนนุสยา สกฺกายนนฺทิยาปิ อนนุสยา อยํ วุจฺจติ ภิกฺขเว ภิกฺขุ นิรนุสโย, อจฺเฉจฺฉิ\footnote{อจฺเฉชฺชิ (สฺยา, กํ, ก)} ตณฺหํ, วิวตฺตยิ\footnote{วาวตฺตยิ (สี)} สํโยชนํ, สมฺมา มานาภิสมยา อนฺตมกาสิ ทุกฺขสฺส.\csromanspacer{} อิมา โข ภิกฺขเว ปญฺจ นิสฺสารณียา ธาตุโยติ.\csromanspacer{}\csromanspacer{}ทสมํ.\csromanspacer{} พฺราหฺมณวคฺโค ปญฺจโม.}

\titlehead{\textbf{ตสฺสุทฺทานํ}}
\csromangathameasure{โสโณ โทโณ สงฺคารโว, การณปาลี จ ปิงฺคิยานี. \\ สุปินา จ วสฺสา วาจา, กุลํ นิสฺสารณีเยน จาติ. \\ \textbf{จตุตฺโถ ปณฺณาสโก สมตฺโต.}}
\csromangathagroup{\csromangathastanza{โสโณ โทโณ สงฺคารโว, การณปาลี จ ปิงฺคิยานี. \\ สุปินา จ วสฺสา วาจา, กุลํ นิสฺสารณีเยน จาติ. \\ \textbf{จตุตฺโถ ปณฺณาสโก สมตฺโต.}}}

\csromanmatikaanchor{mtk.83}
\csromantocmarknopage{chapter}{5. ปญฺจมปณฺณาสก}{mtk.83}
\bookmark[dest={mtk.83},level=0]{5. ปญฺจมปณฺณาสก}
\chapterheadpageread{5. ปญฺจมปณฺณาสก}
\csromanmatikaanchor{mtk.84}
\csromantocmarknopage{section}{(21) 1. กิมิลวคฺค}{mtk.84}
\bookmark[dest={mtk.84},level=1]{(21) 1. กิมิลวคฺค}
\csromanheader{1}{\textbf{(21) 1. กิมิลวคฺค}}
\csromanmatikaanchor{mtk.85}
\csromantocmarkref{subsection}{1. กิมิลสุตฺต}{mtk.85}
\bookmark[dest={mtk.85},level=2]{1. กิมิลสุตฺต}
\csromanheader{2}{1. กิมิลสุตฺต}
\csromanmatikaanchor{mtk.86}
\csromantocmarkref{subsection}{2-4. ธมฺมสฺสวนสุตฺตาทิ}{mtk.86}
\bookmark[dest={mtk.86},level=2]{2-4. ธมฺมสฺสวนสุตฺตาทิ}
\proseitem{201}{\csromanfolio{216}เอกํ สมยํ ภควา กิมิลายํ\footnote{กิมฺพิลายํ (สี, อิ)} วิหรติ เวฬุวเน.\csromanspacer{} อถ โข อายสฺมา กิมิโล\footnote{กิมฺพิโล (สี, สฺยา, กํ, อิ)} เยน ภควา เตนุปสงฺกมิ, อุปสงฺกมิตฺวา ภควนฺตํ อภิวาเทตฺวา เอกมนฺตํ นิสีทิ, เอกมนฺตํ นิสินฺโน โข อายสฺมา กิมิโล ภควนฺตํ เอตทโวจ “โก นุ โข ภนฺเต เหตุ โก ปจฺจโย, เยน ตถาคเต ปรินิพฺพุเต สทฺธมฺโม น จิรฏฺฐิติโก โหตี”ติ.\csromanspacer{} อิธ กิมิล ตถาคเต ปรินิพฺพุเต ภิกฺขู ภิกฺขุนิโย อุปาสกา อุปาสิกาโย สตฺถริ อคารวา วิหรนฺติ อปฺปติสฺสา, ธมฺเม อคารวา วิหรนฺติ อปฺปติสฺสา, สํเฆ อคารวา วิหรนฺติ อปฺปติสฺสา, สิกฺขาย อคารวา วิหรนฺติ อปฺปติสฺสา, อญฺญมญฺญํ อคารวา วิหรนฺติ อปฺปติสฺสา.\csromanspacer{} อยํ โข กิมิล เหตุ อยํ ปจฺจโย, เยน ตถาคเต ปรินิพฺพุเต สทฺธมฺโม น จิรฏฺฐิติโก โหตีติ.}

\prose{โก ปน ภนฺเต เหตุ โก ปจฺจโย, เยน ตถาคเต ปรินิพฺพุเต สทฺธมฺโม จิรฏฺฐิติโก โหตีติ.\csromanspacer{} อิธ กิมิล ตถาคเต ปรินิพฺพุเต ภิกฺขู ภิกฺขุนิโย อุปาสกา อุปาสิกาโย สตฺถริ สคารวา วิหรนฺติ สปฺปติสฺสา, ธมฺเม สคารวา วิหรนฺติ สปฺปติสฺสา, สํเฆ สคารวา วิหรนฺติ สปฺปติสฺสา, สิกฺขาย สคารวา วิหรนฺติ สปฺปติสฺสา, อญฺญมญฺญํ สคารวา วิหรนฺติ สปฺปติสฺสา.\csromanspacer{} อยํ โข กิมิล เหตุ อยํ ปจฺจโย, เยน ตถาคเต ปรินิพฺพุเต สทฺธมฺโม จิรฏฺฐิติโก โหตีติ.\csromanspacer{}\csromanspacer{}ปฐมํ.}

\csromanheader{2}{2. ธมฺมสฺสวนสุตฺต}
\proseitem{202}{ปญฺจิเม ภิกฺขเว อานิสํสา ธมฺมสฺสวเน.\csromanspacer{} กตเม ปญฺจ, อสฺสุตํ สุณาติ, สุตํ ปริโยทาเปติ, กงฺขํ วิตรติ\footnote{วิหนติ (สฺยา, กํ, อิ, ก)}, ทิฏฺฐิํ อุชุํ กโรติ, จิตฺตมสฺส ปสีทติ.\csromanspacer{} อิเม โข ภิกฺขเว ปญฺจ อานิสํสา ธมฺมสฺสวเนติ.\csromanspacer{}\csromanspacer{}ทุติยํ.}

\chapterheadpageread{(21) 1. กิมิลวคฺค \textbf{3. อสฺสาชานียสุตฺต}}
\proseitem{203}{\csromanfolio{217}ปญฺจหิ ภิกฺขเว องฺเคหิ สมนฺนาคโต รญฺโญ ภโทฺร\footnote{ภทฺโท (อิ)} อสฺสาชานีโย ราชารโห โหติ ราชโภคฺโค, รญฺโญ องฺคนฺเตฺวว สงฺขํ คจฺฉติ.}

\prose{กตเมหิ ปญฺจหิ? อชฺชเวน ชเวน มทฺทเวน ขนฺติยา โสรจฺเจน.\csromanspacer{} อิเมหิ โข ภิกฺขเว ปญฺจหิ องฺเคหิ สมนฺนาคโต รญฺโญ ภโทฺร อสฺสาชานีโย ราชารโห โหติ ราชโภคฺโค, รญฺโญ องฺคนฺเตฺวว สงฺขํ คจฺฉติ.\csromanspacer{} เอวเมวํ โข ภิกฺขเว ปญฺจหิ ธมฺเมหิ สมนฺนาคโต ภิกฺขุ อาหุเนยฺโย โหติ ปาหุเนยฺโย ทกฺขิเณยฺโย อญฺชลิกรณีโย อนุตฺตรํ ปุญฺญกฺเขตฺตํ โลกสฺส.}

\prose{กตเมหิ ปญฺจหิ? อชฺชเวน ชเวน มทฺทเวน ขนฺติยา โสรจฺเจน.\csromanspacer{} อิเมหิ โข ภิกฺขเว ปญฺจหิ ธมฺเมหิ สมนฺนาคโต ภิกฺขุ อาหุเนยฺโย โหติ ปาหุเนยฺโย ทกฺขิเณยฺโย อญฺชลิกรณีโย อนุตฺตรํ ปุญฺญกฺเขตฺตํ โลกสฺสาติ.\csromanspacer{}\csromanspacer{}ตติยํ.}

\csromanheader{2}{4. พลสุตฺต}
\proseitem{204}{ปญฺจิมานิ ภิกฺขเว พลานิ.\csromanspacer{} กตมานิ ปญฺจ? สทฺธาพลํ หิริพลํ โอตฺตปฺปพลํ วีริยพลํ ปญฺญาพลํ.\csromanspacer{} อิมานิ โข ภิกฺขเว ปญฺจ พลานีติ.\csromanspacer{}\csromanspacer{}จตุตฺถํ.}

\csromanmatikaanchor{mtk.87}
\csromantocmarkref{subsection}{5. เจโตขิลสุตฺต}{mtk.87}
\bookmark[dest={mtk.87},level=2]{5. เจโตขิลสุตฺต}
\csromanheader{2}{5. เจโตขิลสุตฺต}
\proseitem{205}{\csromansymbolfootnote{*}{อํ 3. 252; ม 1. 145; ที 3. 198 ปิฏฺเฐสุปิ.}ปญฺจิเม ภิกฺขเว เจโตขิลา.\csromanspacer{} กตเม ปญฺจ? อิธ ภิกฺขเว ภิกฺขุ สตฺถริ กงฺขติ วิจิกิจฺฉติ นาธิมุจฺจติ น สมฺปสีทติ.\csromanspacer{} โย โส ภิกฺขเว ภิกฺขุ สตฺถริ กงฺขติ วิจิกิจฺฉติ นาธิมุจฺจติ น สมฺปสีทติ, ตสฺส จิตฺตํ น นมติ อาตปฺปาย อนุโยคาย สาตจฺจาย ปธานาย.\csromanspacer{} ยสฺส จิตฺตํ น นมติ อาตปฺปาย อนุโยคาย สาตจฺจาย ปธานาย.\csromanspacer{} อยํ ปฐโม เจโตขิโล.}

\prose{ปุน จปรํ ภิกฺขเว ภิกฺขุ ธมฺเม กงฺขติ ฯเปฯ.\csromanspacer{} สํเฆ กงฺขติ ฯเปฯ.\csromanspacer{} สิกฺขาย กงฺขติ ฯเปฯ.\csromanspacer{} สพฺรหฺมจารีสุ กุปิโต โหติ อนตฺตมโน อาหตจิตฺโต ขิลชาโต.\csromanspacer{} โย โส ภิกฺขเว ภิกฺขุ สพฺรหฺมจารีสุ กุปิโต \csromanfolio{218}โหติ อนตฺตมโน อาหตจิตฺโต ขิลชาโต, ตสฺส จิตฺตํ น นมติ อาตปฺปาย อนุโยคาย สาตจฺจาย ปธานาย.\csromanspacer{} ยสฺส จิตฺตํ น นมติ อาตปฺปาย อนุโยคาย สาตจฺจาย ปธานาย.\csromanspacer{} อยํ ปญฺจโม เจโตขิโล.\csromanspacer{} อิเม โข ภิกฺขเว ปญฺจ เจโตขิลาติ.\csromanspacer{}\csromanspacer{}ปญฺจมํ.}

\csromanmatikaanchor{mtk.88}
\csromantocmarkref{subsection}{6. วินิพนฺธสุตฺต}{mtk.88}
\bookmark[dest={mtk.88},level=2]{6. วินิพนฺธสุตฺต}
\csromanheader{2}{6. วินิพนฺธสุตฺต}
\proseitem{206}{ปญฺจิเม ภิกฺขเว เจตโสวินิพนฺธา\footnote{เจโตวินิพทฺธา (สารตฺถทีปนีฏีกายํ)}.\csromanspacer{} กตเม ปญฺจ, อิธ ภิกฺขเว ภิกฺขุ กาเมสุ อวีตราโค\footnote{อวิคตราโค (ก)} โหติ อวิคตจฺฉนฺโท อวิคตเปโม อวิคตปิปาโส อวิคตปริฬาโห อวิคตตโณฺห.\csromanspacer{} โย โส ภิกฺขเว ภิกฺขุ กาเมสุ อวีตราโค โหติ อวิคต จฺฉนฺโท อวิคตเปโม อวิคตปิปาโส อวิคตปริฬาโห อวิคตตโณฺห, ตสฺส จิตฺตํ น นมติ อาตปฺปาย อนุโยคาย สาตจฺจาย ปธานาย.\csromanspacer{} ยสฺส จิตฺตํ น นมติ อาตปฺปาย อนุโยคาย สาตจฺจาย ปธานาย.\csromanspacer{} อยํ ปฐโม เจตโสวินิพนฺโธ.}

\prose{ปุน จปรํ ภิกฺขเว ภิกฺขุ กาเย อวีตราโค โหติ ฯเปฯ รูเป อวีตราโค โหติ ฯเปฯ ยาวทตฺถํ อุทราวเทหกํ ภุญฺชิตฺวา เสยฺยสุขํ ปสฺสสุขํ มิทฺธสุขํ อนุยุตฺโต วิหรติ ฯเปฯ อญฺญตรํ เทวนิกายํ ปณิธาย พฺรหฺมจริยํ จรติ “อิมินาหํ สีเลน วา วเตน วา ตเปน วา พฺรหฺมจริเยน วา เทโว วา ภวิสฺสามิ เทวญฺญตโร วา”ติ.\csromanspacer{} โย โส ภิกฺขเว ภิกฺขุ อญฺญตรํ เทวนิกายํ ปณิธาย พฺรหฺมจริยํ จรติ “อิมินาหํ สีเลน วา วเตน วา ตเปน วา พฺรหฺมจริเยน วา เทโว วา ภวิสฺสามิ เทวญฺญตโร วา”ติ.\csromanspacer{} ตสฺส จิตฺตํ น นมติ อาตปฺปาย อนุโยคาย สาตจฺจาย ปธานาย.\csromanspacer{} ยสฺส จิตฺตํ น นมติ อาตปฺปาย อนุโยคาย สาตจฺจาย ปธานาย.\csromanspacer{} อยํ ปญฺจโม เจตโสวินิพนฺโธ.\csromanspacer{} อิเม โข ภิกฺขเว ปญฺจ เจตโสวินิพนฺธาติ.\csromanspacer{}\csromanspacer{}ฉฏฺฐํ.}

\csromanmatikaanchor{mtk.89}
\csromantocmarkref{subsection}{7. ยาคุสุตฺต}{mtk.89}
\bookmark[dest={mtk.89},level=2]{7. ยาคุสุตฺต}
\csromanheader{2}{7. ยาคุสุตฺต}
\proseitem{207}{ปญฺจิเม ภิกฺขเว อานิสํสา ยาคุยา.\csromanspacer{} กตเม ปญฺจ, ขุทฺทํ\footnote{ขุทํ (สี, อิ) ขุธาติ สกฺกตานุโลมํ.} ปฏิหนติ, ปิปาสํ ปฏิวิเนติ, วาตํ อนุโลเมติ, วตฺถิํ โสเธติ,}

\vspace{\tipitakaparskip}
\csromancenter{(21) 1.\csromanspacer{} กิมิลวคฺค อามาวเสสํ ปาเจติ.\csromanspacer{} อิเม โข ภิกฺขเว ปญฺจ อานิสํสา ยาคุยาติ.\csromanspacer{}\csromanspacer{}สตฺตมํ.}
\csromanmatikaanchor{mtk.90}
\csromantocmarkref{subsection}{8. ทนฺตกฏฺฐสุตฺต}{mtk.90}
\bookmark[dest={mtk.90},level=2]{8. ทนฺตกฏฺฐสุตฺต}
\csromanheader{2}{8. ทนฺตกฏฺฐสุตฺต}
\proseitem{208}{\csromanfolio{219}ปญฺจิเม ภิกฺขเว อาทีนวา ทนฺตกฏฺฐสฺส อขาทเน.\csromanspacer{} กตเม ปญฺจ, อจกฺขุสฺสํ, มุขํ ทุคฺคนฺธํ โหติ, รสหรณิโย น วิสุชฺฌนฺติ, ปิตฺตํ เสมฺหํ ภตฺตํ ปริโยนนฺธติ\footnote{ปริโยนทฺธนฺติ (สี, อิ), ปริโยนทฺธติ (สฺยา, กํ, ก) วิ 4. 278 ปิฏฺเฐ ปสฺสิตพฺพํ.}, ภตฺตมสฺส นจฺฉาเทติ.\csromanspacer{} อิเม โข ภิกฺขเว ปญฺจ อาทีนวา ทนฺตกฏฺฐสฺส อขาทเน.}

\prose{ปญฺจิเม ภิกฺขเว อานิสํสา ทนฺตกฏฺฐสฺส ขาทเน.\csromanspacer{} กตเม ปญฺจ, จกฺขุสฺสํ, มุขํ น ทุคฺคนฺธํ โหติ, รสหรณิโย วิสุชฺฌนฺติ, ปิตฺตํ เสมฺหํ ภตฺตํ น ปริโยนนฺธติ, ภตฺตมสฺส ฉาเทติ.\csromanspacer{} อิเม โข ภิกฺขเว ปญฺจ อานิสํสา ทนฺตกฏฺฐสฺส ขาทเนติ.\csromanspacer{}\csromanspacer{}อฏฺฐมํ.}

\csromanmatikaanchor{mtk.91}
\csromantocmarkref{subsection}{9. คีตสฺสรสุตฺต}{mtk.91}
\bookmark[dest={mtk.91},level=2]{9. คีตสฺสรสุตฺต}
\csromanheader{2}{9. คีตสฺสรสุตฺต}
\proseitem{209}{ปญฺจิเม ภิกฺขเว อาทีนวา อายตเกน คีตสฺสเรน ธมฺมํ ภณนฺตสฺส.\csromanspacer{} กตเม ปญฺจ, อตฺตนาปิ ตสฺมิํ สเร สารชฺชติ, ปเรปิ ตสฺมิํ สเร สารชฺชนฺติ, คหปติกาปิ อุชฺฌายนฺติ “ยเถว มยํ คายาม, เอวเมว โข สมณา สกฺยปุตฺติยา คายนฺตี”ติ, สรกุตฺติมฺปิ นิกามยมานสฺส สมาธิสฺส ภงฺโค โหติ, ปจฺฉิมา ชนตา ทิฏฺฐานุคติํ อาปชฺชติ.\csromanspacer{} อิเม โข ภิกฺขเว ปญฺจ อาทีนวา อายตเกน คีตสฺสเรน ธมฺมํ ภณนฺตสฺสาติ.\csromanspacer{}\csromanspacer{}นวมํ.}

\csromanmatikaanchor{mtk.92}
\csromantocmarkref{subsection}{10. มุฏฺฐสฺสติสุตฺต}{mtk.92}
\bookmark[dest={mtk.92},level=2]{10. มุฏฺฐสฺสติสุตฺต}
\csromanheader{2}{10. มุฏฺฐสฺสติสุตฺต}
\proseitem{210}{ปญฺจิเม ภิกฺขเว อาทีนวา มุฏฺฐสฺสติสฺส อสมฺปชานสฺส นิทฺทํ โอกฺกมยโต.\csromanspacer{} กตเม ปญฺจ, ทุกฺขํ สุปติ, ทุกฺขํ ปฏิพุชฺฌติ, ปาปกํ สุปินํ ปสฺสติ, เทวตา น รกฺขนฺติ, อสุจิ มุจฺจติ.\csromanspacer{} อิเม โข ภิกฺขเว ปญฺจ อาทีนวา มุฏฺฐสฺสติสฺส อสมฺปชานสฺส นิทฺทํ โอกฺกมยโต.}

\csromanmatikaanchor{mtk.93}
\csromanmatikaanchor{mtk.94}
\csromantocmarknopage{section}{(22) 2. อกฺโกสกวคฺค}{mtk.93}
\bookmark[dest={mtk.93},level=1]{(22) 2. อกฺโกสกวคฺค}
\csromantocmarkref{subsection}{1-2. อกฺโกสกสุตฺตาทิ}{mtk.94}
\bookmark[dest={mtk.94},level=2]{1-2. อกฺโกสกสุตฺตาทิ}
\prose{ปญฺจิเม ภิกฺขเว อานิสํสา อุปฏฺฐิตสฺสติสฺส สมฺปชานสฺส นิทฺทํ โอกฺกมยโต.\csromanspacer{} กตเม ปญฺจ, สุขํ สุปติ, สุขํ ปฏิพุชฺฌติ, น ปาปกํ สุปินํ \csromanfolio{220}ปสฺสติ, เทวตา รกฺขนฺติ, อสุจิ น มุจฺจติ.\csromanspacer{} อิเม โข ภิกฺขเว ปญฺจ อานิสํสา อุปฏฺฐิตสฺสติสฺส สมฺปชานสฺส นิทฺทํ โอกฺกมยโตติ.\csromanspacer{}\csromanspacer{}ทสมํ.\csromanspacer{} กิมิลวคฺโค ปฐโม.}

\titlehead{\textbf{ตสฺสุทฺทานํ}}
\csromangathameasure{กิมิโล ธมฺมสฺสวนํ, อาชานีโย พลํ ขิลํ. \\ วินิพนฺธํ ยาคุ กฏฺฐํ, คีตํ มุฏฺฐสฺสตินา จาติ.}
\csromangathagroup{\csromangathastanza{กิมิโล ธมฺมสฺสวนํ, อาชานีโย พลํ ขิลํ. \\ วินิพนฺธํ ยาคุ กฏฺฐํ, คีตํ มุฏฺฐสฺสตินา จาติ.}}

\csromanheader{1}{\textbf{(22) 2. อกฺโกสกวคฺค}}
\csromanheader{2}{1. อกฺโกสกสุตฺต}
\proseitem{211}{โย โส ภิกฺขเว ภิกฺขุ อกฺโกสกปริภาสโก อริยูปวาที สพฺรหฺม จรีนํ, ตสฺส ปญฺจ อาทีนวา ปาฏิกงฺขา.\csromanspacer{} กตเม ปญฺจ? ปาราชิโก วา โหติ ฉินฺน ปริปนฺโถ\footnote{ฉินฺนปริพนฺโธ (ก)}, อญฺญตรํ วา สํกิลิฏฺฐํ อาปตฺติํ อาปชฺชติ, พาฬฺหํ วา โรคาตงฺกํ ผุสติ, สมฺมูโฬฺห กาลํ กโรติ, กายสฺส เภทา ปรํ มรณา อปายํ ทุคฺคติํ วินิปาตํ นิรยํ อุปปชฺชติ.\csromanspacer{} โย โส ภิกฺขเว ภิกฺขุ อกฺโกสกปริภาสโก อริยูปวาที สพฺรหฺมจารีนํ, ตสฺส อิเม ปญฺจ อาทีนวา ปาฏิกงฺขาติ.\csromanspacer{}\csromanspacer{}ปฐมํ.}

\csromanheader{2}{2. ภณฺฑนการกสุตฺต}
\proseitem{212}{โย โส ภิกฺขเว ภิกฺขุ ภณฺฑนการโก กลหการโก วิวาทการโก ภสฺสการโก สํเฆ อธิกรณการโก, ตสฺส ปญฺจ อาทีนวา ปาฏิกงฺขา.\csromanspacer{} กตเม ปญฺจ? อนธิคตํ นาธิคจฺฉติ, อธิคตา\footnote{อธิคตํ (สพฺพตฺถ), อุปริ 222 สุตฺเต ปน อธิคตาติปิ ทิสฺสติ.} ปริหายติ, ปาปโก กิตฺติสทฺโท อพฺภุคฺคจฺฉติ, สมฺมูโฬฺห กาลํ กโรติ, กายสฺส เภทา ปรํ มรณา อปายํ ทุคฺคติํ วินิปาตํ นิรยํ อุปปชฺชติ.\csromanspacer{} โย โส ภิกฺขเว ภิกฺขุ ภณฺฑนการโก กลหการโก วิวาทการโก ภสฺสการโก สํเฆ อธิกรณการโก, ตสฺส อิเม ปญฺจ อาทีนวา ปาฏิกงฺขาติ.\csromanspacer{}\csromanspacer{}ทุติยํ.}

\csromanmatikaanchor{mtk.95}
\csromantocmarkref{subsection}{3. สีลสุตฺต}{mtk.95}
\bookmark[dest={mtk.95},level=2]{3. สีลสุตฺต}
\csromanheader{2}{(22) 2. อกฺโกสกวคฺค \textbf{3. สีลสุตฺต}}
\proseitem{213}{\csromanfolio{221}ปญฺจิเม ภิกฺขเว อาทีนวา ทุสฺสีลสฺส สีลวิปตฺติยา.\csromanspacer{} กตเม ปญฺจ, อิธ ภิกฺขเว ทุสฺสีโล สีลวิปนฺโน ปมาทาธิกรณํ มหติํ โภคชานิํ นิคจฺฉติ.\csromanspacer{} อยํ ภิกฺขเว ปฐโม อาทีนโว ทุสฺสีลสฺส สีลวิปตฺติยา.}

\prose{ปุน จปรํ ภิกฺขเว ทุสฺสีลสฺส สีลวิปนฺนสฺส ปาปโก กิตฺติสทฺโท อพฺภุคฺคจฺฉติ.\csromanspacer{} อยํ ภิกฺขเว ทุติโย อาทีนโว ทุสฺสีลสฺส สีลวิปตฺติยา.}

\prose{ปุน จปรํ ภิกฺขเว ทุสฺสีโล สีลวิปนฺโน ยญฺญเทว ปริสํ อุปสงฺกมติ ยทิ ขตฺติยปริสํ ยทิ พฺราหฺมณปริสํ ยทิ คหปติปริสํ ยทิ สมณปริสํ, อวิสารโท อุปสงฺกมติ มงฺกุภูโต.\csromanspacer{} อยํ ภิกฺขเว ตติโย อาทีนโว ทุสฺสีลสฺส สีลวิปตฺติยา.}

\prose{ปุน จปรํ ภิกฺขเว ทุสฺสีโล สีลวิปนฺโน สมฺมูโฬฺห กาลํ กโรติ.\csromanspacer{} อยํ ภิกฺขเว จตุตฺโถ อาทีนโว ทุสฺสีลสฺส สีลวิปตฺติยา.}

\prose{ปุน จปรํ ภิกฺขเว ทุสฺสีโล สีลวิปนฺโน กายสฺส เภทา ปรํ มรณา อปายํ ทุคฺคติํ วินิปาตํ นิรยํ อุปปชฺชติ.\csromanspacer{} อยํ ภิกฺขเว ปญฺจโม อาทีนโว ทุสฺสีลสฺส สีลวิปตฺติยา.\csromanspacer{} อิเม โข ภิกฺขเว ปญฺจ อาทีนวา ทุสฺสีลสฺส สีลวิปตฺติยา.}

\prose{ปญฺจิเม ภิกฺขเว อานิสํสา สีลวโต สีลสมฺปทาย.\csromanspacer{} กตเม ปญฺจ, อิธ ภิกฺขเว สีลวา สีลสมฺปนฺโน อปฺปมาทาธิกรณํ มหนฺตํ โภคกฺขนฺธํ อธิคจฺฉติ.\csromanspacer{} อยํ ภิกฺขเว ปฐโม อานิสํโส สีลวโต สีลสมฺปทาย.}

\prose{ปุน จปรํ ภิกฺขเว สีลวโต สีลสมฺปนฺนสฺส กลฺยาโณ กิตฺติสทฺโท อพฺภุคฺคจฺฉติ.\csromanspacer{} อยํ ภิกฺขเว ทุติโย อานิสํโส สีลวโต สีลสมฺปทาย.}

\prose{ปุน จปรํ ภิกฺขเว สีลวา สีลสมฺปนฺโน ยญฺญเทว ปริสํ อุปสงฺกมติ ยทิ ขตฺติยปริสํ ยทิ พฺราหฺมณปริสํ ยทิ คหปติปริสํ ยทิ สมณปริสํ, วิสารโท อุปสงฺกมติ อมงฺกุภูโต.\csromanspacer{} อยํ ภิกฺขเว ตติโย อานิสํโส สีลวโต สีลสมฺปทาย.}

\csromanmatikaanchor{mtk.96}
\csromanmatikaanchor{mtk.97}
\csromantocmarkref{subsection}{4. พหุภาณิสุตฺต}{mtk.96}
\bookmark[dest={mtk.96},level=2]{4. พหุภาณิสุตฺต}
\csromantocmarkref{subsection}{5-6. อกฺขนฺติสุตฺตทฺวย}{mtk.97}
\bookmark[dest={mtk.97},level=2]{5-6. อกฺขนฺติสุตฺตทฺวย}
\prose{\csromanfolio{222}ปุน จปรํ ภิกฺขเว สีลวา สีลสมฺปนฺโน อสมฺมูโฬฺห กาลํ กโรติ.\csromanspacer{} อยํ ภิกฺขเว จตุตฺโถ อานิสํโส สีลวโต สีลสมฺปทาย.}

\prose{ปุน จปรํ ภิกฺขเว สีลวา สีลสมฺปนฺโน กายสฺส เภทา ปรํ มรณา สุคติํ สคฺคํ โลกํ อุปปชฺชติ.\csromanspacer{} อยํ ภิกฺขเว ปญฺจโม อานิสํโส สีลวโต สีลสมฺปทาย.\csromanspacer{} อิเม โข ภิกฺขเว ปญฺจ อานิสํสา สีลวโต สีลสมฺปทายาติ.\csromanspacer{}\csromanspacer{}ตติยํ.}

\csromanheader{2}{4. พหุภาณิสุตฺต}
\proseitem{214}{ปญฺจิเม ภิกฺขเว อาทีนวา พหุภาณิสฺมิํ ปุคฺคเล.\csromanspacer{} กตเม ปญฺจ? มุสา ภณติ, ปิสุณํ ภณติ, ผรุสํ ภณติ, สมฺผปฺปลาปํ ภณติ, กายสฺส เภทา ปรํ มรณา อปายํ ทุคฺคติํ วินิปาตํ นิรยํ อุปปชฺชติ.\csromanspacer{} อิเม โข ภิกฺขเว ปญฺจ อาทีนวา พหุภาณิสฺมิํ ปุคฺคเล.}

\prose{ปญฺจิเม ภิกฺขเว อานิสํสา มนฺตภาณิสฺมิํ ปุคฺคเล.\csromanspacer{} กตเม ปญฺจ? น มุสา ภณติ, น ปิสุณํ ภณติ, น ผรุสํ ภณติ, น สมฺผปฺปลาปํ ภณติ, กายสฺส เภทา ปรํ มรณา สุคติํ สคฺคํ โลกํ อุปปชฺชติ.\csromanspacer{} อิเม โข ภิกฺขเว ปญฺจ อานิสํสา มนฺตภาณิสฺมิํ ปุคฺคเลติ.\csromanspacer{}\csromanspacer{}จตุตฺถํ.}

\csromanheader{2}{5. ปฐม อกฺขนฺติสุตฺต}
\proseitem{215}{ปญฺจิเม ภิกฺขเว อาทีนวา อกฺขนฺติยา.\csromanspacer{} กตเม ปญฺจ? พหุโน ชนสฺส อปฺปิโย โหติ อมนาโป, เวรพหุโล จ โหติ, วชฺชพหุโล จ, สมฺมูโฬฺห กาลํ กโรติ, กายสฺส เภทา ปรํ มรณา อปายํ ทุคฺคติํ วินิปาตํ นิรยํ อุปปชฺชติ.\csromanspacer{} อิเม โข ภิกฺขเว ปญฺจ อาทีนวา อกฺขนฺติยา.}

\prose{ปญฺจิเม ภิกฺขเว อานิสํสา ขนฺติยา.\csromanspacer{} กตเม ปญฺจ? พหุโน ชนสฺส ปิโย โหติ มนาโป, น เวรพหุโล โหติ, น วชฺชพหุโล, อสมฺมูโฬฺห กาลํ กโรติ, กายสฺส เภทา ปรํ มรณา สุคติํ สคฺคํ โลกํ อุปปชฺชติ.\csromanspacer{} อิเม โข ภิกฺขเว ปญฺจ อานิสํสา ขนฺติยาติ.\csromanspacer{}\csromanspacer{}ปญฺจมํ.}

\csromanheader{2}{6. ทุติย อกฺขนฺติสุตฺต}
\proseitem{216}{ปญฺจิเม ภิกฺขเว อาทีนวา อกฺขนฺติยา.\csromanspacer{} กตเม ปญฺจ? พหุโน ชนสฺส อปฺปิโย โหติ อมนาโป, ลุทฺโท จ โหติ, วิปฺปฏิสารี}

\vspace{\tipitakaparskip}
\csromancenter{(22) 2.\csromanspacer{} อกฺโกสกวคฺค จ, สมฺมูโฬฺห กาลํ กโรติ, กายสฺส เภทา ปรํ มรณา อปายํ ทุคฺคติํ วินิปาตํ นิรยํ อุปปชฺชติ.\csromanspacer{} อิเม โข ภิกฺขเว ปญฺจ อาทีนวา อกฺขนฺติยา.}
\prose{\csromanfolio{223}ปญฺจิเม ภิกฺขเว อานิสํสา ขนฺติยา.\csromanspacer{} กตเม ปญฺจ? พหุโน ชนสฺส ปิโย โหติ มนาโป, อลุทฺโท จ โหติ, อวิปฺปฏิสารี จ, อสมฺมูโฬฺห กาลํ กโรติ, กายสฺส เภทา ปรํ มรณา สุคติํ สคฺคํ โลกํ อุปปชฺชติ.\csromanspacer{} อิเม โข ภิกฺขเว ปญฺจ อานิสํสา ขนฺติยาติ.\csromanspacer{}\csromanspacer{}ฉฏฺฐํ.}

\csromanmatikaanchor{mtk.98}
\csromantocmarkref{subsection}{7-8. อปาสาทิกสุตฺต}{mtk.98}
\bookmark[dest={mtk.98},level=2]{7-8. อปาสาทิกสุตฺต}
\csromanheader{2}{7. ปฐม อปาสาทิกสุตฺต}
\proseitem{217}{ปญฺจิเม ภิกฺขเว อาทีนวา อปาสาทิเก.\csromanspacer{} กตเม ปญฺจ? อตฺตาปิ อตฺตานํ อุปวทติ, อนุวิจฺจ วิญฺญู ครหนฺติ, ปาปโก กิตฺติสทฺโท อพฺภุคฺคจฺฉติ, สมฺมูโฬฺห กาลํ กโรติ, กายสฺส เภทา ปรํ มรณา อปายํ ทุคฺคติํ วินิปาตํ นิรยํ อุปปชฺชติ.\csromanspacer{} อิเม โข ภิกฺขเว ปญฺจ อาทีนวา อปาสาทิเก.}

\prose{ปญฺจิเม ภิกฺขเว อานิสํสา ปาสาทิเก.\csromanspacer{} กตเม ปญฺจ? อตฺตาปิ อตฺตานํ น อุปวทติ, อนุวิจฺจ วิญฺญู ปสํสนฺติ, กลฺยาโณ กิตฺติสทฺโท อพฺภุคฺคจฺฉติ, อสมฺมูโฬฺห กาลํ กโรติ, กายสฺส เภทา ปรํ มรณา สุคติํ สคฺคํ โลกํ อุปปชฺชติ.\csromanspacer{} อิเม โข ภิกฺขเว ปญฺจ อานิสํสา ปาสาทิเกติ.\csromanspacer{}\csromanspacer{}สตฺตมํ.}

\csromanheader{2}{8. ทุติย อปาสาทิกสุตฺต}
\proseitem{218}{ปญฺจิเม ภิกฺขเว อาทีนวา อปาสาทิเก.\csromanspacer{} กตเม ปญฺจ? อปฺปสนฺนา นปฺปสีทนฺติ, ปสนฺนานญฺจ เอกจฺจานํ อญฺญถตฺตํ โหติ, สตฺถุสาสนํ อกตํ\footnote{น กตํ (สฺยา, กํ, ก)} โหติ, ปจฺฉิมา ชนตา ทิฏฺฐานุคติํ อาปชฺชติ, จิตฺตมสฺส นปฺปสีทติ.\csromanspacer{} อิเม โข ภิกฺขเว ปญฺจ อาทีนวา อปาสาทิเก.}

\prose{ปญฺจิเม ภิกฺขเว อานิสํสา ปาสาทิเก.\csromanspacer{} กตเม ปญฺจ? อปฺปสนฺนา ปสีทนฺติ, ปสนฺนานญฺจ ภิยฺโยภาโว โหติ, สตฺถุสาสนํ กตํ โหติ, ปจฺฉิมา ชนตา ทิฏฺฐานุคติํ อาปชฺชติ, จิตฺตมสฺส ปสีทติ.\csromanspacer{} อิเม โข ภิกฺขเว ปญฺจ อานิสํสา ปาสาทิเกติ.\csromanspacer{}\csromanspacer{}อฏฺฐมํ.}

\csromanmatikaanchor{mtk.99}
\csromantocmarkref{subsection}{9. อคฺคิสุตฺต}{mtk.99}
\bookmark[dest={mtk.99},level=2]{9. อคฺคิสุตฺต}
\csromanheader{2}{9. อคฺคิสุตฺต}
\csromanmatikaanchor{mtk.100}
\csromanmatikaanchor{mtk.102}
\csromantocmarkref{subsection}{10. มธุราสุตฺต}{mtk.100}
\bookmark[dest={mtk.100},level=2]{10. มธุราสุตฺต}
\proseitem{219}{\csromanfolio{224}ปญฺจิเม ภิกฺขเว อาทีนวา อคฺคิสฺมิํ.\csromanspacer{} กตเม ปญฺจ? อจกฺขุสฺโส ทุพฺพณฺณกรโณ ทุพฺพลกรโณ สงฺคณิกาปวฑฺฒโน\footnote{สงฺคณิกาปวทฺธโน (สี), สงฺคณิการามพนฺธโน (ก)} ติรจฺฉานกถาปวตฺตนิโก โหติ.\csromanspacer{} อิเม โข ภิกฺขเว ปญฺจ อาทีนวา อคฺคิสฺมินฺติ.\csromanspacer{}\csromanspacer{}นวมํ.}

\csromanheader{2}{10. มธุราสุตฺต}
\proseitem{220}{ปญฺจิเม ภิกฺขเว อาทีนวา มธุรายํ.\csromanspacer{} กตเม ปญฺจ? วิสมา พหุรชา จณฺฑสุนขา วาฬยกฺขา ทุลฺลภปิณฺฑา.\csromanspacer{} อิเม โข ภิกฺขเว ปญฺจ อาทีนวา มธุรายนฺติ.\csromanspacer{}\csromanspacer{}ทสมํ.\csromanspacer{} อกฺโกสกวคฺโค ทุติโย.}

\titlehead{\textbf{ตสฺสุทฺทานํ}}
\csromangathameasure{อกฺโกสภณฺฑนสีลํ, พหุภาณี เทฺว อขนฺติโย. \\ อปาสาทิกา เทฺว วุตฺตา, อคฺคสฺมิํ มธุเรน จาติ.}
\csromangathagroup{\csromangathastanza{อกฺโกสภณฺฑนสีลํ, พหุภาณี เทฺว อขนฺติโย. \\ อปาสาทิกา เทฺว วุตฺตา, อคฺคสฺมิํ มธุเรน จาติ.}}

\csromanmatikaanchor{mtk.101}
\csromantocmarknopage{section}{(23) 3. ทีฆจาริกวคฺค}{mtk.101}
\bookmark[dest={mtk.101},level=1]{(23) 3. ทีฆจาริกวคฺค}
\csromantocmarkref{subsection}{1-2. ทีฆจาริกสุตฺตทฺวย}{mtk.102}
\bookmark[dest={mtk.102},level=2]{1-2. ทีฆจาริกสุตฺตทฺวย}
\csromanheader{1}{\textbf{(23) 3. ทีฆจาริกวคฺค}}
\csromanheader{2}{1. ปฐมทีฆจาริกสุตฺต}
\proseitem{221}{ปญฺจิเม ภิกฺขเว อาทีนวา ทีฆจาริกํ อนวตฺถจาริกํ อนุยุตฺตสฺส วิหรโต.\csromanspacer{} กตเม ปญฺจ? อสฺสุตํ น สุณาติ, สุตํ น ปริโยทาเปติ, สุเตเนกจฺเจน อวิสารโท โหติ, คาฬฺหํ\footnote{พาฬฺหํ (สฺยา)} โรคาตงฺกํ ผุสติ, น จ มิตฺตวา โหติ.\csromanspacer{} อิเม โข ภิกฺขเว ปญฺจ อาทีนวา ทีฆจาริกํ อนวตฺถจาริกํ อนุยุตฺตสฺส วิหรโต.}

\prose{ปญฺจิเม ภิกฺขเว อานิสํสา สมวตฺถจาเร.\csromanspacer{} กตเม ปญฺจ? อสฺสุตํ สุณาติ, สุตํ ปริโยทาเปติ, สุเตเนกจฺเจน วิสารโท โหติ, น คาฬฺหํ โรคาตงฺกํ ผุสติ, มิตฺตวา จ โหติ.\csromanspacer{} อิเม โข ภิกฺขเว ปญฺจ อานิสํสา สมวตฺถจาเรติ.\csromanspacer{}\csromanspacer{}ปฐมํ.}

\chapterheadpageread{(23) 3. ทีฆจาริกวคฺค \textbf{2. ทุติยทีฆจาริกสุตฺต}}
\proseitem{222}{\csromanfolio{225}ปญฺจิเม ภิกฺขเว อาทีนวา ทีฆจาริกํ อนวตฺถจาริกํ อนุยุตฺตสฺส วิหรโต.\csromanspacer{} กตเม ปญฺจ? อนธิคตํ นาธิคจฺฉติ, อธิคตา ปริหายติ, อธิคเตเนกจฺเจน อวิสารโท โหติ, คาฬฺหํ โรคาตงฺกํ ผุสติ, น จ มิตฺตวา โหติ.\csromanspacer{} อิเม โข ภิกฺขเว ปญฺจ อาทีนวา ทีฆจาริกํ อนวตฺถจาริกํ อนุยุตฺตสฺส วิหรโต.}

\prose{ปญฺจิเม ภิกฺขเว อานิสํสา สมวตฺถจาเร.\csromanspacer{} กตเม ปญฺจ? อนธิคตํ อธิคจฺฉติ, อธิคตา น ปริหายติ, อธิคเตเนกจฺเจน วิสารโท โหติ, น คาฬฺหํ โรคาตงฺกํ ผุสติ, มิตฺตวา จ โหติ.\csromanspacer{} อิเม โข ภิกฺขเว ปญฺจ อานิสํสา สมวตฺถจาเรติ.\csromanspacer{}\csromanspacer{}ทุติยํ.}

\csromanmatikaanchor{mtk.103}
\csromantocmarkref{subsection}{3-4. อตินิวาสสุตฺตทฺวย}{mtk.103}
\bookmark[dest={mtk.103},level=2]{3-4. อตินิวาสสุตฺตทฺวย}
\csromanheader{2}{3. อตินิวาสสุตฺต}
\proseitem{223}{ปญฺจิเม ภิกฺขเว อาทีนวา อตินิวาเส.\csromanspacer{} กตเม ปญฺจ? พหุภณฺโฑ โหติ พหุภณฺฑสนฺนิจโย, พหุเภสชฺโช โหติ พหุเภสชฺชสนฺนิจโย, พหุกิจฺโจ โหติ พหุกรณีโย, พฺยตฺโต กิํกรณีเยสุ สํสฏฺโฐ วิหรติ คหฏฺฐปพฺพชิเตหิ อนนุโลมิเกน คิหิสํสคฺเคน, ตมฺหา จ อาวาสา ปกฺกมนฺโต สาเปกฺโข ปกฺกมติ.\csromanspacer{} อิเม โข ภิกฺขเว ปญฺจ อาทีนวา อตินิวาเส.}

\prose{ปญฺจิเม ภิกฺขเว อานิสํสา สมวตฺถวาเส.\csromanspacer{} กตเม ปญฺจ? น พหุภณฺโฑ โหติ น พหุภณฺฑสนฺนิจโย, น พหุเภสชฺโช โหติ น พหุเภสชฺชสนฺนิจโย, น พหุกิจฺโจ โหติ น พหุกรณีโย, น พฺยตฺโต กิํกรณีเยสุ อสํสฏฺโฐ วิหรติ คหฏฺฐปพฺพชิเตหิ อนนุโลมิเกน คิหิสํสคฺเคน, ตมฺหา จ อาวาสา ปกฺกมนฺโต อนเปกฺโข ปกฺกมติ.\csromanspacer{} อิเม โข ภิกฺขเว ปญฺจ อานิสํสา สมวตฺถวาเสติ.\csromanspacer{}\csromanspacer{}ตติยํ.}

\csromanheader{2}{4. มจฺฉรีสุตฺต}
\proseitem{224}{ปญฺจิเม ภิกฺขเว อาทีนวา อตินิวาเส.\csromanspacer{} กตเม ปญฺจ? อาวาสมจฺฉรี โหติ, กุลมจฺฉรี โหติ, ลาภมจฺฉรี โหติ, วณฺณมจฺฉรี โหติ, ธมฺมมจฺฉรี โหติ.\csromanspacer{} อิเม โข ภิกฺขเว ปญฺจ อาทีนวา อตินิวาเส.}

\csromanmatikaanchor{mtk.104}
\csromantocmarkref{subsection}{5-6. กุลูปกสุตฺตทฺวย}{mtk.104}
\bookmark[dest={mtk.104},level=2]{5-6. กุลูปกสุตฺตทฺวย}
\prose{\csromanfolio{226}ปญฺจิเม ภิกฺขเว อานิสํสา สมวตฺถวาเส.\csromanspacer{} กตเม ปญฺจ? น อาวาสมจฺฉรี โหติ, น กุลมจฺฉรี โหติ, น ลาภมจฺฉรี โหติ, น วณฺณมจฺฉรี โหติ, น ธมฺมมจฺฉรี โหติ.\csromanspacer{} อิเม โข ภิกฺขเว ปญฺจ อานิสํสา สมวตฺถวาเสติ.\csromanspacer{}\csromanspacer{}จตุตฺถํ.}

\csromanheader{2}{5. ปฐมกุลูปกสุตฺต}
\proseitem{225}{ปญฺจิเม ภิกฺขเว อาทีนวา กุลูปเก\footnote{กุลุปเก (สฺยา, อิ), กุลูปเค (สี)}.\csromanspacer{} กตเม ปญฺจ? อนามนฺตจาเร อาปชฺชติ, รโห นิสชฺชาย อาปชฺชติ, ปฏิจฺฉนฺเน อาสเน อาปชฺชติ, มาตุคามสฺส อุตฺตริ ฉปฺปญฺจวาจาหิ ธมฺมํ เทเสนฺโต อาปชฺชติ, กามสงฺกปฺปพหุโล วิหรติ.\csromanspacer{} อิเม โข ภิกฺขเว ปญฺจ อาทีนวา กุลูปเกติ.\csromanspacer{}\csromanspacer{}ปญฺจมํ.}

\csromanheader{2}{6. ทุติยกุลูปกสุตฺต}
\proseitem{226}{ปญฺจิเม ภิกฺขเว อาทีนวา กุลูปกสฺส ภิกฺขุโน อติเวลํ กุเลสุ สํสฏฺฐสฺส วิหรโต.\csromanspacer{} กตเม ปญฺจ? มาตุคามสฺส อภิณฺหทสฺสนํ, ทสฺสเน สติ สํสคฺโค, สํสคฺเค สติ วิสฺสาโส, วิสฺสาเส สติ โอตาโร, โอติณฺณจิตฺตสฺเสตํ ปาฏิกงฺขํ “อนภิรโต วา พฺรหฺมจริยํ จริสฺสติ, อญฺญตรํ วา สํกิลิฏฺฐํ อาปตฺติํ อาปชฺชิสฺสติ, สิกฺขํ วา ปจฺจกฺขาย หีนายาวตฺติสฺสติ”.\csromanspacer{} อิเม โข ภิกฺขเว ปญฺจ อาทีนวา กุลูปกสฺส ภิกฺขุโน อติเวลํ กุเลสุ สํสฏฺฐสฺส วิหรโตติ.\csromanspacer{}\csromanspacer{}ฉฏฺฐํ.}

\csromanmatikaanchor{mtk.105}
\csromantocmarkref{subsection}{7. โภคสุตฺต}{mtk.105}
\bookmark[dest={mtk.105},level=2]{7. โภคสุตฺต}
\csromanheader{2}{7. โภคสุตฺต}
\proseitem{227}{ปญฺจิเม ภิกฺขเว อาทีนวา โภเคสุ.\csromanspacer{} กตเม ปญฺจ? อคฺคิสาธารณา โภคา, อุทกสาธารณา โภคา, ราชสาธารณา โภคา, โจรสาธารณา โภคา, อปฺปิเยหิ ทายาเทหิ สาธารณา โภคา.\csromanspacer{} อิเม โข ภิกฺขเว ปญฺจ อาทีนวา โภเคสุ.}

\prose{ปญฺจิเม ภิกฺขเว อานิสํสา โภเคสุ.\csromanspacer{} กตเม ปญฺจ? โภเค นิสฺสาย อตฺตานํ สุเขติ ปีเณติ สมฺมา สุขํ ปริหรติ, มาตาปิตโร สุเขติ ปีเณติ สมฺมา สุขํ ปริหรติ, ปุตฺตทารทาสกมฺมกรโปริเส}

\vspace{\tipitakaparskip}
\csromancenter{(23) 3.\csromanspacer{} ทีฆจาริกวคฺค สุเขติ ปีเณติ สมฺมา สุขํ ปริหรติ, มิตฺตามจฺเจ สุเขติ ปีเณติ สมฺมา สุขํ ปริหรติ, สมณพฺราหฺมเณสุ อุทฺธคฺคิกํ ทกฺขิณํ ปติฏฺฐาเปติ โสวคฺคิกํ สุขวิปากํ สคฺคสํวตฺตนิกํ.\csromanspacer{} อิเม โข ภิกฺขเว ปญฺจ อานิสํสา โภเคสูติ.\csromanspacer{}\csromanspacer{}สตฺตมํ.}
\csromanmatikaanchor{mtk.106}
\csromantocmarkref{subsection}{8. อุสฺสูรภตฺตสุตฺต}{mtk.106}
\bookmark[dest={mtk.106},level=2]{8. อุสฺสูรภตฺตสุตฺต}
\csromanheader{2}{8. อุสฺสูรภตฺตสุตฺต}
\csromanmatikaanchor{mtk.107}
\csromantocmarkref{subsection}{9-10. กณฺหสปฺปสุตฺต}{mtk.107}
\bookmark[dest={mtk.107},level=2]{9-10. กณฺหสปฺปสุตฺต}
\proseitem{228}{\csromanfolio{227}ปญฺจิเม ภิกฺขเว อาทีนวา อุสฺสูรภตฺเต กุเล.\csromanspacer{} กตเม ปญฺจ? เย เต อติถี ปาหุนา, เต น กาเลน ปฏิปูเชนฺติ.\csromanspacer{} ยา ตา พลิปฏิคฺคาหิกา เทวตา, ตา น กาเลน ปฏิปูเชนฺติ.\csromanspacer{} เย เต สมณพฺราหฺมณา เอกภตฺติกา รตฺตูปรตา วิรตา วิกาลโภชนา, เต น กาเลน ปฏิปูเชนฺติ.\csromanspacer{} ทาสกมฺมกรโปริสา วิมุขา กมฺมํ กโรนฺติ.\csromanspacer{} ตาวตกํเยว อสมเยน ภุตฺตํ อโนชวนฺตํ โหติ.\csromanspacer{} อิเม โข ภิกฺขเว ปญฺจ อาทีนวา อุสฺสูรภตฺเต กุเล.}

\prose{ปญฺจิเม ภิกฺขเว อานิสํสา สมยภตฺเต กุเล.\csromanspacer{} กตเม ปญฺจ? เย เต อติถี ปาหุนา, เต กาเลน ปฏิปูเชนฺติ.\csromanspacer{} ยา ตา พลิปฏิคฺคาหิกา เทวตา, ตา กาเลน ปฏิปูเชนฺติ.\csromanspacer{} เย เต สมณพฺราหฺมณา เอกภตฺติกา รตฺตูปรตา วิรตา วิกาลโภชนา, เต กาเลน ปฏิปูเชนฺติ.\csromanspacer{} ทาสกมฺมกรโปริสา อวิมุขา กมฺมํ กโรนฺติ.\csromanspacer{} ตาวตกํเยว สมเยน ภุตฺตํ โอชวนฺตํ โหติ.\csromanspacer{} อิเม โข ภิกฺขเว ปญฺจ อานิสํสา สมยภตฺเต กุเลติ.\csromanspacer{}\csromanspacer{}อฏฺฐมํ.}

\csromanheader{2}{9. ปฐมกณฺหสปฺปสุตฺต}
\proseitem{229}{ปญฺจิเม ภิกฺขเว อาทีนวา กณฺหสปฺเป.\csromanspacer{} กตเม ปญฺจ? อสุจิ ทุคฺคนฺโธ สภีรุ สปฺปฏิภโย มิตฺตทุพฺภี.\csromanspacer{} อิเม โข ภิกฺขเว ปญฺจ อาทีนวา กณฺหสปฺเป.\csromanspacer{} เอวเมวํ โข ภิกฺขเว ปญฺจิเม อาทีนวา มาตุคาเม.\csromanspacer{} กตเม ปญฺจ? อสุจิ ทุคฺคนฺโธ สภีรุ สปฺปฏิภโย มิตฺตทุพฺภี.\csromanspacer{} อิเม โข ภิกฺขเว ปญฺจ อาทีนวา มาตุคาเมติ.\csromanspacer{}\csromanspacer{}นวมํ.}

\csromanheader{2}{10. ทุติยกณฺหสปฺปสุตฺต}
\proseitem{230}{ปญฺจิเม ภิกฺขเว อาทีนวา กณฺหสปฺเป.\csromanspacer{} กตเม ปญฺจ? โกธโน อุปนาหี โฆรวิโส ทุชฺชิโวฺห มิตฺตทุพฺภี.\csromanspacer{} อิเม โข ภิกฺขเว ปญฺจ อาทีนวา กณฺหสปฺเป.}

\prose{\csromanfolio{228}เอวเมวํ โข ภิกฺขเว ปญฺจิเม อาทีนวา มาตุคาเม.\csromanspacer{} กตเม ปญฺจ? โกธโน อุปนาหี โฆรวิโส ทุชฺโชโวฺห มิตฺตทุพฺภี.\csromanspacer{} ตตฺริทํ ภิกฺขเว มาตุคามสฺส โฆรวิสตา, เยภุยฺเยน ภิกฺขเว มาตุคาโม ติพฺพราโค.\csromanspacer{} ตตฺริทํ ภิกฺขเว มาตุคามสฺส ทุชฺชิวฺหตา, เยภุยฺเยน ภิกฺขเว มาตุคาโม ปิสุณวาโจ.\csromanspacer{} ตตฺริทํ ภิกฺขเว มาตุคามสฺส มิตฺตทุพฺภิตา, เยภุยฺเยน ภิกฺขเว มาตุคาโม อติจารินี.\csromanspacer{} อิเม โข ภิกฺขเว ปญฺจ อาทีนวา มาตุคาเมติ.\csromanspacer{}\csromanspacer{}ทสมํ.\csromanspacer{} ทีฆจาริกวคฺโค ตติโย.}

\titlehead{\textbf{ตสฺสุทฺทานํ}}
\csromangathameasure{เทฺว ทีฆจาริกา วุตฺตา, อตินิวาสมจฺฉรี. \\ เทฺว จ กุลูปกา โภคา, ภตฺตํ สปฺปาปเร ทุเวติ.}
\csromangathagroup{\csromangathastanza{เทฺว ทีฆจาริกา วุตฺตา, อตินิวาสมจฺฉรี. \\ เทฺว จ กุลูปกา โภคา, ภตฺตํ สปฺปาปเร ทุเวติ.}}

\csromanmatikaanchor{mtk.108}
\csromantocmarknopage{section}{(24) 4. อาวาสิกวคฺค}{mtk.108}
\bookmark[dest={mtk.108},level=1]{(24) 4. อาวาสิกวคฺค}
\csromanheader{1}{\textbf{(24) 4. อาวาสิกวคฺค}}
\csromanmatikaanchor{mtk.109}
\csromantocmarkref{subsection}{1. อาวาสิกสุตฺต}{mtk.109}
\bookmark[dest={mtk.109},level=2]{1. อาวาสิกสุตฺต}
\csromanheader{2}{1. อาวาสิกสุตฺต}
\proseitem{231}{ปญฺจหิ ภิกฺขเว ธมฺเมหิ สมนฺนาคโต อาวาสิโก ภิกฺขุ อภาวนีโย โหติ.\csromanspacer{} กตเมหิ ปญฺจหิ? น อากปฺปสมฺปนฺโน โหติ น วตฺตสมฺปนฺโน, น พหุสฺสุโต โหติ น สุตธโร, น ปฏิสลฺเลขิตา\footnote{สลฺเลขิตา (ก-สี)} โหติ น ปฏิสลฺลานาราโม, น กลฺยาณวาโจ โหติ น กลฺยาณวากฺกรโณ, ทุปฺปญฺโญ โหติ ชโฬ เอฬมูโค.\csromanspacer{} อิเมหิ โข ภิกฺขเว ปญฺจหิ ธมฺเมหิ สมนฺนาคโต อาวาสิโก ภิกฺขุ อภาวนีโย โหติ.}

\prose{ปญฺจหิ ภิกฺขเว ธมฺเมหิ สมนฺนาคโต อาวาสิโก ภิกฺขุ ภาวนีโย โหติ.\csromanspacer{} กตเมหิ ปญฺจหิ? อากปฺปสมฺปนฺโน โหติ วตฺตสมฺปนฺโน, พหุสฺสุโต โหติ สุตธโร, ปฏิสลฺเลขิตา โหติ ปฏิสลฺลานาราโม, กลฺยาณวาโจ โหติ กลฺยาณวากฺกรโณ, ปญฺญวา โหติ}

\vspace{\tipitakaparskip}
\csromancenter{(24) 4.\csromanspacer{} อาวาสิกวคฺค อชโฬ อเนฬมูโค.\csromanspacer{} อิเมหิ โข ภิกฺขเว ปญฺจหิ ธมฺเมหิ สมนฺนาคโต อาวาสิโก ภิกฺขุ ภาวนีโย โหตีติ.\csromanspacer{}\csromanspacer{}ปฐมํ.}
\csromanheader{2}{2. ปิยสุตฺต}
\csromanmatikaanchor{mtk.110}
\csromantocmarkref{subsection}{2-10. ปิยสุตฺตาทิ}{mtk.110}
\bookmark[dest={mtk.110},level=2]{2-10. ปิยสุตฺตาทิ}
\proseitem{232}{\csromanfolio{229}ปญฺจหิ ภิกฺขเว ธมฺเมหิ สมนฺนาคโต อาวาสิโก ภิกฺขุ สพฺรหฺม จารีนํ ปิโย จ โหติ มนาโป จ ครุ จ ภาวนีโย จ.}

\prose{กตเมหิ ปญฺจหิ? สีลวา โหติ ปาติโมกฺขสํวรสํวุโต วิหรติ อาจาร โคจรสมฺปนฺโน, อณุมตฺเตสุ วชฺเชสุ ภยทสฺสาวี สมาทาย สิกฺขติ สิกฺขาปเทสุ.\csromanspacer{} พหุสฺสุโต โหติ สุตธโร สุตสนฺนิจโย, เย เต ธมฺมา อาทิกลฺยาณา มชฺเฌกลฺยาณา ปริโยสานกลฺยาณา สาตฺถํ สพฺยญฺชนํ เกวลปริปุณฺณํ ปริสุทฺธํ พฺรหฺม จริยํ อภิวทนฺติ, ตถารูปาสฺส ธมฺมา พหุสฺสุตา โหนฺติ ธาตา วจสา ปริจิตา มนสานุเปกฺขิตา ทิฏฺฐิยา สุปฺปฏิวิทฺธา.\csromanspacer{} กลฺยาณวาโจ โหติ กลฺยาณวากฺกรโณ โปริยา วาจาย สมนฺนาคโต วิสฺสฏฺฐาย อเนลคลาย อตฺถสฺส วิญฺญาปนิยา.\csromanspacer{} จตุนฺนํ ฌานานํ อาภิเจตสิกานํ ทิฏฺฐธมฺมสุขวิหารานํ นิกามลาภี โหติ อกิจฺฉลาภี อกสิรลาภี.\csromanspacer{} อาสวานํ ขยา อนาสวํ เจโตวิมุตฺติํ ปญฺญาวิมุตฺติํ ทิฏฺเฐว ธมฺเม สยํ อภิญฺญา สจฺฉิกตฺวา อุปสมฺปชฺช วิหรติ.\csromanspacer{} อิเมหิ โข ภิกฺขเว ปญฺจหิ ธมฺเมหิ สมนฺนาคโต อาวาสิโก ภิกฺขุ สพฺรหฺมจารีนํ ปิโย จ โหติ มนาโป จ ครุ จ ภาวนีโย จาติ.\csromanspacer{}\csromanspacer{}ทุติยํ.}

\csromanheader{2}{3. โสภนสุตฺต}
\proseitem{233}{ปญฺจหิ ภิกฺขเว ธมฺเมหิ สมนฺนาคโต อาวาสิโก ภิกฺขุ อาวาสํ โสเภติ.\csromanspacer{} กตเมหิ ปญฺจหิ? สีลวา โหติ ฯเปฯ สมาทาย สิกฺขติ สิกฺขาปเทสุ.\csromanspacer{} พหุสฺสุโต โหติ ฯเปฯ ทิฏฺฐิยา สุปฺปฏิวิทฺธา.\csromanspacer{} กลฺยาณวาโจ โหติ กลฺยาณวากฺกรโณ โปริยา วาจาย สมนฺนาคโต วิสฺสฏฺฐาย อเนลคลาย อตฺถสฺส วิญฺญาปนิยา.\csromanspacer{} ปฏิพโล โหติ อุปสงฺกมนฺเต ธมฺมิยา กถาย สนฺทสฺเสตุํ สมาทเปตุํ สมุตฺเตเชตุํ สมฺปหํเสตุํ.\csromanspacer{} จตุนฺนํ ฌานานํ อาภิเจกสิกานํ ทิฏฺฐธมฺม สุขวิหารานํ นิกามลาภี โหติ อกิจฺฉลาภี อกสิรลาภี.\csromanspacer{} อิเมหิ โข ภิกฺขเว ปญฺจหิ ธมฺเมหิ สมนฺนาคโต อาวาสิโก ภิกฺขุ อาวาสํ โสเภตีติ.\csromanspacer{}\csromanspacer{}ตติยํ.}

\csromanheader{2}{4. พหูปการสุตฺต}
\proseitem{234}{\csromanfolio{230}ปญฺจหิ ภิกฺขเว ธมฺเมหิ สมนฺนาคโต อาวาสิโก ภิกฺขุ อาวาสสฺส พหูปกาโร โหติ.\csromanspacer{} กตเมหิ ปญฺจหิ? สีลวา โหติ ฯเปฯ สมาทาย สิกฺขติ สิกฺขา ปเทสุ.\csromanspacer{} พหุสฺสุโต โหติ ฯเปฯ ทิฏฺฐิยา สุปฺปฏิวิทฺธา.\csromanspacer{} ขณฺฑผุลฺลํ ปฏิสงฺขโรติ.\csromanspacer{} มหา โข ปน ภิกฺขุสํโฆ อภิกฺกนฺโต นานาเวรชฺชกา ภิกฺขู, คิหีนํ อุปสงฺกมิตฺวา อาโรเจติ “มหา โข อาวุโส ภิกฺขุสํโฆ อภิกฺกนฺโต นานาเวรชฺชกา ภิกฺขู, กโรถ ปุญฺญานิ, สมโย ปุญฺญานิ กาตุนฺ”ติ.\csromanspacer{} จตุนฺนํ ฌานานํ อาภิเจตสิกานํ ทิฏฺฐธมฺมสุขวิหารานํ นิกามลาภี โหติ อกิจฺฉลาภี อกสิรลาภี.\csromanspacer{} อิเมหิ โข ภิกฺขเว ปญฺจหิ ธมฺเมหิ สมนฺนาคโต อาวาสิโก ภิกฺขุ อาวาสสฺส พหูปกาโร โหตีติ.\csromanspacer{}\csromanspacer{}จตุตฺถํ.}

\csromanheader{2}{5. อนุกมฺปสุตฺต}
\proseitem{235}{ปญฺจหิ ภิกฺขเว ธมฺเมหิ สมนฺนาคโต อาวาสิโก ภิกฺขุ คิหีนํ\footnote{คิหิํ (สฺยา), คิหี (กตฺถจิ)} อนุกมฺปติ.\csromanspacer{} กตเมหิ ปญฺจหิ? อธิสีเล\footnote{อธิสีเลสุ (สฺยา)} สมาทเปติ.\csromanspacer{} ธมฺมทสฺสเน นิเวเสติ.\csromanspacer{} คิลานเก อุปสงฺกมิตฺวา สติํ อุปฺปาเทติ “อรหคฺคตํ อายสฺมนฺโต สติํ อุปฏฺฐาเปถา”ติ.\csromanspacer{} มหา โข ปน ภิกฺขุสํโฆ อภิกฺกนฺโต นานาเวรชฺชกา ภิกฺขู, คิหีนํ อุปสงฺกมิตฺวา อาโรเจติ “มหา โข อาวุโส ภิกฺขุสํโฆ อภิกฺกนฺโต นานาเวรชฺชกา ภิกฺขู, กโรถ ปุญฺญานิ, สมโย ปุญฺญานิ กาตุนฺ”ติ.\csromanspacer{} ยํ โข ปนสฺส โภชนํ เทนฺติ ลูขํ วา ปณีตํ วา, ตํ อตฺตนา ปริภุญฺชติ, สทฺธาเทยฺยํ น วินิปาเตติ.\csromanspacer{} อิเมหิ โข ภิกฺขเว ปญฺจหิ ธมฺเมหิ สมนฺนาคโต อาวาสิโก ภิกฺขุ คิหีนํ อนุกมฺปตีติ.\csromanspacer{}\csromanspacer{}ปญฺจมํ.}

\csromanheader{2}{6. ปฐม อวณฺณารหสุตฺต}
\proseitem{236}{ปญฺจหิ ภิกฺขเว ธมฺเมหิ สมนฺนาคโต อาวาสิโก ภิกฺขุ ยถาภตํ นิกฺขิตฺโต เอวํ นิรเย.\csromanspacer{} กตเมหิ ปญฺจหิ? อนนุวิจฺจ อปริโยคาเหตฺวา อวณฺณารหสฺส วณฺณํ ภาสติ, อนนุวิจฺจ อปริโยคาเหตฺวา วณฺณารหสฺส อวณฺณํ ภาสติ, อนนุวิจฺจ อปริโยคาเหตฺวา อปฺปสาทนีเย ฐาเน ปสาทํ อุปทํเสติ, อนนุวิจฺจ}

\vspace{\tipitakaparskip}
\csromancenter{(24) 4.\csromanspacer{} อาวาสิกวคฺค อปริโยคาเหตฺวา ปสาทนีเย ฐาเน อปฺปสาทํ อุปทํเสติ, สทฺธาเทยฺยํ วินิปาเตติ.\csromanspacer{} อิเมหิ โข ภิกฺขเว ปญฺจหิ ธมฺเมหิ สมนฺนาคโต อาวาสิโก ภิกฺขุ ยถาภตํ นิกฺขิตฺโต เอวํ นิรเย.}
\prose{\csromanfolio{231}ปญฺจหิ ภิกฺขเว ธมฺเมหิ สมนฺนาคโต อาวาสิโก ภิกฺขุ ยถาภตํ นิกฺขิตฺโต เอวํ สคฺเค.\csromanspacer{} กตเมหิ ปญฺจหิ? อนุวิจฺจ ปริโยคาเหตฺวา อวณฺณารหสฺส อวณฺณํ ภาสติ, อนุวิจฺจ ปริโยคาเหตฺวา วณฺณารหสฺส วณฺณํ ภาสติ, อนุวิจฺจ ปริโยคาเหตฺวา อปฺปสาทนีเย ฐาเน อปฺปสาทํ อุปทํเสติ, อนุวิจฺจ ปริโยคาเหตฺวา ปสาทนีเย ฐาเน ปสาทํ อุปทํเสติ, สทฺธาเทยฺยํ น วินิปาเตติ.\csromanspacer{} อิเมหิ โข ภิกฺขเว ปญฺจหิ ธมฺเมหิ สมนฺนาคโต อาวาสิโก ภิกฺขุ ยถาภตํ นิกฺขิตฺโต เอวํ สคฺเคติ.\csromanspacer{}\csromanspacer{}ฉฏฺฐํ.}

\csromanheader{2}{7. ทุติย อวณฺณารหสุตฺต}
\proseitem{237}{ปญฺจหิ ภิกฺขเว ธมฺเมหิ สมนฺนาคโต อาวาสิโก ภิกฺขุ ยถาภตํ นิกฺขิตฺโต เอวํ นิรเย.\csromanspacer{} กตเมหิ ปญฺจหิ? อนนุวิจฺจ อปริโยคาเหตฺวา อวณฺณารหสฺส วณฺณํ ภาสติ, อนนุวิจฺจ อปริโยคาเหตฺวา วณฺณารหสฺส อวณฺณํ ภาสติ, อาวา สมจฺฉรี โหติ อาวาสปลิเคธี, กุลมจฺฉรี โหติ กุลปลิเคธี, สทฺธาเทยฺยํ วินิปาเตติ.\csromanspacer{} อิเมหิ โข ภิกฺขเว ปญฺจหิ ธมฺเมหิ สมนฺนาคโต อาวาสิโก ภิกฺขุ ยถาภตํ นิกฺขิตฺโต เอวํ นิรเย.}

\prose{ปญฺจหิ ภิกฺขเว ธมฺเมหิ สมนฺนาคโต อาวาสิโก ภิกฺขุ ยถาภตํ นิกฺขิตฺโต เอวํ สคฺเค.\csromanspacer{} กตเมหิ ปญฺจหิ? อนุวิจฺจ ปริโยคาเหตฺวา อวณฺณารหสฺส อวณฺณํ ภาสติ, อนุวิจฺจ ปริโยคาเหตฺวา วณฺณารหสฺส วณฺณํ ภาสติ, น อาวาสมจฺฉรี โหติ น อาวาสปลิเคธี, น กุลมจฺฉรี โหติ น กุลปลิเคธี, สทฺธาเทยฺยํ น วินิปาเตติ.\csromanspacer{} อิเมหิ โข ภิกฺขเว ปญฺจหิ ธมฺเมหิ สมนฺนาคโต อาวาสิโก ภิกฺขุ ยถาภตํ นิกฺขิตฺโต เอวํ สคฺเคติ.\csromanspacer{}\csromanspacer{}สตฺตมํ.}

\csromanheader{2}{8. ตติย อวณฺณารหสุตฺต}
\proseitem{238}{ปญฺจหิ ภิกฺขเว ธมฺเมหิ สมนฺนาคโต อาวาสิโก ภิกฺขุ ยถาภตํ นิกฺขิตฺโต เอวํ นิรเย.\csromanspacer{} กตเมหิ ปญฺจหิ? อนนุวิจฺจ \csromanfolio{232}อปริโยคาเหตฺวา อวณฺณารหสฺส วณฺณํ ภาสติ, อนนุวิจฺจ อปริโยคาเหตฺวา วณฺณา รหสฺส อวณฺณํ ภาสติ, อาวาสมจฺฉรี โหติ, กุลมจฺฉรี โหติ, ลาภมจฺฉรี โหติ.\csromanspacer{} อิเมหิ โข ภิกฺขเว ปญฺจหิ ธมฺเมหิ สมนฺนาคโต อาวาสิโก ภิกฺขุ ยถาภตํ นิกฺขิตฺโต เอวํ นิรเย.}

\prose{ปญฺจหิ ภิกฺขเว ธมฺเมหิ สมนฺนาคโต อาวาสิโก ภิกฺขุ ยถาภตํ นิกฺขิตฺโต เอวํ สคฺเค.\csromanspacer{} กตเมหิ ปญฺจหิ? อนุวิจฺจ ปริโยคาเหตฺวา อวณฺณารหสฺส อวณฺณํ ภาสติ, อนุวิจฺจ ปริโยคาเหตฺวา วณฺณารหสฺส วณฺณํ ภาสติ, น อาวาสมจฺฉรี โหติ, น กุลมจฺฉรี โหติ, น ลาภมจฺฉรี โหติ.\csromanspacer{} อิเมหิ โข ภิกฺขเว ปญฺจหิ ธมฺเมหิ สมนฺนาคโต อาวาสิโก ภิกฺขุ ยถาภตํ นิกฺขิตฺโต เอวํ สคฺเคติ.\csromanspacer{}\csromanspacer{}อฏฺฐมํ.}

\csromanheader{2}{9. ปฐมมจฺฉริยสุตฺต}
\proseitem{239}{ปญฺจหิ ภิกฺขเว ธมฺเมหิ สมนฺนาคโต อาวาสิโก ภิกฺขุ ยถาภตํ นิกฺขิตฺโต เอวํ นิรเย.\csromanspacer{} กตเมหิ ปญฺจหิ? อาวาสมจฺฉรี โหติ, กุลมจฺฉรี โหติ, ลาภมจฺฉรี โหติ, วณฺณมจฺฉรี\footnote{ธมฺมมจฺฉรี (ก)} โหติ, สทฺธาเทยฺยํ วินิปาเตติ.\csromanspacer{} อิเมหิ โข ภิกฺขเว ปญฺจหิ ธมฺเมหิ สมนฺนาคโต อาวาสิโก ภิกฺขุ ยถาภตํ นิกฺขิตฺโต เอวํ นิรเย.}

\prose{ปญฺจหิ ภิกฺขเว ธมฺเมหิ สมนฺนาคโต อาวาสิโก ภิกฺขุ ยถาภตํ นิกฺขิตฺโต เอวํ สคฺเค.\csromanspacer{} กตเมหิ ปญฺจหิ? น อาวาสมจฺฉรี โหติ, น กุลมจฺฉรี โหติ, น ลาภมจฺฉรี โหติ, น วณฺณมจฺฉรี โหติ, สทฺธาเทยฺยํ น วินิปาเตติ.\csromanspacer{} อิเมหิ โข ภิกฺขเว ปญฺจหิ ธมฺเมหิ สมนฺนาคโต อาวาสิโก ภิกฺขุ ยถาภตํ นิกฺขิตฺโต เอวํ สคฺเคติ.\csromanspacer{}\csromanspacer{}นวมํ.}

\csromanheader{2}{10. ทุติยมจฺฉริยสุตฺต}
\proseitem{240}{ปญฺจหิ ภิกฺขเว ธมฺเมหิ สมนฺนาคโต อาวาสิโก ภิกฺขุ ยถาภตํ นิกฺขิตฺโต เอวํ นิรเย.\csromanspacer{} กตเมหิ ปญฺจหิ? อาวาสมจฺฉรี โหติ, กุลมจฺฉรี โหติ, ลาภมจฺฉรี โหติ, วณฺณมจฺฉรี โหติ, ธมฺมมจฺฉรี โหติ.\csromanspacer{} อิเมหิ โข ภิกฺขเว ปญฺจหิ ธมฺเมหิ สมนฺนาคโต อาวาสิโก ภิกฺขุ ยถาภตํ นิกฺขิตฺโต เอวํ นิรเย.}

\csromanmatikaanchor{mtk.111}
\csromantocmarknopage{section}{(25) 5. ทุจฺจริตวคฺค}{mtk.111}
\bookmark[dest={mtk.111},level=1]{(25) 5. ทุจฺจริตวคฺค}
\csromanheader{1}{\textbf{(25) 5. ทุจฺจริตวคฺค}}
\csromanmatikaanchor{mtk.112}
\csromantocmarkref{subsection}{1-8. ทุจฺจริตสุตฺตานิ}{mtk.112}
\bookmark[dest={mtk.112},level=2]{1-8. ทุจฺจริตสุตฺตานิ}
\prose{\csromanfolio{233}ปญฺจหิ ภิกฺขเว ธมฺเมหิ สมนฺนาคโต อาวาสิโก ภิกฺขุ ยถาภตํ นิกฺขิตฺโต เอวํ สคฺเค.\csromanspacer{} กตเมหิ ปญฺจหิ? น อาวาสมจฺฉรี โหติ, น กุลมจฺฉรี โหติ, น ลาภมจฺฉรี โหติ, น วณฺณมจฺฉรี โหติ, น ธมฺมมจฺฉรี โหติ.\csromanspacer{} อิเมหิ โข ภิกฺขเว ปญฺจหิ ธมฺเมหิ สมนฺนาคโต อาวาสิโก ภิกฺขุ ยถาภตํ นิกฺขิตฺโต เอวํ สคฺเคติ.\csromanspacer{}\csromanspacer{}ทสมํ.\csromanspacer{} อาวาสิกวคฺโค จตุตฺโถ.}

\titlehead{\textbf{ตสฺสุทฺทานํ}}
\csromangathameasure{อาวาสิโก ปิโย จ โสภโน, พหูปกาโร อนุกมฺปโก จ. \\ ตโย อวณฺณารหา เจว, มจฺฉริยา ทุเวปิ จาติ.}
\csromangathagroup{\csromangathastanza{อาวาสิโก ปิโย จ โสภโน, พหูปกาโร อนุกมฺปโก จ. \\ ตโย อวณฺณารหา เจว, มจฺฉริยา ทุเวปิ จาติ\footnote{ยถาภตํ จาปิ อวณฺณเคธา, จตุกฺกมจฺเฉร ปญฺจเมน จาติ (สี, สฺยา) ยถาภตํ อวณฺณญฺจ, จตุโก มจฺฉริเยน จาติ (ก)}.}}

\chapterhead{\textbf{(25) 5. ทุจฺจริตวคฺค}}
\csromanheader{2}{1. ปฐมทุจฺจริตสุตฺต}
\proseitem{241}{ปญฺจิเม ภิกฺขเว อาทีนวา ทุจฺจริเต.\csromanspacer{} กตเม ปญฺจ? อตฺตาปิ อตฺตานํ อุปวทติ, อนุวิจฺจ วิญฺญู ครหนฺติ, ปาปโก กิตฺติสทฺโท อพฺภุคฺคจฺฉติ, สมฺมูโฬฺห กาลํ กโรติ, กายสฺส เภทา ปรํ มรณา อปายํ ทุคฺคติํ วินิปาตํ นิรยํ อุปปชฺชติ.\csromanspacer{} อิเม โข ภิกฺขเว ปญฺจ อาทีนวา ทุจฺจริเต.}

\prose{ปญฺจิเม ภิกฺขเว อานิสํสา สุจริเต.\csromanspacer{} กตเม ปญฺจ? อตฺตาปิ อตฺตานํ น อุปวทติ, อนุวิจฺจ วิญฺญู ปสํสนฺติ, กลฺยาโณ กิตฺติสทฺโท อพฺภุคฺคจฺฉติ, อสมฺมูโฬฺห กาลํ กโรติ, กายสฺส เภทา ปรํ มรณา สุคติํ สคฺคํ โลกํ อุปปชฺชติ.\csromanspacer{} อิเม โข ภิกฺขเว ปญฺจ อานิสํสา สุจริเตติ.\csromanspacer{}\csromanspacer{}ปฐมํ.}

\csromanheader{2}{2. ปฐมกายทุจฺจริตสุตฺต}
\proseitem{242}{ปญฺจิเม ภิกฺขเว อาทีนวา กายทุจฺจริเต ฯเปฯ อานิสํสา กายสุจริเต ฯเปฯ.\csromanspacer{} ทุติยํ.}

\csromanheader{2}{3. ปฐมวจีทุจฺจริตสุตฺต}
\proseitem{243}{\csromanfolio{234}ปญฺจิเม ภิกฺขเว อาทีนวา วจีทุจฺจริเต ฯเปฯ อานิสํสา วจีสุจริเต ฯเปฯ.\csromanspacer{}\csromanspacer{}ตติยํ.}

\csromanheader{2}{4. ปฐมมโนทุจฺจริตสุตฺต}
\proseitem{244}{ปญฺจิเม ภิกฺขเว อาทีนวา มโนทุจฺจริเต ฯเปฯ อานิสํสา มโนสุจริเต.\csromanspacer{} กตเม ปญฺจ? อตฺตาปิ อตฺตานํ น อุปวทติ, อนุวิจฺจ วิญฺญู ปสํสนฺติ, กลฺยาโณ กิตฺติสทฺโท อพฺภุคฺคจฺฉติ, อสมฺมูโฬฺห กาลํ กโรติ, กายสฺส เภทา ปรํ มรณา สุคติํ สคฺคํ โลกํ อุปปชฺชติ.\csromanspacer{} อิเม โข ภิกฺขเว ปญฺจ อานิสํสา มโนสุจริเตติ.\csromanspacer{}\csromanspacer{}จตุตฺถํ.}

\csromanheader{2}{5. ทุติยทุจฺจริตสุตฺต}
\proseitem{245}{ปญฺจิเม ภิกฺขเว อาทีนวา ทุจฺจริเต.\csromanspacer{} กตเม ปญฺจ? อตฺตาปิ อตฺตานํ อุปวทติ, อนุวิจฺจ วิญฺญู ครหนฺติ, ปาปโก กิตฺติสทฺโท อพฺภุคฺคจฺฉติ, สทฺธมฺมา วุฏฺฐาติ, อสทฺธมฺเม ปติฏฺฐาติ.\csromanspacer{} อิเม โข ภิกฺขเว ปญฺจ อาทีนวา ทุจฺจริเต.}

\prose{ปญฺจิเม ภิกฺขเว อานิสํสา สุจริเต.\csromanspacer{} กตเม ปญฺจ? อตฺตาปิ อตฺตานํ น อุปวทติ, อนุวิจฺจ วิญฺญู ปสํสนฺติ, กลฺยาโณ กิตฺติสทฺโท อพฺภุคฺคจฺฉติ, อสทฺธมฺมา วุฏฺฐาติ, สทฺธมฺเม ปติฏฺฐาติ.\csromanspacer{} อิเม โข ภิกฺขเว ปญฺจ อานิสํสา สุจริเตติ.\csromanspacer{}\csromanspacer{}ปญฺจมํ.}

\csromanheader{2}{6. ทุติยกายทุจฺจริตสุตฺต}
\proseitem{246}{ปญฺจิเม ภิกฺขเว อาทีนวา กายทุจฺจริเต ฯเปฯ อานิสํสา กายสุจริเต ฯเปฯ.\csromanspacer{} ฉฏฺฐํ.}

\csromanheader{2}{7. ทุติยวจีทุจฺจริตสุตฺต}
\proseitem{247}{ปญฺจิเม ภิกฺขเว อาทีนวา วจีทุจฺจริเต ฯเปฯ อานิสํสา วจีสุจริเต ฯเปฯ.\csromanspacer{}\csromanspacer{}สตฺตมํ.}

\chapterheadpageread{(25) 5. ทุจฺจริตวคฺค \textbf{8. ทุติยมโนทุจฺจริตสุตฺต}}
\proseitem{248}{\csromanfolio{235}ปญฺจิเม ภิกฺขเว อาทีนวา มโนทุจฺจริเต ฯเปฯ อานิสํสา มโนสุจริเต.\csromanspacer{} กตเม ปญฺจ? อตฺตาปิ อตฺตานํ น อุปวทติ, อนุวิจฺจ วิญฺญู ปสํสนฺติ, กลฺยาโณ กิตฺติสทฺโท อพฺภุคฺคจฺฉติ, อสทฺธมฺมา วุฏฺฐาติ, สทฺธมฺเม ปติฏฺฐาติ.\csromanspacer{} อิเม โข ภิกฺขเว ปญฺจ อานิสํสา มโนสุจริเตติ.\csromanspacer{}\csromanspacer{}อฏฺฐมํ.}

\csromanmatikaanchor{mtk.113}
\csromantocmarkref{subsection}{9. สิวถิกสุตฺต}{mtk.113}
\bookmark[dest={mtk.113},level=2]{9. สิวถิกสุตฺต}
\csromanheader{2}{9. สิวถิกสุตฺต}
\proseitem{249}{ปญฺจิเม ภิกฺขเว อาทีนวา สิวถิกาย\footnote{สีวถิกาย (สี, สฺยา, กํ, อิ)}.\csromanspacer{} กตเม ปญฺจ? อสุจิ, ทุคฺคนฺธา, สปฺปฏิภยา, วาฬานํ อมนุสฺสานํ อาวาโส, พหุโน ชนสฺส อาโรทนา.\csromanspacer{} อิเม โข ภิกฺขเว ปญฺจ อาทีนวา สิวถิกาย.}

\prose{เอวเมวํ โข ภิกฺขเว ปญฺจิเม อาทีนวา สิวถิกูปเม ปุคฺคเล.\csromanspacer{} กตเม ปญฺจ? อิธ ภิกฺขเว เอกจฺโจ ปุคฺคโล อสุจินา กายกมฺเมน สมนฺนาคโต โหติ, อสุจินา วจีกมฺเมน สมนฺนาคโต โหติ, อสุจินา มโนกมฺเมน สมนฺนาคโต โหติ.\csromanspacer{} อิทมสฺส อสุจิตาย วทามิ.\csromanspacer{} เสยฺยถาปิ สา ภิกฺขเว สิวถิกา อสุจิ, ตถูปมาหํ ภิกฺขเว อิมํ ปุคฺคลํ วทามิ.}

\prose{ตสฺส อสุจินา กายกมฺเมน สมนฺนาคตสฺส อสุจินา วจีกมฺเมน สมนฺนาคตสฺส อสุจินา มโนกมฺเมน สมนฺนาคตสฺส ปาปโก กิตฺติสทฺโท อพฺภุคฺคจฺฉติ.\csromanspacer{} อิทมสฺส ทุคฺคนฺธตาย วทามิ.\csromanspacer{} เสยฺยถาปิ สา ภิกฺขเว สิวถิกา ทุคฺคนฺธา, ตถูปมาหํ ภิกฺขเว อิมํ ปุคฺคลํ วทามิ.}

\prose{ตเมนํ อสุจินา กายกมฺเมน สมนฺนาคตํ อสุจินา วจีกมฺเมน สมนฺนาคตํ อสุจินา มโนกมฺเมน สมนฺนาคตํ เปสลา สพฺรหฺมจารี อารกา ปริวชฺชนฺติ.\csromanspacer{} อิทมสฺส สปฺปฏิภยสฺมิํ วทามิ.\csromanspacer{} เสยฺยถาปิ สา ภิกฺขเว สิวถิกา สปฺปฏิภยา, ตถูปมาหํ ภิกฺขเว อิมํ ปุคฺคลํ วทามิ.}

\prose{โส อสุจินา กายกมฺเมน สมนฺนาคโต อสุจินา วจีกมฺเมน สมนฺนาคโต อสุจินา มโนกมฺเมน สมนฺนาคโต สภาเคหิ ปุคฺคเลหิ สทฺธิํ สํวสติ.\csromanspacer{} อิทมสฺส วาฬาวาสสฺมิํ วทามิ.\csromanspacer{} เสยฺยถาปิ \csromanfolio{236}สา ภิกฺขเว สิวถิกา วาฬานํ อมนุสฺสานํ อาวาโส, ตถูปมาหํ ภิกฺขเว อิมํ ปุคฺคลํ วทามิ.}

\prose{ตเมนํ อสุจินา กายกมฺเมน สมนฺนาคตํ อสุจินา วจีกมฺเมน สมนฺนาคตํ อสุจินา มโนกมฺเมน สมนฺนาคตํ เปสลา สพฺรหฺมจารี ทิสฺวา ขียธมฺมํ\footnote{ขียนธมฺมํ (สี)} อาปชฺชนฺติ “อโห วต โน ทุกฺขํ, เย มยํ เอวรูเปหิ ปุคฺคเลหิ สทฺธิํ สํวสามา”ติ.\csromanspacer{} อิทมสฺส อาโรทนาย วทามิ.\csromanspacer{} เสยฺยถาปิ สา ภิกฺขเว สิวถิกา พหุโน ชนสฺส อาโรทนา, ตถูปมาหํ ภิกฺขเว อิมํ ปุคฺคลํ วทามิ.\csromanspacer{} อิเม โข ภิกฺขเว ปญฺจ อาทีนวา สิวถิกูปเม ปุคฺคเลติ.\csromanspacer{}\csromanspacer{}นวมํ.}

\csromanmatikaanchor{mtk.114}
\csromantocmarkref{subsection}{10. ปุคฺคลปฺปสาทสุตฺต}{mtk.114}
\bookmark[dest={mtk.114},level=2]{10. ปุคฺคลปฺปสาทสุตฺต}
\csromanheader{2}{10. ปุคฺคลปฺปสาทสุตฺต}
\proseitem{250}{ปญฺจิเม ภิกฺขเว อาทีนวา ปุคฺคลปฺปสาเท.\csromanspacer{} กตเม ปญฺจ? ยสฺมิํ ภิกฺขเว ปุคฺคเล ปุคฺคโล อภิปฺปสนฺโน โหติ, โส ตถารูปํ อาปตฺติํ อาปนฺโน โหติ, ยถารูปาย อาปตฺติยา สํโฆ อุกฺขิปติ.\csromanspacer{} ตสฺส เอวํ โหติ “โย โข มฺยายํ ปุคฺคโล ปิโย มนาโป โส สํเฆน อุกฺขิตฺโต”ติ, ภิกฺขูสุ อปฺปสาทพหุโล โหติ.\csromanspacer{} ภิกฺขูสุ อปฺปสาทพหุโล สมาโน อญฺเญ ภิกฺขู น ภชติ, อญฺเญ ภิกฺขู อภชนฺโต สทฺธมฺมํ น สุณาติ, สทฺธมฺมํ อสุณนฺโต สทฺธมฺมา ปริหายติ.\csromanspacer{} อยํ ภิกฺขเว ปฐโม อาทีนโว ปุคฺคลปฺปสาเท.}

\prose{ปุน จปรํ ภิกฺขเว ยสฺมิํ ปุคฺคเล ปุคฺคโล อภิปฺปสนฺโน โหติ, โส ตถารูปํ อาปตฺติํ อาปนฺโน โหติ, ยถารูปาย อาปตฺติยา สํโฆ อนฺเต นิสีทาเปติ.\csromanspacer{} ตสฺส เอวํ โหติ “โย โข มฺยายํ ปุคฺคโล ปิโย มนาโป, โส สํเฆน อนฺเต นิสีทาปิโต”ติ, ภิกฺขูสุ อปฺปสาทพหุโล โหติ.\csromanspacer{} ภิกฺขูสุ อปฺปสาทพหุโล สมาโน อญฺเญ ภิกฺขู น ภชติ, อญฺเญ ภิกฺขู อภชนฺโต สทฺธมฺมํ น สุณาติ, สทฺธมฺมํ อสุณนฺโต สทฺธมฺมา ปริหายติ.\csromanspacer{} อยํ ภิกฺขเว ทุติโย อาทีนโว ปุคฺคลปฺปสาเท.}

\prose{ปุน จปรํ ภิกฺขเว ยสฺมิํ ปุคฺคเล ปุคฺคโล อภิปฺปสนฺโน โหติ, โส ทิสาปกฺกนฺโต โหติ ฯเปฯ โส วิพฺภนฺโต โหติ ฯเปฯ.}

\vspace{\tipitakaparskip}
\csromancenter{(25) 5.\csromanspacer{} ทุจฺจริตวคฺค โส กาลงฺกโต โหติ.\csromanspacer{} ตสฺส เอวํ โหติ “โย โข มฺยายํ ปุคฺคโล ปิโย มนาโป, โส กาลงฺกโต”ติ, อญฺเญ ภิกฺขู น ภชติ, อญฺเญ ภิกฺขู อภชนฺโต สทฺธมฺมํ น สุณาติ, สทฺธมฺมํ อสุณนฺโต สทฺธมฺมา ปริหายติ, อยํ ภิกฺขเว ปญฺจโม อาทีนโว ปุคฺคลปฺปสาเท.\csromanspacer{} อิเม โข ภิกฺขเว ปญฺจ อาทีนวา ปุคฺคลปฺปสาเทติ.\csromanspacer{}\csromanspacer{}ทสมํ.\csromanspacer{} ทุจฺจริตวคฺโค ปญฺจโม.}
\titlehead{\textbf{ตสฺสุทฺทานํ}}
\csromangathameasure{ทุจฺจริตํ กายทุจฺจริตํ, วจีทุจฺจริตํ มโนทุจฺจริตํ. \\ จตูหิ ปเร เทฺว สิวถิกา, ปุคฺคลปฺปสาเทน จาติ. \\ \textbf{ปญฺจโม ปณฺณาสโก สมตฺโต.}}
\csromangathagroup{\csromangathastanza{\csromanfolio{237}ทุจฺจริตํ กายทุจฺจริตํ, วจีทุจฺจริตํ มโนทุจฺจริตํ. \\ จตูหิ ปเร เทฺว สิวถิกา, ปุคฺคลปฺปสาเทน จาติ. \\ \textbf{ปญฺจโม ปณฺณาสโก สมตฺโต.}}}

\csromanmatikaanchor{mtk.115}
\csromantocmarknopage{section}{(26) 6. อุปสมฺปทาวคฺค}{mtk.115}
\bookmark[dest={mtk.115},level=1]{(26) 6. อุปสมฺปทาวคฺค}
\csromanheader{1}{\textbf{(26) 6. อุปสมฺปทาวคฺค}}
\csromanheader{2}{1. อุปสมฺปาเทตพฺพสุตฺต}
\csromanmatikaanchor{mtk.116}
\csromanmatikaanchor{mtk.117}
\csromantocmarkref{subsection}{1-3. อุปสมฺปาเทตพฺพสุตฺตาทิ}{mtk.116}
\bookmark[dest={mtk.116},level=2]{1-3. อุปสมฺปาเทตพฺพสุตฺตาทิ}
\csromantocmarkref{subsection}{4-5. มจฺฉริยสุตฺตทฺวย}{mtk.117}
\bookmark[dest={mtk.117},level=2]{4-5. มจฺฉริยสุตฺตทฺวย}
\proseitem{251}{\csromanfolio{238}\csromansymbolfootnote{*}{วิ 3. 90 ปิฏฺเฐปิ.}ปญฺจหิ ภิกฺขเว ธมฺเมหิ สมนฺนาคเตน ภิกฺขุนา อุปสมฺปาเทตพฺพํ.\csromanspacer{} กตเมหิ ปญฺจหิ? อิธ ภิกฺขเว ภิกฺขุ อเสเขน สีลกฺขนฺเธน สมนฺนาคโต โหติ, อเสเขน สมาธิกฺขนฺเธน สมนฺนาคโต โหติ, อเสเขน ปญฺญากฺขนฺเธน สมนฺนาคโต โหติ, อเสเขน วิมุตฺติกฺขนฺเธน สมนฺนาคโต โหติ, อเสเขน วิมุตฺติญาณทสฺสนกฺขนฺเธน สมนฺนาคโต โหติ.\csromanspacer{} อิเมหิ โข ภิกฺขเว ปญฺจหิ ธมฺเมหิ สมนฺนาคเตน ภิกฺขุนา อุปสมฺปาเทตพฺพนฺติ.\csromanspacer{}\csromanspacer{}ปฐมํ.}

\csromanheader{2}{2. นิสฺสยสุตฺต}
\proseitem{252}{\csromansymbolmark{*}ปญฺจหิ ภิกฺขเว ธมฺเมหิ สมนฺนาคเตน ภิกฺขุนา นิสฺสโย ทาตพฺโพ.\csromanspacer{} กตเมหิ ปญฺจหิ? อิธ ภิกฺขเว ภิกฺขุ อเสเขน สีลกฺขนฺเธน สมนฺนาคโต โหติ ฯเปฯ อเสเขน วิมุตฺติญาณทสฺสนกฺขนฺเธน สมนฺนาคโต โหติ.\csromanspacer{} อิเมหิ ฯเปฯ นิสฺสโย ทาตพฺโพติ.\csromanspacer{}\csromanspacer{}ทุติยํ.}

\csromanheader{2}{3. สามเณรสุตฺต}
\proseitem{253}{\csromansymbolmark{*}ปญฺจหิ ภิกฺขเว ธมฺเมหิ สมนฺนาคเตน ภิกฺขุนา สามเณโร อุปฏฺฐาเปตพฺโพ.\csromanspacer{} กตเมหิ ปญฺจหิ? อิธ ภิกฺขเว ภิกฺขุ อเสเขน สีลกฺขนฺเธน สมนฺนาคโต โหติ, อเสเขน สมาธิกฺขนฺเธน.\csromanspacer{} อเสเขน ปญฺญากฺขนฺเธน.\csromanspacer{} อเสเขน วิมุตฺติกฺขนฺเธน.\csromanspacer{} อเสเขน วิมุตฺติญาณทสฺสนกฺขนฺเธน สมนฺนาคโต โหติ.\csromanspacer{} อิเมหิ โข ภิกฺขเว ปญฺจหิ ธมฺเมหิ สมนฺนาคเตน ภิกฺขุนา สามเณโร อุปฏฺฐาเปตพฺโพติ.\csromanspacer{}\csromanspacer{}ตติยํ.}

\csromanheader{2}{4. ปญฺจมจฺฉริยสุตฺต}
\proseitem{254}{ปญฺจิมานิ ภิกฺขเว มจฺฉริยานิ.\csromanspacer{} กตมานิ ปญฺจ? อาวาสมจฺฉริยํ กุลมจฺฉริยํ ลาภมจฺฉริยํ วณฺณมจฺฉริยํ ธมฺมมจฺฉริยํ.\csromanspacer{} อิมานิ โข ภิกฺขเว ปญฺจ มจฺฉริยานิ.\csromanspacer{} อิเมสํ โข ภิกฺขเว ปญฺจนฺนํ มจฺฉริยานํ เอตํ ปฏิกุฏฺฐํ\footnote{ปติกิฏฺฐํ (สี, อิ), ปฏิกฺกิฏฺฐํ (สฺยา, กํ), ปฏิกิฏฺฐํ (ก)}, ยทิทํ ธมฺมมจฺฉริยนฺติ.\csromanspacer{}\csromanspacer{}จตุตฺถํ.}

\chapterheadpageread{(26) 6. อุปสมฺปทาวคฺค}
\csromanmatikaanchor{mtk.118}
\csromantocmarkref{subsection}{6-21. ปฐมฌานสุตฺตาทิ}{mtk.118}
\bookmark[dest={mtk.118},level=2]{6-21. ปฐมฌานสุตฺตาทิ}
\proseitem{255}{\csromanfolio{239}ปญฺจนฺนํ ภิกฺขเว มจฺฉริยานํ ปหานาย สมุจฺเฉทาย พฺรหฺมจริยํ วุสฺสติ.\csromanspacer{} กตเมสํ ปญฺจนฺนํ? อาวาสมจฺฉริยสฺส ปหานาย สมุจฺเฉทาย พฺราหฺมจริยํ วุสฺสติ, กุลมจฺฉริยสฺส.\csromanspacer{} ลาภมจฺฉริยสฺส.\csromanspacer{} วณฺณมจฺฉริยสฺส.\csromanspacer{} ธมฺมมจฺฉริยสฺส ปหานาย สมุจฺเฉทาย พฺรหฺมจริยํ วุสฺสติ.\csromanspacer{} อิเมสํ โข ภิกฺขเว ปญฺจนฺนํ มจฺฉริยานํ ปหานาย สมุจฺเฉทาย พฺรหฺมจริยํ วุสฺสตีติ.\csromanspacer{}\csromanspacer{}ปญฺจมํ.}

\csromanheader{2}{6. ปฐมฌานสุตฺต}
\proseitem{256}{ปญฺจิเม ภิกฺขเว ธมฺเม อปฺปหาย อภพฺโพ ปฐมํ ฌานํ อุปสมฺปชฺช วิหริตุํ.\csromanspacer{} กตเม ปญฺจ? อาวาสมจฺฉริยํ กุลมจฺฉริยํ ลาภมจฺฉริยํ วณฺณมจฺฉริยํ ธมฺมมจฺฉริยํ.\csromanspacer{} อิเม โข ภิกฺขเว ปญฺจ ธมฺเม อปฺปหาย อภพฺโพ ปฐมํ ฌานํ อุปสมฺปชฺช วิหริตุํ.}

\prose{ปญฺจิเม ภิกฺขเว ธมฺเม ปหาย ภพฺโพ ปฐมํ ฌานํ อุปสมฺปชฺช วิหริตุํ.\csromanspacer{} กตเม ปญฺจ? อาวาสมจฺฉริยํ กุลมจฺฉริยํ ลาภมจฺฉริยํ วณฺณมจฺฉริยํ ธมฺมมจฺฉริยํ.\csromanspacer{} อิเม โข ภิกฺขเว ปญฺจ ธมฺเม ปหาย ภพฺโพ ปฐมํ ฌานํ อุปสมฺปชฺช วิหริตุนฺติ.\csromanspacer{}\csromanspacer{}ฉฏฺฐํ.}

\csromanheader{2}{7-13. ทุติยฌานสุตฺตาทิสตฺตก}
\proseitem{257-263}{ปญฺจิเม ภิกฺขเว ธมฺเม อปฺปหาย อภพฺโพ ทุติยํ ฌานํ ฯเปฯ อภพฺโพ ตติยํ ฌานํ.\csromanspacer{} อภพฺโพ จตุตฺถํ ฌานํ.\csromanspacer{} อภพฺโพ โสตาปตฺติผลํ.\csromanspacer{} อภพฺโพ สกทาคามิผลํ.\csromanspacer{} อภพฺโพ อนาคามิผลํ.\csromanspacer{} อภพฺโพ อรหตฺตํ\footnote{อรหตฺตผลํ (สี)} สจฺฉิกาตุํ.\csromanspacer{} กตเม ปญฺจ? อาวาสมจฺฉริยํ กุลมจฺฉริยํ ลาภม จฺฉริยํ วณฺณมจฺฉริยํ ธมฺมมจฺฉริยํ.\csromanspacer{} อิเม โข ภิกฺขเว ปญฺจ ธมฺเม อปฺปหาย อภพฺโพ อรหตฺตํ สจฺฉิกาตุํ.}

\prose{ปญฺจิเม ภิกฺขเว ธมฺเม ปหาย ภพฺโพ ทุติยํ ฌานํ ฯเปฯ ภพฺโพ ตติยํ ฌานํ.\csromanspacer{} ภพฺโพ จตุตฺถํ ฌานํ.\csromanspacer{} ภพฺโพ โสตาปตฺติผลํ.\csromanspacer{} ภพฺโพ สกทาคามิผลํ.\csromanspacer{} ภพฺโพ อนาคามิผลํ.\csromanspacer{} ภพฺโพ อรหตฺตํ สจฺฉิกาตุํ.\csromanspacer{} กตเม ปญฺจ? อาวาสมจฺฉริยํ กุลมจฺฉริยํ ลาภมจฺฉริยํ วณฺณมจฺฉริยํ ธมฺมมจฺฉริยํ.\csromanspacer{} อิเม โข ภิกฺขเว ปญฺจ ธมฺเม ปหาย ภพฺโพ อรหตฺตํ สจฺฉิกาตุนฺติ.\csromanspacer{}\csromanspacer{}เตรสมํ.}

\csromanheader{2}{14. อปรปฐมฌานสุตฺต}
\csromanmatikaanchor{mtk.119}
\csromanmatikaanchor{mtk.120}
\csromantocmarknopage{section}{(27) 1. สมฺมุติเปยฺยาล}{mtk.119}
\bookmark[dest={mtk.119},level=1]{(27) 1. สมฺมุติเปยฺยาล}
\csromantocmarkref{subsection}{1-14. ภตฺตุทฺเทสกสุตฺตาทิ}{mtk.120}
\bookmark[dest={mtk.120},level=2]{1-14. ภตฺตุทฺเทสกสุตฺตาทิ}
\proseitem{264}{\csromanfolio{240}ปญฺจิเม ภิกฺขเว ธมฺเม อปฺปหาย อภพฺโพ ปฐมํ ฌานํ อุปสมฺปชฺช วิหริตุํ.\csromanspacer{} กตเม ปญฺจ? อาวาสมจฺฉริยํ กุลมจฺฉริยํ ลาภมจฺฉริยํ วณฺณมจฺฉริยํ อกตญฺญุตํ อกตเวทิตํ.\csromanspacer{} อิเม โข ภิกฺขเว ปญฺจ ธมฺเม อปฺปหาย อภพฺโพ ปฐมํ ฌานํ อุปสมฺปชฺช วิหริตุํ.}

\prose{ปญฺจิเม ภิกฺขเว ธมฺเม ปหาย ภพฺโพ ปฐมํ ฌานํ อุปสมฺปชฺช วิหริตุํ.\csromanspacer{} กตเม ปญฺจ? อาวาสมจฺฉริยํ กุลมจฺฉริยํ ลาภมจฺฉริยํ วณฺณมจฺฉริยํ อกตญฺญุตํ อกตเวทิตํ.\csromanspacer{} อิเม โข ภิกฺขเว ปญฺจ ธมฺเม ปหาย ภพฺโพ ปฐมํ ฌานํ อุปสมฺปชฺช วิหริตุนฺติ.\csromanspacer{}\csromanspacer{}จุทฺทสมํ.}

\csromanheader{2}{15-21. อปรทุติยฌานสุตฺตาทิ}
\proseitem{265-271}{ปญฺจิเม ภิกฺขเว ธมฺเม อปฺปหาย อภพฺโพ ทุติยํ ฌานํ.\csromanspacer{} ตติยํ ฌานํ.\csromanspacer{} จตุตฺถํ ฌานํ.\csromanspacer{} โสตาปตฺติผลํ.\csromanspacer{} สกทาคามิผลํ.\csromanspacer{} อนาคามิผลํ.\csromanspacer{} อรหตฺตํ สจฺฉิกาตุํ.\csromanspacer{} กตเม ปญฺจ? อาวาสมจฺฉริยํ กุลมจฺฉริยํ ลาภมจฺฉริยํ วณฺณมจฺฉริยํ อกตญฺญุตํ อกตเวทิตํ.\csromanspacer{} อิเม โข ภิกฺขเว ปญฺจ ธมฺเม อปฺปหาย อภพฺโพ อรหตฺตํ สจฺฉิกาตุํ.}

\prose{ปญฺจิเม ภิกฺขเว ธมฺเม ปหาย ภพฺโพ ทุติยํ ฌานํ.\csromanspacer{} ตติยํ ฌานํ.\csromanspacer{} จตุตฺถํ ฌานํ.\csromanspacer{} โสตาปตฺติผลํ.\csromanspacer{} สกทาคามิผลํ.\csromanspacer{} อนาคามิผลํ.\csromanspacer{} อรหตฺตํ สจฺฉิกาตุํ.\csromanspacer{} กตเม ปญฺจ? อาวาสมจฺฉริยํ กุลมจฺฉริยํ ลาภมจฺฉริยํ วณฺณมจฺฉริยํ อกตญฺญุตํ อกตเวทิตํ.\csromanspacer{} อิเม โข ภิกฺขเว ปญฺจ ธมฺเม ปหาย ภพฺโพ อรหตฺตํ สจฺฉิกาตุนฺติ.\csromanspacer{}\csromanspacer{}เอกวีสติมํ.\csromanspacer{} อุปสมฺปทาวคฺโค ฉฏฺโฐ.}

\csromanheader{1}{1. สมฺมุติเปยฺยาล}
\csromanheader{2}{1. ภตฺตุทฺเทสกสุตฺต}
\proseitem{272}{ปญฺจหิ ภิกฺขเว ธมฺเมหิ สมนฺนาคโต ภตฺตุทฺเทสโก น สมฺมนฺนิตพฺโพ\footnote{น สมฺมนิตพฺโพ (ก) วิ 4. 330 ปิฏฺเฐปิ ปสฺสิตพฺพํ.}.\csromanspacer{} กตเมหิ ปญฺจหิ? ฉนฺทาคติํ คจฺฉติ, โทสาคติํ คจฺฉติ.}

\vspace{\tipitakaparskip}
\csromancenter{สมฺมุติเปยฺยาล โมหาคติํ คจฺฉติ, ภยาคติํ คจฺฉติ, อุทฺทิฏฺฐานุทฺทิฏฺฐํ น ชานาติ.\csromanspacer{} อิเมหิ โข ภิกฺขเว ปญฺจหิ ธมฺเมหิ สมนฺนาคโต ภตฺตุทฺเทสโก น สมฺมนฺนิตพฺโพ.}
\prose{\csromanfolio{241}ปญฺจหิ ภิกฺขเว ธมฺเมหิ สมนฺนาคโต ภตฺตุทฺเทสโก สมฺมนฺนิตพฺโพ.\csromanspacer{} กตเมหิ ปญฺจหิ? น ฉนฺทาคติํ คจฺฉติ, น โทสาคติํ คจฺฉติ, น โมหาคติํ คจฺฉติ, น ภยาคติํ คจฺฉติ, อุทฺทิฏฺฐานุทฺทิฏฺฐํ ชานาติ.\csromanspacer{} อิเมหิ โข ภิกฺขเว ปญฺจหิ ธมฺเมหิ สมนฺนาคโต ภตฺตุทฺเทสโก สมฺมนฺนิตพฺโพติ.}

\prose{ปญฺจหิ ภิกฺขเว ธมฺเมหิ สมนฺนาคโต ภตฺตุทฺเทสโก สมฺมโต\footnote{สมฺมโตปิ (สี)} น เปเสตพฺโพ ฯเปฯ สมฺมโต เปเสตพฺโพ.\csromanspacer{} พาโล เวทิตพฺโพ.\csromanspacer{} ปณฺฑิโต เวทิตพฺโพ.\csromanspacer{} ขตํ อุปหตํ อตฺตานํ ปริหรติ.\csromanspacer{} อกฺขตํ อนุปหตํ อตฺตานํ ปริหรติ.\csromanspacer{} ยถาภตํ นิกฺขิตฺโต เอวํ นิรเย.\csromanspacer{} ยถาภตํ นิกฺขิตฺโต เอวํ สคฺเค.\csromanspacer{} กตเมหิ ปญฺจหิ? น ฉนฺทาคติํ คจฺฉติ, น โทสาคติํ คจฺฉติ, น โมหาคติํ คจฺฉติ, น ภยาคติํ คจฺฉติ, อุทฺทิฏฺฐานุทฺทิฏฺฐํ ชานาติ.\csromanspacer{} อิเมหิ โข ภิกฺขเว ปญฺจหิ ธมฺเมหิ สมนฺนาคโต ภตฺตุทฺเทสโก ยถาภตํ นิกฺขิตฺโต เอวํ สคฺเคติ.\csromanspacer{}\csromanspacer{}ปฐมํ.}

\csromanheader{2}{2-14. เสนาสนปญฺญาปกสุตฺตาทิ}
\proseitem{273-285}{ปญฺจหิ ภิกฺขเว ธมฺเมหิ สมนฺนาคโต เสนาสนปญฺญาปโก น สมฺมนฺนิตพฺโพ ฯเปฯ ปญฺญตฺตาปญฺญตฺตํ น ชานาติ ฯเปฯ เสนาสนปญฺญาปโก สมฺมนฺนิตพฺโพ ฯเปฯ ปญฺญตฺตาปญฺญตฺตํ ชานาติ ฯเปฯ.}

\prose{เสนาสนคาหาปโก น สมฺมนฺนิตพฺโพ ฯเปฯ คหิตาคหิตํ\footnote{ปญฺญตฺตาปญฺญตฺตํ (สี, สฺยา, กํ)} น ชานาติ ฯเปฯ เสนาสนคาหาปโก สมฺมนฺนิตพฺโพ ฯเปฯ คหิตาคหิตํ2 ชานาติ ฯเปฯ.}

\prose{ภณฺฑาคาริโก น สมฺมนฺนิตพฺโพ ฯเปฯ คุตฺตาคุตฺตํ น ชานาติ.\csromanspacer{} ภณฺฑาคาริโก สมฺมนฺนิตพฺโพ ฯเปฯ คุตฺตาคุตฺตํ ชานาติ.}

\prose{จีวรปฏิคฺคาหโก น สมฺมนฺนิตพฺโพ ฯเปฯ คหิตาคหิตํ น ชานาติ.\csromanspacer{} จีวรปฏิคฺคาหโก สมฺมนฺนิตพฺโพ ฯเปฯ คหิตาคหิตํ ชานาติ.}

\prose{\csromanfolio{242}จีวรภาชโก น สมฺมนฺนิตพฺโพ ฯเปฯ ภาชิตาภาชิตํ น ชานาติ.\csromanspacer{} จีวรภาชโก สมฺมนฺนิตพฺโพ ฯเปฯ ภาชิตาภาชิตํ ชานาติ.}

\prose{ยาคุภาชโก น สมฺมนฺนิตพฺโพ ฯเปฯ ยาคุภาชโก สมฺมนฺนิตพฺโพ ฯเปฯ.}

\prose{ผลภาชโก น สมฺมนฺนิตพฺโพ ฯเปฯ ผลภาชโก สมฺมนฺนิตพฺโพ ฯเปฯ.}

\prose{ขชฺชกภาชโก น สมฺมนฺนิตพฺโพ ฯเปฯ ภาชิตาภาชิตํ น ชานาติ.\csromanspacer{} ขชฺชกภาชโก สมฺมนฺนิตพฺโพ ฯเปฯ ภาชิตาภาชิตํ ชานาติ.}

\prose{อปฺปมตฺตกวิสฺสชฺชโก น สมฺมนฺนิตพฺโพ ฯเปฯ วิสฺสชฺชิตาวิสฺสชฺชิตํ น ชานาติ.\csromanspacer{} อปฺปมตฺตกวิสฺสชฺชโก สมฺมนฺนิตพฺโพ ฯเปฯ วิสฺสชฺชิตาวิสฺสชฺชิตํ ชานาติ.}

\prose{สาฏิยคฺคาหาปโก น สมฺมนฺนิตพฺโพ ฯเปฯ คหิตาคหิตํ น ชานาติ.\csromanspacer{} สาฏิยคฺคาหาปโก สมฺมนฺนิตพฺโพ ฯเปฯ คหิตาคหิตํ ชานาติ.}

\prose{ปตฺตคฺคาหาปโก น สมฺมนฺนิตพฺโพ ฯเปฯ คหิตาคหิตํ น ชานาติ.\csromanspacer{} ปตฺตคฺคาหาปโก สมฺมนฺนิตพฺโพ ฯเปฯ คหิตาคหิตํ ชานาติ.}

\prose{อารามิกเปสโก น สมฺมนฺนิตพฺโพ ฯเปฯ อารามิกเปสโก สมฺมนฺนิตพฺโพ ฯเปฯ.}

\prose{สามเณรเปสโก น สมฺมนฺนิตพฺโพ ฯเปฯ สามเณรเปสโก สมฺมนฺนิตพฺโพ ฯเปฯ.}

\prose{สมฺมโต น เปเสตพฺโพ ฯเปฯ สมฺมโต เปเสตพฺโพ ฯเปฯ.}

\prose{สามเณรเปสโก พาโล เวทิตพฺโพ ฯเปฯ ปณฺฑิโต เวทิตพฺโพ.\csromanspacer{} ขตํ อุปหตํ อตฺตานํ ปริหรติ.\csromanspacer{} อกฺขตํ อนุปหตํ อตฺตานํ ปริหรติ.\csromanspacer{} ยถาภตํ นิกฺขิตฺโต เอวํ นิรเย.\csromanspacer{} ยถาภตํ นิกฺขิตฺโต เอวํ สคฺเค.\csromanspacer{} กตเมหิ ปญฺจหิ? น ฉนฺทาคติํ คจฺฉติ, น โทสาคติํ คจฺฉติ, น โมหาคติํ คจฺฉติ, น ภยาคติํ คจฺฉติ, เปสิตาเปสิตํ ชานาติ.\csromanspacer{} อิเมหิ โข ภิกฺขเว ปญฺจหิ ธมฺเมหิ สมนฺนาคโต สามเณรเปสโก ยถาภตํ นิกฺขิตฺโต เอวํ สคฺเคติ.\csromanspacer{}\csromanspacer{}จุทฺทสมํ.\csromanspacer{} สมฺมุติเปยฺยาลํ นิฏฺฐิตํ.}

\csromanmatikaanchor{mtk.121}
\csromantocmarknopage{section}{(28) 2. สิกฺขาปทเปยฺยาล}{mtk.121}
\bookmark[dest={mtk.121},level=1]{(28) 2. สิกฺขาปทเปยฺยาล}
\csromanheader{1}{2. สิกฺขาปทเปยฺยาล \textbf{2. สิกฺขาปทเปยฺยาล}}
\csromanheader{2}{1. ภิกฺขุสุตฺต}
\csromanmatikaanchor{mtk.122}
\csromantocmarkref{subsection}{1-17. ภิกฺขุสุตฺตาทิ}{mtk.122}
\bookmark[dest={mtk.122},level=2]{1-17. ภิกฺขุสุตฺตาทิ}
\proseitem{286}{\csromanfolio{243}ปญฺจหิ ภิกฺขเว ธมฺเมหิ สมนฺนาคโต ภิกฺขุ ยถาภตํ นิกฺขิตฺโต เอวํ นิรเย.\csromanspacer{} กตเมหิ ปญฺจหิ? ปาณาติปาตี โหติ, อทินฺนาทายี โหติ, อพฺรหฺมจารี โหติ, มุสาวาที โหติ, สุราเมรยมชฺชปมาทฏฺฐายี โหติ.\csromanspacer{} อิเมหิ โข ภิกฺขเว ปญฺจหิ ธมฺเมหิ สมนฺนาคโต ภิกฺขุ ยถาภตํ นิกฺขิตฺโต เอวํ นิรเย.}

\prose{ปญฺจหิ ภิกฺขเว ธมฺเมหิ สมนฺนาคโต ภิกฺขุ ยถาภตํ นิกฺขิตฺโต เอวํ สคฺเค.\csromanspacer{} กตเมหิ ปญฺจหิ? ปาณาติปาตา ปฏิวิรโต โหติ, อทินฺนาทานา ปฏิวิรโต โหติ, อพฺรหฺมจริยา ปฏิวิรโต โหติ, มุสาวาทา ปฏิวิรโต โหติ, สุราเมรยมชฺชปมาทฏฺฐานา ปฏิวิรโต โหติ.\csromanspacer{} อิเมหิ โข ภิกฺขเว ปญฺจหิ ธมฺเมหิ สมนฺนาคโต ภิกฺขุ ยถาภตํ นิกฺขิตฺโต เอวํ สคฺเคติ.\csromanspacer{}\csromanspacer{}ปฐมํ.}

\csromanheader{2}{2-7. ภิกฺขุนีสุตฺตาทิ}
\proseitem{287-292}{ปญฺจหิ ภิกฺขเว ธมฺเมหิ สมนฺนาคตา ภิกฺขุนี ฯเปฯ สิกฺขมานา, สามเณโร, สามเณรี, อุปาสโก, อุปาสิกา ยถาภตํ นิกฺขิตฺตา เอวํ นิรเย.\csromanspacer{} กตเมหิ ปญฺจหิ? ปาณาติปาตินี โหติ, อทินฺนาทายินี โหติ, กาเมสุมิจฺฉาจารินี โหติ, มุสาวาทินี โหติ, สุราเมรยมชฺชปมาทฏฺฐายินี โหติ.\csromanspacer{} อิเมหิ โข ภิกฺขเว ปญฺจหิ ธมฺเมหิ สมนฺนาคตา อุปาสิกา ยถาภตํ นิกฺขิตฺตา เอวํ นิรเย.}

\prose{ปญฺจหิ ภิกฺขเว ธมฺเมหิ สมนฺนาคตา อุปาสิกา ยถาภตํ นิกฺขิตฺตา เอวํ สคฺเค.\csromanspacer{} กตเมหิ ปญฺจหิ? ปาณาติปาตา ปฏิวิรตา โหติ, อทินฺนาทานา ปฏิวิรตา โหติ, กาเมสุมิจฺฉาจารา ปฏิวิรตา โหติ, มุสาวาทา ปฏิวิรตา โหติ, สุราเมรยมชฺชปมาทฏฺฐานา ปฏิวิรตา โหติ.\csromanspacer{} อิเมหิ โข ภิกฺขเว ปญฺจหิ ธมฺเมหิ สมนฺนาคตา อุปาสิกา ยถาภตํ นิกฺขิตฺตา เอวํ สคฺเคติ.\csromanspacer{}\csromanspacer{}สตฺตมํ.}

\csromanheader{2}{8. อาชีวกสุตฺต}
\csromanmatikaanchor{mtk.123}
\csromanmatikaanchor{mtk.124}
\csromantocmarknopage{section}{(29) 3. ราคเปยฺยาล}{mtk.123}
\bookmark[dest={mtk.123},level=1]{(29) 3. ราคเปยฺยาล}
\csromantocmarkref{section}{ปฐมสุตฺตาทิ}{mtk.124}
\bookmark[dest={mtk.124},level=1]{ปฐมสุตฺตาทิ}
\proseitem{293}{\csromanfolio{244}ปญฺจหิ ภิกฺขเว ธมฺเมหิ สมนฺนาคโต อาชีวโก ยถาภตํ นิกฺขิตฺโต เอวํ นิรเย.\csromanspacer{} กตเมหิ ปญฺจหิ? ปาณาติปาตี โหติ, อทินฺนาทายี โหติ, อพฺรหฺมจารี โหติ, มุสาวาที โหติ, สุราเมรยมชฺชปมาทฏฺฐายี โหติ.\csromanspacer{} อิเมหิ โข ภิกฺขเว ปญฺจหิ ธมฺเมหิ สมนฺนาคโต อาชีวโก ยถาภตํ นิกฺขิตฺโต เอวํ นิรเยติ.\csromanspacer{}\csromanspacer{}อฏฺฐมํ.}

\csromanheader{2}{9-17. นิคณฺฐสุตฺตาทิ}
\proseitem{294-302}{ปญฺจหิ ภิกฺขเว ธมฺเมหิ สมนฺนาคโต นิคณฺโฐ.\csromanspacer{} มุณฺฑสาวโก.\csromanspacer{} ชฏิลโก.\csromanspacer{} ปริพฺพาชโก.\csromanspacer{} มาคณฺฑิโก.\csromanspacer{} เตทณฺฑิโก.\csromanspacer{} อารุทฺธโก.\csromanspacer{} โคตมโก.\csromanspacer{} เทวธมฺมิโก ยถาภตํ นิกฺขิตฺโต เอวํ นิรเย.\csromanspacer{} กตเมหิ ปญฺจหิ? ปาณาติปาตี โหติ, อทินฺนาทายี โหติ ฯเปฯ สุราเมรยมชฺชปมาทฏฺฐายี โหติ.\csromanspacer{} อิเมหิ โข ภิกฺขเว ปญฺจหิ ธมฺเมหิ สมนฺนาคโต เทวธมฺมิโก ยถาภตํ นิกฺขิตฺโต เอวํ นิรเยติ.\csromanspacer{}\csromanspacer{}สตฺตรสมํ.\csromanspacer{} สิกฺขาปทเปยฺยาลํ นิฏฺฐิตํ.}

\csromanheader{2}{3. ราคเปยฺยาล}
\proseitem{303}{ราคสฺส ภิกฺขเว อภิญฺญาย ปญฺจ ธมฺมา ภาเวตพฺพา.\csromanspacer{} กตเม ปญฺจ? อสุภสญฺญา มรณสญฺญา อาทีนวสญฺญา อาหาเร ปฏิกูลสญฺญา สพฺพโลเก อนภิรตสญฺญา\footnote{สพฺพตฺถปิ เอวเมว ทิสฺสติ.}.\csromanspacer{} ราคสฺส ภิกฺขเว อภิญฺญาย อิเม ปญฺจ ธมฺมา ภาเวตพฺพาติ. (1)}

\proseitem{304}{ราคสฺส ภิกฺขเว อภิญฺญาย ปญฺจ ธมฺมา ภาเวตพฺพา.\csromanspacer{} กตเม ปญฺจ? อนิจฺจสญฺญา อนตฺตสญฺญา มรณสญฺญา อาหาเร ปฏิกูลสญฺญา สพฺพโลเก อนภิรตสญฺญา.\csromanspacer{} ราคสฺส ภิกฺขเว อภิญฺญาย อิเม ปญฺจ ธมฺมา ภาเวตพฺพาติ. (2)}

\csromanheader{2}{3. ราคเปยฺยาล}
\proseitem{305}{\csromanfolio{245}ราคสฺส ภิกฺขเว อภิญฺญาย ปญฺจ ธมฺมา ภาเวตพฺพา.\csromanspacer{} กตเม ปญฺจ? อนิจฺจสญฺญา, อนิจฺเจ ทุกฺขสญฺญา, ทุกฺเข อนตฺตสญฺญา, ปหานสญฺญา วิราคสญฺญา.\csromanspacer{} ราคสฺส ภิกฺขเว อภิญฺญาย อิเม ปญฺจ ธมฺมา ภาเวตพฺพาติ. (3)}

\proseitem{306}{ราคสฺส ภิกฺขเว อภิญฺญาย ปญฺจ ธมฺมา ภาเวตพฺพา.\csromanspacer{} กตเม ปญฺจ? สทฺธินฺทฺริยํ วีริยินฺทฺริยํ สตินฺทฺริยํ สมาธินฺทฺริยํ ปญฺญินฺทฺริยํ.\csromanspacer{} ราคสฺส ภิกฺขเว อภิญฺญาย อิเม ปญฺจ ธมฺมา ภาเวตพฺพาติ. (4)}

\proseitem{307}{ราคสฺส ภิกฺขเว อภิญฺญาย ปญฺจ ธมฺมา ภาเวตพฺพา.\csromanspacer{} กตเม ปญฺจ? สทฺธาพลํ วีริยพลํ สติพลํ สมาธิพลํ ปญฺญาพลํ.\csromanspacer{} ราคสฺส ภิกฺขเว อภิญฺญาย อิเม ปญฺจ ธมฺมา ภาเวตพฺพาติ. (5)}

\proseitem{308-1151}{ราคสฺส ภิกฺขเว ปริญฺญาย.\csromanspacer{} ปริกฺขยาย.\csromanspacer{} ปหานาย.\csromanspacer{} ขยาย.\csromanspacer{} วยาย.\csromanspacer{} วิราคาย.\csromanspacer{} นิโรธาย.\csromanspacer{} จาคาย.\csromanspacer{} ปฏินิสฺสคฺคาย ปญฺจ ธมฺมา ภาเวตพฺพา.\csromanspacer{} โทสสฺส.\csromanspacer{} โมหสฺส.\csromanspacer{} โกธสฺส.\csromanspacer{} อุปนาหสฺส.\csromanspacer{} มกฺขสฺส.\csromanspacer{} ปฬาสสฺส.\csromanspacer{} อิสฺสาย.\csromanspacer{} มจฺฉริยสฺส.\csromanspacer{} มายาย.\csromanspacer{} สาเฐยฺยสฺส.\csromanspacer{} ถมฺภสฺส.\csromanspacer{} สารมฺภสฺส.\csromanspacer{} มานสฺส.\csromanspacer{} อติมานสฺส.\csromanspacer{} มทสฺส.\csromanspacer{} ปมาทสฺส อภิญฺญาย.\csromanspacer{} ปริญฺญาย.\csromanspacer{} ปริกฺขยาย.\csromanspacer{} ปหานาย.\csromanspacer{} ขยาย.\csromanspacer{} วยาย.\csromanspacer{} วิราคาย.\csromanspacer{} นิโรธาย.\csromanspacer{} จาคาย.\csromanspacer{} ปฏินิสฺสคฺคาย ปญฺจ ธมฺมา ภาเวตพฺพา. (6-849)}

\prose{กตเม ปญฺจ? สทฺธาพลํ วีริยพลํ สติพลํ สมาธิพลํ ปญฺญาพลํ.\csromanspacer{} ปมาทสฺส ภิกฺขเว ปฏินิสฺสคฺคาย อิเม ปญฺจ ธมฺมา ภาเวตพฺพาติ. (6-850) ราคเปยฺยาลํ นิฏฺฐิตํ.}

\titlehead{\textbf{ตสฺสุทฺทานํ}}
\csromangathameasure{อภิญฺญาย ปริญฺญาย ปริกฺขยาย, ปหานาย ขยาย วเยน จ. \\ วิราคนิโรธา จาคญฺจ, ปฏินิสฺสคฺโค อิเม ทสาติ.}
\csromangathagroup{\csromangathastanza{อภิญฺญาย ปริญฺญาย ปริกฺขยาย, ปหานาย ขยาย วเยน จ. \\ วิราคนิโรธา จาคญฺจ, ปฏินิสฺสคฺโค อิเม ทสาติ.}}

\nitthitam{ปญฺจกนิปาโต นิฏฺฐิโต.}
\csromanheader{2}{\textbf{ตตฺริทํ วคฺคุทฺทานํ}}
\csromangathameasure{เสขพลํ พลญฺเจว, ปญฺจงฺคิกญฺจ สุมนํ. \\ มุณฺฑนีวรณญฺจ สญฺญญฺจ, โยธาชีวญฺจ อฏฺฐมํ. \\ เถรํ กกุธผาสุญฺจ, อนฺธกวินฺท ทฺวาทสํ. \\ คิลานราชติกณฺฑํ, สทฺธมฺมาฆาตุปาสกํ. \\ อรญฺญพฺราหฺมณญฺเจว, กิมิลกฺโกสกํ ตถา. \\ ทีฆาจาราวาสิกญฺจ, ทุจฺจริตูปสมฺปทนฺติ.}
\csromangathagroup{\csromangathastanza{\csromanfolio{246}เสขพลํ พลญฺเจว, ปญฺจงฺคิกญฺจ สุมนํ. \\ มุณฺฑนีวรณญฺจ สญฺญญฺจ, โยธาชีวญฺจ อฏฺฐมํ.}\csromangathastanza{เถรํ กกุธผาสุญฺจ, อนฺธกวินฺท ทฺวาทสํ. \\ คิลานราชติกณฺฑํ, สทฺธมฺมาฆาตุปาสกํ.}\csromangathastanza{อรญฺญพฺราหฺมณญฺเจว, กิมิลกฺโกสกํ ตถา. \\ ทีฆาจาราวาสิกญฺจ, ทุจฺจริตูปสมฺปทนฺติ.}}

\nitthitam{\textbf{ปญฺจกนิปาตปาฬิ นิฏฺฐิตา.}}
\csromanheader{2}{\textbf{องฺคุตฺตรนิกาย}}
\csromanmatikaanchor{mtk.125}
\csromantocmarknopage{chapter}{ฉกฺกนิปาตปาฬิ}{mtk.125}
\bookmark[dest={mtk.125},level=0]{ฉกฺกนิปาตปาฬิ}
\csromantocmarknopage{section}{มติกา}{mtk.126}
\bookmark[dest={mtk.126},level=1]{มติกา}
\csromantocmarknopage{chapter}{1. ปฐมปณฺณาสก}{mtk.127}
\bookmark[dest={mtk.127},level=0]{1. ปฐมปณฺณาสก}
\gambhira{\textbf{ฉกฺกนิปาตปาฬิ}}
\namakkaram{นโม ตสฺส ภควโต อรหโต สมฺมาสมฺพุทฺธสฺส.}
\csromanheader{2}{1. ปฐมปณฺณาสก}
\chapterhead{1. อาหุเนยฺยวคฺค}
\csromanheader{2}{1. ปฐม อาหุเนยฺยสุตฺต}
\csromanmatikaanchor{mtk.128}
\csromanmatikaanchor{mtk.129}
\csromantocmarknopage{section}{1. อาหุเนยฺยวคฺค}{mtk.128}
\bookmark[dest={mtk.128},level=1]{1. อาหุเนยฺยวคฺค}
\csromantocmarkref{subsection}{1-2. อาหุเนยฺยสุตฺตทฺวย}{mtk.129}
\bookmark[dest={mtk.129},level=2]{1-2. อาหุเนยฺยสุตฺตทฺวย}
\proseitem{1}{\csromanfolio{247}เอวํ เม สุตํ\csromandash{}เอกํ สมยํ ภควา สาวตฺถิยํ วิหรติ เชตวเน อนาถปิณฺฑิกสฺส อาราเม.\csromanspacer{} ตตฺร โข ภควา ภิกฺขู อามนฺเตสิ “ภิกฺขโว”ติ. “ภทนฺเต”ติ เต ภิกฺขู ภควโต ปจฺจสฺโสสุํ.\csromanspacer{} ภควา เอตทโวจ\csromandash{}}

\prose{ฉหิ ภิกฺขเว ธมฺเมหิ สมนฺนาคโต ภิกฺขุ อาหุเนยฺโย โหติ ปาหุเนยฺโย ทกฺขิเณยฺโย อญฺชลิกรณีโย อนุตฺตรํ ปุญฺญกฺเขตฺตํ โลกสฺส.\csromanspacer{} กตเมหิ ฉหิ? * อิธ ภิกฺขเว ภิกฺขุ จกฺขุนา รูปํ ทิสฺวา เนว สุมโน โหติ น ทุมฺมโน, อุเปกฺขโก วิหรติ สโต สมฺปชาโน.\csromanspacer{} โสเตน สทฺทํ สุตฺวา.\csromanspacer{} ฆาเนน คนฺธํ ฆายิตฺวา.\csromanspacer{} ชิวฺหาย รสํ สายิตฺวา.\csromanspacer{} กาเยน โผฏฺฐพฺพํ ผุสิตฺวา.\csromanspacer{} มนสา ธมฺมํ วิญฺญาย เนว สุมโน โหติ น ทุมฺมโน, อุเปกฺขโก วิหรติ สโต สมฺปชาโน.\csromanspacer{} อิเมหิ โข ภิกฺขเว ฉหิ ธมฺเมหิ สมนฺนาคโต ภิกฺขุ อาหุเนยฺโย โหติ ปาหุเนยฺโย ทกฺขิเณยฺโย อญฺชลิกรณีโย อนุตฺตรํ ปุญฺญกฺเขตฺตํ โลกสฺสาติ.}

\prose{อิทมโวจ ภควา, อตฺตมนา เต ภิกฺขู ภควโต ภาสิตํ อภินนฺทุนฺติ.\csromanspacer{}\csromanspacer{}ปฐมํ.}

\gambhira{ฉกฺกนิปาตปาฬิ \textbf{2. ทุติย อาหุเนยฺยสุตฺต}}
\proseitem{2}{\csromanfolio{248}ฉหิ ภิกฺขเว ธมฺเมหิ สมนฺนาคโต ภิกฺขุ อาหุเนยฺโย โหติ ฯเปฯ อนุตฺตรํ ปุญฺญกฺเขตฺตํ โลกสฺส.\csromanspacer{} กตเมหิ ฉหิ? * อิธ ภิกฺขเว ภิกฺขุ อเนกวิหิตํ อิทฺธิวิธํ ปจฺจนุโภติ, เอโกปิ หุตฺวา พหุธา โหติ, พหุธาปิ หุตฺวา เอโก โหติ, อาวิภาวํ ติโรภาวํ ติโรกุฏฺฏํ ติโรปาการํ ติโรปพฺพตํ อสชฺชมาโน คจฺฉติ เสยฺยถาปิ อากาเส.\csromanspacer{} ปถวิยาปิ อุมฺมุชฺชนิมุชฺชํ กโรติ เสยฺยถาปิ อุทเก.\csromanspacer{} อุทเกปิ อภิชฺชมาเน คจฺฉติ เสยฺยถาปิ ปถวิยํ.\csromanspacer{} อากาเสปิ ปลฺลงฺเกน กมติ เสยฺยถาปิ ปกฺขี สกุโณ.\csromanspacer{} อิเมปิ จนฺทิมสูริเย เอวํมหิทฺธิเก เอวํมหานุภาเว ปาณินา ปริมสติ\footnote{ปรามสติ (ก)} ปริมชฺชติ, ยาว พฺรหฺมโลกาปิ กาเยน วสํ วตฺเตติ. (1)}

\prose{ทิพฺพาย โสตธาตุยา วิสุทฺธาย อติกฺกนฺตมานุสิกาย อุโภ สทฺเท สุณาติ ทิพฺเพ จ มานุเส จ เย ทูเร สนฺติเก จ. (2)}

\prose{ปรสตฺตานํ ปรปุคฺคลานํ เจตสา เจโต ปริจฺจ ปชานาติ, สราคํ วา จิตฺตํ “สราคํ จิตฺตนฺ”ติ ปชานาติ, วีตราคํ วา จิตฺตํ.\csromanspacer{} สโทสํ วา จิตฺตํ.\csromanspacer{} วีตโทสํ วา จิตฺตํ.\csromanspacer{} สโมหํ วา จิตฺตํ.\csromanspacer{} วีตโมหํ วา จิตฺตํ.\csromanspacer{} สํขิตฺตํ วา จิตฺตํ.\csromanspacer{} วิกฺขิตฺตํ วา จิตฺตํ.\csromanspacer{} มหคฺคตํ วา จิตฺตํ.\csromanspacer{} อมหคฺคตํ วา จิตฺตํ.\csromanspacer{} ส-อุตฺตรํ วา จิตฺตํ.\csromanspacer{} อนุตฺตรํ วา จิตฺตํ.\csromanspacer{} สมาหิตํ วา จิตฺตํ.\csromanspacer{} อสมาหิตํ วา จิตฺตํ.\csromanspacer{} วิมุตฺตํ วา จิตฺตํ.\csromanspacer{} อวิมุตฺตํ วา จิตฺตํ “อวิมุตฺตํ จิตฺตนฺ”ติ ปชานาติ. (3)}

\prose{อเนกวิหิตํ ปุพฺเพนิวาสํ อนุสฺสรติ.\csromanspacer{} เสยฺยถิทํ, เอกมฺปิ ชาติํ เทฺวปิ ชาติโย ฯเปฯ อิติ สาการํ ส-อุทฺเทสํ อเนกวิหิตํ ปุพฺเพนิวาสํ อนุสฺสรติ. (4)}

\prose{ทิพฺเพน จกฺขุนา วิสุทฺเธน อติกฺกนฺตมานุสเกน สตฺเต ปสฺสติ จวมาเน อุปปชฺชมาเน หีเน ปณีเต สุวณฺเณ ทุพฺพณฺเณ สุคเต ทุคฺคเต, ยถากมฺมูปเค สตฺเต ปชานาติ “อิเม วต โภนฺโต สตฺตา กายทุจฺจริเตน สมนฺนาคตา วจีทุจฺจริเตน สมนฺนาคตา มโนทุจฺจริเตน สมนฺนาคตา อริยานํ อุปวาทกา มิจฺฉาทิฏฺฐิกา มิจฺฉาทิฏฺฐิกมฺมสมาทานา, เต กายสฺส เภทา ปรํ มรณา อปายํ}

\vspace{\tipitakaparskip}
\csromancenter{อาหุเนยฺยวคฺค ทุคฺคติํ วินิปาตํ นิรยํ อุปปนฺนา.\csromanspacer{} อิเม วา ปน โภนฺโต สตฺตา กายสุจริเตน สมนฺนาคตา วจีสุจริเตน สมนฺนาคตา มโนสุจริเตน สมนฺนาคตา อริยานํ อนุปวาทกา สมฺมาทิฏฺฐิกา สมฺมาทิฏฺฐิกมฺมสมาทานา, เต กายสฺส เภทา ปรํ มรณา สุคติํ สคฺคํ โลกํ อุปปนฺนา”ติ.\csromanspacer{} อิติ ทิพฺเพน จกฺขุนา วิสุทฺเธน อติกฺกนฺตมานุสเกน สตฺเต ปสฺสติ จวมาเน อุปปชฺชมาเน หีเน ปณีเต สุวณฺเณ ทุพฺพณฺเณ สุคเต ทุคฺคเต, ยถากมฺมูปเค สตฺเต ปชานาติ. (5)}
\csromanmatikaanchor{mtk.130}
\csromantocmarkref{subsection}{3-4. อินฺทฺริยพลสุตฺตานิ}{mtk.130}
\bookmark[dest={mtk.130},level=2]{3-4. อินฺทฺริยพลสุตฺตานิ}
\prose{\csromanfolio{249}อาสวานํ ขยา อนาสวํ เจโตวิมุตฺติํ ปญฺญาวิมุตฺติํ ทิฏฺเฐว ธมฺเม สยํ อภิญฺญา สจฺฉิกตฺวา อุปสมฺปชฺช วิหรติ. (6)}

\prose{อิเมหิ โข ภิกฺขเว ฉหิ ธมฺเมหิ สมนฺนาคโต ภิกฺขู อาหุเนยฺโย โหติ ฯเปฯ อนุตฺตรํ ปุญฺญกฺเขตฺตํ โลกสฺสาติ.\csromanspacer{}\csromanspacer{}ทุติยํ.}

\csromanheader{2}{3. อินฺทฺริยสุตฺต}
\proseitem{3}{ฉหิ ภิกฺขเว ธมฺเมหิ สมนฺนาคโต ภิกฺขุ อาหุเนยฺโย โหติ ฯเปฯ อนุตฺตรํ ปุญฺญกฺเขตฺตํ โลกสฺส.\csromanspacer{} กตเมหิ ฉหิ? สทฺธินฺทฺริเยน วีริยินฺทฺริเยน สตินฺทฺริเยน สมาธินฺทฺริเยน ปญฺญินฺทฺริเยน, อาสวานํ ขยา อนาสวํ เจโตวิมุตฺติํ ปญฺญาวิมุตฺติํ ทิฏฺเฐว ธมฺเม สยํ อภิญฺญา สจฺฉิกตฺวา อุปสมฺปชฺช วิหรติ.\csromanspacer{} อิเมหิ โข ภิกฺขเว ฉหิ ธมฺเมหิ สมนฺนาคโต ภิกฺขุ อาหุเนยฺโย โหติ ปาหุเนยฺโย ทกฺขิเณยฺโย อญฺชลิกรณีโย อนุตฺตรํ ปุญฺญกฺเขตฺตํ โลกสฺสาติ.\csromanspacer{}\csromanspacer{}ตติยํ.}

\csromanheader{2}{4. พลสุตฺต}
\proseitem{4}{ฉหิ ภิกฺขเว ธมฺเมหิ สมนฺนาคโต ภิกฺขุ อาหุเนยฺโย โหติ ฯเปฯ อนุตฺตรํ ปุญฺญกฺเขตฺตํ โลกสฺส.\csromanspacer{} กตเมหิ ฉหิ? สทฺธาพเลน วีริยพเลน สติพเลน สมาธิพเลน ปญฺญาพเลน, อาสวานํ ขยา อนาสวํ เจโตวิมุตฺติํ ปญฺญาวิมุตฺติํ ทิฏฺเฐว ธมฺเม สยํ อภิญฺญา สจฺฉิกตฺวา อุปสมฺปชฺช วิหรติ.\csromanspacer{} อิเมหิ โข ภิกฺขเว ฉหิ ธมฺเมหิ สมนฺนาคโต ภิกฺขุ อาหุเนยฺโย โหติ ฯเปฯ อนุตฺตรํ ปุญฺญกฺเขตฺตํ โลกสฺสาติ.\csromanspacer{}\csromanspacer{}จตุตฺถํ.}

\csromanmatikaanchor{mtk.131}
\csromantocmarkref{subsection}{5-7. ปฐม-อาชานียสุตฺตานิ}{mtk.131}
\bookmark[dest={mtk.131},level=2]{5-7. ปฐม-อาชานียสุตฺตานิ}
\csromanheader{2}{ฉกฺกนิปาตปาฬิ \textbf{5. ปฐม อาชานียสุตฺต}}
\proseitem{5}{\csromanfolio{250}ฉหิ ภิกฺขเว องฺเคหิ สมนฺนาคโต รญฺโญ ภโทฺร อสฺสาชานีโย ราชารโห โหติ ราชโภคฺโค รญฺโญ องฺคนฺเตฺวว สงฺขํ คจฺฉติ.}

\prose{กตเมหิ ฉหิ? อิธ ภิกฺขเว รญฺโญ ภโทฺร อสฺสาชานีโย ขโม โหติ รูปานํ, ขโม สทฺทานํ, ขโม คนฺธานํ, ขโม รสานํ, ขโม โผฏฺฐพฺพานํ, วณฺณสมฺปนฺโน จ โหติ.\csromanspacer{} อิเมหิ โข ภิกฺขเว ฉหิ องฺเคหิ สมนฺนาคโต รญฺโญ ภโทฺร อสฺสาชานีโย ราชารโห โหติ ราชโภคฺโค รญฺโญ องฺคนฺเตฺวว สงฺขํ คจฺฉติ.}

\prose{เอวเมวํ โข ภิกฺขเว ฉหิ ธมฺเมหิ สมนฺนาคโต ภิกฺขุ อาหุเนยฺโย โหติ ฯเปฯ อนุตฺตรํ ปุญฺญกฺเขตฺตํ โลกสฺส.\csromanspacer{} กตเมหิ ฉหิ? อิธ ภิกฺขเว ภิกฺขุ ขโม โหติ รูปานํ, ขโม สทฺทานํ, ขโม คนฺธานํ, ขโม รสานํ, ขโม โผฏฺฐพฺพานํ, ขโม ธมฺมานํ.\csromanspacer{} อิเมหิ โข ภิกฺขเว ฉหิ ธมฺเมหิ สมนฺนาคโต ภิกฺขุ อาหุเนยฺโย โหติ ฯเปฯ อนุตฺตรํ ปุญฺญกฺเขตฺตํ โลกสฺสาติ.\csromanspacer{}\csromanspacer{}ปญฺจมํ.}

\csromanheader{2}{6. ทุติย อาชานียสุตฺต}
\proseitem{6}{ฉหิ ภิกฺขเว องฺเคหิ สมนฺนาคโต รญฺโญ ภโทฺร อสฺสาชานีโย ราชารโห โหติ ราชโภคฺโค, รญฺโญ องฺคนฺเตฺวว สงฺขํ คจฺฉติ.\csromanspacer{} กตเมหิ ฉหิ? อิธ ภิกฺขเว รญฺโญ ภโทฺร อสฺสาชานีโย ขโม โหติ รูปานํ, ขโม สทฺทานํ, ขโม คนฺธานํ, ขโม รสานํ, ขโม โผฏฺฐพฺพานํ, พลสมฺปนฺโน จ โหติ.\csromanspacer{} อิเมหิ โข ภิกฺขเว ฉหิ องฺเคหิ สมนฺนาคโต รญฺโญ ภโทฺร อสฺสาชานีโย ราชารโห โหติ ราชโภคฺโค, รญฺโญ องฺคนฺเตฺวว สงฺขํ คจฺฉติ.}

\prose{เอวเมวํ โข ภิกฺขเว ฉหิ ธมฺเมหิ สมนฺนาคโต ภิกฺขุ อาหุเนยฺโย ฯเปฯ อนุตฺตรํ ปุญฺญกฺเขตฺตํ โลกสฺส.\csromanspacer{} กตเมหิ ฉหิ? อิธ ภิกฺขเว ภิกฺขุ ขโม โหติ รูปานํ ฯเปฯ ขโม ธมฺมานํ.\csromanspacer{} อิเมหิ โข ภิกฺขเว ฉหิ ธมฺเมหิ สมนฺนาคโต ภิกฺขุ อาหุเนยฺโย โหติ ฯเปฯ อนุตฺตรํ ปุญฺญกฺเขตฺตํ โลกสฺสาติ.\csromanspacer{}\csromanspacer{}ฉฏฺฐํ.}

\csromanmatikaanchor{mtk.132}
\csromanmatikaanchor{mtk.278}
\csromantocmarkref{subsection}{8-9. อนุตฺตริยสุตฺตาทิ}{mtk.132}
\bookmark[dest={mtk.132},level=2]{8-9. อนุตฺตริยสุตฺตาทิ}
\chapterheadpageread{1. อาหุเนยฺยวคฺค \textbf{7. ตติย อาชานียสุตฺต}}
\proseitem{7}{\csromanfolio{251}ฉหิ ภิกฺขเว องฺเคหิ สมนฺนาคโต รญฺโญ ภโทฺร อสฺสาชานีโย ราชารโห โหติ ราชโภคฺโค, รญฺโญ องฺคนฺเตฺวว สงฺขํ คจฺฉติ.\csromanspacer{} กตเมหิ ฉหิ? อิธ ภิกฺขเว รญฺโญ ภโทฺร อสฺสาชานีโย ขโม โหติ รูปานํ, ขโม สทฺทานํ, ขโม คนฺธานํ, ขโม รสานํ, ขโม โผฏฺฐพฺพานํ, ชวสมฺปนฺโน จ โหติ.\csromanspacer{} อิเมหิ โข ภิกฺขเว ฉหิ องฺเคหิ สมนฺนาคโต รญฺโญ ภโทฺร อสฺสาชานีโย ราชารโห โหติ ราชโภคฺโค, รญฺโญ องฺคนฺเตฺวว สงฺขํ คจฺฉติ.}

\prose{เอวเมวํ โข ภิกฺขเว ฉหิ ธมฺเมหิ สมนฺนาคโต ภิกฺขุ อาหุเนยฺโย โหติ ฯเปฯ อนุตฺตรํ ปุญฺญกฺเขตฺตํ โลกสฺส.\csromanspacer{} กตเมหิ ฉหิ? อิธ ภิกฺขเว ภิกฺขุ ขโม โหติ รูปานํ ฯเปฯ ขโม ธมฺมานํ.\csromanspacer{} อิเมหิ โข ภิกฺขเว ฉหิ ธมฺเมหิ สมนฺนาคโต ภิกฺขุ อาหุเนยฺโย โหติ ฯเปฯ อนุตฺตรํ ปุญฺญกฺเขตฺตํ โลกสฺสาติ.\csromanspacer{}\csromanspacer{}สตฺตมํ.}

\csromanheader{2}{8. อนุตฺตริยสุตฺต}
\proseitem{8}{\csromansymbolfootnote{*}{ที 3. 207 ปิฏฺเฐปิ.}ฉยิมานิ ภิกฺขเว อนุตฺตริยานิ.\csromanspacer{} กตมานิ ฉ? ทสฺสนานุตฺตริยํ สวนานุตฺตริยํ ลาภานุตฺตริยํ สิกฺขานุตฺตริยํ ปาริจริยานุตฺตริยํ อนุสฺสตานุตฺตริยํ.\csromanspacer{} อิมานิ โข ภิกฺขเว ฉ อนุตฺตริยานีติ.\csromanspacer{}\csromanspacer{}อฏฺฐมํ.}

\csromanheader{2}{9. อนุสฺสติฏฺฐานสุตฺต}
\proseitem{9}{\csromansymbolmark{*}ฉยิมานิ ภิกฺขเว อนุสฺสติฏฺฐานานิ.\csromanspacer{} กตมานิ ฉ? พุทฺธานุสฺสติ ธมฺมานุสฺสติ สํฆานุสฺสติ สีลานุสฺสติ จาคานุสฺสติ เทวตานุสฺสติ.\csromanspacer{} อิมานิ โข ภิกฺขเว ฉ อนุสฺสติฏฺฐานานีติ.\csromanspacer{}\csromanspacer{}นวมํ.}

\csromanmatikaanchor{mtk.133}
\csromantocmarkref{subsection}{10. มหานามสุตฺต}{mtk.133}
\bookmark[dest={mtk.133},level=2]{10. มหานามสุตฺต}
\csromanheader{2}{10. มหานามสุตฺต}
\proseitem{10}{เอกํ สมยํ ภควา สกฺเกสุ วิหรติ กปิลวตฺถุสฺมิํ นิโคฺรธาราเม.\csromanspacer{} อถ โข มหานาโม สกฺโก เยน ภควา เตนุปสงฺกมิ, อุปสงฺกมิตฺวา ภควนฺตํ อภิวาเทตฺวา เอกมนฺตํ นิสีทิ, เอกมนฺตํ นิสินฺโน โข มหานาโม สกฺโก ภควนฺตํ เอตทโวจ “โย โส ภนฺเต อริยสาวโก อาคตผโล วิญฺญาตสาสโน, โส กตเมน วิหาเรน พหุลํ วิหรตี”ติ.}

\gambhira{ฉกฺกนิปาตปาฬิ}
\prose{\csromanfolio{252}โย โส มหานาม อริยสาวโก อาคตผโล วิญฺญาตสาสโน, โส อิมินา วิหาเรน พหุลํ วิหรติ. * อิธ มหานาม อริยสาวโก ตถาคตํ อนุสฺสรติ “อิติปิ โส ภควา อรหํ สมฺมาสมฺพุทฺโธ วิชฺชาจรณสมฺปนฺโน สุคโต โลกวิทู อนุตฺตโร ปุริสทมฺมสารถิ สตฺถา เทวมนุสฺสานํ พุทฺโธ ภควา”ติ.\csromanspacer{} ยสฺมิํ มหานาม สมเย อริยสาวโก ตถาคตํ อนุสฺสรติ, เนวสฺส ตสฺมิํ สมเย ราคปริยุฏฺฐิตํ จิตฺตํ โหติ, น โทสปริยุฏฺฐิตํ จิตฺตํ โหติ, น โมหปริยุฏฺฐิตํ จิตฺตํ โหติ, อุชุคตเมวสฺส ตสฺมิํ สมเย จิตฺตํ โหติ ตถาคตํ อารพฺภ.\csromanspacer{} อุชุคตจิตฺโต โข ปน มหานาม อริยสาวโก ลภติ อตฺถเวทํ, ลภติ ธมฺมเวทํ, ลภติ ธมฺมูปสํหิตํ ปาโมชฺชํ, ปมุทิตสฺส ปีติ ชายติ, ปีติมนสฺส กาโย ปสฺสมฺภติ, ปสฺสทฺธกาโย สุขํ เวทิยติ, สุขิโน จิตฺตํ สมาธิยติ.\csromanspacer{} อยํ วุจฺจติ มหานาม อริยสาวโก วิสมคตาย ปชาย สมปฺปตฺโต วิหรติ, สพฺยาปชฺชาย ปชาย อพฺยาปชฺโช วิหรติ, ธมฺมโสตํ สมาปนฺโน พุทฺธานุสฺสติํ ภาเวติ. (1)}

\prose{ปุน จปรํ มหานาม อริยสาวโก ธมฺมํ อนุสฺสรติ “สฺวากฺขาโต ภควตา ธมฺโม สนฺทิฏฺฐิโก อกาลิโก เอหิปสฺสิโก โอปเนยฺยิโก ปจฺจตฺตํ เวทิตพฺโพ วิญฺญูหี”ติ.\csromanspacer{} ยสฺมิํ มหานาม สมเย อริยสาวโก ธมฺมํ อนุสฺสรติ, เนวสฺส ตสฺมิํ สมเย ราคปริยุฏฺฐิตํ จิตฺตํ โหติ, น โทสปริยุฏฺฐิตํ จิตฺตํ โหติ, น โมหปริยุฏฺฐิตํ จิตฺตํ โหติ, อุชุคตเมวสฺส ตสฺมิํ สมเย จิตฺตํ โหติ ธมฺมํ อารพฺภ.\csromanspacer{} อุชุคตจิตฺโต โข ปน มหานาม อริยสาวโก ลภติ อตฺถเวทํ, ลภติ ธมฺมเวทํ, ลภติ ธมฺมูปสํหิตํ ปาโมชฺชํ, ปมุทิตสฺส ปีติ ชายติ, ปีติมนสฺส กาโย ปสฺสมฺภติ, ปสฺสทฺธกาโย สุขํ เวทิยติ, สุขิโน จิตฺตํ สมาธิยติ.\csromanspacer{} อยํ วุจฺจติ มหานาม อริยสาวโก วิสมคตาย ปชาย สมปฺปตฺโต วิหรติ, สพฺยาปชฺชาย ปชาย อพฺยาปชฺโช วิหรติ, ธมฺมโสตํ สมาปนฺโน ธมฺมานุสฺสติํ ภาเวติ. (2)}

\prose{ปุน จปรํ มหานาม อริยสาวโก สํฆํ อนุสฺสรติ “สุปฺปฏิปนฺโน ภควโต สาวกสํโฆ, อุชุปฺปฏิปนฺโน ภควโต สาวกสํโฆ, ญายปฺปฏิปนฺโน ภควโต สาวกสํโฆ, สามีจิปฺปฏิปนฺโน ภควโต สาวกสํโฆ, ยทิทํ จตฺตาริ ปุริสยุคานิ อฏฺฐ ปุริสปุคฺคลา, เอส}

\vspace{\tipitakaparskip}
\csromancenter{อาหุเนยฺยวคฺค ภควโต สาวกสํโฆ อาหุเนยฺโย ปาหุเนยฺโย ทกฺขิเณยฺโย อญฺชลิกรณีโย อนุตฺตรํ ปุญฺญกฺเขตฺตํ โลกสฺสา”ติ.\csromanspacer{} ยสฺมิํ มหานาม สมเย อริยสาวโก สํฆํ อนุสฺสรติ, เนวสฺส ตสฺมิํ สมเย ราคปริยุฏฺฐิตํ จิตฺตํ โหติ, น โทสปริยุฏฺฐิตํ จิตฺตํ โหติ, น โมหปริยุฏฺฐิตํ จิตฺตํ โหติ, อุชุคตเมวสฺส ตสฺมิํ สมเย จิตฺตํ โหติ สํฆํ อารพฺภ.\csromanspacer{} อุชุคตจิตฺโต โข ปน มหานาม อริยสาวโก ลภติ อตฺถเวทํ, ลภติ ธมฺมเวทํ, ลภติ ธมฺมูปสํหิตํ ปาโมชฺชํ, ปมุทิตสฺส ปีติ ชายติ, ปีติมนสฺส กาโย ปสฺสมฺภติ, ปสฺสทฺธกาโย สุขํ เวทิยติ, สุขิโน จิตฺตํ สมาธิยติ.\csromanspacer{} อยํ วุจฺจติ มหานาม อริยสาวโก วิสมคตาย ปชาย สมปฺปตฺโต วิหรติ, สพฺยาปชฺชาย ปชาย อพฺยาปชฺโช วิหรติ, ธมฺมโสตํ สมาปนฺโน สํฆานุสฺสติํ ภาเวติ. (3)}
\prose{\csromanfolio{253}ปุน จปรํ มหานาม อริยสาวโก อตฺตโน สีลานิ อนุสฺสรติ อขณฺฑานิ อจฺฉิทฺทานิ อสพลานิ อกมฺมาสานิ ภุชิสฺสานิ วิญฺญุปฺปสตฺถานิ อปรามฏฺฐานิ สมาธิสํวตฺตนิกานิ.\csromanspacer{} ยสฺมิํ มหานาม สมเย อริยสาวโก สีลํ อนุสฺสรติ, เนวสฺส ตสฺมิํ สมเย ราคปริยุฏฺฐิตํ จิตฺตํ โหติ, น โทสปริยุฏฺฐิตํ จิตฺตํ โหติ, น โมหปริยุฏฺฐิตํ จิตฺตํ โหติ, อุชุคตเมวสฺส ตสฺมิํ สมเย จิตฺตํ โหติ สีลํ อารพฺภ.\csromanspacer{} อุชุคตจิตฺโต โข ปน มหานาม อริยสาวโก ลภติ อตฺถเวทํ, ลภติ ธมฺมเวทํ, ลภติ ธมฺมูปสํหิตํ ปาโมชฺชํ, ปมุทิตสฺส ปีติ ชายติ, ปีติมนสฺส กาโย ปสฺสมฺภติ, ปสฺสทฺธกาโย สุขํ เวทิยติ, สุขิโน จิตฺตํ สมาธิยติ.\csromanspacer{} อยํ วุจฺจติ มหานาม อริยสาวโก วิสมคตาย ปชาย สมปฺปตฺโต วิหรติ.\csromanspacer{} สพฺยาปชฺชาย ปชาย อพฺยาปชฺโช วิหรติ, ธมฺมโสตํ สมาปนฺโน สีลานุสฺสติํ ภาเวติ. (4)}

\prose{ปุน จปรํ มหานาม อริยสาวโก อตฺตโน จาคํ อนุสฺสรติ “ลาภา วต เม, สุลทฺธํ วต เม, โยหํ มจฺเฉรมลปริยุฏฺฐิตาย ปชาย วิคตมลมจฺเฉเรน เจตสา อคารํ อชฺฌาวสามิ มุตฺตจาโค ปยภปาณิ โวสคฺครโต ยาจโยโค ทานสํ วิภาครโต”ติ.\csromanspacer{} ยสฺมิํ มหานาม สมเย อริยสาวโก จาคํ อนุสฺสรติ, เนวสฺส ตสฺมิํ สมเย ราคปริยุฏฺฐิตํ จิตฺตํ โหติ, น โทสปริยุฏฺฐิตํ จิตฺตํ}

\vspace{\tipitakaparskip}
\csromancenter{ฉกฺกนิปาตปาฬิ โหติ, น โมหปริยุฏฺฐิตํ จิตฺตํ โหติ, อุชุคตเมวสฺส ตสฺมิํ สมเย จิตฺตํ โหติ จาคํ อารพฺภ.\csromanspacer{} อุชุคตจิตฺโต โข ปน มหานาม อริยสาวโก ลภติ อตฺถเวทํ, ลภติ ธมฺมเวทํ, ลภติ ธมฺมูปสํหิตํ ปาโมชฺชํ, ปมุทิตสฺส ปีติ ชายติ, ปีติมนสฺส กาโย ปสฺสมฺภติ, ปสฺสทฺธกาโย สุขํ เวทิยติ, สุขิโน จิตฺตํ สมาธิยติ.\csromanspacer{} อยํ วุจฺจติ มหานาม อริยสาวโก วิสมคตาย ปชาย สมปฺปตฺโต วิหรติ.\csromanspacer{} สพฺยาปชฺชาย ปชาย อพฺยาปชฺโช วิหรติ, ธมฺมโสตํ สมาปนฺโน จาคานุสฺสติํ ภาเวติ. (5)}
\prose{\csromanfolio{254}ปุน จปรํ มหานาม อริยสวโก เทวตานุสฺสติํ ภาเวติ “สนฺติ เทวา จาตุมหาราชิกา\footnote{จาตุมฺมหาราชิกา (สี, สฺยา, กํ, อิ)}, สนฺติ เทวา ตาวติํสา, สนฺติ เทวา ยามา, สนฺติ เทวา ตุสิตา, สนฺติ เทวา นิมฺมานรติโน, สนฺติ เทวา ปรนิมฺมิตวสวตฺติโน, สนฺติ เทวา พฺรหฺมกายิกา, สนฺติ เทวา ตตุตฺตริ\footnote{ตตุตฺตริํ (สี, สฺยา, กํ, อิ), ตทุตฺตริ (ก) อุปริ 276; วิสุทฺธิ 1. 218 ปิฏฺเฐสุปิ ปสฺสิตพฺพํ.}, ยถารูปาย สทฺธาย สมนฺนาคตา ตา เทวตา อิโต จุตา ตตฺถ อุปปนฺนา\footnote{ตตฺถ อุปฺปนฺนา (สี), ตตฺถูปปนฺนา (สฺยา, กํ), ตตฺถุปปนฺนา (อํ 1. 210)}, มยฺหมฺปิ ตถารูปา สทฺธา สํวิชฺชติ.\csromanspacer{} ยถารูเปน สีเลน สมนฺนาคตา ตา เทวตา อิโต จุตา ตตฺถ อุปปนฺนา, มยฺหมฺปิ ตถารูปํ สีลํ สํวิชฺชติ.\csromanspacer{} ยถารูเปน สุเตน สมนฺนาคตา ตา เทวตา อิโต จุตา ตตฺถ อุปปนฺนา, มยฺหมฺปิ ตถารูปํ สุตํ สํวิชฺชติ.\csromanspacer{} ยถารูเปน จาเคน สมนฺนาคตา ตา เทวตา อิโต จุตา ตตฺถ อุปปนฺนา, มยฺหมฺปิ ตถารูโป จาโค สํวิชฺชติ.\csromanspacer{} ยถารูปาย ปญฺญาย สมนฺนาคตา ตา เทวตา อิโต จุตา ตตฺถ อุปปนฺนา, มยฺหมฺปิ ตถารูปา ปญฺญา สํวิชฺชตี”ติ.\csromanspacer{} ยสฺมิํ มหานาม สมเย อริยสาวโก อตฺตโน จ ตาสญฺจ เทวตานํ สทฺธญฺจ สีลญฺจ สุตญฺจ จาคญฺจ ปญฺญญฺจ อนุสฺสรติ, เนวสฺส ตสฺมิํ สมเย ราคปริยุฏฺฐิตํ จิตฺตํ โหติ, น โทสปริยุฏฺฐิตํ จิตฺตํ โหติ, น โมหปริยุฏฺฐิตํ จิตฺตํ โหติ, อุชุคตเมวสฺส ตสฺมิํ สมเย จิตฺตํ โหติ ตา เทวตา อารพฺภ.\csromanspacer{} อุชุคตจิตฺโต โข ปน มหานาม อริยสาวโก ลภติ อตฺถเวทํ, ลภติ ธมฺมเวทํ, ลภติ ธมฺมูปสํหิตํ ปาโมชฺชํ, ปมุทิตสฺส ปีติ ชายติ, ปีติมนสฺส กาโย ปสฺสมฺภติ, ปสฺสทฺธกาโย สุขํ เวทิยติ, สุขิโน จิตฺตํ สมาธิยติ.\csromanspacer{} อยํ วุจฺจติ มหานาม อริยสาวโก วิสมคตาย ปชาย สมปฺปตฺโต}

\vspace{\tipitakaparskip}
\csromancenter{สารณียวคฺค วิหรติ, สพฺยาปชฺชาย\footnote{สพฺยาปชฺฌาย ... อพฺยาปชฺโฌ (ก)} ปชาย อพฺยาปชฺโช1 วิหรติ, ธมฺมโสตํ สมาปนฺโน เทวตานุสฺสติํ ภาเวติ. (6)}
\csromanmatikaanchor{mtk.134}
\csromanmatikaanchor{mtk.135}
\csromantocmarknopage{section}{2. สารณียวคฺค}{mtk.134}
\bookmark[dest={mtk.134},level=1]{2. สารณียวคฺค}
\csromantocmarkref{subsection}{1-2. สารณียสุตฺตทฺวย}{mtk.135}
\bookmark[dest={mtk.135},level=2]{1-2. สารณียสุตฺตทฺวย}
\prose{\csromanfolio{255}โย โส มหานาม อริยสาวโก อาคตผโล วิญฺญาตสาสโน, โส อิมินา วิหาเรน พหุลํ วิหรตีติ.\csromanspacer{}\csromanspacer{}ทสมํ.}

\vspace{\tipitakaparskip}
\csromancenter{อาหุเนยฺยวคฺโค ปฐโม.}
\titlehead{\textbf{ตสฺสุทฺทานํ}}
\csromangathameasure{เทฺว อาหุเนยฺยา อินฺทฺริย, พลานิ ตโย อาชานียา. \\ อนุตฺตริย อนุสฺสตี, มหานาเมน เต ทสาติ.}
\csromangathagroup{\csromangathastanza{เทฺว อาหุเนยฺยา อินฺทฺริย, พลานิ ตโย อาชานียา. \\ อนุตฺตริย อนุสฺสตี, มหานาเมน เต ทสาติ.}}

\csromanheader{1}{2. สารณียวคฺค}
\csromanheader{2}{1. ปฐมสารณียสุตฺต}
\proseitem{11}{ฉยิเม ภิกฺขเว ธมฺมา สารณียา\footnote{สาราณียา (สี, สฺยา, กํ, อิ)}.\csromanspacer{} กตเม ฉ? อิธ ภิกฺขเว ภิกฺขุโน เมตฺตํ กายกมฺมํ ปจฺจุปฏฺฐิตํ โหติ สพฺรหฺมจารีสุ อาวิ เจว รโห จ.\csromanspacer{} อยมฺปิ ธมฺโม สารณีโย.}

\prose{ปุน จปรํ ภิกฺขเว ภิกฺขุโน เมตฺตํ วจีกมฺมํ ปจฺจุปฏฺฐิตํ โหติ สพฺรหฺมจารีสุ อาวิ เจว รโห จ.\csromanspacer{} อยมฺปิ ธมฺโม สารณีโย.}

\prose{ปุน จปรํ ภิกฺขเว ภิกฺขุโน เมตฺตํ มโนกมฺมํ ปจฺจุปฏฺฐิตํ โหติ สพฺรหฺมจารีสุ อาวิ เจว รโห จ.\csromanspacer{} อยมฺปิ ธมฺโม สารณีโย.}

\prose{ปุน จปรํ ภิกฺขเว ภิกฺขุ เย เต ลาภา ธมฺมิกา ธมฺมลทฺธา อนฺตมโส ปตฺตปริยาปนฺนมตฺตมฺปิ ตถารูเปหิ ลาเภหิ อปฺปฏิวิภตฺตโภคี โหติ สีลวนฺเตหิ สพฺรหฺมจารีหิ สาธารณโภคี.\csromanspacer{} อยมฺปิ ธมฺโม สารณีโย.}

\prose{ปุน จปรํ ภิกฺขเว ภิกฺขุ ยานิ ตานิ สีลานิ อขณฺฑานิ อจฺฉิทฺทานิ อสพลานิ อกมฺมาสานิ ภุชิสฺสานิ วิญฺญุปฺปสตฺถานิ อปรามฏฺฐานิ}

\vspace{\tipitakaparskip}
\csromancenter{ฉกฺกนิปาตปาฬิ สมาธิสํวตฺตนิกานิ, ตถารูเปหิ สีเลหิ สีลสามญฺญคโต วิหรติ สพฺรหฺมจารีหิ อาวิ เจว รโห จ.\csromanspacer{} อยมฺปิ ธมฺโม สารณีโย.}
\prose{\csromanfolio{256}ปุน จปรํ ภิกฺขเว ภิกฺขุ ยายํ ทิฏฺฐิ อริยา นิยฺยานิกา นิยฺยาติ ตกฺกรสฺส สมฺมา ทุกฺขกฺขยาย, ตถารูปาย ทิฏฺฐิยา ทิฏฺฐิสามญฺญคโต วิหรติ สพฺรหฺมจารีหิ อาวิ เจว รโห จ.\csromanspacer{} อยมฺปิ ธมฺโม สารณีโย.\csromanspacer{} อิเม โข ภิกฺขเว ฉ ธมฺมา สารณียาติ.\csromanspacer{}\csromanspacer{}ปฐมํ.}

\csromanheader{2}{2. ทุติยสารณียสุตฺต}
\proseitem{12}{ฉยิเม ภิกฺขเว ธมฺมา สารณียา ปิยกรณา ครุกรณา สงฺคหาย อวิวาทาย สามคฺคิยา เอกีภาวาย สํวตฺตนฺติ.\csromanspacer{} กตเม ฉ? อิธ ภิกฺขเว ภิกฺขุโน เมตฺตํ กายกมฺมํ ปจฺจุปฏฺฐิตํ โหติ สพฺรหฺมจารีสุ อาวิ เจว รโห จ.\csromanspacer{} อยมฺปิ ธมฺโม สารณีโย ปิยกรโณ ครุกรโณ สงฺคหาย อวิวาทาย สามคฺคิยา เอกีภาวาย สํวตฺตติ.}

\prose{ปุน จปรํ ภิกฺขเว ภิกฺขุโน เมตฺตํ วจีกมฺมํ ปจฺจุปฏฺฐิตํ โหติ ฯเปฯ เมตฺตํ มโนกมฺมํ ปจฺจุปฏฺฐิตํ โหติ สพฺรหฺมจารีสุ อาวิ เจว รโห จ.\csromanspacer{} อยมฺปิ ธมฺโม สารณีโย ปิยกรโณ ครุกรโณ สงฺคหาย อวิวาทาย สามคฺคิยา เอกีภาวาย สํวตฺตติ.}

\prose{ปุน จปรํ ภิกฺขเว ภิกฺขุ เย เต ลาภา ธมฺมิกา ธมฺมลทฺธา อนฺตมโส ปตฺตปริยาปนฺนมตฺตมฺปิ, ตถารูเปหิ ลาเภหิ อปฺปฏิวิภตฺตโภคี โหติ สีลวนฺเตหิ สพฺรหฺมจารีหิ สาธารณโภคี.\csromanspacer{} อยมฺปิ ธมฺโม สารณีโย ปิยกรโณ ครุกรโณ สงฺคหาย อวิวาทาย สามคฺคิยา เอกีภาวาย สํวตฺตติ.}

\prose{ปุน จปรํ ภิกฺขเว ภิกฺขุ ยานิ ตานิ สีลานิ อขณฺฑานิ อจฺฉิทฺทานิ อสพลานิ อกมฺมาสานิ ภุชิสฺสานิ วิญฺญุปฺปสตฺถานิ อปรามฏฺฐานิ สมาธิ สํวตฺตนิกานิ.\csromanspacer{} ตถารูเปหิ สีเลหิ สีลสามญฺญคโต วิหรติ สพฺรหฺมจารีหิ อาวิ เจว รโห จ.\csromanspacer{} อยมฺปิ ธมฺโม สารณีโย ปิยกรโณ ครุกรโณ สงฺคหาย อวิวาทาย สามคฺคิยา เอกีภาวาย สํวตฺตติ.}

\chapterheadpageread{2. สารณียวคฺค}
\prose{\csromanfolio{257}ปุน จปรํ ภิกฺขเว ภิกฺขุ ยายํ ทิฏฺฐิ อริยา นิยฺยานิกา นิยฺยาติ ตกฺกรสฺส สมฺมา ทุกฺขกฺขยาย, ตถารูปาย ทิฏฺฐิยา ทิฏฺฐิสามญฺญคโต วิหรติ สพฺรหฺมจารีหิ อาวิ เจว รโห จ.\csromanspacer{} อยมฺปิ ธมฺโม สารณีโย ปิยกรโณ ครุกรโณ สงฺคหาย อวิวาทาย สามคฺคิยา เอกีภาวาย สํวตฺตติ.\csromanspacer{} อิเม โข ภิกฺขเว ฉ ธมฺมา สารณียา ปิยกรณา ครุกรณา สงฺคหาย อวิวาทาย สามคฺคิยา เอกีภาวาย สํวตฺตนฺตีติ.\csromanspacer{}\csromanspacer{}ทุติยํ.}

\csromanmatikaanchor{mtk.136}
\csromantocmarkref{subsection}{3. นิสฺสารณียสุตฺต}{mtk.136}
\bookmark[dest={mtk.136},level=2]{3. นิสฺสารณียสุตฺต}
\csromanheader{2}{3. นิสฺสารณียสุตฺต}
\proseitem{13}{ฉยิมา ภิกฺขเว นิสฺสารณียา ธาตุโย.\csromanspacer{} กตมา ฉ? อิธ ภิกฺขเว ภิกฺขุ เอวํ วเทยฺย “เมตฺตา หิ โข เม เจโตวิมุตฺติ ภาวิตา พหุลีกตา ยานีกตา วตฺถุกตา อนุฏฺฐิตา ปริจิตา สุสมารทฺธา, อถ จ ปน เม พฺยาปาโท จิตฺตํ ปริยาทาย ติฏฺฐตี”ติ.\csromanspacer{} โส “มา เหวนฺ”ติสฺส วจนีโย “มายสฺมา เอวํ อวจ, มา ภควนฺตํ อพฺภาจิกฺขิ, น หิ สาธุ ภควโต อพฺภกฺขานํ, น หิ ภควา เอวํ วเทยฺย.\csromanspacer{} อฏฺฐานเมตํ อาวุโส อนวกาโส, ยํ เมตฺตาย เจโตวิมุตฺติยา ภาวิตาย พหุลีกตาย ยานีกตาย วตฺถุกตาย อนุฏฺฐิตาย ปริจิตาย สุสมารทฺธาย, อถ จ ปนสฺส พฺยาปาโท จิตฺตํ ปริยาทาย ฐสฺสติ\footnote{ฐสฺสตีติ (สพฺพตฺถ) ที 3. 205 ปิฏฺเฐ ปสฺสิตพฺพํ.}, เนตํ ฐานํ วิชฺชติ, นิสฺสรณํ เหตํ อาวุโส พฺยาปาทสฺส ยทิทํ เมตฺตาเจโตวิมุตฺตี”ติ\footnote{เมตฺตาเจโตวิมุตฺติ (สพฺพตฺถ)}.}

\prose{อิธ ปน ภิกฺขเว ภิกฺขุ เอวํ วเทยฺย “กรุณา หิ โข เม เจโตวิมุตฺติ ภาวิตา พหุลีกตา ยานีกตา วตฺถุกตา อนุฏฺฐิตา ปริจิตา สุสมารทฺธา, อถ จ ปน เม วิเหสา จิตฺตํ ปริยาทาย ติฏฺฐตี”ติ.\csromanspacer{} โส “มา เหวนฺ”ติสฺส วจนีโย “มายสฺมา เอวํ อวจ, มา ภควนฺตํ อพฺภาจิกฺขิ, น หิ สาธุ ภควโต อพฺภกฺขานํ, น หิ ภควา เอวํ วเทยฺย.\csromanspacer{} อฏฺฐานเมตํ อาวุโส อนวกาโส, ยํ กรุณาย เจโตวิมุตฺติยา ภาวิตาย พหุลีกตาย ยานีกตาย วตฺถุกตาย อนุฏฺฐิตาย ปริจิตาย สุสมารทฺธาย, อถ จ ปนสฺส วิเหสา จิตฺตํ ปริยาทาย ฐสฺสติ, เนตํ ฐานํ วิชฺชติ, นิสฺสรณํ เหตํ อาวุโส วิเหสาย ยทิทํ กรุณาเจโตวิมุตฺตี”ติ.}

\gambhira{ฉกฺกนิปาตปาฬิ}
\prose{\csromanfolio{258}อิธ ปน ภิกฺขเว ภิกฺขุ เอวํ วเทยฺย “มุทิตา หิ โข เม เจโตวิมุตฺติ ภาวิตา พหุลีกตา ยานีกตา วตฺถุกตา อนุฏฺฐิตา ปริจิตา สุสมารทฺธา, อถ จ ปน เม อรติ จิตฺตํ ปริยาทาย ติฏฺฐตี”ติ.\csromanspacer{} โส “มา เหวนฺ”ติสฺส วจนีโย “มายสฺมา เอวํ อวจ, มา ภควนฺตํ อพฺภาจิกฺขิ, น หิ สาธุ ภควโต อพฺภกฺขานํ, น หิ ภควา เอวํ วเทยฺย.\csromanspacer{} อฏฺฐานเมตํ อาวุโส อนวกาโส, ยํ มุทิตาย เจโตวิมุตฺติยา ภาวิตาย พหุลีกตาย ยานีกตาย วตฺถุกตาย อนุฏฺฐิตาย ปริจิตาย สุสมารทฺธาย, อถ จ ปนสฺส อรติ จิตฺตํ ปริยาทาย ฐสฺสติ, เนตํ ฐานํ วิชฺชติ, นิสฺสรณํ เหตํ อาวุโส อรติยา ยทิทํ มุทิตาเจโตวิมุตฺตี”ติ.}

\prose{อิธ ปน ภิกฺขเว ภิกฺขุ เอวํ วเทยฺย “อุเปกฺขา หิ โข เม เจโตวิมุตฺติ ภาวิตา พหุลีกตา ยานีกตา วตฺถุกตา อนุฏฺฐิตา ปริจิตา สุสมารทฺธา, อถ จ ปน เม ราโค จิตฺตํ ปริยาทาย ติฏฺฐตี”ติ.\csromanspacer{} โส “มา เหวนฺ”ติสฺส วจนีโย “มายสฺมา เอวํ อวจ, มา ภควนฺตํ อพฺภาจิกฺขิ, น หิ สาธุ ภควโต อพฺภกฺขานํ, น หิ ภควา เอวํ วเทยฺย.\csromanspacer{} อฏฺฐานเมตํ อาวุโส อนวกาโส, ยํ อุเปกฺขาย เจโต วิมุตฺติยา ภาวิตาย พหุลีกตาย ยานีกตาย วตฺถุกตาย อนุฏฺฐิตาย ปริจิตาย สุสมารทฺธาย, อถ จ ปนสฺส ราโค จิตฺตํ ปริยาทาย ฐสฺสติ, เนตํ ฐานํ วิชฺชติ, นิสฺสรณํ เหตํ อาวุโส ราคสฺส ยทิทํ อุเปกฺขาเจโตวิมุตฺตี”ติ.}

\prose{อิธ ปน ภิกฺขเว ภิกฺขุ เอวํ วเทยฺย “อนิมิตฺตา หิ โข เม เจโตวิมุตฺติ ภาวิตา พหุลีกตา ยานีกตา วตฺถุกตา อนุฏฺฐิตา ปริจิตา สุสมารทฺธา, อถ จ ปน เม นิมิตฺตานุสาริ วิญฺญาณํ โหตี”ติ.\csromanspacer{} โส “มา เหวนฺ”ติสฺส วจนีโย “มายสฺมา เอวํ อวจ, มา ภควนฺตํ อพฺภาจิกฺขิ, น หิ สาธุ ภควโต อพฺภกฺขานํ, น หิ ภควา เอวํ วเทยฺย.\csromanspacer{} อฏฺฐานเมตํ อาวุโส อนวกาโส, ยํ อนิมิตฺตาย เจโต วิมุตฺติยา ภาวิตาย พหุลีกตาย ยานีกตาย วตฺถุกตาย อนุฏฺฐิตาย ปริจิตาย สุสมารทฺธาย, อถ จ ปนสฺส นิมิตฺตานุสาริ วิญฺญาณํ ภวิสฺสติ, เนตํ ฐานํ วิชฺชติ, นิสฺสรณํ เหตํ อาวุโส สพฺพนิมิตฺตานํ ยทิทํ อนิมิตฺตาเจโตวิมุตฺตี”ติ.}

\chapterheadpageread{2. สารณียวคฺค}
\prose{\csromanfolio{259}อิธ ปน ภิกฺขเว ภิกฺขุ เอวํ วเทยฺย “อสฺมีติ โข เม วิคตํ\footnote{วิคเต (สฺยา)}, อยมหมสฺมีติ จ\footnote{อยํ จกาโร ที 3. 207 ปิฏฺเฐ นตฺถิ.} น สมนุปสฺสามิ, อถ จ ปน เม วิจิกิจฺฉากถํกถาสลฺลํ จิตฺตํ ปริยาทาย ติฏฺฐตี”ติ.\csromanspacer{} โส “มา เหวนฺ”ติสฺส วจนีโย “มายสฺมา เอวํ อวจ, มา ภควนฺตํ อพฺภาจิกฺขิ, น หิ สาธุ ภควโต อพฺภกฺขานํ, น หิ ภควา เอวํ วเทยฺย.\csromanspacer{} อฏฺฐานเมตํ อาวุโส อนวกาโส, ยํ อสฺมีติ วิคเต ‘อยมหมสฺมี’ติ จ น สมนุปสฺสโต, อถ จ ปนสฺส วิจิกิจฺฉากถํกถาสลฺลํ จิตฺตํ ปริยาทาย ฐสฺสติ, เนตํ ฐานํ วิชฺชติ, นิสฺสรณํ เหตํ อาวุโส วิจิกิจฺฉากถํกถาสลฺลสฺส ยทิทํ ‘อสฺมี’ติ มานสมุคฺฆาโต”ติ.\csromanspacer{} อิมา โข ภิกฺขเว ฉ นิสฺสารณียา ธาตุโยติ.\csromanspacer{}\csromanspacer{}ตติยํ.}

\csromanmatikaanchor{mtk.137}
\csromantocmarkref{subsection}{4. ภทฺทกสุตฺต}{mtk.137}
\bookmark[dest={mtk.137},level=2]{4. ภทฺทกสุตฺต}
\csromanheader{2}{4. ภทฺทกสุตฺต}
\proseitem{14}{ตตฺร โข อายสฺมา สาริปุตฺโต ภิกฺขู อามนฺเตสิ “อาวุโส ภิกฺขเว”ติ. “อาวุโส”ติ โข เต ภิกฺขู อายสฺมโต สาริปุตฺตสฺส ปจฺจสฺโสสุํ.\csromanspacer{} อายสฺมา สาริปุตฺโต เอตทโวจ\csromandash{}}

\prose{ตถา ตถา อาวุโส ภิกฺขุ วิหารํ กปฺเปติ, ยถา ยถาสฺส วิหารํ กปฺปยโต น ภทฺทกํ มรณํ โหติ, น ภทฺทิกา กาลกิริยา\footnote{กาลํกิริยา (ก) อํ 1. 264 ปิฏฺเฐปิ. 4. น ปหาสิ (สี, สฺยา, กํ, อิ)}.\csromanspacer{} กถญฺจาวุโส ภิกฺขุ ตถา ตถา วิหารํ กปฺเปติ, ยถา ยถาสฺส วิหารํ กปฺปยโต น ภทฺทกํ มรณํ โหติ, น ภทฺทิกา กาลกิริยา.}

\prose{อิธาวุโส ภิกฺขุ กมฺมาราโม โหติ กมฺมรโต กมฺมารามตํ อนุยุตฺโต, ภสฺสาราโม โหติ ภสฺสรโต ภสฺสารามตํ อนุยุตฺโต, นิทฺทาราโม โหติ นิทฺทารโต นิทฺทารามตํ อนุยุตฺโต, สงฺคณิการาโม โหติ สงฺคณิกรโต สงฺคณิการามตํ อนุยุตฺโต สํสคฺคาราโม โหติ สํสคฺครโต สํสคฺคารามตํ อนุยุตฺโต ปปญฺจาราโม โหติ ปปญฺจรโต ปปญฺจารามตํ อนุยุตฺโต.\csromanspacer{} เอวํ โข อาวุโส ภิกฺขุ ตถา ตถา วิหารํ กปฺเปติ, ยถา ยถาสฺส วิหารํ กปฺปยโต น ภทฺทกํ มรณํ โหติ, น ภทฺทิกา กาลกิริยา.\csromanspacer{} อยํ วุจฺจตาวุโส ภิกฺขุ สกฺกายาภิรโต นปฺปชหาสิ4 สกฺกายํ สมฺมา ทุกฺขสฺส อนฺตกิริยาย.}

\gambhira{ฉกฺกนิปาตปาฬิ}
\prose{\csromanfolio{260}ตถา ตถาวุโส ภิกฺขุ วิหารํ กปฺเปติ, ยถา ยถาสฺส วิหารํ กปฺปยโต ภทฺทกํ มรณํ โหติ, ภทฺทิกา กาลกิริยา.\csromanspacer{} กถญฺจาวุโส ภิกฺขุ ตถา ตถา วิหารํ กปฺเปติ, ยถา ยถาสฺส วิหารํ กปฺปยโต ภทฺทกํ มรณํ โหติ, ภทฺทิกา กาลกิริยา.}

\prose{อิธาวุโส ภิกฺขุ น กมฺมาราโม โหติ น กมฺมรโต น กมฺมารามตํ อนุยุตฺโต, น ภสฺสาราโม โหติ น ภสฺสรโต น ภสฺสรามตํ อนุยุตฺโต, น นิทฺทาราโม โหติ น นิทฺทารโต น นิทฺทารามตํ อนุยุตฺโต, น สงฺคณิการาโม โหติ น สงฺคณิกรโต น สงฺคณิการามตํ อนุยุตฺโต, น สํสคฺคาราโม โหติ น สํสคฺครโต น สํสคฺคา รามตํ อนุยุตฺโต น ปปญฺจาราโม โหติ น ปปญฺจรโต น ปปญฺจารามตํ อนุยุตฺโต.\csromanspacer{} เอวํ โข อาวุโส ภิกฺขุ ตถา ตถา วิหารํ กปฺเปติ, ยถา ยถาสฺส วิหารํ กปฺปยโต ภทฺทกํ มรณํ โหติ, ภทฺทิกา กาลกิริยา.\csromanspacer{} อยํ วุจฺจตาวุโส ภิกฺขุ นิพฺพานาภิรโต ปชหาสิ สกฺกายํ สมฺมา ทุกฺขสฺส อนฺตกิริยายาติ.}

\csromangathameasure{โย ปปญฺจมนุยุตฺโต, ปปญฺจาภิรโต มโค. \\ วิราธยี โส นิพฺพานํ, โยคกฺเขมํ อนุตฺตรํ. \\ โย จ ปปญฺจํ หิตฺวาน, นิปฺปปญฺจปเท รโต. \\ อาราธยี โส นิพฺพานํ, โยคกฺเขมํ อนุตฺตรนฺติ.จตุตฺถํ.}
\csromangathagroup{\csromangathastanza{โย ปปญฺจมนุยุตฺโต, ปปญฺจาภิรโต มโค. \\ วิราธยี โส นิพฺพานํ, โยคกฺเขมํ อนุตฺตรํ.}\csromangathastanza{โย จ ปปญฺจํ หิตฺวาน, นิปฺปปญฺจปเท รโต. \\ อาราธยี โส นิพฺพานํ, โยคกฺเขมํ อนุตฺตรนฺติ.\csromanspacer{}\csromanspacer{}จตุตฺถํ.}}

\csromanmatikaanchor{mtk.138}
\csromantocmarkref{subsection}{5. อนุตปฺปิยสุตฺต}{mtk.138}
\bookmark[dest={mtk.138},level=2]{5. อนุตปฺปิยสุตฺต}
\csromanheader{2}{5. อนุตปฺปิยสุตฺต}
\proseitem{15}{ตตฺร โข อายสฺมา สาริปุตฺโต ภิกฺขู อามนฺเตสิ\csromandash{}ตถา ตถาวุโส ภิกฺขุ วิหารํ กปฺเปติ, ยถา ยถาสฺส วิหารํ กปฺปยโต กาลกิริยา อนุตปฺปา โหติ.\csromanspacer{} กถญฺจาวุโส ภิกฺขุ ตถา ตถา วิหารํ กปฺเปติ, ยถา ยถาสฺส วิหารํ กปฺปยโต กาลกิริยา อนุตปฺปา โหติ.}

\prose{อิธาวุโส ภิกฺขุ กมฺมาราโม โหติ กมฺมรโต กมฺมารามตํ อนุยุตฺโต.\csromanspacer{} ภสฺสาราโม โหติ.\csromanspacer{} นิทฺทาราโม โหติ.\csromanspacer{} สงฺคณิการาโม โหติ.\csromanspacer{} สํสคฺคาราโม โหติ.\csromanspacer{} ปปญฺจาราโม โหติ ปปญฺจรโต ปปญฺจารามตํ อนุยุตฺโต.\csromanspacer{} เอวํ โข อาวุโส ภิกฺขุ ตถา ตถา วิหารํ กปฺเปติ, ยถา ยถาสฺส วิหารํ กปฺปยโต กาลกิริยา อนุตปฺปา}

\vspace{\tipitakaparskip}
\csromancenter{สารณียวคฺค โหติ.\csromanspacer{} อยํ วุจฺจตาวุโส ภิกฺขุ สกฺกายาภิรโต นปฺปชหาสิ สกฺกายํ สมฺมา ทุกฺขสฺส อนฺตกิริยาย.}
\prose{\csromanfolio{261}ตถา ตถาวุโส ภิกฺขุ วิหารํ กปฺเปติ, ยถา ยถาสฺส วิหารํ กปฺปยโต กาลกิริยา อนนุตปฺปา โหติ.\csromanspacer{} กถญฺจาวุโส ภิกฺขุ ตถา ตถา วิหารํ กปฺเปติ, ยถา ยถาสฺส วิหารํ กปฺปยโต กาลกิริยา อนนุตปฺปา}

\prose{อิธาวุโส ภิกฺขุ น กมฺมาราโม โหติ น กมฺมรโต น กมฺมารามตํ อนุยุตฺโต.\csromanspacer{} น ภสฺสาราโม โหติ.\csromanspacer{} น นิทฺทาราโม โหติ.\csromanspacer{} น สงฺคณิการาโม โหติ.\csromanspacer{} น สํสคฺคาราโม โหติ.\csromanspacer{} น ปปญฺจาราโม โหติ น ปปญฺจรโต น ปปญฺจารามตํ อนุยุตฺโต.\csromanspacer{} เอวํ โข อาวุโส ภิกฺขุ ตถา ตถา วิหารํ กปฺเปติ, ยถา ยถาสฺส วิหารํ กปฺปยโต กาลกิริยา อนนุตปฺปา โหติ.\csromanspacer{} อยํ วุจฺจตาวุโส ภิกฺขุ นิพฺพานาภิรโต ปชหาสิ สกฺกายํ สมฺมา ทุกฺขสฺส อนฺตกิริยายาติ.}

\csromangathameasure{โย ปปญฺจมนุยุตฺโต, ปปญฺจาภิรโต มโค. \\ วิราธยี โส นิพฺพานํ, โยคกฺเขมํ อนุตฺตรํ. \\ โย จ ปปญฺจํ หิตฺวาน, นิปฺปปญฺจปเท รโต. \\ อาราธยี โส นิพฺพานํ, โยคกฺเขมํ อนุตฺตรนฺติ.ปญฺจมํ.}
\csromangathagroup{\csromangathastanza{โย ปปญฺจมนุยุตฺโต, ปปญฺจาภิรโต มโค. \\ วิราธยี โส นิพฺพานํ, โยคกฺเขมํ อนุตฺตรํ.}\csromangathastanza{โย จ ปปญฺจํ หิตฺวาน, นิปฺปปญฺจปเท รโต. \\ อาราธยี โส นิพฺพานํ, โยคกฺเขมํ อนุตฺตรนฺติ.\csromanspacer{}\csromanspacer{}ปญฺจมํ.}}

\csromanmatikaanchor{mtk.139}
\csromantocmarkref{subsection}{6. นกุลปิตุสุตฺต}{mtk.139}
\bookmark[dest={mtk.139},level=2]{6. นกุลปิตุสุตฺต}
\csromanheader{2}{6. นกุลปิตุสุตฺต}
\proseitem{16}{เอกํ สมยํ ภควา ภคฺเคสุ วิหรติ สุสุมารคิเร\footnote{สุํสุมารคิเร (สี, อิ), สํสุมารคิเร (กตฺถจิ)} เภสกฬาวเน มิคทาเย.\csromanspacer{} เตน โข ปน สมเยน นกุลปิตา คหปติ อาพาธิโก โหติ ทุกฺขิโต พาฬฺหคิลาโน.\csromanspacer{} อถ โข นกุลมาตา คหปตานี นกุลปิตรํ คหปติํ เอตทโวจ\csromandash{}}

\prose{มา โข ตฺวํ คหปติ สาเปกฺโข\footnote{สาเปโข (อิ, ก)} กาลมกาสิ.\csromanspacer{} ทุกฺขา คหปติ สาเปกฺขสฺส กาลกิริยา, ครหิตา จ ภควตา สาเปกฺขสฺส กาลกิริยา.\csromanspacer{} สิยา โข ปน เต คหปติ เอวมสฺส “น นกุลมาตา คหปตานี มมจฺจเยน สกฺขิสฺสติ\footnote{น สกฺขิสฺสติ (สี, สฺยา, กํ), สกฺโกติ (อิ, ก)} ทารเก โปเสตุํ ฆราวาสํ สนฺถริตุนฺ”ติ\footnote{สนฺธริตุนฺติ (ก), สณฺฐริตุํ (สฺยา)}.\csromanspacer{} น โข ปเนตํ คหปติ เอวํ ทฏฺฐพฺพํ.\csromanspacer{} กุสลาหํ คหปติ}

\vspace{\tipitakaparskip}
\csromancenter{ฉกฺกนิปาตปาฬิ กปฺปาสํ กนฺติตุํ เวณิํ โอลิขิตุํ, สกฺโกมหํ คหปติ ตวจฺจเยน ทารเก โปเสตุํ ฆราวาสํ สนฺถริตุํ.\csromanspacer{} ตสฺมาติห ตฺวํ คหปติ มา สาเปกฺโข กาลมกาสิ, ทุกฺขา คหปติ สาเปกฺขสฺส กาลกิริยา, ครหิตา จ ภควตา สาเปกฺขสฺส กาลํกิริยา. (1)}
\prose{\csromanfolio{262}สิยา โข ปน เต คหปติ เอวมสฺส “นกุลมาตา คหปตานี มมจฺจเยน อญฺญํ ฆรํ\footnote{ภตฺตารํ (สฺยา, กํ), วีรํ (สี)} คมิสฺสตี”ติ.\csromanspacer{} น โข ปเนตํ คหปติ เอวํ ทฏฺฐพฺพํ.\csromanspacer{} ตฺวํ เจว โข คหปติ ชานาสิ อหญฺจ, ยํ โน\footnote{ยทา เต (สี), ยถา (สฺยา), ยถา โน (อิ)} โสฬสวสฺสานิ คหฏฺฐกํ พฺรหฺมจริยํ สมาจิณฺณํ\footnote{สมาทินฺนํ (สี)}.\csromanspacer{} ตสฺมาติห ตฺวํ คหปติ มา สาเปกฺโข กาลมกาสิ, ทุกฺขา คหปติ สาเปกฺขสฺส กาลกิริยา, ครหิตา จ ภควตา สาเปกฺขสฺส กาลกิริยา. (2)}

\prose{สิยา โข ปน เต คหปติ เอวมสฺส “นกุลมาตา คหปตานี มมจฺจเยน น ทสฺสนกามา ภวิสฺสติ ภควโต, น ทสฺสนกามา ภิกฺขุสํฆสฺสา”ติ.\csromanspacer{} น โข ปเนตํ คหปติ เอวํ ทฏฺฐพฺพํ.\csromanspacer{} อหํ หิ คหปติ ตวจฺจเยน ทสฺสนกามตรา เจว ภวิสฺสามิ ภควโต, ทสฺสนกามตรา จ ภิกฺขุสํฆสฺส.\csromanspacer{} ตสฺมาติห ตฺวํ คหปติ มา สาเปกฺโข กาลมกาสิ, ทุกฺขา คหปติ สาเปกฺขสฺส กาลกิริยา, ครหิตา จ ภควตา สาเปกฺขสฺส กาลกิริยา. (3)}

\prose{สิยา โข ปน เต คหปติ เอวมสฺส “น นกุลมาตา คหปตานี มมจฺจเยน สีเลสุ\footnote{นกุลมาตา ... น สีเลสุ (สี, อิ)} ปริปูรการินี”ติ.\csromanspacer{} น โข ปเนตํ คหปติ เอวํ ทฏฺฐพฺพํ.\csromanspacer{} ยาวตา โข คหปติ ตสฺส ภควโต สาวิกา คิหี โอทาตวสนา สีเลสุ ปริปูรการินิโย อหํ ตาสํ อญฺญตรา.\csromanspacer{} ยสฺส โข ปนสฺส กงฺขา วา วิมติ วา, อยํ โส ภควา อรหํ สมฺมาสมฺพุทฺโธ ภคฺเคสุ วิหรติ สุสุมารคิเร เภสกฬาวเน มิคทาเย, ตํ ภควนฺตํ อุปสงฺกมิตฺวา ปุจฺฉตุ.\csromanspacer{} ตสฺมาติห ตฺวํ คหปติ มา สาเปกฺโข กาลมกาสิ, ทุกฺขา คหปติ สาเปกฺขสฺส กาลกิริยา, ครหิตา จ ภควตา สาเปกฺขสฺส กาลกิริยา. (4)}

\chapterheadpageread{2. สารณียวคฺค}
\prose{\csromanfolio{263}สิยา โข ปน เต คหปติ เอวมสฺส “น นกุลมาตา คหปตานี ลาภินี\footnote{นกุลมาตา คหปตานี น ลาภินี (อิ)} อชฺฌตฺตํ เจโตสมถสฺสา”ติ.\csromanspacer{} น โข ปเนตํ คหปติ เอวํ ทฏฺฐพฺพํ.\csromanspacer{} ยาวตา โข คหปติ ตสฺส ภควโต สาวิกา คิหี โอทาตวสนา ลาภินิโย อชฺฌตฺตํ เจโตสมถสฺส, อหํ ตาสํ อญฺญตรา.\csromanspacer{} ยสฺส โข ปนสฺส กงฺขา วา วิมติ วา, อยํ โส ภควา อรหํ สมฺมาสมฺพุทฺโธ ภคฺเคสุ วิหรติ สุสุมารคิเร เภสกฬาวเน มิคทาเย, ตํ ภควนฺตํ อุปสงฺกมิตฺวา ปุจฺฉตุ.\csromanspacer{} ตสฺมาติห ตฺวํ คหปติ มา สาเปกฺโข กาลมกาสิ, ทุกฺขา คหปติ สาเปกฺขสฺส กาลกิริยา, ครหิตา จ ภควตา สาเปกฺขสฺส กาลกิริยา. (5)}

\prose{สิยา โข ปน เต คหปติ เอวมสฺส “น นกุลมาตา คหปตานี อิมสฺมิํ ธมฺมวินเย โอคาธปฺปตฺตา ปติคาธปฺปตฺตา อสฺสาสปฺปตฺตา ติณฺณวิจิกิจฺฉา วิคตกถํกถา เวสารชฺชปฺปตฺตา อปรปฺปจฺจยา สตฺถุสาสเน วิหรตี”ติ.\csromanspacer{} น โข ปเนตํ คหปติ เอวํ ทฏฺฐพฺพํ.\csromanspacer{} ยาวตา โข คหปติ ตสฺส ภควโต สาวิกา คิหี โอทาตวสนา อิมสฺมิํ ธมฺมวินเย โอคาธปฺปตฺตา ปติคาธปฺปตฺตา อสฺสาสปฺปตฺตา ติณฺณวิจิกิจฺฉา วิคตกถํกถา เวสารชฺชปฺปตฺตา อปรปฺปจฺจยา สตฺถุสาสเน วิหรนฺติ, อหํ ตาสํ อญฺญตรา.\csromanspacer{} ยสฺส โข ปนสฺส กงฺขา วา วิมติ วา, อยํ โส ภควา อรหํ สมฺมาสมฺพุทฺโธ ภคฺเคสุ วิหรติ สุสุมารคิเร เภสกฬาวเน มิคทาเย, ตํ ภควนฺตํ อุปสงฺกมิตฺวา ปุจฺฉตุ.\csromanspacer{} ตสฺมาติห ตฺวํ คหปติ มา สาเปกฺโข กาลมกาสิ, ทุกฺขา คหปติ สาเปกฺขสฺส กาลกิริยา, ครหิตา จ ภควตา สาเปกฺขสฺส กาลกิริยาติ. (6)}

\prose{อถ โข นกุลปิตุโน คหปติสฺส นกุลมาตรา\footnote{นกุลมาตาย (สี, สฺยา), นกุลมาตุยา (ก)} คหปตานิยา อิมินา โอวาเทน โอวทิยมานสฺส โส อาพาโธ ฐานโส ปฏิปฺปสฺสมฺภิ.\csromanspacer{} วุฏฺฐหิ\footnote{วุฏฺฐาติ (ก)} จ นกุลปิตา คหปติ ตมฺหา อาพาธา, ตถาปหีโน จ ปน นกุลปิตุโน คหปติสฺส โส อาพาโธ อโหสิ.\csromanspacer{} อถ โข นกุลปิตา คหปติ คิลานา วุฏฺฐิโต\footnote{“คิลานภาวโต วุฏฺฐาย ฐิโต, ภาวปฺปธาโน หิ อยํ นิทฺเทโส”ติ ฏีกาสํวณฺณนา.} อจิรวุฏฺฐิโต เคลญฺญา}

\vspace{\tipitakaparskip}
\csromancenter{ฉกฺกนิปาตปาฬิ ทณฺฑโมลุพฺภ เยน ภควา เตนุปสงฺกมิ, อุปสงฺกมิตฺวา ภควนฺตํ อภิวาเทตฺวา เอกมนฺตํ นิสีทิ, เอกมนฺตํ นิสินฺนํ โข นกุลปิตรํ คหปติํ ภควา เอตทโวจ\csromandash{}}
\prose{\csromanfolio{264}ลาภา เต คหปติ, สุลทฺธํ เต คหปติ, ยสฺส เต นกุลมาตา คหปตานี อนุกมฺปิกา อตฺถกามา โอวาทิกา อนุสาสิกา.\csromanspacer{} ยาวตา โข คหปติ มม สาวิกา คิหี โอทาตวสนา สีเลสุ ปริปูรการินิโย, นกุลมาตา คหปตานี ตาสํ อญฺญตรา.\csromanspacer{} ยาวตา โข คหปติ มม สาวิกา คิหี โอทาตวสนา ลาภินิโย อชฺฌตฺตํ เจโตสมถสฺส, นกุลมาตา คหปตานี ตาสํ อญฺญตรา.\csromanspacer{} ยาวตา โข คหปติ มม สาวิกา คิหี โอทาตวสนา อิมสฺมิํ ธมฺมวินเย โอคาธปฺปตฺตา ปติคาธปฺปตฺตา อสฺสาสปฺปตฺตา ติณฺณวิจิกิจฺฉา วิคตกถํกถา เวสารชฺชปฺปตฺตา อปรปฺปจฺจยา สตฺถุสาสเน วิหรนฺติ, นกุลมาตา คหปตานี ตาสํ อญฺญตรา.\csromanspacer{} ลาภา เต คหปติ, สุลทฺธํ เต คหปติ, ยสฺส เต นกุลมาตา คหปตานี อนุกมฺปิกา อตฺถกามา โอวาทิกา อนุสาสิกาติ.\csromanspacer{}\csromanspacer{}ฉฏฺฐํ.}

\csromanmatikaanchor{mtk.140}
\csromantocmarkref{subsection}{7. โสปฺปสุตฺต}{mtk.140}
\bookmark[dest={mtk.140},level=2]{7. โสปฺปสุตฺต}
\csromanheader{2}{7. โสปฺปสุตฺต}
\proseitem{17}{เอกํ สมยํ ภควา สาวตฺถิยํ วิหรติ เชตวเน อนาถปิณฺฑิกสฺส อาราเม.\csromanspacer{} อถ โข ภควา สายนฺหสมยํ ปฏิสลฺลานา วุฏฺฐิโต เยนุปฏฺฐานสาลา เตนุปสงฺกมิ, อุปสงฺกมิตฺวา ปญฺญตฺเต อาสเน นิสีทิ.\csromanspacer{} อายสฺมาปิ โข สาริปุตฺโต สายนฺหสมยํ ปฏิสลฺลานา วุฏฺฐิโต เยนุปฏฺฐานสาลา เตนุปสงฺกมิ, อุปสงฺกมิตฺวา ภควนฺตํ อภิวาเทตฺวา เอกมนฺตํ นิสีทิ.\csromanspacer{} อายสฺมาปิ โข มหาโมคฺคลฺลาโน.\csromanspacer{} อายสฺมาปิ โข มหากสฺสโป.\csromanspacer{} อายสฺมาปิ โข มหากจฺจาโน.\csromanspacer{} อายสฺมาปิ โข มหาโกฏฺฐิโก\footnote{มหาโกฏฺฐิโต (สี, อิ)}.\csromanspacer{} อายสฺมาปิ โข มหาจุนฺโท.\csromanspacer{} อายสฺมาปิ โข มหากปฺปิโน.\csromanspacer{} อายสฺมาปิ โข อนุรุทฺโธ.\csromanspacer{} อายสฺมาปิ โข เรวโต.\csromanspacer{} อายสฺมาปิ โข อานนฺโท สายนฺหสมยํ ปฏิสลฺลานา วุฏฺฐิโต เยนุปฏฺฐานสาลา เตนุปสงฺกมิ, อุปสงฺกมิตฺวา ภควนฺตํ อภิวาเทตฺวา เอกมนฺตํ นิสีทิ.\csromanspacer{} อถ โข ภควา พหุเทว รตฺติํ นิสชฺชาย วีตินาเมตฺวา อุฏฺฐายาสนา วิหารํ ปาวิสิ.\csromanspacer{} เตปิ โข}

\vspace{\tipitakaparskip}
\csromancenter{สารณียวคฺค อายสฺมนฺโต อจิรปกฺกนฺตสฺส ภควโต อุฏฺฐายาสนา ยถาวิหารํ อคมํสุ.\csromanspacer{} เย ปน ตตฺถ ภิกฺขู นวา อจิรปพฺพชิตา อธุนาคตา อิมํ ธมฺมวินยํ, เต ยาว สูริยุคฺคมนา กากจฺฉมานา สุปิํสุ.\csromanspacer{} อทฺทสา โข ภควา ทิพฺเพน จกฺขุนา วิสุทฺเธน อติกฺกนฺตมานุสเกน เต ภิกฺขู ยาว สูริยุคฺคมนา กากจฺฉมาเน สุปนฺเต, ทิสฺวา เยนุปฏฺฐานสาลา เตนุปสงฺกมิ, อุปสงฺกมิตฺวา ปญฺญตฺเต อาสเน นิสีทิ, นิสชฺช โข ภควา ภิกฺขู อามนฺเตสิ\csromandash{}}
\prose{\csromanfolio{265}กหํ นุ โข ภิกฺขเว สาริปุตฺโต, กหํ มหาโมคฺคลฺลาโน, กหํ มหากสฺสโป, กหํ มหากจฺจาโน, กหํ มหาโกฏฺฐิโก, กหํ มหาจุนฺโท, กหํ มหากปฺปิโน, กหํ อนุรุทฺโธ, กหํ เรวโต, กหํ อานนฺโท, กหํ นุ โข เต ภิกฺขเว เถรา สาวกา คตาติ.\csromanspacer{} เตปิ โข ภนฺเต อายสฺมนฺโต อจิรปกฺกนฺตสฺส ภควโต อุฏฺฐายาสนา ยถาวิหารํ อคมํสูติ.\csromanspacer{} เตน โน\footnote{เกน โน (ก), เก นุ (กตฺถจิ)} ตุเมฺห ภิกฺขเว “เถรา ภิกฺขู นาคตา”ติ\footnote{ภิกฺขู นวา (สี, สฺยา, กํ, อิ), ภิกฺขู คตาติ (?)} ยาว สูริยุคฺคมนา กากจฺฉมานา สุปถ.\csromanspacer{} ตํ กิํ มญฺญถ ภิกฺขเว, อปิ นุ ตุเมฺหหิ ทิฏฺฐํ วา สุตํ วา “ราชา ขตฺติโย มุทฺธาวสิตฺโต\footnote{มุทฺธาภิสิตฺโต (ก)} ยาวทตฺถํ เสยฺยสุขํ ปสฺสสุขํ มิทฺธสุขํ อนุยุตฺโต วิหรนฺโต ยาวชีวํ รชฺชํ กาเรนฺโต ชนปทสฺส วา ปิโย มนาโป”ติ.\csromanspacer{} โน เหตํ ภนฺเต.\csromanspacer{} สาธุ ภิกฺขเว, มยาปิ โข เอตํ ภิกฺขเว เนว ทิฏฺฐํ น สุตํ “ราชา ขตฺติโย มุทฺธาวสิตฺโต ยาวทตฺถํ เสยฺยสุขํ ปสฺสสุขํ มิทฺธสุขํ อนุยุตฺโต วิหรนฺโต ยาวชีวํ รชฺชํ กาเรนฺโต ชนปทสฺส วา ปิโย มนาโป”ติ.}

\prose{ตํ กิํ มญฺญถ ภิกฺขเว, อปิ นุ ตุเมฺหหิ ทิฏฺฐํ วา สุตํ วา “รฏฺฐิโก ฯเปฯ เปตฺตณิโก.\csromanspacer{} เสนาปติโก.\csromanspacer{} คามคามณิโก\footnote{คามคามิโก (สี, อิ)}.\csromanspacer{} ปูคคามณิโก ยาวทตฺถํ เสยฺยสุขํ ปสฺสสุขํ มิทฺธสุขํ อนุยุตฺโต วิหรนฺโต ยาวชีวํ ปูคคามณิกตฺตํ กาเรนฺโต ปูคสฺส วา ปิโย มนาโป”ติ.\csromanspacer{} โน เหตํ ภนฺเต.\csromanspacer{} สาธุ ภิกฺขเว, มยาปิ โข เอตํ ภิกฺขเว เนว ทิฏฺฐํ น สุตํ “ปูคคามณิโก ยาวทตฺถํ เสยฺยสุขํ ปสฺสสุขํ มิทฺธสุขํ อนุยุตฺโต วิหรนฺโต ยาวชีวํ ปูคคามณิกตฺตํ วา กาเรนฺโต ปูคสฺส วา ปิโย มนาโป”ติ.}

\gambhira{ฉกฺกนิปาตปาฬิ}
\prose{\csromanfolio{266}ตํ กิํ มญฺญถ ภิกฺขเว, อปิ นุ ตุเมฺหหิ ทิฏฺฐํ วา สุตํ วา “สมโณ วา พฺราหฺมโณ วา ยาวทตฺถํ เสยฺยสุขํ ปสฺสสุขํ มิทฺธสุขํ อนุยุตฺโต อินฺทฺริเยสุ อคุตฺตทฺวาโร โภชเน อมตฺตญฺญู ชาคริยํ อนนุยุตฺโต อวิปสฺสโก กุสลานํ ธมฺมานํ ปุพฺพรตฺตาปรรตฺตํ โพธิปกฺขิยานํ\footnote{โพธปกฺขิยานํ (สี), โพธปกฺขิกานํ (อิ)} ธมฺมานํ ภาวนานุ โยคํ อนนุยุตฺโต อาสวานํ ขยา อนาวสํ เจโตวิมุตฺติํ ปญฺญาวิมุตฺติํ ทิฏฺเฐว ธมฺเม สยํ อภิญฺญา สจฺฉิกตฺวา อุปสมฺปชฺช วิหรนฺโต”ติ.\csromanspacer{} โน เหตํ ภนฺเต.\csromanspacer{} สาธุ ภิกฺขเว, มยาปิ โข เอตํ ภิกฺขเว เนว ทิฏฺฐํ น สุตํ “สมโณ วา พฺราหฺมโณ วา ยาวทตฺถํ เสยฺยสุขํ ปสฺสสุขํ มิทฺธสุขํ อนุยุตฺโต อินฺทฺริเยสุ อคุตฺตทฺวาโร โภชเน อมตฺตญฺญู ชาคริยํ อนนุยุตฺโต อวิปสฺสโก กุสลานํ ธมฺมานํ ปุพฺพรตฺตาปรรตฺตํ โพธิปกฺขิยานํ ธมฺมานํ ภาวนานุโยคํ อนนุยุตฺโต อาสวานํ ขยา อนาสวํ เจโตวิมุตฺติํ ปญฺญาวิมุตฺติํ ทิฏฺเฐว ธมฺเม สยํ อภิญฺญา สจฺฉิกตฺวา อุปสมฺปชฺช วิหรนฺโต”ติ.}

\prose{ตสฺมาติห ภิกฺขเว เอวํ สิกฺขิตพฺพํ “อินฺทฺริเยสุ คุตฺตทฺวารา ภวิสฺสาม, โภชเน มตฺตญฺญุโน ชาคริยํ อนุยุตฺตา วิปสฺสกา กุสลานํ ธมฺมานํ ปุพฺพรตฺตาปรรตฺตํ โพธิปกฺขิยานํ ธมฺมานํ ภาวนานุโยคมนุยุตฺตา วิหริสฺสามา”ติ.\csromanspacer{} เอวํ หิ โว ภิกฺขเว สิกฺขิตพฺพนฺติ.\csromanspacer{}\csromanspacer{}สตฺตมํ.}

\csromanmatikaanchor{mtk.141}
\csromantocmarkref{subsection}{8. มจฺฉพนฺธสุตฺต}{mtk.141}
\bookmark[dest={mtk.141},level=2]{8. มจฺฉพนฺธสุตฺต}
\csromanheader{2}{8. มจฺฉพนฺธสุตฺต}
\proseitem{18}{เอกํ สมยํ ภควา โกสเลสุ จาริกํ จรติ มหตา ภิกฺขุสํเฆน สทฺธิํ.\csromanspacer{} อทฺทสา โข ภควา อทฺธานมคฺคปฺปฏิปนฺโน อญฺญตรสฺมิํ ปเทเส มจฺฉิกํ มจฺฉพนฺธํ มจฺเฉ วธิตฺวา วธิตฺวา วกฺกิณมานํ, ทิสฺวา มคฺคา โอกฺกมฺม อญฺญตรสฺมิํ รุกฺขมูเล ปญฺญตฺเต อาสเน นิสีทิ, นิสชฺช โข ภควา ภิกฺขู อามนฺเตสิ “ปสฺสถ โน ตุเมฺห ภิกฺขเว อมุํ มจฺฉิกํ มจฺฉพนฺธํ มจฺเฉ วธิตฺวา วธิตฺวา วิกฺกิณมานนฺ”ติ.\csromanspacer{} เอวํ ภนฺเต.}

\prose{ตํ กิํ มญฺญถ ภิกฺขเว, อปิ นุ ตุเมฺหหิ ทิฏฺฐํ วา สุตํ วา “มจฺฉิโก มจฺฉพนฺโธ มจฺเฉ วธิตฺวา วธิตฺวา วิกฺกิณมาโน เตน กมฺเมน เตน อาชีเวน หตฺถิยายี วา อสฺสยายี วา รถยายี วา ยานยายี วา โภคโภคี วา มหนฺตํ วา โภคกฺขนฺธํ อชฺฌาวสนฺโต”ติ.\csromanspacer{} โน เหตํ}

\vspace{\tipitakaparskip}
\csromancenter{สารณียวคฺค ภนฺเต.\csromanspacer{} สาธุ ภิกฺขเว, มยาปิ โข เอตํ ภิกฺขเว เนว ทิฏฺฐํ น สุตํ “มจฺฉิโก มจฺฉพนฺโธ มจฺเฉ วธิตฺวา วธิตฺวา วิกฺกิณมาโน เตน กมฺเมน เตน อาชีเวน หตฺถิยายี วา อสฺสยายี วา รถยายี วา ยานยายี วา โภคโภคี วา มหนฺตํ วา โภคกฺขนฺธํ อชฺฌาวสนฺโต”ติ.\csromanspacer{} ตํ กิสฺส เหตุ? เต หิ โส ภิกฺขเว มจฺเฉ วชฺเฌ วธายุปนีเต\footnote{วธายานีเต (สฺยา, กํ), วธาย นีเต (ก)} ปาปเกน มนสานุเปกฺขติ, ตสฺมา โส เนว หตฺถิยายี โหติ น อสฺสยายี น รถยายี น ยานยายี น โภคโภคี น มหนฺตํ โภคกฺขนฺธํ อชฺฌาวสติ.}
\prose{\csromanfolio{267}ตํ กิํ มญฺญถ ภิกฺขเว, อปิ นุ ตุเมฺหหิ ทิฏฺฐํ วา สุตํ วา “โคฆาตโก คาโว วธิตฺวา วธิตฺวา วิกฺกิณมาโน เตน กมฺเมน เตน อาชีเวน หตฺถิยายี วา อสฺสยายี วา รถยายี วา ยานยายี วา โภคโภคี วา มหนฺตํ วา โภคกฺขนฺธํ อชฺฌาวสนฺโต”ติ.\csromanspacer{} โน เหตํ ภนฺเต.\csromanspacer{} สาธุ ภิกฺขเว, มยาปิ โข เอตํ ภิกฺขเว เนว ทิฏฺฐํ น สุตํ “โคฆาตโก คาโว วธิตฺวา วธิตฺวา วิกฺกิณมาโน เตน กมฺเมน เตน อาชีเวน หตฺถิยายี วา อสฺสยายี วา รถยายี วา ยานยายี วา โภคโภคี วา มหนฺตํ วา โภคกฺขนฺธํ อชฺฌาวสนฺโต”ติ.\csromanspacer{} ตํ กิสฺส เหตุ? เต หิ โส ภิกฺขเว คาโว วชฺเฌ วธายุปนีเต ปาปเกน มนสานุเปกฺขติ.\csromanspacer{} ตสฺมา โส เนว หตฺถิยายี โหติ น อสฺสยายี น รถยายี น ยานยายี น โภคโภคี น มหนฺตํ โภคกฺขนฺธํ อชฺฌาวสติ.}

\prose{ตํ กิํ มญฺญถ ภิกฺขเว, อปิ นุ ตุเมฺหหิ ทิฏฺฐํ วา สุตํ วา “โอรพฺภิโก ฯเปฯ สูกริโก\footnote{โสกริโก (สฺยา)} ฯเปฯ สากุณิโก ฯเปฯ มาควิโก มเค\footnote{มิเค (สฺยา, กํ)} วธิตฺวา วธิตฺวา วิกฺกิณมาโน เตน กมฺเมน เตน อาชีเวน หตฺถิยายี วา อสฺสยายี วา รถยายี วา ยานยายี วา โภคโภคี วา มหนฺตํ วา โภคกฺขนฺธํ อชฺฌาวสนฺโต”ติ.\csromanspacer{} โน เหตํ ภนฺเต.\csromanspacer{} สาธุ ภิกฺขเว, มยาปิ โข เอตํ ภิกฺขเว เนว ทิฏฺฐํ น สุตํ “มาควิโก มเค วธิตฺวา วธิตฺวา วิกฺกิณมาโน เตน กมฺเมน เตน อาชีเวน หตฺถิยายี วา อสฺสยายี วา รถยายี วา ยานยายี วา โภคโภคี วา มหนฺตํ วา โภคกฺขนฺธํ อชฺฌาวสนฺโต”ติ.\csromanspacer{} ตํ กิสฺส เหตุ? เต หิ โส ภิกฺขเว มเค วชฺเฌ วธายุปนีเต ปาปเกน มนสานุเปกฺขติ.\csromanspacer{} ตสฺมา โส เนว หตฺถิยายี โหติ น อสฺสยายี น รถยายี น}

\vspace{\tipitakaparskip}
\csromancenter{ฉกฺกนิปาตปาฬิ ยานยายี น โภคโภคี น มหนฺตํ โภคกฺขนฺธํ อชฺฌาวสติ.\csromanspacer{} เตหิ (นาม)\footnote{( ) พหูสุ นตฺถิ.} โส ภิกฺขเว ติรจฺฉานคเต ปาเณ วชฺเฌ วธายุปนีเต ปาปเกน มนสานุเปกฺขมาโน\footnote{มนสานุเปกฺขติ, ตสฺมา โส (สฺยา, ก)} เนว หตฺถิยายี ภวิสฺสติ\footnote{โหติ (สฺยา, ก)} น อสฺสยายี น รถยายี น ยานยายี น โภคโภคี น มหนฺตํ โภคกฺขนฺธํ อชฺฌาวสิสฺสติ\footnote{อชฺฌาวสติ (สฺยา, ก)}, โก ปน วาโท ยํ มนุสฺสภูตํ วชฺฌํ วธายุปนีตํ ปาปเกน มนสานุเปกฺขติ.\csromanspacer{} ตํ หิ ตสฺส\footnote{ตํ หิสฺส (อิ, ก)} ภิกฺขเว โหติ ทีฆรตฺตํ อหิตาย ทุกฺขาย.\csromanspacer{} กายสฺส เภทา ปรํ มรณา อปายํ ทุคฺคติํ วินิปาตํ นิรยํ อุปปชฺชตีติ.\csromanspacer{}\csromanspacer{}อฏฺฐมํ.}
\csromanheader{2}{9. ปฐมมรณสฺสติสุตฺต}
\csromanmatikaanchor{mtk.142}
\csromantocmarkref{subsection}{9-10. มรณสฺสติสุตฺต}{mtk.142}
\bookmark[dest={mtk.142},level=2]{9-10. มรณสฺสติสุตฺต}
\proseitem{19}{\csromanfolio{268}เอกํ สมยํ ภควา นาติเก\footnote{นาทิเก (สี, สฺยา, กํ, อิ) อํ 3. 135 ปิฏฺเฐปิ.} วิหรติ คิญฺชกาวสเถ.\csromanspacer{} ตตฺร โข ภควา ภิกฺขู อามนฺเตสิ “ภิกฺขโว”ติ. “ภทนฺเต”ติ เต ภิกฺขู ภควโต ปจฺจสฺโสสุํ.\csromanspacer{} ภควา เอตทโวจ “มรณสฺสติ ภิกฺขเว ภาวิตา พหุลีกตา มหปฺผลา โหติ มหานิสํสา อมโตคธา อมตปริโยสานา, ภาเวถ โน ตุเมฺห ภิกฺขเว มรณสฺสตินฺ”ติ.}

\prose{เอวํ วุตฺเต อญฺญตโร ภิกฺขุ ภควนฺตํ เอตทโวจ “อหํ โข ภนฺเต ภาเวมิ มรณสฺสตินฺ”ติ.\csromanspacer{} ยถา กถํ ปน ตฺวํ ภิกฺขุ ภาเวสิ มรณสฺสตินฺติ.\csromanspacer{} อิธ มยฺหํ ภนฺเต เอวํ โหติ “อโห วตาหํ รตฺตินฺทิวํ ชีเวยฺยํ, ภควโต สาสนํ มนสิ กเรยฺยํ, พหุ วต เม กตํ อสฺสา”ติ.\csromanspacer{} เอวํ โข อหํ ภนฺเต ภาเวมิ มรณสฺสตินฺติ.}

\prose{อญฺญตโรปิ โข ภิกฺขุ ภควนฺตํ เอตทโวจ “อหมฺปิ โข ภนฺเต ภาเวมิ มรณสฺสตินฺ”ติ.\csromanspacer{} ยถา กถํ ปน ตฺวํ ภิกฺขุ ภาเวสิ มรณสฺสตินฺติ.\csromanspacer{} อิธ มยฺหํ ภนฺเต เอวํ โหติ “อโห วตาหํ ทิวสํ ชีเวยฺยํ, ภควโต สาสนํ มนสิ กเรยฺยํ, พหุ วต เม กตํ อสฺสา”ติ.\csromanspacer{} เอวํ โข อหํ ภนฺเต ภาเวมิ มรณสฺสตินฺติ.}

\prose{อญฺญตโรปิ โข ภิกฺขุ ภควนฺตํ เอตทโวจ “อหมฺปิ โข ภนฺเต ภาเวมิ มรณสฺสตินฺ”ติ.\csromanspacer{} ยถา กถํ ปน ตฺวํ ภิกฺขุ ภาเวสิ}

\vspace{\tipitakaparskip}
\csromancenter{สารณียวคฺค มรณสฺสตินฺติ.\csromanspacer{} อิธ มยฺหํ ภนฺเต เอวํ โหติ “อโห วตาหํ ตทนฺตรํ ชีเวยฺยํ, ยทนฺตรํ เอกปิณฺฑปาตํ ภุญฺชามิ, ภควโต สาสนํ มนสิ กเรยฺยํ, พหุ วต เม กตํ อสฺสา”ติ.\csromanspacer{} เอวํ โข อหํ ภนฺเต ภาเวมิ มรณสฺสตินฺติ.}
\prose{\csromanfolio{269}อญฺญตโรปิ โข ภิกฺขุ ภควนฺตํ เอตทโวจ “อหมฺปิ โข ภนฺเต ภาเวมิ มรณสฺสตินฺ”ติ.\csromanspacer{} ยถา กถํ ปน ตฺวํ ภิกฺขุ ภาเวสิ มรณสฺสตินฺติ.\csromanspacer{} อิธ มยฺหํ ภนฺเต เอวํ โหติ “อโห วตาหํ ตทนฺตรํ ชีเวยฺยํ, ยทนฺตรํ จตฺตาโร ปญฺจ อาโลเป สงฺขาทิตฺวา\footnote{สงฺขริตฺวา (ก)} อชฺโฌหรามิ, ภควโต สาสนํ มนสิ กเรยฺยํ, พหุ วต เม กตํ อสฺสา”ติ.\csromanspacer{} เอวํ โข อหํ ภนฺเต ภาเวมิ มรณสฺสตินฺติ.}

\prose{อญฺญตโรปิ โข ภิกฺขุ ภควนฺตํ เอตทโวจ “อหมฺปิ โข ภนฺเต ภาเวมิ มรณสฺสตินฺ”ติ.\csromanspacer{} ยถา กถํ ปน ตฺวํ ภิกฺขุ ภาเวสิ มรณสฺสตินฺติ.\csromanspacer{} อิธ มยฺหํ ภนฺเต เอวํ โหติ “อโห วตาหํ ตทนฺตรํ ชีเวยฺยํ, ยทนฺตรํ เอกํ อาโลปํ สงฺขาทิตฺวา\footnote{สํหริตฺวา (ก)} อชฺโฌหรามิ, ภควโต สาสนํ มนสิ กเรยฺยํ.\csromanspacer{} พหุ วต เม กตํ อสฺสา”ติ.\csromanspacer{} เอวํ โข อหํ ภนฺเต ภาเวมิ มรณสฺสตินฺติ.}

\prose{อญฺญตโรปิ โข ภิกฺขุ ภควนฺตํ เอตทโวจ “อหมฺปิ โข ภนฺเต ภาเวมิ มรณสฺสตินฺ”ติ.\csromanspacer{} ยถา กถํ ปน ตฺวํ ภิกฺขุ ภาเวสิ มรณสฺสตินฺติ.\csromanspacer{} อิธ มยฺหํ ภนฺเต เอวํ โหติ “อโห วตาหํ ตทนฺตรํ ชีเวยฺยํ, ยทนฺตรํ อสฺสสิตฺวา วา ปสฺสสามิ, ปสฺสสิตฺวา วา อสฺสสามิ, ภควโต สาสนํ มนสิ กเรยฺยํ, พหุ วต เม กตํ อสฺสา”ติ.\csromanspacer{} เอวํ โข อหํ ภนฺเต ภาเวมิ มรณสฺสตินฺติ.}

\prose{เอวํ วุตฺเต ภควา เต ภิกฺขู เอตทโวจ “โย จายํ\footnote{ยฺวายํ (อิ, ก)} ภิกฺขเว ภิกฺขุ เอวํ มรณสฺสติํ ภาเวติ ‘อโห วตาหํ รตฺตินฺทิวํ ชีเวยฺยํ, ภควโต สาสนํ มนสิ กเรยฺยํ, พหุ วต เม กถํ อสฺสา”ติ.}

\prose{โย จายํ\footnote{โยปายํ (ก)} ภิกฺขเว ภิกฺขุ เอวํ มรณสฺสติํ ภาเวติ “อโห วตาหํ ทิวสํ ชีเวยฺยํ, ภควโต สาสนํ มนสิ กเรยฺยํ, พหุ วต เม กตํ อสฺสา”ติ.}

\gambhira{ฉกฺกนิปาตปาฬิ}
\prose{\csromanfolio{270}โย จายํ ภิกฺขเว ภิกฺขุ เอวํ มรณสฺสติํ ภาเวติ “อโห วตาหํ ตทนฺตรํ ชีเวยฺยํ, ยทนฺตรํ เอกปิณฺฑปาตํ ภุญฺชามิ, ภควโต สาสนํ มนสิ กเรยฺยํ, พหุ วต เม กตํ อสฺสา”ติ.}

\prose{โย จายํ ภิกฺขเว ภิกฺขุ เอวํ มรณสฺสติํ ภาเวติ “อโห วตาหํ ตทนฺตรํ ชีเวยฺยํ, ยทนฺตรํ จตฺตาโร ปญฺจ อาโลเป สงฺขาทิตฺวา อชฺโฌหรามิ, ภควโต สาสนํ มนสิ กเรยฺยํ, พหุ วต เม กตํ อสฺสา”ติ.\csromanspacer{} อิเม วุจฺจนฺติ ภิกฺขเว ภิกฺขู ปมตฺตา วิหรนฺติ, ทนฺธํ มรณสฺสติํ ภาเวนฺติ อาสวานํ ขยาย.}

\prose{โย จ ขฺวายํ\footnote{โยปายํ (ก)} ภิกฺขเว ภิกฺขุ เอวํ มรณสฺสติํ ภาเวติ “อโห วตาหํ ตทนฺตรํ ชีเวยฺยํ, ยทนฺตรํ เอกํ อาโลปํ สงฺขาทิตฺวา อชฺโฌหรามิ, ภควโต สาสนํ มนสิ กเรยฺยํ, พหุ วต เม กตํ อสฺสา”ติ.}

\prose{โย จายํ ภิกฺขเว ภิกฺขุ เอวํ มรณสฺสติํ ภาเวติ “อโห วตาหํ ตทนฺตรํ ชีเวยฺยํ, ยทนฺตรํ อสฺสสิตฺวา วา ปสฺสสามิ, ปสฺสสิตฺวา วา อสฺสสามิ, ภควโต สาสนํ มนสิ กเรยฺยํ, พหุ วต เม กตํ อสฺสา”ติ.\csromanspacer{} อิเม วุจฺจนฺติ ภิกฺขเว ภิกฺขู อปฺปมตฺตา วิหรนฺติ, ติกฺขํ มรณสฺสติํ ภาเวนฺติ อาสวานํ ขยาย.}

\prose{ตสฺมาติห ภิกฺขเว เอวํ สิกฺขิตพฺพํ “อปฺปมตฺตา วิหริสฺสาม, ติกฺขํ มรณสฺสติํ ภาเวสฺสาม อาสวานํ ขยายา”ติ.\csromanspacer{} เอวํ หิ โว ภิกฺขเว สิกฺขิตพฺพนฺติ.\csromanspacer{}\csromanspacer{}นวมํ.}

\csromanheader{2}{10. ทุติยมรณสฺสติสุตฺต}
\proseitem{20}{เอกํ สมยํ ภควา นาติเก วิหรติ คิญฺชกาวสเถ.\csromanspacer{} ตตฺร โข ภควา ภิกฺขู อามนฺเตสิ\csromandash{}มรณสฺสติ ภิกฺขเว ภาวิตา พหุลีกตา มหปฺผลา โหติ มหานิสํสา อมโตคธา อมตปริโยสานา.\csromanspacer{} กถํ ภาวิตา จ ภิกฺขเว มรณสฺสติ กถํ พหุลีกตา มหปฺผลา โหติ มหานิสํสา อมโตคธา อมตปริโยสานา.}

\prose{อิธ ภิกฺขเว ภิกฺขุ ทิวเส นิกฺขนฺเต รตฺติยา ปติหิตาย\footnote{ปฏิคตาย (ก) อํ 3. 138 ปิฏฺเฐปิ.} อิติ ปฏิสญฺจิกฺขติ “พหุกา โข เม ปจฺจยา มรณสฺส, อหิ วา มํ ฑํเสยฺย, วิจฺฉิโก วา มํ ฑํเสยฺย, สตปที วา มํ ฑํเสยฺย.\csromanspacer{} เตน เม อสฺส}

\vspace{\tipitakaparskip}
\csromancenter{สารณียวคฺค กาลกิริยา, โส มมสฺส อนฺตราโย.\csromanspacer{} อุปกฺขลิตฺวา วา ปปเตยฺยํ, ภตฺตํ วา เม ภุตฺตํ พฺยาปชฺเชยฺย, ปิตฺตํ วา เม กุปฺเปยฺย, เสมฺหํ วา เม กุปฺเปยฺย, สตฺถกา วา เม วาตา กุปฺเปยฺยุํ, เตน เม อสฺส กาลกิริยา, โส มมสฺส อนฺตราโย”ติ.\csromanspacer{} เตน ภิกฺขเว ภิกฺขุนา อิติ ปฏิสญฺจิกฺขิตพฺพํ “อตฺถิ นุ โข เม ปาปกา อกุสลา ธมฺมา อปฺปหีนา, เย เม อสฺสุ รตฺติํ กาลํ กโรนฺตสฺส อนฺตรายายา”ติ.}
\prose{\csromanfolio{271}สเจ ภิกฺขเว ภิกฺขุ ปจฺจเวกฺขมาโน เอวํ ชานาติ “อตฺถิ เม ปาปกา อกุสลา ธมฺมา อปฺปหีนา, เย เม อสฺสุ รตฺติํ กาลํ กโรนฺตสฺส อนฺตรายายา”ติ.\csromanspacer{} เตน ภิกฺขเว ภิกฺขุนา เตสํเยว ปาปกานํ อกุสลานํ ธมฺมานํ ปหานาย อธิมตฺโต ฉนฺโท จ วายาโม จ อุสฺสาโห จ อุสฺโสฬฺหี จ อปฺปฏิวานี จ สติ จ สมฺปชญฺญญฺจ กรณียํ.\csromanspacer{} เสยฺยถาปิ ภิกฺขเว อาทิตฺตเจโล วา อาทิตฺตสีโส วา ตสฺเสว เจลสฺส วา สีสสฺส วา นิพฺพาปนาย อธิมตฺตํ ฉนฺทญฺจ วายามญฺจ อุสฺสาหญฺจ อุสฺโสฬฺหิญฺจ อปฺปฏิวานิญฺจ สติญฺจ สมฺปชญฺญญฺจ กเรยฺย.\csromanspacer{} เอวเมวํ โข ภิกฺขเว เตน ภิกฺขุนา เตสํเยว ปาปกานํ อกุสลานํ ธมฺมานํ ปหานาย อธิมตฺโต ฉนฺโท จ วายาโม จ อุสฺสาโห จ อุสฺโสฬฺหี จ อปฺปฏิวานี จ สติ จ สมฺปชญฺญญฺจ กรณียํ.}

\prose{สเจ ปน ภิกฺขเว ภิกฺขุ ปจฺจเวกฺขมาโน เอวํ ชานาติ “นตฺถิ เม ปาปกา อกุสลา ธมฺมา อปฺปหีนา, เย เม อสฺสุ รตฺติํ กาลํ กโรนฺตสฺส อนฺตรายายา”ติ.\csromanspacer{} เตน ภิกฺขเว ภิกฺขุนา เตเนว ปีติปาโมชฺเชน วิหาตพฺพํ อโหรตฺตานุสิกฺขินา กุสเลสุ ธมฺเมสุ.}

\prose{อิธ ปน ภิกฺขเว ภิกฺขุ รตฺติยา นิกฺขนฺตาย ทิวเส ปติหิเต อิติ ปฏิสญฺจิกฺขติ “พหุกา โข เม ปจฺจยา มรณสฺส อหิ วา มํ ฑฺฑํเสยฺย, วิจฺฉิโก วา มํ ฑํเสยฺย, สตปที วา มํ ฑํเสยฺย, เตน เม อสฺส กาลกิริยา, โส มมสฺส อนฺตราโย.\csromanspacer{} อุปกฺขลิตฺวา วา ปปเตยฺยํ, ภตฺตํ วา เม ภุตฺตํ พฺยาปชฺเชยฺย, ปิตฺตํ วา เม กุปฺเปยฺย, เสมฺหํ วา เม กุปฺเปยฺย, สตฺถกา วา เม วาตา กุปฺเปยฺยุํ, เตน เม อสฺส กาลกิริยา, โส มมสฺส อนฺตราโย”ติ.\csromanspacer{} เตน ภิกฺขเว ภิกฺขุนา อิติ ปฏิสญฺจิกฺขิตพฺพํ “อตฺถิ นุ โข เม ปาปกา อกุสลา ธมฺมา อปฺปหีนา, เย เม อสฺสุ ทิวา กาลํ กโรนฺตสฺส อนฺตรายายา”ติ.}

\gambhira{ฉกฺกนิปาตปาฬิ}
\prose{\csromanfolio{272}สเจ ภิกฺขเว ภิกฺขุ ปจฺจเวกฺขมาโน เอวํ ชานาติ “อตฺถิ เม ปาปกา อกุสลา ธมฺมา อปฺปหีนา, เย เม อสฺสุ ทิวา กาลํ กโรนฺตสฺส อนฺตรายายา”ติ.\csromanspacer{} เตน ภิกฺขเว ภิกฺขุนา เตสํเยว ปาปกานํ อกุสลานํ ธมฺมานํ ปหานาย อธิมตฺโต ฉนฺโท จ วายาโม จ อุสฺสาโห จ อุสฺโสฬฺหี จ อปฺปฏิวานี จ สติ จ สมฺปชญฺญญฺจ กรณียํ.\csromanspacer{} เสยฺยถาปิ ภิกฺขเว อาทิตฺตเจโล วา อาทิตฺตสีโส วา ตสฺเสว เจลสฺส วา สีสสฺส วา นิพฺพาปนาย อธิมตฺตํ ฉนฺทญฺจ วายามญฺจ อุสฺสาหญฺจ อุสฺโสฬฺหิญฺจ อปฺปฏิวานิญฺจ สติญฺจ สมฺปชญฺญญฺจ กเรยฺย.\csromanspacer{} เอวเมวํ โข ภิกฺขเว เตน ภิกฺขุนา เตสํเยว ปาปกานํ อกุสลานํ ธมฺมานํ ปหานาย อธิมตฺโต ฉนฺโท จ วายาโม จ อุสฺสาโห จ อุสฺโสฬฺหี จ อปฺปฏิวานี จ สติ จ สมฺปชญฺญญฺจ กรณียํ.}

\prose{สเจ ปน ภิกฺขเว ภิกฺขุ ปจฺจเวกฺขมาโน เอวํ ชานาติ “นตฺถิ เม ปาปกา อกุสลา ธมฺมา อปฺปหีนา, เย เม อสฺสุ ทิวา กาลํ กโรนฺตสฺส อนฺตรายายา”ติ.\csromanspacer{} เตน ภิกฺขเว ภิกฺขุนา เตเนว ปีติปาโมชฺเชน วิหาตพฺพํ อโหรตฺตานุสิกฺขินา กุสเลสุ ธมฺเมสุ.\csromanspacer{} เอวํ ภาวิตา โข ภิกฺขเว มรณสฺสติ เอวํ พหุลีกตา มหปฺผลา โหติ มหานิสํสา อมโตคธา อมตปริโยสานาติ.\csromanspacer{}\csromanspacer{}ทสฺสมํ.\csromanspacer{} สารณียวคฺโค ทุติโย.}

\titlehead{\textbf{ตสฺสุทฺทานํ}}
\csromangathameasure{เทฺว สารณี นิสฺสารณียํ, ภทฺทกํ อนุตปฺปิยํ. \\ นกุลํ โสปฺปมจฺฉา จ, เทฺว โหนฺติ มรณสฺสตีติ.}
\csromangathagroup{\csromangathastanza{เทฺว สารณี นิสฺสารณียํ, ภทฺทกํ อนุตปฺปิยํ. \\ นกุลํ โสปฺปมจฺฉา จ, เทฺว โหนฺติ มรณสฺสตีติ.}}

\csromanmatikaanchor{mtk.143}
\csromantocmarknopage{section}{3. อนุตฺตริยวคฺค}{mtk.143}
\bookmark[dest={mtk.143},level=1]{3. อนุตฺตริยวคฺค}
\csromanheader{1}{3. อนุตฺตริยวคฺค}
\csromanmatikaanchor{mtk.144}
\csromantocmarkref{subsection}{1. สามกสุตฺต}{mtk.144}
\bookmark[dest={mtk.144},level=2]{1. สามกสุตฺต}
\csromanheader{2}{1. สามกสุตฺต}
\proseitem{21}{เอกํ สมยํ ภควา สกฺเกสุ วิหรติ สามคามเก โปกฺขรณิยายํ.\csromanspacer{} อถ โข อญฺญตรา เทวตา อภิกฺกนฺตาย รตฺติยา อภิกฺกนฺตวณฺณา เกวลกปฺปํ โปกฺขรณิยํ โอภาเสตฺวา เยน ภควา}

\vspace{\tipitakaparskip}
\csromancenter{อนุตฺตริยวคฺค เตนุปสงฺกมิ, อุปสงฺกมิตฺวา ภควนฺตํ อภิวาเทตฺวา เอกมนฺตํ อฏฺฐาสิ, เอกมนฺตํ ฐิตา โข สา เทวตา ภควนฺตํ เอตทโวจ.}
\prose{\csromanfolio{273}ตโยเม ภนฺเต ธมฺมา ภิกฺขุโน ปริหานาย สํวตฺตนฺติ.\csromanspacer{} กตเม ตโย? กมฺมารามตา ภสฺสารามตา นิทฺทารามตา.\csromanspacer{} อิเม โข ภนฺเต ตโย ธมฺมา ภิกฺขุโน ปริหานาย สํวตฺตนฺตีติ.\csromanspacer{} อิทมโวจ สา เทวตา.\csromanspacer{} สมนุญฺโญ สตฺถา อโหสิ.\csromanspacer{} อถ โข สา เทวตา “สมนุญฺโญ เม สตฺถา”ติ ภควนฺตํ อภิวาเทตฺวา ปทกฺขิณํ กตฺวา ตตฺเถวนฺตรธายิ.}

\prose{อถ โข ภควา ตสฺสา รตฺติยา อจฺจเยน ภิกฺขู อามนฺเตสิ\csromandash{}อิมํ ภิกฺขเว รตฺติํ อญฺญตรา เทวตา อภิกฺกนฺตาย รตฺติยา อภิกฺกนฺตวณฺณา เกวลกปฺปํ โปกฺขรณิยํ โอภาเสตฺวา เยนาหํ เตนุปสงฺกมิ, อุปสงฺกมิตฺวา มํ อภิวาเทตฺวา เอกมนฺตํ อฏฺฐาสิ, เอกมนฺตํ ฐิตา โข ภิกฺขเว สา เทวตา มํ เอตทโวจ “ตโยเม ภนฺเต ธมฺมา ภิกฺขุโน ปริหานาย สํวตฺตนฺติ.\csromanspacer{} กตเม ตโย? กมฺมารามตา ภสฺสารามตา นิทฺทารามตา.\csromanspacer{} อิเม โข ภนฺเต ตโย ธมฺมา ภิกฺขุโน ปริหานาย สํวตฺตนฺตี”ติ.\csromanspacer{} อิทมโวจ ภิกฺขเว สา เทวตา, อิทํ วตฺวา มํ อภิวาเทตฺวา ปทกฺขิณํ กตฺวา ตตฺเถวนฺตรธายิ.\csromanspacer{} เตสํ โว\footnote{โข (ก)} ภิกฺขเว อลาภา, เตสํ ทุลฺลทฺธํ, เย โว เทวตาปิ ชานนฺติ กุสเลหิ ธมฺเมหิ ปริหายมาเน.}

\prose{อปเรปิ ภิกฺขเว ตโย ปริหานิเย ธมฺเม เทเสสฺสามิ, ตํ สุณาถ สาธุกํ มนสิ กโรถ, ภาสิสฺสามีติ. “เอวํ ภนฺเต”ติ โข เต ภิกฺขู ภควโต ปจฺจสฺโสสุํ.\csromanspacer{} ภควา เอตทโวจ\csromandash{}กตเม จ ภิกฺขเว ตโย ปริหานิยา ธมฺมา? สงฺคณิการามตา โทวจสฺสตา ปาปมิตฺตตา.\csromanspacer{} อิเม โข ภิกฺขเว ตโย ปริหานิยา ธมฺมา.}

\prose{เย หิ เกจิ ภิกฺขเว อตีตมทฺธานํ ปริหายิํสุ กุสเลหิ ธมฺเมหิ, สพฺเพเต อิเมเหว ฉหิ ธมฺเมหิ ปริหายิํสุ กุสเลหิ ธมฺเมหิ.\csromanspacer{} เยปิ หิ เกจิ ภิกฺขเว อนาคตมทฺธานํ ปริหายิสฺสนฺติ กุสเลหิ ธมฺเมหิ, สพฺเพเต อิเมเหว ฉหิ ธมฺเมหิ ปริหายิสฺสนฺติ}

\vspace{\tipitakaparskip}
\csromancenter{ฉกฺกนิปาตปาฬิ กุสเลหิ ธมฺเมหิ.\csromanspacer{} เยปิ หิ เกจิ ภิกฺขเว เอตรหิ ปริหายนฺติ กุสเลหิ ธมฺเมหิ, สพฺเพเต อิเมเหว ฉหิ ธมฺเมหิ ปริหายนฺติ กุสเลหิ ธมฺเมหีติ.\csromanspacer{}\csromanspacer{}ปฐมํ.}
\csromanheader{2}{2. อปริหานิยสุตฺต}
\csromanmatikaanchor{mtk.145}
\csromantocmarkref{subsection}{2-4. อปริหานิยสุตฺตาทิ}{mtk.145}
\bookmark[dest={mtk.145},level=2]{2-4. อปริหานิยสุตฺตาทิ}
\proseitem{22}{\csromanfolio{274}ฉยิเม ภิกฺขเว อปริหานิเย ธมฺเม เทเสสฺสามิ, ตํ สุณาถ ฯเปฯ.\csromanspacer{} กตเม จ ภิกฺขเว ฉ อปริหานิยา ธมฺมา, น กมฺมารามตา น ภสฺสารามตา น นิทฺทารามตา น สงฺคณิการามตา โสวจสฺสตา กลฺยาณมิตฺตตา.\csromanspacer{} อิเม โข ภิกฺขเว ฉ อปริหานิยา ธมฺมา.}

\prose{เย หิ เกจิ ภิกฺขเว อตีตมทฺธานํ น ปริหายิํสุ กุสเลหิ ธมฺเมหิ, สพฺเพเต อิเมเหว ฉหิ ธมฺเมหิ น ปริหายิํสุ กุสเลหิ ธมฺเมหิ.\csromanspacer{} เยปิ หิ เกจิ ภิกฺขเว อนาคตมทฺธานํ น ปริหายิสฺสนฺติ กุสเลหิ ธมฺเมหิ, สพฺเพเต อิเมเหว ฉหิ ธมฺเมหิ น ปริหายิสฺสนฺติ กุสเลหิ ธมฺเมหิ.\csromanspacer{} เยปิ หิ เกจิ ภิกฺขเว เอตรหิ น ปริหายนฺติ กุสเลหิ ธมฺเมหิ, สพฺเพเต อิเมเหว ฉหิ ธมฺเมหิ น ปริหายนฺติ กุสเลหิ ธมฺเมหีติ.\csromanspacer{}\csromanspacer{}ทุติยํ.}

\csromanheader{2}{3. ภยสุตฺต}
\proseitem{23}{‘ภยนฺ’ติ ภิกฺขเว กามานเมตํ อธิวจนํ. ‘ทุกฺขนฺ’ติ ภิกฺขเว กามานเมตํ อธิวจนํ. ‘โรโค’ติ ภิกฺขเว กามานเมตํ อธิวจนํ. ‘คณฺโฑ’ติ ภิกฺขเว กามานเมตํ อธิวจนํ. ‘สงฺโค’ติ ภิกฺขเว กามานเมตํ อธิวจนํ. ‘ปงฺโก’ติ ภิกฺขเว กามานเมตํ อธิวจนํ.}

\prose{กสฺมา จ ภิกฺขเว ‘ภยนฺ’ติ กามานเมตํ อธิวจนํ, กามราครตฺตายํ ภิกฺขเว ฉนฺทราควินิพทฺโธ ทิฏฺฐธมฺมิกาปิ ภยา น ปริมุจฺจติ, สมฺปรายิกาปิ ภยา น ปริมุจฺจติ.\csromanspacer{} ตสฺมา ‘ภยนฺ’ติ กามานเมตํ อธิวจนํ.\csromanspacer{} กสฺมา จ ภิกฺขเว ‘ทุกฺขนฺ’ติ ฯเปฯ. ‘โรโค’ติ. ‘คณฺโฑ’ติ. ‘สงฺโค’ติ. ‘ปงฺโก’ติ กามานเมตํ อธิวจนํ, กามราครตฺตายํ ภิกฺขเว ฉนฺทราควินิพทฺโธ ทิฏฺฐธมฺมิกาปิ ปงฺกา น ปริมุจฺจติ, สมฺปรายิกาปิ ปงฺกา น ปริมุจฺจติ.\csromanspacer{} ตสฺมา ‘ปงฺโก’ติ กามานเมตํ อธิวจนนฺติ.}

\csromangathameasure{ภยํ ทุกฺขํ โรโค คณฺโฑ, สงฺโค ปงฺโก จ อุภยํ. \\ เอเต กามา ปวุจฺจนฺติ, ยตฺถ สตฺโต ปุถุชฺชโน.}
\csromangathagroup{\csromangathastanza{ภยํ ทุกฺขํ โรโค คณฺโฑ, สงฺโค ปงฺโก จ อุภยํ. \\ เอเต กามา ปวุจฺจนฺติ, ยตฺถ สตฺโต ปุถุชฺชโน.}}

\chapterheadpageread{3. อนุตฺตริยวคฺค}
\csromangathameasure{อุปาทาเน ภยํ ทิสฺวา, ชาติมรณสมฺภเว. \\ อนุปาทา วิมุจฺจนฺติ, ชาติมรณสงฺขเย. \\ เต เขมปฺปตฺตา สุขิโน, ทิฏฺฐธมฺมาภินิพฺพุตา. \\ สพฺพเวรภยาตีตา, สพฺพทุกฺขํ อุปจฺจคุนฺติ.ตติยํ.}
\csromangathagroup{\csromangathastanza{\csromanfolio{275}อุปาทาเน ภยํ ทิสฺวา, ชาติมรณสมฺภเว. \\ อนุปาทา วิมุจฺจนฺติ, ชาติมรณสงฺขเย.}\csromangathastanza{เต เขมปฺปตฺตา สุขิโน, ทิฏฺฐธมฺมาภินิพฺพุตา. \\ สพฺพเวรภยาตีตา\footnote{สพฺเพ เวรภยาตีตา (สฺยา)}, สพฺพทุกฺขํ อุปจฺจคุนฺติ.\csromanspacer{}\csromanspacer{}ตติยํ.}}

\csromanheader{2}{4. หิมวนฺตสุตฺต}
\proseitem{24}{ฉหิ ภิกฺขเว ธมฺเมหิ สมนฺนาคโต ภิกฺขุ หิมวนฺตํ ปพฺพตราชํ ปทาเลยฺย, โก ปน วาโท ฉวาย อวิชฺชาย.\csromanspacer{} กตเมหิ ฉหิ? อิธ ภิกฺขเว ภิกฺขุ สมาธิสฺส สมาปตฺติกุสโล โหติ, สมาธิสฺส ฐิติกุสโล โหติ, สมาธิสฺส วุฏฺฐานกุสโล โหติ, สมาธิสฺส กลฺลิตกุสโล\footnote{กลฺลตากุสโล (สฺยา, กํ, ก) สํ 2. 225 ปิฏฺเฐปิ ปสฺสิตพฺพํ.} โหติ, สมาธิสฺส โคจรกุสโล โหติ, สมาธิสฺส อภินีหารกุสโล โหติ.\csromanspacer{} อิเมหิ โข ภิกฺขเว ฉหิ ธมฺเมหิ สมนฺนาคโต ภิกฺขุ หิมวนฺตํ ปพฺพตราชํ ปทาเลยฺย, โก ปน วาโท ฉวาย อวิชฺชายาติ.\csromanspacer{}\csromanspacer{}จตุตฺถํ.}

\csromanmatikaanchor{mtk.146}
\csromantocmarkref{subsection}{5. อนุสฺสติฏฺฐานสุตฺต}{mtk.146}
\bookmark[dest={mtk.146},level=2]{5. อนุสฺสติฏฺฐานสุตฺต}
\csromanheader{2}{5. อนุสฺสติฏฺฐานสุตฺต}
\proseitem{25}{ฉยิมานิ ภิกฺขเว อนุสฺสติฏฺฐานานิ.\csromanspacer{} กตมานิ ฉ? อิธ ภิกฺขเว อริยสาวโก ตถาคตํ อนุสฺสรติ “อิติปิ โส ภควา ฯเปฯ สตฺถา เทวมนุสฺสานํ พุทฺโธ ภควา”ติ.\csromanspacer{} ยสฺมิํ ภิกฺขเว สมเย อริยสาวโก ตถาคตํ อนุสฺสรติ, เนวสฺส ตสฺมิํ สมเย ราคปริยุฏฺฐิตํ จิตฺตํ โหติ, น โทสปริยุฏฺฐิตํ จิตฺตํ โหติ, น โมหปริยุฏฺฐิตํ จิตฺตํ โหติ.\csromanspacer{} อุชุคตเมวสฺส ตสฺมิํ สมเย จิตฺตํ โหติ นิกฺขนฺตํ มุตฺตํ วุฏฺฐิตํ เคธมฺหา. ‘เคโธ’ติ โข ภิกฺขเว ปญฺจนฺเนตํ กามคุณานํ อธิวจนํ.\csromanspacer{} อิทมฺปิ โข ภิกฺขเว อารมฺมณํ กริตฺวา เอวมิเธกจฺเจ สตฺตา วิสุชฺฌนฺติ.}

\prose{ปุน จปรํ ภิกฺขเว อริยสาวโก ธมฺมํ อนุสฺสรติ “สฺวากฺขาโต ภควตา ธมฺโม ฯเปฯ ปจฺจตฺตํ เวทิตพฺโพ วิญฺญูหี”ติ.\csromanspacer{} ยสฺมิํ ภิกฺขเว สมเย อริยสาวโก ธมฺมํ อนุสฺสรติ, เนวสฺส ตสฺมิํ สมเย ราคปริยุฏฺฐิตํ จิตฺตํ โหติ, น โทสปริยุฏฺฐิตํ จิตฺตํ โหติ, น โมหปริยุฏฺฐิตํ จิตฺตํ}

\vspace{\tipitakaparskip}
\csromancenter{ฉกฺกนิปาตปาฬิ โหติ.\csromanspacer{} อุชุคตเมวสฺส ตสฺมิํ สมเย จิตฺตํ โหติ นิกฺขนฺตํ มุตฺตํ วุฏฺฐิตํ เคธมฺหา. ‘เคโธ’ติ โข ภิกฺขเว ปญฺจนฺเนตํ กามคุณานํ อธิวจนํ.\csromanspacer{} อิทมฺปิ โข ภิกฺขเว อารมฺมณํ กริตฺวา เอวมิเธกจฺเจ สตฺตา วิสุชฺฌนฺติ.}
\prose{\csromanfolio{276}ปุน จปรํ ภิกฺขเว อริยสาวโก สํฆํ อนุสฺสรติ “สุปฺปฏิปนฺโน ภควโต สาวกสํโฆ ฯเปฯ อนุตฺตรํ ปุญฺญกฺเขตฺตํ โลกสฺสา”ติ.\csromanspacer{} ยสฺมิํ ภิกฺขเว สมเย อริยสาวโก สํฆํ อนุสฺสรติ, เนวสฺส ตสฺมิํ สมเย ราคปริยุฏฺฐิตํ จิตฺตํ โหติ, น โทสปริยุฏฺฐิตํ จิตฺตํ โหติ, น โมหปริยุฏฺฐิตํ จิตฺตํ โหติ.\csromanspacer{} อุชุคตเมวสฺส ตสฺมิํ สมเย จิตฺตํ โหติ นิกฺขนฺตํ มุตฺตํ วุฏฺฐิตํ เคธมฺหา. ‘เคโธ’ติ โข ภิกฺขเว ปญฺจนฺเนตํ กามคุณานํ อธิวจนํ.\csromanspacer{} อิทมฺปิ โข ภิกฺขเว อารมฺมณํ กริตฺวา เอวมิเธกจฺเจ สตฺตา วิสุชฺฌนฺติ.}

\prose{ปุน จปรํ ภิกฺขเว อริยสาวโก อตฺตโน สีลานิ อนุสฺสรติ อขณฺฑานิ ฯเปฯ สมาธิสํวตฺตนิกานิ.\csromanspacer{} ยสฺมิํ ภิกฺขเว สมเย อริยสาวโก สีลํ อนุสฺสรติ, เนวสฺส ตสฺมิํ สมเย ราคปริยุฏฺฐิตํ จิตฺตํ โหติ, น โทสปริยุฏฺฐิตํ จิตฺตํ โหติ, น โมหปริยุฏฺฐิตํ จิตฺตํ โหติ.\csromanspacer{} อุชุคตเมวสฺส ตสฺมิํ สมเย จิตฺตํ โหติ นิกฺขนฺตํ มุตฺตํ วุฏฺฐิตํ เคธมฺหา. ‘เคโธ’ติ โข ภิกฺขเว ปญฺจนฺเนตํ กามคุณานํ อธิวจนํ.\csromanspacer{} อิทมฺปิ โข ภิกฺขเว อารมฺมณํ กริตฺวา เอวมิเธกจฺเจ สตฺตา วิสุชฺฌนฺติ.}

\prose{ปุน จปรํ ภิกฺขเว อริยสาวโก อตฺตโน จาคํ อนุสฺสรติ “ลาภา วต เม, สุลทฺธํ วต เม ฯเปฯ ยาจโยโค ทานสํวิภาครโต”ติ ฯเปฯ เอวมิเธกจฺเจ สตฺตา วิสุชฺฌนฺติ.}

\prose{ปุน จปรํ ภิกฺขเว อริยสาวโก เทวตา อนุสฺสรติ “สนฺติ เทวา จาตุมหาราชิกา, สนฺติ เทวา ตาวติํสา, สนฺติ เทวา ยามา, สนฺติ เทวา ตุสิตา, สนฺติ เทวา นิมฺมานรติโน, สนฺติ เทวา ปรนิมฺมิตวสวตฺติโน, สนฺติ เทวา พฺรหฺมกายิกา, สนฺติ เทวา ตตุตฺตริ, ยถารูปาย สทฺธาย สมนฺนาคตา ตา เทวตา อิโต จุตา ตตฺถ อุปปนฺนา, มยฺหมฺปิ ตถารูปา สทฺธา สํวิชฺชติ.\csromanspacer{} ยถารูเปน สีเลน.\csromanspacer{} สุเตน.\csromanspacer{} จาเคน.\csromanspacer{} ปญฺญาย สมนฺนาคตา ตา เทวตา อิโต จุตา ตตฺถ อุปปนฺนา, มยฺหมฺปิ ตถารูปา ปญฺญา สํวิชฺชตี”ติ.}

\chapterheadpageread{3. อนุตฺตริยวคฺค}
\prose{\csromanfolio{277}ยสฺมิํ ภิกฺขเว สมเย อริยสาวโก อตฺตโน จ ตาสญฺจ เทวตานํ สทฺธญฺจ สีลญฺจ สุตญฺจ จาคญฺจ ปญฺญญฺจ อนุสฺสรติ, เนวสฺส ตสฺมิํ สมเย ราคปริยุฏฺฐิตํ จิตฺตํ โหติ, น โทสปริยุฏฺฐิตํ จิตฺตํ โหติ, น โมหปริยุฏฺฐิตํ จิตฺตํ โหติ.\csromanspacer{} อุชุคตเมวสฺส ตสฺมิํ สมเย จิตฺตํ โหติ นิกฺขนฺตํ มุตฺตํ วุฏฺฐิตํ เคธมฺหา. ‘เคโธ’ติ โข ภิกฺขเว ปญฺจนฺเนตํ กามคุณานํ อธิวจนํ.\csromanspacer{} อิทมฺปิ โข ภิกฺขเว อารมฺมณํ กริตฺวา เอวมิเธกจฺเจ สตฺตา วิสุชฺฌนฺติ.\csromanspacer{} อิมานิ โข ภิกฺขเว ฉ อนุสฺสติฏฺฐานานีติ.\csromanspacer{}\csromanspacer{}ปญฺจมํ.}

\csromanmatikaanchor{mtk.147}
\csromantocmarkref{subsection}{6. มหากจฺจานสุตฺต}{mtk.147}
\bookmark[dest={mtk.147},level=2]{6. มหากจฺจานสุตฺต}
\csromanheader{2}{6. มหากจฺจานสุตฺต}
\proseitem{26}{ตตฺร โข อายสฺมา มหากจฺจาโน ภิกฺขู อามนฺเตสิ “อาวุโส ภิกฺขเว”ติ. “อาวุโส”ติ โข เต ภิกฺขู อายสฺมโต มหากจฺจานสฺส ปจฺจสฺโสสุํ.\csromanspacer{} อายสฺมา มหากจฺจาโน เอตทโวจ\csromandash{}อจฺฉริยํ อาวุโส, อพฺภุตํ อาวุโส, ยาวญฺจิทํ เตน ภควตา ชานตา ปสฺสตา อรหตา สมฺมาสมฺพุทฺเธน สมฺพาเธ โอกาสาธิคโม อนุพุทฺโธ สตฺตานํ วิสุทฺธิยา โสกปริเทวานํ สมติกฺกมาย ทุกฺขโทมนสฺสานํ อตฺถงฺคมาย ญายสฺส อธิคมาย นิพฺพานสฺส สจฺฉิกิริยาย ยทิทํ ฉ อนุสฺสติฏฺฐานานิ.}

\prose{กตมานิ ฉ? อิธาวุโส อริยสาวโก ตถาคตํ อนุสฺสรติ “อิติปิ โส ภควา ฯเปฯ สตฺถา เทวมนุสฺสานํ พุทฺโธ ภควา”ติ.\csromanspacer{} ยสฺมิํ อาวุโส สมเย อริยสาวโก ตถาคตํ อนุสฺสรติ, เนวสฺส ตสฺมิํ สมเย ราคปริยุฏฺฐิตํ จิตฺตํ โหติ, น โทสปริยุฏฺฐิตํ จิตฺตํ โหติ, น โมหปริยุฏฺฐิตํ จิตฺตํ โหติ.\csromanspacer{} อุชุคตเมวสฺส ตสฺมิํ สมเย จิตฺตํ โหติ นิกฺขนฺตํ มุตฺตํ วุฏฺฐิตํ เคธมฺหา. ‘เคโธ’ติ โข อาวุโส ปญฺจนฺเนตํ กามคุณานํ อธิวจนํ.\csromanspacer{} ส โข โส อาวุโส อริยสาวโก สพฺพโส อากาสสเมน เจตสา วิหรติ วิปุเลน มหคฺคเตน อปฺปมาเณน อเวเรน อพฺยาปชฺเชน.\csromanspacer{} อิทมฺปิ โข อาวุโส อารมฺมณํ กริตฺวา เอวมิเธกจฺเจ สตฺตา วิสุทฺธิธมฺมา ภวนฺติ.}

\prose{ปุน จปรํ อาวุโส อริยสาวโก ธมฺมํ อนุสฺสรติ “สฺวากฺขาโต ภควตา ธมฺโม ฯเปฯ ปจฺจตฺตํ เวทิตพฺโพ วิญฺญูหี”ติ.\csromanspacer{} ยสฺมิํ อาวุโส สมเย อริยสาวโก ธมฺมํ อนุสฺสรติ, เนวสฺส ตสฺมิํ สมเย ราคปริยุฏฺฐิตํ จิตฺตํ โหติ, น โทสปริยุฏฺฐิตํ จิตฺตํ โหติ, น}

\vspace{\tipitakaparskip}
\csromancenter{ฉกฺกนิปาตปาฬิ โมหปริยุฏฺฐิตํ จิตฺตํ โหติ.\csromanspacer{} อุชุคตเมวสฺส ตสฺมิํ สมเย จิตฺตํ โหติ นิกฺขนฺตํ มุตฺตํ วุฏฺฐิตํ เคธมฺหา. ‘เคโธ’ติ โข อาวุโส ปญฺจนฺเนตํ กามคุณานํ อธิวจนํ.\csromanspacer{} ส โข โส อาวุโส อริยสาวโก สพฺพโส อากาสสเมน เจตสา วิหรติ วิปุเลน มหคฺคเตน อปฺปมาเณน อเวเรน อพฺยาปชฺเชน.\csromanspacer{} อิทมฺปิ โข อาวุโส อารมฺมณํ กริตฺวา เอวมิเธกจฺเจ สตฺตา วิสุทฺธิธมฺมา ภวนฺติ.}
\prose{\csromanfolio{278}ปุน จปรํ อาวุโส อริยสาวโก สํฆํ อนุสฺสรติ “สุปฺปฏิปนฺโน ภควโต สาวกสํโฆ ฯเปฯ อนุตฺตรํ ปุญฺญกฺเขตฺตํ โลกสฺสา”ติ.\csromanspacer{} ยสฺมิํ อาวุโส สมเย อริยสาวโก สํฆํ อนุสฺสรติ, เนวสฺส ตสฺมิํ สมเย ราคปริยุฏฺฐิตํ จิตฺตํ โหติ, น โทสปริยุฏฺฐิตํ จิตฺตํ โหติ, น โมหปริยุฏฺฐิตํ จิตฺตํ โหติ.\csromanspacer{} อุชุคตเมวสฺส ตสฺมิํ สมเย จิตฺตํ โหติ นิกฺขนฺตํ มุตฺตํ วุฏฺฐิตํ เคธมฺหา. ‘เคโธ’ติ โข อาวุโส ปญฺจนฺเนตํ กามคุณานํ อธิวจนํ.\csromanspacer{} ส โข โส อาวุโส อริยสาวโก สพฺพโส อากาสสเมน เจตสา วิหรติ วิปุเลน มหคฺคเตน อปฺปมาเณน อเวเรน อพฺยาปชฺเชน.\csromanspacer{} อิทมฺปิ โข อาวุโส อารมฺมณํ กริตฺวา เอวมิเธกจฺเจ สตฺตา วิสุทฺธิธมฺมา ภวนฺติ.}

\prose{ปุน จปรํ อาวุโส อริยสาวโก อตฺตโน สีลานิ อนุสฺสรติ อขณฺฑานิ ฯเปฯ สมาธิสํวตฺตนิกานิ.\csromanspacer{} ยสฺมิํ อาวุโส สมเย อริยสาวโก อตฺตโน สีลํ อนุสฺสรติ, เนวสฺส ตสฺมิํ สมเย ราคปริยุฏฺฐิตํ จิตฺตํ โหติ, น โทสปริยุฏฺฐิตํ จิตฺตํ โหติ, น โมหปริยุฏฺฐิตํ จิตฺตํ โหติ.\csromanspacer{} อุชุคตเมวสฺส ตสฺมิํ สมเย จิตฺตํ โหติ นิกฺขนฺตํ มุตฺตํ วุฏฺฐิตํ เคธมฺหา. ‘เคโธ’ติ โข อาวุโส ปญฺจนฺเนตํ กามคุณานํ อธิวจนํ.\csromanspacer{} ส โข โส อาวุโส อริยสาวโก สพฺพโส อากาสสเมน เจตสา วิหรติ วิปุเลน มหคฺคเตน อปฺปมาเณน อเวเรน อพฺยาปชฺเชน.\csromanspacer{} อิทมฺปิ โข อาวุโส อารมฺมณํ กริตฺวา เอวมิเธกจฺเจ สตฺตา วิสุทฺธิธมฺมา ภวนฺติ.}

\prose{ปุน จปรํ อาวุโส อริยสาวโก อตฺตโน จาคํ อนุสฺสรติ “ลาภา วต เม, สุลทฺธํ วต เม ฯเปฯ ยาจโยโค ทานสํวิภาครโต”ติ.\csromanspacer{} ยสฺมิํ อาวุโส สมเย อริยสาวโก อตฺตโน จาคํ อนุสฺสรติ, เนวสฺส ตสฺมิํ สมเย ราคปริยุฏฺฐิตํ จิตฺตํ โหติ,}

\vspace{\tipitakaparskip}
\csromancenter{อนุตฺตริยวคฺค น โทสปริยุฏฺฐิตํ จิตฺตํ โหติ, น โมหปริยุฏฺฐิตํ จิตฺตํ โหติ.\csromanspacer{} อุชุคตเมวสฺส ตสฺมิํ สมเย จิตฺตํ โหติ นิกฺขนฺตํ มุตฺตํ วุฏฺฐิตํ เคธมฺหา. ‘เคโธ’ติ โข อาวุโส ปญฺจนฺเนตํ กามคุณานํ อธิวจนํ.\csromanspacer{} ส โข โส อาวุโส อริยสาวโก สพฺพโส อากาสสเมน เจตสา วิหรติ วิปุเลน มหคฺคเตน อปฺปมาเณน อเวเรน อพฺยาปชฺเชน.\csromanspacer{} อิทมฺปิ โข อาวุโส อารมฺมณํ กริตฺวา เอวมิเธกจฺเจ สตฺตา วิสุทฺธิธมฺมา ภวนฺติ.}
\csromanmatikaanchor{mtk.148}
\csromantocmarkref{subsection}{7-8. สมยสุตฺต}{mtk.148}
\bookmark[dest={mtk.148},level=2]{7-8. สมยสุตฺต}
\prose{\csromanfolio{279}ปุน จปรํ อาวุโส อริยสาวโก เทวตา อนุสฺสรติ “สนฺติ เทวา จาตุมหาราชิกา, สนฺติ เทวา ฯเปฯ ตตุตฺตริ, ยถารูปาย สทฺธาย สมนฺนาคตา ตา เทวตา อิโต จุตา ตตฺถ อุปปนฺนา, มยฺหมฺปิ ตถารูปา สทฺธา สํวิชฺชติ.\csromanspacer{} ยถารูเปน สีเลน.\csromanspacer{} สุเตน.\csromanspacer{} จาเคน.\csromanspacer{} ปญฺญาย สมนฺนาคตา ตา เทวตา อิโต จุตา ตตฺถ อุปปนฺนา, มยฺหมฺปิ ตถารูปา ปญฺญา สํวิชฺชตี”ติ.\csromanspacer{} ยสฺมิํ อาวุโส สมเย อริยสาวโก อตฺตโน จ ตาสญฺจ เทวตานํ สทฺธญฺจ สีลญฺจ สุตญฺจ จาคญฺจ ปญฺญญฺจ อนุสฺสรติ, เนวสฺส ตสฺมิํ สมเย ราคปริยุฏฺฐิตํ จิตฺตํ โหติ, น โทสปริยุฏฺฐิตํ จิตฺตํ โหติ, น โมหปริยุฏฺฐิตํ จิตฺตํ โหติ.\csromanspacer{} อุชุคตเมวสฺส ตสฺมิํ สมเย จิตฺตํ โหติ นิกฺขนฺตํ มุตฺตํ วุฏฺฐิตํ เคธมฺหา. ‘เคโธ’ติ โข อาวุโส ปญฺจนฺเนตํ กามคุณานํ อธิวจนํ.\csromanspacer{} ส โข โส อาวุโส อริยสาวโก สพฺพโส อากาสสเมน เจตสา วิหรติ วิปุเลน มหคฺคเตน อปฺปมาเณน อเวเรน อพฺยาปชฺเชน.\csromanspacer{} อิทมฺปิ โข อาวุโส อารมฺมณํ กริตฺวา เอวมิเธกจฺเจ สตฺตา วิสุทฺธิธมฺมา ภวนฺติ.}

\prose{อจฺฉริยํ อาวุโส, อพฺภุตํ อาวุโส, ยาวญฺจิทํ เตน ภควตา ชานตา ปสฺสตา อรหตา สมฺมาสมฺพุทฺเธน สมฺพาเธ โอกาสาธิคโม อนุพุทฺโธ สตฺตานํ วิสุทฺธิยา โสกปริเทวานํ สมติกฺกมาย ทุกฺขโทมนสฺสานํ อตฺถงฺคมาย ญายสฺส อธิคมาย นิพฺพานสฺส สจฺฉิกิริยาย ยทิทํ ฉ อนุสฺสติฏฺฐานานีติ.\csromanspacer{}\csromanspacer{}ฉฏฺฐํ.}

\csromanheader{2}{7. ปฐมสมยสุตฺต}
\proseitem{27}{อถ โข อญฺญตโร ภิกฺขุ เยน ภควา เตนุปสงฺกมิ, อุปสงฺกมิตฺวา ภควนฺตํ อภิวาเทตฺวา เอกมนฺตํ นิสีทิ, เอกมนฺตํ นิสินฺโน}

\vspace{\tipitakaparskip}
\csromancenter{ฉกฺกนิปาตปาฬิ โข โส ภิกฺขุ ภควนฺตํ เอตทโวจ “กติ นุ โข ภนฺเต สมยา มโนภาวนียสฺส ภิกฺขุโน ทสฺสนาย อุปสงฺกมิตุนฺ”ติ.\csromanspacer{} ฉยิเม ภิกฺขุ สมยา มโนภาวนียสฺส ภิกฺขุโน ทสฺสนาย อุปสงฺกมิตุํ.}
\prose{\csromanfolio{280}กตเม ฉ? อิธ ภิกฺขุ ยสฺมิํ สมเย ภิกฺขุ กามราคปริยุฏฺฐิเตน เจตสา วิหรติ กามราคปเรเตน, อุปฺปนฺนสฺส จ กามราคสฺส นิสฺสรณํ ยถาภูตํ นปฺปชานาติ, ตสฺมิํ สมเย มโนภาวนีโย ภิกฺขุ อุปสงฺกมิตฺวา เอวมสฺส วจนีโย “อหํ โข อาวุโส กามราคปริยุฏฺฐิเตน เจตสา วิหรามิ กามราคปเรเตน, อุปฺปนฺนสฺส จ กามราคสฺส นิสฺสรณํ ยถาภูตํ นปฺปชานามิ, สาธุ วต เม อายสฺมา กามราคสฺส ปหานาย ธมฺมํ เทเสตู”ติ.\csromanspacer{} ตสฺส มโนภาวนีโย ภิกฺขุ กามราคสฺส ปหานาย ธมฺมํ เทเสติ.\csromanspacer{} อยํ ภิกฺขุ ปฐโม สมโย มโนภาวนียสฺส ภิกฺขุโน ทสฺสนาย อุปสงฺกมิตุํ.}

\prose{ปุน จปรํ ภิกฺขุ ยสฺมิํ สมเย ภิกฺขุ พฺยาปาทปริยุฏฺฐิเตน เจตสา วิหรติ พฺยาปาทปเรเตน, อุปฺปนฺนสฺส จ พฺยาปาทสฺส นิสฺสรณํ ยถาภูตํ นปฺปชานาติ, ตสฺมิํ สมเย มโนภาวนีโย ภิกฺขุ อุปสงฺกมิตฺวา เอวมสฺส วจนีโย “อหํ โข อาวุโส พฺยาปาทปริยุฏฺฐิเตน เจตสา วิหรามิ พฺยาปาทปเรเตน, อุปฺปนฺนสฺส จ พฺยาปาทสฺส นิสฺสรณํ ยถาภูตํ นปฺปชานามิ, สาธุ วต เม อายสฺมา พฺยาปาทสฺส ปหานาย ธมฺมํ เทเสตู”ติ.\csromanspacer{} ตสฺส มโนภาวนีโย ภิกฺขุ พฺยาปาทสฺส ปหานาย ธมฺมํ เทเสติ.\csromanspacer{} อยํ ภิกฺขุ ทุติโย สมโย มโนภาวนียสฺส ภิกฺขุโน ทสฺสนาย อุปสงฺกมิตุํ.}

\prose{ปุน จปรํ ภิกฺขุ ยสฺมิํ สมเย ภิกฺขุ ถินมิทฺธปริยุฏฺฐิเตน เจตสา วิหรติ ถินมิทฺธปเรเตน, อุปฺปนฺนสฺส จ ถินมิทฺธสฺส นิสฺสรณํ ยถาภูตํ นปฺปชานาติ, ตสฺมิํ สมเย มโนภาวนีโย ภิกฺขุ อุปสงฺกมิตฺวา เอวมสฺส วจนีโย “อหํ โข อาวุโส ถินมิทฺธปริยุฏฺฐิเตน เจตสา วิหรามิ ถินมิทฺธปเรเตน, อุปฺปนฺนสฺส จ ถินมิทฺธสฺส นิสฺสรณํ ยถาภูตํ นปฺปชานามิ, สาธุ วต เม อายสฺมา ถินมิทฺธสฺส ปหานาย ธมฺมํ เทเสตู”ติ.\csromanspacer{} ตสฺส มโนภาวนีโย ภิกฺขุ ถินมิทฺธสฺส ปหานาย}

\vspace{\tipitakaparskip}
\csromancenter{อนุตฺตริยวคฺค ธมฺมํ เทเสติ.\csromanspacer{} อยํ ภิกฺขุ ตติโย สมโย มโนภาวนียสฺส ภิกฺขุโน ทสฺสนาย อุปสงฺกมิตุํ.}
\prose{\csromanfolio{281}ปุน จปรํ ภิกฺขุ ยสฺมิํ สมเย ภิกฺขุ อุทฺธจฺจกุกฺกุจฺจ ปริยุฏฺฐิเตน เจตสา วิหรติ อุทฺธจฺจกุกฺกุจฺจปเรเตน, อุปฺปนฺนสฺส จ อุทฺธจฺจกุกฺกุจฺจสฺส นิสฺสรณํ ยถาภูตํ นปฺปชานาติ, ตสฺมิํ สมเย มโนภาวนีโย ภิกฺขุ อุปสงฺกมิตฺวา เอวมสฺส วจนีโย “อหํ โข อาวุโส อุทฺธจฺจกุกฺกุจฺจปริยุฏฺฐิเตน เจตสา วิหรามิ อุทฺธจฺจกุกฺกุจฺจปเรเตน, อุปฺปนฺนสฺส จ อุทฺธจฺจกุกฺกุจฺจสฺส นิสฺสรณํ ยถาภูตํ นปฺปชานามิ, สาธุ วต เม อายสฺมา อุทฺธจฺจกุกฺกุจฺจสฺส ปหานาย ธมฺมํ เทเสตู”ติ.\csromanspacer{} ตสฺส มโนภาวนีโย ภิกฺขุ อุทฺธจฺจกุกฺกุจฺจสฺส ปหานาย ธมฺมํ เทเสติ.\csromanspacer{} อยํ ภิกฺขุ จตุตฺโถ สมโย มโนภาวนียสฺส ภิกฺขุโน ทสฺสนาย อุปสงฺกมิตุํ.}

\prose{ปุน จปรํ ภิกฺขุ ยสฺมิํ สมเย ภิกฺขุ วิจิกิจฺฉาปริยุฏฺฐิเตน เจตสา วิหรติ วิจิกิจฺฉาปเรเตน, อุปฺปนฺนาย จ วิจิกิจฺฉาย นิสฺสรณํ ยถาภูตํ นปฺปชานาติ, ตสฺมิํ สมเย มโนภาวนีโย ภิกฺขุ อุปสงฺกมิตฺวา เอวมสฺส วจนีโย “อหํ อาวุโส วิจิกิจฺฉาปริยุฏฺฐิเตน เจตสา วิหรามิ วิจิกิจฺฉาปเรเตน, อุปฺปนฺนาย จ วิจิกิจฺฉาย นิสฺสรณํ ยถาภูตํ นปฺปชานามิ, สาธุ วต เม อายสฺมา วิจิกิจฺฉาย ปหานาย ธมฺมํ เทเสตู”ติ.\csromanspacer{} ตสฺส มโนภาวนีโย ภิกฺขุ วิจิกิจฺฉาย ปหานาย ธมฺมํ เทเสติ.\csromanspacer{} อยํ ภิกฺขุ ปญฺจโม สมโย มโนภาวนียสฺส ภิกฺขุโน ทสฺสนาย อุปสงฺกมิตุํ.}

\prose{ปุน จปรํ ภิกฺขุ ยสฺมิํ สมเย ภิกฺขุ ยํ นิมิตฺตํ อาคมฺม ยํ นิมิตฺตํ มนสิกโรโต อนนฺตรา อาสวานํ ขโย โหติ, ตํ นิมิตฺตํ นปฺปชานาติ, ตสฺมิํ สมเย มโนภาวนีโย ภิกฺขุ อุปสงฺกมิตฺวา เอวมสฺส วจนีโย “อหํ โข อาวุโส ยํ นิมิตฺตํ อาคมฺม ยํ นิมิตฺตํ มนสิกโรโต อนนฺตรา อาสวานํ ขโย โหติ, ตํ นิมิตฺตํ นปฺปชานามิ, สาธุ วต เม อายสฺมา อาสวานํ ขยาย ธมฺมํ เทเสตู”ติ.\csromanspacer{} ตสฺส มโนภาวนีโย ภิกฺขุ อาสวานํ ขยาย ธมฺมํ เทเสติ.\csromanspacer{} อยํ ภิกฺขุ ฉฏฺโฐ สมโย มโนภาวนียสฺส ภิกฺขุโน ทสฺสนาย อุปสงฺกมิตุํ.\csromanspacer{} อิเม โข ภิกฺขุ ฉ สมยา มโนภาวนียสฺส ภิกฺขุโน ทสฺสนาย อุปสงฺกมิตุนฺติ.\csromanspacer{}\csromanspacer{}สตฺตมํ.}

\gambhira{ฉกฺกนิปาตปาฬิ \textbf{8. ทุติยสมยสุตฺต}}
\proseitem{28}{\csromanfolio{282}เอกํ สมยํ สมฺพหุลา เถรา ภิกฺขู พาราณสิยํ วิหรนฺติ อิสิปตเน มิคทาเย.\csromanspacer{} อถ โข เตสํ เถรานํ ภิกฺขูนํ ปจฺฉาภตฺตํ ปิณฺฑปาตปฏิกฺกนฺตานํ มณฺฑลมาเฬ สนฺนิสินฺนานํ สนฺนิปติตานํ อยมนฺตรากถา อุทปาทิ “โก นุ โข อาวุโส สมโย มโนภาวนียสฺส ภิกฺขุโน ทสฺสนาย อุปสงฺกมิตุนฺ”ติ.}

\prose{เอวํ วุตฺเต อญฺญตโร ภิกฺขุ เถเร ภิกฺขู เอตทโวจ “ยสฺมิํ อาวุโส สมเย มโนภาวนีโย ภิกฺขุ ปจฺฉาภตฺตํ ปิณฺฑปาตปฏิกฺกนฺโต ปาเท ปกฺขาเลตฺวา นิสินฺโน โหติ ปลฺลงฺกํ อาภุชิตฺวา อุชุํ กายํ ปณิธาย ปริมุขํ สติํ อุปฏฺฐเปตฺวา, โส สมโย มโนภาวนียสฺส ภิกฺขุโน ทสฺสนาย อุปสงฺกมิตุนฺ”ติ.}

\prose{เอวํ วุตฺเต อญฺญตโร ภิกฺขุ ตํ ภิกฺขุํ เอตทโวจ “น โข อาวุโส โส สมโย มโนภาวนียสฺส ภิกฺขุโน ทสฺสนาย อุปสงฺกมิตุํ.\csromanspacer{} ยสฺมิํ อาวุโส สมเย มโนภาวนีโย ภิกฺขุ ปจฺฉาภตฺตํ ปิณฺฑปาตปฏิกฺกนฺโต ปาเท ปกฺขาเลตฺวา นิสินฺโน โหติ ปลฺลงฺกํ อาภุชิตฺวา อุชุํ กายํ ปณิธาย ปริมุขํ สติํ อุปฏฺฐเปตฺวา, จาริตฺตกิลมโถปิสฺส ตสฺมิํ สมเย อปฺปฏิปฺปสฺสทฺโธ โหติ, ภตฺตกิลมโถปิสฺส ตสฺมิํ สมเย อปฺปฏิปฺปสฺสทฺโธ โหติ, ตสฺมา โส อสมโย มโนภาวนียสฺส ภิกฺขุโน ทสฺสนาย อุปสงฺกมิตุํ.\csromanspacer{} ยสฺมิํ อาวุโส สมเย มโนภาวนีโย ภิกฺขุ สายนฺหสมยํ ปฏิสลฺลานา วุฏฺฐิโต วิหารปจฺฉายายํ นิสินฺโน โหติ ปลฺลงฺกํ อาภุชิตฺวา อุชุํ กายํ ปณิธาย ปริมุขํ สติํ อุปฏฺฐเปตฺวา, โส สมโย มโนภาวนียสฺส ภิกฺขุโน ทสฺสนาย อุปสงฺกมิตุนฺ”ติ.}

\prose{เอวํ วุตฺเต อญฺญตโร ภิกฺขุ ตํ ภิกฺขุํ เอตทโวจ “น โข อาวุโส โส สมโย มโนภาวนียสฺส ภิกฺขุโน ทสฺสนาย อุปสงฺกมิตุํ.\csromanspacer{} ยสฺมิํ อาวุโส สมเย มโนภาวนีโย ภิกฺขุ สายนฺหสมยํ ปฏิสลฺลานา วุฏฺฐิโต วิหารปจฺฉายายํ นิสินฺโน โหติ ปลฺลงฺกํ อาภุชิตฺวา อุชุํ กายํ ปณิธาย ปริมุขํ สติํ อุปฏฺฐเปตฺวา.\csromanspacer{} ยเทวสฺส ทิวา สมาธินิมิตฺตํ มนสิกตํ โหติ, ตเทวสฺส ตสฺมิํ สมเย สมุทาจรติ.\csromanspacer{} ตสฺมา โส อสมโย มโนภาวนียสฺส ภิกฺขุโน}

\vspace{\tipitakaparskip}
\csromancenter{อนุตฺตริยวคฺค ทสฺสนาย อุปสงฺกมิตุํ.\csromanspacer{} ยสฺมิํ อาวุโส สมเย มโนภาวนีโย ภิกฺขุ รตฺติยา ปจฺจูสสมยํ ปจฺจุฏฺฐาย นิสินฺโน โหติ ปลฺลงฺกํ อาภุชิตฺวา อุชุํ กายํ ปณิธาย ปริมุขํ สติํ อุปฏฺฐเปตฺวา, โส สมโย มโนภาวนียสฺส ภิกฺขุโน ทสฺสนาย อุปสงฺกมิตุนฺ”ติ.}
\prose{\csromanfolio{283}เอวํ วุตฺเต อญฺญตโร ภิกฺขุ ตํ ภิกฺขุํ เอตทโวจ “น โข อาวุโส โส สมโย มโนภาวนียสฺส ภิกฺขุโน ทสฺสนาย อุปสงฺกมิตุํ.\csromanspacer{} ยสฺมิํ อาวุโส สมเย มโนภาวนีโย ภิกฺขุ รตฺติยา ปจฺจูสสมยํ ปจฺจุฏฺฐาย นิสินฺโน โหติ ปลฺลงฺกํ อาภุชิตฺวา อุชุํ กายํ ปณิธาย ปริมุขํ สติํ อุปฏฺฐเปตฺวา, โอชฏฺฐายิสฺส ตสฺมิํ สมเย กาโย โหติ, ผาสุสฺส โหติ พุทฺธานํ สาสนํ มนสิ กาตุํ.\csromanspacer{} ตสฺมา โส อสมโย มโน-อภาวนียสฺส ภิกฺขุโน ทสฺสนาย อุปสงฺกมิตุนฺ”ติ.}

\prose{เอวํ วุตฺเต อายสฺมา มหากจฺจาโน เถเร ภิกฺขู เอตทโวจ\csromandash{}สมฺมุขา เมตํ อาวุโส ภควโต สุตํ สมฺมุขา ปฏิคฺคหิตํ “ฉยิเม ภิกฺขุ สมยา มโนภาวนียสฺส ภิกฺขุโน ทสฺสนาย อุปสงฺกมิตุํ.}

\prose{กตเม ฉ? อิธ ภิกฺขุ ยสฺมิํ สมเย ภิกฺขุ กามราคปริยุฏฺฐิเตน เจตสา วิหรติ กามราคปเรเตน, อุปฺปนฺนสฺส จ กามราคสฺส นิสฺสรณํ ยถาภูตํ นปฺปชานาติ.\csromanspacer{} ตสฺมิํ สมเย มโนภาวนีโย ภิกฺขุ อุปสงฺกมิตฺวา เอวมสฺส วจนีโย ‘อหํ โข อาวุโส กามราคปริยุฏฺฐิเตน เจตสา วิหรามิ กามราคปเรเตน, อุปฺปนฺนสฺส จ กามราคสฺส นิสฺสรณํ ยถาภูตํ นปฺปชานามิ, สาธุ วต เม อายสฺมา กามราคสฺส ปหานาย ธมฺมํ เทเสตู”ติ.\csromanspacer{} ตสฺส มโนภาวนีโย ภิกฺขุ กามราคสฺส ปหานาย ธมฺมํ เทเสติ.\csromanspacer{} อยํ ภิกฺขุ ปฐโม สมโย มโนภาวนียสฺส ภิกฺขุโน ทสฺสนาย อุปสงฺกมิตุํ.}

\prose{ปุน จปรํ ภิกฺขุ ยสฺมิํ สมเย ภิกฺขุ พฺยาปาทปริยุฏฺฐิเตน เจตสา วิหรติ ฯเปฯ.\csromanspacer{} ถินมิทฺธปริยุฏฺฐิเตน เจตสา วิหรติ.\csromanspacer{} อุทฺธจฺจกุกฺกุจฺจ ปริยุฏฺฐิเตน เจตสา วิหรติ.\csromanspacer{} วิจิกิจฺฉาปริยุฏฺฐิเตน เจตสา วิหรติ.\csromanspacer{} ยํ นิมิตฺตํ อาคมฺม ยํ นิมิตฺตํ มนสิกโรโต อนนฺตรา อาสวานํ ขโย โหติ, ตํ นิมิตฺตํ น ชานาติ น ปสฺสติ.\csromanspacer{} ตสฺมิํ สมเย มโนภาวนีโย ภิกฺขุ อุปสงฺกมิตฺวา เอวมสฺส วจนีโย ‘อหํ โข อาวุโส ยํ นิมิตฺตํ}

\vspace{\tipitakaparskip}
\csromancenter{ฉกฺกนิปาตปาฬิ อาคมฺม ยํ นิมิตฺตํ มนสิกโรโต อนนฺตรา อาสวานํ ขโย โหติ, ตํ นิมิตฺตํ น ชานามิ น ปสฺสามิ, สาธุ วต เม อายสฺมา อาสวานํ ขยาย ธมฺมํ เทเสตู’ติ.\csromanspacer{} ตสฺส มโนภาวนีโย ภิกฺขุ อาสวานํ ขยาย ธมฺมํ เทเสติ.\csromanspacer{} อยํ ภิกฺขุ ฉฏฺโฐ สมโย มโนภาวนียสฺส ภิกฺขุโน ทสฺสนาย อุปสงฺกมิตุํ.}
\prose{\csromanfolio{284}สมฺมุขา เมตํ อาวุโส ภควโต สุตํ สมฺมุขา ปฏิคฺคหิตํ.\csromanspacer{} อิเม โข ภิกฺขุ ฉ สมยา มโนภาวนียสฺส ภิกฺขุโน ทสฺสนาย อุปสงฺกมิตุนฺ”ติ.\csromanspacer{}\csromanspacer{}อฏฺฐมํ.}

\csromanmatikaanchor{mtk.149}
\csromantocmarkref{subsection}{9. อุทายีสุตฺต}{mtk.149}
\bookmark[dest={mtk.149},level=2]{9. อุทายีสุตฺต}
\csromanheader{2}{9. อุทายีสุตฺต}
\proseitem{29}{อถ โข ภควา อายสฺมนฺตํ อุทายิํ อามนฺเตสิ “กติ นุ โข อุทายิ อนุสฺสติฏฺฐานานี”ติ.\csromanspacer{} เอวํ วุตฺเต อายสฺมา อุทายี ตุณฺหี อโหสิ.\csromanspacer{} ทุติยมฺปิ โข ภควา อายสฺมนฺตํ อุทายิํ อามนฺเตสิ “กติ นุ โข อุทายิ อนุสฺสติฏฺฐานานี”ติ.\csromanspacer{} ทุติยมฺปิ โข อายสฺมา อุทายี ตุณฺหี อโหสิ.\csromanspacer{} ตติยมฺปิ โข ภควา อายสฺมนฺตํ อุทายิํ อามนฺเตสิ “กติ นุ โข อุทายิ อนุสฺสติฏฺฐานานี”ติ.\csromanspacer{} ตติยมฺปิ โข อายสฺมา อุทายี ตุณฺหี อโหสิ.}

\prose{อถ โข อายสฺมา อานนฺโท อายสฺมนฺตํ อุทายิํ เอตทโวจ “สตฺถา ตํ อาวุโส อุทายิ อามนฺเตสี”ติ.\csromanspacer{} สุโณมหํ อาวุโส อานนฺท ภควโต.\csromanspacer{} อิธ ภนฺเต ภิกฺขุ อเนกวิหิตํ ปุพฺเพนิวาสํ อนุสฺสรติ.\csromanspacer{} เสยฺยถิทํ, เอกมฺปิ ชาติํ เทฺวปิ ชาติโย ฯเปฯ อิติ สาการํ ส-อุทฺเทสํ อเนกวิหิตํ ปุพฺเพนิวาสํ อนุสฺสรติ, อิทํ ภนฺเต อนุสฺสติฏฺฐานนฺติ.}

\prose{อถ โข ภควา อายสฺมนฺตํ อานนฺทํ อามนฺเตสิ “อญฺญาสิํ โข อหํ อานนฺท ‘เนวายํ อุทายี โมฆปุริโส อธิจิตฺตํ อนุยุตฺโต วิหรตี’ติ”.\csromanspacer{} กติ นุ โข อานนฺท อนุสฺสติฏฺฐานานีติ.}

\prose{ปญฺจ ภนฺเต อนุสฺสติฏฺฐานานิ.\csromanspacer{} กตมานิ ปญฺจ? อิธ ภนฺเต ภิกฺขุ วิวิจฺเจว กาเมหิ ฯเปฯ ตติยํ ฌานํ อุปสมฺปชฺช วิหรติ.\csromanspacer{} อิทํ ภนฺเต อนุสฺสติฏฺฐานํ เอวํ ภาวิตํ เอวํ พหุลีกตํ ทิฏฺฐธมฺมสุขวิหาราย สํวตฺตติ.}

\prose{ปุน จปรํ ภนฺเต ภิกฺขุ อาโลกสญฺญํ มนสิ กโรติ, ทิวา สญฺญํ อธิฏฺฐาติ, ยถา ทิวา ตถา รตฺติํ ยถา รตฺติํ ตถา ทิวา, อิติ วิวเฏน}

\vspace{\tipitakaparskip}
\csromancenter{อนุตฺตริยวคฺค เจตสา อปริโยนทฺเธน สปฺปภาสํ จิตฺตํ ภาเวติ.\csromanspacer{} อิทํ ภนฺเต อนุสฺสติฏฺฐานํ เอวํ ภาวิตํ เอวํ พหุลีกตํ ญาณทสฺสนปฺปฏิลาภาย สํวตฺตติ.}
\prose{\csromanfolio{285}ปุน จปรํ ภนฺเต ภิกฺขุ อิมเมว กายํ อุทฺธํ ปาทตลา อโธ เกสมตฺถกา ตจปริยนฺตํ ปูรํ นานปฺปการสฺส อสุจิโน ปจฺจเวกฺขติ “อตฺถิ อิมสฺมิํ กาเย เกสา โลมา นขา ทนฺตา ตโจ, มํสํ นฺหารุ\footnote{นหารุ (สี, อิ) ที 2. 233; ม 1. 72 ปิฏฺฐาทีสุปิ.} อฏฺฐิ อฏฺฐิมิญฺชํ วกฺกํ, หทยํ ยกนํ กิโลมกํ ปิหกํ ปปฺผาสํ, อนฺตํ อนฺตคุณํ อุทริยํ กรีสํ, ปิตฺตํ เสมฺหํ ปุพฺโพ โลหิตํ เสโท เมโท, อสฺสุ วสา เขโฬ สิงฺฆาณิกา ลสิกา มุตฺตนฺ”ติ.\csromanspacer{} อิทํ ภนฺเต อนุสฺสติฏฺฐานํ เอวํ ภาวิตํ เอวํ พหุลีกตํ กามราคปฺปหานาย สํวตฺตติ.}

\prose{ปุน จปรํ ภนฺเต ภิกฺขุ เสยฺยถาปิ ปสฺเสยฺย สรีรํ สิวถิกาย ฉฏฺฏิตํ\footnote{ฉฑฺฑิตํ (สี, สฺยา, อิ)} เอกาหมตํ วา ทฺวีหมตํ วา ตีหมตํ วา อุทฺธุมาตกํ วินีลกํ วิปุพฺพกชาตํ.\csromanspacer{} โส อิมเมว กายํ เอวํ\footnote{เอวนฺติ อิทํ สติปฏฺฐานสุตฺตาทีสุ นตฺถิ.} อุปสํหรติ “อยมฺปิ โข กาโย เอวํธมฺโม เอวํภาวี เอวํอนตีโต”ติ\footnote{เอตํ อนตีโตติ (สี)}.}

\prose{เสยฺยถาปิ วา ปน\footnote{เสยฺยถา วา ปน (สฺยา)} ปสฺเสยฺย สรีรํ สีวถิกาย ฉฏฺฏิตํ กาเกหิ วา ขชฺชมานํ กุลเลหิ วา ขชฺชมานํ คิชฺเฌหิ วา ขชฺชมานํ สุนเขหิ วา ขชฺชมานํ สิงฺคาเลหิ\footnote{สิคาเลหิ (สี)} วา ขชฺชมานํ วิวิเธหิ วา ปาณกชาเตหิ ขชฺชมานํ.\csromanspacer{} โส อิมเมว กายํ เอวํ อุปสํหรติ “อยมฺปิ โข กาโย เอวํธมฺโม เอวํภาวี เอวํอนตีโต”ติ.}

\prose{เสยฺยถาปิ วา ปน ปสฺเสยฺย สรีรํ สีวถิกาย ฉฏฺฏิตํ อฏฺฐิกสงฺขลิกํ สมํสโลหิตํ นฺหารุสมฺพนฺธํ ฯเปฯ.\csromanspacer{} อฏฺฐิกสงฺขลิกํ นิมฺมํสโลหิตมกฺขิตํ นฺหารุสมฺพนฺธํ.\csromanspacer{} อฏฺฐิกสงฺขลิกํ อปคตมํสโลหิตํ นฺหารุสมฺพนฺธํ.\csromanspacer{} อฏฺฐิกานิ อปคตสมฺพนฺธานิ ทิสาวิทิสาวิกฺขิตฺตานิ\footnote{ทิสาวิทิสาสุ วิกฺขิตฺตานิ (สี)}, อญฺเญน หตฺถฏฺฐิกํ อญฺเญน ปาทฏฺฐิกํ อญฺเญน ชงฺฆฏฺฐิกํ อญฺเญน อูรุฏฺฐิกํ อญฺเญน กฏิฏฺฐิกํ\footnote{กฏฏฺฐิกํ (สี)} อญฺเญน}

\vspace{\tipitakaparskip}
\csromancenter{ฉกฺกนิปาตปาฬิ 1ผาสุกฏฺฐิกํ อญฺเญน ปิฏฺฐิกณฺฏกฏฺฐิกํ อญฺเญน ขนฺธฏฺฐิกํ อญฺเญน คีวฏฺฐิกํ อญฺเญน หนุกฏฺฐิกํ อญฺเญน ทนฺตกฏฺฐิกํ อญฺเญน สีสกฏาหํ1.\csromanspacer{} อฏฺฐิกานิ เสตานิ สงฺขวณฺณปฺปฏิภาคานิ\footnote{สงฺขวณฺณูปนิภานิ (สี, สฺยา, อิ)} อฏฺฐิกานิ ปุญฺชกิตานิ\footnote{ปุญฺชกตานิ (อิ)} เตโรวสฺสิกานิ.\csromanspacer{} อฏฺฐิกานิ ปูตีนิ จุณฺณกชาตานิ.\csromanspacer{} โส อิมเมว กายํ เอวํ อุปสํหรติ. “อยมฺปิ โข กาโย เอวํธมฺโม เอวํภาวี เอวํอนตีโต”ติ.\csromanspacer{} อิทํ ภนฺเต อนุสฺสติฏฺฐานํ เอวํ ภาวิตํ เอวํ พหุลีกตํ อสฺมิมานสมุคฺฆาตาย สํวตฺตติ.}
\prose{\csromanfolio{286}ปุน จปรํ ภนฺเต ภิกฺขุ สุขสฺส จ ปหานา ฯเปฯ จตุตฺถํ ฌานํ อุปสมฺปชฺช วิหรติ.\csromanspacer{} อิทํ ภนฺเต อนุสฺสติฏฺฐานํ เอวํ ภาวิตํ เอวํ พหุลีกตํ อเนกธาตุปฏิเวธาย สํวตฺตติ.\csromanspacer{} อิมานิ โข ภนฺเต ปญฺจ อนุสฺสติฏฺฐานานีติ.}

\prose{สาธุ สาธุ อานนฺท, เตน หิ ตฺวํ อานนฺท อิทมฺปิ ฉฏฺฐํ อนุสฺสติฏฺฐานํ ธาเรหิ.\csromanspacer{} อิธานนฺท ภิกฺขุ สโตว อภิกฺกมติ, สโตว ปฏิกฺกมติ, สโตว ติฏฺฐติ, สโตว นิสีทติ, สโตว เสยฺยํ กปฺเปติ, สโตว กมฺมํ อธิฏฺฐาติ.\csromanspacer{} อิทํ อานนฺท อนุสฺสติฏฺฐานํ เอวํ ภาวิตํ เอวํ พหุลีกตํ สติสมฺปชญฺญาย สํวตฺตตีติ.\csromanspacer{}\csromanspacer{}นวมํ.}

\csromanmatikaanchor{mtk.150}
\csromantocmarkref{subsection}{10. อนุตฺตริยสุตฺต}{mtk.150}
\bookmark[dest={mtk.150},level=2]{10. อนุตฺตริยสุตฺต}
\csromanheader{2}{10. อนุตฺตริยสุตฺต}
\proseitem{30}{ฉยิมานิ ภิกฺขเว อนุตฺตริยานิ.\csromanspacer{} กตมานิ ฉ? ทสฺสนานุตฺตริยํ สวนานุตฺตริยํ ลาภานุตฺตริยํ สิกฺขานุตฺตริยํ ปาริจริยานุตฺตริยํ อนุสฺสตานุตฺตริยนฺติ.}

\prose{กตมญฺจ ภิกฺขเว ทสฺสนานุตฺตริยํ? อิธ ภิกฺขเว เอกจฺโจ หตฺถิรตนมฺปิ ทสฺสนาย คจฺฉติ, อสฺสรตนมฺปิ ทสฺสนาย คจฺฉติ, มณิรตนมฺปิ ทสฺสนาย คจฺฉติ, อุจฺจาวจํ วา ปน ทสฺสนาย คจฺฉติ, สมณํ วา พฺราหฺมณํ วา มิจฺฉาทิฏฺฐิกํ มิจฺฉาปฏิปนฺนํ ทสฺสนาย คจฺฉติ.\csromanspacer{} อตฺเถตํ ภิกฺขเว ทสฺสนํ, เนตํ นตฺถีติ วทามิ.\csromanspacer{} ตญฺจ โข เอตํ ภิกฺขเว ทสฺสนํ หีนํ คมฺมํ โปถุชฺชนิกํ อนริยํ อนตฺถสํหิตํ น นิพฺพิทาย น วิราคาย น นิโรธาย น อุปสมาย น อภิญฺญาย น สมฺโพธาย น นิพฺพานาย สํวตฺตติ.\csromanspacer{} โย จ โข ภิกฺขเว ตถาคตํ วา ตถาคตสาวกํ วา ทสฺสนาย}

\vspace{\tipitakaparskip}
\csromancenter{อนุตฺตริยวคฺค คจฺฉติ นิวิฏฺฐสทฺโธ นิวิฏฺฐเปโม เอกนฺตคโต อภิปฺปสนฺโน.\csromanspacer{} เอตทานุตฺตริยํ ภิกฺขเว ทสฺสนานํ สตฺตานํ วิสุทฺธิยา โสกปริเทวานํ\footnote{โสกปริทฺทวานํ (สี)} สมติกฺกมาย ทุกฺขโทมนสฺสานํ อตฺถงฺคมาย\footnote{อตฺถคมาย (สี)} ญายสฺส อธิคมาย นิพฺพานสฺส สจฺฉิกิริยาย, ยทิทํ ตถาคตํ วา ตถาคตสาวกํ วา ทสฺสนาย คจฺฉติ นิวิฏฺฐสทฺโธ นิวิฏฺฐเปโม เอกนฺตคโต อภิปฺปสนฺโน.\csromanspacer{} อิทํ วุจฺจติ ภิกฺขเว ทสฺสนานุตฺตริยํ.\csromanspacer{} อิติ ทสฺสนานุตฺตริยํ.}
\prose{\csromanfolio{287}สวนานุตฺตริยญฺจ กถํ โหติ? อิธ ภิกฺขเว เอกจฺโจ เภริสทฺทมฺปิ\footnote{เภริสทฺทสฺสปิ (ก) เอวํ วีณาสทฺทมฺปิ-อิจฺจาทีสุปิ.} สวนาย คจฺฉติ, วีณาสทฺทมฺปิ สวนาย คจฺฉติ, คีตสทฺทมฺปิ สวนาย คจฺฉติ, อุจฺจาวจํ วา ปน สวนาย คจฺฉติ, สมณสฺส วา พฺราหฺมณสฺส วา มิจฺฉาทิฏฺฐิกสฺส มิจฺฉาปฏิปนฺนสฺส ธมฺมสฺสวนาย คจฺฉติ.\csromanspacer{} อตฺเถตํ ภิกฺขเว สวนํ, เนตํ นตฺถีติ วทามิ.\csromanspacer{} ตญฺจ โข เอตํ ภิกฺขเว สวนํ หีนํ คมฺมํ โปถุชฺชนิกํ อนริยํ อนตฺถสํหิตํ น นิพฺพิทาย น วิราคาย น นิโรธาย น อุปสมาย น อภิญฺญาย น สมฺโพธาย น นิพฺพานาย สํวตฺตติ.\csromanspacer{} โย จ โข ภิกฺขเว ตถาคตสฺส วา ตถาคตสาวกสฺส วา ธมฺมสฺสวนาย คจฺฉติ นิวิฏฺฐสทฺโธ นิวิฏฺฐเปโม เอกนฺตคโต อภิปฺปสนฺโน.\csromanspacer{} เอตทานุตฺตริยํ ภิกฺขเว สวนานํ สตฺตานํ วิสุทฺธิยา โสกปริเทวานํ สมติกฺกมาย ทุกฺขโทมนสฺสานํ อตฺถงฺคมาย ญายสฺส อธิคมาย นิพฺพานสฺส สจฺฉิกิริยาย, ยทิทํ ตถาคตสฺส วา ตถาคตสาวกสฺส วา ธมฺมสฺสวนาย คจฺฉติ นิวิฏฺฐสทฺโธ นิวิฏฺฐเปโม เอกนฺตคโต อภิปฺปสนฺโน.\csromanspacer{} อิทํ วุจฺจติ ภิกฺขเว สวนานุตฺตริยํ.\csromanspacer{} อิติ ทสฺสนานุตฺตริยํ สวนานุตฺตริยํ.}

\prose{ลาภานุตฺตริยญฺจ กถํ โหติ? อิธ ภิกฺขเว เอกจฺโจ ปุตฺตลาภมฺปิ ลภติ, ทารลาภมฺปิ ลภติ, ธนลาภมฺปิ ลภติ, อุจฺจาวจํ วา ปน ลาภํ ลภติ, สมเณ วา พฺราหฺมเณ วา มิจฺฉาทิฏฺฐิเก มิจฺฉาปฏิปนฺเน สทฺธํ ปฏิลภติ.\csromanspacer{} อตฺเถโส ภิกฺขเว ลาโภ, เนโส นตฺถีติ วทามิ.\csromanspacer{} โส จ โข เอโส ภิกฺขเว ลาโภ หีโน คมฺโม โปถุชฺชนิโก อนริโย อนตฺถสํหิโต น นิพฺพิทาย น วิราคาย น นิโรธาย น อุปสมาย น อภิญฺญาย น สมฺโพธาย น นิพฺพานาย สํวตฺตติ.\csromanspacer{} โย จ โข ภิกฺขเว ตถาคเต วา ตถาคตสาวเก}

\vspace{\tipitakaparskip}
\csromancenter{ฉกฺกนิปาตปาฬิ วา สทฺธํ ปฏิลภติ นิวิฏฺฐสทฺโธ นิวิฏฺฐเปโม เอกนฺตคโต อภิปฺปสนฺโน.\csromanspacer{} เอตทานุตฺตริยํ ภิกฺขเว ลาภานํ สตฺตานํ วิสุทฺธิยา โสกปริเทวานํ สมติกฺกมาย ทุกฺขโทมนสฺสานํ อตฺถงฺคมาย ญายสฺส อธิคมาย นิพฺพานสฺส สจฺฉิกิริยาย, ยทิทํ ตถาคเต วา ตถาคตสาวเก วา สทฺธํ ปฏิลภติ นิวิฏฺฐสทฺโธ นิวิฏฺฐเปโม เอกนฺตคโต อภิปฺปสนฺโน.\csromanspacer{} อิทํ วุจฺจติ ภิกฺขเว ลาภานุตฺตริยํ, อิติ ทสฺสนานุตฺตริยํ สวนานุตฺตริยํ ลาภานุตฺตริยํ.}
\prose{\csromanfolio{288}สิกฺขานุตฺตริยญฺจ กถํ โหติ? อิธ ภิกฺขเว เอกจฺโจ หตฺถิสฺมิมฺปิ สิกฺขติ, อสฺสสฺมิมฺปิ สิกฺขติ, รถสฺมิมฺปิ สิกฺขติ, ธนุสฺมิมฺปิ สิกฺขติ, ถรุสฺมิมฺปิ สิกฺขติ, อุจฺจาวจํ วา ปน สิกฺขติ, สมณสฺส วา พฺราหฺมณสฺส วา มิจฺฉาทิฏฺฐิกสฺส มิจฺฉาปฏิปนฺนสฺส\footnote{มิจฺฉาปฏิปตฺติํ (ก)} สิกฺขติ.\csromanspacer{} อตฺเถสา ภิกฺขเว สิกฺขา, เนสา นตฺถีติ วทามิ.\csromanspacer{} สา จ โข เอสา ภิกฺขเว สิกฺขา หีนา คมฺมา โปถุชฺชนิกา อนริยา อนตฺถสํหิตา น นิพฺพิทาย น วิราคาย น นิโรธาย น อุปสมาย น อภิญฺญาย น สมฺโพธาย น นิพฺพานาย สํวตฺตติ.\csromanspacer{} โย จ โข ภิกฺขเว ตถาคตปฺปเวทิโต ธมฺมวินเย อธิสีลมฺปิ สิกฺขติ, อธิจิตฺตมฺปิ สิกฺขติ, อธิปญฺญมฺปิ สิกฺขติ, นิวิฏฺฐสทฺโธ นิวิฏฺฐเปโม เอกนฺตคโต อภิปฺปสนฺโน.\csromanspacer{} เอตทานุตฺตริยํ ภิกฺขเว สิกฺขานํ สตฺตานํ วิสุทฺธิยา โสกปริเทวานํ สมติกฺกมาย ทุกฺขโทมนสฺสานํ อตฺถงฺคมาย ญายสฺส อธิคมาย นิพฺพานสฺส สจฺฉิกิริยาย, ยทิทํ ตถาคตปฺปเวทิเต ธมฺมวินเย อธิสีลมฺปิ สิกฺขติ, อธิจิตฺตมฺปิ สิกฺขติ, อธิปญฺญมฺปิ สิกฺขติ, นิวิฏฺฐสทฺโธ นิวิฏฺฐเปโม เอกนฺตคโต อภิปฺปสนฺโน.\csromanspacer{} อิทํ วุจฺจติ ภิกฺขเว สิกฺขานุตฺตริยํ.\csromanspacer{} อิติ ทสฺสนานุตฺตริยํ สวนานุตฺตริยํ ลาภานุตฺตริยํ สิกฺขานุตฺตริยํ.}

\prose{ปาริจริยานุตฺตริยญฺจ กถํ โหติ? อิธ ภิกฺขเว เอกจฺโจ ขตฺติยมฺปิ ปริจรติ, พฺราหฺมณมฺปิ ปริจรติ, คหปติมฺปิ ปริจรติ, อุจฺจาวจํ วา ปน ปริจรติ, สมณํ วา พฺราหฺมณํ วา มิจฺฉาทิฏฺฐิกํ มิจฺฉาปฏิปนฺนํ ปริจรติ.\csromanspacer{} อตฺเถสา ภิกฺขเว ปาริจริยา, เนสา นตฺถีติ วทามิ.\csromanspacer{} สา จ โข เอสา ภิกฺขเว ปาริจริยา หีนา คมฺมา โปถุชฺชนิกา อนริยา อนตฺถสํหิตา น นิพฺพิทาย ฯเปฯ น นิพฺพานาย สํวตฺตติ.\csromanspacer{} โย จ โข ภิกฺขเว ตถาคตํ วา ตถาคตสาวกํ วา ปริจรติ นิวิฏฺฐสทฺโธ นิวิฏฺฐเปโม เอกนฺตคโต}

\vspace{\tipitakaparskip}
\csromancenter{อนุตฺตริยวคฺค อภิปฺปสนฺโน.\csromanspacer{} เอตทานุตฺตริยํ ภิกฺขเว ปาริจริยานํ สตฺตานํ วิสุทฺธิยา โสกปริเทวานํ สมติกฺกมาย ทุกฺขโทมนสฺสานํ อตฺถงฺคมาย ญายสฺส อธิคมาย นิพฺพานสฺส สจฺฉิกิริยาย ยทิทํ ตถาคตํ วา ตถาคตสาวกํ วา ปริจรติ นิวิฏฺฐสทฺโธ นิวิฏฺฐเปโม เอกนฺตคโต อภิปฺปสนฺโน.\csromanspacer{} อิทํ วุจฺจติ ภิกฺขเว ปาริจริยานุตฺตริยํ.\csromanspacer{} อิติ ทสฺสนานุตฺตริยํ สวนานุตฺตริยํ ลาภานุตฺตริยํ สิกฺขานุตฺตริยํ ปาริจริยานุตฺตริยํ.}
\prose{\csromanfolio{289}อนุสฺสตานุตฺตริยญฺจ กถํ โหติ? อิธ ภิกฺขเว เอกจฺโจ ปุตฺตลาภมฺปิ อนุสฺสรติ, ทารลาภมฺปิ อนุสฺสรติ, ธนลาภมฺปิ อนุสฺสรติ, อุจฺจาวจํ วา ปน ลาภํ อนุสฺสรติ, สมณํ วา พฺราหฺมณํ วา มิจฺฉาทิฏฺฐิกํ มิจฺฉาปฏิปนฺนํ อนุสฺสรติ.\csromanspacer{} อตฺเถสา ภิกฺขเว อนุสฺสติ, เนสา นตฺถีติ วทามิ.\csromanspacer{} สา จ โข เอสา ภิกฺขเว อนุสฺสติ หีนา คมฺมา โปถุชฺชนิกา อนริยา อนตฺถสํหิตา น นิพฺพิทาย น วิราคาย น นิโรธาย น อุปสมาย น อภิญฺญาย น สมฺโพธาย น นิพฺพานาย สํวตฺตติ.\csromanspacer{} โย จ โข ภิกฺขเว ตถาคตํ วา ตถาคตสาวกํ วา อนุสฺสรติ นิวิฏฺฐสทฺโธ นิวิฏฺฐเปโม เอกนฺตคโต อภิปฺปสนฺโน.\csromanspacer{} เอตทานุตฺตริยํ ภิกฺขเว อนุสฺสตีนํ สตฺตานํ วิสุทฺธิยา โสกปริเทวานํ สมติกฺกมาย ทุกฺขโทมนสฺสานํ อตฺถงฺคมาย ญายสฺส อธิคมาย นิพฺพานสฺส สจฺฉิกิริยาย, ยทิทํ ตถาคตํ วา ตถาคตสาวกํ วา อนุสฺสรติ นิวิฏฺฐสทฺโธ นิวิฏฺฐเปโม เอกนฺตคโต อภิปฺปสนฺโน.\csromanspacer{} อิทํ วุจฺจติ ภิกฺขเว อนุสฺสตานุตฺตริยํ.\csromanspacer{} อิมานิ โข ภิกฺขเว ฉ อนุตฺตริยานีติ.}

\csromangathameasure{เย ทสฺสนานุตฺตรํ ลทฺธา, สวนญฺจ อนุตฺตรํ. \\ ลาภานุตฺตริยํ ลทฺธา, สิกฺขานุตฺตริเย รตา. \\ อุปฏฺฐิตา ปาริจริยา, ภาวยนฺติ อนุสฺสติํ. \\ วิเวกปฺปฏิสํยุตฺตํ, เขมํ อมตคามินิํ. \\ อปฺปมาเท ปมุทิตา, นิปกา สีลสํวุตา. \\ เต เว กาเลน ปจฺเจนฺติ, ยตฺถ ทุกฺขํ นิรุชฺฌตีติ.ทสมํ. อนุตฺตริยวคฺโค ตติโย.}
\csromangathagroup{\csromangathastanza{เย ทสฺสนานุตฺตรํ ลทฺธา\footnote{เย ทสฺสนวรํ ลทฺธา (สี, อิ), ทสฺสนานุตฺตริยํ ลทฺธา (สฺยา, กํ)}, สวนญฺจ อนุตฺตรํ. \\ ลาภานุตฺตริยํ ลทฺธา, สิกฺขานุตฺตริเย รตา\footnote{เย ทสฺสนวรํ ลทฺธา (สี, อิ), ทสฺสนานุตฺตริยํ ลทฺธา (สฺยา, กํ)}.}\csromangathastanza{อุปฏฺฐิตา ปาริจริยา, ภาวยนฺติ อนุสฺสติํ. \\ วิเวกปฺปฏิสํยุตฺตํ, เขมํ อมตคามินิํ.}\csromangathastanza{อปฺปมาเท ปมุทิตา, นิปกา สีลสํวุตา. \\ เต เว กาเลน ปจฺเจนฺติ\footnote{ปจฺจนฺติ (สฺยา, ก)}, ยตฺถ ทุกฺขํ นิรุชฺฌตีติ.\csromanspacer{}\csromanspacer{}ทสมํ.\csromanspacer{} อนุตฺตริยวคฺโค\footnote{สามกวคฺโค (ก)} ตติโย.}}

\gambhira{ฉกฺกนิปาตปาฬิ \textbf{ตสฺสุทฺทานํ}}
\csromanmatikaanchor{mtk.151}
\csromanmatikaanchor{mtk.153}
\csromantocmarknopage{section}{4. เทวตาวคฺค}{mtk.151}
\bookmark[dest={mtk.151},level=1]{4. เทวตาวคฺค}
\csromangathameasure{สามโก อปริหานิโย, ภยํ หิมวานุสฺสติ. \\ กจฺจาโน เทฺว จ สมยา, อุทายี อนุตฺตริเยนาติ.}
\csromangathagroup{\csromangathastanza{\csromanfolio{290}สามโก อปริหานิโย, ภยํ หิมวานุสฺสติ. \\ กจฺจาโน เทฺว จ สมยา, อุทายี อนุตฺตริเยนาติ.}}

\chapterhead{4. เทวตาวคฺค}
\csromanmatikaanchor{mtk.152}
\csromantocmarkref{subsection}{1. เสขสุตฺต}{mtk.152}
\bookmark[dest={mtk.152},level=2]{1. เสขสุตฺต}
\csromantocmarkref{subsection}{2-3. อปริหานสุตฺตาทิ}{mtk.153}
\bookmark[dest={mtk.153},level=2]{2-3. อปริหานสุตฺตาทิ}
\csromanheader{2}{1. เสขสุตฺต}
\proseitem{31}{ฉยิเม ภิกฺขเว ธมฺมา เสขสฺส ภิกฺขุโน ปริหานาย สํวตฺตนฺติ.\csromanspacer{} กตเม ฉ? กมฺมารามตา ภสฺสารามตา นิทฺทารามตา สงฺคณิการามตา อินฺทฺริเยสุ อคุตฺตทฺวารตา โภชเน อมตฺตญฺญุตา.\csromanspacer{} อิเม โข ภิกฺขเว ฉ ธมฺมา เสขสฺส ภิกฺขุโน ปริหานาย สํวตฺตนฺติ.}

\prose{ฉยิเม ภิกฺขเว ธมฺมา เสขสฺส ภิกฺขุโน อปริหานาย สํวตฺตนฺติ.\csromanspacer{} กตเม ฉ? น กมฺมารามตา น ภสฺสารามตา น นิทฺทารามตา น สงฺคณิการามตา อินฺทฺริเยสุ คุตฺตทฺวารตา โภชเน มตฺตญฺญุตา.\csromanspacer{} อิเม โข ภิกฺขเว ฉ ธมฺมา เสขสฺส ภิกฺขุโน อปริหานาย สํวตฺตนฺตีติ.\csromanspacer{}\csromanspacer{}ปฐมํ.}

\csromanheader{2}{2. ปฐม อปริหานสุตฺต}
\proseitem{32}{อถ โข อญฺญตรา เทวตา อภิกฺกนฺตาย รตฺติยา อภิกฺกนฺตวณฺณา เกวลกปฺปํ เชตวนํ โอภาเสตฺวา เยน ภควา เตนุปสงฺกมิ, อุปสงฺกมิตฺวา ภควนฺตํ อภิวาเทตฺวา เอกมนฺตํ อฏฺฐาสิ, เอกมนฺตํ ฐิตา โข สา เทวตา ภควนฺตํ เอตทโวจ\csromandash{}}

\prose{ฉยิเม ภนฺเต ธมฺมา ภิกฺขุโน อปริหานาย สํวตฺตนฺติ.\csromanspacer{} กตเม ฉ? สตฺถุคารวตา ธมฺมคารวตา สํฆคารวตา สิกฺขาคารวตา อปฺปมาทคารวตา ปฏิสนฺถารคารวตา\footnote{ปฏิสนฺธารคารวตา (ก)}.\csromanspacer{} อิเม โข ภนฺเต ฉ ธมฺมา ภิกฺขุโน อปริหานาย สํวตฺตนฺตีติ.\csromanspacer{} อิทมโวจ สา เทวตา, สมนุญฺโญ สตฺถา อโหสิ.\csromanspacer{} อถ โข สา เทวตา “สมนุญฺโญ เม สตฺถา”ติ ภควนฺตํ อภิวาเทตฺวา ปทกฺขิณํ กตฺวา ตตฺเถวนฺตรธายิ.}

\chapterheadpageread{4. เทวตาวคฺค}
\csromanmatikaanchor{mtk.154}
\csromantocmarkref{subsection}{4. มหาโมคฺคลฺลานสุตฺตาทิ}{mtk.154}
\bookmark[dest={mtk.154},level=2]{4. มหาโมคฺคลฺลานสุตฺตาทิ}
\prose{\csromanfolio{291}อถ โข ภควา ตสฺสา รตฺติยา อจฺจเยน ภิกฺขู อามนฺเตสิ “อิมํ ภิกฺขเว รตฺติํ อญฺญตรา เทวตา อภิกฺกนฺตาย รตฺติยา อภิกฺกนฺตวณฺณา เกวลกปฺปํ เชตวนํ โอภาเสตฺวา เยนาหํ เตนุปสงฺกมิ, อุปสงฺกมิตฺวา มํ อภิวาเทตฺวา เอกมนฺตํ อฏฺฐาสิ, เอกมนฺตํ ฐิตา โข ภิกฺขเว สา เทวตา มํ เอตทโวจ ‘ฉยิเม ภนฺเต ธมฺมา ภิกฺขุโน อปริหานาย สํวตฺตนฺติ.\csromanspacer{} กตเม ฉ? สตฺถุคารวตา ธมฺมคารวตา สํฆคารวตา สิกฺขาคารวตา อปฺปมาทคารวตา ปฏิสนฺถารคารวตา.\csromanspacer{} อิเม โข ภนฺเต ฉ ธมฺมา ภิกฺขุโน อปริหานาย สํวตฺตนฺตี’ติ.\csromanspacer{} อิทมโวจ ภิกฺขเว สา เทวตา, อิทํ วตฺวา มํ อภิวาเทตฺวา ปทกฺขิณํ กตฺวา ตตฺเถวนฺตรธายี”ติ.}

\csromangathameasure{สตฺถุครุ ธมฺมครุ, สํเฆ จ ติพฺพคารโว. \\ อปฺปมาทครุ ภิกฺขุ, ปฏิสนฺถารคารโว. \\ อภพฺโพ ปริหานาย, นิพฺพานสฺเสว สนฺติเกติ.ทุติยํ.}
\csromangathagroup{\csromangathastanza{สตฺถุครุ ธมฺมครุ, สํเฆ จ ติพฺพคารโว. \\ อปฺปมาทครุ ภิกฺขุ, ปฏิสนฺถารคารโว. \\ อภพฺโพ ปริหานาย, นิพฺพานสฺเสว สนฺติเกติ.\csromanspacer{}\csromanspacer{}ทุติยํ.}}

\csromanheader{2}{3. ทุติย อปริหานสุตฺต}
\proseitem{33}{อิมํ ภิกฺขเว รตฺติํ อญฺญตรา เทวตา อภิกฺกนฺตาย รตฺติยา อภิกฺกนฺตวณฺณา เกวลกปฺปํ เชตวนํ โอภาเสตฺวา เยนาหํ เตนุปสงฺกมิ, อุปสงฺกมิตฺวา มํ อภิวาเทตฺวา เอกมนฺตํ อฏฺฐาสิ, เอกมนฺตํ ฐิตา โข ภิกฺขเว สา เทวตา มํ เอตทโวจ “ฉยิเม ภนฺเต ธมฺมา ภิกฺขุโน อปริหานาย สํวตฺตนฺติ.\csromanspacer{} กตเม ฉ? สตฺถุคารวตา ธมฺมคารวตา สํฆคารวตา สิกฺขาคารวตา หิริคารวตา โอตฺตปฺปคารวตา.\csromanspacer{} อิเม โข ภนฺเต ฉ ธมฺมา ภิกฺขุโน อปริหานาย สํวตฺตนฺตี”ติ.\csromanspacer{} อิทมโวจ ภิกฺขเว สา เทวตา, อิทํ วตฺวา มํ อภิวาเทตฺวา ปทกฺขิณํ กตฺวา ตตฺเถวนฺตรธายีติ.}

\csromangathameasure{สตฺถุครุ ธมฺมครุ, สํเฆ จ ติพฺพคารโว. \\ หิริ-โอตฺตปฺปสมฺปนฺโน, สปฺปติสฺโส สคารโว. \\ อภพฺโพ ปริหานาย, นิพฺพานสฺเสว สนฺติเกติ.ตติยํ.}
\csromangathagroup{\csromangathastanza{สตฺถุครุ ธมฺมครุ, สํเฆ จ ติพฺพคารโว. \\ หิริ-โอตฺตปฺปสมฺปนฺโน, สปฺปติสฺโส สคารโว. \\ อภพฺโพ ปริหานาย, นิพฺพานสฺเสว สนฺติเกติ.\csromanspacer{}\csromanspacer{}ตติยํ.}}

\csromanheader{2}{4. มหาโมคฺคลฺลานสุตฺต}
\proseitem{34}{เอกํ สมยํ ภควา สาวตฺถิยํ วิหรติ เชตวเน อนาถปิณฺฑิกสฺส อาราเม.\csromanspacer{} อถ โข อายสฺมโต มหาโมคฺคลฺลานสฺส}

\vspace{\tipitakaparskip}
\csromancenter{ฉกฺกนิปาตปาฬิ รโหคตสฺส ปฏิสลฺลีนสฺส เอวํ เจตโส ปริวิตกฺโก อุทปาทิ “กตเมสานํ เทวานํ เอวํ ญาณํ โหติ ‘โสตาปนฺนา นาม\footnote{โสตาปนฺนามฺหา (สี), โสตาปนฺนามฺห (สฺยา, กํ, อิ)} อวินิปาตธมฺมา นิยตา สมฺโพธิปรายณา’ติ”.\csromanspacer{} เตน โข ปน สมเยน ติสฺโส นาม ภิกฺขุ อธุนากาลงฺกโต อญฺญตรํ พฺรหฺมโลกํ อุปปนฺโน โหติ.\csromanspacer{} ตตฺรปิ นํ เอวํ ชานนฺติ “ติสฺโส พฺรหฺมา มหิทฺธิโก มหานุภาโว”ติ.}
\prose{\csromanfolio{292}อถ โข อายสฺมา มหาโมคฺคลฺลาโน เสยฺยถาปิ นาม พลวา ปุริโส สมิญฺชิตํ\footnote{สมฺมิญฺชิตํ (สี, สฺยา, กํ, อิ)} วา พาหํ ปสาเรยฺย, ปสาริตํ วา พาหํ สมิญฺเชยฺย, เอวเมวํ เชตวเน อนฺตรหิโต ตสฺมิํ พฺรหฺมโลเก ปาตุรโหสิ.\csromanspacer{} อทฺทสา โข ติสฺโส พฺรหฺมา อายสฺมนฺตํ มหาโมคฺคลฺลานํ ทูรโตว อาคจฺฉนฺตํ, ทิสฺวาน อายสฺมนฺตํ มหาโมคฺคลฺลานํ เอตทโวจ “เอหิ โข มาริส โมคฺคลฺลาน, สฺวาคตํ\footnote{สาคตํ (สี)} มาริส โมคฺคลฺลาน, จิรสฺสํ โข มาริส โมคฺคลฺลาน อิมํ ปริยายมกาสิ ยทิทํ อิธาคมนาย, นิสีท มาริส โมคฺคลฺลาน, อิทมาสนํ ปญฺญตฺตนฺ”ติ.\csromanspacer{} นิสีทิ โข อายสฺมา มหาโมคฺคลฺลาโน ปญฺญตฺเต อาสเน.\csromanspacer{} ติสฺโสปิ โข พฺรหฺมา อายสฺมนฺตํ มหาโมคฺคลฺลานํ อภิวาเทตฺวา เอกมนฺตํ นิสีทิ, เอกมนฺตํ นิสินฺนํ โข ติสฺสํ พฺรหฺมานํ อายสฺมา มหาโมคฺคลฺลาโน เอตทโวจ\csromandash{}}

\prose{กตเมสานํ โข ติสฺส เทวานํ เอวํ ญาณํ โหติ “โสตาปนฺนา นาม อวินิปาตธมฺมา นิยตา สมฺโพธิปรายณา”ติ.\csromanspacer{} จตุมหาราชิกานํ โข มาริส โมคฺคลฺลาน เทวานํ เอวํ ญาณํ โหติ “โสตาปนฺนา นาม อวินิปาตธมฺมา นิยตา สมฺโพธิปรายณา”ติ.}

\prose{สพฺเพสญฺเญว นุ โข ติสฺส จาตุมหาราชิกานํ เทวานํ เอวํ ญาณํ โหติ “โสตาปนฺนา นาม อวินิปาตธมฺมา นิยตา สมฺโพธิปรายณา”ติ.\csromanspacer{} น โข มาริส โมคฺคลฺลาน สพฺเพสํ จตุมหาราชิกานํ เทวานํ เอวํ ญาณํ โหติ “โสตาปนฺนา นาม อวินิปาตธมฺมา นิยตา สมฺโพธิปรายณา”ติ.\csromanspacer{} เย โข เต มาริส โมคฺคลฺลาน จาตุมหาราชิกา เทวา พุทฺเธ อเวจฺจปฺปสาเทน อสมนฺนาคตา, ธมฺเม อเวจฺจปฺปสาเทน อสมนฺนาคตา, สํเฆ อเวจฺจปฺปสาเทน อสมนฺนาคตา, อริยกนฺเตหิ}

\vspace{\tipitakaparskip}
\csromancenter{เทวตาวคฺค สีเลหิ อสมนฺนาคตา, น เตสํ เทวานํ เอวํ ญาณํ โหติ “โสตาปนฺนา นาม อวินิปาตธมฺมา นิยตา สมฺโพธิปรายณา”ติ.\csromanspacer{} เย จ โข เต มาริส โมคฺคลฺลาน จาตุมหาราชิกา เทวา พุทฺเธ อเวจฺจปฺปสาเทน สมนฺนาคตา, ธมฺเม อเวจฺจปฺปสาเทน สมนฺนาคตา, สํเฆ อเวจฺจปฺปสาเทน สมนฺนาคตา, อริยกนฺเตหิ สีเลหิ สมนฺนาคตา, เตสํ เอวํ ญาณํ โหติ “โสตาปนฺนา นาม อวินิปาตธมฺมา นิยตา สมฺโพธิปรายณา”ติ.}
\prose{\csromanfolio{293}จาตุมหาราชิกานญฺเญว นุ โข ติสฺส เทวานํ เอวํ ญาณํ โหติ “โสตาปนฺนา นาม อวินิปาตธมฺมา นิยตา สมฺโพธิปรายณา”ติ.\csromanspacer{} อุทาหุ ตาวติํสานมฺปิ เทวานํ ฯเปฯ ยามานมฺปิ เทวานํ.\csromanspacer{} ตุสิตานมฺปิ เทวานํ.\csromanspacer{} นิมฺมานรตีนมฺปิ เทวานํ.\csromanspacer{} ปรนิมฺมิตวสวตฺตีนมฺปิ เทวานํ เอวํ ญาณํ โหติ “โสตาปนฺนา นาม อวินิปาตธมฺมา นิยตา สมฺโพธิปรายณา”ติ.\csromanspacer{} ปรนิมฺมิตวสวตฺตีนมฺปิ โข มาริส โมคฺคลฺลาน เทวานํ เอวํ ญาณํ โหติ “โสตาปนฺนา นาม อวินิปาตธมฺมา นิยตา สมฺโพธิปรายณา”ติ.}

\prose{สพฺเพสญฺเญว นุ โข ติสฺส ปรนิมฺมิตวสวตฺตีนํ เทวานํ เอวํ ญาณํ โหติ “โสตาปนฺนา นาม อวินิปาตธมฺมา นิยตา สมฺโพธิปรายณา”ติ.\csromanspacer{} น โข มาริส โมคฺคลฺลาน สพฺเพสํ ปรนิมฺมิตวสวตฺตีนํ เทวานํ เอวํ ญาณํ โหติ “โสตาปนฺนา นาม อวินิปาตธมฺมา นิยตา สมฺโพธิปรายณา”ติ.\csromanspacer{} เย โข เต มาริส โมคฺคลฺลาน ปรนิมฺมิตสวตฺตี เทวา พุทฺเธ อเวจฺจปฺปสาเทน อสมนฺนาคตา, ธมฺเม อเวจฺจปฺปสาเทน อสมนฺนาคตา, สํเฆ อเวจฺจปฺปสาเทน อสมนฺนาคตา, อริยกนฺเตหิ สีเลหิ อสมนฺนาคตา, น เตสํ เทวานํ เอวํ ญาณํ โหติ “โสตาปนฺนา นาม อวินิปาตธมฺมา นิยตา สมฺโพธิปรายณา”ติ.\csromanspacer{} เย จ โข เต มาริส โมคฺคลฺลาน ปรนิมฺมิตสวตฺตี เทวา พุทฺเธ อเวจฺจปฺปสาเทน สมนฺนาคตา, ธมฺเม อเวจฺจปฺปสาเทน สมนฺนาคตา, สํเฆ อเวจฺจปฺปสาเทน สมนฺนาคตา, อริยกนฺเตหิ สีเลหิ สมนฺนาคตา, เตสํ เอวํ ญาณํ โหติ “โสตาปนฺนา นาม อวินิปาตธมฺมา นิยตา สมฺโพธิปรายณา”ติ.}

\prose{อถ โข อายสฺมา มหาโมคฺคลฺลาโน ติสฺสสฺส พฺรหฺมุโน ภาสิตํ อภินนฺทิตฺวา อนุโมทิตฺวา เสยฺยถาปิ นาม พลวา ปุริโส สมิญฺชิตํ วา พาหํ ปสาเรยฺย, ปสาริตํ วา พาหํ สมิญฺเชยฺย, เอวเมวํ พฺรหฺมโลเก อนฺตรหิโต เชตวเน ปาตุรโหสีติ.\csromanspacer{}\csromanspacer{}จตุตฺถํ.}

\csromanmatikaanchor{mtk.155}
\csromantocmarkref{subsection}{5. วิชฺชาภาคิยสุตฺต}{mtk.155}
\bookmark[dest={mtk.155},level=2]{5. วิชฺชาภาคิยสุตฺต}
\csromanheader{2}{ฉกฺกนิปาตปาฬิ \textbf{5. วิชฺชาภาคิยสุตฺต}}
\proseitem{35}{\csromanfolio{294}ฉยิเม ภิกฺขเว ธมฺมา วิชฺชาภาคิยา.\csromanspacer{} กตเม ฉ? อนิจฺจสญฺญา, อนิจฺเจ ทุกฺขสญฺญา, ทุกฺเข อนตฺตสญฺญา, ปหานสญฺญา, วิราคสญฺญา, นิโรธสญฺญา.\csromanspacer{} อิเม โข ภิกฺขเว ฉ ธมฺมา วิชฺชาภาคิยาติ.\csromanspacer{}\csromanspacer{}ปญฺจมํ.}

\csromanmatikaanchor{mtk.156}
\csromantocmarkref{subsection}{6. วิวาทมูลสุตฺต}{mtk.156}
\bookmark[dest={mtk.156},level=2]{6. วิวาทมูลสุตฺต}
\csromanheader{2}{6. วิวาทมูลสุตฺต}
\proseitem{36}{\csromansymbolfootnote{*}{ที 3. 204; ม 3. 33; วิ 4. 212 ปิฏฺเฐสุปิ.}ฉยิมานิ ภิกฺขเว วิวาทมูลานิ.\csromanspacer{} กตมานิ ฉ? อิธ ภิกฺขเว ภิกฺขุ โกธโน โหติ อุปนาหี.\csromanspacer{} โย โส ภิกฺขเว ภิกฺขุ โกธโน โหติ อุปนาหี, โส สตฺถริปิ อคารโว วิหรติ อปฺปติสฺโส, ธมฺเมปิ อคารโว วิหรติ อปฺปติสฺโส, สํเฆปิ อคารโว วิหรติ อปฺปติสฺโส, สิกฺขายปิ น ปริปูรการี โหติ.\csromanspacer{} โย โส ภิกฺขเว ภิกฺขุ สตฺถริ อคารโว วิหรติ อปฺปติสฺโส, ธมฺเม อคารโว วิหรติ อปฺปติสฺโส, สํเฆ อคารโว วิหรติ อปฺปติสฺโส, สิกฺขาย น ปริปูรการี, โส สํเฆ วิวาทํ ชเนติ.\csromanspacer{} โย โหติ วิวาโท พหุชนาหิตาย พหุชนาสุขาย พหุโน ชนสฺส อนตฺถาย อหิตาย ทุกฺขาย เทวมนุสฺสานํ.\csromanspacer{} เอวรูปญฺเจ ตุเมฺห ภิกฺขเว วิวาทมูลํ อชฺฌตฺตํ วา พหิทฺธา วา สมนุปสฺเสยฺยาถ, ตตฺร ตุเมฺห ภิกฺขเว ตสฺเสว ปาปกสฺส วิวาทมูลสฺส ปหานาย วายเมยฺยาถ.\csromanspacer{} เอวรูปญฺเจ ตุเมฺห ภิกฺขเว วิวาทมูลํ อชฺฌตฺตํ วา พหิทฺธา วา น สมนุปสฺเสยฺยาถ, ตตฺร ตุเมฺห ภิกฺขเว ตสฺเสว ปาปกสฺส วิวาทมูลสฺส อายติํ อนวสฺสวาย ปฏิปชฺเชยฺยาถ.\csromanspacer{} เอวเมตสฺส ปาปกสฺส วิวาทมูลสฺส ปหานํ โหติ, เอวเมตสฺส ปาปกสฺส วิวาทมูลสฺส อายติํ อนวสฺสโว โหติ.}

\prose{ปุน จปรํ ภิกฺขเว ภิกฺขุ มกฺขี โหติ ปฬาสี ฯเปฯ.\csromanspacer{} อิสฺสุกี โหติ มจฺฉรี.\csromanspacer{} สโฐ โหติ มายาวี.\csromanspacer{} ปาปิจฺโฉ โหติ มิจฺฉาทิฏฺฐิ.\csromanspacer{} สนฺทิฏฺฐิปรามาสี โหติ อาธานคฺคาหี ทุปฺปฏินิสฺสคฺคี.\csromanspacer{} โย โส ภิกฺขเว ภิกฺขุ สนฺทิฏฺฐิปรามาสี โหติ อาธานคฺคาหี ทุปฺปฏินิสฺสคฺคี, โส สตฺถริปิ อคารโว วิหรติ อปฺปติสฺโส, ธมฺเมปิ อคารโว วิหรติ อปฺปติสฺโส, สํเฆปิ อคารโว วิหรติ อปฺปติสฺโส, สิกฺขายปิ น ปริปูรการี โหติ.\csromanspacer{} โย โส ภิกฺขเว ภิกฺขุ สตฺถริ อคารโว วิหรติ อปฺปติสฺโส.\csromanspacer{} ธมฺเม.\csromanspacer{} สํเฆ อคารโว วิหรติ อปฺปติสฺโส, สิกฺขาย น ปริปูรการี, โส สํเฆ วิวาทํ ชเนติ.\csromanspacer{} โย โหติ วิวาโท}

\vspace{\tipitakaparskip}
\csromancenter{เทวตาวคฺค พหุชนาหิตาย พหุชนาสุขาย พหุโน ชนสฺส อนตฺถาย อหิตาย ทุกฺขาย เทวมนุสฺสานํ.\csromanspacer{} เอวรูปญฺเจ ตุเมฺห ภิกฺขเว วิวาทมูลํ อชฺฌตฺตํ วา พหิทฺธา วา สมนุปสฺเสยฺยาถ, ตตฺร ตุเมฺห ภิกฺขเว ตสฺเสว ปาปกสฺส วิวาทมูลสฺส ปหานาย วายเมยฺยาถ.\csromanspacer{} เอวรูปญฺเจ ตุเมฺห ภิกฺขเว วิวาทมูลํ อชฺฌตฺตํ วา พหิทฺธา วา น สมนุปสฺเสยฺยาถ, ตตฺร ตุเมฺห ภิกฺขเว ตสฺเสว ปาปกสฺส วิวาทมูลสฺส อายติํ อนวสฺสวาย ปฏิปชฺเชยฺยาถ.\csromanspacer{} เอวเมตสฺส ปาปกสฺส วิวาทมูลสฺส ปหานํ โหติ, เอวเมตสฺส ปาปกสฺส วิวาทมูลสฺส อายติํ อนวสฺสโว โหติ.\csromanspacer{} อิมานิ โข ภิกฺขเว ฉ วิวาทมูลานีติ.\csromanspacer{}\csromanspacer{}ฉฏฺฐํ.}
\csromanmatikaanchor{mtk.157}
\csromantocmarkref{subsection}{7. ฉฬงฺคทานสุตฺต}{mtk.157}
\bookmark[dest={mtk.157},level=2]{7. ฉฬงฺคทานสุตฺต}
\csromanheader{2}{7. ฉฬงฺคทานสุตฺต}
\proseitem{37}{\csromanfolio{295}เอกํ สมยํ ภควา สาวตฺถิยํ วิหรติ เชตวเน อนาถปิณฺฑิกสฺส อาราเม.\csromanspacer{} เตน โข ปน สมเยน เวฬุกณฺฑกี\footnote{เวฬุกณฺฑกิยา (อุปริ 444; อํ 1. 88; สํ 1. 432 ปิฏฺเฐปิ.)} นนฺทมาตา อุปาสิกา สาริปุตฺตโมคฺคลฺลานปฺปมุเข ภิกฺขุสํเฆ ฉฬงฺคสมนฺนาคตํ ทกฺขิณํ ปติฏฺฐาเปติ.\csromanspacer{} อทฺทสา โข ภควา ทิพฺเพน จกฺขุนา วิสุทฺเธน อติกฺกนฺตมานุสเกน เวฬุกณฺฑกิํ นนฺทมาตรํ อุปาสิกํ สาริปุตฺตโมคฺคลฺลานปฺปมุเข ภิกฺขุสํเฆ ฉฬงฺคสมนฺนาคตํ ทกฺขิณํ ปติฏฺฐาเปนฺติํ, ทิสฺวา ภิกฺขู อามนฺเตสิ “เอสา ภิกฺขเว เวฬุกณฺฑกี นนฺทมาตา อุปาสิกา สาริปุตฺต- โมคฺคลฺลานปฺปมุเข ภิกฺขุสํเฆ ฉฬงฺคสมนฺนาคตํ ทกฺขิณํ ปติฏฺฐาเปติ.}

\prose{กถญฺจ ภิกฺขเว ฉฬงฺคสมนฺนาคตา ทกฺขิณา โหติ? อิธ ภิกฺขเว ทายกสฺส ตีณงฺคานิ โหนฺติ, ปฏิคฺคาหกานํ ตีณงฺคานิ.\csromanspacer{} กตมานิ ทายกสฺส ตีณงฺคานิ? อิธ ภิกฺขเว ทายโก ปุพฺเพว ทานา สุมโน โหติ, ททํ จิตฺตํ ปสาเทติ, ทตฺวา อตฺตมโน โหติ.\csromanspacer{} อิมานิ ทายกสฺส ตีณงฺคานิ.}

\prose{กตมานิ ปฏิคฺคาหกานํ ตีณงฺคานิ? อิธ ภิกฺขเว ปฏิคฺคาหกา วีตราคา วา โหนฺติ ราควินยาย วา ปฏิปนฺนา, วีตโทสา วา โหนฺติ โทสวินยาย วา ปฏิปนฺนา, วีตโมหา วา โหนฺติ โมหวินยาย วา ปฏิปนฺนา.\csromanspacer{} อิมานิ ปฏิคฺคาหกานํ ตีณงฺคานิ.\csromanspacer{} อิติ ทายกสฺส ตีณงฺคานิ, ปฏิคฺคาหกานํ ตีณงฺคานิ.\csromanspacer{} เอวํ โข ภิกฺขเว ฉฬงฺคสมนฺนาคตา ทกฺขิณา โหติ.}

\gambhira{ฉกฺกนิปาตปาฬิ}
\prose{\csromanfolio{296}เอวํ ฉฬงฺคสมนฺนาคตาย ภิกฺขเว ทกฺขิณาย น สุกรํ ปุญฺญสฺส ปมาณํ คเหตุํ “เอตฺตโก ปุญฺญาภิสนฺโท กุสลาภิสนฺโท สุขสฺสาหาโร โสวคฺคิโก สุขวิปาโก สคฺคสํวตฺตนิโก อิฏฺฐาย กนฺตาย มนาปาย หิตาย สุขาย สํวตฺตตี”ติ.\csromanspacer{} อถ โข อสงฺเขฺยยฺโย\footnote{อสงฺเขยฺโย (สี, สฺยา, กํ, อิ)} อปฺปเมยฺโย มหาปุญฺญกฺขนฺโธเตฺวว สงฺขํ คจฺฉติ.}

\prose{เสยฺยถาปิ ภิกฺขเว มหาสมุทฺเท น สุกรํ อุทกสฺส ปมาณํ คเหตุํ “เอตฺตกานิ อุทกาฬฺหกานี”ติ วา “เอตฺตกานิ อุทกาฬฺหกสตานี”ติ วา “เอตฺตกานิ อุทกาฬฺหกสหสฺสานี”ติ วา “เอตฺตกานิ อุทกาฬฺหกสตสหสฺสานี”ติ วา, อถ โข อสงฺเขฺยยฺโย อปฺปเมยฺโย มหา-อุทกกฺขนฺโธเตฺวว สงฺขํ คจฺฉติ.\csromanspacer{} เอวเมวํ โข ภิกฺขเว เอวํ ฉฬงฺคสมนฺนาคตาย ทกฺขิณาย น สุกรํ ปุญฺญสฺส ปมาณํ คเหตุํ “เอตฺตโก ปุญฺญาภิสนฺโท กุสลาภิสนฺโท สุขสฺสาหาโร โสวคฺคิโก สุขวิปาโก สคฺคสํวตฺตนิโก อิฏฺฐาย กนฺตาย มนาปาย หิตาย สุขาย สํวตฺตตี”ติ.\csromanspacer{} อถ โข อสงฺเขฺยยฺโย อปฺปเมยฺโย มหาปุญฺญกฺขนฺโธเตฺวว สงฺขํคจฺฉตีติ.}

\prose{\csromansymbolfootnote{*}{ขุ 2. 159 เปตวตฺถุมฺหิปิ.}ปุพฺเพว ทานา สุมโน, ททํ จิตฺตํ ปสาทเย.}

\csromangathameasure{ทตฺวา อตฺตมโน โหติ, เอสา ยญฺญสฺส สมฺปทา. \\ วีตราคา วีตโทสา, วีตโมหา อนาสวา. \\ เขตฺตํ ยญฺญสฺส สมฺปนฺนํ, สญฺญตา พฺรหฺมจารโย. \\ สยํ อาจมยิตฺวาน, ทตฺวา สเกหิ ปาณิภิ. \\ อตฺตโน ปรโต เจโส, ยญฺโญ โหติ มหปฺผโล.}
\csromangathagroup{\csromangathastanza{ทตฺวา อตฺตมโน โหติ, เอสา ยญฺญสฺส\footnote{ปุญฺญสฺส (ก)} สมฺปทา. \\ วีตราคา\footnote{ปุญฺญสฺส (ก)} วีตโทสา, วีตโมหา อนาสวา.}\csromangathastanza{เขตฺตํ ยญฺญสฺส สมฺปนฺนํ, สญฺญตา พฺรหฺมจารโย\footnote{พฺรหฺมจาริโน (สฺยา, กํ)}. \\ สยํ อาจมยิตฺวาน, ทตฺวา สเกหิ ปาณิภิ. \\ อตฺตโน ปรโต เจโส, ยญฺโญ โหติ มหปฺผโล.}}

\prose{\csromansymbolfootnote{+}{อํ 1. 353 ปิฏฺเฐปิ.}เอวํ ยชิตฺวา เมธาวี, สทฺโธ มุตฺเตน เจตสา.}

\csromangathameasure{อพฺยาพชฺฌํ สุขํ โลกํ, ปณฺฑิโต อุปปชฺชตีติ.สตฺตมํ.}
\csromangathagroup{\csromangathastanza{อพฺยาพชฺฌํ สุขํ โลกํ, ปณฺฑิโต อุปปชฺชตีติ.\csromanspacer{}\csromanspacer{}สตฺตมํ.}}

\csromanmatikaanchor{mtk.158}
\csromantocmarkref{subsection}{8. อตฺตการีสุตฺต}{mtk.158}
\bookmark[dest={mtk.158},level=2]{8. อตฺตการีสุตฺต}
\csromanheader{2}{8. อตฺตการีสุตฺต}
\proseitem{38}{อถ โข อญฺญตโร พฺราหฺมโณ เยน ภควา เตนุปสงฺกมิ, อุปสงฺกมิตฺวา ภควตา สทฺธิํ สมฺโมทิ, สมฺโมทนียํ กถํ สารณียํ วีติสาเรตฺวา เอกมนฺตํ นิสีทิ, เอกมนฺตํ นิสินฺโน โข โส พฺราหฺมโณ}

\vspace{\tipitakaparskip}
\csromancenter{เทวตาวคฺค ภควนฺตํ เอตทโวจ “อหํ หิ โภ โคตม เอวํวาที เอวํทิฏฺฐิ ‘นตฺถิ อตฺตกาโร, นตฺถิ ปรกาโร’ติ”.\csromanspacer{} มาหํ พฺรหฺมณ เอวํวาทิํ เอวํทิฏฺฐิํ อทฺทสํ วา อสฺโสสิํ วา, กถํ หิ นาม สยํ อภิกฺกมนฺโต สยํ ปฏิกฺกมนฺโต เอวํ วกฺขติ “นตฺถิ อตฺตกาโร, นตฺถิ ปรกาโร”ติ.}
\prose{\csromanfolio{297}ตํ กิํ มญฺญสิ พฺราหฺมณ, อตฺถิ อารพฺภธาตูติ.\csromanspacer{} เอวํ โภ.\csromanspacer{} อารพฺภธาตุยา สติ อารพฺภวนฺโต สตฺตา ปญฺญายนฺตีติ.\csromanspacer{} เอวํ โภ.\csromanspacer{} ยํ โข พฺราหฺมณ อารพฺภธาตุยา สติ อารพฺภวนฺโต สตฺตา ปญฺญายนฺติ.\csromanspacer{} อยํ สตฺตานํ อตฺตกาโร อยํ ปรกาโร.}

\prose{ตํ กิํ มญฺญสิ พฺราหฺมณ, อตฺถิ นิกฺกมธาตุ ฯเปฯ อตฺถิ ปรกฺกมธาตุ.\csromanspacer{} อตฺถิ ถามธาตุ.\csromanspacer{} อตฺถิ ฐิติธาตุ.\csromanspacer{} อตฺถิ อุปกฺกมธาตูติ.\csromanspacer{} เอวํ โภ.\csromanspacer{} อุปกฺกมธาตุยา สติ อุปกฺกมวนฺโต สตฺตา ปญฺญายนฺตีติ.\csromanspacer{} เอวํ โภ.\csromanspacer{} ยํ โข พฺราหฺมณ อุปกฺกมธาตุยา สติ อุปกฺกมวนฺโต สตฺตา ปญฺญายนฺติ.\csromanspacer{} อยํ สตฺตานํ อตฺตกาโร อยํ ปรกาโร.}

\prose{มาหํ พฺราหฺมณ\footnote{ตํ กิํ มญฺญสิ พฺราหฺมณ, มาหํ (ก)} เอวํวาทิํ เอวํทิฏฺฐิํ อทฺทสํ วา อสฺโสสิํ วา, กถํ หิ นาม สยํ อภิกฺกมนฺโต สยํ ปฏิกฺกมนฺโต เอวํ วกฺขติ “นตฺถิ อตฺตกาโร นตฺถิ ปรกาโร”ติ.}

\prose{อภิกฺกนฺตํ โภ โคตม ฯเปฯ อชฺชตคฺเค ปาณุเปตํ สรณํ คตนฺติ.\csromanspacer{}\csromanspacer{}อฏฺฐมํ.}

\csromanmatikaanchor{mtk.159}
\csromantocmarkref{subsection}{9. นิทานสุตฺต}{mtk.159}
\bookmark[dest={mtk.159},level=2]{9. นิทานสุตฺต}
\csromanheader{2}{9. นิทานสุตฺต}
\proseitem{39}{ตีณิมานิ ภิกฺขเว นิทานานิ กมฺมานํ สมุทยาย.\csromanspacer{} กตมานิ ตีณิ? โลโภ นิทานํ กมฺมานํ สมุทยาย, โทโส นิทานํ กมฺมานํ สมุทยาย, โมโห นิทานํ กมฺมานํ สมุทยาย.\csromanspacer{} น ภิกฺขเว โลภา อโลโภ สมุเทติ, อถ โข ภิกฺขเว โลภา โลโภว สมุเทติ.\csromanspacer{} น ภิกฺขเว โทสา อโทโส สมุเทติ, อถ โข ภิกฺขเว โทสา โทโสว สมุเทติ.\csromanspacer{} น ภิกฺขเว โมหา อโมโห สมุเทติ, อถ โข ภิกฺขเว โมหา โมโหว สมุเทติ.\csromanspacer{} น ภิกฺขเว โลภเชน กมฺเมน โทสเชน กมฺเมน โมหเชน กมฺเมน เทวา ปญฺญายนฺติ, มนุสฺสา ปญฺญายนฺติ, ยา วา ปนญฺญาปิ กาจิ สุคติโย, อถ โข ภิกฺขเว โลภเชน กมฺเมน โทสเชน กมฺเมน โมหเชน}

\vspace{\tipitakaparskip}
\csromancenter{ฉกฺกนิปาตปาฬิ กมฺเมน นิรโย ปญฺญายติ, ติรจฺฉานโยนิ ปญฺญายติ, เปตฺติวิสโย ปญฺญายติ, ยา วา ปนญฺญาปิ กาจิ ทุคฺคติโย.\csromanspacer{} อิมานิ โข ภิกฺขเว ตีณิ นิทานานิ กมฺมานํ สมุทยาย.}
\prose{\csromanfolio{298}ตีณิมานิ ภิกฺขเว นิทานานิ กมฺมานํ สมุทยาย.\csromanspacer{} กตมานิ ตีณิ? อโลโภ นิทานํ กมฺมานํ สมุทยาย, อโทโส นิทานํ กมฺมานํ สมุทยาย, อโมโห นิทานํ กมฺมานํ สมุทยาย.\csromanspacer{} น ภิกฺขเว อโลภา โลโภ สมุเทติ, อถ โข ภิกฺขเว อโลภา อโลโภว สมุเทติ.\csromanspacer{} น ภิกฺขเว อโทสา โทโส สมุเทติ, อถ โข ภิกฺขเว อโทสา อโทโสว สมุเทติ.\csromanspacer{} น ภิกฺขเว อโมหา โมโห สมุเทติ.\csromanspacer{} อถ โข ภิกฺขเว อโมหา อโมโหว สมุเทติ.\csromanspacer{} น ภิกฺขเว อโลภเชน กมฺเมน อโทสเชน กมฺเมน อโมหเชน กมฺเมน นิรโย ปญฺญายติ, ติรจฺฉานโยนิ ปญฺญายติ, เปตฺติวิสโย ปญฺญายติ, ยา วา ปนญฺญาปิ กาจิ ทุคฺคติโย.\csromanspacer{} อถ โข ภิกฺขเว อโลภเชน กมฺเมน อโทสเชน กมฺเมน อโมหเชน กมฺเมน เทวา ปญฺญายนฺติ, มนุสฺสา ปญฺญายนฺติ, ยา วา ปนญฺญาปิ กาจิ สุคติโย.\csromanspacer{} อิมานิ โข ภิกฺขเว ตีณิ นิทานานิ กมฺมานํ สมุทยายาติ.\csromanspacer{}\csromanspacer{}นวมํ.}

\csromanmatikaanchor{mtk.160}
\csromantocmarkref{subsection}{10. กิมิลสุตฺต}{mtk.160}
\bookmark[dest={mtk.160},level=2]{10. กิมิลสุตฺต}
\csromanheader{2}{10. กิมิลสุตฺต}
\proseitem{40}{เอวํ เม สุตํ\csromandash{}เอกํ สมยํ ภควา กิมิลายํ วิหรติ นิจุลวเน.\csromanspacer{} อถ โข อายสฺมา กิมิโล เยน ภควา เตนุปสงฺกมิ, อุปสงฺกมิตฺวา ภควนฺตํ อภิวาเทตฺวา เอกมนฺตํ นิสีทิ, เอกมนฺตํ นิสินฺโน โข อายสฺมา กิมิโล ภควนฺตํ เอตทโวจ “โก นุ โข ภนฺเต เหตุ โก ปจฺจโย, เยน ตถาคเตปรินิพฺพุเต สทฺธมฺโม นจิรฏฺฐิติโก โหตี”ติ.\csromanspacer{} อิธ กิมิล ตถาคเต ปรินิพฺพุเต ภิกฺขู ภิกฺขุนิโย อุปาสกา อุปาสิกาโย สตฺถริ อคารวา วิหรนฺติ อปฺปติสฺสา, ธมฺเม อคารวา วิหรนฺติ อปฺปติสฺสา, สํเฆ อคารวา วิหรนฺติ อปฺปติสฺสา, สิกฺขาย อคารวา วิหรนฺติ อปฺปติสฺสา, อปฺปมาเท อคารวา วิหรนฺติ อปฺปติสฺสา, ปฏิสนฺถาเร\footnote{ปฏิสนฺธาเร (ก), เหฏฺฐา 216 ปิฏฺเฐ, อุปริ 460 ปิฏฺเฐสุปิ.} อคารวา วิหรนฺติ อปฺปติสฺสา.\csromanspacer{} อยํ โข กิมิล เหตุ อยํ ปจฺจโย, เยน ตถาคเต ปรินิพฺพุเต สทฺธมฺโมน จิรฏฺฐิติโก โหติ.}

\chapterheadpageread{4. เทวตาวคฺค}
\csromanmatikaanchor{mtk.161}
\csromantocmarkref{subsection}{11. ทารุกฺขนฺทสุตฺต}{mtk.161}
\bookmark[dest={mtk.161},level=2]{11. ทารุกฺขนฺทสุตฺต}
\prose{\csromanfolio{299}โก ปน ภนฺเต เหตุ โก ปจฺจโย, เยน ตถาคเต ปรินิพฺพุเต สทฺธมฺโม จิรฏฺฐิติโก โหตีติ.\csromanspacer{} อิธ กิมิล ตถาคเต ปรินิพฺพุเต ภิกฺขู ภิกฺขุนิโย อุปาสกา อุปาสิกาโย สตฺถริ สคารวา วิหรนฺติ สปฺปติสฺสา, ธมฺเม สคารวา วิหรนฺติ สปฺปติสฺสา, สํเฆ สคารวา วิหรนฺติ สปฺปติสฺสา, สิกฺขาย สคารวา วิหรนฺติ สปฺปติสฺสา, อปฺปมาเท สคารวา วิหรนฺติ สปฺปติสฺสา, ปฏิสนฺถาเร สคารวา วิหรนฺติ สปฺปติสฺสา.\csromanspacer{} อยํ โข กิมิล เหตุ อยํ ปจฺจโย, เยน ตถาคเต ปรินิพฺพุเต สทฺธมฺโม จิรฏฺฐิติโก โหตีติ.\csromanspacer{}\csromanspacer{}ทสมํ.}

\csromanheader{2}{11. ทารุกฺขนฺธสุตฺต}
\proseitem{41}{เอวํ เม สุตํ\csromandash{}เอกํ สมยํ อายสฺมา สาริปุตฺโต ราชคเห วิหรติ คิชฺฌกูเฏ ปพฺพเต.\csromanspacer{} อถ โข อายสฺมา สาริปุตฺโต ปุพฺพณฺหสมยํ นิวาเสตฺวา ปตฺตจีวรมาทาย สมฺพหุเลหิ ภิกฺขูหิ สทฺธิํ คิชฺฌกูฏา ปพฺพตา โอโรหนฺโต อทฺทส อญฺญตรสฺมิํ ปเทเส มหนฺตํ ทารุกฺขนฺธํ, ทิสฺวา ภิกฺขู อามนฺเตสิ “ปสฺสถ โน อาวุโส ตุเมฺห อมุํ มหนฺตํ ทารุกฺขนฺธนฺ”ติ.\csromanspacer{} เอวมาวุโสติ.}

\prose{อากงฺขมาโน อาวุโส ภิกฺขุ อิทฺธิมา เจโตวสิปฺปตฺโต อมุํ ทารุกฺขนฺธํ ปถวีเตฺวว อธิมุจฺเจยฺย.\csromanspacer{} ตํ กิสฺส เหตุ? อตฺถิ อาวุโส อมุมฺหิ ทารุกฺขนฺเธ ปถวีธาตุ, ยํ นิสฺสาย ภิกฺขุ อิทฺธิมา เจโตวสิปฺปตฺโต อมุํ ทารุกฺขนฺธํ ปถวีเตฺวว อธิมุจฺเจยฺย.\csromanspacer{} อากงฺขมาโน อาวุโส ภิกฺขุ อิทฺธิมา เจโตวสิปฺปตฺโต อมุํ ทารุกฺขนฺธํ อาโปเตฺวว อธิมุจฺเจยฺย ฯเปฯ เตโชเตฺวว อธิมุจฺเจยฺย.\csromanspacer{} วาโยเตฺวว อธิมุจฺเจยฺย.\csromanspacer{} สุภนฺเตฺวว อธิมุจฺเจยฺย.\csromanspacer{} อสุภนฺเตฺวว อธิมุจฺเจยฺย.\csromanspacer{} ตํ กิสฺส เหตุ? อตฺถิ อาวุโส อมุมฺหิ ทารุกฺขนฺเธ อสุภธาตุ, ยํ นิสฺสาย ภิกฺขุ อิทฺธิมา เจโตวสิปฺปตฺโต อมุํ ทารุกฺขนฺธํ อสุภนฺเตฺวว อธิมุจฺเจยฺยาติ.\csromanspacer{}\csromanspacer{}เอกาทสมํ.}

\csromanmatikaanchor{mtk.162}
\csromantocmarkref{subsection}{12. นาคิตสุตฺต}{mtk.162}
\bookmark[dest={mtk.162},level=2]{12. นาคิตสุตฺต}
\csromanheader{2}{12. นาคิตสุตฺต}
\proseitem{42}{เอวํ เม สุตํ\csromandash{}เอกํ สมยํ ภควา โกสเลสุ จาริกํ จรมาโน มหตา ภิกฺขุสํเฆน สทฺธิํ เยน อิจฺฉานงฺคลํ นาม โกสลานํ พฺราหฺมณคาโม, ตทวสริ.\csromanspacer{} ตตฺร สุทํ ภควา อิจฺฉานงฺคเล}

\vspace{\tipitakaparskip}
\csromancenter{ฉกฺกนิปาตปาฬิ วิหรติ อิจฺฉานงฺคลวนสณฺเฑ.\csromanspacer{} อสฺโสสุํ โข อิจฺฉานงฺคลกา พฺราหฺมณคหปติกา “สมโณ ขลุ โภ โคตโม สกฺยปุตฺโต สกฺยกุลา ปพฺพชิโต อิจฺฉานงฺคลํ อนุปฺปตฺโต อิจฺฉานงฺคเล วิหรติ อิจฺฉานงฺคลวนสณฺเฑ.\csromanspacer{} ตํ โข ปน ภวนฺตํ โคตมํ เอวํ กลฺยาโณ กิตฺติสทฺโท อพฺภุคฺคโต ‘อิติปิ โส ภควา อรหํ สมฺมาสมฺพุทฺโธ วิชฺชาจรณสมฺปนฺโน ฯเปฯ พุทฺโธ ภควา’ติ.\csromanspacer{} โส อิมํ โลกํ สเทวกํ ฯเปฯ อรหตํ ทสฺสนํ โหตี”ติ.\csromanspacer{} อถ โข อิจฺฉานงฺคลกา พฺราหฺมณคหปติกา ตสฺสา รตฺติยา อจฺจเยน ปหูตํ ขาทนียํ โภชนียํ อาทาย เยน อิจฺฉานงฺคลวนสณฺโฑ เตนุปสงฺกมิํสุ, อุปสงฺกมิตฺวา พหิทฺวารโกฏฺฐเก อฏฺฐํสุ อุจฺจาสทฺทา มหาสทฺทา.}
\prose{\csromanfolio{300}เตน โข ปน สมเยน อายสฺมา นาคิโต ภควโต อุปฏฺฐาโก โหติ.\csromanspacer{} อถ โข ภควา อายสฺมนฺตํ นาคิตํ อามนฺเตสิ “เก ปน เต นาคิต อุจฺจาสทฺทา มหาสทฺทา เกวฏฺฏา มญฺเญ มจฺฉวิโลเป”ติ.\csromanspacer{} เอเต ภนฺเต อิจฺฉานงฺคลกา พฺราหฺมณคหปติกา ปหูตํ ขาทนียํ โภชนียํ อาทาย พหิทฺวารโกฏฺฐเก ฐิตา ภควนฺตํเยว อุทฺทิสฺส ภิกฺขุสํฆญฺจาติ.\csromanspacer{} มาหํ นาคิต ยเสน สมาคมํ, มา จ มยา ยโส.\csromanspacer{} โย โข นาคิต นยิมสฺส เนกฺขมฺมสุขสฺส ปวิเวกสุขสฺส อุปสมสุขสฺส สมฺโพธสุขสฺส นิกามลาภี อสฺส อกิจฺฉลาภี อกสิรลาภี, ยสฺสาหํ เนกฺขมฺมสุขสฺส ปวิเวกสุขสฺส อุปสมสุขสฺส สมฺโพธสุขสฺส นิกามลาภี อกิจฺฉลาภี อกสิรลาภี, โส ตํ มีฬฺหสุขํ มิทฺธสุขํ ลาภสกฺการสิโลกสุขํ สาทิเยยฺยาติ.}

\prose{อธิวาเสตุ ทานิ ภนฺเต ภควา, อธิวาเสตุ สุคโต, อธิวาสนกาโล ทานิ ภนฺเต ภควโต, เยน เยเนว ทานิ ภนฺเต ภควา คมิสฺสติ, ตนฺนินฺนาว ภวิสฺสนฺติ พฺราหฺมณคหปติกา เนคมา เจว ชานปทา จ.\csromanspacer{} เสยฺยถาปิ ภนฺเต ถุลฺลผุสิตเก เทเว วสนฺเต ยถานินฺนํ อุทกานิ ปวตฺตนฺติ.\csromanspacer{} เอวเมวํ โข ภนฺเต เยน เยเนว ทานิ ภควา คมิสฺสติ, ตนฺนินฺนาว ภวิสฺสนฺติ พฺราหฺมณคหปติกา เนคมา เจว ชานปทา จ.\csromanspacer{} ตํ กิสฺส เหตุ? ตถา หิ ภนฺเต ภควโต สีลปญฺญาณนฺติ. มาหํ นาคิต ยเสน สมาคมํ, มา จ มยา ยโส.\csromanspacer{} โย โข นาคิต นยิมสฺส เนกฺขมฺมสุขสฺส ปวิเวกสุขสฺส อุปสมสุขสฺส}

\vspace{\tipitakaparskip}
\csromancenter{เทวตาวคฺค สมฺโพธสุขสฺส นิกามลาภี อสฺส อกิจฺฉลาภี อกสิรลาภี, ยสฺสาหํ เนกฺขมฺมสุขสฺส ปวิเวกสุขสฺส อุปสมสุขสฺส สมฺโพธสุขสฺส นิกามลาภี อกิจฺฉลาภี อกสิรลาภี, โส ตํ มีฬฺหสุขํ มิทฺธสุขํ ลาภสกฺการสิโลกสุขํ สาทิเยยฺย.}
\prose{\csromanfolio{301}อิธาหํ นาคิต ภิกฺขุํ ปสฺสามิ คามนฺตวิหาริํ สมาหิตํ นิสินฺนํ, ตสฺส มยฺหํ นาคิต เอวํ โหติ “อิทานิมํ อายสฺมนฺตํ1 อารามิโก วา อุปฏฺฐหิสฺสติ สมณุทฺเทโส วา ตํ ตมฺหา1 สมาธิมฺหา จาเวสฺสตี”ติ เตนาหํ นาคิต ตสฺส ภิกฺขุโน น อตฺตมโน โหมิ คามนฺตวิหาเรน.\csromanspacer{} อิธ ปนาหํ นาคิต ภิกฺขุํ ปสฺสามิ อารญฺญิกํ อรญฺเญ ปจลายมานํ นิสินฺนํ, ตสฺส มยฺหํ นาคิต เอวํ โหติ “อิทานิ อยมายสฺมา อิมํ นิทฺทากิลมถํ ปฏิวิโนเทตฺวา อรญฺญสญฺญํเยว มนสิ กริสฺสติ เอกตฺตนฺ”ติ.\csromanspacer{} เตนาหํ นาคิต ตสฺส ภิกฺขุโน อตฺตมโน โหมิ อรญฺญวิหาเรน. อิธ ปนาหํ นาคิต ภิกฺขุํ ปสฺสามิ อารญฺญิกํ อรญฺเญ อสมาหิตํ นิสินฺนํ, ตสฺส มยฺหํ นาคิต เอวํ โหติ “อิทานิ อยมายสฺมา อสมาหิตํ วา จิตฺตํ สมาทหิสฺสติ, สมาหิตํ วา จิตฺตํ อนุรกฺขิสฺสตี”ติ.\csromanspacer{} เตนาหํ นาคิต ตสฺส ภิกฺขุโน อตฺตมโน โหมิ อรญฺญวิหาเรน. อิธ ปนาหํ นาคิต ภิกฺขุํ ปสฺสามิ อารญฺญิกํ อรญฺเญ สมาหิตํ นิสินฺนํ, ตสฺส มยฺหํ นาคิต เอวํ โหติ “อิทานิ อยมายสฺมา อวิมุตฺตํ วา จิตฺตํ วิโมเจสฺสติ, วิมุตฺตํ วา จิตฺตํ อนุรกฺขิสฺสตี”ติ.\csromanspacer{} เตนาหํ นาคิต ตสฺส ภิกฺขุโน อตฺตมโน โหมิ อรญฺญวิหาเรน. อิธ ปนาหํ นาคิต ภิกฺขุํ ปสฺสามิ คามนฺตวิหาริํ ลาภิํ จีวรปิณฺฑปาตเสนาสนคิลานปฺปจฺจยเภสชฺชปริกฺขารานํ.\csromanspacer{} โส ตํ ลาภสกฺการสิโลกํ นิกามยมาโน ริญฺจติ ปฏิสลฺลานํ, ริญฺจติ อรญฺญวนปตฺถานิ ปนฺตานิ เสนาสนานิ, คามนิคมราชธานิํ โอสริตฺวา วาสํ กปฺเปติ.\csromanspacer{} เตนาหํ นาคิต ตสฺส ภิกฺขุโน น อตฺตมโน โหมิ คามนฺตวิหาเรน.}

\gambhira{ฉกฺกนิปาตปาฬิ}
\prose{\csromanfolio{302}อิธ ปนาหํ นาคิต ภิกฺขุํ ปสฺสามิ อารญฺญิกํ ลาภิํ จีวร- ปิณฺฑปาตเสนาสนคิลานปฺปจฺจยเภสชฺชปริกฺขารานํ, โส ตํ ลาภสกฺการสิโลกํ ปฏิปณาเมตฺวา น ริญฺจติ ปฏิสลฺลานํ, น ริญฺจติ อรญฺญวนปตฺถานิ ปนฺตานิ เสนาสนานิ.\csromanspacer{} เตนาหํ นาคิต ตสฺส ภิกฺขุโน อตฺตมโน โหมิ อรญฺญวิหาเรน.\csromanspacer{} ยสฺมาหํ นาคิต สมเย อทฺธานมคฺคปฺปฏิปนฺโน น กญฺจิ\footnote{น กิญฺจิ (สี, อิ, ก)} ปสฺสามิ ปุรโต วา ปจฺฉโต วา, ผาสุ เม นาคิต ตสฺมิํ สมเย โหติ อนฺตมโส อุจฺจารปสฺสาวกมฺมายาติ.\csromanspacer{}\csromanspacer{}ทฺวาทสมํ.\csromanspacer{} เทวตาวคฺโค\footnote{เสกฺขปริหานิยวคฺโค (สฺยา)} จตุตฺโถ.}

\titlehead{\textbf{ตสฺสุทฺทานํ}}
\csromangathameasure{เสขา เทฺว อปริหานิ, โมคฺคลฺลานวิชฺชาภาคิยา. \\ วิวาททานตฺตการี นิทานํ, กิมิลทารุกฺขนฺเธน นาคิโตติ.}
\csromangathagroup{\csromangathastanza{เสขา เทฺว อปริหานิ, โมคฺคลฺลานวิชฺชาภาคิยา. \\ วิวาททานตฺตการี นิทานํ, กิมิลทารุกฺขนฺเธน นาคิโตติ.}}

\csromanmatikaanchor{mtk.163}
\csromantocmarknopage{section}{5. ธมฺมิกวคฺค}{mtk.163}
\bookmark[dest={mtk.163},level=1]{5. ธมฺมิกวคฺค}
\csromanheader{1}{5. ธมฺมิกวคฺค}
\csromanmatikaanchor{mtk.164}
\csromantocmarkref{subsection}{1. นาคสุตฺต}{mtk.164}
\bookmark[dest={mtk.164},level=2]{1. นาคสุตฺต}
\csromanheader{2}{1. นาคสุตฺต}
\proseitem{43}{เอกํ สมยํ ภควา สาวตฺถิยํ วิหรติ เชตวเน อนาถปิณฺฑิกสฺส อาราเม.\csromanspacer{} อถ โข ภควา ปุพฺพณฺหสมยํ นิวาเสตฺวา ปตฺตจีวรมาทาย สาวตฺถิยํ ปิณฺฑาย ปาวิสิ, สาวตฺถิยํ ปิณฺฑาย จริตฺวา ปจฺฉาภตฺตํ ปิณฺฑปาตปฏิกฺกนฺโต อายสฺมนฺตํ อานนฺทํ อามนฺเตสิ “อายามานนฺท เยน ปุพฺพาราโม มิคารมาตุปาสาโท เตนุปสงฺกมิสฺสาม ทิวาวิหารายา”ติ. “เอวํ ภนฺเต”ติ โข อายสฺมา อานนฺโท ภควโต ปจฺจสฺโสสิ.}

\prose{อถ โข ภควา อายสฺมตา อานนฺเทน สทฺธิํ เยน ปุพฺพาราโม มิคารมาตุปาสาโท เตนุปสงฺกมิ.\csromanspacer{} อถ โข ภควา สายนฺหสมยํ ปฏิสลฺลานา วุฏฺฐิโต อายสฺมนฺตํ อานนฺทํ อามนฺเตสิ “อายามานนฺท เยน ปุพฺพโกฏฺฐโก เตนุปสงฺกมิสฺสาม คตฺตานิ ปริสิญฺจิตุนฺ”ติ. “เอวํ}

\vspace{\tipitakaparskip}
\csromancenter{ธมฺมิกวคฺค ภนฺเต”ติ โข อายสฺมา อานนฺโท ภควโต ปจฺจสฺโสสิ.\csromanspacer{} อถ โข ภควา อายสฺมตา อานนฺเทน สทฺธิํ เยน ปุพฺพโกฏฺฐโก เตนุปสงฺกมิ คตฺตานิ ปริสิญฺจิตุํ, ปุพฺพโกฏฺฐเก คตฺตานิ ปริสิญฺจิตฺวา ปจฺจุตฺตริตฺวา เอกจีวโร อฏฺฐาสิ คตฺตานิ ปุพฺพาปยมาโน.}
\prose{\csromanfolio{303}เตน โข ปน สมเยน รญฺโญ ปเสนทิสฺส โกสลสฺส เสโต นาม นาโค มหาตูริย\footnote{มหาตุริย (สี, สฺยา, กํ, อิ)} ตาฬิตวาทิเตน ปุพฺพโกฏฺฐกา ปจฺจุตฺตรติ, อปิสฺสุ ตํ ชโน ทิสฺวา เอวมาห “อภิรูโป วต โภ รญฺโญ นาโค, ทสฺสนีโย วต โภ รญฺโญ นาโค, ปาสาทิโก วต โภ รญฺโญ นาโค, กายุปปนฺโน วต โภ รญฺโญ นาโค”ติ.\csromanspacer{} เอวํ วุตฺเต อายสฺมา อุทายี ภควนฺตํ เอตทโวจ “หตฺถิเมว นุ โข ภนฺเต มหนฺตํ พฺรหนฺตํ\footnote{มหนฺตํ พฺรุหนฺตํ (สี), มหตฺตํ พฺรหฺมตฺตํ (ก)} กายุปปนฺนํ ชโน ทิสฺวา เอวมาห ‘นาโค วต โภ นาโค’ติ.\csromanspacer{} อุทาหุ อญฺญมฺปิ กญฺจิ\footnote{กิญฺจิ (ก)} มหนฺตํ พฺรหนฺตํ กายุปปนฺนํ ชโน ทิสฺวา เอวมาห ‘นาโค วต โภ นาโค’ติ”.\csromanspacer{} หตฺถิมฺปิ โข อุทายิ มหนฺตํ พฺรหนฺตํ กายุปปนฺนํ ชโน ทิสฺวา เอวมาห “นาโค วต โภ นาโค”ติ.\csromanspacer{} อสฺสมฺปิ โข อุทายิ มหนฺตํ พฺรหนฺตํ ฯเปฯ.\csromanspacer{} โคณมฺปิ โข อุทายิ มหนฺตํ พฺรหนฺตํ ฯเปฯ.\csromanspacer{} อุรคมฺปิ\footnote{นาคมฺปิ (ก)} โข อุทายิ มหนฺตํ พฺรหนฺตํ ฯเปฯ.\csromanspacer{} รุกฺขมฺปิ โข อุทายิ มหนฺตํ พฺรหนฺตํ ฯเปฯ.\csromanspacer{} มนุสฺสมฺปิ โข อุทายิ มหนฺตํ พฺรหนฺตํ กายุปปนฺนํ ชโน ทิสฺวา เอวมาห “นาโค วต โภ นาโค”ติ.\csromanspacer{} อปิ จ อุทายิ โย สเทวเก โลเก สมารเก สพฺรหฺมเก สสฺสมณพฺราหฺมณิยา ปชาย สเทวมนุสฺสาย อาคุํ น กโรติ กาเยน วาจาย มนสา, ตมหํ “นาโค”ติ พฺรูมีติ.}

\prose{อจฺฉริยํ ภนฺเต, อพฺภุตํ ภนฺเต, ยาว สุภาสิตํ จิทํ ภนฺเต ภควตา.\csromanspacer{} อปิ จ อุทายิ โย สเทวเก โลเก สมารเก สพฺรหฺมเก สสฺสมณ พฺราหฺมณิยา ปชาย สเทวมนุสฺสาย อาคุํ น กโรติ กาเยน วาจาย มนสา, ตมหํ “นาโค”ติ พฺรูมีติ.\csromanspacer{} อิทญฺจ ปนาหํ ภนฺเต ภควตา สุภาสิตํ อิมาหิ คาถาหิ อนุโมทามิ\csromandash{}}

\csromangathameasure{“มนุสฺสภูตํ สมฺพุทฺธํ, อตฺตทนฺตํ สมาหิตํ. \\ อิริยมานํ พฺรหฺมปเถ, จิตฺตสฺสูปสเม รตํ.}
\csromangathagroup{\csromangathastanza{“มนุสฺสภูตํ สมฺพุทฺธํ, อตฺตทนฺตํ สมาหิตํ. \\ อิริยมานํ พฺรหฺมปเถ, จิตฺตสฺสูปสเม รตํ.}}

\gambhira{ฉกฺกนิปาตปาฬิ}
\csromangathameasure{ยํ มนุสฺสา นมสฺสนฺติ, สพฺพธมฺมาน ปารคุํ. \\ เทวาปิ ตํ นมสฺสนฺติ, อิติ เม อรหโต สุตํ. \\ สพฺพสํโยชนาตีตํ, วนา นิพฺพน2มาคตํ. \\ กาเมหิ เนกฺขมฺมรตํ3, มุตฺตํ เสลาว กญฺจนํ. \\ สพฺเพ อจฺจรุจี นาโค, หิมวาญฺเญ สิลุจฺจเย. \\ สพฺเพสํ นาคนามานํ, สจฺจนาโม อนุตฺตโร. \\ นาคํ โว กิตฺตยิสฺสามิ, น หิ อาคุํ กโรติ โส. \\ โสรจฺจํ อวิหิํสา จ, ปาทา นาคสฺส เต ทุเว. \\ ตโป จ พฺรหฺมจริยํ, จรณา นาคสฺส ตฺยาปเร. \\ สทฺธาหตฺโถ มหานาโค, อุเปกฺขาเสตทนฺตวา. \\ สติ คีวา สิโร ปญฺญา, วีมํสา ธมฺมจินฺตนา. \\ ธมฺมกุจฺฉิสมาตโป, วิเวโก ตสฺส วาลธิ. \\ โส ฌายี อสฺสาสรโต, อชฺฌตฺตํ สุสมาหิโต. \\ คจฺฉํ สมาหิโต นาโค, ฐิโต นาโค สมาหิโต. \\ เสยฺยํ สมาหิโต นาโค, นิสินฺโนปิ สมาหิโต. \\ สพฺพตฺถ สํวุโต นาโค, เอสา นาคสฺส สมฺปทา. \\ ภุญฺชติ อนวชฺชานิ, สาวชฺชานิ น ภุญฺชติ. \\ ฆาสมจฺฉาทนํ ลทฺธา, สนฺนิธิํ ปริวชฺชยํ. \\ สํโยชนํ อณุํถูลํ, สพฺพํ เฉตฺวาน พนฺธนํ. \\ เยน เยเนว คจฺฉติ, อนเปกฺโขว คจฺฉติ. \\ ยถาปิ อุทเก ชาตํ, ปุณฺฑรีกํ ปวฑฺฒติ. \\ นุปลิปฺปติ โตเยน, สุจิคนฺธํ มโนรมํ. \\ ตเถว โลเก สุชาโต, พุทฺโธ โลเก วิหรติ. \\ นุปลิปฺปติ โลเกน, โตเยน ปทุมํ ยถา.}
\csromangathagroup{\csromangathastanza{\csromanfolio{304}ยํ มนุสฺสา นมสฺสนฺติ, สพฺพธมฺมาน ปารคุํ. \\ เทวาปิ ตํ\footnote{นํ (สี, อิ)} นมสฺสนฺติ, อิติ เม อรหโต สุตํ.}\csromangathastanza{สพฺพสํโยชนาตีตํ, วนา นิพฺพน2มาคตํ. \\ กาเมหิ เนกฺขมฺมรตํ3, มุตฺตํ เสลาว กญฺจนํ.}\csromangathastanza{สพฺเพ อจฺจรุจี นาโค, หิมวาญฺเญ สิลุจฺจเย. \\ สพฺเพสํ นาคนามานํ, สจฺจนาโม อนุตฺตโร.}\csromangathastanza{นาคํ โว\footnote{เต (ก)} กิตฺตยิสฺสามิ, น หิ อาคุํ กโรติ โส. \\ โสรจฺจํ อวิหิํสา จ, ปาทา นาคสฺส เต ทุเว.}\csromangathastanza{ตโป จ พฺรหฺมจริยํ, จรณา นาคสฺส ตฺยาปเร. \\ สทฺธาหตฺโถ มหานาโค, อุเปกฺขาเสตทนฺตวา.}\csromangathastanza{สติ คีวา สิโร ปญฺญา, วีมํสา ธมฺมจินฺตนา. \\ ธมฺมกุจฺฉิสมาตโป, วิเวโก ตสฺส วาลธิ.}\csromangathastanza{โส ฌายี อสฺสาสรโต, อชฺฌตฺตํ สุสมาหิโต\footnote{อชฺฌตฺตุปสมาหิโต (สฺยา, ก)}. \\ คจฺฉํ สมาหิโต นาโค, ฐิโต นาโค สมาหิโต.}\csromangathastanza{เสยฺยํ สมาหิโต นาโค, นิสินฺโนปิ สมาหิโต. \\ สพฺพตฺถ สํวุโต นาโค, เอสา นาคสฺส สมฺปทา.}\csromangathastanza{ภุญฺชติ อนวชฺชานิ, สาวชฺชานิ น ภุญฺชติ. \\ ฆาสมจฺฉาทนํ ลทฺธา, สนฺนิธิํ ปริวชฺชยํ.}\csromangathastanza{สํโยชนํ อณุํถูลํ, สพฺพํ เฉตฺวาน พนฺธนํ. \\ เยน เยเนว คจฺฉติ, อนเปกฺโขว คจฺฉติ.}\csromangathastanza{ยถาปิ อุทเก ชาตํ, ปุณฺฑรีกํ ปวฑฺฒติ. \\ นุปลิปฺปติ\footnote{น อุปลิปฺปติ (สี, สฺยา, กํ, อิ), นุปลิมฺปติ (ก)} โตเยน, สุจิคนฺธํ มโนรมํ.}\csromangathastanza{ตเถว โลเก สุชาโต, พุทฺโธ โลเก วิหรติ. \\ นุปลิปฺปติ โลเกน, โตเยน ปทุมํ ยถา.}}

\chapterheadpageread{5. ธมฺมิกวคฺค}
\csromangathameasure{มหาคินีว ชลิโต, อนาหารูปสมฺมติ. \\ สงฺขาเรสูปสนฺเตสุ, “นิพฺพุโต”ติ ปวุจฺจติ. \\ อตฺถสฺสายํ วิญฺญาปนี, อุปมา วิญฺญูหิ เทสิตา. \\ วิญฺญสฺสนฺติ มหานาคา, นาคํ นาเคน เทสิตํ. \\ วีตราโค วีตโทโส, วีตโมโห อนาสโว. \\ สรีรํ วิชหํ นาโค, ปรินิพฺพิสฺสติ อนาสโว”ติ.ปฐมํ.}
\csromangathagroup{\csromangathastanza{\csromanfolio{305}มหาคินีว ชลิโต\footnote{มหาคฺคินิ ปชฺชลิโต (สี, สฺยา, กํ)}, อนาหารูปสมฺมติ. \\ สงฺขาเรสูปสนฺเตสุ\footnote{มหาคฺคินิ ปชฺชลิโต (สี, สฺยา, กํ)}, “นิพฺพุโต”ติ ปวุจฺจติ.}\csromangathastanza{อตฺถสฺสายํ วิญฺญาปนี, อุปมา วิญฺญูหิ เทสิตา. \\ วิญฺญสฺสนฺติ\footnote{วิญฺญิสฺสนฺติ (ก)} มหานาคา, นาคํ นาเคน เทสิตํ.}\csromangathastanza{วีตราโค วีตโทโส, วีตโมโห อนาสโว. \\ สรีรํ วิชหํ นาโค, ปรินิพฺพิสฺสติ\footnote{ปรินิพฺพาติ (อิ, ก)} อนาสโว”ติ.\csromanspacer{}\csromanspacer{}ปฐมํ.}}

\csromanmatikaanchor{mtk.165}
\csromantocmarkref{subsection}{2. มิคสาลาสุตฺต}{mtk.165}
\bookmark[dest={mtk.165},level=2]{2. มิคสาลาสุตฺต}
\csromanheader{2}{2. มิคสาลาสุตฺต}
\proseitem{44}{อถ โข อายสฺมา อานนฺโท ปุพฺพณฺหสมยํ นิวาเสตฺวา ปตฺตจีวรมาทาย เยน มิคสาลาย อุปาสิกาย นิเวสนํ เตนุปสงฺกมิ, อุปสงฺกมิตฺวา ปญฺญตฺเต อาสเน นิสีทิ.\csromanspacer{} อถ โข มิคสาลา\footnote{มิคสาณา (ก) อํ 3. 365 ปิฏฺเฐปิ.} อุปาสิกา เยนายสฺมา อานนฺโท เตนุปสงฺกมิ, อุปสงฺกมิตฺวา อายสฺมนฺตํ อานนฺทํ อภิวาเทตฺวา เอกมนฺตํ นิสีทิ, เอกมนฺตํ นิสินฺนา โข มิคสาลา อุปาสิกา อายสฺมนฺตํ อานนฺทํ เอตทโวจ\csromandash{}}

\prose{กถํ กถํ นามายํ ภนฺเต อานนฺท ภควตา ธมฺโม เทสิโต อญฺเญยฺโย.\csromanspacer{} ยตฺร หิ นาม พฺรหฺมจารี จ อพฺรหฺมจารี จ อุโภ สมสมคติกา ภวิสฺสนฺติ อภิสมฺปรายํ.\csromanspacer{} ปิตา เม ภนฺเต ปุราโณ พฺรหฺมจารี อโหสิ อาราจารี วิรโต เมถุนา คามธมฺมา, โส กาลงฺกโต ภควตา พฺยากโต “สกทาคามี สตฺโต\footnote{สกทาคามิปตฺโต (ก, สฺยา, อิ)} ตุสิตํ กายํ อุปปนฺโน”ติ.\csromanspacer{} เปตฺเตยฺโยปิ\footnote{เปตฺเตยฺโย ปิโย (สี, อิ, ก), ปิตุ ปิโย (สฺยา, กํ)} เม ภนฺเต อิสิทตฺโต อพฺรหฺมจารี อโหสิ สทารสนฺตุฏฺโฐ, โสปิ กาลงฺกโต ภควตา พฺยากโต “สกทาคามิปตฺโต ตุสิตํ กายํ อุปปนฺโน”ติ.\csromanspacer{} กถํ กถํ นามายํ ภนฺเต อานนฺท ภควตา ธมฺโม เทสิโต อญฺเญยฺโย.\csromanspacer{} ยตฺร หิ นาม พฺรหฺมจารี จ อพฺรหฺมจารี จ อุโภ สมสมคติกา ภวิสฺสนฺติ อภิสมฺปรายนฺติ.\csromanspacer{} เอวํ โข ปเนตํ ภคินิ ภควตา พฺยากตนฺติ.}

\gambhira{ฉกฺกนิปาตปาฬิ}
\prose{\csromanfolio{306}อถ โข อายสฺมา อานนฺโท มิคสาลาย อุปาสิกาย นิเวสเน ปิณฺฑปาตํ คเหตฺวา อุฏฺฐายาสนา ปกฺกามิ.\csromanspacer{} อถ โข อายสฺมา อานนฺโท ปจฺฉาภตฺตํ ปณฺฑปาตปฏิกฺกนฺโต เยน ภควา เตนุปสงฺกมิ, อุปสงฺกมิตฺวา ภควนฺตํ อภิวาเทตฺวา เอกมนฺตํ นิสีทิ, เอกมนฺตํ นิสินฺโน โข อายสฺมา อานนฺโท ภควนฺตํ เอตทโวจ\csromandash{}}

\prose{อิธาหํ ภนฺเต ปุพฺพณฺหสมยํ นิวาเสตฺวา ปตฺตจีวรมาทาย เยน มิคสาลาย อุปาสิกาย นิเวสนํ เตนุปสงฺกมิํ, อุปสงฺกมิตฺวา ปญฺญตฺเต อาสเน นิสีทิํ.\csromanspacer{} อถ โข ภนฺเต มิคสาลา อุปาสิกา เยนาหํ เตนุปสงฺกมิ, อุปสงฺกมิตฺวา มํ อภิวาเทตฺวา เอกมนฺตํ นิสีทิ, เอกมนฺตํ นิสินฺนา โข ภนฺเต มิคสาลา อุปาสิกา มํ เอตทโวจ “กถํ กถํ นามายํ ภนฺเต อานนฺท ภควตา ธมฺโม เทสิโต อญฺเญยฺโย.\csromanspacer{} ยตฺร หิ นาม พฺรหฺมจารี จ อพฺรหฺมจารี จ อุโภ สมสมคติกา ภวิสฺสนฺติ อภิสมฺปรายํ.\csromanspacer{} ปิตา เม ภนฺเต ปุราโณ พฺรหฺมจารี อโหสิ อาราจารี วิรโต เมถุนา คามธมฺมา, โส กาลงฺกโต ภควตา พฺยากโต ‘สกทาคามิปตฺโต ตุสิตํ กายํ อุปปนฺโน’ติ.\csromanspacer{} เปตฺเตยฺโยปิ เม ภนฺเต อิสิทตฺโต อพฺรหฺมจารี อโหสิ สทารสนฺตุฏฺโฐ, โสปิ กาลงฺกโต ภควตา พฺยากโต ‘สกทาคามิปตฺโต ตุสิตํ กายํ อุปปนฺโน’ติ.\csromanspacer{} กถํ กถํ นามายํ ภนฺเต อานนฺท ภควตา ธมฺโม เทสิโต อญฺเญยฺโย.\csromanspacer{} ยตฺร หิ นาม พฺรหฺมจารี จ อพฺรหฺมจารี จ อุโภ สมสมคติกา ภวิสฺสนฺติ อภิสมฺปรายนฺ”ติ.\csromanspacer{} เอวํ วุตฺเต อหํ ภนฺเต มิคสาลํ อุปาสิกํ เอตทโวจํ “เอวํ โข ปเนตํ ภคินิ ภควตา พฺยากตนฺ”ติ.}

\prose{กา จานนฺท มิคสาลา อุปาสิกา พาลา อพฺยตฺตา อมฺมกา อมฺมกปญฺญา\footnote{อมฺพกา อมฺพกปญฺญา (สี, อิ), อมฺพกา อมฺพกสญฺญา (สฺยา, กํ) อํ 3. 367 ปิฏฺเฐ ปสฺสิตพฺพํ.}, เก จ ปุริสปุคฺคลปโรปริยญาเณ.\csromanspacer{} ฉยิเม อานนฺท ปุคฺคลา สนฺโต สํวิชฺชมานา โลกสฺมิํ.}

\prose{กตเม ฉ, อิธานนฺท เอกจฺโจ ปุคฺคโล โสรโต โหติ สุขสํวาโส, อภินนฺทนฺติ สพฺรหฺมจารี เอกตฺตวาเสน.\csromanspacer{} ตสฺส สวเนนปิ อกตํ โหติ, พาหุสจฺเจนปิ อกตํ โหติ, ทิฏฺฐิยาปิ อปฺปฏิวิทฺธํ โหติ, สามายิกมฺปิ วิมุตฺติํ น ลภติ.\csromanspacer{} โส กายสฺส เภทา}

\vspace{\tipitakaparskip}
\csromancenter{ธมฺมิกวคฺค ปรํ มรณา หานาย ปเรติ โน วิเสสาย, หานคามีเยว โหติ โน วิเสสคามี. (1)}
\prose{\csromanfolio{307}อิธ ปนานนฺท เอกจฺโจ ปุคฺคโล โสรโต โหติ สุขสํวาโส, อภินนฺทนฺติ สพฺรหฺมจารี เอกตฺตวาเสน.\csromanspacer{} ตสฺส สวเนนปิ กตํ โหติ, พาหุสจฺเจนปิ กตํ โหติ, ทิฏฺฐิยาปิ ปฏิวิทฺธํ โหโต, สามายิกมฺปิ วิมุตฺติํ ลภติ.\csromanspacer{} โส กายสฺส เภทา ปรํ มรณา วิเสสาย ปเรติ โน หานาย, วิเสสคามีเยว โหติ โน หานคามี. (2)}

\prose{ตตฺรานนฺท ปมาณิกา ปมิณนฺติ “อิมสฺสปิ เตว ธมฺมา อปรสฺสปิ เตว ธมฺมา, กสฺมา เตสํ เอโก หีโน เอโก ปณีโต”ติ.\csromanspacer{} ตํ หิ เตสํ อานนฺท โหติ ทีฆรตฺตํ อหิตาย ทุกฺขาย.}

\prose{ตตฺรานนฺท ยฺวายํ ปุคฺคโล โสรโต โหติ สุขสํวาโส, อภินนฺทนฺติ สพฺรหฺมจารี เอกตฺตวาเสน.\csromanspacer{} ตสฺส สวเนนปิ กตํ โหติ, พาหุสจฺเจนปิ กตํ โหติ, ทิฏฺฐิยาปิ ปฏิวิทฺธํ โหติ, สามายิกมฺปิ วิมุตฺติํ ลภติ.\csromanspacer{} อยํ อานนฺท ปุคฺคโล อมุนา ปุริเมน ปุคฺคเลน อภิกฺกนฺตตโร จ ปณีตตโร จ.\csromanspacer{} ตํ กิสฺส เหตุ? อิมํ หานนฺท ปุคฺคลํ ธมฺมโสโต นิพฺพหติ, ตทนฺตรํ โก ชาเนยฺย อญฺญตฺร ตถาคเตน.\csromanspacer{} ตสฺมาติหานนฺท มา ปุคฺคเลสุ ปมาณิกา อหุวตฺถ, มา ปุคฺคเลสุ ปมาณํ คณฺหิตฺถ.\csromanspacer{} ขญฺญติ หานนฺท ปุคฺคเลสุ ปมาณํ คณฺหนฺโต, อหํ วา อานนฺท ปุคฺคเลสุ ปมาณํ คเณฺหยฺยํ, โย วา ปนสฺส มาทิโส.}

\prose{อิธ ปนานนฺท เอกจฺจสฺส ปุคฺคลสฺส โกธมาโน อธิคโต\footnote{อวิคโต (ก)} โหติ, สมเยน สมยญฺจสฺส โลภธมฺมา อุปฺปชฺชนฺติ.\csromanspacer{} ตสฺส สวเนนปิ อกตํ โหติ, พาหุสจฺเจนปิ อกตํ โหติ, ทิฏฺฐิยาปิ อปฺปฏิวิทฺธํ โหติ, สามายิกมฺปิ วิมุตฺติํ น ลภติ.\csromanspacer{} โส กายสฺส เภทา ปรํ มรณา หานาย ปเรติ โน วิเสสาย, หานคามีเยว โหติ โน วิเสสคามี. (3)}

\prose{อิธ ปนานนฺท เอกจฺจสฺส ปุคฺคลสฺส โกธมาโน อธิคโต โหติ, สมเยน สมยญฺจสฺส โลภธมฺมา อุปฺปชฺชนฺติ.\csromanspacer{} ตสฺส สวเนนปิ กตํ โหติ ฯเปฯ โน หานคามี. (4)}

\gambhira{ฉกฺกนิปาตปาฬิ}
\prose{\csromanfolio{308}ตตฺรานนฺท ปมาณิกา ปมิณนฺติ ฯเปฯ โย วา ปนสฺส มาทิโส.}

\prose{อิธ ปนานนฺท เอกจฺจสฺส ปุคฺคลสฺส โกธมาโน อธิคโต โหติ, สมเยน สมยญฺจสฺส วจีสงฺขารา อุปฺปชฺชนฺติ.\csromanspacer{} ตสฺส สวเนนปิ อกตํ โหติ ฯเปฯ สามายิกมฺปิ วิมุตฺติํ น ลภติ.\csromanspacer{} โส กายสฺส เภทา ปรํ มรณา หานาย ปเรติ โน วิเสสาย, หานคามีเยว โหติ โน วิเสสาคามี. (5)}

\prose{อิธ ปนานนฺท เอกจฺจสฺส ปุคฺคลสฺส โกธมาโน อธิคโต โหติ, สมเยน สมยญฺจสฺส วจีสงฺขารา อุปฺปชฺชนฺติ.\csromanspacer{} ตสฺส สวเนนปิ กตํ โหติ, พาหุสจฺเจนปิ กตํ โหติ, ทิฏฺฐิยาปิ ปฏิวิทฺธํ โหติ, สามายิกมฺปิ วิมุตฺติํ ลภติ.\csromanspacer{} โส กายสฺส เภทา ปรํ มรณา วิเสสาย ปเรติ โน หานาย, วิเสสคามีเยว โหติ โน หานคามี. (6)}

\prose{ตตฺรานนฺท ปมาณิกา ปมิณนฺติ “อิมสฺสปิ เตว ธมฺมา อปรสฺสปิ เตว ธมฺมา, กสฺมา เตสํ เอโก หีโน เอโก ปณีโต”ติ.\csromanspacer{} ตํ หิ เตสํ อานนฺท โหติ ทีฆรตฺตํ อหิตาย ทุกฺขาย.}

\prose{ตตฺรานนฺท ยสฺส ปุคฺคลสฺส โกธมาโน อธิคโต โหติ, สมเยน สมยญฺจสฺส วจีสงฺขารา อุปฺปชฺชนฺติ.\csromanspacer{} ตสฺส สวเนนปิ กตํ โหติ, พาหุสจฺเจนปิ กตํ โหติ, ทิฏฺฐิยาปิ ปฏิวิทฺธํ โหติ, สามายิกมฺปิ วิมุตฺติํ ลภติ.\csromanspacer{} อยํ อานนฺท ปุคฺคโล อมุนา ปุริเมน ปุคฺคเลน อภิกฺกนฺตตโร จ ปณีตตโร จ.\csromanspacer{} ตํ กิสฺส เหตุ? อิมํ หานนฺท ปุคฺคลํ ธมฺมโสโต นิพฺพหติ, ตทนฺตรํ โก ชาเนยฺย อญฺญตฺร ตถาคเตน.\csromanspacer{} ตสฺมาติหานนฺท มา ปุคฺคเลสุ ปมาณิกา อหุวตฺถ, มา ปุคฺคเลสุ ปมาณํ คณฺหิตฺถ.\csromanspacer{} ขญฺญติ หานนฺท ปุคฺคเลสุ ปมาณํ คณฺหนฺโต, อหํ วา อานนฺท ปุคฺคเลสุ ปมาณํ คเณฺหยฺยํ, โย วา ปนสฺส มาทิโส.}

\prose{กา จานนฺท มิคสาลา อุปาสิกา พาลา อพฺยตฺตา อมฺมกา อมฺมกปญฺญา, เก จ ปุริสปุคฺคลปโรปริยญาเณ.\csromanspacer{} อิเม โข อานนฺท ฉ ปุคฺคลา สนฺโต สํวิชฺชมานา โลกสฺมิํ.}

\prose{ยถารูเปน อานนฺท สีเลน ปุราโณ สมนฺนาคโต อโหสิ, ตถารูเปน สีเลน อิสิทตฺโต สมนฺนาคโต อภวิสฺส.\csromanspacer{} นยิธ ปุราโณ อิสิทตฺตสฺส คติมฺปิ อญฺญสฺส.\csromanspacer{} ยถารูปาย จ อานนฺท ปญฺญาย อิสิทตฺโต สมนฺนาคโต อโหสิ, ตถารูปาย ปญฺญาย ปุราโณ}

\vspace{\tipitakaparskip}
\csromancenter{ธมฺมิกวคฺค สมนฺนาคโต อภวิสฺส.\csromanspacer{} นยิธ อิสิทตฺโต ปุราณสฺส คติมฺปิ อญฺญสฺส.\csromanspacer{} อิติ โข อานนฺท อิเม ปุคฺคลา อุโภ เอกงฺคหีนาติ.\csromanspacer{}\csromanspacer{}ทุติยํ.}
\csromanmatikaanchor{mtk.166}
\csromantocmarkref{subsection}{3. อิณสุตฺต}{mtk.166}
\bookmark[dest={mtk.166},level=2]{3. อิณสุตฺต}
\csromanheader{2}{3. อิณสุตฺต}
\proseitem{45}{\csromanfolio{309}ทาลิทฺทิยํ\footnote{ทาฬิทฺทิยํ (สี)} ภิกฺขเว ทุกฺขํ โลกสฺมิํ กามโภคิโนติ.\csromanspacer{} เอวํ ภนฺเต.\csromanspacer{} ยมฺปิ ภิกฺขเว ทลิทฺโท\footnote{ทฬิทฺโท (สี)} อสฺสโก อนาฬฺหิโก\footnote{อนทฺธิโก (สฺยา, กํ)} อิณํ อาทิยติ.\csromanspacer{} อิณาทานมฺปิ ภิกฺขเว ทุกฺขํ โลกสฺมิํ กามโภคิโนติ.\csromanspacer{} เอวํ ภนฺเต.\csromanspacer{} ยมฺปิ ภิกฺขเว ทลิทฺโท อสฺสโก อนาฬฺหิโก อิณํ อาทิยิตฺวา วฑฺฒิํ ปฏิสฺสุณาติ, วฑฺฒิปิ ภิกฺขเว ทุกฺขา โลกสฺมิํ กามโภคิโนติ.\csromanspacer{} เอวํ ภนฺเต.\csromanspacer{} ยมฺปิ ภิกฺขเว ทลิทฺโท อสฺสโก อนาฬฺหิโก วฑฺฒิํ ปฏิสฺสุณิตฺวา กาลาภตํ\footnote{กาลคตํ (ก)} วฑฺฒิํ น เทติ, โจเทนฺติปิ นํ.\csromanspacer{} โจทนาปิ ภิกฺขเว ทุกฺขา โลกสฺมิํ กามโภคิโนติ.\csromanspacer{} เอวํ ภนฺเต.\csromanspacer{} ยมฺปิ ภิกฺขเว ทลิทฺโท อสฺสโก อนาฬฺหิโก โจทิยมาโน น เทติ, อนุจรนฺติปิ นํ.\csromanspacer{} อนุจริยาปิ ภิกฺขเว ทุกฺขา โลกสฺมิํ กามโภคิโนติ.\csromanspacer{} เอวํ ภนฺเต.\csromanspacer{} ยมฺปิ ภิกฺขเว ทลิทฺโท อสฺสโก อนาฬฺหิโก อนุจริยมาโน น เทติ, พนฺธนฺติปิ นํ.\csromanspacer{} พนฺธนมฺปิ ภิกฺขเว ทุกฺขํ โลกสฺมิํ กามโภคิโนติ.\csromanspacer{} เอวํ ภนฺเต.}

\prose{อิติ โข ภิกฺขเว ทาลิทฺทิยมฺปิ ทุกฺขํ โลกสฺมิํ กามโภคิโน, อิณาทานมฺปิ ทุกฺขํ โลกสฺมิํ กามโภคิโน, วฑฺฒิปิ ทุกฺขา โลกสฺมิํ กามโภคิโน, โจทนาปิ ทุกฺขา โลกสฺมิํ กามโภคิโน, อนุจริยาปิ ทุกฺขา โลกสฺมิํ กามโภคิโน, พนฺธนมฺปิ ทุกฺขํ โลกสฺมิํ กามโภคิโน.\csromanspacer{} เอวเมวํ โข ภิกฺขเว ยสฺส กสฺสจิ สทฺธา นตฺถิ กุสเลสุ ธมฺเมสุ, หิรี นตฺถิ กุสเลสุ ธมฺเมสุ, โอตฺตปฺปํ นตฺถิ กุสเลสุ ธมฺเมสุ, วีริยํ นตฺถิ กุสเลสุ ธมฺเมสุ, ปญฺญา นตฺถิ กุสเลสุ ธมฺเมสุ.\csromanspacer{} อยํ วุจฺจติ ภิกฺขเว อริยสฺส วินเย ทลิทฺโท อสฺสโก อนาฬฺหิโก.}

\prose{ส โข โส ภิกฺขเว ทลิทฺโท อสฺสโก อนาฬฺหิโก สทฺธาย อสติ กุสเลสุ ธมฺเมสุ, หิริยา อสติ กุสเลสุ ธมฺเมสุ, โอตฺตปฺเป อสติ กุสเลสุ ธมฺเมสุ, วีริเย อสติ กุสเลสุ ธมฺเมสุ, ปญฺญาย}

\vspace{\tipitakaparskip}
\csromancenter{ฉกฺกนิปาตปาฬิ อสติ กุสเลสุ ธมฺเมสุ, กาเยน ทุจฺจริตํ จรติ, วาจาย ทุจฺจริตํ จรติ, มนสา ทุจฺจริตํ จรติ.\csromanspacer{} อิทมสฺส อิณาทานสฺมิํ วทามิ.}
\prose{\csromanfolio{310}โส ตสฺส กายทุจฺจริตสฺส ปฏิจฺฉาทนเหตุ ปาปิกํ อิจฺฉํ ปณิทหติ\footnote{ปทหติ (ก)}, “มา มํ ชญฺญู”ติ อิจฺฉติ, “มา มํ ชญฺญู”ติ สงฺกปฺปติ, “มา มํ ชญฺญู”ติ วาจํ ภาสติ, “มา มํ ชญฺญู”ติ กาเยนปรกฺกมติ.\csromanspacer{} โส ตสฺส วจีทุจฺจริตสฺส ปฏิจฺฉาทนเหตุ ฯเปฯ โส ตสฺส มโนทุจฺจริตสฺส ปฏิจฺฉาทนเหตุ ฯเปฯ “มา มํ ชญฺญู”ติ กาเยน ปรกฺกมติ.\csromanspacer{} อิทมสฺส วฑฺฒิยา วทามิ.}

\prose{ตเมนํ เปสลา สพฺรหฺมจารี เอวมาหํสุ “อยญฺจ โส อายสฺมา เอวํการี เอวํสมาจาโร”ติ.\csromanspacer{} อิทมสฺส โจทนาย วทามิ.}

\prose{ตเมนํ อรญฺญคตํ วา รุกฺขมูลคตํ วา สุญฺญาคารคตํ วา วิปฺปฏิสารสหคตา ปาปกา อกุสลวิตกฺกา สมุทาจรนฺติ.\csromanspacer{} อิทมสฺส อนุจริยาย วทามิ.}

\prose{ส โข โส ภิกฺขเว ทลิทฺโท อสฺสโก อนาฬฺหิโก กาเยน ทุจฺจริตํ จริตฺวา วาจาย ทุจฺจริตํ จริตฺวา มนสา ทุจฺจริตํ จริตฺวา กายสฺส เภทา ปรํ มรณา นิรยพนฺธเน วา พชฺฌติ ติรจฺฉานโยนิพนฺธเน วา.\csromanspacer{} นาหํ ภิกฺขเว อญฺญํ เอกพนฺธนมฺปิ สมนุปสฺสามิ เอวํทารุณํ เอวํกฏุกํ\footnote{เอวํทุกฺขํ (สฺยา, กํ, ก)} เอวํอนฺตรายกรํ อนุตฺตรสฺส โยคกฺเขมสฺส อธิคมาย, ยถยิทํ ภิกฺขเว นิรยพนฺธนํ วา ติรจฺฉานโยนิพนฺธนํ วาติ.}

\csromangathameasure{ทาลิทฺทิยํ ทุกฺขํ โลเก, อิณาทานญฺจ วุจฺจติ. \\ ทลิทฺโท อิณมาทาย, ภุญฺชมาโน วิหญฺญติ. \\ ตโต อนุจรนฺติ นํ, พนฺธนมฺปิ นิคจฺฉติ. \\ เอตํ หิ พนฺธนํ ทุกฺขํ, กามลาภาภิชปฺปินํ. \\ ตเถว อริยวินเย, สทฺธา ยสฺส น วิชฺชติ. \\ อหิรีโก อโนตฺตปฺปี, ปาปกมฺมวินิพฺพโย. \\ กายทุจฺจริตํ กตฺวา, วจีทุจฺจริตานิ จ. \\ มโนทุจฺจริตํ กตฺวา, “มา มํ ชญฺญู”ติ อิจฺฉติ.}
\csromangathagroup{\csromangathastanza{ทาลิทฺทิยํ ทุกฺขํ โลเก, อิณาทานญฺจ วุจฺจติ. \\ ทลิทฺโท อิณมาทาย, ภุญฺชมาโน วิหญฺญติ.}\csromangathastanza{ตโต อนุจรนฺติ นํ, พนฺธนมฺปิ นิคจฺฉติ. \\ เอตํ หิ พนฺธนํ ทุกฺขํ, กามลาภาภิชปฺปินํ.}\csromangathastanza{ตเถว อริยวินเย, สทฺธา ยสฺส น วิชฺชติ. \\ อหิรีโก อโนตฺตปฺปี, ปาปกมฺมวินิพฺพโย.}\csromangathastanza{กายทุจฺจริตํ กตฺวา, วจีทุจฺจริตานิ จ. \\ มโนทุจฺจริตํ กตฺวา, “มา มํ ชญฺญู”ติ อิจฺฉติ.}}

\chapterheadpageread{5. ธมฺมิกวคฺค}
\csromangathameasure{โส สํสปฺปติ กาเยน, วาจาย อุท เจตสา. \\ ปาปกมฺมํ ปวฑฺเฒนฺโต, ตตฺถ ตตฺถ ปุนปฺปุนํ. \\ โส ปาปกมฺโม ทุมฺเมโธ, ชานํ ทุกฺกฏมตฺตโน. \\ ทลิทฺโท อิณมาทาย, ภุญฺชมาโน วิหญฺญติ. \\ ตโต อนุจรนฺติ นํ, สงฺกปฺปา มานสา ทุขา. \\ คาเม วา ยทิ วารญฺเญ, ยสฺส วิปฺปฏิสารชา. \\ โส ปาปกมฺโม ทุมฺเมโธ, ชานํ ทุกฺกฏมตฺตโน. \\ โยนิมญฺญตรํ คนฺตฺวา, นิรเย วาปิ พชฺฌติ. \\ เอตํ หิ พนฺธนํ ทุกฺขํ, ยมฺหา ธีโร ปมุจฺจติ. \\ ธมฺมลทฺเธหิ โภเคหิ, ททํ จิตฺตํ ปสาทยํ. \\ อุภยตฺถ กฏคฺคาโห, สทฺธสฺส ฆรเมสิโน. \\ ทิฏฺฐธมฺมหิตตฺถาย, สมฺปรายสุขาย จ. \\ เอวเมตํ คหฏฺฐานํ, จาโค ปุญฺญํ ปวฑฺฒติ. \\ ตเถว อริยวินเย, สทฺธา ยสฺส ปติฏฺฐิตา. \\ หิรีมโน จ โอตฺตปฺปี, ปญฺญวา สีลสํวุโต. \\ เอโส โข อริยวินเย, “สุขชีวี”ติ วุจฺจติ. \\ นิรามิสํ สุขํ ลทฺธา, อุเปกฺขํ อธิติฏฺฐติ. \\ ปญฺจ นีวรเณ หิตฺวา, นิจฺจํ อารทฺธวีริโย. \\ ฌานานิ อุปสมฺปชฺช, เอโกทิ นิปโก สโต. \\ เอวํ ญตฺวา ยถาภูตํ, สพฺพสํโยชนกฺขเย. \\ สพฺพโส อนุปาทาย, สมฺมา จิตฺตํ วิมุจฺจติ. \\ ตสฺส สมฺมาวิมุตฺตสฺส, ญาณํ เจ โหติ ตาทิโน. \\ “อกุปฺปา เม วิมุตฺตี”ติ, ภวสํโยชนกฺขเย. \\ เอตํ โข ปรมํ ญาณํ, เอตํ สุขมนุตฺตรํ. \\ อโสกํ วิรชํ เขมํ, เอตํ อานณฺยมุตฺตมนฺติ.ตติยํ.}
\csromangathagroup{\csromangathastanza{\csromanfolio{311}โส สํสปฺปติ\footnote{สงฺกปฺปติ (ก)} กาเยน, วาจาย อุท เจตสา. \\ ปาปกมฺมํ ปวฑฺเฒนฺโต, ตตฺถ ตตฺถ ปุนปฺปุนํ.}\csromangathastanza{โส ปาปกมฺโม ทุมฺเมโธ, ชานํ ทุกฺกฏมตฺตโน. \\ ทลิทฺโท อิณมาทาย, ภุญฺชมาโน วิหญฺญติ.}\csromangathastanza{ตโต อนุจรนฺติ นํ, สงฺกปฺปา มานสา ทุขา. \\ คาเม วา ยทิ วารญฺเญ, ยสฺส วิปฺปฏิสารชา.}\csromangathastanza{โส ปาปกมฺโม ทุมฺเมโธ, ชานํ ทุกฺกฏมตฺตโน. \\ โยนิมญฺญตรํ คนฺตฺวา, นิรเย วาปิ พชฺฌติ.}\csromangathastanza{เอตํ หิ พนฺธนํ ทุกฺขํ, ยมฺหา ธีโร ปมุจฺจติ. \\ ธมฺมลทฺเธหิ โภเคหิ, ททํ จิตฺตํ ปสาทยํ.}\csromangathastanza{อุภยตฺถ กฏคฺคาโห, สทฺธสฺส ฆรเมสิโน. \\ ทิฏฺฐธมฺมหิตตฺถาย, สมฺปรายสุขาย จ.}\csromangathastanza{เอวเมตํ คหฏฺฐานํ, จาโค ปุญฺญํ ปวฑฺฒติ. \\ ตเถว อริยวินเย, สทฺธา ยสฺส ปติฏฺฐิตา.}\csromangathastanza{หิรีมโน จ โอตฺตปฺปี, ปญฺญวา สีลสํวุโต. \\ เอโส โข อริยวินเย, “สุขชีวี”ติ วุจฺจติ.}\csromangathastanza{นิรามิสํ สุขํ ลทฺธา, อุเปกฺขํ อธิติฏฺฐติ. \\ ปญฺจ นีวรเณ หิตฺวา, นิจฺจํ อารทฺธวีริโย.}\csromangathastanza{ฌานานิ อุปสมฺปชฺช, เอโกทิ นิปโก สโต. \\ เอวํ ญตฺวา ยถาภูตํ, สพฺพสํโยชนกฺขเย.}\csromangathastanza{สพฺพโส อนุปาทาย, สมฺมา จิตฺตํ วิมุจฺจติ. \\ ตสฺส สมฺมาวิมุตฺตสฺส, ญาณํ เจ โหติ ตาทิโน.}\csromangathastanza{“อกุปฺปา เม วิมุตฺตี”ติ, ภวสํโยชนกฺขเย. \\ เอตํ โข ปรมํ ญาณํ, เอตํ สุขมนุตฺตรํ. \\ อโสกํ วิรชํ เขมํ, เอตํ อานณฺยมุตฺตมนฺติ.\csromanspacer{}\csromanspacer{}ตติยํ.}}

\csromanmatikaanchor{mtk.167}
\csromantocmarkref{subsection}{4. มหาจุนฺทสุตฺต}{mtk.167}
\bookmark[dest={mtk.167},level=2]{4. มหาจุนฺทสุตฺต}
\csromanheader{2}{ฉกฺกนิปาตปาฬิ \textbf{4. มหาจุนฺทสุตฺต}}
\proseitem{46}{\csromanfolio{312}เอวํ เม สุตํ\csromandash{}เอกํ สมยํ อายสฺมา มหาจุนฺโท เจตีสุ วิหรติ สยํชาติยํ\footnote{สหชาติยํ (สี, อิ), สญฺชาติยํ (สฺยา, กํ)}.\csromanspacer{} ตตฺร โข อายสฺมา มหาจุนฺโท ภิกฺขู อามนฺเตสิ “อาวุโส ภิกฺขเว”ติ. “อาวุโส”ติ โข เต ภิกฺขู อายสฺมโต มหาจุนฺทสฺส ปจฺจสฺโสสุํ.\csromanspacer{} อายสฺมา มหาจุนฺโท เอตทโวจ\csromandash{}}

\prose{อิธาวุโส ธมฺมโยคา ภิกฺขู ฌายี ภิกฺขู อปสาเทนฺติ “อิเม ปน ‘ฌายิโนมฺหา ฌายิโนมฺหา’ติ ฌายนฺติ ปชฺฌายนฺติ นิชฺฌายนฺติ อวชฺฌายนฺติ\footnote{อปชฺฌายนฺติ (ม 1. 410 ปิฏฺเฐ) 3. กิํ หิเม (สี, สฺยา, กํ, อิ)}, กิมิเม3 ฌายนฺติ, กินฺติเม ฌายนฺติ, กถํ อิเม ฌายนฺตี”ติ.\csromanspacer{} ตตฺถ ธมฺมโยคา จ ภิกฺขู นปฺปสีทนฺติ, ฌายี จ ภิกฺขู นปฺปสีทนฺติ, น จ พหุชนหิตาย ปฏิปนฺนา โหนฺติ พหุชนสุขาย พหุโน ชนสฺส อตฺถาย หิตาย สุขาย เทวมนุสฺสานํ.}

\prose{อิธ ปนาวุโส ฌายี ภิกฺขู ธมฺมโยเค ภิกฺขู อปสาเทนฺติ “อิเม ปน ‘ธมฺมโยคมฺหา ธมฺมโยคมฺหา’ติ อุทฺธตา อุนฺนฬา จปลา มุขรา วิกิณฺณวาจา มุฏฺฐสฺสตี อสมฺปชานา อสมาหิตา วิพฺภนฺตจิตฺตา ปากตินฺทฺริยา, กิมิเม ธมฺมโยคา, กินฺติเม ธมฺมโยคา, กถํ อิเม ธมฺมโยคา”ติ.\csromanspacer{} ตตฺถ ฌายี จ ภิกฺขู นปฺปสีทนฺติ, ธมฺมโยคา จ ภิกฺขู นปฺปสีทนฺติ, น จ พหุชนหิตาย ปฏิปนฺนา โหนฺติ พหุชนสุขาย พหุโน ชนสฺส อตฺถาย หิตาย สุขาย เทวมนุสฺสานํ.}

\prose{อิธ ปนาวุโส ธมฺมโยคา ภิกฺขู ธมฺมโยคานญฺเญว ภิกฺขูนํ วณฺณํ ภาสนฺติ, โน ฌายีนํ ภิกฺขูนํ วณฺณํ ภาสนฺติ.\csromanspacer{} ตตฺถ ธมฺมโยคา จ ภิกฺขู นปฺปสีทนฺติ, ฌายี จ ภิกฺขู นปฺปสีทนฺติ, น จ พหุชนหิตาย ปฏิปนฺนา โหนฺติ พหุชนสุขาย พหุโน ชนสฺส อตฺถาย หิตาย สุขาย เทวมนุสฺสานํ.}

\prose{อิธ ปนาวุโส ฌายี ภิกฺขู ฌายีนญฺเญว ภิกฺขูนํ วณฺณํ ภาสนฺติ, โน ธมฺมโยคานํ ภิกฺขูนํ วณฺณํ ภาสนฺติ.\csromanspacer{} ตตฺถ ฌายี จ ภิกฺขู นปฺปสีทนฺติ, ธมฺมโยคา จ ภิกฺขู นปฺปสีทนฺติ, น จ พหุชนหิตาย ปฏิปนฺนา โหนฺติ}

\vspace{\tipitakaparskip}
\csromancenter{ธมฺมิกวคฺค พหุชนสุขาย พหุโน ชนสฺส อตฺถาย หิตาย สุขาย เทวมนุสฺสานํ.}
\csromanmatikaanchor{mtk.168}
\csromantocmarkref{subsection}{5-6. สนฺทิฏฺฐิกสุตฺตทฺวย}{mtk.168}
\bookmark[dest={mtk.168},level=2]{5-6. สนฺทิฏฺฐิกสุตฺตทฺวย}
\prose{\csromanfolio{313}ตสฺมาติหาวุโส เอวํ สิกฺขิตพฺพํ “ธมฺมโยคา สมานา ฌายีนํ ภิกฺขูนํ วณฺณํ ภาสิสฺสามา”ติ, เอวํ หิ โว อาวุโส สิกฺขิตพฺพํ.\csromanspacer{} ตํ กิสฺส เหตุ? อจฺฉริยา เหเต อาวุโส ปุคฺคลา ทุลฺลภา โลกสฺมิํ, เย อมตํ ธาตุํ กาเยน ผุสิตฺวา วิหรนฺติ.\csromanspacer{} ตสฺมาติหาวุโส เอวํ สิกฺขิตพฺพํ “ฌายี สมานา ธมฺมโยคานํ ภิกฺขูนํ วณฺณํ ภาสิสฺสามา”ติ, เอวํ หิ โว อาวุโส สิกฺขิตพฺพํ.\csromanspacer{} ตํ กิสฺส เหตุ? อจฺฉริยา เหเต อาวุโส ปุคฺคลา ทุลฺลภา โลกสฺมิํ, เย คมฺภีรํ อตฺถปทํ ปญฺญาย อติวิชฺฌ ปสฺสนฺตีติ.\csromanspacer{}\csromanspacer{}จตุตฺถํ.}

\csromanheader{2}{5. ปฐมสนฺทิฏฺฐิกสุตฺต}
\proseitem{47}{อถ โข โมฬิยสีวโก\footnote{โมลิยสีวโก (สี, อิ), โมฬิสีวโก (ก)} ปริพฺพาชโก เยน ภควา เตนุปสงฺกมิ, อุปสงฺกมิตฺวา ภควตา สทฺธิํ สมฺโมทิ, สมฺโมทนียํ กถํ สารณียํ วีติสาเรตฺวา เอกมนฺตํ นิสีทิ, เอกมนฺตํ นิสินฺโน โข โมฬิยสีวโก ปริพฺพาชโก ภควนฺตํ เอตทโวจ “สนฺทิฏฺฐิโก ธมฺโม สนฺทิฏฺฐิโก ธมฺโมติ ภนฺเต วุจฺจติ, กิตฺตาวตา นุ โข ภนฺเต สนฺทิฏฺฐิโก ธมฺโม โหติ อกาลิโก เอหิปสฺสิโก โอปเนยฺยิโก ปจฺจตฺตํ เวทิตพฺโพ วิญฺญูหี”ติ.}

\prose{เตน หิ สีวก ตญฺเญเวตฺถ ปฏิปุจฺฉามิ, ยถา เต ขเมยฺย, ตถา นํ พฺยากเรยฺยาสิ.\csromanspacer{} ตํ กิํ มญฺญสิ สีวก, สนฺตํ วา อชฺฌตฺตํ โลภํ “อตฺถิ เม อชฺฌตฺตํ โลโภ”ติ ปชานาสิ, อสนฺตํ วา อชฺฌตฺตํ โลภํ “นตฺถิ เม อชฺฌตฺตํ โลโภ”ติ ปชานาสีติ.\csromanspacer{} เอวํ ภนฺเต.\csromanspacer{} ยํ โข ตฺวํ สีวก สนฺตํ วา อชฺฌตฺตํ โลภํ “อตฺถิ เม อชฺฌตฺตํ โลโภ”ติ ปชานาสิ, อสนฺตํ วา อชฺฌตฺตํ โลภํ “นตฺถิ เม อชฺฌตฺตํ โลโภ”ติ ปชานาสิ.\csromanspacer{} เอวมฺปิ โข สีวก สนฺทิฏฺฐิโก ธมฺโม โหติ ฯเปฯ. ตํ กิํ มญฺญสิ สีวก, สนฺตํ วา อชฺฌตฺตํ โทสํฯเปฯ สนฺตํ วา อชฺฌตฺตํ โมหํ ฯเปฯ สนฺตํ วา อชฺฌตฺตํ โลภธมฺมํ ฯเปฯ สนฺตํ วา อชฺฌตฺตํ โทสธมฺมํ ฯเปฯ สนฺตํ วา อชฺฌตฺตํ โมหธมฺมํ “อตฺถิ เม อชฺฌตฺตํ โมหธมฺโม”ติ ปชานาสิ,}

\vspace{\tipitakaparskip}
\csromancenter{ฉกฺกนิปาตปาฬิ อสนฺตํ วา อชฺฌตฺตํ โมหธมฺมํ “นตฺถิ เม อชฺฌตฺตํ โมหธมฺโม”ติ ปชานาสีติ.\csromanspacer{} เอวํ ภนฺเต.\csromanspacer{} ยํ โข ตฺวํ สีวก สนฺตํ วา อชฺฌตฺตํ โมหธมฺมํ “อตฺถิ เม อชฺฌตฺตํ โมหธมฺโม”ติ ปชานาสิ, อสนฺตํ วา อชฺฌตฺตํ โมหธมฺมํ “นตฺถิ เม อชฺฌตฺตํ โมหธมฺโม”ติ ปชานาสิ.\csromanspacer{} เอวํ โข สีวก สนฺทิฏฺฐิโก ธมฺโม โหติ อกาลิโก เอหิปสฺสิโก โอปเนยฺยิโก ปจฺจตฺตํ เวทิตพฺโพ วิญฺญูหีติ.}
\prose{\csromanfolio{314}อภิกฺกนฺตํ ภนฺเต, อภิกฺกนฺตํ ภนฺเต ฯเปฯ อุปาสกํ มํ ภนฺเต ภควา ธาเรตุ อชฺชตคฺเค ปาณุเปตํ สรณํ คตนฺติ.\csromanspacer{}\csromanspacer{}ปญฺจมํ.}

\csromanheader{2}{6. ทุติยสนฺทิฏฺฐิกสุตฺต}
\proseitem{48}{อถ โข อญฺญตโร พฺราหฺมโณ เยน ภควา เตนุปสงฺกมิ, อุปสงฺกมิตฺวา ภควตา สทฺธิํ สมฺโมทิ, สมฺโมทนียํ กถํ สารณียํ วีติสาเรตฺวา เอกมนฺตํ นิสีทิ, เอกมนฺตํ นิสินฺโน โข โส พฺราหฺมโณ ภควนฺตํ เอตทโวจ “สนฺทิฏฺฐิโก ธมฺโม สนฺทิฏฺฐิโก ธมฺโมติ โภ โคตม วุจฺจติ, กิตฺตาวตา นุ โข โภ โคตม สนฺทิฏฺฐิโก ธมฺโม โหติ อกาลิโก เอหิปสฺสิโก โอปเนยฺยิโก ปจฺจตฺตํ เวทิตพฺโพ วิญฺญูหี”ติ.}

\prose{เตน หิ พฺราหฺมณ ตญฺเญเวตฺถ ปฏิปุจฺฉิสฺสามิ, ยถา เต ขเมยฺย, ตถา นํ พฺยากเรยฺยาสิ.\csromanspacer{} ตํ กิํ มญฺญสิ พฺราหฺมณ, สนฺตํ วา อชฺฌตฺตํ ราคํ “อตฺถิ เม อชฺฌตฺตํ ราโค”ติ ปชานาสิ, อสนฺตํ วา อชฺฌตฺตํ ราคํ “นตฺถิ เม อชฺฌตฺตํ ราโคติ ปชานาสี”ติ.\csromanspacer{} เอวํ โภ.\csromanspacer{} ยํ โข ตฺวํ พฺราหฺมณ สนฺตํ วา อชฺฌตฺตํ ราคํ “อตฺถิ เม อชฺฌตฺตํ ราโค”ติ ปชานาสิ, อสนฺตํ วา อชฺฌตฺตํ ราคํ “นตฺถิ เม อชฺฌตฺตํ ราโค”ติ ปชานาสิ.\csromanspacer{} เอวมฺปิ โข พฺราหฺมณ สนฺทิฏฺฐิโก ธมฺโม โหติ ฯเปฯ.}

\prose{ตํ กิํ มญฺญสิ พฺราหฺมณ, สนฺตํ วา อชฺฌตฺตํ โทสํ ฯเปฯ สนฺตํ วา อชฺฌตฺตํ โมหํ ฯเปฯสนฺตํ วา อชฺฌตฺตํ กายสนฺโทสํ ฯเปฯ สนฺตํ วา อชฺฌตฺตํ วจีสนฺโทสํ ฯเปฯ สนฺตํ วา อชฺฌตฺตํ มโนสนฺโทสํ “อตฺถิ เม อชฺฌตฺตํ มโนสนฺโทโส”ติ ปชานาสิ, อสนฺตํ วา อชฺฌตฺตํ มโนสนฺโทสํ “นตฺถิ เม อชฺฌตฺตํ มโนสนฺโทโส”ติ ปชานาสีติ.\csromanspacer{} เอวํ โภ.\csromanspacer{} ยํ โข ตฺวํ พฺราหฺมณ สนฺตํ วา อชฺฌตฺตํ มโนสนฺโทสํ “อตฺถิ เม อชฺฌตฺตํ มโนสนฺโทโส”ติ ปชานาสิ, อสนฺตํ วา อชฺฌตฺตํ มโนสนฺโทสํ “นตฺถิ เม อชฺฌตฺตํ}

\vspace{\tipitakaparskip}
\csromancenter{ธมฺมิกวคฺค มโนสนฺโทโส”ติ ปชานาสิ.\csromanspacer{} เอวํ โข พฺราหฺมณ สนฺทิฏฺฐิโก ธมฺโม โหติ อกาลิโก เอหิปสฺสิโก โอปเนยฺยิโก ปจฺจตฺตํ เวทิตพฺโพ วิญฺญูหีติ.}
\prose{\csromanfolio{315}อภิกฺกนฺตํ โภ โคตม, อภิกฺกนฺตํ โภ โคตม ฯเปฯ อุปาสกํ มํ ภวํ โคตโม ธาเรตุ อชฺชตคฺเค ปาณุเปตํ สรณํ คตนฺติ.\csromanspacer{}\csromanspacer{}ฉฏฺฐํ.}

\csromanmatikaanchor{mtk.169}
\csromantocmarkref{subsection}{7. เขมสุตฺต}{mtk.169}
\bookmark[dest={mtk.169},level=2]{7. เขมสุตฺต}
\csromanheader{2}{7. เขมสุตฺต}
\proseitem{49}{เอกํ สมยํ ภควา สาวตฺถิยํ วิหรติ เชตวเน อนาถปิณฺฑิกสฺส อาราเม.\csromanspacer{} เตน โข ปน สมเยน อายสฺมา จ เขโม อายสฺมา จ สุมโน สาวตฺถิยํ วิหรนฺติ อนฺธวนสฺมิํ.\csromanspacer{} อถ โข อายสฺมา จ เขโม อายสฺมา จ สุมโน เยน ภควา เตนุปสงฺกมิํสุ, อุปสงฺกมิตฺวา ภควนฺตํ อภิวาเทตฺวา เอกมนฺตํ นิสีทิํสุ, เอกมนฺตํ นิสินฺโน โข อายสฺมา เขโม ภควนฺตํ เอตทโวจ\csromandash{}}

\prose{โย โส ภนฺเต ภิกฺขุ อรหํ ขีณาสโว วุสิตวา กตกรณีโย โอหิตภาโร อนุปฺปตฺตสทตฺโถ ปริกฺขีณภวสํโยชโน สมฺมทญฺญา วิมุตฺโต.\csromanspacer{} ตสฺส น เอวํ โหติ “อตฺถิ เม เสยฺโยติ วา อตฺถิ เม สทิโสติ วา อตฺถิ เม หีโนติ วา”ติ.\csromanspacer{} อิทมโวจายสฺมา เขโม, สมนุญฺโญ สตฺถา อโหสิ.\csromanspacer{} อถ โข อายสฺมา เขโม “สมนุญฺโญ เม สตฺถา”ติ อุฏฺฐายาสนา ภควนฺตํ อภิวาเทตฺวา ปทกฺขิณํ กตฺวา ปกฺกามิ.}

\prose{อถ โข อายสฺมา สุมโน อจิรปกฺกนฺเต อายสฺมนฺเต เขเม ภควนฺตํ เอตทโวจ “โย โส ภนฺเต ภิกฺขุ อรหํ ขีณาสโว วุสิตวา กตกรณีโย โอหิตภาโร อนุปฺปตฺตสทตฺโถ ปริกฺขีณภวสํโยชโน สมฺมทญฺญา วิมุตฺโต.\csromanspacer{} ตสฺส น เอวํ โหติ ‘นตฺถิ เม เสยฺโยติ วา นตฺถิ เม สทิโสติ วา นตฺถิ เม หีโนติ วา’ติ”.\csromanspacer{} อิทมโว จายสฺมา สุมโน, สมนุญฺโญ สตฺถา อโหสิ.\csromanspacer{} อถ โข อายสฺมา สุมโน “สมนุญฺโญ เม สตฺถา”ติ อุฏฺฐายาสนา ภควนฺตํ อภิวาเทตฺวา ปทกฺขิณํ กตฺวา ปกฺกามิ.}

\prose{อถ โข ภควา อจิรปกฺกนฺเตสุ อายสฺมนฺเต จ เขเม อายสฺมนฺเต จ สุมเน ภิกฺขู อามนฺเตสิ “เอวํ โข ภิกฺขเว กุลปุตฺตา อญฺญํ}

\vspace{\tipitakaparskip}
\csromancenter{ฉกฺกนิปาตปาฬิ พฺยากโรนฺติ, อตฺโถ จ วุตฺโต, อตฺตา จ อนุปนีโต, อถ จ ปน อิเธกจฺเจ โมฆปุริสา หสมานกา\footnote{หสมานกํ (ก) วิ 3. 272 ปิฏฺเฐปิ.} มญฺเญ อญฺญํ พฺยากโรนฺติ, เต ปจฺฉา วิฆาตํ อาปชฺชนฺตี”ติ.}
\prose{\csromanfolio{316}น อุสฺเสสุ น โอเมสุ, สมตฺเต โนปนียเร\footnote{โนปนิยฺยเร (สฺยา, อิ, ก)}.}

\prose{ขีณา ชาติ วุสิตํ พฺรหฺมจริยํ, จรนฺติ สํโยชนวิปฺปมุตฺตาติ.\csromanspacer{} สตฺตมํ.}

\csromanmatikaanchor{mtk.170}
\csromantocmarkref{subsection}{8. อินฺทฺริยสํวรสุตฺต}{mtk.170}
\bookmark[dest={mtk.170},level=2]{8. อินฺทฺริยสํวรสุตฺต}
\csromanheader{2}{8. อินฺทฺริยสํวรสุตฺต}
\proseitem{50}{\csromansymbolfootnote{*}{เหฏฺฐา 16, 175 ปิฏฺเฐสุ; อุปริ 472 ปิฏฺเฐปิ.}อินฺทฺริยสํวเร ภิกฺขเว อสติ อินฺทฺริยสํวรวิปนฺนสฺส หตูปนิสํ โหติ สีลํ, สีเล อสติ สีลวิปนฺนสฺส หตูปนิโส โหติ สมฺมาสมาธิ, สมฺมาสมาธิมฺหิ อสติ สมฺมาสมาธิวิปนฺนสฺส หตูปนิสํ โหติ ยถาภูตญาณทสฺสนํ, ยถาภูตญาณทสฺสเน อสติ ยถาภูตญาณทสฺสนวิปนฺนสฺส หตูปนิโส โหติ นิพฺพิทาวิราโค, นิพฺพิทาวิราเค อสติ นิพฺพิทาวิราควิปนฺนสฺส หตูปนิสํ โหติ วิมุตฺติญาณทสฺสนํ.\csromanspacer{} เสยฺยถาปิ ภิกฺกเว รุกฺโข สาขาปลาสวิปนฺโน, ตสฺส ปปฏิกาปิ น ปาริปูริํ คจฺฉติ, ตโจปิ น ปาริปูริํ คจฺฉติ, เผคฺคุปิ น ปาริปูริํ คจฺฉติ, สาโรปิ น ปาริปูริํ คจฺฉติ.\csromanspacer{} เอวเมวํ โข ภิกฺขเว อินฺทฺริยสํวเร อสติ อินฺทฺริยสํวรวิปนฺนสฺส หตูปนิสํ โหติ สีลํ ฯเปฯ วิมุตฺติญาณทสฺสนํ.}

\prose{อินฺทฺริยสํวเร ภิกฺขเว สติ อินฺทฺริยสํวรสมฺปนฺนสฺส อุปนิสสมฺปนฺนํ โหติ สีลํ, สีเล สติ สีลสมฺปนฺนสฺส อุปนิสสมฺปนฺโน โหติ สมฺมาสมาธิ, สมฺมาสมาธิมฺหิ สติ สมฺมาสมาธิสมฺปนฺนสฺส อุปนิสสมฺปนฺนํ โหติ ยถาภูต ญาณทสฺสนํ, ยถาภูตญาณทสฺสเน สติ ยถาภูตญาณทสฺสนสมฺปนฺนสฺส อุปนิสสมฺปนฺโน โหติ นิพฺพิทาวิราโค, นิพฺพิทาวิราเค สติ นิพฺพิทาวิราคสมฺปนฺนสฺส อุปนิสสมฺปนฺนํ โหติ วิมุตฺติญาณทสฺสนํ.\csromanspacer{} เสยฺยถาปิ ภิกฺขเว รุกฺโข สาขาปลาสสมฺปนฺโน, ตสฺส ปปฏิกาปิ ปาริปูริํ คจฺฉติ, ตโจปิ ปาริปูริํ คจฺฉติ, เผคฺคุปิ ปาริปูริํ คจฺฉติ, สาโรปิ ปาริปูริํ}

\vspace{\tipitakaparskip}
\csromancenter{ธมฺมิกวคฺค คจฺฉติ.\csromanspacer{} เอวเมวํ โข ภิกฺขเว อินฺทฺริยสํวเร สติ อินฺทฺริยสํวรสมฺปนฺนสฺส อุปนิสสมฺปนฺนํ โหติ สีลํ ฯเปฯ วิมุตฺติญาณทสฺสนนฺติ.\csromanspacer{}\csromanspacer{}อฏฺฐมํ.}
\csromanmatikaanchor{mtk.171}
\csromantocmarkref{subsection}{9. อานนฺทสุตฺต}{mtk.171}
\bookmark[dest={mtk.171},level=2]{9. อานนฺทสุตฺต}
\csromanheader{2}{9. อานนฺทสุตฺต}
\proseitem{51}{\csromanfolio{317}อถ โข อายสฺมา อานนฺโท เยนายสฺมา สาริปุตฺโต เตนุปสงฺกมิ, อุปสงฺกมิตฺวา อายสฺมตา สาริปุตฺเตน สทฺธิํ สมฺโมทิ, สมฺโมทนียํ กถํ สารณียํ วีติสาเรตฺวา เอกมนฺตํ นิสีทิ, เอกมนฺตํ นิสินฺโน โข อายสฺมา อานนฺโท อายสฺมนฺตํ สาริปุตฺตํ เอตทโวจ\csromandash{}}

\prose{กิตฺตาวตา นุ โข อาวุโส สาริปุตฺต ภิกฺขุ อสฺสุตญฺเจว ธมฺมํ สุณาติ, สุตา จสฺส ธมฺมา น สมฺโมสํ คจฺฉนฺติ, เย จสฺส ธมฺมา ปุพฺเพ เจตสา สมฺผุฏฺฐปุพฺพา, เต จ สมุทาจรนฺติ, อวิญฺญาตญฺจ วิชานาตีติ.\csromanspacer{} อายสฺมา โข อานนฺโท พหุสฺสุโต ปฏิภาตุ อายสฺมนฺตํเยว อานนฺทนฺติ.\csromanspacer{} เตนหาวุโส สาริปุตฺต สุณาหิ, สาธุกํ มนสิ กโรหิ ภาสิสฺสามีติ. “เอวมาวุโส”ติ โข อายสฺมา สาริปุตฺโต อายสฺมโต อานนฺทสฺส ปจฺจสฺโสสิ.\csromanspacer{} อายสฺมา อานนฺโท เอตทโวจ\csromandash{}}

\prose{อิธาวุโส สาริปุตฺต ภิกฺขุ ธมฺมํ ปริยาปุณาติ สุตฺตํ เคยฺยํ เวยฺยากรณํ คาถํ อุทานํ อิติวุตฺตกํ ชาตกํ อพฺภุตธมฺมํ เวทลฺลํ.\csromanspacer{} โส ยถาสุตํ ยถาปริยตฺตํ ธมฺมํ วิตฺถาเรน ปเรสํ เทเสติ.\csromanspacer{} ยถาสุตํ ยถาปริยตฺตํ ธมฺมํ วิตฺถาเรน ปเรสํ วาเจติ.\csromanspacer{} ยถาสุตํ ยถาปริยตฺตํ ธมฺมํ วิตฺถาเรน สชฺฌายํ กโรติ.\csromanspacer{} ยถาสุตํ ยถาปริยตฺตํ ธมฺมํ เจตสา อนุวิตกฺเกติ อนุวิจาเรติ มนสานุเปกฺขติ.\csromanspacer{} ยสฺมิํ อาวาเส เถรา ภิกฺขู วิหรนฺติ พหุสฺสุตา อาคตาคมา ธมฺมธรา วินยธรา มาติกาธรา, ตสฺมิํ อาวาเส วสฺสํ อุเปติ.\csromanspacer{} เต กาเลน กาลํ อุปสงฺกมิตฺวา ปริปุจฺฉติ ปริปญฺหติ “อิทํ ภนฺเต กถํ อิมสฺส กฺวตฺโถ”ติ, เต ตสฺส อายสฺมโต อวิวฏํ เจว วิวรนฺติ, อนุตฺตานีกตญฺจ อุตฺตานีกโรนฺติ, อเนกวิหิเตสุ จ กงฺขาฐานิเยสุ ธมฺเมสุ กงฺขํ ปฏิวิโนเทนฺติ.\csromanspacer{} เอตฺตาวตา โข อาวุโส สาริปุตฺต ภิกฺขุ อสฺสุตญฺเจว ธมฺมํ สุณาติ, สุตา จสฺส ธมฺมา น สมฺโมสํ คจฺฉนฺติ, เย จสฺส ธมฺมา ปุพฺเพ เจตสา สมฺผุฏฺฐปุพฺพา, เต จ สมุทาจรนฺติ, อวิญฺญาตญฺจ วิชานาตีติ.}

\gambhira{ฉกฺกนิปาตปาฬิ}
\prose{\csromanfolio{318}อจฺฉริยํ อาวุโส, อพฺภุตํ อาวุโส, ยาว สุภาสิตํ จิทํ อายสฺมตา อานนฺเทน.\csromanspacer{} อิเมหิ จ มยํ ฉหิ ธมฺเมหิ สมนฺนาคตํ อายสฺมนฺตํ อานนฺทํ ธาเรม.\csromanspacer{} อายสฺมา หิ อานนฺโท ธมฺมํ ปริยาปุณาติ สุตฺตํ เคยฺยํ เวยฺยากรณํ คาถํ อุทานํ อิติวุตฺตกํ ชาตกํ อพฺภุตธมฺมํ เวทลฺลํ.\csromanspacer{} อายสฺมา อานนฺโท ยถาสุตํ ยถาปริยตฺตํ ธมฺมํ วิตฺถาเรน ปเรสํ เทเสติ.\csromanspacer{} อายสฺมา อานนฺโท ยถาสุตํ ยถาปริยตฺตํ ธมฺมํ วิตฺถาเรน ปเรสํ วาเจติ.\csromanspacer{} อายสฺมา อานนฺโท ยถาสุตํ ยถาปริยตฺตํ ธมฺมํ วิตฺถาเรน สชฺฌายํ กโรติ.\csromanspacer{} อายสฺมา อานนฺโท ยถาสุตํ ยถาปริยตฺตํ ธมฺมํ เจตสา อนุวิตกฺเกติ อนุวิจาเรติ มนสานุ เปกฺขติ.\csromanspacer{} อายสฺมา อานนฺโท ยสฺมิํ อาวาเส เถรา ภิกฺขู วิหรนฺติ พหุสฺสุตา อาคตาคมา ธมฺมธรา วินยธรา มาติกาธรา, ตสฺมิํ อาวาเส วสฺสํ อุเปติ.\csromanspacer{} เต อายสฺมา อานนฺโท กาเลน กาลํ อุปสงฺกมิตฺวา ปริปุจฺฉติ ปริปญฺหติ “อิทํ ภนฺเต กถํ อิมสฺส กฺวตฺโถ”ติ, เต อายสฺมโต อานนฺทสฺส อวิวฏญฺเจว วิวรนฺติ, อนุตฺตานีกตญฺจ อุตฺตานีกโรนฺติ, อเนกวิหิเตสุ จ กงฺขาฐานิเยสุ ธมฺเมสุ กงฺขํ ปฏิวิโนเทนฺตีติ.\csromanspacer{}\csromanspacer{}นวมํ.}

\csromanmatikaanchor{mtk.172}
\csromantocmarkref{subsection}{10. ขตฺติยสุตฺต}{mtk.172}
\bookmark[dest={mtk.172},level=2]{10. ขตฺติยสุตฺต}
\csromanheader{2}{10. ขตฺติยสุตฺต}
\proseitem{52}{อถ โข ชาณุสฺโสณิ\footnote{ชาณุโสณิ (ก)} พฺราหฺมโณ เยน ภควา เตนุปสงฺกมิ, อุปสงฺก มิตฺวา ภควตา สทฺธิํ สมฺโมทิ, สมฺโมทนียํ กถํ สารณียํ วีติสาเรตฺวา เอกมนฺตํ นิสีทิ, เอกมนฺตํ นิสินฺโน โข ชาณุสฺโสณิ พฺราหฺมโณ ภควนฺตํ เอตทโวจ\csromandash{}}

\prose{ขตฺติยา โภ โคตม กิํอธิปฺปายา กิํอุปวิจารา กิํอธิฏฺฐานา กิํอภินิเวสา กิํปริโยสานาติ.\csromanspacer{} ขตฺติยา โข พฺราหฺมณ โภคาธิปฺปายา ปญฺญูปวิจารา พลาธิฏฺฐานา ปถวีภินิเวสา อิสฺสริยปริโยสานาติ.}

\prose{พฺราหฺมณา ปน โภ โคตม กิํอธิปฺปายา กิํอุปวิจารา กิํอธิฏฺฐานา กิํอภินิเวสา กิํปริโยสานาติ.\csromanspacer{} พฺราหฺมณา โข พฺราหฺมณ โภคาธิปฺปายา ปญฺญูปวิจารา มนฺตาธิฏฺฐานา ยญฺญาภินิเวสา พฺรหฺมโลกปริโยสานาติ.}

\chapterheadpageread{5. ธมฺมิกวคฺค}
\prose{\csromanfolio{319}คหปติกา ปน โภ โคตม กิํอธิปฺปายา กิํอุปวิจารา กิํอภินิเวสา กิํปริโยสานาติ.\csromanspacer{} คหปติกา โข พฺราหฺมณ โภคาธิปฺปายา ปญฺญูปวิจารา สิปฺปาธิฏฺฐานา กมฺมนฺตาภินิเวสา นิฏฺฐิตกมฺมนฺตปริโยสานาติ.}

\prose{อิตฺถี ปน โภ โคตม กิํอธิปฺปายา กิํอุปวิจารา กิํอธิฏฺฐานา กิํอภินิเวสา กิํปริโยสานาติ.\csromanspacer{} อิตฺถี โข พฺราหฺมณ ปุริสาธิปฺปายา อลงฺการูปวิจารา ปุตฺตาธิฏฺฐานา อสปตฺตีภินิเวสา อิสฺสริยปริโยสานาติ.}

\prose{โจรา ปน โภ โคตม กิํอธิปฺปายา กิํอุปวิจารา กิํอธิฏฺฐานา กิํอภินิเวสา กิํปริโยสานาติ.\csromanspacer{} โจรา โข พฺราหฺมณ อาทานาธิปฺปายา คหนูปวิจารา สตฺถาธิฏฺฐานา อนฺธการาภินิเวสา อทสฺสนปริโยสานาติ.}

\prose{สมณา ปน โภ โคตม กิํอธิปฺปายา กิํอุปวิจารา กิํอธิฏฺฐานา กิํอภินิเวสา กิํปริโยสานาติ.\csromanspacer{} สมณา โข พฺราหฺมณ ขนฺติโสรจฺจาธิปฺปายา ปญฺญูปวิจารา สีลาธิฏฺฐานา อากิญฺจญฺญาภินิเวสา\footnote{อกิญฺจนาภินิเวสา (สฺยา, ก)} นิพฺพานปริโยสานาติ.}

\prose{อจฺฉริยํ โภ โคตม, อพฺภุตํ โภ โคตม, ขตฺติยานมฺปิ ภวํ โคตโม ชานาติ อธิปฺปายญฺจ อุปวิจารญฺจ อธิฏฺฐานญฺจ อภินิเวสญฺจ ปริโยสานญฺจ.\csromanspacer{} พฺราหฺมณานมฺปิ ภวํ โคตโม ชานาติ ฯเปฯ คหปตีนมฺปิ ภวํ โคตโม ชานาติ.\csromanspacer{} อิตฺถีนมฺปิ ภวํ โคตโม ชานาติ.\csromanspacer{} โจรานมฺปิ ภวํ โคตโม ชานาติ.\csromanspacer{} สมณานมฺปิ ภวํ โคตโม ชานาติ อธิปฺปายญฺจ อุปวิจารญฺจ อธิฏฺฐานญฺจ อภินิเวสญฺจ ปริโยสานญฺจ.\csromanspacer{} อภิกฺกนฺตํ โภ โคตม ฯเปฯ อุปาสกํ มํ ภวํ โคตโม ธาเรตุ อชฺชตคฺเค ปาณุเปตํ สรณํ คตนฺติ.\csromanspacer{}\csromanspacer{}ทสมํ.}

\csromanmatikaanchor{mtk.173}
\csromantocmarkref{subsection}{11. อปฺปมาทสุตฺต}{mtk.173}
\bookmark[dest={mtk.173},level=2]{11. อปฺปมาทสุตฺต}
\csromanheader{2}{11. อปฺปมาทสุตฺต}
\proseitem{53}{อถ โข อญฺญตโร พฺราหฺมโณ เยน ภควา เตนุปสงฺกมิ, อุปสงฺกมิตฺวา ภควตา สทฺธิํ สมฺโมทิ, สมฺโมทนียํ กถํ}

\vspace{\tipitakaparskip}
\csromancenter{ฉกฺกนิปาตปาฬิ สารณียํ วีติสาเรตฺวา เอกมนฺตํ นิสีทิ, เอกมนฺตํ นิสินฺโน โข โส พฺราหฺมโณ ภควนฺตํ เอตทโวจ\csromandash{}}
\prose{\csromanfolio{320}อตฺถิ นุ โข โภ โคตม เอโก ธมฺโม ภาวิโต พหุลีกโต, โย อุโภ อตฺเถ สมธิคฺคยฺห ติฏฺฐติ ทิฏฺฐธมฺมิกญฺเจว อตฺถํ โย จ อตฺโถ สมฺปรายิโกติ.\csromanspacer{} อตฺถิ โข พฺราหฺมณ เอโก ธมฺโม ภาวิโต พหุลีกโต, โย อุโภ อตฺเถ สมธิคฺคยฺห ติฏฺฐติ ทิฏฺฐธมฺมิกญฺเจว อตฺถํ โย จ อตฺโถ สมฺปรายิโกติ.}

\prose{กตโม ปน โภ โคตม เอโก ธมฺโม ภาวิโต พหุลีกโต, โย อุโภ อตฺเถ สมธิคฺคยฺห ติฏฺฐติ ทิฏฺฐธมฺมิกญฺเจว อตฺถํ โย จ อตฺโถ สมฺปรายิโกติ.\csromanspacer{} อปฺปมาโท โข พฺราหฺมณ เอโก ธมฺโม ภาวิโต พหุลีกโต อุโภ อตฺเถ สมธิคฺคยฺห ติฏฺฐติ ทิฏฺฐธมฺมิกญฺเจว อตฺถํ โย จ อตฺโถ สมฺปรายิโก.}

\prose{เสยฺยถาปิ พฺราหฺมณ ยานิ กานิจิ ชงฺคลานํ\footnote{ชงฺคมานํ (สี, อิ) อํ 3. 273; ม 1. 242 ปิฏฺเฐสุปิ.} ปาณานํ ปทชาตานิ, สพฺพานิ ตานิ หตฺถิปเท สโมธานํ คจฺฉนฺติ, หตฺถิปทํ เตสํ อคฺคมกฺขายติ ยทิทํ มหนฺตตฺเตน.\csromanspacer{} เอวเมวํ โข พฺราหฺมณ อปฺปมาโท เอโก ธมฺโม ภาวิโต พหุลีกโต อุโภ อตฺเถ สมธิคฺคยฺห ติฏฺฐติ ทิฏฺฐธมฺมิกญฺเจว อตฺถํ โย จ อตฺโถ สมฺปรายิโก.}

\prose{เสยฺยถาปิ พฺราหฺมณ กูฏาคารสฺส ยา กาจิ โคปานสิโย, สพฺพา ตา กูฏงฺคมา กูฏนินฺนา กูฏสโมสรณา, กูฏํ ตาสํ อคฺคมกฺขายติ.\csromanspacer{} เอวเมวํ โข พฺราหฺมณ ฯเปฯ.}

\prose{เสยฺยถาปิ พฺราหฺมณ ปพฺพชลายโก ปพฺพชํ\footnote{พพฺพชลายโก พพฺพชํ (สี, อิ)} ลายิตฺวา อคฺเค คเหตฺวา โอธุนาติ นิธุนาติ นิจฺฉาเทติ.\csromanspacer{} เอวเมวํ โข พฺราหฺมณ ฯเปฯ.}

\prose{เสยฺยถาปิ พฺราหฺมณ อมฺพปิณฺฑิยา วณฺฏจฺฉินฺนาย ยานิ กานิจิ อมฺพานิ วณฺฏูปนิพนฺธนานิ, สพฺพานิ ตานิ ตทนฺวยานิ ภวนฺติ.\csromanspacer{} เอวเมวํ โข พฺราหฺมณ ฯเปฯ.}

\prose{เสยฺยาถาปิ พฺราหฺมณ เย เกจิ ขุทฺทราชาโน\footnote{กุฑฺฑราชาโน (สี, สฺยา-ฏฺฐ), กุทฺทราชาโน (อิ) อํ 3. 274.}, สพฺเพเต รญฺโญ จกฺกวตฺติสฺส อนุยนฺตา\footnote{อนุยุตฺตา (สี, สฺยา, กํ, อิ)} ภวนฺติ, ราชา เตสํ จกฺกวตฺตี อคฺคมกฺขายติ.\csromanspacer{} เอวเมวํ โข พฺราหฺมณ ฯเปฯ.}

\chapterheadpageread{5. ธมฺมิกวคฺค}
\prose{\csromanfolio{321}เสยฺยถาปิ พฺราหฺมณ ยา กาจิ ตารกรูปานํ ปภา, สพฺพา ตา จนฺทสฺส ปภาย กลํ นาคฺฆนฺติ โสฬสิํ, จนฺทปฺปภา ตาสํ อคฺคมกฺขายติ.\csromanspacer{} เอวเมวํ โข พฺราหฺมณ อปฺปมาโท เอโก ธมฺโม ภาวิโต พหุลีกโต อุโภ อตฺเถ สมธิคฺคยฺห ติฏฺฐติ ทิฏฺฐธมฺมิกญฺเจว อตฺถํ โย จ อตฺโถ สมฺปรายิโก.}

\prose{อยํ โข พฺราหฺมณ เอโก ธมฺโม ภาวิโต พหุลีกโต อุโภ อตฺเถ สมธิคฺคยฺห ติฏฺฐติ ทิฏฺฐธมฺมิกญฺเจว อตฺถํ โย จ อตฺโถ สมฺปรายิโกติ.}

\prose{อภิกฺกนฺตํ โภ โคตม, อภิกฺกนฺตํ โภ โคตม ฯเปฯ อุปาสกํ มํ ภวํ โคตโม ธาเรตุ อชฺชตคฺเค ปาณุเปตํ สรณํ คตนฺติ.\csromanspacer{}\csromanspacer{}เอกาทสมํ.}

\csromanmatikaanchor{mtk.174}
\csromantocmarkref{subsection}{12. ธมฺมิกสุตฺต}{mtk.174}
\bookmark[dest={mtk.174},level=2]{12. ธมฺมิกสุตฺต}
\csromantocmarknopage{chapter}{2. ทุติยปณฺณาสก}{mtk.175}
\bookmark[dest={mtk.175},level=0]{2. ทุติยปณฺณาสก}
\csromanheader{2}{12. ธมฺมิกสุตฺต}
\proseitem{54}{เอกํ สมยํ ภควา ราชคเห วิหรติ คิชฺฌกูเฏ ปพฺพเต.\csromanspacer{} เตน โข ปน สมเยน อายสฺมา ธมฺมิโก ชาติภูมิยํ อาวาสิโก โหติ สพฺพโส ชาติภูมิยํ สตฺตสุ อาวาเสสุ.\csromanspacer{} ตตฺร สุทํ อายสฺมา ธมฺมิโก อาคนฺตุเก ภิกฺขู อกฺโกสติ ปริภาสติ วิหิํสติ วิตุทติ โรเสติ วาจาย.\csromanspacer{} เต จ อาคนฺตุกา ภิกฺขู อายสฺมตา ธมฺมิเกน อกฺโกสิยมานา ปริภาสิยมานา วิเหสิยมานา วิตุทิยมานา โรสิยมานา วาจาย ปกฺกมนฺติ น สณฺฐนฺติ\footnote{น สณฺฐหนฺติ (สี)}, ริญฺจนฺติ อาวาสํ.}

\prose{อถ โข ชาติภูมกานํ\footnote{ชาติภูมิกานํ (สฺยา, อิ, ก)} อุปาสกานํ เอตทโหสิ “มยํ โข ภิกฺขุสํฆํ ปจฺจุปฏฺฐิตา จีวรปิณฺฑปาตเสนาสนคิลานปฺปจฺจย- เภสชฺชปริกฺขาเรน, อถ จ ปน อาคนฺตุกา ภิกฺขู ปกฺกมนฺติ น สณฺฐนฺติ, ริญฺจนฺติ อาวาสํ.\csromanspacer{} โก นุ โข เหตุ โก ปจฺจโย, เยน อาคนฺตุกา ภิกฺขู ปกฺกมนฺติ น สณฺฐนฺติ, ริญฺจนฺติ อาวาสนฺ”ติ.\csromanspacer{} อถ โข ชาติภูมกานํ อุปาสกานํ เอตทโหสิ “อยํ โข อายสฺมา ธมฺมิโก อาคนฺตุเก ภิกฺขู อกฺโกสติ ปริภาสติ วิหิํสติ วิตุทติ โรเสติ วาจาย, เต จ อาคนฺตุกา ภิกฺขู อายสฺมตา ธมฺมิเกน อกฺโกสิยมานา ปริภาสิยมานา วิเหสิยมานา วิตุทิยมานา โรสิยมานา วาจาย}

\vspace{\tipitakaparskip}
\csromancenter{ฉกฺกนิปาตปาฬิ ปกฺกมนฺติ น สณฺฐนฺติ, ริญฺจนฺติ อาวาสํ.\csromanspacer{} ยํนูน มยํ อายสฺมนฺตํ ธมฺมิกํ ปพฺพาเชยฺยามา”ติ.}
\prose{\csromanfolio{322}อถ โข ชาติภูมกา อุปาสกา เยน อายสฺมา ธมฺมิโก เตนุปสงฺกมิํสุ, อุปสงฺกมิตฺวา อายสฺมนฺตํ ธมฺมิกํ เอตทโวจุํ “ปกฺกมตุ ภนฺเต อายสฺมา ธมฺมิโก อิมมฺหา อาวาสา, อลํ เต อิธ วาเสนา”ติ.\csromanspacer{} อถ โข อายสฺมา ธมฺมิโก ตมฺหา อาวาสา อญฺญํ อาวาสํ อคมาสิ.\csromanspacer{} ตตฺรปิ สุทํ อายสฺมา ธมฺมิโก อาคนฺตุเก ภิกฺขู อกฺโกสติ ปริภาสติ วิหิํสติ วิตุทติ โรเสติ วาจาย, เต จ อาคนฺตุกา ภิกฺขู อายสฺมตา ธมฺมิเกน อกฺโกสิยมานา ปริภาสิยมานา วิเหสิยมานา วิตุทิยมานา โรสิยมานา วาจาย ปกฺกมนฺติ น สณฺฐนฺติ, ริญฺจนฺติ อาวาสํ.}

\prose{อถ โข ชาติภูมกานํ อุปาสกานํ เอตทโหสิ “มยํ โข ภิกฺขุสํฆํ ปจฺจุปฏฺฐิตา จีวรปิณฺฑปาตเสนาสนคิลานปฺปจฺจย- เภสชฺชปริกฺขาเรน, อถ จ ปน อาคนฺตุกา ภิกฺขู ปกฺกมนฺติ น สณฺฐนฺติ, ริญฺจนฺติ อาวาสํ.\csromanspacer{} โก นุ โข เหตุ โก ปจฺจโย, เยน อาคนฺตุกา ภิกฺขู ปกฺกมนฺติ น สณฺฐนฺติ, ริญฺจนฺติ อาวาสนฺติ.\csromanspacer{} อถ โข ชาติภูมกานํ อุปาสกานํ เอตทโหสิ “อยํ โข อายสฺมา ธมฺมิโก อาคนฺตุเก ภิกฺขู อกฺโกสติ ปริภาสติ วิหิํสติ วิตุทติ โรเสติ วาจาย, เต จ อาคนฺตุกา ภิกฺขู อายสฺมตา ธมฺมิเกน อกฺโกสิยมานา ปริภาสิยมานา วิเหสิยมานา วิตุทิยมานา โรสิยมานา วาจาย ปกฺกมนฺติ น สณฺฐนฺติ, ริญฺจนฺติ อาวาสํ.\csromanspacer{} ยํนูน มยํ อายสฺมนฺตํ ธมฺมิกํ ปพฺพาเชยฺยามา”ติ.}

\prose{อถ โข ชาติภูมกา อุปาสกา เยนายสฺมา ธมฺมิโก เตนุปสงฺกมิํสุ, อุปสงฺกมิตฺวา อายสฺมนฺตํ ธมฺมิกํ เอตทโวจุํ “ปกฺกมตุ ภนฺเต อายสฺมา ธมฺมิโก อิมมฺหาปิ อาวาสา, อลํ เต อิธ วาเสนา”ติ.\csromanspacer{} อถ โข อายสฺมา ธมฺมิโก ตมฺหาปิ อาวาสา อญฺญํ อาวาสํ อคมาสิ.\csromanspacer{} ตตฺรปิ สุทํ อายสฺมา ธมฺมิโก อาคนฺตุเก ภิกฺขู อกฺโกสติ ปริภาสติ วิหิํสติ วิตุทติ โรเสติ วาจาย, เต จ อาคนฺตุกา ภิกฺขู อายสฺมตา ธมฺมิเกน อกฺโกสิยมานา ปริภาสิยมานา วิเหสิยมานา วิตุทิยมานา โรสิยมานา วาจาย ปกฺกมนฺติ น สณฺฐนฺติ, ริญฺจนฺติ อาวาสํ.}

\chapterheadpageread{5. ธมฺมิกวคฺค}
\prose{\csromanfolio{323}อถ โข ชาติภูมกานํ อุปาสกานํ เอตทโหสิ “มยํ โข ภิกฺขุสํฆํ ปจฺจุปฏฺฐิตา จีวร ปิณฺฑปาต เสนาสน คิลานปฺปจฺจย เภสชฺชปริกฺขาเรน, อถ จ ปน อาคนฺตุกา ภิกฺขู ปกฺกมนฺติ น สณฺฐนฺติ, ริญฺจนฺติ อาวาสํ.\csromanspacer{} โก นุ โข เหตุ โก ปจฺจโย, เยน อาคนฺตุกา ภิกฺขู ปกฺกมนฺติ น สณฺฐนฺติ, ริญฺจนฺติ อาวาสนฺ”ติ.\csromanspacer{} อถ โข ชาติภูมกานํ อุปาสกานํ เอตทโหสิ “อยํ โข อายสฺมา ธมฺมิโก อาคนฺตุเก ภิกฺขู อกฺโกสติ ฯเปฯ.\csromanspacer{} ยํนูน มยํ อายสฺมนฺตํ ธมฺมิกํ ปพฺพาเชยฺยาม สพฺพโส ชาติภูมิยํ สตฺตหิ อาวาเสหี”ติ.\csromanspacer{} อถ โข ชาติภูมกา อุปาสกา เยนายสฺมา ธมฺมิโก เตนุปสงฺกมิํสุ, อุปสงฺกมิตฺวา อายสฺมนฺตํ ธมฺมิกํ เอตทโวจุํ “ปกฺกมตุ ภนฺเต อายสฺมา ธมฺมิโก สพฺพโส ชาติภูมิยํ สตฺตหิ อาวาเสหี”ติ.\csromanspacer{} อถ โข อายสฺมโต ธมฺมิกสฺส เอตทโหสิ “ปพฺพาชิโต โขมฺหิ ชาติภูมเกหิ อุปาสเกหิ สพฺพโส ชาติภูมิยํ สตฺตหิ อาวาเสหิ, กหํ นุ โข ทานิ คจฺฉามี”ติ.\csromanspacer{} อถ โข อายสฺมโต ธมฺมิกสฺส เอตทโหสิ “ยํนูนาหํ เยน ภควา เตนุปสงฺกเมยฺยนฺ”ติ.}

\prose{อถ โข อายสฺมา ธมฺมิโก ปตฺตจีวรมาทาย เยน ราชคหํ เตน ปกฺกามิ, อนุปุพฺเพน เยน ราชคหํ คิชฺฌกูโฏ ปพฺพโต เยน ภควา เตนุปสงฺกมิ, อุปสงฺกมิตฺวา ภควนฺตํ อภิวาเทตฺวา เอกมนฺตํ นิสีทิ, เอกมนฺตํ นิสินฺนํ โข อายสฺมนฺตํ ธมฺมิกํ ภควา เอตทโวจ “หนฺท กุโต นุ ตฺวํ พฺราหฺมณ ธมฺมิก อาคจฺฉสี”ติ.\csromanspacer{} ปพฺพาชิโต อหํ ภนฺเต ชาติภูมเกหิ อุปาสเกหิ สพฺพโส ชาติภูมิยํ สตฺตหิ อาวาเสหีติ.\csromanspacer{} อลํ พฺราหฺมณ ธมฺมิก, กิํ เต อิมินา, ยํ ตํ ตโต ตโต ปพฺพาเชนฺติ, โส ตฺวํ ตโต ตโต ปพฺพาชิโต มเมว สนฺติเก อาคจฺฉสิ.}

\prose{ภูตปุพฺพํ พฺราหฺมณ ธมฺมิก สามุทฺทิกา วาณิชา ตีรทสฺสิํ สกุณํ คเหตฺวา นาวาย สมุทฺทํ อชฺโฌคาหนฺติ, เต อตีรทกฺขิณิยา\footnote{อตีรทสฺสนิยา (สฺยา), อตีรทสฺสิยา (ก)} นาวาย ตีรทสฺสิํ สกุณํ มุญฺจนฺติ, โส คจฺฉเตว ปุรตฺถิมํ ทิสํ, คจฺฉติ ปจฺฉิมํ ทิสํ, คจฺฉติ อุตฺตรํ ทิสํ, คจฺฉติ ทกฺขิณํ ทิสํ, คจฺฉติ อุทฺธํ, คจฺฉติ อนุทิสํ.\csromanspacer{} สเจ โส สมนฺตา ตีรํ ปสฺสติ, ตถาคตโกว\footnote{ตถาคโต (ก) ที 1. 213 ปิฏฺเฐ ปสฺสิตพฺพํ.} โหติ.\csromanspacer{} สเจ ปน โส}

\vspace{\tipitakaparskip}
\csromancenter{ฉกฺกนิปาตปาฬิ สมนฺตา ตีรํ น ปสฺสติ, ตเมว นาวํ ปจฺจาคจฺฉติ.\csromanspacer{} เอวเมวํ โข พฺราหฺมณ ธมฺมิก ยํ ตํ ตโต ตโต ปพฺพาเชนฺติ, โส ตฺวํ ตโต ตโต ปพฺพาชิโต มเมว สนฺติเก อาคจฺฉสิ.}
\prose{\csromanfolio{324}ภูตปุพฺพํ พฺราหฺมณ ธมฺมิก รญฺโญ โกรพฺยสฺส สุปฺปติฏฺโฐ นาม นิโคฺรธราชา อโหสิ ปญฺจสาโข สีตจฺฉาโย มโนรโม.\csromanspacer{} สุปฺปติฏฺฐสฺส โข ปน พฺราหฺมณ ธมฺมิก นิโคฺรธราชสฺส ทฺวาทสโยชนานิ อภินิเวโส อโหสิ, ปญฺจ โยชนานิ มูลสนฺตานกานํ.\csromanspacer{} สุปฺปติฏฺฐสฺส โข ปน พฺราหฺมณ ธมฺมิก นิโคฺรธราชสฺส ตาว มหนฺตานิ ผลานิ อเหสุํ เสยฺยถาปิ นาม อาฬฺหกถาลิกา.\csromanspacer{} เอวมสฺส สาทูนิ ผลานิ อเหสุํ เสยฺยถาปิ นาม ขุทฺทํ มธุํ อเนลกํ.\csromanspacer{} สุปฺปติฏฺฐสฺส โข ปน พฺราหฺมณ ธมฺมิก นิโคฺรธราชสฺส เอกํ ขนฺธํ ราชา ปริภุญฺชติ สทฺธิํ อิตฺถาคาเรน, เอกํ ขนฺธํ พลกาโย ปริภุญฺชติ, เอกํ ขนฺธํ เนคมชานปทา ปริภุญฺชนฺติ, เอกํ ขนฺธํ สมณพฺราหฺมณา ปริภุญฺชนฺติ, เอกํ ขนฺธํ มิคา\footnote{มิคปกฺขิโน (สี, สฺยา, อิ)} ปริภุญฺชนฺติ.\csromanspacer{} สุปฺปติฏฺฐสฺส โข ปน พฺราหฺมณ ธมฺมิก นิโคฺรธราชสฺส น โกจิ ผลานิ รกฺขติ, น จ สุทํ\footnote{น จ ปุน (ก) 3. ยาว ปาปมนุสฺโส (สฺยา), ยาวตา ปาปมนุสฺโส (ก)} อญฺญมญฺญสฺส ผลานิ หิํสนฺติ.}

\prose{อถ โข พฺราหฺมณ ธมฺมิก อญฺญตโร ปุริโส สุปฺปติฏฺฐสฺส นิโคฺรธราชสฺส ยาวทตฺถํ ผลานิ ภกฺขิตฺวา สาขํ ภญฺชิตฺวา ปกฺกามิ.\csromanspacer{} อถ โข พฺราหฺมณ ธมฺมิก สุปฺปติฏฺเฐ นิโคฺรธราเช อธิวตฺถาย เทวตาย เอตทโหสิ “อจฺฉริยํ วต โภ, อพฺภุตํ วต โภ, ยาว ปาโป มนุสฺโส3, ยตฺร หิ นาม สุปฺปติฏฺฐสฺส นิโคฺรธราชสฺส ยาวทตฺถํ ผลานิ ภกฺขิตฺวา สาขํ ภญฺชิตฺวา ปกฺกมิสฺสติ, ยํนูน สุปฺปติฏฺโฐ นิโคฺรธราชา อายติํ ผลํ น ทเทยฺยา”ติ.\csromanspacer{} อถ โข พฺราหฺมณ ธมฺมิก สุปฺปติฏฺโฐ นิโคฺรธราชา อายติํ ผลํ น อทาสิ.}

\prose{อถ โข พฺราหฺมณ ธมฺมิก ราชา โกรโพฺย เยน สกฺโก เทวานมินฺโท เตนุปสงฺกมิ, อุปสงฺกมิตฺวา สกฺกํ เทวานมินฺทํ เอตทโวจ “ยคฺเฆ มาริส ชาเนยฺยาสิ, สุปฺปติฏฺโฐ นิโคฺรธราชา ผลํ น เทตี”ติ.\csromanspacer{} อถ โข พฺราหฺมณ ธมฺมิก สกฺโก เทวานมินฺโท ตถารูปํ อิทฺธาภิสงฺขารํ อภิสงฺขาสิ\footnote{อภิสงฺขาเรสิ (สฺยา, ก)}, ยถา ภุสา วาตวุฏฺฐิ อาคนฺตฺวา สุปฺปติฏฺฐํ นิโคฺรธราชํ}

\vspace{\tipitakaparskip}
\csromancenter{ธมฺมิกวคฺค ปวตฺเตสิ\footnote{ปาเตสิ (สี, อิ)}, อุมฺมูลมกาสิ.\csromanspacer{} อถ โข พฺราหฺมณ ธมฺมิก สุปฺปติฏฺเฐ นิโคฺรธราเช อธิวตฺถา เทวตา ทุกฺขี ทุมฺมนา อสฺสุมุขี รุทมานา เอกมนฺตํ อฏฺฐาสิ.}
\prosecont{\csromanfolio{325}อถ โข พฺราหฺมณ ธมฺมิก สกฺโก เทวานมินฺโท เยน สุปฺปติฏฺเฐ นิโคฺรธราเช อธิวตฺถา เทวตา เตนุปสงฺกมิ, อุปสงฺกมิตฺวา สุปฺปติฏฺเฐ นิโคฺรธราเช อธิวตฺถํ เทวตํ เอตทโวจ “กิํ นุ ตฺวํ เทวเต ทุกฺขี ทุมฺมนา อสฺสุมุขี รุทมานา เอกมนฺตํ ฐิตา”ติ.\csromanspacer{} ตถา หิ ปน เม มาริส ภุสา วาตวุฏฺฐิ อาคนฺตฺวา ภวนํ ปวตฺเตสิ, อุมฺมูลมกาสีติ.\csromanspacer{} อปิ นุ ตฺวํ เทวเต รุกฺขธมฺเม ฐิตาย ภุสา วาตวุฏฺฐิ อาคนฺตฺวา ภวนํ ปวตฺเตสิ, อุมฺมูลมกาสีติ.\csromanspacer{} กถํ ปน มาริส รุกฺโข รุกฺขธมฺเม ฐิโต โหตีติ.\csromanspacer{} อิธ เทวเต รุกฺขสฺส มูลํ มูลตฺถิกา หรนฺติ, ตจํ ตจตฺถิกา หรนฺติ, ปตฺตํ ปตฺตตฺถิกา หรนฺติ, ปุปฺผํ ปุปฺผตฺถิกา หรนฺติ, ผลํ ผลตฺถิกา หรนฺติ, น จ เตน เทวตาย อนตฺตมนตา วา อนภินนฺทิ\footnote{อนภิรทฺธิ (สี)} วา กรณียา.\csromanspacer{} เอวํ โข เทวเต รุกฺโข รุกฺขธมฺเม ฐิโต โหตีติ.\csromanspacer{} อฏฺฐิตาเยว โข เม มาริส รุกฺขธมฺเม ภุสา วาตวุฏฺฐิ อาคนฺตฺวา ภวนํ ปวตฺเตสิ, อุมฺมูลมกาสีติ.\csromanspacer{} สเจ โข ตฺวํ เทวเต รุกฺขธมฺเม ติฏฺเฐยฺยาสิ, สิยา\footnote{สิยาปิ (สี, อิ)} เต ภวนํ ยถา ปุเรติ.\csromanspacer{} ฐสฺสามหํ\footnote{ติฏฺเฐยฺยามหํ (สฺยา)} มาริส รุกฺขธมฺเม, โหตุ เม ภวนํ ยถา ปุเรติ. อถ โข พฺราหฺมณ ธมฺมิก สกฺโก เทวานมินฺโท ตถารูปํ อิทฺธาภิสงฺขารํ อภิสงฺขาสิ\footnote{อภิสงฺขาริ (สฺยา, ก)}, ยถา ภุสา วาตวุฏฺฐิ อาคนฺตฺวา สุปฺปติฏฺฐํ นิโคฺรธราชํ อุสฺสาเปสิ, สจฺฉวีนิ มูลานิ อเหสุํ.\csromanspacer{} เอวเมวํ โข พฺราหฺมณ ธมฺมิก อปิ นุ ตํ สมณธมฺเม ฐิตํ ชาติภูมกา อุปาสกา ปพฺพาเชสุํ สพฺพโส ชาติภูมิยํ สตฺตหิ อาวาเสหีติ.\csromanspacer{} กถํ ปน ภนฺเต สมโณ สมณธมฺเม ฐิโต โหตีติ.\csromanspacer{} อิธ พฺราหฺมณ ธมฺมิก สมโณ อกฺโกสนฺตํ น ปจฺจกฺโกสติ, โรสนฺตํ น ปฏิโรสติ, ภณฺฑนฺตํ น ปฏิภณฺฑติ.\csromanspacer{} เอวํ โข พฺราหฺมณ ธมฺมิก สมโณ สมณธมฺเม ฐิโต โหตีติ.\csromanspacer{} อฏฺฐิตํเยว มํ ภนฺเต สมณธมฺเม ชาติภูมกา อุปาสกา ปพฺพาเชสุํ สพฺพโส ชาติภูมิยํ สตฺตหิ อาวาเสหีติ.}

\gambhira{ฉกฺกนิปาตปาฬิ}
\prose{\csromanfolio{326}\csromansymbolfootnote{*}{อุปริ 476, 500 ปิฏฺเฐสุปิ.}ภูตปุพฺพํ พฺราหฺมณ ธมฺมิก สุเนตฺโต นาม สตฺถา อโหสิ ติตฺถกโร กาเมสุ วีตราโค, สุเนตฺตสฺส โข ปน พฺราหฺมณ ธมฺมิก สตฺถุโน อเนกานิ สาวกสตานิ อเหสุํ.\csromanspacer{} สุเนตฺโต สตฺถา สาวกานํ พฺรหฺมโลกสหพฺยตาย ธมฺมํ เทเสสิ.\csromanspacer{} เย โข ปน พฺราหฺมณ ธมฺมิก สุเนตฺตสฺส สตฺถุโน พฺรหฺมโลกสหพฺยตาย ธมฺมํ เทเสนฺตสฺส จิตฺตานิ น ปสาเทสุํ.\csromanspacer{} เต กายสฺส เภทา ปรํ มรณา อปายํ ทุคฺคติํ วินิปาตํ นิรยํ อุปปชฺชิํสุ.\csromanspacer{} เย โข ปน พฺราหฺมณ ธมฺมิก สุเนตฺตสฺส สตฺถุโน พฺรหฺมโลกสหพฺยตาย ธมฺมํ เทเสนฺตสฺส จิตฺตานิ ปสาเทสุํ, เต กายสฺส เภทา ปรํ มรณา สุคติํ สคฺคํ โลกํ อุปปชฺชิํสุ.}

\prose{ภูตปุพฺพํ พฺราหฺมณ ธมฺมิก มูคปกฺโข นาม สตฺถา อโหสิ ฯเปฯ.\csromanspacer{} อรเนมิ นาม สตฺถา อโหสิ.\csromanspacer{} กุทฺทาลโก นาม สตฺถา อโหสิ.\csromanspacer{} หตฺถิปาโล นาม สตฺถา อโหสิ.\csromanspacer{} โชติปาโล นาม สตฺถา อโหสิ ติตฺถกโร กาเมสุ วีตราโค โชติปาลสฺส โข ปน พฺราหฺมณ ธมฺมิก สตฺถุโน อเนกานิ สาวกสตานิ อเหสุํ.\csromanspacer{} โชติปาโล สตฺถา สาวกานํ พฺรหฺมโลกสหพฺยตาย ธมฺมํ เทเสสิ.\csromanspacer{} เย โข ปน พฺราหฺมณ ธมฺมิก โชติปาลสฺส สตฺถุโน พฺรหฺมโลกสหพฺยตาย ธมฺมํ เทเสนฺตสฺส จิตฺตานิ น ปสาเทสุํ, เต กายสฺส เภทา ปรํ มรณา อปายํ ทุคฺคติํ วินิปาตํ นิรยํ อุปปชฺชิํสุ.\csromanspacer{} เย โข ปน พฺราหฺมณ ธมฺมิก โชติปาลสฺส สตฺถุโน พฺรหฺมโลกสหพฺยตาย ธมฺมํ เทเสนฺตสฺส จิตฺตานิ ปสาเทสุํ, เต กายสฺส เภทา ปรํ มรณา สุคติํ สคฺคํ โลกํ อุปปชฺชิํสุ.}

\prose{ตํ กิํ มญฺญสิ พฺราหฺมณ ธมฺมิก, โย อิเม ฉ สตฺถาเร ติตฺถกเร กาเมสุ วีตราเค อเนกสตปริวาเร สสาวกสํเฆ ปทุฏฺฐจิตฺโต อกฺโกเสยฺย ปริภาเสยฺย, พหุํ โส อปุญฺญํ ปสเวยฺยาติ.\csromanspacer{} เอวํ ภนฺเต, โย โข พฺราหฺมณ ธมฺมิก อิเม ฉ สตฺถาเร ติตฺถกเร กาเมสุ วีตราเค อเนกสตปริวาเร สสาวกสํเฆ ปทุฏฺฐจิตฺโต อกฺโกเสยฺย ปริภาเสยฺย, พหุํ โส อปุญฺญํ ปสเวยฺย.\csromanspacer{} โย เอกํ ทิฏฺฐิสมฺปนฺนํ ปุคฺคลํ ปทุฏฺฐจิตฺโต อกฺโกสติ ปริภาสติ, อยํ ตโต พหุตรํ อปุญฺญํ ปสวติ.\csromanspacer{} ตํ กิสฺส เหตุ? นาหํ พฺราหฺมณ ธมฺมิก อิโต พหิทฺธา เอวรูปิํ ขนฺติํ\footnote{เอวรูปํ ขนฺตํ (สฺยา)} วทามิ, ยถา’มํ สพฺรหฺมจารีสุ.\csromanspacer{} ตสฺมาติห พฺราหฺมณ}

\vspace{\tipitakaparskip}
\csromancenter{ธมฺมิกวคฺค ธมฺมิก เอวํ สิกฺขิตพฺพํ “น โน สมสพฺรหฺมจารีสุ\footnote{น โน อามสพฺรหฺมจารีสุ (สฺยา), น โน สพฺรหฺมจารีสุ (สี, อิ)} จิตฺตานิ ปทุฏฺฐานิ ภวิสฺสนฺตี”ติ.\csromanspacer{} เอวํ หิ เต พฺราหฺมณ ธมฺมิก สิกฺขิตพฺพนฺติ.}
\vspace{0.5\baselineskip}
\csromangathameasure{สุเนตฺโต มูคปกฺโข จ, อรเนมิ จ พฺราหฺมโณ. \\ กุทฺทาลโก อหุ สตฺถา, หตฺถิปาโล จ มาณโว. \\ โชติปาโล จ โควินฺโท, อหุ สตฺตปุโรหิโต. \\ อหิํสกา อตีตํเส, ฉ สตฺถาโร ยสสฺสิโน. \\ นิรามคนฺธา กรุเณธิมุตฺตา, กามสํโยชนาติคา. \\ กามราคํ วิราเชตฺวา, พฺรหฺมโลกูปคา อหุํ. \\ อเหสุํ สาวกา เตสํ, อเนกานิ สตานิปิ. \\ นิรามคนฺธา กรุเณธิมุตฺตา, กามสํโยชนาติคา. \\ กามราคํ วิราเชตฺวา, พฺรหฺมโลกูปคา อหุํ5. \\ เยเต อิสี พาหิรเก, วีตราเค สมาหิเต. \\ ปทุฏฺฐมนสงฺกปฺโป, โย นโร ปริภาสติ. \\ พหุญฺจ โส ปสวติ, อปุญฺญํ ตาทิโส นโร. \\ โย เจกํ ทิฏฺฐิสมฺปนฺนํ, ภิกฺขุํ พุทฺธสฺส สาวกํ. \\ ปทุฏฺฐมนสงฺกปฺโป, โย นโร ปริภาสติ. \\ อยํ ตโต พหุตรํ, อปุญฺญํ ปสเว นโร. \\ น สาธุรูปํ อาสีเท, ทิฏฺฐิฏฺฐานปฺปหายินํ. \\ สตฺตโม ปุคฺคโล เอโส, อริยสํฆสฺส วุจฺจติ. \\ อวีตราโค กาเมสุ, ยสฺส ปญฺจินฺทฺริยา มุทู. \\ สทฺธา สติ จ วีริยํ, สมโถ จ วิปสฺสนา. \\ ตาทิสํ ภิกฺขุมาสชฺช, ปุพฺเพว อุปหญฺญติ. \\ อตฺตานํ อุปหนฺตฺวาน, ปจฺฉา อญฺญํ วิหิํสติ.}
\csromangathagroup{\csromangathastanza{\csromanfolio{327}สุเนตฺโต มูคปกฺโข จ, อรเนมิ จ พฺราหฺมโณ. \\ กุทฺทาลโก อหุ สตฺถา, หตฺถิปาโล จ มาณโว.}\csromangathastanza{โชติปาโล จ โควินฺโท, อหุ สตฺตปุโรหิโต. \\ อหิํสกา\footnote{อภิเสกา (สฺยา)} อตีตํเส, ฉ สตฺถาโร ยสสฺสิโน.}\csromangathastanza{นิรามคนฺธา กรุเณธิมุตฺตา\footnote{วิมุตฺตา (สี, สฺยา, อิ)}, กามสํโยชนาติคา\footnote{กามสํโยชนาติตา (สฺยา)}. \\ กามราคํ วิราเชตฺวา, พฺรหฺมโลกูปคา อหุํ\footnote{วิมุตฺตา (สี, สฺยา, อิ)}.}\csromangathastanza{อเหสุํ สาวกา เตสํ, อเนกานิ สตานิปิ. \\ นิรามคนฺธา กรุเณธิมุตฺตา, กามสํโยชนาติคา.}\csromangathastanza{กามราคํ วิราเชตฺวา, พฺรหฺมโลกูปคา อหุํ5. \\ เยเต อิสี พาหิรเก, วีตราเค สมาหิเต.}\csromangathastanza{ปทุฏฺฐมนสงฺกปฺโป, โย นโร ปริภาสติ. \\ พหุญฺจ โส ปสวติ, อปุญฺญํ ตาทิโส นโร.}\csromangathastanza{โย เจกํ ทิฏฺฐิสมฺปนฺนํ, ภิกฺขุํ พุทฺธสฺส สาวกํ. \\ ปทุฏฺฐมนสงฺกปฺโป, โย นโร ปริภาสติ.}\csromangathastanza{อยํ ตโต พหุตรํ, อปุญฺญํ ปสเว นโร. \\ น สาธุรูปํ อาสีเท, ทิฏฺฐิฏฺฐานปฺปหายินํ.}\csromangathastanza{สตฺตโม ปุคฺคโล เอโส, อริยสํฆสฺส วุจฺจติ. \\ อวีตราโค กาเมสุ, ยสฺส ปญฺจินฺทฺริยา มุทู.}\csromangathastanza{สทฺธา สติ จ วีริยํ, สมโถ จ วิปสฺสนา. \\ ตาทิสํ ภิกฺขุมาสชฺช, ปุพฺเพว อุปหญฺญติ. \\ อตฺตานํ อุปหนฺตฺวาน, ปจฺฉา อญฺญํ วิหิํสติ.}}

\gambhira{ฉกฺกนิปาต ปาฬิ}
\csromangathameasure{โย จ รกฺขติ อตฺตานํ, รกฺขิโต ตสฺส พาหิโร. \\ ตสฺมา รกฺเขยฺย อตฺตานํ, อกฺขโต ปณฺฑิโต สทาติ. ทฺวาทสมํ. ธมฺมิกวคฺโค ปญฺจโม.}
\csromangathagroup{\csromangathastanza{\csromanfolio{328}โย จ รกฺขติ อตฺตานํ, รกฺขิโต ตสฺส พาหิโร. \\ ตสฺมา รกฺเขยฺย อตฺตานํ, อกฺขโต ปณฺฑิโต สทาติ.\csromanspacer{} ทฺวาทสมํ.\csromanspacer{} ธมฺมิกวคฺโค ปญฺจโม.}}

\titlehead{\textbf{ตสฺสุทฺทานํ}}
\csromangathameasure{นาคมิคสาลา อิณํ, จุนฺทํ เทฺว สนฺทิฏฺฐิกา ทุเว. \\ เขม-อินฺทฺริย อานนฺท, ขตฺติยา อปฺปมาเทน ธมฺมิโกติ. \\ \textbf{ปฐโม ปณฺณาสโก สมตฺโต.}}
\csromangathagroup{\csromangathastanza{นาคมิคสาลา อิณํ, จุนฺทํ เทฺว สนฺทิฏฺฐิกา ทุเว. \\ เขม-อินฺทฺริย อานนฺท, ขตฺติยา อปฺปมาเทน ธมฺมิโกติ. \\ \textbf{ปฐโม ปณฺณาสโก สมตฺโต.}}}

\csromanheader{2}{2. ทุติยปณฺณาสก}
\chapterhead{6. มหาวคฺค}
\vspace{\tipitakaparskip}
\csromancenter{โสณสุตฺต 55. * เอวํ เม สุตํ\csromandash{}เอกํ สมยํ ภควา ราชคเห วิหรติ คิชฺฌกูเฏ ปพฺพเต.\csromanspacer{} เตน โข ปน สมเยน อายสฺมา โสโณ ราชคเห วิหรติ สีตวนสฺมิํ.\csromanspacer{} อถ โข อายสฺมโต โสณสฺส รโหคตสฺส ปฏิสลฺลีนสฺส เอวํ เจตโส ปริวิตกฺโก อุทปาทิ “เย โข เกจิ ภควโต สาวกา อารทฺธวีริยา วิหรนฺติ, อหํ เตสํ อญฺญตโร, อถ จ ปน เม น อนุปาทาย อาสเวหิ จิตฺตํ วิมุจฺจติ, สํวิชฺชนฺติ โข ปน เม กุเล โภคา, สกฺกา โภคา จ ภุญฺชิตุํ ปุญฺญานิ จ กาตุํ.\csromanspacer{} ยํนูนาหํ สิกฺขํ ปจฺจกฺขาย หีนายาวตฺติตฺวา โภเค จ ภุญฺเชยฺยํ, ปุญฺญานิ จ กเรยฺยนฺ”ติ.}
\csromanmatikaanchor{mtk.176}
\csromanmatikaanchor{mtk.177}
\csromantocmarknopage{section}{6. มหาวคฺค}{mtk.176}
\bookmark[dest={mtk.176},level=1]{6. มหาวคฺค}
\csromantocmarkref{subsection}{1. โสณสุตฺต}{mtk.177}
\bookmark[dest={mtk.177},level=2]{1. โสณสุตฺต}
\prose{\csromanfolio{329}อถ โข ภควา อายสฺมโต โสณสฺส เจตสา เจโตปริวิตกฺกมญฺญาย เสยฺยถาปิ นาม พลวา ปุริโส สมิญฺชิตํ วา พหํ ปสาเรยฺย, ปสาริตํ วา พาหํ สมิญฺเชยฺย.\csromanspacer{} เอวเมวํ โข คิชฺฌกูเฏ ปพฺพเต อนฺตรหิโต สีตวเน อายสฺมโต โสณสฺส สมฺมุเข ปาตุรโหสิ.\csromanspacer{} นิสีทิ ภควา ปญฺญตฺเต อาสเน.\csromanspacer{} อายสฺมาปิ โข โสโณ ภควนฺตํ อภิวาเทตฺวา เอกมนฺตํ นิสีทิ, เอกมนฺตํ นิสินฺนํ โข อายสฺมนฺตํ โสณํ ภควา เอตทโวจ\csromandash{}}

\prose{นนุ เต โสณ รโหคตสฺส ปฏิสลฺลีนสฺส เอวํ เจตโส ปริวิตกฺโก อุทปาทิ “เย โข เกจิ ภควโต สาวกา อารทฺธวีริยา วิหรนฺติ, อหํ เตสํ อญฺญตโร, อถ จ ปน เม น อนุปาทาย อาสเวหิ จิตฺตํ วิมุจฺจติ, สํวิชฺชนฺติ โข ปน เม กุเล โภคา, สกฺกา โภคา\footnote{โภเค (วิ 3. 268; สฺยา, ก)} จ ภุญฺชิตุํ ปุญฺญานิ จ กาตุํ.\csromanspacer{} ยํนูนาหํ สิกฺขํ ปจฺจกฺขาย หีนายาวตฺติตฺวา โภเค จ ภุญฺเชยฺยํ, ปุญฺญานิ จ กเรยฺยนฺ”ติ.\csromanspacer{} เอวํ ภนฺเต.}

\gambhira{ฉกฺกนิปาตปาฬิ}
\prose{\csromanfolio{330}ตํ กิํ มญฺญสิ โสณ, กุสโล ตฺวํ ปุพฺเพ อคาริยภูโต\footnote{อาคาริกภูโต (สฺยา), อคาริกภูโต (วิ 3. 269)} วีณาย ตนฺติสฺสเรติ.\csromanspacer{} เอวํ ภนฺเต.\csromanspacer{} ตํ กิํ มญฺญสิ โสณ, ยทา เต วีณาย ตนฺติโย อจฺจายตา โหนฺติ, อปิ นุ เต วีณา ตสฺมิํ สมเย สรวตี วา โหติ กมฺมญฺญา วาติ.\csromanspacer{} โน เหตํ ภนฺเต.}

\prose{ตํ กิํ มญฺญสิ โสณ, ยทา เต วีณาย ตนฺติโย อติสิถิลา โหนฺติ, อปิ นุ เต วีณา ตสฺมิํ สมเย สรวตี วา โหติ กมฺมญฺญา วาติ.\csromanspacer{} โน เหตํ ภนฺเต.}

\prose{ยทา ปน เต โสณ วีณาย ตนฺติโย น อจฺจายตา โหนฺติ, นาติสิถิลา สเม คุเณ ปติฏฺฐิตา, อปิ นุ เต วีณา ตสฺมิํ สมเย สรวตี วา โหติ กมฺมญฺญา วาติ.\csromanspacer{} เอวํ ภนฺเต.}

\prose{เอวเมวํ โข โสณ อจฺจารทฺธวีริยํ อุทฺธจฺจาย สํวตฺตติ, อติสิถิลวีริยํ โกสชฺชาย สํวตฺตติ.\csromanspacer{} ตสฺมาติห ตฺวํ โสณ วีริยสมถํ อธิฏฺฐห, อินฺทฺริยานญฺจ สมตํ ปฏิวิชฺฌ, ตตฺถ จ นิมิตฺตํ คณฺหาหีติ. “เอวํ ภนฺเต”ติ โข อายสฺมา โสโณ ภควโต ปจฺจสฺโสสิ.\csromanspacer{} อถ โข ภควา อายสฺมนฺตํ โสณํ อิมินา โอวาเทน โอวทิตฺวา เสยฺยถาปิ นาม พลวา ปุริโส สมิญฺชิตํ วา พาหํ ปสาเรยฺย, ปสาริตํ วา พาหํ สมิญฺเชยฺย.\csromanspacer{} เอวเมวํ โข สีตวเน อนฺตรหิโต คิชฺฌกูเฏ ปพฺพเต ปาตุรโหสิ.}

\prose{อถ โข อายสฺมา โสโณ อปเรน สมเยน วีริยสมถํ อธิฏฺฐาสิ, อินฺทฺริยานญฺจ สมตํ ปฏิวิชฺฌิ, ตตฺถ จ นิมิตฺตํ อคฺคเหสิ.\csromanspacer{} อถ โข อายสฺมา โสโณ เอโก วูปกฏฺโฐ อปฺปมตฺโต อาตาปี ปหิตตฺโต วิหรนฺโต นจิรสฺเสว ยสฺสตฺถาย กุลปุตฺตา สมฺมเทว อคารสฺมา อนาคาริยํ ปพฺพชนฺติ, ตทนุตฺตรํ พฺรหฺมจริยปริโยสานํ ทิฏฺเฐว ธมฺเม สยํ อภิญฺญา สจฺฉิกตฺวา อุปสมฺปชฺช วิหาสิ, “ขีณา ชาติ, วุสิตํ พฺรหฺมจริยํ, กตํ กรณียํ, นาปรํ อิตฺถตฺตายา”ติ อพฺภญฺญาสิ, อญฺญตโร จ ปนายสฺมา โสโณ อรหตํ อโหสิ.}

\prose{อถ โข อายสฺมโต โสณสฺส อรหตฺตปฺปตฺตสฺส เอตทโหสิ “ยํนูนาหํ เยน ภควา เตนุปสงฺกเมยฺยํ, อุปสงฺกมิตฺวา ภควโต สนฺติเก อญฺญํ พฺยากเรยฺยนฺ”ติ.\csromanspacer{} อถ โข อายสฺมา โสโณ เยน}

\vspace{\tipitakaparskip}
\csromancenter{มหาวคฺค ภควา เตนุปสงฺกมิ, อุปสงฺกมิตฺวา ภควนฺตํ อภิวาเทตฺวา เอกมนฺตํ นิสีทิ, เอกมนฺตํ นิสินฺโน โข อายสฺมา โสโณ ภควนฺตํ เอตทโวจ\csromandash{}}
\prose{\csromanfolio{331}โย โส ภนฺเต ภิกฺขุ อรหํ ขีณาสโว วุสิตวา กตกรณีโย โอหิตภาโร อนุปฺปตฺตสทตฺโถ ปริกฺขีณภวสํโยชโน สมฺมทญฺญาวิมุตฺโต, โส ฉ ฐานานิ อธิมุตฺโต โหติ.\csromanspacer{} เนกฺขมฺมาธิมุตฺโต โหติ, ปวิเวกาธิมุตฺโต โหติ, อพฺยาปชฺชาธิมุตฺโต\footnote{อพฺยาปชฺฌาธิมุตฺโต (ก) วิ 3. 270 ปิฏฺเฐ ปสฺสิตพฺพํ.} โหติ, ตณฺหากฺขยาธิมุตฺโต โหติ, อุปาทานกฺขยาธิมุตฺโต โหติ, อสมฺโมหาธิมุตฺโต โหติ.}

\prose{สิยา โข ปน ภนฺเต อิเธกจฺจสฺส อายสฺมโต เอวมสฺส “เกวลํ สทฺธามตฺตกํ นูน อยมายสฺมา นิสฺสาย เนกฺขมฺมาธิมุตฺโต”ติ.\csromanspacer{} น โข ปเนตํ ภนฺเต เอวํ ทฏฺฐพฺพํ.\csromanspacer{} ขีณาสโว ภนฺเต ภิกฺขุ วุสิตวา กตกรณีโย กรณียํ อตฺตโน อสมนุปสฺสนฺโต กตสฺส วา ปฏิจยํ ขยา ราคสฺส วีตราคตฺตา เนกฺขมฺมาธิมุตฺโต โหติ, ขยา โทสสฺส วีตโทสตฺตา เนกฺขมฺมาธิมุตฺโต โหติ, ขยา โมหสฺส วีตโมหตฺตา เนกฺขมฺมาธิมุตฺโต โหติ.}

\prose{สิยา โข ปน ภนฺเต อิเธกจฺจสฺส อายสฺมโต เอวมสฺส “ลาภสกฺการสิโลกํ นูน อยมายสฺมา นิกามยมาโน ปวิเวกาธิมุตฺโต”ติ.\csromanspacer{} น โข ปเนตํ ภนฺเต เอวํ ทฏฺฐพฺพํ.\csromanspacer{} ขีณาสโว ภนฺเต ภิกฺขุ วุสิตวา กตกรณีโย กรณียํ อตฺตโน อสมนุปสฺสนฺโต กตสฺส วา ปฏิจยํ ขยา ราคสฺส วีตราคตฺตา ปวิเวกาธิมุตฺโต โหติ, ขยา โทสสฺส วีตโทสตฺตา ปวิเวกาธิมุตฺโต โหติ, ขยา โมหสฺส วีตโมหตฺตา ปวิเวกาธิมุตฺโต โหติ.}

\prose{สิยา โข ปน ภนฺเต อิเธกจฺจสฺส อายสฺมโต เอวมสฺส “สีลพฺพตปรามาสํ นูน อยมายสฺมา สารโต ปจฺจาคจฺฉนฺโต อพฺยาปชฺชาธิมุตฺโต”ติ.\csromanspacer{} น โข ปเนตํ ภนฺเต เอวํ ทฏฺฐพฺพํ.\csromanspacer{} ขีณาสโว ภนฺเต ภิกฺขุ วุสิตวา กตกรณีโย กรณียํ อตฺตโน อสมนุปสฺสนฺโต กตสฺส วา ปฏิจยํ ขยา ราคสฺส วีตราคตฺตา อพฺยาปชฺชาธิมุตฺโต โหติ, ขยา โทสสฺส วีตโทสตฺตา อพฺยาปชฺชาธิมุตฺโต โหติ, ขยา โมหสฺส วีตโมหตฺตา อพฺยาปชฺชาธิมุตฺโต โหติ.}

\gambhira{ฉกฺกนิปาตปาฬิ}
\prose{\csromanfolio{332}ขยา ราคสฺส วีตราคตฺตา ตณฺหากฺขยาธิมุตฺโต โหติ, ขยา โทสสฺส วีตโทสตฺตา ตณฺหากฺขยาธิมุตฺโต โหติ, ขยา โมหสฺส วีตโมหตฺตา ตณฺหากฺขยาธิมุตฺโต โหติ.}

\prose{ขยา ราคสฺส วีตราคตฺตา อุปาทานกฺขยาธิมุตฺโต โหติ, ขยา โทสสฺส วีตโทสตฺตา อุปาทานกฺขยาธิมุตฺโต โหติ, ขยา โมหสฺส วีตโมหตฺตา อุปาทานกฺขยาธิมุตฺโต โหติ.}

\prose{ขยา ราคสฺส วีตราคตฺตา อสมฺโมหาธิมุตฺโต โหติ, ขยา โทสสฺส วีตโทสตฺตา อสมฺโมหาธิมุตฺโต โหติ, ขยา โมหสฺส วีตโมหตฺตา อสมฺโมหาธิมุตฺโต โหติ.}

\prose{เอวํ สมฺมา วิมุตฺตจิตฺตสฺส ภนฺเต ภิกฺขุโน ภุสา เจปิ จกฺขุวิญฺเญยฺยา รูปา จกฺขุสฺส อาปาถํ\footnote{อาปาตํ (ก)} อาคจฺฉนฺติ, เนวสฺส จิตฺตํ ปริยาทิยนฺติ, อมิสฺสีกตเมวสฺส จิตฺตํ โหติ ฐิตํ อาเนญฺชปฺปตฺตํ, วยญฺจสฺสานุปสฺสติ.\csromanspacer{} ภุสา เจปิ โสตวิญฺเญยฺยา สทฺทา.\csromanspacer{} ฆานวิญฺเญยฺยา คนฺธา.\csromanspacer{} ชิวฺหาวิญฺเญยฺยา รสา.\csromanspacer{} กายวิญฺเญยฺยา โผฏฺฐพฺพา.\csromanspacer{} มโนวิญฺเญยฺยา ธมฺมา มนสฺส อาปาถํ อาคจฺฉนฺติ, เนวสฺส จิตฺตํ ปริยาทิยนฺติ, อมิสฺสีกตเมวสฺส จิตฺตํ โหติ ฐิตํ อาเนญฺชปฺปตฺตํ, วยญฺจสฺสานุปสฺสติ.\csromanspacer{} เสยฺยถาปิ ภนฺเต เสโล ปพฺพโต อจฺฉิทฺโท อสุสิโร เอกคฺฆโน, อถ ปุรตฺถิมาย เจปิ ทิสาย อาคจฺเฉยฺย ภุสา วาตวุฏฺฐิ, เนว นํ สงฺกมฺเปยฺย น สมฺปกมฺเปยฺย น สมฺปเวเธยฺย.\csromanspacer{} อถ ปจฺฉิมาย เจปิ ทิสาย อาคจฺเฉยฺย ภุสา วาตวุฏฺฐิ.\csromanspacer{} อถ อุตฺตราย เจปิ ทิสาย อาคจฺเฉยฺย ภุสา วาตวุฏฺฐิ.\csromanspacer{} อถ ทกฺขิณาย เจปิ ทิสาย อาคจฺเฉยฺย ภุสา วาตวุฏฺฐิ, เนว นํ สงฺกมฺเปยฺย น สมฺปกมฺเปยฺย น สมฺปเวเธยฺย.\csromanspacer{} เอวเมวํ โข ภนฺเต เอวํ สมฺมาวิมุตฺตจิตฺตสฺส ภิกฺขุโน ภุสา เจปิ จกฺขุวิญฺเญยฺยา รูปา จกฺขุสฺส อาปาถํ อาคจฺฉนฺติ, เนวสฺส จิตฺตํ ปริยาทิยนฺติ, อมิสฺสีกตเมวสฺส จิตฺตํ โหติ, ฐิตํ อาเนญฺชปฺปตฺตํ, วยญฺจสฺสานุปสฺสติ.\csromanspacer{} ภุสา เจปิ โสตวิญฺเญยฺยา สทฺทา.\csromanspacer{} ฆานวิญฺเญยฺยา คนฺธา.\csromanspacer{} ชิวฺหาวิญฺเญยฺยา รสา.\csromanspacer{} กายวิญฺเญยฺยา โผฏฺฐพฺพา.\csromanspacer{} มโนวิญฺเญยฺยา ธมฺมา มนสฺส อาปาถํ อาคจฺฉนฺติ, เนวสฺส จิตฺตํ ปริยาทิยนฺติ, อมิสฺสีกตเมวสฺส จิตฺตํ โหติ, ฐิตํ อาเนญฺชปฺปตฺตํ วยญฺจสฺสานุปสฺสตีติ.}

\chapterheadpageread{6. มหาวคฺค}
\csromangathameasure{เนกฺขมฺมํ อธิมุตฺตสฺส, ปวิเวกญฺจ เจตโส. \\ อพฺยาปชฺชาธิมุตฺตสฺส, อุปาทานกฺขยสฺส จ. \\ ตณฺหากฺขยาธิมุตฺตสฺส, อสมฺโมหญฺจ เจตโส. \\ ทิสฺวา อายตนุปฺปาทํ, สมฺมา จิตฺตํ วิมุจฺจติ. \\ ตสฺส สมฺมา วิมุตฺตสฺส, สนฺตจิตฺตสฺส ภิกฺขุโน. \\ กตสฺส ปฏิจโย นตฺถิ, กรณียํ น วิชฺชติ. \\ เสโล ยถา เอกคฺฆโน, วาเตน น สมีรติ. \\ เอวํ รูปา รสา สทฺทา, คนฺธา ผสฺสา จ เกวลา. \\ อิฏฺฐา ธมฺมา อนิฏฺฐา จ, นปฺปเวเธนฺติ ตาทิโน. \\ ฐิตํ จิตฺตํ วิปฺปมุตฺตํ, วยญฺจสฺสานุปสฺสตีติ.ปฐมํ.}
\csromangathagroup{\csromangathastanza{\csromanfolio{333}เนกฺขมฺมํ อธิมุตฺตสฺส, ปวิเวกญฺจ เจตโส. \\ อพฺยาปชฺชาธิมุตฺตสฺส, อุปาทานกฺขยสฺส จ.}\csromangathastanza{ตณฺหากฺขยาธิมุตฺตสฺส, อสมฺโมหญฺจ เจตโส. \\ ทิสฺวา อายตนุปฺปาทํ, สมฺมา จิตฺตํ วิมุจฺจติ.}\csromangathastanza{ตสฺส สมฺมา วิมุตฺตสฺส, สนฺตจิตฺตสฺส ภิกฺขุโน. \\ กตสฺส ปฏิจโย นตฺถิ, กรณียํ น วิชฺชติ.}\csromangathastanza{เสโล ยถา เอกคฺฆโน, วาเตน น สมีรติ. \\ เอวํ รูปา รสา สทฺทา, คนฺธา ผสฺสา จ เกวลา.}\csromangathastanza{อิฏฺฐา ธมฺมา อนิฏฺฐา จ, นปฺปเวเธนฺติ ตาทิโน. \\ ฐิตํ จิตฺตํ วิปฺปมุตฺตํ\footnote{วิมุตฺตญฺจ (ก) วิ 3. 271; อภิ 4. 77 ปิฏฺเฐสุปิ.}, วยญฺจสฺสานุปสฺสตีติ.\csromanspacer{}\csromanspacer{}ปฐมํ.}}

\csromanmatikaanchor{mtk.178}
\csromantocmarkref{subsection}{2. ผคฺคุนสุตฺต}{mtk.178}
\bookmark[dest={mtk.178},level=2]{2. ผคฺคุนสุตฺต}
\csromanheader{2}{2. ผคฺคุนสุตฺต}
\proseitem{56}{เตน โข ปน สมเยน อายสฺมา ผคฺคุโน\footnote{เผคฺคุโน (ก), ผคฺคุโณ (สี, สฺยา)} อาพาธิโก โหติ ทุกฺขิโต พาฬฺหคิลาโน.\csromanspacer{} อถ โข อายสฺมา อานนฺโท เยน ภควา เตนุปสงฺกมิ, อุปสงฺกมิตฺวา ภควนฺตํ อภิวาเทตฺวา เอกมนฺตํ นิสีทิ, เอกมนฺตํ นิสินฺโน โข อายสฺมา อานนฺโท ภควนฺตํ เอตทโวจ “อายสฺมา ภนฺเต ผคฺคุโน อาพาธิโก ทุกฺขิโต พาฬฺหคิลาโน, สาธุ ภนฺเต ภควา เยนายสฺมา ผคฺคุโน เตนุปสงฺกมตุ อนุกมฺปํ อุปาทายา”ติ.\csromanspacer{} อธิวาเสสิ ภควา ตุณฺหีภาเวน.\csromanspacer{} อถ โข ภควา สายนฺหสมยํ ปฏิสลฺลานา วุฏฺฐิโต เยนายสฺมา ผคฺคุโน เตนุปสงฺกมิ.\csromanspacer{} อทฺทสา โข อายสฺมา ผคฺคุโน ภควนฺตํ ทูรโตว อาคจฺฉนฺตํ, ทิสฺวาน มญฺจเก สมโธสิ\footnote{สมญฺโจปิ (สี, สฺยา, อิ), สํ + ธู + อี = สมโธสิ.}.\csromanspacer{} อถ โข ภควา อายสฺมนฺตํ ผคฺคุนํ เอตทโวจ “อลํ ผคฺคุน มา ตฺวํ มญฺจเก สมโธสิ, สนฺติมานิ อาสนานิ ปเรหิ ปญฺญตฺตานิ, ตตฺถาหํ นิสีทิสฺสามี”ติ.\csromanspacer{} นิสีทิ ภควา ปญฺญตฺเต อาสเน, นิสชฺช โข ภควา อายสฺมนฺตํ ผคฺคุนํ เอตทโวจ\csromandash{}}

\prose{กจฺจิ เต ผคฺคุน ขมนียํ, กจฺจิ ยาปนียํ, กจฺจิ เต ทุกฺขา เวทนา ปฏิกฺกมนฺติ โน อภิกฺกมนฺติ, ปฏิกฺกโมสานํ ปญฺญายติ โน อภิกฺกโมติ.\csromanspacer{} น เม}

\vspace{\tipitakaparskip}
\csromancenter{ฉกฺกนิปาตปาฬิ ภนฺเต ขมนียํ น ยาปนียํ, พาฬฺหา เม ทุกฺขา เวทนา อภิกฺกมนฺติ โน ปฏิกฺกมนฺติ, อภิกฺกโมสานํ ปญฺญายติ โน ปฏิกฺกโม.}
\prose{\csromanfolio{334}เสยฺยถาปิ ภนฺเต พลวา ปุริโส ติเณฺหน สิขเรน มุทฺธนิ\footnote{มุทฺธานํ (สี, อิ)} อภิมตฺเถยฺย\footnote{อภิมฏฺเฐยฺย (สฺยา)}.\csromanspacer{} เอวเมวํ โข เม ภนฺเต อธิมตฺตา วาตา มุทฺธนิ1 อูหนนฺติ\footnote{หนนฺติ (สี, อิ), โอหนนฺติ (สฺยา)}, น เม ภนฺเต ขมนียํ น ยาปนียํ, พาฬฺหา เม ทุกฺขา เวทนา อภิกฺกมนฺติ โน ปฏิกฺกมนฺติ, อภิกฺกโมสานํ ปญฺญายติ โน ปฏิกฺกโม.}

\prose{เสยฺยถาปิ ภนฺเต พลวา ปุริโส ทเฬฺหน วรตฺตกฺขณฺเฑน สีสเวฐนํ ทเทยฺย.\csromanspacer{} เอวเมวํ โข เม ภนฺเต อธิมตฺตา สีเส สีสเวทนา, น เม ภนฺเต ขมนียํ น ยาปนียํ, พาฬฺหา เม ทุกฺขา เวทนา อภิกฺกมนฺติ โน ปฏิกฺกมนฺติ, อภิกฺกโมสานํ ปญฺญายติ โน ปฏิกฺกโม.}

\prose{เสยฺยถาปิ ภนฺเต ทกฺโข โคฆาตโก วา โคฆาตกนฺเตวาสี วา ติเณฺหน โควิกนฺตเนน กุจฺฉิํ ปริกนฺเตยฺย.\csromanspacer{} เอวเมวํ โข เม ภนฺเต อธิมตฺตา วาตา กุจฺฉิํ ปริกนฺตนฺติ.\csromanspacer{} น เม ภนฺเต ขมนียํ น ยาปนียํ, พาฬฺหา เม ทุกฺขา เวทนา อภิกฺกมนฺติ โน ปฏิกฺกมนฺติ, อภิกฺกโมสานํ ปญฺญายติ โน ปฏิกฺกโม.}

\prose{เสยฺยถาปิ ภนฺเต เทฺว พลวนฺโต ปุริสา ทุพฺพลตรํ ปุริสํ นานาพาหาสุ คเหตฺวา องฺคารกาสุยา สนฺตาเปยฺยุํ สมฺปริตาเปยฺยุํ.\csromanspacer{} เอวเมวํ โข เม ภนฺเต อธิมตฺโต กายสฺมิํ ฑาโห, น เม ภนฺเต ขมนียํ น ยาปนียํ, พาฬฺหา เม ทุกฺขา เวทนา อภิกฺกมนฺติ โน ปฏิกฺกมนฺติ, อภิกฺกโมสานํ ปญฺญายติ โน ปฏิกฺกโมติ.\csromanspacer{} อถ โข ภควา อายสฺมนฺตํ ผคฺคุนํ ธมฺมิยา กถาย สนฺทสฺเสตฺวา สมาทเปตฺวา สมุตฺเตเชตฺวา สมฺปหํเสตฺวา อุฏฺฐายาสนา ปกฺกามิ.}

\prose{อถ โข อายสฺมา ผคฺคุโน อจิรปกฺกนฺตสฺส ภควโต กาลมกาสิ, ตมฺหิ จสฺส สมเย มรณกาเล อินฺทฺริยานิ วิปฺปสีทิํสุ.\csromanspacer{} อถ โข อายสฺมา อานนฺโท เยน ภควา เตนุปสงฺกมิ, อุปสงฺกมิตฺวา ภควนฺตํ อภิวาเทตฺวา เอกมนฺตํ นิสีทิ, เอกมนฺตํ นิสินฺโน โข อายสฺมา อานนฺโท ภควนฺตํ เอตทโวจ “อายสฺมา ภนฺเต ผคฺคุโน อจิรปกฺกนฺตสฺส}

\vspace{\tipitakaparskip}
\csromancenter{มหาวคฺค ภควโต กาลมกาสิ, ตมฺหิ จสฺส สมเย มรณกาเล อินฺทฺริยานิ วิปฺปสีทิํสู”ติ.}
\prose{\csromanfolio{335}กิํ หานนฺท ผคฺคุนสฺส\footnote{เผคฺคุนสฺส อานนฺท (ก)} ภิกฺขุโน อินฺทฺริยานิ น วิปฺปสีทิสฺสนฺติ.\csromanspacer{} ผคฺคุนสฺส อานนฺท ภิกฺขุโน ปญฺจหิ โอรมฺภาคิเยหิ สํโยชเนหิ จิตฺตํ อวิมุตฺตํ อโหสิ.\csromanspacer{} ตสฺส ตํ ธมฺมเทสนํ สุตฺวา ปญฺจหิ โอรมฺภาคิเยหิ สํโยชเนหิ จิตฺตํ วิมุตฺตํ.}

\prose{ฉยิเม อานนฺท อานิสํสา กาเลน ธมฺมสฺสวเน\footnote{ธมฺมสวเณ (สี)} กาเลน อตฺถุปปริกฺขาย.\csromanspacer{} กตเม ฉ? อิธานนฺท ภิกฺขุโน ปญฺจหิ โอรมฺภาคิเยหิ สํโยชเนหิ จิตฺตํ อวิมุตฺตํ โหติ, โส ตมฺหิ สมเย มรณกาเล ลภติ ตถาคตํ ทสฺสนาย.\csromanspacer{} ตสฺส ตถาคโต ธมฺมํ เทเสติ อาทิกลฺยาณํ มชฺเฌกลฺยาณํ ปริโยสานกลฺยาณํ สาตฺถํ สพฺยญฺชนํ, เกวลปริปุณฺณํ ปริสุทฺธํ พฺรหฺมจริยํ ปกาเสติ.\csromanspacer{} ตสฺส ตํ ธมฺมเทสนํ สุตฺวา ปญฺจหิ โอรมฺภาคิเยหิ สํโยชเนหิ จิตฺตํ วิมุจฺจติ.\csromanspacer{} อยํ อานนฺท ปฐโม อานิสํโส กาเลน ธมฺมสฺสวเน.}

\prose{ปุน จปรํ อานนฺท ภิกฺขุโน ปญฺจหิ โอรมฺภาคิเยหิ สํโยชเนหิ จิตฺตํ อวิมุตฺตํ โหติ, โส ตมฺหิ สมเย มรณกาเล น เหว โข\footnote{โน จ โข (ก)} ลภติ ตถาคตํ ทสฺสนาย, อปิ จ โข ตถาคตสาวกํ ลภติ ทสฺสนาย.\csromanspacer{} ตสฺส ตถาคตสาวโก ธมฺมํ เทเสติ อาทิกลฺยาณํ มชฺเฌกลฺยาณํ ปริโยสานกลฺยาณํ สาตฺถํ สพฺยญฺชนํ, เกวลปริปุณฺณํ ปริสุทฺธํ พฺรหฺมจริยํ ปกาเสติ.\csromanspacer{} ตสฺส ตํ ธมฺมเทสนํ สุตฺวา ปญฺจหิ โอรมฺภาคิเยหิ สํโยชเนหิ จิตฺตํ วิมุจฺจติ.\csromanspacer{} อยํ อานนฺท ทุติโย อานิสํโส กาเลน ธมฺมสฺสวเน.}

\prose{ปุน จปรํ อานนฺท ภิกฺขุโน ปญฺจหิ โอรมฺภาคิเยหิ สํโยชเนหิ จิตฺตํ อวิมุตฺตํ โหติ, โส ตมฺหิ สมเย มรณกาเล น เหว โข ลภติ ตถาคตํ ทสฺสนาย, นปิ ตถาคตสาวกํ ลภติ ทสฺสนาย, อปิ จ โข ยถาสุตํ ยถาปริยตฺตํ ธมฺมํ เจตสา อนุวิตกฺเกติ อนุวิจาเรติ มนสานุเปกฺขติ.\csromanspacer{} ตสฺส ยถาสุตํ}

\vspace{\tipitakaparskip}
\csromancenter{ฉกฺกนิปาตปาฬิ ยถาปริยตฺตํ ธมฺมํ เจตสา อนุวิตกฺกยโต อนุวิจารยโต มนสานุเปกฺขโต ปญฺจหิ โอรมฺภาคิเยหิ สํโยชเนหิ จิตฺตํ วิมุจฺจติ.\csromanspacer{} อยํ อานนฺท ตติโย อานิสํโส กาเลน อตฺถุปปริกฺขาย.}
\prose{\csromanfolio{336}อิธานนฺท ภิกฺขุโน ปญฺจหิ โอรมฺภาคิเยหิ สํโยชเนหิ จิตฺตํ วิมุตฺตํ โหติ, อนุตฺตเร จ โข อุปธิสงฺขเย จิตฺตํ อวิมุตฺตํ โหติ, โส ตมฺหิ สมเย มรณกาเล ลภติ ตถาคตํ ทสฺสนาย.\csromanspacer{} ตสฺส ตถาคโต ธมฺมํ เทเสติ อาทิกลฺยาณํ มชฺเฌกลฺยาณํ ฯเปฯ พฺรหฺมจริยํ ปกาเสติ.\csromanspacer{} ตสฺส ตํ ธมฺมเทสนํ สุตฺวา อนุตฺตเร อุปธิสงฺขเย จิตฺตํ วิมุจฺจติ.\csromanspacer{} อยํ อานนฺท จตุตฺโถ อานิสํโส กาเลน ธมฺมสฺสวเน.}

\prose{ปุน จปรํ อานนฺท ภิกฺขุโน ปญฺจหิ โอรมฺภาคิเยหิ สํโยชเนหิ จิตฺตํ วิมุตฺตํ โหติ, อนุตฺตเร จ โข อุปธิสงฺขเย จิตฺตํ อวิมุตฺตํ โหติ, โส ตมฺหิ สมเย มรณกาเล น เหว โข ลภติ ตถาคตํ ทสฺสนาย, อปิ จ โข ตถาคตสาวกํ ลภติ ทสฺสนาย.\csromanspacer{} ตสฺส ตถาคตสาวโก ธมฺมํ เทเสติ อาทิกลฺยาณํ ฯเปฯ ปริสุทฺธํ พฺรหฺมจริยํ ปกาเสติ.\csromanspacer{} ตสฺส ตํ ธมฺมเทสนํ สุตฺวา อนุตฺตเร อุปธิสงฺขเย จิตฺตํ วิมุจฺจติ.\csromanspacer{} อยํ อานนฺท ปญฺจโม อานิสํโส กาเลน ธมฺมสฺสวเน.}

\prose{ปุน จปรํ อานนฺท ภิกฺขุโน ปญฺจหิ โอรมฺภาคิเยหิ สํโยชเนหิ จิตฺตํ วิมุตฺตํ โหติ, อนุตฺตเร จ โข อุปธิสงฺขเย จิตฺตํ อวิมุตฺตํ โหติ, โส ตมฺหิ สมเย มรณกาเล น เหว โข ลภติ ตถาคตํ ทสฺสนาย, นปิ ตถาคตสาวกํ ลภติ ทสฺสนาย, อปิ จ โข ยถาสุตํ ยถาปริยตฺตํ ธมฺมํ เจตสา อนุวิตกฺเกติ อนุวิจาเรติ มนสานุเปกฺขติ.\csromanspacer{} ตสฺส ยถาสุตํ ยถาปริยตฺตํ ธมฺมํ เจตสา อนุวิตกฺกยโต อนุวิจารยโต มนสานุเปกฺขโต อนุตฺตเร อุปธิสงฺขเย จิตฺตํ วิมุจฺจติ.\csromanspacer{} อยํ อานนฺท ฉฏฺโฐ อานิสํโส กาเลน อตฺถุปปริกฺขาย.\csromanspacer{}\csromanspacer{}อิเม โข อานนฺท ฉ อานิสํสา กาเลน ธมฺมสฺสวเน กาเลน อตฺถุปปริกฺขายาติ.\csromanspacer{}\csromanspacer{}ทุติยํ.}

\csromanmatikaanchor{mtk.179}
\csromantocmarkref{subsection}{3. ฉฬภิชาติสุตฺต}{mtk.179}
\bookmark[dest={mtk.179},level=2]{3. ฉฬภิชาติสุตฺต}
\csromanheader{2}{3. ฉฬภิชาติสุตฺต}
\proseitem{57}{เอกํ สมยํ ภควา ราชคเห วิหรติ คิชฺฌกูเฏ ปพฺพเต.\csromanspacer{} อถ โข อายสฺมา อานนฺโท เยน ภควา เตนุปสงฺกมิ, อุปสงฺกมิตฺวา}

\vspace{\tipitakaparskip}
\csromancenter{มหาวคฺค ภควนฺตํ อภิวาเทตฺวา เอกมนฺตํ นิสีทิ, เอกมนฺตํ นิสินฺโน โข อายสฺมา อานนฺโท ภควนฺตํ เอตทโวจ\csromandash{}ปูรเณน ภนฺเต กสฺสเปน ฉฬภิชาติโย ปญฺญตฺตา, กณฺหาภิชาติ ปญฺญตา, นีลาภิชาติ ปญฺญตฺตา, โลหิตาภิชาติ ปญฺญตฺตา, หลิทฺทาภิชาติ ปญฺญตฺตา, สุกฺกาภิชาติ ปญฺญตฺตา, ปรมสุกฺกาภิชาติ ปญฺญตฺตา.}
\prose{\csromanfolio{337}ตตฺริทํ ภนฺเต ปูรเณน กสฺสเปน กณฺหาภิชาติ ปญฺญตฺตา โอรพฺภิกา สูกริกา สากุณิกา มาควิกา ลุทฺทา มจฺฉฆาตกา โจรา โจรฆาตกา พนฺธนาคาริกา เย วา ปนญฺเญปิ เกจิ กุรุรกมฺมนฺตา.}

\prose{ตตฺริทํ ภนฺเต ปูรเณน กสฺสเปน นีลาภิชาติ ปญฺญตฺตา ภิกฺขู กณฺฏกวุตฺติกา เย วา ปนญฺเญปิ เกจิ กมฺมวาทา กฺริยวาทา.}

\prose{ตตฺริทํ ภนฺเต ปูรเณน กสฺสเปน โลหิตาภิชาติ ปญฺญตฺตา นิคณฺฐา เอกสาฏกา.}

\prose{ตตฺริทํ ภนฺเต ปูรเณน กสฺสเปน หลิทฺทาภิชาติ ปญฺญตฺตา คิหี โอทาตวสนา อเจลกสาวกา.}

\prose{ตตฺริทํ ภนฺเต ปูรเณน กสฺสเปน สุกฺกาภิชาติ ปญฺญตฺตา อาชีวกา อาชีวกินิโย.}

\prose{ตตฺริทํ ภนฺเต ปูรเณน กสฺสเปน ปรมสุกฺกาภิชาติ ปญฺญตฺตา นนฺโท วจฺโฉ, กิโส สํกิจฺโจ, มกฺขลิ โคสาโล.\csromanspacer{} ปูรเณน ภนฺเต กสฺสเปน อิมา ฉฬภิชาติโย ปญฺญตฺตาติ.}

\prose{กิํ ปนานนฺท ปูรณสฺส กสฺสปสฺส สพฺโพ โลโก เอตทพฺภนุชานาติ อิมา ฉฬภิชาติโย ปญฺญาเปตุนฺติ.\csromanspacer{} โน เหตํ ภนฺเต.\csromanspacer{} เสยฺยถาปิ อานนฺท ปุริโส ทลิทฺโท อสฺสโก อนาฬฺหิโก, ตสฺส อกามกสฺส พิลํ โอลคฺเคยฺยุํ, อิทํ เต อมฺโภ ปุริส มํสญฺจ ขาทิตพฺพํ, มูลญฺจ อนุปฺปทาตพฺพนฺติ.\csromanspacer{} เอวเมวํ โข อานนฺท ปูรเณน กสฺสเปน อปฺปฏิญฺญาย เอเตสํ สมณพฺราหฺมณานํ อิมา ฉฬภิชาติโย ปญฺญตฺตา, ยถา ตํ พาเลน อพฺยตฺเตน อเขตฺตญฺญุนา อกุสเลน.}

\prose{อหํ โข ปนานนฺท ฉฬภิชาติโย ปญฺญาเปมิ, ตํ สุณาหิ สาธุกํ มนสิ กโรหิ, ภาสิสฺสามีติ. “เอวํ ภนฺเต”ติ โข อายสฺมา อานนฺโท ภควโต ปจฺจสฺโสสิ.\csromanspacer{} ภควา เอตทโวจ\csromandash{}กตมา จานนฺท}

\vspace{\tipitakaparskip}
\csromancenter{ฉกฺกนิปาตปาฬิ ฉฬภิชาติโย, อิธานนฺท เอกจฺโจ กณฺหาภิชาติโย สมาโน กณฺหํ ธมฺมํ อภิชายติ, อิธ ปนานนฺท เอกจฺโจ กณฺหาภิชาติโย สมาโน สุกฺกํ ธมฺมํ อภิชายติ, อิธ ปนานนฺท เอกจฺโจ กณฺหาภิชาติโย สมาโน อกณฺหํ อสุกฺกํ นิพฺพานํ อภิชายติ, อิธ ปนานนฺท เอกจฺโจ สุกฺกาภิชาติโย สมาโน กณฺหํ ธมฺมํ อภิชายติ, อิธ ปนานนฺท เอกจฺโจ สุกฺกาภิชาติโย สมาโน สุกฺกํ ธมฺมํ อภิชายติ, อิธ ปนานนฺท เอกจฺโจ สุกฺกาภิชาติโย สมาโน อกณฺหํ อสุกฺกํ นิพฺพานํ อภิชายติ.}
\prose{\csromanfolio{338}กถญฺจานนฺท กณฺหาภิชาติโย สมาโน กณฺหํ ธมฺมํ อภิชายติ? อิธานนฺท เอกจฺโจ นีเจ กุเล ปจฺจาชาโต โหติ จณฺฑาลกุเล วา เนสาทกุเล วา เวนกุเล\footnote{เวณกุเล (สพฺพตฺถ)} วา รถการกุเล วา ปุกฺกุสกุเล วา ทลิทฺเท อปฺปนฺนปานโภชเน กสิรวุตฺติเก, ยตฺถ กสิเรน ฆาสจฺฉาโท ลพฺภติ, โส จ โหติ ทุพฺพณฺโณ ทุทฺทสิโก โอโกฏิมโก พวฺหาพาโธ กาโณ วา กุณี วา ขญฺโช วา ปกฺขหโต วา, น ลาภี อนฺนสฺส ปานสฺส วตฺถสฺส ยานสฺส มาลาคนฺธวิเลปนสฺส เสยฺยาวสถปทีเปยฺยสฺส.\csromanspacer{} โส กาเยน ทุจฺจริตํ จรติ, วาจาย ทุจฺจริตํ จรติ, มนสา ทุจฺจริตํ จรติ, โส กาเยน ทุจฺจริตํ จริตฺวา วาจาย ทุจฺจริตํ จริตฺวา มนสา ทุจฺจริตํ จริตฺวา กายสฺส เภทา ปรํ มรณา อปายํ ทุคฺคติํ วินิปาตํ นิรยํ อุปปชฺชติ.\csromanspacer{} เอวํ โข อานนฺท กณฺหาภิชาติโย สมาโน กณฺหํ ธมฺมํ อภิชายติ.}

\prose{กถญฺจานนฺท กณฺหาภิชาติโย สมาโน สุกฺกํ ธมฺมํ อภิชายติ? อิธานนฺท เอกจฺโจ นีเจ กุเล ปจฺจาชาโต โหติ จณฺฑาลกุเล วา ฯเปฯ เสยฺยาวสถปทีเปยฺยสฺส.\csromanspacer{} โส กาเยน สุจริตํ จรติ, วาจาย สุจริตํ จรติ, มนสา สุจริตํ จรติ, โส กาเยน สุจริตํ จริตฺวา วาจาย สุจริตํ จริตฺวา มนสา สุจริตํ จริตฺวา กายสฺส เภทา ปรํ มรณา สุคติํ สคฺคํ โลกํ อุปปชฺชติ.\csromanspacer{} เอวํ โข อานนฺท กณฺหาภิชาติโย สมาโน สุกฺกํ ธมฺมํ อภิชายติ.}

\prose{กถญฺจานนฺท กณฺหาภิชาติโย สมาโน อกณฺหํ อสุกฺกํ นิพฺพานํ อภิชายติ? อิธานนฺท เอกจฺโจ นีเจ กุเล ปจฺจาชาโต โหติ}

\vspace{\tipitakaparskip}
\csromancenter{มหาวคฺค จณฺฑาลกุเล วา ฯเปฯ โส จ โหติ ทุพฺพณฺโณ ทุทฺทสิโก โอโกฏิมโก.\csromanspacer{} โส เกสมสฺสุํ โอหาเรตฺวา กาสายานิ วตฺถานิ อจฺฉาเทตฺวา อคารสฺมา อนคาริยํ ปพฺพชติ, โส เอวํ ปพฺพชิโต สมาโน ปญฺจ นีวรเณ ปหาย เจตโส อุปกฺกิเลเส ปญฺญาย ทุพฺพลีกรเณ จตูสุ สติปฏฺฐาเนสุ สุปฺปติฏฺฐิตจิตฺโต สตฺต โพชฺฌงฺเค ยถาภูตํ ภาเวตฺวา อกณฺหํ อสุกฺกํ นิพฺพานํ อภิชายติ.\csromanspacer{} เอวํ โข อานนฺท กณฺหาภิชาติโย สมาโน อกณฺหํ อสุกฺกํ นิพฺพานํ อภิชายติ.}
\prose{\csromanfolio{339}กถญฺจานนฺท สุกฺกาภิชาติโย สมาโน กณฺหํ ธมฺมํ อภิชายติ? อิธานนฺท เอกจฺโจ อุจฺเจ กุเล ปจฺจาชาโต โหติ ขตฺติยมหาสาลกุเล วา พฺราหฺมณมหาสาลกุเล วา คหปติมหาสาลกุเล วา อฑฺเฒ มหทฺธเน มหาโภเค ปหูตชาตรูปรชเต ปหูตวิตฺตูปกรเณ ปหูตธนธญฺเญ, โส จ โหติ อภิรูโป ทสฺสนีโย ปาสาทิโก ปรมาย วณฺณโปกฺขรตาย สมนฺนาคโต, ลาภี อนฺนสฺส ปานสฺส วตฺถสฺส ยานสฺส มาลาคนฺธวิเลปนสฺส เสยฺยาวสถปทีเปยฺยสฺส.\csromanspacer{} โส กาเยน ทุจฺจริตํ จรติ, วาจาย ทุจฺจริตํ จรติ, มนสา ทุจฺจริตํ จรติ, โส กาเยน ทุจฺจริตํ จริตฺวา วาจาย ทุจฺจริตํ จริตฺวา มนสา ทุจฺจริตํ จริตฺวา กายสฺส เภทา ปรํ มรณา อปายํ ทุคฺคติํ วินิปาตํ นิรยํ อุปปชฺชติ.\csromanspacer{} เอวํ โข อานนฺท สุกฺกาภิชาติโย สมาโน กณฺหํ ธมฺมํ อภิชายติ.}

\prose{กถญฺจานนฺท สุกฺกาภิชาติโย สมาโน สุกฺกํ ธมฺมํ อภิชายติ? อิธานนฺท เอกจฺโจ อุจฺเจ กุเล ปจฺจาชาโต โหติ ขตฺติยมหาสาลกุเล วา ฯเปฯ เสยฺยาวสถปทีเปยฺยสฺส.\csromanspacer{} โส กาเยน สุจริตํ จรติ, วาจาย สุจริตํ จรติ, มนสา สุจริตํ จรติ, โส กาเยน สุจริตํ จริตฺวา วาจาย สุจริตํ จริตฺวา มนสา สุจริตํ จริตฺวา กายสฺส เภทา ปรํ มรณา สุคติํ สคฺคํ โลกํ อุปปชฺชติ.\csromanspacer{} เอวํ โข อานนฺท สุกฺกาภิชาติโย สมาโน สุกฺกํ ธมฺมํ อภิชายติ.}

\prose{กถญฺจานนฺท สุกฺกาภิชาติโย สมาโน อกณฺหํ อสุกฺกํ นิพฺพานํ อภิชายติ? อิธานนฺท เอกจฺโจ อุจฺเจ กุเล ปจฺจาชาโต โหติ}

\vspace{\tipitakaparskip}
\csromancenter{ฉกฺกนิปาตปาฬิ ขตฺติยมหาสาลกุเล วา พฺราหฺมณมหาสาลกุเล วา คหปติมหาสาลกุเล วา อฑฺเฒ มหทฺธเน มหาโภเค ปหูตชาตรูปรชเต ปหูตวิตฺตูปกรเณ ปหูตธนธญฺเญ, โส จ โหติ อภิรูโป ทสฺสนีโย ปาสาทิโก ปรมาย วณฺณโปกฺขร ตาย สมนฺนาคโต, ลาภี อนฺนสฺส ปานสฺส วตฺถสฺส ยานสฺส มาลาคนฺธวิเลปนสฺส เสยฺยาวสถปทีเปยฺยสฺส.\csromanspacer{} โส เกสมสฺสุํ โอหาเรตฺวา กาสายานิ วตฺถานิ อจฺฉาเทตฺวา อคารสฺมา อนคาริยํ ปพฺพชติ, โส เอวํ ปพฺพชิโต สมาโน ปญฺจ นีวรเณ ปหาย เจตโส อุปกฺกิเลเส ปญฺญาย ทุพฺพลีกรเณ จตูสุ สติปฏฺฐาเนสุ สุปฺปติฏฺฐิตจิตฺโต สตฺต โพชฺฌงฺเค ยถาภูตํ ภาเวตฺวา อกณฺหํ อสุกฺกํ นิพฺพานํ อภิชายติ.\csromanspacer{} เอวํ โข อานนฺท สุกฺกาภิชาติโย สมาโน อกณฺหํ อสุกฺกํ นิพฺพานํ อภิชายติ.\csromanspacer{} อิมา โข อานนฺท ฉฬภิชาติโยติ.\csromanspacer{}\csromanspacer{}ตติยํ.}
\csromanmatikaanchor{mtk.180}
\csromantocmarkref{subsection}{4. อาสวสุตฺต}{mtk.180}
\bookmark[dest={mtk.180},level=2]{4. อาสวสุตฺต}
\csromanheader{2}{4. อาสวสุตฺต}
\proseitem{58}{\csromanfolio{340}ฉหิ ภิกฺขเว ธมฺเมหิ สมนฺนาคโต ภิกฺขุ อาหุเนยฺโย โหติ ปาหุเนยฺโย ทกฺขิเณยฺโย อญฺชลิกรณีโย อนุตฺตรํ ปุญฺญกฺเขตฺตํ โลกสฺส.}

\prose{กตเมหิ ฉหิ? อิธ ภิกฺขเว ภิกฺขุโน เย อาสวา สํวรา ปหาตพฺพา, เต สํวเรน ปหีนา โหนฺติ.\csromanspacer{} เย อาสวา ปฏิเสวนา ปหาตพฺพา, เต ปฏิเสวนาย ปหีนา โหนฺติ.\csromanspacer{} เย อาสวา อธิวาสนา ปหาตพฺพา, เต อธิวาสนาย ปหีนา โหนฺติ.\csromanspacer{} เย อาสวา ปริวชฺชนา ปหาตพฺพา, เต ปริวชฺชนาย ปหีนา โหนฺติ.\csromanspacer{} เย อาสวา วิโนทนา ปหาตพฺพา, เต วิโนทนาย ปหีนา โหนฺติ.\csromanspacer{} เย อาสวา ภาวนา ปหาตพฺพา, เต ภาวนาย ปหีนา โหนฺติ.}

\prose{กตเม จ ภิกฺขเว อาสวา สํวรา ปหาตพฺพา? เย สํวเรน ปหีนา โหนฺติ, อิธ ภิกฺขเว ภิกฺขุ ปฏิสงฺขา โยนิโส จกฺขุนฺทฺริยสํวรสํวุโต วิหรติ.\csromanspacer{} ยํ หิสฺส ภิกฺขเว จกฺขุนฺทฺริยสํวรํ อสํวุตสฺส วิหรโต อุปฺปชฺเชยฺยุํ อาสวา วิฆาตปริฬาหา, จกฺขุนฺทฺริยสํวรํ สํวุตสฺส วิหรโต เอวํส เต อาสวา วิฆาตปริฬาหา น โหนฺติ.\csromanspacer{} ปฏิสงฺขา โยนิโส โสตินฺทฺริย ฯเปฯ.\csromanspacer{} ฆานินฺทฺริย.\csromanspacer{} ชิวฺหินฺทฺริย.\csromanspacer{} กายินฺทฺริย.}

\vspace{\tipitakaparskip}
\csromancenter{มหาวคฺค มนินฺทฺริยสํวรสํวุโต วิหรติ.\csromanspacer{} ยํ หิสฺส ภิกฺขเว มนินฺทฺริยสํวรํ อสํวุตสฺส วิหรโต อุปฺปชฺเชยฺยุํ อาสวา วิฆาตปริฬาหา, มนินฺทฺริยสํวรํ สํวุตสฺส วิหรโต เอวํส เต อาสวา วิฆาตปริฬาหา น โหนฺติ.\csromanspacer{} อิเม วุจฺจนฺติ ภิกฺขเว อาสวา สํวรา ปหาตพฺพา, เย สํวเรน ปหีนา โหนฺติ.}
\prose{\csromanfolio{341}กตเม จ ภิกฺขเว อาสวา ปฏิเสวนา ปหาตพฺพา? เย ปฏิเสวนาย ปหีนา โหนฺติ, อิธ ภิกฺขเว ภิกฺขุ ปฏิสงฺขา โยนิโส จีวรํ ปฏิเสวติ, ยาวเทว สีตสฺส ปฏิฆาตาย อุณฺหสฺส ปฏิฆาตาย ฑํส มกส วาตาตป สรีสป สมฺผสฺสานํ\footnote{qอํส...สิริํสปสมฺผสฺสานํ (สี, สฺยา, กํ, อิ) ม 1. 12 ปิฏฺเฐปิ.} ปฏิฆาตาย ยาวเทว หิริโกปีนปฏิจฺฉาทนตฺถํ.\csromanspacer{} ปฏิสงฺขา โยนิโส ปิณฺฑปาตํ ปฏิเสวติ, เนว ทวาย น มทาย น มณฺฑนาย น วิภูสนาย, ยาวเทว อิมสฺส กายสฺส ฐิติยา ยาปนาย วิหิํสูปรติยา พฺรหฺมจริยานุคฺคหาย, อิติ ปุราณญฺจ เวทนํ ปฏิหงฺขามิ, นวญฺจ เวทนํ น อุปฺปาเทสฺสามิ, ยาตฺรา จ เม ภวิสฺสติ อนวชฺชตา จ ผาสุวิหาโร จ\footnote{ผาสุวิหาโร จาติ (สี, สฺยา, กํ, อิ)}.\csromanspacer{} ปฏิสงฺขา โยนิโส เสนาสนํ ปฏิเสวติ, ยาวเทว สีตสฺส ปฏิฆาตาย อุณฺหสฺส ปฏิฆาตาย ฑํส มกส วาตาตป สรีสป สมฺผสฺสานํ ปฏิฆาตาย ยาวเทว อุตุปริสฺสยวิโนทนปฏิสลฺลานารามตฺถํ.\csromanspacer{} ปฏิสงฺขา โยนิโส คิลานปฺปจฺจยเภสชฺชปริกฺขารํ ปฏิเสวติ, ยาวเทว อุปฺปนฺนานํ เวยฺยาพาธิกานํ เวทนานํ ปฏิฆาตาย อพฺยาพชฺฌปรมตายาติ.\csromanspacer{} ยํ หิสฺส ภิกฺขเว อปฺปฏิเสวโต อุปฺปชฺเชยฺยุํ อาสวา วิฆาตปริฬาหา, ปฏิเสวโต เอวํส เต อาสวา วิฆาตปริฬาหา น โหนฺติ.\csromanspacer{} อิเม วุจฺจนฺติ ภิกฺขเว อาสวา ปฏิเสวนา ปหาตพฺพา, เย ปฏิเสวนาย ปหีนา โหนฺติ.}

\prose{กตเม จ ภิกฺขเว อาสวา อธิวาสนา ปหาตพฺพา? เย อธิวาสนาย ปหีนา โหนฺติ, อิธ ภิกฺขเว ภิกฺขุ ปฏิสงฺขา โยนิโส ขโม โหติ สีตสฺส อุณฺหสฺส ชิฆจฺฉาย ปิปาสาย ฑํส มกสวาตาตป สรีสป สมฺผสฺสานํ ทุรุตฺตานํ ทุราคตานํ วจนปถานํ, อุปฺปนฺนานํ สารีริกานํ เวทนานํ ทุกฺขานํ ติพฺพานํ ขรานํ กฏุกานํ อสาตานํ อมนาปานํ ปาณหรานํ อธิวาสกชาติโก โหติ.\csromanspacer{} ยํ หิสฺส}

\vspace{\tipitakaparskip}
\csromancenter{ฉกฺกนิปาตปาฬิ ภิกฺขเว อนธิวาสโต อุปฺปชฺเชยฺยุํ อาสวา วิฆาตปริฬาหา, อธิวาสโต เอวํส เต อาสวา วิฆาตปริฬาหา น โหนฺติ.\csromanspacer{} อิเม วุจฺจนฺติ ภิกฺขเว อาสวา อธิวาสนา ปหาตพฺพา, เย อธิวาสนาย ปหีนา โหนฺติ.}
\prose{\csromanfolio{342}กตเม จ ภิกฺขเว อาสวา ปริวชฺชนา ปหาตพฺพา? เย ปริวชฺชนาย ปหีนา โหนฺติ, อิธ ภิกฺขเว ภิกฺขุ ปฏิสงฺขา โยนิโส จณฺฑํ หตฺถิํ ปริวชฺเชติ, จณฺฑํ อสฺสํ ปริวชฺเชติ, จณฺฑํ โคณํ ปริวชฺเชติ, จณฺฑํ กุกฺกุรํ ปริวชฺเชติ, อหิํ ขาณุํ กณฺฏกฏฺฐานํ โสพฺภํ ปปาตํ จนฺทนิกํ โอฬิคลฺลํ, ยถารูเป อนาสเน นิสินฺนํ ยถารุเป อโคจเร จรนฺตํ ยถารูเป ปาปเก มิตฺเต ภชนฺตํ วิญฺญู สพฺรหฺมจารี ปาปเกสุ ฐาเนสุ โอกปฺเปยฺยุํ, โส ตญฺจ อนาสนํ ตญฺจ อโคจรํ เต จ ปาปเก มิตฺเต ปฏิสงฺขา โยนิโส ปริวชฺเชติ.\csromanspacer{} ยํ หิสฺส ภิกฺขเว อปริวชฺชยโต อุปฺปชฺเชยฺยุํ อาสวา วิฆาตปริฬาหา, ปริวชฺชยโต เอวํส เต อาสวา วิฆาตปริฬาหา น โหนฺติ.\csromanspacer{} อิเม วุจฺจนฺติ ภิกฺขเว อาสวา ปริวชฺชนา ปหาตพฺพา, เย ปริวชฺชนาย ปหีนา โหนฺติ.}

\prose{กตเม จ ภิกฺขเว อาสวา วิโนทนา ปหาตพฺพา? เย วิโนทนาย ปหีนา โหนฺติ, อิธ ภิกฺขเว ภิกฺขุ ปฏิสงฺขา โยนิโส อุปฺปนฺนํ กามวิตกฺกํ นาธิวาเสติ ปชหติ วิโนเทติ พฺยนฺตีกโรติ อนภาวํ คเมติ, ปฏิสงฺขา โยนิโส อุปฺปนฺนํ พฺยาปาทวิตกฺกํ, อุปฺปนฺนํ วิหิํสาวิตกฺกํ, อุปฺปนฺนุปฺปนฺเน ปาปเก อกุสเล ธมฺเม นาธิวาเสติ ปชหติ วิโนเทติ พฺยนฺตีกโรติ อนภาวํ คเมติ.\csromanspacer{} ยํ หิสฺส ภิกฺขเว อวิโนทยโต อุปฺปชฺเชยฺยุํ อาสวา วิฆาตปริฬาหา, วิโนทยโต เอวํส เต อาสวา วิฆาตปริฬาหา น โหนฺติ.\csromanspacer{} อิเม วุจฺจนฺติ ภิกฺขเว อาสวา วิโนทนา ปหาตพฺพา, เย วิโนทนาย ปหีนา โหนฺติ.}

\prose{กตเม จ ภิกฺขเว อาสวา ภาวนา ปหาตพฺพา? เย ภาวนาย ปหีนา โหนฺติ, อิธ ภิกฺขเว ภิกฺขุ ปฏิสงฺขา โยนิโส สติสมฺโพชฺฌงฺคํ ภาเวติ วิเวกนิสฺสิตํ วิราคนิสฺสิตํ นิโรธนิสฺสิตํ โวสคฺคปริณามิํ.\csromanspacer{} ปฏิสงฺขา โยนิโส ธมฺมวิจย สมฺโพชฺฌงฺคํ ภาเวติ, วีริยสมฺโพชฺฌงฺคํ ภาเวติ, ปีติสมฺโพชฺฌงฺคํ ภาเวติ, ปสฺสทฺธิสมฺโพชฺฌงฺคํ ภาเวติ, สมาธิสมฺโพชฺฌงฺคํ ภาเวติ, อุเปกฺขาสมฺโพชฺฌงฺคํ ภาเวติ วิเวกนิสฺสิตํ วิราคนิสฺสิตํ นิโรธนิสฺสิตํ โวสคฺคปริณามิํ.\csromanspacer{} ยํ หิสฺส ภิกฺขเว}

\vspace{\tipitakaparskip}
\csromancenter{มหาวคฺค อภาวยโต อุปฺปชฺเชยฺยุํ อาสวา วิฆาตปริฬาหา, ภาวยโต เอวํส เต อาสวา วิฆาตปริฬาหา น โหนฺติ.\csromanspacer{} อิเม วุจฺจนฺติ ภิกฺขเว อาสวา ภาวนา ปหาตพฺพา, เย ภาวนาย ปหีนา โหนฺติ.}
\prose{\csromanfolio{343}อิเมหิ โข ภิกฺขเว ฉหิ ธมฺเมหิ สมนฺนาคโต ภิกฺขุ อาหุเนยฺโย โหติ ปาหุเนยฺโย ทกฺขิเณยฺโย อญฺชลิกรณีโย อนุตฺตรํ ปุญฺญกฺเขตฺตํ โลกสฺสาติ.\csromanspacer{}\csromanspacer{}จตุตฺถํ.}

\csromanmatikaanchor{mtk.181}
\csromantocmarkref{subsection}{5. ทารุกมฺมิกสุตฺต}{mtk.181}
\bookmark[dest={mtk.181},level=2]{5. ทารุกมฺมิกสุตฺต}
\csromanheader{2}{5. ทารุกมฺมิกสุตฺต}
\proseitem{59}{เอวํ เม สุตํ\csromandash{}เอกํ สมยํ ภควา นาติเก วิหรติ คิญฺชกาวสเถ.\csromanspacer{} อถ โข ทารุกมฺมิโก คหปติ เยน ภควา เตนุปสงฺกมิ, อุปสงฺกมิตฺวา ภควนฺตํ อภิวาเทตฺวา เอกมนฺตํ นิสีทิ.\csromanspacer{} เอกมนฺตํ นิสินฺนํ โข ทารุกมฺมิกํ คหปติํ ภควา เอตทโวจ “อปิ นุ เต คหปติ กุเล ทานํ ทียตี”ติ.\csromanspacer{} ทียติ เม ภนฺเต กุเล ทานํ, ตญฺจ โข เย เต ภิกฺขู อารญฺญิกา ปิณฺฑปาติกา ปํสุกูลิกา อรหนฺโต วา อรหตฺตมคฺคํ วา สมาปนฺนา, ตถารูเปสุ เม ภนฺเต ภิกฺขูสุ ทานํ ทียตีติ.}

\prose{ทุชฺชานํ โข เอตํ คหปติ ตยา คิหินา กามโภคินา ปุตฺตสมฺพาธสยนํ อชฺฌาวสนฺเตน กาสิกจนฺทนํ ปจฺจนุโภนฺเตน มาลาคนฺธวิเลปนํ ธารยนฺเตน ชาตรูปรชตํ สาทิยนฺเตน อิเม วา อรหนฺโต อิเม วา อรหตฺตมคฺคํ สมาปนฺนาติ.}

\prose{อารญฺญิโก เจปิ คหปติ ภิกฺขุ โหติ อุทฺธโต อุนฺนโฬ จปโล มุขโร วิกิณฺณวาโจ มุฏฺฐสฺสติ อสมฺปชาโน อสมาหิโต วิพฺภนฺตจิตฺโต ปากตินฺทฺริโย, เอวํ โส เตนงฺเคน คารโยฺห.\csromanspacer{} อารญฺญิโก เจปิ คหปติ ภิกฺขุ โหติ อนุทฺธโต อนุนฺนโฬ อจปโล อมุขโร อวิกิณฺณวาโจ อุปฏฺฐิตสฺสติ สมฺปชาโน สมาหิโต เอกคฺคจิตฺโต สํวุตินฺทฺริโย, เอวํ โส เตนงฺเคน ปาสํโส.}

\prose{คามนฺตวิหารี เจปิ คหปติ ภิกฺขุ โหติ อุทฺธโต ฯเปฯ เอวํ โส เตนงฺเคน คารโยฺห.\csromanspacer{} คามนฺตวิหารี เจปิ คหปติ ภิกฺขุ โหติ อนุทฺธโต ฯเปฯ เอวํ โส เตนงฺเคน ปาสํโส.}

\gambhira{ฉกฺกนิปาตปาฬิ}
\prose{\csromanfolio{344}ปิณฺฑปาติโก เจปิ คหปติ ภิกฺขุ โหติ อุทฺธโต ฯเปฯ เอวํ โส เตนงฺเคน คารโยฺห.\csromanspacer{} ปิณฺฑปาติโก เจปิ คหปติ ภิกฺขุ โหติ อนุทฺธโต ฯเปฯ เอวํ โส เตนงฺเคน ปาสํโส.}

\prose{เนมนฺตนิโก เจปิ คหปติ ภิกฺขุ โหติ อุทฺธโต ฯเปฯ เอวํ โส เตนงฺเคน คารโยฺห.\csromanspacer{} เนมนฺตนิโก เจปิ คหปติ ภิกฺขุ โหติ อนุทฺธโต ฯเปฯ เอวํ โส เตนงฺเคน ปาสํโส.}

\prose{ปํสุกูลิโก เจปิ คหปติ ภิกฺขุ โหติ อุทฺธโต ฯเปฯ เอวํ โส เตนงฺเคน คารโยฺห.\csromanspacer{} ปํสุกูลิโก เจปิ คหปติ ภิกฺขุ โหติ อนุทฺธโต ฯเปฯ เอวํ โส เตนงฺเคน ปาสํโส.}

\prose{คหปติจีวรธโร เจปิ คหปติ ภิกฺขุ โหติ อุทฺธโต อุนฺนโฬ จปโล มุขโร วิกิณฺณวาโจ มุฏฺฐสฺสติ อสมฺปชาโน อสมาหิโต วิพฺภนฺตจิตฺโต ปากตินฺทฺริโย, เอวํ โส เตนงฺเคน คารโยฺห.\csromanspacer{} คหปติจีวรธโร เจปิ คหปติ ภิกฺขุ โหติ อนุทฺธโต อนุนฺนโฬ อจปโล อมุขโร อวิกิณฺณวาโจ อุปฏฺฐิตสฺสติ สมฺปชาโน สมาหิโต เอกคฺคจิตฺโต สํวุตินฺทฺริโย, เอวํ โส เตนงฺเคน ปาสํโส.}

\prose{อิงฺฆ ตฺวํ คหปติ สํเฆ ทานํ\footnote{ทานานิ (ก)} เทหิ, สํเฆ เต ทานํ ททโต จิตฺตํ ปสีทิสฺสติ, โส ตฺวํ ปสนฺนจิตฺโต กายสฺส เภทา ปรํ มรณา สุคติํ สคฺคํ โลกํ อุปปชฺชิสฺสสีติ.\csromanspacer{} เอสาหํ ภนฺเต อชฺชตคฺเค สํเฆ ทานํ ทสฺสามีติ.\csromanspacer{}\csromanspacer{}ปญฺจมํ.}

\csromanmatikaanchor{mtk.182}
\csromantocmarkref{subsection}{6. หตฺถิสาริปุตฺตสุตฺต}{mtk.182}
\bookmark[dest={mtk.182},level=2]{6. หตฺถิสาริปุตฺตสุตฺต}
\csromanheader{2}{6. หตฺถิสาริปุตฺตสุตฺต}
\proseitem{60}{เอวํ เม สุตํ\csromandash{}เอกํ สมยํ ภควา พาราณสิยํ วิหรติ อิสิปตเน มิคทาเย.\csromanspacer{} เตน โข ปน สมเยน สมฺพหุลา เถรา ภิกฺขู ปจฺฉาภตฺตํ ปิณฺฑปาตปฏิกฺกนฺตา มณฺฑลมาเฬ สนฺนิสินฺนา สนฺนิปติตา อภิธมฺมกถํ กเถนฺติ.\csromanspacer{} ตตฺร สุทํ อายสฺมา จิตฺโต หตฺถิสาริปุตฺโต เถรานํ ภิกฺขูนํ อภิธมฺมกถํ กเถนฺตานํ อนฺตรนฺตรา กถํ โอปาเตติ.\csromanspacer{} อถ โข อายสฺมา มหาโกฏฺฐิโก อายสฺมนฺตํ จิตฺตํ หตฺถิสาริปุตฺตํ เอตทโวจ “มายสฺมา จิตฺโต หตฺถิสาริปุตฺโต เถรานํ ภิกฺขูนํ}

\vspace{\tipitakaparskip}
\csromancenter{มหาวคฺค อภิธมฺมกถํ กเถนฺตานํ อนฺตรนฺตรา กถํ โอปาเตสิ, ยาว กถาปริโยสานํ อายสฺมา จิตฺโต อาคเมตู”ติ.\csromanspacer{} เอวํ วุตฺเต อายสฺมโต จิตฺตสฺส หตฺถิสาริปุตฺตสฺส สหายกา ภิกฺขู อายสฺมนฺตํ มหาโกฏฺฐิกํ เอตทโวจุํ “มายสฺมา มหาโกฏฺฐิโก อายสฺมนฺตํ จิตฺตํ หตฺถิสาริปุตฺตํ อปสาเทสิ, ปณฺฑิโต อายสฺมา จิตฺโต หตฺถิสาริปุตฺโต, ปโหติ จายสฺมา จิตฺโต หตฺถิสาริปุตฺโต เถรานํ ภิกฺขูนํ อภิธมฺมกถํ กเถตุนฺ”ติ.}
\prose{\csromanfolio{345}ทุชฺชานํ โข เอตํ อาวุโส ปรสฺส เจโตปริยายํ อชานนฺเตหิ.\csromanspacer{} อิธาวุโส เอกจฺโจ ปุคฺคโล ตาวเทว โสรตโสรโต โหติ, นิวาตนิวาโต โหติ, อุปสนฺตุปสนฺโต โหติ, ยาว สตฺถารํ อุปนิสฺสาย วิหรติ อญฺญตรํ วา ครุฏฺฐานิยํ สพฺรหฺมจาริํ.\csromanspacer{} ยโต จ โข โส วปกสฺสเตว สตฺถารา, วปกสฺสติ ครุฏฺฐานิเยหิ สพฺรหฺมจารีหิ, โส สํสฏฺโฐ วิหรติ ภิกฺขูหิ ภิกฺขุนีหิ อุปาสเกหิ อุปาสิกาหิ รญฺญา ราชมหามตฺเตหิ ติตฺถิเยหิ ติตฺถิยสาวเกหิ.\csromanspacer{} ตสฺส สํสฏฺฐสฺส วิสฺสตฺถสฺส ปากตสฺส ภสฺสมนุยุตฺตสฺส วิหรโต ราโค จิตฺตํ อนุทฺธํเสติ, โส ราคานุทฺธํสิเตน จิตฺเตน สิกฺขํ ปจฺจกฺขาย หีนายาวตฺตติ.}

\prose{เสยฺยาถาปิ อาวุโส โคโณ กิฏฺฐาโท ทาเมน วา พทฺโธ\footnote{อาราเม วา พนฺโธ (ก)} วเช วา โอรุทฺโธ.\csromanspacer{} โย นุ โข อาวุโส เอวํ วเทยฺย “น ทานายํ โคโณ กิฏฺฐาโท ปุนเทว กิฏฺฐํ โอตริสฺสตี”ติ.\csromanspacer{} สมฺมา นุ โข โส อาวุโส วทมาโน วเทยฺยาติ.\csromanspacer{} โน หิทํ อาวุโส, ฐานํ เหตํ อาวุโส วิชฺชติ, ยํ โส โคโณ กิฏฺฐาโท ทามํ วา เฉตฺวา วชํ วา ภินฺทิตฺวา อถ ปุนเทว กิฏฺฐํ โอตเรยฺยาติ.\csromanspacer{} เอวเมวํ โข อาวุโส อิเธกจฺโจ ปุคฺคโล ตาวเทว โสรตโสรโต โหติ, นิวาตนิวาโต โหติ, อุปสนฺตุปสนฺโต โหติ, ยาว สตฺถารํ อุปนิสฺสาย วิหรติ อญฺญตรํ วา ครุฏฺฐานิยํ สพฺรหฺมจาริํ.\csromanspacer{} ยโต จ โข โส วปกสฺสเตว สตฺถารา, วปกสฺสติ ครุฏฺฐานิเยหิ สพฺรหฺมจารีหิ, โส สํสฏฺโฐ วิหรติ ภิกฺขูหิ ภิกฺขุนีหิ อุปาสเกหิ อุปาสิกาหิ รญฺญา ราชมหามตฺเตหิ ติตฺถิเยหิ ติตฺถิยสาวเกหิ.\csromanspacer{} ตสฺส สํสฏฺฐสฺส วิสฺสตฺถสฺส ปากตสฺส ภสฺสมนุยุตฺตสฺส วิหรโต ราโค จิตฺตํ}

\vspace{\tipitakaparskip}
\csromancenter{ฉกฺกนิปาตปาฬิ อนุทฺธํเสติ, โส ราคานุทฺธํสิเตน จิตฺเตน สิกฺขํ ปจฺจกฺขาย หีนายาวตฺตติ. (1)}
\prose{\csromanfolio{346}อิธ ปนาวุโส เอกจฺโจ ปุคฺคโล วิวิจฺเจว กาเมหิ ฯเปฯ ปฐมํ ฌานํ อุปสมฺปชฺช วิหรติ, โส “ลาภิมฺหิ ปฐมสฺส ฌานสฺสา”ติ สํสฏฺโฐ วิหรติ ภิกฺขูหิ ฯเปฯ สิกฺขํ ปจฺจกฺขาย หีนายาวตฺตติ.\csromanspacer{} เสยฺยถาปิ อาวุโส จาตุมหาปเถ ถุลฺลผุสิตโก เทโว วสฺสนฺโต\footnote{ถุลฺลผุสิตเก เทเว วสฺสนฺเต (ก)} รชํ อนฺตรธาเปยฺย, จิกฺขลฺลํ ปาตุกเรยฺย.\csromanspacer{} โย นุ โข อาวุโส เอวํ วเทยฺย “น ทานิ อมุสฺมิํ\footnote{อมุกสฺมิํ (ก)} จาตุมหาปเถ ปุนเทว รโช ปาตุภวิสฺสตี”ติ.\csromanspacer{} สมฺมา นุ โข โส อาวุโส วทมาโน วเทยฺยาติ.\csromanspacer{} โน หิทํ อาวุโส, ฐานํ เหตํ อาวุโส วิชฺชติ, ยํ อมุสฺมิํ จาตุมหาปเถ มนุสฺสา วา อติกฺกเมยฺยุํ, โคปสู วา อติกฺกเมยฺยุํ, วาตาตโป วา สฺเนหคตํ ปริยาทิเยยฺย, อถ ปุนเทว รโช ปาตุภเวยฺยาติ.\csromanspacer{} เอวเมวํ โข อาวุโส อิเธกจฺโจ ปุคฺคโล วิวิจฺเจว กาเมหิ ฯเปฯ ปฐมํ ฌานํ อุปสมฺปชฺช วิหรติ.\csromanspacer{} โส “ลาภิมฺหิ ปฐมสฺส ฌานสฺสา”ติ สํสฏฺโฐ วิหรติ ภิกฺขูหิ ฯเปฯ สิกฺขํ ปจฺจกฺขาย หีนายาวตฺตติ. (2)}

\prose{อิธ ปนาวุโส เอกจฺโจ ปุคฺคโล วิตกฺกวิจารานํ วูปสมา ฯเปฯ ทุติยํ ฌานํ อุปสมฺปชฺช วิหรติ, โส “ลาภิมฺหิ ทุติยสฺส ฌานสฺสา”ติ สํสฏฺโฐ วิหรติ ภิกฺขูหิ ฯเปฯ สิกฺขํ ปจฺจกฺขาย หีนายาวตฺตติ.\csromanspacer{} เสยฺยถาปิ อาวุโส คามสฺส วา นิคมสฺส วา อวิทูเร มหนฺตํ ตฬากํ, ตตฺถ ถุลฺลผุสิตโก เทโว วุฏฺโฐ สิปฺปิสมฺพุกมฺปิ สกฺขรกฐลมฺปิ อนฺตรธาเปยฺย.\csromanspacer{} โย นุ โข อาวุโส เอวํ วเทยฺย “น ทานิ อมุสฺมิํ ตฬาเก ปุนเทว สิปฺปิสมฺพุกา วา สกฺขรกฐลา วา ปาตุภวิสฺสนฺตี”ติ.\csromanspacer{} สมฺมา นุ โข โส อาวุโส วทมาโน วเทยฺยาติ.\csromanspacer{} โน หิทํ อาวุโส, ฐานํ เหตํ อาวุโส วิชฺชติ, ยํ อมุสฺมิํ ตฬาเก มนุสฺสา วา ปิเวยฺยุํ, โคปสู วา ปิเวยฺยุํ, วาตาตโป วา สฺเนหคตํ ปริยาทิเยยฺย, อถ ปุนเทว สิปฺปิสมฺพุกาปิ สกฺขรกฐลาปิ ปาตุภเวยฺยุนฺติ.\csromanspacer{} เอวเมวํ โข อาวุโส อิเธกจฺโจ ปุคฺคโล วิตกฺกวิจารานํ วูปสมา ฯเปฯ ทุติยํ ฌานํ อุปสมฺปชฺช วิหรติ.\csromanspacer{} โส “ลาภิมฺหิ ทุติยสฺส ฌานสฺสา”ติ สํสฏฺโฐ วิหรติ ภิกฺขูหิ ฯเปฯ สิกฺขํ ปจฺจกฺขาย หีนายาวตฺตติ. (3)}

\chapterheadpageread{6. มหาวคฺค}
\prose{\csromanfolio{347}อิธ ปนาวุโส เอกจฺโจ ปุคฺคโล ปีติยา จ วิราคา ฯเปฯ ตติยํ ฌานํ อุปสมฺปชฺช วิหรติ, โส “ลาภิมฺหิ ตติยสฺส ฌานสฺสา”ติ สํสฏฺโฐ วิหรติ ภิกฺขูหิ ฯเปฯ สิกฺขํ ปจฺจกฺขาย หีนายาวตฺตติ.\csromanspacer{} เสยฺยถาปิ อาวุโส ปุริสํ ปณีตโภชนํ ภุตฺตาวิํ อาภิโทสิกํ โภชนํ นจฺฉาเทยฺย.\csromanspacer{} โย นุ โข อาวุโส เอวํ วเทยฺย “น ทานิ อมุํ ปุริสํ ปุนเทว โภชนํ ฉาเทสฺสตี”ติ.\csromanspacer{} สมฺมา นุ โข โส อาวุโส วทมาโน วเทยฺยาติ.\csromanspacer{} โน หิทํ อาวุโส, ฐานํ เหตํ อาวุโส วิชฺชติ, อมุํ ปุริสํ ปณีตโภชนํ ภุตฺตาวิํ ยาวสฺส สา โอชา กาเย ฐสฺสติ, ตาว น อญฺญํ โภชนํ ฉาเทสฺสติ.\csromanspacer{} ยโต จ ขฺวสฺส สา โอชา อนฺตรธายิสฺสติ, อถ ปุนเทว ตํ โภชนํ ฉาเทยฺยาติ.\csromanspacer{} เอวเมวํ โข อาวุโส อิเธกจฺโจ ปุคฺคโล ปีติยา จ วิราคา ฯเปฯ ตติยํ ฌานํ อุปสมฺปชฺช วิหรติ.\csromanspacer{} โส “ลาภิมฺหิ ตติยสฺส ฌานสฺสา”ติ สํสฏฺโฐ วิหรติ ภิกฺขูหิ ฯเปฯ สิกฺขํ ปจฺจกฺขาย หีนายาวตฺตติ. (4)}

\prose{อิธ ปนาวุโส เอกจฺโจ ปุคฺคโล สุขสฺส จ ปหานา ทุกฺขสฺส จ ปหานา ฯเปฯ จตุตฺถํ ฌานํ อุปสมฺปชฺช วิหรติ, โส “ลาภิมฺหิ จตุตฺถสฺส ฌานสฺสา”ติ สํสฏฺโฐ วิหรติ ภิกฺขูหิ ฯเปฯ สิกฺขํ ปจฺจกฺขาย หีนายาวตฺตติ.\csromanspacer{} เสยฺยถาปิ อาวุโส ปพฺพตสงฺเขเป อุทกรหโท นิวาโต วิคต-อูมิโก.\csromanspacer{} โย นุ โข อาวุโส เอวํ วเทยฺย “น ทานิ อมุสฺมิํ อุทกรหเท ปุนเทว อูมิ ปาตุภวิสฺสตี”ติ.\csromanspacer{} สมฺมา นุ โข โส อาวุโส วทมาโน วเทยฺยาติ.\csromanspacer{} โน หิทํ อาวุโส, ฐานํ เหตํ อาวุโส วิชฺชติ, ยา ปุรตฺถิมาย ทิสาย อาคจฺเฉยฺย ภุสา วาตวุฏฺฐิ, สา ตสฺมิํ อุทกรหเท อูมิํ ชเนยฺย.\csromanspacer{} ยา ปจฺฉิมาย ทิสาย อาคจฺเฉยฺย ฯเปฯ ยา อุตฺตราย ทิสาย อาคจฺเฉยฺย.\csromanspacer{} ยา ทกฺขิณาย ทิสาย อาคจฺเฉยฺย ภุสา วาตวุฏฺฐิ, สา ตสฺมิํ อุทกรหเท อูมิํ ชเนยฺยาติ.\csromanspacer{} เอวเมวํ โข อาวุโส อิเธกจฺโจ ปุคฺคโล สุขสฺส จ ปหานา ทุกฺขสฺส จ ปหานา ฯเปฯ จตุตฺถํ ฌานํ อุปสมฺปชฺช วิหรติ.\csromanspacer{} โส “ลาภิมฺหิ จตุตฺถสฺส ฌานสฺสา”ติ สํสฏฺโฐ วิหรติ ภิกฺขูหิ ฯเปฯ สิกฺขํ ปจฺจกฺขาย หีนายาวตฺตติ. (5)}

\prose{อิธ ปนาวุโส เอกจฺโจ ปุคฺคโล สพฺพนิมิตฺตานํ อมนสิการา อนิมิตฺตํ เจโตสมาธิํ อุปสมฺปชฺช วิหรติ, โส “ลาภิมฺหิ อนิมิตฺตสฺส เจโตสมาธิสฺสา”ติ สํสฏฺโฐ วิหรติ ภิกฺขูหิ ภิกฺขุนีหิ อุปาสเกหิ อุปาสิกาหิ รญฺญา ราชมหามตฺเตหิ ติตฺถิเยหิ ติตฺถิย}

\vspace{\tipitakaparskip}
\csromancenter{ฉกฺกนิปาตปาฬิ สาวเกหิ.\csromanspacer{} ตสฺส สํสฏฺฐสฺส วิสฺสตฺถสฺส ปากตสฺส ภสฺสมนุยุตฺตสฺส วิหรโต ราโค จิตฺตํ อนุทฺธํเสติ.\csromanspacer{} โส ราคานุทฺธํสิเตน จิตฺเตน สิกฺขํ ปจฺจกฺขาย หีนายาวตฺตติ, เสยฺยถาปิ อาวุโส ราชา วา ราชมหามตฺโต วา จตุรงฺคินิยา เสนาย อทฺธานมคฺคปฺปฏิปนฺโน อญฺญตรสฺมิํ วนสณฺเฑ เอกรตฺติํ วาสํ อุปคจฺเฉยฺย.\csromanspacer{} ตตฺร\footnote{ตตฺถ (สี, อิ)} หตฺถิสทฺเทน อสฺสสทฺเทน รถสทฺเทน ปตฺติสทฺเทน เภริปณวสงฺขติณว นินฺนาทสทฺเทน จีริกสทฺโท\footnote{จิริฬิกาสทฺโท (สี, สฺยา, กํ, อิ)} อนฺตรธาเยยฺย\footnote{อนฺตรธาเปยฺย (สฺยา, อิ, ก)}.\csromanspacer{} โย นุ โข อาวุโส เอวํ วเทยฺย “น ทานิ อมุสฺมิํ วนสณฺเฑ ปุนเทว จีริกสทฺโท ปาตุภวิสฺสตี”ติ.\csromanspacer{} สมฺมา นุ โข โส อาวุโส วทมาโน วเทยฺยาติ.\csromanspacer{} โน หิทํ อาวุโส, ฐานํ เหตํ อาวุโส วิชฺชติ, ยํ โส ราชา วา ราชมหามตฺโต วา ตมฺหา วนสณฺฑา ปกฺกเมยฺย, อถ ปุนเทว จีริกสทฺโท ปาตุภเวยฺยาติ.\csromanspacer{} เอวเมวํ โข อาวุโส อิเธกจฺโจ ปุคฺคโล สพฺพนิมิตฺตานํ อมนสิการา อนิมิตฺตํ เจโตสมาธิํ อุปสมฺปชฺช วิหรติ.\csromanspacer{} โส “ลาภิมฺหิ อนิมิตฺตสฺส เจโตสมาธิสฺสา”ติ สํสฏฺโฐ วิหรติ ภิกฺขูหิ ภิกฺขุนีหิ อุปาสเกหิ อุปาสิกาหิ รญฺญา ราชมหามตฺเตหิ ติตฺถิเยหิ ติตฺถิยสาวเกหิ.\csromanspacer{} ตสฺส สํสฏฺฐสฺส วิสฺสตฺถสฺส ปากตสฺส ภสฺสมนุยุตฺตสฺส วิหรโต ราโค จิตฺตํ อนุทฺธํเสติ, โส ราคานุทฺธํสิเตน จิตฺเตน สิกฺขํ ปจฺจกฺขาย หีนายาวตฺตตีติ. (6)}
\prose{\csromanfolio{348}อถ โข อายสฺมา จิตฺโต หตฺถิสาริปุตฺโต อปเรน สมเยน สิกฺขํ ปจฺจกฺขาย หีนายาวตฺตติ.\csromanspacer{} อถ โข จิตฺตสฺส หตฺถิสาริปุตฺตสฺส สหายกา ภิกฺขู เยนายสฺมา มหาโกฏฺฐิโก เตนุปสงฺกมิํสุ, อุปสงฺกมิตฺวา อายสฺมนฺตํ มหาโกฏฺฐิกํ เอตทโวจุํ “กิํ นุ โข อายสฺมตา มหาโกฏฺฐิเกน จิตฺโต หตฺถิสาริปุตฺโต เจตสา เจโต ปริจฺจ วิทิโต ‘อิมาสญฺจ อิมาสญฺจ วิหารสมาปตฺตีนํ จิตฺโต หตฺถิสาริปุตฺโต ลาภี, อถ จ ปน สิกฺขํ ปจฺจกฺขาย หีนายาวตฺติสฺสตี’ติ.\csromanspacer{} อุทาหุ เทวตา เอตมตฺถํ อาโรเจสุํ ‘จิตฺโต ภนฺเต หตฺถิสาริปุตฺโต อิมาสญฺจ อิมาสญฺจ วิหารสมาปตฺตีนํ ลาภี, อถ จ ปน สิกฺขํ ปจฺจกฺขาย หีนายาวตฺติสฺสตี’ติ”.\csromanspacer{} เจตสา เจโต ปริจฺจ วิทิโต เม อาวุโส จิตฺโต หตฺถิสาริปุตฺโต “อิมาสญฺจ อิมาสญฺจ วิหารสมาปตฺตีนํ ลาภี, อถ จ ปน สิกฺขํ ปจฺจกฺขาย หีนายาวตฺติสฺสตี”ติ.\csromanspacer{} เทวตาปิ เม}

\vspace{\tipitakaparskip}
\csromancenter{มหาวคฺค เอตมตฺถํ อาโรเจสุํ “จิตฺโต ภนฺเต หตฺถิสาริปุตฺโต อิมาสญฺจ อิมาสญฺจ วิหารสมาปตฺตีนํ ลาภี, อถ จ ปน สิกฺขํ ปจฺจกฺขาย หีนายาวตฺติสฺสตี”ติ.}
\prose{\csromanfolio{349}อถ โข จิตฺตสฺส หตฺถิสาริปุตฺตสฺส สหายกา ภิกฺขู เยน ภควา เตนุปสงฺกมิํสุ, อุปสงฺกมิตฺวา ภควนฺตํ อภิวาเทตฺวา เอกมนฺตํ นิสีทิํสุ, เอกมนฺตํ นิสินฺนา โข เต ภิกฺขู ภควนฺตํ เอตทโวจุํ “จิตฺโต ภนฺเต หตฺถิสาริปุตฺโต อิมาสญฺจ อิมาสญฺจ วิหารสมาปตฺตีนํ ลาภี, อถ จ ปน สิกฺขํ ปจฺจกฺขาย หีนายาวตฺตตี”ติ.\csromanspacer{} น ภิกฺขเว จิตฺโต จิรํ สริสฺสติ\footnote{ปทิสฺสติ (ก)} เนกฺขมฺมสฺสาติ.}

\prose{อถ โข จิตฺโต หตฺถิสาริปุตฺโต นจิรสฺเสว เกสมสฺสุํ โอหาเรตฺวา กาสายานิ วตฺถานิ อจฺฉาเทตฺวา อคารสฺมา อนคาริยํ ปพฺพชิ.\csromanspacer{} อถ โข อายสฺมา จิตฺโต หตฺถิสาริปุตฺโต เอโก วูปกฏฺโฐ อปฺปมตฺโต อาตาปี ปหิตตฺโต วิหรนฺโต นจิรสฺเสว ยสฺสตฺถาย กุลปุตฺตา สมฺมเทว อคารสฺมา อนคาริยํ ปพฺพชนฺติ, ตทนุตฺตรํ พฺรหฺมจริยปริโยสานํ ทิฏฺเฐว ธมฺเม สยํ อภิญฺญา สจฺฉิกตฺวา อุปสมฺปชฺช วิหาสิ, “ขีณา ชาติ, วุสิตํ พฺรหฺมจริยํ, กตํ กรณียํ, นาปรํ อิตฺถตฺตายา”ติ อพฺภญฺญาสิ.\csromanspacer{} อญฺญตโร จ ปนายสฺมา จิตฺโต หตฺถิสาริปุตฺโต อรหตํ อโหสีติ.\csromanspacer{}\csromanspacer{}ฉฏฺฐํ.}

\csromanmatikaanchor{mtk.183}
\csromantocmarkref{subsection}{7. มชฺเฌสุตฺต}{mtk.183}
\bookmark[dest={mtk.183},level=2]{7. มชฺเฌสุตฺต}
\csromanheader{2}{7. มชฺเฌสุตฺต}
\proseitem{61}{เอวํ เม สุตํ\csromandash{}เอกํ สมยํ ภควา พาราณสิยํ วิหรติ อิสิปตเน มิคทาเย.\csromanspacer{} เตน โข ปน สมเยน สมฺพหุลานํ เถรานํ ภิกฺขูนํ ปจฺฉาภตฺตํ ปิณฺฑปาตปฏิกฺกนฺตานํ มณฺฑลมาเฬ สนฺนิสินฺนานํ สนฺนิปติตานํ อยมนฺตรากถา อุทปาทิ “วุตฺตมิทํ อาวุโส ภควตา ปารายเน เมตฺเตยฺยปเญฺห\csromandash{}}

\csromangathameasure{*‘โย อุโภนฺเต วิทิตฺวาน, มชฺเฌ มนฺตา น ลิปฺปติ. \\ ตํ พฺรูมิ มหาปุริโสติ, โสธ สิพฺพินิ มจฺจคา’ติ.}
\csromangathagroup{\csromangathastanza{\csromansymbolfootnote{*}{ขุ 1. 435; ขุ 8. 7 ปิฏฺเฐสุ.}‘โย อุโภนฺเต วิทิตฺวาน, มชฺเฌ มนฺตา น ลิปฺปติ\footnote{น ลิมฺปติ (ก)}. \\ ตํ พฺรูมิ มหาปุริโสติ, โสธ สิพฺพินิ\footnote{น ลิมฺปติ (ก)} มจฺจคา’ติ.}}

\prose{กตโม นุ โข อาวุโส เอโก อนฺโต, กตโม ทุติโย อนฺโต, กิํ มชฺเฌ.\csromanspacer{} กา สิพฺพินี”ติ.\csromanspacer{} เอวํ วุตฺเต อญฺญตโร ภิกฺขุ เถเร}

\vspace{\tipitakaparskip}
\csromancenter{ฉกฺกนิปาตปาฬิ ภิกฺขู เอตทโวจ “ผสฺโส โข อาวุโส เอโก อนฺโต, ผสฺสสมุทโย ทุติโย อนฺโต, ผสฺสนิโรโธ มชฺเฌ, ตณฺหา สิพฺพินี, ตณฺหา หิ นํ สิพฺพติ ตสฺส ตสฺเสว ภวสฺส อภินิพฺพตฺติยา.\csromanspacer{} เอตฺตาวตา โข อาวุโส ภิกฺขุ อภิญฺเญยฺยํ อภิชานาติ, ปริญฺเญยฺยํ ปริชานาติ, อภิญฺเญยฺยํ อภิชานนฺโต\footnote{อภิชานิตฺวา (ก)} ปริญฺเญยฺยํ ปริชานนฺโต\footnote{ปริชานิตฺวา (ก)} ทิฏฺเฐว ธมฺเม ทุกฺขสฺสนฺตกโร โหตี”ติ.}
\prose{\csromanfolio{350}เอวํ วุตฺเต อญฺญตโร ภิกฺขุ เถเร ภิกฺขู เอตทโวจ “อตีตํ โข อาวุโส เอโก อนฺโต, อนาคตํ ทุติโย อนฺโต, ปจฺจุปฺปนฺนํ มชฺเฌ, ตณฺหา สิพฺพินี, ตณฺหา หิ นํ สิพฺพติ ตสฺส ตสฺเสว ภวสฺส อภินิพฺพตฺติยา.\csromanspacer{} เอตฺตาวตา โข อาวุโส ภิกฺขุ อภิญฺเญยฺยํ อภิชานาติ, ปริญฺเญยฺยํ ปริชานาติ, อภิญฺเญยฺยํ อภิชานนฺโต ปริญฺเญยฺยํ ปริชานนฺโต ทิฏฺเฐว ธมฺเม ทุกฺขสฺสนฺตกโร โหตี”ติ.}

\prose{เอวํ วุตฺเต อญฺญตโร ภิกฺขุ เถเร ภิกฺขู เอตทโวจ “สุขา อาวุโส เวทนา เอโก อนฺโต, ทุกฺขา เวทนา ทุติโย อนฺโต, อทุกฺขมสุขา เวทนา มชฺเฌ, ตณฺหา สิพฺพินี, ตณฺหา หิ นํ สิพฺพติ ตสฺส ตสฺเสว ภวสฺส อภินิพฺพตฺติยา.\csromanspacer{} เอตฺตาวตา โข อาวุโส ภิกฺขุ อภิญฺเญยฺยํ อภิชานาติ, ปริญฺเญยฺยํ ปริชานาติ, อภิญฺเญยฺยํ อภิชานนฺโต ปริญฺเญยฺยํ ปริชานนฺโต ทิฏฺเฐว ธมฺเม ทุกฺขสฺสนฺตกโร โหตี”ติ.}

\prose{เอวํ วุตฺเต อญฺญตโร ภิกฺขุ เถเร ภิกฺขู เอตทโวจ “นามํ โข อาวุโส เอโก อนฺโต, รูปํ ทุติโย อนฺโต, วิญฺญาณํ มชฺเฌ, ตณฺหา สิพฺพินี, ตณฺหา หิ นํ สิพฺพติ ตสฺส ตสฺเสว ภวสฺส อภินิพฺพตฺติยา.\csromanspacer{} เอตฺตาวตา โข อาวุโส ภิกฺขุ อภิญฺเญยฺยํ อภิชานาติ, ปริญฺเญยฺยํ ปริชานาติ, อภิญฺเญยฺยํ อภิชานนฺโต ปริญฺเญยฺยํ ปริชานนฺโต ทิฏฺเฐว ธมฺเม ทุกฺขสฺสนฺตกโร โหตี”ติ.}

\prose{เอวํ วุตฺเต อญฺญตโร ภิกฺขุ เถเร ภิกฺขู เอตทโวจ “ฉ โข อาวุโส อชฺฌตฺติกานิ อายตนานิ เอโก อนฺโต, ฉ พาหิรานิ อายตนานิ ทุติโย อนฺโต, วิญฺญาณํ มชฺเฌ, ตณฺหา สิพฺพินี, ตณฺหา หิ นํ สิพฺพติ ตสฺส ตสฺเสว ภวสฺส อภินิพฺพตฺติยา.\csromanspacer{} เอตฺตาวตา โข}

\vspace{\tipitakaparskip}
\csromancenter{มหาวคฺค อาวุโส ภิกฺขุ อภิญฺเญยฺยํ อภิชานาติ, ปริญฺเญยฺยํ ปริชานาติ, อภิญฺเญยฺยํ อภิชานนฺโต ปริญฺเญยฺยํ ปริชานนฺโต ทิฏฺเฐว ธมฺเม ทุกฺขสฺสนฺตกโร โหตี”ติ.}
\prose{\csromanfolio{351}เอวํ วุตฺเต อญฺญตโร ภิกฺขุ เถเร ภิกฺขู เอตทโวจ “สกฺกาโย โข อาวุโส เอโก อนฺโต, สกฺกายสมุทโย ทุติโย อนฺโต, สกฺกายนิโรโธ มชฺเฌ, ตณฺหา สิพฺพินี, ตณฺหา หิ นํ สิพฺพติ ตสฺส ตสฺเสว ภวสฺส อภินิพฺพตฺติยา.\csromanspacer{} เอตฺตาวตา โข อาวุโส ภิกฺขุ อภิญฺเญยฺยํ อภิชานาติ, ปริญฺเญยฺยํ ปริชานาติ, อภิญฺเญยฺยํ อภิชานนฺโต\footnote{สพฺพตฺถปิ เอวเมว ทิสฺสติ.} ปริญฺเญยฺยํ ปริชานนฺโต1 ทิฏฺเฐว ธมฺเม ทุกฺขสฺสนฺตกโร โหตี”ติ.}

\prose{เอวํ วุตฺเต อญฺญตโร ภิกฺขุ เถเร ภิกฺขู เอตทโวจ “พฺยากตํ โข อาวุโส อเมฺหหิ สพฺเพเหว ยถาสกํ ปฏิภานํ, อายามาวุโส เยน ภควา เตนุปสงฺกมิสฺสาม, อุปสงฺกมิตฺวา ภควโต เอตมตฺถํ อาโรเจสฺสาม.\csromanspacer{} ยถา โน ภควา พฺยากริสฺสติ, ตถา นํ ธาเรสฺสามา”ติ. “เอวมาวุโส”ติ โข เถรา ภิกฺขู ตสฺส ภิกฺขุโน ปจฺจสฺโสสุํ.\csromanspacer{} อถ โข เถรา ภิกฺขู เยน ภควา เตนุป สงฺกมิํสุ, อุปสงฺกมิตฺวา ภควนฺตํ อภิวาเทตฺวา เอกมนฺตํ นิสีทิํสุ, เอกมนฺตํ นิสินฺนา โข เถรา ภิกฺขู ยาวตโก อโหสิ สพฺเพเหว สทฺธิํ กถาสลฺลาโป, ตํ สพฺพํ ภควโต อาโรเจสุํ “กสฺส นุ โข ภนฺเต สุภาสิตนฺ”ติ.\csromanspacer{} สพฺเพสํ โว ภิกฺขเว สุภาสิตํ ปริยาเยน, อปิ จ ยํ มยา สนฺธาย ภาสิตํ ปารายเน เมตฺเตยฺยปเญฺห\csromandash{}}

\csromangathameasure{“โย อุโภนฺเต วิทิตฺวาน, มชฺเฌ มนฺตา น ลิปฺปติ. \\ ตํ พฺรูมิ ‘มหาปุริโส’ติ, โสธ สิพฺพินิมจฺจคา”ติ.}
\csromangathagroup{\csromangathastanza{“โย อุโภนฺเต วิทิตฺวาน, มชฺเฌ มนฺตา น ลิปฺปติ. \\ ตํ พฺรูมิ ‘มหาปุริโส’ติ, โสธ สิพฺพินิมจฺจคา”ติ.}}

\prose{ตํ สุณาถ สาธุกํ มนสิ กโรถ, ภาสิสฺสามีติ. “เอวํ ภนฺเต”ติ โข เถรา ภิกฺขู ภควโต ปจฺจสฺโสสุํ.\csromanspacer{} ภควา เอตทโวจ “ผสฺโส โข ภิกฺขเว เอโก อนฺโต, ผสฺสสมุทโย ทุติโย อนฺโต, ผสฺสนิโรโธ มชฺเฌ, ตณฺหา สิพฺพินี, ตณฺหา หิ นํ สิพฺพติ ตสฺส ตสฺเสว ภวสฺส อภินิพฺพตฺติยา.\csromanspacer{} เอตฺตาวตา โข ภิกฺขเว ภิกฺขุ อภิญฺเญยฺยํ อภิชานาติ, ปริญฺเญยฺยํ ปริชานาติ, อภิญฺเญยฺยํ}

\vspace{\tipitakaparskip}
\csromancenter{ฉกฺกนิปาตปาฬิ อภิชานนฺโต ปริญฺเญยฺยํ ปริชานนฺโต ทิฏฺเฐว ธมฺเม ทุกฺขสฺสนฺตกโร โหตี”ติ.\csromanspacer{}\csromanspacer{}สตฺตมํ.}
\csromanmatikaanchor{mtk.184}
\csromantocmarkref{subsection}{8. ปุริสินฺทฺริยญาณสุตฺต}{mtk.184}
\bookmark[dest={mtk.184},level=2]{8. ปุริสินฺทฺริยญาณสุตฺต}
\csromanheader{2}{8. ปุริสินฺทฺริยญาณสุตฺต}
\proseitem{62}{\csromanfolio{352}เอวํ เม สุตํ\csromandash{}เอกํ สมยํ ภควา โกสเลสุ จาริกํ จรมาโน มหตา ภิกฺขุสํเฆน สทฺธิํ เยน ทณฺฑกปฺปกํ นาม โกสลานํ นิคโม ตทวสริ.\csromanspacer{} อถ โข ภควา มคฺคา โอกฺกมฺม อญฺญตรสฺมิํ รุกฺขมูเล ปญฺญตฺเต อาสเน นิสีทิ.\csromanspacer{} เต จ ภิกฺขู ทณฺฑกปฺปกํ ปวิสิํสุ อาวสถํ ปริเยสิตุํ.}

\prose{อถ โข อายสฺมา อานนฺโท สมฺพหุเลหิ ภิกฺขูหิ สทฺธิํ เยน อจิรวตี นที เตนุปสงฺกมิ คตฺตานิ ปริสิญฺจิตุํ, อจิรวติยา นทิยา คตฺตานิ ปริสิญฺจิตฺวา ปจฺจุตฺตริตฺวา เอกจีวโร อฏฺฐาสิ คตฺตานิ ปุพฺพาปยมาโน.\csromanspacer{} อถ โข อญฺญตโร ภิกฺขุ เยนายสฺมา อานนฺโท เตนุปสงฺกมิ, อุปสงฺกมิตฺวา อายสฺมนฺตํ อานนฺทํ เอตทโวจ “กิํ นุ โข อาวุโส อานนฺท สพฺพํ เจตโส สมนฺนาหริตฺวา นุ โข เทวทตฺโต ภควตา พฺยากโต ‘อาปายิโก เทวทตฺโต เนรยิโก กปฺปฏฺโฐ อเตกิจฺโฉ’ติ, * อุทาหุ เกนจิเทว ปริยาเยนา”ติ.\csromanspacer{} เอวํ โข ปเนตํ อาวุโส ภควตา พฺยากตนฺติ.}

\prose{อถ โข อายสฺมา อานนฺโท เยน ภควา เตนุปสงฺกมิ, อุปสงฺกมิตฺวา ภควนฺตํ อภิวาเทตฺวา เอกมนฺตํ นิสีทิ, เอกมนฺตํ นิสินฺโน โข อายสฺมา อานนฺโท ภควนฺตํ เอตทโวจ “อิธาหํ ภนฺเต สมฺพหุเลหิ ภิกฺขูหิ สทฺธิํ เยน อจิรวตี นที เตนุปสงฺกมิํ คตฺตานิ ปริสิญฺจิตุํ, อจิรวติยา นทิยา คตฺตานิ ปริสิญฺจิตฺวา ปจฺจุตฺตริตฺวา เอกจีวโร อฏฺฐาสิํ คตฺตานิ ปุพฺพาปยมาโน.\csromanspacer{} อถ โข ภนฺเต อญฺญตโร ภิกฺขุ เยนาหํ เตนุปสงฺกมิ, อุปสงฺกมิตฺวา มํ เอตทโวจ “กิํ นุ โข อาวุโส อานนฺท สพฺพํ เจตโส สมนฺนาหริตฺวา นุ โข เทวทตฺโต ภควตา พฺยากโต “อาปายิโก เทวทตฺโต เนรยิโก กปฺปฏฺโฐ อเตกิจฺโฉ’ติ, อุทาหุ เกนจิเทว ปริยาเยนา”ติ.\csromanspacer{} เอวํ วุตฺเต อหํ ภนฺเต ตํ ภิกฺขุํ เอตทโวจํ “เอวํ โข ปเนตํ อาวุโส ภควตา พฺยากตนฺ”ติ.}

\chapterheadpageread{6. มหาวคฺค}
\prose{\csromanfolio{353}โส วา\footnote{โส จ (สฺยา)} โข อานนฺท ภิกฺขุ นโว ภวิสฺสติ อจิรปพฺพชิโต, เถโร วา ปน พาโล อพฺยตฺโต.\csromanspacer{} กถํ หิ นาม ยํ มยา เอกํเสน พฺยากตํ ตตฺถ เทฺวชฺฌํ อาปชฺชิสฺสติ, นาหํ อานนฺท อญฺญํ เอกปุคฺคลมฺปิ สมนุปสฺสามิ, โย เอวํ มยา สพฺพํ เจตโส สมนฺนาหริตฺวา พฺยากโต ยถยิทํ เทวทตฺโต.\csromanspacer{} ยาวกีวญฺจาหํ อานนฺท เทวทตฺตสฺส วาลคฺคโกฏินิตฺตุทนมตฺตมฺปิ สุกฺกธมฺมํ อทฺทสํ, เนว ตาวาหํ เทวทตฺตํ พฺยากาสิํ “อาปายิโก เทวทตฺโต เนรยิโก กปฺปฏฺโฐ อเตกิจฺโฉ”ติ.\csromanspacer{} ยโต จ โข อหํ อานนฺท เทวทตฺตสฺส วาลคฺคโกฏินิตฺตุทนมตฺตมฺปิ สุกฺกธมฺมํ น อทฺทสํ, อถาหํ เทวทตฺตํ พฺยากาสิํ “อาปายิโก เทวทตฺโต เนรยิโก กปฺปฏฺโฐ อเตกิจฺโฉ”ติ.}

\prose{เสยฺยาถาปิ อานนฺท คูถกูโป สาธิกโปริโส ปูโร คูถสฺส สมติตฺติโก, ตตฺร ปุริโส สสีสโก นิมุคฺโค อสฺส, ตสฺส โกจิเทว ปุริโส อุปฺปชฺเชยฺย อตฺถกาโม หิตกาโม โยคกฺเขมกาโม ตมฺหา คูถกูปา อุทฺธริตุกาโม.\csromanspacer{} โส ตํ คูถกูปํ สมนฺตานุปริคจฺฉนฺโต เนว ปสฺเสยฺย ตสฺส ปุริสสฺส วาลคฺคโกฏินิตฺตุทนมตฺตมฺปิ คูเถน อมกฺขิตํ, ยตฺถ ตํ คเหตฺวา อุทฺธเรยฺย.\csromanspacer{} เอวเมวํ โข อหํ อานนฺท ยโต เทวทตฺตสฺส วาลคฺคโกฏินิตฺตุทนมตฺตมฺปิ สุกฺกธมฺมํ น อทฺทสํ, อถาหํ เทวทตฺตํ พฺยากาสิํ “อาปายิโก เทวทตฺโต เนรยิโก กปฺปฏฺโฐ อเตกิจฺโฉ”ติ.\csromanspacer{} สเจ ตุเมฺห อานนฺท สุเณยฺยาถ ตถาคตสฺส ปุริสินฺทฺริยญาณานิ วิภชิสฺสามีติ\footnote{วิภชนฺตสฺสาติ (สี, สฺยา, อิ)}.}

\prose{เอตสฺส ภควา กาโล, เอตสฺส สุคต กาโล, ยํ ภควา ปุริสินฺทฺริยญาณานิ วิภเชยฺย, ภควโต สุตฺวา ภิกฺขู ธาเรสฺสนฺตีติ.\csromanspacer{} เตนหานนฺท สุณาหิ สาธุกํ มนสิ กโรหิ, ภาสิสฺสามีติ. “เอวํ ภนฺเต”ติ โข อายสฺมา อานนฺโท ภควโต ปจฺจสฺโสสิ.\csromanspacer{} ภควา เอตทโวจ\csromandash{}}

\prose{อิธาหํ อานนฺท เอกจฺจํ ปุคฺคลํ เอวํ เจตสา เจโต ปริจฺจ ปชานามิ “อิมสฺส โข ปุคฺคลสฺส วิชฺชมานา กุสลาปิ ธมฺมา อกุสลาปิ ธมฺมา”ติ.\csromanspacer{} ตเมนํ อปเรน สมเยน เอวํ เจตสา เจโต ปริจฺจ ปชานามิ “อิมสฺส โข ปุคฺคลสฺส กุสลา ธมฺมา อนฺตรหิตา, อกุสลา ธมฺมา}

\vspace{\tipitakaparskip}
\csromancenter{ฉกฺกนิปาตปาฬิ สมฺมุขีภูตา.\csromanspacer{} อตฺถิ จ ขฺวสฺส กุสลมูลํ อสมุจฺฉินฺนํ, ตมฺหา ตสฺส กุสลา กุสลํ ปาตุภวิสฺสติ.\csromanspacer{} เอวมยํ ปุคฺคโล อายติํ อปริหานธมฺโม ภวิสฺสตี”ติ.\csromanspacer{} เสยฺยถาปิ อานนฺท พีชานิ อขณฺฑานิ อปูตีนิ อวาตาตปหตานิ สารทานิ สุขสยิตานิ สุเขตฺเต สุปริกมฺมกตาย ภูมิยา นิกฺขิตฺตานิ, ชาเนยฺยาสิ ตฺวํ อานนฺท “อิมานิ พีชานิ วุทฺธิํ วิรูฬฺหิํ เวปุลฺลํ อาปชฺชิสฺสนฺตี”ติ.\csromanspacer{} เอวํ ภนฺเต.\csromanspacer{} เอวเมวํ โข อหํ อานนฺท อิเธกจฺจํ ปุคฺคลํ เอวํ เจตสา เจโต ปริจฺจ ปชานามิ “อิมสฺส โข ปุคฺคลสฺส วิชฺชมานา กุสลาปิ ธมฺมา อกุสลาปิ ธมฺมา”ติ.\csromanspacer{} ตเมนํ อปเรน สมเยน เอวํ เจตสา เจโต ปริจฺจ ปชานามิ “อิมสฺส โข ปุคฺคลสฺส กุสลา ธมฺมา อนฺตรหิตา, อกุสลา ธมฺมา สมฺมุขีภูตา.\csromanspacer{} อตฺถิ จ ขฺวสฺส กุสลมูลํ อสมุจฺฉินฺนํ, ตมฺหา ตสฺส กุสลา กุสลํ ปาตุภวิสฺสติ.\csromanspacer{} เอวมยํ ปุคฺคโล อายติํ อปริหานธมฺโม ภวิสฺสตี”ติ.\csromanspacer{} เอวมฺปิ โข อานนฺท ตถาคตสฺส ปุริสปุคฺคโล เจตสา เจโต ปริจฺจ วิทิโต โหติ, เอวมฺปิ โข อานนฺท ตถาคตสฺส ปุริสินฺทฺริยญาณํ เจตสา เจโต ปริจฺจ วิทิตํ โหติ, เอวมฺปิ โข อานนฺท ตถาคตสฺส อายติํ ธมฺมสมุปฺปาโท เจตสา เจโต ปริจฺจ วิทิโต โหติ. (1)}
\prose{\csromanfolio{354}อิธ ปนาหํ อานนฺท เอกจฺจํ ปุคฺคลํ เอวํ เจตสา เจโต ปริจฺจ ปชานามิ “อิมสฺส โข ปุคฺคลสฺส วิชฺชมานา กุสลาปิ ธมฺมา อกุสลาปิ ธมฺมา”ติ.\csromanspacer{} ตเมนํ อปเรน สมเยน เอวํ เจตสา เจโต ปริจฺจ ปชานามิ “อิมสฺส โข ปุคฺคลสฺส อกุสลา ธมฺมา อนฺตรหิตา, กุสลา ธมฺมา สมฺมุขีภูตา.\csromanspacer{} อตฺถิ จ ขฺวสฺส อกุสลมูลํ อสมุจฺฉินฺนํ, ตมฺหา ตสฺส อกุสลา อกุสลํ ปาตุภวิสฺสติ.\csromanspacer{} เอวมยํ ปุคฺคโล อายติํ ปริหานธมฺโม ภวิสฺสตี”ติ.\csromanspacer{} เสยฺยถาปิ อานนฺท พีชานิ อขณฺฑานิ อปูตีนิ อวาตาตปหตานิ สารทานิ สุขสยิตานิ ปุถุสิลาย นิกฺขิตฺตานิ, ชาเนยฺยาสิ ตฺวํ อานนฺท “นยิมานิ พีชานิ วุทฺธิํ วิรูฬฺหิํ เวปุลฺลํ อาปชฺชิสฺสนฺตี”ติ.\csromanspacer{} เอวํ ภนฺเต.\csromanspacer{} เอวเมวํ โข อหํ อานนฺท อิเธกจฺจํ ปุคฺคลํ เอวํ เจตสา เจโต ปริจฺจ ปชานามิ “อิมสฺส โข ปุคฺคลสฺส วิชฺชมานา กุสลาปิ ธมฺมา อกุสลาปิ ธมฺมา”ติ.\csromanspacer{} ตเมนํ อปเรน สมเยน เอวํ เจตสา เจโต ปริจฺจ ปชานามิ “อิมสฺส โข ปุคฺคลสฺส อกุสลา ธมฺมา อนฺตรหิตา, กุสลา ธมฺมา สมฺมุขีภูตา.\csromanspacer{} อตฺถิ จ ขฺวสฺส อกุสลมูลํ อสมุจฺฉินฺนํ, ตมฺหา ตสฺส อกุสลา อกุสลํ}

\vspace{\tipitakaparskip}
\csromancenter{มหาวคฺค ปาตุภวิสฺสติ.\csromanspacer{} เอวมยํ ปุคฺคโล อายติํ ปริหานธมฺโม ภวิสฺสตี”ติ.\csromanspacer{} เอวมฺปิ โข อานนฺท ตถาคตสฺส ปุริสปุคฺคโล เจตสา เจโต ปริจฺจ วิทิโต โหติ, เอวมฺปิ โข อานนฺท ตถาคตสฺส ปุริสินฺทฺริยญาณํ เจตสา เจโต ปริจฺจ วิทิตํ โหติ, เอวมฺปิ โข อานนฺท ตถาคตสฺส อายติํ ธมฺมสมุปฺปาโท เจตสา เจโต ปริจฺจ วิทิโต โหติ. (2)}
\prose{\csromanfolio{355}อิธ ปนาหํ อานนฺท เอกจฺจํ ปุคฺคลํ เอวํ เจตสา เจโต ปริจฺจ ปชานามิ “อิมสฺส โข ปุคฺคลสฺส วิชฺชมานา กุสลาปิ ธมฺมา อกุสลาปิ ธมฺมา”ติ.\csromanspacer{} ตเมนํ อปเรน สมเยน เอวํ เจตสา เจโต ปริจฺจ ปชานามิ “นตฺถิ อิมสฺส ปุคฺคลสฺส วาลคฺคโกฏินิตฺตุทนมตฺโตปิ สุกฺโก ธมฺโม, สมนฺนาคโตยํ ปุคฺคโล เอกนฺตกาฬเกหิ อกุสเลหิ ธมฺเมหิ, กายสฺส เภทา ปรํ มรณา อปายํ ทุคฺคติํ วินิปาตํ นิรยํ อุปปชฺชิสฺสตี”ติ.\csromanspacer{} เสยฺยาถาปิ อานนฺท พีชานิ ขณฺฑานิ ปูตีนิ วาตาตปหตานิ สุเขตฺเต สุปริกมฺมกตาย ภูมิยา นิกฺขิตฺตานิ, ชาเนยฺยาสิ ตฺวํ อานนฺท “นยิมานิ พีชานิ วุทฺธิํ วิรูฬฺหิํ เวปุลฺลํ อาปชฺชิสฺสนฺตี”ติ.\csromanspacer{} เอวํ ภนฺเต.\csromanspacer{} เอวเมวํ โข อหํ อานนฺท อิเธกจฺจํ ปุคฺคลํ เอวํ เจตสา เจโต ปริจฺจ ปชานามิ “อิมสฺส โข ปุคฺคลสฺส วิชฺชมานา กุสลาปิ ธมฺมา อกุสลาปิ ธมฺมา”ติ.\csromanspacer{} ตเมนํ อปเรน สมเยน เอวํ เจตสา เจโต ปริจฺจ ปชานามิ “นตฺถิ อิมสฺส ปุคฺคลสฺส วาลคฺคโกฏิ นิตฺตุทนมตฺโตปิ สุกฺโก ธมฺโม, สมนฺนาคโตยํ ปุคฺคโล เอกนฺตกาฬเกหิ อกุสเลหิ ธมฺเมหิ, กายสฺส เภทา ปรํ มรณา อปายํ ทุคฺคติํ วินิปาตํ นิรยํ อุปปชฺชิสฺสตี”ติ.\csromanspacer{} เอวมฺปิ โข อานนฺท ตถาคตสฺส ปุริสปุคฺคโล เจตสา เจโต ปริจฺจ วิทิโต โหติ, เอวมฺปิ โข อานนฺท ตถาคตสฺส ปุริสินฺทฺริยญาณํ เจตสา เจโต ปริจฺจ วิทิตํ โหติ, เอวมฺปิ โข อานนฺท ตถาคตสฺส อายติํ ธมฺมสมุปฺปาโท เจตสา เจโต ปริจฺจ วิทิโต โหตีติ. (3)}

\prose{เอวํ วุตฺเต อายสฺมา อานนฺโท ภควนฺตํ เอตทโวจ\csromandash{}สกฺกา นุ โข ภนฺเต อิเมสํ ติณฺณํ ปุคฺคลานํ อปเรปิ ตโย ปุคฺคลา สปฺปฏิภาคา ปญฺญาเปตุนฺติ. “สกฺกา อานนฺทา”ติ ภควา อโวจ.\csromanspacer{} อิธาหํ อานนฺท เอกจฺจํ ปุคฺคลํ เอวํ เจตสา เจโต ปริจฺจ ปชานามิ “อิมสฺส โข ปุคฺคลสฺส วิชฺชมานา กุสลาปิ ธมฺมา อกุสลาปิ ธมฺมา”ติ.\csromanspacer{} ตเมนํ อปเรน สมเยน เอวํ เจตสา เจโต ปริจฺจ ปชานามิ “อิมสฺส โข ปุคฺคลสฺส กุสลา ธมฺมา อนฺตรหิตา, อกุสลา ธมฺมา สมฺมุขีภูตา.\csromanspacer{} อตฺถิ จ ขฺวสฺส}

\vspace{\tipitakaparskip}
\csromancenter{ฉกฺกนิปาตปาฬิ กุสลมูลํ อสมุจฺฉินฺนํ, ตมฺปิ สพฺเพน สพฺพํ สมุคฺฆาตํ คจฺฉติ.\csromanspacer{} เอวมยํ ปุคฺคโล อายติํ ปริหานธมฺโม ภวิสฺสตี”ติ.\csromanspacer{} เสยฺยถาปิ อานนฺท องฺคารานิ อาทิตฺตานิ สมฺปชฺชลิตานิ สโชติภูตานิ ปุถุสิลาย นิกฺขิตฺตานิ, ชาเนยฺยาสิ ตฺวํ อานนฺท “นยิมานิ องฺคารานิ วุทฺธิํ วิรูฬฺหิํ เวปุลฺลํ อาปชฺชิสฺสนฺตี”ติ.\csromanspacer{} เอวํ ภนฺเต.\csromanspacer{} เสยฺยถาปิ วา ปน อานนฺท สายนฺหสมยํ\footnote{สายนฺหสมเย (สฺยา, ก)} สูริเย โอคจฺฉนฺเต ชาเนยฺยาสิ ตฺวํ อานนฺท “อาโลโก อนฺตรธายิสฺสติ, อนฺธกาโร ปาตุภวิสฺสตี”ติ.\csromanspacer{} เอวํ ภนฺเต.\csromanspacer{} เสยฺยถาปิ วา ปนานนฺท อภิโท อทฺธรตฺตํ ภตฺตกาลสมเย ชาเนยฺยาสิ ตฺวํ อานนฺท “อาโลโก อนฺตรหิโต, อนฺธกาโร ปาตุภูโต”ติ.\csromanspacer{} เอวํ ภนฺเต.\csromanspacer{} เอวเมวํ โข อหํ อานนฺท อิเธกจฺจํ ปุคฺคลํ เอวํ เจตสา เจโต ปริจฺจ ปชานามิ “อิมสฺส โข ปุคฺคลสฺส วิชฺชมานา กุสลาปิ ธมฺมา อกุสลาปิ ธมฺมา”ติ.\csromanspacer{} ตเมนํ อปเรน สมเยน เอวํ เจตสา เจโต ปริจฺจ ปชานามิ “อิมสฺส โข ปุคฺคลสฺส กุสลา ธมฺมา อนฺตรหิตา, อกุสลา ธมฺมา สมฺมุขีภูตา.\csromanspacer{} อตฺถิ จ ขฺวสฺส กุสลมูลํ อสมุจฺฉินฺนํ, ตมฺปิ สพฺเพน สพฺพํ สมุคฺฆาตํ คจฺฉติ.\csromanspacer{} เอวมยํ ปุคฺคโล อายติํ ปริหานธมฺโม ภวิสฺสตี”ติ.\csromanspacer{} เอวมฺปิ โข อานนฺท ตถาคตสฺส ปุริสปุคฺคโล เจตสา เจโต ปริจฺจ วิทิโต โหติ, เอวมฺปิ โข อานนฺท ตถาคตสฺส ปุริสินฺทฺริยญาณํ เจตสา เจโต ปริจฺจ วิทิตํ โหติ, เอวมฺปิ โข อานนฺท ตถาคตสฺส อายติํ ธมฺมสมุปฺปาโท เจตสา เจโต ปริจฺจ วิทิโต โหติ. (4)}
\prose{\csromanfolio{356}อิธ ปนาหํ อานนฺท เอกจฺจํ ปุคฺคลํ เอวํ เจตสา เจโต ปริจฺจ ปชานามิ “อิมสฺส โข ปุคฺคลสฺส วิชฺชมานา กุสลาปิ ธมฺมา อกุสลาปิ ธมฺมา”ติ.\csromanspacer{} ตเมนํ อปเรน สมเยน เอวํ เจตสา เจโต ปริจฺจ ปชานามิ “อิมสฺส โข ปุคฺคลสฺส อกุสลา ธมฺมา อนฺตรหิตา, กุสลา ธมฺมา สมฺมุขีภูตา.\csromanspacer{} อตฺถิ จ ขฺวสฺส อกุสลมูลํ อสมุจฺฉินฺนํ, ตมฺปิ สพฺเพน สพฺพํ สมุคฺฆาตํ คจฺฉติ.\csromanspacer{} เอวมยํ ปุคฺคโล อายติํ อปริหานธมฺโม ภวิสฺสตี”ติ.\csromanspacer{} เสยฺยถาปิ อานนฺท องฺคารานิ อาทิตฺตานิ สมฺปชฺชลิตานิ สโชติภูตานิ สุกฺเข ติณปุญฺเช วา กฏฺฐปุญฺเช วา นิกฺขิตฺตานิ, ชาเนยฺยาสิ ตฺวํ อานนฺท “อิมานิ องฺคารานิ วุฑฺฒิํ วิรูฬฺหิํ เวปุลฺลํ อาปชฺชิสฺสนฺตี”ติ.\csromanspacer{} เอวํ ภนฺเต.\csromanspacer{} เสยฺยถาปิ}

\vspace{\tipitakaparskip}
\csromancenter{มหาวคฺค วา ปนานนฺท รตฺติยา ปจฺจูสสมยํ\footnote{รตฺติปจฺจูสสมเย (ก)} สูริเย อุคฺคจฺฉนฺเต ชาเนยฺยาสิ ตฺวํ อานนฺท “อนฺธกาโร อนฺตรธายิสฺสติ, อาโลโก ปาตุภวิสฺสตี”ติ.\csromanspacer{} เอวํ ภนฺเต.\csromanspacer{} เสยฺยถาปิ วา ปนานนฺท อภิโท มชฺฌนฺหิเก ภตฺตกาลสมเย ชาเนยฺยาสิ ตฺวํ อานนฺท “อนฺธกาโร อนฺตรหิโต, อาโลโก ปาตุภูโต”ติ.\csromanspacer{} เอวํ ภนฺเต.\csromanspacer{} เอวเมวํ โข อหํ อานนฺท อิเธกจฺจํ ปุคฺคลํ เอวํ เจตสา เจโต ปริจฺจ ปชานามิ “อิมสฺส โข ปุคฺคลสฺส วิชฺชมานา กุสลาปิ ธมฺมา อกุสลาปิ ธมฺมา”ติ.\csromanspacer{} ตเมนํ อปเรน สมเยน เอวํ เจตสา เจโต ปริจฺจ ปชานามิ “อิมสฺส โข ปุคฺคลสฺส อกุสลา ธมฺมา อนฺตรหิตา, กุสลา ธมฺมา สมฺมุขีภูตา.\csromanspacer{} อตฺถิ จ ขฺวสฺส อกุสลมูลํ อสมุจฺฉินฺนํ, ตมฺปิ สพฺเพน สพฺพํ สมุคฺฆาตํ คจฺฉติ.\csromanspacer{} เอวมยํ ปุคฺคโล อายติํ อปริหานธมฺโม ภวิสฺสตี”ติ.\csromanspacer{} เอวมฺปิ โข อานนฺท ตถาคตสฺส ปุริสปุคฺคโล เจตสา เจโต ปริจฺจ วิทิโต โหติ, เอวมฺปิ โข อานนฺท ตถาคตสฺส ปุริสินฺทฺริยญาณํ เจตสา เจโต ปริจฺจ วิทิตํ โหติ, เอวมฺปิ โข อานนฺท ตถาคตสฺส อายติํ ธมฺมสมุปฺปาโท เจตสา เจโต ปริจฺจ วิทิโต โหติ. (5)}
\prose{\csromanfolio{357}อิธ ปนาหํ อานนฺท เอกจฺจํ ปุคฺคลํ เจตสา เจโต ปริจฺจ ปชานามิ “อิมสฺส โข ปุคฺคลสฺส วิชฺชมานา กุสลาปิ ธมฺมา อกุสลาปิ ธมฺมา”ติ.\csromanspacer{} ตเมนํ อปเรน สมเยน เอวํ เจตสา เจโต ปริจฺจ ปชานามิ “นตฺถิ อิมสฺส ปุคฺคลสฺส วาลคฺคโกฏิ นิตฺตุทนมตฺโตปิ อกุสโล ธมฺโม, สมนฺนาคโตยํ ปุคฺคโล เอกนฺตสุกฺเกหิ อนวชฺเชหิ ธมฺเมหิ, ทิฏฺเฐว ธมฺเม ปรินิพฺพายิสฺสตี”ติ.\csromanspacer{} เสยฺยถาปิ อานนฺท องฺคารานิ สีตานิ นิพฺพุตานิ สุกฺเข ติณปุญฺเช วา กฏฺฐปุญฺเช วา นิกฺขิตฺตานิ, ชาเนยฺยาสิ ตฺวํ อานนฺท “นยิมานิ องฺคารานิ วุทฺธิํ วิรูฬฺหิํ เวปุลฺลํ อาปชฺชิสฺสนฺตี”ติ.\csromanspacer{} เอวํ ภนฺเต.\csromanspacer{} เอวเมวํ โข อหํ อานนฺท อิเธกจฺจํ ปุคฺคลํ เอวํ เจตสา เจโต ปริจฺจ ปชานามิ “อิมสฺส โข ปุคฺคลสฺส วิชฺชมานา กุสลาปิ ธมฺมา อกุสลาปิ ธมฺมา”ติ.\csromanspacer{} ตเมนํ อปเรน สมเยน เอวํ เจตสา เจโต ปริจฺจ ปชานามิ “นตฺถิ อิมสฺส ปุคฺคลสฺส วาลคฺคโกฏินิตฺตุทนมตฺโตปิ อกุสโล ธมฺโม, สมนฺนาคโตยํ ปุคฺคโล เอกนฺตสุกฺเกหิ อนวชฺเชหิ ธมฺเมหิ, ทิฏฺเฐว ธมฺเม ปรินิพฺพายิสฺสตี”ติ.\csromanspacer{} เอวมฺปิ โข อานนฺท ตถาคตสฺส ปุริสปุคฺคโล เจตสา เจโต ปริจฺจ วิทิโต โหติ, เอวมฺปิ โข อานนฺท ตถาคตสฺส}

\vspace{\tipitakaparskip}
\csromancenter{ฉกฺกนิปาตปาฬิ ปุริสินฺทฺริยญาณํ เจตสา เจโต ปริจฺจ วิทิตํ โหติ, เอวมฺปิ โข อานนฺท ตถาคตสฺส อายติํ ธมฺมสมุปฺปาโท เจตสา เจโต ปริจฺจ วิทิโต โหติ. (6)}
\prose{\csromanfolio{358}ตตฺรานนฺท เย เต ปุริมา ตโย ปุคฺคลา, เตสํ ติณฺณํ ปุคฺคลานํ เอโก อปริหานธมฺโม, เอโก ปริหานธมฺโม, เอโก อาปายิโก เนรยิโก.\csromanspacer{} ตตฺรานนฺท เยเม ปจฺฉิมา ตโย ปุคฺคลา, อิเมสํ ติณฺณํ ปุคฺคลานํ เอโก ปริหานธมฺโม, เอโก อปริหานธมฺโม, เอโก ปรินิพฺพานธมฺโมติ.\csromanspacer{}\csromanspacer{}อฏฺฐมํ.}

\csromanmatikaanchor{mtk.185}
\csromantocmarkref{subsection}{9. นิพฺเพธิกสุตฺต}{mtk.185}
\bookmark[dest={mtk.185},level=2]{9. นิพฺเพธิกสุตฺต}
\csromanheader{2}{9. นิพฺเพธิกสุตฺต}
\proseitem{63}{นิพฺเพธิกปริยายํ โว ภิกฺขเว ธมฺมปริยายํ เทเสสฺสามิ, ตํ สุณาถ สาธุกํ มนสิ กโรถ, ภาสิสฺสามีติ. “เอวํ ภนฺเต”ติ โข เต ภิกฺขู ภควโต ปจฺจสฺโสสุํ.\csromanspacer{} ภควา เอตทโวจ\csromandash{}}

\prose{กตโม จ โส ภิกฺขเว นิพฺเพธิกปริยาโย ธมฺมปริยาโย? กามา ภิกฺขเว เวทิตพฺพา, กามานํ นิทานสมฺภโว เวทิตพฺโพ, กามานํ เวมตฺตตา เวทิตพฺพา, กามานํ วิปาโก เวทิตพฺโพ, กามนิโรโธ\footnote{กามานํ นิโรโธ (ก) เอวํ เวทนานิโรโธ-อิจฺจาทีสุปิ.} เวทิตพฺโพ, กามนิโรธคามินี\footnote{กามานํ นิโรธคามินี (ก) เอวํ เวทนานิโรธคามินี-อิจฺจาทีสุปิ.} ปฏิปทา เวทิตพฺพา. (1)}

\prose{เวทนา ภิกฺขเว เวทิตพฺพา, เวทนานํ นิทานสมฺภโว เวทิตพฺโพ, เวทนานํ เวมตฺตตา เวทิตพฺพา, เวทนานํ วิปาโก เวทิตพฺโพ, เวทนานิโรโธ เวทิตพฺโพ, เวทนานิโรธคามินี ปฏิปทา เวทิตพฺพา. (2)}

\prose{สญฺญา ภิกฺขเว เวทิตพฺพา, สญฺญานํ นิทานสมฺภโว เวทิตพฺโพ, สญฺญานํ เวมตฺตตา เวทิตพฺพา, สญฺญานํ วิปาโก เวทิตพฺโพ, สญฺญานิโรโธ เวทิตพฺโพ, สญฺญานิโรธคามินี ปฏิปทา เวทิตพฺพา. (3)}

\prose{อาสวา ภิกฺขเว เวทิตพฺพา, อาสวานํ นิทานสมฺภโว เวทิตพฺโพ, อาสวานํ เวมตฺตตา เวทิตพฺพา, อาสวานํ วิปาโก เวทิตพฺโพ, อาสวนิโรโธ เวทิตพฺโพ, อาสวนิโรธคามินี ปฏิปทา เวทิตพฺพา. (4)}

\chapterheadpageread{6. มหาวคฺค}
\prose{\csromanfolio{359}กมฺมํ ภิกฺขเว เวทิตพฺพํ, กมฺมานํ นิทานสมฺภโว เวทิตพฺโพ, กมฺมานํ เวมตฺตตา เวทิตพฺพา, กมฺมานํ วิปาโก เวทิตพฺโพ, กมฺมนิโรโธ เวทิตพฺโพ, กมฺมนิโรธคามินี ปฏิปทา เวทิตพฺพา. (5)}

\prose{ทุกฺขํ ภิกฺขเว เวทิตพฺพํ, ทุกฺขสฺส นิทานสมฺภโว เวทิตพฺโพ, ทุกฺขสฺส เวมตฺตตา เวทิตพฺพา, ทุกฺขสฺส วิปาโก เวทิตพฺโพ, ทุกฺขนิโรโธ เวทิตพฺโพ, ทุกฺขนิโรธคามินี ปฏิปทา เวทิตพฺพา. (6)}

\prose{“กามา ภิกฺขเว เวทิตพฺพา, กามานํ นิทานสมฺภโว เวทิตพฺโพ, กามานํ เวมตฺตตา เวทิตพฺพา, กามานํ วิปาโก เวทิตพฺโพ, กามนิโรโธ เวทิตพฺโพ, กามนิโรธคามินี ปฏิปทา เวทิตพฺพา”ติ อิติ โข ปเนตํ วุตฺตํ, กิญฺเจตํ ปฏิจฺจ วุตฺตํ. + ปญฺจิเม ภิกฺขเว กามคุณา จกฺขุวิญฺเญยฺยา รูปา อิฏฺฐา กนฺตา มนาปา ปิยรูปา กามูปสํหิตา รชนียา.\csromanspacer{} โสตวิญฺเญยฺยา สทฺทา.\csromanspacer{} ฆานวิญฺเญยฺยา คนฺธา.\csromanspacer{} ชิวฺหาวิญฺเญยฺยา รสา.\csromanspacer{} กายวิญฺเญยฺยา โผฏฺฐพฺพา อิฏฺฐา กนฺตา มนาปา ปิยรูปา กามูปสํหิตา รชนียา.\csromanspacer{} อปิ จ โข ภิกฺขเว เนเต กามา, กามคุณา นาเมเต\footnote{เต กามคุณา นาม เนเต กามา (ก)} อริยสฺส วินเย วุจฺจนฺติ.}

\csromangathameasure{*สงฺกปฺปราโค ปุริสสฺส กาโม, \\ เนเต กามา ยานิ จิตฺรานิ โลเก. \\ สงฺกปฺปราโค ปุริสสฺส กาโม, \\ ติฏฺฐนฺติ จิตฺรานิ ตเถว โลเก. \\ อเถตฺถ ธีรา วินยนฺติ ฉนฺทนฺติ.}
\csromangathagroup{\csromangathastanza{\csromansymbolfootnote{*}{อภิ 4. 275 ปิฏฺเฐ.}สงฺกปฺปราโค ปุริสสฺส กาโม, \\ เนเต\footnote{น เต (สฺยา)} กามา ยานิ จิตฺรานิ โลเก. \\ สงฺกปฺปราโค ปุริสสฺส กาโม, \\ ติฏฺฐนฺติ จิตฺรานิ ตเถว โลเก.}\csromangathastanza{อเถตฺถ ธีรา วินยนฺติ ฉนฺทนฺติ.}}

\prose{กตโม จ ภิกฺขเว กามานํ นิทานสมฺภโว? ผสฺโส ภิกฺขเว กามานํ นิทานสมฺภโว.}

\prose{กตมา จ ภิกฺขเว กามานํ เวมตฺตตา? อญฺโญ ภิกฺขเว กาโม รูเปสุ.\csromanspacer{} อญฺโญ กาโม สทฺเทสุ.\csromanspacer{} อญฺโญ กาโม คนฺเธสุ.\csromanspacer{} อญฺโญ กาโม รเสสุ.\csromanspacer{} อญฺโญ กาโม โผฏฺฐพฺเพสุ.\csromanspacer{} อยํ วุจฺจติ ภิกฺขเว กามานํ เวมตฺตตา.}

\gambhira{ฉกฺกนิปาตปาฬิ}
\prose{\csromanfolio{360}กตโม จ ภิกฺขเว กามานํ วิปาโก? ยํ โข ภิกฺขเว กามยมาโน ตชฺชํ ตชฺชํ อตฺตภาวํ อภินิพฺพตฺเตติ ปุญฺญภาคิยํ วา อปุญฺญภาคิยํ วา.\csromanspacer{} อยํ วุจฺจติ ภิกฺขเว กามานํ วิปาโก.}

\prose{กตโม จ ภิกฺขเว กามนิโรโธ? ผสฺสนิโรโธ\footnote{ผสฺสนิโรธา (สฺยา)} ภิกฺขเว กามนิโรโธ.\csromanspacer{} อยเมว อริโย อฏฺฐงฺคิโก มคฺโค กามนิโรธคามินี ปฏิปทา.\csromanspacer{} เสยฺยถิทํ, สมฺมาทิฏฺฐิ สมฺมาสงฺกปฺโป สมฺมาวาจา สมฺมากมฺมนฺโต สมฺมา-อาชีโว สมฺมาวายาโม สมฺมาสติ สมฺมาสมาธิ.}

\prose{ยโต โข\footnote{ยโต จ โข (พหูสุ) 3. สพฺพตฺถปิ เอวเมว ทิสฺสติ.} ภิกฺขเว อริยสาวโก เอวํ กาเม ปชานาติ, เอวํ กามานํ นิทานสมฺภวํ ปชานาติ, เอวํ กามานํ เวมตฺตตํ ปชานาติ, เอวํ กามานํ วิปากํ ปชานาติ, เอวํ กามนิโรธํ ปชานาติ, เอวํ กามนิโรธคามินิํ ปฏิปทํ ปชานาติ.\csromanspacer{} โส อิมํ นิพฺเพธิกํ พฺรหฺมจริยํ ปชานาติ กามนิโรธํ. “กามา ภิกฺขเว เวทิตพฺพา ฯเปฯ กามนิโรธคามินี3 ปฏิปทา เวทิตพฺพา”ติ อิติ ยํ ตํ วุตฺตํ, อิทเมตํ ปฏิจฺจ วุตฺตํ. (1)}

\prose{“เวทนา ภิกฺขเว เวทิตพฺพา ฯเปฯ เวทนานิโรธคามินี ปฏิปทา เวทิตพฺพา”ติ อิติ โข ปเนตํ วุตฺตํ, กิญฺเจตํ ปฏิจฺจ วุตฺตํ.\csromanspacer{} ติสฺโส อิมา ภิกฺขเว เวทนา\csromandash{}สุขา เวทนา ทุกฺขา เวทนา อทุกฺขมสุขา เวทนา.}

\prose{กตโม จ ภิกฺขเว เวทนานํ นิทานสมฺภโว? ผสฺโส ภิกฺขเว เวทนานํ นิทานสมฺภโว.}

\prose{กตมา จ ภิกฺขเว เวทนานํ เวมตฺตตา? อตฺถิ ภิกฺขเว สามิสา สุขา เวทนา, อตฺถิ นิรามิสา สุขา เวทนา, อตฺถิ สามิสา ทุกฺขา เวทนา, อตฺถิ นิรามิสา ทุกฺขา เวทนา, อตฺถิ สามิสา อทุกฺขมสุขา เวทนา, อตฺถิ นิรามิสา อทุกฺขมสุขา เวทนา.\csromanspacer{} อยํ วุจฺจติ ภิกฺขเว เวทนานํ เวมตฺตตา.}

\prose{กตโม จ ภิกฺขเว เวทนานํ วิปาโก? ยํ โข ภิกฺขเว เวทิยมาโน\footnote{เวทยมาโน (สฺยา, กํ) อํ 1. 554 ปิฏฺเฐปิ.} ตชฺชํ ตชฺชํ อตฺตภาวํ อภินิพฺพตฺเตติ ปุญฺญภาคิยํ วา อปุญฺญภาคิยํ วา.\csromanspacer{} อยํ วุจฺจติ ภิกฺขเว เวทนานํ วิปาโก.}

\chapterheadpageread{6. มหาวคฺค}
\prose{\csromanfolio{361}กตโม จ ภิกฺขเว เวทนานิโรโธ? ผสฺสนิโรโธ\footnote{ผสฺสนิโรธา (สฺยา, กํ, ก)} ภิกฺขเว เวทนานิโรโธ.\csromanspacer{} อยเมว อริโย อฏฺฐงฺคิโก มคฺโค เวทนานิโรธคามินี ปฏิปทา.\csromanspacer{} เสยฺยถิทํ, สมฺมาทิฏฺฐิ ฯเปฯ สมฺมาสมาธิ.}

\prose{ยโต โข ภิกฺขเว อริยสาวโก เอวํ เวทนํ ปชานาติ, เอวํ เวทนานํ นิทานสมฺภวํ ปชานาติ, เอวํ เวทนานํ เวมตฺตตํ ปชานาติ, เอวํ เวทนานํ วิปากํ ปชานาติ, เอวํ เวทนานิโรธํ ปชานาติ, เอวํ เวทนานิโรธคามินิํ ปฏิปทํ ปชานาติ.\csromanspacer{} โส อิมํ นิพฺเพธิกํ พฺรหฺมจริยํ ปชานาติ เวทนานิโรธํ. “เวทนา ภิกฺขเว เวทิตพฺพา ฯเปฯ เวทนานิโรธคามินี ปฏิปทา เวทิตพฺพา”ติ อิติ ยํ ตํ วุตฺตํ, อิทเมตํ ปฏิจฺจ วุตฺตํ. (2)}

\prose{“สญฺญา ภิกฺขเว เวทิตพฺพา ฯเปฯ สญฺญานิโรธคามินี ปฏิปทา เวทิตพฺพา”ติ อิติ โข ปเนตํ วุตฺตํ, กิญฺเจตํ ปฏิจฺจ วุตฺตํ, ฉยิมา ภิกฺขเว สญฺญา รูปสญฺญา สทฺทสญฺญา คนฺธสญฺญา รสสญฺญา โผฏฺฐพฺพสญฺญา ธมฺมสญฺญา.}

\prose{กตโม จ ภิกฺขเว สญฺญานํ นิทานสมฺภโว? ผสฺโส ภิกฺขเว สญฺญานํ นิทานสมฺภโว.}

\prose{กตมา จ ภิกฺขเว สญฺญานํ เวมตฺตตา? อญฺญา ภิกฺขเว สญฺญา รูเปสุ.\csromanspacer{} อญฺญา สญฺญา สทฺเทสุ\footnote{อญฺญา ภิกฺขเว รูเปสุ สญฺญา อญฺญา สทฺเทสุ สญฺญา (ก) เอวํ เสเสสุปิ.}.\csromanspacer{} อญฺญา สญฺญา คนฺเธสุ.\csromanspacer{} อญฺญา สญฺญา รเสสุ.\csromanspacer{} อญฺญา สญฺญา โผฏฺฐพฺเพสุ.\csromanspacer{} อญฺญา สญฺญา ธมฺเมสุ.\csromanspacer{} อยํ วุจฺจติ ภิกฺขเว สญฺญานํ เวมตฺตตา.}

\prose{กตโม จ ภิกฺขเว สญฺญานํ วิปาโก? โวหารเวปกฺกํ\footnote{โวหารเวปกฺกาหํ (สฺยา, อิ), โวหารปกฺกาหํ (สี)} ภิกฺขเว สญฺญํ\footnote{สญฺญา (สฺยา, อิ)} วทามิ.\csromanspacer{} ยถา ยถา นํ สญฺชานาติ, ตถา ตถา โวหรติ “เอวํ สญฺญี อโหสินฺ”ติ\footnote{อโหสีติ (ก)}.\csromanspacer{} อยํ วุจฺจติ ภิกฺขเว สญฺญานํ วิปาโก.}

\prose{กตโม จ ภิกฺขเว สญฺญานิโรโธ? ผสฺสนิโรโธ\footnote{ผสฺสนิโรธา (สฺยา, ก)} ภิกฺขเว สญฺญานิโรโธ.\csromanspacer{} อยเมว อริโย อฏฺฐงฺคิโก มคฺโค สญฺญานิโรธคามินี ปฏิปทา.\csromanspacer{} เสยฺยถิทํ, สมฺมาทิฏฺฐิ ฯเปฯ สมฺมาสมาธิ.}

\gambhira{ฉกฺกนิปาตปาฬิ}
\prose{\csromanfolio{362}ยโต โข ภิกฺขเว อริยสาวโก เอวํ สญฺญํ ปชานาติ, เอวํ สญฺญานํ นิทานสมฺภวํ ปชานาติ, เอวํ สญฺญานํ เวมตฺตตํ ปชานาติ, เอวํ สญฺญานํ วิปากํ ปชานาติ, เอวํ สญฺญานิโรธํ ปชานาติ, เอวํ สญฺญานิโรธคามินิํ ปฏิปทํ ปชานาติ.\csromanspacer{} โส อิมํ นิพฺเพธิกํ พฺรหฺมจริยํ ปชานาติ สญฺญานิโรธํ. “สญฺญา ภิกฺขเว เวทิตพฺพา ฯเปฯ สญฺญานิโรธคามินี ปฏิปทา เวทิตพฺพา”ติ อิติ ยํ ตํ วุตฺตํ, อิทเมตํ ปฏิจฺจ วุตฺตํ. (3)}

\prose{"อาสวา ภิกฺขเว เวทิตพฺพา ฯเปฯ อาสวนิโรธคามินี\footnote{สพฺพตฺถปิ เอวเมว ทิสฺสติ.} ปฏิปทา เวทิตพฺพา”ติ อิติ โข ปเนตํ วุตฺตํ, กิญฺเจตํ ปฏิจฺจ วุตฺตํ.\csromanspacer{} ตโยเม ภิกฺขเว อาสวา กามาสโว ภวาสโว อวิชฺชาสโว.}

\prose{กตโม จ ภิกฺขเว อาสวานํ นิทานสมฺภโว? อวิชฺชา ภิกฺขเว อาสวานํ นิทานสมฺภโว.}

\prose{กตมา จ ภิกฺขเว อาสวานํ เวมตฺตตา? อตฺถิ ภิกฺขเว อาสวา นิรยคมนียา\footnote{นิรยคามินิยา (สี, ก)}, อตฺถิ อาสวา ติรจฺฉานโยนิคมนียา, อตฺถิ อาสวา เปตฺติวิสยคมนียา, อตฺถิ อาสวา มนุสฺสโลกคมนียา, อตฺถิ อาสวา เทวโลกคมนียา.\csromanspacer{} อยํ วุจฺจติ ภิกฺขเว อาสวานํ เวมตฺตตา.}

\prose{กตโม จ ภิกฺขเว อาสวานํ วิปาโก? ยํ โข ภิกฺขเว อวิชฺชาคโต ตชฺชํ ตชฺชํ อตฺตภาวํ อภินิพฺพตฺเตติ ปุญฺญภาคิยํ วา อปุญฺญภาคิยํ วา.\csromanspacer{} อยํ วุจฺจติ ภิกฺขเว อาสวานํ วิปาโก.}

\prose{กตโม จ ภิกฺขเว อาสวนิโรโธ? อวิชฺชานิโรโธ ภิกฺขเว อาสวนิโรโธ.\csromanspacer{} อยเมว อริโย อฏฺฐงฺคิโก มคฺโค อาสวนิโรธคามินี1 ปฏิปทา.\csromanspacer{} เสยฺยถิทํ, สมฺมาทิฏฺฐิ ฯเปฯ สมฺมาสมาธิ.}

\prose{ยโต โข ภิกฺขเว อริยสาวโก เอวํ อาสเว ปชานาติ, เอวํ อาสวานํ นิทานสมฺภวํ ปชานาติ, เอวํ อาสวานํ เวมตฺตตํ ปชานาติ, เอวํ อาสวานํ วิปากํ ปชานาติ, เอวํ อาสวานํ นิโรธํ ปชานาติ, เอวํ อาสวานํ นิโรธคามินิํ ปฏิปทํ ปชานาติ.\csromanspacer{} โส อิมํ นิพฺเพธิกํ พฺรหฺมจริยํ ปชานาติ อาสวนิโรธํ. “อาสวา ภิกฺขเว เวทิตพฺพา ฯเปฯ.}

\vspace{\tipitakaparskip}
\csromancenter{มหาวคฺค อาสวนิโรธคามินี\footnote{สพฺพตฺถปิ เอวเมว ทิสฺสติ.} ปฏิปทา เวทิตพฺพา”ติ อิติ ยํ ตํ วุตฺตํ, อิทเมตํ ปฏิจฺจ วุตฺตํ. (4)}
\prose{\csromanfolio{363}“กมฺมํ ภิกฺขเว เวทิตพฺพํ ฯเปฯ กมฺมนิโรธคามินี1 ปฏิปทา เวทิตพฺพา”ติ อิติ โข ปเนตํ วุตฺตํ, กิญฺเจตํ ปฏิจฺจ วุตฺตํ. * เจตนาหํ ภิกฺขเว กมฺมํ วทามิ, เจตยิตฺวา กมฺมํ กโรติ กาเยน วาจาย มนสา.}

\prose{กตโม จ ภิกฺขเว กมฺมานํ นิทานสมฺภโว? ผสฺโส ภิกฺขเว กมฺมานํ นิทานสมฺภโว.}

\prose{กตมา จ ภิกฺขเว กมฺมานํ เวมตฺตตา, อตฺถิ ภิกฺขเว กมฺมํ นิรยเวทนียํ, อตฺถิ กมฺมํ ติรจฺฉานโยนิเวทนียํ, อตฺถิ กมฺมํ เปตฺติวิสยเวทนียํ, อตฺถิ กมฺมํ มนุสฺสโลกเวทนียํ, อตฺถิ กมฺมํ เทวโลกเวทนียํ.\csromanspacer{} อยํ วุจฺจติ ภิกฺขเว กมฺมานํ เวมตฺตตา.}

\prose{กตโม จ ภิกฺขเว กมฺมานํ วิปาโก? ติวิธาหํ\footnote{อิมาหํ (ก)} ภิกฺขเว กมฺมานํ วิปากํ วทามิ ทิฏฺเฐว\footnote{ทิฏฺเฐ วา (สี)} ธมฺเม อุปปชฺเช วา\footnote{อุปปชฺชํ วา (ก-สี, อํ 3. 497), อุปปชฺชวา (?), ม 3. 257 ปิฏฺเฐ ปาฬิยา ตทตฺถวณฺณนาย จ สํสนฺเทตพฺพํ.} อปเร วา ปริยาเย.\csromanspacer{} อยํ วุจฺจติ ภิกฺขเว กมฺมานํ วิปาโก.}

\prose{กตโม จ ภิกฺขเว กมฺมนิโรโธ? ผสฺสนิโรโธ\footnote{ผสฺสนิโรธา (ก-สี, สฺยา, ก)} ภิกฺขเว กมฺมนิโรโธ.\csromanspacer{} อยเมว อริโย อฏฺฐงฺคิโก มคฺโค กมฺมนิโรธคามินี\footnote{สพฺพตฺถปิ เอวเมว ทิสฺสติ.} ปฏิปทา.\csromanspacer{} เสยฺยถิทํ, สมฺมาทิฏฺฐิ ฯเปฯ สมฺมาสมาธิ.}

\prose{ยโต โข ภิกฺขเว อริยสาวโก เอวํ กมฺมํ ปชานาติ, เอวํ กมฺมานํ นิทานสมฺภวํ ปชานาติ, เอวํ กมฺมานํ เวมตฺตตํ ปชานาติ, เอวํ กมฺมานํ วิปากํ ปชานาติ, เอวํ กมฺมนิโรธํ ปชานาติ, เอวํ กมฺมนิโรธคามินิํ ปฏิปทํ ปชานาติ.\csromanspacer{} โส อิมํ นิพฺเพธิกํ พฺรหฺมจริยํ ปชานาติ กมฺมนิโรธํ. “กมฺมํ ภิกฺขเว เวทิตพฺพํ ฯเปฯ กมฺมนิโรธคามินี ปฏิปทา เวทิตพฺพา”ติ อิติ ยํ ตํ วุตฺตํ, อิทเมตํ ปฏิจฺจ วุตฺตํ. (5)}

\prose{“ทุกฺขํ ภิกฺขเว เวทิตพฺพํ, ทุกฺขสฺส นิทานสมฺภโว เวทิตพฺโพ, ทุกฺขสฺส เวมตฺตตา เวทิตพฺพา, ทุกฺขสฺส วิปาโก เวทิตพฺโพ, ทุกฺขนิโรโธ เวทิตพฺโพ, ทุกฺขนิโรธคามินี ปฏิปทา เวทิตพฺพา”ติ อิติ โข ปเนตํ}

\vspace{\tipitakaparskip}
\csromancenter{ฉกฺกนิปาตปาฬิ วุตฺตํ, กิญฺเจตํ ปฏิจฺจ วุตฺตํ.\csromanspacer{} ชาติปิ ทุกฺขา, ชราปิ ทุกฺขา, พฺยาธิปิ ทุกฺโข,\footnote{พฺยาธิปิ ทุกฺขา (สฺยา, อิ, ก)} มรณมฺปิ ทุกฺขํ, โสกปริเทวทุกฺขโทมนสฺสุปายาสาปิ ทุกฺขา, ยมฺปิจฺฉํ น ลภติ ตมฺปิ ทุกฺขํ, สํขิตฺเตน ปญฺจุปาทานกฺขนฺธา\footnote{ปญฺจุปาทานกฺขนฺธาปิ (ก)} ทุกฺขา.}
\prose{\csromanfolio{364}กตโม จ ภิกฺขเว ทุกฺขสฺส นิทานสมฺภโว? ตณฺหา ภิกฺขเว ทุกฺขสฺส นิทานสมฺภโว.}

\prose{กตมา จ ภิกฺขเว ทุกฺขสฺส เวมตฺตตา? อตฺถิ ภิกฺขเว ทุกฺขํ อธิมตฺตํ, อตฺถิ ปริตฺตํ, อตฺถิ ทนฺธวิราคิ, อตฺถิ ขิปฺปวิราคิ.\csromanspacer{} อยํ วุจฺจติ ภิกฺขเว ทุกฺขสฺส เวมตฺตตา.}

\prose{กตโม จ ภิกฺขเว ทุกฺขสฺส วิปาโก? อิธ ภิกฺขเว เอกจฺโจ เยน ทุกฺเขน อภิภูโต ปริยาทินฺนจิตฺโต\footnote{ปริยาทิณฺณจิตฺโต (ก)} โสจติ กิลมติ ปริเทวติ อุรตฺตาฬิํ กนฺทติ สมฺโมหํ อาปชฺชติ.\csromanspacer{} เยน วา ปน ทุกฺเขน อภิภูโต ปริยาทินฺนจิตฺโต พหิทฺธา ปริเยฏฺฐิํ อาปชฺชติ “โก\footnote{โส น (ก)} เอกปทํ ทฺวิปทํ ชานาติ\footnote{ปชานาติ (ก)} อิมสฺส ทุกฺขสฺส นิโรธายา”ติ.\csromanspacer{} สมฺโมหเวปกฺกํ วาหํ ภิกฺขเว ทุกฺขํ วทามิ ปริเยฏฺฐิเวปกฺกํ วา.\csromanspacer{} อยํ วุจฺจติ ภิกฺขเว ทุกฺขสฺส วิปาโก.}

\prose{กตโม จ ภิกฺขเว ทุกฺขนิโรโธ? ตณฺหานิโรโธ\footnote{ตณฺหานิโรธา (ก-สี, สฺยา, ก)} ภิกฺขเว ทุกฺขนิโรโธ.\csromanspacer{} อยเมว อริโย อฏฺฐงฺคิโก มคฺโค ทุกฺขสฺส นิโรธคามินี ปฏิปทา.\csromanspacer{} เสยฺยถิทํ, สมฺมาทิฏฺฐิ ฯเปฯ สมฺมาสมาธิ.}

\prose{ยโต โข ภิกฺขเว อริยสาวโก เอวํ ทุกฺขํ ปชานาติ, เอวํ ทุกฺขสฺส นิทานสมฺภวํ ปชานาติ, เอวํ ทุกฺขสฺส เวมตฺตตํ ปชานาติ, เอวํ ทุกฺขสฺส วิปากํ ปชานาติ, เอวํ ทุกฺขนิโรธํ ปชานาติ, เอวํ ทุกฺขนิโรธคามินิํ ปฏิปทํ ปชานาติ.\csromanspacer{} โส อิมํ นิพฺเพธิกํ พฺรหฺมจริยํ ปชานาติ ทุกฺขนิโรธํ. “ทุกฺขํ ภิกฺขเว เวทิตพฺพํ, ทุกฺขสฺส นิทานสมฺภโว เวทิตพฺโพ, ทุกฺขสฺส เวมตฺตตา เวทิตพฺพา, ทุกฺขสฺส วิปาโก เวทิตพฺโพ, ทุกฺขนิโรโธ เวทิตพฺโพ, ทุกฺขนิโรธคามินี ปฏิปทา เวทิตพฺพา”ติ อิติ ยํ ตํ วุตฺตํ, อิทเมตํ ปฏิจฺจ วุตฺตํ. (6)}

\prose{อยํ โข โส ภิกฺขเว นิพฺเพธิกปริยาโย ธมฺมปริยาโยติ.\csromanspacer{}\csromanspacer{}นวมํ.}

\csromanmatikaanchor{mtk.186}
\csromanmatikaanchor{mtk.260}
\csromantocmarkref{subsection}{10. สีหนาทสุตฺต}{mtk.186}
\bookmark[dest={mtk.186},level=2]{10. สีหนาทสุตฺต}
\csromanheader{2}{6. มหาวคฺค \textbf{10. สีหนาทสุตฺต}}
\proseitem{64}{\csromanfolio{365}\csromansymbolfootnote{*}{ม 1. 99; อภิ 2. 328; อํ 3. 283; ขุ 9. 356 ปิฏฺเฐสุ.}ฉยิมานิ ภิกฺขเว ตถาคตสฺส ตถาคตพลานิ, เยหิ พเลหิ สมนฺนาคโต ตถาคโต อาสภํ ฐานํ ปฏิชานาติ ปริสาสุ สีหนาทํ นทติ, พฺรหฺมจกฺกํ ปวตฺเตติ.\csromanspacer{} กตมานิ ฉ? อิธ ภิกฺขเว ตถาคโต ฐานญฺจ ฐานโต อฏฺฐานญฺจ อฏฺฐานโต ยถาภูตํ ปชานาติ.\csromanspacer{} ยมฺปิ ภิกฺขเว ตถาคโต ฐานญฺจ ฐานโต อฏฺฐานญฺจ อฏฺฐานโต ยถาภูตํ ปชานาติ, อิทมฺปิ ภิกฺขเว ตถาคตสฺส ตถาคตพลํ โหติ, ยํ พลํ อาคมฺม ตถาคโต อาสภํ ฐานํ ปฏิชานาติ, ปริสาสุ สีหนาทํ นทติ, พฺรหฺมจกฺกํ ปวตฺเตติ. (1)}

\prose{ปุน จปรํ ภิกฺขเว ตถาคโต อตีตานาคตปจฺจุปฺปนฺนานํ กมฺมสมา ทานานํ ฐานโส เหตุโส วิปากํ ยถาภูตํ ปชานาติ.\csromanspacer{} ยมฺปิ ภิกฺขเว ตถาคโต อตีตานาคตปจฺจุปฺปนฺนานํ กมฺมสมาทานานํ ฐานโส เหตุโส วิปากํ ยถาภูตํ ปชานาติ, อิทมฺปิ ภิกฺขเว ตถาคตสฺส ตถาคตพลํ โหติ, ยํ พลํ อาคมฺม ตถาคโต อาสภํ ฐานํ ปฏิชานาติ, ปริสาสุ สีหนาทํ นทติ, พฺรหฺมจกฺกํ ปวตฺเตติ. (2)}

\prose{ปุน จปรํ ภิกฺขเว ตถาคโต ฌานวิโมกฺขสมาธิ- สมาปตฺตีนํ สํกิเลสํ โวทานํ วุฏฺฐานํ ยถาภูตํ ปชานาติ.\csromanspacer{} ยมฺปิ ภิกฺขเว ตถาคโต ฯเปฯ อิทมฺปิ ภิกฺขเว ตถาคตสฺส ตถาคตพลํ โหติ, ยํ พลํ อาคมฺม ตถาคโต อาสภํ ฐานํ ปฏิชานาติ, ปริสาสุ สีหนาทํ นทติ, พฺรหฺมจกฺกํ ปวตฺเตติ. (3)}

\prose{ปุน จปรํ ภิกฺขเว ตถาคโต อเนกวิหิตํ ปุพฺเพนิวาสํ อนุสฺสรติ.\csromanspacer{} เสยฺยถิทํ, เอกมฺปิ ชาติํ เทฺวปิ ชาติโย ฯเปฯ อิติ สาการํ ส-อุทฺเทสํ อเนกวิหิตํ ปุพฺเพนิวาสํ อนุสฺสรติ.\csromanspacer{} ยมฺปิ ภิกฺขเว ตถาคโต อเนกวิหิตํ ปุพฺเพนิวาสํ อนุสฺสรติ.\csromanspacer{} เสยฺยถิทํ, เอกมฺปิ ชาติํ เทฺวปิ ชาติโย ฯเปฯ อิติ สาการํ ส-อุทฺเทสํ อเนกวิหิตํ ปุพฺเพนิวาสํ อนุสฺสรติ.\csromanspacer{} อิทมฺปิ ภิกฺขเว ตถาคตสฺส ตถาคตพลํ โหติ, ยํ พลํ อาคมฺม ตถาคโต อาสภํ ฐานํ ปฏิชานาติ, ปริสาสุ สีหนาทํ นทติ, พฺรหฺมจกฺกํ ปวตฺเตติ. (4)}

\prose{ปุน จปรํ ภิกฺขเว ตถาคโต ทิพฺเพน จกฺขุนา วิสุทฺเธน อติกฺกนฺต มานุสเกน ฯเปฯ ยถากมฺมูปเค สตฺเต ปชานาติ.\csromanspacer{} ยมฺปิ ภิกฺขเว}

\vspace{\tipitakaparskip}
\csromancenter{ฉกฺกนิปาตปาฬิ ตถาคโต ทิพฺเพน จกฺขุนา วิสุทฺเธน อติกฺกนฺตมานุสเกน ฯเปฯ ยถากมฺมูปเค สตฺเต ปชานาติ.\csromanspacer{} อิทมฺปิ ภิกฺขเว ตถาคตสฺส ตถาคตพลํ โหติ, ยํ พลํ อาคมฺม ตถาคโต อาสภํ ฐานํ ปฏิชานาติ, ปริสาสุ สีหนาทํ นทติ, พฺรหฺมจกฺกํ ปวตฺเตติ. (5)}
\prose{\csromanfolio{366}ปุน จปรํ ภิกฺขเว ตถาคโต อาสวานํ ขยา ฯเปฯ สจฺฉิกตฺวา อุปสมฺปชฺช วิหรติ.\csromanspacer{} ยมฺปิ ภิกฺขเว ตถาคโต อาสวานํ ขยา ฯเปฯ สจฺฉิกตฺวา อุปสมฺปชฺช วิหรติ.\csromanspacer{} อิทมฺปิ ภิกฺขเว ตถาคตสฺส ตถาคตพลํ โหติ, ยํ พลํ อาคมฺม ตถาคโต อาสภํ ฐานํ ปฏิชานาติ, ปริสาสุ สีหนาทํ นทติ, พฺรหฺมจกฺกํ ปวตฺเตติ.\csromanspacer{} อิมานิ โข ภิกฺขเว ฉ ตถาคตสฺส ตถาคตพลานิ, เยหิ พเลหิ สมนฺนาคโต ตถาคโต อาสภํ ฐานํ ปฏิชานาติ, ปริสาสุ สีหนาทํ นทติ, พฺรหฺมจกฺกํ ปวตฺเตติ. (6)}

\prose{ตตฺร เจ ภิกฺขเว ปเร ตถาคตํ ฐานญฺจ ฐานโต อฏฺฐานญฺจ อฏฺฐานโต ยถาภูตํ ญาเณน อุปสงฺกมิตฺวา ปญฺหํ ปุจฺฉนฺติ.\csromanspacer{} ยถา ยถา ภิกฺขเว ตถาคตสฺส ฐานญฺจ ฐานโต อฏฺฐานญฺจ อฏฺฐานโต ยถาภูตํ ญาณํ วิทิตํ, ตถา ตถา เตสํ ตถาคโต ฐานญฺจ ฐานโต อฏฺฐานญฺจ อฏฺฐานโต ยถาภูตํ ญาเณน ปญฺหํ ปุฏฺโฐ พฺยากโรติ. (1)}

\prose{ตตฺร เจ ภิกฺขเว ปเร ตถาคตํ อตีตานาคตปจฺจุปฺปนฺนานํ กมฺมสมาทานานํ ฐานโส เหตุโส วิปากํ ยถาภูตํ ญาเณน อุปสงฺกมิตฺวา ปญฺหํ ปุจฺฉนฺติ.\csromanspacer{} ยถา ยถา ภิกฺขเว ตถาคตสฺส อตีตานาคตปจฺจุปฺปนฺนานํ กมฺมสมาทานานํ ฐานโส เหตุโส วิปากํ ยถาภูตํ ญาณํ วิทิตํ, ตถา ตถา เตสํ ตถาคโต อตีตานาคตปจฺจุปฺปนฺนานํ กมฺมสมาทานานํ ฐานโส เหตุโส วิปากํ ยถาภูตํ ญาเณน ปญฺหํ ปุฏฺโฐ พฺยากโรติ. (2)}

\prose{ตตฺร เจ ภิกฺขเว ปเร ตถาคตํ ฌานวิโมกฺขสมาธิสมาปตฺตีนํ สํกิเลสํ โวทานํ วุฏฺฐานํ ยถาภูตํ ญาเณน อุปสงฺกมิตฺวา ปญฺหํ ปุจฺฉนฺติ.\csromanspacer{} ยถา ยถา ภิกฺขเว ตถาคตสฺส ฌานวิโมกฺขสมาธิสมาปตฺตีนํ สํกิเลสํ โวทานํ วุฏฺฐานํ ยถาภูตํ ญาณํ วิทิตํ, ตถา ตถา เตสํ ตถาคโต ฌานวิโมกฺขสมาธิสมาปตฺตีนํ สํกิเลสํ โวทานํ วุฏฺฐานํ ยถาภูตํ ญาเณน ปญฺหํ ปุฏฺโฐ พฺยากโรติ. (3)}

\chapterheadpageread{6. มหาวคฺค}
\prose{\csromanfolio{367}ตตฺร เจ ภิกฺขเว ปเร ตถาคตํ ปุพฺเพนิวาสานุสฺสติํ ยถาภูตํ ญาเณน อุปสงฺกมิตฺวา ปญฺหํ ปุจฺฉนฺติ.\csromanspacer{} ยถา ยถา ภิกฺขเว ตถาคตสฺส ปุพฺเพนิวาสานุสฺสติํ ยถาภูตํ ญาณํ วิทิตํ, ตถา ตถา เตสํ ตถาคโต ปุพฺเพนิวาสานุสฺสติํ ยถาภูตํ ญาเณน ปญฺหํ ปุฏฺโฐ พฺยากโรติ. (4)}

\prose{ตตฺร เจ ภิกฺขเว ปเร ตถาคตํ สตฺตานํ จุตูปปาตํ ยถาภูตํ ญาเณน อุปสงฺกมิตฺวา ปญฺหํ ปุจฺฉนฺติ.\csromanspacer{} ยถา ยถา ภิกฺขเว ตถาคตสฺส สตฺตานํ จุตูปปาตํ ยถาภูตํ ญาณํ วิทิตํ, ตถา ตถา เตสํ ตถาคโต สตฺตานํ จุตูปปาตํ ยถาภูตํ ญาเณน ปญฺหํ ปุฏฺโฐ พฺยากโรติ. (5)}

\prose{ตตฺร เจ ภิกฺขเว ปเร ตถาคตํ อาสวานํ ขยา ฯเปฯ ยถาภูตํ ญาเณน อุปสงฺกมิตฺวา ปญฺหํ ปุจฺฉนฺติ.\csromanspacer{} ยถา ยถา ภิกฺขเว ตถาคตสฺส อาสวานํ ขยา ฯเปฯ ยถาภูตํ ญาณํ วิทิตํ, ตถา ตถา เตสํ ตถาคโต อาสวานํ ขยา ฯเปฯ ยถาภูตํ ญาเณน ปญฺหํ ปุฏฺโฐ พฺยากโรติ. (6)}

\prose{ตตฺร ภิกฺขเว ยมฺปิทํ\footnote{ยมิทํ (สี, อิ), ยทิทํ (ก)} ฐานญฺจ ฐานโต อฏฺฐานญฺจ อฏฺฐานโต ยถาภูตํ ญาณํ, ตมฺปิ สมาหิตสฺส วทามิ โน อสมาหิตสฺส.\csromanspacer{} ยมฺปิทํ\footnote{ยทิทํ (ก)} อตีตานาคตปจฺจุปฺปนฺนานํ กมฺมสมาทานานํ ฐานโส เหตุโส วิปากํ ยถาภูตํ ญาณํ, ตมฺปิ สมาหิตสฺส วทามิ โน อสมาหิตสฺส.\csromanspacer{} ยมฺปิทํ2 ฌานวิโมกฺขสมาธิสมาปตฺตีนํ สํกิเลสํ โวทานํ วุฏฺฐานํ ยถาภูตํ ญาณํ, ตมฺปิ สมาหิตสฺส วทามิ โน อสมาหิตสฺส.\csromanspacer{} ยมฺปิทํ2 ปุพฺเพนิวาสานุสฺสติํ ยถาภูตํ ญาณํ, ตมฺปิ สมาหิตสฺส วทามิ โน อสมาหิตสฺส.\csromanspacer{} ยมฺปิทํ2 สตฺตานํ จุตูปปาตํ ยถาภูตํ ญาณํ, ตมฺปิ สมาหิตสฺส วทามิ โน อสมาหิตสฺส.\csromanspacer{} ยมฺปิทํ2 อาสวานํ ขยา ฯเปฯ ยถาภูตํ ญาณํ, ตมฺปิ สมาหิตสฺส วทามิ โน อสมาหิตสฺส.\csromanspacer{} อิติ โข ภิกฺขเว สมาธิ มคฺโค, อสมาธิ กุมฺมคฺโคติ.\csromanspacer{}\csromanspacer{}ทสมํ.\csromanspacer{} มหาวคฺโค ฉฏฺโฐ\footnote{ปฐโม (สฺยา, ก)}.}

\gambhira{ฉกฺกนิปาตปาฬิ \textbf{ตสฺสุทฺทานํ}}
\csromanmatikaanchor{mtk.187}
\csromanmatikaanchor{mtk.188}
\csromanmatikaanchor{mtk.189}
\csromantocmarknopage{section}{7. เทวตาวคฺค}{mtk.187}
\bookmark[dest={mtk.187},level=1]{7. เทวตาวคฺค}
\csromantocmarkref{subsection}{1-2. อนาคามิผลสุตฺตาทิ}{mtk.188}
\bookmark[dest={mtk.188},level=2]{1-2. อนาคามิผลสุตฺตาทิ}
\csromantocmarkref{subsection}{3-8. มิตฺตสุตฺตาทิ}{mtk.189}
\bookmark[dest={mtk.189},level=2]{3-8. มิตฺตสุตฺตาทิ}
\csromangathameasure{โสโณ ผคฺคุโน ภิชาติ, อาสวา ทารุหตฺถิ จ. \\ มชฺเช ญาณํ นิพฺเพธิกํ, สีหนาโทติ เต ทสาติ.}
\csromangathagroup{\csromangathastanza{\csromanfolio{368}โสโณ ผคฺคุโน ภิชาติ, อาสวา ทารุหตฺถิ จ. \\ มชฺเช ญาณํ นิพฺเพธิกํ, สีหนาโทติ เต ทสาติ.}}

\chapterhead{7. เทวตาวคฺค}
\csromanheader{2}{1. อนาคามิผลสุตฺต}
\proseitem{65}{ฉ ภิกฺขเว ธมฺเม อปฺปหาย อภพฺโพ อนาคามิผลํ สจฺฉิกาตุํ.\csromanspacer{} กตเม ฉ? อสฺสทฺธิยํ อหิริกํ อโนตฺตปฺปํ โกสชฺชํ มุฏฺฐสฺสจฺจํ ทุปฺปญฺญตํ.\csromanspacer{} อิเม โข ภิกฺขเว ฉ ธมฺเม อปฺปหาย อภพฺโพ อนาคามิผลํ สจฺฉิกาตุํ.}

\prose{ฉ ภิกฺขเว ธมฺเม ปหาย ภพฺโพ อนาคามิผลํ สจฺฉิกาตุํ.\csromanspacer{} กตเม ฉ? อสฺสทฺธิยํ อหิริกํ อโนตฺตปฺปํ โกสชฺชํ มุฏฺฐสฺสจฺจํ ทุปฺปญฺญตํ.\csromanspacer{} อิเม โข ภิกฺขเว ฉ ธมฺเม ปหาย ภพฺโพ อนาคามิผลํ สจฺฉิกาตุนฺติ.\csromanspacer{} ปฐมํ. \textbf{2.}\csromanspacer{}\textbf{ อรหตฺตสุตฺต}}

\proseitem{66}{ฉ ภิกฺขเว ธมฺเม อปฺปหาย อภพฺโพ อรหตฺตํ สจฺฉิกาตุํ.\csromanspacer{} กตเม ฉ? ถินํ\footnote{ถีนํ (สี, สฺยา, กํ, อิ)} มิทฺธํ อุทฺธจฺจํ กุกฺกุจฺจํ อสฺสทฺธิยํ ปมาทํ.\csromanspacer{} อิเม โข ภิกฺขเว ฉ ธมฺเม อปฺปหาย อภพฺโพ อรหตฺตํ สจฺฉิกาตุํ.}

\prose{ฉ ภิกฺขเว ธมฺเม ปหาย ภพฺโพ อรหตฺตํ สจฺฉิกาตุํ.\csromanspacer{} กตเม ฉ? ถินํ มิทฺธํ อุทฺธจฺจํ กุกฺกุจฺจํ อสฺสทฺธิยํ ปมาทํ.\csromanspacer{} อิเม โข ภิกฺขเว ฉ ธมฺเม ปหาย ภพฺโพ อรหตฺตํ สจฺฉิกาตุนฺติ.\csromanspacer{}\csromanspacer{}ทุติยํ.}

\csromanheader{2}{3. มิตฺตสุตฺต}
\proseitem{67}{โส วต ภิกฺขเว ภิกฺขุ ปาปมิตฺโต ปาปสหาโย ปาปสมฺปวงฺโก ปาปมิตฺเต\footnote{ปาปเก มิตฺเต (ก)} เสวมาโน ภชมาโน ปยิรุปาสมาโน เตสญฺจ ทิฏฺฐานุคติํ อาปชฺชมาโน อาภิสมาจาริกํ ธมฺมํ ปริปูเรสฺสตีติ เนตํ ฐานํ วิชฺชติ, อาภิสมาจาริกํ ธมฺมํ อปริปูเรตฺวา เสขํ}

\vspace{\tipitakaparskip}
\csromancenter{เทวตาวคฺค ธมฺมํ ปริปูเรสฺสตีติ เนตํ ฐานํ วิชฺชติ, เสขํ ธมฺมํ อปริปูเรตฺวา สีลานิ ปริปูเรสฺสตีติ เนตํ ฐานํ วิชฺชติ, สีลานิ อปริปูเรตฺวา กามราคํ วา รูปราคํ วา อรูปราคํ วา ปชหิสฺสตีติ เนตํ ฐานํ วิชฺชติ.}
\prose{\csromanfolio{369}โส วต ภิกฺขเว ภิกฺขุ กลฺยาณมิตฺโต กลฺยาณสหาโย กลฺยาณสมฺปวงฺโก กลฺยาณมิตฺเต เสวมาโน ภชมาโน ปยิรุปาสมาโน เตสญฺจ ทิฏฺฐานุคติํ อาปชฺชมาโน อาภิสมาจาริกํ ธมฺมํ ปริปูเรสฺสตีติ ฐานเมตํ วิชฺชติ, อาภิสมาจาริกํ ธมฺมํ ปริปูเรตฺวา เสขํ ธมฺมํ ปริปูเรสฺสตีติ ฐานเมตํ วิชฺชติ, เสขํ ธมฺมํ ปริปูเรตฺวา สีลานิ ปริปูเรสฺสตีติ ฐานเมตํ วิชฺชติ, สีลานิ ปริปูเรตฺวา กามราคํ วา รูปราคํ วา อรูปราคํ วา ปชหิสฺสตีติ ฐานเมตํ วิชฺชตีติ.\csromanspacer{}\csromanspacer{}ตติยํ.}

\csromanheader{2}{4. สงฺคณิการามสุตฺต}
\proseitem{68}{โส วต ภิกฺขเว ภิกฺขุ สงฺคณิการาโม สงฺคณิกรโต สงฺคณิการามตํ อนุยุตฺโต คณาราโม คณรโต คณารามตํ อนุยุตฺโต เอโก ปวิเวเก อภิรมิสฺสตีติ เนตํ ฐานํ วิชฺชติ, เอโก ปวิเวเก อนภิรมนฺโต จิตฺตสฺส นิมิตฺตํ คเหสฺสตีติ เนตํ ฐานํ วิชฺชติ, จิตฺตสฺส นิมิตฺตํ อคณฺหนฺโต สมฺมาทิฏฺฐิํ ปริปูเรสฺสตีติ เนตํ ฐานํ วิชฺชติ, สมฺมาทิฏฺฐิํ อปริปูเรตฺวา สมฺมาสมาธิํ ปริปูเรสฺสตีติ เนตํ ฐานํ วิชฺชติ, สมฺมาสมาธิํ อปริปูเรตฺวา สํโยชนานิ ปชหิสฺสตีติ เนตํ ฐานํ วิชฺชติ, สํโยชนานิ อปฺปหาย นิพฺพานํ สจฺฉิกริสฺสตีติ เนตํ ฐานํ วิชฺชติ.}

\prose{โส วต ภิกฺขเว ภิกฺขุ น สงฺคณิการาโม น สงฺคณิกรโต น สงฺคณิการามตํ อนุยุตฺโต น คณาราโม น คณรโต น คณารามตํ อนุยุตฺโต เอโก ปวิเวเก อภิรมิสฺสตีติ ฐานเมตํ วิชฺชติ, เอโก ปวิเวเก อภิรมนฺโต จิตฺตสฺส นิมิตฺตํ คเหสฺสตีติ ฐานเมตํ วิชฺชติ, จิตฺตสฺส นิมิตฺตํ คณฺหนฺโต สมฺมาทิฏฺฐิํ ปริปูเรสฺสตีติ ฐานเมตํ วิชฺชติ, สมฺมาทิฏฺฐิํ ปริปูเรตฺวา สมฺมาสมาธิํ ปริปูเรสฺสตีติ ฐานเมตํ วิชฺชติ, สมฺมาสมาธิํ ปริปูเรตฺวา สํโยชนานิ ปชหิสฺสตีติ ฐานเมตํ วิชฺชติ, สํโยชนานิ ปหาย นิพฺพานํ สจฺฉิกริสฺสตีติ ฐานเมตํ วิชฺชตีติ.\csromanspacer{}\csromanspacer{}จตุตฺถํ.}

\gambhira{ฉกฺกนิปาตปาฬิ \textbf{5. เทวตาสุตฺต}}
\proseitem{69}{\csromanfolio{370}อถ โข อญฺญตรา เทวตา อภิกฺกนฺตาย รตฺติยา อภิกฺกนฺตวณฺณา เกวลกปฺปํ เชตวนํ โอภาเสตฺวา เยน ภควา เตนุปสงฺกมิ, อุปสงฺกมิตฺวา ภควนฺตํ อภิวาเทตฺวา เอกมนฺตํ อฏฺฐาสิ, เอกมนฺตํ ฐิตา โข สา เทวตา ภควนฺตํ เอตทโวจ “ฉยิเม ภนฺเต ธมฺมา ภิกฺขุโน อปริหานาย สํวตฺตนฺติ.\csromanspacer{} กตเม ฉ? สตฺถุคารวตา ธมฺมคารวตา สํฆคารวตา สิกฺขาคารวตา โสวจสฺสตา กลฺยาณมิตฺตตา.\csromanspacer{} อิเม โข ภนฺเต ฉ ธมฺมา ภิกฺขุโน อปริหานาย สํวตฺตนฺตี”ติ.\csromanspacer{} อิทมโวจ สา เทวตา, สมนุญฺโญ สตฺถา อโหสิ.\csromanspacer{} อถ โข สา เทวตา “สมนุญฺโญ เม สตฺถา”ติ ภควนฺตํ อภิวาเทตฺวา ปทกฺขิณํ กตฺวา ตตฺเถวนฺตรธายิ.}

\prose{อถ โข ภควา ตสฺสา รตฺติยา อจฺจเยน ภิกฺขู อามนฺเตสิ\csromandash{}อิมํ ภิกฺขเว รตฺติํ อญฺญตรา เทวตา อภิกฺกนฺตาย รตฺติยา อภิกฺกนฺตวณฺณา เกวลกปฺปํ เชตวนํ โอภาเสตฺวา เยนาหํ เตนุปสงฺกมิ, อุปสงฺกมิตฺวา มํ อภิวาเทตฺวา เอกมนฺตํ อฏฺฐาสิ, เอกมนฺตํ ฐิตา โข ภิกฺขเว สา เทวตา มํ เอตทโวจ “ฉยิเม ภนฺเต ธมฺมา ภิกฺขุโน อปริหานาย สํวตฺตนฺติ.\csromanspacer{} กตเม ฉ? สตฺถุคารวตา ธมฺมคารวตา สํฆคารวตา สิกฺขาคารวตา โสวจสฺสตา กลฺยาณมิตฺตตา.\csromanspacer{} อิเม โข ภนฺเต ฉ ธมฺมา ภิกฺขุโน อปริหานาย สํวตฺตนฺตี”ติ.\csromanspacer{} อิทมโวจ ภิกฺขเว สา เทวตา, อิทํ วตฺวา มํ อภิวาเทตฺวา ปทกฺขิณํ กตฺวา ตตฺเถวนฺตรธายีติ.}

\prose{เอวํ วุตฺเต อายสฺมา สาริปุตฺโต ภควนฺตํ อภิวาเทตฺวา เอตทโวจ “อิมสฺส โข อหํ ภนฺเต ภควตา สํขิตฺเตน ภาสิตสฺส เอวํ วิตฺถาเรน อตฺถํ อาชานามิ.\csromanspacer{} อิธ ภนฺเต ภิกฺขุ อตฺตนา จ สตฺถุคารโว โหติ สตฺถุคารวตาย จ วณฺณวาที.\csromanspacer{} เย จญฺเญ ภิกฺขู น สตฺถุคารวา, เต จ สตฺถุคารวตาย สมาทเปติ.\csromanspacer{} เย จญฺเญ ภิกฺขู สตฺถุคารวา, เตสญฺจ วณฺณํ ภณติ ภูตํ ตจฺฉํ กาเลน.\csromanspacer{} อตฺตนา จ ธมฺมคารโว โหติ.\csromanspacer{} สํฆคารโว โหติ.\csromanspacer{} สิกฺขาคารโว โหติ.\csromanspacer{} สุวโจ โหติ.\csromanspacer{} กลฺยาณมิตฺโต โหติ กลฺยาณมิตฺตตาย จ วณฺณวาที.\csromanspacer{} เย จญฺเญ ภิกฺขู น กลฺยาณมิตฺตา, เต จ กลฺยาณมิตฺตตาย สมาทเปติ.\csromanspacer{} เย จญฺเญ ภิกฺขู กลฺยาณมิตฺตา, เตสญฺจ วณฺณํ ภณติ ภูตํ ตจฺฉํ}

\vspace{\tipitakaparskip}
\csromancenter{เทวตาวคฺค กาเลน.\csromanspacer{} อิมสฺส โข อหํ ภนฺเต ภควตา สํขิตฺเตน ภาสิตสฺส เอวํ วิตฺถาเรน อตฺถํ อาชานามี”ติ.}
\prose{\csromanfolio{371}สาธุ สาธุ สาริปุตฺต, สาธุ โข ตฺวํ สาริปุตฺต อิมสฺส มยา สํขิตฺเตน ภาสิตสฺส เอวํ วิตฺถาเรน อตฺถํ อาชานาสิ.\csromanspacer{} อิธ สาริปุตฺต ภิกฺขุ อตฺตนา จ สตฺถุคารโว โหติ, สตฺถุคารวตาย จ วณฺณวาที.\csromanspacer{} เย จญฺเญ ภิกฺขู น สตฺถุคารวา, เต จ สตฺถุคารวตาย สมาทเปติ.\csromanspacer{} เย จญฺเญ ภิกฺขู สตฺถุคารวา, เตสญฺจ วณฺณํ ภณติ ภูตํ ตจฺฉํ กาเลน.\csromanspacer{} อตฺตนา จ ธมฺมคารโว โหติ.\csromanspacer{} สํฆคารโว โหติ.\csromanspacer{} สิกฺขาคารโว โหติ.\csromanspacer{} สุวโจ โหติ.\csromanspacer{} กลฺยาณมิตฺโต โหติ กลฺยาณมิตฺตตาย จ วณฺณวาที.\csromanspacer{} เย จญฺเญ ภิกฺขู น กลฺยาณมิตฺตา, เต จ กลฺยาณมิตฺตตาย สมาทเปติ.\csromanspacer{} เย จญฺเญ ภิกฺขู กลฺยาณมิตฺตา, เตสญฺจ วณฺณํ ภณติ ภูตํ ตจฺฉํ กาเลน.\csromanspacer{} อิมสฺส โข สาริปุตฺต มยา สํขิตฺเตน ภาสิตสฺส เอวํ วิตฺถาเรน อตฺโถ ทฏฺฐพฺโพติ.\csromanspacer{}\csromanspacer{}ปญฺจมํ.}

\csromanheader{2}{6. สมาธิสุตฺต}
\proseitem{70}{โส วต ภิกฺขเว ภิกฺขุ น สนฺเตน สมาธินา น ปณีเตน น ปฏิปฺปสฺสทฺธิลทฺเธน\footnote{น ปฏิปฺปสฺสทฺธลทฺเธน (สี)} น เอโกทิภาวาธิคเตน อเนกวิหิตํ อิทฺธิวิธํ ปจฺจนุภวิสฺสติ.\csromanspacer{} เอโกปิ หุตฺวา พหุธา ภวิสฺสติ, พหุธาปิ หุตฺวา เอโก ภวิสฺสติ ฯเปฯ ยาว พฺรหฺมโลกาปิ กาเยน วสํ วตฺเตสฺสตีติ เนตํ ฐานํ วิชฺชติ.\csromanspacer{} ทิพฺพาย โสตธาตุยา วิสุทฺธาย อติกฺกนฺตมานุสิกาย อุโภ สทฺเท สุณิสฺสติ ทิพฺเพ จ มานุเส จ เย ทูเร สนฺติเก จาติ เนตํ ฐานํ วิชฺชติ.\csromanspacer{} ปรสตฺตานํ ปรปุคฺคลานํ เจตสา เจโต ปริจฺจ ปชานิสฺสติ.\csromanspacer{} สราคํ วา จิตฺตํ “สราคํ จิตฺตนฺ”ติ ปชานิสฺสติ ฯเปฯ วิมุตฺตํ วา จิตฺตํ “วิมุตฺตํ จิตฺตนฺ”ติ ปชานิสฺสตีติ เนตํ ฐานํ วิชฺชติ.\csromanspacer{} อเนกวิหิตํ ปุพฺเพนิวาสํ อนุสฺสริสฺสติ.\csromanspacer{} เสยฺยถิทํ, เอกมฺปิ ชาติํ เทฺวปิ ชาติโย ฯเปฯ อิติ สาการํ ส-อุทฺเทสํ อเนกวิหิตํ ปุพฺเพนิวาสํ อนุสฺสริสฺสตีติ เนตํ ฐานํ วิชฺชติ.\csromanspacer{} ทิพฺเพน จกฺขุนา วิสุทฺเธน อติกฺกนฺตมานุสเกน สตฺเต ปสฺสิสฺสติ ฯเปฯ ยถากมฺมูปเค สตฺเต ปชานิสฺสตีติ เนตํ ฐานํ วิชฺชติ.\csromanspacer{} อาสวานํ ขยา ฯเปฯ สจฺฉิกตฺวา อุปสมฺปชฺช วิหริสฺสตีติ เนตํ ฐานํ วิชฺชติ.}

\gambhira{ฉกฺกนิปาตปาฬิ}
\prose{\csromanfolio{372}โส วต ภิกฺขเว ภิกฺขุ สนฺเตน สมาธินา ปณีเตน ปฏิปฺปสฺสทฺธิ ลทฺเธน เอโกทิภาวาธิคเตน อเนกวิหิตํ อิทฺธิวิธํ ปจฺจนุภวิสฺสติ ฯเปฯ ยาว พฺรหฺมโลกาปิ กาเยน วสํ วตฺเตสฺสตีติ ฐานเมตํ วิชฺชติ.\csromanspacer{} ทิพฺพาย โสตธาตุยา วิสุทฺธาย อติกฺกนฺตมานุสิกาย อุโภ สทฺเท สุณิสฺสติ ทิพฺเพ จ มานุเส จ เย ทูเร สนฺติเก จาติ ฐานเมตํ วิชฺชติ.\csromanspacer{} ปรสตฺตานํ ปรปุคฺคลานํ เจตสา เจโต ปริจฺจ ปชานิสฺสติ.\csromanspacer{} สราคํ วา จิตฺตํ “สราคํ จิตฺตนฺ”ติ ปชานิสฺสติ ฯเปฯ วิมุตฺตํ วา จิตฺตํ “วิมุตฺตํ จิตฺตนฺ”ติ ปชานิสฺสตีติ ฐานเมตํ วิชฺชติ.\csromanspacer{} อเนกวิหิตํ ปุพฺเพนิวาสํ อนุสฺสริสฺสติ.\csromanspacer{} เสยฺยถิทํ, เอกมฺปิ ชาติํ เทฺวปิ ชาติโย ฯเปฯ อิติ สาการํ ส-อุทฺเทสํ อเนกวิหิตํ ปุพฺเพนิวาสํ อนุสฺสริสฺสตีติ ฐานเมตํ วิชฺชติ.\csromanspacer{} ทิพฺเพน จกฺขุนา วิสุทฺเธน อติกฺกนฺตมานุสเกน สตฺเต ปสฺสิสฺสติ จวมาเน อุปปชฺชมาเน หีเน ปณีเต สุวณฺเณ ทุพฺพณฺเณ สุคเต ทุคฺคเต, ยถากมฺมูปเค สตฺเต ปชานิสฺสตีติ ฐานเมตํ วิชฺชติ.\csromanspacer{} อาสวานํ ขยา อนาสวํ เจโตวิมุตฺติํ ฯเปฯ สจฺฉิกตฺวา อุปสมฺปชฺช วิหริสฺสตีติ ฐานเมตํ วิชฺชตีติ.\csromanspacer{}\csromanspacer{}ฉฏฺฐํ.}

\csromanheader{2}{7. สกฺขิภพฺพสุตฺต}
\proseitem{71}{ฉหิ ภิกฺขเว ธมฺเมหิ สมนฺนาคโต ภิกฺขุ อภพฺโพ ตตฺร ตเตฺรว สกฺขิภพฺพตํ ปาปุณิตุํ สติ สติ อายตเน.\csromanspacer{} กตเมหิ ฉหิ? อิธ ภิกฺขเว ภิกฺขุ “อิเม หานภาคิยา ธมฺมา”ติ ยถาภูตํ นปฺปชานาติ. “อิเม ฐิติภาคิยา ธมฺมา”ติ ยถาภูตํ นปฺปชานาติ. “อิเม วิเสสภาคิยา ธมฺมา”ติ ยถาภูตํ นปฺปชานาติ. “อิเม นิพฺเพธภาคิยา ธมฺมา”ติ ยถาภูตํ นปฺปชานาติ.\csromanspacer{} อสกฺกจฺจการี จ โหติ อสปฺปายการี จ.\csromanspacer{} อิเมหิ โข ภิกฺขเว ฉหิ ธมฺเมหิ สมนฺนาคโต ภิกฺขุ อภพฺโพ ตตฺร ตเตฺรว สกฺขิภพฺพตํ ปาปุณิตุํ สติ สติ อายตเน.}

\prose{ฉหิ ภิกฺขเว ธมฺเมหิ สมนฺนาคโต ภิกฺขุ ภพฺโพ ตตฺร ตเตฺรว สกฺขิภพฺพตํ ปาปุณิตุํ สติ สติ อายตเน.\csromanspacer{} กตเมหิ ฉหิ? อิธ ภิกฺขเว ภิกฺขุ “อิเม หานภาคิยา ธมฺมา”ติ ยถาภูตํ ปชานาติ. “อิเม ฐิติภาคิยา ธมฺมา”ติ ยถาภูตํ ปชานาติ. “อิเม วิเสสภาคิยา ธมฺมา”ติ ยถาภูตํ ปชานาติ. “อิเม นิพฺเพธภาคิยา ธมฺมา”ติ ยถาภูตํ ปชานาติ.\csromanspacer{} สกฺกจฺจการี จ โหติ สปฺปายการี จ.}

\vspace{\tipitakaparskip}
\csromancenter{เทวตาวคฺค อิเมหิ โข ภิกฺขเว ฉหิ ธมฺเมหิ สมนฺนาคโต ภิกฺขุ ภพฺโพ ตตฺร ตเตฺรว สกฺขิภพฺพตํ ปาปุณิตุํ สติ สติ อายตเนติ.\csromanspacer{}\csromanspacer{}สตฺตมํ.}
\csromanheader{2}{8. พลสุตฺต}
\csromanmatikaanchor{mtk.190}
\csromantocmarkref{subsection}{9-10. ตชฺฌานสุตฺตาทิ}{mtk.190}
\bookmark[dest={mtk.190},level=2]{9-10. ตชฺฌานสุตฺตาทิ}
\proseitem{72}{\csromanfolio{373}ฉหิ ภิกฺขเว ธมฺเมหิ สมนฺนาคโต ภิกฺขุ อภพฺโพ สมาธิสฺมิํ\footnote{สมาธิมฺหิ (ก)} พลตํ ปาปุณิตุํ.\csromanspacer{} กตเมหิ ฉหิ? อิธ ภิกฺขเว ภิกฺขุ น สมาธิสฺส สมาปตฺติกุสโล โหติ, น สมาธิสฺส ฐิติกุสโล โหติ, น สมาธิสฺส\footnote{น สมาธิมฺหา (ก) อุปริสตฺตกนิปาเต เทวตาวคฺเค ปน สพฺพตฺถปิ “สมาธิสฺส” อิเตฺวว ทิสฺสติ.} วุฏฺฐานกุสโล โหติ, อสกฺกจฺจการี จ โหติ อสาตจฺจการี จ อสปฺปายการี จ.\csromanspacer{} อิเมหิ โข ภิกฺขเว ฉหิ ธมฺเมหิ สมนฺนาคโต ภิกฺขุ อภพฺโพ สมาธิสฺมิํ พลตํ ปาปุณิตุํ.}

\prose{ฉหิ ภิกฺขเว ธมฺเมหิ สมนฺนาคโต ภิกฺขุ ภพฺโพ สมาธิสฺมิํ พลตํ ปาปุณิตุํ.\csromanspacer{} กตเมหิ ฉหิ? อิธ ภิกฺขเว ภิกฺขุ สมาธิสฺส สมาปตฺติกุสโล โหติ, สมาธิสฺส ฐิติกุสโล โหติ, สมาธิสฺส วุฏฺฐานกุสโล โหติ, สกฺกจฺจการี จ โหติ สาตจฺจการี จ สปฺปายการี จ.\csromanspacer{} อิเมหิ โข ภิกฺขเว ฉหิ ธมฺเมหิ สมนฺนาคโต ภิกฺขุ ภพฺโพ สมาธิสฺมิํ พลตํ ปาปุณิตุนฺติ.\csromanspacer{}\csromanspacer{}อฏฺฐมํ.}

\csromanheader{2}{9. ปฐมตชฺฌานสุตฺต}
\proseitem{73}{ฉ ภิกฺขเว ธมฺเม อปฺปหาย อภพฺโพ ปฐมํ ฌานํ อุปสมฺปชฺช วิหริตุํ.\csromanspacer{} กตเม ฉ? กามจฺฉนฺทํ พฺยาปาทํ ถินมิทฺธํ อุทฺธจฺจกุกฺกุจฺจํ วิจิกิจฺฉํ, กาเมสุ โข ปนสฺส อาทีนโว น ยถาภูตํ สมฺมปฺปญฺญาย สุทิฏฺโฐ โหติ.\csromanspacer{} อิเม โข ภิกฺขเว ฉ ธมฺเม อปฺปหาย อภพฺโพ ปฐมํ ฌานํ อุปสมฺปชฺช วิหริตุํ.}

\prose{ฉ ภิกฺขเว ธมฺเม ปหาย ภพฺโพ ปฐมํ ฌานํ อุปสมฺปชฺช วิหริตุํ.\csromanspacer{} กตเม ฉ? กามจฺฉนฺทํ พฺยาปาทํ ถินมิทฺธํ อุทฺธจฺจกุกฺกุจฺจํ วิจิกิจฺฉํ, กาเมสุ โข ปนสฺส อาทีนโว น ยถาภูตํ สมฺมปฺปญฺญาย สุทิฏฺโฐ โหติ.\csromanspacer{} อิเม โข ภิกฺขเว ฉ ธมฺเม ปหาย ภพฺโพ ปฐมํ ฌานํ อุปสมฺปชฺช วิหริตุนฺติ.\csromanspacer{}\csromanspacer{}นวมํ.}

\gambhira{ฉกฺกนิปาตปาฬิ \textbf{10. ทุติยตชฺฌานสุตฺต}}
\csromanmatikaanchor{mtk.191}
\csromanmatikaanchor{mtk.192}
\csromantocmarknopage{section}{8. อรหตฺตวคฺค}{mtk.191}
\bookmark[dest={mtk.191},level=1]{8. อรหตฺตวคฺค}
\csromantocmarkref{subsection}{1-9. ทุกฺขสุตฺตาทิ}{mtk.192}
\bookmark[dest={mtk.192},level=2]{1-9. ทุกฺขสุตฺตาทิ}
\proseitem{74}{\csromanfolio{374}ฉ ภิกฺขเว ธมฺเม อปฺปหาย อภพฺโพ ปฐมํ ฌานํ อุปสมฺปชฺช วิหริตุํ.\csromanspacer{} กตเม ฉ? กามวิตกฺกํ พฺยาปาทวิตกฺกํ วิหิํสาวิตกฺกํ กามสญฺญํ พฺยาปาทสญฺญํ วิหิํสาสญฺญํ.\csromanspacer{} อิเม โข ภิกฺขเว ฉ ธมฺเม อปฺปหาย อภพฺโพ ปฐมํ ฌานํ อุปสมฺปชฺช วิหริตุํ.}

\prose{ฉ ภิกฺขเว ธมฺเม ปหาย ภพฺโพ ปฐมํ ฌานํ อุปสมฺปชฺช วิหริตุํ.\csromanspacer{} กตเม ฉ? กามวิตกฺกํ พฺยาปาทวิตกฺกํ วิหิํสาวิตกฺกํ กามสญฺญํ พฺยาปาทสญฺญํ วิหิํสาสญฺญํ.\csromanspacer{} อิเม โข ภิกฺขเว ฉ ธมฺเม ปหาย ภพฺโพ ปฐมํ ฌานํ อุปสมฺปชฺช วิหริตุนฺติ.\csromanspacer{}\csromanspacer{}ทสมํ.\csromanspacer{} เทวตาวคฺโค สตฺตโม\footnote{ทุติโย (สฺยา, ก)}.}

\titlehead{\textbf{ตสฺสุทฺทานํ}}
\csromangathameasure{อนาคามิ อรหํ มิตฺตา, สงฺคณิการามเทวตา. \\ สมาธิ สกฺขิภพฺพํ พลํ, ตชฺฌานา อปเร ทุเวติ.}
\csromangathagroup{\csromangathastanza{อนาคามิ อรหํ มิตฺตา, สงฺคณิการามเทวตา. \\ สมาธิ สกฺขิภพฺพํ พลํ, ตชฺฌานา อปเร ทุเวติ.}}

\chapterhead{8. อรหตฺตวคฺค}
\csromanheader{2}{1. ทุกฺขสุตฺต}
\proseitem{75}{ฉหิ ภิกฺขเว ธมฺเมหิ สมนฺนาคโต ภิกฺขุ ทิฏฺเฐว ธมฺเม ทุกฺขํ วิหรติ สวิฆาตํ ส-อุปายาสํ สปริฬาหํ, กายสฺส เภทา ปรํ มรณา ทุคฺคติ ปาฏิกงฺขา.\csromanspacer{} กตเมหิ ฉหิ? กามวิตกฺเกน พฺยาปาทวิตกฺเกน วิหิํสาวิตกฺเกน กามสญฺญาย พฺยาปาทสญฺญาย วิหิํสาสญฺญาย.\csromanspacer{} อิเมหิ โข ภิกฺขเว ฉหิ ธมฺเมหิ สมนฺนาคโต ภิกฺขุ ทิฏฺเฐว ธมฺเม ทุกฺขํ วิหรติ สวิฆาตํ ส-อุปายาสํ สปริฬาหํ.\csromanspacer{} กายสฺส เภทา ปรํ มรณา ทุคฺคติ ปาฏิกงฺขา.}

\prose{ฉหิ ภิกฺขเว ธมฺเมหิ สมนฺนาคโต ภิกฺขุ ทิฏฺเฐว ธมฺเม สุขํ วิหรติ อวิฆาตํ อนุปายาสํ อปริฬาหํ, กายสฺส เภทา}

\vspace{\tipitakaparskip}
\csromancenter{อรหตฺตวคฺค ปรํ มรณา สุคติ ปาฏิกงฺขา.\csromanspacer{} กตเมหิ ฉหิ? เนกฺขมฺมวิตกฺเกน อพฺยาปาทวิตกฺเกน อวิหิํสาวิตกฺเกน เนกฺขมฺมสญฺญาย อพฺยาปาทสญฺญาย อวิหิํสาสญฺญาย.\csromanspacer{} อิเมหิ โข ภิกฺขเว ฉหิ ธมฺเมหิ สมนฺนาคโต ภิกฺขุ ทิฏฺเฐว ธมฺเม สุขํ วิหรติ อวิฆาตํ อนุปายาสํ อปริฬาหํ.\csromanspacer{} กายสฺส เภทา ปรํ มรณา สุคติ ปาฏิกงฺขาติ.\csromanspacer{}\csromanspacer{}ปฐมํ.}
\csromanheader{2}{2. อรหตฺตสุตฺต}
\proseitem{76}{\csromanfolio{375}ฉ ภิกฺขเว ธมฺเม อปฺปหาย อภพฺโพ อรหตฺตํ สจฺฉิกาตุํ.\csromanspacer{} กตเม ฉ? มานํ โอมานํ อติมานํ อธิมานํ ถมฺภํ อตินิปาตํ.\csromanspacer{} อิเม โข ภิกฺขเว ฉ ธมฺเม อปฺปหาย อภพฺโพ อรหตฺตํ สจฺฉิกาตุํ.}

\prose{ฉ ภิกฺขเว ธมฺเม ปหาย ภพฺโพ อรหตฺตํ สจฺฉิกาตุํ.\csromanspacer{} กตเม ฉ? มานํ โอมานํ อติมานํ อธิมานํ ถมฺภํ อตินิปาตํ.\csromanspacer{} อิเม โข ภิกฺขเว ฉ ธมฺเม ปหาย ภพฺโพ อรหตฺตํ สจฺฉิกาตุนฺติ.\csromanspacer{}\csromanspacer{}ทุติยํ.}

\csromanheader{2}{3. อุตฺตริมนุสฺสธมฺมสุตฺต}
\proseitem{77}{ฉ ภิกฺขเว ธมฺเม อปฺปหาย อภพฺโพ อุตฺตริมนุสฺสธมฺมํ อลมริยญาณทสฺสนวิเสสํ สจฺฉิกาตุํ.\csromanspacer{} กตเม ฉ? มุฏฺฐสฺสจฺจํ อสมฺปชญฺญํ อินฺทฺริเยสุ อคุตฺตทฺวารตํ โภชเน อมตฺตญฺญุตํ กุหนํ ลปนํ.\csromanspacer{} อิเม โข ภิกฺขเว ฉ ธมฺเม อปฺปหาย อภพฺโพ อุตฺตริมนุสฺสธมฺมํ อลมริยญาณทสฺสนวิเสสํ สจฺฉิกาตุํ.}

\prose{ฉ ภิกฺขเว ธมฺเม ปหาย ภพฺโพ อุตฺตริมนุสฺสธมฺมํ อลมริยญาณทสฺสนวิเสสํ สจฺฉิกาตุํ.\csromanspacer{} กตเม ฉ? มุฏฺฐสฺสจฺจํ อสมฺปชญฺญํ อินฺทฺริเยสุ อคุตฺตทฺวารตํ โภชเน อมตฺตญฺญุตํ กุหนํ ลปนํ.\csromanspacer{} อิเม โข ภิกฺขเว ฉ ธมฺเม ปหาย ภพฺโพ อุตฺตริมนุสฺสธมฺมํ อลมริยญาณทสฺสนวิเสสํ สจฺฉิกาตุนฺติ.\csromanspacer{}\csromanspacer{}ตติยํ.}

\csromanheader{2}{4. สุขโสมนสฺสสุตฺต}
\proseitem{78}{ฉหิ ภิกฺขเว ธมฺเมหิ สมนฺนาคโต ภิกฺขุ ทิฏฺเฐว ธมฺเม สุขโสมนสฺสพหุโล วิหรติ, โยนิ จสฺส อารทฺธา โหติ อาสวานํ ขยาย.\csromanspacer{} กตเมหิ ฉหิ? อิธ ภิกฺขเว ภิกฺขุ ธมฺมาราโม โหติ, ภาวนาราโม โหติ, ปหานาราโม โหติ, ปวิเวการาโม}

\vspace{\tipitakaparskip}
\csromancenter{ฉกฺกนิปาตปาฬิ โหติ, อพฺยาปชฺชาราโม โหติ, นิปฺปปญฺจาราโม โหติ.\csromanspacer{} อิเมหิ โข ภิกฺขเว ฉหิ ธมฺเมหิ สมนฺนาคโต ภิกฺขุ ทิฏฺเฐว ธมฺเม สุขโสมนสฺสพหุโล วิหรติ, โยนิ จสฺส อารทฺธา โหติ อาสวานํ ขยายาติ.\csromanspacer{}\csromanspacer{}จตุตฺถํ.}
\csromanheader{2}{5. อธิคมสุตฺต}
\proseitem{79}{\csromanfolio{376}ฉหิ ภิกฺขเว ธมฺเมหิ สมนฺนาคโต ภิกฺขุ อภพฺโพ อนธิคตํ วา กุสลํ ธมฺมํ อธิคนฺตุํ, อธิคตํ วา กุสลํ ธมฺมํ ผาติํ กาตุํ\footnote{ผาติกตฺตุํ (สี), ผาติกาตุํ (สฺยา, กํ, อิ)}.\csromanspacer{} กตเมหิ ฉหิ? อิธ ภิกฺขเว ภิกฺขุ น อายกุสโล จ โหติ, น อปายกุสโล จ โหติ, น อุปายกุสโล จ โหติ, อนธิคตานํ กุสลานํ ธมฺมานํ อธิคมาย น ฉนฺทํ ชเนติ, อธิคเต กุสเล ธมฺเม น อารกฺขติ\footnote{สารกฺขติ (สี, สฺยา, กํ, อิ)}, สาตจฺจกิริยาย น สมฺปาเทติ.\csromanspacer{} อิเมหิ โข ภิกฺขเว ฉหิ ธมฺเมหิ สมนฺนาคโต ภิกฺขุ อภพฺโพ อนธิคตํ วา กุสลํ ธมฺมํ อธิคนฺตุํ, อธิคตํ วา กุสลํ ธมฺมํ ผาติํ กาตุํ.}

\prose{ฉหิ ภิกฺขเว ธมฺเมหิ สมนฺนาคโต ภิกฺขุ ภพฺโพ อนธิคตํ วา กุสลํ ธมฺมํ อธิคนฺตุํ, อธิคตํ วา กุสลํ ธมฺมํ ผาติํ กาตุํ.\csromanspacer{} กตเมหิ ฉหิ? อิธ ภิกฺขเว ภิกฺขุ อายกุสโล จ โหติ, อปายกุสโล จ โหติ, อุปายกุสโล จ โหติ, อนธิคตานํ กุสลานํ ธมฺมานํ อธิคมาย ฉนฺทํ ชเนติ, อธิคเต กุสเล ธมฺเม อารกฺขติ, สาตจฺจกิริยาย สมฺปาเทติ.\csromanspacer{} อิเมหิ โข ภิกฺขเว ฉหิ ธมฺเมหิ สมนฺนาคโต ภิกฺขุ ภพฺโพ อนธิคตํ วา กุสลํ ธมฺมํ อธิคนฺตุํ อธิคตํ วา กุสลํ ธมฺมํ ผาติํ กาตุนฺติ.\csromanspacer{}\csromanspacer{}ปญฺจมํ.}

\csromanheader{2}{6. มหนฺตตฺตสุตฺต}
\proseitem{80}{ฉหิ ภิกฺขเว ธมฺเมหิ สมนฺนาคโต ภิกฺขุ นจิรสฺเสว มหนฺตตฺตํ\footnote{มหตฺตํ (สฺยา, กํ)} เวปุลฺลตฺตํ ปาปุณาติ ธมฺเมสุ.\csromanspacer{} กตเมหิ ฉหิ? อิธ ภิกฺขเว ภิกฺขุ อาโลกพหุโล จ โหติ โยคพหุโล จ เวทพหุโล จ อสนฺตุฏฺฐิพหุโล จ อนิกฺขิตฺตธุโร จ กุสเลสุ ธมฺเมสุ อุตฺตริ จ ปตาเรติ\footnote{ปกโรติ (ก)}.\csromanspacer{} อิเมหิ โข ภิกฺขเว ฉหิ ธมฺเมหิ สมนฺนาคโต ภิกฺขุ นจิรสฺเสว มหนฺตตฺตํ เวปุลฺลตฺตํ ปาปุณาติ ธมฺเมสูติ.\csromanspacer{}\csromanspacer{}ฉฏฺฐํ.}

\csromanheader{1}{8. อรหตฺตวคฺค \textbf{7. ปฐมนิรยสุตฺต}}
\proseitem{81}{\csromanfolio{377}ฉหิ ภิกฺขเว ธมฺเมหิ สมนฺนาคโต ยถาภตํ นิกฺขิตฺโต เอวํ นิรเย.\csromanspacer{} กตเมหิ ฉหิ? ปาณาติปาตี โหติ, อทินฺนาทายี โหติ, กาเมสุมิจฺฉาจารี โหติ, มุสาวาที โหติ, ปาปิจฺโฉ จ มิจฺฉาทิฏฺฐิ จ.\csromanspacer{} อิเมหิ โข ภิกฺขเว ฉหิ ธมฺเมหิ สมนฺนาคโต ยถาภตํ นิกฺขิตฺโต เอวํ นิรเย.}

\prose{ฉหิ ภิกฺขเว ธมฺเมหิ สมนฺนาคโต ยถาภตํ นิกฺขิตฺโต เอวํ สคฺเค.\csromanspacer{} กตเมหิ ฉหิ? ปาณาติปาตา ปฏิวิรโต โหติ, อทินฺนาทานา ปฏิวิรโต โหติ, กาเมสุมิจฺฉาจารา ปฏิวิรโต โหติ, มุสาวาทา ปฏิวิรโต โหติ, อปฺปิจฺโฉ จ สมฺมาทิฏฺฐิ จ.\csromanspacer{} อิเมหิ โข ภิกฺขเว ฉหิ ธมฺเมหิ สมนฺนาคโต ยถาภตํ นิกฺขิตฺโต เอวํ สคฺเคติ.\csromanspacer{}\csromanspacer{}สตฺตมํ.}

\csromanheader{2}{8. ทุติยนิรยสุตฺต}
\proseitem{28}{ฉหิ ภิกฺขเว ธมฺเมหิ สมนฺนาคโต ยถาภตํ นิกฺขิตฺโต เอวํ นิรเย.\csromanspacer{} กตเมหิ ฉหิ?1 ปาณาติปาตี โหติ, อทินฺนาทายี โหติ, กาเมสุมิจฺฉาจารี โหติ, มุสาวาที โหติ1, ลุทฺโธ จ ปคพฺโภ จ.\csromanspacer{} อิเมหิ โข ภิกฺขเว ฉหิ ธมฺเมหิ สมนฺนาคโต ยถาภตํ นิกฺขิตฺโต เอวํ นิรเย.}

\prose{ฉหิ ภิกฺขเว ธมฺเมหิ สมนฺนาคโต ยถาภตํ นิกฺขิตฺโต เอวํ สคฺเค.\csromanspacer{} กตเมหิ ฉหิ? ปาณาติปาตา ปฏิวิรโต โหติ, อทินฺนาทานา ปฏิวิรโต โหติ, กาเมสุ มิจฺฉาจารา ปฏิวิรโต โหติ, มุสาวาทา ปฏิวิรโต โหติ, อลุทฺโธ จ อปฺปคพฺโภ จ.\csromanspacer{} อิเมหิ โข ภิกฺขเว ฉหิ ธมฺเมหิ สมนฺนาคโต ยถาภตํ นิกฺขิตฺโต เอวํ สคฺเคติ.\csromanspacer{}\csromanspacer{}อฏฺฐมํ.}

\csromanheader{2}{9. อคฺคธมฺมสุตฺต}
\proseitem{83}{ฉหิ ภิกฺขเว ธมฺเมหิ สมนฺนาคโต ภิกฺขุ อภพฺโพ อคฺคํ ธมฺมํ อรหตฺตํ สจฺฉิกาตุํ.\csromanspacer{} กตเมหิ ฉหิ? อิธ ภิกฺขเว ภิกฺขุ อสฺสทฺโธ โหติ, อหิริโก โหติ, อโนตฺตปฺปี โหติ, กุสีโต โหติ,}

\vspace{\tipitakaparskip}
\csromancenter{ฉกฺกนิปาตปาฬิ ทุปฺปญฺโญ โหติ, กาเย จ ชีวิเต จ สาเปกฺโข โหติ.\csromanspacer{} อิเมหิ โข ภิกฺขเว ฉหิ ธมฺเมหิ สมนฺนาคโต ภิกฺขุ อภพฺโพ อคฺคํ ธมฺมํ อรหตฺตํ สจฺฉิกาตุํ.}
\prose{\csromanfolio{378}ฉหิ ภิกฺขเว ธมฺเมหิ สมนฺนาคโต ภิกฺขุ ภพฺโพ อคฺคํ ธมฺมํ อรหตฺตํ สจฺฉิกาตุํ.\csromanspacer{} กตเมหิ ฉหิ? อิธ ภิกฺขเว ภิกฺขุ สทฺโธ โหติ, หิริมา โหติ, โอตฺตปฺปี โหติ, อารทฺธวีริโย โหติ, ปญฺญวา โหติ, กาเย จ ชีวิเต จ อนเปกฺโข โหติ.\csromanspacer{} อิเมหิ โข ภิกฺขเว ฉหิ ธมฺเมหิ สมนฺนาคโต ภิกฺขุ ภพฺโพ อคฺคํ ธมฺมํ อรหตฺตํ สจฺฉิกาตุนฺติ.\csromanspacer{}\csromanspacer{}นวมํ.}

\csromanmatikaanchor{mtk.193}
\csromantocmarkref{subsection}{10. รตฺติทิวสสุตฺต}{mtk.193}
\bookmark[dest={mtk.193},level=2]{10. รตฺติทิวสสุตฺต}
\csromanheader{2}{10. รตฺติทิวสสุตฺต}
\proseitem{84}{ฉหิ ภิกฺขเว ธมฺเมหิ สมนฺนาคตสฺส ภิกฺขุโน ยา รตฺติ วา ทิวโส วา อาคจฺฉติ, หานิเยว ปาฏิกงฺขา กุสเลสุ ธมฺเมสุ โน วุทฺธิ.\csromanspacer{} กตเมหิ ฉหิ? อิธ ภิกฺขเว ภิกฺขุ มหิจฺโฉ โหติ วิฆาตวา อสนฺตุฏฺโฐ อิตรีตรจีวรปิณฺฑปาตเสนาสนคิลานปฺปจฺจยเภสชฺช- ปริกฺขาเรน, อสฺสทฺโธ โหติ, ทุสฺสีโล โหติ, กุสีโต โหติ, มุฏฺฐสฺสติ โหติ, ทุปฺปญฺโญ โหติ.\csromanspacer{} อิเมหิ โข ภิกฺขเว ฉหิ ธมฺเมหิ สมนฺนาคตสฺส ภิกฺขุโน ยา รตฺติ วา ทิวโส วา อาคจฺฉติ, หานิเยว ปาฏิกงฺขา กุสเลสุ ธมฺเมสุ โน วุทฺธิ.}

\prose{ฉหิ ภิกฺขเว ธมฺเมหิ สมนฺนาคตสฺส ภิกฺขุโน ยา รตฺติ วา ทิวโส วา อาคจฺฉติ, วุทฺธิเยว ปาฏิกงฺขา กุสเลสุ ธมฺเมสุ โน ปริหานิ.\csromanspacer{} กตเมหิ ฉหิ? อิธ ภิกฺขเว ภิกฺขุ น มหิจฺโฉ โหติ อวิฆาตวา สนฺตุฏฺโฐ อิตรีตรจีวรปิณฺฑปาตเสนาสนคิลานปฺปจฺจย- เภสชฺชปริกฺขาเรน, สทฺโธ โหติ, สีลวา โหติ, อารทฺธวีริโย โหติ, สติมา โหติ, ปญฺญวา โหติ.\csromanspacer{} อิเมหิ โข ภิกฺขเว ฉหิ ธมฺเมหิ สมนฺนาคตสฺส ภิกฺขุโน ยา รตฺติ วา ทิวโส วา อาคจฺฉติ, วุทฺธิเยว ปาฏิกงฺขา กุสเลสุ ธมฺเมสุ, โน ปริหานีติ.\csromanspacer{}\csromanspacer{}ทสมํ.\csromanspacer{} อรหตฺตวคฺโค อฏฺฐโม\footnote{ตติโย (สฺยา, ก)}.}

\chapterheadpageread{9. สีติวคฺค \textbf{ตสฺสุทฺทานํ}}
\csromangathameasure{ทุกฺขํ อรหตฺตํ อุตฺตริ จ, สุขํ อธิคเมน จ. \\ มหนฺตตฺตํ ทฺวยํ นิรเย, อคฺคธมฺมญฺจ รตฺติโยติ.}
\csromangathagroup{\csromangathastanza{\csromanfolio{379}ทุกฺขํ อรหตฺตํ อุตฺตริ จ, สุขํ อธิคเมน จ. \\ มหนฺตตฺตํ ทฺวยํ นิรเย\footnote{มหตฺตทฺวยนิรเย (สฺยา)}, อคฺคธมฺมญฺจ รตฺติโยติ.}}

\csromanmatikaanchor{mtk.194}
\csromantocmarknopage{section}{9. สีติวคฺค}{mtk.194}
\bookmark[dest={mtk.194},level=1]{9. สีติวคฺค}
\csromanheader{1}{9. สีติวคฺค}
\csromanmatikaanchor{mtk.195}
\csromantocmarkref{subsection}{1. สีติภาวสุตฺต}{mtk.195}
\bookmark[dest={mtk.195},level=2]{1. สีติภาวสุตฺต}
\csromanheader{2}{1. สีติภาวสุตฺต}
\proseitem{85}{ฉหิ ภิกฺขเว ธมฺเมหิ สมนฺนาคโต ภิกฺขุ อภพฺโพ อนุตฺตรํ สีติภาวํ สจฺฉิกาตุํ.\csromanspacer{} กตเมหิ ฉหิ? * อิธ ภิกฺขเว ภิกฺขุ ยสฺมิํ สมเย จิตฺตํ นิคฺคเหตพฺพํ, ตสฺมิํ สมเย จิตฺตํ น นิคฺคณฺหาติ.\csromanspacer{} ยสฺมิํ สมเย จิตฺตํ ปคฺคเหตพฺพํ, ตสฺมิํ สมเย จิตฺตํ น ปคฺคณฺหาติ.\csromanspacer{} ยสฺมิํ สมเย จิตฺตํ สมฺปหํสิตพฺพํ, ตสฺมิํ สมเย จิตฺตํ น สมฺปหํเสติ.\csromanspacer{} ยสฺมิํ สมเย จิตฺตํ อชฺฌุเปกฺขิตพฺพํ, ตสฺมิํ สมเย จิตฺตํ น อชฺฌุเปกฺขติ.\csromanspacer{} หีนาธิมุตฺติโก จ โหติ สกฺกายาภิรโต จ.\csromanspacer{} อิเมหิ โข ภิกฺขเว ฉหิ ธมฺเมหิ สมนฺนาคโต ภิกฺขุ อภพฺโพ อนุตฺตรํ สีติภาวํ สจฺฉิกาตุํ.}

\prose{ฉหิ ภิกฺขเว ธมฺเมหิ สมนฺนาคโต ภิกฺขุ ภพฺโพ อนุตฺตรํ สีติภาวํ สจฺฉิกาตุํ.\csromanspacer{} กตเมหิ ฉหิ? อิธ ภิกฺขเว ภิกฺขุ ยสฺมิํ สมเย จิตฺตํ นิคฺคเหตพฺพํ, ตสฺมิํ สมเย จิตฺตํ นิคฺคณฺหาติ.\csromanspacer{} ยสฺมิํ สมเย จิตฺตํ ปคฺคเหตพฺพํ, ตสฺมิํ สมเย จิตฺตํ ปคฺคณฺหาติ.\csromanspacer{} ยสฺมิํ สมเย จิตฺตํ สมฺปหํสิตพฺพํ, ตสฺมิํ สมเย จิตฺตํ สมฺปหํเสติ.\csromanspacer{} ยสฺมิํ สมเย จิตฺตํ อชฺฌุเปกฺขิตพฺพํ, ตสฺมิํ สมเย จิตฺตํ อชฺฌุเปกฺขติ, ปณีตาธิมุตฺติโก จ โหติ นิพฺพานาภิรโต จ.\csromanspacer{} อิเมหิ โข ภิกฺขเว ฉหิ ธมฺเมหิ สมนฺนาคโต ภิกฺขุ ภพฺโพ อนุตฺตรํ สีติภาวํ สจฺฉิกาตุนฺติ.\csromanspacer{}\csromanspacer{}ปฐมํ.}

\csromanmatikaanchor{mtk.196}
\csromantocmarkref{subsection}{2. อาวรณสุตฺต}{mtk.196}
\bookmark[dest={mtk.196},level=2]{2. อาวรณสุตฺต}
\csromanheader{2}{2. อาวรณสุตฺต}
\proseitem{86}{ฉหิ ภิกฺขเว ธมฺเมหิ สมนฺนาคโต สุณนฺโตปิ สทฺธมฺมํ อภพฺโพ นิยามํ โอกฺกมิตุํ กุสเลสุ ธมฺเมสุ สมฺมตฺตํ.\csromanspacer{} กตเมหิ}

\vspace{\tipitakaparskip}
\csromancenter{ฉกฺกนิปาตปาฬิ ฉหิ? กมฺมาวรณตาย สมนฺนาคโต โหติ, กิเลสาวรณตาย สมนฺนาคโต โหติ, วิปากาวรณตาย สมนฺนาคโต โหติ, อสฺสทฺโธ จ โหติ อจฺฉนฺทิโก จ ทุปฺปญฺโญ จ.\csromanspacer{} อิเมหิ โข ภิกฺขเว ฉหิ ธมฺเมหิ สมนฺนาคโต สุณนฺโตปิ สทฺธมฺมํ อภพฺโพ นิยามํ โอกฺกมิตุํ กุสเลสุ ธมฺเมสุ สมฺมตฺตํ.}
\prose{\csromanfolio{380}ฉหิ ภิกฺขเว ธมฺเมหิ สมนฺนาคโต สุณนฺโต สทฺธมฺมํ ภพฺโพ นิยามํ โอกฺกมิตุํ กุสเลสุ ธมฺเมสุ สมฺมตฺตํ.\csromanspacer{} กตเมหิ ฉหิ? น กมฺมาวรณตาย สมนฺนาคโต โหติ, น กิเลสาวรณตาย สมนฺนาคโต โหติ, น วิปากาวรณตาย สมนฺนาคโต โหติ, สทฺโธ จ โหติ ฉนฺทิโก จ ปญฺญวา จ.\csromanspacer{} อิเมหิ โข ภิกฺขเว ฉหิ ธมฺเมหิ สมนฺนาคโต สุณนฺโต สทฺธมฺมํ ภพฺโพ นิยามํ โอกฺกมิตุํ กุสเลสุ ธมฺเมสุ สมฺมตฺตนฺติ.\csromanspacer{}\csromanspacer{}ทุติยํ.}

\csromanmatikaanchor{mtk.197}
\csromantocmarkref{subsection}{3. โวโรปิตสุตฺต}{mtk.197}
\bookmark[dest={mtk.197},level=2]{3. โวโรปิตสุตฺต}
\csromanheader{2}{3. โวโรปิตสุตฺต}
\proseitem{87}{ฉหิ ภิกฺขเว ธมฺเมหิ สมนฺนาคโต สุณนฺโตปิ สทฺธมฺมํ อภพฺโพ นิยามํ โอกฺกมิตุํ กุสเลสุ ธมฺเมสุ สมฺมตฺตํ.\csromanspacer{} กตเมหิ ฉหิ? มาตา ชีวิตา โวโรปิตา โหติ, ปิตา ชีวิตา โวโรปิโต โหติ, อรหํ\footnote{อรหา (สฺยา, กํ), อรหนฺโต (ก)} ชีวิตา โวโรปิโต โหติ, ตถาคตสฺส ทุฏฺเฐน จิตฺเตน โลหิตํ อุปฺปาทิตํ โหติ, สํโฆ ภินฺโน โหติ, ทุปฺปญฺโญ โหติ ชโฬ เอฬมูโค.\csromanspacer{} อิเมหิ โข ภิกฺขเว ฉหิ ธมฺเมหิ สมนฺนาคโต สุณนฺโตปิ สทฺธมฺมํ อภพฺโพ นิยามํ โอกฺกมิตุํ กุสเลสุ ธมฺเมสุ สมฺมตฺตํ.}

\prose{ฉหิ ภิกฺขเว ธมฺเมหิ สมนฺนาคโต สุณนฺโต สทฺธมฺมํ ภพฺโพ นิยามํ โอกฺกมิตุํ กุสเลสุ ธมฺเมสุ สมฺมตฺตํ.\csromanspacer{} กตเมหิ ฉหิ? น มาตา ชีวิตา โวโรปิตา โหติ, น ปิตา ชีวิตา โวโรปิโต โหติ, น อรหํ ชีวิตา โวโรปิโต โหติ, น ตถาคตสฺส ทุฏฺเฐน จิตฺเตน โลหิตํ อุปฺปาทิตํ โหติ, น สํโฆ ภินฺโน โหติ, ปญฺญวา โหติ อชโฬ อเนฬมูโค.\csromanspacer{} อิเมหิ โข ภิกฺขเว ฉหิ ธมฺเมหิ}

\vspace{\tipitakaparskip}
\csromancenter{สีติวคฺค สมนฺนาคโต สุณนฺโต สทฺธมฺมํ ภพฺโพ นิยามํ โอกฺกมิตุํ กุสเลสุ ธมฺเมสุ สมฺมตฺตนฺติ.\csromanspacer{}\csromanspacer{}ตติยํ.}
\csromanheader{2}{4. สุสฺสูสติสุตฺต}
\csromanmatikaanchor{mtk.198}
\csromantocmarkref{subsection}{4-6. สุสฺสูสติสุตฺตาทิ}{mtk.198}
\bookmark[dest={mtk.198},level=2]{4-6. สุสฺสูสติสุตฺตาทิ}
\proseitem{88}{\csromanfolio{381}ฉหิ ภิกฺขเว ธมฺเมหิ สมนฺนาคโต สุณนฺโตปิ สทฺธมฺมํ อภพฺโพ นิยามํ โอกฺกมิตุํ กุสเลสุ ธมฺเมสุ สมฺมตฺตํ.\csromanspacer{} กตเมหิ ฉหิ? ตถาคตปฺปเวทิเต ธมฺมวินเย เทสิยมาเน น สุสฺสูสติ, น โสตํ โอทหติ, น อญฺญา จิตฺตํ อุปฏฺฐาเปติ\footnote{อุปฏฺฐเปติ (สี, สฺยา, กํ, อิ)}, อนตฺถํ คณฺหาติ, อตฺถํ ริญฺจติ, อนนุโลมิกาย ขนฺติยา สมนฺนาคโต โหติ.\csromanspacer{} อิเมหิ โข ภิกฺขเว ฉหิ ธมฺเมหิ สมนฺนาคโต สุณนฺโตปิ สทฺธมฺมํ อภพฺโพ นิยามํ โอกฺกมิตุํ กุสเลสุ ธมฺเมสุ สมฺมตฺตํ.}

\prose{ฉหิ ภิกฺขเว ธมฺเมหิ สมนฺนาคโต สุณนฺโต สทฺธมฺมํ ภพฺโพ นิยามํ โอกฺกมิตุํ กุสเลสุ ธมฺเมสุ สมฺมตฺตํ.\csromanspacer{} กตเมหิ ฉหิ? ตถาคตปฺปเวทิเต ธมฺมวินเย เทสิยมาเน สุสฺสูสติ, โสตํ โอทหติ, อญฺญา จิตฺตํ อุปฏฺฐาเปติ, อตฺถํ คณฺหาติ, อนตฺถํ ริญฺจติ, อนุโลมิกาย ขนฺติยา สมนฺนาคโต โหติ.\csromanspacer{} อิเมหิ โข ภิกฺขเว ฉหิ ธมฺเมหิ สมนฺนาคโต สุณนฺโต สทฺธมฺมํ ภพฺโพ นิยามํ โอกฺกมิตุํ กุสเลสุ ธมฺเมสุ สมฺมตฺตนฺติ.\csromanspacer{}\csromanspacer{}จตุตฺถํ.}

\csromanheader{2}{5. อปฺปหายสุตฺต}
\proseitem{89}{ฉ ภิกฺขเว ธมฺเม อปฺปหาย อภพฺโพ ทิฏฺฐิสมฺปทํ สจฺฉิกาตุํ.\csromanspacer{} กตเม ฉ? สกฺกายทิฏฺฐิํ วิจิกิจฺฉํ สีลพฺพตปรามาสํ อปายคมนียํ ราคํ อปายคมนียํ โทสํ อปายคมนียํ โมหํ.\csromanspacer{} อิเม โข ภิกฺขเว ฉ ธมฺเม อปฺปหาย อภพฺโพ ทิฏฺฐิสมฺปทํ สจฺฉิกาตุํ.}

\prose{ฉ ภิกฺขเว ธมฺเม ปหาย ภพฺโพ ทิฏฺฐิสมฺปทํ สจฺฉิกาตุํ.\csromanspacer{} กตเม ฉ? สกฺกายทิฏฺฐิํ วิจิกิจฺฉํ สีลพฺพตปรามาสํ อปายคมนียํ ราคํ อปายคมนียํ โทสํ อปายคมนียํ โมหํ.\csromanspacer{} อิเม โข ภิกฺขเว ฉ ธมฺเม ปหาย ภพฺโพ ทิฏฺฐิสมฺปทํ สจฺฉิกาตุนฺติ.\csromanspacer{}\csromanspacer{}ปญฺจมํ.}

\gambhira{ฉกฺกนิปาตปาฬิ \textbf{6. ปหีนสุตฺต}}
\proseitem{90}{\csromanfolio{382}ฉยิเม ภิกฺขเว ธมฺมา ทิฏฺฐิสมฺปนฺนสฺส ปุคฺคลสฺส ปหีนา.\csromanspacer{} กตเม ฉ? สกฺกายทิฏฺฐิ วิจิกิจฺฉา สีลพฺพตปรามาโส อปายคมนีโย ราโค อปายคมนีโย โทโส อปายคมนีโย โมโห.\csromanspacer{} อิเม โข ภิกฺขเว ฉ ธมฺมา ทิฏฺฐิสมฺปนฺนสฺส ปุคฺคลสฺส ปหีนาติ.\csromanspacer{}\csromanspacer{}ฉฏฺฐํ.}

\csromanmatikaanchor{mtk.199}
\csromantocmarkref{subsection}{7. อภพฺพสุตฺต}{mtk.199}
\bookmark[dest={mtk.199},level=2]{7. อภพฺพสุตฺต}
\csromanheader{2}{7. อภพฺพสุตฺต}
\proseitem{91}{ฉ ภิกฺขเว ธมฺเม อภพฺโพ ทิฏฺฐิสมฺปนฺโน ปุคฺคโล อุปฺปาเทตุํ.\csromanspacer{} กตเม ฉ? สกฺกายทิฏฺฐิํ วิจิกิจฺฉํ สีลพฺพตปรามาสํ อปายคมนียํ ราคํ อปายคมนียํ โทสํ อปายคมนียํ โมหํ.\csromanspacer{} อิเม โข ภิกฺขเว ฉ ธมฺเม อภพฺโพ ทิฏฺฐิสมฺปนฺโน ปุคฺคโล อุปฺปาเทตุนฺติ.\csromanspacer{}\csromanspacer{}สตฺตมํ.}

\csromanmatikaanchor{mtk.200}
\csromantocmarkref{subsection}{8-11. อภพฺพฏฺฐานสุตฺต}{mtk.200}
\bookmark[dest={mtk.200},level=2]{8-11. อภพฺพฏฺฐานสุตฺต}
\csromanheader{2}{8. ปฐม อภพฺพฏฺฐานสุตฺต}
\proseitem{92}{ฉยิมานิ ภิกฺขเว อภพฺพฏฺฐานานิ.\csromanspacer{} กตมานิ ฉ? อภพฺโพ ทิฏฺฐิสมฺปนฺโน ปุคฺคโล สตฺถริ อคารโว วิหริตุํ อปฺปติสฺโส, อภพฺโพ ทิฏฺฐิสมฺปนฺโน ปุคฺคโล ธมฺเม อคารโว วิหริตุํ อปฺปติสฺโส, อภพฺโพ ทิฏฺฐิสมฺปนฺโน ปุคฺคโล สํเฆ อคารโว วิหริตุํ อปฺปติสฺโส, อภพฺโพ ทิฏฺฐิสมฺปนฺโน ปุคฺคโล สิกฺขาย อคารโว วิหริตุํ อปฺปติสฺโส, อภพฺโพ ทิฏฺฐิสมฺปนฺโน ปุคฺคโล อนาคมนียํ วตฺถุํ ปจฺจาคนฺตุํ, อภพฺโพ ทิฏฺฐิสมฺปนฺโน ปุคฺคโล อฏฺฐมํ ภวํ นิพฺพตฺเตตุํ.\csromanspacer{} อิมานิ โข ภิกฺขเว ฉ อภพฺพฏฺฐานานีติ.\csromanspacer{}\csromanspacer{}อฏฺฐมํ.}

\csromanheader{2}{9. ทุติย อภพฺพฏฺฐานสุตฺต}
\proseitem{93}{ฉยิมานิ ภิกฺขเว อภพฺพฏฺฐานานิ.\csromanspacer{} กตมานิ ฉ? อภพฺโพ ทิฏฺฐิสมฺปนฺโน ปุคฺคโล กญฺจิ\footnote{กิญฺจิ (ก) อภิ 2. 348; ม 3. 109 ปิฏฺเฐสุปิ. 2. อานนฺตริยกมฺมํ (สี), อนนฺตริยกมฺมํ (สฺยา, อิ) อํ 1. 468 ปิฏฺเฐ ปสฺสิตพฺพํ.} สงฺขารํ นิจฺจโต อุปคนฺตุํ, อภพฺโพ ทิฏฺฐิสมฺปนฺโน ปุคฺคโล กญฺจิ สงฺขารํ สุขโต อุปคนฺตุํ, อภพฺโพ ทิฏฺฐิสมฺปนฺโน ปุคฺคโล กญฺจิ ธมฺมํ อตฺตโต อุปคนฺตุํ, อภพฺโพ ทิฏฺฐิสมฺปนฺโน ปุคฺคโล อานนฺตริยํ กมฺมํ2 กาตุํ, อภพฺโพ ทิฏฺฐิสมฺปนฺโน ปุคฺคโล โกตูหลมงฺคเลน สุทฺธิํ}

\vspace{\tipitakaparskip}
\csromancenter{สีติวคฺค ปจฺจาคนฺตุํ, อภพฺโพ ทิฏฺฐิสมฺปนฺโน ปุคฺคโล อิโต พหิทฺธา ทกฺขิเณยฺยํ คเวสิตุํ.\csromanspacer{} อิมานิ โข ภิกฺขเว ฉ อภพฺพฏฺฐานานีติ.\csromanspacer{}\csromanspacer{}นวมํ.}
\csromanheader{2}{10. ตติย อภพฺพฏฺฐานสุตฺต}
\proseitem{94}{\csromanfolio{383}ฉยิมานิ ภิกฺขเว อภพฺพฏฺฐานานิ.\csromanspacer{} กตมานิ ฉ? อภพฺโพ ทิฏฺฐิสมฺปนฺโน ปุคฺคโล มาตรํ ชีวิตา โวโรเปตุํ, อภพฺโพ ทิฏฺฐิสมฺปนฺโน ปุคฺคโล ปิตรํ ชีวิตา โวโรเปตุํ, อภพฺโพ ทิฏฺฐิสมฺปนฺโน ปุคฺคโล อรหนฺตํ ชีวิตา โวโรเปตุํ, อภพฺโพ ทิฏฺฐิสมฺปนฺโน ปุคฺคโล ตถาคตสฺส ทุฏฺเฐน จิตฺเตน โลหิตํ อุปฺปาเทตุํ, อภพฺโพ ทิฏฺฐิสมฺปนฺโน ปุคฺคโล สํฆํ ภินฺทิตุํ, อภพฺโพ ทิฏฺฐิสมฺปนฺโน ปุคฺคโล อญฺญํ สตฺถารํ อุทฺทิสิตุํ.\csromanspacer{} อิมานิ โข ภิกฺขเว ฉ อภพฺพฏฺฐานานีติ.\csromanspacer{}\csromanspacer{}ทสมํ.}

\csromanheader{2}{11. จตุตฺถ อภพฺพฏฺฐานสุตฺต}
\proseitem{95}{ฉยิมานิ ภิกฺขเว อภพฺพฏฺฐานานิ.\csromanspacer{} กตมานิ ฉ? อภพฺโพ ทิฏฺฐิสมฺปนฺโน ปุคฺคโล สยํกตํ สุขทุกฺขํ ปจฺจาคนฺตุํ, อภพฺโพ ทิฏฺฐิสมฺปนฺโน ปุคฺคโล ปรํกตํ\footnote{ปรกตํ (สี, สฺยา)} สุขทุกฺขํ ปจฺจาคนฺตุํ, อภพฺโพ ทิฏฺฐิสมฺปนฺโน ปุคฺคโล สยํกตญฺจ ปรํกตญฺจ สุขทุกฺขํ ปจฺจาคนฺตุํ, อภพฺโพ ทิฏฺฐิสมฺปนฺโน ปุคฺคโล อสยํการํ อธิจฺจสมุปฺปนฺนํ สุขทุกฺขํ ปจฺจาคนฺตุํ, อภพฺโพ ทิฏฺฐิสมฺปนฺโน ปุคฺคโล อปรํการํ อธิจฺจสมุปฺปนฺนํ สุขทุกฺขํ ปจฺจาคนฺตุํ, อภพฺโพ ทิฏฺฐิสมฺปนฺโน ปุคฺคโล อสยํการญฺจ อปรํการญฺจ อธิจฺจสมุปฺปนฺนํ สุขทุกฺขํ ปจฺจาคนฺตุํ.\csromanspacer{} ตํ กิสฺส เหตุ? ตถา หิสฺส ภิกฺขเว ทิฏฺฐิสมฺปนฺนสฺส ปุคฺคลสฺส เหตุ จ สุทิฏฺโฐ เหตุสมุปฺปนฺนา จ ธมฺมา.\csromanspacer{} อิมานิ โข ภิกฺขเว ฉ อภพฺพฏฺฐานานีติ.\csromanspacer{}\csromanspacer{}เอกาทสมํ.\csromanspacer{} สีติวคฺโค นวโม\footnote{จตุตฺโถ (สฺยา, ก)}.}

\titlehead{\textbf{ตสฺสุทฺทานํ}}
\csromangathameasure{สีติภาวํ อาวรณํ, โวโรปิตา สุสูสติ. \\ อปฺปหาย ปหีนา’ภพฺโพ, ตฏฺฐานา จตุโรปิ จาติ.}
\csromangathagroup{\csromangathastanza{สีติภาวํ อาวรณํ, โวโรปิตา สุสูสติ. \\ อปฺปหาย ปหีนา’ภพฺโพ, ตฏฺฐานา จตุโรปิ จาติ.}}

\csromanmatikaanchor{mtk.201}
\csromantocmarknopage{section}{10. อานิสํสวคฺค}{mtk.201}
\bookmark[dest={mtk.201},level=1]{10. อานิสํสวคฺค}
\csromanheader{1}{ฉกฺกนิปาตปาฬิ \textbf{10. อานิสํสวคฺค}}
\csromanmatikaanchor{mtk.202}
\csromantocmarkref{subsection}{1. ปาตุภาวสุตฺต}{mtk.202}
\bookmark[dest={mtk.202},level=2]{1. ปาตุภาวสุตฺต}
\csromanheader{2}{1. ปาตุภาวสุตฺต}
\csromanmatikaanchor{mtk.203}
\csromanmatikaanchor{mtk.204}
\csromantocmarkref{subsection}{2. อานิสํสสุตฺต}{mtk.203}
\bookmark[dest={mtk.203},level=2]{2. อานิสํสสุตฺต}
\csromantocmarkref{subsection}{3-6. อนิจฺจสุตฺตาทิ}{mtk.204}
\bookmark[dest={mtk.204},level=2]{3-6. อนิจฺจสุตฺตาทิ}
\proseitem{96}{\csromanfolio{384}ฉนฺนํ ภิกฺขเว ปาตุภาโว ทุลฺลโภ โลกสฺมิํ.\csromanspacer{} กตเมสํ ฉนฺนํ? ตถาคตสฺส อรหโต สมฺมาสมฺพุทฺธสฺส ปาตุภาโว ทุลฺลโภ โลกสฺมิํ, ตถาคตปฺปเวทิตสฺส ธมฺมวินยสฺส เทเสตา ปุคฺคโล ทุลฺลโภ โลกสฺมิํ, อริยายตเน ปจฺจาชาติ ทุลฺลภา\footnote{ปจฺจาชาโต ทุลฺลโภ (สฺยา)} โลกสฺมิํ, อินฺทฺริยานํ อเวกลฺลตา ทุลฺลภา โลกสฺมิํ, อชฬตา อเนฬมูคตา ทุลฺลภา โลกสฺมิํ, กุสเล ธมฺเม ฉนฺโท\footnote{กุสลธมฺมจฺฉนฺโท (สี, สฺยา, อิ)} ทุลฺลโภ โลกสฺมิํ.\csromanspacer{} อิเมสํ โข ภิกฺขเว ฉนฺนํ ปาตุภาโว ทุลฺลโภ โลกสฺมินฺติ.\csromanspacer{}\csromanspacer{}ปฐมํ.}

\csromanheader{2}{2. อานิสํสสุตฺต}
\proseitem{97}{ฉยิเม ภิกฺขเว อานิสํสา โสตาปตฺติผลสจฺฉิกิริยาย.\csromanspacer{} กตเม ฉ? สทฺธมฺมนิยโต โหติ, อปริหานธมฺโม โหติ, ปริยนฺตกต’สฺส ทุกฺขํ โหติ\footnote{ทุกฺขํ น โหติ (สฺยา, อิ, ก)}, อสาธารเณน ญาเณน สมนฺนาคโต โหติ, เหตุ จสฺส สุทิฏฺโฐ เหตุสมุปฺปนฺนา จ ธมฺมา.\csromanspacer{} อิเม โข ภิกฺขเว ฉ อานิสํสา โสตาปตฺติผลสจฺฉิกิริยายาติ.\csromanspacer{}\csromanspacer{}ทุติยํ.}

\csromanheader{2}{3. อนิจฺจสุตฺต}
\proseitem{98}{โส วต ภิกฺขเว ภิกฺขุ กญฺจิ สงฺขารํ นิจฺจโต สมนุปสฺสนฺโต อนุโลมิกาย ขนฺติยา สมนฺนาคโต ภวิสฺสตีติ เนตํ ฐานํ วิชฺชติ, อนุโลมิกาย ขนฺติยา อสมนฺนาคโต สมฺมตฺตนิยามํ โอกฺกมิสฺสตีติ เนตํ ฐานํ วิชฺชติ, สมฺมตฺตนิยามํ อโนกฺกมมาโน โสตาปตฺติผลํ วา สกทาคามิผลํ วา อนาคามิผลํ วา อรหตฺตํ\footnote{อรหตฺตผลํ (ก) ขุ 9. 409 ปิฏฺเฐปิ.} วา สจฺฉิกริสฺสตีติ เนตํ ฐานํ วิชฺชติ.}

\prose{โส วต ภิกฺขเว ภิกฺขุ สพฺพสงฺขาเร\footnote{สพฺพสงฺขารํ (สี, อิ)} อนิจฺจโต สมนุปสฺสนฺโต อนุโลมิกาย ขนฺติยา สมนฺนาคโต ภวิสฺสตีติ ฐานเมตํ วิชฺชติ, อนุโลมิกาย ขนฺติยา สมนฺนาคโต สมฺมตฺตนิยามํ โอกฺกมิสฺสตีติ ฐานเมตํ วิชฺชติ, สมฺมตฺตนิยามํ โอกฺกมมาโน โสตาปตฺติผลํ วา}

\vspace{\tipitakaparskip}
\csromancenter{อานิสํสวคฺค สกทาคามิผลํ วา อนาคามิผลํ วา อรหตฺตํ วา สจฺฉิกริสฺสตีติ ฐานเมตํ วิชฺชตีติ.\csromanspacer{}\csromanspacer{}ตติยํ.}
\csromanheader{2}{4. ทุกฺขสุตฺต}
\csromanmatikaanchor{mtk.205}
\csromantocmarkref{subsection}{7-11. อนวตฺถิตสุตฺตาทิ}{mtk.205}
\bookmark[dest={mtk.205},level=2]{7-11. อนวตฺถิตสุตฺตาทิ}
\proseitem{99}{\csromanfolio{385}โส วต ภิกฺขเว ภิกฺขุ กญฺจิ สงฺขารํ สุขโต สมนุปสฺสนฺโต ฯเปฯ สพฺพสงฺขาเร ทุกฺขโต สมนุปสฺสนฺโต ฯเปฯ ฐานเมตํ วิชฺชติ.\csromanspacer{}\csromanspacer{}จตุตฺถํ.}

\csromanheader{2}{5. อนตฺตสุตฺต}
\proseitem{100}{โส วต ภิกฺขเว ภิกฺขุ กญฺจิ ธมฺมํ อตฺตโต สมนุปสฺสนฺโต ฯเปฯ สพฺพธมฺเม\footnote{สพฺพธมฺมํ (สี, อิ), กิญฺจิธมฺมํ (ก)} อนตฺตโต สมนุปสฺสนฺโต ฯเปฯ ฐานเมตํ วิชฺชติ.\csromanspacer{}\csromanspacer{}ปญฺจมํ.}

\csromanheader{2}{6. นิพฺพานสุตฺต}
\proseitem{101}{โส วต ภิกฺขเว ภิกฺขุ นิพฺพานํ ทุกฺขโต สมนุปสฺสนฺโต อนุโลมิกาย ขนฺติยา สมนฺนาคโต ภวิสฺสตีติ เนตํ ฐานํ วิชฺชติ, อนุโลมิกาย ขนฺติยา อสมนฺนาคโต สมฺมตฺตนิยามํ โอกฺกมิสฺสตีติ เนตํ ฐานํ วิชฺชติ, สมฺมตฺตนิยามํ อโนกฺกมมาโน โสตาปตฺติผลํ วา สกทาคามิผลํ วา อนาคามิผลํ วา อรหตฺตํ วา สจฺฉิกริสฺสตีติ เนตํ ฐานํ วิชฺชติ.}

\prose{โส วต ภิกฺขเว ภิกฺขุ นิพฺพานํ สุขโต สมนุปสฺสนฺโต อนุโลมิกาย ขนฺติยา สมนฺนาคโต ภวิสฺสตีติ ฐานเมตํ วิชฺชติ, อนุโลมิกาย ขนฺติยา สมนฺนาคโต สมฺมตฺตนิยามํ โอกฺกมิสฺสตีติ ฐานเมตํ วิชฺชติ, สมฺมตฺตนิยามํ โอกฺกมมาโน โสตาปตฺติผลํ วา สกทาคามิผลํ วา อนาคามิผลํ วา อรหตฺตํ วา สจฺฉิกริสฺสตีติ ฐานเมตํ วิชฺชตีติ.\csromanspacer{}\csromanspacer{}ฉฏฺฐํ.}

\csromanheader{2}{7. อนวตฺถิตสุตฺต}
\proseitem{102}{ฉ ภิกฺขเว อานิสํเส สมฺปสฺสมาเนน อลเมว ภิกฺขุนา สพฺพสงฺขาเรสุ อโนธิํ กริตฺวา อนิจฺจสญฺญํ อุปฏฺฐาเปตุํ.\csromanspacer{} กตเม ฉ, สพฺพสงฺขารา จ เม อนวตฺถิตา\footnote{อนวฏฺฐิตโต (สี, สฺยา, อิ)} ขายิสฺสนฺติ, สพฺพโลเก จ เม มโน}

\vspace{\tipitakaparskip}
\csromancenter{ฉกฺกนิปาตปาฬิ นาภิรมิสฺสติ\footnote{น รมิสฺสติ (ก)}, สพฺพโลกา จ เม มโน วุฏฺฐหิสฺสติ, นิพฺพานโปณญฺจ เม มานสํ ภวิสฺสติ.\csromanspacer{} สํโยชนา จ เม ปหานํ คจฺฉิสฺสนฺติ\footnote{คจฺฉนฺติ (สฺยา, อิ, ก)}, ปรเมน จ สามญฺเญน สมนฺนาคโต ภวิสฺสามีติ.\csromanspacer{} อิเม โข ภิกฺขเว ฉ อานิสํเส สมฺปสฺสมาเนน อลเมว ภิกฺขุนา สพฺพสงฺขาเรสุ อโนธิํ กริตฺวา อนิจฺจสญฺญํ อุปฏฺฐาเปตุนฺติ.\csromanspacer{}\csromanspacer{}สตฺตมํ.}
\csromanheader{2}{8. อุกฺขิตฺตาสิกสุตฺต}
\proseitem{103}{\csromanfolio{386}ฉ ภิกฺขเว อานิสํเส สมฺปสฺสมาเนน อลเมว ภิกฺขุนา สพฺพสงฺขาเรสุ อโนธิํ กริตฺวา ทุกฺขสญฺญํ อุปฏฺฐาเปตุํ.\csromanspacer{} กตเม ฉ, สพฺพสงฺขาเรสุ จ เม นิพฺพิทสญฺญา ปจฺจุปฏฺฐิตา ภวิสฺสติ.\csromanspacer{} เสยฺยถาปิ อุกฺขิตฺตาสิเก วธเก, สพฺพโลกา จ เม มโน วุฏฺฐหิสฺสติ, นิพฺพาเน จ สนฺตทสฺสาวี ภวิสฺสามิ, อนุสยา จ เม สมุคฺฆาตํ คจฺฉิสฺสนฺติ\footnote{คจฺฉนฺติ (อิ, ก)}, กิจฺจการี จ ภวิสฺสามิ, สตฺถา จ เม ปริจิณฺโณ ภวิสฺสติ เมตฺตาวตายาติ.\csromanspacer{} อิเม โข ภิกฺขเว ฉ อานิสํเส สมฺปสฺสมาเนน อลเมว ภิกฺขุนา สพฺพสงฺขาเรสุ อโนธิํ กริตฺวา ทุกฺขสญฺญํ อุปฏฺฐาเปตุนฺติ.\csromanspacer{}\csromanspacer{}อฏฺฐมํ.}

\csromanheader{2}{9. อตมฺมยสุตฺต}
\proseitem{104}{ฉ ภิกฺขเว อานิสํเส สมฺปสฺสมาเนน อลเมว ภิกฺขุนา สพฺพธมฺเมสุ อโนธิํ กริตฺวา อนตฺตสญฺญํ อุปฏฺฐาเปตุํ.\csromanspacer{} กตเม ฉ, สพฺพโลเก จ อตมฺมโย ภวิสฺสามิ, อหงฺการา จ เม อุปรุชฺฌิสฺสนฺติ, มมงฺการา จ เม อุปรุชฺฌิสฺสนฺติ, อสาธารเณน จ ญาเณน สมนฺนาคโต ภวิสฺสามิ, เหตุ จ เม สุทิฏฺโฐ ภวิสฺสติ เหตุสมุปฺปนฺนา จ ธมฺมา.\csromanspacer{} อิเม โข ภิกฺขเว ฉ อานิสํเส สมฺปสฺสมาเนน อลเมว ภิกฺขุนา สพฺพธมฺเมสุ อโนธิํ กริตฺวา อนตฺตสญฺญํ อุปฏฺฐาเปตุนฺติ.\csromanspacer{}\csromanspacer{}นวมํ.}

\csromanheader{2}{10. ภวสุตฺต}
\proseitem{105}{ตโยเม ภิกฺขเว ภวา ปหาตพฺพา, ตีสุ สิกฺขาสุ สิกฺขิตพฺพํ.\csromanspacer{} กตเม ตโย ภวา ปหาตพฺพา, กามภโว รูปภโว}

\vspace{\tipitakaparskip}
\csromancenter{อานิสํสวคฺค อรูปภโว.\csromanspacer{} อิเม ตโย ภวา ปหาตพฺพา.\csromanspacer{} กตมาสุ ตีสุ สิกฺขาสุ สิกฺขิตพฺพํ, อธิสีลสิกฺขาย อธิจิตฺตสิกฺขาย อธิปญฺญาสิกฺขาย.\csromanspacer{} อิมาสุ ตีสุ สิกฺขาสุ สิกฺขิตพฺพํ.\csromanspacer{} ยโต โข ภิกฺขเว ภิกฺขุโน อิเม ตโย ภวา ปหีนา โหนฺติ, อิมาสุ จ ตีสุ สิกฺขาสุ สิขิตสิกฺโข โหติ.\csromanspacer{} อยํ วุจฺจติ ภิกฺขเว ภิกฺขุ อจฺเฉจฺฉิ ตณฺหํ, วิวตฺตยิ สํโยชนํ, สมฺมา มานาภิสมยา อนฺตมกาสิ ทุกฺขสฺสาติ.\csromanspacer{}\csromanspacer{}ทสมํ.}
\csromanheader{2}{11. ตณฺหาสุตฺต}
\proseitem{106}{\csromanfolio{387}ติสฺโส อิมา ภิกฺขเว ตณฺหา ปหาตพฺพา ตโย จ มานา.\csromanspacer{} กตมา ติสฺโส ตณฺหา ปหาตพฺพา, กามตณฺหา ภวตณฺหา วิภวตณฺหา.\csromanspacer{} อิมา ติสฺโส ตณฺหา ปหาตพฺพา.\csromanspacer{} กตเม ตโย มานา ปหาตพฺพา, มาโน โอมาโน อติมาโน.\csromanspacer{} อิเม ตโย มานา ปหาตพฺพา.\csromanspacer{} ยโต โข ภิกฺขเว ภิกฺขุโน อิมา ติสฺโส ตณฺหา ปหีนา โหนฺติ อิเม จ ตโย มานา.\csromanspacer{} อยํ วุจฺจติ ภิกฺขเว ภิกฺขุ อจฺเฉจฺฉิ ตณฺหํ, วิวตฺตยิ สํโยชนํ, สมฺมา มานาภิสมยา อนฺตมกาสิ ทุกฺขสฺสาติ.\csromanspacer{}\csromanspacer{}เอกาทสมํ.\csromanspacer{} อานิสํสวคฺโค ทสโม\footnote{ปญฺจโม (สฺยา, ก)}.}

\titlehead{\textbf{ตสฺสุทฺทานํ}}
\csromangathameasure{ปาตุภาโว อานิสํโส, \\ อนิจฺจทุกฺข-อนตฺตโต. \\ นิพฺพานํ อนวตฺถิ อุกฺขิตฺตาสิ อตมฺมโย, \\ ภวา ตณฺหาเยกา ทสาติ. \\ \textbf{ทุติโย ปณฺณาสโก สมตฺโต.}}
\csromangathagroup{\csromangathastanza{ปาตุภาโว อานิสํโส, \\ อนิจฺจทุกฺข-อนตฺตโต. \\ นิพฺพานํ อนวตฺถิ อุกฺขิตฺตาสิ อตมฺมโย, \\ ภวา ตณฺหาเยกา ทสาติ.}\csromangathastanza{\textbf{ทุติโย ปณฺณาสโก สมตฺโต.}}}

\csromanmatikaanchor{mtk.206}
\csromantocmarknopage{section}{11. ติกวคฺค}{mtk.206}
\bookmark[dest={mtk.206},level=1]{11. ติกวคฺค}
\csromanheader{1}{ฉกฺกนิปาตปาฬิ \textbf{11. ติกวคฺค}}
\csromanheader{2}{1. ราคสุตฺต}
\csromanmatikaanchor{mtk.207}
\csromantocmarkref{subsection}{1-10. ราคสุตฺตาทิ}{mtk.207}
\bookmark[dest={mtk.207},level=2]{1-10. ราคสุตฺตาทิ}
\proseitem{107}{\csromanfolio{388}ตโยเม ภิกฺขเว ธมฺมา.\csromanspacer{} กตเม ตโย, ราโค โทโส โมโห.\csromanspacer{} อิเม โข ภิกฺขเว ตโย ธมฺมา.\csromanspacer{} อิเมสํ โข ภิกฺขเว ติณฺณํ ธมฺมานํ ปหานาย ตโย ธมฺมา ภาเวตพฺพา.\csromanspacer{} กตเม ตโย, ราคสฺส ปหานาย อสุภา ภาเวตพฺพา, โทสสฺส ปหานาย เมตฺตา ภาเวตพฺพา, โมหสฺส ปหานาย ปญฺญา ภาเวตพฺพา.\csromanspacer{} อิเมสํ โข ภิกฺขเว ติณฺณํ ธมฺมานํ ปหานาย อิเม ตโย ธมฺมา ภาเวตพฺพาติ.\csromanspacer{}\csromanspacer{}ปฐมํ.}

\csromanheader{2}{2. ทุจฺจริตสุตฺต}
\proseitem{108}{ตโยเม ภิกฺขเว ธมฺมา.\csromanspacer{} กตเม ตโย, กายทุจฺจริตํ วจีทุจฺจริตํ มโนทุจฺจริตํ.\csromanspacer{} อิเม โข ภิกฺขเว ตโย ธมฺมา.\csromanspacer{} อิเมสํ โข ภิกฺขเว ติณฺณํ ธมฺมานํ ปหานาย ตโย ธมฺมา ภาเวตพฺพา.\csromanspacer{} กตเม ตโย, กายทุจฺจริตสฺส ปหานาย กายสุจริตํ ภาเวตพฺพํ, วจีทุจฺจริตสฺส ปหานาย วจีสุจริตํ ภาเวตพฺพํ, มโนทุจฺจริตสฺส ปหานาย มโนสุจริตํ ภาเวตพฺพํ.\csromanspacer{} อิเมสํ โข ภิกฺขเว ติณฺณํ ธมฺมานํ ปหานาย อิเม ตโย ธมฺมา ภาเวตพฺพาติ.\csromanspacer{}\csromanspacer{}ทุติยํ.}

\csromanheader{2}{3. วิตกฺกสุตฺต}
\proseitem{109}{ตโยเม ภิกฺขเว ธมฺมา.\csromanspacer{} กตเม ตโย, กามวิตกฺโก พฺยาปาทวิตกฺโก วิหิํสาวิตกฺโก.\csromanspacer{} อิเม โข ภิกฺขเว ตโย ธมฺมา.\csromanspacer{} อิเมสํ โข ภิกฺขเว ติณฺณํ ธมฺมานํ ปหานาย ตโย ธมฺมา ภาเวตพฺพา.\csromanspacer{} กตเม ตโย, กามวิตกฺกสฺส ปหานาย เนกฺขมฺมวิตกฺโก ภาเวตพฺโพ, พฺยาปาทวิตกฺกสฺส ปหานาย อพฺยาปาท วิตกฺโก ภาเวตพฺโพ, วิหิํสาวิตกฺกสฺส ปหานาย อวิหิํสาวิตกฺโก ภาเวตพฺโพ.\csromanspacer{} อิเมสํ โข ภิกฺขเว ติณฺณํ ธมฺมานํ ปหานาย อิเม ตโย ธมฺมา ภาเวตพฺพาติ.\csromanspacer{}\csromanspacer{}ตติยํ.}

\chapterheadpageread{11. ติกวคฺค \textbf{4. สญฺญาสุตฺต}}
\proseitem{110}{\csromanfolio{389}ตโยเม ภิกฺขเว ธมฺมา.\csromanspacer{} กตเม ตโย, กามสญฺญา พฺยาปาทสญฺญา วิหิํสาสญฺญา.\csromanspacer{} อิเม โข ภิกฺขเว ตโย ธมฺมา.\csromanspacer{} อิเมสํ โข ภิกฺขเว ติณฺณํ ธมฺมานํ ปหานาย ตโย ธมฺมา ภาเวตพฺพา.\csromanspacer{} กตเม ตโย, กามสญฺญาย ปหานาย เนกฺขมฺมสญฺญา ภาเวตพฺพา, พฺยาปาทสญฺญาย ปหานาย อพฺยาปาทสญฺญา ภาเวตพฺพา, วิหิํสาสญฺญาย ปหานาย อวิหิํสาสญฺญา ภาเวตพฺพา.\csromanspacer{} อิเมสํ โข ภิกฺขเว ติณฺณํ ธมฺมานํ ปหานาย อิเม ตโย ธมฺมา ภาเวตพฺพาติ.\csromanspacer{}\csromanspacer{}จตุตฺถํ.}

\csromanheader{2}{5. ธาตุสุตฺต}
\proseitem{111}{ตโยเม ภิกฺขเว ธมฺมา.\csromanspacer{} กตเม ตโย, กามธาตุ พฺยาปาทธาตุ วิหิํสาธาตุ.\csromanspacer{} อิเม โข ภิกฺขเว ตโย ธมฺมา.\csromanspacer{} อิเมสํ โข ภิกฺขเว ติณฺณํ ธมฺมานํ ปหานาย ตโย ธมฺมา ภาเวตพฺพา.\csromanspacer{} กตเม ตโย, กามธาตุยา ปหานาย เนกฺขมฺมธาตุ ภาเวตพฺพา, พฺยาปาทธาตุยา ปหานาย อพฺยาปาทธาตุ ภาเวตพฺพา, วิหิํสาธาตุยา ปหานาย อวิหิํสาธาตุ ภาเวตพฺพา.\csromanspacer{} อิเมสํ โข ภิกฺขเว ติณฺณํ ธมฺมานํ ปหานาย อิเม ตโย ธมฺมา ภาเวตพฺพาติ.\csromanspacer{}\csromanspacer{}ปญฺจมํ.}

\csromanheader{2}{6. อสฺสาทสุตฺต}
\proseitem{112}{ตโยเม ภิกฺขเว ธมฺมา.\csromanspacer{} กตเม ตโย, อสฺสาททิฏฺฐิ อตฺตานุทิฏฺฐิ มิจฺฉาทิฏฺฐิ.\csromanspacer{} อิเม โข ภิกฺขเว ตโย ธมฺมา.\csromanspacer{} อิเมสํ โข ภิกฺขเว ติณฺณํ ธมฺมานํ ปหานาย ตโย ธมฺมา ภาเวตพฺพา.\csromanspacer{} กตเม ตโย, อสฺสาททิฏฺฐิยา ปหานาย อนิจฺจสญฺญา ภาเวตพฺพา, อตฺตานุทิฏฺฐิยา ปหานาย อนตฺตสญฺญา ภาเวตพฺพา, มิจฺฉาทิฏฺฐิยา ปหานาย สมฺมาทิฏฺฐิ ภาเวตพฺพา.\csromanspacer{} อิเมสํ โข ภิกฺขเว ติณฺณํ ธมฺมานํ ปหานาย อิเม ตโย ธมฺมา ภาเวตพฺพาติ.\csromanspacer{}\csromanspacer{}ฉฏฺฐํ.}

\csromanheader{2}{7. อรติสุตฺต}
\proseitem{113}{ตโยเม ภิกฺขเว ธมฺมา.\csromanspacer{} กตเม ตโย, อรติ วิหิํสา\footnote{วิเหสา (ก)} อธมฺมจริยา.\csromanspacer{} อิเม โข ภิกฺขเว ตโย ธมฺมา.\csromanspacer{} อิเมสํ โข ภิกฺขเว}

\vspace{\tipitakaparskip}
\csromancenter{ฉกฺกนิปาตปาฬิ ติณฺณํ ธมฺมานํ ปหานาย ตโย ธมฺมา ภาเวตพฺพา.\csromanspacer{} กตเม ตโย, อรติยา ปหานาย มุทิตา ภาเวตพฺพา, วิเหสาย ปหานาย อวิหิํสา ภาเวตพฺพา, อธมฺมจริยาย ปหานาย ธมฺมจริยา ภาเวตพฺพา.\csromanspacer{} อิเมสํ โข ภิกฺขเว ติณฺณํ ธมฺมานํ ปหานาย อิเม ตโย ธมฺมา ภาเวตพฺพาติ.\csromanspacer{}\csromanspacer{}สตฺตมํ.}
\csromanheader{2}{8. สนฺตุฏฺฐิตาสุตฺต}
\proseitem{114}{\csromanfolio{390}ตโยเม ภิกฺขเว ธมฺมา.\csromanspacer{} กตเม ตโย, อสนฺตุฏฺฐิตา อสมฺปชญฺญํ มหิจฺฉตา.\csromanspacer{} อิเม โข ภิกฺขเว ตโย ธมฺมา.\csromanspacer{} อิเมสํ โข ภิกฺขเว ติณฺณํ ธมฺมานํ ปหานาย ตโย ธมฺมา ภาเวตพฺพา.\csromanspacer{} กตเม ตโย, อสนฺตุฏฺฐิตาย ปหานาย สนฺตุฏฺฐิตา ภาเวตพฺพา, อสมฺปชญฺญสฺส ปหานาย สมฺปชญฺญํ ภาเวตพฺพํ, มหิจฺฉตาย ปหานาย อปฺปิจฺฉตา ภาเวตพฺพา.\csromanspacer{} อิเมสํ โข ภิกฺขเว ติณฺณํ ธมฺมานํ ปหานาย อิเม ตโย ธมฺมา ภาเวตพฺพาติ.\csromanspacer{}\csromanspacer{}อฏฺฐมํ.}

\csromanheader{2}{9. โทวจสฺสตาสุตฺต}
\proseitem{115}{ตโยเม ภิกฺขเว ธมฺมา.\csromanspacer{} กตเม ตโย, โทวจสฺสตา ปาปมิตฺตตา เจตโส วิกฺเขโป.\csromanspacer{} อิเม โข ภิกฺขเว ตโย ธมฺมา.\csromanspacer{} อิเมสํ โข ภิกฺขเว ติณฺณํ ธมฺมานํ ปหานาย ตโย ธมฺมา ภาเวตพฺพา.\csromanspacer{} กตเม ตโย, โทวจสฺสตาย ปหานาย โสวจสฺสตา ภาเวตพฺพา, ปาปมิตฺตตาย ปหานาย กลฺยาณมิตฺตตา ภาเวตพฺพา, เจตโส วิกฺเขปสฺส ปหานาย อานาปานสฺสติ ภาเวตพฺพา.\csromanspacer{} อิเมสํ โข ภิกฺขเว ติณฺณํ ธมฺมานํ ปหานาย อิเม ตโย ธมฺมา ภาเวตพฺพาติ.\csromanspacer{}\csromanspacer{}นวมํ.}

\csromanheader{2}{10. อุทฺธจฺจสุตฺต}
\proseitem{116}{ตโยเม ภิกฺขเว ธมฺมา.\csromanspacer{} กตเม ตโย, อุทฺธจฺจํ อสํวโร ปมาโท.\csromanspacer{} อิเม โข ภิกฺขเว ตโย ธมฺมา.\csromanspacer{} อิเมสํ โข ภิกฺขเว ติณฺณํ ธมฺมานํ ปหานาย ตโย ธมฺมา ภาเวตพฺพา.\csromanspacer{} กตเม ตโย, อุทฺธจฺจสฺส ปหานาย สมโถ ภาเวตพฺโพ, อสํวรสฺส ปหานาย สํวโร ภาเวตพฺโพ, ปมาทสฺส ปหานาย อปฺปมาโท ภาเวตพฺโพ.}

\vspace{\tipitakaparskip}
\csromancenter{สามญฺญวคฺค อิเมสํ โข ภิกฺขเว ติณฺณํ ธมฺมานํ ปหานาย อิเม ตโย ธมฺมา ภาเวตพฺพาติ.\csromanspacer{}\csromanspacer{}ทสมํ.\csromanspacer{} ติกวคฺโค เอกาทสโม\footnote{ปฐโม (สฺยา)}.}
\titlehead{\textbf{ตสฺสุทฺทานํ}}
\csromangathameasure{ราคทุจฺจริตวิตกฺก, สญฺญา ธาตูติ วุจฺจติ. \\ อสฺสาท-อรติตุฏฺฐิ, ทุเว จ อุทฺธจฺเจน วคฺโคติ.}
\csromangathagroup{\csromangathastanza{\csromanfolio{391}ราคทุจฺจริตวิตกฺก, สญฺญา ธาตูติ วุจฺจติ. \\ อสฺสาท-อรติตุฏฺฐิ, ทุเว จ อุทฺธจฺเจน วคฺโคติ.}}

\csromanmatikaanchor{mtk.208}
\csromantocmarknopage{section}{12. สามญฺญวคฺค}{mtk.208}
\bookmark[dest={mtk.208},level=1]{12. สามญฺญวคฺค}
\csromanheader{1}{12. สามญฺญวคฺค}
\csromanmatikaanchor{mtk.209}
\csromantocmarkref{subsection}{1. กายานุปสฺสีสุตฺต}{mtk.209}
\bookmark[dest={mtk.209},level=2]{1. กายานุปสฺสีสุตฺต}
\csromanheader{2}{1. กายานุปสฺสีสุตฺต}
\proseitem{117}{ฉ ภิกฺขเว ธมฺเม อปฺปหาย อภพฺโพ กาเย กายานุปสฺสี วิหริตุํ.\csromanspacer{} กตเม ฉ? กมฺมารามตํ ภสฺสารามตํ นิทฺทารามตํ สงฺคณิการามตํ อินฺทฺริเยสุ อคุตฺตทฺวารตํ โภชเน อมตฺตญฺญุตํ.\csromanspacer{} อิเม โข ภิกฺขเว ฉ ธมฺเม อปฺปหาย อภพฺโพ กาเย กายานุปสฺสี วิหริตุํ.}

\prose{ฉ ภิกฺขเว ธมฺเม ปหาย ภพฺโพ กาเย กายานุปสฺสี วิหริตุํ.\csromanspacer{} กตเม ฉ? กมฺมารามตํ ภสฺสารามตํ นิทฺทารามตํ สงฺคณิการามตํ อินฺทฺริเยสุ อคุตฺตทฺวารตํ โภชเน อมตฺตญฺญุตํ.\csromanspacer{} อิเม โข ภิกฺขเว ฉ ธมฺเม ปหาย ภพฺโพ กาเย กายานุปสฺสี วิหริตุนฺติ.\csromanspacer{}\csromanspacer{}ปฐมํ.}

\csromanmatikaanchor{mtk.210}
\csromantocmarkref{subsection}{2. ธมฺมานุปสฺสีสุตฺต}{mtk.210}
\bookmark[dest={mtk.210},level=2]{2. ธมฺมานุปสฺสีสุตฺต}
\csromanheader{2}{2. ธมฺมานุปสฺสีสุตฺต}
\proseitem{118}{ฉ ภิกฺขเว ธมฺเม อปฺปหาย อภพฺโพ อชฺฌตฺตํ กาเย ฯเปฯ พหิทฺธา กาเย ฯเปฯ อชฺฌตฺตพหิทฺธา กาเย ฯเปฯ อชฺฌตฺตํ เวทนาสุ ฯเปฯ พหิทฺธา เวทนาสุ ฯเปฯ อชฺฌตฺตพหิทฺธา เวทนาสุ ฯเปฯ อชฺฌตฺตํ จิตฺเต ฯเปฯ พหิทฺธา จิตฺเต ฯเปฯ อชฺฌตฺตพหิทฺธา จิตฺเต ฯเปฯ อชฺฌตฺตํ ธมฺเมสุ ฯเปฯ พหิทฺธา ธมฺเมสุ ฯเปฯ อชฺฌตฺตพหิทฺธา ธมฺเมสุ ธมฺมานุปสฺสี วิหริตุํ.\csromanspacer{} กตเม ฉ? กมฺมารามตํ ภสฺสารามตํ นิทฺทารามตํ สงฺคณิการามตํ อินฺทฺริเยสุ}

\vspace{\tipitakaparskip}
\csromancenter{ฉกฺกนิปาตปาฬิ อคุตฺตทฺวารตํ โภชเน อมตฺตญฺญุตํ.\csromanspacer{} อิเม โข ภิกฺขเว ฉ ธมฺเม ปหาย ภพฺโพ อชฺฌตฺตพหิทฺธา ธมฺเมสุ ธมฺมานุปสฺสี วิหริตุนฺติ.\csromanspacer{}\csromanspacer{}ทุติยํ.}
\csromanheader{2}{3. ตปุสฺสสุตฺต}
\csromanmatikaanchor{mtk.211}
\csromantocmarkref{subsection}{3-23. ตปุสฺสสุตฺตาทิ}{mtk.211}
\bookmark[dest={mtk.211},level=2]{3-23. ตปุสฺสสุตฺตาทิ}
\proseitem{119}{\csromanfolio{392}ฉหิ ภิกฺขเว ธมฺเมหิ สมนฺนาคโต ตปุสฺโส\footnote{ตปสฺโส (อิ) อํ 1. 27 ปิฏฺเฐปิ.} คหปติ ตถาคเต นิฏฺฐงฺคโต อมตทฺทโส อมตํ สจฺฉิกตฺวา อิริยติ.\csromanspacer{} กตเมหิ ฉหิ? พุทฺเธ อเวจฺจปฺปสาเทน ธมฺเม อเวจฺจปฺปสาเทน สํเฆ อเวจฺจปฺปสาเทน อริเยน สีเลน อริเยน ญาเณน อริยาย วิมุตฺติยา.\csromanspacer{} อิเมหิ โข ภิกฺขเว ฉหิ ธมฺเมหิ สมนฺนาคโต ตปุสฺโส คหปติ ตถาคเต นิฏฺฐงฺคโต อมตทฺทโส อมตํ สจฺฉิกตฺวา อิริยตีติ.\csromanspacer{}\csromanspacer{}ตติยํ.}

\csromanheader{2}{4-23. ภลฺลิกาทิสุตฺต}
\proseitem{120-139}{ฉหิ ภิกฺขเว ธมฺเมหิ สมนฺนาคโต ภลฺลิโก คหปติ ฯเปฯ.\csromanspacer{} สุทตฺโต คหปติ อนาถปิณฺฑิโก.\csromanspacer{} จิตฺโต คหปติ มจฺฉิกาสณฺฑิโก.\csromanspacer{} หตฺถโก อาฬวโก.\csromanspacer{} มหานาโม สกฺโก.\csromanspacer{} อุคฺโค คหปติ เวสาลิโก.\csromanspacer{} อุคฺคโต คหปติ.\csromanspacer{} สูรมฺพฏฺโฐ\footnote{สูโร อมฺพฏฺโฐ (ก)}.\csromanspacer{} ชีวโก โกมารภจฺโจ.\csromanspacer{} นกุลปิตา คหปติ.\csromanspacer{} ตวกณฺณิโก คหปติ.\csromanspacer{} ปูรโณ คหปติ.\csromanspacer{} อิสิทตฺโต คหปติ.\csromanspacer{} สนฺธาโน\footnote{สนฺตาโน (ก) 4. วิชโย (สี, สฺยา, อิ)} คหปติ.\csromanspacer{} วิจโย4 คหปติ.\csromanspacer{} วิชยมาหิโก\footnote{วชฺชิยมหิโต (สี, สฺยา, อิ)} คหปติ.\csromanspacer{} เมณฺฑโก คหปติ.\csromanspacer{} วาเสฏฺโฐ อุปาสโก.\csromanspacer{} อริฏฺโฐ อุปาสโก.\csromanspacer{} สารคฺโค\footnote{สาทตฺโต (สฺยา)} อุปาสโก ตถาคเต นิฏฺฐงฺคโต อมตทฺทโส อมตํ สจฺฉิกตฺวา อิริยติ.\csromanspacer{} กตเมหิ ฉหิ? พุทฺเธ อเวจฺจปฺปสาเทน ธมฺเม อเวจฺจปฺปสาเทน สํเฆ อเวจฺจปฺปสาเทน อริเยน สีเลน อริเยน ญาเณน อริยาย วิมุตฺติยา.\csromanspacer{} อิเมหิ โข ภิกฺขเว ฉหิ ธมฺเมหิ สมนฺนาคโต สารคฺโค อุปาสโก ตถาคเต นิฏฺฐงฺคโต อมตทฺทโส อมตํ สจฺฉิกตฺวา อิริยตีติ.\csromanspacer{}\csromanspacer{}เตวีสติมํ.\csromanspacer{} สามญฺญวคฺโค ทฺวาทสโม.}

\csromanmatikaanchor{mtk.212}
\csromanmatikaanchor{mtk.280}
\csromantocmarkref{subsection}{24. ราคเปยฺยาล}{mtk.212}
\bookmark[dest={mtk.212},level=2]{24. ราคเปยฺยาล}
\csromanheader{2}{24. ราคเปยฺยาล \textbf{24. ราคเปยฺยาล}}
\proseitem{140}{\csromanfolio{393}ราคสฺส ภิกฺขเว อภิญฺญาย ฉ ธมฺมา ภาเวตพฺพา.\csromanspacer{} กตเม ฉ? ทสฺสนานุตฺตริยํ สวนานุตฺตริยํ ลาภานุตฺตริยํ สิกฺขานุตฺตริยํ ปาริจริยานุตฺตริยํ อนุสฺสตานุตฺตริยํ.\csromanspacer{} ราคสฺส ภิกฺขเว อภิญฺญาย อิเม ฉ ธมฺมา ภาเวตพฺพาติ. (1)}

\proseitem{141}{ราคสฺส ภิกฺขเว อภิญฺญาย ฉ ธมฺมา ภาเวตพฺพา.\csromanspacer{} กตเม ฉ? พุทฺธานุสฺสติ ธมฺมานุสฺสติ สํฆานุสฺสติ สีลานุสฺสติ จาคานุสฺสติ เทวตานุสฺสติ.\csromanspacer{} ราคสฺส ภิกฺขเว อภิญฺญาย อิเม ฉ ธมฺมา ภาเวตพฺพาติ. (2)}

\proseitem{142}{ราคสฺส ภิกฺขเว อภิญฺญาย ฉ ธมฺมา ภาเวตพฺพา.\csromanspacer{} กตเม ฉ? อนิจฺจสญฺญา อนิจฺเจ ทุกฺขสญฺญา ทุกฺเข อนตฺตสญฺญา ปหานสญฺญา วิราคสญฺญา นิโรธสญฺญา.\csromanspacer{} ราคสฺส ภิกฺขเว อภิญฺญาย อิเม ฉ ธมฺมา ภาเวตพฺพาติ. (3)}

\proseitem{143-169}{ราคสฺส ภิกฺขเว ปริญฺญาย ฯเปฯ ปริกฺขยาย.\csromanspacer{} ปหานาย.\csromanspacer{} ขยาย.\csromanspacer{} วยาย.\csromanspacer{} วิราคาย.\csromanspacer{} นิโรธาย.\csromanspacer{} จาคาย.\csromanspacer{} ปฏินิสฺสคฺคาย ฉ ธมฺมา ภาเวตพฺพา. (4-30)}

\proseitem{170-649}{โทสสฺส ฯเปฯ.\csromanspacer{} โมหสฺส.\csromanspacer{} โกธสฺส.\csromanspacer{} อุปนาหสฺส.\csromanspacer{} มกฺขสฺส.\csromanspacer{} ปฬาสสฺส.\csromanspacer{} อิสฺสาย.\csromanspacer{} มจฺฉริยสฺส.\csromanspacer{} มายาย.\csromanspacer{} สาเฐยฺยสฺส.\csromanspacer{} ถมฺภสฺส.\csromanspacer{} สารมฺภสฺส.\csromanspacer{} มานสฺส.\csromanspacer{} อติมานสฺส.\csromanspacer{} มทสฺส.\csromanspacer{} ปมาทสฺส อภิญฺญาย ฯเปฯ ปริญฺญาย.\csromanspacer{} ปริกฺขยาย.\csromanspacer{} ปหานาย.\csromanspacer{} ขยาย.\csromanspacer{} วยาย.\csromanspacer{} วิราคาย.\csromanspacer{} นิโรธาย.\csromanspacer{} จาคาย ปฏินิสฺสคฺคาย ฯเปฯ อิเม ฉ ธมฺมา ภาเวตพฺพาติ.\csromanspacer{} อิทมโวจ ภควา.\csromanspacer{} อตฺตมนา เต ภิกฺขู ภควโต ภาสิตํ อภินนฺทุนฺติ. (31-510) ราคเปยฺยาลํ นิฏฺฐิตํ.}

\nitthitam{\textbf{ฉกฺกนิปาตปาฬิ นิฏฺฐิตา.}}
\csromanmatikaanchor{mtk.213}
\csromantocmarknopage{chapter}{สตฺตกนิปาตปาฬิ}{mtk.213}
\bookmark[dest={mtk.213},level=0]{สตฺตกนิปาตปาฬิ}
\gambhira{\textbf{องฺคุตฺตรนิกาย} \textbf{สตฺตกนิปาตปาฬิ}}
\namakkaram{นโม ตสฺส ภควโต อรหโต สมฺมาสมฺพุทฺธสฺส.}
\csromanmatikaanchor{mtk.214}
\csromantocmarknopage{section}{ปณฺณาสก}{mtk.214}
\bookmark[dest={mtk.214},level=1]{ปณฺณาสก}
\csromanheader{1}{\textbf{ปณฺณาสก}}
\csromanmatikaanchor{mtk.215}
\csromantocmarknopage{chapter}{1. ธนวคฺค}{mtk.215}
\bookmark[dest={mtk.215},level=0]{1. ธนวคฺค}
\chapterhead{1. ธนวคฺค}
\csromanheader{2}{1. ปฐมปิยสุตฺต}
\csromanmatikaanchor{mtk.216}
\csromantocmarkref{subsection}{1-2. ปิยสุตฺตทฺวย}{mtk.216}
\bookmark[dest={mtk.216},level=2]{1-2. ปิยสุตฺตทฺวย}
\proseitem{1}{\csromanfolio{395}เอวํ เม สุตํ\csromandash{}เอกํ สมยํ ภควา สาวตฺถิยํ วิหรติ เชตวเน อนาถปิณฺฑิกสฺส อาราเม.\csromanspacer{} ตตฺร โข ภควา ภิกฺขู อามนฺเตสิ “ภิกฺขโว”ติ. “ภทนฺเต”ติ เต ภิกฺขู ภควโต ปจฺจสฺโสสุํ ภควา เอตทโวจ\csromandash{}}

\prose{สตฺตหิ ภิกฺขเว ธมฺเมหิ สมนฺนาคโต ภิกฺขุ สพฺรหฺมจารีนํ อปฺปิโย จ โหติ อมนาโป จ อครุ จ อภาวนีโย จ.\csromanspacer{} กตเมหิ สตฺตหิ? อิธ ภิกฺขเว ภิกฺขุ ลาภกาโม จ โหติ, สกฺการกาโม จ โหติ, อนวญฺญตฺติกาโม จ โหติ, อหิริโก จ โหติ, อโนตฺตปฺปี จ ปาปิจฺโฉ จ มิจฺฉาทิฏฺฐิ จ.\csromanspacer{} อิเมหิ โข ภิกฺขเว สตฺตหิ ธมฺเมหิ สมนฺนาคโต ภิกฺขุ สพฺรหฺมจารีนํ อปฺปิโย จ โหติ อมนาโป จ อครุ จ อภาวนีโย จ.}

\prose{สตฺตหิ ภิกฺขเว ธมฺเมหิ สมนฺนาคโต ภิกฺขุ สพฺรหฺมจารีนํ ปิโย จ โหติ มนาโป จ ครุ จ ภาวนีโย จ.\csromanspacer{} กตเมหิ สตฺตหิ? อิธ ภิกฺขเว ภิกฺขุ น ลาภกาโม จ โหติ, น สกฺการกาโม จ โหติ, น อนวญฺญตฺติกาโม จ โหติ, หิรีมา จ โหติ, โอตฺตปฺปี จ}

\vspace{\tipitakaparskip}
\csromancenter{สตฺตกนิปาตปาฬิ อปฺปิจฺโฉ จ สมฺมาทิฏฺฐิ จ.\csromanspacer{} อิเมหิ โข ภิกฺขเว สตฺตหิ ธมฺเมหิ สมนฺนาคโต ภิกฺขุ สพฺรหฺมจารีนํ ปิโย จ โหติ มนาโป จ ครุ จ ภาวนีโย จาติ.\csromanspacer{}\csromanspacer{}ปฐมํ.}
\csromanheader{2}{2. ทุติยปิยสุตฺต}
\csromanmatikaanchor{mtk.217}
\csromantocmarkref{subsection}{3-4. พลสุตฺตทฺวย}{mtk.217}
\bookmark[dest={mtk.217},level=2]{3-4. พลสุตฺตทฺวย}
\proseitem{2}{\csromanfolio{396}สตฺตหิ ภิกฺขเว ธมฺเมหิ สมนฺนาคโต ภิกฺขุ สพฺรหฺมจารีนํ อปฺปิโย จ โหติ อมนาโป จ อครุ จ อภาวนีโย จ.\csromanspacer{} กตเมหิ สตฺตหิ? อิธ ภิกฺขเว ภิกฺขุ ลาภกาโม จ โหติ, สกฺการกาโม จ โหติ, อนวญฺญตฺติกาโม จ โหติ, อหิริโก จ โหติ อโนตฺตปฺปี จ อิสฺสุกี จ มจฺฉรี จ.\csromanspacer{} อิเมหิ โข ภิกฺขเว สตฺตหิ ธมฺเมหิ สมนฺนาคโต ภิกฺขุ สพฺรหฺมจารีนํ อปฺปิโย จ โหติ อมนาโป จ อครุ จ อภาวนีโย จ.}

\prose{สตฺตหิ ภิกฺขเว ธมฺเมหิ สมนฺนาคโต ภิกฺขุ สพฺรหฺมจารีนํ ปิโย จ โหติ มนาโป จ ครุ จ ภาวนีโย จ.\csromanspacer{} กตเมหิ สตฺตหิ? อิธ ภิกฺขเว ภิกฺขุ น ลาภกาโม จ โหติ, น สกฺการกาโม จ โหติ, น อนวญฺญตฺติกาโม จ โหติ, หิรีมา จ โหติ โอตฺตปฺปี จ อนิสฺสุกี จ อมจฺฉรี จ.\csromanspacer{} อิเมหิ โข ภิกฺขเว สตฺตหิ ธมฺเมหิ สมนฺนาคโต ภิกฺขุ สพฺรหฺมจารีนํ ปิโย จ โหติ มนาโป จ ครุ จ ภาวนีโย จาติ.\csromanspacer{} ทุติยํ.}

\csromanheader{2}{3. สํขิตฺตพลสุตฺต}
\proseitem{3}{เอวํ เม สุตํ\csromandash{}เอกํ สมยํ ภควา สาวตฺถิยํ วิหรติ เชตวเน อนาถปิณฺฑิกสฺส อาราเม ฯเปฯ.\csromanspacer{} สตฺติมานิ ภิกฺขเว พลานิ.\csromanspacer{} กตมานิ สตฺต? สทฺธาพลํ วีริยพลํ หิรีพลํ โอตฺตปฺปพลํ สติพลํ สมาธิพลํ ปญฺญาพลํ.\csromanspacer{} อิมานิ โข ภิกฺขเว สตฺต พลานีติ.}

\csromangathameasure{สทฺธาพลํ วีริยญฺจ, หิรี โอตฺตปฺปิยํ พลํ. \\ สติพลํ สมาธิ จ, ปญฺญา เว สตฺตมํ พลํ. \\ เอเตหิ พลวา ภิกฺขุ, สุขํ ชีวติ ปณฺฑิโต.}
\csromangathagroup{\csromangathastanza{สทฺธาพลํ วีริยญฺจ, หิรี\footnote{หิริ (สี, อิ, ก)} โอตฺตปฺปิยํ พลํ. \\ สติพลํ สมาธิ จ, ปญฺญา เว สตฺตมํ พลํ. \\ เอเตหิ พลวา ภิกฺขุ, สุขํ ชีวติ ปณฺฑิโต.}}

\chapterheadpageread{1. ธนวคฺค}
\csromangathameasure{โยนิโส วิจิเน ธมฺมํ, ปญฺญายตฺถํ วิปสฺสติ. \\ ปชฺโชตสฺเสว นิพฺพานํ, วิโมกฺโข โหติ เจตโสติ. ตติยํ.}
\csromangathagroup{\csromangathastanza{\csromanfolio{397}โยนิโส วิจิเน ธมฺมํ, ปญฺญายตฺถํ วิปสฺสติ. \\ ปชฺโชตสฺเสว นิพฺพานํ, วิโมกฺโข โหติ เจตโสติ.\csromanspacer{} ตติยํ.}}

\csromanheader{2}{4. วิตฺถตพลสุตฺต}
\proseitem{4}{สตฺติมานิ ภิกฺขเว พลานิ.\csromanspacer{} กตมานิ สตฺต? สทฺธาพลํ วีริยพลํ หิรีพลํ โอตฺตปฺปพลํ สติพลํ สมาธิพลํ ปญฺญาพลํ.}

\prose{กตมญฺจ ภิกฺขเว สทฺธาพลํ? อิธ ภิกฺขเว อริยสาวโก สทฺโธ โหติ, สทฺทหติ ตถาคตสฺส โพธิํ “อิติปิ โส ภควา อรหํ สมฺมาสมฺพุทฺโธ ฯเปฯ สตฺถา เทวมนุสฺสานํ พุทฺโธ ภควา”ติ.\csromanspacer{} อิทํ วุจฺจติ ภิกฺขเว สทฺธาพลํ.}

\prose{กตมญฺจ ภิกฺขเว วีริยพลํ? อิธ ภิกฺขเว อริยสาวโก อารทฺธวีริโย วิหรติ อกุสลานํ ธมฺมานํ ปหานาย กุสลานํ ธมฺมานํ อุปสมฺปทาย, ถามวา ทฬฺหปรกฺกโม อนิกฺขิตฺตธุโร กุสเลสุ ธมฺเมสุ.\csromanspacer{} อิทํ วุจฺจติ ภิกฺขเว วีริยพลํ.}

\prose{กตมญฺจ ภิกฺขเว หิรีพลํ? อิธ ภิกฺขเว อริยสาวโก หิรีมา โหติ, หิรียติ กายทุจฺจริเตน วจีทุจฺจริเตน มโนทุจฺจริเตน.\csromanspacer{} หิรียติ ปาปกานํ อกุสลานํ ธมฺมานํ สมาปตฺติยา.\csromanspacer{} อิทํ วุจฺจติ ภิกฺขเว หิรีพลํ.}

\prose{กตมญฺจ ภิกฺขเว โอตฺตปฺปพลํ? อิธ ภิกฺขเว อริยสาวโก โอตฺตปฺปี โหติ, โอตฺตปฺปติ กายทุจฺจริเตน วจีทุจฺจริเตน มโนทุจฺจริเตน, โอตฺตปฺปติ ปาปกานํ อกุสลานํ ธมฺมานํ สมาปตฺติยา.\csromanspacer{} อิทํ วุจฺจติ ภิกฺขเว โอตฺตปฺปพลํ.}

\prose{กตมญฺจ ภิกฺขเว สติพลํ? อิธ ภิกฺขเว อริยสาวโก สติมา โหติ, ปรเมน สติเนปกฺเกน สมานฺนาคโต จิรกตมฺปิ จิรภาสิตมฺปิ สริตา อนุสฺสริตา.\csromanspacer{} อิทํ วุจฺจติ ภิกฺขเว สติพลํ.}

\prose{กตมญฺจ ภิกฺขเว สมาธิพลํ? อิธ ภิกฺขเว อริยสาวโก วิวิจฺเจว กาเมหิ ฯเปฯ จตุตฺถํ ฌานํ อุปสมฺปชฺช วิหรติ.\csromanspacer{} อิทํ วุจฺจติ ภิกฺขเว สมาธิพลํ.}

\gambhira{สตฺตกนิปาตปาฬิ}
\csromanmatikaanchor{mtk.218}
\csromantocmarkref{subsection}{5-6. ธนสุตฺตทฺวย}{mtk.218}
\bookmark[dest={mtk.218},level=2]{5-6. ธนสุตฺตทฺวย}
\prose{\csromanfolio{398}กตมญฺจ ภิกฺขเว ปญฺญาพลํ, อิธ ภิกฺขเว อริยสาวโก ปญฺญวา โหติ, อุทยตฺถคามินิยา ปญฺญาย สมนฺนาคโต อริยาย นิพฺเพธิกาย สมฺมา ทุกฺขกฺขยคามินิยา.\csromanspacer{} อิทํ วุจฺจติ ภิกฺขเว ปญฺญาพลํ.\csromanspacer{} อิมานิ โข ภิกฺขเว สตฺต พลานีติ.}

\csromangathameasure{สทฺธาพลํ วีริยญฺจ, หิรี โอตฺตปฺปิยํ พลํ. \\ สติพลํ สมาธิ จ, ปญฺญา เว สตฺตมํ พลํ. \\ เอเตหิ พลวา ภิกฺขุ, สุขํ ชีวติ ปณฺฑิโต. \\ โยนิโส วิจิเน ธมฺมํ, ปญฺญายตฺถํ วิปสฺสติ. \\ ปชฺโชตสฺเสว นิพฺพานํ, วิโมกฺโข โหติ เจตโสติ. จตุตฺถํ.}
\csromangathagroup{\csromangathastanza{สทฺธาพลํ วีริยญฺจ, หิรี โอตฺตปฺปิยํ พลํ. \\ สติพลํ สมาธิ จ, ปญฺญา เว สตฺตมํ พลํ.}\csromangathastanza{เอเตหิ พลวา ภิกฺขุ, สุขํ ชีวติ ปณฺฑิโต. \\ โยนิโส วิจิเน ธมฺมํ, ปญฺญายตฺถํ วิปสฺสติ. \\ ปชฺโชตสฺเสว นิพฺพานํ, วิโมกฺโข โหติ เจตโสติ.\csromanspacer{} จตุตฺถํ.}}

\csromanheader{2}{5. สํขิตฺตธนสุตฺต}
\proseitem{5}{สตฺติมานิ ภิกฺขเว ธนานิ.\csromanspacer{} กตมานิ สตฺต, สทฺธาธนํ สีลธนํ หิรีธนํ โอตฺตปฺปธนํ สุตธนํ จาคธนํ ปญฺญาธนํ.\csromanspacer{} อิมานิ โข ภิกฺขเว สตฺต ธนานีติ.}

\csromangathameasure{สทฺธาธนํ สีลธนํ, หิรี โอตฺตปฺปิยํ ธนํ. \\ สุตธนญฺจ จาโค จ, ปญฺญา เว สตฺตมํ ธนํ. \\ ยสฺส เอเต ธนา อตฺถิ, อิตฺถิยา ปุริสสฺส วา. \\ อทลิทฺโทติ ตํ อาหุ, อโมฆํ ตสฺส ชีวิตํ. \\ ตสฺมา สทฺธญฺจ สีลญฺจ, ปสาทํ ธมฺมทสฺสนํ. \\ อนุยุญฺเชถ เมธาวี, สรํ พุทฺธาน สาสนนฺติ.ปญฺจมํ.}
\csromangathagroup{\csromangathastanza{สทฺธาธนํ สีลธนํ, หิรี โอตฺตปฺปิยํ ธนํ. \\ สุตธนญฺจ จาโค จ, ปญฺญา เว สตฺตมํ ธนํ.}\csromangathastanza{ยสฺส เอเต ธนา อตฺถิ, อิตฺถิยา ปุริสสฺส วา. \\ อทลิทฺโทติ ตํ อาหุ, อโมฆํ ตสฺส ชีวิตํ.}\csromangathastanza{ตสฺมา สทฺธญฺจ สีลญฺจ, ปสาทํ ธมฺมทสฺสนํ. \\ อนุยุญฺเชถ เมธาวี, สรํ\footnote{วรํ (ก)} พุทฺธาน สาสนนฺติ.\csromanspacer{}\csromanspacer{}ปญฺจมํ.}}

\csromanheader{2}{6. วิตฺถตธนสุตฺต}
\proseitem{6}{สตฺติมานิ ภิกฺขเว ธนานิ.\csromanspacer{} กตมานิ สตฺต, สทฺธาธนํ สีลธนํ หิรีธนํ โอตฺตปฺปธนํ สุตธนํ จาคธนํ ปญฺญาธนํ.}

\prose{กตมญฺจ ภิกฺขเว สทฺธาธนํ, อิธ ภิกฺขเว อริยสาวโก สทฺโธ โหติ, สทฺทหติ ตถาคตสฺส โพธิํ “อิติปิ โส ภควา อรหํ สมฺมาสมฺพุทฺโธ ฯเปฯ พุทฺโธ ภควา”ติ.\csromanspacer{} อิทํ วุจฺจติ ภิกฺขเว สทฺธาธนํ.}

\chapterheadpageread{1. ธนวคฺค}
\prose{\csromanfolio{399}กตมญฺจ ภิกฺขเว สีลธนํ? อิธ ภิกฺขเว อริยสาวโก ปาณาติปาตา ปฏิวิรโต โหติ ฯเปฯ สุราเมรยมชฺชปมาทฏฺฐานา ปฏิวิรโต โหติ.\csromanspacer{} อิทํ วุจฺจติ ภิกฺขเว สีลธนํ.}

\prose{กตมญฺจ ภิกฺขเว หิรีธนํ? อิธ ภิกฺขเว อริยสาวโก หิรีมา โหติ, หิรียติ กายทุจฺจริเตน วจีทุจฺจริเตน มโนทุจฺจริเตน, หิรียติ ปาปกานํ อกุสลานํ ธมฺมานํ สมาปตฺติยา.\csromanspacer{} อิทํ วุจฺจติ ภิกฺขเว หิรีธนํ.}

\prose{กตมญฺจ ภิกฺขเว โอตฺตปฺปธนํ? อิธ ภิกฺขเว อริยสาวโก โอตฺตปฺปี โหติ, โอตฺตปฺปติ กายทุจฺจริเตน วจีทุจฺจริเตน มโนทุจฺจริเตน, โอตฺตปฺปติ ปาปกานํ อกุสลานํ ธมฺมานํ สมาปตฺติยา.\csromanspacer{} อิทํ วุจฺจติ ภิกฺขเว โอตฺตปฺปธนํ.}

\prose{กตมญฺจ ภิกฺขเว สุตธนํ? อิธ ภิกฺขเว อริยสาวโก พหุสฺสุโต โหติ สุตธโร สุตสนฺนิจโย, เย เต ธมฺมา อาทิกลฺยาณา มชฺเฌกลฺยาณา ปริโยสานกลฺยาณา สาตฺถํ สพฺยญฺชนํ เกวลปริปุณฺณํ ปริสุทฺธํ พฺรหฺมจริยํ อภิวทนฺติ, ตถารูปาสฺส ธมฺมา พหุสฺสุตา โหนฺติ ธาตา วจสา ปริจิตา มนสานุเปกฺขิตา ทิฏฺฐิยา สุปฺปฏิวิทฺธา.\csromanspacer{} อิทํ วุจฺจติ ภิกฺขเว สุตธนํ.}

\prose{กตมญฺจ ภิกฺขเว จาคธนํ? อิธ ภิกฺขเว อริยสาวโก วิคตมลมจฺเฉเรน เจตสา อคารํ อชฺฌาวสติ มุตฺตจาโค ปยตปาณิ โวสคฺครโต ยาจโยโค ทานสํวิภาครโต.\csromanspacer{} อิทํ วุจฺจติ ภิกฺขเว จาคธนํ.}

\prose{กตมญฺจ ภิกฺขเว ปญฺญาธนํ? อิธ ภิกฺขเว อริยสาวโก ปญฺญวา โหติ ฯเปฯ สมฺมา ทุกฺขกฺขยคามินิยา.\csromanspacer{} อิทํ วุจฺจติ ภิกฺขเว ปญฺญาธนํ.\csromanspacer{} อิมานิ โข ภิกฺขเว สตฺตธนานีติ.}

\csromangathameasure{สทฺธาธนํ สีลธนํ, หิรี โอตฺตปฺปิยํ ธนํ. \\ สุตธนญฺจ จาโค จ, ปญฺญา เว สตฺตมํ ธนํ. \\ ยสฺส เอเต ธนา อตฺถิ, อิตฺถิยา ปุริสสฺส วา. \\ “อทลิทฺโท”ติ ตํ อาหุ, อโมฆํ ตสฺส ชีวิตํ. \\ ตสฺมา สทฺธญฺจ สีลญฺจ, ปสาทํ ธมฺมทสฺสนํ. \\ อนุยุญฺเชถ เมธาวี, สรํ พุทฺธาน สาสนนฺติ.ฉฏฺฐํ.}
\csromangathagroup{\csromangathastanza{สทฺธาธนํ สีลธนํ, หิรี โอตฺตปฺปิยํ ธนํ. \\ สุตธนญฺจ จาโค จ, ปญฺญา เว สตฺตมํ ธนํ.}\csromangathastanza{ยสฺส เอเต ธนา อตฺถิ, อิตฺถิยา ปุริสสฺส วา. \\ “อทลิทฺโท”ติ ตํ อาหุ, อโมฆํ ตสฺส ชีวิตํ.}\csromangathastanza{ตสฺมา สทฺธญฺจ สีลญฺจ, ปสาทํ ธมฺมทสฺสนํ. \\ อนุยุญฺเชถ เมธาวี, สรํ พุทฺธาน สาสนนฺติ.\csromanspacer{}\csromanspacer{}ฉฏฺฐํ.}}

\csromanmatikaanchor{mtk.219}
\csromantocmarkref{subsection}{7. อุคฺคสุตฺต}{mtk.219}
\bookmark[dest={mtk.219},level=2]{7. อุคฺคสุตฺต}
\csromanheader{2}{สตฺตกนิปาตปาฬิ \textbf{7. อุคฺคสุตฺต}}
\csromanmatikaanchor{mtk.220}
\csromantocmarkref{subsection}{8-10. สํโยชนสุตฺตาทิ}{mtk.220}
\bookmark[dest={mtk.220},level=2]{8-10. สํโยชนสุตฺตาทิ}
\proseitem{7}{\csromanfolio{400}อถ โข อุคฺโค ราชมหามตฺโต เยน ภควา เตนุปสงฺกมิ, อุปสงฺกมิตฺวา ภควนฺตํ อภิวาเทตฺวา เอกมนฺตํ นิสีทิ, เอกมนฺตํ นิสินฺโน โข อุคฺโค ราชมหามตฺโต ภควนฺตํ เอตทโวจ\csromandash{}}

\prose{อจฺฉริยํ ภนฺเต, อพฺภุตํ ภนฺเต, ยาว อฑฺโฒ จายํ ภนฺเต มิคาโร โรหเณยฺโย ยาว มหทฺธโน ยาว มหาโภโคติ.\csromanspacer{} กีว อฑฺโฒ ปนุคฺค มิคาโร โรหเณยฺโย กีว มหทฺธโน กีว มหาโภโคติ.\csromanspacer{} สตํ ภนฺเต สตสหสฺสานํ\footnote{สหสฺสานํ (สี), สหสฺสานิ (สฺยา), สตสหสฺสานิ (?)} หิรญฺญสฺส, โก ปน วาโท รูปิยสฺสาติ.\csromanspacer{} อตฺถิ โข เอตํ อุคฺค ธนํ, เนตํ นตฺถีติ วทามีติ, ตญฺจ โข เอตํ อุคฺค ธนํ, สาธารณํ อคฺคินา อุทเกน ราชูหิ โจเรหิ อปฺปิเยหิ ทายาเทหิ.\csromanspacer{} สตฺต โข อิมานิ อุคฺค ธนานิ อสาธารณานิ อคฺคินา อุทเกน ราชูหิ โจเรหิ อปฺปิเยหิ ทายาเทหิ.\csromanspacer{} กตมานิ สตฺต, สทฺธาธนํ สีลธนํ หิรีธนํ โอตฺตปฺปธนํ สุตธนํ จาคธนํ ปญฺญาธนํ.\csromanspacer{} อิมานิ โข อุคฺค สตฺต ธนานิ อสาธารณานิ อคฺคินา อุทเกน ราชูหิ โจเรหิ อปฺปิเยหิ ทายาเทหีติ.}

\csromangathameasure{สทฺธาธนํ สีลธนํ, หิรี โอตฺตปฺปิยํ ธนํ. \\ สุตธนญฺจ จาโค จ, ปญฺญา เว สตฺตมํ ธนํ. \\ ยสฺส เอเต ธนา อตฺถิ, อิตฺถิยา ปุริสสฺส วา. \\ ส เว มหทฺธโน โลเก, อเชยฺโย เทวมานุเส. \\ ตสฺมา สทฺธญฺจ สีลญฺจ, ปสาทํ ธมฺมทสฺสนํ. \\ อนุยุญฺเชถ เมธาวี, สรํ พุทฺธาน สาสนนฺติ.สตฺตมํ.}
\csromangathagroup{\csromangathastanza{สทฺธาธนํ สีลธนํ, หิรี โอตฺตปฺปิยํ ธนํ. \\ สุตธนญฺจ จาโค จ, ปญฺญา เว สตฺตมํ ธนํ.}\csromangathastanza{ยสฺส เอเต ธนา อตฺถิ, อิตฺถิยา ปุริสสฺส วา. \\ ส เว มหทฺธโน โลเก, อเชยฺโย เทวมานุเส.}\csromangathastanza{ตสฺมา สทฺธญฺจ สีลญฺจ, ปสาทํ ธมฺมทสฺสนํ. \\ อนุยุญฺเชถ เมธาวี, สรํ พุทฺธาน สาสนนฺติ.\csromanspacer{}\csromanspacer{}สตฺตมํ.}}

\csromanheader{2}{8. สํโยชนสุตฺต}
\proseitem{8}{สตฺติมานิ ภิกฺขเว สํโยชนานิ.\csromanspacer{} กตมานิ สตฺต, อนุนยสํโยชนํ ปฏิฆสํโยชนํ ทิฏฺฐิสํโยชนํ วิจิกิจฺฉาสํโยชนํ มานสํโยชนํ ภวราคสํโยชนํ อวิชฺชาสํโยชนํ.\csromanspacer{} อิมานิ โข ภิกฺขเว สตฺต สํโยชนานีติ.\csromanspacer{}\csromanspacer{}อฏฺฐมํ.}

\csromanmatikaanchor{mtk.221}
\csromantocmarknopage{chapter}{2. อนุสยวคฺค}{mtk.221}
\bookmark[dest={mtk.221},level=0]{2. อนุสยวคฺค}
\chapterheadpageread{2. อนุสยวคฺค \textbf{9. ปหานสุตฺต}}
\csromanmatikaanchor{mtk.222}
\csromantocmarkref{subsection}{1-2. อนุสยสุตฺตทฺวย}{mtk.222}
\bookmark[dest={mtk.222},level=2]{1-2. อนุสยสุตฺตทฺวย}
\proseitem{9}{\csromanfolio{401}สตฺตนฺนํ ภิกฺขเว สํโยชนานํ ปหานาย สมุจฺเฉทาย พฺรหฺมจริยํ วุสฺสติ.\csromanspacer{} กตเมสํ สตฺตนฺนํ? อนุนยสํโยชนสฺส ปหานาย สมุจฺเฉทาย พฺรหฺมจริยํ วุสฺสติ.\csromanspacer{} ปฏิฆสํโยชนสฺส ฯเปฯ.\csromanspacer{} ทิฏฺฐิสํโยชนสฺส.\csromanspacer{} วิจิกิจฺฉาสํโยชนสฺส.\csromanspacer{} มานสํโยชนสฺส.\csromanspacer{} ภวราคสํโยชนสฺส.\csromanspacer{} อวิชฺชาสํโยชนสฺส ปหานาย สมุจฺเฉทาย พฺรหฺมจริยํ วุสฺสติ.\csromanspacer{} อิเมสํ โข ภิกฺขเว สตฺตนฺนํ สํโยชนานํ ปหานาย สมุจฺเฉทาย พฺรหฺมจริยํ วุสฺสติ.\csromanspacer{} ยโต จ โข ภิกฺขเว ภิกฺขุโน อนุนยสํโยชนํ ปหีนํ โหติ อุจฺฉินฺนมูลํ ตาลาวตฺถุกตํ อนภาวํกตํ อายติํ อนุปฺปาทธมฺมํ.\csromanspacer{} ปฏิฆสํโยชนํ ฯเปฯ.\csromanspacer{} ทิฏฺฐิสํโยชนํ.\csromanspacer{} วิจิกิจฺฉาสํโยชนํ.\csromanspacer{} มานสํโยชนํ.\csromanspacer{} ภวราคสํโยชนํ.\csromanspacer{} อวิชฺชาสํโยชนํ ปหีนํ โหติ อุจฺฉินฺนมูลํ ตาลาวตฺถุกตํ อนภาวํกตํ อายติํ อนุปฺปาท- ธมฺมํ.\csromanspacer{} อยํ วุจฺจติ ภิกฺขเว ภิกฺขุ อจฺเฉจฺฉิ ตณฺหํ, วิวตฺตยิ สํโยชนํ, สมฺมา มานาภิสมยา อนฺตมกาสิ ทุกฺขสฺสาติ.\csromanspacer{}\csromanspacer{}นวมํ.}

\csromanheader{2}{10. มจฺฉริยสุตฺต}
\proseitem{10}{สตฺติมานิ ภิกฺขเว สํโยชนานิ.\csromanspacer{} กตมานิ สตฺต? อนุนยสํโยชนํ ปฏิฆสํโยชนํ ทิฏฺฐิสํโยชนํ วิจิกิจฺฉาสํโยชนํ มานสํโยชนํ อิสฺสาสํโยชนํ มจฺฉริยสํโยชนํ.\csromanspacer{} อิมานิ โข ภิกฺขเว สตฺต สํโยชนานีติ.\csromanspacer{}\csromanspacer{}ทสมํ.\csromanspacer{} ธนวคฺโค ปฐโม.}

\titlehead{\textbf{ตสฺสุทฺทานํ}}
\csromangathameasure{เทฺว ปิยานิ พลํ ธนํ, สํขิตฺตญฺเจว วิตฺถตํ. \\ อุคฺคํ สํโยชนญฺเจว, ปหานํ มจฺฉริเยน จาติ.}
\csromangathagroup{\csromangathastanza{เทฺว ปิยานิ พลํ ธนํ, สํขิตฺตญฺเจว วิตฺถตํ. \\ อุคฺคํ สํโยชนญฺเจว, ปหานํ มจฺฉริเยน จาติ.}}

\chapterhead{2. อนุสยวคฺค}
\csromanheader{2}{1. ปฐม อนุสยสุตฺต}
\proseitem{11}{สตฺติเม ภิกฺขเว อนุสยา.\csromanspacer{} กตเม สตฺต? กามราคานุสโย ปฏิฆานุสโย ทิฏฺฐานุสโย วิจิกิจฺฉานุสโย มานานุสโย}

\vspace{\tipitakaparskip}
\csromancenter{สตฺตกนิปาตปาฬิ ภวราคานุสโย อวิชฺชานุสโย.\csromanspacer{} อิเม โข ภิกฺขเว สตฺต อนุสยาติ.\csromanspacer{}\csromanspacer{}ปฐมํ.}
\csromanheader{2}{2. ทุติย อนุสยสุตฺต}
\csromanmatikaanchor{mtk.223}
\csromantocmarkref{subsection}{3-4. กุลสุตฺตาทิ}{mtk.223}
\bookmark[dest={mtk.223},level=2]{3-4. กุลสุตฺตาทิ}
\proseitem{12}{\csromanfolio{402}สตฺตนฺนํ ภิกฺขเว อนุสยานํ ปหานาย สมุจฺเฉทาย พฺรหฺมจริยํ วุสฺสติ.\csromanspacer{} กตเมสํ สตฺตนฺนํ? กามราคานุสยสฺส ปหานาย สมุจฺเฉทาย พฺรหฺมจริยํ วุสฺสติ.\csromanspacer{} ปฏิฆานุสยสฺส ฯเปฯ.\csromanspacer{} ทิฏฺฐานุสยสฺส.\csromanspacer{} วิจิกิจฺฉานุสยสฺส.\csromanspacer{} มานานุสยสฺส.\csromanspacer{} ภวราคานุสยสฺส.\csromanspacer{} อวิชฺชานุสยสฺส ปหานาย สมุจฺเฉทาย พฺรหฺมจริยํ วุสฺสติ.\csromanspacer{} อิเมสํ โข ภิกฺขเว สตฺตนฺนํ อนุสยานํ ปหานาย สมุจฺเฉทาย พฺรหฺมจริยํ วุสฺสติ.}

\prose{ยโต จ โข ภิกฺขเว ภิกฺขุโน กามราคานุสโย ปหีโน โหติ อุจฺฉินฺนมูโล ตาลาวตฺถุกโต อนภาวํกโต อายติํ อนุปฺปาทธมฺโม.\csromanspacer{} ปฏิฆานุสโย ฯเปฯ.\csromanspacer{} ทิฏฺฐานุสโย.\csromanspacer{} วิจิกิจฺฉานุสโย.\csromanspacer{} มานานุสโย.\csromanspacer{} ภวราคานุสโย.\csromanspacer{} อวิชฺชานุสโย ปหีโน โหติ อุจฺฉินฺนมูโล ตาลาวตฺถุกโต อนภาวํกโต อายติํ อนุปฺปาทธมฺโม.\csromanspacer{} อยํ วุจฺจติ ภิกฺขเว ภิกฺขุ อจฺเฉจฺฉิ ตณฺหํ, วิวตฺตยิ สํโยชนํ, สมฺมา มานาภิสมยา อนฺตมกาสิ ทุกฺขสฺสาติ.\csromanspacer{}\csromanspacer{}ทุติยํ.}

\csromanheader{2}{3. กุลสุตฺต}
\proseitem{13}{สตฺตหิ ภิกฺขเว องฺเคหิ สมนฺนาคตํ กุลํ อนุปคนฺตฺวา วา นาลํ อุปคนฺตุํ, อุปคนฺตฺวา วา นาลํ อุปนิสีทิตุํ.\csromanspacer{} กตเมหิ สตฺตหิ? น มนาเปน ปจฺจุฏฺเฐนฺติ, น มนาเปน อภิวาเทนฺติ, น มนาเปน อาสนํ เทนฺติ, สนฺตมสฺส ปริคุหนฺติ, พหุกมฺปิ โถกํ เทนฺติ, ปณีตมฺปิ ลูขํ เทนฺติ, อสกฺกจฺจํ เทนฺติ โน สกฺกจฺจํ.\csromanspacer{} อิเมหิ โข ภิกฺขเว สตฺตหิ องฺเคหิ สมนฺนาคตํ กุลํ อนุปคนฺตฺวา วา นาลํ อุปคนฺตุํ, อุปคนฺตฺวา วา นาลํ อุปนิสีทิตุํ.}

\prose{สตฺตหิ ภิกฺขเว องฺเคหิ สมนฺนาคตํ กุลํ อนุปคนฺตฺวา วา อลํ อุปคนฺตุํ, อุปคนฺตฺวา วา อลํ อุปนิสีทิตุํ.\csromanspacer{} กตเมหิ สตฺตหิ? มนาเปน ปจฺจุฏฺเฐนฺติ, มนาเปน อภิวาเทนฺติ, มนาเปน อาสนํ เทนฺติ, สนฺตมสฺส น ปริคุหนฺติ, พหุกมฺปิ พหุกํ เทนฺติ, ปณีตมฺปิ ปณีตํ เทนฺติ, สกฺกจฺจํ เทนฺติ โน อสกฺกจฺจํ.\csromanspacer{} อิเมหิ โข ภิกฺขเว สตฺตหิ องฺเคหิ สมนฺนาคตํ กุลํ อนุปคนฺตฺวา วา อลํ อุปคนฺตุํ, อุปคนฺตฺวา วา อลํ อุปนิสีทิตุนฺติ.\csromanspacer{}\csromanspacer{}ตติยํ.}

\chapterheadpageread{2. อนุสยวคฺค \textbf{4. ปุคฺคลสุตฺต}}
\proseitem{14}{\csromanfolio{403}สตฺติเม ภิกฺขเว ปุคฺคลา อาหุเนยฺยา ปาหุเนยฺยา ทกฺขิเณยฺยา อญฺชลิกรณียา อนุตฺตรํ ปุญฺญกฺเขตฺตํ โลกสฺส.\csromanspacer{} กตเม สตฺต? อุภโตภาควิมุตฺโต ปญฺญาวิมุตฺโต กายสกฺขี ทิฏฺฐิปฺปตฺโต\footnote{ทิฏฺฐปฺปตฺโต (ก)} สทฺธาวิมุตฺโต ธมฺมานุสารี สทฺธานุสารี.\csromanspacer{} อิเม โข ภิกฺขเว สตฺต ปุคฺคลา อาหุเนยฺยา ปาหุเนยฺยา ทกฺขิเณยฺยา อญฺชลิกรณียา อนุตฺตรํ ปุญฺญกฺเขตฺตํ โลกสฺสาติ.\csromanspacer{}\csromanspacer{}จตุตฺถํ.}

\csromanmatikaanchor{mtk.224}
\csromantocmarkref{subsection}{5. อุทกูปมาสุตฺต}{mtk.224}
\bookmark[dest={mtk.224},level=2]{5. อุทกูปมาสุตฺต}
\csromanheader{2}{5. อุทกูปมาสุตฺต}
\proseitem{15}{สตฺติเม ภิกฺขเว อุทกูปมา ปุคฺคลา สนฺโต สํวิชฺชมานา โลกสฺมิํ.\csromanspacer{} กตเม สตฺต? * อิธ ภิกฺขเว เอกจฺโจ ปุคฺคโล สกิํ นิมุคฺโค นิมุคฺโคว โหติ, อิธ ปน ภิกฺขเว เอกจฺโจ ปุคฺคโล อุมฺมุชฺชิตฺวา นิมุชฺชติ, อิธ ปน ภิกฺขเว เอกจฺโจ ปุคฺคโล อุมฺมุชฺชิตฺวา ฐิโต โหติ, อิธ ปน ภิกฺขเว เอกจฺโจ ปุคฺคโล อุมฺมุชฺชิตฺวา วิปสฺสติ วิโลเกติ, อิธ ปน ภิกฺขเว เอกจฺโจ ปุคฺคโล อุมฺมุชฺชิตฺวา ปตรติ, อิธ ปน ภิกฺขเว เอกจฺโจ ปุคฺคโล อุมฺมุชฺชิตฺวา ปติคาธปฺปตฺโต โหติ, อิธ ปน ภิกฺขเว เอกจฺโจ ปุคฺคโล อุมฺมุชฺชิตฺวา ติณฺโณ โหติ ปารงฺคโต\footnote{ปารคโต (สี, สฺยา, กํ)} ถเล ติฏฺฐติ พฺราหฺมโณ.}

\prose{กถญฺจ ภิกฺขเว ปุคฺคโล สกิํ นิมุคฺโค นิมุคฺโคว โหติ? อิธ ภิกฺขเว เอกจฺโจ ปุคฺคโล สมนฺนาคโต โหติ เอกนฺตกาฬเกหิ อกุสเลหิ ธมฺเมหิ.\csromanspacer{} เอวํ โข ภิกฺขเว ปุคฺคโล สกิํ นิมุคฺโค นิมุคฺโคว โหติ. (1)}

\prose{กถญฺจ ภิกฺขเว ปุคฺคโล อุมฺมุชฺชิตฺวา นิมุชฺชติ? อิธ ภิกฺขเว เอกจฺโจ ปุคฺคโล อุมฺมุชฺชติ, สาธุ สทฺธา กุสเลสุ ธมฺเมสุ.\csromanspacer{} สาธุ หิรี ฯเปฯ.\csromanspacer{} สาธุ โอตฺตปฺปํ.\csromanspacer{} สาธุ วีริยํ\footnote{วิริยํ (สี, สฺยา, กํ, อิ)}.\csromanspacer{} สาธุ ปญฺญา กุสเลสุ ธมฺเมสูติ.\csromanspacer{} ตสฺส สา สทฺธา เนว ติฏฺฐติ โน วฑฺฒติ หายติเยว.\csromanspacer{} ตสฺส สา หิรี ฯเปฯ.\csromanspacer{} ตสฺส ตํ โอตฺตปฺปํ.\csromanspacer{} ตสฺส ตํ วีริยํ.\csromanspacer{} ตสฺส สา ปญฺญา เนว ติฏฺฐติ โน วฑฺฒติ หายติเยว.\csromanspacer{} เอวํ โข ภิกฺขเว ปุคฺคโล อุมฺมุชฺชิตฺวา นิมุชฺชติ. (2)}

\gambhira{สตฺตกนิปาตปาฬิ}
\prose{\csromanfolio{404}กถญฺจ ภิกฺขเว ปุคฺคโล อุมฺมุชฺชิตฺวา ฐิโต โหติ? อิธ ภิกฺขเว เอกจฺโจ ปุคฺคโล อุมฺมุชฺชติ, สาธุ สทฺธา กุสเลสุ ธมฺเมสุ.\csromanspacer{} สาธุ หิรี ฯเปฯ.\csromanspacer{} สาธุ โอตฺตปฺปํ.\csromanspacer{} สาธุ วีริยํ.\csromanspacer{} สาธุ ปญฺญา กุสเลสุ ธมฺเมสูติ.\csromanspacer{} ตสฺส สา สทฺธา เนว หายติ โน วฑฺฒติ ฐิตา โหติ.\csromanspacer{} ตสฺส สา หิรี ฯเปฯ.\csromanspacer{} ตสฺส ตํ โอตฺตปฺปํ.\csromanspacer{} ตสฺส ตํ วีริยํ.\csromanspacer{} ตสฺส สา ปญฺญา เนว หายติ โน วฑฺฒติ ฐิตา โหติ.\csromanspacer{} เอวํ โข ภิกฺขเว ปุคฺคโล อุมฺมุชฺชิตฺวา ฐิโต โหติ. (3)}

\prose{กถญฺจ ภิกฺขเว ปุคฺคโล อุมฺมุชฺชิตฺวา วิปสฺสติ วิโลเกติ? อิธ ภิกฺขเว เอกจฺโจ ปุคฺคโล อุมฺมุชฺชติ, สาธุ สทฺธา กุสเลสุ ธมฺเมสุ.\csromanspacer{} สาธุ หิรี ฯเปฯ.\csromanspacer{} สาธุ โอตฺตปฺปํ.\csromanspacer{} สาธุ วีริยํ.\csromanspacer{} สาธุ ปญฺญา กุสเลสุ ธมฺเมสูติ.\csromanspacer{} โส ติณฺณํ สํโยชนานํ ปริกฺขยา โสตาปนฺโน โหติ อวินิปาตธมฺโม นิยโต สมฺโพธิปรายโณ.\csromanspacer{} เอวํ โข ภิกฺขเว ปุคฺคโล อุมฺมุชฺชิตฺวา วิปสฺสติ วิโลเกติ. (4)}

\prose{กถญฺจ ภิกฺขเว ปุคฺคโล อุมฺมุชฺชิตฺวา ปตรติ? อิธ ภิกฺขเว เอกจฺโจ ปุคฺคโล อุมฺมุชฺชติ สาธุ สทฺธา กุสเลสุ ธมฺเมสุ.\csromanspacer{} สาธุ หิรี ฯเปฯ.\csromanspacer{} สาธุ โอตฺตปฺปํ.\csromanspacer{} สาธุ วีริยํ.\csromanspacer{} สาธุ ปญฺญา กุสเลสุ ธมฺเมสูติ.\csromanspacer{} โส ติณฺณํ สํโยชนานํ ปริกฺขยา ราคโทสโมหานํ ตนุตฺตา สกทาคามี โหติ สกิเทว\footnote{สกิํเทว (ก)} อิมํ โลกํ อาคนฺตฺวา ทุกฺขสฺสนฺตํ กโรติ\footnote{ทุกฺขสฺสนฺตกโร โหติ (ก)}.\csromanspacer{} เอวํ โข ภิกฺขเว ปุคฺคโล อุมฺมุชฺชิตฺวา ปตรติ. (5)}

\prose{กถญฺจ ภิกฺขเว ปุคฺคโล อุมฺมุชฺชิตฺวา ปติคาธปฺปตฺโต โหติ? อิธ ภิกฺขเว เอกจฺโจ ปุคฺคโล อุมฺมุชฺชติ, สาธุ สทฺธา กุสเลสุ ธมฺเมสุ.\csromanspacer{} สาธุ หิรี ฯเปฯ.\csromanspacer{} สาธุ โอตฺตปฺปํ.\csromanspacer{} สาธุ วีริยํ.\csromanspacer{} สาธุ ปญฺญา กุสเลสุ ธมฺเมสูติ.\csromanspacer{} โส ปญฺจนฺนํ โอรมฺภาคิยานํ สํโยชนานํ ปริกฺขยา โอปปาติโก โหติ ตตฺถ ปรินิพฺพายี อนาวตฺติธมฺโม ตสฺมา โลกา.\csromanspacer{} เอวํ โข ภิกฺขเว ปุคฺคโล อุมฺมุชฺชิตฺวา ปติคาธปฺปตฺโต โหติ. (6)}

\prose{กถญฺจ ภิกฺขเว ปุคฺคโล อุมฺมุชฺชิตฺวา ติณฺโณ โหติ ปารงฺคโต, ถเล ติฏฺฐติ พฺราหฺมโณ? อิธ ภิกฺขเว เอกจฺโจ ปุคฺคโล อุมฺมุชฺชติ, สาธุ สทฺธา กุสเลสุ ธมฺเมสุ.\csromanspacer{} สาธุ หิรี ฯเปฯ.\csromanspacer{} สาธุ โอตฺตปฺปํ.\csromanspacer{} สาธุ วีริยํ.\csromanspacer{} สาธุ}

\vspace{\tipitakaparskip}
\csromancenter{อนุสยวคฺค ปญฺญา กุสเลสุ ธมฺเมสูติ.\csromanspacer{} โส อาสวานํ ขยา อนาสวํ เจโตวิมุตฺติํ ปญฺญาวิมุตฺติํ ทิฏฺเฐว ธมฺเม สยํ อภิญฺญา สจฺฉิกตฺวา อุปสมฺปชฺช วิหรติ.\csromanspacer{} เอวํ โข ภิกฺขเว ปุคฺคโล อุมฺมุชฺชิตฺวา ติณฺโณ โหติ ปารงฺคโต, ถเล ติฏฺฐติ พฺราหฺมโณ. (7)}
\csromanmatikaanchor{mtk.225}
\csromantocmarkref{subsection}{6-9. อนิจฺจานุปสฺสีสุตฺตาทิ}{mtk.225}
\bookmark[dest={mtk.225},level=2]{6-9. อนิจฺจานุปสฺสีสุตฺตาทิ}
\prose{\csromanfolio{405}อิเม โข ภิกฺขเว สตฺต อุทกูปมา ปุคฺคลา สนฺโต สํวิชฺชมานา โลกสฺมินฺติ.\csromanspacer{}\csromanspacer{}ปญฺจมํ.}

\csromanheader{2}{6. อนิจฺจานุปสฺสีสุตฺต}
\proseitem{16}{สตฺติเม ภิกฺขเว ปุคฺคลา อาหุเนยฺยา ปาหุเนยฺยา ทกฺขิเณยฺยา อญฺชลิกรณียา อนุตฺตรํ ปุญฺญกฺเขตฺตํ โลกสฺส.\csromanspacer{} กตเม สตฺต? อิธ ภิกฺขเว เอกจฺโจ ปุคฺคโล สพฺพสงฺขาเรสุ อนิจฺจานุปสฺสี วิหรติ อนิจฺจสญฺญี อนิจฺจปฏิสํเวที สตตํ สมิตํ อพฺโพกิณฺณํ เจตสา อธิมุจฺจมาโน ปญฺญาย ปริโยคาหมาโน.\csromanspacer{} โส อาสวานํ ขยา ฯเปฯ สจฺฉิกตฺวา อุปสมฺปชฺช วิหรติ.\csromanspacer{} อยํ ภิกฺขเว ปฐโม ปุคฺคโล อาหุเนยฺโย ปาหุเนยฺโย ทกฺขิเณยฺโย อญฺชลิกรณีโย อนุตฺตรํ ปุญฺญกฺเขตฺตํ โลกสฺส.}

\prose{ปุน จปรํ ภิกฺขเว อิเธกจฺโจ ปุคฺคโล สพฺพสงฺขาเรสุ อนิจฺจานุปสฺสี วิหรติ อนิจฺจสญฺญี อนิจฺจปฏิสํเวที สตตํ สมิตํ อพฺโพกิณฺณํ เจตสา อธิมุจฺจมาโน ปญฺญาย ปริโยคาหมาโน.\csromanspacer{} ตสฺส อปุพฺพํ อจริมํ อาสวปริยาทานญฺจ โหติ ชีวิตปริยาทานญฺจ.\csromanspacer{} อยํ ภิกฺขเว ทุติโย ปุคฺคโล อาหุเนยฺโย ฯเปฯ อนุตฺตรํ ปุญฺญกฺเขตฺตํ โลกสฺส.}

\prose{ปุน จปรํ ภิกฺขเว อิเธกจฺโจ ปุคฺคโล สพฺพสงฺขาเรสุ อนิจฺจานุปสฺสี วิหรติ อนิจฺจสญฺญี อนิจฺจปฏิสํเวที สตตํ สมิตํ อพฺโพกิณฺณํ เจตสา อธิมุจฺจมาโน ปญฺญาย ปริโยคาหมาโน.\csromanspacer{} โส ปญฺจนฺนํ โอรมฺภาคิยานํ สํโยชนานํ ปริกฺขยา อนฺตราปรินิพฺพายี โหติ ฯเปฯ อุปหจฺจปรินิพฺพายี โหติ ฯเปฯ อสงฺขารปรินิพฺพายี โหติ ฯเปฯ สสงฺขารปรินิพฺพายี โหติ ฯเปฯ อุทฺธํโสโต โหติ อกนิฏฺฐคามี.\csromanspacer{} อยํ ภิกฺขเว สตฺตโม ปุคฺคโล อาหุเนยฺโย ปาหุเนยฺโย ทกฺขิเณยฺโย อญฺชลิกรณีโย อนุตฺตรํ ปุญฺญกฺเขตฺตํ โลกสฺส.\csromanspacer{} อิเม โข ภิกฺขเว สตฺต ปุคฺคลา อาหุเนยฺยา ปาหุเนยฺยา ทกฺขิเณยฺยา อญฺชลิกรณียา อนุตฺตรํ ปุญฺญกฺเขตฺตํ โลกสฺสาติ.\csromanspacer{}\csromanspacer{}ฉฏฺฐํ.}

\gambhira{สตฺตกนิปาตปาฬิ \textbf{7. ทุกฺขานุปสฺสีสุตฺต}}
\proseitem{17}{\csromanfolio{406}สตฺติเม ภิกฺขเว ปุคฺคลา อาหุเนยฺยา ฯเปฯ อนุตฺตรํ ปุญฺญกฺเขตฺตํ โลกสฺส.\csromanspacer{} กตเม สตฺต? อิธ ภิกฺขเว เอกจฺโจ ปุคฺคโล สพฺพสงฺขาเรสุ ทุกฺขานุปสฺสี วิหรติ ฯเปฯ.\csromanspacer{}\csromanspacer{}สตฺตมํ.}

\csromanheader{2}{8. อนตฺตานุปสฺสีสุตฺต}
\vspace{\tipitakaparskip}
\csromancenter{สพฺเพสุ ธมฺเมสุ อนตฺตานุปสฺสี วิหรติ ฯเปฯ.\csromanspacer{}\csromanspacer{}อฏฺฐมํ.}
\csromanheader{2}{9. นิพฺพานสุตฺต}
\proseitem{19}{\csromansymbolfootnote{*}{อภิ 4. 295 ปิฏฺเฐปิ.}นิพฺพาเน สุขานุปสฺสี วิหรติ สุขสญฺญี สุขปฏิสํเวที สตตํ สมิตํ อพฺโพกิณฺณํ เจตสา อธิมุจฺจมาโน ปญฺญาย ปริโยคาหมาโน.\csromanspacer{} โส อาสวานํ ขยา ฯเปฯ สจฺฉิกตฺวา อุปสมฺปชฺช วิหรติ.\csromanspacer{} อยํ ภิกฺขเว ปฐโม ปุคฺคโล อาหุเนยฺโย ฯเปฯ ปุญฺญกฺเขตฺตํ โลกสฺส.}

\prose{ปุน จปรํ ภิกฺขเว อิเธกจฺโจ ปุคฺคโล นิพฺพาเน สุขานุปสฺสี วิหรติ สุขสญฺญี สุขปฏิสํเวที สตตํ สมิตํ อพฺโพกิณฺณํ เจตสา อธิมุจฺจมาโน ปญฺญาย ปริโยคาหมาโน.\csromanspacer{} ตสฺส อปุพฺพํ อจริมํ อาสวปริยาทานญฺจ โหติ ชีวิตปริยาทานญฺจ.\csromanspacer{} อยํ ภิกฺขเว ทุติโย ปุคฺคโล อาหุเนยฺโย ฯเปฯ อนุตฺตรํ ปุญฺญกฺเขตฺตํ โลกสฺส.}

\prose{ปุน จปรํ ภิกฺขเว อิเธกจฺโจ ปุคฺคโล นิพฺพาเน สุขานุปสฺสี วิหรติ สุขสญฺญี สุขปฏิสํเวที สตตํ สมิตํ อพฺโพกิณฺณํ เจตสา อธิมุจฺจมาโน ปญฺญาย ปริโยคาหมาโน.\csromanspacer{} โส ปญฺจนฺนํ โอรมฺภาคิยานํ สํโยชนานํ ปริกฺขยา อนฺตราปรินิพฺพายี โหติ ฯเปฯ อุปหจฺจปรินิพฺพายี โหติ ฯเปฯ อสงฺขารปรินิพฺพายี โหติ ฯเปฯ สสงฺขารปรินิพฺพายี โหติ ฯเปฯ อุทฺธํโสโต โหติ อกนิฏฺฐคามี.\csromanspacer{} อยํ ภิกฺขเว สตฺตโม ปุคฺคโล อาหุเนยฺโย ฯเปฯ อนุตฺตรํ ปุญฺญกฺเขตฺตํ โลกสฺส.\csromanspacer{} อิเม โข ภิกฺขเว สตฺต ปุคฺคลา อาหุเนยฺยา ฯเปฯ อนุตฺตรํ ปุญฺญกฺเขตฺตํ โลกสฺสาติ.\csromanspacer{}\csromanspacer{}นวมํ.}

\csromanmatikaanchor{mtk.226}
\csromanmatikaanchor{mtk.227}
\csromantocmarkref{subsection}{10. นิทฺทสวตฺถุสุตฺต}{mtk.226}
\bookmark[dest={mtk.226},level=2]{10. นิทฺทสวตฺถุสุตฺต}
\csromantocmarknopage{chapter}{3. วชฺชิสตฺตกวคฺค}{mtk.227}
\bookmark[dest={mtk.227},level=0]{3. วชฺชิสตฺตกวคฺค}
\csromanheader{2}{3. วชฺชิสตฺตกวคฺค \textbf{10. นิทฺทสวตฺถุสุตฺต}}
\proseitem{20}{\csromanfolio{407}สตฺติมานิ ภิกฺขเว นิทฺทสวตฺถูนิ.\csromanspacer{} กตมานิ สตฺต? อิธ ภิกฺขเว ภิกฺขุ สิกฺขาสมาทาเน ติพฺพจฺฉนฺโท โหติ, อายติํ จ สิกฺขาสมาทาเน อวิคตเปโม\footnote{อธิคตเปโม (สฺยา), อุปริ 425; ที 3. 209 ปิฏฺเฐสุปิ.}.\csromanspacer{} ธมฺมนิสนฺติยา ติพฺพจฺฉนฺโท โหติ, อายติํ จ ธมฺมนิสนฺติยา อวิคตเปโม.\csromanspacer{} อิจฺฉาวินเย ติพฺพจฺฉนฺโท โหติ, อายติํ จ อิจฺฉาวินเย อวิคตเปโม.\csromanspacer{} ปฏิสลฺลาเน ติพฺพจฺฉนฺโท โหติ, อายติํ จ ปฏิสลฺลาเน อวิคตเปโม.\csromanspacer{} วีริยารมฺเภ\footnote{วีริยารพฺเภ (ก)} ติพฺพจฺฉนฺโท โหติ, อายติํ จ วีริยารมฺเภ อวิคตเปโม.\csromanspacer{} สติเนปกฺเก ติพฺพจฺฉนฺโท โหติ, อายติํ จ สติเนปกฺเก อวิคตเปโม.\csromanspacer{} ทิฏฺฐิปฏิเวเธ ติพฺพจฺฉนฺโท โหติ, อายติํ จ ทิฏฺฐิปฏิเวเธ อวิคตเปโม.\csromanspacer{} อิมานิ โข ภิกฺขเว สตฺต นิทฺทสวตฺถูนีติ.\csromanspacer{}\csromanspacer{}ทสมํ.\csromanspacer{} อนุสยวคฺโค ทุติโย.}

\titlehead{\textbf{ตสฺสุทฺทานํ}}
\csromangathameasure{ทุเว อนุสยา กุลํ, ปุคฺคลํ อุทกูปมํ. \\ อนิจฺจํ ทุกฺขํ อนตฺตา จ, นิพฺพานํ นิทฺทสวตฺถุ จาติ.}
\csromangathagroup{\csromangathastanza{ทุเว อนุสยา กุลํ, ปุคฺคลํ อุทกูปมํ. \\ อนิจฺจํ ทุกฺขํ อนตฺตา จ, นิพฺพานํ นิทฺทสวตฺถุ จาติ.}}

\chapterhead{3. วชฺชิสตฺตกวคฺค}
\csromanmatikaanchor{mtk.228}
\csromantocmarkref{subsection}{1. สารนฺททสุตฺต}{mtk.228}
\bookmark[dest={mtk.228},level=2]{1. สารนฺททสุตฺต}
\csromanheader{2}{1. สารนฺททสุตฺต}
\proseitem{21}{เอวํ เม สุตํ\csromandash{}เอกํ สมยํ ภควา เวสาลิยํ วิหรติ สารนฺทเท เจติเย.\csromanspacer{} อถ โข สมฺพหุลา ลิจฺฉวี เยน ภควา เตนุปสงฺกมิํสุ, อุปสงฺกมิตฺวา ภควนฺตํ อภิวาเทตฺวา เอกมนฺตํ นิสีทิํสุ, เอกมนฺตํ นิสินฺเน โข เต ลิจฺฉวี ภควา เอตทโวจ “สตฺต โว ลิจฺฉวี อปริหานิเย\footnote{อปริหานีเย (ก)} ธมฺเม เทเสสฺสามิ, ตํ สุณาถ สาธุกํ มนสิ กโรถ, ภาสิสฺสามีติ. “เอวํ ภนฺเต”ติ โข เต ลิจฺฉวี ภควโต ปจฺจสฺโสสุํ.\csromanspacer{} ภควา เอตทโวจ\csromandash{}}

\gambhira{สตฺตกนิปาตปาฬิ}
\prose{\csromanfolio{408}กตเม จ ลิจฺฉวี สตฺต อปริหานิยา ธมฺมา? ยาวกีวญฺจ ลิจฺฉวี วชฺชี อภิณฺหํ สนฺนิปาตา ภวิสฺสนฺติ สนฺนิปาตพหุลา, วุทฺธิเยว ลิจฺฉวี วชฺชีนํ ปาฏิกงฺขา โน ปริหานิ. (1)}

\prose{ยาวกีวญฺจ ลิจฺฉวี วชฺชี สมคฺคา สนฺนิปติสฺสนฺติ, สมคฺคา วุฏฺฐหิสฺสนฺติ, สมคฺคา วชฺชิกรณียานิ กริสฺสนฺติ, วุทฺธิเยว ลิจฺฉวี วชฺชีนํ ปาฏิกงฺขา โน ปริหานิ. (2)}

\prose{ยาวกีวญฺจ ลิจฺฉวี วชฺชี อปญฺญตฺตํ น ปญฺญาเปสฺสนฺติ, ปญฺญตฺตํ น สมุจฺฉินฺทิสฺสนฺติ, ยถาปญฺญตฺเต โปราเณ วชฺชิธมฺเม สมาทาย วตฺติสฺสนฺติ, วุทฺธิเยว ลิจฺฉวี วชฺชีนํ ปาฏิกงฺขา โน ปริหานิ. (3)}

\prose{ยาวกีวญฺจ ลิจฺฉวี วชฺชี เย เต วชฺชีนํ วชฺชิมหลฺลกา, เต สกฺกริสฺสนฺติ, ครุํ กริสฺสนฺติ มาเนสฺสนฺติ ปูเชสฺสนฺติ, เตสญฺจ โสตพฺพํ มญฺญิสฺสนฺติ, วุทฺธิเยว ลิจฺฉวี วชฺชีนํ ปาฏิกงฺขา โน ปริหานิ. (4)}

\prose{ยาวกีวญฺจ ลิจฺฉวี วชฺชี ยา ตา กุลิตฺถิโย กุลกุมาริโย, ตา น โอกสฺส ปสยฺห วาเสสฺสนฺติ, วุทฺธิเยว ลิจฺฉวี วชฺชีนํ ปาฏิกงฺขา โน ปริหานิ. (5)}

\prose{ยาวกีวญฺจ ลิจฺฉวี วชฺชี ยานิ ตานิ วชฺชีนํ วชฺชิเจติยานิ อพฺภนฺตรานิ เจว พาหิรานิ จ, ตานิ สกฺกริสฺสนฺติ ครุํ กริสฺสนฺติ มาเนสฺสนฺติ ปูเชสฺสนฺติ, เตสญฺจ ทินฺนปุพฺพํ กตปุพฺพํ ธมฺมิกํ พลิํ โน ปริหาเปสฺสนฺติ, วุทฺธิเยว ลิจฺฉวี วชฺชีนํ ปาฏิกงฺขา โน ปริหานิ. (6)}

\prose{ยาวกีวญฺจ ลิจฺฉวี วชฺชีนํ อรหนฺเตสุ ธมฺมิกา รกฺขาวรณคุตฺติ สุสํวิหิตา ภวิสฺสติ “กินฺติ อนาคตา จ อรหนฺโต วิชิตํ อาคจฺเฉยฺยุํ, อาคตา จ อรหนฺโต วิชิเต ผาสุํ วิหเรยฺยุนฺ”ติ, วุทฺธิเยว ลิจฺฉวี วชฺชีนํ ปาฏิกงฺขา โน ปริหานิ. (7)}

\prose{ยาวกีวญฺจ ลิจฺฉวี อิเม สตฺต อปริหานิยา ธมฺมา วชฺชีสุ ฐสฺสนฺติ\footnote{วตฺติสฺสนฺติ (ก)}, อิเมสุ จ สตฺตสุ อปริหานิเยสุ ธมฺเมสุ วชฺชี สนฺทิสฺสิสฺสนฺติ\footnote{สนฺทิสฺสนฺติ (สี, อิ, ก)}.\csromanspacer{} วุทฺธิเยว ลิจฺฉวี วชฺชีนํ ปาฏิกงฺขา โน ปริหานีติ.\csromanspacer{}\csromanspacer{}ปฐมํ.}

\csromanmatikaanchor{mtk.229}
\csromantocmarkref{subsection}{2. วสฺสการสุตฺต}{mtk.229}
\bookmark[dest={mtk.229},level=2]{2. วสฺสการสุตฺต}
\csromanheader{2}{3. วชฺชิสตฺตกวคฺค \textbf{2. วสฺสการสุตฺต}}
\proseitem{22}{\csromanfolio{409}เอวํ เม สุตํ\csromandash{}เอกํ สมยํ ภควา ราชคเห วิหรติ คิชฺฌกูเฏ ปพฺพเต.\csromanspacer{} เตน โข ปน สมเยน ราชา มาคโธ อชาตสตฺตุ เวเทหิปุตฺโต วชฺชี อภิยาตุกาโม โหติ.\csromanspacer{} โส เอวมาห “อหํ หิเม วชฺชี เอวํมหิทฺธิเก เอวํมหานุภาเว อุจฺเฉจฺฉามิ\footnote{อุจฺเฉชฺชิสฺสามิ (สฺยา), อุจฺฉิชฺชิสฺสามิ (ก)}, วชฺชี วินาเสสฺสามิ, วชฺชี อนยพฺยสนํ อาปาเทสฺสามี”ติ\footnote{อาปาเทสฺสามิ วชฺชีติ (ก) ที 2. 61 ปิฏฺเฐปิ.}.}

\prose{อถ โข ราชา มาคโธ อชาตสตฺตุ เวเทหิปุตฺโต วสฺสการํ พฺราหฺมณํ มาคธมหามตฺตํ อามนฺเตสิ “เอหิ ตฺวํ พฺราหฺมณ เยน ภควา เตนุปสงฺกม, อุปสงฺกมิตฺวา มม วจเนน ภควโต ปาเท สิรสา วนฺทาหิ อปฺปาพาธํ อปฺปาตงฺกํ ลหุฏฺฐานํ พลํ ผาสุวิหารํ ปุจฺฉ ‘ราชา ภนฺเต มาคโธ อชาตสตฺตุ เวเทหิปุตฺโต ภควโต ปาเท สิรสา วนฺทติ, อปฺปาพาธํ อปฺปาตงฺกํ ลหุฏฺฐานํ พลํ ผาสุวิหารํ ปุจฺฉตี’ติ.\csromanspacer{} เอวญฺจ วเทหิ ‘ราชา ภนฺเต มาคโธ อชาตสตฺตุ เวเทหิปุตฺโต วชฺชี อภิยาตุกาโม, โส เอวมาห อหํ หิเม วชฺชี เอวํมหิทฺธิเก เอวํมหานุภาเว อุจฺเฉจฺฉามิ, วชฺชี วินาเสสฺสามิ, วชฺชี อนยพฺยสนํ อาปาเทสฺสามี’ติ.\csromanspacer{} ยถา เต ภควา พฺยากโรติ, ตํ สาธุกํ อุคฺคเหตฺวา มม อาโรเจยฺยาสิ, น หิ ตถาคตา วิตถํ ภณนฺตี”ติ.}

\prose{“เอวํ โภ”ติ โข วสฺสกาโร พฺราหฺมโณ มาคธมหามตฺโต รญฺโญ มาคธสฺส อชาตสตฺตุสฺส เวเทหิปุตฺตสฺส ปฏิสฺสุตฺวา เยน ภควา เตนุปสงฺกมิ, อุปสงฺกมิตฺวา ภควตา สทฺธิํ สมฺโมทิ, สมฺโมทนียํ กถํ สารณียํ วีติสาเรตฺวา เอกมนฺตํ นิสีทิ, เอกมนฺตํ นิสินฺโน โข วสฺสกาโร พฺราหฺมโณ มาคธมหามตฺโต ภควนฺตํ เอตทโวจ “ราชา โภ โคตม มาคโธ อชาตสตฺตุ เวเทหิปุตฺโต โภโต โคตมสฺส ปาเท สิรสา วนฺทติ, อปฺปาพาธํ อปฺปาตงฺกํ ลหุฏฺฐานํ พลํ ผาสุวิหารํ ปุจฺฉติ, ราชา\footnote{เอวญฺจ วเทติ ราชา (สี, ก)} โภ โคตม มาคโธ อชาตสตฺตุ เวเทหิปุตฺโต วชฺชี อภิยาตุกาโม, โส เอวมาห ‘อหํ หิเม วชฺชี เอวํมหิทฺธิเก เอวํมหานุภาเว อุจฺเฉจฺฉามิ, วชฺชี วินาเสสฺสามิ, วชฺชี อนยพฺยสนํ อาปาเทสฺสามี’ติ”.}

\gambhira{สตฺตกนิปาตปาฬิ}
\prose{\csromanfolio{410}เตน โข ปน สมเยน อายสฺมา อานนฺโท ภควโต ปิฏฺฐิโต ฐิโต โหติ ภควนฺตํ พีชยมาโน\footnote{วีชมาโน (สี, สฺยา)}.\csromanspacer{} อถ โข ภควา อายสฺมนฺตํ อานนฺทํ อามนฺเตสิ “กินฺติ เต อานนฺท สุตํ วชฺชี อภิณฺหํ สนฺนิปาตา สนฺนิปาตพหุลา”ติ.\csromanspacer{} สุตํ เมตํ ภนฺเต “วชฺชี อภิณฺหํ สนฺนิปาตา สนฺนิปาตพหุลา”ติ.\csromanspacer{} ยาวกีวญฺจ อานนฺท วชฺชี อภิณฺหํ สนฺนิปาตา ภวิสฺสนฺติ สนฺนิปาตพหุลา, วุทฺธิเยว อานนฺท วชฺชีนํ ปาฏิกงฺขา โน ปริหานิ. (1)}

\prose{กินฺติ เต อานนฺท สุตํ “วชฺชี สมคฺคา สนฺนิปตนฺติ, สมคฺคา วุฏฺฐหนฺติ, สมคฺคา วชฺชิกรณียานิ กโรนฺตี”ติ.\csromanspacer{} สุตํ เมตํ ภนฺเต “วชฺชี สมคฺคา สนฺนิปตนฺติ, สมคฺคา วุฏฺฐหนฺติ, สมคฺคา วชฺชิกรณียานิ กโรนฺตี”ติ.\csromanspacer{} ยาวกีวญฺจ อานนฺท วชฺชี สมคฺคา สนฺนิปติสฺสนฺติ, สมคฺคา วุฏฺฐหิสฺสนฺติ, สมคฺคา วชฺชิกรณียานิ กริสฺสนฺติ, วุทฺธิเยว อานนฺท วชฺชีนํ ปาฏิกงฺขา โน ปริหานิ. (2)}

\prose{กินฺติ เต อานนฺท สุตํ “วชฺชี อปญฺญตฺตํ น ปญฺญาเปนฺติ, ปญฺญตฺตํ น สมุจฺฉินฺทนฺติ, ยถาปญฺญตฺเต โปราเณ วชฺชิธมฺเม สมาทาย วตฺตนฺตี”ติ.\csromanspacer{} สุตํ เมตํ ภนฺเต “วชฺชี อปญฺญตฺตํ น ปญฺญาเปนฺติ, ปญฺญตฺตํ น สมุจฺฉินฺทนฺติ, ยถาปญฺญตฺเต โปราเณ วชฺชิธมฺเม สมาทาย วตฺตนฺตี”ติ.\csromanspacer{} ยาวกีวญฺจ อานนฺท วชฺชี อปญฺญตฺตํ น ปญฺญาเปสฺสนฺติ, ปญฺญตฺตํ น สมุจฺฉินฺทิสฺสนฺติ, ยถาปญฺญตฺเต โปราเณ วชฺชิธมฺเม สมาทาย วตฺติสฺสนฺติ, วุทฺธิเยว อานนฺท วชฺชีนํ ปาฏิกงฺขา โน ปริหานิ. (3)}

\prose{กินฺติ เต อานนฺท สุตํ “วชฺชี เย เต วชฺชีนํ วชฺชิมหลฺลกา, เต สกฺกโรนฺติ ครุํ กโรนฺติ มาเนนฺติ ปูเชนฺติ, เตสญฺจ โสตพฺพํ มญฺญนฺตี”ติ.\csromanspacer{} สุตํ เมตํ ภนฺเต “วชฺชี เย เต วชฺชีนํ วชฺชิมหลฺลกา, เต สกฺกโรนฺติ ครุํ กโรนฺติ มาเนนฺติ ปูเชนฺติ, เตสญฺจ โสตพฺพํ มญฺญนฺตี”ติ.\csromanspacer{} ยาวกีวญฺจ อานนฺท วชฺชี เย เต วชฺชีนํ วชฺชิมหลฺลกา, เต สกฺกริสฺสนฺติ ครุํ กริสฺสนฺติ มาเนสฺสนฺติ ปูเชสฺสนฺติ, เตสญฺจ โสตพฺพํ มญฺญิสฺสนฺติ, วุทฺธิเยว อานนฺท วชฺชีนํ ปาฏิกงฺขา โน ปริหานิ. (4)}

\prose{กินฺติ เต อานนฺท สุตํ “วชฺชี ยา ตา กุลิตฺถิโย กุลกุมาริโย, ตา น โอกสฺส ปสยฺห วาเสนฺตี”ติ.\csromanspacer{} สุตํ เมตํ ภนฺเต “วชฺชี ยา ตา}

\vspace{\tipitakaparskip}
\csromancenter{วชฺชิสตฺตกวคฺค กุลิตฺถิโย กุลกุมาริโย, ตา น โอกสฺส ปสยฺห วาเสนฺตี”ติ.\csromanspacer{} ยาวกีวญฺจ อานนฺท วชฺชี ยา ตา กุลิตฺถิโย กุลกุมาริโย, ตา น โอกสฺส ปสยฺห วาเสสฺสนฺติ, วุทฺธิเยว อานนฺท วชฺชีนํ ปาฏิกงฺขา โน ปริหานิ. (5)}
\prose{\csromanfolio{411}กินฺติ เต อานนฺท สุตํ “วชฺชี ยานิ ตานิ วชฺชีนํ วชฺชิเจติยานิ อพฺภนฺตรานิ เจว พาหิรานิ จ, ตานิ สกฺกโรนฺติ ครุํ กโรนฺติ มาเนนฺติ ปูเชนฺติ, เตสญฺจ ทินฺนปุพฺพํ กตปุพฺพํ ธมฺมิกํ พลิํ โน ปริหาเปนฺตี”ติ.\csromanspacer{} สุตํ เมตํ ภนฺเต “วชฺชี ยานิ ตานิ วชฺชีนํ วชฺชิเจติยานิ อพฺภนฺตรานิ เจว พาหิรานิ จ, ตานิ สกฺกโรนฺติ ครุํ กโรนฺติ มาเนนฺติ ปูเชนฺติ, เตสญฺจ ทินฺนปุพฺพํ กตปุพฺพํ ธมฺมิกํ พลิํ โน ปริหาเปนฺตี”ติ.\csromanspacer{} ยาวกีวญฺจ อานนฺท วชฺชี ยานิ ตานิ วชฺชีนํ วชฺชิเจติยานิ อพฺภนฺตรานิ เจว พาหิรานิ จ, ตานิ สกฺกริสฺสนฺติ ครุํ กริสฺสนฺติ มาเนสฺสนฺติ ปูเชสฺสนฺติ, เตสญฺจ ทินฺนปุพฺพํ กตปุพฺพํ ธมฺมิกํ พลิํ โน ปริหาเปสฺสนฺติ, วุทฺธิเยว อานนฺท วชฺชีนํ ปาฏิกงฺขา โน ปริหานิ. (6)}

\prose{กินฺติ เต อานนฺท สุตํ “วชฺชีนํ อรหนฺเตสุ ธมฺมิกา รกฺขาวรณคุตฺติ สุสํวิหิตา ‘กินฺติ อนาคตา จ อรหนฺโต วิชิตํ อาคจฺเฉยฺยุํ, อาคตา จ อรหนฺโต วิชิเต ผาสุํ วิหเรยฺยุนฺ’ติ”.\csromanspacer{} สุตํ เมตํ ภนฺเต วชฺชีนํ อรหนฺเตสุ ธมฺมิกา รกฺขาวรณคุตฺติ สุสํวิหิตา ภวิสฺสติ “กินฺติ อนาคตา จ อรหนฺโต วิชิตํ อาคจฺเฉยฺยุํ, อาคตา จ อรหนฺโต วิชิเต ผาสุํ วิหเรยฺยุนฺ”ติ.\csromanspacer{} ยาวกีวญฺจ อานนฺท วชฺชีนํ อรหนฺเตสุ ธมฺมิกา รกฺขาวรณคุตฺติ สุสํวิหิตา ภวิสฺสติ “กินฺติ อนาคตา จ อรหนฺโต วิชิตํ อาคจฺเฉยฺยุํ, อาคตา จ อรหนฺโต วิชิเต ผาสุํ วิหเรยฺยุนฺ”ติ, วุทฺธิเยว อานนฺท วชฺชีนํ ปาฏิกงฺขา โน ปริหานีติ. (7)}

\prose{อถ โข ภควา วสฺสการํ พฺราหฺมณํ มาคธมหามตฺตํ อามนฺเตสิ “เอกมิทาหํ พฺราหฺมณ สมยํ เวสาลิยํ วิหรามิ สารนฺทเท เจติเย.\csromanspacer{} ตตฺราหํ พฺราหฺมณ วชฺชีนํ อิเม สตฺต อปริหานิเย ธมฺเม เทเสสิํ.\csromanspacer{} ยาวกีวญฺจ พฺราหฺมณ อิเม สตฺต อปริหานิยา ธมฺมา วชฺชีสุ ฐสฺสนฺติ, อิเมสุ จ สตฺตสุ อปริหานิเยสุ ธมฺเมสุ วชฺชี สนฺทิสฺสิสฺสนฺติ, วุทฺธิเยว พฺราหฺมณ วชฺชีนํ ปาฏิกงฺขา โน ปริหานี”ติ.}

\gambhira{สตฺตกนิปาตปาฬิ}
\csromanmatikaanchor{mtk.230}
\csromantocmarkref{subsection}{3-7. (อปริหานิย) สตฺตกสุตฺตาทิ}{mtk.230}
\bookmark[dest={mtk.230},level=2]{3-7. (อปริหานิย) สตฺตกสุตฺตาทิ}
\prose{\csromanfolio{412}เอกเมเกนปิ\footnote{เอกเมเกนปิ เตน โข (ก) ที 2. 64 ปิฏฺเฐปิ.} โภ โคตม อปริหานิเยน ธมฺเมน สมนฺนาคตานํ วชฺชีนํ วุทฺธิเยว ปาฏิกงฺขา โน ปริหานิ, โก ปน วาโท สตฺตหิ อปริหานิเยหิ ธมฺเมหิ.\csromanspacer{} อกรณียา จ โภ โคตม วชฺชี รญฺญา มาคเธน อชาตสตฺตุนา เวเทหิปุตฺเตน ยทิทํ ยุทฺธสฺส อญฺญตฺร อุปลาปนาย\footnote{อุปลาปนา (ก-สี, ก)} อญฺญตฺร มิถุเภทา, หนฺท จ ทานิ มยํ โภ โคตม คจฺฉาม, พหุกิจฺจา มยํ พหุกรณียาติ.\csromanspacer{} ยสฺสทานิ ตฺวํ พฺราหฺมณ กาลํ มญฺญสีติ.\csromanspacer{} อถ โข วสฺสกาโร พฺราหฺมโณ มาคธมหามตฺโต ภควโต ภาสิตํ อภินนฺทิตฺวา อนุโมทิตฺวา อุฏฺฐายาสนา ปกฺกามีติ.\csromanspacer{}\csromanspacer{}ทุติยํ.}

\csromanheader{2}{3. ปฐมสตฺตกสุตฺต}
\proseitem{23}{เอวํ เม สุตํ\csromandash{}เอกํ สมยํ ภควา ราชคเห วิหรติ คิชฺฌกูเฏ ปพฺพเต.\csromanspacer{} ตตฺร โข ภควา ภิกฺขู อามนฺเตสิ “สตฺต โว ภิกฺขเว อปริหานิเย ธมฺเม เทเสสฺสามิ, ตํ สุณาถ สาธุกํ มนสิ กโรถ, ภาสิสฺสามี”ติ. “เอวํ ภนฺเต”ติ โข เต ภิกฺขู ภควโต ปจฺจสฺโสสุํ.\csromanspacer{} ภควา เอตทโวจ\csromandash{}}

\prose{กตเม จ ภิกฺขเว สตฺต อปริหานิยา ธมฺมา? ยาวกีวญฺจ ภิกฺขเว ภิกฺขู อภิณฺหํ สนฺนิปาตา ภวิสฺสนฺติ สนฺนิปาตพหุลา, วุทฺธิเยว ภิกฺขเว ภิกฺขูนํ ปาฏิกงฺขา โน ปริหานิ. (1)}

\prose{ยาวกีวญฺจ ภิกฺขเว ภิกฺขู สมคฺคา สนฺนิปติสฺสนฺติ, สมคฺคา วุฏฺฐหิสฺสนฺติ, สมคฺคา สํฆกรณียานิ กริสฺสนฺติ, วุทฺธิเยว ภิกฺขเว ภิกฺขูนํ ปาฏิกงฺขา โน ปริหานิ. (2)}

\prose{ยาวกีวญฺจ ภิกฺขเว ภิกฺขู อปญฺญตฺตํ น ปญฺญาเปสฺสนฺติ, ปญฺญตฺตํ น สมุจฺฉินฺทิสฺสนฺติ, ยถาปญฺญตฺเตสุ สิกฺขาปเทสุ สมาทาย วตฺติสฺสนฺติ, วุทฺธิเยว ภิกฺขเว ภิกฺขูนํ ปาฏิกงฺขา โน ปริหานิ. (3)}

\prose{ยาวกีวญฺจ ภิกฺขเว ภิกฺขู เย เต ภิกฺขู เถรา รตฺตญฺญู จิรปพฺพชิตา สํฆปิตโร สํฆปริณายกา, เต สกฺกริสฺสนฺติ ครุํ กริสฺสนฺติ มาเนสฺสนฺติ ปูเชสฺสนฺติ, เตสญฺจ โสตพฺพํ มญฺญิสฺสนฺติ, วุทฺธิเยว ภิกฺขเว ภิกฺขูนํ ปาฏิกงฺขา โน ปริหานิ. (4)}

\chapterheadpageread{3. วชฺชิสตฺตกวคฺค}
\prose{\csromanfolio{413}ยาวกีวญฺจ ภิกฺขเว ภิกฺขู อุปฺปนฺนาย ตณฺหาย โปโนภวิกาย น วสํ คจฺฉิสฺสนฺติ, วุทฺธิเยว ภิกฺขเว ภิกฺขูนํ ปาฏิกงฺขา โน ปริหานิ. (5)}

\prose{ยาวกีวญฺจ ภิกฺขเว ภิกฺขู อารญฺญเกสุ เสนาสเนสุ สาเปกฺขา ภวิสฺสนฺติ, วุทฺธิเยว ภิกฺขเว ภิกฺขูนํ ปาฏิกงฺขา โน ปริหานิ. (6)}

\prose{ยาวกีวญฺจ ภิกฺขเว ภิกฺขู ปจฺจตฺตญฺเญว สติํ อุปฏฺฐาเปสฺสนฺติ “กินฺติ อนาคตา จ เปสลา สพฺรหฺมจารี อาคจฺเฉยฺยุํ, อาคตา จ เปสลา สพฺรหฺมจารี ผาสุํ วิหเรยฺยุนฺ”ติ, วุทฺธิเยว ภิกฺขเว ภิกฺขูนํ ปาฏิกงฺขา โน ปริหานิ. (7)}

\prose{ยาวกีวญฺจ ภิกฺขเว อิเม สตฺต อปริหานิยา ธมฺมา ภิกฺขูสุ ฐสฺสนฺติ, อิเมสุ จ สตฺตสุ อปริหานิเยสุ ธมฺเมสุ ภิกฺขู สนฺทิสฺสิสฺสนฺติ, วุทฺธิเยว ภิกฺขเว ภิกฺขูนํ ปาฏิกงฺขา โน ปริหานีติ.\csromanspacer{}\csromanspacer{}ตติยํ.}

\csromanheader{2}{4. ทุติยสตฺตกสุตฺต}
\proseitem{24}{\csromansymbolfootnote{*}{ที 2. 66 ปิฏฺเฐปิ.}สตฺต โว ภิกฺขเว อปริหานิเย ธมฺเม เทเสสฺสามิ, ตํ สุณาถ สาธุกํ มนสิ กโรถ ฯเปฯ.\csromanspacer{} กตเม จ ภิกฺขเว สตฺต อปริหานิยา ธมฺมา?}

\prose{ยาวกีวญฺจ ภิกฺขเว ภิกฺขู น กมฺมารามา ภวิสฺสนฺติ น กมฺมรตา น กมฺมารามตํ อนุยุตฺตา, วุทฺธิเยว ภิกฺขเว ภิกฺขูนํ ปาฏิกงฺขา โน ปริหานิ.}

\prose{ยาวกีวญฺจ ภิกฺขเว ภิกฺขู น ภสฺสารามา ภวิสฺสนฺติ ฯเปฯ น นิทฺทารามา ภวิสฺสนฺติ.\csromanspacer{} น สงฺคณิการามา ภวิสฺสนฺติ.\csromanspacer{} น ปาปิจฺฉา ภวิสฺสนฺติ น ปาปิกานํ อิจฺฉานํ วสํ คตา.\csromanspacer{} น ปาปมิตฺตา ภวิสฺสนฺติ น ปาปสหายา น ปาปสมฺปวงฺกา.\csromanspacer{} น โอรมตฺตเกน วิเสสาธิคเมน อนฺตราโวสานํ อาปชฺชิสฺสนฺติ.\csromanspacer{} วุทฺธิเยว ภิกฺขเว ภิกฺขูนํ ปาฏิกงฺขา โน ปริหานิ.}

\prose{ยาวกีวญฺจ ภิกฺขเว อิเม สตฺต อปริหานิยา ธมฺมา ภิกฺขูสุ ฐสฺสนฺติ, อิเมสุ จ สตฺตสุ อปริหานิเยสุ ธมฺเมสุ ภิกฺขู สนฺทิสฺสิสฺสนฺติ, วุทฺธิเยว ภิกฺขเว ภิกฺขูนํ ปาฏิกงฺขา โน ปริหานีติ.\csromanspacer{}\csromanspacer{}จตุตฺถํ.}

\gambhira{สตฺตกนิปาตปาฬิ \textbf{5. ตติยสตฺตกสุตฺต}}
\proseitem{25}{\csromanfolio{414}สตฺต โว ภิกฺขเว อปริหานิเย ธมฺเม เทเสสฺสามิ, ตํ สุณาถ สาธุกํ มนสิ กโรถ ฯเปฯ.\csromanspacer{} กตเม จ ภิกฺขเว สตฺต อปริหานิยา ธมฺมา? ยาวกีวญฺจ ภิกฺขเว ภิกฺขู สทฺธา ภวิสฺสนฺติ, วุทฺธิเยว ภิกฺขเว ภิกฺขูนํ ปาฏิกงฺขา โน ปริหานิ. ยาวกีวญฺจ ภิกฺขเว ภิกฺขู หิริมนฺโต\footnote{หิรีมา (สี), หิริมนา (ที 2. 66 ปิฏฺเฐ)} ภวิสฺสนฺติ ฯเปฯ โอตฺตปฺปิโน\footnote{โอตฺตาปีโน (สี)} ภวิสฺสนฺติ.\csromanspacer{} พหุสฺสุตา ภวิสฺสนฺติ.\csromanspacer{} อารทฺธวีริยา ภวิสฺสนฺติ.\csromanspacer{} สติมนฺโต ภวิสฺสนฺติ.\csromanspacer{} ปญฺญวนฺโต ภวิสฺสนฺติ, วุทฺธิเยว ภิกฺขเว ภิกฺขูนํ ปาฏิกงฺขา โน ปริหานิ.\csromanspacer{} ยาวกีวญฺจ ภิกฺขเว อิเม สตฺต อปริหานิยา ธมฺมา ภิกฺขูสุ ฐสฺสนฺติ, อิเมสุ จ สตฺตสุ อปริหานิเยสุ ธมฺเมสุ ภิกฺขู สนฺทิสฺสิสฺสนฺติ, วุทฺธิเยว ภิกฺขเว ภิกฺขูนํ ปาฏิกงฺขา โน ปริหานีติ.\csromanspacer{}\csromanspacer{}ปญฺจมํ.}

\csromanheader{2}{6. โพชฺฌงฺคสุตฺต}
\proseitem{26}{สตฺต โว ภิกฺขเว อปริหานิเย ธมฺเม เทเสสฺสามิ, ตํ สุณาถ สาธุกํ มนสิ กโรถ ฯเปฯ.\csromanspacer{} กตเม จ ภิกฺขเว สตฺต อปริหานิยา ธมฺมา? ยาวกีวญฺจ ภิกฺขเว ภิกฺขู สติสมฺโพชฺฌงฺคํ ภาเวสฺสนฺติ, วุทฺธิเยว ภิกฺขเว ภิกฺขูนํ ปาฏิกงฺขา โน ปริหานิ. ยาวกีวญฺจ ภิกฺขเว ภิกฺขู ธมฺมวิจยสมฺโพชฺฌงฺคํ ภาเวสฺสนฺติ ฯเปฯ วีริยสมฺโพชฺฌงฺคํ ภาเวสฺสนฺติ.\csromanspacer{} ปีติสมฺโพชฺฌงฺคํ ภาเวสฺสนฺติ.\csromanspacer{} ปสฺสทฺธิสมฺโพชฺฌงฺคํ ภาเวสฺสนฺติ.\csromanspacer{} สมาธิสมฺโพชฺฌงฺคํ ภาเวสฺสนฺติ.\csromanspacer{} อุเปกฺขาสมฺโพชฺฌงฺคํ ภาเวสฺสนฺติ, วุทฺธิเยว ภิกฺขเว ภิกฺขูนํ ปาฏิกงฺขา โน ปริหานิ.\csromanspacer{} ยาวกีวญฺจ ภิกฺขเว อิเม สตฺต อปริหานิยา ธมฺมา ภิกฺขูสุ ฐสฺสนฺติ, อิเมสุ จ สตฺตสุ อปริหานิเยสุ ธมฺเมสุ ภิกฺขู สนฺทิสฺสิสฺสนฺติ, วุทฺธิเยว ภิกฺขเว ภิกฺขูนํ ปาฏิกงฺขา โน ปริหานีติ.\csromanspacer{}\csromanspacer{}ฉฏฺฐํ.}

\csromanheader{2}{7. สญฺญาสุตฺต}
\proseitem{27}{สตฺต โว ภิกฺขเว อปริหานิเย ธมฺเม เทเสสฺสามิ, ตํ สุณาถ สาธุกํ มนสิ กโรถ ฯเปฯ.\csromanspacer{} กตเม จ ภิกฺขเว สตฺต}

\vspace{\tipitakaparskip}
\csromancenter{วชฺชิสตฺตกวคฺค อปริหานิยา ธมฺมา? ยาวกีวญฺจ ภิกฺขเว ภิกฺขู อนิจฺจสญฺญํ ภาเวสฺสนฺติ, วุทฺธิเยว ภิกฺขเว ภิกฺขูนํ ปาฏิกงฺขา โน ปริหานิ.}
\csromanmatikaanchor{mtk.231}
\csromantocmarkref{subsection}{8-11. ปริหานิสุตฺตาทิ}{mtk.231}
\bookmark[dest={mtk.231},level=2]{8-11. ปริหานิสุตฺตาทิ}
\prosecont{\csromanfolio{415}ยาวกีวญฺจ ภิกฺขเว ภิกฺขู อนตฺตสญฺญํ ภาเวสฺสนฺติ ฯเปฯ อสุภสญฺญํ ภาเวสฺสนฺติ.\csromanspacer{} อาทีนวสญฺญํ ภาเวสฺสนฺติ.\csromanspacer{} ปหานสญฺญํ ภาเวสฺสนฺติ.\csromanspacer{} วิราคสญฺญํ ภาเวสฺสนฺติ.\csromanspacer{} นิโรธสญฺญํ ภาเวสฺสนฺติ, วุทฺธิเยว ภิกฺขเว ภิกฺขูนํ ปาฏิกงฺขา โน ปริหานิ.\csromanspacer{} ยาวกีวญฺจ ภิกฺขเว อิเม สตฺต อปริหานิยา ธมฺมา ภิกฺขูสุ ฐสฺสนฺติ, อิเมสุ จ สตฺตสุ อปริหานิเยสุ ธมฺเมสุ ภิกฺขู สนฺทิสฺสิสฺสนฺติ, วุทฺธิเยว ภิกฺขเว ภิกฺขูนํ ปาฏิกงฺขา โน ปริหานีติ*.\csromanspacer{}\csromanspacer{}สตฺตมํ.}

\csromanheader{2}{8. ปฐมปริหานิสุตฺต}
\proseitem{28}{เอวํ เม สุตํ\csromandash{}เอกํ สมยํ ภควา สาวตฺถิยํ วิหรติ เชตวเน อนาถปิณฺฑิกสฺส อาราเม.\csromanspacer{} ตตฺร โข ภควา ภิกฺขู อามนฺเตสิ “สตฺติเม ภิกฺขเว ธมฺมา เสขสฺส ภิกฺขุโน ปริหานาย สํวตฺตนฺติ.\csromanspacer{} กตเม สตฺต? กมฺมารามตา ภสฺสารามตา นิทฺทารามตา สงฺคณิการามตา อินฺทฺริเยสุ อคุตฺตทฺวารตา โภชเน อมตฺตญฺญุตา.\csromanspacer{} สนฺติ โข ปน สํเฆ สํฆกรณียานิ, ตตฺร เสโข ภิกฺขุ\footnote{ตตฺร ภิกฺขุ (สี, สฺยา)} อิติ ปฏิสญฺจิกฺขติ ‘สนฺติ โข ปน สํเฆ เถรา\footnote{โข สํฆตฺเถรา (ก)} รตฺตญฺญู จิรปพฺพชิตา ภารวาหิโน, เต\footnote{น เต (ก)} เตน ปญฺญายิสฺสนฺตี’ติ อตฺตนา เตสุ โยคํ\footnote{อตฺตนา โวโยคํ (สี, สฺยา)} อาปชฺชติ.\csromanspacer{} อิเม โข ภิกฺขเว สตฺต ธมฺมา เสขสฺส ภิกฺขุโน ปริหานาย สํวตฺตนฺติ.}

\prose{สตฺติเม ภิกฺขเว ธมฺมา เสขสฺส ภิกฺขุโน อปริหานาย สํวตฺตนฺติ.\csromanspacer{} กตเม สตฺต? น กมฺมารามตา น ภสฺสารามตา น นิทฺทารามตา น สงฺคณิการามตา อินฺทฺริเยสุ คุตฺตทฺวารตา โภชเน มตฺตญฺญุตา.\csromanspacer{} สนฺติ โข ปน สํเฆ สํฆกรณียานิ, ตตฺร เสโข ภิกฺขุ อิติ ปฏิสญฺจิกฺขติ ‘สนฺติ โข ปน สํเฆ เถรา รตฺตญฺญู จิรปพฺพชิตา ภารวาหิโน, เต เตน ปญฺญายิสฺสนฺตี’ติ อตฺตนา น เตสุ โยคํ อาปชฺชติ.\csromanspacer{} อิเม โข ภิกฺขเว สตฺต ธมฺมา เสขสฺส ภิกฺขุโน อปริหานาย สํวตฺตนฺตี”ติ.\csromanspacer{}\csromanspacer{}อฏฺฐมํ.}

\gambhira{สตฺตกนิปาตปาฬิ \textbf{9. ทุติยปริหานิสุตฺต}}
\proseitem{29}{\csromanfolio{416}สตฺติเม ภิกฺขเว ธมฺมา อุปาสกสฺส ปริหานาย สํวตฺตนฺติ.\csromanspacer{} กตเม สตฺต? ภิกฺขุทสฺสนํ หาเปติ, สทฺธมฺมสฺสวนํ ปมชฺชติ, อธิสีเล น สิกฺขติ, อปฺปสาทพหุโล โหติ ภิกฺขูสุ เถเรสุ เจว นเวสุ จ มชฺฌิเมสุ จ, อุปารมฺภจิตฺโต ธมฺมํ สุณาติ รนฺธคเวสี, อิโต พหิทฺธา ทกฺขิเณยฺยํ คเวสติ, ตตฺถ จ ปุพฺพการํ กโรติ.\csromanspacer{} อิเม โข ภิกฺขเว สตฺต ธมฺมา อุปาสกสฺส ปริหานาย สํวตฺตนฺติ.}

\prose{สตฺติเม ภิกฺขเว ธมฺมา อุปาสกสฺส อปริหานาย สํวตฺตนฺติ.\csromanspacer{} กตเม สตฺต? ภิกฺขุทสฺสนํ น หาเปติ, สทฺธมฺมสฺสวนํ นปฺปมชฺชติ, อธิสีเล สิกฺขติ, ปสาทพหุโล โหติ ภิกฺขูสุ เถเรสุ เจว นเวสุ จ มชฺฌิเมสุ จ, อนุปารมฺภจิตฺโต ธมฺมํ สุณาติ น รนฺธคเวสี, น อิโต พหิทฺธา ทกฺขิเณยฺยํ คเวสติ, อิธ จ ปุพฺพการํ กโรติ.\csromanspacer{} อิเม โข ภิกฺขเว สตฺต ธมฺมา อุปาสกสฺส อปริหานาย สํวตฺตนฺตีติ.\csromanspacer{} อิทมโวจ ภควา, อิทํ วตฺวาน สุคโต, อถาปรํ เอตทโวจ สตฺถา\csromandash{}}

\csromangathameasure{“ทสฺสนํ ภาวิตตฺตานํ, โย หาเปติ อุปาสโก. \\ สวนญฺจ อริยธมฺมานํ, อธิสีเล น สิกฺขติ. \\ อปฺปสาโท จ ภิกฺขูสุ, ภิยฺโย ภิยฺโย ปวฑฺฒติ. \\ อุปารมฺภกจิตฺโต จ, สทฺธมฺมํ โสตุมิจฺฉติ. \\ อิโต จ พหิทฺธา อญฺญํ, ทกฺขิเณยฺยํ คเวสติ. \\ ตตฺเถว จ ปุพฺพการํ, โย กโรติ อุปาสโก. \\ เอเต โข ปริหานิเย, สตฺต ธมฺเม สุเทสิเต. \\ อุปาสโก เสวมาโน, สทฺธมฺมา ปริหายติ. \\ ทสฺสนํ ภาวิตตฺตานํ, โย น หาเปติ อุปาสโก. \\ สวนญฺจ อริยธมฺมานํ, อธิสีเล จ สิกฺขติ. \\ ปสาโท จสฺส ภิกฺขูสุ, ภิยฺโย ภิยฺโย ปวฑฺฒติ. \\ อนุปารมฺภจิตฺโต จ, สทฺธมฺมํ โสตุมิจฺฉติ. \\ น อิโต พหิทฺธา อญฺญํ, ทกฺขิเณยฺยํ คเวสติ. \\ อิเธว จ ปุพฺพการํ, โย กโรติ อุปาสโก.}
\csromangathagroup{\csromangathastanza{“ทสฺสนํ ภาวิตตฺตานํ, โย หาเปติ อุปาสโก. \\ สวนญฺจ อริยธมฺมานํ, อธิสีเล น สิกฺขติ.}\csromangathastanza{อปฺปสาโท จ ภิกฺขูสุ, ภิยฺโย ภิยฺโย ปวฑฺฒติ. \\ อุปารมฺภกจิตฺโต จ, สทฺธมฺมํ โสตุมิจฺฉติ.}\csromangathastanza{อิโต จ พหิทฺธา อญฺญํ, ทกฺขิเณยฺยํ คเวสติ. \\ ตตฺเถว จ ปุพฺพการํ, โย กโรติ อุปาสโก.}\csromangathastanza{เอเต โข ปริหานิเย, สตฺต ธมฺเม สุเทสิเต. \\ อุปาสโก เสวมาโน, สทฺธมฺมา ปริหายติ.}\csromangathastanza{ทสฺสนํ ภาวิตตฺตานํ, โย น หาเปติ อุปาสโก. \\ สวนญฺจ อริยธมฺมานํ, อธิสีเล จ สิกฺขติ.}\csromangathastanza{ปสาโท จสฺส ภิกฺขูสุ, ภิยฺโย ภิยฺโย ปวฑฺฒติ. \\ อนุปารมฺภจิตฺโต จ, สทฺธมฺมํ โสตุมิจฺฉติ.}\csromangathastanza{น อิโต พหิทฺธา อญฺญํ, ทกฺขิเณยฺยํ คเวสติ. \\ อิเธว จ ปุพฺพการํ, โย กโรติ อุปาสโก.}}

\chapterheadpageread{3. วชฺชิสตฺตกวคฺค}
\csromangathameasure{เอเต โข อปริหานิเย, สตฺต ธมฺเม สุเทสิเต. \\ อุปาสโก เสวมาโน, สทฺธมฺมา น ปริหายตี”ติ. นวมํ.}
\csromangathagroup{\csromangathastanza{\csromanfolio{417}เอเต โข อปริหานิเย, สตฺต ธมฺเม สุเทสิเต. \\ อุปาสโก เสวมาโน, สทฺธมฺมา น ปริหายตี”ติ.\csromanspacer{} นวมํ.}}

\csromanheader{2}{10. วิปตฺติสุตฺต}
\proseitem{30}{สตฺติมา ภิกฺขเว อุปาสกสฺส วิปตฺติโย ฯเปฯ.\csromanspacer{} สตฺติมา ภิกฺขเว อุปาสกสฺส สมฺปทา ฯเปฯ.\csromanspacer{}\csromanspacer{}ทสมํ.}

\csromanheader{2}{11. ปราภวสุตฺต}
\proseitem{31}{สตฺติเม ภิกฺขเว อุปาสกสฺส ปราภวา ฯเปฯ.\csromanspacer{} สตฺติเม ภิกฺขเว อุปาสกสฺส สมฺภวา.\csromanspacer{} กตเม สตฺต? ภิกฺขุทสฺสนํ น หาเปติ, สทฺธมฺมสฺสวนํ นปฺปมชฺชติ, อธิสีเล สิกฺขติ, ปสาทพหุโล โหติ ภิกฺขูสุ เถเรสุ เจว นเวสุ จ มชฺฌิเมสุ จ, อนุปารมฺภจิตฺโต ธมฺมํ สุณาติ น รนฺธคเวสี, น อิโต พหิทฺธา ทกฺขิเณยฺยํ คเวสติ, อิธ จ ปุพฺพการํ กโรติ.\csromanspacer{} อิเม โข ภิกฺขเว สตฺต อุปาสกสฺส สมฺภวาติ.}

\csromangathameasure{ทสฺสนํ ภาวิตตฺตานํ, โย หาเปติ อุปาสโก. \\ สวนญฺจ อริยธมฺมานํ, อธิสีเล น สิกฺขติ. \\ อปฺปสาโท จ ภิกฺขูสุ, ภิยฺโย ภิยฺโย ปวฑฺฒติ. \\ อุปารมฺภกจิตฺโต จ, สทฺธมฺมํ โสตุมิจฺฉติ. \\ อิโต จ พหิทฺธา อญฺญํ, ทกฺขิเณยฺยํ คเวสติ. \\ ตตฺเถว จ ปุพฺพการํ, โย กโรติ อุปาสโก. \\ เอเต โข ปริหานิเย, สตฺต ธมฺเม สุเทสิเต. \\ อุปาสโก เสวมาโน, สทฺธมฺมา ปริหายติ. \\ ทสฺสนํ ภาวิตตฺตานํ, โย น หาเปติ อุปาสโก. \\ สวนญฺจ อริยธมฺมานํ, อธิสีเล จ สิกฺขติ. \\ ปสาโท จสฺส ภิกฺขูสุ, ภิยฺโย ภิยฺโย ปวฑฺฒติ. \\ อนุปารมฺภจิตฺโต จ, สทฺธมฺมํ โสตุมิจฺฉติ.}
\csromangathagroup{\csromangathastanza{ทสฺสนํ ภาวิตตฺตานํ, โย หาเปติ อุปาสโก. \\ สวนญฺจ อริยธมฺมานํ, อธิสีเล น สิกฺขติ.}\csromangathastanza{อปฺปสาโท จ ภิกฺขูสุ, ภิยฺโย ภิยฺโย ปวฑฺฒติ. \\ อุปารมฺภกจิตฺโต จ, สทฺธมฺมํ โสตุมิจฺฉติ.}\csromangathastanza{อิโต จ พหิทฺธา อญฺญํ, ทกฺขิเณยฺยํ คเวสติ. \\ ตตฺเถว จ ปุพฺพการํ, โย กโรติ อุปาสโก.}\csromangathastanza{เอเต โข ปริหานิเย, สตฺต ธมฺเม สุเทสิเต. \\ อุปาสโก เสวมาโน, สทฺธมฺมา ปริหายติ.}\csromangathastanza{ทสฺสนํ ภาวิตตฺตานํ, โย น หาเปติ อุปาสโก. \\ สวนญฺจ อริยธมฺมานํ, อธิสีเล จ สิกฺขติ.}\csromangathastanza{ปสาโท จสฺส ภิกฺขูสุ, ภิยฺโย ภิยฺโย ปวฑฺฒติ. \\ อนุปารมฺภจิตฺโต จ, สทฺธมฺมํ โสตุมิจฺฉติ.}}

\gambhira{สตฺตกนิปาตปาฬิ}
\csromanmatikaanchor{mtk.232}
\csromanmatikaanchor{mtk.233}
\csromantocmarknopage{chapter}{4. เทวตาวคฺค}{mtk.232}
\bookmark[dest={mtk.232},level=0]{4. เทวตาวคฺค}
\csromantocmarkref{subsection}{1-2. อปฺปมาทคารวสุตฺตาทิ}{mtk.233}
\bookmark[dest={mtk.233},level=2]{1-2. อปฺปมาทคารวสุตฺตาทิ}
\csromangathameasure{น อิโต พหิทฺธา อญฺญํ, ทกฺขิเณยฺยํ คเวสติ. \\ อิเธว จ ปุพฺพการํ, โย กโรติ อุปาสโก. \\ เอเต โข อปริหานิเย, สตฺต ธมฺเม สุเทสิเต. \\ อุปาสโก เสวมาโน, สทฺธมฺมา น ปริหายตีติ. เอกาทสมํ. วชฺชิสตฺตกวคฺโค ตติโย.}
\csromangathagroup{\csromangathastanza{\csromanfolio{418}น อิโต พหิทฺธา อญฺญํ, ทกฺขิเณยฺยํ คเวสติ. \\ อิเธว จ ปุพฺพการํ, โย กโรติ อุปาสโก.}\csromangathastanza{เอเต โข อปริหานิเย, สตฺต ธมฺเม สุเทสิเต. \\ อุปาสโก เสวมาโน, สทฺธมฺมา น ปริหายตีติ.\csromanspacer{} เอกาทสมํ.\csromanspacer{} วชฺชิสตฺตกวคฺโค\footnote{วชฺชิวคฺโค (สฺยา)} ตติโย.}}

\titlehead{\textbf{ตสฺสุทฺทานํ}}
\csromangathameasure{สารนฺท วสฺสกาโร จ, ติสตฺตกานิ ภิกฺขุกา. \\ โพธิสญฺญา เทฺว จ หานิ, วิปตฺติ จ ปราภโวติ.}
\csromangathagroup{\csromangathastanza{สารนฺท วสฺสกาโร จ, ติสตฺตกานิ ภิกฺขุกา. \\ โพธิสญฺญา เทฺว จ หานิ, วิปตฺติ จ ปราภโวติ.}}

\chapterhead{4. เทวตาวคฺค}
\csromanheader{2}{1. อปฺปมาทคารวสุตฺต}
\proseitem{32}{อถ โข อญฺญตรา เทวตา อภิกฺกนฺตาย รตฺติยา อภิกฺกนฺตวณฺณา เกวลกปฺปํ เชตวนํ โอภาเสตฺวา เยน ภควา เตนุปสงฺกมิ, อุปสงฺกมิตฺวา ภควนฺตํ อภิวาเทตฺวา เอกมนฺตํ อฏฺฐาสิ, เอกมนฺตํ ฐิตา โข สา เทวตา ภควนฺตํ เอตทโวจ\csromandash{}}

\prose{สตฺติเม ภนฺเต ธมฺมา ภิกฺขุโน อปริหานาย สํวตฺตนฺติ.\csromanspacer{} กตเม สตฺต? สตฺถุคารวตา ธมฺมคารวตา สํฆคารวตา สิกฺขาคารวตา สมาธิคารวตา อปฺปมาทคารวตา ปฏิสนฺถารคารวตา.\csromanspacer{} อิเม โข ภนฺเต สตฺต ธมฺมา ภิกฺขุโน อปริหานาย สํวตฺตนฺตีติ.\csromanspacer{} อิทมโวจ สา เทวตา, สมนุญฺโญ สตฺถา อโหสิ.\csromanspacer{} อถ โข สา เทวตา “สมนุญฺโญ เม สตฺถา”ติ ภควนฺตํ อภิวาเทตฺวา ปทกฺขิณํ กตฺวา ตตฺเถวนฺตรธายิ.}

\prose{อถ โข ภควา ตสฺสา รตฺติยา อจฺจเยน ภิกฺขู อามนฺเตสิ “อิมํ ภิกฺขเว รตฺติํ อญฺญตรา เทวตา อภิกฺกนฺตาย รตฺติยา อภิกฺกนฺตวณฺณา เกวลกปฺปํ เชตวนํ โอภาเสตฺวา เยนาหํ เตนุปสงฺกมิ,}

\vspace{\tipitakaparskip}
\csromancenter{เทวตาวคฺค อุปสงฺกมิตฺวา มํ อภิวาเทตฺวา เอกมนฺตํ อฏฺฐาสิ, เอกมนฺตํ ฐิตา โข ภิกฺขเว สา เทวตา มํ เอตทโวจ ‘สตฺติเม ภนฺเต ธมฺมา ภิกฺขุโน อปริหานาย สํวตฺตนฺติ.\csromanspacer{} กตเม สตฺต? สตฺถุคารวตา ธมฺมคารวตา สํฆคารวตา สิกฺขาคารวตา สมาธิคารวตา อปฺปมาทคารวตา ปฏิสนฺถารคารวตา.\csromanspacer{} อิเม โข ภนฺเต สตฺต ธมฺมา ภิกฺขุโน อปริหานาย สํวตฺตนฺตี’ติ.\csromanspacer{} อิทมโวจ ภิกฺขเว สา เทวตา, อิทํ วตฺวา มํ อภิวาเทตฺวา ปทกฺขิณํ กตฺวา ตตฺเถวนฺตรธายี”ติ.}
\vspace{0.5\baselineskip}
\csromangathameasure{สตฺถุครุ ธมฺมครุ, สํเฆ จ ติพฺพคารโว. \\ สมาธิครุ อาตาปี, สิกฺขาย ติพฺพคารโว. \\ อปฺปมาทครุ ภิกฺขุ, ปฏิสนฺถารคารโว. \\ อภพฺโพ ปริหานาย, นิพฺพานสฺเสว สนฺติเกติ.ปฐมํ.}
\csromangathagroup{\csromangathastanza{\csromanfolio{419}สตฺถุครุ ธมฺมครุ, สํเฆ จ ติพฺพคารโว. \\ สมาธิครุ อาตาปี\footnote{สมาธิคารวตาปิ จ (ก)}, สิกฺขาย ติพฺพคารโว.}\csromangathastanza{อปฺปมาทครุ ภิกฺขุ, ปฏิสนฺถารคารโว. \\ อภพฺโพ ปริหานาย, นิพฺพานสฺเสว สนฺติเกติ.\csromanspacer{}\csromanspacer{}ปฐมํ.}}

\csromanheader{2}{2. หิรีคารวสุตฺต}
\proseitem{33}{อิมํ ภิกฺขเว รตฺติํ อญฺญตรา เทวตา อภิกฺกนฺตาย รตฺติยา อภิกฺกนฺตวณฺณา เกวลกปฺปํ เชตวนํ โอภาเสตฺวา เยนาหํ เตนุปสงฺกมิ, อุปสงฺกมิตฺวา มํ อภิวาเทตฺวา เอกมนฺตํ อฏฺฐาสิ, เอกมนฺตํ ฐิตา โข ภิกฺขเว สา เทวตา มํ เอตทโวจ “สตฺติเม ภนฺเต ธมฺมา ภิกฺขุโน อปริหานาย สํวตฺตนฺติ.\csromanspacer{} กตเม สตฺต? สตฺถุคารวตา ธมฺมคารวตา สํฆคารวตา สิกฺขาคารวตา สมาธิคารวตา หิริคารวตา โอตฺตปฺปคารวตา.\csromanspacer{} อิเม โข ภนฺเต สตฺต ธมฺมา ภิกฺขุโน อปริหานาย สํวตฺตนฺตี”ติ.\csromanspacer{} อิทมโวจ ภิกฺขเว สา เทวตา, อิทํ วตฺวา มํ อภิวาเทตฺวา ปทกฺขิณํ กตฺวา ตตฺเถวนฺตรธายีติ.}

\csromangathameasure{สตฺถุครุ ธมฺมครุ, สํเฆ จ ติพฺพคารโว. \\ สมาธิครุ อาตาปี, สิกฺขาย ติพฺพคารโว. \\ หิริ-โอตฺตปฺปสมฺปนฺโน, สปฺปติสฺโส สคารโว. \\ อภพฺโพ ปริหานาย, นิพฺพานสฺเสว สนฺติเกติ.ทุติยํ.}
\csromangathagroup{\csromangathastanza{สตฺถุครุ ธมฺมครุ, สํเฆ จ ติพฺพคารโว. \\ สมาธิครุ อาตาปี, สิกฺขาย ติพฺพคารโว.}\csromangathastanza{หิริ-โอตฺตปฺปสมฺปนฺโน, สปฺปติสฺโส สคารโว. \\ อภพฺโพ ปริหานาย, นิพฺพานสฺเสว สนฺติเกติ.\csromanspacer{}\csromanspacer{}ทุติยํ.}}

\csromanheader{2}{สตฺตกนิปาตปาฬิ \textbf{3. ปฐมโสวจสฺสตาสุตฺต}}
\csromanmatikaanchor{mtk.234}
\csromantocmarkref{subsection}{3-4. โสวจสฺสตาสุตฺตทฺวย}{mtk.234}
\bookmark[dest={mtk.234},level=2]{3-4. โสวจสฺสตาสุตฺตทฺวย}
\proseitem{34}{\csromanfolio{420}อิมํ ภิกฺขเว รตฺติํ อญฺญตรา เทวตา ฯเปฯ มํ เอตทโวจ “สตฺติเม ภนฺเต ธมฺมา ภิกฺขุโน อปริหานาย สํวตฺตนฺติ.\csromanspacer{} กตเม สตฺต? สตฺถุคารวตา ธมฺมคารวตา สํฆคารวตา สิกฺขาคารวตา สมาธิคารวตา โสวจสฺสตา กลฺยาณมิตฺตตา.\csromanspacer{} อิเม โข ภนฺเต สตฺต ธมฺมา ภิกฺขุโน อปริหานาย สํวตฺตนฺตี”ติ.\csromanspacer{} อิทมโวจ ภิกฺขเว สา เทวตา, อิทํ วตฺวา มํ อภิวาเทตฺวา ปทกฺขิณํ กตฺวา ตตฺเถวนฺตรธายีติ.}

\csromangathameasure{สตฺถุครุ ธมฺมครุ, สํเฆ จ ติพฺพคารโว. \\ สมาธิครุ อาตาปี, สิกฺขาย ติพฺพคารโว. \\ กลฺยาณมิตฺโต สุวโจ, สปฺปติสฺโส สคารโว. \\ อภพฺโพ ปริหานาย, นิพฺพานสฺเสว สนฺติเกติ.ตติยํ.}
\csromangathagroup{\csromangathastanza{สตฺถุครุ ธมฺมครุ, สํเฆ จ ติพฺพคารโว. \\ สมาธิครุ อาตาปี, สิกฺขาย ติพฺพคารโว.}\csromangathastanza{กลฺยาณมิตฺโต สุวโจ, สปฺปติสฺโส สคารโว. \\ อภพฺโพ ปริหานาย, นิพฺพานสฺเสว สนฺติเกติ.\csromanspacer{}\csromanspacer{}ตติยํ.}}

\csromanheader{2}{4. ทุติยโสวจสฺสตาสุตฺต}
\proseitem{35}{อิมํ ภิกฺขเว รตฺติํ อญฺญตรา เทวตา อภิกฺกนฺตาย รตฺติยา อภิกฺกนฺตวณฺณา ฯเปฯ สตฺติเม ภนฺเต ธมฺมา ภิกฺขุโน อปริหานาย สํวตฺตนฺติ.\csromanspacer{} กตเม สตฺต? สตฺถุคารวตา ธมฺมคารวตา สํฆคารวตา สิกฺขาคารวตา สมาธิคารวตา โสวจสฺสตา กลฺยาณมิตฺตตา.\csromanspacer{} อิเม โข ภนฺเต สตฺต ธมฺมา ภิกฺขุโน อปริหานาย สํวตฺตนฺตีติ.\csromanspacer{} อิทมโวจ ภิกฺขเว สา เทวตา, อิทํ วตฺวา มํ อภิวาเทตฺวา ปทกฺขิณํ กตฺวา ตตฺเถวนฺตรธายีติ.}

\prose{เอวํ วุตฺเต อายสฺมา สาริปุตฺโต ภควนฺตํ เอตทโวจ “อิมสฺส โข อหํ ภนฺเต ภควตา สํขิตฺเตน ภาสิตสฺส เอวํ วิตฺถาเรน อตฺถํ อาชานามิ.\csromanspacer{} อิธ ภนฺเต ภิกฺขุ อตฺตนา จ สตฺถุคารโว โหติ, สตฺถุคารวตาย จ วณฺณวาที.\csromanspacer{} เย จญฺเญ ภิกฺขู น สตฺถุคารวา, เต จ สตฺถุคารวตาย สมาทเปติ.\csromanspacer{} เย จญฺเญ ภิกฺขู สตฺถุคารวา, เตสญฺจ วณฺณํ ภณติ ภูตํ ตจฺฉํ กาเลน.\csromanspacer{} อตฺตนา จ ธมฺมคารโว โหติ ฯเปฯ สํฆคารโว โหติ.\csromanspacer{} สิกฺขาคารโว โหติ.\csromanspacer{} สมาธิคารโว โหติ.\csromanspacer{} สุวโจ โหติ.\csromanspacer{} กลฺยาณมิตฺโต โหติ, กลฺยาณมิตฺตตาย จ วณฺณวาที.\csromanspacer{} เย จญฺเญ ภิกฺขู น กลฺยาณมิตฺตา, เต จ กลฺยาณมิตฺตตาย}

\vspace{\tipitakaparskip}
\csromancenter{เทวตาวคฺค สมาทเปติ.\csromanspacer{} เย จญฺเญ ภิกฺขู กลฺยาณมิตฺตา, เตสญฺจ วณฺณํ ภณติ ภูตํ ตจฺฉํ กาเลนาติ.\csromanspacer{} อิมสฺส โข อหํ ภนฺเต ภควตา สํขิตฺเตน ภาสิตสฺส เอวํ วิตฺถาเรน อตฺถํ อาชานามี”ติ.}
\csromanmatikaanchor{mtk.235}
\csromantocmarkref{subsection}{5-6. มิตฺตสุตฺตทฺวย}{mtk.235}
\bookmark[dest={mtk.235},level=2]{5-6. มิตฺตสุตฺตทฺวย}
\prose{\csromanfolio{421}สาธุ สาธุ สาริปุตฺต, สาธุ โข ตฺวํ สาริปุตฺต อิมสฺส มยา สํขิตฺเตน ภาสิตสฺส เอวํ วิตฺถาเรน อตฺถํ อาชานาสิ.\csromanspacer{} อิธ สาริปุตฺต ภิกฺขุ อตฺตนา จ สตฺถุคารโว โหติ, สตฺถุคารวตาย จ วณฺณวาที.\csromanspacer{} เย จญฺเญ ภิกฺขู น สตฺถุคารวา, เต จ สตฺถุคารวตาย สมาทเปติ.\csromanspacer{} เย จญฺเญ ภิกฺขู สตฺถุคารวา, เตสญฺจ วณฺณํ ภณติ ภูตํ ตจฺฉํ กาเลน.\csromanspacer{} อตฺตนา จ ธมฺมคารโว โหติ ฯเปฯ สํฆคารโว โหติ.\csromanspacer{} สิกฺขาคารโว โหติ.\csromanspacer{} สมาธิคารโว โหติ.\csromanspacer{} สุวโจ โหติ.\csromanspacer{} กลฺยาณมิตฺโต โหติ, กลฺยาณมิตฺตตาย จ วณฺณวาที.\csromanspacer{} เย จญฺเญ ภิกฺขู น กลฺยาณมิตฺตา, เต จ กลฺยาณมิตฺตตาย สมาทเปติ.\csromanspacer{} เย จญฺเญ ภิกฺขู กลฺยาณมิตฺตา, เตสญฺจ วณฺณํ ภณติ ภูตํ ตจฺฉํ กาเลนาติ.\csromanspacer{} อิมสฺส โข สาริปุตฺต มยา สํขิตฺเตน ภาสิตสฺส เอวํ วิตฺถาเรน อตฺโถ ทฏฺฐพฺโพติ.\csromanspacer{}\csromanspacer{}จตุตฺถํ.}

\csromanheader{2}{5. ปฐมมิตฺตสุตฺต}
\proseitem{36}{สตฺตหิ ภิกฺขเว องฺเคหิ สมนฺนาคโต มิตฺโต เสวิตพฺโพ.\csromanspacer{} กตเมหิ สตฺตหิ? ทุทฺททํ ททาติ, ทุกฺกรํ กโรติ, ทุกฺขมํ ขมติ, คุยฺหมสฺส\footnote{คุยฺหสฺส (ก)} อาวิ กโรติ, คุยฺหมสฺส\footnote{คุยฺหํ อสฺส (สี), คุยฺหสฺส (ก)} ปริคุหติ\footnote{ปริคูหติ (สี, สฺยา), ปริคุยฺหติ (ก)}, อาปทาสุ น ชหติ, ขีเณน\footnote{ขีเณ (ก)} นาติมญฺญติ.\csromanspacer{} อิเมหิ โข ภิกฺขเว สตฺตหิ องฺเคหิ สมนฺนาคโต มิตฺโต เสวิตพฺโพติ.}

\csromangathameasure{ทุทฺททํ ททาติ มิตฺโต, ทุกฺกรญฺจาปิ กุพฺพติ. \\ อโถปิสฺส ทุรุตฺตานิ, ขมติ ทุกฺขมานิ จ. \\ คุยฺหญฺจ ตสฺส อกฺขาติ, คุยฺหสฺส ปริคูหติ. \\ อาปทาสุ น ชหาติ, ขีเณน นาติมญฺญติ. \\ ยมฺหิ เอตานิ ฐานานิ, สํวิชฺชนฺตีธ ปุคฺคเล. \\ โส มิตฺโต มิตฺตกาเมน, ภชิตพฺโพ ตถาวิโธติ.ปญฺจมํ.}
\csromangathagroup{\csromangathastanza{ทุทฺททํ ททาติ มิตฺโต, ทุกฺกรญฺจาปิ กุพฺพติ. \\ อโถปิสฺส ทุรุตฺตานิ, ขมติ ทุกฺขมานิ จ\footnote{ทุกฺขมานิปิ (สี, สฺยา)}.}\csromangathastanza{คุยฺหญฺจ ตสฺส\footnote{คุยฺหมสฺส จ (สฺยา)} อกฺขาติ, คุยฺหสฺส ปริคูหติ. \\ อาปทาสุ น ชหาติ, ขีเณน นาติมญฺญติ.}\csromangathastanza{ยมฺหิ เอตานิ ฐานานิ, สํวิชฺชนฺตีธ\footnote{สํวิชฺชนฺติ จ (ก)} ปุคฺคเล. \\ โส มิตฺโต มิตฺตกาเมน, ภชิตพฺโพ ตถาวิโธติ.\csromanspacer{}\csromanspacer{}ปญฺจมํ.}}

\gambhira{สตฺตกนิปาตปาฬิ \textbf{6. ทุติยมิตฺตสุตฺต}}
\csromanmatikaanchor{mtk.236}
\csromantocmarkref{subsection}{7-8. ปฏิสมฺภิทาสุตฺตทฺวย}{mtk.236}
\bookmark[dest={mtk.236},level=2]{7-8. ปฏิสมฺภิทาสุตฺตทฺวย}
\proseitem{37}{\csromanfolio{422}สตฺตหิ ภิกฺขเว ธมฺเมหิ สมนฺนาคโต ภิกฺขุ มิตฺโต เสวิตพฺโพ ภชิตพฺโพ ปยิรุปาสิตพฺโพ อปิ ปนุชฺชมาเนนปิ\footnote{ปณุชฺชมาเนนปิ (สี)}.\csromanspacer{} กตเมหิ สตฺตหิ? ปิโย จ โหติ มนาโป จ ครุ จ ภาวนีโย จ วตฺตา จ วจนกฺขโม จ คมฺภีรญฺจ กถํ กตฺตา โหติ, โน จ อฏฺฐาเน นิโยเชติ.\csromanspacer{} อิเมหิ โข ภิกฺขเว สตฺตหิ ธมฺเมหิ สมนฺนาคโต ภิกฺขุ มิตฺโต เสวิตพฺโพ ภชิตพฺโพ ปยิรุปาสิตพฺโพ อปิ ปนุชฺชมาเนนปีติ.}

\csromangathameasure{ปิโย ครุ ภาวนีโย, วตฺตา จ วจนกฺขโม. \\ คมฺภีรญฺจ กถํ กตฺตา, โน จฏฺฐาเน นิโยชโก. \\ ยมฺหิ เอตานิ ฐานานิ, สํวิชฺชนฺตีธ ปุคฺคเล. \\ โส มิตฺโต มิตฺตกาเมน, อตฺถกามานุกมฺปโต. \\ อปิ นาสิยมาเนน, ภชิตพฺโพ ตถาวิโธติ.ฉฏฺฐํ.}
\csromangathagroup{\csromangathastanza{ปิโย ครุ ภาวนีโย, วตฺตา จ วจนกฺขโม. \\ คมฺภีรญฺจ กถํ กตฺตา, โน จฏฺฐาเน นิโยชโก\footnote{นิโยชเย (สี, สฺยา)}.}\csromangathastanza{ยมฺหิ เอตานิ ฐานานิ, สํวิชฺชนฺตีธ ปุคฺคเล. \\ โส มิตฺโต มิตฺตกาเมน, อตฺถกามานุกมฺปโต. \\ อปิ นาสิยมาเนน, ภชิตพฺโพ ตถาวิโธติ.\csromanspacer{}\csromanspacer{}ฉฏฺฐํ.}}

\csromanheader{2}{7. ปฐมปฏิสมฺภิทาสุตฺต}
\proseitem{38}{สตฺตหิ ภิกฺขเว ธมฺเมหิ สมนฺนาคโต ภิกฺขุ นจิรสฺเสว จตสฺโส ปฏิสมฺภิทา สยํ อภิญฺญา สจฺฉิกตฺวา อุปสมฺปชฺช วิหเรยฺย.\csromanspacer{} กตเมหิ สตฺตหิ? อิธ ภิกฺขเว ภิกฺขุ “อิทํ เม เจตโส ลีนตฺตนฺ”ติ ยถาภูตํ ปชานาติ.\csromanspacer{} อชฺฌตฺตํ สํขิตฺตํ วา จิตฺตํ “อชฺฌตฺตํ เม สํขิตฺตํ จิตฺตนฺ”ติ ยถาภูตํ ปชานาติ.\csromanspacer{} พหิทฺธา วิกฺขิตฺตํ วา จิตฺตํ “พหิทฺธา เม วิกฺขิตฺตํ จิตฺตนฺ”ติ ยถาภูตํ ปชานาติ.\csromanspacer{} ตสฺส วิทิตา เวทนา อุปฺปชฺชนฺติ, วิทิตา อุปฏฺฐหนฺติ, วิทิตา อพฺภตฺถํ คจฺฉนฺติ.\csromanspacer{} วิทิตา สญฺญา อุปฺปชฺชนฺติ, วิทิตา อุปฏฺฐหนฺติ, วิทิตา อพฺภตฺถํ คจฺฉนฺติ.\csromanspacer{} วิทิตา วิตกฺกา อุปฺปชฺชนฺติ, วิทิตา อุปฏฺฐหนฺติ, วิทิตา อพฺภตฺถํ คจฺฉนฺติ.\csromanspacer{} สปฺปายาสปฺปาเยสุ โข ปนสฺส ธมฺเมสุ หีนปฺปณีเตสุ กณฺหสุกฺกสปฺปติภาเคสุ นิมิตฺตํ สุคฺคหิตํ โหติ สุมนสิกตํ สูปธาริตํ สุปฺปฏิวิทฺธํ ปญฺญาย.\csromanspacer{} อิเมหิ โข ภิกฺขเว สตฺตหิ ธมฺเมหิ สมนฺนาคโต ภิกฺขุ นจิรสฺเสว จตสฺโส ปฏิสมฺภิทา สยํ อภิญฺญา สจฺฉิกตฺวา อุปสมฺปชฺช วิหเรยฺยาติ.\csromanspacer{}\csromanspacer{}สตฺตมํ.}

\chapterheadpageread{4. เทวตาวคฺค \textbf{8. ทุติยปฏิสมฺภิทาสุตฺต}}
\csromanmatikaanchor{mtk.237}
\csromantocmarkref{subsection}{9-10. วสสุตฺตทฺวย}{mtk.237}
\bookmark[dest={mtk.237},level=2]{9-10. วสสุตฺตทฺวย}
\proseitem{39}{\csromanfolio{423}สตฺตหิ ภิกฺขเว ธมฺเมหิ สมนฺนาคโต สาริปุตฺโต จตสฺโส ปฏิสมฺภิทา สยํ อภิญฺญา สจฺฉิกตฺวา อุปสมฺปชฺช วิหรติ.\csromanspacer{} กตเมหิ สตฺตหิ? อิธ ภิกฺขเว สาริปุตฺโต “อิทํ เม เจตโส ลีนตฺตนฺ”ติ ยถาภูตํ ปชานาติ.\csromanspacer{} อชฺฌตฺตํ สํขิตฺตํ วา จิตฺตํ “อชฺฌตฺตํ เม สํขิตฺตํ จิตฺตนฺ”ติ ยถาภูตํ ปชานาติ.\csromanspacer{} พหิทฺธา วิกฺขิตฺตํ วา จิตฺตํ “พหิทฺธา เม วิกฺขิตฺตํ จิตฺตนฺ”ติ ยถาภูตํ ปชานาติ.\csromanspacer{} ตสฺส วิทิตา เวทนา อุปฺปชฺชนฺติ, วิทิตา อุปฏฺฐหนฺติ, วิทิตา อพฺภตฺถํ คจฺฉนฺติ.\csromanspacer{} วิทิตา สญฺญา ฯเปฯ วิตกฺกา อุปฺปชฺชนฺติ, วิทิตา อุปฏฺฐหนฺติ, วิทิตา อพฺภตฺถํ คจฺฉนฺติ.\csromanspacer{} สปฺปายา สปฺปาเยสุ โข ปนสฺส ธมฺเมสุ หีนปฺปณีเตสุ กณฺหสุกฺกสปฺปติภาเคสุ นิมิตฺตํ สุคฺคหิตํ สุมนสิกตํ สูปธาริตํ สุปฺปฏิวิทฺธํ ปญฺญาย.\csromanspacer{} อิเมหิ โข ภิกฺขเว สตฺตหิ ธมฺเมหิ สมนฺนาคโต สาริปุตฺโต จตสฺโส ปฏิสมฺภิทา สยํ อภิญฺญา สจฺฉิกตฺวา อุปสมฺปชฺช วิหรตีติ.\csromanspacer{}\csromanspacer{}อฏฺฐมํ.}

\csromanheader{2}{9. ปฐมวสสุตฺต}
\proseitem{40}{สตฺตหิ ภิกฺขเว ธมฺเมหิ สมนฺนาคโต ภิกฺขุ จิตฺตํ วเส\footnote{วสํ (ก)} วตฺเตติ, โน จ ภิกฺขุ จิตฺตสฺส วเสน วตฺตติ.\csromanspacer{} กตเมหิ สตฺตหิ? อิธ ภิกฺขเว ภิกฺขุ สมาธิกุสโล โหติ, สมาธิสฺส สมาปตฺติกุสโล โหติ, สมาธิสฺส ฐิติกุสโล โหติ, สมาธิสฺส วุฏฺฐานกุสโล โหติ, สมาธิสฺส กลฺยาณกุสโล โหติ, สมาธิสฺส โคจรกุสโล โหติ, สมาธิสฺส อภินีหารกุสโล โหติ.\csromanspacer{} อิเมหิ โข ภิกฺขเว สตฺตหิ ธมฺเมหิ สมนฺนาคโต ภิกฺขุ จิตฺตํ วเส วตฺเตติ, โน จ ภิกฺขุ จิตฺตสฺส วเสน วตฺตตีติ.\csromanspacer{}\csromanspacer{}นวมํ.}

\csromanheader{2}{10. ทุติยวสสุตฺต}
\proseitem{41}{สตฺตหิ ภิกฺขเว ธมฺเมหิ สมนฺนาคโต สาริปุตฺโต จิตฺตํ วเส วตฺเตติ, โน จ สาริปุตฺโต จิตฺตสฺส วเสน วตฺตติ.\csromanspacer{} กตเมหิ สตฺตหิ? อิธ ภิกฺขเว สาริปุตฺโต สมาธิกุสโล โหติ, สมาธิสฺส สมาปตฺติกุสโล, สมาธิสฺส ฐิติกุสโล, สมาธิสฺส วุฏฺฐานกุสโล, สมาธิสฺส กลฺยาณกุสโล, สมาธิสฺส โคจรกุสโล,}

\vspace{\tipitakaparskip}
\csromancenter{สตฺตกนิปาตปาฬิ สมาธิสฺส อภินีหารกุสโล โหติ.\csromanspacer{} อิเมหิ โข ภิกฺขเว สตฺตหิ ธมฺเมหิ สมนฺนาคโต สาริปุตฺโต จิตฺตํ วเส วตฺเตติ, โน จ สาริปุตฺโต จิตฺตสฺส วเสน วตฺตตีติ.\csromanspacer{}\csromanspacer{}ทสมํ.}
\csromanheader{2}{11. ปฐมนิทฺทสสุตฺต}
\csromanmatikaanchor{mtk.238}
\csromantocmarkref{subsection}{11-12. นิทฺทสสุตฺตทฺวย}{mtk.238}
\bookmark[dest={mtk.238},level=2]{11-12. นิทฺทสสุตฺตทฺวย}
\proseitem{42}{\csromanfolio{424}อถ โข อายสฺมา สาริปุตฺโต ปุพฺพณฺหสมยํ นิวาเสตฺวา ปตฺตจีวรมาทาย สาวตฺถิํ ปิณฺฑาย ปาวิสิ.\csromanspacer{} อถ โข อายสฺมโต สาริปุตฺตสฺส เอตทโหสิ “อติปฺปโค โข ตาว สาวตฺถิยํ ปิณฺฑาย จริตุํ.\csromanspacer{} ยํนูนาหํ เยน อญฺญติตฺถิยานํ ปริพฺพาชกานํ อาราโม เตนุปสงฺกเมยฺยนฺ”ติ.\csromanspacer{} อถ โข อายสฺมา สาริปุตฺโต เยน อญฺญติตฺถิยานํ ปริพฺพาชกานํ อาราโม เตนุปสงฺกมิ, อุปสงฺกมิตฺวา เตหิ อญฺญติตฺถิเยหิ ปริพฺพาชเกหิ สทฺธิํ สมฺโมทิ, สมฺโมทนียํ กถํ สารณียํ วีติสาเรตฺวา เอกมนฺตํ นิสีทิ.\csromanspacer{} เตน โข ปน สมเยน เตสํ อญฺญติตฺถิยานํ ปริพฺพาชกานํ สนฺนิสินฺนานํ สนฺนิปติตานํ อยมนฺตรากถา อุทปาทิ “โย หิ โกจิ อาวุโส ทฺวาทสวสฺสานิ ปริปุณฺณํ ปริสุทฺธํ พฺราหฺมจริยํ จรติ ‘นิทฺทโส ภิกฺขู’ติ อลํ วจนายา”ติ.}

\prose{อถ โข อายสฺมา สาริปุตฺโต เตสํ อญฺญติตฺถิยานํ ปริพฺพาชกานํ ภาสิตํ เนว อภินนฺทิ นปฺปฏิกฺโกสิ, อนภินนฺทิตฺวา อปฺปฏิกฺโกสิตฺวา อุฏฺฐายาสนา ปกฺกามิ “ภควโต สนฺติเก เอตสฺส ภาสิตสฺส อตฺถํ อาชานิสฺสามี”ติ.\csromanspacer{} อถ โข อายสฺมา สาริปุตฺโต สาวตฺถิยํ ปิณฺฑาย จริตฺวา ปจฺฉาภตฺตํ ปิณฺฑปาตปฏิกฺกนฺโต เยน ภควา เตนุปสงฺกมิ, อุปสงฺกมิตฺวา ภควนฺตํ อภิวาเทตฺวา เอกมนฺตํ นิสีทิ, เอกมนฺตํ นิสินฺโน โข อายสฺมา สาริปุตฺโต ภควนฺตํ เอตทโวจ\csromandash{}}

\prose{อิธาหํ ภนฺเต ปุพฺพณฺหสมยํ นิวาเสตฺวา ปตฺตจีวรมาทาย สาวตฺถิํ ปิณฺฑาย ปาวิสิํ.\csromanspacer{} ตสฺส มยฺหํ ภนฺเต เอตทโหสิ “อติปฺปโค โข ตาว สาวตฺถิยํ ปิณฺฑาย จริตุํ, ยํนูนาหํ เยน อญฺญติตฺถิยานํ ปริพฺพาชกานํ อาราโม เตนุปสงฺกเมยฺยนฺ”ติ.\csromanspacer{} อถ ขฺวาหํ ภนฺเต เยน อญฺญติตฺถิยานํ ปริพฺพาชกานํ อาราโม เตนุปสงฺกมิํ, อุปสงฺกมิตฺวา เตหิ อญฺญติตฺถิเยหิ ปริพฺพาชเกหิ สทฺธิํ สมฺโมทิํ, สมฺโมทนียํ กถํ สารณียํ วีติสาเรตฺวา เอกมนฺตํ นิสีทิํ.\csromanspacer{} เตน โข ปน ภนฺเต สมเยน เตสํ อญฺญติตฺถิยานํ ปริพฺพาชกานํ สนฺนิสินฺนานํ สนฺนิปติตานํ อยมนฺตรากถา}

\vspace{\tipitakaparskip}
\csromancenter{เทวตาวคฺค อุทาปาทิ “โย หิ โกจิ อาวุโส ทฺวาทสวสฺสานิ ปริปุณฺณํ ปริสุทฺธํ พฺรหฺมจริยํ จรติ ‘นิทฺทโส ภิกฺขู’ติ อลํ วจนายา”ติ.\csromanspacer{} อถ ขฺวาหํ ภนฺเต เตสํ อญฺญติตฺถิยานํ ปริพฺพาชกานํ ภาสิตํ เนว อภินนฺทิํ นปฺปฏิกฺโกสิํ, อนภินนฺทิตฺวา อปฺปฏิกฺโกสิตฺวา อุฏฺฐายาสนา ปกฺกมิํ\footnote{ปกฺกามิํ (สี, สฺยา)} “ภควโต สนฺติเก เอตสฺส อตฺถํ อาชานิสฺสามี”ติ.\csromanspacer{} สกฺกา นุ โข ภนฺเต อิมสฺมิํ ธมฺมวินเย เกวลํ วสฺสคณนมตฺเตน นิทฺทโส ภิกฺขุ ปญฺญาเปตุนฺติ.}
\prose{\csromanfolio{425}น โข สาริปุตฺต สกฺกา อิมสฺมิํ ธมฺมวินเย เกวลํ วสฺสคณนมตฺเตน นิทฺทโส ภิกฺขุ ปญฺญาเปตุํ.\csromanspacer{} สตฺต โข อิมานิ สาริปุตฺต นิทฺทสวตฺถูนิ มยา สยํ อภิญฺญา สจฺฉิกตฺวา ปเวทิตานิ.}

\prose{\csromansymbolfootnote{*}{เหฏฺฐา 407; ที 3. 209 ปิฏฺเฐสุปิ.}กตมานิ สตฺต? อิธ สาริปุตฺต ภิกฺขุ สิกฺขาสมาทาเน ติพฺพจฺฉนฺโท โหติ, อายติํ จ สิกฺขาสมาทาเน อวิคตเปโม.\csromanspacer{} ธมฺมนิสนฺติยา ติพฺพจฺฉนฺโท โหติ, อายติํ จ ธมฺมนิสนฺติยา อวิคตเปโม.\csromanspacer{} อิจฺฉาวินเย ติพฺพจฺฉนฺโท โหติ, อายติํ จ อิจฺฉาวินเย อวิคตเปโม.\csromanspacer{} ปฏิสลฺลาเน ติพฺพจฺฉนฺโท โหติ, อายติํ จ ปฏิสลฺลาเน อวิคตเปโม.\csromanspacer{} วีริยารมฺเภ ติพฺพจฺฉนฺโท โหติ, อายติํ จ วีริยารมฺเภ อวิคตเปโม.\csromanspacer{} สติเนปกฺเก ติพฺพจฺฉนฺโท โหติ, อายติํ จ สติเนปกฺเก อวิคตเปโม.\csromanspacer{} ทิฏฺฐิปฏิเวเธ ติพฺพจฺฉนฺโท โหติ, อายติํ จ ทิฏฺฐิปฏิเวเธ อวิคตเปโม.\csromanspacer{} อิมานิ โข สาริปุตฺต สตฺต นิทฺทสวตฺถูนิ มยา สยํ อภิญฺญา สจฺฉิกตฺวา ปเวทิตานิ.\csromanspacer{} อิเมหิ โข สาริปุตฺต สตฺตหิ นิทฺทสวตฺถูหิ สมนฺนาคโต ภิกฺขุ ทฺวาทส เจปิ วสฺสานิ ปริปุณฺณํ ปริสุทฺธํ พฺรหฺมจริยํ จรติ “นิทฺทโส ภิกฺขู”ติ อลํ วจนาย, จตุพฺพีสติ เจปิ วสฺสานิ ปริปุณฺณํ ปริสุทฺธํ พฺรหฺมจริยํ จรติ “นิทฺทโส ภิกฺขู”ติ อลํ วจนาย, ฉตฺติํสติ เจปิ วสฺสานิ ปริปุณฺณํ ปริสุทฺธํ พฺรหฺมจริยํ จรติ “นิทฺทโส ภิกฺขู”ติ อลํ วจนาย, อฏฺฐจตฺตารีสญฺเจปิ วสฺสานิ ปริปุณฺณํ ปริสุทฺธํ พฺรหฺมจริยํ จรติ “นิทฺทโส ภิกฺขู”ติ อลํ วจนายาติ.\csromanspacer{}\csromanspacer{}เอกาทสมํ.}

\csromanheader{2}{12. ทุติยนิทฺทสสุตฺต}
\proseitem{43}{เอวํ เม สุตํ\csromandash{}เอกํ สมยํ ภควา โกสมฺพิยํ วิหรติ โฆสิตาราเม.\csromanspacer{} อถ โข อายสฺมา อานนฺโท ปุพฺพณฺหสมยํ นิวาเสตฺวา}

\vspace{\tipitakaparskip}
\csromancenter{สตฺตกนิปาตปาฬิ ปตฺตจีวรมาทาย โกสมฺพิํ ปิณฺฑาย ปาวิสิ.\csromanspacer{} อถ โข อายสฺมโต อานนฺทสฺส เอตทโหสิ “อติปฺปโค โข ตาว โกสมฺพิยํ ปิณฺฑาย จริตุํ, ยํนูนาหํ เยน อญฺญติตฺถิยานํ ปริพฺพาชกานํ อาราโม เตนุปสงฺกเมยฺยนฺ”ติ.\csromanspacer{} อถ โข อายสฺมา อานนฺโท เยน อญฺญติตฺถิยานํ ปริพฺพาชกานํ อาราโม เตนุปสงฺกมิ, อุปสงฺกมิตฺวา เตหิ อญฺญติตฺถิเยหิ ปริพฺพาชเกหิ สทฺธิํ สมฺโมทิ, สมฺโมทนียํ กถํ สารณียํ วีติสาเรตฺวา เอกมนฺตํ นิสีทิ.}
\prose{\csromanfolio{426}เตน โข ปน สมเยน เตสํ อญฺญติตฺถิยานํ ปริพฺพาชกานํ สนฺนิสินฺนานํ สนฺนิปติตานํ อยมนฺตรากถา อุทปาทิ “โย หิ โกจิ อาวุโส ทฺวาทส วสฺสานิ ปริปุณฺณํ ปริสุทฺธํ พฺรหฺมจริยํ จรติ ‘นิทฺทโส ภิกฺขู’ติ อลํ วจนายา”ติ.}

\prose{อถ โข อายสฺมา อานนฺโท เตสํ อญฺญติตฺถิยานํ ปริพฺพาชกานํ ภาสิตํ เนว อภินนฺทิ นปฺปฏิกฺโกสิ, อนภินนฺทิตฺวา อปฺปฏิกฺโกสิตฺวา อุฏฺฐายาสนา ปกฺกามิ “ภควโต สนฺติเก เอตสฺส ภาสิตสฺส อตฺถํ อาชานิสฺสามี”ติ.\csromanspacer{} อถ โข อายสฺมา อานนฺโท โกสมฺพิยํ ปิณฺฑาย จริตฺวา ปจฺฉาภตฺตํ ปิณฺฑปาตปฏิกฺกนฺโต เยน ภควา เตนุปสงฺกมิ, อุปสงฺกมิตฺวา ภควนฺตํ อภิวาเทตฺวา เอกมนฺตํ นิสีทิ, เอกมนฺตํ นิสินฺโน โข อายสฺมา อานนฺโท ภควนฺตํ เอตทโวจ\csromandash{}}

\prose{อิธาหํ ภนฺเต ปุพฺพณฺหสมยํ นิวาเสตฺวา ปตฺตจีวรมาทาย โกสมฺพิํ ปิณฺฑาย ปาวิสิํ.\csromanspacer{} ตสฺส มยฺหํ ภนฺเต เอตทโหสิ “อติปฺปโค โข ตาว โกสมฺพิยํ ปิณฺฑาย จริตุํ, ยํนูนาหํ เยน อญฺญติตฺถิยานํ ปริพฺพาชกานํ อาราโม เตนุปสงฺกเมยฺยนฺ”ติ ฯเปฯ เตหิ สทฺธิํ สมฺโมทิํ, สมฺโมทนียํ กถํ สารณียํ วีติสาเรตฺวา เอกมนฺตํ นิสีทิํ.}

\prose{เตน โข ปน ภนฺเต สมเยน เตสํ อญฺญติตฺถิยานํ ปริพฺพาชกานํ สนฺนิสินฺนานํ สนฺนิปติตานํ อยมนฺตรากถา อุทปาทิ “โย หิ โกจิ อาวุโส ทฺวาทสวสฺสานิ ปริปุณฺณํ ปริสุทฺธํ พฺรหฺมจริยํ จรติ ‘นิทฺทโส ภิกฺขู’ติ อลํ วจนายา”ติ.\csromanspacer{} อถ ขฺวาหํ ภนฺเต เตสํ อญฺญติตฺถิยานํ ปริพฺพาชกานํ ภาสิตํ เนว อภินนฺทิํ นปฺปฏิกฺโกสิํ, อนภินนฺทิตฺวา อปฺปฏิกฺโกสิตฺวา อุฏฺฐายาสนา ปกฺกมิํ “ภควโต สนฺติเก เอตสฺส}

\vspace{\tipitakaparskip}
\csromancenter{มหายญฺญวคฺค ภาสิตสฺส อตฺถํ อาชานิสฺสามี”ติ.\csromanspacer{} สกฺกา นุ โข ภนฺเต อิมสฺมิํ ธมฺมวินเย เกวลํ วสฺสคณนมตฺเตน นิทฺทโส ภิกฺขุ ปญฺญาเปตุนฺติ.}
\prose{\csromanfolio{427}น โข อานนฺท สกฺกา อิมสฺมิํ ธมฺมวินเย เกวลํ วสฺสคณนมตฺเตน นิทฺทโส ภิกฺขุ ปญฺญาเปตุํ.\csromanspacer{} สตฺต โข อิมานิ อานนฺท นิทฺทสวตฺถูนิ มยา สยํ อภิญฺญา สจฺฉิกตฺวา ปเวทิตานิ.}

\prose{กตมานิ สตฺต? อิธานนฺท ภิกฺขุ สทฺโธ โหติ, หิรีมา โหติ, โอตฺตปฺปี โหติ, พหุสฺสุโต โหติ, อารทฺธวีริโย โหติ, สติมา โหติ, ปญฺญวา โหติ.\csromanspacer{} อิมานิ โข อานนฺท สตฺต นิทฺทสวตฺถูนิ มยา สยํ อภิญฺญา สจฺฉิกตฺวา ปเวทิตานิ.\csromanspacer{} อิเมหิ โข อานนฺท สตฺตหิ นิทฺทสวตฺถูหิ สมนฺนาคโต ภิกฺขุ ทฺวาทส เจปิ วสฺสานิ ปริปุณฺณํ ปริสุทฺธํ พฺรหฺมจริยํ จรติ “นิทฺทโส ภิกฺขู”ติ อลํ วจนาย, จตุพฺพีสติ เจปิ วสฺสานิ ปริปุณฺณํ ปริสุทฺธํ พฺรหฺมจริยํ จรติ “นิทฺทโส ภิกฺขู”ติ อลํ วจนาย, ฉตฺติํสติ เจปิ วสฺสานิ ปริปุณฺณํ ปริสุทฺธํ พฺรหฺมจริยํ จรติ “นิทฺทโส ภิกฺขู”ติ อลํ วจนาย, อฏฺฐจตฺตารีสญฺเจปิ วสฺสานิ ปริปุณฺณํ ปริสุทฺธํ พฺรหฺมจริยํ จรติ “นิทฺทโส ภิกฺขู”ติ อลํ วจนายาติ.\csromanspacer{}\csromanspacer{}ทฺวาทสมํ.\csromanspacer{} เทวตาวคฺโค จตุตฺโถ.}

\titlehead{\textbf{ตสฺสุทฺทานํ}}
\csromangathameasure{อปฺปมาโท หิรี เจว, เทฺว สุวจา ทุเว มิตฺตา. \\ เทฺว ปฏิสมฺภิทา เทฺว วสา, ทุเว นิทฺทสวตฺถุนาติ.}
\csromangathagroup{\csromangathastanza{อปฺปมาโท หิรี เจว, เทฺว สุวจา ทุเว มิตฺตา. \\ เทฺว ปฏิสมฺภิทา เทฺว วสา, ทุเว นิทฺทสวตฺถุนาติ.}}

\csromanmatikaanchor{mtk.239}
\csromantocmarknopage{chapter}{5. มหายญฺญวคฺค}{mtk.239}
\bookmark[dest={mtk.239},level=0]{5. มหายญฺญวคฺค}
\chapterhead{5. มหายญฺญวคฺค}
\csromanmatikaanchor{mtk.240}
\csromantocmarkref{subsection}{1. สตฺตวิญฺญาณฏฺฐิติสุตฺต}{mtk.240}
\bookmark[dest={mtk.240},level=2]{1. สตฺตวิญฺญาณฏฺฐิติสุตฺต}
\csromanheader{2}{1. สตฺตวิญฺญาณฏฺฐิติสุตฺต}
\proseitem{44}{\csromansymbolfootnote{*}{ที 3. 209, 241; ขุ 8. 167 ปิฏฺเฐสุปิ.}สตฺติมา ภิกฺขเว วิญฺญาณฏฺฐิติโย.\csromanspacer{} กตมา สตฺต? สนฺติ ภิกฺขเว สตฺตา นานตฺตกายา นานตฺตสญฺญิโน.\csromanspacer{} เสยฺยถาปิ มนุสฺสา เอกจฺเจ จ เทวา เอกจฺเจ จ วินิปาติกา.\csromanspacer{} อยํ ปฐมา วิญฺญาณฏฺฐิติ.}

\prose{สนฺติ ภิกฺขเว สตฺตา นานตฺตกายา เอกตฺตสญฺญิโน.\csromanspacer{} เสยฺยถาปิ เทวา พฺรหฺมกายิกา ปฐมาภินิพฺพตฺตา.\csromanspacer{} อยํ ทุติยา วิญฺญาณฏฺฐิติ.}

\vspace{\tipitakaparskip}
\csromancenter{สตฺตกนิปาตปาฬิ สนฺติ ภิกฺขเว สตฺตา เอกตฺตกายา นานตฺตสญฺญิโน.\csromanspacer{} เสยฺยถาปิ เทวา อาภสฺสรา.\csromanspacer{} อยํ ตติยา วิญฺญาณฏฺฐิติ.}
\csromanmatikaanchor{mtk.241}
\csromanmatikaanchor{mtk.242}
\csromantocmarkref{subsection}{2. สมาธิปริกฺขารสุตฺต}{mtk.241}
\bookmark[dest={mtk.241},level=2]{2. สมาธิปริกฺขารสุตฺต}
\csromantocmarkref{subsection}{3-4. อคฺคิสุตฺตทฺวย}{mtk.242}
\bookmark[dest={mtk.242},level=2]{3-4. อคฺคิสุตฺตทฺวย}
\prosecont{\csromanfolio{428}สนฺติ ภิกฺขเว สตฺตา เอกตฺตกายา เอกตฺตสญฺญิโน.\csromanspacer{} เสยฺยถาปิ เทวา สุภกิณฺหา.\csromanspacer{} อยํ จตุตฺถา วิญฺญาณฏฺฐิติ.}

\prose{สนฺติ ภิกฺขเว สตฺตา สพฺพโส รูปสญฺญานํ สมติกฺกมา ปฏิฆสญฺญานํ อตฺถงฺคมา นานตฺตสญฺญานํ อมนสิการา “อนนฺโต อากาโส”ติ อากาสานญฺจายตนูปคา.\csromanspacer{} อยํ ปญฺจมา วิญฺญาณฏฺฐิติ.}

\prose{สนฺติ ภิกฺขเว สตฺตา สพฺพโส อากาสานญฺจายตนํ สมติกฺกมฺม “อนนฺตํ วิญฺญาณนฺ”ติ วิญฺญาณญฺจายตนูปคา.\csromanspacer{} อยํ ฉฏฺฐา วิญฺญาณฏฺฐิติ.}

\prose{สนฺติ ภิกฺขเว สตฺตา สพฺพโส วิญฺญาณญฺจายตนํ สมติกฺกมฺม “นตฺถิ กิญฺจี”ติ อากิญฺจญฺญายตนูปคา.\csromanspacer{} อยํ สตฺตมา วิญฺญาณฏฺฐิติ.\csromanspacer{} อิมา โข ภิกฺขเว สตฺต วิญฺญาณฏฺฐิติโยติ.\csromanspacer{}\csromanspacer{}ปฐมํ.}

\csromanheader{2}{2. สมาธิปริกฺขารสุตฺต}
\proseitem{45}{สตฺติเม ภิกฺขเว สมาธิปริกฺขารา.\csromanspacer{} กตเม สตฺต? สมฺมาทิฏฺฐิ สมฺมาสงฺกปฺโป สมฺมาวาจา สมฺมากมฺมนฺโต สมฺมา-อาชีโว สมฺมาวายาโม สมฺมาสติ.\csromanspacer{} ยา โข ภิกฺขเว อิเมหิ สตฺตหงฺเคหิ จิตฺตสฺเสกคฺคตา ปริกฺขตา\footnote{‘สปริกฺขารตา’ติปิ, ‘สปริกฺขตา’ติปิ (สํ 3. 17 ปิฏฺเฐ)}, อยํ วุจฺจติ ภิกฺขเว “อริโย สมฺมาสมาธิ\footnote{สมาธิ (สฺยา)} ส-อุปนิโส” อิติปิ “สปริกฺขาโร” อิติปีติ.\csromanspacer{}\csromanspacer{}ทุติยํ.}

\csromanheader{2}{3. ปฐม อคฺคิสุตฺต}
\proseitem{46}{สตฺติเม ภิกฺขเว อคฺคี.\csromanspacer{} กตเม สตฺต? ราคคฺคิ โทสคฺคิ โมหคฺคิ อาหุเนยฺยคฺคิ คหปตคฺคิ ทกฺขิเณยฺยคฺคิ กฏฺฐคฺคิ.\csromanspacer{} อิเม โข ภิกฺขเว สตฺต อคฺคีติ.\csromanspacer{}\csromanspacer{}ตติยํ.}

\csromanheader{2}{4. ทุติย อคฺคิสุตฺต}
\proseitem{47}{เตน โข ปน สมเยน อุคฺคตสรีรสฺส พฺราหฺมณสฺส มหายญฺโญ อุปกฺขโฏ โหติ.\csromanspacer{} ปญฺจ อุสภสตานิ ถูณูปนีตานิ โหนฺติ ยญฺญตฺถาย, ปญฺจ วจฺฉตรสตานิ ถูณูปนีตานิ โหนฺติ}

\vspace{\tipitakaparskip}
\csromancenter{มหายญฺญวคฺค ยญฺญตฺถาย, ปญฺจ วจฺฉตริสตานิ ถูณูปนีตานิ โหนฺติ ยญฺญตฺถาย, ปญฺจ อชสตานิ ถูณูปนีตานิ โหนฺติ ยญฺญตฺถาย, ปญฺจ อุรพฺภสตานิ ถูณูปนีตานิ โหนฺติ ยญฺญตฺถาย.\csromanspacer{} อถ โข อุคฺคตสรีโร พฺราหฺมโณ เยน ภควา เตนุปสงฺกมิ, อุปสงฺกมิตฺวา ภควตา สทฺธิํ สมฺโมทิ, สมฺโมทนียํ กถํ สารณียํ วีติสาเรตฺวา เอกมนฺตํ นิสีทิ, เอกมนฺตํ นิสินฺโน โข อุคฺคตสรีโร พฺราหฺมโณ ภควนฺตํ เอตทโวจ\csromandash{}}
\prose{\csromanfolio{429}สุตํ เมตํ โภ โคตม อคฺคิสฺส อาทานํ\footnote{อาธานํ (สี, สฺยา)} ยูปสฺส อุสฺสาปนํ มหปฺผลํ โหติ มหานิสํสนฺติ.\csromanspacer{} มยาปิ โข เอตํ พฺราหฺมณ สุตํ อคฺคิสฺส อาทานํ ยูปสฺส อุสฺสาปนํ มหปฺผลํ โหติ มหานิสํสนฺติ.\csromanspacer{} ทุติยมฺปิ โข อุคฺคตสรีโร พฺราหฺมโณ ฯเปฯ.\csromanspacer{} ตติยมฺปิ โข อุคฺคตสรีโร พฺราหฺมโณ ภควนฺตํ เอตทโวจ “สุตํ เมตํ โภ โคตม อคฺคิสฺส อาทานํ ยูปสฺส อุสฺสาปนํ มหปฺผลํ โหติ มหานิสํสนฺ”ติ.\csromanspacer{} มยาปิ โข เอตํ พฺราหฺมณ สุตํ “อคฺคิสฺส อาทานํ ยูปสฺส อุสฺสาปนํ มหปฺผลํ โหติ มหานิสํสนฺ”ติ.\csromanspacer{} ตยิทํ โภ โคตม สเมติ โภโต เจว โคตมสฺส อมฺหากญฺจ, ยทิทํ สพฺเพน สพฺพํ.}

\prose{เอวํ วุตฺเต อายสฺมา อานนฺโท อุคฺคตสรีรํ พฺราหฺมณํ เอตทโวจ “น โข พฺราหฺมณ ตถาคตา เอวํ ปุจฺฉิตพฺพา สุตํ เมตํ โภ โคตม อคฺคิสฺส อาทานํ ยูปสฺส อุสฺสาปนํ มหปฺผลํ โหติ มหานิสํสนฺ”ติ.\csromanspacer{} เอวํ โข พฺราหฺมณ ตถาคตา ปุจฺฉิตพฺพา “อหํ หิ ภนฺเต อคฺคิํ\footnote{อคฺคิสฺส (ก)} อาทาตุกาโม\footnote{อาธาตุกาโม (สี, สฺยา)} ยูปํ อุสฺสาเปตุกาโม, โอวทตุ มํ ภนฺเต ภควา, อนุสาสตุ มํ ภนฺเต ภควา, ยํ มม อสฺส ทีฆรตฺตํ หิตาย สุขายา”ติ.}

\prose{อถ โข อุคฺคตสรีโร พฺราหฺมโณ ภควนฺตํ เอตทโวจ “อหํ หิ โภ โคตม อคฺคิํ อาทาตุกาโม ยูปํ อุสฺสาเปตุกาโม, โอวทตุ มํ ภวํ โคตโม, อนุสาสตุ มํ ภวํ โคตโม, ยํ มม อสฺส ทีฆรตฺตํ หิตาย สุขายา”ติ.}

\prose{อคฺคิํ พฺราหฺมณ อาเทนฺโต\footnote{อาเธนฺโต (สี, สฺยา)} ยูปํ อุสฺสาเปนฺโต ปุพฺเพว ยญฺญา ตีณิ สตฺถานิ อุสฺสาเปติ อกุสลานิ ทุกฺขุทฺรยานิ\footnote{ทุกฺขุทฺทยานิ (สี)} ทุกฺขวิปากานิ.\csromanspacer{} กตมานิ}

\vspace{\tipitakaparskip}
\csromancenter{สตฺตกนิปาตปาฬิ ตีณิ? กายสตฺถํ วจีสตฺถํ มโนสตฺถํ, อคฺคิํ พฺราหฺมณ อาเทนฺโต ยูปํ อุสฺสาเปนฺโต ปุพฺเพว ยญฺญา เอวํ จิตฺตํ อุปฺปาเทสิ “เอตฺตกา อุสภา หญฺญนฺตุ ยญฺญตฺถาย, เอตฺตกา วจฺฉตรา หญฺญนฺตุ ยญฺญตฺถาย, เอตฺตกา วจฺฉตริโย หญฺญนฺตุ ยญฺญตฺถาย, เอตฺตกา อชา หญฺญนฺตุ ยญฺญตฺถาย, เอตฺตกา อุรพฺภา หญฺญนฺตุ ยญฺญตฺถายา”ติ.\csromanspacer{} โส “ปุญฺญํ กโรมี”ติ อปุญฺญํ กโรติ, “กุสลํ กโรมี”ติ อกุสลํ กโรติ, “สุคติยา มคฺคํ ปริเยสามี”ติ ทุคฺคติยา มคฺคํ ปริเยสติ.\csromanspacer{} อคฺคิํ พฺราหฺมณ อาเทนฺโต ยูปํ อุสฺสาเปนฺโต ปุพฺเพว ยญฺญา อิทํ ปฐมํ มโนสตฺถํ อุสฺสาเปติ อกุสลํ ทุกฺขุทฺรยํ ทุกฺขวิปากํ.}
\prose{\csromanfolio{430}ปุน จปรํ พฺราหฺมณ อคฺคิํ อาเทนฺโต ยูปํ อุสฺสาเปนฺโต ปุพฺเพว ยญฺญา เอวํ วาจํ ภาสติ “เอตฺตกา อุสภา หญฺญนฺตุ ยญฺญตฺถาย, เอตฺตกา วจฺฉตรา หญฺญนฺตุ ยญฺญตฺถาย, เอตฺตกา วจฺฉตริโย หญฺญนฺตุ ยญฺญตฺถาย, เอตฺตกา อชา หญฺญนฺตุ ยญฺญตฺถาย, เอตฺตกา อุรพฺภา หญฺญนฺตุ ยญฺญตฺถายา”ติ.\csromanspacer{} โส “ปุญฺญํ กโรมี”ติ อปุญฺญํ กโรติ, “กุสลํ กโรมี”ติ อกุสลํ กโรติ, “สุคติยา มคฺคํ ปริเยสามี”ติ ทุคฺคติยา มคฺคํ ปริเยสติ.\csromanspacer{} อคฺคิํ พฺราหฺมณ อาเทนฺโต ยูปํ อุสฺสาเปนฺโต ปุพฺเพว ยญฺญา อิทํ ทุติยํ วจีสตฺถํ อุสฺสาเปติ อกุสลํ ทุกฺขุทฺรยํ ทุกฺขวิปากํ.}

\prose{ปุน จปรํ พฺราหฺมณ อคฺคิํ อาเทนฺโต ยูปํ อุสฺสาเปนฺโต ปุพฺเพว ยญฺญา สยํ ปฐมํ สมารมฺภติ\footnote{สมารพฺภติ (สี, ก), สมารภติ (สฺยา)} อุสภา หนฺตุํ\footnote{หญฺญนฺตุ (สี, สฺยา)} ยญฺญตฺถาย, สยํ ปฐมํ สมารมฺภติ วจฺฉตรา หนฺตุํ ยญฺญตฺถาย, สยํ ปฐมํ สมารมฺภติ วจฺฉตริโย หนฺตุํ ยญฺญตฺถาย, สยํ ปฐมํ สมารมฺภติ อชา หนฺตุํ ยญฺญตฺถาย, สยํ ปฐมํ สมารมฺภติ อุรพฺภา หนฺตุํ ยญฺญตฺถาย\footnote{หญฺญนฺตุ ยญฺญตฺถายาติ (สี, สฺยา)}.\csromanspacer{} โส “ปุญฺญํ กโรมี”ติ อปุญฺญํ กโรติ, “กุสลํ กโรมี”ติ อกุสลํ กโรติ, “สุคติยา มคฺคํ ปริเยสามี”ติ ทุคฺคติยา มคฺคํ ปริเยสติ.\csromanspacer{} อคฺคิํ พฺราหฺมณ อาเทนฺโต ยูปํ อุสฺสาเปนฺโต ปุพฺเพว ยญฺญา อิทํ ตติยํ กายสตฺถํ อุสฺสาเปติ อกุสลํ ทุกฺขุทฺรยํ ทุกฺขวิปากํ.\csromanspacer{} อคฺคิํ พฺราหฺมณ อาเทนฺโต ยูปํ อุสฺสาเปนฺโต ปุพฺเพว ยญฺญา อิมานิ ตีณิ สตฺถานิ อุสฺสาเปติ อกุสลานิ ทุกฺขุทฺรยานิ ทุกฺขวิปากานิ.}

\chapterheadpageread{5. มหายญฺญวคฺค}
\prose{\csromanfolio{431}ตโยเม พฺราหฺมณ อคฺคี ปหาตพฺพา ปริวชฺเชตพฺพา น เสวิตพฺพา.\csromanspacer{} กตเม ตโย? ราคคฺคิ โทสคฺคิ โมหคฺคิ.}

\prose{กสฺมา จายํ พฺราหฺมณ ราคคฺคิ ปหาตพฺโพ ปริวชฺเชตพฺโพ น เสวิตพฺโพ? รตฺโต โข พฺราหฺมณ ราเคน อภิภูโต ปริยาทินฺนจิตฺโต กาเยน ทุจฺจริตํ จรติ, วาจาย ทุจฺจริตํ จรติ, มนสา ทุจฺจริตํ จรติ.\csromanspacer{} โส กาเยน ทุจฺจริตํ จริตฺวา วาจาย ทุจฺจริตํ จริตฺวา มนสา ทุจฺจริตํ จริตฺวา กายสฺส เภทา ปรํ มรณา อปายํ ทุคฺคติํ วินิปาตํ นิรยํ อุปปชฺชติ.\csromanspacer{} ตสฺมายํ ราคคฺคิ ปหาตพฺโพ ปริวชฺเชตพฺโพ น เสวิตพฺโพ.}

\prose{กสฺมา จายํ พฺราหฺมณ โทสคฺคิ ปหาตพฺโพ ปริวชฺเชตพฺโพ น เสวิตพฺโพ? ทุฏฺโฐ โข พฺราหฺมณ โทเสน อภิภูโต ปริยาทินฺนจิตฺโต กาเยน ทุจฺจริตํ จรติ, วาจาย ทุจฺจริตํ จรติ, มนสา ทุจฺจริตํ จรติ.\csromanspacer{} โส กาเยน ทุจฺจริตํ จริตฺวา วาจาย ทุจฺจริตํ จริตฺวา มนสา ทุจฺจริตํ จริตฺวา กายสฺส เภทา ปรํ มรณา อปายํ ทุคฺคติํ วินิปาตํ นิรยํ อุปปชฺชติ.\csromanspacer{} ตสฺมายํ โทสคฺคิ ปหาตพฺโพ ปริวชฺเชตพฺโพ น เสวิตพฺโพ.}

\prose{กสฺมา จายํ พฺราหฺมณ โมหคฺคิ ปหาตพฺโพ ปริวชฺเชตพฺโพ น เสวิตพฺโพ? มูโฬฺห โข พฺราหฺมณ โมเหน อภิภูโต ปริยาทินฺนจิตฺโต กาเยน ทุจฺจริตํ จรติ, วาจาย ทุจฺจริตํ จรติ, มนสา ทุจฺจริตํ จรติ.\csromanspacer{} โส กาเยน ทุจฺจริตํ จริตฺวา วาจาย ทุจฺจริตํ จริตฺวา มนสา ทุจฺจริตํ จริตฺวา กายสฺส เภทา ปรํ มรณา อปายํ ทุคฺคติํ วินิปาตํ นิรยํ อุปปชฺชติ.\csromanspacer{} ตสฺมายํ โมหคฺคิ ปหาตพฺโพ ปริวชฺเชตพฺโพ น เสวิตพฺโพ.\csromanspacer{} อิเม โข ตโย พฺราหฺมณ อคฺคี ปหาตพฺพา ปริวชฺเชตพฺพา น เสวิตพฺพา.}

\prose{ตโย โข พฺราหฺมณ อคฺคี สกฺกตฺวา ครุํ กตฺวา มาเนตฺวา ปูเชตฺวา สมฺมา สุขํ ปริหาตพฺพา.\csromanspacer{} กตเม ตโย? อาหุเนยฺยคฺคิ คหปตคฺคิ ทกฺขิเณยฺยคฺคิ.}

\prose{กตโม จ พฺราหฺมณ อาหุเนยฺยคฺคิ? อิธ พฺราหฺมณ ยสฺส เต โหนฺติ มาตาติ วา ปิตาติ วา, อยํ วุจฺจติ พฺราหฺมณ อาหุเนยฺยคฺคิ.\csromanspacer{} ตํ กิสฺส}

\vspace{\tipitakaparskip}
\csromancenter{สตฺตกนิปาตปาฬิ เหตุ? อโต หยํ\footnote{อิโตปายํ (ก)} พฺราหฺมณ อาหุโต สมฺภูโต, ตสฺมายํ อาหุเนยฺยคฺคิ สกฺกตฺวา ครุํ กตฺวา มาเนตฺวา ปูเชตฺวา สมฺมา สุขํ ปริหาตพฺโพ.}
\csromanmatikaanchor{mtk.243}
\csromantocmarkref{subsection}{5-6. สญฺญาสุตฺตทฺวย}{mtk.243}
\bookmark[dest={mtk.243},level=2]{5-6. สญฺญาสุตฺตทฺวย}
\prose{\csromanfolio{432}กตโม จ พฺราหฺมณ คหปตคฺคิ? อิธ พฺราหฺมณ ยสฺส เต โหนฺติ ปุตฺตาติ วา ทาราติ วา ทาสาติ วา เปสฺสาติ วา กมฺมกราติ วา, อยํ วุจฺจติ พฺราหฺมณ คหปตคฺคิ, ตสฺมายํ คหปตคฺคิ สกฺกตฺวา ครุํ กตฺวา มาเนตฺวา ปูเชตฺวา สมฺมา สุขํ ปริหาตพฺโพ.}

\prose{กตโม จ พฺราหฺมณ ทกฺขิเณยฺยคฺคิ? อิธ พฺราหฺมณ เย เต สมณพฺราหฺมณา ปรปฺปวาทา ปฏิวิรตา ขนฺติโสรจฺเจ นิวิฏฺฐา เอกมตฺตานํ ทเมนฺติ, เอกมตฺตานํ สเมนฺติ, เอกมตฺตานํ ปรินิพฺพาเปนฺติ.\csromanspacer{} อยํ วุจฺจติ พฺราหฺมณ ทกฺขิเณยฺยคฺคิ, ตสฺมายํ ทกฺขิเณยฺยคฺคิ สกฺกตฺวา ครุํ กตฺวา มาเนตฺวา ปูเชตฺวา สมฺมา สุขํ ปริหาตพฺโพ.\csromanspacer{} อิเม โข พฺราหฺมณ ตโย อคฺคี สกฺกตฺวา ครุํ กตฺวา มาเนตฺวา ปูเชตฺวา สมฺมา สุขํ ปริหาตพฺพา.}

\prose{อยํ โข ปน พฺราหฺมณ กฏฺฐคฺคิ กาเลน กาลํ อุชฺชเลตพฺโพ, กาเลน กาลํ อชฺฌุเปกฺขิตพฺโพ, กาเลน กาลํ นิพฺพาเปตพฺโพ, กาเลน กาลํ นิกฺขิปิตพฺโพติ.}

\prose{เอวํ วุตฺเต อุคฺคตสรีโร พฺราหฺมโณ ภควนฺตํ เอตทโวจ “อภิกฺกนฺตํ โภ โคตม, อภิกฺกนฺตํ โภ โคตม ฯเปฯ อุปาสกํ มํ ภวํ โคตโม ธาเรตุ อชฺชตคฺเค ปาณุเปตํ สรณํ คตนฺติ.\csromanspacer{} เอสาหํ โภ โคตม ปญฺจ อุสภสตานิ มุญฺจามิ, ชีวิตํ เทมิ.\csromanspacer{} ปญฺจ วจฺฉตรสตานิ มุญฺจามิ, ชีวิตํ เทมิ.\csromanspacer{} ปญฺจ วจฺฉตริสตานิ มุญฺจามิ, ชีวิตํ เทมิ.\csromanspacer{} ปญฺจ อชสตานิ มุญฺจามิ, ชีวิตํ เทมิ.\csromanspacer{} ปญฺจ อุรพฺภสตานิ มุญฺจามิ, ชีวิตํ เทมิ.\csromanspacer{} หริตานิ เจว ติณานิ ขาทนฺตุ, สีตานิ จ ปานียานิ ปิวนฺตุ, สีโต จ เนสํ วาโต อุปวายตนฺติ\footnote{อุปวายตูติ (ก)}.\csromanspacer{}\csromanspacer{}จตุตฺถํ.}

\csromanheader{2}{5. ปฐมสญฺญาสุตฺต}
\proseitem{48}{สตฺติมา ภิกฺขเว สญฺญา ภาวิตา พหุลีกตา มหปฺผลา โหนฺติ มหานิสํสา อมโตคธา อมตปริโยสานา.}

\chapterheadpageread{5. มหายญฺญวคฺค}
\prose{\csromanfolio{433}กตมา สตฺต? อสุภสญฺญา มรณสญฺญา อาหาเร ปฏิกูลสญฺญา สพฺพโลเก อนภิรตสญฺญา อนิจฺจสญฺญา อนิจฺเจ ทุกฺขสญฺญา ทุกฺเข อนตฺตสญฺญา.\csromanspacer{} อิมา โข ภิกฺขเว สตฺต สญฺญา ภาวิตา พหุลีกตา มหปฺผลา โหนฺติ มหานิสํสา อมโตคธา อมตปริโยสานาติ.\csromanspacer{}\csromanspacer{}ปญฺจมํ.}

\csromanheader{2}{6. ทุติยสญฺญาสุตฺต}
\proseitem{49}{สตฺติมา ภิกฺขเว สญฺญา ภาวิตา พหุลีกตา มหปฺผลา โหนฺติ มหานิสํสา อมโตคธา อมตปริโยสานา.\csromanspacer{} กตมา สตฺต? อสุภสญฺญา มรณสญฺญา อาหาเร ปฏิกูลสญฺญา สพฺพโลเก อนภิรตสญฺญา อนิจฺจสญฺญา อนิจฺเจ ทุกฺขสญฺญา ทุกฺเข อนตฺตสญฺญา.\csromanspacer{} อิมา โข ภิกฺขเว สตฺต สญฺญา ภาวิตา พหุลีกตา มหปฺผลา โหนฺติ มหานิสํสา อมโตคธา อมตปริโยสานาติ.}

\prose{“อสุภสญฺญา ภิกฺขเว ภาวิตา พหุลีกตา มหปฺผลา โหติ มหานิสํสา อมโตคธา อมตปริโยสานา”ติ อิติ โข ปเนตํ วุตฺตํ, กิญฺเจตํ ปฏิจฺจ วุตฺตํ.\csromanspacer{} อสุภสญฺญาปริจิเตน ภิกฺขเว ภิกฺขุโน เจตสา พหุลํ วิหรโต เมถุนธมฺมสมาปตฺติยา จิตฺตํ ปติลียติ ปติกุฏติ ปติวตฺตติ น สมฺปสาริยติ, อุเปกฺขา วา ปาฏิกุลฺยตา วา สณฺฐาติ.\csromanspacer{} เสยฺยถาปิ ภิกฺขเว กุกฺกุฏปตฺตํ วา นฺหารุททฺทุลํ วา อคฺคิมฺหิ ปกฺขิตฺตํ ปติลียติ ปติกุฏติ ปติวตฺตติ น สมฺปสาริยติ.\csromanspacer{} เอวเมวํ โข ภิกฺขเว ภิกฺขุโน อสุภสญฺญาปริจิเตน เจตสา พหุลํ วิหรโต เมถุนธมฺมสมาปตฺติยา จิตฺตํ ปติลียติ ปติกุฏติ ปติวตฺตติ น สมฺปสาริยติ, อุเปกฺขา วา ปาฏิกุลฺยตา วา สณฺฐาติ.}

\prose{สเจ ภิกฺขเว ภิกฺขุโน อสุภสญฺญาปริจิเตน เจตสา พหุลํ วิหรโต เมถุนธมฺมสมาปตฺติยา จิตฺตํ อนุสนฺทหติ\footnote{อนุสณฺฐาติ (สี)}, อปฺปฏิกุลฺยตา สณฺฐาติ.\csromanspacer{} เวทิตพฺพเมตํ ภิกฺขเว ภิกฺขุนา “อภาวิตา เม อสุภสญฺญา, นตฺถิ เม ปุพฺเพนาปรํ วิเสโส, อปฺปตฺตํ เม ภาวนาพลนฺ”ติ, อิติห ตตฺถ สมฺปชาโน โหติ.\csromanspacer{} สเจ ปน ภิกฺขเว ภิกฺขุโน อสุภสญฺญาปริจิเตน เจตสา พหุลํ วิหรโต เมถุนธมฺมสมาปตฺติยา จิตฺตํ ปติลียติ ปติกุฏติ ปติวตฺตติ น สมฺปสาริยติ, อุเปกฺขา วา ปาฏิกุลฺยตา วา สณฺฐาติ.}

\vspace{\tipitakaparskip}
\csromancenter{สตฺตกนิปาตปาฬิ เวทิตพฺพเมตํ ภิกฺขเว ภิกฺขุนา “สุภาวิตา เม อสุภสญฺญา, อตฺถิ เม ปุพฺเพนาปรํ วิเสโส, ปตฺตํ เม ภาวนาพลนฺ”ติ, อิติห ตตฺถ สมฺปชาโน โหติ. “อสุภสญฺญา ภิกฺขเว ภาวิตา พหุลีกตา มหปฺผลา โหติ มหานิสํสา อมโตคธา อมตปริโยสานา”ติ อิติ ยํ ตํ วุตฺตํ, อิทเมตํ ปฏิจฺจ วุตฺตํ. (1)}
\prose{\csromanfolio{434}“มรณสญฺญา ภิกฺขเว ภาวิตา พหุลีกตา มหปฺผลา โหติ มหานิสํสา อมโตคธา อมตปริโยสานา”ติ อิติ โข ปเนตํ วุตฺตํ, กิญฺเจตํ ปฏิจฺจ วุตฺตํ.\csromanspacer{} มรณสญฺญาปริจิเตน ภิกฺขเว ภิกฺขุโน เจตสา พหุลํ วิหรโต ชีวิตนิกนฺติยา จิตฺตํ ปติลียติ ปติกุฏติ ปติวตฺตติ น สมฺปสาริยติ, อุเปกฺขา วา ปาฏิกุลฺยตา วา สณฺฐาติ.\csromanspacer{} เสยฺยถาปิ ภิกฺขเว กุกฺกุฏปตฺตํ วา นฺหารุททฺทุลํ วา อคฺคิมฺหิ ปกฺขิตฺตํ ปติลียติ ปติกุฏติ ปติวตฺตติ น สมฺปสาริยติ.\csromanspacer{} เอวเมวํ โข ภิกฺขเว ภิกฺขุโน มรณสญฺญาปริจิเตน เจตสา พหุลํ วิหรโต ชีวิตนิกนฺติยา จิตฺตํ ปติลียติ ปติกุฏติ ปติวตฺตติ น สมฺปสาริยติ, อุเปกฺขา วา ปาฏิกุลฺยตา วา สณฺฐาติ.}

\prose{สเจ ภิกฺขเว ภิกฺขุโน มรณสญฺญาปริจิเตน เจตสา พหุลํ วิหรโต ชีวิตนิกนฺติยา จิตฺตํ อนุสนฺทหติ, อปฺปฏิกุลฺยตา สณฺฐาติ.\csromanspacer{} เวทิตพฺพเมตํ ภิกฺขเว ภิกฺขุนา “อภาวิตา เม มรณสญฺญา, นตฺถิ เม ปุพฺเพนาปรํ วิเสโส, อปฺปตฺตํ เม ภาวนาพลนฺ”ติ, อิติห ตตฺถ สมฺปชาโน โหติ.\csromanspacer{} สเจ ปน ภิกฺขเว ภิกฺขุโน มรณสญฺญาปริจิเตน เจตสา พหุลํ วิหรโต ชีวิตนิกนฺติยา จิตฺตํ ปติลียติ ปติกุฏติ ปติวตฺตติ น สมฺปสาริยติ, อุเปกฺขา วา ปาฏิกุลฺยตา วา สณฺฐาติ.\csromanspacer{} เวทิตพฺพเมตํ ภิกฺขเว ภิกฺขุนา “สุภาวิตา เม มรณสญฺญา, อตฺถิ เม ปุพฺเพนาปรํ วิเสโส, ปตฺตํ เม ภาวนาพลนฺ”ติ, อิติห ตตฺถ สมฺปชาโน โหติ. “มรณสญฺญา ภิกฺขเว ภาวิตา พหุลีกตา มหปฺผลา โหติ มหานิสํสา อมโตคธา อมตปริโยสานา”ติ อิติ ยํ ตํ วุตฺตํ, อิทเมตํ ปฏิจฺจ วุตฺตํ. (2)}

\prose{"อาหาเร ปฏิกูลสญฺญา ภิกฺขเว ภาวิตา พหุลีกตา มหปฺผลา โหติ มหานิสํสา อมโตคธา อมตปริโยสานา”ติ อิติ โข ปเนตํ วุตฺตํ, กิญฺเจตํ ปฏิจฺจ วุตฺตํ.\csromanspacer{} อาหาเร ปฏิกูลสญฺญาปริจิเตน ภิกฺขเว ภิกฺขุโน เจตสา พหุลํ วิหรโต รสตณฺหาย จิตฺตํ}

\vspace{\tipitakaparskip}
\csromancenter{มหายญฺญวคฺค ปติลียติ ฯเปฯ อุเปกฺขา วา ปาฏิกุลฺยตา วา สณฺฐาติ.\csromanspacer{} เสยฺยถาปิ ภิกฺขเว กุกฺกุฏปตฺตํ วา นฺหารุททฺทุลํ วา อคฺคิมฺหิ ปกฺขิตฺตํ ปติลียติ ปติกุฏติ ปติวตฺตติ น สมฺปสาริยติ.\csromanspacer{} เอวเมวํ โข ภิกฺขเว ภิกฺขุโน อาหาเร ปฏิกูลสญฺญาปริจิเตน เจตสา พหุลํ วิหรโต รสตณฺหาย จิตฺตํ ปติลียติ ฯเปฯ อุเปกฺขา วา ปาฏิกุลฺยตา สณฺฐาติ.}
\prose{\csromanfolio{435}สเจ ภิกฺขเว ภิกฺขุโน อาหาเร ปฏิกูลสญฺญาปริจิเตน เจตสา พหุลํ วิหรโต รสตณฺหาย จิตฺตํ อนุสนฺทหติ, อปฺปฏิกุลฺยตา สณฺฐาติ.\csromanspacer{} เวทิตพฺพเมตํ ภิกฺขเว ภิกฺขุนา “อภาวิตา เม อาหาเร ปฏิกูลสญฺญา, นตฺถิ เม ปุพฺเพนาปรํ วิเสโส, อปฺปตฺตํ เม ภาวนาพลนฺ”ติ, อิติห ตตฺถ สมฺปชาโน โหติ.\csromanspacer{} สเจ ปน ภิกฺขเว ภิกฺขุโน อาหาเร ปฏิกูลสญฺญาปริจิเตน เจตสา พหุลํ วิหรโต รสตณฺหาย จิตฺตํ ปติลียติ ฯเปฯ อุเปกฺขา วา ปาฏิกุลฺยตา วา สณฺฐาติ.\csromanspacer{} เวทิตพฺพเมตํ ภิกฺขเว ภิกฺขุนา “สุภาวิตา เม อาหาเร ปฏิกูลสญฺญา, อตฺถิ เม ปุพฺเพนาปรํ วิเสโส, ปตฺตํ เม ภาวนาพลนฺ”ติ, อิติห ตตฺถ สมฺปชาโน โหติ. “อาหาเร ปฏิกูลสญฺญา ภิกฺขเว ภาวิตา พหุลีกตา มหปฺผลา โหติ มหานิสํสา อมโตคธา อมตปริโยสานา”ติ อิติ ยํ ตํ วุตฺตํ, อิทเมตํ ปฏิจฺจ วุตฺตํ. (3)}

\prose{“สพฺพโลเก อนภิรตสญฺญา ภิกฺขเว ภาวิตา พหุลีกตา มหปฺผลา โหติ มหานิสํสา อมโตคธา อมตปริโยสานา”ติ อิติ โข ปเนตํ วุตฺตํ, กิญฺเจตํ ปฏิจฺจ วุตฺตํ.\csromanspacer{} สพฺพโลเก อนภิรตสญฺญาปริจิเตน ภิกฺขเว ภิกฺขุโน เจตสา พหุลํ วิหรโต โลกจิเตฺรสุ จิตฺตํ ปติลียติ ฯเปฯ.\csromanspacer{} เสยฺยถาปิ ภิกฺขเว ฯเปฯ ปติลียติ ปติกุฏติ ปติวตฺตติ น สมฺปสาริยติ.\csromanspacer{} เอวเมวํ โข ภิกฺขเว ภิกฺขุโน สพฺพโลเก อนภิรตสญฺญาปริจิเตน เจตสา พหุลํ วิหรโต โลกจิเตฺรสุ จิตฺตํ ปติลียติ ปติกุฏติ ปติวตฺตติ น สมฺปสาริยติ, อุเปกฺขา วา ปาฏิกุลฺยตา วา สณฺฐาติ.}

\prose{สเจ ภิกฺขเว ภิกฺขุโน สพฺพโลเก อนภิรตสญฺญาปริจิเตน เจตสา พหุลํ วิหรโต โลกจิเตฺรสุ จิตฺตํ อนุสนฺทหติ, อปฺปฏิกุลฺยตา สณฺฐาติ.\csromanspacer{} เวทิตพฺพเมตํ ภิกฺขเว ภิกฺขุนา “อภาวิตา เม สพฺพโลเก อนภิรตสญฺญา, นตฺถิ เม ปุพฺเพนาปรํ วิเสโส, อปฺปตฺตํ}

\vspace{\tipitakaparskip}
\csromancenter{สตฺตกนิปาตปาฬิ เม ภาวนาพลนฺ”ติ, อิติห ตตฺถ สมฺปชาโน โหติ.\csromanspacer{} สเจ ปน ภิกฺขเว ภิกฺขุโน สพฺพโลเก อนภิรตสญฺญาปริจิเตน เจตสา พหุลํ วิหรโต โลกจิเตฺรสุ จิตฺตํ ปติลียติ ฯเปฯ อุเปกฺขา วา ปาฏิกุลฺยตา วา สณฺฐาติ.\csromanspacer{} เวทิตพฺพเมตํ ภิกฺขเว ภิกฺขุนา “สุภาวิตา เม สพฺพโลเก อนภิรตสญฺญา, อตฺถิ เม ปุพฺเพนาปรํ วิเสโส, ปตฺตํ เม ภาวนาพลนฺ”ติ, อิติห ตตฺถ สมฺปชาโน โหติ. “สพฺพโลเก อนภิรตสญฺญา ภิกฺขเว ภาวิตา พหุลีกตา มหปฺผลา โหติ มหานิสํสา อมโตคธา อมตปริโยสานา”ติ อิติ ยํ ตํ วุตฺตํ, อิทเมตํ ปฏิจฺจ วุตฺตํ. (4)}
\prose{\csromanfolio{436}“อนิจฺจสญฺญา ภิกฺขเว ภาวิตา พหุลีกตา มหปฺผลา โหติ มหานิสํสา อมโตคธา อมตปริโยสานา”ติ อิติ โข ปเนตํ วุตฺตํ, กิญฺเจตํ ปฏิจฺจ วุตฺตํ.\csromanspacer{} อนิจฺจสญฺญาปริจิเตน ภิกฺขเว ภิกฺขุโน เจตสา พหุลํ วิหรโต ลาภสกฺการสิโลเก จิตฺตํ ปติลียติ ฯเปฯ อุเปกฺขา วา ปาฏิกุลฺยตา วา สณฺฐาติ.\csromanspacer{} เสยฺยถาปิ ภิกฺขเว กุกฺกุฏปตฺตํ วา นฺหารุททฺทุลํ วา อคฺคิมฺหิ ปกฺขิตฺตํ ปติลียติ ปติกุฏติ ปติวตฺตติ น สมฺปสาริยติ.\csromanspacer{} เอวเมวํ โข ภิกฺขเว ภิกฺขุโน อนิจฺจสญฺญาปริจิเตน เจตสา พหุลํ วิหรโต ลาภสกฺการสิโลเก จิตฺตํ ปติลียติ ฯเปฯ อุเปกฺขา วา ปาฏิกุลฺยตา วา สณฺฐาติ.}

\prose{สเจ ภิกฺขเว ภิกฺขุโน อนิจฺจสญฺญาปริจิเตน เจตสา พหุลํ วิหรโต ลาภสกฺการสิโลเก จิตฺตํ อนุสนฺทหติ, อปฺปฏิกุลฺยตา สณฺฐาติ.\csromanspacer{} เวทิตพฺพเมตํ ภิกฺขเว ภิกฺขุนา “อภาวิตา เม อนิจฺจสญฺญา, นตฺถิ เม ปุพฺเพนาปรํ วิเสโส, อปฺปตฺตํ เม ภาวนาพลนฺ”ติ, อิติห ตตฺถ สมฺปชาโน โหติ.\csromanspacer{} สเจ ปน ภิกฺขเว ภิกฺขุโน อนิจฺจสญฺญาปริจิเตน เจตสา พหุลํ วิหรโต ลาภสกฺการสิโลเก จิตฺตํ ปติลียติ ปติกุฏติ ปติวตฺตติ น สมฺปสาริยติ, อุเปกฺขา วา ปาฏิกุลฺยตา วา สณฺฐาติ.\csromanspacer{} เวทิตพฺพเมตํ ภิกฺขเว ภิกฺขุนา “สุภาวิตา เม อนิจฺจสญฺญา, อตฺถิ เม ปุพฺเพนาปรํ วิเสโส, ปตฺตํ เม ภาวนาพลนฺ”ติ, อิติห ตตฺถ สมฺปชาโน โหติ. “อนิจฺจสญฺญา ภิกฺขเว ภาวิตา พหุลีกตา มหปฺผลา โหติ มหานิสํสา อมโตคธา อมตปริโยสานา”ติ อิติ ยํ ตํ วุตฺตํ, อิทเมตํ ปฏิจฺจ วุตฺตํ. (5)}

\chapterheadpageread{5. มหายญฺญวคฺค}
\prose{\csromanfolio{437}“อนิจฺเจ ทุกฺขสญฺญา ภิกฺขเว ภาวิตา พหุลีกตา มหปฺผลา โหติ มหานิสํสา อมโตคธา อมตปริโยสานา”ติ อิติ โข ปเนตํ วุตฺตํ, กิญฺเจตํ ปฏิจฺจ วุตฺตํ.\csromanspacer{} อนิจฺเจ ทุกฺขสญฺญาปริจิเตน ภิกฺขเว ภิกฺขุโน เจตสา พหุลํ วิหรโต อาลเสฺย โกสชฺเช วิสฺสฏฺฐิเย ปมาเท อนนุโยเค อปจฺจเวกฺขณาย ติพฺพา ภยสญฺญา ปจฺจุปฏฺฐิตา โหติ.\csromanspacer{} เสยฺยถาปิ ภิกฺขเว อุกฺขิตฺตาสิเก วธเก.}

\prose{สเจ ภิกฺขเว ภิกฺขุโน อนิจฺเจ ทุกฺขสญฺญาปริจิเตน เจตสา พหุลํ วิหรโต อาลเสฺย โกสชฺเช วิสฺสฏฺฐิเย ปมาเท อนนุโยเค อปจฺจเวกฺขณาย ติพฺพา ภยสญฺญา น ปจฺจุปฏฺฐิตา โหติ.\csromanspacer{} เสยฺยถาปิ ภิกฺขเว อุกฺขิตฺตาสิเก วธเก.\csromanspacer{} เวทิตพฺพเมตํ ภิกฺขเว ภิกฺขุนา “อภาวิตา เม อนิจฺเจ ทุกฺขสญฺญา, นตฺถิ เม ปุพฺเพนาปรํ วิเสโส, อปฺปตฺตํ เม ภาวนาพลนฺ”ติ, อิติห ตตฺถ สมฺปชาโน โหติ.\csromanspacer{} สเจ ปน ภิกฺขเว ภิกฺขุโน อนิจฺเจ ทุกฺขสญฺญาปริจิเตน เจตสา พหุลํ วิหรโต อาลเสฺย โกสชฺเช วิสฺสฏฺฐิเย ปมาเท อนนุโยเค อปจฺจเวกฺขณาย ติพฺพา ภยสญฺญา ปจฺจุปฏฺฐิตา โหติ.\csromanspacer{} เสยฺยถาปิ ภิกฺขเว อุกฺขิตฺตาสิเก วธเก.\csromanspacer{} เวทิตพฺพเมตํ ภิกฺขเว ภิกฺขุนา “สุภาวิตา เม อนิจฺเจ ทุกฺขสญฺญา, อตฺถิ เม ปุพฺเพนาปรํ วิเสโส, ปตฺตํ เม ภาวนาพลนฺ”ติ, อิติห ตตฺถ สมฺปชาโน โหติ. “อนิจฺเจ ทุกฺขสญฺญา ภิกฺขเว ภาวิตา พหุลีกตา มหปฺผลา โหติ มหานิสํสา อมโตคธา อมตปริโยสานา”ติ อิติ ยํ ตํ วุตฺตํ, อิทเมตํ ปฏิจฺจ วุตฺตํ. (6)}

\prose{“ทุกฺเข อนตฺตสญฺญา ภิกฺขเว ภาวิตา พหุลีกตา มหปฺผลา โหติ มหานิสํสา อมโตคธา อมตปริโยสานา”ติ อิติ โข ปเนตํ วุตฺตํ, กิญฺเจตํ ปฏิจฺจ วุตฺตํ, ทุกฺเข อนตฺตสญฺญาปริจิเตน ภิกฺขเว ภิกฺขุโน เจตสา พหุลํ วิหรโต อิมสฺมิํ จ สวิญฺญาณเก กาเย พหิทฺธา จ สพฺพนิมิตฺเตสุ อหงฺการมมงฺการมานาปคตํ มานสํ โหติ วิธาสมติกฺกนฺตํ สนฺตํ สุวิมุตฺตํ.}

\prose{สเจ ภิกฺขเว ภิกฺขุโน ทุกฺเข อนตฺตสญฺญาปริจิเตน เจตสา พหุลํ วิหรโต อิมสฺมิํ จ สวิญฺญาณเก กาเย พหิทฺธา จ สพฺพนิมิตฺเตสุ น อหงฺการมมงฺการมานาปคตํ มานสํ โหติ วิธาสมติกฺกนฺตํ สนฺตํ สุวิมุตฺตํ.\csromanspacer{} เวทิตพฺพเมตํ ภิกฺขเว ภิกฺขุนา “อภาวิตา เม ทุกฺเข}

\vspace{\tipitakaparskip}
\csromancenter{สตฺตกนิปาตปาฬิ อนตฺตสญฺญา, นตฺถิ เม ปุพฺเพนาปรํ วิเสโส, อปฺปตฺตํ เม ภาวนาพลนฺ”ติ, อิติห ตตฺถ สมฺปชาโน โหติ.}
\prose{\csromanfolio{438}สเจ ปน ภิกฺขเว ภิกฺขุโน ทุกฺเข อนตฺตสญฺญาปริจิเตน เจตสา พหุลํ วิหรโต อิมสฺมิํ จ สวิญฺญาณเก กาเย พหิทฺธา จ สพฺพนิมิตฺเตสุ อหงฺการมมงฺการมานาปคตํ มานสํ โหติ วิธาสมติกฺกนฺตํ สนฺตํ สุวิมุตฺตํ.\csromanspacer{} เวทิตพฺพเมตํ ภิกฺขเว ภิกฺขุนา “สุภาวิตา เม ทุกฺเข อนตฺตสญฺญา, อตฺถิ เม ปุพฺเพนาปรํ วิเสโส, ปตฺตํ เม ภาวนาพลนฺ”ติ, อิติห ตตฺถ สมฺปชาโน โหติ. “ทุกฺเข อนตฺตสญฺญา ภิกฺขเว ภาวิตา พหุลีกตา มหปฺผลา โหติ มหานิสํสา อมโตคธา อมตปริโยสานา”ติ อิติ ยํ ตํ วุตฺตํ, อิทเมตํ ปฏิจฺจ วุตฺตํ. (7)}

\prose{อิมา โข ภิกฺขเว สตฺต สญฺญา ภาวิตา พหุลีกตา มหปฺผลา โหนฺติ มหานิสํสา อมโตคธา อมตปริโยสานาติ.\csromanspacer{}\csromanspacer{}ฉฏฺฐํ.}

\csromanmatikaanchor{mtk.244}
\csromantocmarkref{subsection}{7. เมถุนสุตฺต}{mtk.244}
\bookmark[dest={mtk.244},level=2]{7. เมถุนสุตฺต}
\csromanheader{2}{7. เมถุนสุตฺต}
\proseitem{50}{อถ โข ชาณุสฺโสณิ พฺราหฺมโณ เยน ภควา เตนุปสงฺกมิ, อุปสงฺกมิตฺวา ภควตา สทฺธิํ สมฺโมทิ, สมฺโมทนียํ กถํ สารณียํ วีติสาเรตฺวา เอกมนฺตํ นิสีทิ, เอกมนฺตํ นิสินฺโน โข ชาณุสฺโสณิ พฺราหฺมโณ ภควนฺตํ เอตทโวจ “ภวมฺปิ โน โคตโม พฺรหฺมจารี ปฏิชานาตี”ติ.\csromanspacer{} ยํ หิ ตํ พฺราหฺมณ สมฺมา วทมาโน วเทยฺย “อขณฺฑํ อจฺฉิทฺทํ อสพลํ อกมฺมาสํ ปริปุณฺณํ ปริสุทฺธํ พฺรหฺมจริยํ จรตี”ติ.\csromanspacer{} มเมว ตํ พฺราหฺมณ สมฺมา วทมาโน วเทยฺย “อหํ หิ พฺราหฺมณ อขณฺฑํ อจฺฉิทฺทํ อสพลํ อกมฺมาสํ ปริปุณฺณํ ปริสุทฺธํ พฺรหฺมจริยํ จรามี”ติ.\csromanspacer{} กิํ ปน โภ โคตม พฺรหฺมจริยสฺส ขณฺฑมฺปิ ฉิทฺทมฺปิ สพลมฺปิ กมฺมาสมฺปีติ.}

\prose{อิธ พฺราหฺมณ เอกจฺโจ สมโณ วา พฺราหฺมโณ วา สมฺมา พฺรหฺมจารี ปฏิชานมาโน น เหว โข มาตุคาเมน สทฺธิํ ทฺวยํทฺวยสมาปตฺติํ สมาปชฺชติ, อปิ จ โข มาตุคามสฺส อุจฺฉาทน ปริมทฺทน นฺหาปน สมฺพาหนํ สาทิยติ, โส ตํ อสฺสาเทติ\footnote{โส ตทสฺสาเทติ (สี)}, ตํ นิกาเมติ, เตน จ วิตฺติํ อาปชฺชติ, อิทมฺปิ โข พฺราหฺมณ พฺรหฺมจริยสฺส ขณฺฑมฺปิ ฉิทฺทมฺปิ สพลมฺปิ กมฺมาสมฺปิ, อยํ วุจฺจติ พฺราหฺมณ อปริสุทฺธํ พฺรหฺมจริยํ จรติ สํยุตฺโต เมถุเนน สํโยเคน, น ปริมุจฺจติ ชาติยา ชราย มรเณน\footnote{ชรามรเณน (สี, สฺยา)}}

\vspace{\tipitakaparskip}
\csromancenter{มหายญฺญวคฺค โสเกหิ ปริเทเวหิ ทุกฺเขหิ โทมนสฺเสหิ อุปายาเสหิ, น ปริมุจฺจติ ทุกฺขสฺมาติ วทามิ. (1)}
\prose{\csromanfolio{439}ปุน จปรํ พฺราหฺมณ อิเธกจฺโจ สมโณ วา พฺราหฺมโณ วา สมฺมา พฺรหฺมจารี ปฏิชานมาโน น เหว โข มาตุคาเมน สทฺธิํ ทฺวยํทฺวยสมาปตฺติํ สมาปชฺชติ, นปิ มาตุคามสฺส อุจฺฉาทนปริมทฺทนนฺหาปนสมฺพาหนํ สาทิยติ, อปิ จ โข มาตุคาเมน สทฺธิํ สญฺชคฺฆติ สํกีฬติ สํเกลายติ ฯเปฯ.\csromanspacer{} นปิ มาตุคาเมน สทฺธิํ สญฺชคฺฆติ สํกีฬติ สํเกลายติ, อปิ จ โข มาตุคามสฺส จกฺขุนา จกฺขุํ อุปนิชฺฌายติ เปกฺขติ ฯเปฯ.\csromanspacer{} นปิ มาตุคามสฺส จกฺขุนา จกฺขุํ อุปนิชฺฌายติ เปกฺขติ, อปิ จ โข มาตุคามสฺส สทฺทํ สุณาติ ติโรกุฏฺฏํ วา ติโรปาการํ วา หสนฺติยา วา ภณนฺติยา วา คายนฺติยา วา โรทนฺติยา วา ฯเปฯ.\csromanspacer{} นปิ มาตุคามสฺส สทฺทํ สุณาติ ติโรกุฏฺฏํ วา ติโรปาการํ วา หสนฺติยา วา ภณนฺติยา วา คายนฺติยา วา โรทนฺติยา วา, อปิ จ โข ยานิสฺส ตานิ ปุพฺเพ มาตุคาเมน สทฺธิํ หสิตลปิตกีฬิตานิ ตานิ อนุสฺสรติ ฯเปฯ.\csromanspacer{} นปิ ยานิสฺส ตานิ ปุพฺเพ มาตุคาเมน สทฺธิํ หสิตลปิตกีฬิตานิ ตานิ อนุสฺสรติ, อปิ จ โข ปสฺสติ คหปติํ วา คหปติปุตฺตํ วา ปญฺจหิ กามคุเณหิ สมปฺปิตํ สมงฺคีภูตํ ปริจารยมานํ ฯเปฯ.\csromanspacer{} นปิ ปสฺสติ คหปติํ วา คหปติปุตฺตํ วา ปญฺจหิ กามคุเณหิ สมปฺปิตํ สมงฺคีภูตํ ปริจารยมานํ, อปิ จ โข อญฺญตรํ เทวนิกายํ ปณิธาย พฺรหฺมจริยํ จรติ “อิมินาหํ สีเลน วา วเตน วา ตเปน วา พฺรหฺมจริเยน วา เทโว วา ภวิสฺสามิ เทวญฺญตโร วา”ติ.\csromanspacer{} โส ตํ อสฺสาเทติ, ตํ นิกาเมติ, เตน จ วิตฺติํ อาปชฺชติ.\csromanspacer{} อิทมฺปิ โข พฺรหฺมณ พฺรหฺมจริยสฺส ขณฺฑมฺปิ ฉิทฺทมฺปิ สพลมฺปิ กมฺมาสมฺปิ.\csromanspacer{} อยํ วุจฺจติ พฺราหฺมณ อปริสุทฺธํ พฺรหฺมจริยํ จรติ สํยุตฺโต เมถุเนน สํโยเคน, น ปริมุจฺจติ ชาติยา ชราย มรเณน โสเกหิ ปริเทเวหิ ทุกฺเขหิ โทมนสฺเสหิ อุปายาเสหิ, น ปริมุจฺจติ ทุกฺขสฺมาติ วทามิ. (\footnote{อภิสมฺพุทฺโธ ปจฺจญฺญาสิํ (สี, สฺยา)}-7)}

\prose{ยาวกีวญฺจาหํ พฺราหฺมณ อิเมสํ สตฺตนฺนํ เมถุนสํโยคานํ อญฺญตรญฺญตรเมถุนสํโยคํ\footnote{อญฺญตรํ เมถุนสํโยคํ (สี, สฺยา)} อตฺตนิ อปฺปหีนํ สมนุปสฺสิํ, เนว ตาวาหํ พฺราหฺมณ สเทวเก โลเก สมารเก สพฺรหฺมเก สสฺสมณพฺราหฺมณิยา ปชาย สเทวมนุสฺสาย อนุตฺตรํ สมฺมาสมฺโพธิํ อภิสมฺพุทฺโธติ ปจฺจญฺญาสิํ2.}

\gambhira{สตฺตกนิปาตปาฬิ}
\prose{\csromanfolio{440}ยโต จ โขหํ พฺราหฺมณ อิเมสํ สตฺตนฺนํ เมถุนสํโยคานํ อญฺญตรญฺญตรเมถุนสํโยคํ อตฺตนิ อปฺปหีนํ น สมนุปสฺสิํ, อถาหํ พฺราหฺมณ สเทวเก โลเก สมารเก สพฺรหฺมเก สสฺสมณพฺราหฺมณิยา ปชาย สเทวมนุสฺสาย อนุตฺตรํ สมฺมาสมฺโพธิํ อภิสมฺพุทฺโธติ ปจฺจญฺญาสิํ.\csromanspacer{} ญาณญฺจ ปน เม ทสฺสนํ อุทปาทิ, อกุปฺปา เม วิมุตฺติ\footnote{เจโตวิมุตฺติ (สี, ก)}, อยมนฺติมา ชาติ, นตฺถิ ทานิ ปุนพฺภโวติ.}

\prose{เอวํ วุตฺเต ชาณุสฺโสณิ พฺราหฺมโณ ภควนฺตํ เอตทโวจ “อภิกฺกนฺตํ โภ โคตม, อภิกฺกนฺตํ โภ โคตม ฯเปฯ อุปาสกํ มํ ภวํ โคตโม ธาเรตุ อชฺชตคฺเค ปาณุเปตํ สรณํ คตนฺติ.\csromanspacer{}\csromanspacer{}สตฺตมํ.}

\csromanmatikaanchor{mtk.245}
\csromantocmarkref{subsection}{8. สํโยคสุตฺต}{mtk.245}
\bookmark[dest={mtk.245},level=2]{8. สํโยคสุตฺต}
\csromanheader{2}{8. สํโยคสุตฺต}
\proseitem{51}{สํโยควิสํโยคํ โว ภิกฺขเว ธมฺมปริยายํ เทเสสฺสามิ, ตํ สุณาถ ฯเปฯ.\csromanspacer{} กตโม จ โส ภิกฺขเว สํโยโค วิสํโยโค ธมฺมปริยาโย.}

\prose{อิตฺถี ภิกฺขเว อชฺฌตฺตํ อิตฺถินฺทฺริยํ มนสิ กโรติ อิตฺถิกุตฺตํ อิตฺถากปฺปํ อิตฺถิวิธํ อิตฺถิจฺฉนฺทํ อิตฺถิสฺสรํ อิตฺถาลงฺการํ.\csromanspacer{} สา ตตฺถ รชฺชติ, ตตฺราภิรมติ.\csromanspacer{} สา ตตฺถ รตฺตา ตตฺราภิรตา พหิทฺธา ปุริสินฺทฺริยํ มนสิ กโรติ ปุริสกุตฺตํ ปุริสากปฺปํ ปุริสวิธํ ปุริสจฺฉนฺทํ ปุริสสฺสรํ ปุริสาลงฺการํ.\csromanspacer{} สา ตตฺถ รชฺชติ, ตตฺราภิรมติ.\csromanspacer{} สา ตตฺถ รตฺตา ตตฺราภิรตา พหิทฺธา สํโยคํ อากงฺขติ, ยญฺจสฺสา สํโยคปจฺจยา อุปฺปชฺชติ สุขํ โสมนสฺสํ, ตญฺจ อากงฺขติ.\csromanspacer{} อิตฺถตฺเต ภิกฺขเว อภิรตา สตฺตา ปุริเสสุ สํโยคํ คตา.\csromanspacer{} เอวํ โข ภิกฺขเว อิตฺถี อิตฺถตฺตํ นาติวตฺตติ.}

\prose{ปุริโส ภิกฺขเว อชฺฌตฺตํ ปุริสินฺทฺริยํ มนสิ กโรติ ปุริสกุตฺตํ ปุริสากปฺปํ ปุริสวิธํ ปุริสจฺฉนฺทํ ปุริสสฺสรํ ปุริสาลงฺการํ.\csromanspacer{} โส ตตฺถ รชฺชติ, ตตฺราภิรมติ.\csromanspacer{} โส ตตฺถ รตฺโต ตตฺราภิรโต พหิทฺธา อิตฺถินฺทฺริยํ มนสิ กโรติ อิตฺถิกุตฺตํ อิตฺถากปฺปํ อิตฺถิวิธํ อิตฺถิจฺฉนฺทํ อิตฺถิสฺสรํ อิตฺถาลงฺการํ.\csromanspacer{} โส ตตฺถ รชฺชติ, ตตฺราภิรมติ.\csromanspacer{} โส ตตฺถ รตฺโต ตตฺราภิรโต พหิทฺธา สํโยคํ อากงฺขติ, ยญฺจสฺส สํโยคปจฺจยา อุปฺปชฺชติ สุขํ โสมนสฺสํ, ตญฺจ อากงฺขติ.\csromanspacer{} ปุริสตฺเต ภิกฺขเว อภิรตา}

\vspace{\tipitakaparskip}
\csromancenter{มหายญฺญวคฺค สตฺตา อิตฺถีสุ สํโยคํ คตา.\csromanspacer{} เอวํ โข ภิกฺขเว ปุริโส ปุริสตฺตํ นาติวตฺตติ.\csromanspacer{} เอวํ โข ภิกฺขเว สํโยโค โหติ.}
\prose{\csromanfolio{441}กถญฺจ ภิกฺขเว วิสํโยโค โหติ.\csromanspacer{} อิตฺถี ภิกฺขเว อชฺฌตฺตํ อิตฺถินฺทฺริยํ น มนสิ กโรติ อิตฺถิกุตฺตํ อิตฺถากปฺปํ อิตฺถิวิธํ อิตฺถิจฺฉนฺทํ อิตฺถิสฺสรํ อิตฺถาลงฺการํ.\csromanspacer{} สา ตตฺถ น รชฺชติ, สา ตตฺร นาภิรมติ.\csromanspacer{} สา ตตฺถ อรตฺตา ตตฺร อนภิรตา พหิทฺธา ปุริสินฺทฺริยํ น มนสิ กโรติ ปุริสกุตฺตํ ปุริสากปฺปํ ปุริสวิธํ ปุริสจฺฉนฺทํ ปุริสสฺสรํ ปุริสาลงฺการํ.\csromanspacer{} สา ตตฺถ น รชฺชติ, ตตฺร นาภิรมติ.\csromanspacer{} สา ตตฺถ อรตฺตา ตตฺร อนภิรตา พหิทฺธา สํโยคํ นากงฺขติ, ยญฺจสฺสา สํโยคปจฺจยา อุปฺปชฺชติ สุขํ โสมนสฺสํ, ตญฺจ นากงฺขติ.\csromanspacer{} อิตฺถตฺเต ภิกฺขเว อนภิรตา สตฺตา ปุริเสสุ วิสํโยคํ คตา.\csromanspacer{} เอวํ โข ภิกฺขเว อิตฺถี อิตฺถตฺตํ อติวตฺตติ.}

\prose{ปุริโส ภิกฺขเว อชฺฌตฺตํ ปุริสินฺทฺริยํ น มนสิ กโรติ ปุริสกุตฺตํ ปุริสากปฺปํ ปุริสวิธํ ปุริสจฺฉนฺทํ ปุริสสฺสรํ ปุริสาลงฺการํ.\csromanspacer{} โส ตตฺถ น รชฺชติ, โส ตตฺร นาภิรมติ.\csromanspacer{} โส ตตฺถ อรตฺโต ตตฺร อนภิรโต พหิทฺธา อิตฺถินฺทฺริยํ น มนสิ กโรติ อิตฺถิกุตฺตํ อิตฺถากปฺปํ อิตฺถิวิธํ อิตฺถิจฺฉนฺทํ อิตฺถิสฺสรํ อิตฺถาลงฺการํ.\csromanspacer{} โส ตตฺถ น รชฺชติ, ตตฺร นาภิรมติ.\csromanspacer{} โส ตตฺถ อรตฺโต ตตฺร อนภิรโต พหิทฺธา สํโยคํ นากงฺขติ, ยญฺจสฺส สํโยคปจฺจยา อุปฺปชฺชติ สุขํ โสมนสฺสํ, ตญฺจ นากงฺขติ.\csromanspacer{} ปุริสตฺเต ภิกฺขเว อนภิรตา สตฺตา อิตฺถีสุ วิสํโยคํ คตา.\csromanspacer{} เอวํ โข ภิกฺขเว ปุริโส ปุริสตฺตํ อติวตฺตติ.\csromanspacer{} เอวํ โข ภิกฺขเว วิสํโยโค โหติ.\csromanspacer{} อยํ โข ภิกฺขเว สํโยโค วิสํโยโค ธมฺมปริยาโยติ.\csromanspacer{}\csromanspacer{}อฏฺฐมํ.}

\csromanmatikaanchor{mtk.246}
\csromantocmarkref{subsection}{9. ทานมหปฺผลสุตฺต}{mtk.246}
\bookmark[dest={mtk.246},level=2]{9. ทานมหปฺผลสุตฺต}
\csromanheader{2}{9. ทานมหปฺผลสุตฺต}
\proseitem{52}{เอกํ สมยํ ภควา จมฺปายํ วิหรติ คคฺคราย โปกฺขรณิยา ตีเร.\csromanspacer{} อถ โข สมฺพหุลา จมฺเปยฺยกา อุปาสกา เยน อายสฺมา สาริปุตฺโต เตนุปสงฺกมิํสุ, อุปสงฺกมิตฺวา อายสฺมนฺตํ สาริปุตฺตํ อภิวาเทตฺวา เอกมนฺตํ นิสีทิํสุ, เอกมนฺตํ นิสินฺนา โข จมฺเปยฺยกา อุปาสกา อายสฺมนฺตํ สาริปุตฺตํ เอตทโวจุํ “จิรสฺสุตา โน ภนฺเต\footnote{ภนฺเต สาริปุตฺต (สี)}}

\vspace{\tipitakaparskip}
\csromancenter{สตฺตกนิปาตปาฬิ ภควโต สมฺมุขา ธมฺมีกถา, สาธุ มยํ ภนฺเต ลเภยฺยาม ภควโต สมฺมุขา ธมฺมิํ กถํ\footnote{ภควโต สนฺติกา ธมฺมิํ กถํ (สี), ภควโต ธมฺมิํ กถํ (สฺยา)} สวนายา”ติ.\csromanspacer{} เตนหาวุโส ตทหุโปสเถ อาคจฺเฉยฺยาถ, อปฺเปว นาม ลเภยฺยาถ ภควโต สมฺมุขา\footnote{ภควโต สนฺติเก (สฺยา)} ธมฺมิํ กถํ สวนายาติ. “เอวํ ภนฺเต”ติ โข จมฺเปยฺยกา อุปาสกา อายสฺมโต สาริปุตฺตสฺส ปฏิสฺสุตฺวา อุฏฺฐายาสนา อายสฺมนฺตํ สาริปุตฺตํ อภิวาเทตฺวา ปทกฺขิณํ กตฺวา ปกฺกมิํสุ.}
\prose{\csromanfolio{442}อถ โข จมฺเปยฺยกา อุปาสกา ตทหุโปสเถ เยนายสฺมา สาริปุตฺโต เตนุปสงฺกมิํสุ, อุปสงฺกมิตฺวา อายสฺมนฺตํ สาริปุตฺตํ อภิวาเทตฺวา เอกมนฺตํ อฏฺฐํสุ.\csromanspacer{} อถ โข อายสฺมา สาริปุตฺโต เตหิ จมฺเปยฺยเกหิ อุปาสเกหิ สทฺธิํ เยน ภควา เตนุปสงฺกมิ, อุปสงฺกมิตฺวา ภควนฺตํ อภิวาเทตฺวา เอกมนฺตํ นิสีทิ, เอกมนฺตํ นิสินฺโน โข อายสฺมา สาริปุตฺโต ภควนฺตํ เอตทโวจ\csromandash{}}

\prose{สิยา นุ โข ภนฺเต อิเธกจฺจสฺส ตาทิสํเยว ทานํ ทินฺนํ น มหปฺผลํ โหติ น มหานิสํสํ.\csromanspacer{} สิยา ปน ภนฺเต อิเธกจฺจสฺส ตาทิสํเยว ทานํ ทินฺนํ มหปฺผลํ โหติ มหานิสํสนฺติ.\csromanspacer{} สิยา สาริปุตฺต อิเธกจฺจสฺส ตาทิสํเยว ทานํ ทินฺนํ น มหปฺผลํ โหติ น มหานิสํสํ.\csromanspacer{} สิยา ปน สาริปุตฺต อิเธกจฺจสฺส ตาทิสํเยว ทานํ ทินฺนํ มหปฺผลํ โหติ มหานิสํสนฺติ.\csromanspacer{} โก นุ โข ภนฺเต เหตุ โก ปจฺจโย, เยน มิเธกจฺจสฺส ตาทิสํเยว ทานํ ทินฺนํ น มหปฺผลํ โหติ น มหานิสํสํ.\csromanspacer{} โก นุ โข ภนฺเต เหตุ โก ปจฺจโย, เยน มิเธกจฺจสฺส ตาทิสํเยว ทานํ ทินฺนํ มหปฺผลํ โหติ มหานิสํสนฺติ.}

\prose{อิธ สาริปุตฺต เอกจฺโจ สาเปโข\footnote{สาเปโข (สฺยา)} ทานํ เทติ, ปติพทฺธจิตฺโต\footnote{ปติพนฺธจิตฺโต (ก)} ทานํ เทติ, สนฺนิธิเปโข ทานํ เทติ, “อิมํ เปจฺจ ปริภุญฺชิสฺสามี”ติ ทานํ เทติ.\csromanspacer{} โส ตํ ทานํ เทติ สมณสฺส วา พฺราหฺมณสฺส วา อนฺนํ ปานํ วตฺถํ ยานํ มาลาคนฺธวิเลปนํ เสยฺยาวสถปทีเปยฺยํ.\csromanspacer{} ตํ กิํ มญฺญสิ สาริปุตฺต ทเทยฺย อิเธกจฺโจ เอวรูปํ ทานนฺติ.\csromanspacer{} เอวํ ภนฺเต.}

\prose{ตตฺร สาริปุตฺต ยฺวายํ สาเปโข ทานํ เทติ, ปติพทฺธจิตฺโต ทานํ เทติ, สนฺนิธิเปโข ทานํ เทติ, “อิมํ เปจฺจ ปริภุญฺชิสฺสามี”ติ ทานํ เทติ.}

\vspace{\tipitakaparskip}
\csromancenter{มหายญฺญวคฺค โส ตํ ทานํ ทตฺวา กายสฺส เภทา ปรํ มรณา จาตุมหาราชิกานํ เทวานํ สหพฺยตํ อุปปชฺชติ.\csromanspacer{} โส ตํ กมฺมํ เขเปตฺวา ตํ อิทฺธิํ ตํ ยสํ ตํ อาธิปจฺจํ อาคามี โหติ อาคนฺตา อิตฺถตฺตํ.}
\prose{\csromanfolio{443}อิธ ปน สาริปุตฺต เอกจฺโจ น เหว โข สาเปโข ทานํ เทติ, น ปติพทฺธจิตฺโต ทานํ เทติ, น สนฺนิธิเปโข ทานํ เทติ, น “อิมํ เปจฺจ ปริภุญฺชิสฺสามี”ติ ทานํ เทติ.\csromanspacer{} อปิ จ โข “สาหุ ทานนฺ”ติ ทานํ เทติ ฯเปฯ.\csromanspacer{} นปิ “สาหุ ทานนฺ”ติ ทานํ เทติ.\csromanspacer{} อปิ จ โข “ทินฺนปุพฺพํ กตปุพฺพํ ปิตุปิตามเหหิ, น อรหามิ โปราณํ กุลวํสํ หาเปตุนฺ”ติ ทานํ เทติ ฯเปฯ.\csromanspacer{} นปิ “ทินฺนปุพฺพํ กตปุพฺพํ ปิตุปิตามเหหิ, น อรหามิ โปราณํ กุลวํสํ หาเปตุนฺ”ติ ทานํ เทติ.\csromanspacer{} อปิ จ โข “อหํ ปจามิ, อิเม น ปจนฺติ, นารหามิ ปจนฺโต อปจนฺตานํ ทานํ อทาตุนฺ”ติ ทานํ เทติ ฯเปฯ.\csromanspacer{} นปิ “อหํ ปจามิ, อิเม น ปจนฺติ, นารหามิ ปจนฺโต อปจนฺตานํ ทานํ อทาตุนฺ”ติ ทานํ เทติ.\csromanspacer{} อปิ จ โข “ยถา เตสํ ปุพฺพกานํ อิสีนํ ตานิ มหายญฺญานิ อเหสุํ.\csromanspacer{} เสยฺยถิทํ, อฏฺฐกสฺส วามกสฺส วามเทวสฺส เวสฺสามิตฺตสฺส ยมทคฺคิโน องฺคีรสสฺส ภารทฺวาชสฺส วาเสฏฺฐสฺส กสฺสปสฺส ภคุโน.\csromanspacer{} เอวํ เม อยํ ทานสํวิภาโค ภวิสฺสตี”ติ ทานํ เทติ ฯเปฯ.\csromanspacer{} นปิ “ยถา เตสํ ปุพฺพกานํ อิสีนํ ตานิ มหายญฺญานิ อเหสุํ.\csromanspacer{} เสยฺยถิทํ, อฏฺฐกสฺส วามกสฺส วามเทวสฺส เวสฺสามิตฺตสฺส ยมทคฺคิโน องฺคีรสสฺส ภารทฺวาชสฺส วาเสฏฺฐสฺส กสฺสปสฺส ภคุโน.\csromanspacer{} เอวํ เม อยํ ทานสํวิภาโค ภวิสฺสตี”ติ ทานํ เทติ.\csromanspacer{} อปิ จ โข “อิมํ เม ทานํ ททโต จิตฺตํ ปสีทติ, อตฺตมนตา โสมนสฺสํ อุปชายตี”ติ ทานํ เทติ ฯเปฯ.\csromanspacer{} นปิ “อิมํ เม ทานํ ททโต จิตฺตํ ปสีทติ, อตฺตมนตา โสมนสฺสํ อุปชายตี”ติ ทานํ เทติ.\csromanspacer{} อปิ จ โข จิตฺตาลงฺการจิตฺตปริกฺขารํ ทานํ เทติ.\csromanspacer{} โส ตํ ทานํ เทติ สมณสฺส วา พฺราหฺมณสฺส วา อนฺนํ ปานํ วตฺถํ ยานํ มาลาคนฺธวิเลปนํ เสยฺยาวสถปทีเปยฺยํ.\csromanspacer{} ตํ กิํ มญฺญสิ สาริปุตฺต, ทเทยฺย อิเธกจฺโจ เอวรูปํ ทานนฺติ.\csromanspacer{} เอวํ ภนฺเต.}

\prose{ตตฺร สาริปุตฺต ยฺวายํ น เหว\footnote{นเหว โข (สี, สฺยา)} สาเปโข ทานํ เทติ, น ปติพทฺธจิตฺโต ทานํ เทติ, น สนฺนิธิเปโข ทานํ เทติ, น “อิมํ ปจฺจ ปริภุญฺชิสฺสามี”ติ ทานํ เทติ, นปิ “สาหุ ทานนฺ”ติ ทานํ เทติ.\csromanspacer{} นปิ “ทินฺนปุพฺพํ กตปุพฺพํ}

\vspace{\tipitakaparskip}
\csromancenter{สตฺตกนิปาตปาฬิ ปิตุปิตามเหหิ น อรหามิ โปราณํ กุลวํสํ หาเปตุนฺ”ติ ทานํ เทติ.\csromanspacer{} นปิ “อหํ ปจามิ, อิเม น ปจนฺติ, นารหามิ ปจนฺโต อปจนฺตานํ ทานํ อทาตุนฺ”ติ ทานํ เทติ.\csromanspacer{} นปิ “ยถา เตสํ ปุพฺพกานํ อิสีนํ, ตานิ มหายญฺญานิ อเหสุํ.\csromanspacer{} เสยฺยถิทํ, อฏฺฐกสฺส วามกสฺส วามเทวสฺส เวสฺสามิตฺตสฺส ยมทคฺคิโน องฺคีรสสฺส ภารทฺวาชสฺส วาเสฏฺฐสฺส กสฺสปสฺส ภคุโน.\csromanspacer{} เอวํ เม อยํ ทานสํวิภาโค ภวิสฺสตี”ติ ทานํ เทติ.\csromanspacer{} นปิ “อิมํ เม ทานํ ททโต จิตฺตํ ปสีทติ, อตฺตมนตา โสมนสฺสํ อุปชายตี”ติ ทานํ เทติ.\csromanspacer{} อปิ จ โข จิตฺตาลงฺการจิตฺตปริกฺขารํ ทานํ เทติ.\csromanspacer{} โส ตํ ทานํ ทตฺวา กายสฺส เภทา ปรํ มรณา พฺรหฺมกายิกานํ เทวานํ สหพฺยตํ อุปปชฺชติ, โส ตํ กมฺมํ เขเปตฺวา ตํ อิทฺธิํ ตํ ยสํ ตํ อาธิปจฺจํ อนาคามี โหติ อนาคนฺตา อิตฺถตฺตํ.\csromanspacer{} อยํ โข สาริปุตฺต เหตุ อยํ ปจฺจโย, เยน มิเธกจฺจสฺส ตาทิสํเยว ทานํ ทินฺนํ น มหปฺผลํ โหติ น มหานิสํสํ.\csromanspacer{} อยํ ปน สาริปุตฺต เหตุ อยํ ปจฺจโย, เยน มิเธกจฺจสฺส ตาทิสํเยว ทานํ ทินฺนํ มหปฺผลํ โหติ มหานิสํสนฺติ.\csromanspacer{}\csromanspacer{}นวมํ.}
\csromanmatikaanchor{mtk.247}
\csromantocmarkref{subsection}{10. นนฺทมาตาสุตฺต}{mtk.247}
\bookmark[dest={mtk.247},level=2]{10. นนฺทมาตาสุตฺต}
\csromanheader{2}{10. นนฺทมาตาสุตฺต}
\proseitem{53}{\csromanfolio{444}เอวํ เม สุตํ\csromandash{}เอกํ สมยํ อายสฺมา จ สาริปุตฺโต อายสฺมา จ มหาโมคฺคลฺลาโน ทกฺขิณาคิริสฺมิํ จาริกํ จรนฺติ มหตา ภิกฺขุสํเฆน สทฺธิํ.\csromanspacer{} เตน โข ปน สมเยน เวฬุกณฺฑกี\footnote{เวฬุกณฺฏกี (สฺยา) เหฏฺฐา 295; อํ 1. 88; สํ 1. 432 ปิฏฺเฐสุปิ ปสฺสิตพฺพํ.} นนฺทมาตา อุปาสิกา รตฺติยา ปจฺจูสสมยํ ปจฺจุฏฺฐาย ปารายนํ\footnote{ขุ 1. 428; ขุ 8. 1 ปิฏฺฐาทีสุ.} สเรน ภาสติ.}

\prose{เตน โข ปน สมเยน เวสฺสวโณ มหาราชา อุตฺตราย ทิสาย ทกฺขิณํ ทิสํ คจฺฉติ เกนจิเทว กรณีเยน.\csromanspacer{} อสฺโสสิ โข เวสฺสวโณ มหาราชา นนฺทมาตาย อุปาสิกาย ปารายนํ สเรน ภาสนฺติยา, สุตฺวา กถาปริโยสานํ อาคมยมาโน อฏฺฐาสิ.}

\prose{อถ โข นนฺทมาตา อุปาสิกา ปารายนํ สเรน ภาสิตฺวา ตุณฺหี อโหสิ.\csromanspacer{} อถ โข เวสฺสวโณ มหาราชา นนฺทมาตาย อุปาสิกาย กถาปริโยสานํ วิทิตฺวา อพฺภานุโมทิ “สาธุ ภคินิ สาธุ ภคินี”ติ.\csromanspacer{} โก ปเนโส ภทฺรมุขาติ.\csromanspacer{} อหํ เต ภคินิ ภาตา เวสฺสวโณ มหาราชาติ.\csromanspacer{} สาธุ ภทฺรมุข เตน หิ โย เม อยํ ธมฺมปริยาโย}

\vspace{\tipitakaparskip}
\csromancenter{มหายญฺญวคฺค ภณิโต, อิทํ เต โหตุ อาติเถยฺยนฺติ.\csromanspacer{} สาธุ ภคินิ เอตญฺเจว เม โหตุ อาติเถยฺยํ.\csromanspacer{} เสฺวว\footnote{เสฺว จ (สี)} สาริปุตฺตโมคฺคลฺลานปฺปมุโข ภิกฺขุสํโฆ อกตปาตราโส เวฬุกณฺฑกํ อาคมิสฺสติ, ตญฺจ ภิกฺขุสํฆํ ปริวิสิตฺวา มม ทกฺขิณํ อาทิเสยฺยาสิ.\csromanspacer{} เอตญฺเจว\footnote{เอวญฺจ (สี, สฺยา), เอตญฺจ (?)} เม ภวิสฺสติ อาติเถยฺยนฺติ.}
\prose{\csromanfolio{445}อถ โข นนฺทมาตา อุปาสิกา ตสฺสา รตฺติยา อจฺจเยน สเก นิเวสเน ปณีตํ ขาทนียํ โภชนียํ ปฏิยาทาเปสิ.\csromanspacer{} อถ โข สาริปุตฺตโมคฺคลฺลานปฺปมุโข ภิกฺขุสํโฆ อกตปาตราโส เยน เวฬุกณฺฑโก ตทวสริ.\csromanspacer{} อถ โข นนฺทมาตา อุปาสิกา อญฺญตรํ ปุริสํ อามนฺเตสิ “เอหิ ตฺวํ อมฺโภ ปุริส อารามํ คนฺตฺวา ภิกฺขุสํฆสฺส กาลํ อาโรเจหิ ‘กาโล ภนฺเต อยฺยาย นนฺทมาตุยา นิเวสเน นิฏฺฐิตํ ภตฺตนฺ’ติ”. “เอวํ อยฺเย”ติ โข โส ปุริโส นนฺทมาตาย อุปาสิกาย ปฏิสฺสุตฺวา อารามํ คนฺตฺวา ภิกฺขุสํฆสฺส กาลํ อาโรเจสิ “กาโล ภนฺเต อยฺยาย นนฺทมาตุยา นิเวสเน นิฏฺฐิตํ ภตฺตนฺ”ติ.\csromanspacer{} อถ โข สาริปุตฺตโมคฺคลฺลานปฺปมุโข ภิกฺขุสํโฆ ปุพฺพณฺหาสมยํ นิวาเสตฺวา ปตฺตจีวรมาทาย เยน นนฺทมาตาย อุปาสิกาย นิเวสนํ เตนุปสงฺกมิ, อุปสงฺกมิตฺวา ปญฺญตฺเต อาสเน นิสีทิ.\csromanspacer{} อถ โข นนฺทมาตา อุปาสิกา สาริปุตฺตโมคฺคลฺลานปฺปมุขํ ภิกฺขุสํฆํ ปณีเตน ขาทนีเยน โภชนีเยน สหตฺถา สนฺตปฺเปสิ สมฺปวาเรสิ.}

\prose{อถ โข นนฺทมาตา อุปาสิกา อายสฺมนฺตํ สาริปุตฺตํ ภุตฺตาวิํ โอนีตปตฺตปาณิํ เอกมนฺตํ นิสีทิ, เอกมนฺตํ นิสินฺนํ โข นนฺทมาตรํ อุปาสิกํ อายสฺมา สาริปุตฺโต เอตทโวจ “โก ปน เต นนฺทมาเต ภิกฺขุสํฆสฺส อพฺภาคมนํ อาโรเจสี”ติ.}

\prose{อิธาหํ ภนฺเต รตฺติยา ปจฺจูสสมยํ ปจฺจุฏฺฐาย ปารายนํ สเรน ภาสิตฺวา ตุณฺหี อโหสิํ.\csromanspacer{} อถ โข ภนฺเต เวสฺสวโณ มหาราชา มม กถาปริโยสานํ วิทิตฺวา อพฺภานุโมทิ “สาธุ ภคินิ สาธุ ภคินี”ติ.\csromanspacer{} โก ปเนโส ภทฺรมุขาติ.\csromanspacer{} อหํ เต ภคินิ ภาตา เวสฺสวโณ มหาราชาติ.\csromanspacer{} สาธุ ภทฺรมุข เตน หิ โย เม อยํ}

\vspace{\tipitakaparskip}
\csromancenter{สตฺตกนิปาตปาฬิ ธมฺมปริยาโย ภณิโต, อิทํ เต โหตุ อาติเถยฺยนฺติ.\csromanspacer{} สาธุ ภคินิ เอตญฺเจว เม โหตุ อาติเถยฺยํ.\csromanspacer{} เสฺวว สาริปุตฺตโมคฺคลฺลานปฺปมุโข ภิกฺขุสํโฆ อกตปาตราโส เวฬุกณฺฑกํ อาคมิสฺสติ, ตญฺจ ภิกฺขุสํฆํ ปริวิสิตฺวา มม ทกฺขิณํ อาทิเสยฺยาสิ.\csromanspacer{} เอตญฺเจว\footnote{เอตญฺจ (สี), เอวญฺจ (สฺยา)} เม ภวิสฺสติ อาติเถยฺยนฺติ, ยทิทํ\footnote{ยมิทํ (ม 1. 299 ปิฏฺเฐ) 3-3. ปุญฺญํ หิ ตํ (สี), ปุญฺญํ ปุญฺญมหิตํ (สฺยา), ปุญฺญํ วา ปุญฺญมหํ วา (อิ), ปุญฺญํ วา ปุญฺญมหี วา (ก)} ภนฺเต ทาเน3 ปุญฺญญฺจ ปุญฺญมหี จ, ตํ3 เวสฺสวณสฺส มหาราชสฺส สุขาย โหตูติ.}
\prose{\csromanfolio{446}อจฺฉริยํ นนฺทมาเต, อพฺภุตํ นนฺทมาเต, ยตฺร หิ นาม เวสฺสวเณน มหาราเชน เอวํมหิทฺธิเกน เอวํมเหสกฺเขน เทวปุตฺเตน สมฺมุขา สลฺลปิสฺสสีติ. (1)}

\prose{น โข เม ภนฺเต เอเสว อจฺฉริโย อพฺภุโต ธมฺโม, อตฺถิ เม อญฺโญปิ อจฺฉริโย อพฺภุโต ธมฺโม.\csromanspacer{} อิธ เม ภนฺเต นนฺโท นาม เอกปุตฺตโก ปิโย มนาโป, ตํ ราชาโน กิสฺมิญฺจิเทว ปกรเณ โอกสฺส ปสยฺห ชีวิตา โวโรเปสุํ.\csromanspacer{} ตสฺมิํ โข ปนาหํ ภนฺเต ทารเก คหิเต วา คยฺหมาเน วา วเธ วา วชฺฌมาเน วา หเต วา หญฺญมาเน วา นาภิชานามิ จิตฺตสฺส อญฺญถตฺตนฺติ.\csromanspacer{} อจฺฉริยํ นนฺทมาเต, อพฺภุตํ นนฺทมาเต.\csromanspacer{} ยตฺร หิ นาม จิตฺตุปฺปาทมฺปิ\footnote{จิตฺตุปฺปาทมตฺตมฺปิ (สฺยา)} ปริโสเธสฺสสีติ. (2)}

\prose{น โข เม ภนฺเต เอเสว อจฺฉริโย อพฺภุโต ธมฺโม, อตฺถิ เม อญฺโญปิ อจฺฉริโย อพฺภุโต ธมฺโม.\csromanspacer{} อิธ เม ภนฺเต สามิโก กาลงฺกโต อญฺญตรํ ยกฺขโยนิํ อุปปนฺโน, โส เม เตเนว ปุริเมน อตฺตภาเวน อุทฺทสฺเสสิ.\csromanspacer{} น โข ปนาหํ ภนฺเต อภิชานามิ ตโตนิทานํ จิตฺตสฺส อญฺญถตฺตนฺติ.\csromanspacer{} อจฺฉริยํ นนฺทมาเต, อพฺภุตํ นนฺทมาเต, ยตฺร หิ นาม จิตฺตุปฺปาทมฺปิ ปริโสเธสฺสสีติ. (3)}

\prose{น โข เม ภนฺเต เอเสว อจฺฉริโย อพฺภุโต ธมฺโม, อตฺถิ เม อญฺโญปิ อจฺฉริโย อพฺภุโต ธมฺโม.\csromanspacer{} ยโตหํ ภนฺเต สามิกสฺส ทหรสฺเสว ทหรา อานีตา.\csromanspacer{} นาภิชานามิ สามิกํ มนสาปิ อติจริตา\footnote{อติจริตุํ (สฺยา), อติจาริตฺตํ (ก)} กุโต ปน กาเยนาติ.\csromanspacer{} อจฺฉริยํ นนฺทมาเต, อพฺภุตํ นนฺทมาเต, ยตฺร หิ นาม จิตฺตุปฺปาทมฺปิ ปริโสเธสฺสสีติ. (4)}

\chapterheadpageread{5. มหายญฺญวคฺค}
\prose{\csromanfolio{447}น โข เม ภนฺเต เอเสว อจฺฉริโย อพฺภุโต ธมฺโม, อตฺถิ เม อญฺโญปิ อจฺฉริโย อพฺภุโต ธมฺโม.\csromanspacer{} ยทาหํ ภนฺเต อุปาสิกา ปฏิเทสิตา, นาภิชานามิ กิญฺจิ สิกฺขาปทํ สญฺจิจฺจ วีติกฺกมิตาติ.\csromanspacer{} อจฺฉริยํ นนฺทมาเต, อพฺภุตํ นนฺทมาเตติ. (5)}

\prose{น โข เม ภนฺเต เอเสว อจฺฉริโย อพฺภุโต ธมฺโม, อตฺถิ เม อญฺโญปิ อจฺฉริโย อพฺภุโต ธมฺโม.\csromanspacer{} อิธาหํ ภนฺเต ยาวเท\footnote{ยาวเทว (สี, สฺยา) สํ 1. 412 ปิฏฺเฐ ปาฬิ จ อฏฺฐกถาฏีกา จ ปสฺสิตพฺพา.} อากงฺขามิ วิวิจฺเจว กาเมหิ วิวิจฺจ อกุสเลหิ ธมฺเมหิ สวิตกฺกํ สวิจารํ วิเวกชํ ปีติสุขํ ปฐมํ ฌานํ อุปสมฺปชฺช วิหรามิ.\csromanspacer{} วิตกฺกวิจารานํ วูปสมา อชฺฌตฺตํ สมฺปสาทนํ เจตโส เอโกทิภาวํ อวิตกฺกํ อวิจารํ วิเวกชํ ปีติสุขํ ทุติยํ ฌานํ อุปสมฺปชฺช วิหรามิ.\csromanspacer{} ปีติยา จ วิราคา อุเปกฺขิกา จ วิหรามิ สตา จ สมฺปชานา, สุขญฺจ กาเยน ปฏิสํเวเทมิ, ยํ ตํ อริยา อาจิกฺขนฺติ “อุเปกฺขโก สติมา สุขวิหารี”ติ ตติยํ ฌานํ อุปสมฺปชฺช วิหรามิ.\csromanspacer{} สุขสฺส จ ปหานา ทุกฺขสฺส จ ปหานา ปุพฺเพว โสมนสฺสโทมนสฺสานํ อตฺถงฺคมา อทุกฺขมสุขํ อุเปกฺขาสติปาริสุทฺธิํ จตุตฺถํ ฌานํ อุปสมฺปชฺช วิหรามีติ.\csromanspacer{} อจฺฉริยํ นนฺทมาเต, อพฺภุตํ นนฺทมาเตติ. (6)}

\prose{น โข เม ภนฺเต เอเสว อจฺฉริโย อพฺภุโต ธมฺโม, อตฺถิ เม อญฺโญปิ อจฺฉริโย อพฺภุโต ธมฺโม.\csromanspacer{} ยานิมานิ ภนฺเต ภควตา เทสิตานิ ปญฺโจรมฺภาคิยานิ สํโยชนานิ, นาหํ เตสํ กิญฺจิ อตฺตนิ อปฺปหีนํ สมนุปสฺสามีติ.\csromanspacer{} อจฺฉริยํ นนฺทมาเต, อพฺภุตํ นนฺทมาเตติ. (7)}

\prose{อถ โข อายสฺมา สาริปุตฺโต นนฺทมาตรํ อุปาสิกํ ธมฺมิยา กถาย สนฺทสฺเสตฺวา สมาทเปตฺวา สมุตฺเตเชตฺวา สมฺปหํเสตฺวา อุฏฺฐายาสนา ปกฺกามีติ.\csromanspacer{}\csromanspacer{}ทสมํ.\csromanspacer{} มหายญฺญวคฺโค ปญฺจโม.}

\titlehead{\textbf{ตสฺสุทฺทานํ}}
\csromangathameasure{ฐิติ จ ปริกฺขารํ เทฺว, อคฺคี สญฺญา จ เทฺว ปรา. \\ เมถุนา สํโยโค ทานํ, นนฺทมาเตน เต ทสาติ. \\ \textbf{ปณฺณาสโก สมตฺโต.}}
\csromangathagroup{\csromangathastanza{ฐิติ จ ปริกฺขารํ เทฺว, อคฺคี สญฺญา จ เทฺว ปรา. \\ เมถุนา สํโยโค ทานํ, นนฺทมาเตน เต ทสาติ. \\ \textbf{ปณฺณาสโก สมตฺโต.}}}

\csromanmatikaanchor{mtk.248}
\csromantocmarknopage{chapter}{6. อพฺยากตวคฺค}{mtk.248}
\bookmark[dest={mtk.248},level=0]{6. อพฺยากตวคฺค}
\chapterheadpageread{สตฺตกนิปาตปาฬิ \textbf{6. อพฺยากตวคฺค}}
\csromanmatikaanchor{mtk.249}
\csromantocmarkref{subsection}{1. อพฺยากตสุตฺต}{mtk.249}
\bookmark[dest={mtk.249},level=2]{1. อพฺยากตสุตฺต}
\csromanheader{2}{1. อพฺยากตสุตฺต}
\proseitem{54}{\csromanfolio{448}อถ โข อญฺญตโร ภิกฺขุ เยน ภควา เตนุปสงฺกมิ, อุปสงฺกมิตฺวา ภควนฺตํ อภิวาเทตฺวา เอกมนฺตํ นิสีทิ, เอกมนฺตํ นิสินฺโน โข โส ภิกฺขุ ภควนฺตํ เอตทโวจ “โก นุ โข ภนฺเต เหตุ โก ปจฺจโย, เยน สุตวโต อริยสาวกสฺส วิจิกิจฺฉา นุปฺปชฺชติ อพฺยากตวตฺถูสู”ติ.}

\prose{ทิฏฺฐินิโรธา โข ภิกฺขุ สุตวโต อริยสาวกสฺส วิจิกิจฺฉา นุปฺปชฺชติ อพฺยากตวตฺถูสุ, “โหติ ตถาคโต ปรํ มรณา”ติ โข ภิกฺขุ ทิฏฺฐิคตเมตํ, “น โหติ ตถาคโต ปรํ มรณา”ติ โข ภิกฺขุ ทิฏฺฐิคตเมตํ, “โหติ จ น จ โหติ ตถาคโต ปรํ มรณา”ติ โข ภิกฺขุ ทิฏฺฐิคตเมตํ, “เนว โหติ น น โหติ ตถาคโต ปรํ มรณา”ติ โข ภิกฺขุ ทิฏฺฐิคตเมตํ.\csromanspacer{} อสฺสุตวา ภิกฺขุ ปุถุชฺชโน ทิฏฺฐิํ นปฺปชานาติ, ทิฏฺฐิสมุทยํ นปฺปชานาติ, ทิฏฺฐินิโรธํ นปฺปชานาติ, ทิฏฺฐินิโรธคามินิํ ปฏิปทํ นปฺปชานาติ.\csromanspacer{} ตสฺส สา ทิฏฺฐิ ปวฑฺฒติ, โส น ปริมุจฺจติ ชาติยา ชราย มรเณน โสเกหิ ปริเทเวหิ ทุกฺเขหิ โทมนสฺเสหิ อุปายาเสหิ, น ปริมุจฺจติ ทุกฺขสฺมาติ วทามิ.}

\prose{สุตวา จ โข ภิกฺขุ อริยสาวโก ทิฏฺฐิํ ปชานาติ, ทิฏฺฐิสมุทยํ ปชานาติ, ทิฏฺฐินิโรธํ ปชานาติ, ทิฏฺฐินิโรธคามินิํ ปฏิปทํ ปชานาติ.\csromanspacer{} ตสฺส สา ทิฏฺฐิ นิรุชฺฌติ, โส ปริมุจฺจติ ชาติยา ชราย มรเณน โสเกหิ ปริเทเวหิ ทุกฺเขหิ โทมนสฺเสหิ อุปายาเสหิ, ปริมุจฺจติ ทุกฺขสฺมาติ วทามิ.\csromanspacer{} เอวํ ชานํ โข ภิกฺขุ สุตวา อริยสาวโก เอวํ ปสฺสํ “โหติ ตถาคโต ปรํ มรณา”ติปิ น พฺยากโรติ, “น โหติ ตถาคโต ปรํ มรณา”ติปิ น พฺยากโรติ, “โหติ จ น จ โหติ ตถาคโต ปรํ มรณา”ติปิ น พฺยากโรติ, “เนว โหติ น น โหติ ตถาคโต ปรํ มรณา”ติปิ น พฺยากโรติ.\csromanspacer{} เอวํ ชานํ โข ภิกฺขุ สุตวา อริยสาวโก เอวํ ปสฺสํ เอวํ อพฺยากรณธมฺโม โหติ อพฺยากตวตฺถูสุ.\csromanspacer{} เอวํ ชานํ โข ภิกฺขุ สุตวา อริยสาวโก เอวํ ปสฺสํ น ฉมฺภติ น กมฺปติ น เวธติ น สนฺตาสํ อาปชฺชติ อพฺยากตวตฺถูสุ.}

\chapterheadpageread{6. อพฺยากตวคฺค}
\prose{\csromanfolio{449}“โหติ ตถาคโต ปรํ มรณา”ติ โข ภิกฺขุ ตณฺหาคตเมตํ ฯเปฯ.\csromanspacer{} สญฺญาคตเมตํ ฯเปฯ.\csromanspacer{} มญฺญิตเมตํ ฯเปฯ.\csromanspacer{} ปปญฺจิตเมตํ ฯเปฯ.\csromanspacer{} อุปาทานคตเมตํ ฯเปฯ. “โหติ ตถาคโต ปรํ มรณา”ติ โข ภิกฺขุ วิปฺปฏิสาโร เอโส, “น โหติ ตถาคโต ปรํ มรณา”ติ โข ภิกฺขุ วิปฺปฏิสาโร เอโส, “โหติ จ น จ โหติ ตถาคโต ปรํ มรณา”ติ โข ภิกฺขุ วิปฺปฏิสาโร เอโส, “เนว โหติ น น โหติ ตถาคโต ปรํ มรณา”ติ โข ภิกฺขุ วิปฺปฏิสาโร เอโส.\csromanspacer{} อสฺสุตวา ภิกฺขุ ปุถุชฺชโน วิปฺปฏิสารํ นปฺปชานาติ, วิปฺปฏิสารสมุทยํ นปฺปชานาติ, วิปฺปฏิสารนิโรธํ นปฺปชานาติ, วิปฺปฏิสารนิโรธคามินิํ ปฏิปทํ นปฺปชานาติ.\csromanspacer{} ตสฺส โส วิปฺปฏิสาโร ปวฑฺฒติ, โส น ปริมุจฺจติ ชาติยา ชราย มรเณน โสเกหิ ปริเทเวหิ ทุกฺเขหิ โทมนสฺเสหิ อุปายาเสหิ, น ปริมุจฺจติ ทุกฺขสฺมาติ วทามิ.}

\prose{สุตวา จ โข ภิกฺขุ อริยสาวโก วิปฺปฏิสารํ ปชานาติ, วิปฺปฏิสารสมุทยํ ปชานาติ, วิปฺปฏิสารนิโรธํ ปชานาติ, วิปฺปฏิสารนิโรธคามินิํ ปฏิปทํ ปชานาติ.\csromanspacer{} ตสฺส โส วิปฺปฏิสาโร นิรุชฺฌติ, โส ปริมุจฺจติ ชาติยา ฯเปฯ ทุกฺขสฺมาติ วทามิ.\csromanspacer{} เอวํ ชานํ โข ภิกฺขุ สุตวา อริยสาวโก เอวํ ปสฺสํ “โหติ ตถาคโต ปรํ มรณา”ติปิ น พฺยากโรติ ฯเปฯ “เนว โหติ น น โหติ ตถาคโต ปรํ มรณา”ติปิ น พฺยากโรติ.\csromanspacer{} เอวํ ชานํ โข ภิกฺขุ สุตวา อริยสาวโก เอวํ ปสฺสํ เอวํ อพฺยากรณธมฺโม โหติ อพฺยากตวตฺถูสุ.\csromanspacer{} เอวํ ชานํ โข ภิกฺขุ สุตวา อริยสาวโก เอวํ ปสฺสํ น ฉมฺภติ น กมฺปติ น เวธติ น สนฺตาสํ อาปชฺชติ อพฺยากตวตฺถูสุ.\csromanspacer{} อยํ โข ภิกฺขุ เหตุ อยํ ปจฺจโย, เยน สุตวโต อริยสาวกสฺส วิจิกิจฺฉา นปฺปชฺชติ อพฺยากตวตฺถูสูติ.\csromanspacer{}\csromanspacer{}ปฐมํ.}

\csromanmatikaanchor{mtk.250}
\csromantocmarkref{subsection}{2. ปุริสคติสุตฺต}{mtk.250}
\bookmark[dest={mtk.250},level=2]{2. ปุริสคติสุตฺต}
\csromanheader{2}{2. ปุริสคติสุตฺต}
\proseitem{55}{สตฺต จ\footnote{สตฺต (สี), สตฺต จ โข (ก)} ภิกฺขเว ปุริสคติโย เทเสสฺสามิ อนุปาทา จ ปรินิพฺพานํ\footnote{ปรินิพฺพาณํ (สี)}, ตํ สุณาถ สาธุกํ มนสิ กโรถ, ภาสิสฺสามีติ. “เอวํ ภนฺเต”ติ โข เต ภิกฺขู ภควโต ปจฺจสฺโสสุํ.\csromanspacer{} ภควา เอตทโวจ\csromandash{} กตมา จ ภิกฺขเว สตฺต ปุริสคติโย.}

\gambhira{สตฺตกนิปาตปาฬิ}
\prose{\csromanfolio{450}อิธ ภิกฺขเว ภิกฺขุ เอวํ ปฏิปนฺโน โหติ “โน จสฺส โน จ เม สิยา, น ภวิสฺสติ น เม ภวิสฺสติ, ยทตฺถิ ยํ ภูตํ ตํ ปชหามี”ติ อุเปกฺขํ ปฏิลภติ.\csromanspacer{} โส ภเว น รชฺชติ, สมฺภเว น รชฺชติ, อตฺถุตฺตริ ปทํ สนฺตํ สมฺมปฺปญฺญาย ปสฺสติ.\csromanspacer{} ตญฺจ ขฺวสฺส ปทํ น สพฺเพน สพฺพํ สจฺฉิกตํ โหติ, ตสฺส น สพฺเพน สพฺพํ มานานุสโย ปหีโน โหติ, น สพฺเพน สพฺพํ ภวราคานุสโย ปหีโน โหติ, น สพฺเพน สพฺพํ อวิชฺชานุสโย ปหีโน โหติ.\csromanspacer{} โส ปญฺจนฺนํ โอรมฺภาคิยานํ สํโยชนานํ ปริกฺขยา อนฺตราปรินิพฺพายี โหติ.\csromanspacer{} เสยฺยถาปิ ภิกฺขเว ทิวสํสนฺตตฺเต\footnote{ทิวสสนฺตตฺเต (สี, สฺยา) ม 2. 117 ปิฏฺเฐปิ.} อโยกปาเล หญฺญมาเน ปปฏิกา นิพฺพตฺติตฺวา นิพฺพาเยยฺย.\csromanspacer{} เอวเมวํ โข ภิกฺขเว ภิกฺขุ เอวํ ปฏิปนฺโน โหติ “โน จสฺส โน จ เม สิยา, น ภวิสฺสติ น เม ภวิสฺสติ, ยทตฺถิ ยํ ภูตํ ตํ ปชหามี”ติ อุเปกฺขํ ปฏิลภติ.\csromanspacer{} โส ภเว น รชฺชติ, สมฺภเว น รชฺชติ, อตฺถุตฺตริ ปทํ สนฺตํ สมฺมปฺปญฺญาย ปสฺสติ.\csromanspacer{} ตญฺจ ขฺวสฺส ปทํ น สพฺเพน สพฺพํ สจฺฉิกตํ โหติ, ตสฺส น สพฺเพน สพฺพํ มานานุสโย ปหีโน โหติ, น สพฺเพน สพฺพํ ภวราคานุสโย ปหีโน โหติ, น สพฺเพน สพฺพํ อวิชฺชานุสโย ปหีโน โหติ.\csromanspacer{} โส ปญฺจนฺนํ โอรมฺภาคิยานํ สํโยชนานํ ปริกฺขยา อนฺตราปรินิพฺพายี โหติ. (1)}

\prose{อิธ ปน ภิกฺขเว ภิกฺขุ เอวํ ปฏิปนฺโน โหติ “โน จสฺส โน จ เม สิยา, น ภวิสฺสติ น เม ภวิสฺสติ, ยทตฺถิ ยํ ภูตํ ตํ ปชหามี”ติ อุเปกฺขํ ปฏิลภติ.\csromanspacer{} โส ภเว น รชฺชติ, สมฺภเว น รชฺชติ, อตฺถุตฺตริ ปทํ สนฺตํ สมฺมปฺปญฺญาย ปสฺสติ.\csromanspacer{} ตญฺจ ขฺวสฺส ปทํ น สพฺเพน สพฺพํ สจฺฉิกตํ โหติ, ตสฺส น สพฺเพน สพฺพํ มานานุสโย ปหีโน โหติ, น สพฺเพน สพฺพํ ภวราคานุสโย ปหีโน โหติ, น สพฺเพน สพฺพํ อวิชฺชานุสโย ปหีโน โหติ.\csromanspacer{} โส ปญฺจนฺนํ โอรมฺภาคิยานํ สํโยชนานํ ปริกฺขยา อนฺตราปรินิพฺพายี โหติ.\csromanspacer{} เสยฺยถาปิ ภิกฺขเว ทิวสํสนฺตตฺเต อโยกปาเล หญฺญมาเน ปปฏิกา นิพฺพตฺติตฺวา อุปฺปติตฺวา นิพฺพาเยยฺย.\csromanspacer{} เอวเมวํ โข ภิกฺขเว ภิกฺขุ เอวํ ปฏิปนฺโน โหติ “โน จสฺส, โน จ เม สิยา ฯเปฯ.\csromanspacer{} โส ปญฺจนฺนํ โอรมฺภาคิยานํ สํโยชนานํ ปริกฺขยา อนฺตราปรินิพฺพายี โหติ. (2)}

\chapterheadpageread{6. อพฺยากตวคฺค}
\prose{\csromanfolio{451}อิธ ปน ภิกฺขเว ภิกฺขุ เอวํ ปฏิปนฺโน โหติ “โน จสฺส, โน จ เม สิยา ฯเปฯ.\csromanspacer{} โส ปญฺจนฺนํ โอรมฺภาคิยานํ สํโยชนานํ ปริกฺขยา อนฺตราปรินิพฺพายี โหติ.\csromanspacer{} เสยฺยถาปิ ภิกฺขเว ทิวสํสนฺตตฺเต อาโยกปาเล หญฺญมาเน ปปฏิกา นิพฺพตฺติตฺวา อุปฺปติตฺวา อนุปหจฺจ ตลํ นิพฺพาเยยฺย.\csromanspacer{} เอวเมวํ โข ภิกฺขเว ภิกฺขุ เอวํ ปฏิปนฺโน โหติ “โน จสฺส, โน จ เม สิยา ฯเปฯ.\csromanspacer{} โส ปญฺจนฺนํ โอรมฺภาคิยานํ สํโยชนานํ ปริกฺขยา อนฺตราปรินิพฺพยี โหติ. (3)}

\prose{อิธ ปน ภิกฺขเว ภิกฺขุ เอวํ ปฏิปนฺโน โหติ “โน จสฺส, โน จ เม สิยา ฯเปฯ.\csromanspacer{} โส ปญฺจนฺนํ โอรมฺภาคิยานํ สํโยชนานํ ปริกฺขยา อุปหจฺจปรินิพฺพายี โหติ.\csromanspacer{} เสยฺยถาปิ ภิกฺขเว ทิวสํสนฺตตฺเต อโยกปาเล หญฺญมาเน ปปฏิกา นิพฺพตฺติตฺวา อุปฺปติตฺวา อุปหจฺจ ตลํ นิพฺพาเยยฺย.\csromanspacer{} เอวเมวํ โข ภิกฺขเว ภิกฺขุ เอวํ ปฏิปนฺโน โหติ “โน จสฺส โน จ เม สิยา ฯเปฯ.\csromanspacer{} โส ปญฺจนฺนํ โอรมฺภาคิยานํ สํโยชนานํ ปริกฺขยา อุปหจฺจปรินิพฺพายี โหติ. (4)}

\prose{อิธ ปน ภิกฺขเว ภิกฺขุ เอวํ ปฏิปนฺโน โหติ “โน จสฺส โน จ เม สิยา ฯเปฯ.\csromanspacer{} โส ปญฺจนฺนํ โอรมฺภาคิยานํ สํโยชนานํ ปริกฺขยา อสงฺขารปรินิพฺพายี โหติ.\csromanspacer{} เสยฺยถาปิ ภิกฺขเว ทิวสํสนฺตตฺเต อโยกปาเล หญฺญมาเน ปปฏิกา นิพฺพตฺติตฺวา อุปฺปติตฺวา ปริตฺเต ติณปุญฺเช วา กฏฺฐปุญฺเช วา นิปเตยฺย, สา ตตฺถ อคฺคิมฺปิ ชเนยฺย, ธูมมฺปิ ชเนยฺย, อคฺคิมฺปิ ชเนตฺวา ธูมมฺปิ ชเนตฺวา ตเมว ปริตฺตํ ติณปุญฺชํ วา กฏฺฐปุญฺชํ วา ปริยาทิยิตฺวา อนาหารา นิพฺพาเยยฺย.\csromanspacer{} เอวเมวํ โข ภิกฺขเว ภิกฺขุ เอวํ ปฏิปนฺโน โหติ “โน จสฺส โน จ เม สิยา ฯเปฯ.\csromanspacer{} โส ปญฺจนฺนํ โอรมฺภาคิยานํ สํโยชนานํ ปริกฺขยา อสงฺขารปรินิพฺพายี โหติ. (5)}

\prose{อิธ ปน ภิกฺขเว ภิกฺขุ เอวํ ปฏิปนฺโน โหติ “โน จสฺส, โน จ เม สิยา ฯเปฯ.\csromanspacer{} โส ปญฺจนฺนํ โอรมฺภาคิยานํ สํโยชนานํ ปริกฺขยา สสงฺขารปรินิพฺพายี โหติ.\csromanspacer{} เสยฺยถาปิ ภิกฺขเว ทิวสํสนฺตตฺเต อโยกปาเล หญฺญมาเน ปปฏิกา นิพฺพตฺติตฺวา อุปฺปติตฺวา วิปุเล ติณปุญฺเช วา กฏฺฐปุญฺเช วา นิปเตยฺย, สา ตตฺถ อคฺคิมฺปิ ชเนยฺย, ธูมมฺปิ ชเนยฺย, อคฺคิมฺปิ ชเนตฺวา ธูมมฺปิ ชเนตฺวา ตเมว วิปุลํ ติณปุญฺชํ วา กฏฺฐปุญฺชํ วา ปริยาทิยิตฺวา อนาหารา นิพฺพาเยยฺย.\csromanspacer{} เอวเมวํ โข ภิกฺขเว}

\vspace{\tipitakaparskip}
\csromancenter{สตฺตกนิปาตปาฬิ ภิกฺขุ เอวํ ปฏิปนฺโน โหติ “โน จสฺส, โน จ เม สิยา ฯเปฯ.\csromanspacer{} โส ปญฺจนฺนํ โอรมฺภาคิยานํ สํโยชนานํ ปริกฺขยา สสงฺขารปรินิพฺพายี โหติ. (6)}
\prose{\csromanfolio{452}อิธ ปน ภิกฺขเว ภิกฺขุ เอวํ ปฏิปนฺโน โหติ “โน จสฺส, โน จ เม สิยา, น ภวิสฺสติ, น เม ภวิสฺสติ, ยทตฺถิ ยํ ภูตํ, ตํ ปชหามี”ติ อุเปกฺขํ ปฏิลภติ.\csromanspacer{} โส ภเว น รชฺชติ, สมฺภเว น รชฺชติ, อตฺถุตฺตริ ปทํ สนฺตํ สมฺมปญฺญาย ปสฺสติ.\csromanspacer{} ตญฺจ ขฺวสฺส ปทํ น สพฺเพน สพฺพํ สจฺฉิกตํ โหติ, ตสฺส น สพฺเพน สพฺพํ มานานุสโย ปหีโน โหติ, น สพฺเพน สพฺพํ ภวราคานุสโย ปหีโน โหติ, น สพฺเพน สพฺพํ อวิชฺชานุสโย ปหีโน โหติ.\csromanspacer{} โส ปญฺจนฺนํ โอรมฺภาคิยานํ สํโยชนานํ ปริกฺขยา อุทฺธํโสโต โหติ อกนิฏฺฐคามี.\csromanspacer{} เสยฺยถาปิ ภิกฺขเว ทิวสํสนฺตตฺเต อโยกปาเล หญฺญมาเน ปปฏิกา นิพฺพตฺติตฺวา อุปฺปติตฺวา มหนฺเต ติณปุญฺเช วา กฏฺฐปุญฺเช วา นิปเตยฺย, สา ตตฺถ อคฺคิมฺปิ ชเนยฺย, ธูมมฺปิ ชเนยฺย, อคฺคิมฺปิ ชเนตฺวา ธูมมฺปิ ชเนตฺวา ตเมว มหนฺตํ ติณปุญฺชํ วา กฏฺฐปุญฺชํ วา ปริยาทิยิตฺวา คจฺฉมฺปิ ทเหยฺย\footnote{qอเหยฺย (อญฺญตฺถ)}, ทายมฺปิ ทเหยฺย, คจฺฉมฺปิ ทหิตฺวา ทายมฺปิ ทหิตฺวา หริตนฺตํ วา ปถนฺตํ วา\footnote{ปนฺถนฺตํ วา (สี) สฺยามโปตฺถเก อิทํ น ทิสฺสติ.} เสลนฺตํ วา อุทกนฺตํ วา รมณียํ วา ภูมิภาคํ อาคมฺม อนาหารา นิพฺพาเยยฺย.\csromanspacer{} เอวเมวํ โข ภิกฺขเว ภิกฺขุ เอวํ ปฏิปนฺโน โหติ “โน จสฺส, โน จ เม สิยา ฯเปฯ.\csromanspacer{} โส ปญฺจนฺนํ โอรมฺภาคิยานํ สํโยชนานํ ปริกฺขยา อุทฺธํโสโต โหติ อกนิฏฺฐคามี.\csromanspacer{} อิมา โข ภิกฺขเว สตฺต ปุริสคติโย.}

\prose{กตมญฺจ ภิกฺขเว อนุปาทาปรินิพฺพานํ, อิธ ภิกฺขเว ภิกฺขุ เอวํ ปฏิปนฺโน โหติ “โน จสฺส, โน จ เม สิยา, น ภวิสฺสติ, น เม ภวิสฺสติ, ยทตฺถิ ยํ ภูตํ, ตํ ปชหามี”ติ อุเปกฺขํ ปฏิลภติ.\csromanspacer{} โส ภเว น รชฺชติ, สมฺภเว น รชฺชติ, อตฺถุตฺตริ ปทํ สนฺตํ สมฺมปฺปญฺญาย ปสฺสติ.\csromanspacer{} ตญฺจ ขฺวสฺส ปทํ สพฺเพน สพฺพํ สจฺฉิกตํ โหติ, ตสฺส สพฺเพน สพฺพํ มานานุสโย ปหีโน โหติ, สพฺเพน สพฺพํ ภวราคานุสโย ปหีโน โหติ, สพฺเพน สพฺพํ อวิชฺชานุสโย ปหีโน โหติ, โส อาสวานํ ขยา ฯเปฯ สจฺฉิกตฺวา อุปสมฺปชฺช วิหรติ.\csromanspacer{} อิทํ วุจฺจติ ภิกฺขเว อนุปาทาปรินิพฺพานํ.\csromanspacer{} อิมา โข ภิกฺขเว สตฺต ปุริสคติโย อนุปาทา จ ปรินิพฺพานนฺติ.\csromanspacer{}\csromanspacer{}ทุติยํ.}

\csromanmatikaanchor{mtk.251}
\csromantocmarkref{subsection}{3. ติสฺสพฺรหฺมาสุตฺต}{mtk.251}
\bookmark[dest={mtk.251},level=2]{3. ติสฺสพฺรหฺมาสุตฺต}
\csromanheader{2}{6. อพฺยากตวคฺค \textbf{3. ติสฺสพฺรหฺมาสุตฺต}}
\proseitem{56}{\csromanfolio{453}เอวํ เม สุตํ\csromandash{}เอกํ สมยํ ภควา ราชคเห วิหรติ คิชฺฌกูเฏ ปพฺพเต.\csromanspacer{} อถ โข เทฺว เทวตา อภิกฺกนฺตาย รตฺติยา อภิกฺกนฺตวณฺณา เกวลกปฺปํ คิชฺฌกูฏํ โอภาเสตฺวา เยน ภควา เตนุปสงฺกมิํสุ, อุปสงฺกมิตฺวา ภควนฺตํ อภิวาเทตฺวา เอกมนฺตํ อฏฺฐํสุ, เอกมนฺตํ ฐิตา โข เอกา เทวตา ภควนฺตํ เอตทโวจ “เอตา ภนฺเต ภิกฺขุนิโย วิมุตฺตา”ติ.\csromanspacer{} อปรา เทวตา ภควนฺตํ เอตทโวจ “เอตา ภนฺเต ภิกฺขุนิโย อนุปาทิเสสา สุวิมุตฺตา”ติ.\csromanspacer{} อิทมโวจุํ ตา เทวตา, สมนุญฺโญ สตฺถา อโหสิ.\csromanspacer{} อถ โข ตา เทวตา “สมนุญฺโญ สตฺถา”ติ ภควนฺตํ อภิวาเทตฺวา ปทกฺขิณํ กตฺวา ตตฺเถวนฺตรธายิํสุ.}

\prose{อถ โข ภควา ตสฺสา รตฺติยา อจฺจเยน ภิกฺขู อามนฺเตสิ\csromandash{}อิมํ ภิกฺขเว รตฺติํ เทฺว เทวตา อภิกฺกนฺตาย รตฺติยา อภิกฺกนฺตวณฺณา เกวลกปฺปํ คิชฺฌกูฏํ โอภาเสตฺวา เยนาหํ เตนุปสงฺกมิํสุ, อุปสงฺกมิตฺวา มํ อภิวาเทตฺวา เอกมนฺตํ อฏฺฐํสุ, เอกมนฺตํ ฐิตา โข ภิกฺขเว เอกา เทวตา มํ เอตทโวจ “เอตา ภนฺเต ภิกฺขุนิโย วิมุตฺตา”ติ.\csromanspacer{} อปรา เทวตา มํ เอตทโวจ “เอตา ภนฺเต ภิกฺขุนิโย อนุปาทิเสสา สุวิมุตฺตา”ติ.\csromanspacer{} อิทมโวจุํ ภิกฺขเว ตา เทวตา, อิทํ วตฺวา มํ อภิวาเทตฺวา ปทกฺขิณํ กตฺวา ตตฺเถวนฺตรธายิํสูติ.}

\prose{เตน โข ปน สมเยน อายสฺมา มหาโมคฺคลฺลาโน ภควโต อวิทูเร นิสินฺโน โหติ.\csromanspacer{} อถ โข อายสฺมโต มหาโมคฺคลฺลานสฺส เอตทโหสิ “กตเมสานํ โข เทวานํ เอวํ ญาณํ โหติ ส-อุปาทิเสเส วา ‘ส-อุปาทิเสโส’ติ, อนุปาทิเสเส วา ‘อนุปาทิเสโส’ติ”.\csromanspacer{} เตน โข ปน สมเยน ติสฺโส นาม ภิกฺขุ อธุนากาลงฺกโต อญฺญตรํ พฺรหฺมโลกํ อุปปนฺโน โหติ.\csromanspacer{} ตตฺราปิ นํ เอวํ ชานนฺติ “ติสฺโส พฺรหฺมา มหิทฺธิโก มหานุภาโว”ติ.}

\prose{อถ โข อายสฺมา มหาโมคฺคลฺลาโน เสยฺยถาปิ นาม พลวา ปุริโส สมิญฺชิตํ วา พาหํ ปสาเรยฺย, ปสาริตํ วา พาหํ สมิญฺเชยฺย, เอวเมวํ คิชฺฌกูเฏ ปพฺพเต อนฺตรหิโต ตสฺมิํ พฺรหฺมโลเก ปาตุรโหสิ.\csromanspacer{} อทฺทสา โข ติสฺโส พฺรหฺมา อายสฺมนฺตํ มหาโมคฺคลฺลานํ}

\vspace{\tipitakaparskip}
\csromancenter{สตฺตกนิปาตปาฬิ ทูรโตว อาคจฺฉนฺตํ, ทิสฺวา อายสฺมนฺตํ มหาโมคฺคลฺลานํ เอตทโวจ “เอหิ โข มาริส โมคฺคลฺลาน, สฺวาคตํ มาริส โมคฺคลฺลาน, จิรสฺสํ โข มาริส โมคฺคลฺลาน อิมํ ปริยายมกาสิ, ยทิทํ อิธาคมนาย, นิสีท มาริส โมคฺคลฺลาน, อิทมาสนํ ปญฺญตฺตนฺ”ติ.\csromanspacer{} นิสีทิ โข อายสฺมา มหาโมคฺคลฺลาโน ปญฺญตฺเต อาสเน.\csromanspacer{} ติสฺโสปิ โข พฺรหฺมา อายสฺมนฺตํ มหาโมคฺคลฺลานํ อภิวาเทตฺวา เอกมนฺตํ นิสีทิ, เอกมนฺตํ นิสินฺนํ โข ติสฺสํ พฺรหฺมานํ อายสฺมา มหาโมคฺคลฺลาโน เอตทโวจ “กตเมสานํ โข ติสฺส เทวานํ เอวํ ญาณํ โหติ ส-อุปาทิเสเส วา ‘ส-อุปาทิเสโส’ติ, อนุปาทิเสเส วา ‘อนุปาทิเสโส’ติ”.\csromanspacer{} พฺรหฺมกายิกานํ โข มาริส โมคฺคลฺลาน เทวานํ เอวํ ญาณํ โหติ ส-อุปาทิเสเส วา “ส-อุปาทิเสโส”ติ, อนุปาทิเสเส วา “อนุปาทิเสโส”ติ.}
\prose{\csromanfolio{454}สพฺเพสญฺเญว โข ติสฺส พฺรหฺมกายิกานํ เทวานํ เอวํ ญาณํ โหติ ส-อุปาทิเสเส วา “ส-อุปาทิเสโส”ติ, “อนุปาทิเสเส วา “อนุปาทิเสโส”ติ.\csromanspacer{} น โข มาริส โมคฺคลฺลาน สพฺเพสํ พฺรหฺมกายิกานํ เทวานํ เอวํ ญาณํ โหติ ส-อุปาทิเสเส วา “ส-อุปาทิเสโส”ติ, อนุปาทิเสเส วา “อนุปาทิเสโส”ติ.}

\prose{เย โข เต มาริส โมคฺคลฺลาน พฺรหฺมกายิกา เทวา พฺรหฺเมน อายุนา สนฺตุฏฺฐา พฺรหฺเมน วณฺเณน พฺรหฺเมน สุเขน พฺรหฺเมน ยเสน พฺรหฺเมน อาธิปเตยฺเยน สนฺตุฏฺฐา.\csromanspacer{} เต อุตฺตริ นิสฺสรณํ ยถาภูตํ นปฺปชานนฺติ.\csromanspacer{} เตสํ น เอวํ ญาณํ โหติ ส-อุปาทิเสเส วา “ส-อุปาทิเสโส”ติ, อนุปาทิเสเส วา “อนุปาทิเสโส”ติ.\csromanspacer{} เย จ โข เต มาริส โมคฺคลฺลาน พฺรหฺมกายิกา เทวา พฺรหฺเมน อายุนา อสนฺตุฏฺฐา พฺรหฺเมน วณฺเณน พฺรหฺเมน สุเขน พฺรหฺเมน ยเสน พฺรหฺเมน อาธิปเตยฺเยน อสนฺตุฏฺฐา.\csromanspacer{} เต จ อุตฺตริ นิสฺสรณํ ยถาภูตํ ปชานนฺติ.\csromanspacer{} เตสํ เอวํ ญาณํ โหติ ส-อุปาทิเสเส วา “ส-อุปาทิเสโส”ติ, อนุปาทิเสเส วา “อนุปาทิเสโส”ติ.}

\prose{อิธ มาริส โมคฺคลฺลาน ภิกฺขุ อุภโตภาควิมุตฺโต โหติ, ตเมนํ เต เทวา เอวํ ชานนฺติ “อยํ โข อายสฺมา อุภโตภาควิมุตฺโต.\csromanspacer{} ยาวสฺส กาโย ฐสฺสติ, ตาว นํ ทกฺขนฺติ เทวมนุสฺสา.\csromanspacer{} กายสฺส เภทา น นํ ทกฺขนฺติ เทวมนุสฺสา”ติ.\csromanspacer{} เอวมฺปิ โข มาริส}

\vspace{\tipitakaparskip}
\csromancenter{อพฺยากตวคฺค โมคฺคลฺลาน เตสํ เทวานํ ญาณํ โหติ ส-อุปาทิเสเส วา “ส-อุปาทิเสโส”ติ, อนุปาทิเสเส วา “อนุปาทิเสโส”ติ.}
\prose{\csromanfolio{455}อิธ ปน มาริส โมคฺคลฺลาน ภิกฺขุ ปญฺญาวิมุตฺโต โหติ, ตเมนํ เต เทวา เอวํ ชานนฺติ “อยํ โข อายสฺมา ปญฺญาวิมุตฺโต ยาวสฺส กาโย ฐสฺสติ, ตาว นํ ทกฺขนฺติ เทวมนุสฺสา, กายสฺส เภทา น นํ ทกฺขนฺติ เทวมนุสฺสา”ติ.\csromanspacer{} เอวมฺปิ โข มาริส โมคฺคลฺลาน เตสํ เทวานํ ญาณํ โหติ ส-อุปาทิเสเส วา “ส-อุปาทิเสโส”ติ, อนุปาทิเสเส วา “อนุปาทิเสโส”ติ.}

\prose{อิธ ปน มาริส โมคฺคลฺลาน ภิกฺขุ กายสกฺขี โหติ, ตเมนํ เทวา เอวํ ชานนฺติ “อยํ โข อายสฺมา กายสกฺขี, อปฺเปว นาม อยมายสฺมา อนุโลมิกานิ เสนาสนานิ ปฏิเสวมาโน กลฺยาณมิตฺเต ภชมาโน อินฺทฺริยานิ สมนฺนานยมาโน ยสฺสตฺถาย กุลปุตฺตา สมฺมเทว อคารสฺมา อนคาริยํ ปพฺพชนฺติ, ตทนุตฺตรํ พฺรหฺมจริยปริโยสานํ ทิฏฺเฐว ธมฺเม สยํ อภิญฺญา สจฺฉิกตฺวา อุปสมฺปชฺช วิหเรยฺยา”ติ.\csromanspacer{} เอวมฺปิ โข มาริส โมคฺคลฺลาน เตสํ เทวานํ ญาณํ โหติ ส-อุปาทิเสเส วา “ส-อุปาทิเสโส”ติ อนุปาทิเสเส วา “อนุปาทิเสโส”ติ.}

\prose{อิธ ปน มาริส โมคฺคลฺลาน ภิกฺขุ ทิฏฺฐิปฺปตฺโต โหติ ฯเปฯ.\csromanspacer{} สทฺธาวิมุตฺโต โหติ ฯเปฯ.\csromanspacer{} ธมฺมานุสารี โหติ, ตเมนํ เต เทวา เอวํ ชานนฺติ “อยํ โข อายสฺมา ธมฺมานุสารี, อปฺเปว นาม อยมายสฺมา อนุโลมิกานิ เสนาสนานิ ปฏิเสวมาโน กลฺยาณมิตฺเต ภชมาโน อินฺทฺริยานิ สมนฺนานยมาโน ยสฺสตฺถาย กุลปุตฺตา สมฺมเทว อคารสฺมา อนคาริยํ ปพฺพชนฺติ, ตทนุตฺตรํ พฺรหฺมจริยปริโยสานํ ทิฏฺเฐว ธมฺเม สยํ อภิญฺญา สจฺฉิกตฺวา อุปสมฺปชฺช วิหเรยฺยา”ติ.\csromanspacer{} เอวมฺปิ โข มาริส โมคฺคลฺลาน เตสํ เทวานํ ญาณํ โหติ ส-อุปาทิเสเส วา “ส-อุปาทิเสโส”ติ, อนุปาทิเสเส วา “อนุปาทิเสโส”ติ.}

\prose{อถ โข อายสฺมา มหาโมคฺคลฺลาโน ติสฺสสฺส พฺรหฺมุโน ภาสิตํ อภินนฺทิตฺวา อนุโมทิตฺวา เสยฺยถาปิ นาม พลวา ปุริโส สมิญฺชิตํ วา พาหํ ปสาเรยฺย, ปสาริตํ วา พาหํ สมิญฺเชยฺย.\csromanspacer{} เอวเมวํ พฺรหฺมโลเก อนฺตรหิโต คิชฺฌกูเฏ ปพฺพเต ปาตุรโหสิ.\csromanspacer{} อถ โข อายสฺมา มหาโมคฺคลฺลาโน เยน ภควา เตนุปสงฺกมิ, อุปสงฺกมิตฺวา ภควนฺตํ}

\vspace{\tipitakaparskip}
\csromancenter{สตฺตกนิปาตปาฬิ อภิวาเทตฺวา เอกมนฺตํ นิสีทิ, เอกมนฺตํ นิสินฺโน โข อายสฺมา มหาโมคฺคลฺลาโน ยาวตโก อโหสิ ติสฺเสน พฺรหฺมุนา สทฺธิํ กถาสลฺลาโป, ตํ สพฺพํ ภควโต อาโรเจสิ.}
\prose{\csromanfolio{456}น หิ ปน เต โมคฺคลฺลาน ติสฺโส พฺรหฺมา สตฺตมํ อนิมิตฺตวิหาริํ ปุคฺคลํ เทเสติ.\csromanspacer{} เอตสฺส ภควา กาโล, เอตสฺส สุคต กาโล, ยํ ภควา สตฺตมํ อนิมิตฺตวิหาริํ ปุคฺคลํ เทเสยฺย.\csromanspacer{} ภควโต สุตฺวา ภิกฺขู ธาเรสฺสนฺตีติ.\csromanspacer{} เตน หิ โมคฺคลฺลาน สุณาหิ สาธุกํ มนสิ กโรหิ, ภาสิสฺสามีติ. “เอวํ ภนฺเต”ติ โข อายสฺมา มหาโมคฺคลฺลาโน ภควโต ปจฺจสฺโสสิ.\csromanspacer{} ภควา เอตทโวจ\csromandash{}}

\prose{อิธ โมคฺคลฺลาน ภิกฺขุ สพฺพนิมิตฺตานํ อมนสิการา อนิมิตฺตํ เจโตสมาธิํ อุปสมฺปชฺช วิหรติ, ตเมนํ เต เทวา เอวํ ชานนฺติ “อยํ โข อายสฺมา สพฺพนิมิตฺตานํ อมนสิการา อนิมิตฺตํ เจโตสมาธิํ อุปสมฺปชฺช วิหรติ, อปฺเปว นาม อยมายสฺมา อนุโลมิกานิ เสนาสนานิ ปฏิเสวมาโน กลฺยาณมิตฺเต ภชมาโน อินฺทฺริยานิ สมนฺนานยมาโน ยสฺสตฺถาย กุลปุตฺตา สมฺมเทว อคารสฺมา อนคาริยํ ปพฺพชนฺติ, ตทนุตฺตรํ พฺรหฺมจริยปริโยสานํ ทิฏฺเฐว ธมฺเม สยํ อภิญฺญา สจฺฉิกตฺวา อุปสมฺปชฺช วิหเรยฺยา”ติ.\csromanspacer{} เอวํ โข โมคฺคลฺลาน เตสํ เทวานํ ญาณํ โหติ ส-อุปาทิเสเส วา “ส-อุปาทิเสโส”ติ, อนุปาทิเสเส วา “อนุปาทิเสโส”ติ.\csromanspacer{}\csromanspacer{}ตติยํ.}

\csromanmatikaanchor{mtk.252}
\csromantocmarkref{subsection}{4. สีหเสนาปติสุตฺต}{mtk.252}
\bookmark[dest={mtk.252},level=2]{4. สีหเสนาปติสุตฺต}
\csromanheader{2}{4. สีหเสนาปติสุตฺต}
\proseitem{57}{เอวํ เม สุตํ\csromandash{}เอกํ สมยํ ภควา เวสาลิยํ วิหรติ มหาวเน กูฏาคารสาลายํ.\csromanspacer{} อถ โข สีโห เสนาปติ เยน ภควา เตนุปสงฺกมิ, อุปสงฺกมิตฺวา ภควนฺตํ อภิวาเทตฺวา เอกมนฺตํ นิสีทิ, เอกมนฺตํ นิสินฺโน โข สีโห เสนาปติ ภควนฺตํ เอตทโวจ “สกฺกา นุ โข ภนฺเต สนฺทิฏฺฐิกํ ทานผลํ ปญฺญาเปตุนฺ”ติ.}

\prose{เตน หิ สีห ตญฺเญเวตฺถ ปฏิปุจฺฉิสฺสามิ, ยถา เต ขเมยฺย, ตถา นํ พฺยากเรยฺยาสิ.\csromanspacer{} ตํ กิํ มญฺญสิ สีห, อิธ เทฺว ปุริสา เอโก ปุริโส อสฺสทฺโธ มจฺฉรี กทริโย ปริภาสโก, เอโก ปุริโส สทฺโธ}

\vspace{\tipitakaparskip}
\csromancenter{อพฺยากตวคฺค ทานปติ อนุปฺปทานรโต.\csromanspacer{} ตํ กิํ มญฺญสิ สีห, กํ นุ โข\footnote{กิํ นุ โข (ก)} อรหนฺโต ปฐมํ อนุกมฺปนฺตา อนุกมฺเปยฺยุํ, โย วา โส ปุริโส อสฺสทฺโธ มจฺฉรี กทริโย ปริภาสโก, โย วา โส ปุริโส สทฺโธ ทานปติ อนุปฺปทานรโตติ.}
\prose{\csromanfolio{457}โย โส ภนฺเต ปุริโส อสฺสทฺโธ มจฺฉรี กทริโย ปริภาสโก, กินฺตํ\footnote{กินฺติ (ก)} อรหนฺโต ปฐมํ อนุกมฺปนฺตา อนุกมฺปิสฺสนฺติ.\csromanspacer{} โย จ โข โส ภนฺเต ปุริโส สทฺโธ ทานปติ อนุปฺปทานรโต, ตํเยว อรหนฺโต ปฐมํ อนุกมฺปนฺตา อนุกมฺเปยฺยุํ. ตํ กิํ มญฺญสิ สีห, กํ นุ โข อรหนฺโต ปฐมํ อุปสงฺกมนฺตา อุปสงฺกเมยฺยุํ, โย วา โส ปุริโส อสฺสทฺโธ มจฺฉรี กทริโย ปริภาสโก, โย วา โส ปุริโส สทฺโธ ทานปติ อนุปฺปทานรโตติ.\csromanspacer{} โย โส ภนฺเต ปุริโส อสฺสทฺโธ มจฺฉรี กทริโย ปริภาสโก, กินฺตํ อรหนฺโต ปฐมํ อุปสงฺกมนฺตา อุปสงฺกมิสฺสนฺติ.\csromanspacer{} โย จ โข โส ภนฺเต ปุริโส สทฺโธ ทานปติ อนุปฺปทานรโต, ตํเยว อรหนฺโต ปฐมํ อุปสงฺกมนฺตา อุปสงฺกเมยฺยุํ.}

\prose{ตํ กิํ มญฺญสิ สีห, กสฺส นุ โข อรหนฺโต ปฐมํ ปฏิคฺคณฺหนฺตา ปฏิคฺคเณฺหยฺยุํ, โย วา โส ปุริโส อสฺสทฺโธ มจฺฉรี กทริโย ปริภาสโก, โย วา โส ปุริโส สทฺโธ ทานปติ อนุปฺปทานรโตติ.\csromanspacer{} โย โส ภนฺเต ปุริโส อสฺสทฺโธ มจฺฉรี กทริโย ปริภาสโก, กินฺตํ ตสฺส อรหนฺโต ปฐมํ ปฏิคฺคณฺหนฺตา ปฏิคฺคณฺหิสฺสนฺติ.\csromanspacer{} โย จ โข โส ภนฺเต ปุริโส สทฺโธ ทานปติ อนุปฺปทานรโต, ตสฺเสว อรหนฺโต ปฐมํ ปฏิคฺคณฺหนฺตา ปฏิคฺคเณฺหยฺยุํ.}

\prose{ตํ กิํ มญฺญสิ สีห, กสฺส นุ โข อรหนฺโต ปฐมํ ธมฺมํ เทเสนฺตา เทเสยฺยุํ, โย วา โส ปุริโส อสฺสทฺโธ มจฺฉรี กทริโย ปริภาสโก, โย วา โส ปุริโส สทฺโธ ทานปติ อนุปฺปทานรโตติ.\csromanspacer{} โย โส ภนฺเต ปุริโส อสฺสทฺโธ มจฺฉรี กทริโย ปริภาสโก, กินฺตํ ตสฺส อรหนฺโต ปฐมํ ธมฺมํ เทเสนฺตา เทเสสฺสนฺติ.\csromanspacer{} โย จ โข โส ภนฺเต ปุริโส สทฺโธ ทานปติ อนุปฺปทานรโต, ตสฺเสว อรหนฺโต ปฐมํ ธมฺมํ เทเสนฺตา เทเสยฺยุํ.}

\gambhira{สตฺตกนิปาตปาฬิ}
\prose{\csromanfolio{458}ตํ กิํ มญฺญสิ สีห, กสฺส นุ โข กลฺยาโณ กิตฺติสทฺโท อพฺภุคฺคจฺเฉยฺย, โย วา โส ปุริโส อสฺสทฺโธ มจฺฉรี กทริโย ปริภาสโก, โย วา โส ปุริโส สทฺโธ ทานปติ อนุปฺปทานรโตติ.\csromanspacer{} โย โส ภนฺเต ปุริโส อสฺสทฺโธ มจฺฉรี กทริโย ปริภาสโก, กินฺตํ ตสฺส กลฺยาโณ กิตฺติสทฺโท อพฺภุคฺคจฺฉิสฺสติ.\csromanspacer{} โย จ โข โส ภนฺเต ปุริโส สทฺโธ ทานปติ อนุปฺปทานรโต, ตสฺเสว กลฺยาโณ กิตฺติสทฺโท อพฺภุคฺคจฺเฉยฺย.}

\prose{ตํ กิํ มญฺญสิ สีห, โก นุ โข ยํยเทว ปริสํ อุปสงฺกเมยฺย ยทิ ขตฺติยปริสํ ยทิ พฺราหฺมณปริสํ ยทิ คหปติปริสํ ยทิ สมณปริสํ, วิสารโท อุปสงฺกเมยฺย อมงฺกุภูโต, โย วา โส ปุริโส, อสฺสทฺโธ มจฺฉรี กทริโย ปริภาสโก, โย วา โส ปุริโส สทฺโธ ทานปติ อนุปฺปทานรโตติ.\csromanspacer{} โย โส ภนฺเต ปุริโส อสฺสทฺโธ มจฺฉรี กทริโย ปริภาสโก, กิํ โส ยํยเทว ปริสํ อุปสงฺกมิสฺสติ ยทิ ขตฺติยปริสํ ยทิ พฺราหฺมณปริสํ ยทิ คหปติปริสํ ยทิ สมณปริสํ วิสารโท อุปสงฺกมิสฺสติ อมงฺกุภูโต.\csromanspacer{} โย จ โข โส ภนฺเต ปุริโส สทฺโธ ทานปติ อนุปฺปทานรโต, โส ยํยเทว ปริสํ อุปสงฺกเมยฺย ยทิ ขตฺติยปริสํ ยทิ พฺราหฺมณปริสํ ยทิ คหปติปริสํ ยทิ สมณปริสํ, วิสารโท อุปสงฺกเมยฺย อมงฺกุภูโต.}

\prose{ตํ กิํ มญฺญสิ สีห, โก นุ โข กายสฺส เภทา ปรํ มรณา สุคติํ สคฺคํ โลกํ อุปปชฺเชยฺย, โย วา โส ปุริโส อสฺสทฺโธ มจฺฉรี กทริโย ปริภาสโก, โย วา โส ปุริโส สทฺโธ ทานปติ อนุปฺปทานรโตติ.\csromanspacer{} โย โส ภนฺเต ปุริโส อสฺสทฺโธ มจฺฉรี กทริโย ปริภาสโก, กิํ โส กายสฺส เภทา ปรํ มรณา สุคติํ สคฺคํ โลกํ อุปปชฺชิสฺสติ.\csromanspacer{} โย จ โข โส ภนฺเต ปุริโส สทฺโธ ทานปติ อนุปฺปทานรโต, โส กายสฺส เภทา ปรํ มรณา สุคติํ สคฺคํ โลกํ อุปปชฺเชยฺย.}

\prose{ยานิมานิ ภนฺเต ภควตา สนฺทิฏฺฐิกานิ ทานผลานิ อกฺขาตานิ, นาหํ เอตฺถ ภควโต สทฺธาย คจฺฉามิ.\csromanspacer{} อหมฺปิ เอตานิ ชานามิ, อหํ ภนฺเต ทายโก ทานปติ, มํ อรหนฺโต ปฐมํ อนุกมฺปนฺตา อนุกมฺปนฺติ.\csromanspacer{} อหํ ภนฺเต ทายโก ทานปติ, มํ อรหนฺโต ปฐมํ อุปสงฺกมนฺตา อุปสงฺกมนฺติ.\csromanspacer{} อหํ ภนฺเต ทายโก ทานปติ, มยฺหํ อรหนฺโต ปฐมํ ปฏิคฺคณฺหนฺตา}

\vspace{\tipitakaparskip}
\csromancenter{อพฺยากตวคฺค ปฏิคฺคณฺหนฺติ.\csromanspacer{} อหํ ภนฺเต ทายโก ทานปติ, มยฺหํ อรหนฺโต ปฐมํ ธมฺมํ เทเสนฺตา เทเสนฺติ.\csromanspacer{} อหํ ภนฺเต ทายโก ทานปติ, มยฺหํ กลฺยาโณ กิตฺติสทฺโท อพฺภุคฺคโต “สีโห เสนาปติ ทายโก การโก สํฆุปฏฺฐาโก”ติ.\csromanspacer{} อหํ ภนฺเต ทายโก ทานปติ, ยํยเทว ปริสํ อุปสงฺกมามิ ยทิ ขตฺติยปริสํ ฯเปฯ ยทิ สมณปริสํ, วิสารโท อุปสงฺกมามิ อมงฺกุภูโต.\csromanspacer{} ยานิมานิ ภนฺเต ภควตา สนฺทิฏฺฐิกานิ ทานผลานิ อกฺขาตานิ, นาหํ เอตฺถ ภควโต สทฺธาย คจฺฉามิ.\csromanspacer{} อหมฺปิ เอตานิ ชานามิ.\csromanspacer{} ยญฺจ โข มํ ภนฺเต ภควา เอวมาห “ทายโก สีห ทานปติ กายสฺส เภทา ปรํ มรณา สุคติํ สคฺคํ โลกํ อุปปชฺชตี”ติ.\csromanspacer{} เอตาหํ น ชานามิ, เอตฺถ จ ปนาหํ ภควโต สทฺธาย คจฺฉามีติ.\csromanspacer{} เอวเมตํ สีห, เอวเมตํ สีห, ทายโก สีห ทานปติ กายสฺส เภทา ปรํ มรณา สุคติํ สคฺคํ โลกํ อุปปชฺชตีติ.\csromanspacer{}\csromanspacer{}จตุตฺถํ.}
\csromanmatikaanchor{mtk.253}
\csromantocmarkref{subsection}{5. อรกฺเขยฺยสุตฺต}{mtk.253}
\bookmark[dest={mtk.253},level=2]{5. อรกฺเขยฺยสุตฺต}
\csromanheader{2}{5. อรกฺเขยฺยสุตฺต}
\proseitem{58}{\csromanfolio{459}จตฺตาริมานิ ภิกฺขเว ตถาคตสฺส อรกฺเขยฺยานิ, ตีหิ จ อนุปวชฺโช.\csromanspacer{} กตมานิ จตฺตาริ ตถาคตสฺส อรกฺเขยฺยานิ? ปริสุทฺธกายสมาจาโร ภิกฺขเว ตถาคโต, นตฺถิ ตถาคตสฺส กายทุจฺจริตํ, ยํ ตถาคโต รกฺเขยฺย “มา เม อิทํ ปโร อญฺญาสี”ติ.\csromanspacer{} ปริสุทฺธวจีสมาจาโร ภิกฺขเว ตถาคโต, นตฺถิ ตถาคตสฺส วจีทุจฺจริตํ, ยํ ตถาคโต รกฺเขยฺย “มา เม อิทํ ปโร อญฺญาสี”ติ.\csromanspacer{} ปริสุทฺธมโนสมาจาโร ภิกฺขเว ตถาคโต, นตฺถิ ตถาคตสฺส มโนทุจฺจริตํ, ยํ ตถาคโต รกฺเขยฺย “มา เม อิทํ ปโร อญฺญาสี”ติ.\csromanspacer{} ปริสุทฺธาชีโว ภิกฺขเว ตถาคโต, นตฺถิ ตถาคตสฺส มิจฺฉา-อาชีโว, ยํ ตถาคโต รกฺเขยฺย “มา เม อิทํ ปโร อญฺญาสี”ติ.\csromanspacer{} อิมานิ จตฺตาริ ตถาคตสฺส อรกฺเขยฺยานิ.}

\prose{กตเมหิ ตีหิ อนุปวชฺโช? สฺวากฺขาตธมฺโม ภิกฺขเว ตถาคโต.\csromanspacer{} ตตฺร วต มํ สมโณ วา พฺราหฺมโณ วา เทโว วา มาโร วา พฺรหฺมา วา โกจิ วา โลกสฺมิํ สหธมฺเมน ปฏิโจเทสฺสติ “อิติปิ ตฺวํ น สฺวากฺขาตธมฺโม”ติ, นิมิตฺตเมตํ ภิกฺขเว น สมนุปสฺสามิ, เอตมหํ\footnote{เอตํปหํ (สี, สฺยา)}}

\vspace{\tipitakaparskip}
\csromancenter{สตฺตกนิปาตปาฬิ ภิกฺขเว นิมิตฺตํ อสมนุปสฺสนฺโต เขมปฺปตฺโต อภยปฺปตฺโต เวสารชฺชปฺปตฺโต วิหรามิ.}
\prose{\csromanfolio{460}สุปญฺญตฺตา โข ปน เม ภิกฺขเว สาวกานํ นิพฺพานคามินี ปฏิปทา, ยถาปฏิปนฺนา มม สาวกา อาสวานํ ขยา อนาสวํ เจโตวิมุตฺติํ ปญฺญาวิมุตฺติํ ทิฏฺเฐว ธมฺเม สยํ อภิญฺญา สจฺฉิกตฺวา อุปสมฺปชฺช วิหรนฺติ.\csromanspacer{} ตตฺร วต มํ สมโณ วา พฺราหฺมโณ วา เทโว วา มาโร วา พฺรหฺมา วา โกจิ วา โลกสฺมิํ สหธมฺเมน ปฏิโจเทสฺสติ “อิติปิ เต น สุปญฺญตฺตา สาวกานํ นิพฺพานคามินี ปฏิปทา, ยถาปฏิปนฺนา ตว สาวกา อาสวานํ ขยา ฯเปฯ สจฺฉิกตฺวา อุปสมฺปชฺช วิหรนฺตี”ติ, นิมิตฺตเมตํ ภิกฺขเว น สมนุปสฺสามิ, เอตมหํ ภิกฺขเว นิมิตฺตํ อสมนุปสฺสนฺโต เขมปฺปตฺโต อภยปฺปตฺโต เวสารชฺชปฺปโต วิหรามิ.}

\prose{อเนกสตา โข ปน เม ภิกฺขเว สาวกปริสา อาสวานํ ขยา ฯเปฯ สจฺฉิกตฺวา อุปสมฺปชฺช วิหรนฺติ.\csromanspacer{} ตตฺร วต มํ สมโณ วา พฺราหฺมโณ วา เทโว วา มาโร วา พฺรหฺมา วา โกจิ วา โลกสฺมิํ สหธมฺเมน ปฏิโจเทสฺสติ “อิติปิ เต น อเนกสตา สาวกปริสา อาสวานํ ขยา อนาสวํ เจโตวิมุตฺติํ ปญฺญาวิมุตฺติํ ทิฏฺเฐว ธมฺเม สยํ อภิญฺญา สจฺฉิกตฺวา อุปสมฺปชฺช วิหรนฺตี”ติ, นิมิตฺตเมตํ ภิกฺขเว น สมนุปสฺสามิ.\csromanspacer{} เอตมหํ ภิกฺขเว นิมิตฺตํ อสมนุปสฺสนฺโต เขมปฺปตฺโต อภยปฺปตฺโต เวสารชฺชปฺปตฺโต วิหรามิ.\csromanspacer{} อิเมหิ ตีหิ อนุปวชฺโช.}

\prose{อิมานิ โข ภิกฺขเว.\csromanspacer{} จตฺตาริ ตถาคตสฺส อรกฺเขยฺยานิ, อิเมหิ จ ตีหิ อนุปวชฺโชติ.\csromanspacer{}\csromanspacer{}ปญฺจมํ}

\csromanmatikaanchor{mtk.254}
\csromantocmarkref{subsection}{6. กิมิลสุตฺต}{mtk.254}
\bookmark[dest={mtk.254},level=2]{6. กิมิลสุตฺต}
\csromanheader{2}{6. กิมิลสุตฺต}
\proseitem{59}{เอวํ เม สุตํ\csromandash{}เอกํ สมยํ ภควา กิมิลายํ วิหรติ นิจุลวเน\footnote{เวฬุวเน (สี, สฺยา, กํ, อิ) เหฏฺฐา 216, 298 ปิฏฺเฐสุปิ.}.\csromanspacer{} อถ โข อายสฺมา กิมิโล เยน ภควา เตนุปสงฺกมิ, อุปสงฺกมิตฺวา ภควนฺตํ อภิวาเทตฺวา เอกมนฺตํ นิสีทิ, เอกมนฺตํ นิสินฺโน โข อายสฺมา กิมิโล ภควนฺตํ เอตทโวจ “โก นุ โข ภนฺเต เหตุ โก ปจฺจโย, เยน ตถาคเต ปรินิพฺพุเต สทฺธมฺโม น จิรฏฺฐิติโก โหตี”ติ.}

\chapterheadpageread{6. อพฺยากตวคฺค}
\prose{\csromanfolio{461}อิธ กิมิล ตถาคเต ปรินิพฺพุเต ภิกฺขู ภิกฺขุนิโย อุปาสกา อุปาสิกาโย สตฺถริ อคารวา วิหรนฺติ อปฺปติสฺสา, ธมฺเม อคารวา วิหรนฺติ อปฺปติสฺสา, สํเฆ อคารวา วิหรนฺติ อปฺปติสฺสา, สิกฺขาย อคารวา วิหรนฺติ อปฺปติสฺสา, สมาธิสฺมิํ อคารวา วิหรนฺติ อปฺปติสฺสา, อปฺปมาเท อคารวา วิหรนฺติ อปฺปติสฺสา, ปฏิสนฺถาเร อคารวา วิหรนฺติ อปฺปติสฺสา, อยํ โข กิมิล เหตุ อยํ ปจฺจโย, เยน ตถาคเต ปรินิพฺพุเต สทฺธมฺโม น จิรฏฺฐิติโก โหตีติ.}

\prose{โก ปน ภนฺเต เหตุ โก ปจฺจโย, เยน ตถาคเต ปรินิพฺพุเต สทฺธมฺโม จิรฏฺฐิติโก โหตีติ.\csromanspacer{} อิธ กิมิล ตถาคเต ปรินิพฺพุเต ภิกฺขู ภิกฺขุนิโย อุปาสกา อุปาสิกาโย สตฺถริ สคารวา วิหรนฺติ สปฺปติสฺสา, ธมฺเม สคารวา วิหรนฺติ สปฺปติสฺสา, สํเฆ สคารวา วิหรนฺติ สปฺปติสฺสา, สิกฺขาย สคารวา วิหรนฺติ สปฺปติสฺสา, สมาธิสฺมิํ สคารวา วิหรนฺติ สปฺปติสฺสา, อปฺปมาเท สคารวา วิหรนฺติ สปฺปติสฺสา, ปฏิสนฺถาเร สคารวา วิหรนฺติ สปฺปติสฺสา.\csromanspacer{} อยํ โข กิมิล เหตุ อยํ ปจฺจโย, เยน ตถาคเต ปรินิพฺพุเต สทฺธมฺโม จิรฏฺฐิติโก โหตีติ.\csromanspacer{}\csromanspacer{}ฉฏฺฐํ.}

\csromanmatikaanchor{mtk.255}
\csromantocmarkref{subsection}{7. สตฺตธมฺมสุตฺต}{mtk.255}
\bookmark[dest={mtk.255},level=2]{7. สตฺตธมฺมสุตฺต}
\csromanheader{2}{7. สตฺตธมฺมสุตฺต}
\proseitem{60}{สตฺตหิ ภิกฺขเว ธมฺเมหิ สมนฺนาคโต ภิกฺขุ นจิรสฺเสว อาสวานํ ขยา ฯเปฯ สจฺฉิกตฺวา อุปสมฺปชฺช วิหเรยฺย.\csromanspacer{} กตเมหิ สตฺตหิ? อิธ ภิกฺขเว ภิกฺขุ สทฺโธ โหติ, สีลวา โหติ, พหุสฺสุโต โหติ, ปฏิสลฺลีโน โหติ, อารทฺธวีริโย โหติ, สติมา โหติ, ปญฺญวา โหติ.\csromanspacer{} อิเมหิ โข ภิกฺขเว สตฺตหิ ธมฺเมหิ สมนฺนาคโต ภิกฺขุ นจิรสฺเสว อาสวานํ ขยา ฯเปฯ สจฺฉิกตฺวา อุปสมฺปชฺช วิหเรยฺยาติ.\csromanspacer{}\csromanspacer{}สตฺตมํ.}

\csromanmatikaanchor{mtk.256}
\csromantocmarkref{subsection}{8. ปจลายมานสุตฺต}{mtk.256}
\bookmark[dest={mtk.256},level=2]{8. ปจลายมานสุตฺต}
\csromanheader{2}{8. ปจลายมานสุตฺต}
\proseitem{61}{เอวํ เม สุตํ\csromandash{}เอกํ สมยํ ภควา ภคฺเคสุ วิหรติ สุสุมารคิเร เภสกฬาวเน มิคทาเย.\csromanspacer{} เตน โข ปน สมเยน อายสฺมา มหาโมคฺคลฺลาโน มคเธสุ กลฺลวาฬปุตฺตคาเม\footnote{กลฺลวาลมุตฺตคาเม (สฺยา)} ปจลายมาโน}

\vspace{\tipitakaparskip}
\csromancenter{สตฺตกนิปาตปาฬิ นิสินฺโน โหติ.\csromanspacer{} อทฺทสา โข ภควา ทิพฺเพน จกฺขุนา วิสุทฺเธน อติกฺกนฺตมานุสเกน อายสฺมนฺตํ มหาโมคฺคลฺลานํ มคเธสุ กลฺลวาฬปุตฺตคาเม ปจลายมานํ นิสินฺนํ, ทิสฺวา เสยฺยถาปิ นาม พลวา ปุริโส สมิญฺชิตํ วา พาหํ ปสาเรยฺย, ปสาริตํ วา พาหํ สมิญฺเชยฺย.\csromanspacer{} เอวเมวํ ภคฺเคสุ สุสุมารคิเร เภสกฬาวเน มิคทาเย อนฺตรหิโต มคเธสุ กลฺลวาฬปุตฺตคาเม อายสฺมโต มหาโมคฺคลฺลานสฺส สมฺมุเข ปาตุรโหสิ.\csromanspacer{} นิสีทิ ภควา ปญฺญตฺเต อาสเน, นิสชฺช โข ภควา อายสฺมนฺตํ มหาโมคฺคลฺลานํ เอตทโวจ\csromandash{}}
\prose{\csromanfolio{462}ปจลายสิ โน ตฺวํ โมคฺคลฺลาน, ปจลายสิ โน ตฺวํ โมคฺคลฺลานาติ.\csromanspacer{} เอวํ ภนฺเต.\csromanspacer{} ตสฺมาติห โมคฺคลฺลาน ยถาสญฺญิสฺส เต วิหรโต ตํ มิทฺธํ โอกฺกมติ, ตํ สญฺญํ มา มนสากาสิ\footnote{มา มนสิกาสิ (สี), มนสิ กเรยฺยาสิ (สฺยา), มนสากาสิ (ก)}, ตํ สญฺญํ มา พหุลมกาสิ\footnote{ตํ สญฺญํ พหุลํ กเรยฺยาสิ (สฺยา), ตํ สญฺญํ พหุลมกาสิ (ก)}, ฐานํ โข ปเนตํ โมคฺคลฺลาน วิชฺชติ, ยํ เต เอวํ วิหรโต ตํ มิทฺธํ ปหีเยถ. (1)}

\prose{โน เจ เต เอวํ วิหรโต ตํ มิทฺธํ ปหีเยถ, ตโต ตฺวํ โมคฺคลฺลาน ยถาสุตํ ยถาปริยตฺตํ ธมฺมํ เจตสา อนุวิตกฺเกยฺยาสิ อนุวิจาเรยฺยาสิ มนสา อนุเปกฺเขยฺยาสิ.\csromanspacer{} ฐานํ โข ปเนตํ วิชฺชติ, ยํ เต เอวํ วิหรโต ตํ มิทฺธํ ปหีเยถ. (2)}

\prose{โน เจ เต เอวํ วิหรโต ตํ มิทฺธํ ปหีเยถ, ตโต ตฺวํ โมคฺคลฺลาน ยถาสุตํ ยถาปริยตฺตํ ธมฺมํ วิตฺถาเรน สชฺฌายํ กเรยฺยาสิ.\csromanspacer{} ฐานํ โข ปเนตํ วิชฺชติ, ยํ เต เอวํ วิหรโต ตํ มิทฺธํ ปหีเยถ. (3)}

\prose{โน เจ เต เอวํ วิหรโต ตํ มิทฺธํ ปหีเยถ, ตโต ตฺวํ โมคฺคลฺลาน อุโภ กณฺณโสตานิ อาวิญฺเฉยฺยาสิ\footnote{อาวิญฺเชยฺยาสิ (สี, สฺยา)}, ปาณินา คตฺตานิ อนุมชฺเชยฺยาสิ.\csromanspacer{} ฐานํ โข ปเนตํ วิชฺชติ, ยํ เต เอวํ วิหรโต ตํ มิทฺธํ ปหีเยถ. (4)}

\prose{โน เจ เต เอวํ วิหรโต ตํ มิทฺธํ ปหีเยถ, ตโต ตฺวํ โมคฺคลฺลาน อุฏฺฐายาสนา อุทเกน อกฺขีนิ อนุมชฺชิตฺวา\footnote{ปนิญฺชิตฺวา (ก)} ทิสา}

\vspace{\tipitakaparskip}
\csromancenter{อพฺยากตวคฺค อนุวิโลเกยฺยาสิ, นกฺขตฺตานิ ตารกรูปานิ อุลฺโลเกยฺยาสิ.\csromanspacer{} ฐานํ โข ปเนตํ วิชฺชติ, ยํ เต เอวํ วิหรโต ตํ มิทฺธํ ปหีเยถ. (5)}
\prose{\csromanfolio{463}โน เจ เต เอวํ วิหรโต ตํ มิทฺธํ ปหีเยถ, ตโต ตฺวํ โมคฺคลฺลาน อาโลกสญฺญํ มนสิ กเรยฺยาสิ, ทิวาสญฺญํ อธิฏฺฐเหยฺยาสิ, ยถา ทิวา ตถา รตฺติํ ยถา รตฺติํ ตถา ทิวา, อิติ วิวเฏน\footnote{วิวฏฺเฏน (สฺยา), มิทฺธวิคเตน (ก)} เจตสา อปริโยนทฺเธน สปฺปภาสํ จิตฺตํ ภาเวยฺยาสิ.\csromanspacer{} ฐานํ โข ปเนตํ วิชฺชติ, ยํ เต เอวํ วิหรโต ตํ มิทฺธํ ปหีเยถ. (6)}

\prose{โน เจ เต เอวํ วิหรโต ตํ มิทฺธํ ปหีเยถ, ตโต ตฺวํ โมคฺคลฺลาน ปจฺฉาปุเรสญฺญี\footnote{ปจฺฉาปุเร ตถาสญฺญี (กตฺถจิ)} จงฺกมํ อธิฏฺฐเหยฺยาสิ อนฺโตคเตหิ อินฺทฺริเยหิ อพหิคเตน มานเสน.\csromanspacer{} ฐานํ โข ปเนตํ วิชฺชติ, ยํ เต เอวํ วิหรโต ตํ มิทฺธํ ปหีเยถ. (7)}

\prose{โน เจ เต เอวํ วิหรโต ตํ มิทฺธํ ปหีเยถ, ตโต ตฺวํ โมคฺคลฺลาน ทกฺขิเณน ปสฺเสน สีหเสยฺยํ กปฺเปยฺยาสิ ปาเท ปาทํ อจฺจาธาย สโต สมฺปชาโน อุฏฺฐานสญฺญํ มนสิ กริตฺวา.\csromanspacer{} ปฏิพุทฺเธน จ\footnote{ปฏิพุทฺเธเนว (สฺยา)} เต โมคฺคลฺลาน ขิปฺปญฺเญว ปจฺจุฏฺฐาตพฺพํ น เสยฺยสุขํ น ปสฺสสุขํ น มิทฺธสุขํ อนุยุตฺโต วิหริสฺสามีติ.\csromanspacer{} เอวํ หิ เต โมคฺคลฺลาน สิกฺขิตพฺพํ.}

\prose{ตสฺมาติห โมคฺคลฺลาน เอวํ สิกฺขิตพฺพํ “น อุจฺจาโสณฺฑํ ปคฺคเหตฺวา กุลานิ อุปสงฺกมิสฺสามี”ติ.\csromanspacer{} เอวํ หิ เต โมคฺคลฺลาน สิกฺขิตพฺพํ, สเจ โมคฺคลฺลาน ภิกฺขุ อุจฺจาโสณฺฑํ ปคฺคเหตฺวา กุลานิ อุปสงฺกมติ.\csromanspacer{} สนฺติ หิ โมคฺคลฺลาน กุเลสุ กิจฺจกรณียานิ, เยหิ มนุสฺสา อาคตํ ภิกฺขุํ น มนสิ กโรนฺติ, ตตฺร ภิกฺขุสฺส เอวํ โหติ “โกสุ นาม อิทานิ มํ อิมสฺมิํ กุเล ปริภินฺทิ, วิรตฺตรูปา ทานิเม มยิ มนุสฺสา”ติ, อิติสฺส อลาเภน มงฺกุภาโว, มงฺกุภูตสฺส อุทฺธจฺจํ, อุทฺธตสฺส อสํวโร, อสํวุตสฺส อารา จิตฺตํ สมาธิมฺหา.}

\prose{ตสฺมาติห โมคฺคลฺลาน เอวํ สิกฺขิตพฺพํ “น วิคฺคาหิกกถํ กเถสฺสามี”ติ.\csromanspacer{} เอวํ หิ เต โมคฺคลฺลาน สิกฺขิตพฺพํ, วิคฺคาหิกาย โมคฺคลฺลาน กถาย}

\vspace{\tipitakaparskip}
\csromancenter{สตฺตกนิปาตปาฬิ สติ กถาพาหุลฺลํ ปาฏิกงฺขํ, กถาพาหุลฺเล สติ อุทฺธจฺจํ, อุทฺธตสฺส อสํวโร, อสํวุตสฺส อารา จิตฺตํ สมาธิมฺหา.1 นาหํ โมคฺคลฺลาน สพฺเพเหว สํสคฺคํ วณฺณยามิ, น ปนาหํ โมคฺคลฺลาน สพฺเพเหว สํสคฺคํ น วณฺณยามิ, สคหฏฺฐปพฺพชิเตหิ โข อหํ โมคฺคลฺลาน สํสคฺคํ น วณฺณยามิ1, ยานิ จ โข ตานิ เสนาสนานิ อปฺปสทฺทานิ อปฺปนิคฺโฆสานิ วิชนวาตานิ มนุสฺสราหสฺเสยฺยกานิ\footnote{มนุสฺสราหเสยฺยกานิ (สี, สฺยา)} ปฏิสลฺลานสารุปฺปานิ, ตถารูเปหิ เสนาสเนหิ สํสคฺคํ\footnote{สมคฺคํ (ก)} วณฺณยามีติ.}
\prose{\csromanfolio{464}เอวํ วุตฺเต อายสฺมา มหาโมคฺคลฺลาโน ภควนฺตํ เอตทโวจ “กิตฺตาวตา นุ โข ภนฺเต ภิกฺขุ สํขิตฺเตน ตณฺหาสงฺขยวิมุตฺโต โหติ, อจฺจนฺตนิฏฺโฐ อจฺจนฺตโยคกฺเขมี อจฺจนฺตพฺรหฺมจารี อจฺจนฺตปริโยสาโน เสฏฺโฐ เทวมนุสฺสานนฺ”ติ.}

\prose{อิธ โมคฺคลฺลาน ภิกฺขุโน สุตํ โหติ “สพฺเพ ธมฺมา นาลํ อภินิเวสายา”ติ, เอวญฺเจตํ โมคฺคลฺลาน ภิกฺขุโน สุตํ โหติ “สพฺเพ ธมฺมา นาลํ อภินิเวสายา”ติ.\csromanspacer{} โส สพฺพํ ธมฺมํ อภิชานาติ, สพฺพํ ธมฺมํ อภิญฺญาย สพฺพํ ธมฺมํ ปริชานาติ.\csromanspacer{} สพฺพํ ธมฺมํ ปริญฺญาย ยํกิญฺจิ เวทนํ เวทิยติ สุขํ วา ทุกฺขํ วา อทุกฺขมสุขํ วา.\csromanspacer{} โส ตาสุ เวทนาสุ อนิจฺจานุปสฺสี วิหรติ, วิราคานุปสฺสี วิหรติ, นิโรธานุปสฺสี วิหรติ, ปฏินิสฺสคฺคานุปสฺสี วิหรติ.\csromanspacer{} โส ตาสุ เวทนาสุ อนิจฺจานุปสฺสี วิหรนฺโต วิราคานุปสฺสี วิหรนฺโต นิโรธานุปสฺสี วิหรนฺโต ปฏินิสฺสคฺคานุปสฺสี วิหรนฺโต น กิญฺจิ\footnote{น จ กิญฺจิ (สี, สฺยา, ก) ม 1. 318 ปิฏฺเฐ ปสฺสิตพฺพํ.} โลเก อุปาทิยติ, อนุปาทิยํ น ปริตสฺสติ, อปริตสฺสํ ปจฺจตฺตํเยว ปรินิพฺพายติ, “ขีณา ชาติ, วุสิตํ พฺรหฺมจริยํ, กตํ กรณียํ, นาปรํ อิตฺถตฺตายา”ติ ปชานาติ, เอตฺตาวตา โข โมคฺคลฺลาน ภิกฺขุ สํขิตฺเตน ตณฺหาสงฺขยวิมุตฺโต โหติ, อจฺจนฺตนิฏฺโฐ อจฺจนฺตโยคกฺเขมี อจฺจนฺตพฺรหฺมจารี อจฺจนฺตปริโยสาโน เสฏฺโฐ เทวมนุสฺสานนฺติ.\csromanspacer{}\csromanspacer{}อฏฺฐมํ.}

\csromanmatikaanchor{mtk.257}
\csromantocmarkref{subsection}{9. เมตฺตสุตฺต}{mtk.257}
\bookmark[dest={mtk.257},level=2]{9. เมตฺตสุตฺต}
\csromanheader{2}{6. อพฺยากตวคฺค \textbf{9. เมตฺตสุตฺต}}
\proseitem{62}{\csromanfolio{465}\csromansymbolfootnote{*}{ขุ 1. 205 ปิฏฺเฐ อิติวุตฺตเกปิ.}มา ภิกฺขเว ปุญฺญานํ ภายิตฺถ, สุขสฺเสตํ ภิกฺขเว อธิวจนํ ยทิทํ ปุญฺญานิ\footnote{ยทิทํ ปุญฺญนฺติ (สี), ยทิทํ ปุญฺญานิ (ก) 2-2. ภิกฺขเว ทีฆรตฺตํ อิฏฺฐํ (สฺยา), ภิกฺขเว ทีฆรตฺตํ ปุญฺญานํ อิฏฺฐํ (?)}.\csromanspacer{} อภิชานามิ โข ปนาหํ2 ภิกฺขเว ทีฆรตฺตํ กตานํ ปุญฺญานํ ทีฆรตฺตํ อิฏฺฐํ2 กนฺตํ มนาปํ วิปากํ ปจฺจนุภูตํ.\csromanspacer{} สตฺต วสฺสานิ เมตฺตจิตฺตํ ภาเวสิํ, สตฺต วสฺสานิ เมตฺตจิตฺตํ ภาเวตฺวา สตฺต สํวฏฺฏ วิวฏฺฏกปฺเป นยิมํ โลกํ ปุนาคมาสิํ, สํวฏฺฏมาเน สุทาหํ\footnote{สํวฏฺฏมานสฺสุทาหํ (ก)} ภิกฺขเว โลเก อาภสฺสรูปโค โหมิ, วิวฏฺฏมาเน โลเก สุญฺญํ พฺรหฺมวิมานํ อุปปชฺชามิ.}

\prose{ตตฺร สุทํ ภิกฺขเว พฺรหฺมา โหมิ มหาพฺรหฺมา อภิภู อนภิภูโต อญฺญทตฺถุทโส วสวตฺตี.\csromanspacer{} ฉตฺติํสกฺขตฺตุํ โข ปนาหํ ภิกฺขเว สกฺโก อโหสิํ เทวานมินฺโท, อเนกสตกฺขตฺตุํ ราชา อโหสิํ จกฺกวตฺตี ธมฺมิโก ธมฺมราชา จาตุรนฺโต วิชิตาวี ชนปทตฺถาวริยปฺปตฺโต สตฺตรตนสมนฺนาคโต.\csromanspacer{} ตสฺส มยฺหํ ภิกฺขเว อิมานิ สตฺต รตนานิ อเหสุํ.\csromanspacer{} เสยฺยถิทํ, จกฺกรตนํ หตฺถิรตนํ อสฺสรตนํ มณิรตนํ อิตฺถิรตนํ คหปติรตนํ ปริณายกรตนเมว สตฺตมํ.\csromanspacer{} ปโรสหสฺสํ โข ปน เม ภิกฺขเว ปุตฺตา อเหสุํ สูรา วีรงฺครูปา ปรเสนปฺปมทฺทนา.\csromanspacer{} โส อิมํ ปถวิํ สาครปริยนฺตํ อทณฺเฑน อสตฺเถน ธมฺเมน อภิวิชิย อชฺฌาวสินฺ’ติ\footnote{อชฺฌาวสนฺติ (สฺยา), อชฺฌาวสติ (สี, ก)}.}

\csromangathameasure{ปสฺส ปุญฺญานํ วิปากํ, กุสลานํ สุเขสิโน. \\ เมตฺตํ จิตฺตํ วิภาเวตฺวา, สตฺต วสฺสานิ ภิกฺขโว. \\ สตฺตสํวฏฺฏวิวฏฺฏกปฺเป, นยิมํ โลกํ ปุนาคมิํ. \\ สํวฏฺฏมาเน โลกมฺหิ, โหมิ อาภสฺสรูปโค. \\ วิวฏฺฏมาเน โลกสฺมิํ, สุญฺญพฺรหฺมูปโค อหุํ. \\ สตฺตกฺขตฺตุํ มหาพฺรหฺมา, วสวตฺตี ตทา อหุํ. \\ ฉตฺติํสกฺขตฺตุํ เทวินฺโท, เทวรชฺชมการยิํ.}
\csromangathagroup{\csromangathastanza{ปสฺส ปุญฺญานํ วิปากํ, กุสลานํ สุเขสิโน\footnote{สุเขสินํ (สี)}. \\ เมตฺตํ จิตฺตํ วิภาเวตฺวา, สตฺต วสฺสานิ ภิกฺขโว\footnote{สุเขสินํ (สี)}.}\csromangathastanza{สตฺตสํวฏฺฏวิวฏฺฏกปฺเป, นยิมํ โลกํ ปุนาคมิํ\footnote{ปุนาคมํ (สฺยา)}. \\ สํวฏฺฏมาเน โลกมฺหิ, โหมิ อาภสฺสรูปโค.}\csromangathastanza{วิวฏฺฏมาเน โลกสฺมิํ, สุญฺญพฺรหฺมูปโค อหุํ. \\ สตฺตกฺขตฺตุํ มหาพฺรหฺมา, วสวตฺตี ตทา อหุํ. \\ ฉตฺติํสกฺขตฺตุํ เทวินฺโท, เทวรชฺชมการยิํ.}}

\gambhira{สตฺตกนิปาตปาฬิ}
\csromangathameasure{จกฺกวตฺตี อหุํ ราชา, ชมฺพุมณฺฑสฺส อิสฺสโร. \\ มุทฺธาวสิตฺโต ขตฺติโย, มนุสฺสาธิปตี อหุํ. \\ อทณฺเฑน อสตฺเถน, วิเชยฺย ปถวิํ อิมํ. \\ อสาหเสน กมฺเมน, สเมน มนุสาสิ ตํ. \\ ธมฺเมน รชฺชํ กาเรตฺวา, อสฺมิํ ปถวิมณฺฑเล. \\ มหทฺธเน มหาโภเค, อฑฺเฒ อชายิหํ กุเล. \\ สพฺพกาเมหิ สมฺปนฺเน, รตเนหิ จ สตฺตหิ. \\ พุทฺธา สงฺคาหกา โลเก, เตหิ เอตํ สุเทสิตํ. \\ เอโส เหตุ มหนฺตสฺส, ปถโพฺย เม น วิปชฺชติ. \\ ปหูตวิตฺตูปกรโณ, ราชา โหติ6 ปตาปวา. \\ อิทฺธิมา ยสวา โหติ6, ชมฺพุมณฺฑสฺส อิสฺสโร. \\ โก สุตฺวา นปฺปสีเทยฺย, อปิ กณฺหาภิชาติโย. \\ ตสฺมา หิ อตฺตกาเมน, มหตฺตมภิกงฺขตา. \\ สทฺธมฺโม ครุกาตพฺโพ, สรํ พุทฺธานสาสนนฺติ.นวมํ.}
\csromangathagroup{\csromangathastanza{\csromanfolio{466}จกฺกวตฺตี อหุํ ราชา, ชมฺพุมณฺฑสฺส\footnote{ชมฺพุทีปสฺส (สี), ชมฺพุสณฺฑสฺส (สฺยา)} อิสฺสโร. \\ มุทฺธาวสิตฺโต\footnote{ชมฺพุทีปสฺส (สี), ชมฺพุสณฺฑสฺส (สฺยา)} ขตฺติโย, มนุสฺสาธิปตี อหุํ.}\csromangathastanza{อทณฺเฑน อสตฺเถน, วิเชยฺย ปถวิํ อิมํ. \\ อสาหเสน กมฺเมน\footnote{ธมฺเมน (สี, สฺยา)}, สเมน มนุสาสิ ตํ.}\csromangathastanza{ธมฺเมน รชฺชํ กาเรตฺวา, อสฺมิํ ปถวิมณฺฑเล. \\ มหทฺธเน มหาโภเค, อฑฺเฒ อชายิหํ กุเล.}\csromangathastanza{สพฺพกาเมหิ สมฺปนฺเน\footnote{สมฺปุณฺเณ (ก)}, รตเนหิ จ สตฺตหิ. \\ พุทฺธา สงฺคาหกา โลเก, เตหิ เอตํ สุเทสิตํ.}\csromangathastanza{เอโส เหตุ มหนฺตสฺส, ปถโพฺย เม น วิปชฺชติ\footnote{เอส เหตุ มหนฺตสฺส, ปุถโพฺย เยน วุจฺจติ (สี, สฺยา) 6. โหมิ (สี, สฺยา)}. \\ ปหูตวิตฺตูปกรโณ, ราชา โหติ6 ปตาปวา.}\csromangathastanza{อิทฺธิมา ยสวา โหติ6, ชมฺพุมณฺฑสฺส\footnote{ชมฺพุสณฺฑสฺส (สี, สฺยา)} อิสฺสโร. \\ โก สุตฺวา นปฺปสีเทยฺย, อปิ กณฺหาภิชาติโย.}\csromangathastanza{ตสฺมา หิ อตฺตกาเมน\footnote{อตฺถกาเมน (สฺยา, ก)}, มหตฺตมภิกงฺขตา. \\ สทฺธมฺโม ครุกาตพฺโพ, สรํ พุทฺธานสาสนนฺติ.\csromanspacer{}\csromanspacer{}นวมํ.}}

\csromanmatikaanchor{mtk.258}
\csromantocmarkref{subsection}{10. ภริยาสุตฺต}{mtk.258}
\bookmark[dest={mtk.258},level=2]{10. ภริยาสุตฺต}
\csromanheader{2}{10. ภริยาสุตฺต}
\proseitem{63}{อถ โข ภควา ปุพฺพณฺหสมยํ นิวาเสตฺวา ปตฺตจีวรมาทาย เยน อนาถปิณฺฑิกสฺส คหปติสฺส นิเวสนํ เตนุปสงฺกมิ, อุปสงฺกมิตฺวา ปญฺญตฺเต อาสเน นิสีทิ.\csromanspacer{} เตน โข ปน สมเยน อนาถปิณฺฑิกสฺส คหปติสฺส นิเวสเน มนุสฺสา อุจฺจาสทฺทา มหาสทฺทา โหนฺติ.\csromanspacer{} อถ โข อนาถปิณฺฑิโก คหปติ เยน ภควา เตนุปสงฺกมิ, อุปสงฺกมิตฺวา ภควนฺตํ อภิวาเทตฺวา เอกมนฺตํ นิสีทิ, เอกมนฺตํ นิสินฺนํ โข อนาถปิณฺฑิกํ คหปติํ ภควา เอตทโวจ\csromandash{}}

\chapterheadpageread{6. อพฺยากตวคฺค}
\prose{\csromanfolio{467}กิํ นุ เต คหปติ นิเวสเน มนุสฺสา อุจฺจาสทฺทา มหาสทฺทา เกวฏฺฏา มญฺเญ มจฺฉวิโลเปติ.\csromanspacer{} อยํ ภนฺเต สุชาตา ฆรสุณฺหา อฑฺฒกุลา อานีตา, สา เนว สสฺสุํ อาทิยติ.\csromanspacer{} น สสุรํ อาทิยติ, น สามิกํ อาทิยติ, ภควนฺตมฺปิ น สกฺกโรติ น ครุํ กโรติ น มาเนติ น ปูเชตีติ.}

\prose{อถ โข ภควา สุชาตํ ฆรสุณฺหํ อามนฺเตสิ “เอหิ สุชาเต”ติ. “เอวํ ภนฺเต”ติ โข สุชาตา ฆรสุณฺหา ภควโต ปฏิสฺสุตฺวา เยน ภควา เตนุปสงฺกมิ, อุปสงฺกมิตฺวา ภควนฺตํ อภิวาเทตฺวา เอกมนฺตํ นิสีทิ, เอกมนฺตํ นิสินฺนํ โข สุชาตํ ฆรสุณฺหํ ภควา เอตทโวจ\csromandash{}}

\prose{สตฺต โข อิมา สุชาเต ปุริสสฺส ภริยาโย.\csromanspacer{} กตมา สตฺต? วธกสมา โจรีสมา อยฺยสมา มาตาสมา ภคินีสมา สขีสมา ทาสีสมา.\csromanspacer{} อิมา โข สุชาเต สตฺต ปุริสสฺส ภริยาโย, ตาสํ ตฺวํ กตมาติ.\csromanspacer{} น โข อหํ\footnote{นาหํ (สฺยา)} ภนฺเต อิมสฺส ภควตา สํขิตฺเตน ภาสิตสฺส วิตฺถาเรน อตฺถํ อาชานามิ, สาธุ เม ภนฺเต ภควา ตถา ธมฺมํ เทเสตุ, ยถาหํ อิมสฺส ภควตา สํขิตฺเตน ภาสิตสฺส วิตฺถาเรน อตฺถํ ชาเนยฺยนฺติ.\csromanspacer{} เตน หิ สุชาเต สุณาหิ สาธุกํ มนสิ กโรหิ, ภาสิสฺสามีติ. “เอวํ ภนฺเต”ติ โข สุชาตา ฆรสุณฺหา ภควโต ปจฺจสฺโสสิ.\csromanspacer{} ภควา เอตทโวจ\csromandash{}}

\csromangathameasure{“ปทุฏฺฐจิตฺตา อหิตานุกมฺปินี, \\ อญฺเญสุ รตฺตา อติมญฺญเต ปติํ. \\ ธเนน กีตสฺส วธาย อุสฺสุกา, \\ ยา เอวรูปา ปุริสสฺส ภริยา.}
\csromangathagroup{\csromangathastanza{“ปทุฏฺฐจิตฺตา อหิตานุกมฺปินี, \\ อญฺเญสุ รตฺตา อติมญฺญเต ปติํ. \\ ธเนน กีตสฺส วธาย อุสฺสุกา, \\ ยา เอวรูปา ปุริสสฺส ภริยา.}}

\csromanheader{2}{วธา จ ภริยาติ จ สา ปวุจฺจติ. (1)}
\csromangathameasure{ยํ อิตฺถิยา วินฺทติ สามิโก ธนํ, \\ สิปฺปํ วณิชฺชญฺจ กสิํ อธิฏฺฐหํ. \\ อปฺปมฺปิ ตสฺส อปหาตุมิจฺฉติ, \\ ยา เอวรูปา ปุริสสฺส ภริยา.}
\csromangathagroup{\csromangathastanza{ยํ อิตฺถิยา วินฺทติ สามิโก ธนํ, \\ สิปฺปํ วณิชฺชญฺจ กสิํ อธิฏฺฐหํ. \\ อปฺปมฺปิ ตสฺส อปหาตุมิจฺฉติ, \\ ยา เอวรูปา ปุริสสฺส ภริยา.}}

\csromanheader{2}{โจรี จ ภริยาติ จ สา ปวุจฺจติ. (2)}
\gambhira{สตฺตกนิปาตปาฬิ}
\csromangathameasure{อกมฺมกามา อลสา มหคฺฆสา, \\ ผรุสา จ จณฺฑี ทุรุตฺตวาทินี. \\ อุฏฺฐายกานํ อภิภุยฺย วตฺตติ, \\ ยา เอวรูปา ปุริสสฺส ภริยา.}
\csromangathagroup{\csromangathastanza{\csromanfolio{468}อกมฺมกามา อลสา มหคฺฆสา, \\ ผรุสา จ จณฺฑี ทุรุตฺตวาทินี. \\ อุฏฺฐายกานํ อภิภุยฺย วตฺตติ, \\ ยา เอวรูปา ปุริสสฺส ภริยา.}}

\csromanheader{2}{อยฺยา จ ภริยาติ จ สา ปวุจฺจติ. (3)}
\csromangathameasure{ยา สพฺพทา โหติ หิตานุกมฺปินี, \\ มาตาว ปุตฺตํ อนุรกฺขเต ปติํ. \\ ตโต ธนํ สมฺภตมสฺส รกฺขติ, \\ ยา เอวรูปา ปุริสสฺส ภริยา.}
\csromangathagroup{\csromangathastanza{ยา สพฺพทา โหติ หิตานุกมฺปินี, \\ มาตาว ปุตฺตํ อนุรกฺขเต ปติํ. \\ ตโต ธนํ สมฺภตมสฺส รกฺขติ, \\ ยา เอวรูปา ปุริสสฺส ภริยา.}}

\csromanheader{2}{มาตา จ ภริยาติ จ สา ปวุจฺจติ. (4)}
\csromangathameasure{ยถาปิ เชฏฺฐา ภคินี กนิฏฺฐกา, \\ สคารวา โหติ สกมฺหิ สามิเก. \\ หิรีมนา ภตฺตุวสานุวตฺตินี, \\ ยา เอวรูปา ปุริสสฺส ภริยา.}
\csromangathagroup{\csromangathastanza{ยถาปิ เชฏฺฐา ภคินี กนิฏฺฐกา\footnote{กณิฏฺฐา (สี), กนิฏฺฐา (สฺยา)}, \\ สคารวา โหติ สกมฺหิ สามิเก. \\ หิรีมนา ภตฺตุวสานุวตฺตินี, \\ ยา เอวรูปา ปุริสสฺส ภริยา.}}

\csromanheader{2}{ภคินี จ ภริยาติ จ สา ปวุจฺจติ. (5)}
\csromangathameasure{ยาจีธ ทิสฺวาน ปติํ ปโมทติ, \\ สขี สขารํว จิรสฺส มาคตํ. \\ โกเลยฺยกา สีลวตี ปติพฺพตา, \\ ยา เอวรูปา ปุริสสฺส ภริยา.}
\csromangathagroup{\csromangathastanza{ยาจีธ ทิสฺวาน ปติํ ปโมทติ, \\ สขี สขารํว จิรสฺส มาคตํ. \\ โกเลยฺยกา สีลวตี ปติพฺพตา, \\ ยา เอวรูปา ปุริสสฺส ภริยา.}}

\csromanheader{2}{สขี จ ภริยาติ จ สา ปวุจฺจติ. (6)}
\csromangathameasure{อกฺกุทฺธสนฺตา วธทณฺฑตชฺชิตา, \\ อทุฏฺฐจิตฺตา ปติโน ติติกฺขติ. \\ อกฺโกธนา ภตฺตุวสานุวตฺตินี, \\ ยา เอวรูปา ปุริสสฺส ภริยา.}
\csromangathagroup{\csromangathastanza{อกฺกุทฺธสนฺตา วธทณฺฑตชฺชิตา, \\ อทุฏฺฐจิตฺตา ปติโน ติติกฺขติ. \\ อกฺโกธนา ภตฺตุวสานุวตฺตินี, \\ ยา เอวรูปา ปุริสสฺส ภริยา.}}

\csromanheader{2}{ทาสี จ ภริยาติ จ สา ปวุจฺจติ. (7)}
\csromangathameasure{ยาจีธ ภริยา วธกาติ วุจฺจติ, \\ โจรี จ อยฺยาติ จ ยา ปวุจฺจติ. \\ ทุสฺสีลรูปา ผรุสา อนาทรา,}
\csromangathagroup{\csromangathastanza{ยาจีธ ภริยา วธกาติ วุจฺจติ, \\ โจรี จ อยฺยาติ จ ยา ปวุจฺจติ. \\ ทุสฺสีลรูปา ผรุสา อนาทรา,}}

\csromanheader{2}{กายสฺส เภทา นิรยํ วชนฺติ ตา. (1-3)}
\chapterheadpageread{6. อพฺยากตวคฺค}
\csromangathameasure{ยาจีธ มาตา ภคินี สขีติ จ, \\ ทาสี จ ภริยาติ จ สา ปวุจฺจติ. \\ สีเล ฐิตตฺตา จิรรตฺตสํวุตา,}
\csromangathagroup{\csromangathastanza{\csromanfolio{469}ยาจีธ มาตา ภคินี สขีติ จ, \\ ทาสี จ ภริยาติ จ สา ปวุจฺจติ. \\ สีเล ฐิตตฺตา จิรรตฺตสํวุตา,}}

\csromanheader{2}{กายสฺส เภทา สุคติํ วชนฺติ ตา”ติ. (4-7)}
\prose{อิมา โข สุชาเต สตฺต ปุริสสฺส ภริยาโย, ตาสํ ตฺวํ กตมาติ.\csromanspacer{} อชฺชตคฺเค มํ ภนฺเต ภควา ทาสีสมํ สามิกสฺส ภริยํ ธาเรตูติ.\csromanspacer{}\csromanspacer{}ทสมํ.}

\csromanmatikaanchor{mtk.259}
\csromantocmarkref{subsection}{11. โกธนสุตฺต}{mtk.259}
\bookmark[dest={mtk.259},level=2]{11. โกธนสุตฺต}
\csromantocmarknopage{chapter}{7. มหาวคฺค}{mtk.260}
\bookmark[dest={mtk.260},level=0]{7. มหาวคฺค}
\csromanheader{2}{11. โกธนสุตฺต}
\proseitem{64}{สตฺติเม ภิกฺขเว ธมฺมา สปตฺตกนฺตา สปตฺตกรณา โกธนํ อาคจฺฉนฺติ อิตฺถิํ วา ปุริสํ วา.\csromanspacer{} กตเม สตฺต? อิธ ภิกฺขเว สปตฺโต สปตฺตสฺส เอวํ อิจฺฉติ “อโห วตายํ ทุพฺพณฺโณ อสฺสา”ติ.\csromanspacer{} ตํ กิสฺส เหตุ? น ภิกฺขเว สปตฺโต สปตฺตสฺส วณฺณวตาย นนฺทติ.\csromanspacer{} โกธโนยํ\footnote{โกธนายํ (ก)} ภิกฺขเว ปุริสปุคฺคโล โกธาภิภูโต โกธปเรโต กิญฺจาปิ โส โหติ สุนฺหาโต สุวิลิตฺโต กปฺปิตเกสมสฺสุ โอทาตวตฺถวสโน\footnote{โอทาตวสโน (ก)}, อถ โข โส ทุพฺพณฺโณว โหติ โกธาภิภูโต.\csromanspacer{} อยํ ภิกฺขเว ปฐโม ธมฺโม สปตฺตกนฺโต สปตฺตกรโณ โกธนํ อาคจฺฉติ อิตฺถิํ วา ปุริสํ วา.}

\prose{ปุน จปรํ ภิกฺขเว สปตฺโต สปตฺตสฺส เอวํ อิจฺฉติ “อโห วตายํ ทุกฺขํ สเยยฺยา”ติ.\csromanspacer{} ตํ กิสฺส เหตุ? น ภิกฺขเว สปตฺโต สปตฺตสฺส สุขเสยฺยาย นนฺทติ.\csromanspacer{} โกธโนยํ ภิกฺขเว ปุริสปุคฺคโล โกธาภิภูโต โกธปเรโต กิญฺจาปิ โส ปลฺลงฺเก เสติ โคนกตฺถเต ปฏลิกตฺถเต กทลิมิคปวรปจฺจตฺถรเณ ส-อุตฺตรจฺฉเท อุภโตโลหิตกูปธาเน.\csromanspacer{} อถ โข โส ทุกฺขญฺเญว เสติ โกธาภิภูโต.\csromanspacer{} อยํ ภิกฺขเว ทุติโย ธมฺโม สปตฺตกนฺโต สปตฺตกรโณ โกธนํ อาคจฺฉติ อิตฺถิํ วา ปุริสํ วา.}

\prose{ปุน จปรํ ภิกฺขเว สปตฺโต สปตฺตสฺส เอวํ อิจฺฉติ “อโห วตายํ น ปจุรตฺโถ อสฺสา”ติ.\csromanspacer{} ตํ กิสฺส เหตุ? น ภิกฺขเว สปตฺโต สปตฺตสฺส ปจุรตฺถตาย นนฺทติ.\csromanspacer{} โกธโนยํ ภิกฺขเว ปุริสปุคฺคโล โกธาภิภูโต โกธปเรโต อนตฺถมฺปิ คเหตฺวา “อตฺโถ เม คหิโต”ติ}

\vspace{\tipitakaparskip}
\csromancenter{สตฺตกนิปาตปาฬิ มญฺญติ.\csromanspacer{} อตฺถมฺปิ คเหตฺวา “อนตฺโถ เม คหิโต”ติ มญฺญติ.\csromanspacer{} ตสฺสิเม ธมฺมา อญฺญมญฺญํ วิปจฺจนีกา คหิตา ทีฆรตฺตํ อหิตาย ทุกฺขาย สํวตฺตนฺติ โกธาภิภูตสฺส.\csromanspacer{} อยํ ภิกฺขเว ตติโย ธมฺโม สปตฺตกนฺโต สปตฺตกรโณ โกธนํ อาคจฺฉติ อิตฺถิํ วา ปุริสํ วา.}
\prose{\csromanfolio{470}ปุน จปรํ ภิกฺขเว สปตฺโต สปตฺตสฺส เอวํ อิจฺฉติ “อโห วตายํ น โภควา อสฺสา”ติ.\csromanspacer{} ตํ กิสฺส เหตุ? น ภิกฺขเว สปตฺโต สปตฺตสฺส โภควตาย นนฺทติ.\csromanspacer{} โกธนสฺส ภิกฺขเว ปุริสปุคฺคลสฺส โกธาภิภูตสฺส โกธปเรตสฺส เยปิสฺส เต โหนฺติ โภคา อุฏฺฐานวีริยาธิคตา พาหาพลปริจิตา เสทาวกฺขิตฺตา ธมฺมิกา ธมฺมลทฺธา, เตปิ ราชาโน ราชโกสํ ปเวเสนฺติ โกธาภิภูตสฺส.\csromanspacer{} อยํ ภิกฺขเว จตุตฺโถ ธมฺโม สปตฺตกนฺโต สปตฺตกรโณ โกธนํ อาคจฺฉติ อิตฺถิํ วา ปุริสํ วา.}

\prose{ปุน จปรํ ภิกฺขเว สปตฺโต สปตฺตสฺส เอวํ อิจฺฉติ “อโห วตายํ น ยสวา อสฺสา”ติ.\csromanspacer{} ตํ กิสฺส เหตุ? น ภิกฺขเว สปตฺโต สปตฺตสฺส ยสวตาย นนฺทติ.\csromanspacer{} โกธโนยํ ภิกฺขเว ปุริสปุคฺคโล โกธาภิภูโต โกธปเรโต.\csromanspacer{} โยปิสฺส โส โหติ ยโส อปฺปมาทาธิคโต ตมฺหาปิ ธํสติ โกธาภิภูโต.\csromanspacer{} อยํ ภิกฺขเว ปญฺจโม ธมฺโม สปตฺตกนฺโต สปตฺตกรโณ โกธนํ อาคจฺฉติ อิตฺถิํ วา ปุริสํ วา.}

\prose{ปุน จปรํ ภิกฺขเว สปตฺโต สปตฺตสฺส เอวํ อิจฺฉติ “อโห วตายํ น มิตฺตวา อสฺสา”ติ.\csromanspacer{} ตํ กิสฺส เหตุ? น ภิกฺขเว สปตฺโต สปตฺตสฺส มิตฺตวตาย นนฺทติ.\csromanspacer{} โกธนํ ภิกฺขเว ปุริสปุคฺคลํ โกธาภิภูตํ โกธปเรตํ เยปิสฺส เต โหนฺติ มิตฺตามจฺจา ญาติสาโลหิตา, เตปิ อารกา ปริวชฺชนฺติ โกธาภิภูตํ.\csromanspacer{} อยํ ภิกฺขเว ฉฏฺโฐ ธมฺโม สปตฺตกนฺโต สปตฺตกรโณ โกธนํ อาคจฺฉติ อิตฺถิํ วา ปุริสํ วา.}

\prose{ปุน จปรํ ภิกฺขเว สปตฺโต สปตฺตสฺส เอวํ อิจฺฉติ “อโห วตายํ กายสฺส เภทา ปรํ มรณา อปายํ ทุคฺคติํ วินิปาตํ นิรยํ อุปปชฺเชยฺยา”ติ.\csromanspacer{} ตํ กิสฺส เหตุ? น ภิกฺขเว สปตฺโต สปตฺตสฺส สุคติคมเน นนฺทติ.\csromanspacer{} โกธโนยํ ภิกฺขเว ปุริสปุคฺคโล โกธาภิภูโต โกธปเรโต กาเยน ทุจฺจริตํ จรติ, วาจาย ทุจฺจริตํ จรติ, มนสา ทุจฺจริตํ จรติ.\csromanspacer{} โส กาเยน ทุจฺจริตํ จริตฺวา วาจาย ฯเปฯ}

\vspace{\tipitakaparskip}
\csromancenter{อพฺยากตวคฺค กายสฺส เภทา ปรํ มรณา อปายํ ทุคฺคติํ วินิปาตํ นิรยํ อุปปชฺชติ โกธาภิภูโต.\csromanspacer{} อยํ ภิกฺขเว สตฺตโม ธมฺโม สปตฺตกนฺโต สปตฺตกรโณ โกธนํ อาคจฺฉติ อิตฺถิํ วา ปุริสํ วา.\csromanspacer{} อิเม โข ภิกฺขเว สตฺต ธมฺมา สปตฺตกนฺตา สปตฺตกรณา โกธนํ อาคจฺฉนฺติ อิตฺถิํ วา ปุริสํ วาติ.}
\prose{\csromanfolio{471}โกธโน ทุพฺพณฺโณ โหติ, อโถ ทุกฺขมฺปิ เสติ โส.}

\prose{อโถ อตฺถํ คเหตฺวาน, อนตฺถํ อธิปชฺชติ\footnote{อธิคจฺฉติ (สี), ปฏิปชฺชติ (สฺยา)}.}

\prose{ตโต กาเยน วาจาย, วธํ กตฺวาน โกธโน.}

\prose{โกธาภิภูโต ปุริโส, ธนชานิํ นิคจฺฉติ.}

\prose{โกธสมฺมทสมฺมตฺโต, อายสกฺยํ\footnote{อายสกฺขํ (สฺยา)} นิคจฺฉติ.}

\prose{ญาติมิตฺตา สุหชฺชา จ, ปริวชฺชนฺติ โกธนํ. * อนตฺถชนโน โกโธ, โกโธ จิตฺตปฺปโกปโน.}

\prose{ภยมนฺตรโต ชาตํ, ตํ ชโน นาวพุชฺฌติ.}

\prose{กุทฺโธ อตฺถํ น ชานาติ, กุทฺโธ ธมฺมํ น ปสฺสติ.}

\prose{อนฺธตมํ ตทา โหติ, ยํ โกโธ สหเต นรํ.}

\prose{ยํ กุทฺโธ อุปโรเธติ, สุกรํ วิย ทุกฺกรํ.}

\prose{ปจฺฉา โส วิคเต โกเธ, อคฺคิทฑฺโฒว ตปฺปติ.}

\prose{ทุมฺมงฺกุยํ ปทสฺเสติ\footnote{สทสฺเสติ (สี), ปฐมํ ทสฺเสติ (สฺยา)}, ธูมํ ธูมีว ปาปโก.}

\prose{ยโต ปตายติ โกโธ, เยน กุชฺฌนฺติ มานวา.}

\prose{นาสฺส หิรี น โอตฺตปฺปํ\footnote{น อสฺส หิรี โอตฺตปฺปญฺจ (ก)}, น วาโจ โหติ คารโว.}

\prose{โกเธน อภิภูตสฺส, น ทีปํ โหติ กิญฺจนํ.}

\prose{ตปนียานิ กมฺมานิ, ยานิ ธมฺเมหิ อารกา.}

\prose{ตานิ อาโรจยิสฺสามิ, ตํ สุณาถ ยถา ตถํ.}

\prose{กุทฺโธ หิ ปิตรํ หนฺติ, หนฺติ กุทฺโธ สมาตรํ.}

\prose{กุทฺโธ หิ พฺราหฺมณํ หนฺติ, หนฺติ กุทฺโธ ปุถุชฺชนํ.}

\gambhira{สตฺตกนิปาตปาฬิ}
\prose{\csromanfolio{472}ยาย มาตุ ภโต โปโส, อิมํ โลกํ อเวกฺขติ.}

\prose{ตมฺปิ ปาณททิํ สนฺติํ, หนฺติ กุทฺโธ ปุถุชฺชโน.}

\prose{อตฺตูปมา หิ เต สตฺตา, อตฺตา หิ ปรโม\footnote{ปรํ (สี, สฺยา)} ปิโย.}

\prose{หนฺติ กุทฺโธ ปุถุตฺตานํ, นานารูเปสุ มุจฺฉิโต.}

\prose{อสินา หนฺติ อตฺตานํ, วิสํ ขาทนฺติ มุจฺฉิตา.}

\prose{รชฺชุยา พชฺฌ มียนฺติ, ปพฺพตามปิ กนฺทเร.}

\prose{ภูนหจฺจานิ\footnote{ภูตหจฺจานิ (สี, สฺยา)} กมฺมานิ, อตฺตมารณิยานิ จ.}

\prose{กโรนฺตา นาวพุชฺฌนฺติ\footnote{กโรนฺโต นาวพุชฺฌติ (ก)}, โกธชาโต ปราภโว.}

\prose{อิตายํ โกธรูเปน, มจฺจุปาโส คุหาสโย.}

\prose{ตํ ทเมน สมุจฺฉินฺเท, ปญฺญาวีริเยน ทิฏฺฐิยา.}

\prose{ยถา เมตํ\footnote{เอกเมกํ (สฺยา), เอกเมตํ (สี)} อกุสลํ, สมุจฺฉินฺเทถ ปณฺฑิโต.}

\prose{ตเถว ธมฺเม สิกฺเขถ, มา โน ทุมฺมงฺกุยํ อหุ.}

\prose{วีตโกธา อนายาสา, วีตโลภา อนุสฺสุกา\footnote{อนิสฺสุกา (สี, สฺยา) ตทฏฺฐกถาสุ ปน “อนุสฺสุกา”เตฺวว ทิสฺสติ.}.}

\prose{ทนฺตา โกธํ ปหนฺตฺวาน, ปรินิพฺพนฺติ อนาสวาติ\footnote{ปรินิพฺพิสฺสถนาสวาติ (สฺยา), ปรินิพฺพิํสุ อนาสวาติ (ก)}.\csromanspacer{}\csromanspacer{}เอกาทสมํ.\csromanspacer{} อพฺยากตวคฺโค ฉฏฺโฐ.}

\titlehead{\textbf{ตสฺสุทฺทานํ}}
\csromangathameasure{อพฺยากโต ปุริสคติ, ติสฺส สีห อรกฺขิยํ. \\ กิมิลํ สตฺต ปจลา, เมตฺตา ภริยา โกเธกาทสาติ.}
\csromangathagroup{\csromangathastanza{อพฺยากโต ปุริสคติ, ติสฺส สีห อรกฺขิยํ. \\ กิมิลํ สตฺต ปจลา, เมตฺตา ภริยา โกเธกาทสาติ.}}

\csromanmatikaanchor{mtk.261}
\csromantocmarkref{subsection}{1. หิรี-โอตฺตปฺปสุตฺต}{mtk.261}
\bookmark[dest={mtk.261},level=2]{1. หิรี-โอตฺตปฺปสุตฺต}
\csromanheader{2}{7. มหาวคฺค \textbf{1. หิรี โอตฺตปฺปสุตฺต}}
\proseitem{65}{\csromansymbolfootnote{*}{เหฏฺฐา 16, 175, 316 ปิฏฺเฐสุปิ.}หิโรตฺตปฺเป ภิกฺขเว อสติ หิโรตฺตปฺปวิปนฺนสฺส หตูปนิโส โหติ อินฺทฺริยสํวโร, อินฺทฺริยสํวเร อสติ อินฺทฺริยสํวรวิปนฺนสฺส}

\vspace{\tipitakaparskip}
\csromancenter{มหาวคฺค หตูปนิสํ โหติ สีลํ, สีเล อสติ สีลวิปนฺนสฺส หตูปนิโส โหติ สมฺมาสมาธิ, สมฺมาสมาธิมฺหิ อสติ สมฺมาสมาธิวิปนฺนสฺส หตูปนิสํ โหติ ยถาภูตญาณทสฺสนํ, ยถาภูตญาณทสฺสเน อสติ ยถาภูตญาณ ทสฺสนวิปนฺนสฺส หตูปนิโส โหติ นิพฺพิทาวิราโค, นิพฺพิทาวิราเค อสติ นิพฺพิทาวิราควิปนฺนสฺส หตูปนิสํ โหติ วิมุตฺติญาณทสฺสนํ.\csromanspacer{} เสยฺยถาปิ ภิกฺขเว รุกฺโข สาขาปลาสวิปนฺโน, ตสฺส ปปฏิกาปิ น ปาริปูริํ คจฺฉติ.\csromanspacer{} ตโจปิ เผคฺคุปิ สาโรปิ น ปาริปูริํ คจฺฉติ.\csromanspacer{} เอวเมวํ โข ภิกฺขเว หิโรตฺตปฺเป อสติ หิโรตฺตปฺปวิปนฺนสฺส หตูปนิโส โหติ อินฺทฺริยสํวโร, อินฺทฺริยสํวเร อสติ อินฺทฺริยสํวรวิปนฺนสฺส หตูปนิสํ โหติ สีลํ, สีเล อสติ สีลวิปนฺนสฺส หตูปนิโส โหติ สมฺมาสมาธิ, สมฺมาสมาธิมฺหิ อสติ สมฺมาสมาธิวิปนฺนสฺส หตูปนิสํ โหติ ยถาภูตญาณทสฺสนํ, ยถาภูตญาณทสฺสเน อสติ ยถาภูตญาณ ทสฺสนวิปนฺนสฺส หตูปนิโส โหติ นิพฺพิทาวิราโค, นิพฺพิทาวิราเค อสติ นิพฺพิทาวิราควิปนฺนสฺส หตูปนิสํ โหติ วิมุตฺติญาณทสฺสนํ.}
\prose{\csromanfolio{473}หิโรตฺตปฺเป ภิกฺขเว สติ หิโรตฺตปฺปสมฺปนฺนสฺส อุปนิสสมฺปนฺโน โหติ อินฺทฺริยสํวโร, อินฺทฺริยสํวเร สติ อินฺทฺริยสํวรสมฺปนฺนสฺส อุปนิสสมฺปนฺนํ โหติ สีลํ, สีเล สติ สีลสมฺปนฺนสฺส อุปนิสสมฺปนฺโน โหติ สมฺมาสมาธิ, สมฺมาสมาธิมฺหิ สติ สมฺมาสมาธิสมฺปนฺนสฺส อุปนิสสมฺปนฺนํ โหติ ยถาภูตญาณทสฺสนํ, ยถาภูตญาณทสฺสเน สติ ยถาภูตญาณ ทสฺสนสมฺปนฺนสฺส อุปนิสสมฺปนฺโน โหติ นิพฺพิทาวิราโค, นิพฺพิทาวิราเค สติ นิพฺพิทาวิราคสมฺปนฺนสฺส อุปนิสสมฺปนฺนํ โหติ วิมุตฺติญาณทสฺสนํ.\csromanspacer{} เสยฺยถาปิ ภิกฺขเว รุกฺโข สาขาปลาสสมฺปนฺโน, ตสฺส ปปฏิกาปิ ปาริปูริํ คจฺฉติ, ตโจปิ เผคฺคุปิ สาโรปิ ปาริปูริํ คจฺฉติ.\csromanspacer{} เอวเมวํ โข ภิกฺขเว หิโรตฺตปฺเป สติ หิโรตฺตปฺปสมฺปนฺนสฺส อุปนิสสมฺปนฺโน โหติ ฯเปฯ วิมุตฺติญาณทสฺสนนฺติ.\csromanspacer{}\csromanspacer{}ปฐมํ.}

\csromanmatikaanchor{mtk.262}
\csromantocmarkref{subsection}{2. สตฺตสูริยสุตฺต}{mtk.262}
\bookmark[dest={mtk.262},level=2]{2. สตฺตสูริยสุตฺต}
\csromanheader{2}{2. สตฺตสูริยสุตฺต}
\proseitem{66}{เอวํ เม สุตํ\csromandash{}เอกํ สมยํ ภควา เวสาลิยํ วิหรติ อมฺพปาลิวเน.\csromanspacer{} ตตฺร โข ภควา ภิกฺขู อามนฺเตสิ “ภิกฺขโว”ติ. “ภทนฺเต”ติ เต ภิกฺขู ภควโต ปจฺจสฺโสสุํ.\csromanspacer{} ภควา เอตทโวจ\csromandash{}}

\gambhira{สตฺตกนิปาตปาฬิ}
\prose{\csromanfolio{474}อนิจฺจา ภิกฺขเว สงฺขารา, อธุวา ภิกฺขเว สงฺขารา, อนสฺสาสิกา ภิกฺขเว สงฺขารา, ยาวญฺจิทํ ภิกฺขเว อลเมว สพฺพสงฺขาเรสุ นิพฺพินฺทิตุํ อลํ วิรชฺชิตุํ อลํ วิมุจฺจิตุํ.}

\prose{สิเนรุ ภิกฺขเว ปพฺพตราชา จตุราสีติโยชนสหสฺสานิ อายาเมน, จตุราสีติโยชนสหสฺสานิ วิตฺถาเรน, จตุราสีติโยชนสหสฺสานิ มหาสมุทฺเท อชฺโฌคาโฬฺห, จตุราสีติโยชนสหสฺสานิ มหาสมุทฺทา อจฺจุคฺคโต.\csromanspacer{} โหติ โข โส ภิกฺขเว สมโย, ยํ กทาจิ กรหจิ ทีฆสฺส อทฺธุโน อจฺจเยน พหูนิ วสฺสานิ พหูนิ วสฺสสตานิ พหูนิ วสฺสสหสฺสานิ พหูนิ วสฺสสตสหสฺสานิ เทโว น วสฺสติ เทเว โข ปน ภิกฺขเว อวสฺสนฺเต เย เกจิเม พีชคามภูตคามา โอสธิติณวนปฺปตโย, เต อุสฺสุสฺสนฺติ วิสุสฺสนฺติ น ภวนฺติ.\csromanspacer{} เอวํ อนิจฺจา ภิกฺขเว สงฺขารา, เอวํ อธุวา ภิกฺขเว สงฺขารา ฯเปฯ อลํ วิมุจฺจิตุํ.}

\prose{โหติ โข โส ภิกฺขเว สมโย, ยํ กทาจิ กรหจิ ทีฆสฺส อทฺธุโน อจฺจเยน ทุติโย สูริโย ปาตุภวติ, ทุติยสฺส ภิกฺขเว สูริยสฺส ปาตุภาวา ยา กาจิ กุนฺนทิโย กุโสพฺภา\footnote{กุสฺสุพฺภา (สี), กุสฺโสพฺภา (สฺยา)}, ตา อุสฺสุสฺสนฺติ วิสุสฺสนฺติ น ภวนฺติ.\csromanspacer{} เอวํ อนิจฺจา ภิกฺขเว สงฺขารา ฯเปฯ อลํ วิมุจฺจิตุํ.}

\prose{โหติ โข โส ภิกฺขเว สมโย, ยํ กทาจิ กรหจิ ทีฆสฺส อทฺธุโน อจฺจเยน ตติโย สูริโย ปาตุภวติ, ตติยสฺส ภิกฺขเว สูริยสฺส ปาตุภาวา ยา กาจิ มหานทิโย.\csromanspacer{} เสยฺยถิทํ, คงฺคา ยมุนา อจิรวตี สรภู มหี, ตา อุสฺสุสฺสนฺติ วิสุสฺสนฺติ น ภวนฺติ.\csromanspacer{} เอวํ อนิจฺจา ภิกฺขเว สงฺขารา ฯเปฯ อลํ วิมุจฺจิตุํ.}

\prose{โหติ โข โส ภิกฺขเว สมโย, ยํ กทาจิ กรหจิ ทีฆสฺส อทฺธุโน อจฺจเยน จตุตฺโถ สูริโย ปาตุภวติ, จตุตฺถสฺส ภิกฺขเว สูริยสฺส ปาตุภาวา เย เต มหาสรา, ยโต อิมา มหานทิโย ปวตฺตนฺติ.\csromanspacer{} เสยฺยถิทํ, อโนตตฺตา สีหปปาตา รถการา กณฺณมุณฺฑา กุณาลา ฉทฺทนฺตา มนฺทากินิยา, ตา อุสฺสุสฺสนฺติ วิสุสฺสนฺติ น ภวนฺติ.\csromanspacer{} เอวํ อนิจฺจา ภิกฺขเว สงฺขารา ฯเปฯ อลํ วิมุจฺจิตุํ.}

\prose{โหติ โข โส ภิกฺขเว สมโย, ยํ กทาจิ กรหจิ ทีฆสฺส อทฺธุโน อจฺจเยน ปญฺจโม สูริโย ปาตุภวติ, ปญฺจมสฺส ภิกฺขเว}

\vspace{\tipitakaparskip}
\csromancenter{มหาวคฺค สูริยสฺส ปาตุภาวา โยชนสติกานิปิ มหาสมุทฺเท อุทกานิ โอคจฺฉนฺติ, ทฺวิโยชนสติกานิปิ มหาสมุทฺเท อุทกานิ โอคจฺฉนฺติ, ติโยชนสติกานิปิ จตุโยชนสติกานิปิ ปญฺจโยชนสติกานิปิ ฉโยชนสติกานิปิ สตฺตโยชนสติกานิปิ มหาสมุทฺเท อุทกานิ โอคจฺฉนฺติ.\csromanspacer{} สตฺตตาลมฺปิ มหาสมุทฺเท อุทกํ สณฺฐาติ, ฉตาลมฺปิ ปญฺจตาลมฺปิ จตุตาลมฺปิ ติตาลมฺปิ ทฺวิตาลมฺปิ ตาลมตฺตมฺปิ มหาสมุทฺเท อุทกํ สณฺฐาติ, สตฺตโปริสมฺปิ มหาสมุทฺเท อุทกํ สณฺฐาติ, ฉโปริสมฺปิ ปญฺจโปริสมฺปิ จตุโปริสมฺปิ ติโปริสมฺปิ ทฺวิโปริสมฺปิ โปริสมฺปิ\footnote{โปริสมตฺตมฺปิ (สฺยา)} อฑฺฒโปริสมฺปิ กฏิมตฺตมฺปิ ชณฺณุกามตฺตมฺปิ โคปฺผกมตฺตมฺปิ มหาสมุทฺเท อุทกํ สณฺฐาติ.\csromanspacer{} เสยฺยถาปิ ภิกฺขเว สรทสมเย ถุลฺลผุสิตเก เทเว วสฺสนฺเต ตตฺถ ตตฺถ โคปเทสุ\footnote{โคปฺผกปเทเสสุ (ก)} อุทกานิ ฐิตานิ โหนฺติ.\csromanspacer{} เอวเมวํ โข ภิกฺขเว ตตฺถ ตตฺถ โคปฺผกมตฺตานิ\footnote{โคปทมตฺตานิ (สี, สฺยา) 4. ธูปายนฺติ สนฺธูปายนฺติ สมฺปธูปายนฺติ (สี, สฺยา)} มหาสมุทฺเท อุทกานิ ฐิตานิ โหนฺติ, ปญฺจมสฺส ภิกฺขเว สูริยสฺส ปาตุภาวา องฺคุลิปพฺพมตฺตมฺปิ มหาสมุทฺเท อุทกํ น โหติ.\csromanspacer{} เอวํ อนิจฺจา ภิกฺขเว สงฺขารา ฯเปฯ อลํ วิมุจฺจิตุํ.}
\prose{\csromanfolio{475}โหติ โข โส ภิกฺขเว สมโย, ยํ กทาจิ กรหจิ ทีฆสฺส อทฺธุโน อจฺจเยน ฉฏฺโฐ สูริโย ปาตุภวติ, ฉฏฺฐสฺส ภิกฺขเว สูริยสฺส ปาตุภาวา อยญฺจ มหาปถวี สิเนรุ จ ปพฺพตราชา ธูมายนฺติ สํธูมายนฺติ สมฺปธูมายนฺติ4.\csromanspacer{} เสยฺยถาปิ ภิกฺขเว กุมฺภการปาโก อาเลปิโต\footnote{อาลิมฺปิโต (สี, สฺยา)} ปฐมํ ธูเมติ สํธูเมติ สมฺปธูเมติ.\csromanspacer{} เอวเมวํ โข ภิกฺขเว ฉฏฺฐสฺส สูริยสฺส ปาตุภาวา อยญฺจ มหาปถวี สิเนรุ จ ปพฺพตราชา ธูมายนฺติ สํธูมายนฺติ สมฺปธูมายนฺติ.\csromanspacer{} เอวํ อนิจฺจา ภิกฺขเว สงฺขารา ฯเปฯ อลํ วิมุจฺจิตุํ.}

\prose{โหติ โข โส ภิกฺขเว สมโย, ยํ กทาจิ กรหจิ ทีฆสฺส อทฺธุโน อจฺจเยน สตฺตโม สูริโย ปาตุภวติ, สตฺตมสฺส ภิกฺขเว สูริยสฺส ปาตุภาวา อยญฺจ มหาปถวี สิเนรุ จ ปพฺพตราชา อาทิปฺปนฺติ ปชฺชลนฺติ เอกชาลา ภวนฺติ.\csromanspacer{} อิมิสฺสา จ ภิกฺขเว มหาปถวิยา}

\vspace{\tipitakaparskip}
\csromancenter{สตฺตกนิปาตปาฬิ สิเนรุสฺส จ ปพฺพตราชสฺส ฌายมานานํ ทยฺหมานานํ อจฺจิ วาเตน ขิตฺตา ยาว พฺรหฺมโลกาปิ คจฺฉติ, สิเนรุสฺส ภิกฺขเว ปพฺพตราชสฺส ฌายมานสฺส ทยฺหมานสฺส วินสฺสมานสฺส มหตา เตโชขนฺเธน อภิภูตสฺส โยชนสติกานิปิ กูฏานิ ปลุชฺชนฺติ ทฺวิโยชนสติกานิปิ ติโยชนสติกานิปิ จตุโยชนสติกานิปิ ปญฺจโยชนสติกานิปิ กูฏานิ ปลุชฺชนฺติ.\csromanspacer{} อิมิสฺสา จ ภิกฺขเว มหาปถวิยา สิเนรุสฺส จ ปพฺพตราชสฺส ฌายมานานํ ทยฺหมานานํ เนว ฉาริกา ปญฺญายติ น มสิ, เสยฺยถาปิ ภิกฺขเว สปฺปิสฺส วา เตลสฺส วา ฌายมานสฺส ทยฺหมานสฺส เนว ฉาริกา ปญฺญายติ น มสิ.\csromanspacer{} เอวเมวํ โข ภิกฺขเว อิมิสฺสา จ มหาปถวิยา สิเนรุสฺส จ ปพฺพตราชสฺส ฌายมานานํ ทยฺหมานานํ เนว ฉาริกา ปญฺญายติ น มสิ.\csromanspacer{} เอวํ อนิจฺจา ภิกฺขเว สงฺขารา, เอวํ อธุวา ภิกฺขเว สงฺขารา, เอวํ อนสฺสาสิกา ภิกฺขเว สงฺขารา, ยาวญฺจิทํ ภิกฺขเว อลเมว สพฺพสงฺขาเรสุ นิพฺพินฺทิตุํ อลํ วิรชฺชิตุํ อลํ วิมุจฺจิตุํ.}
\prose{\csromanfolio{476}ตตฺร ภิกฺขเว โก มนฺตา โก สทฺธาตา “อยญฺจ ปถวี สิเนรุ จ ปพฺพตราชา ทยฺหิสฺสนฺติ วินสฺสิสฺสนฺติ น ภวิสฺสนฺตี”ติ อญฺญตฺร ทิฏฺฐปเทหิ.}

\prose{\csromansymbolfootnote{*}{เหฏฺฐา 326, อุปริ 500 ปิฏฺเฐปิ.}ภูตปุพฺพํ ภิกฺขเว สุเนตฺโต นาม สตฺถา อโหสิ ติตฺถกโร กาเมสุ วีตราโค, สุเนตฺตสฺส โข ปน ภิกฺขเว สตฺถุโน อเนกานิ สาวกสตานิ อเหสุํ, สุเนตฺโต ภิกฺขเว สตฺถา สาวกานํ พฺรหฺมโลกสหพฺยตาย ธมฺมํ เทเสสิ.\csromanspacer{} เย โข ปน ภิกฺขเว สุเนตฺตสฺส สตฺถุโน พฺรหฺมโลกสหพฺยตาย ธมฺมํ เทเสนฺตสฺส สพฺเพน สพฺพํ สาสนํ อาชานิํสุ.\csromanspacer{} เต กายสฺส เภทา ปรํ มรณา สุคติํ พฺรหฺมโลกํ อุปปชฺชิํสุ.\csromanspacer{} เย น สพฺเพน สพฺพํ สาสนํ อาชานิํสุ.\csromanspacer{} เต กายสฺส เภทา ปรํ มรณา อปฺเปกจฺเจ ปรนิมฺมิตวสวตฺตีนํ เทวานํ สหพฺยตํ อุปปชฺชิํสุ, อปฺเปกจฺเจ นิมฺมานรตีนํ เทวานํ สหพฺยตํ อุปปชฺชิํสุ, อปฺเปกจฺเจ ตุสิตานํ เทวานํ สหพฺยตํ อุปปชฺชิํสุ, อปฺเปกจฺเจ ยามานํ เทวานํ สหพฺยตํ อุปปชฺชิํสุ, อปฺเปกจฺเจ ตาวติํสานํ เทวานํ สหพฺยตํ อุปปชฺชิํสุ, อปฺเปกจฺเจ จาตุมหาราชิกานํ เทวานํ สหพฺยตํ อุปปชฺชิํสุ, อปฺเปกจฺเจ ขตฺติยมหาสาลานํ สหพฺยตํ อุปปชฺชิํสุ, อปฺเปกจฺเจ พฺราหฺมณมหาสาลานํ สหพฺยตํ อุปปชฺชิํสุ, อปฺเปกจฺเจ คหปติมหาสาลานํ สหพฺยตํ อุปปชฺชิํสุ.}

\chapterheadpageread{7. มหาวคฺค}
\prose{\csromanfolio{477}อถ โข ภิกฺขเว สุเนตฺตสฺส สตฺถุโน เอตทโหสิ “น โข เมตํ ปติรูปํ, โยหํ สาวกานํ สมสมคติโย อสฺสํ อภิสมฺปรายํ, ยํนูนาหํ อุตฺตริ เมตฺตํ\footnote{อุตฺตริ มคฺคํ (ก)} ภาเวยฺยนฺ”ติ.}

\prose{อถ โข ภิกฺขเว สุเนตฺโต สตฺถา สตฺต วสฺสานิ เมตฺตํ จิตฺตํ ภาเวสิ, สตฺต วสฺสานิ เมตฺตํ จิตฺตํ ภาเวตฺวา สตฺต สํวฏฺฏวิวฏฺฏกปฺเป นยิมํ โลกํ ปุนราคมาสิ.\csromanspacer{} สํวฏฺฏมาเน สุทํ ภิกฺขเว โลเก อาภสฺสรูปโค โหติ, วิวฏฺฏมาเน โลเก สุญฺญํ พฺรหฺมวิมานํ อุปปชฺชติ.\csromanspacer{} ตตฺร สุทํ ภิกฺขเว พฺรหฺมา โหติ มหาพฺรหฺมา อภิภู อนภิภูโต อญฺญทตฺถุทโส วสวตฺตี.\csromanspacer{} ฉตฺติํสกฺขตฺตุํ โข ปน ภิกฺขเว สกฺโก อโหสิ เทวานมินฺโท, อเนกสตกฺขตฺตุํ ราชา อโหสิ จกฺกวตฺตี ธมฺมิโก ธมฺมราชา จาตุรนฺโต วิชิตาวี ชนปทตฺถาวริยปฺปตฺโต สตฺตรตนสมนฺนาคโต, ปโรสหสฺสํ โข ปนสฺส ปุตฺตา อเหสุํ สูรา วีรงฺครูปา ปรเสนปฺปมทฺทนา.\csromanspacer{} โส อิมํ ปถวิํ สาครปริยนฺตํ อทณฺเฑน อสตฺเถน ธมฺเมน อภิวิชิย อชฺฌาวสิ.\csromanspacer{} โส หิ นาม ภิกฺขเว สุเนตฺโต สตฺถา เอวํ ทีฆายุโก สมาโน เอวํ จิรฏฺฐิติโก อปริมุตฺโต อโหสิ ชาติยา ชราย มรเณน โสเกหิ ปริเทเวหิ ทุกฺเขหิ โทมนสฺเสหิ อุปายาเสหิ, อปริมุตฺโต ทุกฺขสฺมาติ วทามิ.}

\prose{ตํ กิสฺส เหตุ? จตุนฺนํ ธมฺมานํ อนนุโพธา อปฺปฏิเวธา.\csromanspacer{} กตเมสํ จตุนฺนํ? อริยสฺส ภิกฺขเว สีลสฺส อนนุโพธา อปฺปฏิเวธา อริยสฺส สมาธิสฺส อนนุโพธา อปฺปฏิเวธา อริยาย ปญฺญาย อนนุโพธา อปฺปฏิเวธา อริยาย วิมุตฺติยา อนนุโพธา อปฺปฏิเวธา, ตยิทํ ภิกฺขเว อริยํ สีลํ อนุพุทฺธํ ปฏิวิทฺธํ, อริโย สมาธิ อนุโพโธ ปฏิวิทฺโธ, อริยา ปญฺญา อนุโพธา ปฏิวิทฺธา, อริยา วิมุตฺติ อนุโพธา ปฏิวิทฺธา, อุจฺฉินฺนา ภวตณฺหา, ขีณา ภวเนตฺติ, นตฺถิ ทานิ ปุนพฺภโวติ.\csromanspacer{} อิทมโวจ ภควา, อิทํ วตฺวาน สุคโต, อถาปรํ เอตทโวจ สตฺถา\csromandash{}}

\csromangathameasure{“สีลํ สมาธิ ปญฺญา จ, วิมุตฺติ จ อนุตฺตรา. \\ อนุพุทฺธา อิเม ธมฺมา, โคตเมน ยสสฺสินา. \\ อิติ พุทฺโธ อภิญฺญาย, ธมฺมมกฺขาสิ ภิกฺขุนํ. \\ ทุกฺขสฺสนฺตกโร สตฺถา, จกฺขุมา ปรินิพฺพุโต”ติ.ทุติยํ.}
\csromangathagroup{\csromangathastanza{“สีลํ สมาธิ ปญฺญา จ, วิมุตฺติ จ อนุตฺตรา. \\ อนุพุทฺธา อิเม ธมฺมา, โคตเมน ยสสฺสินา.}\csromangathastanza{อิติ พุทฺโธ อภิญฺญาย, ธมฺมมกฺขาสิ ภิกฺขุนํ. \\ ทุกฺขสฺสนฺตกโร สตฺถา, จกฺขุมา ปรินิพฺพุโต”ติ.\csromanspacer{}\csromanspacer{}ทุติยํ.}}

\csromanmatikaanchor{mtk.263}
\csromantocmarkref{subsection}{3. นคโรปมสุตฺต}{mtk.263}
\bookmark[dest={mtk.263},level=2]{3. นคโรปมสุตฺต}
\csromanheader{2}{สตฺตกนิปาตปาฬิ \textbf{3. นคโรปมสุตฺต}}
\proseitem{67}{\csromanfolio{478}ยโต โข ภิกฺขเว รญฺโญ ปจฺจนฺติมํ นครํ สตฺตหิ นครปริกฺขาเรหิ สุปริกฺขตํ\footnote{สุปริกฺขิตฺตํ (ก)} โหติ, จตุนฺนญฺจ อาหารานํ นิกามลาภี โหติ อกิจฺฉลาภี อกสิรลาภี.\csromanspacer{} อิทํ วุจฺจติ ภิกฺขเว รญฺโญ ปจฺจนฺติมํ นครํ อกรณียํ พาหิเรหิ ปจฺจตฺถิเกหิ ปจฺจามิตฺเตหิ.}

\prose{กตเมหิ สตฺตหิ นครปริกฺขาเรหิ สุปริกฺขตํ โหติ? อิธ ภิกฺขเว รญฺโญ ปจฺจนฺติเม นคเร เอสิกา โหติ คมฺภีรเนมา\footnote{คมฺภีรเนมิ (ก)} สุนิขาตา อจลา อสมฺปเวธี\footnote{อสมฺปเวธิ (สี, สฺยา)}.\csromanspacer{} อิมินา ปฐเมน นครปริกฺขาเรน สุปริกฺขตํ โหติ รญฺโญ ปจฺจนฺติมํ นครํ อพฺภนฺตรานํ คุตฺติยา พาหิรานํ ปฏิฆาตาย. (1)}

\prose{ปุน จปรํ ภิกฺขเว รญฺโญ ปจฺจนฺติเม นคเร ปริกฺขา โหติ คมฺภีรา เจว วิตฺถตา จ.\csromanspacer{} อิมินา ทุติเยน นครปริกฺขาเรน สุปริกฺขตํ โหติ รญฺโญ ปจฺจนฺติมํ นครํ อพฺภนฺตรานํ คุตฺติยา พาหิรานํ ปฏิฆาตาย. (2)}

\prose{ปุน จปรํ ภิกฺขเว รญฺโญ ปจฺจนฺติเม นคเร อนุปริยายปโถ โหติ อุจฺโจ เจว วิตฺถโต จ.\csromanspacer{} อิมินา ตติเยน นครปริกฺขาเรน สุปริกฺขตํ โหติ รญฺโญ ปจฺจนฺติมํ นครํ อพฺภนฺตรานํ คุตฺติยา พาหิรานํ ปฏิฆาตาย. (3)}

\prose{ปุน จปรํ ภิกฺขเว รญฺโญ ปจฺจนฺติเม นคเร พหุํ อาวุธํ สนฺนิจิตํ โหติ สลากญฺเจว เชวนิกญฺจ\footnote{เชวนิยญฺจ (สี-ฏฺฐ)}.\csromanspacer{} อิมินา จตุตฺเถน นครปริกฺขาเรน สุปริกฺขตํ โหติ รญฺโญ ปจฺจนฺติมํ นครํ อพฺภนฺตรานํ คุตฺติยา พาหิรานํ ปฏิฆาตาย. (4)}

\prose{ปุน จปรํ ภิกฺขเว รญฺโญ ปจฺจนฺติเม นคเร พหุพลกาโย ปฏิวสติ.\csromanspacer{} เสยฺยถิทํ, หตฺถาโรหา อสฺสาโรหา รถิกา ธนุคฺคหา เจลกา จลกา ปิณฺฑทายกา อุคฺคา ราชปุตฺตา ปกฺขนฺทิโน มหานาคา สูรา จมฺมโยธิโน ทาสกปุตฺตา.\csromanspacer{} อิมินา ปญฺจเมน นครปริกฺขาเรน สุปริกฺขตํ โหติ รญฺโญ ปจฺจนฺติมํ นครํ อพฺภนฺตรานํ คุตฺติยา พาหิรานํ ปฏิฆาตาย. (5)}

\prose{ปุน จปรํ ภิกฺขเว รญฺโญ ปจฺจนฺติเม นคเร โทวาริโก โหติ ปณฺฑิโต พฺยตฺโต เมธาวี อญฺญาตานํ นิวาเรตา ญาตานํ ปเวเสตา.}

\vspace{\tipitakaparskip}
\csromancenter{มหาวคฺค อิมินา ฉฏฺเฐน นครปริกฺขาเรน สุปริกฺขตํ โหติ รญฺโญ ปจฺจนฺติมํ นครํ อพฺภนฺตรานํ คุตฺติยา พาหิรานํ ปฏิฆาตาย. (6)}
\prose{\csromanfolio{479}ปุน จปรํ ภิกฺขเว รญฺโญ ปจฺจนฺติเม นคเร ปากาโร โหติ อุจฺโจ เจว วิตฺถโต จ วาสนเลปนสมฺปนฺโน จ.\csromanspacer{} อิมินา สตฺตเมน นครปริกฺขาเรน สุปริกฺขตํ โหติ รญฺโญ ปจฺจนฺติมํ นครํ อพฺภนฺตรานํ คุตฺติยา พาหิรานํ ปฏิฆาตาย.\csromanspacer{} อิเมหิ สตฺตหิ นครปริกฺขาเรหิ สุปริกฺขตํ โหติ. (7)}

\prose{กตเมสํ จตุนฺนํ อาหารานํ นิกามลาภี โหติ อกิจฺฉลาภี อกสิรลาภี? อิธ ภิกฺขเว รญฺโญ ปจฺจนฺติเม นคเร พหุํ ติณกฏฺโฐทกํ สนฺนิจิตํ โหติ อพฺภนฺตรานํ รติยา อปริตสฺสาย ผาสุวิหาราย พาหิรานํ ปฏิฆาตาย. (1)}

\prose{ปุน จปรํ ภิกฺขเว รญฺโญ ปจฺจนฺติเม นคเร พหุํ สาลิยวกํ สนฺนิจิตํ โหติ อพฺภนฺตรานํ รติยา อปริตสฺสาย ผาสุวิหาราย พาหิรานํ ปฏิฆาตาย. (2)}

\prose{ปุน จปรํ ภิกฺขเว รญฺโญ ปจฺจนฺติเม นคเร พหุํ ติลมุคฺคมาสาปรณฺณํ สนฺนิจิตํ โหติ อพฺภนฺตรานํ รติยา อปริตสฺสาย ผาสุวิหาราย พาหิรานํ ปฏิฆาตาย. (3)}

\prose{ปุน จปรํ ภิกฺขเว รญฺโญ ปจฺจนฺติเม นคเร พหุํ เภสชฺชํ สนฺนิจิตํ โหติ.\csromanspacer{} เสยฺยถิทํ, สปฺปิ นวนีตํ เตลํ มธุ ผาณิตํ โลณํ, อพฺภนฺตรานํ รติยา อปริตสฺสาย ผาสุวิหาราย พาหิรานํ ปฏิฆาตาย.\csromanspacer{} อิเมสํ โข ภิกฺขเว จตุนฺนํ อาหารานํ นิกามลาภี โหติ อกิจฺฉลาภี อกสิรลาภี. (4)}

\prose{ยโต โข ภิกฺขเว รญฺโญ ปจฺจนฺติมํ นครํ อิเมหิ สตฺตหิ นครปริกฺขาเรหิ สุปริกฺขตํ โหติ, อิเมสญฺจ จตุนฺนํ อาหารานํ นิกามลาภี โหติ อกิจฺฉลาภี อกสิรลาภี.\csromanspacer{} อิทํ วุจฺจติ ภิกฺขเว รญฺโญ ปจฺจนฺติมํ นครํ อกรณียํ พาหิเรหิ ปจฺจตฺถิเกหิ ปจฺจามิตฺเตหิ.\csromanspacer{} เอวเมวํ โข ภิกฺขเว ยโต อริยสาวโก สตฺตหิ สทฺธมฺเมหิ สมนฺนาคโต โหติ จตุนฺนญฺจ ฌานานํ อาภิเจตสิกานํ ทิฏฺฐธมฺมสุขวิหารานํ นิกามลาภี}

\vspace{\tipitakaparskip}
\csromancenter{สตฺตกนิปาตปาฬิ โหติ อกิจฺฉลาภี อกสิรลาภี.\csromanspacer{} อยํ วุจฺจติ ภิกฺขเว อริยสาวโก อกรณีโย มารสฺส อกรณีโย ปาปิมโต.\csromanspacer{} กตเมหิ สตฺตหิ สทฺธมฺเมหิ สมนฺนาคโต โหติ?}
\prose{\csromanfolio{480}เสยฺยถาปิ ภิกฺขเว รญฺโญ ปจฺจนฺติเม นคเร เอสิกา โหติ คมฺภีรเนมา สุนิขาตา อจลา อสมฺปเวธี อพฺภนฺตรานํ คุตฺติยา พาหิรานํ ปฏิฆาตาย.\csromanspacer{} เอวเมวํ โข ภิกฺขเว อริยสาวโก สทฺโธ โหติ, สทฺทหติ ตถาคตสฺส โพธิํ “อิติปิ โส ฯเปฯ พุทฺโธ ภควา”ติ, สทฺเธสิโก ภิกฺขเว อริยสาวโก อกุสลํ ปชหติ, กุสลํ ภาเวติ, สาวชฺชํ ปชหติ, อนวชฺชํ ภาเวติ, สุทฺธํ อตฺตานํ ปริหรติ.\csromanspacer{} อิมินา ปฐเมน สทฺธมฺเมน สมนฺนาคโต โหติ. (1)}

\prose{เสยฺยถาปิ ภิกฺขเว รญฺโญ ปจฺจนฺติเม นคเร ปริกฺขา โหติ คมฺภีรา เจว วิตฺถตา จ อพฺภนฺตรานํ คุตฺติยา พาหิรานํ ปฏิฆาตาย.\csromanspacer{} เอวเมวํ โข ภิกฺขเว อริยสาวโก หิรีมา โหติ, หิรียติ กายทุจฺจริเตน วจีทุจฺจริเตน มโนทุจฺจริเตน, หิรียติ ปาปกานํ อกุสลานํ ธมฺมานํ สมาปตฺติยา.\csromanspacer{} หิรีปริกฺโข โข ภิกฺขเว อริยสาวโก อกุสลํ ปชหติ, กุสลํ ภาเวติ, สาวชฺชํ ปชหติ, อนวชฺชํ ภาเวติ, สุทฺธํ อตฺตานํ ปริหรติ.\csromanspacer{} อิมินา ทุติเยน สทฺธมฺเมน สมนฺนาคโต โหติ. (2)}

\prose{เสยฺยถาปิ ภิกฺขเว รญฺโญ ปจฺจนฺติเม นคเร อนุปริยายปโถ โหติ อุจฺโจ เจว วิตฺถโต จ อพฺภนฺตรานํ คุตฺติยา พาหิรานํ ปฏิฆาตาย.\csromanspacer{} เอวเมว โข ภิกฺขเว อริยสาวโก โอตฺตปฺปี โหติ, โอตฺตปฺปติ กายทุจฺจริเตน วจีทุจฺจริเตน มโนทุจฺจริเตน, โอตฺตปฺปติ ปาปกานํ อกุสลานํ ธมฺมานํ สมาปตฺติยา.\csromanspacer{} โอตฺตปฺปปริยายปโถ ภิกฺขเว อริยสาวโก อกุสลํ ปชหติ, กุสลํ ภาเวติ, สาวชฺชํ ปชหติ, อนวชฺชํ ภาเวติ, สุทฺธํ อตฺตานํ ปริหรติ.\csromanspacer{} อิมินา ตติเยน สทฺธมฺเมน สมนฺนาคโต โหติ. (3)}

\prose{เสยฺยถาปิ ภิกฺขเว รญฺโญ ปจฺจนฺติเม นคเร พหุํ อาวุธํ สนฺนิจิตํ โหติ สลากญฺเจว เชวนิกญฺจ อพฺภนฺตรานํ คุตฺติยา พาหิรานํ ปฏิฆาตาย.\csromanspacer{} เอวเมวํ โข ภิกฺขเว อริยสาวโก พหุสฺสุโต โหติ ฯเปฯ ทิฏฺฐิยา สุปฺปฏิวิทฺธา.\csromanspacer{} สุตาวุโธ ภิกฺขเว อริยสาวโก}

\vspace{\tipitakaparskip}
\csromancenter{มหาวคฺค อกุสลํ ปชหติ, กุสลํ ภาเวติ, สาวชฺชํ ปชหติ, อนวชฺชํ ภาเวติ, สุทฺธํ อตฺตานํ ปริหรติ.\csromanspacer{} อิมินา จตุตฺเถน สทฺธมฺเมน สมนฺนาคโต โหติ. (4)}
\prose{\csromanfolio{481}เสยฺยถาปิ ภิกฺขเว รญฺโญ ปจฺจนฺติเม นคเร พหุพลกาโย ปฏิวสติ.\csromanspacer{} เสยฺยถิทํ, หตฺถาโรหา อสฺสาโรหา รถิกา ธนุคฺคหา เจลกา จลกา ปิณฺฑทายกา อุคฺคา ราชปุตฺตา ปกฺขนฺทิโน มหานาคา สูรา จมฺมโยธิโน ทาสกปุตฺตา อพฺภนฺตรานํ คุตฺติยา พาหิรานํ ปฏิฆาตาย.\csromanspacer{} เอวเมวํ โข ภิกฺขเว อริยสาวโก อารทฺธวีริโย วิหรติ อกุสลานํ ธมฺมานํ ปหานาย กุสลานํ ธมฺมานํ อุปสมฺปทาย ถามวา ทฬฺหปรกฺกโม อนิกฺขิตฺตธุโร กุสเลสุ ธมฺเมสุ.\csromanspacer{} วีริยพลกาโย ภิกฺขเว อริยสาวโก อกุสลํ ปชหติ, กุสลํ ภาเวติ, สาวชฺชํ ปชหติ, อนวชฺชํ ภาเวติ, สุทฺธํ อตฺตานํ ปริหรติ.\csromanspacer{} อิมินา ปญฺจเมน สทฺธมฺเมน สมนฺนาคโต โหติ. (5)}

\prose{เสยฺยถาปิ ภิกฺขเว รญฺโญ ปจฺจนฺติเม นคเร โทวาริโก โหติ ปณฺฑิโต พฺยตฺโต เมธาวี อญฺญาตานํ นิวาเรตา ญาตานํ ปเวเสตา อพฺภนฺตรานํ คุตฺติยา พาหิรานํ ปฏิฆาตาย.\csromanspacer{} เอวเมวํ โข ภิกฺขเว อริยสาวโก สติมา โหติ, ปรเมน สติเนปกฺเกน สมนฺนาคโต จิรกตมฺปิ จิรภาสิตมฺปิ สริตา อนุสฺสริตา.\csromanspacer{} สติโทวาริโก ภิกฺขเว อริยสาวโก อกุสลํ ปชหติ, กุสลํ ภาเวติ, สาวชฺชํ ปชหติ, อนวชฺชํ ภาเวติ, สุทฺธํ อตฺตานํ ปริหรติ.\csromanspacer{} อิมินา ฉฏฺเฐน สทฺธมฺเมน สมนฺนาคโต โหติ. (6)}

\prose{เสยฺยถาปิ ภิกฺขเว รญฺโญ ปจฺจนฺติเม นคเร ปากาโร โหติ อุจฺโจ เจว วิตฺถโต จ วาสนเลปนสมฺปนฺโน จ อพฺภนฺตรานํ คุตฺติยา พาหิรานํ ปฏิฆาตาย.\csromanspacer{} เอวเมวํ โข ภิกฺขเว อริยสาวโก ปญฺญวา โหติ อุทยตฺถคามินิยา ปญฺญาย สมนฺนาคโต อริยาย นิพฺเพธิกาย สมฺมา ทุกฺขกฺขยคามินิยา, ปญฺญาวาสนเลปนสมฺปนฺโน ภิกฺขเว อริยสาวโก อกุสลํ ปชหติ, กุสลํ ภาเวติ, สาวชฺชํ ปชหติ, อนวชฺชํ ภาเวติ, สุทฺธํ อตฺตานํ ปริหรติ.\csromanspacer{} อิมินา สตฺตเมน สทฺธมฺเมน สมนฺนาคโต โหติ.\csromanspacer{} อิเมหิ สตฺตหิ สทฺธมฺเมหิ สมนฺนาคโต โหติ. (7)}

\gambhira{สตฺตกนิปาตปาฬิ}
\prose{\csromanfolio{482}กตเมสํ จตุนฺนํ ฌานานํ อาภิเจตสิกานํ ทิฏฺฐธมฺมสุข วิหารานํ นิกามลาภี โหติ อกิจฺฉลาภี อกสิรลาภี.\csromanspacer{} เสยฺยถาปิ ภิกฺขเว รญฺโญ ปจฺจนฺติเม นคเร พหุํ ติณกฏฺโฐทกํ สนฺนิจิตํ โหติ อพฺภนฺตรานํ รติยา อปริตสฺสาย ผาสุวิหาราย พาหิรานํ ปฏิฆาตาย.\csromanspacer{} เอวเมวํ โข ภิกฺขเว อริยสาวโก วิวิจฺเจว กาเมหิ ฯเปฯ ปฐมํ ฌานํ อุปสมฺปชฺช วิหรติ อตฺตโน รติยา อปริตสฺสาย ผาสุวิหาราย โอกฺกมนาย นิพฺพานสฺส. (1)}

\prose{เสยฺยถาปิ ภิกฺขเว รญฺโญ ปจฺจนฺติเม นคเร พหุํ สาลิ ยวกํ สนฺนิจิตํ โหติ อพฺภนฺตรานํ รติยา อปริตสฺสาย ผาสุวิหาราย พาหิรานํ ปฏิฆาตาย.\csromanspacer{} เอวเมวํ โข ภิกฺขเว อริยสาวโก วิตกฺกวิจารานํ วูปสมา ฯเปฯ ทุติยํ ฌานํ อุปสมฺปชฺช วิหรติ อตฺตโน รติยา อปริตสฺสาย ผาสุวิหาราย โอกฺกมนาย นิพฺพานสฺส. (2)}

\prose{เสยฺยถาปิ ภิกฺขเว รญฺโญ ปจฺจนฺติเม นคเร พหุํ ติลมุคฺคมาสาปรณฺณํ สนฺนิจิตํ โหติ อพฺภนฺตรานํ รติยา อปริตสฺสาย ผาสุวิหาราย พาหิรานํ ปฏิฆาตาย.\csromanspacer{} เอวเมวํ โข ภิกฺขเว อริยสาวโก ปีติยา จ วิราคา ฯเปฯ ตติยํ ฌานํ อุปสมฺปชฺช วิหรติ อตฺตโน รติยา อปริตสฺสาย ผาสุวิหาราย โอกฺกมนาย นิพฺพานสฺส. (3)}

\prose{เสยฺยถาปิ ภิกฺขเว รญฺโญ ปจฺจนฺติเม นคเร พหุํ เภสชฺชํ สนฺนิจิตํ โหติ.\csromanspacer{} เสยฺยถิทํ, สปฺปิ นวนีตํ เตลํ มธุ ผาณิตํ โลณํ, อพฺภนฺตรานํ รติยา อปริตสฺสาย ผาสุวิหาราย พาหิรานํ ปฏิฆาตาย.\csromanspacer{} เอวเมวํ โข ภิกฺขเว อริยสาวโก สุขสฺส จ ปหานา ทุกฺขสฺส จ ปหานา ปุพฺเพว โสมนสฺสโทมนสฺสานํ อตฺถงฺคมา อทุกฺขมสุขํ อุเปกฺขาสติปาริสุทฺธิํ จตุตฺถํ ฌานํ อุปสมฺปชฺช วิหรติ อตฺตโน รติยา อปริตสฺสาย ผาสุวิหาราย โอกฺกมนาย นิพฺพานสฺส.\csromanspacer{} อิเมสํ จตุนฺนํ ฌานานํ อาภิเจตสิกานํ ทิฏฺฐธมฺมสุขวิหารานํ นิกามลาภี โหติ อกิจฺฉลาภี อกสิรลาภี. (4)}

\prose{ยโต โข ภิกฺขเว อริยสาวโก อิเมหิ สตฺตหิ สทฺธมฺเมหิ สมนฺนาคโต โหติ, อิเมสญฺจ จตุนฺนํ ฌานานํ อาภิเจตสิกานํ ทิฏฺฐธมฺเมสุขวิหารานํ นิกามลาภี โหติ อกิจฺฉลาภี อกสิรลาภี.}

\vspace{\tipitakaparskip}
\csromancenter{มหาวคฺค อยํ วุจฺจติ ภิกฺขเว อริยสาวโก อกรณีโย มารสฺส อกรณีโย ปาปิมโตติ.\csromanspacer{}\csromanspacer{}ตติยํ.}
\csromanmatikaanchor{mtk.264}
\csromantocmarkref{subsection}{4. ธมฺมญฺญูสุตฺต}{mtk.264}
\bookmark[dest={mtk.264},level=2]{4. ธมฺมญฺญูสุตฺต}
\csromanheader{2}{4. ธมฺมญฺญูสุตฺต}
\proseitem{68}{\csromanfolio{483}สตฺตหิ ภิกฺขเว ธมฺเมหิ สมนฺนาคโต ภิกฺขุ อาหุเนยฺโย โหติ ฯเปฯ อนุตฺตรํ ปุญฺญกฺเขตฺตํ โลกสฺส.\csromanspacer{} กตเมหิ สตฺตหิ? อิธ ภิกฺขเว ภิกฺขุ ธมฺมญฺญู จ โหติ อตฺถญฺญู จ อตฺตญฺญู จ มตฺตญฺญู จ กาลญฺญู จ ปริสญฺญู จ ปุคฺคลปโรปรญฺญู จ.}

\prose{กถญฺจ ภิกฺขเว ภิกฺขุ ธมฺมญฺญู โหติ? อิธ ภิกฺขเว ภิกฺขุ ธมฺมํ ชานาติ สุตฺตํ เคยฺยํ เวยฺยากรณํ คาถํ อุทานํ อิติวุตฺตกํ ชาตกํ อพฺภุตธมฺมํ เวทลฺลํ.\csromanspacer{} โน เจ ภิกฺขุ ธมฺมํ ชาเนยฺย สุตฺตํ เคยฺยํ ฯเปฯ อพฺภุตธมฺมํ เวทลฺลํ, นยิธ “ธมฺมญฺญู”ติ วุจฺเจยฺย.\csromanspacer{} ยสฺมา จ โข ภิกฺขเว ภิกฺขุ ธมฺมํ ชานาติ สุตฺตํ เคยฺยํ ฯเปฯ อพฺภุตธมฺมํ เวทลฺลํ, ตสฺมา “ธมฺมญฺญู”ติ วุจฺจติ.\csromanspacer{} อิติ ธมฺมญฺญู. (1)}

\prose{อตฺถญฺญู จ กถํ โหติ? อิธ ภิกฺขเว ภิกฺขุ ตสฺส ตสฺเสว ภาสิตสฺส อตฺถํ ชานาติ “อยํ อิมสฺส ภาสิตสฺส อตฺโถ, อยํ อิมสฺส ภาสิตสฺส อตฺโถ”ติ.\csromanspacer{} โน เจ ภิกฺขเว ภิกฺขุ ตสฺส ตสฺเสว ภาสิตสฺส อตฺถํ ชาเนยฺย “อยํ อิมสฺส ภาสิตสฺส อตฺโถ, อยํ อิมสฺส ภาสิตสฺส อตฺโถ”ติ, นยิธ “อตฺถญฺญู”ติ วุจฺเจยฺย.\csromanspacer{} ยสฺมา จ โข ภิกฺขเว ภิกฺขุ ตสฺส ตสฺเสว ภาสิตสฺส อตฺถํ ชานาติ “อยํ อิมสฺส ภาสิตสฺส อตฺโถ, อยํ อิมสฺส ภาสิตสฺส อตฺโถ”ติ, ตสฺมา “อตฺถญฺญู”ติ วุจฺจติ.\csromanspacer{} อิติ ธมฺมญฺญู อตฺถญฺญู. (2)}

\prose{อตฺตญฺญู จ กถํ โหติ? อิธ ภิกฺขเว ภิกฺขุ อตฺตานํ ชานาติ “เอตฺตโกมฺหิ สทฺธาย สีเลน สุเตน จาเคน ปญฺญาย ปฏิภาเนนา”ติ.\csromanspacer{} โน เจ ภิกฺขเว ภิกฺขุ อตฺตานํ ชาเนยฺย “เอตฺตโกมฺหิ สทฺธาย สีเลน สุเตน จาเคน ปญฺญาย ปฏิภาเนนา”ติ, นยิธ “อตฺตญฺญู”ติ วุจฺเจยฺย.\csromanspacer{} ยสฺมา จ ภิกฺขเว ภิกฺขุ อตฺตานํ ชานาติ “เอตฺตโกมฺหิ สทฺธาย สีเลน สุเตน จาเคน ปญฺญาย ปฏิภาเนนา”ติ, ตสฺมา “อตฺตญฺญู”ติ วุจฺจติ.\csromanspacer{} อิติ ธมฺมญฺญู อตฺถญฺญู อตฺตญฺญู. (3)}

\gambhira{สตฺตกนิปาตปาฬิ}
\prose{\csromanfolio{484}มตฺตญฺญู จ กถํ โหติ? อิธ ภิกฺขเว ภิกฺขุ มตฺตํ ชานาติ จีวรปิณฺฑปาตเสนาสนคิลานปฺปจฺจยเภสชฺชปริกฺขารานํ ปฏิคฺคหณาย.\csromanspacer{} โน เจ ภิกฺขเว ภิกฺขุ มตฺตํ ชาเนยฺย จีวรปิณฺฑปาตเสนาสนคิลานปฺปจฺจย เภสชฺชปริกฺขารานํ ปฏิคฺคหณาย, นยิธ “มตฺตญฺญู”ติ วุจฺเจยฺย.\csromanspacer{} ยสฺมา จ โข ภิกฺขเว ภิกฺขุ มตฺตํ ชานาติ จีวรปิณฺฑปาตเสนาสนคิลานปฺปจฺจย เภสชฺชปริกฺขารานํ ปฏิคฺคหณาย, ตสฺมา “มตฺตญฺญู”ติ วุจฺจติ.\csromanspacer{} อิติ ธมฺมญฺญู อตฺถญฺญู อตฺตญฺญู มตฺตญฺญู. (4)}

\prose{กาลญฺญู จ กถํ โหติ? อิธ ภิกฺขเว ภิกฺขุ กาลํ ชานาติ “อยํ กาโล อุทฺเทสสฺส, อยํ กาโล ปริปุจฺฉาย, อยํ กาโล โยคสฺส, อยํ กาโล ปฏิสลฺลานสฺสา”ติ.\csromanspacer{} โน เจ ภิกฺขเว ภิกฺขุ กาลํ ชาเนยฺย “อยํ กาโล อุทฺเทสสฺส, อยํ กาโล ปริปุจฺฉาย, อยํ กาโล โยคสฺส, อยํ กาโล ปฏิสลฺลานสฺสา”ติ, นยิธ “กาลญฺญู”ติ วุจฺเจยฺย.\csromanspacer{} ยสฺมา จ โข ภิกฺขเว ภิกฺขุ กาลํ ชานาติ “อยํ กาโล อุทฺเทสสฺส, อยํ กาโล ปริปุจฺฉาย, อยํ กาโล โยคสฺส, อยํ กาโล ปฏิสลฺลานสฺสา”ติ, ตสฺมา “กาลญฺญู”ติ วุจฺจติ.\csromanspacer{} อิติ ธมฺมญฺญู อตฺถญฺญู อตฺตญฺญู มตฺตญฺญู กาลญฺญู. (5)}

\prose{ปริสญฺญู จ กถํ โหติ? อิธ ภิกฺขเว ภิกฺขุ ปริสํ ชานาติ “อยํ ขตฺติยปริสา, อยํ พฺราหฺมณปริสา, อยํ คหปติปริสา, อยํ สมณปริสา.\csromanspacer{} ตตฺถ เอวํ อุปสงฺกมิตพฺพํ, เอวํ ฐาตพฺพํ, เอวํ กตฺตพฺพํ, เอวํ นิสีทิตพฺพํ, เอวํ ภาสิตพฺพํ, เอวํ ตุณฺหี ภวิตพฺพนฺ”ติ.\csromanspacer{} โน เจ ภิกฺขเว ภิกฺขุ ปริสํ ชาเนยฺย “อยํ ขตฺติยปริสา ฯเปฯ เอวํ ตุณฺหี ภวิตพฺพนฺ”ติ, นยิธ “ปริสญฺญู”ติ วุจฺเจยฺย.\csromanspacer{} ยสฺมา จ โข ภิกฺขเว ภิกฺขุ ปริสํ ชานาติ “อยํ ขตฺติยปริสา, อยํ พฺราหฺมณปริสา, อยํ คหปติปริสา, อยํ สมณปริสา.\csromanspacer{} ตตฺถ เอวํ อุปสงฺกมิตพฺพํ, เอวํ ฐาตพฺพํ, เอวํ กตฺตพฺพํ, เอวํ นิสีทิตพฺพํ, เอวํ ภาสิตพฺพํ, เอวํ ตุณฺหี ภวิตพฺพนฺ”ติ, ตสฺมา “ปริสญฺญู”ติ วุจฺจติ.\csromanspacer{} อิติ ธมฺมญฺญู อตฺถญฺญู อตฺตญฺญู มตฺตญฺญู กาลญฺญู ปริสญฺญู. (6)}

\prose{ปุคฺคลปโรปรญฺญู จ กถํ โหติ? อิธ ภิกฺขเว ภิกฺขุโน\footnote{สพฺพตฺถปิ อิธ อารมฺเภ “ภิกฺขุโน”เตฺวว ทิสฺสติ.} ทฺวเยน ปุคฺคลา วิทิตา โหนฺติ.\csromanspacer{} เทฺว ปุคฺคลา เอโก อริยานํ ทสฺสนกาโม, เอโก อริยานํ น ทสฺสนกาโม.\csromanspacer{} ยฺวายํ ปุคฺคโล อริยานํ น}

\vspace{\tipitakaparskip}
\csromancenter{มหาวคฺค ทสฺสนกาโม, เอวํ โส เตนงฺเคน คารโยฺห.\csromanspacer{} ยฺวายํ ปุคฺคโล อริยานํ ทสฺสนกาโม, เอวํ โส เตนงฺเคน ปาสํโส.}
\prose{\csromanfolio{485}เทฺว ปุคฺคลา อริยานํ ทสฺสนกาโม, เอโก สทฺธมฺมํ โสตุกาโม, เอโก สทฺธมฺมํ น โสตุกาโม.\csromanspacer{} ยฺวายํ ปุคฺคโล สทฺธมฺมํ น โสตุกาโม, เอวํ โส เตนงฺเคน คารโยฺห.\csromanspacer{} ยฺวายํ ปุคฺคโล สทฺธมฺมํ โสตุกาโม, เอวํ โส เตนงฺเคน ปาสํโส.}

\prose{เทฺว ปุคฺคลา สทฺธมฺมํ โสตุกามา, เอโก โอหิตโสโต ธมฺมํ สุณาติ, เอโก อโนหิตโสโต ธมฺมํ สุณาติ.\csromanspacer{} ยฺวายํ ปุคฺคโล อโนหิตโสโต ธมฺมํ สุณาติ, เอวํ โส เตนงฺเคน คารโยฺห.\csromanspacer{} ยฺวายํ ปุคฺคโล โอหิตโสโต ธมฺมํ สุณาติ, เอวํ โส เตนงฺเคน ปาสํโส.}

\prose{เทฺว ปุคฺคลา โอหิตโสตา ธมฺมํ สุณนฺติ\footnote{สุณนฺตา (ก)}, เอโก สุตฺวา ธมฺมํ ธาเรติ, เอโก สุตฺวา ธมฺมํ น ธาเรติ.\csromanspacer{} ยฺวายํ ปุคฺคโล สุตฺวา น ธมฺมํ ธาเรติ, เอวํ โส เตนงฺเคน คารโยฺห.\csromanspacer{} ยฺวายํ ปุคฺคโล สุตฺวา ธมฺมํ ธาเรติ, เอวํ โส เตนงฺเคน ปาสํโส.}

\prose{เทฺว ปุคฺคลา สุตฺวา ธมฺมํ ธาเรนฺติ\footnote{ธาเรนฺตา (ก)}, เอโก ธาตานํ\footnote{ธตานํ (สี, สฺยา, กํ, อิ)} ธมฺมานํ อตฺถํ อุปปริกฺขติ, เอโก ธาตานํ ธมฺมานํ อตฺถํ น อุปปริกฺขติ.\csromanspacer{} ยฺวายํ ปุคฺคโล ธาตานํ ธมฺมานํ อตฺถํ น อุปปริกฺขติ, เอวํ โส เตนงฺเคน คารโยฺห.\csromanspacer{} ยฺวายํ ปุคฺคโล ธาตานํ ธมฺมานํ อตฺถํ อุปปริกฺขติ, เอวํ โส เตนงฺเคน ปาสํโส.}

\prose{เทฺว ปุคฺคลา ธาตานํ ธมฺมานํ อตฺถํ อุปปริกฺขนฺติ\footnote{อุปปริกฺขนฺตา (ก)}, เอโก อตฺถมญฺญาย ธมฺมมญฺญาย ธมฺมานุธมฺมปฺปฏิปนฺโน, เอโก อตฺถมญฺญาย ธมฺมมญฺญาย น ธมฺมานุธมฺมปฺปฏิปนฺโน.\csromanspacer{} ยฺวายํ ปุคฺคโล อตฺถมญฺญาย ธมฺมมญฺญาย น ธมฺมานุธมฺมปฺปฏิปนฺโน, เอวํ โส เตนงฺเคน คารโยฺห.\csromanspacer{} ยฺวายํ ปุคฺคโล อตฺถมญฺญาย ธมฺมมญฺญาย ธมฺมานุธมฺมปฺปฏิปนฺโน, เอวํ โส เตนงฺเคน ปาสํโส.}

\prose{เทฺว ปุคฺคลา อตฺถมญฺญาย ธมฺมมญฺญาย ธมฺมานุธมฺมปฺปฏิ- ปนฺนา, เอโก อตฺตหิตาย ปฏิปนฺโน โน ปรหิตาย, เอโก อตฺตหิตาย จ}

\vspace{\tipitakaparskip}
\csromancenter{สตฺตกนิปาตปาฬิ ปฏิปนฺโน ปรหิตาย จ.\csromanspacer{} ยฺวายํ ปุคฺคโล อตฺตหิตาย ปฏิปนฺโน โน ปรหิตาย, เอวํ โส เตนงฺเคน คารโยฺห.\csromanspacer{} ยฺวายํ ปุคฺคโล อตฺตหิตาย จ ปฏิปนฺโน ปรหิตาย จ, เอวํ โส เตนงฺเคน ปาสํโส.\csromanspacer{} เอวํ โข ภิกฺขเว ภิกฺขุโน\footnote{ภิกฺขุนา (สี, สฺยา)} ทฺวเยน ปุคฺคลา วิทิตา โหนฺติ.\csromanspacer{} เอวํ ภิกฺขเว ภิกฺขุ ปุคฺคลปโรปรญฺญู โหติ.\csromanspacer{} อิเมหิ โข ภิกฺขเว สตฺตหิ ธมฺเมหิ สมนฺนาคโต ภิกฺขุ อาหุเนยฺโย โหติ ปาหุเนยฺโย ฯเปฯ อนุตฺตรํ ปุญฺญกฺเขตฺตํ โลกสฺสาติ.\csromanspacer{}\csromanspacer{}จตุตฺถํ.}
\csromanmatikaanchor{mtk.265}
\csromantocmarkref{subsection}{5. ปาริจฺฉตฺตกสุตฺต}{mtk.265}
\bookmark[dest={mtk.265},level=2]{5. ปาริจฺฉตฺตกสุตฺต}
\csromanheader{2}{5. ปาริจฺฉตฺตกสุตฺต}
\proseitem{69}{\csromanfolio{486}ยสฺมิํ ภิกฺขเว สมเย เทวานํ ตาวติํสานํ ปาริจฺฉตฺตโก โกวิฬาโร ปณฺฑุปลาโส โหติ, อตฺตมนา ภิกฺขเว เทวา ตาวติํสา ตสฺมิํ สมเย โหนฺติ “ปณฺฑุปลาโส ทานิ ปาริจฺฉตฺตโก โกวิฬาโร, นจิรสฺเสว ทานิ ปนฺนปลาโส ภวิสฺสตี”ติ. (1)}

\prose{ยสฺมิํ ภิกฺขเว สมเย เทวานํ ตาวติํสานํ ปาริจฺฉตฺตโก โกวิฬาโร ปนฺนปลาโส โหติ, อตฺตมนา ภิกฺขเว เทวา ตาวติํสา ตสฺมิํ สมเย โหนฺติ “ปนฺนปลาโส ทานิ ปาริจฺฉตฺตโก โกวิฬาโร, นจิรสฺเสว ทานิ ชาลกชาโต ภวิสฺสตี”ติ. (2)}

\prose{ยสฺมิํ ภิกฺขเว สมเย เทวานํ ตาวติํสานํ ปาริจฺฉตฺตโก โกวิฬาโร ชาลกชาโต โหติ, อตฺตมนา ภิกฺขเว เทวา ตาวติํสา ตสฺมิํ สมเย โหนฺติ “ชาลกชาโต ทานิ ปาริจฺฉตฺตโก โกวิฬาโร, นจิรสฺเสว ทานิ ขารกชาโต ภวิสฺสตี”ติ. (3)}

\prose{ยสฺมิํ ภิกฺขเว สมเย เทวานํ ตาวติํสานํ ปาริจฺฉตฺตโก โกวิฬาโร ขารกชาโต โหติ, อตฺตมนา ภิกฺขเว เทวา ตาวติํสา ตสฺมิํ สมเย โหนฺติ “ขารกชาโต ทานิ ปาริจฺฉตฺตโก โกวิฬาโร, นจิรสฺเสว ทานิ กุฏุมลกชาโต\footnote{กุฑุมลกชาโต (สี, สฺยา, อิ)} ภวิสฺสตี”ติ. (4)}

\prose{ยสฺมิํ ภิกฺขเว สมเย เทวานํ ตาวติํสานํ ปาริจฺฉตฺตโก โกวิฬาโร กุฏุมลกชาโต โหติ, อตฺตมนา ภิกฺขเว เทวา}

\vspace{\tipitakaparskip}
\csromancenter{มหาวคฺค ตาวติํสา ตสฺมิํ สมเย โหนฺติ “กุฏุมลกชาโต ทานิ ปาริจฺฉตฺตโก โกวิฬาโร, นจิรสฺเสว ทานิ โกรกชาโต\footnote{โกกาสกชาโต (สี, สฺยา), โกสกชาโต (ก)} ภวิสฺสตี”ติ. (5)}
\prose{\csromanfolio{487}ยสฺมิํ ภิกฺขเว สมเย เทวานํ ตาวติํสานํ ปาริจฺฉตฺตโก โกวิฬาโร โกรกชาโต โหติ, อตฺตมนา ภิกฺขเว เทวา ตาวติํสา ตสฺมิํ สมเย โหนฺติ “โกรกชาโต ทานิ ปาริจฺฉตฺตโก โกวิฬาโร, นจิรสฺเสว ทานิ สพฺพผาลิผุลฺโล\footnote{สพฺพปาลิผุลฺโล (สี, อิ)} ภวิสฺสตี”ติ. (6)}

\prose{ยสฺมิํ ภิกฺขเว สมเย เทวานํ ตาวติํสานํ ปาริจฺฉตฺตโก โกวิฬาโร สพฺพผาลิผุลฺโล โหติ, อตฺตมนา ภิกฺขเว เทวา ตาวติํสา ปาริฉตฺตกสฺส โกวิฬารสฺส มูเล ทิพฺเพ จตฺตาโร มาเส ปญฺจหิ กามคุเณหิ สมปฺปิตา สมงฺคีภูตา ปริจาเรนฺติ. (7)}

\prose{สพฺพผาลิผุลฺลสฺส โข ปน ภิกฺขเว ปาริจฺฉตฺตกสฺส โกวิฬารสฺส สมนฺตา ปญฺญาสโยชนานิ อาภาย ผุฏํ โหติ, อนุวาตํ โยชนสตํ คนฺโธ คจฺฉติ, อยมานุภาโว ปาริจฺฉตฺตกสฺส โกวิฬารสฺส.}

\prose{เอวเมวํ โข ภิกฺขเว ยสฺมิํ สมเย อริยสาวโก อคารสฺมา อนคาริยํ ปพฺพชฺชาย เจเตติ, ปณฺฑุปลาโส ภิกฺขเว อริยสาวโก ตสฺมิํ สมเย โหติ เทวานํว ตาวติํสานํ ปาริจฺฉตฺตโก โกวิฬาโร. (1)}

\prose{ยสฺมิํ ภิกฺขเว สมเย อริยสาวโก เกสมสฺสุํ โอหาเรตฺวา กาสายานิ วตฺถานิ อจฺฉาเทตฺวา อคารสฺมา อนคาริยํ ปพฺพชิโต โหติ, ปนฺนปลาโส ภิกฺขเว อริยสาวโก ตสฺมิํ สมเย โหติ เทวานํว ตาวติํสานํ ปาริจฺฉตฺตโก โกวิฬาโร. (2)}

\prose{ยสฺมิํ ภิกฺขเว สมเย อริยสาวโก วิวิจฺเจว กาเมหิ ฯเปฯ ปฐมํ ฌานํ อุปสมฺปชฺช วิหรติ, ชาลกชาโต ภิกฺขเว อริยสาวโก ตสฺมิํ สมเย โหติ เทวานํว ตาวติํสานํ ปาริจฺฉตฺตโก โกวิฬาโร. (3)}

\gambhira{สตฺตกนิปาตปาฬิ}
\prose{\csromanfolio{488}ยสฺมิํ ภิกฺขเว สมเย อริยสาวโก วิตกฺกวิจารานํ วูปสมา ฯเปฯ ทุติยํ ฌานํ อุปสมฺปชฺช วิหรติ, ขารกชาโต ภิกฺขเว อริยสาวโก ตสฺมิํ สมเย โหติ เทวานํว ตาวติํสานํ ปาริจฺฉตฺตโก โกวิฬาโร. (4)}

\prose{ยสฺมิํ ภิกฺขเว สมเย อริยสาวโก ปีติยา จ วิราคา ฯเปฯ ตติยํ ฌานํ อุปสมฺปชฺช วิหรติ, กุฏุมลกชาโต ภิกฺขเว อริยสาวโก ตสฺมิํ สมเย โหติ เทวานํว ตาวติํสานํ ปาริจฺฉตฺตโก โกวิฬาโร. (5)}

\prose{ยสฺมิํ ภิกฺขเว สมเย อริยสาวโก สุขสฺส จ ปหานา ทุกฺขสฺส จ ปหานา ฯเปฯ จตุตฺถํ ฌานํ อุปสมฺปชฺช วิหรติ, โกรกชาโต ภิกฺขเว อริยสาวโก ตสฺมิํ สมเย โหติ เทวานํว ตาวติํสานํ ปาริจฺฉตฺตโก โกวิฬาโร. (6)}

\prose{ยสฺมิํ ภิกฺขเว สมเย อริยสาวโก อาสวานํ ขยา ฯเปฯ สจฺฉิกตฺวา อุปสมฺปชฺช วิหรติ, สพฺพผาลิผุลฺโล ภิกฺขเว อริยสาวโก ตสฺมิํ สมเย โหติ เทวานํว ตาวติํสานํ ปาริจฺฉตฺตโก โกวิฬาโร. (7)}

\prose{ตสฺมิํ ภิกฺขเว สมเย ภุมฺมา เทวา สทฺทมนุสฺสาเวนฺติ “เอโส อิตฺถนฺนาโม อายสฺมา อิตฺถนฺนามสฺส อายสฺมโต สทฺธิวิหาริโก อมุกมฺหา คามา วา นิคมา วา อคารสฺมา อนคาริยํ ปพฺพชิโต อาสวานํ ขยา อนาสวํ เจโตวิมุตฺติํ ปญฺญาวิมุตฺติํ ทิฏฺเฐว ธมฺเม สยํ อภิญฺญา สจฺฉิกตฺวา อุปสมฺปชฺช วิหรตี”ติ.\csromanspacer{} ภุมฺมานํ เทวานํ สทฺทํ สุตฺวา จาตุมหาราชิกา เทวา ฯเปฯ ตาวติํสา เทวา.\csromanspacer{} ยามา เทวา.\csromanspacer{} ตุสิตา เทวา.\csromanspacer{} นิมฺมานรตี เทวา.\csromanspacer{} ปรนิมฺมิตวสวตฺตี เทวา.\csromanspacer{} พฺรหฺมกายิกา เทวา สทฺทมนุสฺสาเวนฺติ “เอโส อิตฺถนฺนาโม อายสฺมา อิตฺถนฺนามสฺส อายสฺมโต สทฺธิวิหาริโก อมุกมฺหา คามา วา นิคมา วา อคารสฺมา อนคาริยํ ปพฺพชิโต อาสวานํ ขยา อนาสวํ เจโตวิมุตฺติํ ปญฺญาวิมุตฺติํ ทิฏฺเฐว ธมฺเม สยํ อภิญฺญา สจฺฉิกตฺวา อุปสมฺปชฺช วิหรตี”ติ.\csromanspacer{} อิติห เตน ขเณน เตน มุหุตฺเตน ยาว พฺรหฺมโลกา สทฺโท\footnote{สาธุการสทฺโท (สี-ฏฺฐ, ก-ฏฺฐ)} อพฺภุคฺคจฺฉติ.\csromanspacer{} อยมานุภาโว ขีณาสวสฺส ภิกฺขุโนติ.\csromanspacer{}\csromanspacer{}ปญฺจมํ.}

\csromanmatikaanchor{mtk.266}
\csromantocmarkref{subsection}{6. สกฺกจฺจสุตฺต}{mtk.266}
\bookmark[dest={mtk.266},level=2]{6. สกฺกจฺจสุตฺต}
\csromanheader{2}{7. มหาวคฺค \textbf{6. สกฺกจฺจสุตฺต}}
\proseitem{70}{\csromanfolio{489}อถ โข อายสฺมโต สาริปุตฺตสฺส รโหคตสฺส ปฏิสลฺลีนสฺส เอวํ เจตโส ปริวิตกฺโก อุทปาทิ “กิํ นุ โข ภิกฺขุ สกฺกตฺวา ครุํ กตฺวา อุปนิสฺสาย วิหรนฺโต อกุสลํ ปชเหยฺย, กุสลํ ภาเวยฺยา”ติ.\csromanspacer{} อถ โข อายสฺมโต สาริปุตฺตสฺส เอตทโหสิ “สตฺถารํ โข ภิกฺขุ สกฺกตฺวา ครุํ กตฺวา อุปนิสฺสาย วิหรนฺโต อกุสลํ ปชเหยฺย, กุสลํ ภาเวยฺย.\csromanspacer{} ธมฺมํ โข ภิกฺขุ ฯเปฯ.\csromanspacer{} สํฆํ โข ภิกฺขุ ฯเปฯ.\csromanspacer{} สิกฺขํ โข ภิกฺขุ ฯเปฯ.\csromanspacer{} สมาธิํ โข ภิกฺขุ ฯเปฯ.\csromanspacer{} อปฺปมาทํ โข ภิกฺขุ ฯเปฯ.\csromanspacer{} ปฏิสนฺถารํ โข ภิกฺขุ สกฺกตฺวา ครุํ กตฺวา อุปนิสฺสาย วิหรนฺโต อกุสลํ ปชเหยฺย, กุสลํ ภาเวยฺยา”ติ.}

\prose{อถ โข อายสฺมโต สาริปุตฺตสฺส เอตทโหสิ “อิเม โข เม ธมฺมา ปริสุทฺธา ปริโยทาตา, ยํนูนาหํ อิเม ธมฺเม คนฺตฺวา\footnote{คเหตฺวา (ก)} ภควโต อาโรเจยฺยํ, เอวํ เม อิเม ธมฺมา ปริสุทฺธา เจว ภวิสฺสนฺติ ปริสุทฺธสงฺขาตตรา จ.\csromanspacer{} เสยฺยถาปิ นาม ปุริโส สุวณฺณนิกฺขํ อธิคจฺเฉยฺย ปริสุทฺธํ ปริโยทาตํ.\csromanspacer{} ตสฺส เอวมสฺส อยํ โข เม สุวณฺณนิกฺโข ปริสุทฺโธ ปริโยทาโต, ยํนูนาหํ อิมํ สุวณฺณนิกฺขํ คนฺตฺวา1 กมฺมารานํ ทสฺเสยฺยํ, เอวํ เม อยํ สุวณฺณนิกฺโข สกมฺมารคโต ปริสุทฺโธ เจว ภวิสฺสติ ปริสุทฺธสงฺขาตตโร จ.\csromanspacer{} เอวเมวํ เม\footnote{โข (ก)} อิเม ธมฺมา ปริสุทฺธา ปริโยทาตา, ยํนูนาหํ อิเม ธมฺเม คนฺตฺวา1 ภควโต อาโรเจยฺยํ, เอวํ เม อิเม ธมฺมา ปริสุทฺธา เจว ภวิสฺสนฺติ ปริสุทฺธสงฺขาตตรา จา”ติ.}

\prose{อถ โข อายสฺมา สาริปุตฺโต สายนฺหสมยํ ปฏิสลฺลานา วุฏฺฐิโต เยน ภควา เตนุปสงฺกมิ, อุปสงฺกมิตฺวา ภควนฺตํ อภิวาเทตฺวา เอกมนฺตํ นิสีทิ, เอกมนฺตํ นิสินฺโน โข อายสฺมา สาริปุตฺโต ภควนฺตํ เอตทโวจ\csromandash{}}

\prose{อิธ มยฺหํ ภนฺเต รโหคตสฺส ปฏิสลฺลีนสฺส เอวํ เจตโส ปริวิตกฺโก อุทปาทิ “กิํ นุ โข ภิกฺขุ สกฺกตฺวา ครุํ กตฺวา อุปนิสฺสาย วิหรนฺโต อกุสลํ ปชเหยฺย, กุสลํ ภาเวยฺยา”ติ.\csromanspacer{} อถ โข ตสฺส มยฺหํ ภนฺเต เอตทโหสิ “สตฺถารํ โข ภิกฺขุ สกฺกตฺวา}

\vspace{\tipitakaparskip}
\csromancenter{สตฺตกนิปาตปาฬิ ครุํ กตฺวา อุปนิสฺสาย วิหรนฺโต อกุสลํ ปชเหยฺย, กุสลํ ภาเวยฺย.\csromanspacer{} ธมฺมํ โข ภิกฺขุ ฯเปฯ.\csromanspacer{} ปฏิสนฺถารํ โข ภิกฺขุ สกฺกตฺวา ฯเปฯ กุสลํ ภาเวยฺยา”ติ.\csromanspacer{} อถ โข ตสฺส มยฺหํ ภนฺเต เอตทโหสิ “อิเม โข เม ธมฺมา ปริสุทฺธา ปริโยทาตา, ยํนูนาหํ อิเม ธมฺเม คนฺตฺวา ภควตา อาโรเจยฺยํ, เอวํ เม อิเม ธมฺมา ปริสุทฺธา เจว ภวิสฺสนฺติ ปริสุทฺธสงฺขาตตรา จ.\csromanspacer{} เสยฺยถาปิ นาม ปุริโส สุวณฺณนิกฺขํ อธิคจฺเฉยฺย ปริสุทฺธํ ปริโยทาตํ.\csromanspacer{} ตสฺส เอวมสฺส อยํ โข เม สุวณฺณนิกฺโข ปริสุทฺโธ ปริโยทาโต, ยํนูนาหํ อิมํ สุวณฺณนิกฺขํ คนฺตฺวา กมฺมารานํ ทสฺเสยฺยํ, เอวํ เม อยํ สุวณฺณนิกฺโข สกมฺมารคโต ปริสุทฺโธ เจว ภวิสฺสติ ปริสุทฺธสงฺขาตตโร จ.\csromanspacer{} เอวเมวํ เม อิเม ธมฺมา ปริสุทฺธา ปริโยทาตา, ยํนูนาหํ อิเม ธมฺเม คนฺตฺวา ภควโต อาโรเจยฺยํ, เอวํ เม อิเม ธมฺมา ปริสุทฺธา เจว ภวิสฺสนฺติ ปริสุทฺธสงฺขาตตรา จา”ติ.}
\prose{\csromanfolio{490}สาธุ สาธุ สาริปุตฺต, สตฺถารํ โข สาริปุตฺต ภิกฺขุ สกฺกตฺวา ครุํ กตฺวา อุปนิสฺสาย วิหรนฺโต อกุสลํ ปชเหยฺย, กุสลํ ภาเวยฺย.\csromanspacer{} ธมฺมํ โข สาริปุตฺต ภิกฺขุ สกฺกตฺวา ครุํ กตฺวา อุปนิสฺสาย วิหรนฺโต อกุสลํ ปชเหยฺย, กุสลํ ภาเวยฺย.\csromanspacer{} สํฆํ โข ฯเปฯ.\csromanspacer{} สิกฺขํ โข.\csromanspacer{} สมาธิํ โข.\csromanspacer{} อปฺปมาทํ โข.\csromanspacer{} ปฏิสนฺถารํ โข สาริปุตฺต ภิกฺขุ สกฺกตฺวา ครุํ กตฺวา อุปนิสฺสาย วิหรนฺโต อกุสลํ ปชเหยฺย, กุสลํ ภาเวยฺยาติ.}

\prose{เอวํ วุตฺเต อายสฺมา สาริปุตฺโต ภควนฺตํ เอตทโวจ “อิมสฺส โข อหํ ภนฺเต ภควตา สํขิตฺเตน ภาสิตสฺส เอวํ วิตฺถาเรน อตฺถํ อาชานามิ.\csromanspacer{} โส วต ภนฺเต ภิกฺขุ สตฺถริ อคารโว “ธมฺเม สคารโว ภวิสฺสตี”ติ เนตํ ฐานํ วิชฺชติ.\csromanspacer{} โย โส ภนฺเต ภิกฺขุ สตฺถริ อคารโว, ธมฺเมปิ โส อคารโว. (1-2) โส วต ภนฺเต ภิกฺขุ สตฺถริ อคารโว ธมฺเม อคารโว “สํเฆ สคารโว ภวิสฺสตี”ติ เนตํ ฐานํ วิชฺชติ.\csromanspacer{} โย โส ภนฺเต ภิกฺขุ สตฺถริ อคารโว ธมฺเม อคารโว, สํเฆปิ โส อคารโว. (3) โส วต ภนฺเต ภิกฺขุ สตฺถริ อคารโว ธมฺเม อคารโว สํเฆ อคารโว “สิกฺขาย สคารโว ภวิสฺสตี”ติ เนตํ ฐานํ วิชฺชติ.\csromanspacer{} โย โส ภนฺเต ภิกฺขุ สตฺถริ อคารโว ธมฺเม อคารโว สํเฆ อคารโว, สิกฺขายปิ โส อคารโว. (4)}

\chapterheadpageread{7. มหาวคฺค}
\prose{\csromanfolio{491}โส วต ภนฺเต ภิกฺขุ สตฺถริ อคารโว ธมฺเม อคารโว สํเฆ อคารโว สิกฺขาย อคารโว “สมาธิสฺมิํ สคารโว ภวิสฺสตี”ติ เนตํ ฐานํ วิชฺชติ.\csromanspacer{} โย โส ภนฺเต ภิกฺขุ สตฺถริ อคารโว ธมฺเม อคารโว สํเฆ อคารโว สิกฺขาย อคารโวสมาธิสฺมิมฺปิ โส อคารโว. (5)}

\prose{โส วต ภนฺเต ภิกฺขุ สตฺถริ อคารโว ธมฺเม อคารโว สํเฆ อคารโว สิกฺขาย อคารโว สมาธิสฺมิํ อคารโว “อปฺปมาเท สคารโว ภวิสฺสตี”ติ เนตํ ฐานํ วิชฺชติ.\csromanspacer{} โย โส ภนฺเต ภิกฺขุ สตฺถริ อคารโว ธมฺเม อคารโว สํเฆ อคารโว สิกฺขาย อคารโว สมาธิสฺมิํ อคารโว, อปฺปมาเทปิ โส อคารโว. (6)}

\prose{โส วต ภนฺเต ภิกฺขุ สตฺถริ อคารโว ธมฺเม อคารโว สํเฆ อคารโว สิกฺขาย อคารโว สมาธิสฺมิํ อคารโว อปฺปมาเท อคารโว “ปฏิสนฺถาเร สคารโว ภวิสฺสตี”ติ เนตํ ฐานํ วิชฺชติ.\csromanspacer{} โย โส ภนฺเต ภิกฺขุ สตฺถริ อคารโว ฯเปฯ อปฺปมาเท อคารโว, ปฏิสนฺถาเรปิ โส อคารโว. (7)}

\prose{โส วต ภนฺเต ภิกฺขุ สตฺถริ สคารโว “ธมฺเม อคารโว ภวิสฺสตี”ติ เนตํ ฐานํ วิชฺชติ.\csromanspacer{} โย โส ภนฺเต ภิกฺขุ สตฺถริ สคารโว, ธมฺเมปิ โส สคารโว ฯเปฯ. (1-6)}

\prose{โส วต ภนฺเต ภิกฺขุ สตฺถริ สคารโว ฯเปฯ อปฺปมาเท สคารโว “ปฏิสนฺถาเร อคารโว ภวิสฺสตี”ติ เนตํ ฐานํ วิชฺชติ.\csromanspacer{} โย โส ภนฺเต ภิกฺขุ สตฺถริ สคารโว ฯเปฯ อปฺปมาเท สคารโว, ปฏิสนฺถาเรปิ โส สคารโว. (7)}

\prose{โส วต ภนฺเต ภิกฺขุ สตฺถริ สคารโว “ธมฺเมปิ สคารโว ภวิสฺสตี”ติ ฐานเมตํ วิชฺชติ.\csromanspacer{} โย โส ภนฺเต ภิกฺขุ สตฺถริ สคารโว, ธมฺเมปิ โส สคารโว ฯเปฯ. (1-6)}

\prose{โส วต ภนฺเต ภิกฺขุ สตฺถริ สคารโว ฯเปฯ อปฺปมาเท สคารโว, “ปฏิสนฺถาเรปิ สคารโว ภวิสฺสตี”ติ ฐานเมตํ วิชฺชติ.\csromanspacer{} โย โส ภนฺเต ภิกฺขุ สตฺถริ สคารโว ธมฺเม สคารโว สํเฆ สคารโว สิกฺขาย สคารโว สมาธิสฺมิํ สคารโว อปฺปมาเท สคารโว, ปฏิสนฺถาเรปิ โส สคารโวติ. (7)}

\gambhira{สตฺตกนิปาตปาฬิ}
\prose{\csromanfolio{492}อิมสฺส โข อหํ ภนฺเต ภควตา สํขิตฺเตน ภาสิตสฺส เอวํ วิตฺถาเรน อตฺถํ อาชานามี”ติ.}

\prose{สาธุ สาธุ สาริปุตฺต, สาธุ โข ตฺวํ สาริปุตฺต อิมสฺส มยา สํขิตฺเตน ภาสิตสฺส เอวํ วิตฺถาเรน อตฺถํ อาชานาสิ.\csromanspacer{} โส วต สาริปุตฺต ภิกฺขุ สตฺถริ อคารโว “ธมฺเม สคารโว ภวิสฺสตี”ติ เนตํ ฐานํ วิชฺชติ ฯเปฯ.\csromanspacer{} โย โส สาริปุตฺต ภิกฺขุ สตฺถริ อคารโว ธมฺเม อคารโว สํเฆ อคารโว สิกฺขาย อคารโว สมาธิสฺมิํ อคารโว, อปฺปมาเทปิ โส อคารโว. (1-6)}

\prose{โส วต สาริปุตฺต ภิกฺขุ สตฺถริ อคารโว ธมฺเม อคารโว สํเฆ อคารโว สิกฺขาย อคารโว สมาธิสฺมิํ อคารโว อปฺปมาเท อคารโว “ปฏิสนฺถาเร สคารโว ภวิสฺสตี”ติ เนตํ ฐานํ วิชฺชติ.\csromanspacer{} โย โส สาริปุตฺต ภิกฺขุ สตฺถริ อคารโว ธมฺเม อคารโว สํเฆ อคารโว สิกฺขาย อคารโว สมาธิสฺมิํ อคารโว อปฺปมาเท อคารโว, ปฏิสนฺถาเรปิ โส อคารโว. (7)}

\prose{โส วต สาริปุตฺต ภิกฺขุ สตฺถริ สคารโว “ธมฺเม อคารโว ภวิสฺสตี”ติ เนตํ ฐานํ วิชฺชติ ฯเปฯ.\csromanspacer{} โย โส สาริปุตฺต ภิกฺขุ สตฺถริ สคารโว, ธมฺเมปิ โส สคารโว ฯเปฯ. (1-6)}

\prose{โส วต สาริปุตฺต ภิกฺขุ สตฺถริ สคารโว ธมฺเม สคารโว ฯเปฯ อปฺปมาเท สคารโว “ปฏิสนฺถาเร อคารโว ภวิสฺสตี”ติ เนตํ ฐานํ วิชฺชติ.\csromanspacer{} โย โส สาริปุตฺต ภิกฺขุ สตฺถริ สคารโว ฯเปฯ อปฺปมาเท สคารโว, ปฏิสนฺถาเรปิ โส สคารโว. (7)}

\prose{โส วต สาริปุตฺต ภิกฺขุ สตฺถริ สคารโว “ธมฺเมปิ สคารโว ภวิสฺสตี”ติ ฐานเมตํ วิชฺชติ.\csromanspacer{} โย โส สาริปุตฺต ภิกฺขุ สตฺถริ สคารโว, ธมฺเมปิ โส สคารโว ฯเปฯ. (1-6)}

\prose{โส วต สาริปุตฺต ภิกฺขุ สตฺถริ สคารโว ฯเปฯ อปฺปมาเท สคารโว “ปฏิสนฺถาเรปิ โส สคารโว ภวิสฺสตี”ติ ฐานเมตํ วิชฺชติ.\csromanspacer{} โย โส สาริปุตฺต ภิกฺขุ สตฺถริ สคารโว ฯเปฯ อปฺปมาเท สคารโว, ปฏิสนฺถาเรปิ โส สคารโวติ. (7)}

\chapterheadpageread{7. มหาวคฺค}
\prose{\csromanfolio{493}อิมสฺส โข สาริปุตฺต มยา สํขิตฺเตน ภาสิตสฺส เอวํ วิตฺถาเรน อตฺโถ ทฏฺฐพฺโพติ.\csromanspacer{}\csromanspacer{}ฉฏฺฐํ.}

\csromanmatikaanchor{mtk.267}
\csromantocmarkref{subsection}{7. ภาวนาสุตฺต}{mtk.267}
\bookmark[dest={mtk.267},level=2]{7. ภาวนาสุตฺต}
\csromanheader{2}{7. ภาวนาสุตฺต}
\proseitem{71}{ภาวนํ อนนุยุตฺตสฺส ภิกฺขเว ภิกฺขุโน วิหรโต กิญฺจาปิ เอวํ อิจฺฉา อุปฺปชฺเชยฺย “อโห วต เม อนุปาทาย อาสเวหิ จิตฺตํ วิมุจฺเจยฺยา”ติ.\csromanspacer{} อถ ขฺวสฺส เนว อนุปาทาย อาสเวหิ จิตฺตํ วิมุจฺจติ.\csromanspacer{} ตํ กิสฺส เหตุ? อภาวิตตฺตาติสฺส วจนียํ.\csromanspacer{} กิสฺส อภาวิตตฺตา? จตุนฺนํ สติปฏฺฐานานํ จตุนฺนํ สมฺมปฺปธานานํ จตุนฺนํ อิทฺธิปาทานํ ปญฺจนฺนํ อินฺทฺริยานํ ปญฺจนฺนํ พลานํ สตฺตนฺนํ โพชฺฌงฺคานํ อริยสฺส อฏฺฐงฺคิกสฺส มคฺคสฺส.}

\prose{เสยฺยถาปิ ภิกฺขเว กุกฺกุฏิยา อณฺฑานิ อฏฺฐ วา ทส วา ทฺวาทส วา, ตานสฺสุ กุกฺกุฏิยา น สมฺมา อธิสยิตานิ น สมฺมา ปริเสทิตานิ น สมฺมา ปริภาวิตานิ.\csromanspacer{} กิญฺจาปิ ตสฺสา กุกฺกุฏิยา เอวํ อิจฺฉา อุปฺปชฺเชยฺย “อโห วต เม กุกฺกุฏโปตกา ปาทนขสิขาย วา มุขตุณฺฑเกน วา อณฺฑโกสํ ปทาเลตฺวา โสตฺถินา อภินิพฺภิชฺเชยฺยุนฺ”ติ.\csromanspacer{} อถ โข อภพฺพาว เต กุกฺกุฏโปตกา ปาทนขสิขาย วา มุขตุณฺฑเกน วา อณฺฑโกสํ ปทาเลตฺวา โสตฺถินา อภินิพฺภิชฺชิตุํ.\csromanspacer{} ตํ กิสฺส เหตุ? ตถา หิ ภิกฺขเว กุกฺกุฏิยา อณฺฑานิ น สมฺมา อธิสยิตานิ น สมฺมา ปริเสทิตานิ น สมฺมา ปริภาวิตานิ.\csromanspacer{} เอวเมวํ โข ภิกฺขเว ภาวนํ อนนุยุตฺตสฺส ภิกฺขุโน วิหรโต กิญฺจาปิ เอวํ อิจฺฉา อุปฺปชฺเชยฺย “อโห วต เม อนุปาทาย อาสเวหิ จิตฺตํ วิมุจฺเจยฺยา”ติ.\csromanspacer{} อถ ขฺวสฺส เนว อนุปาทาย อาสเวหิ จิตฺตํ วิมุจฺจติ.\csromanspacer{} ตํ กิสฺส เหตุ? อภาวิตตฺตาติสฺส วจนียํ.\csromanspacer{} กิสฺส อภาวิตตฺตา? จตุนฺนํ สติปฏฺฐานานํ จตุนฺนํ สมฺมปฺปธานานํ จตุนฺนํ อิทฺธิปาทานํ ปญฺจนฺนํ อินฺทฺริยานํ ปญฺจนฺนํ พลานํ สตฺตนฺนํ โพชฺฌงฺคานํ อริยสฺส อฏฺฐงฺคิกสฺส มคฺคสฺส.}

\prose{ภาวนํ อนุยุตฺตสฺส ภิกฺขเว ภิกฺขุโน วิหรโต กิญฺจาปิ น เอวํ อิจฺฉา อุปฺปชฺเชยฺย “อโห วต เม อนุปาทาย อาสเวหิ จิตฺตํ วิมุจฺเจยฺยา”ติ.\csromanspacer{} อถ ขฺวสฺส อนุปาทาย อาสเวหิ จิตฺตํ วิมุจฺจติ.\csromanspacer{} ตํ กิสฺส เหตุ? ภาวิตตฺตาติสฺส วจนียํ.\csromanspacer{} กิสฺส ภาวิตตฺตา? จตุนฺนํ สติปฏฺฐานานํ จตุนฺนํ สมฺมปฺปธานานํ จตุนฺนํ อิทฺธิปาทานํ ปญฺจนฺนํ อินฺทฺริยานํ ปญฺจนฺนํ พลานํ สตฺตนฺนํ โพชฺฌงฺคานํ อริยสฺส อฏฺฐงฺคิกสฺส มคฺคสฺส.}

\gambhira{สตฺตกนิปาตปาฬิ}
\prose{\csromanfolio{494}เสยฺยถาปิ ภิกฺขเว กุกฺกุฏิยา อณฺฑานิ อฏฺฐ วา ทส วา ทฺวาทส วา, ตานสฺสุ กุกฺกุฏิยา สมฺมา อธิสยิตานิ สมฺมา ปริเสทิตานิ สมฺมา ปริภาวิตานิ.\csromanspacer{} กิญฺจาปิ ตสฺสา กุกฺกุฏิยา น เอวํ อิจฺฉา อุปฺปชฺเชยฺย “อโห วต เม กุกฺกุฏโปตกา ปาทนขสิขาย วา มุขตุณฺฑเกน วา อณฺฑโกสํ ปทาเลตฺวา โสตฺถินา อภินิพฺภิชฺเชยฺยุนฺ”ติ.\csromanspacer{} อถ โข ภพฺพาว เต กุกฺกุฏโปตกา ปาทนขสิขาย วา มุขตุณฺฑเกน วา อณฺฑโกสํ ปทาเลตฺวา โสตฺถินา อภินิพฺภิชฺชิตุํ.\csromanspacer{} ตํ กิสฺส เหตุ? ตถา หิ ภิกฺขเว กุกฺกุฏิยา อณฺฑานิ สมฺมา อธิสยิตานิ สมฺมา ปริเสทิตานิ สมฺมา ปริภาวิตานิ.\csromanspacer{} เอวเมวํ โข ภิกฺขเว ภาวนํ อนุยุตฺตสฺส ภิกฺขุโน วิหรโต กิญฺจาปิ น เอวํ อิจฺฉา อุปฺปชฺเชยฺย “อโห วต เม อนุปาทาย อาสเวหิ จิตฺตํ วิมุจฺเจยฺยา”ติ.\csromanspacer{} อถ ขฺวสฺส อนุปาทาย อาสเวหิ จิตฺตํ วิมุจฺจติ.\csromanspacer{} ตํ กิสฺส เหตุ? ภาวิตตฺตาติสฺส วจนียํ.\csromanspacer{} กิสฺส ภาวิตตฺตา? จตุนฺนํ สติปฏฺฐานานํ ฯเปฯ อริยสฺส อฏฺฐงฺคิกสฺส มคฺคสฺส.}

\prose{เสยฺยถาปิ ภิกฺขเว ผลคณฺฑสฺส\footnote{ปลคณฺฑสฺส (?)} วา ผลคณฺฑนฺเตวาสิกสฺส วา ทิสฺสนฺเตว\footnote{ขียนฺเต (ก) 3-3. ทิสฺสนฺติ องฺคุลิปทานิ (สี), องฺคุลปทานิ ทิสฺสนฺติ องฺคุลปทํ (ก)} วาสิชเฏ3 องฺคุลิปทานิ, ทิสฺสติ\footnote{ทิสฺสนฺติ (สฺยา)} องฺคุฏฺฐปทํ3.\csromanspacer{} โน จ ขฺวสฺส เอวํ ญาณํ โหติ “เอตฺตกํ เม อชฺช วาสิชฏสฺส ขีณํ, เอตฺตกํ หิยฺโย, เอตฺตกํ ปเร”ติ, อถ ขฺวสฺส ขีเณ “ขีณนฺ”เตว ญาณํ โหติ.\csromanspacer{} เอวเมวํ โข ภิกฺขเว ภาวนํ อนุยุตฺตสฺส ภิกฺขุโน วิหรโต กิญฺจาปิ น เอวํ ญาณํ โหติ “เอตฺตกํ เม อชฺช อาสวานํ ขีณํ, เอตฺตกํ หิยฺโย, เอตฺตกํ ปเร”ติ, อถ ขฺวสฺส ขีเณ “ขีณนฺ”เตว ญาณํ โหติ.}

\prose{เสยฺยถาปิ ภิกฺขเว สามุทฺทิกาย นาวาย เวตฺตพนฺธนพทฺธาย\footnote{เวตฺตพนฺธาย (ก)} ฉ มาสานิ อุทเก ปริยาทาย เหมนฺติเกน ถเล อุกฺขิตฺตาย วาตาตปปเรตานิ พนฺธนานิ, ตานิ ปาวุสฺสเกน เมเฆน อภิปฺปวุฏฺฐานิ อปฺปกสิเรเนว ปริหายนฺติ\footnote{ปฏิปฺปสฺสมฺภนฺติ (สี, สฺยา)}, ปูติกานิ ภวนฺติ.\csromanspacer{} เอวเมวํ โข ภิกฺขเว ภาวนํ อนุยุตฺตสฺส ภิกฺขุโน วิหรโต อปฺปกสิเรเนว สํโยชนานิ ปฏิปฺปสฺสมฺภนฺติ, ปูติกานิ ภวนฺตีติ.\csromanspacer{}\csromanspacer{}สตฺตมํ.}

\csromanmatikaanchor{mtk.268}
\csromantocmarkref{subsection}{8. อคฺคิกฺขนฺโธปมสุตฺต}{mtk.268}
\bookmark[dest={mtk.268},level=2]{8. อคฺคิกฺขนฺโธปมสุตฺต}
\csromanheader{2}{7. มหาวคฺค \textbf{8. อคฺคิกฺขนฺโธปมสุตฺต}}
\proseitem{72}{\csromanfolio{495}เอวํ เม สุตํ\csromandash{}เอกํ สมยํ ภควา โกสเลสุ จาริกํ จรติ มหตา ภิกฺขุสํเฆน สทฺธิํ.\csromanspacer{} อทฺทสา โข ภควา อทฺธานมคฺคปฺปฏิปนฺโน อญฺญตรสฺมิํ ปเทเส มหนฺตํ อคฺคิกฺขนฺธํ อาทิตฺตํ สมฺปชฺชลิตํ สโชติภูตํ\footnote{สญฺโชติภูตํ (สฺยา)}, ทิสฺวาน มคฺคา โอกฺกมฺม อญฺญตรสฺมิํ รุกฺขมูเล ปญฺญตฺเต อาสเน นิสีทิ, นิสชฺช โข ภควา ภิกฺขู อามนฺเตสิ “ปสฺสถ โน ตุเมฺห ภิกฺขเว อมุํ มหนฺตํ อคฺคิกฺขนฺธํ อาทิตฺตํ สมฺปชฺชลิตํ สโชติภูตนฺ”ติ.\csromanspacer{} เอวํ ภนฺเตติ.}

\prose{ตํ กิํ มญฺญถ ภิกฺขเว, กตมํ นุ โข วรํ, ยํ อมุํ มหนฺตํ อคฺคิกฺขนฺธํ อาทิตฺตํ สมฺปชฺชลิตํ สโชติภูตํ อาลิงฺเคตฺวา อุปนิสีเทยฺย วา อุปนิปชฺเชยฺย วา, ยํ วา ขตฺติยกญฺญํ วา พฺราหฺมณกญฺญํ วา คหปติกญฺญํ วา มุทุตลุนหตฺถปาทํ อาลิงฺเคตฺวา อุปนิสีเทยฺย วา อุปนิปชฺเชยฺย วาติ.\csromanspacer{} เอตเทว ภนฺเต วรํ, ยํ ขตฺติยกญฺญํ วา พฺราหฺมณกญฺญํ วา คหปติกญฺญํ วา มุทุตลุนหตฺถปาทํ อาลิงฺเคตฺวา อุปนิสีเทยฺย วา อุปนิปชฺเชยฺย วา.\csromanspacer{} ทุกฺขํ เหตํ ภนฺเต, ยํ อมุํ มหนฺตํ อคฺคิกฺขนฺธํ อาทิตฺตํ สมฺปชฺชลิตํ สโชติภูตํ อาลิงฺเคตฺวา อุปนิสีเทยฺย วา อุปนิปชฺเชยฺย วาติ.}

\prose{อาโรจยามิ โว ภิกฺขเว, ปฏิเวทยามิ โว ภิกฺขเว, ยถา เอตเทว ตสฺส วรํ ทุสฺสีลสฺส ปาปธมฺมสฺส อสุจิสงฺกสฺสรสมาจารสฺส ปฏิจฺฉนฺนกมฺมนฺตสฺส อสฺสมณสฺส สมณปฏิญฺญสฺส อพฺรหฺมจาริสฺส พฺราหฺมจาริปฏิญฺญสฺส อนฺโตปูติกสฺส อวสฺสุตสฺส กสมฺพุชาตสฺส, ยํ อมุํ มหนฺตํ อคฺคิกฺขนฺธํ อาทิตฺตํ สมฺปชฺชลิตํ สโชติภูตํ อาลิงฺเคตฺวา อุปนิสีเทยฺย วา อุปนิปชฺเชยฺย วา.\csromanspacer{} ตํ กิสฺส เหตุ? ตโตนิทานํ หิ โส ภิกฺขเว มรณํ วา นิคจฺเฉยฺย มรณมตฺตํ วา ทุกฺขํ, น เตฺวว ตปฺปจฺจยา กายสฺส เภทา ปรํ มรณา อปายํ ทุคฺคติํ วินิปาตํ นิรยํ อุปปชฺเชยฺย.}

\prose{ยญฺจ โข โส ภิกฺขเว ทุสฺสีโล ปาปธมฺโม อสุจิสงฺกสฺสรสมาจาโร ฯเปฯ กสมฺพุชาโต ขตฺติยกญฺญํ วา พฺราหฺมณกญฺญํ วา คหปติกญฺญํ วา มุทุตลุนหตฺถปาทํ อาลิงฺเคตฺวา อุปนิสีทติ วา อุปนิปชฺชติ วา, ตํ หิ ตสฺส\footnote{ตํ หิสฺส (ก)} ภิกฺขเว โหติ ทีฆรตฺตํ อหิตาย ทุกฺขาย, กายสฺส เภทา ปรํ มรณา อปายํ ทุคฺคติํ วินิปาตํ นิรยํ อุปปชฺชติ. (1)}

\gambhira{สตฺตกนิปาตปาฬิ}
\prose{\csromanfolio{496}ตํ กิํ มญฺญถ ภิกฺขเว, กตมํ นุ โข วรํ, ยํ พลวา ปุริโส ทฬฺหาย วาลรชฺชุยา อุโภ ชงฺฆา เวเฐตฺวา ฆํเสยฺย, สา ฉวิํ ฉินฺเทยฺย, ฉวิํ เฉตฺวา จมฺมํ ฉินฺเทยฺย, จมฺมํ เฉตฺวา มํสํ ฉินฺเทยฺย, มํสํ เฉตฺวา นฺหารุํ ฉินฺเทยฺย, นฺหารุํ เฉตฺวา อฏฺฐิํ ฉินฺเทยฺย, อฏฺฐิํ เฉตฺวา อฏฺฐิมิญฺชํ อาหจฺจ ติฏฺเฐยฺย.\csromanspacer{} ยํ วา ขตฺติยมหาสาลานํ วา พฺราหฺมณมหาสาลานํ วา คหปติ มหาสาลานํ วา อภิวาทนํ สาทิเยยฺยาติ.\csromanspacer{} เอตเทว ภนฺเต วรํ, ยํ ขตฺติยมหาสาลานํ วา พฺราหฺมณมหาสาลานํ วา คหปติมหาสาลานํ วา อภิวาทนํ สาทิเยยฺย.\csromanspacer{} ทุกฺขํ เหตํ ภนฺเต, ยํ พลวา ปุริโส ทฬฺหาย วาลรชฺชุยา ฯเปฯ อฏฺฐิมิญฺชํ อาหจฺจ ติฏฺเฐยฺยาติ.}

\prose{อาโรจยามิ โว ภิกฺขเว, ปฏิเวทยามิ โว ภิกฺขเว, ยถา เอตเทว ตสฺส วรํ ทุสฺสีลสฺส ฯเปฯ กสมฺพุชาตสฺส, ยํ พลวา ปุริโส ทฬฺหาย วาลรชฺชุยา อุโภ ชงฺฆา เวเฐตฺวา ฯเปฯ อฏฺฐิมิญฺชํ อาหจฺจ ติฏฺเฐยฺย.\csromanspacer{} ตํ กิสฺส เหตุ? ตโตนิทานํ หิ โส ภิกฺขเว มรณํ วา นิคจฺเฉยฺย มรณมตฺตํ วา ทุกฺขํ, น เตฺวว ตปฺปจฺจยา กายสฺส เภทา ปรํ มรณา อปายํ ทุคฺคติํ วินิปาตํ นิรยํ อุปปชฺเชยฺย.\csromanspacer{} ยญฺจ โข โส ภิกฺขเว ทุสฺสีโล ฯเปฯ กสมฺพุชาโต ขตฺติยมหาสาลานํ วา พฺราหฺมณมหาสาลานํ วา คหปติมหาสาลานํ วา อภิวาทนํ สาทิยติ, ตํ หิ ตสฺส ภิกฺขเว โหติ ทีฆรตฺตํ อหิตาย ทุกฺขาย, กายสฺส เภทา ปรํ มรณา อปายํ วินิปาตํ นิรยํ อุปปชฺชติ. (2)}

\prose{ตํ กิํ มญฺญถ ภิกฺขเว, กตมํ นุ โข วรํ, ยํ พลวา ปุริโส ติณฺหาย สตฺติยา เตลโธตาย ปจฺโจรสฺมิํ ปหเรยฺย, ยํ วา ขตฺติยมหาสาลานํ วา พฺราหฺมณมหาสาลานํ วา คหปติมหาสาลานํ วา อญฺชลิกมฺมํ สาทิเยยฺยาติ.\csromanspacer{} เอตเทว ภนฺเต วรํ, ยํ ขตฺติยมหาสาลานํ วา พฺราหฺมณมหาสาลานํ วา คหปติมหาสาลานํ วา อญฺชลิกมฺมํ สาทิเยยฺย.\csromanspacer{} ทุกฺขํ เหตํ ภนฺเต, ยํ พลวา ปุริโส ติณฺหาย สตฺติยา เตลโธตาย ปจฺโจรสฺมิํ ปหเรยฺยาติ.}

\prose{อาโรจยามิ โว ภิกฺขเว, ปฏิเวทยามิ โว ภิกฺขเว, ยถา เอตเทว ตสฺส วรํ ทุสฺสีลสฺส ฯเปฯ กสมฺพุชาตสฺส, ยํ พลวา ปุริโส ติณฺหาย สตฺติยา เตลโธตาย ปจฺโจรสฺมิํ ปหเรยฺย.\csromanspacer{} ตํ กิสฺส เหตุ? ตโตนิทานํ หิ โส ภิกฺขเว มรณํ วา นิคจฺเฉยฺย มรณมตฺตํ วา ทุกฺขํ, น เตฺวว ตปฺปจฺจยา กายสฺส เภทา ปรํ มรณา อปายํ ทุคฺคติํ}

\vspace{\tipitakaparskip}
\csromancenter{มหาวคฺค วินิปาตํ นิรยํ อุปปชฺเชยฺย.\csromanspacer{} ยญฺจ โข โส ภิกฺขเว ทุสฺสีโล ปาปธมฺโม ฯเปฯ กสมฺพุชาโต ขตฺติยมหาสาลานํ วา พฺราหฺมณมหาสาลานํ วา คหปติมหาสาลานํ วา อญฺชลิกมฺมํ สาทิยติ, ตํ หิ ตสฺส ภิกฺขเว โหติ ทีฆรตฺตํ อหิตาย ทุกฺขาย, กายสฺส เภทา ปรํ มรณา อปายํ ทุคฺคติํ วินิปาตํ นิรยํ อุปปชฺชติ. (3)}
\prose{\csromanfolio{497}ตํ กิํ มญฺญถ ภิกฺขเว, กตมํ นุ โข วรํ, ยํ พลวา ปุริโส ตตฺเตน อโยปฏฺเฏน อาทิตฺเตน สมฺปชฺชลิเตน สโชติภูเตน กายํ สมฺปลิเวเฐยฺย, ยํ วา ขตฺติยมหาสาลานํ วา พฺราหฺมณมหาสาลานํ วา คหปติมหาสาลานํ วา สทฺธาเทยฺยํ จีวรํ ปริภุญฺเชยฺยาติ.\csromanspacer{} เอตเทว ภนฺเต วรํ, ยํ ขตฺติยมหาสาลานํ วา ฯเปฯ สทฺธาเทยฺยํ จีวรํ ปริภุญฺเชยฺย.\csromanspacer{} ทุกฺขํ เหตํ ภนฺเต, ยํ พลวา ปุริโส ตตฺเตน อโยปฏฺเฏน อาทิตฺเตน สมฺปชฺชลิเตน สโชติภูเตน กายํ สมฺปลิเวเฐยฺยาติ.}

\prose{อาโรจยามิ โว ภิกฺขเว, ปฏิเวทยามิ โว ภิกฺขเว, ยถา เอตเทว ตสฺส วรํ ทุสฺสีลสฺส ฯเปฯ กสมฺพุชาตสฺส, ยํ พลวา ปุริโส ตตฺเตน อโยปฏฺเฏน อาทิตฺเตน สมฺปชฺชลิเตน สโชติภูเตน กายํ สมฺปลิเวเฐยฺย.\csromanspacer{} ตํ กิสฺส เหตุ? ตโตนิทานํ หิ โส ภิกฺขเว มรณํ วา นิคจฺเฉยฺย มรณมตฺตํ วา ทุกฺขํ, น เตฺวว ตปฺปจฺจยา กายสฺส เภทา ปรํ มรณา อปายํ ทุคฺคติํ วินิปาตํ นิรยํ อุปปชฺเชยฺย.\csromanspacer{} ยญฺจ โข โส ภิกฺขเว ทุสฺสีโล ฯเปฯ กสมฺพุชาโต ขตฺติยมหาสาลานํ วา พฺราหฺมณมหาสาลานํ วา คหปติมหาสาลานํ วา สทฺธาเทยฺยํ จีวรํ ปริภุญฺชติ, ตํ หิ ตสฺส ภิกฺขเว โหติ ทีฆรตฺตํ อหิตาย ทุกฺขาย, กายสฺส เภทา ปรํ มรณา อปายํ ทุคฺคติํ วินิปาตํ นิรยํ อุปปชฺชติ. (4)}

\prose{ตํ กิํ มญฺญถ ภิกฺขเว, กตมํ นุ โข วรํ, ยํ พลวา ปุริโส ตตฺเตน อโยสงฺกุนา มุขํ วิวริตฺวา ตตฺตํ โลหคุฬํ อาทิตฺตํ สมฺปชฺชลิตํ สโชติภูตํ มุเข ปกฺขิเปยฺย, ตํ ตสฺส โอฏฺฐมฺปิ ทเหยฺย\footnote{qอเหยฺย (กตฺถจิ)}, มุขมฺปิ ทเหยฺย, ชิวฺหมฺปิ ทเหยฺย, กณฺฐมฺปิ ทเหยฺย, อุรมฺปิ\footnote{ม 3. 224 เทวทูตสุตฺเต ปน “อุทรมฺปิ” อิติ วิเทสปาโฐ ทิสฺสติ.} ทเหยฺย, อนฺตมฺปิ อนฺตคุณมฺปิ อาทาย อโธภาคา\footnote{อโธภาคํ (ก)} นิกฺขเมยฺย, ยํ วา ขตฺติยมหาสาลานํ วา พฺราหฺมณมหาสาลานํ วา คหปติมหาสาลานํ วา สทฺธาเทยฺยํ}

\vspace{\tipitakaparskip}
\csromancenter{สตฺตกนิปาตปาฬิ ปิณฺฑปาตํ ปริภุญฺเชยฺยาติ.\csromanspacer{} เอตเทว ภนฺเต วรํ, ยํ ขตฺติยมหาสาลานํ วา พฺราหฺมณมหาสาลานํ วา คหปติมหาสาลานํ วา สทฺธาเทยฺยํ ปิณฺฑปาตํ ปริภุญฺเชยฺย.\csromanspacer{} ทุกฺขํ เหตํ ภนฺเต, ยํ พลวา ปุริโส ตตฺเตน อโยสงฺกุนา มุขํ วิวริตฺวา ตตฺตํ โลหคุฬํ อาทิตฺตํ สมฺปชฺชลิตํ สโชติภูตํ มุเข ปกฺขิเปยฺย, ตํ ตสฺส โอฏฺฐมฺปิ ทเหยฺย, มุขมฺปิ ทเหยฺย, ชิวฺหมฺปิ ทเหยฺย, กณฺฐมฺปิ ทเหยฺย, อุรมฺปิ ทเหยฺย, อนฺตมฺปิ อนฺตคุณมฺปิ อาทาย อโธภาคํ นิกฺขเมยฺยาติ.}
\prose{\csromanfolio{498}อาโรจยามิ โว ภิกฺขเว, ปฏิเวทยามิ โว ภิกฺขเว, ยถา เอตเทว ตสฺส วรํ ทุสฺสีลสฺส ฯเปฯ กสมฺพุชาตสฺส, ยํ พลวา ปุริโส ตตฺเตน อโยสงฺกุนา มุขํ วิวริตฺวา ตตฺตํ โลหคุฬํ อาทิตฺตํ สมฺปชฺชลิตํ สโชติภูตํ มุเข ปกฺขิเปยฺย, ตํ ตสฺส โอฏฺฐมฺปิ ทเหยฺย, มุขมฺปิ ทเหยฺย, ชิวฺหมฺปิ ทเหยฺย, กณฺฐมฺปิ ทเหยฺย, อุรมฺปิ ทเหยฺย, อนฺตมฺปิ อนฺตคุณมฺปิ อาทาย อโธภาคํ นิกฺขเมยฺย.\csromanspacer{} ตํ กิสฺส เหตุ? ตโตนิทานํ หิ โส ภิกฺขเว มรณํ วา นิคจฺเฉยฺย มรณมตฺตํ วา ทุกฺขํ, น เตฺวว ตปฺปจฺจยา กายสฺส เภทา ปรํ มรณา อปายํ ทุคฺคติํ วินิปาตํ นิรยํ อุปปชฺเชยฺย.\csromanspacer{} ยญฺจ โข โส ภิกฺขเว ทุสฺสีโล ปาปธมฺโม ฯเปฯ กสมฺพุชาโต ขตฺติยมหาสาลานํ วา พฺราหฺมณมหาสาลานํ วา คหปติมหาสาลานํ วา สทฺธาเทยฺยํ ปิณฺฑปาตํ ปริภุญฺชติ, ตํ หิ ตสฺส โหติ ทีฆรตฺตํ อหิตาย ทุกฺขาย, กายสฺส เภทา ปรํ มรณา อปายํ ทุคฺคติํ วินิปาตํ นิรยํ อุปปชฺชติ. (5)}

\prose{ตํ กิํ มญฺญถ ภิกฺขเว, กตมํ นุ โข วรํ, ยํ พลวา ปุริโส สีเส วา คเหตฺวา ขนฺเธ วา คเหตฺวา ตตฺตํ อโยมญฺจํ วา อโยปีฐํ วา อภินิสีทาเปยฺย วา อภินิปชฺชาเปยฺย วา, ยํ วา ขตฺติยมหาสาลานํ วา พฺราหฺมณมหาสาลานํ วา คหปติมหาสาลานํ วา สทฺธาเทยฺยํ มญฺจปีฐํ\footnote{มญฺจํ วา ปีฐํ วา (ก)} ปริภุญฺเชยฺยาติ.\csromanspacer{} เอตเทว ภนฺเต วรํ, ยํ ขตฺติยมหาสาลานํ วา พฺราหฺมณมหาสาลานํ วา คหปติมหาสาลานํ วา สทฺธาเทยฺยํ มญฺจปีฐํ ปริภุญฺเชยฺย.\csromanspacer{} ทุกฺขํ เหตํ ภนฺเต, ยํ พลวา ปุริโส สีเส วา คเหตฺวา ขนฺเธ วา คเหตฺวา ตตฺตํ อโยมญฺจํ วา อโยปีฐํ วา อภินิสีทาเปยฺย วา อภินิปชฺชาเปยฺย วาติ.}

\chapterheadpageread{7. มหาวคฺค}
\prose{\csromanfolio{499}อาโรจยามิ โว ภิกฺขเว, ปฏิเวทยามิ โว ภิกฺขเว, ยถา เอตเทว ตสฺส วรํ ทุสฺสีลสฺส ฯเปฯ กสมฺพุชาตสฺส, ยํ พลวา ปุริโส สีเส วา คเหตฺวา ขนฺเธ วา คเหตฺวา ตตฺตํ อโยมญฺจํ วา อโยปีฐํ วา อภินิสีทาเปยฺย วา อภินิปชฺชาเปยฺย วา.\csromanspacer{} ตํ กิสฺส เหตุ? ตโตนิทานํ หิ โส ภิกฺขเว มรณํ วา นิคจฺเฉยฺย มรณมตฺตํ วา ทุกฺขํ, น เตฺวว ตปฺปจฺจยา กายสฺส เภทา ปรํ มรณา อปายํ ทุคฺคติํ วินิปาตํ นิรยํ อุปปชฺเชยฺย.\csromanspacer{} ยญฺจ โข โส ภิกฺขเว ทุสฺสีโล ปาปธมฺโม ฯเปฯ กสมฺพุชาโต ขตฺติยมหาสาลานํ วา พฺราหฺมณมหาสาลานํ วา คหปติมหาสาลานํ วา สทฺธาเทยฺยํ มญฺจปีฐํ ปริภุญฺชติ, ตํ หิ ตสฺส ภิกฺขเว โหติ ทีฆรตฺตํ อหิตาย ทุกฺขาย, กายสฺส เภทา ปรํ มรณา อปายํ ทุคฺคติํ วินิปาตํ นิรยํ อุปปชฺชติ. (6)}

\prose{ตํ กิํ มญฺญถ ภิกฺขเว, กตมํ นุ โข วรํ, ยํ พลวา ปุริโส อุทฺธํปาทํ อโธสิรํ คเหตฺวา ตตฺตาย โลหกุมฺภิยา ปกฺขิเปยฺย อาทิตฺตาย สมฺปชฺชลิตาย สโชติภูตาย, โส ตตฺถ เผณุทฺเทหกํ ปจฺจมาโน สกิมฺปิ อุทฺธํ คจฺเฉยฺย, สกิมฺปิ อโธ คจฺเฉยฺย, สกิมฺปิ ติริยํ คจฺเฉยฺย, ยํ วา ขตฺติยมหาสาลานํ วา พฺราหฺมณมหาสาลานํ วา คหปติมหาสาลานํ วา สทฺธาเทยฺยํ วิหารํ ปริภุญฺเชยฺยาติ.\csromanspacer{} เอตเทว ภนฺเต วรํ, ยํ ขตฺติยมหาสาลานํ วา พฺราหฺมณมหาสาลานํ วา คหปติมหาสาลานํ วา สทฺธาเทยฺยํ วิหารํ ปริภุญฺเชยฺย.\csromanspacer{} ทุกฺขํ เหตํ ภนฺเต, ยํ พลวา ปุริโส อุทฺธํปาทํ อโธสิรํ คเหตฺวา ตตฺตาย โลหกุมฺภิยา ปกฺขิเปยฺย อาทิตฺตาย สมฺปชฺชลิตาย สโชติภูตาย, โส ตตฺถ เผณุทฺเทหกํ ปจฺจมาโน สกิมฺปิ อุทฺธํ คจฺเฉยฺย, สกิมฺปิ อโธ คจฺเฉยฺย, สกิมฺปิ ติริยํ คจฺเฉยฺยาติ.}

\prose{อาโรจยามิ โว ภิกฺขเว, ปฏิเวทยามิ โว ภิกฺขเว, ยถา เอตเทว ตสฺส วรํ ทุสฺสีลสฺส ปาปธมฺมสฺส ฯเปฯ กสมฺพุชาตสฺส, ยํ พลวา ปุริโส อุทฺธํปาทํ อโธสิรํ คเหตฺวา ฯเปฯ สกิมฺปิ ติริยํ คจฺเฉยฺย.\csromanspacer{} ตํ กิสฺส เหตุ? ตโตนิทานํ หิ โส ภิกฺขเว มรณํ วา นิคจฺเฉยฺย มรณมตฺตํ วา ทุกฺขํ, น เตฺวว ตปฺปจฺจยา กายสฺส เภทา ปรํ มรณา อปายํ ทุคฺคติํ วินิปาตํ นิรยํ อุปปชฺเชยฺย.\csromanspacer{} ยญฺจ โข โส ภิกฺขเว ทุสฺสีโล ปาปธมฺโม ฯเปฯ กสมฺพุชาโต ขตฺติยมหาสาลานํ วา พฺราหฺมณมหาสาลานํ วา คหปติมหาสาลานํ วา สทฺธาเทยฺยํ วิหารํ}

\vspace{\tipitakaparskip}
\csromancenter{สตฺตกนิปาตปาฬิ ปริภุญฺชติ, ตํ หิ ตสฺส ภิกฺขเว โหติ ทีฆรตฺตํ อหิตาย ทุกฺขาย, กายสฺส เภทา ปรํ มรณา อปายํ ทุคฺคติํ วินิปาตํ นิรยํ อุปปชฺชติ. (7)}
\prose{\csromanfolio{500}ตสฺมาติห ภิกฺขเว เอวํ สิกฺขิตพฺพํ “เยสญฺจ\footnote{เยสํ (?)} มยํ ปริภุญฺชาม จีวรปิณฺฑปาตเสนาสนคิลานปฺปจฺจยเภสชฺชปริกฺขารํ\footnote{...ปริกฺขารานํ (สี, สฺยา, ก)}, เตสํ เต การา มหปฺผลา ภวิสฺสนฺติ มหานิสํสา, อมฺหากญฺเจวายํ ปพฺพชฺชา อวญฺฌา ภวิสฺสติ สผลา ส-อุทฺรยา”ติ.\csromanspacer{} เอวํ หิ โว ภิกฺขเว สิกฺขิตพฺพํ.\csromanspacer{} อตฺตตฺถํ วา ภิกฺขเว สมฺปสฺสมาเนน อลเมว อปฺปมาเทน สมฺปาเทตุํ, ปรตฺถํ วา ภิกฺขเว สมฺปสฺสมาเนน อลเมว อปฺปมาเทน สมฺปาเทตุํ, อุภยตฺถํ วา ภิกฺขเว สมฺปสฺสมาเนน อลเมว อปฺปมาเทน สมฺปาเทตุนฺติ.}

\prose{อิทมโวจ ภควา\footnote{อิทมโวจ ภควา ฯเปฯ (ก) 4. อุคฺคจฺฉิ (ก)}, อิมสฺมิํ จ ปน เวยฺยากรณสฺมิํ ภญฺญมาเน สฏฺฐิมตฺตานํ ภิกฺขูนํ อุณฺหํ โลหิตํ มุขโต อุคฺคญฺฉิ4, สฏฺฐิมตฺตา ภิกฺขู สิกฺขํ ปจฺจกฺขาย หีนายาวตฺติํสุ “สุทุกฺกรํ ภควา สุทุกฺกรํ ภควา”ติ, สฏฺฐิมตฺตานํ ภิกฺขูนํ อนุปาทาย อาสเวหิ จิตฺตานิ วิมุจฺจิํสูติ.\csromanspacer{}\csromanspacer{}อฏฺฐมํ.}

\csromanmatikaanchor{mtk.269}
\csromantocmarkref{subsection}{9. สุเนตฺตสุตฺต}{mtk.269}
\bookmark[dest={mtk.269},level=2]{9. สุเนตฺตสุตฺต}
\csromanheader{2}{9. สุเนตฺตสุตฺต}
\proseitem{73}{\csromansymbolfootnote{*}{เหฏฺฐา 326, 476 ปิฏฺเฐสุปิ.}ภูตปุพฺพํ ภิกฺขเว สุเนตฺโต นาม สตฺถา อโหสิ ติตฺถกโร กาเมสุ วีตราโค.\csromanspacer{} สุเนตฺตสฺส โข ปน ภิกฺขเว สตฺถุโน อเนกานิ สาวกสตานิ อเหสุํ, สุเนตฺโต สตฺถา สาวกานํ พฺรหฺมโลกสหพฺยตาย ธมฺมํ เทเสสิ.\csromanspacer{} เย โข ปน ภิกฺขเว\footnote{เย โข ภิกฺขเว (สี, สฺยา)} สุเนตฺตสฺส สตฺถุโน พฺรหฺมโลกสหพฺยตาย ธมฺมํ เทเสนฺตสฺส จิตฺตานิ นปฺปสาเทสุํ, เต กายสฺส เภทา ปรํ มรณา อปายํ ทุคฺคติํ วินิปาตํ นิรยํ อุปปชฺชิํสุ.\csromanspacer{} เย โข ปน ภิกฺขเว สุเนตฺตสฺส สตฺถุโน พฺรหฺมโลกสหพฺยตาย ธมฺมํ เทเสนฺตสฺส จิตฺตานิ ปสาเทสุํ, เต กายสฺส เภทา ปรํ มรณา สุคติํ สคฺคํ โลกํ อุปปชฺชิํสุ.}

\prose{ภูตปุพฺพํ ภิกฺขเว มูคปกฺโข นาม สตฺถา อโหสิ ฯเปฯ อรเนมิ นาม สตฺถา อโหสิ ฯเปฯ กุทฺทาลโก\footnote{กุทฺทาโล (สี, สฺยา)} นาม สตฺถา อโหสิ ฯเปฯ หตฺถิปาโล นาม สตฺถา อโหสิ ฯเปฯ โชติปาโล นาม สตฺถา อโหสิ ฯเปฯ อรโก นาม สตฺถา อโหสิ ติตฺถกโร กาเมสุ}

\vspace{\tipitakaparskip}
\csromancenter{มหาวคฺค วีตราโค.\csromanspacer{} อรกสฺส โข ปน ภิกฺขเว สตฺถุโน อเนกานิ สาวกสตานิ อเหสุํ, อรโก นาม สตฺถา สาวกานํ พฺรหฺมโลกสหพฺยตาย ธมฺมํ เทเสสิ.\csromanspacer{} เย โข ปน ภิกฺขเว อรกสฺส สตฺถุโน พฺรหฺมโลกสหพฺยตาย ธมฺมํ เทเสนฺตสฺส จิตฺตานิ นปฺปสาเทสุํ, เต กายสฺส เภทา ปรํ มรณา อปายํ ทุคฺคติํ วินิปาตํ นิรยํ อุปฺปชฺชิํสุ.\csromanspacer{} เย โข ปน ภิกฺขเว อรกสฺส สตฺถุโน พฺรหฺมโลกสหพฺยตาย ธมฺมํ เทเสนฺตสฺส จิตฺตานิ ปสาเทสุํ, เต กายสฺส เภทา ปรํ มรณา สุคติํ สคฺคํ โลกํ อุปปชฺชิํสุ.}
\prose{\csromanfolio{501}ตํ กิํ มญฺญถ ภิกฺขเว, โย อิเม สตฺต สตฺถาเร ติตฺถกเร กาเมสุ วีตราเค อเนกสตปริวาเร สสาวกสํเฆ ปทุฏฺฐจิตฺโต อกฺโกเสยฺย ปริภาเสยฺย, พหุํ โส อปุญฺญํ ปสเวยฺยาติ.\csromanspacer{} เอวํ ภนฺเต.\csromanspacer{} โย ภิกฺขเว อิเม สตฺต สตฺถาเร ติตฺถกเร กาเมสุ วีตราเค อเนกสตปริวาเร สสาวกสํเฆ ปทุฏฺฐจิตฺโต อกฺโกเสยฺย ปริภาเสยฺย, พหุํ โส อปุญฺญํ ปสเวยฺย.\csromanspacer{} โย เอกํ ทิฏฺฐิสมฺปนฺนํ ปุคฺคลํ ปทุฏฺฐจิตฺโต อกฺโกสติ ปริภาสติ, อยํ ตโต พหุตรํ อปุญฺญํ ปสวติ.\csromanspacer{} ตํ กิสฺส เหตุ? นาหํ ภิกฺขเว อิโต พหิทฺธา เอวรูปิํ ขนฺติํ วทามิ, ยถามํ สพฺรหฺมจารีสุ.}

\prose{ตสฺมาติห ภิกฺขเว เอวํ สิกฺขิตพฺพํ “น โน สพฺรหฺมจารีสุ\footnote{น เตฺวว อมฺหํ สพฺรหฺมจารีสุ (สฺยา) อํ 2. 313 ปิฏฺเฐ อญฺญถา ทิสฺสติ.} จิตฺตานิ ปทุฏฺฐานิ ภวิสฺสนฺตี”ติ, เอวํ หิ โว ภิกฺขเว สิกฺขิตพฺพนฺติ.\csromanspacer{}\csromanspacer{}นวมํ.}

\csromanmatikaanchor{mtk.270}
\csromantocmarkref{subsection}{10. อรกสุตฺต}{mtk.270}
\bookmark[dest={mtk.270},level=2]{10. อรกสุตฺต}
\csromanheader{2}{10. อรกสุตฺต}
\proseitem{74}{ภูตปุพฺพํ ภิกฺขเว อรโก นาม สตฺถา อโหสิ ติตฺถกโร กาเมสุ วีตราโค.\csromanspacer{} อรกสฺส โข ปน ภิกฺขเว สตฺถุโน อเนกานิ สาวกสตานิ อเหสุํ, อรโก สตฺถา สาวกานํ เอวํ ธมฺมํ เทเสติ “อปฺปกํ พฺราหฺมณ ชีวิตํ มนุสฺสานํ ปริตฺตํ ลหุกํ\footnote{ลหุสํ (ฏีกา)} พหุทุกฺขํ พหุปายาสํ มนฺตายํ\footnote{มนฺตาย (สพฺพตฺถ) ฏีกา ปสฺสิตพฺพา.} โพทฺธพฺพํ, กตฺตพฺพํ กุสลํ, จริตพฺพํ พฺรหฺมจริยํ, นตฺถิ ชาตสฺส อมรณํ.}

\gambhira{สตฺตกนิปาตปาฬิ}
\prose{\csromanfolio{502}เสยฺยถาปิ พฺราหฺมณ ติณคฺเค อุสฺสาวพินฺทุ สูริเย อุคฺคจฺฉนฺเต ขิปฺปํเยว ปฏิวิคจฺฉติ, น จิรฏฺฐิติกํ โหติ.\csromanspacer{} เอวเมว โข พฺราหฺมณ อุสฺสาวพินฺทูปมํ ชีวิตํ มนุสฺสานํ ปริตฺตํ ลหุกํ พหุทุกฺขํ พหุปายาสํ มนฺตายํ โพทฺธพฺพํ, กตฺตพฺพํ กุสลํ, จริตพฺพํ พฺรหฺมจริยํ, นตฺถิ ชาตสฺส อมรณํ.}

\prose{เสยฺยถาปิ พฺราหฺมณ ถุลฺลผุสิตเก เทเว วสฺสนฺเต อุทกพุพฺพุฬํ\footnote{อุทกปุปฺผุฬํ (ก)} ขิปฺปํเยว ปฏิวิคจฺฉติ, น จิรฏฺฐิติกํ โหติ.\csromanspacer{} เอวเมว โข พฺราหฺมณ อุทกพุพฺพุฬูปมํ ชีวิตํ มนุสฺสานํ ปริตฺตํ ลหุกํ พหุทุกฺขํ พหุปายาสํ มนฺตายํ โพทฺธพฺพํ, กตฺตพฺพํ กุสลํ, จริตพฺพํ พฺรหฺมจริยํ, นตฺถิ ชาตสฺส อมรณํ.}

\prose{เสยฺยถาปิ พฺราหฺมณ อุทเก ทณฺฑราชิ ขิปฺปํเยว ปฏิวิคจฺฉติ, น จิรฏฺฐิติกา โหติ.\csromanspacer{} เอวเมว โข พฺราหฺมณ อุทเก ทณฺฑราชูปมํ ชีวิตํ มนุสฺสานํ ปริตฺตํ ฯเปฯ นตฺถิ ชาตสฺส อมรณํ.}

\prose{เสยฺยถาปิ พฺราหฺมณ นที ปพฺพเตยฺยา ทูรงฺคมา สีฆโสตา หารหารินี.\csromanspacer{} นตฺถิ โส ขโณ วา ลโย วา มุหุตฺโต วา.\csromanspacer{} ยํ สา\footnote{ยาย (ก)} อาวตฺตติ\footnote{ถรติ (สี), ธรติ (สฺยา), อวติฏฺเฐยฺย (?)}, อถ โข สา คจฺฉเตว วตฺตเตว สนฺทเตว.\csromanspacer{} เอวเมว โข พฺราหฺมณ นทีปพฺพเตยฺยูปมํ ชีวิตํ มนุสฺสานํ ปริตฺตํ ลหุกํ ฯเปฯ นตฺถิ ชาตสฺส อมรณํ.}

\prose{เสยฺยถาปิ พฺราหฺมณ พลวา ปุริโส ชิวฺหคฺเค เขฬปิณฺฑํ สํยูหิตฺวา อกสิเรเนว วเมยฺย\footnote{ปตาเปยฺย (ก)}.\csromanspacer{} เอวเมว โข พฺราหฺมณ เขฬปิณฺฑูปมํ ชีวิตํ มนุสฺสานํ ปริตฺตํ ฯเปฯ นตฺถิ ชาตสฺส อมรณํ.}

\prose{เสยฺยถาปิ พฺราหฺมณ ทิวสํสนฺตตฺเต อโยกฏาเห มํสเปสิ\footnote{มํสเปสี (สี, สฺยา)} ปกฺขิตฺตา ขิปฺปํเยว ปฏิวิคจฺฉติ, น จิรฏฺฐิติกา โหติ.\csromanspacer{} เอวเมว โข พฺราหฺมณ มํสเปสูปมํ ชีวิตํ มนุสฺสานํ ปริตฺตํ ฯเปฯ นตฺถิ ชาตสฺส อมรณํ.}

\prose{เสยฺยถาปิ พฺราหฺมณ คาวี วชฺฌา อาฆาตนํ นียมานา ยํ ยเทว ปาทํ อุทฺธรติ, สนฺติเกว โหติ วธสฺส, สนฺติเกว มรณสฺส.\csromanspacer{} เอวเมว โข พฺราหฺมณ โควชฺฌูปมํ\footnote{คาวีวชฺฌูปมํ (สี, สฺยา)} ชีวิตํ มนุสฺสานํ ปริตฺตํ ลหุกํ พหุทุกฺขํ}

\vspace{\tipitakaparskip}
\csromancenter{มหาวคฺค พหุปายาสํ มนฺตายํ โพทฺธพฺพํ, กตฺตพฺพํ กุสลํ, จริตพฺพํ พฺรหฺมจริยํ, นตฺถิ ชาตสฺส อมรณนฺ”ติ.}
\prose{\csromanfolio{503}เตน โข ปน ภิกฺขเว สมเยน มนุสฺสานํ สฏฺฐิวสฺสสหสฺสานิ อายุปฺปมาณํ อโหสิ, ปญฺจวสฺสสติกา กุมาริกา อลํปเตยฺยา อโหสิ.\csromanspacer{} เตน โข ปน ภิกฺขเว สมเยน มนุสฺสานํ ฉเฬว อาพาธา อเหสุํ สีตํ อุณฺหํ ชีฆจฺฉา ปิปาสา อุจฺจาโร ปสฺสาโว.\csromanspacer{} โส หิ นาม ภิกฺขเว อรโก สตฺถา เอวํ ทีฆายุเกสุ มนุสฺเสสุ เอวํ จิรฏฺฐิติเกสุ เอวํ อปฺปาพาเธสุ สาวกานํ เอวํ ธมฺมํ เทเสสฺสติ “อปฺปกํ พฺราหฺมณ ชีวิตํ มนุสฺสานํ ปริตฺตํ ลหุกํ พหุทุกฺขํ พหุปายาสํ มนฺตายํ โพทฺธพฺพํ, กตฺตพฺพํ กุสลํ, จริตพฺพํ พฺรหฺมจริยํ, นตฺถิ ชาตสฺส อมรณนฺ”ติ.}

\prose{เอตรหิ ตํ ภิกฺขเว สมฺมา วทมาโน วเทยฺย “อปฺปกํ ชีวิตํ มนุสฺสานํ ปริตฺตํ ลหุกํ พหุทุกฺขํ พหุปายาสํ มนฺตายํ โพทฺธพฺพํ, กตฺตพฺพํ กุสลํ, จริตพฺพํ พฺรหฺมจริยํ, นตฺถิ ชาตสฺส อมรณนฺ”ติ.\csromanspacer{} เอตรหิ ภิกฺขเว โย จิรํ ชีวติ, โส วสฺสสตํ อปฺปํ วา ภิยฺโย.\csromanspacer{} วสฺสสตํ โข ปน ภิกฺขเว ชีวนฺโต ตีณิเยว อุตุสตานิ ชีวติ, อุตุสตํ เหมนฺตานํ, อุตุสตํ คิมฺหานํ, อุตุสตํ วสฺสานํ.\csromanspacer{} ตีณิ โข ปน ภิกฺขเว อุตุสตานิ ชีวนฺโต ทฺวาทส\footnote{ทฺวาทสํ (สี, ก)} เยว มาสสตานิ ชีวติ.\csromanspacer{} จตฺตาริ มาสสตานิ เหมนฺตานํ, จตฺตาริ มาสสตานิ คิมฺหานํ, จตฺตาริ มาสสตานิ วสฺสานํ.\csromanspacer{} ทฺวาทส โข ปน ภิกฺขเว มาสสตานิ ชีวนฺโต จตุวีสติเยว อทฺธมาสสตานิ ชีวติ, อฏฺฐทฺธมาสสตานิ เหมนฺตานํ, อฏฺฐทฺธมาสสตานิ คิมฺหานํ, อฏฺฐทฺธมาสสตานิ วสฺสานํ.\csromanspacer{} จตุวีสติ โข ปน ภิกฺขเว อทฺธมาสสตานิ ชีวนฺโต ฉตฺติํสํเยว รตฺติสหสฺสานิ ชีวติ, ทฺวาทส รตฺติสหสฺสานิ เหมนฺตานํ, ทฺวาทส รตฺติสหสฺสานิ คิมฺหานํ, ทฺวาทส รตฺติสหสฺสานิ วสฺสานํ.\csromanspacer{} ฉตฺติํสํ โข ปน ภิกฺขเว รตฺติสหสฺสานิ ชีวนฺโต เทฺวสตฺตติเยว\footnote{เทฺวสตฺตติญฺเญว (สฺยา), เทฺวสตฺตติญฺเจว (ก)} ภตฺตสหสฺสานิ ภุญฺชติ, จตุวีสติ ภตฺตสหสฺสานิ เหมนฺตานํ, จตุวีสติ ภตฺตสหสฺสานิ คิมฺหานํ, จตุวีสติ ภตฺตสหสฺสานิ วสฺสานํ สทฺธิํ มาตุถญฺญาย สทฺธิํ ภตฺตนฺตราเยน.}

\prose{ตตฺริเม ภตฺตนฺตรายา, กปิมิทฺโธปิ ภตฺตํ น ภุญฺชติ, ทุกฺขิโตปิ ภตฺตํ น ภุญฺชติ, พฺยาธิโตปิ ภตฺตํ น ภุญฺชติ, อุโปสถิโกปิ ภตฺตํ น}

\vspace{\tipitakaparskip}
\csromancenter{สตฺตกนิปาตปาฬิ ภุญฺชติ, อลาภเกนปิ ภตฺตํ น ภุญฺชติ.\csromanspacer{} อิติ โข ภิกฺขเว มยา วสฺสสตายุกสฺส มนุสฺสสฺส อายุปิ สงฺขาโต\footnote{สงฺขาตํ (?)}, อายุปฺปมาณมฺปิ สงฺขาตํ, อุตูปิ สงฺขาตา, สํวจฺฉราปิ สงฺขาตา, มาสาปิ สงฺขาตา, อทฺธมาสาปิ สงฺขาตา, รตฺติปิ สงฺขาตา, ทิวาปิ สงฺขาตา, ภตฺตาปิ สงฺขาตา, ภตฺตนฺตรายาปิ สงฺขาตา.\csromanspacer{} ยํ ภิกฺขเว สตฺถารา กรณียํ สาวกานํ หิเตสินา อนุกมฺปเกน อนุกมฺปํ อุปาทาย.\csromanspacer{} กตํ โว ตํ มยา, เอตานิ ภิกฺขเว รุกฺขมูลานิ, เอตานิ สุญฺญาคารานิ, ฌายถ ภิกฺขเว มา ปมาทตฺถ, มา ปจฺฉา วิปฺปฏิสาริโน อหุวตฺถ, อยํ โว อมฺหากํ อนุสาสนีติ.\csromanspacer{}\csromanspacer{}ทสมํ.\csromanspacer{} มหาวคฺโค สตฺตโม.}
\titlehead{\textbf{ตสฺสุทฺทานํ}}
\csromanmatikaanchor{mtk.271}
\csromanmatikaanchor{mtk.272}
\csromantocmarknopage{chapter}{8. วินยวคฺค}{mtk.271}
\bookmark[dest={mtk.271},level=0]{8. วินยวคฺค}
\csromantocmarkref{subsection}{1-4. วินยธรสุตฺตานิ}{mtk.272}
\bookmark[dest={mtk.272},level=2]{1-4. วินยธรสุตฺตานิ}
\csromangathameasure{หิรีสูริยํ อุปมา, ธมฺมญฺญู ปาริฉตฺตกํ. \\ สกฺกจฺจํ ภาวนา อคฺคิ, สุเนตฺต-อรเกน จาติ.}
\csromangathagroup{\csromangathastanza{\csromanfolio{504}หิรีสูริยํ อุปมา, ธมฺมญฺญู ปาริฉตฺตกํ. \\ สกฺกจฺจํ ภาวนา อคฺคิ, สุเนตฺต-อรเกน จาติ.}}

\chapterhead{8. วินยวคฺค}
\csromanheader{2}{1. ปฐมวินยธรสุตฺต}
\proseitem{75}{สตฺตหิ ภิกฺขเว ธมฺเมหิ สมนฺนาคโต ภิกฺขุ วินยธโร โหติ.\csromanspacer{} กตเมหิ สตฺตหิ, อาปตฺติํ ชานาติ, อนาปตฺติํ ชานาติ, ลหุกํ อาปตฺติํ ชานาติ, ครุกํ อาปตฺติํ ชานาติ, สีลวา โหติ, ปาติโมกฺขสํวรสํวุโต วิหรติ, อาจารโคจรสมฺปนฺโน อณุมตฺเตสุ วชฺเชสุ ภยทสฺสาวี, สมาทาย สิกฺขติ สิกฺขาปเทสุ, จตุนฺนํ ฌานานํ อาภิเจตสิกานํ ทิฏฺฐธมฺมสุขวิหารานํ นิกามลาภี โหติ อกิจฺฉลาภี อกสิรลาภี, อาสวานํ ขยา อนาสวํ เจโตวิมุตฺติํ ปญฺญาวิมุตฺติํ ทิฏฺเฐว ธมฺเม สยํ อภิญฺญา สจฺฉิกตฺวา อุปสมฺปชฺช วิหรติ.\csromanspacer{} อิเมหิ โข ภิกฺขเว สตฺตหิ ธมฺเมหิ สมนฺนาคโต ภิกฺขุ วินยธโร โหตีติ.\csromanspacer{}\csromanspacer{}ปฐมํ.}

\chapterheadpageread{8. วินยวคฺค \textbf{2. ทุติยวินยธรสุตฺต}}
\proseitem{76}{\csromanfolio{505}สตฺตหิ ภิกฺขเว ธมฺเมหิ สมนฺนาคโต ภิกฺขุ วินยธโร โหติ.\csromanspacer{} กตเมหิ สตฺตหิ, อาปตฺติํ ชานาติ, อนาปตฺติํ ชานาติ, ลหุกํ อาปตฺติํ ชานาติ, ครุกํ อาปตฺติํ ชานาติ, อุภยานิ โข ปนสฺส ปาติโมกฺขานิ วิตฺถาเรน สฺวาคตานิ โหนฺติ สุวิภตฺตานิ สุปฺปวตฺตีนิ, สุวินิจฺฉิตานิ สุตฺตโส อนุพฺยญฺชนโส, จตุนฺนํ ฌานานํ อาภิเจตสิกานํ ทิฏฺฐธมฺมสุข- วิหารานํ นิกามลาภี โหติ อกิจฺฉลาภี อกสิรลาภี, อาสวานํ ขยา อนาสวํ เจโตวิมุตฺติํ ปญฺญาวิมุตฺติํ ทิฏฺเฐว ธมฺเม สยํ อภิญฺญา สจฺฉิกตฺวา อุปสมฺปชฺช วิหรติ.\csromanspacer{} อิเมหิ โข ภิกฺขเว สตฺตหิ ธมฺเมหิ สมนฺนาคโต ภิกฺขุ วินยธโร โหตีติ.\csromanspacer{}\csromanspacer{}ทุติยํ.}

\csromanheader{2}{3. ตติยวินยธรสุตฺต}
\proseitem{77}{สตฺตหิ ภิกฺขเว ธมฺเมหิ สมนฺนาคโต ภิกฺขุ วินยธโร โหติ.\csromanspacer{} กตเมหิ สตฺตหิ, อาปตฺติํ ชานาติ, อนาปตฺติํ ชานาติ, ลหุกํ อาปตฺติํ ชานาติ, ครุกํ อาปตฺติํ ชานาติ, วินเย โข ปน ฐิโต โหติ อสํหีโร, จตุนฺนํ ฌานานํ อาภิเจตสิกานํ ทิฏฺฐธมฺมสุขวิหารานํ นิกามลาภี โหติ อกิจฺฉลาภี อกสิรลาภี, อาสวานํ ขยา อนาสวํ เจโตวิมุตฺติํ ปญฺญาวิมุตฺติํ ทิฏฺเฐว ธมฺเม สยํ อภิญฺญา สจฺฉิกตฺวา อุปสมฺปชฺช วิหรติ.\csromanspacer{} อิเมหิ โข ภิกฺขเว สตฺตหิ ธมฺเมหิ สมนฺนาคโต ภิกฺขุ วินยธโร โหตีติ.\csromanspacer{}\csromanspacer{}ตติยํ.}

\csromanheader{2}{4. จตุตฺถวินยธรสุตฺต}
\proseitem{78}{สตฺตหิ ภิกฺขเว ธมฺเมหิ สมนฺนาคโต ภิกฺขุ วินยธโร โหติ.\csromanspacer{} กตเมหิ สตฺตหิ, อาปตฺติํ ชานาติ, อนาปตฺติํ ชานาติ, ลหุกํ อาปตฺติํ ชานาติ, ครุกํ อาปตฺติํ ชานาติ, อเนกวิหิตํ ปุพฺเพนิวาสํ อนุสฺสรติ.\csromanspacer{} เสยฺยถิทํ, เอกมฺปิ ชาติํ เทฺวปิ ชาติโย ฯเปฯ อิติ สาการํ ส-อุทฺเทสํ อเนกวิหิตํ ปุพฺเพนิวาสํ อนุสฺสรติ, ทิพฺเพน จกฺขุนา วิสุทฺเธน อติกฺกนฺตมานุสเกน, ยถากมฺมูปเค สตฺเต ปชานาติ, อาสวานํ ขยา อนาสวํ เจโตวิมุตฺติํ ปญฺญาวิมุตฺติํ ทิฏฺเฐว ธมฺเม}

\vspace{\tipitakaparskip}
\csromancenter{สตฺตกนิปาตปาฬิ สยํ อภิญฺญา สจฺฉิกตฺวา อุปสมฺปชฺช วิหรติ.\csromanspacer{} อิเมหิ โข ภิกฺขเว สตฺตหิ ธมฺเมหิ สมนฺนาคโต ภิกฺขุ วินยธโร โหตีติ.\csromanspacer{}\csromanspacer{}จตุตฺถํ.}
\csromanheader{2}{5. ปฐมวินยธรโสภนสุตฺต}
\csromanmatikaanchor{mtk.273}
\csromantocmarkref{subsection}{5-8. วินยธรโสภนสุตฺตานิ}{mtk.273}
\bookmark[dest={mtk.273},level=2]{5-8. วินยธรโสภนสุตฺตานิ}
\proseitem{79}{\csromanfolio{506}สตฺตหิ ภิกฺขเว ธมฺเมหิ สมนฺนาคโต\footnote{สมนฺนาคโต ภิกฺขุ (สี, สฺยา, ก) อนนฺตรสุตฺตทฺวเย ปน อิทํ ปาฐนานตฺตํ นตฺถิ.} วินยธโร โสภติ.\csromanspacer{} กตเมหิ สตฺตหิ, อาปตฺติํ ชานาติ, อนาปตฺติํ ชานาติ, ลหุกํ อาปตฺติํ ชานาติ, ครุกํ อาปตฺติํ ชานาติ, สีลวา โหติ ฯเปฯ สมาทาย สิกฺขติ สิกฺขาปเทสุ, จตุนฺนํ ฌานานํ อาภิเจตสิกานํ นิกามลาภี โหติ อกิจฺฉลาภี อกสิรลาภี, อาสวานํ ขยา ฯเปฯ สจฺฉิกตฺวา อุปสมฺปชฺช วิหรติ.\csromanspacer{} อิเมหิ โข ภิกฺขเว สตฺตหิ ธมฺเมหิ สมนฺนาคโต วินยธโร โสภตีติ.\csromanspacer{}\csromanspacer{}ปญฺจมํ.}

\csromanheader{2}{6. ทุติยวินยธรโสภนสุตฺต}
\proseitem{80}{สตฺตหิ ภิกฺขเว ธมฺเมหิ สมนฺนาคโต วินยธโร โสภติ.\csromanspacer{} กตเมหิ สตฺตหิ, อาปตฺติํ ชานาติ, อนาปตฺติํ ชานาติ, ลหุกํ อาปตฺติํ ชานาติ, ครุกํ อาปตฺติํ ชานาติ, อุภยานิ โข ปนสฺส ปาติโมกฺขานิ วิตฺถาเรน สฺวาคตานิ โหนฺติ สุวิภตฺตานิ สุปฺปวตฺตีนิ สุวินิจฺฉิตานิ สุตฺตโส อนุพฺยญฺชนโส, จตุนฺนํ ฌานานํ ฯเปฯ อกสิรลาภี, อาสวานํ ขยา ฯเปฯ สจฺฉิกตฺวา อุปสมฺปชฺช วิหรติ.\csromanspacer{} อิเมหิ โข ภิกฺขเว สตฺตหิ ธมฺเมหิ สมนฺนาคโต วินยธโร โสภตีติ.\csromanspacer{}\csromanspacer{}ฉฏฺฐํ.}

\csromanheader{2}{7. ตติยวินยธรโสภนสุตฺต}
\proseitem{81}{สตฺตหิ ภิกฺขเว ธมฺเมหิ สมนฺนาคโต วินยธโร โสภติ.\csromanspacer{} กตเมหิ สตฺตหิ, อาปตฺติํ ชานาติ, อนาปตฺติํ ชานาติ, ลหุกํ อาปตฺติํ ชานาติ, ครุกํ อาปตฺติํ ชานาติ, วินเย โข ปน ฐิโต โหติ อสํหีโร, จตุนฺนํ ฌานานํ ฯเปฯ อกสิรลาภี, อาสวานํ ขยา ฯเปฯ สจฺฉิกตฺวา อุปสมฺปชฺช วิหรติ.\csromanspacer{} อิเมหิ โข ภิกฺขเว สตฺตหิ ธมฺเมหิ สมนฺนาคโต วินยธโร โสภตีติ.\csromanspacer{}\csromanspacer{}สตฺตมํ.}

\chapterheadpageread{8. วินยวคฺค \textbf{8. จตุตฺถวินยธรโสภนสุตฺต}}
\proseitem{82}{\csromanfolio{507}สตฺตหิ ภิกฺขเว ธมฺเมหิ สมนฺนาคโต\footnote{สมนฺนาคโต ภิกฺขุ (ก)} วินยธโร โสภติ.\csromanspacer{} กตเมหิ สตฺตหิ, อาปตฺติํ ชานาติ, อนาปตฺติํ ชานาติ, ลหุกํ อาปตฺติํ ชานาติ, ครุกํ อาปตฺติํ ชานาติ, อเนกวิหิตํ ปุพฺเพนิวาสํ อนุสฺสรติ.\csromanspacer{} เสยฺยถิทํ, เอกมฺปิ ชาติํ เทฺวปิ ชาติโย ฯเปฯ อิติ สาการํ ส-อุทฺเทสํ อเนกวิหิตํ ปุพฺเพนิวาสํ อนุสฺสรติ, ทิพฺเพน จกฺขุนา วิสุทฺเธน อติกฺกนฺตมานุสเกน ฯเปฯ อาสวานํ ขยา ฯเปฯ สจฺฉิกตฺวา อุปสมฺปชฺช วิหรติ.\csromanspacer{} อิเมหิ โข ภิกฺขเว สตฺตหิ ธมฺเมหิ สมนฺนาคโต วินยธโร โสภตีติ.\csromanspacer{}\csromanspacer{}อฏฺฐมํ.}

\csromanmatikaanchor{mtk.274}
\csromantocmarkref{subsection}{9. สตฺถุสาสนสุตฺต}{mtk.274}
\bookmark[dest={mtk.274},level=2]{9. สตฺถุสาสนสุตฺต}
\csromanheader{2}{9. สตฺถุสาสนสุตฺต}
\proseitem{83}{อถ โข อายสฺมา อุปาลิ เยน ภควา เตนุปสงฺกมิ อุปสงฺกมิตฺวา ภควนฺตํ อภิวาเทตฺวา เอกมนฺตํ นิสีทิ, เอกมนฺตํ นิสินฺโน โข อายสฺมา อุปาลิ ภควนฺตํ เอตทโวจ\csromandash{}}

\prose{สาธุ เม ภนฺเต ภควา สํขิตฺเตน ธมฺมํ เทเสตุ, ยมหํ ภควโต ธมฺมํ สุตฺวา เอโก วูปกฏฺโฐ อปฺปมตฺโต อาตาปี ปหิตตฺโต วิหเรยฺยนฺติ.\csromanspacer{} เย โข ตฺวํ อุปาลิ ธมฺเม ชาเนยฺยาสิ “อิเม ธมฺมา น เอกนฺตนิพฺพิทาย วิราคาย นิโรธาย อุปสมาย อภิญฺญาย สมฺโพธาย นิพฺพานาย สํวตฺตนฺตี”ติ.\csromanspacer{} เอกํเสน อุปาลิ ธาเรยฺยาสิ “เนโส ธมฺโม เนโส วินโย เนตํ สตฺถุสาสนนฺ”ติ.\csromanspacer{} เย จ โข ตฺวํ อุปาลิ ธมฺเม ชาเนยฺยาสิ “อิเม ธมฺมา เอกนฺตนิพฺพิทาย วิราคาย นิโรธาย อุปสมาย อภิญฺญาย สมฺโพธาย นิพฺพานาย สํวตฺตนฺตี”ติ.\csromanspacer{} เอกํเสน อุปาลิ ธาเรยฺยาสิ “เอโส ธมฺโม เอโส วินโย เอตํ สตฺถุสาสนนฺ”ติ.\csromanspacer{}\csromanspacer{}นวมํ.}

\csromanmatikaanchor{mtk.275}
\csromantocmarkref{subsection}{10. อธิกรณสมถสุตฺต}{mtk.275}
\bookmark[dest={mtk.275},level=2]{10. อธิกรณสมถสุตฺต}
\csromanheader{2}{10. อธิกรณสมถสุตฺต}
\proseitem{84}{สตฺติเม ภิกฺขเว อธิกรณสมถา ธมฺมา อุปฺปนฺนุปฺปนฺนานํ อธิกรณานํ สมถาย วูปสมาย.\csromanspacer{} กตเม สตฺต, สมฺมุขาวินโย}

\vspace{\tipitakaparskip}
\csromancenter{สตฺตกนิปาตปาฬิ ทาตพฺโพ, สติวินโย ทาตพฺโพ, อมูฬฺหวินโย ทาตพฺโพ, 1 ปฏิญฺญาตกรณํ ทาตพฺพํ, เยภุยฺยสิกา ทาตพฺพา, ตสฺสปาปิยสิกา ทาตพฺพา, ติณวตฺถารโก ทาตพฺโพ1.\csromanspacer{} อิเม โข ภิกฺขเว สตฺต อธิกรณสมถา ธมฺมา อุปฺปนฺนุปฺปนฺนานํ อธิกรณานํ สมถาย วูปสมายาติ.\csromanspacer{} วินยวคฺโค อฏฺฐโม.}
\titlehead{\textbf{ตสฺสุทฺทานํ}}
\csromanmatikaanchor{mtk.276}
\csromanmatikaanchor{mtk.277}
\csromantocmarknopage{chapter}{9. สมณวคฺค}{mtk.276}
\bookmark[dest={mtk.276},level=0]{9. สมณวคฺค}
\csromantocmarkref{subsection}{1-10. ภิกฺขุสุตฺตาทีนิ}{mtk.277}
\bookmark[dest={mtk.277},level=2]{1-10. ภิกฺขุสุตฺตาทีนิ}
\csromantocmarknopage{chapter}{10. อาหุเนยฺยวคฺค}{mtk.278}
\bookmark[dest={mtk.278},level=0]{10. อาหุเนยฺยวคฺค}
\csromangathameasure{จตุโร วินยธรา, จตุโร เจว โสภนา. \\ สาสนํ อธิกรณ-สมเถนฏฺฐเม ทสาติ.}
\csromangathagroup{\csromangathastanza{\csromanfolio{508}จตุโร วินยธรา, จตุโร เจว โสภนา. \\ สาสนํ อธิกรณ-สมเถนฏฺฐเม ทสาติ.}}

\chapterhead{9. สมณวคฺค}
\csromanheader{2}{1. ภิกฺขุสุตฺต}
\proseitem{85}{\csromansymbolfootnote{*}{ขุ 7. 54; ขุ 8. 39 ปิฏฺเฐสุปิ.}สตฺตนฺนํ ภิกฺขเว ธมฺมานํ ภินฺนตฺตา ภิกฺขุ โหติ.\csromanspacer{} กตเมสํ สตฺตนฺนํ? สกฺกายทิฏฺฐิ ภินฺนา โหติ, วิจิกิจฺฉา ภินฺนา โหติ, สีลพฺพตปรามาโส ภินฺโน โหติ, ราโค ภินฺโน โหติ, โทโส ภินฺโน โหติ, โมโห ภินฺโน โหติ, มาโน ภินฺโน โหติ.\csromanspacer{} อิเมสํ โข ภิกฺขเว สตฺตนฺนํ ธมฺมานํ ภินฺนตฺตา ภิกฺขุ โหตีติ.\csromanspacer{}\csromanspacer{}ปฐมํ.}

\csromanheader{2}{2. สมณสุตฺต}
\vspace{\tipitakaparskip}
\csromancenter{สตฺตนฺนํ ภิกฺขเว ธมฺมานํ สมิตตฺตา สมโณ โหติ ฯเปฯ.\csromanspacer{}\csromanspacer{}ทุติยํ.}
\csromanheader{2}{3. พฺราหฺมณสุตฺต}
\vspace{\tipitakaparskip}
\csromancenter{พาหิตตฺตา พฺราหฺมโณ โหติ ฯเปฯ.\csromanspacer{}\csromanspacer{}ตติยํ.}
\csromanheader{2}{4. โสตฺติยสุตฺต}
\proseitem{88}{นิสฺสุตตฺตา\footnote{นิสฺสุตฺตตฺตา (สฺยา) ม 1. 348 ปิฏฺเฐ ปสฺสิตพฺพํ.} โสตฺติโย\footnote{โสตฺถิโก (สี), โสตฺติโก (สฺยา)} โหติ ฯเปฯ.\csromanspacer{}\csromanspacer{}จตุตฺถํ.}

\chapterheadpageread{9. สมณวคฺค \textbf{5. นฺหาตกสุตฺต}}
\proseitem{89}{\csromanfolio{509}นฺหาตตฺตา นฺหาตโก\footnote{นหาตตฺตา นหาตโก (สี, สฺยา)} โหติ ฯเปฯ.\csromanspacer{}\csromanspacer{}ปญฺจมํ.}

\csromanheader{2}{6. เวทคูสุตฺต}
\vspace{\tipitakaparskip}
\csromancenter{วิทิตตฺตา เวทคู โหติ ฯเปฯ.\csromanspacer{}\csromanspacer{}ฉฏฺฐํ.}
\csromanheader{2}{7. อริยสุตฺต}
\vspace{\tipitakaparskip}
\csromancenter{อารกตฺตา\footnote{อรหตฺตา (สี), อรี หตตฺตา (ก) ม 1. 348 ปิฏฺเฐ ปาฬิ อฏฺฐกถาฏีกา ปสฺสิตพฺพา. สฺยามโปตฺถเก ปน สกลมฺปิ อิทํ สตฺตมสุตฺตํ นตฺถิ.} อริโย โหติ ฯเปฯ.\csromanspacer{}\csromanspacer{}สตฺตมํ.}
\csromanheader{2}{8. อรหาสุตฺต}
\proseitem{92}{อารกตฺตา อรหา โหติ.\csromanspacer{} กตเมสํ สตฺตนฺนํ? สกฺกายทิฏฺฐิ อารกา โหติ, วิจิกิจฺฉา อารกา โหติ, สีลพฺพตปรามาโส อารโก โหติ, ราโค อารโก โหติ, โทโส อารโก โหติ, โมโห อารโก โหติ, มาโน อารโก โหติ.\csromanspacer{} อิเมสํ โข ภิกฺขเว สตฺตนฺนํ ธมฺมานํ อารกตฺตา อรหา โหตีติ.\csromanspacer{}\csromanspacer{}อฏฺฐมํ.}

\csromanheader{2}{9. อสทฺธมฺมสุตฺต}
\proseitem{93}{สตฺติเม ภิกฺขเว อสทฺธมฺมา.\csromanspacer{} กตเม สตฺต? อสฺสทฺโธ โหติ, อหิริโก โหติ, อโนตฺตปฺปี โหติ, อปฺปสฺสุโต โหติ, กุสีโต โหติ, มุฏฺฐสฺสติ โหติ, ทุปฺปญฺโญ โหติ.\csromanspacer{} อิเม โข ภิกฺขเว สตฺต อสทฺธมฺมาติ.\csromanspacer{}\csromanspacer{}นวมํ.}

\csromanheader{2}{10. สทฺธมฺมสุตฺต}
\proseitem{94}{สตฺติเม ภิกฺขเว สทฺธมฺมา.\csromanspacer{} กตเม สตฺต? สทฺโธ โหติ, หิรีมา โหติ, โอตฺตปฺปี โหติ, พหุสฺสุโต โหติ, อารทฺธวีริโย โหติ, สติมา โหติ, ปญฺญวา โหติ.\csromanspacer{} อิเม โข ภิกฺขเว สตฺต สทฺธมฺมาติ.\csromanspacer{}\csromanspacer{}ทสมํ.\csromanspacer{} สมณวคฺโค นวโม.}

\gambhira{สตฺตกนิปาตปาฬิ \textbf{ตสฺสุทฺทานํ}}
\csromanmatikaanchor{mtk.279}
\csromantocmarkref{section}{ปฐมสุตฺตาทีนิ}{mtk.279}
\bookmark[dest={mtk.279},level=1]{ปฐมสุตฺตาทีนิ}
\csromantocmarknopage{chapter}{11. ราคเปยฺยาล}{mtk.280}
\bookmark[dest={mtk.280},level=0]{11. ราคเปยฺยาล}
\csromangathameasure{ภิกฺขุํ สมโณ พฺราหฺมโณ, โสตฺติโย เจว นฺหาตโก. \\ เวทคู อริโย อรหา, อสทฺธมฺมา จ สทฺธมฺมาติ.}
\csromangathagroup{\csromangathastanza{\csromanfolio{510}ภิกฺขุํ สมโณ พฺราหฺมโณ, โสตฺติโย เจว นฺหาตโก. \\ เวทคู อริโย อรหา, อสทฺธมฺมา จ สทฺธมฺมาติ.}}

\chapterhead{10. อาหุเนยฺยวคฺค}
\proseitem{95}{สตฺติเม ภิกฺขเว ปุคฺคลา อาหุเนยฺยา ฯเปฯ ทกฺขิเณยฺยา อญฺชลิกรณียา อนุตฺตรํ ปุญฺญกฺเขตฺตํ โลกสฺส.\csromanspacer{} กตเม สตฺต, อิธ ภิกฺขเว เอกจฺโจ ปุคฺคโล จกฺขุสฺมิํ อนิจฺจานุปสฺสี วิหรติ อนิจฺจสญฺญี อนิจฺจปฏิสํเวที สตตํ สมิตํ อพฺโพกิณฺณํ เจตสา อธิมุจฺจมาโน ปญฺญาย ปริโยคาหมาโน, โส อาสวานํ ขยา อนาสวํ เจโตวิมุตฺติํ ปญฺญาวิมุตฺติํ ทิฏฺเฐว ธมฺเม สยํ อภิญฺญา สจฺฉิกตฺวา อุปสมฺปชฺช วิหรติ.\csromanspacer{} อยํ โข ภิกฺขเว ปฐโม ปุคฺคโล อาหุเนยฺโย ปาหุเนยฺโย ฯเปฯ อนุตฺตรํ ปุญฺญกฺเขตฺตํ โลกสฺส.}

\prose{ปุน จปรํ ภิกฺขเว อิเธกจฺโจ ปุคฺคโล จกฺขุสฺมิํ อนิจฺจานุปสฺสี วิหรติ อนิจฺจสญฺญี อนิจฺจปฏิสํเวที สตตํ สมิตํ อพฺโพกิณฺณํ เจตสา อธิมุจฺจมาโน ปญฺญาย ปริโยคาหมาโน, ตสฺส อปุพฺพํ อจริมํ อาสวปริยาทานญฺจ โหติ ชีวิตปริยาทานญฺจ.\csromanspacer{} อยํ ภิกฺขเว ทุติโย ปุคฺคโล อาหุเนยฺโย ฯเปฯ อนุตฺตรํ ปุญฺญกฺเขตฺตํ โลกสฺส.}

\prose{ปุน จปรํ ภิกฺขเว อิเธกจฺโจ ปุคฺคโล จกฺขุสฺมิํ อนิจฺจานุปสฺสี วิหรติ อนิจฺจสญฺญี อนิจฺจปฏิสํเวที สตตํ สมิตํ อพฺโพกิณฺณํ เจตสา อธิมุจฺจมาโน ปญฺญาย ปริโยคาหมาโน, โส ปญฺจนฺนํ โอรมฺภาคิยานํ สํโยชนานํ ปริกฺขยา อนฺตราปรินิพฺพายี โหติ ฯเปฯ อุปหจฺจปรินิพฺพายี โหติ ฯเปฯ อสงฺขารปรินิพฺพายี โหติ ฯเปฯ สสงฺขารปรินิพฺพายี โหติ ฯเปฯ อุทฺธํโสโต โหติ อกนิฏฺฐคามี.\csromanspacer{} อยํ ภิกฺขเว สตฺตโม ปุคฺคโล อาหุเนยฺโย ฯเปฯ อนุตฺตรํ ปุญฺญกฺเขตฺตํ โลกสฺส.\csromanspacer{} อิเม โข ภิกฺขเว สตฺต ปุคฺคลา อาหุเนยฺยา ปาหุเนยฺยา ทกฺขิเณยฺยา อญฺชลิกรณียา อนุตฺตรํ ปุญฺญกฺเขตฺตํ โลกสฺสาติ. (1)}

\chapterheadpageread{10. อาหุเนยฺยวคฺค}
\proseitem{96-622}{\csromanfolio{511}สตฺติเม ภิกฺขเว ปุคฺคลา อาหุเนยฺยา ปาหุเนยฺยา ฯเปฯ อนุตฺตรํ ปุญฺญกฺเขตฺตํ โลกสฺส.\csromanspacer{} กตเม สตฺต? อิธ ภิกฺขเว เอกจฺโจ ปุคฺคโล จกฺขุสฺมิํ ทุกฺขานุปสฺสี วิหรติ ฯเปฯ จกฺขุสฺมิํ อนตฺตานุปสฺสี วิหรติ ฯเปฯ จกฺขุสฺมิํ ขยานุปสฺสี วิหรติ ฯเปฯ จกฺขุสฺมิํ วยานุปสฺสี วิหรติ ฯเปฯ จกฺขุสฺมิํ วิราคานุปสฺสี วิหรติ ฯเปฯ จกฺขุสฺมิํ นิโรธานุปสฺสี วิหรติ ฯเปฯ จกฺขุสฺมิํ ปฏินิสฺสคฺคานุปสฺสี วิหรติ ฯเปฯ. (2-8)}

\prose{โสตสฺมิํ ฯเปฯ.\csromanspacer{} ฆานสฺมิํ.\csromanspacer{} ชิวฺหาย.\csromanspacer{} กายสฺมิํ.\csromanspacer{} มนสฺมิํ ฯเปฯ. (9- 48)}

\prose{รูเปสุ ฯเปฯ.\csromanspacer{} สทฺเทสุ.\csromanspacer{} คนฺเธสุ.\csromanspacer{} รเสสุ.\csromanspacer{} โผฏฺฐพฺเพสุ.\csromanspacer{} ธมฺเมสุ ฯเปฯ. (49-96)}

\prose{จกฺขุวิญฺญาเณ ฯเปฯ.\csromanspacer{} โสตวิญฺญาเณ.\csromanspacer{} ฆานวิญฺญาเณ.\csromanspacer{} ชิวฺหาวิญฺญาเณ.\csromanspacer{} กายวิญฺญาเณ.\csromanspacer{} มโนวิญฺญาเณ ฯเปฯ. (97-144)}

\prose{จกฺขุสมฺผสฺเส ฯเปฯ.\csromanspacer{} โสตสมฺผสฺเส.\csromanspacer{} ฆานสมฺผสฺเส.\csromanspacer{} ชิวฺหาสมฺผสฺเส.\csromanspacer{} กายสมฺผสฺเส.\csromanspacer{} มโนสมฺผสฺเส ฯเปฯ. (145-192)}

\prose{จกฺขุสมฺผสฺสชาย เวทนาย ฯเปฯ.\csromanspacer{} โสตสมฺผสฺสชาย เวทนาย.\csromanspacer{} ฆานสมฺผสฺสชาย เวทนาย.\csromanspacer{} ชิวฺหาสมฺผสฺสชาย เวทนาย.\csromanspacer{} กายสมฺผสฺสชาย เวทนาย.\csromanspacer{} มโนสมฺผสฺสชาย เวทนาย ฯเปฯ. (193- 240)}

\prose{รูปสญฺญาย ฯเปฯ.\csromanspacer{} สทฺทสญฺญาย.\csromanspacer{} คนฺธสญฺญาย.\csromanspacer{} รสสญฺญาย.\csromanspacer{} โผฏฺฐพฺพสญฺญาย.\csromanspacer{} ธมฺมสญฺญาย ฯเปฯ. (241-288)}

\prose{รูปสญฺเจตนาย ฯเปฯ.\csromanspacer{} สทฺทสญฺเจตนาย.\csromanspacer{} คนฺธสญฺเจตนาย.\csromanspacer{} รสสญฺเจตนาย.\csromanspacer{} โผฏฺฐพฺพสญฺเจตนาย.\csromanspacer{} ธมฺมสญฺเจตนาย ฯเปฯ. (289- 336)}

\prose{รูปตณฺหาย ฯเปฯ.\csromanspacer{} สทฺทตณฺหาย.\csromanspacer{} คนฺธตณฺหาย.\csromanspacer{} รสตณฺหาย.\csromanspacer{} โผฏฺฐพฺพตณฺหาย.\csromanspacer{} ธมฺมตณฺหาย ฯเปฯ. (337-384)}

\prose{รูปวิตกฺเก ฯเปฯ.\csromanspacer{} สทฺทวิตกฺเก.\csromanspacer{} คนฺธวิตกฺเก.\csromanspacer{} รสวิตกฺเก.\csromanspacer{} โผฏฺฐพฺพวิตกฺเก.\csromanspacer{} ธมฺมวิตกฺเก ฯเปฯ. (385-432)}

\prose{รูปวิจาเร ฯเปฯ.\csromanspacer{} สทฺทวิจาเร.\csromanspacer{} คนฺธวิจาเร.\csromanspacer{} รสวิจาเร.\csromanspacer{} โผฏฺฐพฺพวิจาเร.\csromanspacer{} ธมฺมวิจาเร ฯเปฯ. (433-480)}

\vspace{\tipitakaparskip}
\csromancenter{สตฺตกนิปาตปาฬิ (ปญฺจกฺขนฺเธ)\footnote{( ) สี-สฺยา-โปตฺถเกสุ นตฺถิ.} ฯเปฯ.\csromanspacer{} รูปกฺขนฺเธ.\csromanspacer{} เวทนากฺขนฺเธ.\csromanspacer{} สญฺญากฺขนฺเธ.\csromanspacer{} สงฺขารกฺขนฺเธ.\csromanspacer{} วิญฺญาณกฺขนฺเธ อนิจฺจานุปสฺสี วิหรติ ฯเปฯ ทุกฺขานุปสฺสี วิหรติ.\csromanspacer{} อนตฺตานุปสฺสี วิหรติ.\csromanspacer{} ขยานุปสฺสี วิหรติ.\csromanspacer{} วยานุปสฺสี วิหรติ.\csromanspacer{} วิราคานุปสฺสี วิหรติ.\csromanspacer{} นิโรธานุปสฺสี วิหรติ.\csromanspacer{} ปฏินิสฺสคฺคานุปสฺสี วิหรติ ฯเปฯ โลกสฺสาติ. (481-528)}
\vspace{0.5\baselineskip}
\csromanmatikaanchor{mtk.281}
\csromantocmarkref{section}{ปฐมสุตฺตาทีนิ}{mtk.281}
\bookmark[dest={mtk.281},level=1]{ปฐมสุตฺตาทีนิ}
\csromangathameasure{ฉทฺวารารมฺมเณเสฺวตฺถ, วิญฺญาเณสุ จ ผสฺเสสุ. \\ เวทนาสุ จ ทฺวารสฺส, สุตฺตา โหนฺติ วิสุํ อฏฺฐ. \\ สญฺญา สญฺเจตนา ตณฺหา, วิตกฺเกสุ วิจาเร จ. \\ โคจรสฺส วิสุํ อฏฺฐ, ปญฺจกฺขนฺเธ จ ปจฺเจเก. \\ โสฬสเสฺวตฺถ มูเลสุ, อนิจฺจํ ทุกฺขมนตฺตา. \\ ขยา วยา วิราคา จ, นิโรธา ปฏินิสฺสคฺคา. \\ กมํ อฏฺฐานุปสฺสนา, โยเชตฺวาน วิสุํ วิสุํ. \\ สมฺปิณฺฑิเตสุ สพฺเพสุ, โหนฺติ ปญฺจ สตานิ จ. \\ อฏฺฐวีสติ สุตฺตานิ, อาหุเนยฺเย จ วคฺคิเก. อาหุเนยฺยวคฺโค ทสโม.}
\csromangathagroup{\csromangathastanza{\csromanfolio{512}ฉทฺวารารมฺมเณเสฺวตฺถ, วิญฺญาเณสุ จ ผสฺเสสุ. \\ เวทนาสุ จ ทฺวารสฺส, สุตฺตา โหนฺติ วิสุํ อฏฺฐ.}\csromangathastanza{สญฺญา สญฺเจตนา ตณฺหา, วิตกฺเกสุ วิจาเร จ. \\ โคจรสฺส วิสุํ อฏฺฐ, ปญฺจกฺขนฺเธ จ ปจฺเจเก.}\csromangathastanza{โสฬสเสฺวตฺถ มูเลสุ, อนิจฺจํ ทุกฺขมนตฺตา. \\ ขยา วยา วิราคา จ, นิโรธา ปฏินิสฺสคฺคา.}\csromangathastanza{กมํ อฏฺฐานุปสฺสนา, โยเชตฺวาน วิสุํ วิสุํ. \\ สมฺปิณฺฑิเตสุ สพฺเพสุ, โหนฺติ ปญฺจ สตานิ จ. \\ อฏฺฐวีสติ สุตฺตานิ, อาหุเนยฺเย จ วคฺคิเก\footnote{อิมา อุทฺทานคาถาโย สี-สฺยา-โปตฺถเกสุ นตฺถิ.}.\csromanspacer{} อาหุเนยฺยวคฺโค ทสโม.}}

\csromanheader{2}{11. ราคเปยฺยาล}
\proseitem{623}{ราคสฺส ภิกฺขเว อภิญฺญาย สตฺต ธมฺมา ภาเวตพฺพา.\csromanspacer{} กตเม สตฺต? สติสมฺโพชฺฌงฺโค ฯเปฯ อุเปกฺขาสมฺโพชฺฌงฺโค.\csromanspacer{} ราคสฺส ภิกฺขเว อภิญฺญาย อิเม สตฺต ธมฺมา ภาเวตพฺพาติ. (1)}

\proseitem{624}{ราคสฺส ภิกฺขเว อภิญฺญาย สตฺต ธมฺมา ภาเวตพฺพา.\csromanspacer{} กตเม สตฺต? อนิจฺจสญฺญา อนตฺตสญฺญา อสุภสญฺญา อาทีนวสญฺญา ปหานสญฺญา วิราคสญฺญา นิโรธสญฺญา.\csromanspacer{} ราคสฺส ภิกฺขเว อภิญฺญาย อิเม สตฺต ธมฺมา ภาเวตพฺพาติ. (2)}

\csromanheader{2}{11. ราคเปยฺยาล}
\proseitem{625}{\csromanfolio{513}ราคสฺส ภิกฺขเว อภิญฺญาย สตฺต ธมฺมา ภาเวตพฺพา.\csromanspacer{} กตเม สตฺต? อสุภสญฺญา มรณสญฺญา อาหาเร ปฏิกูลสญฺญา สพฺพโลเก อนภิรตสญฺญา อนิจฺจสญฺญา อนิจฺเจ ทุกฺขสญฺญา ทุกฺเข อนตฺตสญฺญา.\csromanspacer{} ราคสฺส ภิกฺขเว อภิญฺญาย อิเม สตฺต ธมฺมา ภาเวตพฺพาติ. (3)}

\proseitem{626-652}{ราคสฺส ภิกฺขเว ปริญฺญาย ฯเปฯ.\csromanspacer{} ปริกฺขยาย.\csromanspacer{} ปหานาย.\csromanspacer{} ขายาย.\csromanspacer{} วยาย.\csromanspacer{} วิราคาย.\csromanspacer{} นิโรธาย.\csromanspacer{} จาคาย ฯเปฯ ปฏินิสฺสคฺคาย อิเม สตฺต ธมฺมา ภาเวตพฺพาติ. (4-30)}

\proseitem{653-1132}{โทสสฺส ฯเปฯ.\csromanspacer{} โมหสฺส.\csromanspacer{} โกธสฺส.\csromanspacer{} อุปนาหสฺส.\csromanspacer{} มกฺขสฺส.\csromanspacer{} ปฬาสสฺส.\csromanspacer{} อิสฺสาย.\csromanspacer{} มจฺฉริยสฺส.\csromanspacer{} มายาย.\csromanspacer{} สาเฐยฺยสฺส.\csromanspacer{} ถมฺภสฺส.\csromanspacer{} สารมฺภสฺส.\csromanspacer{} มานสฺส.\csromanspacer{} อติมานสฺส.\csromanspacer{} มทสฺส.\csromanspacer{} ปมาทสฺส.\csromanspacer{} อภิญฺญาย ฯเปฯ ปริญฺญาย.\csromanspacer{} ปริกฺขยาย.\csromanspacer{} ปหานาย.\csromanspacer{} ขยาย.\csromanspacer{} วยาย.\csromanspacer{} วิราคาย.\csromanspacer{} นิโรธาย.\csromanspacer{} จาคาย.\csromanspacer{} ปฏินิสฺสคฺคาย ฯเปฯ อิเม สตฺต ธมฺมา ภาเวตพฺพาติ. (31-510)}

\prose{อิทมโวจ ภควา.\csromanspacer{} อตฺตมนา เต ภิกฺขู ภควโต ภาสิตํ อภินนฺทุนฺติ.\csromanspacer{} ราคเปยฺยาลํ นิฏฺฐิตํ.}

\nitthitam{\textbf{สตฺตกนิปาตปาฬิ นิฏฺฐิตา.}}
\orphannote{ขุ 5. 131 ปิฏฺเฐ ชาตเกปิ.}

\orphannote{วิ 4. 342, 355 ปิฏฺเฐสุปิ.}

\orphannote{...สิริํสป (สี, สฺยา, กํ, อิ)}

\orphannote{อภิ 3. 174 ปิฏฺเฐปิ.}

\orphannote{เหฏฺฐา 71 ปิฏฺเฐปิ.}

\orphannote{วิ 3. 221; ที 3. 196 ปิฏฺเฐสุ.}

\orphannote{ที 3. 207; ขุ 9. 390 ปิฏฺเฐสุปิ.}

\orphannote{ที 3. 239 ปิฏฺเฐปิ.}

\orphannote{อํ 3. 531 ปิฏฺเฐปิ.}

\orphannote{ปิฏฺฐิกณฺฏกํ อญฺเญน สีสกฏาหํ (สี, อิ), ปิฏฺฐิกณฺฏกฏฺฐิกํ อญฺเญน สีสกฏาหํ (สฺยา, กํ)}

\orphannote{อารามิโก วา ฆฏฺเฏสฺสติ สมณุทฺเทโส วา, ตํ ตมฺหา (ก-สี, อิ) อารามิโก วา ฆฏฺเฏสฺสติ สมณุทฺเทโส วา, โส ตํ ตมฺหา (ก-สี) อํ 3. 156 ปิฏฺเฐ ปสฺสิตพฺพํ.}

\orphannote{นิพฺพาน (สี, สฺยา, กํ, อิ) 3. เนกฺขมฺเม รตํ (ก-สี)}

\orphannote{วิ 3. 267 ปิฏฺเฐปิ อาคตํ.}

\orphannote{วิ 4. 365; อํ 3. 10 ปิฏฺเฐสุ ปสฺสิตพฺพํ.}

\orphannote{ม 1. 119; สํ 2. 427 ปิฏฺเฐสุ.}

\orphannote{อภิ 4. 290 ปิฏฺเฐปิ.}

\orphannote{มุสาวาที โหติ, ปิสุณวาโจ โหติ, ผรุสวาโจ โหติ, สมฺผปฺปลาปี โหติ, (สี, สฺยา, อิ) เอวํ สุกฺกปกฺเขปิ.}

\orphannote{วิสุทฺธิ 1. 127 ปิฏฺฐาทีสุ วิตฺถาโร.}

\orphannote{อภิ 3. 181; อภิ 4. 425 ปิฏฺเฐปิ.}

\orphannote{ที 2. 67 ปิฏฺเฐปิ.}

\orphannote{นาหํ โมคฺคลฺลาน สพฺเพเหว สมคฺคํ วณฺณยามิ คหฏฺเฐหิ, ปพฺพชิเตหิ โข อหํ โมคฺคลฺลาน สมคฺคํ วณฺณยามิ (ก)}

\orphannote{ขุ 1. 252 ปิฏฺเฐ อิติวุตฺตเกปิ.}

\orphannote{วิ 5. 242 ปิฏฺเฐปิ.}

\orphannote{วิ 5. 241 ปิฏฺเฐปิ.}

\orphannote{ปฏิญฺญาตกรณํ, เยภุยฺยสฺสิกา, ตสฺสปาปิยฺยสฺสิกา, ติณวตฺถารโก (สฺยา) ที 3. 210 สงฺคีติสุตฺเตน จ วิ 2. 272 วินเยน จ สํสนฺเทตพฺพํ.}

